\documentclass[11pt,a4paper]{article}
\pdfextension mapfile{+cm.map}
\pdfextension mapfile{+cmextra.map}
\pdfextension mapfile{+symbols.map}
\pdfextension mapfile{+euler.map}
\pdfextension mapfile{+latxfont.map}
\pdfextension mapfile{+rsfs.map}
\usepackage[margin=25mm,headheight=15pt]{geometry}
\usepackage{fontspec}
\setmainfont{Latin Modern Roman}
\usepackage{amsmath,amssymb,amsthm,amscd,mathtools,mathrsfs}
\usepackage{microtype,booktabs,longtable,array,xcolor,graphicx,enumitem,needspace,fancyhdr}
\usepackage[round,authoryear]{natbib}
\usepackage[unicode,colorlinks=true,linkcolor=blue!40!black,citecolor=blue!40!black,urlcolor=blue!40!black]{hyperref}
\usepackage{xurl}
\usepackage{listings}
\usepackage{tabularx}
\lstset{basicstyle=\ttfamily\small,breaklines=true,columns=fullflexible,keepspaces=true}
\usepackage{tikz}
\usetikzlibrary{arrows.meta,calc,patterns,positioning}
\usepackage{aliascnt}
\usepackage[nameinlink,noabbrev]{cleveref}
\allowdisplaybreaks[2]
\setcounter{MaxMatrixCols}{32}
\setlength{\emergencystretch}{3em}
\setlength{\parskip}{2pt}
\newtheorem{theorem}{Theorem}[section]
\newaliascnt{proposition}{theorem}
\newtheorem{proposition}[proposition]{Proposition}
\aliascntresetthe{proposition}
\crefname{proposition}{proposition}{propositions}
\newaliascnt{lemma}{theorem}
\newtheorem{lemma}[lemma]{Lemma}
\aliascntresetthe{lemma}
\crefname{lemma}{lemma}{lemmas}
\newaliascnt{corollary}{theorem}
\newtheorem{corollary}[corollary]{Corollary}
\aliascntresetthe{corollary}
\crefname{corollary}{corollary}{corollarys}
\newaliascnt{claim}{theorem}
\newtheorem{claim}[claim]{Claim}
\aliascntresetthe{claim}
\crefname{claim}{claim}{claims}
\newaliascnt{statement}{theorem}
\newtheorem{statement}[statement]{Statement}
\aliascntresetthe{statement}
\crefname{statement}{statement}{statements}
\theoremstyle{definition}
\newaliascnt{definition}{theorem}
\newtheorem{definition}[definition]{Definition}
\aliascntresetthe{definition}
\crefname{definition}{definition}{definitions}
\newaliascnt{remark}{theorem}
\newtheorem{remark}[remark]{Remark}
\aliascntresetthe{remark}
\crefname{remark}{remark}{remarks}
\newaliascnt{example}{theorem}
\newtheorem{example}[example]{Example}
\aliascntresetthe{example}
\crefname{example}{example}{examples}
\newaliascnt{problem}{theorem}
\newtheorem{problem}[problem]{Open problem}
\aliascntresetthe{problem}
\crefname{problem}{open problem}{open problems}
\newaliascnt{question}{theorem}
\newtheorem{question}[question]{Question}
\aliascntresetthe{question}
\crefname{question}{question}{questions}
\newaliascnt{auditrecord}{theorem}
\newtheorem{auditrecord}[auditrecord]{Verification record}
\aliascntresetthe{auditrecord}
\crefname{auditrecord}{verification record}{verification records}
\newaliascnt{test}{theorem}
\newtheorem{test}[test]{Verification}
\aliascntresetthe{test}
\crefname{test}{verification}{verifications}
\newaliascnt{conjecture}{theorem}
\newtheorem{conjecture}[conjecture]{Conjecture}
\aliascntresetthe{conjecture}
\crefname{conjecture}{conjecture}{conjectures}
\newcommand{\C}{\mathbb C}\newcommand{\R}{\mathbb R}\newcommand{\Z}{\mathbb Z}\newcommand{\Q}{\mathbb Q}\newcommand{\N}{\mathbb N}\newcommand{\F}{\mathbb F}
\newcommand{\CC}{\mathbb C}\newcommand{\RR}{\mathbb R}\newcommand{\ZZ}{\mathbb Z}\newcommand{\QQ}{\mathbb Q}\newcommand{\NN}{\mathbb N}
\newcommand{\PP}{\mathbb P}\newcommand{\Pp}{\mathbb P}\newcommand{\Pone}{\mathbb P^1}\newcommand{\HH}{\mathbb H}
\newcommand{\RP}{\mathbb{RP}}\newcommand{\CP}{\mathbb{CP}}
\newcommand{\OO}{\mathcal O}\newcommand{\cO}{\mathcal O}\newcommand{\cM}{\mathcal M}\newcommand{\cF}{\mathcal F}
\newcommand{\h}{\mathfrak h}\newcommand{\hz}{\mathfrak h_z}\newcommand{\Acon}{\mathfrak A}\newcommand{\wt}[1]{\widehat{#1}}
\newcommand{\red}{\mathrm{red}}\newcommand{\tor}{\mathrm{tor}}\newcommand{\dP}{\mathrm{dP}}
\newcommand{\gh}{\widehat\gamma}\newcommand{\uh}{\widehat u}\newcommand{\wh}{\widehat w}\newcommand{\dhat}{\widehat\delta}\renewcommand{\dh}{\widehat\delta}
\newcommand{\dd}{\mathrm d}\newcommand{\ii}{\mathrm i}\newcommand{\eps}{\varepsilon}\newcommand{\partt}{\mathrm{part}}\newcommand{\Bcirc}{B^\circ}\newcommand{\ddbar}{\partial\bar\partial}
\DeclareMathOperator{\Pic}{Pic}\DeclareMathOperator{\Alb}{Alb}\DeclareMathOperator{\Aut}{Aut}\DeclareMathOperator{\NS}{NS}
\DeclareMathOperator{\Hom}{Hom}\DeclareMathOperator{\End}{End}\DeclareMathOperator{\Ext}{Ext}\DeclareMathOperator{\Tor}{Tor}
\DeclareMathOperator{\Spec}{Spec}\DeclareMathOperator{\Proj}{Proj}\DeclareMathOperator{\Sym}{Sym}\DeclareMathOperator{\Diff}{Diff}
\DeclareMathOperator{\GL}{GL}\DeclareMathOperator{\SL}{SL}\DeclareMathOperator{\PSL}{PSL}\DeclareMathOperator{\Sp}{Sp}
\DeclareMathOperator{\im}{im}\DeclareMathOperator{\coker}{coker}\DeclareMathOperator{\rank}{rank}\DeclareMathOperator{\diag}{diag}
\DeclareMathOperator{\ord}{ord}\DeclareMathOperator{\vol}{vol}\DeclareMathOperator{\pr}{pr}\DeclareMathOperator{\id}{id}
\DeclareMathOperator{\divisor}{div}\DeclareMathOperator{\Tr}{Tr}\DeclareMathOperator{\Nm}{Nm}\DeclareMathOperator{\Sha}{Sha}
\DeclareMathOperator{\Imn}{Im}\DeclareMathOperator{\Ren}{Re}\DeclareMathOperator{\trdeg}{tr.deg}
\newcommand{\Uses}[1]{\par\smallskip\noindent\textit{Uses:} #1.\par\smallskip}
\newcommand{\Verified}{}\newcommand{\AuthorSelfCheck}{}\newcommand{\NeedsCheck}{}\newcommand{\Open}{\textbf{Open.}\ }

\newcommand{\StatementRecord}[4]{\par\smallskip\noindent{\upshape\textbf{Label:} \nolinkurl{#1}.\par\noindent\textbf{Status:} #2.\par\noindent\textbf{Check record:} #3.\par\noindent\textbf{Uses:} #4.}\par\smallskip}
\newcommand{\RungZero}[1]{\par\noindent\textbf{Result and mechanism.} #1\par\smallskip}
\newcommand{\RungOne}{\par\noindent\textbf{Proof at the first rung.}\par\smallskip}
\newcommand{\CurrentChecker}{GPT-6 Astra (OpenAI, ChatGPT Pro mode, effort Pro), 18 November 2026}
\newaliascnt{cited}{theorem}
\newtheorem{cited}[cited]{Cited result}
\aliascntresetthe{cited}
\crefname{cited}{cited result}{cited results}
\newaliascnt{refuted}{theorem}
\newtheorem{refuted}[refuted]{Correction}
\aliascntresetthe{refuted}
\crefname{refuted}{correction}{corrections}
\newaliascnt{computed}{theorem}
\newtheorem{computed}[computed]{Computation}
\aliascntresetthe{computed}
\crefname{computed}{computation}{computations}
\newaliascnt{openq}{theorem}
\newtheorem{openq}[openq]{Open question}
\aliascntresetthe{openq}
\crefname{openq}{open question}{open questions}

\providecommand{\rk}{\operatorname{rk}}
\providecommand{\Ann}{\operatorname{Ann}}
\providecommand{\Log}{\operatorname{Log}}
\providecommand{\vir}{\mathrm{vir}}

\begin{document}
\section{Smooth returns with their physical circle and cusp fan}
\label{sec:IV-results}

We identify $\Theta _7$ with $\mathbb Z/28$ by
$[\Sigma]\mapsto28\mu(\Sigma)$, with outward boundary orientation
and positive generator $\mu=1/28$.  The return is formed from the
specified positive parameter loop and its actual endpoint map.
The physical circle is the circle in the given punctured line bundle
$M^\times$.  A cusp fan includes its lattice basis and the marked
collar identifications.  These choices remain in every formula.

\subsection{The physical transfer and the reference class}
\RungZero{The reference $(3k,4l)$ family has return zero on its
reference fan for every coprime pair.  Its two characteristic terms
cancel on the actual page, and the physical-line comparison fixes
the sign of the nonzero values in the other families.}
\RungOne
\begin{enumerate}[label=\arabic*.]
\item The fixed curve has the specified page surface; in the
  $(3,4)$ case this is the elliptic surface over
  $b^3=a^4+1$, with twelve nodal fibres.  The zero section and the
  specified section remain disjoint, giving the relative terms
  $p_{\rm rel}=-24$ and $s_\delta=-2$
  (\cref{push:src:C.3.24}).
\item The characteristic formula for its six-dimensional return
  is $(p_{\rm rel}-12s_\delta)/24$, so this reference value is zero.
  The spin-c coboundary and the explicitly surgered free-circle
  quotient give the same cancellation, including the coprime
  $(3k,4l)$ reference families
  (\cref{push:src:C.3.27,maps:src:B.6.24}).
\item The meridian-one Thom line is $U=M^{-1}$.  The Hermitian
  map $D_M(v)(w)=\langle v,w\rangle/\langle v,v\rangle$ satisfies
  $D_M(e^{i\theta}v)=e^{-i\theta}D_M(v)$, so $x_U=-x_M$.
  The quaternionic map $(a,b)\mapsto(\bar b,-\bar a)$ intertwines
  the two $SU(2)$ cap systems and preserves their real
  four-dimensional orientation (\cref{night:P3K.68}).
\item With the actual order of the two cap products, the physical
  Hopf transfer is $-\Sigma_{F^{-1}}$ and the dual transfer is
  $\Sigma_{F^{-1}}$.  Opposite collar cancellation gives
  $\Sigma_{F^{-1}}=-\Sigma_F$.  Thus the exact line-and-cap
  comparison gives $A_M=-A_U$, and for the page data $(H,\chi)$
  it gives
  \[
  (k_M,\ell_M)=(3H,24\chi),\qquad
  A_M=(4k_M-\ell_M)/24=H/2-\chi
  \]
  (\cref{night:P3K.68,night:P3K.69}).
\end{enumerate}

\subsection{The actual fan algorithm}
\RungZero{Every lattice-basis fan in the specified scalar families
has an exact integer calculation of its return.  The calculation
starts from a vertical fan, transports the complete source, line,
unit and cap maps through an Euclidean fan word, and solves the
target-prism equations after each flop.  The scalar polynomial
$KL(S-1)(S-3)/8$ from the earlier plan is false.}
\RungOne
\begin{enumerate}[label=\arabic*.]
\item The fan normalization is effected by the displayed lattice
  and toric coordinate maps.  Euclidean subtraction expresses each
  marked fan as a vertical seed
  $T_n=\left(\begin{smallmatrix}1&0\\n&1\end{smallmatrix}\right)$
  followed by the two explicit generators
  $Q_0=\left(\begin{smallmatrix}1&-1\\0&1\end{smallmatrix}\right)$
  and $Q_1=\left(\begin{smallmatrix}1&0\\-1&1\end{smallmatrix}\right)$
  (\cref{night:P3K.71}).
\item For a fixed-source scalar loop $\lambda=r^J$, the literal
  maps $Z_n=Z$, $V_n=Z^{-n}V$, $h_n=hB^n$, $B_n=B$ identify the
  full vertical source and cap geometry.  Its physical integer is
  \[
  A_{M,n}=H/2-\chi-Jn(n-1)/2,
  \]
  with free characteristic pair
  $(3H-3Jn(n-1),24\chi)$
  (\cref{night:P3K.72}).
\item At each actual fan step, the doubled ordinary-product
  matrix $V$ and vector $\delta$ obey the stated target-prism
  equations.  Prescribing the selected face entries
  $V'_{22}=V_{ii}$ and
  $\delta'_2=2V_{ii}-\delta_i$ gives their unique rational
  solution.  The transported geometric classes supply the integral
  solution, proving integrality of every displayed update
  (\cref{night:P3K.73}).
\item The native free mapping-class action is
  $(m,n,s)\mapsto(m+2n,-n,s+n)$.
  It makes the single-loop physical path sum change by
  $(\delta_i-V_{ii})/4$ at the corresponding step.  The vertical
  seed and these updates therefore compute the complete return,
  including the explicitly constructed shorter loops
  (\cref{night:P3K.73}).
\item For the original shorter $(3,4)$ loop, $J=2$ and the
  reference integer is zero.  At $T_2$ the actual value is $-2=26$
  modulo 28.  The old polynomial gives zero there because
  $T_2^{-1}e_2=e_2$.  This is an exact counterexample in the same
  marked family.  The corrected algorithm supplies fan witnesses
  for all 28 classes, including
  $\left(\begin{smallmatrix}1&-2\\0&1\end{smallmatrix}\right)$
  with class three (\cref{night:P3K.74,night:P3K.73}).
\end{enumerate}

\subsection{The fixed-edge flop and the generators}
\RungZero{In the specified fixed-edge comparison the flop sends the
physical index $A_M$ to $7A_M$ modulo 28.  The OS45 reference family
has class three and hence its positive covers realize every class;
the multiplier and sign follow from the same physical transfer and
the unchanged native free generators.}
\RungOne
\begin{enumerate}[label=\arabic*.]
\item For the page the native free pair is
  $(m,n)=(6\chi,3H-6\chi)$.  In the physical orientation
  $n=6A_M$ (\cref{night:P3K.68,night:P3K.69}).
\item The actual two-ruling Cox map conjugates the native
  coordinates by $(m,n,s)\mapsto(m+2n,-n,s+n)$.
  The line dualization does not replace this map: its explicit
  quaternionic intertwiner preserves the native real generators
  (\cref{night:P3K.68}).
\item The central increment is consequently $n=6A_M$; adding it
  to the initial physical class gives $7A_M$.  This proves the
  corrected fixed-edge rule and refutes the former $-5A$ rule
  in its stated physical interpretation
  (\cref{night:P3K.69}).
\item For OS45, $M=\mathcal O(8G-O)$ and $Q=G$ give reference
  physical class three by the constructed page and caps.  A
  degree-$d$ positive parameter cover composes the actual return
  $d$ times and has class $3d$.  Since three is a unit modulo 28,
  these covers realize every class
  (\cref{night:P3K.69}).
\item The fixed-edge flop of that generator has class 21.  The
  degree-four weighted OS47 example has physical class four.
  These values use the same outward-boundary and positive-parameter
  conventions as the reference calculation
  (\cref{night:P3K.69}).
\end{enumerate}

\subsection{The families and their complete marked scopes}
\RungZero{The completed one-cap, two-cap and selected-cusp families
have the characteristic pairs and all-fan algorithms given below.
Each value is transported through an explicit source, line and
endpoint map; the remaining source-moduli attachments are stated
separately in \cref{night:P3K.75}.}
\RungOne
\begin{enumerate}[label=\arabic*.]
\item The Kodaira component equations, root substitutions and
  affine norm conditions construct the actual finite caps and
  determine the complete retained character and winding data.
  The ordinary annular line cochains identify them with the
  marked regular-fibre lattice; their full units are retained
  (the cap statements and \cref{night:P3K.68,night:P3K.69}).
\item The ordinary-product characteristic calculation gives the
  physical pair $(k_M,\ell_M)$ on a seed fan.  The appropriate
  Euclidean fan word and the uniquely solved target-prism equations
  transport it to every lattice-basis fan
  (\cref{night:P3K.71,night:P3K.73,night:P3K.77,night:P3K.87,night:P3K.90}).
\item For the component-supported one-cap OS49 family the actual
  shortest parameter is $w=r^3$.  Its positive reference fan has
  class zero and its indicated negative fan has class $-k$;
  the whole $r$ loop has class $-3k$.  The source-line-cap endpoint
  map is constructed before the degree-three characteristic
  pullback is used (\cref{night:P3K.65,night:P3K.67,night:P3K.76,night:P3K.77}).
\item For the two-cap OS49 family, $m=3k$, $h=6l+1$ and the full
  ordinary clutch obey
  $D=hn_A+mn_B+mh\sigma=\pm1$.
  The shortest loop $w=r^3$, $\lambda=w^{kh}$ has positive-fan
  pair $(2kh,8k(3l+1))$ and integer $kl$.
  The actual endpoint group and all admitted windings are computed
  by the complete cap and annular maps
  (\cref{night:P3K.84,night:P3K.85,night:P3K.87,night:P3K.91}).
\item The source-zero E7 branches
  $(a,b,m)=(1,2,4p)$ and $(2,1,4p-1)$ have
  $N=km$, $\mu=2-a+6pk$ and vertical class
  \[
  A_j=2b\binom{j+1}{3}+(3-b)\binom{j+1}{2}
      +k(a-1)(3p-1)-km\binom j2.
  \]
  Their pair is $(\mu+6A_j,4\mu)$; the source, Cartier, root,
  normal and inverse-closing arrows prove this seed and its
  all-fan transport.  Each fixed-scalar branch has explicit fan
  witnesses for all 28 classes
  (\cref{night:P3K.86,night:P3K.88,night:P3K.89,night:P3K.90}).
\end{enumerate}

\begin{center}
\small
\begin{tabular}{@{}p{.23\linewidth}p{.29\linewidth}p{.20\linewidth}p{.18\linewidth}@{}}
\toprule
Family and marking & Construction & Physical class & Positive covers\\
\midrule
Coprime $(3k,4l)$, reference fan & The specified section page and caps
  & $0$ & $0$\\
Original $(3,4)$, all basis fans & The seed and prism algorithm
  & Every element occurs & Multiples of the selected value\\
OS45, reference fan & $M=\mathcal O(8G-O)$, $Q=G$
  & $3$ & All of $\mathbb Z/28$\\
OS45, specified fixed-edge flop & The two-ruling Cox map
  & $21$ & $\{0,7,14,21\}$\\
Weighted OS47 example & The specified degree-four page
  & $4$ & $4\mathbb Z/28$\\
OS49, one cap, shortest loop & $w=r^3$, indicated negative fan
  & $-k$ & $k\mathbb Z/28$\\
OS49, two caps, positive fan & $m=3k$, $h=6l+1$, $\lambda=w^{kh}$
  & $kl$ & $kl\mathbb Z/28$\\
E7 source-zero, fixed scalar & Each of the two completed branches
  & Every element occurs over fans & Multiples of the selected value\\
\bottomrule
\end{tabular}
\end{center}

\subsection{The source-infinity formula and the completed E6 loops}
\RungZero{The complete E7 four-boundary groupoid gives a closed
source-infinity formula on every fan.  The E6 source, Cartier and
root maps identify its completed selected-cusp loops with the same
selected-cusp algorithm, and extend its specified unselected and
native-infinity loops over proper discs, giving return zero there.}
\RungOne
\begin{enumerate}[label=\arabic*.]
\item The full branch reflection and infinity rescaling, including
  the root endpoint and its inverse scalar path on the caps, give
  \[
  I_a(P)=-Z_a(P)-Z_a(JP)-C(P),\qquad
  J=\operatorname{diag}(-1,1).
  \]
  This is the relation between the actual four marked boundary
  loops (\cref{night:P3K.92}).
\item For $(x,y)^t=P^{-1}e_2$ put $S=x^2+xy+y^2$.
  Combining the two reflected source-zero prism systems with
  twice the selected-cusp system gives winding $-8I$.
  Its unique solution and the common first-step seed yield
  \[
  I_a(P)=1-a-3(a-1)(S-1)
       -\frac{2a+1}{4}(S-1)(S-3)\pmod {28}.
  \]
  Both fan-generator identities are proved by substitution into
  the full prism equations (\cref{night:P3K.93}).
\item A primitive pair makes $S$ odd.  Writing $S=2t+1$ gives
  $-3t(t-1)$ for $a=1$ and $-5t^2-t-1$ for $a=2$.
  Their period modulo 28 and the displayed fan witnesses give
  the respective complete images
  $\{0,6,8,10,14,20,22,24\}$ and
  $\{5,7,9,13,19,21,23,27\}$
  (\cref{night:P3K.93}).
\item For $y^2=x^3+rx^2+tx+t^2$, $G=(0,t)$ and
  $M_6=\mathcal O(6G-O)$, the actual selected-cusp map on the
  ordered $w$ cover has
  $r-1/3=-w^4/[3(w^2-3)^2]$.
  Its full cap comparison gives vertical pair
  $(-6-18j(j-1),0)$ and class $-1-3j(j-1)$, and identifies
  every fan step with \cref{night:P3K.71}
  (\cref{night:P3K.95,night:P3K.98}).
\item The full cap progressions
  $(a,b;m)=(6,1;3l+2)$, $l\geq0$, and
  $(2,3;3l)$, $l\geq1$, extend at the unselected $r=0$
  collision and at each native infinity point $w^2=3$.
  The constructed proper disc families have identity closing,
  so their marked positive $w$ returns vanish.  At infinity
  $1/r=(w^2-3)^2/(3-2w^2)$ has degree two; this exact map is
  retained in the statement of the return
  (\cref{night:P3K.94,night:P3K.96,night:P3K.97}).
\end{enumerate}

% Additional notation transported from the night manuscript.
\providecommand{\LL}{\mathcal L}
\providecommand{\cT}{\mathcal T}
\providecommand{\HP}{\mathbb H\mathbb P}
\providecommand{\tr}{\operatorname{tr}}
\providecommand{\sgn}{\operatorname{sgn}}
\providecommand{\vir}{\mathrm{vir}}

\section{The page and the reference return}



\begin{proposition}[The actual positive-loop clutching map]\label{fam:src:E.2.2}
\StatementRecord{fam:src:E.2.2}{proved}{GPT-6 Astra (OpenAI, ChatGPT Pro mode), effort Pro, 17 November 2026; Claude Fable 5.2 (Anthropic), effort high, 2026-11-12; Claude Fable 5.2 (Anthropic), effort high, 12 November 2026}{the hypotheses, notation and definitions stated here; \cite[Theorem~0 and Corollaries~1--2, p.~8]{Cerf1970}; \cite[pp.~393--394]{Smale61}; \cite[Section~6, equation~(11), pp.~101--103]{EellsKuiper1962}}

For the retained $S^6$ family over the punctured $w$-disc, the
monodromy acts as the identity on $H^*(S^6;\Z)$. The remaining
smooth monodromy is represented explicitly as follows. Set
$c_0=-i(M+1)$ and $c(r)=c_0-2r$, $0\leq r\leq1$. Let
$H:X_{c_0}\to X_{c_0-2}$ be parallel transport in the proper smooth
family over this interval. With the coherent isomorphism from A.3,
\[
       \Phi=F_{2,c_0-2}\circ H:X_{c_0}\longrightarrow X_{c_0}  \tag{E9}\label{famIII:eq:E9}
\]
represents the monodromy of the positively oriented $w$-circle.
Its isotopy class is independent of the smooth connection. The
associated oriented homotopy sphere
$\Sigma_\Phi=D^7\cup_\Phi(-D^7)$ reduces the mapping-class
calculation to the explicit residue
\[
  [\Phi]=28\mu(\Sigma_\Phi)
    =\frac{\langle p_1(W)^2,[W,\partial W]\rangle
                 -4\sigma(W)}{32}\pmod {28},                 \tag{E10}\label{famIII:eq:E10}
\]
where $W$ is a spin eight-manifold with boundary $\Sigma_\Phi$ and
the Pontryagin class is lifted to relative rational cohomology.
The identification of $\Z/28$ is oriented by a sphere with
$\mu=1/28$.
\end{proposition}


\begin{proof}
\begin{enumerate}[label=\arabic*.]
\item The proper submersion has a smooth horizontal distribution; integration
along a compact path exists because the fibres are compact. The
complex orientation varies continuously. Its monodromy therefore
preserves the orientation class in $H^6$ and fixes the constant
class in $H^0$. These are the only nonzero integral cohomology
groups, proving the first assertion with integral coefficients.

\item Since $e^{-\pi i(c_0-2r)}=e^{-\pi ic_0}e^{2\pi ir}$, the path
$c(r)$ is the lift of the positive circle. The descent gluing at
its endpoint is exactly $F_{2,c_0-2}$, proving \eqref{famIII:eq:E9}. A path of
horizontal distributions gives a smooth path of the resulting
monodromy diffeomorphisms, so their isotopy classes agree. Changing
the coherent descent maps by a member of the connected automorphism
group also leaves this class unchanged. The diffeomorphism
$X_{c_0}\simeq S^6$ is chosen with the complex orientation.

\item The glued manifold $\Sigma_\Phi$ is simply connected by van Kampen
and has the integral homology of $S^7$ by Mayer--Vietoris; the
Hurewicz and Whitehead arguments give a homotopy sphere. The
gluing comparison $\pi_0\operatorname{Diff}^+(S^6)=\Theta_7$
sends $\Phi$ to this construction. Here we use Cerf's sphere-gluing comparison in dimension six \cite[Theorem~0 and Corollaries~1--2, p.~8]{Cerf1970}: isotopy classes
are the group of sphere gluing maps modulo extension over a disc.
The extension equivalence is exactly the equivalence of their
two-disc gluings. In dimension seven, the simply connected
h-cobordism comparison identifies these classes with oriented
homotopy spheres; the comparison with oriented diffeomorphism
classes is Smale's Theorem~J \cite[pp.~393--394]{Smale61}.
Thus the map being evaluated is the actual boundary gluing of \eqref{famIII:eq:E9}.
Eells--Kuiper's
seven-dimensional invariant is
$(p_1^2-4\sigma)/896$ modulo one and distinguishes the twenty-eight
oriented homotopy-seven-sphere classes
\cite[Section~6, equation~(11), pp.~101--103]{EellsKuiper1962}. Multiplying that formula by
twenty-eight gives \eqref{famIII:eq:E10}. Its hypotheses hold here: a homotopy
seven-sphere is spin and has vanishing $H^3(-;\Q)$ and
$H^4(-;\Q)$, so the relative lift is unique. The same theorem
and its cobordism construction supply a spin coboundary. Thus
\eqref{famIII:eq:E10} is an invariant of the specific gluing \eqref{famIII:eq:E9}, rather than an
inference from its action on cohomology.

\end{enumerate}
\end{proof}




\begin{remark}[Resolution of the original compact-filling and mapping-class questions]\label{fam:src:E.2.3}
\StatementRecord{fam:src:E.2.3}{proved}{GPT-6 Astra (OpenAI, ChatGPT Pro mode), effort Pro, 17 November 2026; Claude Fable 5.2 (Anthropic), effort high, 2026-11-12; Claude Fable 5.2 (Anthropic), effort high, 2026-11-17; Claude Fable 5.2 (Anthropic), effort high, 12 November 2026}{Definitions and the displayed constructions}

The original question asked for a flat proper filling of the compact $S^6$ family over $w=0$. The specific equivariant compactification is excluded by Paper~I, Paper~I, Theorem~21.58. In particular no global central fibre is named
whose Euler number, fundamental group, or integral homology could
be asserted. Paper~I, Proposition~19.6 constructs and computes the two reduced-surface
degenerations after the stated additional base change; it does not
construct a neighbourhood of the nonreduced fibres $3S_1$ and
$4S_2$ in a compact limit.

The finite-place test in that question was the extension of the
order-three affine deck action over a two-parameter plumbing in
$(t,s)$, with prescribed normal character and central multiplicity
three, followed by descent from $s$ to $w=s^6$. The order-four
analogue must retain the normal character of order four. At infinity
the prescribed periodic gluing was the one specified in Paper~I, Remark~19.5, now
compatible with these finite-place quotients. The computed central
surface homology groups cannot be substituted for the homology of
that unconstructed threefold.

For the mapping-class part, E.2.2 gives the exact clutching map and
its characteristic-number reduction. The complete characteristic and boundary calculations in \cref{push:src:C.3.24} and \cref{push:src:C.3.27} give zero for this particular $\Sigma_\Phi$. The numeric instance to be evaluated was the residue \eqref{famIII:eq:E10} for the single loop based at
$c_0=-i(M+1)$; all other basepoints in the connected parameter disc
give conjugate monodromy, and $\Theta_7$ is abelian. The identity on cohomology alone does not evaluate that residue; the complete cap calculations do.
\end{remark}


The actual positive-loop smooth monodromy is zero. \ref{push:src:C.3.24} proves
this through an equivariant orbit-knot integer; \ref{push:src:C.3.25}--\ref{push:src:C.3.27}
give a second determination through an explicit spin$^c$
eight-dimensional coboundary and direct index additivity. Their
boundary-invariant and one-loop arguments are independent; their
actual annular representative and characteristic inputs are shared
and identified below. Paper~I, Theorem~21.58 excludes the specified descended analytic
filling, so \ref{push:src:C.5.1} records that its requested Chern
data are undefined. In this section $M_\Phi$ is a closed seven-manifold,
whereas $W_8$ denotes a
compact oriented eight-manifold with boundary. The circle parameter
has positive orientation. On $M_\Phi$ we use the orientation
$d\theta\wedge o_{X_c}$; this is the boundary orientation of a
holomorphic family over a positively oriented disc. The statements
below do not assign a characteristic number to that excluded
analytic filling.



The credited statement and complete re-derivation are in \cref{push:src:C.1.1}.
The credited statement and complete re-derivation are in \cref{push:src:C.1.2}.




The credited statement and complete re-derivation are in \cref{push:src:C.2.1}.

\begin{proposition}[the untwisted intersection form]\label{push:src:C.2.2}
\StatementRecord{push:src:C.2.2}{proved}{GPT-6 Astra (OpenAI, ChatGPT Pro mode), effort Pro, 2026-11-17; Claude Fable 5.2 (Anthropic), effort high, 16 November 2026}{the hypotheses, notation and definitions stated here; the explicit constructions and calculations in the proof}

Let $W_A$ be a sufficiently small compact tubular representative of the
smooth total space $\mathcal X^+$ of Paper~I, Theorem~20.4 over a closed $u$-disc.
Its boundary is the mapping torus of Engel's \emph{untwisted} family
$X(u)$. It is not the six-sphere mapping torus $M_\Phi$.
Write $\mathcal L=H^2(S;\Z)$, $F=c_1\mathcal O_S(S_\infty)$,
$m=c_1(L)$, and $v=6P-O-3F$ in the conventions of Paper~I, Proposition~20.11.
Let $H$ be the infinity-section class on $R$ and $A$ the exceptional
Cartier divisor pulled back to $R^+$. Then
\[
 H_4(W_A;\Z)=
 \frac{\mathcal L\oplus\Z H\oplus\Z A}{\Z(m-A)}.
\]
For $\alpha=(a,h,b)\in H^2(R^+;\Z)$ set
\[
 d(\alpha)=(I-T^*)a+h(m-F)+bF.
\]
The middle intersection form of $W_A$ is
\[
 B(\alpha,\beta)=-d(\alpha)\cdot d(\beta).
\]
Its rank is two, its signature is $+2$, and its two nonzero Smith
factors are $1,11$. On the image modulo its radical it is represented,
in the ordered basis $(v,m)$, by
\[
 \begin{pmatrix}6&1\\1&2\end{pmatrix}.
\]
\end{proposition}


\begin{proof}
\begin{enumerate}[label=\arabic*.]
\item The local semistable charts $xy=u$ give a tubular representative which
retracts to $Q^+$; its boundary is the nearby-fibre mapping torus. One
may equivalently choose a regular neighbourhood of $Q^+$ in the smooth
total space and round its boundary. The fibrewise plumbing identifies
that boundary with the inverse image of a small positive $u$-circle.
The normalization mapping cylinder and $H_3(S)=0$ identify its fourth
homology with the cokernel of
$H_4(S)\to H_4(R^+)$. In the divisor coordinates of Paper~I, Proposition~20.12 the two
sections are $D'_0=H+m-A$ and $D_\infty=H$. Their difference is
$m-A$, proving the group presentation.

\item The normalization is an immersed complex hypersurface in $W_A$.
Its normal line has class $-D'_0-D_\infty$, because the whole central
divisor is the divisor of the base coordinate $u$. Put
$r_0\alpha=a+bF$ and $r_\infty\alpha=a-hm$; these are the
restrictions to the two sections. The self-intersection of two
four-cycles pushed forward from the normalization contains the normal
Euler term and the intersections between distinct local branches.
All transverse complex branch intersections have positive sign.
Consequently it is
\begin{align*}
 &-r_0\alpha\cdot r_0\beta-r_\infty\alpha\cdot r_\infty\beta\\
 &\quad+r_0\alpha\cdot T^*r_\infty\beta
       +T^*r_\infty\alpha\cdot r_0\beta\\
 &=-\bigl(r_0\alpha-T^*r_\infty\alpha\bigr)
       \cdot\bigl(r_0\beta-T^*r_\infty\beta\bigr).
\end{align*}
This is an intersection-of-cycles calculation in the actual smooth
fourfold, rather than an assignment of a form to an unconstructed
central fibre. The isometry property of $T$ and $T^*m=m-F$ give
the stated expression for $d$. Also $d(m-A)=0$, so it descends.

\item The translation formula in Paper~I, Proposition~20.11 shows that the image of $d$ is
$\langle F,v,m\rangle$. The surjectivity of the three intersection
coordinates proved there shows equality integrally. Its radical is
$\Z F$, since $F^2=Fv=Fm=0$ and
$v^2=-6$, $m^2=-2$, $vm=-1$. The remaining form is the displayed
positive matrix, of determinant eleven and with an entry equal to
one. Integral row and column reduction gives Smith factors $1,11$.
These computations, including the full eleven-by-eleven form, are
independently checked in \path{code/e_push_C_attempts.py}.

\end{enumerate}
\end{proof}




\begin{proposition}[a relative Pontryagin number for $W_A$]\label{push:src:C.2.3}
\StatementRecord{push:src:C.2.3}{proved}{GPT-6 Astra (OpenAI, ChatGPT Pro mode), effort Pro, 2026-11-17; Claude Fable 5.2 (Anthropic), effort high, 16 November 2026}{the hypotheses, notation and definitions stated here; \cref{push:src:C.2.2}; \cref{push:src:C.2.1}; \citet[Theorem~4.3, pp.~564--565]{AtiyahSinger1968IIIE}}

In the coordinates of \ref{push:src:C.2.2}, $p=p_1(TW_A)$ evaluates as
\[
 p(a,h,b)=-4(aF)-28h.
\]
It vanishes on the radical of the middle form and therefore has a
rational relative lift. The square of any such lift is
\[
 \bar p^{\,2}=\frac{4512}{11}.
\]
These numbers do not compute $\mu(M_\Phi)$. In particular the
canonical determinant class of $W_A$ restricts nontrivially to its
own boundary, so \ref{push:src:C.2.1} cannot be applied with that determinant.
\end{proposition}


\begin{proof}
\begin{enumerate}[label=\arabic*.]
\item Let $h:R^+\to R$ be the modification, $C$ its nodal elliptic-fibre
centre, and $c=c_1(R)=F+2H+m$. The centre is a regular embedding of
codimension two, also at its node: its ideal there is $(X,zv)$.
Its normal bundle has degree zero, and $c\cdot C=0$. The blow-up
and small resolution give
\[
 c_1(R^+)=h^*c-A,\qquad
 A\,h^*D\,h^*E=0,\quad
 A^2h^*D=-C\cdot D,\quad A^3=0.
\]
These are the regular-embedding blow-up identities. The subsequent
small resolution is crepant, and $A$ is the pullback of a Cartier
divisor, so they remain valid on $R^+$. The normal line to the
immersed normalization has class
$h^*N_0+A$, where $N_0=-(2H+m)$.

\item For completeness, the required second Chern pairings do not assume
that $R^+$ is K\"ahler. Analytic Hirzebruch--Riemann--Roch applies to
a compact complex manifold; see \citet[Theorem~4.3, pp.~564--565]{AtiyahSinger1968IIIE}.
The modification has $Rh_*\mathcal O_{R^+}=\mathcal O_R$.
For the blow-up this follows locally from its hypersurface description
in the projective line bundle and the vanishing of the direct images
of $\mathcal O_{\mathbb P^1}(-1)$; for the small resolution it follows
from the formal neighbourhood of its $(-1,-1)$ rational curve.
Thus $\chi(h^*D)=\chi(D)$ for every line bundle $D$ on $R$.
The linear coefficient in the Riemann--Roch polynomial for $kD$
and the preceding intersection identities give
\[
 c_2(R^+)h^*D=c_2(R)D+C\cdot D.
\]
Moreover $\chi(\mathcal O(A))=\chi(\mathcal O)=1$: on the
exceptional projective line bundle over $C$, the restriction is its
relative $\mathcal O(-1)$, which has zero derived direct image.
Pulling through the small resolution does not change this Euler
characteristic. Applying Riemann--Roch to $A$ yields $c_2(R^+)A=0$,
because $A^3$, $c_1(R^+)A^2$, and $c_1(R^+)^2A$ all vanish.

\item The real tangent bundle of $W_A$ restricted to the normalization is
the sum of its real tangent bundle and this normal line. Therefore
\[
 h^*p_1(R)+h^*N_0^2-4[C]
\]
has the same pairings on pulled-back divisors as $p_1(TW_A)$, and
the pairing with $A$ is zero. Here $[C]$ denotes its degree-four
cohomology class on $R$. Since $H^2=-mH$, we have
$(2H+m)^2=m^2=-2[\mathrm{pt}]$. The surface $S$ has
$p_1(S)=-24[\mathrm{pt}]$. The ruled-bundle tangent sequence thus
gives
\[
 p_1(R)+N_0^2=p_1(S)+2(2H+m)^2=-28[\mathrm{pt}],
\]
and $C\cdot(a+hH)=aF$. This proves the covector formula.

\item Modulo $\Z F$ the coordinate of $d(a,h,b)$ along $v$ is $aF$
and the coordinate along $m$ is $h$. Thus the covector on the
nondegenerate quotient is $(-4,-28)$. It annihilates the radical,
which is precisely the condition for a rational relative lift.
Its square is
\[
 (-4,-28)\frac1{11}
 \begin{pmatrix}2&-1\\-1&6\end{pmatrix}
 \binom{-4}{-28}=\frac{4512}{11}.
\]
Different relative lifts differ by a class pairing trivially with
absolute classes, so this value does not depend on the choice.

\item Finally the canonical first Chern class on the nearby untwisted
fibre is $f^*c_1\mathcal O_B(1)$. It is nonzero: the surjectivity of
$a\mapsto(aF,av,am)$ in Paper~I, Proposition~20.11 supplies an integral two-cycle with
evaluation one, and the Mayer--Vietoris relations in Paper~I, Proposition~20.12 do not
alter that evaluation. Hence no boundary trivialization of this
determinant exists. Even the formal expression
$(\bar p^{\,2}-4\sigma(W_A))/896=79/176$ is only an arithmetic
combination for this different boundary, not the invariant $\mu(M_\Phi)$ of the original mapping torus. A cobordism implementing the finite logarithmic transforms
and the cusp closing map would be required to transport the result.

\end{enumerate}
\end{proof}


The credited statement and complete re-derivation are in \cref{push:src:C.3.1}.

\begin{proposition}[a global scalar gauge would be trivial]\label{push:src:C.3.2}
\StatementRecord{push:src:C.3.2}{proved}{GPT-6 Astra (OpenAI, ChatGPT Pro mode), effort Pro, 2026-11-17; Claude Fable 5.2 (Anthropic), effort high, 16 November 2026}{the hypotheses, notation and definitions stated here; the explicit constructions and calculations in the proof}

Let a smooth circle act on a simply connected smooth manifold $Z$.
If $g:Z\to S^1$ is smooth and invariant under the circle, then
$z\mapsto g(z)\cdot z$ is smoothly isotopic to the identity.
The primitive coordinate $\theta$ detecting $\gamma$ on a regular
fibre of $X_c$ has no continuous extension to $X_c$.
\end{proposition}



\begin{proof}
\begin{enumerate}[label=\arabic*.]
\item Since $Z$ is simply connected, $g$ lifts to a smooth real function
$a$ with $g=e^{2\pi ia}$. Invariance implies
$a(\lambda z)-a(z)\in\Z$. This difference is continuous in the
connected circle and is zero at its identity, so it is zero
everywhere. The maps $z\mapsto e^{2\pi ita(z)}\cdot z$ are
diffeomorphisms with inverses obtained by replacing $t$ by $-t$.
They give the isotopy. In the present application a loop representing
$\widehat\gamma$ in a regular torus fibre has degree one under
$\theta$. If $\theta$ extended to the simply connected $S^6$,
that degree would be zero. The contradiction shows why the proposed local scalar twist is not a global scalar-gauge calculation.

\end{enumerate}
\end{proof}


The credited statement and complete re-derivation are in \cref{push:src:C.3.3}.

\begin{proposition}[information lost under the proposed base change]\label{push:src:C.3.4}
\StatementRecord{push:src:C.3.4}{proved}{GPT-6 Astra (OpenAI, ChatGPT Pro mode), effort Pro, 2026-11-17; GPT-6 Astra (OpenAI, ChatGPT Pro mode), S6\_06\_BLIND, effort Pro, 2026-11-17; Claude Fable 5.2 (Anthropic), effort high, 16 November 2026}{the hypotheses, notation and definitions stated here; the explicit constructions and calculations in the proof}

Under $w=r^6$ the smooth monodromy class becomes $6[\Phi]$.
Its value alone leaves two possibilities for $[\Phi]$ in $\Z/28$,
differing by 14. In particular a vanishing value after that base
change would not prove that $[\Phi]=0$.
\end{proposition}


\begin{proof}
\begin{enumerate}[label=\arabic*.]
\item The induced homomorphism on fundamental groups sends a positive
$r$-circle to six positive $w$-circles, so its transport is
$\Phi^6$ up to isotopy. The kernel of multiplication by six on
$\Z/28$ is $\{0,14\}$. The explicit preimages of all elements
of its image are checked in \path{code/c_exact.py}.

\end{enumerate}
\end{proof}


The complete coherent closing translation and its full chartwise proof are in Paper~I, Proposition~18.5.

\begin{proposition}[independence of the equivariant lift]\label{push:src:C.3.6}
\StatementRecord{push:src:C.3.6}{proved}{GPT-6 Astra (OpenAI, ChatGPT Pro mode), effort Pro, 2026-11-17; Claude Fable 5.2 (Anthropic), effort high, 16 November 2026}{the hypotheses, notation and definitions stated here; \cref{push:src:C.3.1}}

Let \(X=X_c\) carry the effective semifree circle action of \ref{push:src:C.3.1},
with fixed surface \(F\cong S^2\). If two smooth equivariant
diffeomorphisms of \(X\) induce the same map of its orbit space,
they are smoothly isotopic. More generally, a smooth isotopy of
orbit maps which is supported in a compact subset of the free
orbit stratum lifts to an equivariant isotopy. Thus, whenever the
annular orbit twist has been constructed with a smooth equivariant
lift, its class in \(\Theta_7\) is independent of that lift.
This assertion does not construct an isotopy of the annular twist.
\end{proposition}


\begin{proof}
\begin{enumerate}[label=\arabic*.]
\item Their quotient is an equivariant orbit-preserving diffeomorphism
\(h\). On \(X\setminus F\) it has a unique expression
\(h(x)=g(x)\cdot x\), with \(g\) a smooth invariant circle-valued
function. We check its extension across \(F\) directly. Choose an
equivariant tubular slice and trivialize its normal bundle locally
over \(F\). Conjugating the negative-weight complex coordinate
identifies the normal action with scalar multiplication on
\(\mathbb C^2\). Since \(h\) preserves each orbit, it fixes the
base variable \(y\) and has normal part
\[
 (f_1(y,z),f_2(y,z))=g(y,z)(z_1,z_2),\qquad z\ne0.
\]
In particular \(z_2f_1=z_1f_2\), an identity of smooth functions
also at zero. In the homogeneous Taylor expansion in the four real
normal variables, each pair \((p_1,p_2)\) therefore satisfies
\(z_2p_1=z_1p_2\). In the complexified polynomial ring with
variables \(z_1,z_2,\bar z_1,\bar z_2\), this implies
\(p_1=z_1q\) and \(p_2=z_2q\): first set \(z_1=0\) and divide
by \(z_1\), then cancel \(z_1\). This coefficient argument works
with coefficients depending smoothly on \(y\). It produces a
common formal Taylor series for \(g\).

\item The formal series can be realized by a smooth complex function
\(G(y,z)\). Here is the parametric Borel construction needed for
this assertion. On a relatively compact base slice, write its
homogeneous terms as \(q_n(y,z)\), choose a smooth cutoff equal to
one near zero, and replace each term by
\(\chi(z/\epsilon_n)q_n(y,z)\). Choose
\(\epsilon_n\downarrow0\) so rapidly that the term and all
mixed base/normal derivatives of total order at most \(n/2\)
are bounded by \(2^{-n}\), away from the finitely many lower
degrees. This is possible because each such derivative retains a
positive power of \(\epsilon_n\). The resulting series converges
in every smooth seminorm, and its Taylor series is the prescribed
one: the cutoff is constant near zero and higher homogeneous
terms have zero lower jets. This proves the required local
parametric extension. The remainders
\(f_j-Gz_j\) are flat in the normal variables, locally uniformly
with all derivatives in \(y\). For \(z\ne0\),
\[
 g-G=\frac{\bar z_1(f_1-Gz_1)+\bar z_2(f_2-Gz_2)}{|z|^2}.
\]
A smooth function flat along \(z=0\), divided by \(|z|^2\),
extends smoothly and flatly by zero: the usual quotient derivative
formula and Taylor bounds of arbitrarily high order prove this for
every derivative. Thus \(g\) extends smoothly across the slice.
Its absolute value is one by continuity. Equivariance implies
invariance on the free part, hence on the slice by continuity.
The local extensions agree because they agree on the dense free
stratum. We obtain a global smooth invariant \(g:X\to S^1\).

\item Since \(X\) is simply connected, the covering
\(\mathbb R\to S^1\), \(a\mapsto e^{2\pi ia}\), gives a smooth
lift \(a\). The integer-valued difference
\(a(\lambda x)-a(x)\) is zero at \(\lambda=1\) and continuous
in the connected circle, so \(a\) is invariant. The maps
\(x\mapsto e^{2\pi isa(x)}\cdot x\), \(0\leq s\leq1\),
are diffeomorphisms with the explicit inverse obtained by changing
\(s\) to \(-s\). They give the required isotopy.

\item For the second assertion the free quotient is a principal circle
bundle. A connection lifts the compactly supported time-dependent
vector field of the orbit isotopy to an invariant vector field.
Its flow exists for the full finite time interval, is equivariant,
and extends by the identity near \(F\). The endpoint differs from
any other smooth equivariant lift of the same orbit map by the
orbit-preserving diffeomorphism just treated.

\item The general orbit-preserving statement is Haefliger and Salem
(1991, Theorem~3.1). We used the direct slice argument because the
non-equivariant Lemma~3.2 in that paper, as originally printed,
requires correction: complex conjugation on \(\mathbb C\)
preserves scalar-circle orbits but its quotient
\(\bar z/z\) has no continuous extension to zero. The theorem
with equivariance is unaffected; the corrected argument is also
explained in Karshon's research slides, which reproduce the original
statement. No conclusion about smooth maps of the topological orbit
space near \(F\) is assumed here.
\nocite{HaefligerSalem1991D,KarshonOrbitSlidesD}

\end{enumerate}
\end{proof}




\begin{proposition}[the central cusp orbit fibre]\label{push:src:C.3.7}
\StatementRecord{push:src:C.3.7}{proved}{GPT-6 Astra (OpenAI, ChatGPT Pro mode), effort Pro, 2026-11-17; Claude Fable 5.2 (Anthropic), effort high, 16 November 2026}{the hypotheses, notation and definitions stated here; \cref{push:src:C.3.1}}

Let \(H=W/\rho\), and let \(F\cong S^2\) be the fixed double
curve. Then
\[
 \pi_1(H)=\mathbb Z^2,\qquad
 H_i(H;\mathbb Z)=
 \begin{cases}
 \mathbb Z&i=0,3,\\
 \mathbb Z^2&i=1,2,\\
 0&i>3.
 \end{cases}
\]
The class of \(F\) is nonzero and primitive in \(H_2(H;\mathbb Z)\).
In particular \(F\) does not bound a three-chain in the central
orbit fibre. This statement does not assert a deformation retraction
of the entire cusp neighbourhood onto \(H\), or compute the class
of the annular orbit twist.
\end{proposition}


\begin{proof}
\begin{enumerate}[label=\arabic*.]
\item Number the six toric boundary curves of the normalization
\(\widetilde W=\mathrm{dP}_6\) cyclically as
\(D_1,\ldots,D_6\), with \(D_3,D_6\) the fixed pair, and put
\(p_i=D_i\cap D_{i+1}\), with indices modulo six. Opposite sides
are identified by the maps already checked in Paper~I, Proposition~4.9 and Paper~I, Proposition~4.10 and Paper~I, Theorem~5.10 and \ref{push:src:C.3.1}.
Their endpoint rule is
\[
 p_{i-1}\longmapsto p_{i+3},\qquad
 p_i\longmapsto p_{i+2}\quad(D_i\longrightarrow D_{i+3}).
\]
It follows either from the two triangle classes of the periodic
fan or directly from these six endpoint equations that the vertex
classes are
\(A=\{p_2,p_4,p_6\}\) and \(B=\{p_1,p_3,p_5\}\).
The multiplicative units in the side maps preserve this endpoint
rule. The fixed curve is their resulting sphere, with endpoints
\(A,B\).

\item Put \(Y=\widetilde W/\rho\). This is an oriented compact
three-manifold whose boundary consists of the two fixed spheres.
Here are the local checks, including the two remaining vertices.
At a free orbit the quotient is an oriented three-dimensional
slice. At a fixed divisor the local model is
\(\mathbb C\times\mathbb C\), with the circle rotating only
the first factor, and its quotient is a half-space. The two toric
vertices not on \(D_3\cup D_6\) are isolated fixed points with
weights \((1,-1)\). Conjugating one coordinate gives the scalar
Hopf action on \(\mathbb C^2\), whose quotient is the cone on
\(\mathbb CP^1=S^2\), namely \(\mathbb R^3\). Thus they are
interior points, not additional boundary components.
All other stabilizers are trivial: the determinants of the circle
cocharacter with the four nonfixed edge rays are \(\pm1\).

\item The quotient is simply connected. Indeed local slices give path
lifts of orbit-space paths, and lifts on successive subintervals
can be joined by paths in their connected circle orbits. A lifted
loop closes after adding a path in its final orbit; its projection
is the original loop with a constant segment. Since
\(\widetilde W\) is simply connected, the resulting loop is
null-homotopic, and projecting the null-homotopy proves the claim.
Poincar\'e--Lefschetz duality and the boundary sequence now give
\[
 H_1(Y)=H_3(Y)=0,\qquad H_2(Y)=\mathbb Z.
\]
More explicitly, \(H_3(Y,\partial Y)=\mathbb Z\) maps to
\(H_2(\partial Y)=\mathbb Z^2\) by the sum of the two boundary
fundamental classes, and \(H_2(Y,\partial Y)=H^1(Y)=0\).
No classification theorem for simply connected three-manifolds
is needed for these conclusions.

\item Let \(\widetilde D\subset Y\) be the quotient of the six
boundary curves, and let \(D\subset H\) be its image. The space
\(\widetilde D\) is two spheres joined by two paths, each
subdivided into two edges. Hence
\[
 H_0(\widetilde D)=\mathbb Z,\quad
 H_1(\widetilde D)=\mathbb Z,\quad
 H_2(\widetilde D)=\mathbb Z^2.
\]
The image \(D\) is the sphere \(F\) together with two arcs
\(E_1,E_2\), each joining \(A\) to \(B\). Consequently
\[
 H_0(D)=\mathbb Z,\quad H_1(D)=\mathbb Z^2,\quad
 H_2(D)=\mathbb Z.
\]
Choose a path \(c\subset F\) from \(A\) to \(B\). The
endpoint rule gives the image of the generator of
\(\pi_1(\widetilde D)\) as
\[
 E_1E_2^{-1}cE_1^{-1}E_2c^{-1}.
\]
Writing \(a=E_1c^{-1}\), \(b=E_2c^{-1}\), this is
\(ab^{-1}a^{-1}b\). In particular the induced map on first
homology is zero.

\item We also fix the second-homology sign. Near a fixed divisor use a
complex normal coordinate \(re^{i\theta}\) and a complex tangent
coordinate. Orient the quotient by putting the positive circle
direction first in the orientation of the surface. At circle
weight \(+1\) the quotient boundary orientation is the complex
curve orientation; at weight \(-1\) it is its negative. The
two normal weights are opposite. Thus, although their boundary
orientations sum to zero in \(H_2(Y)\), the two fixed curves
with their complex orientations represent the same generator
of \(H_2(Y)\). Their maps to \(F\) are holomorphic of degree
one. Therefore the Mayer--Vietoris map
\[
 H_2(\widetilde D)\longrightarrow H_2(Y)\oplus H_2(D)
\]
has matrix
\[
 \begin{pmatrix}1&1\\-1&-1\end{pmatrix}.
\]
It has a primitive rank-one image, rank-one kernel and torsion-free
rank-one cokernel.

\item The normalization conductor is a topological pushout, and taking
a compact group quotient commutes with that pushout. The inclusion
\(\widetilde D\hookrightarrow Y\) is a cofibration: the finite
toric orbit-stratum decomposition can be triangulated with its
boundary spheres and edge arcs as a subcomplex. Thus this pushout
is a homotopy pushout and its Mayer--Vietoris and van Kampen
sequences apply. Van Kampen gives
\[
 \pi_1(H)=\langle a,b\mid ab^{-1}a^{-1}b=1\rangle
            =\mathbb Z^2.
\]
The preceding homology maps give
\[
 H_3(H)=\mathbb Z,\qquad
 0\longrightarrow\mathbb Z\longrightarrow H_2(H)
      \longrightarrow\mathbb Z\longrightarrow0,
 \qquad H_1(H)=\mathbb Z^2.
\]
The middle sequence splits as a sequence of abelian groups. Its
left generator is the image of \((0,[F])\) in the displayed
second-homology cokernel. It is therefore a primitive nonzero
class in \(H_2(H)\). There are no cells above dimension three.
This proves all assertions, including the nonbounding claim.
The endpoint equivalence classes, incidence ranks, commutator
abelianization, and primitive matrix are checked in
\path{code/d_push_cusp_orbit_cells.py}; the orientation argument
and comparison with the geometric pushout are supplied above.

\end{enumerate}
\end{proof}




\begin{lemma}[equivariant retraction of the cusp body]\label{push:src:C.3.8}
\StatementRecord{push:src:C.3.8}{proved}{GPT-6 Astra (OpenAI, ChatGPT Pro mode), effort Pro, 2026-11-17; Claude Fable 5.2 (Anthropic), effort high, 16 November 2026}{the hypotheses, notation and definitions stated here; \cref{push:src:C.3.7}}

For a sufficiently small closed cusp disc \(\overline\Delta\),
the compact manifold with boundary
\(V=f^{-1}(\overline\Delta)\) has a \(\rho\)-equivariant strong
deformation retraction onto \(W=f^{-1}(0)\). Consequently
\(V/\rho\) strongly deformation retracts onto \(H=W/\rho\),
and the fixed sphere \(F\) is primitive and nonzero in
\(H_2(V/\rho;\mathbb Z)\). It cannot bound a three-chain,
and hence cannot bound a three-ball, inside this cusp orbit body.
\end{lemma}


The existence of the torus-equivariant strong deformation retraction is already \citet[Proposition~7.2, p.~41; proof completed on p.~50]{AC26}. Restricting that equivariance to the scalar circle gives the retraction used here. The additional assertion is the primitive nonzero class of the fixed sphere in the quotient body; its identification uses the quotient calculation of \cref{push:src:C.3.7}. The complete retraction argument below is retained.

\begin{proof}
\begin{enumerate}[label=\arabic*.]
\item Write \(q\) for the cusp coordinate and set \(g=|q\circ f|^2\).
Choose a smooth invariant Hermitian metric by averaging any smooth
Hermitian metric over the compact circle. The negative gradient
flow of \(g\) is equivariant, fixes \(W\) pointwise, and points
strictly inward on \(\partial V\). Its flow \(\varphi_t\) exists
for every \(t\geq0\), since \(V\) is compact.

\item We prove the uniform convergence to \(W\) that is needed for a
strong deformation retraction. At every point of \(W\), the
actual semistable cusp charts give holomorphic coordinates with
\(q\circ f=z_1\cdots z_k\), where \(1\leq k\leq3\).
An initial holomorphic unit can be absorbed in one coordinate.
For a point off \(W\), some \(i\) satisfies
\( |z_i|\leq|q\circ f|^{1/k}\), and therefore
\[
 \left|\frac{\partial(q\circ f)}{\partial z_i}\right|
       \geq |q\circ f|^{1-1/k}.
\]
Uniform comparison with the chosen metric on a finite cover of
the compact fibre, and \(g\leq1\) after shrinking, give a
constant \(c>0\) such that
\[
 |\nabla g|\geq c g^{5/6}
\]
near \(W\). On the compact complement of a smaller such
neighbourhood, \(g>0\) and \(\nabla g\ne0\), because
\(f\) is a submersion off its central fibre. Decreasing \(c\)
extends this inequality to all of \(V\setminus W\).

\item Along a trajectory put \(h(t)=g(\varphi_t(x))\). Then
\(h'=-|\nabla g|^2\leq-c^2h^{5/3}\). If
\(G=\max_V g>0\), integration gives the uniform bound
\[
 h(t)\leq\bigl(G^{-2/3}+(2c^2/3)t\bigr)^{-3/2}.
\]
The remaining length of the trajectory starting at time \(t\)
is bounded by
\[
 \int_t^\infty|\nabla g(\varphi_s(x))|\,ds
 =\int_0^{h(t)}\frac{dv}{|\nabla g|}
 \leq\frac6c h(t)^{1/6}.
\]
One may first integrate to a finite time; the displayed uniform
bound and monotone convergence justify the endpoint expression.
Thus every trajectory has a limit in \(W\), and the tail-length
bound tends uniformly to zero. The maps \(\varphi_t\) converge
uniformly to a continuous equivariant retraction \(r:V\to W\).
The homotopy given by \(\varphi_{s/(1-s)}\) for \(s<1\)
and by \(r\) at \(s=1\) is continuous, equivariant, and fixes
\(W\) throughout. This is the required strong deformation
retraction. It descends to the orbit quotient, so the homology
and primitivity statement transfers from \ref{push:src:C.3.7}. The same argument
does not supply a global Seifert surface in the full orbit sphere.

\end{enumerate}
\end{proof}




\begin{lemma}[the primitive circle marking]\label{push:src:C.3.9}
\StatementRecord{push:src:C.3.9}{proved}{GPT-6 Astra (OpenAI, ChatGPT Pro mode), effort Pro, 2026-11-17; Claude Fable 5.2 (Anthropic), effort high, 16 November 2026}{the hypotheses, notation and definitions stated here; \cref{push:src:C.3.7}}

Use the conductor generators \(a,b\) of \ref{push:src:C.3.7}, with
\(E_1\) represented by the edge of direction \(e_1-e_2\),
\(E_2\) by the edge of direction \(e_1\), and the closing path
\(c\) on the fixed edge of direction \(e_2\). Under the
natural period marking of the cusp orbit fibre,
\[
 a\longmapsto\widehat u,\qquad
 b\longmapsto\widehat\gamma+\widehat u.
\]
In particular the circle obtained by gluing the two fixed boundary
spheres of the normalization has oriented generator
\(ab^{-1}=-\widehat\gamma\). It is primitive and has no
additional \(\widehat u\) component.
\end{lemma}


\begin{proof}
\begin{enumerate}[label=\arabic*.]
\item The cusp period action has fan translation
\(\lambda\mapsto B_0\lambda\), with
\[
 B_0=\begin{pmatrix}0&1\\-1&0\end{pmatrix},
 \quad \lambda=(\gamma,u)^t.
\]
This is the actual period matrix
\(Z=sB_0+C(q)\) of the cusp construction, not an unmarked
torus identification. The unit factors \(\exp(2\pi i C(q)\lambda)\)
fix the toric strata and hence do not change their triangle labels.
Take the cyclic fan directions
\[
 (D_1,\ldots,D_6)=
 (e_1-e_2,e_1,e_2,e_2-e_1,-e_1,-e_2).
\]
For \(L_v=(v,v+e_1,v+e_2)\) and
\(U_v=(v+e_1,v+e_2,v+e_1+e_2)\), the vertices of the
normalizing component have the following actual lifts:
\[
 \begin{array}{c|cccccc}
 &p_1&p_2&p_3&p_4&p_5&p_6\\\hline
 \text{triangle}&U_{-e_2}&L_0&U_{-e_1}&L_{-e_1}
                 &U_{-e_1-e_2}&L_{-e_2}.
 \end{array}
\]
For example \(p_1\) is the triangle
\(\{0,e_1-e_2,e_1\}=U_{-e_2}\); the other entries follow
by the same explicit substitution.

\item Choose \(A=L_0\) and the lift of \(c\) from
\(p_2=L_0\) to \(p_3=U_{-e_1}\).
The edge \(D_1\), oriented from \(A\) to \(B\), goes
from \(L_{-e_2}\) to \(U_{-e_2}\). Translate it by
\(+e_2\) to start at the selected lift \(L_0\).
It ends at \(U_0\), which is the chosen \(B\)-lift
translated by \(+e_1\). Thus \(a=E_1c^{-1}\) has
fan displacement \(+e_1=B_0\widehat u\).
The edge \(D_2\) oriented from \(A\) to \(B\)
starts already at \(L_0\) and ends at \(U_{-e_2}\).
Its displacement relative to the chosen \(B\)-lift is
\(e_1-e_2=B_0(\widehat\gamma+\widehat u)\).
This proves the two labels.

\item The real normalization circle is the path along \(D_1\)
and back along \(D_2\), closed by the gluing of the fixed
boundary spheres. Its image is \(E_1E_2^{-1}=ab^{-1}\),
so its label is \(-\widehat\gamma\). Equivalently, the
fixed-edge gluing has fan translation
\(B_0(1,0)^t=-e_2\), corresponding to the period
\(\widehat\gamma\); traversing its fundamental
domain in the stated direction has the opposite sign.
The two computed labels have determinant \(-1\), and hence
give the entire period lattice. A loop lift is legitimate even
at a vertex: a nonzero fan translation cannot preserve the finite
set of component labels incident to that vertex, so the period
action on the toric cover, also after the circle quotient, is free.
The finite triangle and displacement calculation is checked in
\path{code/d_push_cusp_marking.py}.

\end{enumerate}
\end{proof}




\begin{proposition}[reduction to a specified annular orbit map]\label{push:src:C.3.10}
\StatementRecord{push:src:C.3.10}{proved}{GPT-6 Astra (OpenAI, ChatGPT Pro mode), effort Pro, 2026-11-17; GPT-6 Astra (OpenAI, ChatGPT Pro mode), S6\_06\_BLIND, effort Pro, 2026-11-17; Claude Fable 5.2 (Anthropic), effort high, 16 November 2026}{the hypotheses, notation and definitions stated here; \cref{push:src:C.3.1}; \cref{push:src:C.1.1}; \cref{push:src:C.3.6}}

Let $\mathcal Q=X_{c_0}/\rho$ be the orbit space for the circle of
\ref{push:src:C.3.1}. Choose a small cusp disc with coordinate $q$, and a smooth
radial function $\chi$ supported in that disc, equal to one on a
smaller disc. On its punctured annulus, let
$v_\gamma=\operatorname{pr}_{\delta}\Pi_{c_0}\widehat\gamma$
be the projected real period. Define an orbit map $\psi$ by
\[
 \psi([z])=[z+\chi(q)v_\gamma(q)]
\]
on the cusp annulus and by the identity elsewhere. This is a
well-defined diffeomorphism, equal to the identity near the fixed
set. It has a smooth equivariant lift to $X_{c_0}$. Every such lift
has mapping class $6[\Phi]$, for the positive $w$-loop convention
of \ref{push:src:C.1.1} and the closing map of Paper~I, Proposition~18.5.

The assertion concerns the support of the \emph{orbit map}.
It does not assert an identity-ended lift supported in the annulus.
In particular it does not evaluate $6[\Phi]$.
\end{proposition}


\begin{proof}
\begin{enumerate}[label=\arabic*.]
\item We first specify the quotient transport used in the comparison.
Take $c(r)=c_0-12r$, $0\leq r\leq1$. On an ordinary fibre,
the quotient by $\rho$ is obtained by projecting the additive
period coordinates to $\mathbb C\times\mathbb R$, discarding
$\operatorname{Re}\zeta_2$. The period change
$(c(r)-c_0)e\otimes\gamma$ has real coefficient, so its
projected lattice is unchanged. This gives constant quotient
charts over the ordinary base.

\item The finite root quotient charts are also constant: the change of
their period columns and affine deck translations is in the real
$e$-direction and disappears under this projection. Their overlap
with the ordinary chart changes, however. With the notation and
signs of Paper~I, Proposition~18.5,
\[
 \sigma_{j,c(r)}-\sigma_{j,c_0}
       =-12r\,\epsilon_j w_j e,
 \qquad w_j=(2\pi i)^{-1}\log s_j,
 \quad (\epsilon_1,\epsilon_2)=(1,-1).
\]
Let $R_a$ denote the positive real subgroup of the global
$\mathbb C^*$-action, with multiplier $e^a$. The transport in
the root chart which agrees with the ordinary identity is thus
$R_{12r\epsilon_j\log|s_j|}$ on the overlap. Choose a smooth
rotation-invariant real function $f_j(s_j)$ which vanishes near
zero and equals $\log|s_j|$ on that overlap. Replacing this by
$R_{12r\epsilon_j f_j(s_j)}$ extends the transport over the
root centre and commutes with its finite deck group. These
functions are smooth on the free root-cover quotient; no assertion
that $|t|^{1/m_j}$ is a smooth coordinate on the coarse base is
used. At the cusp, the varying multiplicative period unit is a
$\delta$-phase, so the toric quotient charts and their overlaps
are constant. We take the identity transport there.

\item This quotient transport is induced by a smooth equivariant
transport on the actual family. Here is the lifting justification.
On an ordinary or finite root chart the affine family and its
smooth lattice coordinates give local smooth horizontal lifts of
the prescribed quotient vector field. At the cusp use its smooth
toric charts, with the finite base parameter interval retained;
the dependence just removed is a circle phase, so these charts
give local equivariant lifts too, including at fixed points.
Two such lifts have the same quotient field and differ by a
vertical circle field on the free stratum. The difference vanishes
at a fixed point. An invariant smooth partition of unity on the
actual total space patches these local fields. Their parameter
component remains one and their quotient is the specified field.
Properness gives a flow for the whole compact interval. Denote
its endpoint by $H$. This argument uses smooth local lifts in the
given slices, rather than assuming an arbitrary orbit-space
isotopy lifts across the fixed set.

\item Put $G=\widetilde F_{12,c_0-12}\circ H$. By \ref{push:src:C.1.1} this smooth
equivariant diffeomorphism represents $6[\Phi]$. On ordinary
quotient charts, Paper~I, Proposition~18.5 gives the positive real flow
$R_{a(t)}$, where
\[
 a(t)=\log|K(t)|,\qquad K(t)=t^4/(t-1)^3.
\]
On the finite root chart the same map is the flow with time
\[
 a_j(s_j)=\log|K(t)s_j^{-12\epsilon_j}|
                       +12\epsilon_j f_j(s_j).
\]
The first term is the logarithm of a nonvanishing unit and the
second is smooth; hence this time is smooth through the root
centre. On its outer overlap it equals $a(t)$. At the cusp the
quotient map is the bounded toric unit $U=U_{12,c_0-12}$ of
Paper~I, Proposition~18.5. Write $u(q)$ for its additive projected logarithm.
That logarithm is single-valued on the full disc. The exact
period identity of Paper~I, Proposition~18.5 is
\[
 u(q)+v_\gamma(q)=R_{a(q)}
\]
in the additive universal covering vector space; on the right
we mean the additive displacement of that real flow.

\item Multiplication by $U^{-s\chi(q)}$, $0\leq s\leq1$, defines a
smooth equivariant isotopy supported in the cusp. To see this
through every toric stratum, choose the explicit holomorphic
additive logarithm of the two nonzero unit coordinates of $U$.
Exponentiating $-s\chi$ times that logarithm gives a torus
element depending smoothly on $q$; multiplication acts smoothly
in every toric chart and commutes with the period identifications.
It is the identity near the outer edge of the disc.

\item There is a second global equivariant isotopy, the real
$\mathbb C^*$-flow with time $-s(1-\chi)a$, using $a_j$ in
the finite charts. The times agree on overlaps, are smooth on
the actual root quotients, and vanish on an inner cusp disc.
They are constant along the $\mathbb C^*$-orbits, so the inverse
flow is obtained by negating the time. Consequently this is an
isotopy of diffeomorphisms, not just a pointwise multiplication.

\item Compose $G$ with the endpoints of these two isotopies. Away
from the cusp its orbit map becomes the identity. In the cusp
annulus its projected displacement is exactly
\[
 a-\chi u-(1-\chi)a=\chi(a-u)=\chi v_\gamma.
\]
Where $\chi=1$, this is a whole projected lattice period and
the orbit map is again the identity; on the central fibre the
cusp-unit description gives the same identity directly. Thus the
result is an actual smooth equivariant lift of the stated $\psi$,
and has the same mapping class as $G$.

\item Finally $A_0\widehat\gamma=\widehat\gamma-\widehat\delta$
makes $v_\gamma$ single-valued on the punctured cusp after the
$\delta$-quotient. Translation by $\chi v_\gamma$ and its
negative are smooth mutually inverse maps on that free orbit
annulus, and they are the identity on collars of its ends.
This proves the direct well-definedness and smoothness assertion.
\ref{push:src:C.3.6} shows that every other existing smooth equivariant lift has
the same mapping class. This construction supplies no assertion
that the full lift is supported in the annulus: its residual
circle gauge outside that annulus has not been discarded.

\end{enumerate}
\end{proof}




\begin{proposition}[a nonzero normal Euler number]\label{push:src:C.3.11}
\StatementRecord{push:src:C.3.11}{proved}{GPT-6 Astra (OpenAI, ChatGPT Pro mode), effort Pro, 2026-11-17; GPT-6 Astra (OpenAI, ChatGPT Pro mode), S6\_06\_BLIND, effort Pro, 2026-11-17; Claude Fable 5.2 (Anthropic), effort high, 16 November 2026}{the hypotheses, notation and definitions stated here; \cref{push:src:C.3.7}}

Glue the two fixed boundary spheres of \(Y=\mathrm{dP}_6/\rho\)
using their actual conductor identification, and call the resulting
closed oriented three-manifold \(N\). Its natural map to the
smooth cusp orbit five-manifold is an immersion. Its sole
self-intersection circle has two disjoint preimage circles, the
two pairs of nonfixed edge intervals of \ref{push:src:C.3.7}.
Along the fixed sphere \(F\subset N\), the oriented normal
two-plane bundle of this immersion has Euler number of absolute
value two. If its complex orientation is that of the Hopf invariant
\(a=xz\), where \(x,z\) have circle weights \(+1,-1\), the
Euler number is \(-2\).
\end{proposition}


\begin{proof}
\begin{enumerate}[label=\arabic*.]
\item Away from the conductor the normalization quotient is embedded.
At an ordinary nonfixed double point the circle acts freely along
the double curve, and the two smooth three-dimensional branches
intersect transversely along its one-dimensional orbit quotient.
At a fixed double point the two normal coordinates have weights
\(+1,-1\). The Hopf quotient coordinates are
\[
 a=xz\in\mathbb C,\qquad
 b=(|x|^2-|z|^2)/2\in\mathbb R,
\]
together with the coordinate along the fixed curve. The two
halves of \(Y\) glue as the two signs of \(b\), so their
union is smooth. At a triple point the full base equation is
\(q=xyz\), with \(y\) tangent to \(F\) and invariant.
In the quotient it becomes \(q=ay\). Hence the central orbit
fibre is the union of the two smooth three-planes
\(\{a=0\}\) and \(\{y=0\}\), intersecting transversely
along the \(b\)-axis. This verifies immersion and transversality
at both previously distinguished vertices. The endpoint
identifications of \ref{push:src:C.3.7} join the two nonfixed double arcs into
one circle; its preimages in \(N\) are disjoint circles.

\item The two complex weight lines in \(N_{F/X}\) are each
\(\mathcal O_{\mathbb P^1}(-1)\). To see this with the actual
normalization, each is the normal line of one fixed edge in its
\(\mathrm{dP}_6\) branch, and each edge has self-intersection
\(-1\). This remains true at the two vertices because those
branches and their edge normals are smooth there. The circle
characters distinguish the lines globally, so no choice of a
splitting of an extension is required.

\item The Hopf invariant \(a=xz\) is valued in the tensor product
of these two weight lines. It therefore gives the complex normal
line of the branch \(\{a=0\}\) along \(F\) as
\[
 \mathcal O(-1)\otimes\mathcal O(-1)=\mathcal O(-2).
\]
At a triple point \(q=ay\), the coordinate \(a\) still is
a normal coordinate of that branch. The zero of \(y\) does
not alter its normal bundle; using \(q\) instead would lose
precisely this distinction. The stated Euler number is consequently
\(-2\) for the \(a\)-complex orientation and \(+2\) if that
normal orientation is reversed.

\item In particular this immersion cannot simply be compressed along
\(F\) to an immersed three-sphere while retaining its normal
two-plane bundle: extension over a bounding three-disc would force
its Euler pairing on the boundary sphere to vanish. This identifies
a required framing correction; it does not compute that correction
or a Haefliger invariant for the annular twist.

\end{enumerate}
\end{proof}




\begin{proposition}[obstruction to the literal translational Seifert quotient]\label{push:src:C.3.12}
\StatementRecord{push:src:C.3.12}{proved}{GPT-6 Astra (OpenAI, ChatGPT Pro mode), effort Pro, 2026-11-17; Claude Fable 5.2 (Anthropic), effort high, 16 November 2026}{the hypotheses, notation and definitions stated here; \citet[Remark 8.3]{E26}; \cref{push:src:C.3.10}; \cref{push:eq:C12-u-flow}; \cref{push:eq:C12-limiting-identification}}

Fix an admissible parameter and a small closed cusp body $V$, with
central fibre $W$. On every regular fibre let $K_q$ be the real
three-dimensional subtorus marked by $\ker\gamma$, namely by
$\widehat u,\widehat w,\widehat\delta$.
If $Z$ is Hausdorff and a continuous map $h:V\to Z$ is constant
on every translate of $K_q$ for $q\ne0$, then $h|_W$ is constant.

Consequently there is no continuous extension of this literal
regular-fibre translational quotient whose central fibre is a
circle. In particular, the translational quotient in the first
paragraph of \citet[Remark 8.3]{E26} cannot be identified with the
circle-valued cusp moment quotient described there while retaining
those equivalence classes on every regular fibre. A construction
which changes the regular-fibre equivalence relation is not excluded.
This correction concerns the rank-three quotient; it does not affect
the smooth $\widehat\delta$-circle quotient in \ref{push:src:C.3.10}.
\end{proposition}


The final paragraph below also proves nonextension of the raw
invariant-circle coordinate. The proposition has the stronger quantified
scope of every Hausdorff-valued continuous map with the stated regular
translational leaves. 
\begin{proof}
\begin{enumerate}[label=\arabic*.]
\item We use the actual cusp toric coordinates and periods. Write
\[
 q=e^{2\pi i s},\qquad
 Z=sB_0+C(q),\qquad
 B_0=\begin{pmatrix}0&1\\-1&0\end{pmatrix},
 \qquad C_{22}(q)=\mu(q).
\]
The columns correspond to $(\widehat\gamma,\widehat u)$;
the unitary torus directions are $(\widehat w,\widehat\delta)$.
On the dense torus, with multiplicative coordinates $(z_1,z_2)$,
translation by the real multiple $r\widehat u$ is exactly
\begin{equation}\label{push:eq:C12-u-flow}
 (z_1,z_2)\longmapsto
 \bigl(q^r e^{2\pi i r C_{12}(q)}z_1,
             e^{2\pi i r\mu(q)}z_2\bigr).
\end{equation}
Here a branch is chosen on the positive real $q$-ray. This is a
translation in $K_q$ for every real $r$, not merely for integral $r$.
The two remaining circle directions rotate $z_1$ and $z_2$
independently.

\item The limiting constant $\mu_0=\mu(0)$ has nonzero imaginary part.
For the branch used in Paper~I, Proposition~22.16 it is
\[
 \mu_0=\frac12-\frac{i}{2\pi}\log(5+2\sqrt6),\qquad
 |e^{2\pi i\mu_0}|=5+2\sqrt6>1.
\]
Only the nonvanishing of its imaginary part is needed. It can
also be checked directly without that branch choice: the cusp
limit of $\wp(\mu;\tau)=0$ is
$\pi^2\csc^2(\pi\mu_0)-\pi^2/3=0$, so
$\sin^2(\pi\mu_0)=3$, which is impossible for real $\mu_0$.
The Fourier limit and holomorphic finite cusp branch are the
ones checked in Paper~I, Proposition~22.16.

\item Consider the edge between the lattice vertices $0$ and $e_1$
of the Mumford fan. On the open part of its two adjacent
components the toric chart is
\[
 x=z_1,\qquad y=q/z_1,\qquad z=z_2\in\C^*,
 \qquad xy=q.
\]
This chart follows directly from the two primitive height-one
rays $(0,0,1)$ and $(1,0,1)$; its additional invertible character
is $z_2$. It is also the ordinary double-curve chart of the
fixed bottom triangulation. Let $P(x,z)$, for $x,z\in\C^*$,
denote its central point $(x,0,z)$, and let $Q(z)$ denote
$(0,0,z)$ on the double curve. We apply these notations also
to their images in the period quotient. The set of all $P(x,z)$
is the dense torus of the normalizing $\mathrm{dP}_6$ component,
whose image is dense in the irreducible central fibre $W$.

\item First $h(P(x,z))$ is independent of $x$. To check positive
radial changes, fix $\lambda\in\R$, put $q_n=e^{-n}$, and
translate $(x,q_n/x,z)$ by
$r_n\widehat u$, where $r_n=-\lambda/n$. Since $C_{12}$ and
$\mu$ are bounded, the two multiplicative factors in
\eqref{push:eq:C12-u-flow} tend respectively to $e^\lambda$ and
$1$. The translated sequence therefore tends to $P(e^\lambda x,z)$,
whereas the original sequence tends to $P(x,z)$. Their $h$-images
agree term by term. Continuity and the Hausdorff property give
equality of the two limiting images. Rotation by
$\widehat w$ treats the argument of $x$ in the same way.
Thus we may write $H(z)=h(P(1,z))$.

\item Next fix $0<r<1$ in \eqref{push:eq:C12-u-flow} and let $q\downarrow0$
on the positive real ray. Starting at $(x,q/x,z)$, the translated
coordinates satisfy
\[
 x_r=q^r e^{2\pi i r C_{12}(q)}x\longrightarrow0,\qquad
 y_r=q^{1-r} e^{-2\pi i r C_{12}(q)}/x\longrightarrow0,
\]
and $z_r\longrightarrow e^{2\pi i r\mu_0}z$. Hence the same
continuity argument gives the exact equality
\begin{equation}\label{push:eq:C12-limiting-identification}
 H(z)=h\bigl(Q(e^{2\pi i r\mu_0}z)\bigr)
 \qquad(0<r<1).
\end{equation}
For $r_1,r_2\in(0,1)$ apply this equality first to $z$ and
then to $z'=e^{2\pi i(r_1-r_2)\mu_0}z$. Their double-curve
limits are the same point, so
\[
 H(z)=H(e^{2\pi i a\mu_0}z)\qquad(-1<a<1).
\]
Iteration gives the equality for every real $a$. Rotation by
$\widehat\delta$ also gives $H(z)=H(e^{i\vartheta}z)$ for every
$\vartheta$. Since $\operatorname{Im}\mu_0\ne0$, the moduli
$|e^{2\pi i a\mu_0}|$ range over all positive real numbers.
Combining these two invariances proves that $H$ is constant on
$\C^*$. Therefore $h$ is constant on the dense torus in $W$.
The inverse image of a point of the Hausdorff space $Z$ is
closed; density proves constancy on all of $W$.

\item For comparison, the raw real period coordinate already fails
continuity in an elementary fixed-edge chart. Take
$q=e^{-2\pi T}$, $z_1=1$, $z_2=q^r$, with $0<r<1$ and
$T\to+\infty$. In the chart
$(z_2,q/z_2,z_1)$ all these paths tend to the same point
$(0,0,1)$. Their logarithmic modulus vectors are $(0,rT)$.
Since $\operatorname{Im}Z=TB_0+O(1)$, their real period
vectors are
\[
 (\operatorname{Im}Z)^{-1}(0,rT)^t
                  =(-r,0)^t+O(T^{-1}).
\]
Thus their raw $\gamma$ circle coordinates tend to $-r$ modulo
$\Z$, with different limits at that one central point. The
stronger argument \eqref{push:eq:C12-limiting-identification} shows
why an arbitrary base-dependent reparametrization of those
circle coordinates cannot repair the asserted central circle
while keeping the same regular-fibre translational leaves.
The exact limiting multiplier and the inverse-matrix calculation
are independently checked in
\path{code/d_push_cusp_quotient_obstruction.py}; the continuity
and closure argument above, rather than that finite check,
proves the universal assertion.

\end{enumerate}
\end{proof}




\begin{proposition}[a corrected marked Seifert comparison]\label{push:src:C.3.13}
\StatementRecord{push:src:C.3.13}{proved}{GPT-6 Astra (OpenAI, ChatGPT Pro mode), effort Pro, 2026-11-17; Claude Fable 5.2 (Anthropic), effort high, 16 November 2026}{the hypotheses, notation and definitions stated here; \cref{push:eq:C13-weighted-map}; \cref{push:eq:C13-real-overlap}; \cref{push:src:C.3.7}; \cref{push:src:C.3.9}; \cref{push:src:C.3.12}; \cref{push:eq:C13-cusp-gauge}; \cref{push:src:C.3.20}; \cref{push:src:C.3.22}}

There is a smooth map $\mathcal G:X_c\to S^3$, invariant under
the $\widehat\delta$ circle, with the following properties.
It covers the given base coordinate through the explicit weighted
Seifert map
\[
 p:S^3\longrightarrow\PP^1,\qquad
 p(z_1,z_2)=\frac{z_2^3}{z_2^3-z_1^4},\qquad
 \lambda\cdot(z_1,z_2)=(\lambda^3z_1,\lambda^4z_2).
\]
Outside a chosen small cusp body it is the actual quotient by
the regular $\ker\gamma$ subtori, followed by a specified
base-dependent circle gauge and the explicit identification below.
It includes the actual finite root quotients with rotation data
$+1/3$ and $-1/4$. Inside the cusp body the map is modified;
it does not preserve the original translational leaves.

On the outer cusp collar a positively marked $\widehat\gamma$
circle maps orientation-preservingly onto a regular weighted
Seifert fibre. The central fibre of the target is the positive
$(3,4)$ torus knot, and a nearby regular fibre has linking number
$12$ with it, in the standard orientation of $S^3\subset\C^2$.
This is framing data in the constructed three-dimensional target.
No identification of its normal framing with the six-dimensional
orbit-knot framing, and no value of $6[\Phi]$, follows from this
statement.
\end{proposition}


\begin{proof}
\begin{enumerate}[label=\arabic*.]
\item We first fix the real character and all of its overlap corrections.
On an ordinary fibre, let
$\gamma_z:\C^2\to\R$ be the real-linear functional determined by
\[
 \gamma_z(\Pi_c(z)v)=\gamma(v)\qquad(v\in\R^4).
\]
The real period matrix is invertible in the admissible range.
The resulting circle coordinate
$\theta=\gamma_z(\zeta)\pmod\Z$ is invariant under ordinary
monodromy because $\gamma A_j=\gamma$. It is the literal
regular-fibre quotient character, not a new choice of marking.

\item On the $j$th finite root chart write $\theta_j$ for this
character before the affine finite quotient. The exact map from
that chart to the ordinary chart is translation by
$\sigma_{j,c}=w_j\Pi_cv_j$, where
$w_j=(2\pi i)^{-1}\log s_j$ and
$\epsilon_j=\gamma(v_j)$ is $1,-1$, respectively.
Consequently its real circle overlap is exactly
\begin{equation}\label{push:eq:C13-real-overlap}
 \theta=\theta_j+\epsilon_j\operatorname{Re}w_j+a_j^{\mathrm{rad}}(t),
 \qquad
 a_j^{\mathrm{rad}}(t)=\operatorname{Im}w_j\,
                  \gamma_z(i\Pi_cv_j).
\end{equation}
The last expression is real and single-valued on the punctured
base disc. Indeed $\operatorname{Im}w_j$ is single-valued,
$v_j$ is fixed by the local integral monodromy, and its
complex-linear covariance preserves $\gamma_z$.
No logarithmic radial term has been dropped.

\item Choose disjoint finite discs and smooth radial cutoffs equal to
one on smaller discs and zero outside the larger ones. They define
a real smooth function $d(t)$ on $B^\circ$ which is $-a_j^{\mathrm{rad}}(t)$
on the smaller punctured discs and is zero on a cusp collar.
Choose the roots $F^3=t$, $G^4=t-1$ on the uniformizing cover and
the logarithm $\ell=(2\pi i)^{-1}\log(F/G)$ used in Paper~I, Proposition~18.5.
Put
\[
 \Theta=\theta-\operatorname{Re}\ell+d(t).
\]
Let $R_t>0$ be the unique solution of
\[
 R_t^6|G|^2+R_t^8|F|^2=1.
\]
Existence and uniqueness follow from strict increase from zero
to infinity; the derivative in $R_t$ is positive, so the
implicit-function theorem gives the required smooth dependence.
The proposed map on the ordinary region is
\begin{equation}\label{push:eq:C13-weighted-map}
 (z_1,z_2)=
 \bigl(R_t^3G e^{6\pi i\Theta},\,
                R_t^4F e^{8\pi i\Theta}\bigr).
\end{equation}
Its norm is one and direct substitution gives $p(z_1,z_2)=t$.
Increasing $\theta$ by $a$ is exactly the weighted action of
$e^{2\pi ia}$, so the positive $\gamma$ orientation is retained.

\item We check descent and the finite centres explicitly. For the
clockwise first root generator $F$ changes by $e^{-2\pi i/3}$,
$G$ is unchanged, and $\ell$ changes by $-1/3$.
Thus $\Theta$ changes by $+1/3$, and both coordinates in
\eqref{push:eq:C13-weighted-map} remain unchanged. For the second
generator $G$ changes by $e^{-2\pi i/4}$, $F$ is unchanged,
and $\Theta$ changes by $-1/4$. Again both coordinates are
unchanged: their total exponents are respectively $-1$ and
$-1$. These generators and changes of logarithm prove descent
on the entire ordinary base.

\item On a full finite root disc the exact identity from Paper~I, Proposition~18.5 is
$\ell=\epsilon_jw_j+\kappa_j$, where $\kappa_j$ is a
holomorphic logarithm of an invariant nonvanishing unit.
Equations \eqref{push:eq:C13-real-overlap} and $d=-a_j^{\mathrm{rad}}$ therefore give
\[
 \Theta=\theta_j-\operatorname{Re}\kappa_j.
\]
This expression is smooth through the root centre. The vanishing
factor $F$ at the first centre, or $G$ at the second, is the
root coordinate times a holomorphic unit. The other factor is
nonzero. Thus \eqref{push:eq:C13-weighted-map} extends smoothly to
the full root chart. Its transverse derivative is nonzero and
the circle factor has degree one; it identifies the actual
finite quotient solid torus with the corresponding weighted
solid torus. The deck actions on $(s_j,\theta_j)$ are precisely
\[
 (s_1,\theta_1)\mapsto(e^{-2\pi i/3}s_1,\theta_1+1/3),
 \qquad
 (s_2,\theta_2)\mapsto(e^{-2\pi i/4}s_2,\theta_2-1/4).
\]
This verifies the overlap maps, not merely the exceptional
multiplicities of an abstract Seifert space.

\item We next verify the cusp meridian and construct the replacement
map upstairs, without assuming a smooth quotient structure at
the fixed set. Let $V$ be a small smooth compact cusp body and
$\overline V=V/\rho$. \ref{push:src:C.3.7}--\ref{push:src:C.3.9} identify
$H_1(\overline V;\Z)=\Z\widehat\gamma\oplus\Z\widehat u$.
The outer orbit link is the mapping torus of the action
$u\mapsto u+w$ on the $(u,w)$ torus, times the $\gamma$ circle.
Its first homology is therefore freely generated by
$\widehat\gamma,\widehat u$ and a transverse cusp meridian;
$\widehat w$ is a commutator. The two first generators have
the stated images in $H_1(\overline V)$ by the actual fan marking
in \ref{push:src:C.3.9}.

\item The transverse meridian has zero $\gamma$ value and bounds
an actual disc in $V$. To see this without choosing a meridian
up to an unknown fibre twist, take the toric section
$z_1=z_2=1$. In the chart of \ref{push:src:C.3.12} it is
$(x,y,z)=(1,q,1)$, so it extends across $q=0$.
Its additive logarithms are zero and hence $\theta=0$
identically on its boundary. Thus the integral class on
$\overline V$ with values $1,0$ on
$(\widehat\gamma,\widehat u)$ restricts to exactly the raw
character $[d\theta]$ on the orbit link, including the transverse
meridian value. Pull it back to $V$.

\item Represent the pulled-back integral class by a smooth closed real
one-form $\alpha$ with integral periods. Average it over the
$\delta$ circle. Each group element is isotopic to the identity,
so the cohomology class and its integral periods are unchanged.
The averaged form is invariant. If $K$ denotes the circle
generator, Cartan's formula gives $d(\alpha(K))=0$.
The body is connected and has fixed points, where $K=0$;
therefore $\alpha(K)=0$ everywhere.
On an outer collar $\alpha-d\theta$ is exact, since the two
forms have the same periods. Choose its real primitive,
average the primitive over the circle, and extend it to $V$
using an invariant collar cutoff. Subtracting its differential
from $\alpha$ makes the resulting invariant integral closed
form equal to $d\theta$ on a smaller outer collar. Integrating
it modulo $\Z$, with the constant chosen at one collar point,
gives a smooth map
\[
 g:V\longrightarrow\R/\Z
\]
which equals $\theta$ on that collar. Its contraction with $K$
is zero, so $g$ is invariant. This construction is entirely on
the smooth manifold $V$ and yields a continuous factor on
$\overline V$.

\item For completeness, the prescribed orbit boundary map has a
unique relative homotopy class of extensions to $\overline V$.
The restriction of $H^1(\overline V;\Z)$ to the link is
injective by the two period generators just checked, and both
spaces are connected. The relative cohomology exact sequence
therefore gives $H^1(\overline V,\partial\overline V;\Z)=0$.
Maps to $\R/\Z$ with fixed boundary values are classified by
this relative group: a map represents its integral period class,
and a map with zero class has a real lift which supplies the
relative homotopy. This assertion does not classify normal
framings of any embedded trace.

\item Use now the holomorphic cusp coordinate $q=1/t$; changing from
the original Mumford coordinate multiplies it by a unit and
does not change the meridian just verified. Write
$q=e^{2\pi i s}$. Compatible roots satisfy
\[
 F=q^{-1/3},\qquad G=q^{-1/4}(1-q)^{1/4},\qquad
 \ell=-s/12+\kappa_0(q),\qquad
 \kappa_0(q)=-\frac{\log(1-q)}{8\pi i}.
\]
A different compatible root choice only changes the displayed
unit gauge. The smooth fibre coordinate in the weighted model
near its regular cusp knot is
\begin{equation}\label{push:eq:C13-cusp-gauge}
 \vartheta=\Theta-\operatorname{Re}s/12
           =\theta-\operatorname{Re}\kappa_0(q),
\end{equation}
because $d=0$ there. In particular no additional integer
meridian winding is present. This is the required boundary
framing comparison; the finite labels alone would not give it.

\item Let $S(q)>0$ solve
$S(q)^6|1-q|^{1/2}+S(q)^8=1$. The explicit smooth solid-torus
chart of the target is
\[
 \Phi_0(q,\vartheta)=
 \bigl(S(q)^3(1-q)^{1/4}e^{6\pi i\vartheta},
                         S(q)^4e^{8\pi i\vartheta}\bigr).
\]
It is a diffeomorphism onto a neighbourhood of the regular
cusp knot. Indeed its inverse has
$q=1-z_1^4/z_2^3$ and
$e^{2\pi i\vartheta}=z_2(1-q)^{1/4}/(S(q)z_1)$;
all factors are nonzero there. Define the replacement on $V$
by
\[
 \mathcal G(x)=
 \Phi_0\bigl(q(x),g(x)-\operatorname{Re}\kappa_0(q(x))\bigr).
\]
It agrees exactly with \eqref{push:eq:C13-weighted-map} on the outer
collar by \eqref{push:eq:C13-cusp-gauge}, and hence glues with the
ordinary and finite maps already constructed. Smoothness and
$\delta$ invariance hold upstairs on every piece. The first
coordinate of the solid-torus chart proves that it still covers
$f$. By \ref{push:src:C.3.12} its new cusp leaves cannot be the original
translational leaves on every punctured fibre; that change is
an explicit part of the construction.

\item Finally the knot $q=0$ in the target is
$z_2^3=z_1^4$, parametrized with positive orientation by the
weighted circle $(e^{3i\phi},e^{4i\phi})$ with the fixed positive
radii prescribed by $S(0)$. It is the positive $(3,4)$ torus
knot. For a small fixed $q\ne0$, along the nearby fibre
\[
 (z_2^3-z_1^4)\circ\Phi_0(q,\vartheta)
                  =q S(q)^{12}e^{24\pi i\vartheta}.
\]
The argument therefore has degree $12$. The function
$z_2^3-z_1^4$ has a simple complex-normal zero along the
knot in $S^3$, so its argument represents the positive
meridian class in the knot complement. This degree is its
linking number with the knot. The orientation is the one
already matched to positive $\gamma$ in
\eqref{push:eq:C13-weighted-map}. The computation establishes the
three-dimensional push-off datum only. The actual normal-line
comparison and its contribution are established separately in
\ref{push:src:C.3.20}--\ref{push:src:C.3.22}; they are not inferred from this linking number.
The finite phases, radial cancellation, link abelianization and
polynomial identities are checked exactly in
\path{code/d_push_corrected_seifert.py}.

\end{enumerate}
\end{proof}




\begin{proposition}[an actual characteristic-number target]\label{push:src:C.3.14}
\StatementRecord{push:src:C.3.14}{proved}{GPT-6 Astra (OpenAI, ChatGPT Pro mode), effort Pro, 2026-11-17; Claude Fable 5.2 (Anthropic), effort high, 16 November 2026}{the hypotheses, notation and definitions stated here; \cref{push:src:C.3.1}; \cref{push:src:C.3.10}; \cref{push:src:C.1.1}; \citet[Section 1, Propositions 1.3--1.4]{FintushelPao1983D}; \citep{KervaireMilnor1963}; \citep[Corollary 6.6]{Haefliger1966D}; \cref{push:src:C.3.11}; \cref{push:src:C.3.6}; \cref{push:src:C.2.1}; \citet[Section 3, page 1911]{Olbermann2010D}; \citep[Theorem 4, pages 360--361]{Wall1966VE}; \citep{Wall1967CorrE}; \citet[Example 7.14 and Proposition 7.15]{Davis1982D}}

The smooth orbit pair for the circle of \ref{push:src:C.3.1} is diffeomorphic to
the standard pair $(S^5,S^2)$. Its free complement, after deleting
a small fixed tubular neighbourhood, is therefore
$P\simeq S^2\times D^3$. The circle action on $X_c$ is smoothly
equivariantly equivalent, with its circle direction fixed, to the
linear Hopf action on
$\partial(D^4_{\mathbb H}\times D^3)$.

Choose such an identification and regard the actual annular map
$\psi$ of \ref{push:src:C.3.10} as a diffeomorphism of $P$ equal to the identity
near its boundary. Let $\widetilde\psi$ be its extension by the
identity to the double $DP\simeq S^2\times S^3$. Form
\[
 M_\psi=(S^2\times D^4)\cup_{\widetilde\psi}
                  -(S^2\times D^4).
\]
This is a closed simply connected spin six-manifold with
$H^2(M_\psi;\Z)=H^4(M_\psi;\Z)=\Z$ and $H^3(M_\psi;\Z)=0$.
Let $x$ be the generator restricting to the oriented meridian
generator on each $S^2\times D^4$, and put
\[
 k_\psi=\langle x^3,[M_\psi]\rangle,\qquad
 J_\psi=\frac{\langle p_1(M_\psi)x-4x^3,[M_\psi]\rangle}{24}.
\]
Both are integers. The unframed knot coordinate of the orbit
knot of the clutching sphere for $6\Phi$ is $-J_\psi$ in
Wall's surgery convention. Consequently, with one fixed sign
depending on comparison of that convention with the generator
of \ref{push:src:C.1.1},
\[
 \mu(\Sigma_{6\Phi})=\pm\frac{J_\psi}{28}\pmod\Z.
\]
The pure framing subgroup has zero image in $\Theta_7$.
No value of $J_\psi$ or $k_\psi$ is asserted here.
\end{proposition}


\begin{proof}
\begin{enumerate}[label=\arabic*.]
\item The dimension hypothesis in
\citet[Section 1, Propositions 1.3--1.4]{FintushelPao1983D}
is $n>5$. Applied to this semifree action on $S^6$, it gives a
smooth orbit homotopy five-sphere and a smooth embedded fixed
two-sphere. Since $\Theta_5=0$ \citep{KervaireMilnor1963}, the
orbit is the standard $S^5$. The unknotting range
$C_n^q=0$ for $n<2q-3$
\citep[Corollary 6.6]{Haefliger1966D} applies at $(n,q)=(2,3)$.
Thus the fixed embedding is standard and its exterior is
$S^2\times D^3$.

\item Here the orbit identification can be lifted to the claimed
equivariant identification. In the normal bundle of the fixed
curve, conjugating the negative weight line converts the circle
representation to two scalar weight-one complex lines. By
\ref{push:src:C.3.11} their Chern classes are $-1$ and $+1$. The resulting
rank-two complex bundle over $S^2$ has first Chern class zero,
and is trivial, since clutching rank-two complex bundles over
$S^2$ is classified by $\pi_1(U(2))$ and its determinant.
Thus its equivariant tubular neighbourhood is the linear model.
The oriented normal three-plane bundle of the orbit sphere is
also trivial. The ambient pair diffeomorphism can be adjusted
on a tubular neighbourhood to this prescribed normal model:
the relevant change of oriented framing is a map
$S^2\to SO(3)$, and $\pi_2(SO(3))=0$.

\item On the complement, the principal circle bundle has Euler class
equal to a primitive generator of $H^2(P;\Z)$, as is seen on
each normal Hopf sphere. Choose the orbit identification to
match this signed generator with the linear model. Principal
circle bundles are classified by that class, so their complements
are equivariantly isomorphic. On the common boundary the
difference from the prescribed tubular isomorphism is a gauge
$S^2\times S^2\to S^1$. It has a real logarithm because
$H^1(S^2\times S^2;\Z)=0$. Extending that logarithm over a
collar and using a cutoff corrects the bundle isomorphism there.
The two isomorphisms now glue smoothly. This proof uses the
linear tubular model at the fixed set; it does not assume that
an arbitrary quotient diffeomorphism lifts smoothly there.

\item The support of $\psi$ misses the fixed neighbourhood by \ref{push:src:C.3.10}.
Any boundary-fixed diffeomorphism of $P$ preserves the meridian
Euler class, since restriction $H^2(P)\to H^2(\partial P)$ is
injective. It therefore has a circle-bundle lift, which can be
made the identity on a boundary collar by the same logarithmic
gauge correction. Extend that lift by the identity on the fixed
neighbourhood. \ref{push:src:C.3.6} compares it with the lift actually constructed
in \ref{push:src:C.3.10}, so its mapping class is $6[\Phi]$.

\item Clutch two linear seven-discs along this lift. Their fixed discs
clutch by the identity to a fixed $S^3$. Their orbit six-discs
clutch along $\psi$ on $S^5$, giving an orbit homotopy six-sphere,
which is standard because $\Theta_6=0$. Denote the resulting
orbit knot by $K_\psi\subset S^6$. This is the orbit knot of
the actual clutching sphere for $6\Phi$.

\item We spell out its framed surgery, so that the cap is specified.
In each orbit six-disc, the complement of a standard fixed
three-disc neighbourhood is $N=S^2\times D^4$. Its boundary
is divided into an outer and an inner copy of $P=S^2\times D^3$.
The two outer copies glue by $\psi$ to form the knot exterior.
The inner copies give its boundary $S^2\times S^3$ with the
framing inherited from the fixed-disc normal products. Framed
surgery fills that boundary by a third copy of $N$.
Combine this third copy with one of the first two along its
inner $P$. Two four-discs glued along a boundary three-disc
form a four-disc. The combined piece is therefore again $N$,
and the remaining gluing map is $\psi$ on the outer half and
the identity on the inner half. This is exactly the displayed
$M_\psi$, with the boundary orientation convention used there.

\item Van Kampen gives $\pi_1(M_\psi)=0$. Mayer--Vietoris for its
two copies of $S^2\times D^4$, with intersection $S^2\times S^3$,
gives the asserted cohomology groups and no torsion. The
restriction map in degree two is injective, and $w_2$ restricts
to zero on both pieces; thus $w_2(M_\psi)=0$. The spin
structure is unique. If $L_x$ is the line bundle with Chern
class $x$, the twisted spin index formula used in \ref{push:src:C.2.1} gives
\[
 \operatorname{ind}(D\otimes L_x)
 =\left\langle\frac{x^3}{6}-\frac{p_1x}{24},[M_\psi]\right\rangle
 =-J_\psi\in\Z.
\]
This also verifies the required integrality independently of the
classification of six-manifolds. The identical characteristic
integer is recorded in \citet[Section 3, page 1911]{Olbermann2010D}.

\item For a single framed three-knot, Wall's direct surgery calculation
gives
\[
 k=6\varphi+\beta',\qquad \langle p_1x,[M]\rangle=4\beta',
\]
where $\varphi$ is its unframed knot integer and $\beta'$ its
framing integer
\citep[Theorem 4, pages 360--361]{Wall1966VE}.
Hence $J=-\varphi$. The original pages were inspected directly;
the published corrigendum changes Theorem 14, not this formula
\citep{Wall1967CorrE}. Finally the orbit-knot classification and
\citet[Example 7.14 and Proposition 7.15]{Davis1982D} send a
knot generator to the Milnor sphere $M^7_3=S_{2,-1}(\mathbb H)$.
The associated oriented rank-four bundle over $S^4$ has Euler
evaluation one and first Pontryagin evaluation six. Thus \ref{push:src:C.2.1}
gives $(6^2-4)/896=1/28$ with the corresponding
orientation. Comparing the two choices of knot generator gives
the fixed sign in the statement.

\item The framing kernel has a useful direct verification. Represent
its generator on $P$ by
$(a,y)\mapsto(r(y)a,y)$, with
$r:(D^3,\partial D^3)\to(SO(3),1)$. Lift $r$ to
$R:D^3\to SU(2)$, equal to one near the boundary. In the
quaternionic linear model its equivariant lift is
$(v,y)\mapsto(R(y)v,y)$ on
$\partial(D^4_{\mathbb H}\times D^3)$; the circle acts by
right complex multiplication. This formula extends over the
seven-disc. Its clutching sphere is therefore standard.
The framing direction cannot change the $\Theta_7$ class,
in agreement with the formula $J=-\varphi$.

\end{enumerate}
\end{proof}




\begin{corollary}[the specified cap integer vanishes]\label{push:src:C.3.15}
\StatementRecord{push:src:C.3.15}{proved}{GPT-6 Astra (OpenAI, ChatGPT Pro mode), effort Pro, 2026-11-17; GPT-6 Astra (OpenAI, ChatGPT Pro mode), S6\_06\_BLIND, effort Pro, 2026-11-17; Claude Fable 5.2 (Anthropic), effort high, 16 November 2026}{the hypotheses, notation and definitions stated here; \cref{push:src:C.3.18}; \cref{push:src:C.3.22}; \cref{push:src:C.3.23}; \cref{push:src:C.3.24}}

The characteristic integer of this actual annular cap is
$J_\psi=0$ in $\Z$.
\end{corollary}


\begin{proof}
\begin{enumerate}[label=\arabic*.]
\item The complete evaluation is given below in \ref{push:src:C.3.18}--\ref{push:src:C.3.22}.
The actual affine-frame defect is $p_{\mathrm{rel}}=-24$,
the actual delta-line relative square is $s_\delta=-2$, and
the common signed cap formula is
\[
 J_\psi=\epsilon_{\rm cap}
              (p_{\mathrm{rel}}-12s_\delta)/24=0.
\]
This uses the global frame and Picard-line comparisons, not
only the local unipotent derivative or a target linking number.
The fixed sign comparing Wall's nonzero generator with the
positive Eells--Kuiper generator is immaterial for this zero.
\ref{push:src:C.3.23}--\ref{push:src:C.3.24} give the integer-level argument for the original
one-loop monodromy.

\end{enumerate}
\end{proof}




\begin{proposition}[the actual cusp-link inclusion]\label{push:src:C.3.16}
\StatementRecord{push:src:C.3.16}{proved}{GPT-6 Astra (OpenAI, ChatGPT Pro mode), effort Pro, 2026-11-17; Claude Fable 5.2 (Anthropic), effort high, 16 November 2026}{the hypotheses, notation and definitions stated here; \cref{push:src:C.3.13}; \cref{push:src:C.3.14}; \cref{push:src:C.3.7}}

Orient the central orbit cycle $H=W/S^1_\delta\subset Q=X_c/S^1_\delta$
by the complex $q$-normal orientation, where $q=1/t$ near the cusp.
Use the marked positive period $\gamma$ and the actual transverse
$q$-meridian of \ref{push:src:C.3.13}. Then
\[
 H_1(Q\setminus H;\Z)=\Z,
 \qquad \operatorname{lk}(\gamma,H)=12.
\]
More precisely, write the cusp orbit link as
$L=S^1_\gamma\times\operatorname{Nil}^3(u,w,q)$, with the integral
basis $(\gamma,u,q)$ of $H_1(L;\Z)$ fixed in \ref{push:src:C.3.13}. The inclusion
homomorphism is
\[
 H_1(L;\Z)\longrightarrow H_1(Q\setminus H;\Z),
 \qquad (a,b,c)\longmapsto 12a+c,
\]
where the target generator is the positive meridian of $H$.
This is a statement about the actual cusp inclusion, not only the
weighted target three-sphere.
\end{proposition}


\begin{proof}
\begin{enumerate}[label=\arabic*.]
\item By \ref{push:src:C.3.14}, $Q$ is the standard smooth five-sphere. The finite
normalization pushout in \ref{push:src:C.3.7} makes $H$ a compact, locally
contractible space, and gives $H^3(H;\Z)=\Z$. Its oriented
normalization has degree one at a generic smooth point; identification
along the conductor, of dimension at most two, does not change the
top-dimensional fundamental cycle. Thus $H$ has the multiplicity-one
oriented fundamental cycle specified in the statement. Integral
Alexander duality gives
\[
 \widetilde H_1(Q\setminus H;\Z)
 \simeq \widetilde H^3(H;\Z)=\Z.
\]
The local positive meridian at any generic smooth point pairs to one
with that fundamental cycle and therefore represents the indicated
primitive generator.

\item Let $\overline G:Q\to S^3$ be the continuous descent of the smooth
$S^1_\delta$-invariant map constructed in \ref{push:src:C.3.13}, and write its two
complex coordinate functions as $G_1,G_2$. Define the globally
single-valued continuous complex function
\[
 h=G_2^3-G_1^4.
\]
Since the weighted target map is
$p(z_1,z_2)=z_2^3/(z_2^3-z_1^4)$ and
$p\circ\overline G=\overline f$, the zero set of $h$ is exactly $H$.
In the cusp solid-torus formula of \ref{push:src:C.3.13}, it is explicitly
\[
 h=q\,S(q)^{12}\exp(24\pi i\vartheta),\qquad
 \vartheta=g-\operatorname{Re}\kappa_0(q),
\]
where $S(q)>0$, $g$ is circle-valued, and $\kappa_0$ is a
single-valued holomorphic function on the cusp disc. The factor
multiplying $q$ is a continuous nowhere-zero function. Hence on a
small positive normal circle at a generic smooth point of $H$, the
map $h/|h|$ has degree $+1$. It follows that
\[
 [h/|h|]\in H^1(Q\setminus H;\Z)
\]
is exactly the integral meridian, or linking, class. This argument
requires only continuity of the descended map at fixed points; it
does not assert that an arbitrary invariant smooth map descends
smoothly in every orbit-space coordinate system.

\item On the outer cusp collar the extension $g$ equals the raw marked
character $\theta$. Traversing the positive $\gamma$ period changes
$\theta$ by one, while $q$ and the other torus coordinates stay fixed.
The displayed formula for $h$ consequently has argument degree
$12$ on that loop. A traversal of $u$ changes neither $q$ nor
$\theta$, and has degree zero. Finally, the actual positive
transverse meridian is the section $z_1=z_2=1$ of \ref{push:src:C.3.13}. On it
$\theta=0$ identically, so
\[
 \frac{h}{|h|}
 =\frac{q}{|q|}\exp\bigl(-24\pi i\operatorname{Re}\kappa_0(q)\bigr).
\]
The second factor extends over the whole disc and has degree zero.
The meridian therefore has degree one. These three evaluations give
the asserted inclusion homomorphism. The commutator direction $w$
is already zero in $H_1(L;\Z)$.

\end{enumerate}
\end{proof}


\begin{remark}[Scope]
The numerical value $12$ is now an actual linking invariant in the
orbit five-sphere. The function $q$ alone cannot provide it: globally
$q$ is a meromorphic base coordinate with poles, so its argument
compares the zero and pole fibres. The globally defined function $h$
removes that ambiguity. The present result specifies one part of the
global attaching data for the cap of \ref{push:src:C.3.14}; it does not evaluate its
Pontryagin framing term or its characteristic integer $J_\psi$.
\end{remark}

\begin{proposition}[a marked circle extension over the cusp]\label{push:src:C.3.17}
\StatementRecord{push:src:C.3.17}{proved}{GPT-6 Astra (OpenAI, ChatGPT Pro mode), effort Pro, 2026-11-17; Claude Fable 5.2 (Anthropic), effort high, 16 November 2026}{the hypotheses, notation and definitions stated here; \cref{push:src:C.3.14}; \cref{push:src:C.3.10}}

Let $V$ be a sufficiently small closed cusp body and
$\overline V=V/S^1_\delta$. Its smooth orbit structure can be
chosen compatibly with \ref{push:src:C.3.14} so that there is a diffeomorphism
\[
 \overline V\simeq E_{I_1}\times S^1,
\]
where $E_{I_1}$ is a closed neighbourhood of one nodal fibre in
an elliptic surface. More precisely, the actual positive
$\gamma$-translation circle action on the outer cusp collar
extends to a smooth circle action $\rho$ on $\overline V$.
The extension is free. It is not asserted to preserve the
fixed sphere $F\subset\overline V$ or the original interior
map to the $q$-disc.

Consequently the actual annular map $\psi$ of \ref{push:src:C.3.10} is the
endpoint of an ambient isotopy of $Q$ whose restriction to $F$
is the loop of embedded spheres $t\mapsto\rho_t(F)$.
After the marked deformation used in the proof, its trace is
the normalization immersion of a nodal elliptic curve times a
circle: it has one transverse double circle and two distinct
preimage circles. The regular cusp three-torus with period
basis $(\gamma,u,w)$ represents a primitive generator of
$H_3(\overline V;\Z)$; in particular it does not bound a
four-chain inside the cusp body.
\end{proposition}


\begin{proof}
\begin{enumerate}[label=\arabic*.]
\item We first remove the bounded period units by a marked smooth
deformation. On the cusp write, as in Paper~I, Proposition~18.5,
\[
 Z(q)=sB_0+C(q),\qquad
 B_0=\begin{pmatrix}0&1\\-1&0\end{pmatrix},\qquad
 q=e^{2\pi i s}.
\]
For $0\leq a\leq1$ replace $C(q)$ by $(1-a)C(q)$, keeping
the same $A_2$ triangulation and the real periods $w,\delta$.
This is a deformation of the actual toric period quotient,
not a replacement inferred from its homology. The two period
automorphisms still have the same integral translations of the
fan and commute; only their nowhere-zero holomorphic unit
multipliers change. They depend smoothly on $a$, and even
holomorphically on a small complex neighbourhood of this
interval. On a sufficiently small punctured $q$-disc the
imaginary period matrices are invertible uniformly in $a$:
their leading part is $(\operatorname{Im}s)B_0$, while the
remaining matrices are uniformly bounded.

\item The fixed toric charts, with this extra parameter, therefore
give a smooth family of the same Mumford cusp quotients.
There is no new stabilizer at the central fibre: a nonzero
period translates every height-one ray by a nonzero lattice
vector, and cannot preserve any finite cone of the fan.
The usual finite fundamental collection of toric charts gives
proper discontinuity, uniformly for the compact parameter
interval. The central family has the constant compact
normalization $dP_6$ with its smoothly varying opposite-side
identifications. Thus it has a proper representative over a
uniform smaller $q$-disc. One may see the uniformity directly
by taking a relatively compact neighbourhood of this compact
central family: its boundary is disjoint from $q=0$, and after
shrinking $|q|$ no fibre meets that boundary. The nonzero
fibres are connected compact tori, so the selected representative
contains the whole nearby fibre. Denote these compact smooth
cusp bodies by $V_a$.

\item Their $S^1_\delta$ actions form a smooth semifree family, with
a smooth family of fixed two-spheres. Choose an equivariant
Ehresmann transport with the following two prescribed local
forms. On the outer collar keep $q$ and all real coefficients
in the marked period basis fixed. This is a smooth transport
there and carries the actual $\gamma$-translation action to
the corresponding action for every $a$. Near the fixed spheres,
choose family equivariant normal tubular coordinates. After
conjugating the negative-weight line the normal representation
is scalar multiplication on $\mathbb C^2$. Choose Hermitian
metrics and a unitary connection, and make the transport
linear and unitary in these normal coordinates. The collar
and fixed tubes are disjoint. A circle-invariant partition of
unity extends and patches the prescribed horizontal fields;
each has derivative one on $a$, so their patched field does
also. Compactness gives its flow for the entire interval.

\item This particular transport induces a smooth transport of orbit
spaces. On the free locus this follows from the principal-bundle
quotient. In the fixed tubes, unitary normal transport induces
the usual smooth linear rotation in the Hopf $\mathbb R^3$
coordinates. Thus it is smooth there as well. This explicit
choice avoids any assertion that every equivariant smooth map
has a smooth map in an arbitrarily preassigned quotient atlas.
The Hopf smooth structures and tubular choices are legitimate
ones for the orbit construction in \ref{push:src:C.3.14}; the annular map is
supported in the free collar, where these structures coincide.
We obtain a marked orbit diffeomorphism
\[
 (\overline V,\text{boundary }\gamma\text{-action})
 \simeq(\overline V_1,\text{boundary }\gamma\text{-action}).
\]
The transport is not required to preserve $q$ in the interior.

\item At $a=1$ the period matrix is $sB_0$. The complex torus is
the product of two Tate elliptic curves, and the square toric
model is the fibre product of their nodal families. The chosen
$A_2$ triangulation is its chosen small resolution. Project
this resolution to the first Tate family $E_{I_1}$, and let
$S^1_\delta$ act on the second factor. The projection descends
to a proper smooth submersion
\[
 p:\overline V_1\longrightarrow E_{I_1}
\]
with oriented circle fibres. Here is the local check at the
only place where this assertion is not immediate. Write the
fibre product as $xy=uv$, with $q=xy=uv$, and take the small
resolution chart
\[
 u=xA,\qquad y=Av,
\]
whose coordinates are $(x,A,v)$. The weights of $(A,v)$ are
$(+1,-1)$ and $x$ has weight zero. Its Hopf quotient coordinates
are
\[
 \bigl(x,y=Av,d=(|A|^2-|v|^2)/2\bigr),
\]
and the descended projection is simply $(x,y,d)\mapsto(x,y)$.
The second chart, $x=uB$, $v=By$, gives the identical check
with $d'=(|u|^2-|B|^2)/2$. The local orbit fibre at $x=y=0$
consists of the two orbit rays joined by the quotient interval
of the exceptional projective line, and is a smooth real line.
Globally it closes to the second Tate orbit circle. At fixed
points away from this chart the model is likewise
$(p,A,v)\mapsto(p,Av)$, whose quotient is the same smooth
projection. We make the compatibility of these quotient charts
explicit. On overlaps near the same fixed branch the two weight
coordinates have the form
\[
 A_j=a_{ij}(p,z)A_i,\qquad v_j=b_{ij}(p,z)v_i,
 \qquad z=A_iv_i,
\]
where the units and the change of the tangential coordinate
$p$ are pulled back from the projection base. The two fixed
lifts over the node have disjoint fixed tubular neighbourhoods;
they require no overlap identification at the node itself.
Choose Hermitian metrics for the two weight line bundles on
these base charts, retaining the separated lifts of the fixed
normalization. If their positive local coefficients are
$h_{A,i},h_{v,i}$, their transition rules are
\[
 h_{A,j}|a_{ij}|^2=h_{A,i},\qquad
 h_{v,j}|b_{ij}|^2=h_{v,i}.
\]
The weighted Hopf coordinate
\[
 d_i=\tfrac12\bigl(h_{A,i}|A_i|^2-h_{v,i}|v_i|^2\bigr)
\]
therefore agrees with $d_j$ on overlaps. These are legitimate
Hopf coordinates: rescale the two normal variables by
$\sqrt{h_{A,i}}$ and $\sqrt{h_{v,i}}$. Their product is
$\sqrt{h_{A,i}h_{v,i}}\,z$, a smooth local change of the
complex coordinate $z$ with invertible derivative at $z=0$.
Thus $(p,z,d_i)$ are smooth quotient coordinates, their
overlap maps are smooth, and the projection remains
$(p,z,d_i)\mapsto(p,z)$. This specifies the Hopf smoothing
used for the transport above, without a general smoothness
claim for arbitrary equivariant quotient maps.

\item The section of the second Tate family represented by $z_2=1$
avoids its node. It therefore gives a section of the fibre
product which lifts through the small resolution, and then
gives a section of $p$. Hence this oriented circle bundle is
trivial. This proves the asserted product diffeomorphism by
an actual projection and section.

\item We can match the boundary action exactly, not just its homology
class. On the boundary of this circle bundle the positive
$\gamma$ action is the circle translation in the second Tate
orbit fibre. Its invariant vertical length form has fibre
integral one. Extend that length form to any smooth positive
vertical length form on the entire oriented circle bundle,
using a collar cutoff. Divide on each fibre by its integral;
this is a smooth positive function of the base because $p$
is a proper submersion. The dual vertical vector field now has
period exactly one, and agrees with the specified vector field
on the boundary collar. Its flow defines the desired free
circle action. Transport it back to $\overline V$ by the marked
orbit diffeomorphism.

\item Choose an invariant function $\chi$ equal to one in the inner
body and zero outside the outer collar, with the same collar
profile as in \ref{push:src:C.3.10}. Such a function can be chosen because
the action agrees with the actual $q$-fibre translations there;
extend it constantly by one on the remaining interior. The
maps
\[
 H_t(z)=\rho_{t\chi(z)}(z),\qquad 0\leq t\leq1,
\]
are diffeomorphisms, since $\chi$ is invariant. They extend
by the identity to $Q$. Their endpoint is the actual annular
map: on the inner body it is $\rho_1=\id$, and on the collar
it is translation by $\chi\gamma$. On the fixed sphere
$\chi=1$, so this isotopy traces precisely $\rho_t(F)$.
It is not an isotopy relative to the fixed sphere.

\item In the zero-unit model the fixed sphere projects to the
normalization of the first nodal elliptic curve. Over a smooth
point of that curve it meets the orbit circle once, while the
two preimages of its node give two distinct points of the
orbit circle over the node. Its circle sweep is therefore the
normalization immersion of the nodal curve times the circle.
The local form $(x,y,d)$ with $q=xy$ gives its two transverse
branches at the double circle. Finally $E_{I_1}$ retracts onto
its nodal elliptic curve, whose fundamental two-cycle is the
degree-one specialization of a regular elliptic fibre. The
K\"unneth theorem now gives
\[
 H_3(E_{I_1}\times S^1;\Z)
 =H_2(E_{I_1};\Z)\otimes H_1(S^1;\Z)=\Z,
\]
generated by that regular fibre times the marked circle.
Boundary marking transports this three-torus to the actual
$(\gamma,u,w)$ torus, proving the last assertion.

\end{enumerate}
\end{proof}


\begin{remark}[Scope]
The raw fractional period translations need not extend across
the central fibre. The action above modifies their interior
orbits. Its loop moves the fixed sphere, which is exactly the
feature retained by the annular relative mapping class. Thus
the construction does not trivialize that class or evaluate
$J_\psi$. In particular the nonbounding three-torus prevents
discarding its stable framing contribution by a chain lying
entirely in the cusp body.
\end{remark}

\begin{proposition}[an actual bounding four-manifold and its frame defect]\label{push:src:C.3.18}
\StatementRecord{push:src:C.3.18}{proved}{GPT-6 Astra (OpenAI, ChatGPT Pro mode), effort Pro, 2026-11-17; Claude Fable 5.2 (Anthropic), effort high, 16 November 2026}{the hypotheses, notation and definitions stated here; \cref{push:src:C.3.13}; \cref{push:src:C.3.17}}

There is a smooth genus-one Lefschetz fibration
\[
 \mathfrak g:Q\simeq S^5\longrightarrow S^3
\]
whose critical set is a circle mapping bijectively to the positive
$(3,4)$ knot
\[
 K=\{z_2^3-z_1^4=0\}\subset S^3.
\]
It agrees on the exterior and the outer cusp collar with the
actual torus quotient of \ref{push:src:C.3.13}. In particular, the inverse image
of a nearby positive weighted circle $K'$ is the actual marked
three-torus with periods $(\gamma,u,w)$.

This three-torus bounds an embedded compact four-manifold
$M\subset Q$ carrying a positive genus-one Lefschetz fibration
over a Seifert surface of $K'$, with exactly twelve nodal critical
points. Its Euler number is $12$. Give $M$ its Lefschetz complex
orientation and its boundary the induced orientation, which in
the marked period convention is $(\gamma,w,u)$. Let
$\mathfrak f_{\mathrm{aff}}$ be the stable boundary frame consisting
of the three period directions and the two real $q$-normal
directions, with a fixed ordering consistent with the ambient
orientation. Relative to this actual affine frame,
\[
 \left\langle p_1(TM\oplus\R,\mathfrak f_{\mathrm{aff}}),
                   [M,\partial M]\right\rangle=-24.
\]
Equivalently, its transgression is $+24$ when the boundary
three-torus is oriented by the ordered basis $(\gamma,u,w)$.
No value of $\langle p_1(M_\psi)x,[M_\psi]\rangle$ or of
$J_\psi$ is inferred without the cap transgression calculation.
\end{proposition}


\begin{proof}
\begin{enumerate}[label=\arabic*.]
\item On the complement of the cusp body, the descended map of \ref{push:src:C.3.13}
is a smooth two-torus bundle: on an ordinary chart it is the
quotient by the two residual real period directions $u,w$,
and the finite root charts give the same submersion after their
free finite quotients. The explicit finite gauges in \ref{push:src:C.3.13}
ensure that this includes the two exceptional weighted circles
in the target.

\item Use the marked product of \ref{push:src:C.3.17} on the cusp body. In its
zero-unit model the projection is an elliptic $I_1$ neighbourhood
times a circle. Choose the product circle coordinate by the
second Tate identity section $z_2=1$ and the extending circle
action. On the boundary this coordinate is exactly the marked
character $\theta$: the section has value zero and the boundary
action is the actual $\gamma$ translation. The prescribed period
transport in \ref{push:src:C.3.17} preserves all these coordinates on the
outer collar. Consequently the map
\[
 (x,\theta)\longmapsto
 \Phi_0\bigl(q_E(x),\theta-\operatorname{Re}\kappa_0(q_E(x))\bigr)
\]
agrees there with the exterior map, where $q_E:E_{I_1}\to\Delta$
is the nodal elliptic map and $\Phi_0,\kappa_0$ are the exact
solid-torus coordinates of \ref{push:src:C.3.13}. The two maps glue smoothly.
Near the critical circle, after the indicated smooth target
coordinate change, their local expression is
\[
 (x,y,\theta)\longmapsto(xy,\theta).
\]
Thus there is one positive nodal degeneration over each point
of $K$, and no other critical point. This constructs the stated
map from the actual marked collar. It need not equal the old
map $\overline f$ in the interior of the cusp body.

\item Choose $0<\epsilon\ll1$ and put
\[
 K'=\{(1-\epsilon)z_2^3-z_1^4=0\}\subset S^3,
 \qquad F_\epsilon=(1-\epsilon)z_2^3-z_1^4.
\]
Then $K'$ is the positive weighted circle $q=\epsilon$ in
the actual collar coordinates. The weighted circle field
$(z_1,z_2)\mapsto(e^{6\pi it}z_1,e^{8\pi it}z_2)$ satisfies
\[
 F_\epsilon(t\cdot z)=e^{24\pi it}F_\epsilon(z).
\]
Hence $\arg F_\epsilon$ is a submersion off $K'$. The closure
of one of its pages is a smooth compact oriented Seifert
surface $\Sigma'$ with boundary $K'$: the polynomial has a
simple complex normal zero along the link, which gives the
usual smooth half-plane collar of a page. Its fibres are
connected, since the normal meridian has argument degree one.
On the critical knot $K$,
\[
 F_\epsilon=-\epsilon z_2^3.
\]
Its argument increases twelve times along the positive weighted
parameter. Every page therefore meets $K$ transversely in
exactly twelve points, all with positive local intersection.

\item Set $M=\mathfrak g^{-1}(\Sigma')$. The displayed local nodal
model and transversality of $\Sigma'$ to the critical circle
show that $M$ is a smooth embedded compact four-manifold with
boundary $\mathfrak g^{-1}(K')$. The restriction to $M$ is a
positive genus-one Lefschetz fibration over $\Sigma'$ with
twelve critical points. A smooth torus has Euler number zero,
and an $I_1$ fibre has Euler number one. Stratification over
these twelve values therefore gives $\chi(M)=12$.

\item We next specify the boundary framing whose characteristic number
is computed. Choose a fibrewise complex structure on the regular
two-tori, compatible with their affine markings, and the usual
complex structure in the nodal charts. Such choices glue: on
the intervening collars the marked elliptic upper half-plane
is contractible. Choose a complex structure on $\Sigma'$ and
a compatible horizontal distribution. These give an almost
complex structure on $M$, equal to the standard nodal structure
near its critical points. Along $\partial M$, its vertical
complex tangent is framed by the translation vector dual to
the elliptic one-form with $w$-period one. Frame its horizontal
complex line by a lift of the outward normal of $\Sigma'$;
the associated real imaginary direction is the positive
boundary tangent. Denote this complex frame by $\mathfrak f_C$.

\item For completeness, the relative determinant class is pulled
back from the base. On a regular fibre its nonzero holomorphic
one-forms form a line $\lambda$, and evaluation identifies
the relative canonical line with the pullback of $\lambda$.
This remains true through each nodal chart: for $t=xy$ the
relative form is $dx/x=-dy/y$, and wedging with $dt$ gives
a nowhere-zero multiple of $dx\wedge dy$. Thus, as complex
line bundles,
\[
 \det(TM,J)\simeq\pi^*(T\Sigma'\otimes\lambda^{-1}).
\]
The determinant trivialization supplied by $\mathfrak f_C$
is the pullback of the corresponding base trivialization along
$\partial\Sigma'$. Therefore for some
$\beta\in H^2(\Sigma',\partial\Sigma';\Z)$,
\[
 c_1(TM,\mathfrak f_C)=\pi^*\beta,
 \qquad c_1(TM,\mathfrak f_C)^2=0.
\]
This does not require assigning a degree to $\lambda$ or a
holomorphic structure to the entire four-manifold. Its local
evaluation argument concerns the fibrewise elliptic structures
and the positive nodal charts.

\item The horizontal vector in $\mathfrak f_C$ points outwards on
the boundary. The relative top Chern class is therefore the
relative Euler class for an outward-pointing vector, and
Poincar\'e--Hopf gives
\[
 \langle c_2(TM,\mathfrak f_C),[M,\partial M]\rangle
 =\chi(M)=12.
\]
The relative characteristic-class identity now gives
\[
 \langle p_1(TM,\mathfrak f_C),[M,\partial M]\rangle
 =\langle c_1^2-2c_2,[M,\partial M]\rangle=-24.
\]

\item We finally compare this frame with the actual annular one.
On the target collar the exact formulas of \ref{push:src:C.3.13} give
\[
 F_\epsilon=(q-\epsilon)S(q)^{12}
              \exp(24\pi i\vartheta),
 \qquad\vartheta=\theta-\operatorname{Re}\kappa_0(q).
\]
Along $q=\epsilon$, the inward normal ray of the chosen page
therefore has $q$-angle equal to a constant minus $12\theta$.
Its complementary normal line in $S^3$ rotates by the same
angle. Add the pullback of that normal real line to $TM$,
identifying $TM\oplus\R$ with $TQ|_M$ after a horizontal
choice. On the boundary, the complex vertical translation
frame and the positive $\gamma$ tangent are the actual marked
period frame. The two remaining real vectors differ from
the fixed real and imaginary $q$ directions by a rotation
of winding $-12$ as $\gamma$ traverses the circle. Changes
of horizontal lifts add only contractible shear matrices.
Thus the stable frame comparison, up to fixed ordering, is
\[
 \operatorname{diag}\bigl(I_3,R_{-24\pi\theta}\bigr):
 T^3\longrightarrow SO(5).
\]
It factors through the $\gamma$ circle and is nullhomotopic,
since $\pi_1(SO(5))=\Z/2$ and twelve is even. Consequently
the relative first Pontryagin number is unchanged when
$\mathfrak f_C\oplus1$ is replaced by
$\mathfrak f_{\mathrm{aff}}$.

\item For the orientation convention, the elliptic complex orientation
is $(w,u)$, because its marked periods are $1,\tau$ with
$\operatorname{Im}\tau>0$. The outward-first boundary convention
for the Lefschetz orientation therefore orients the boundary
by $(\gamma,w,u)$. Reversing this to $(\gamma,u,w)$ reverses
the transgression sign, giving the equivalent $+24$ convention
stated above.

\end{enumerate}
\end{proof}


\begin{remark}[Scope]
The four-manifold leaves the cusp body; its existence is
consistent with the local nonbounding assertion in \ref{push:src:C.3.17}.
The value computed here is a relative Pontryagin defect of
the specified affine frame, not yet the Pontryagin pairing
of the closed cap in \ref{push:src:C.3.14}. That latter conversion must
retain the actual gluing and its line-bundle transition.
\end{remark}

\begin{proposition}[the cap Pontryagin pairing]\label{push:src:C.3.19}
\StatementRecord{push:src:C.3.19}{proved}{GPT-6 Astra (OpenAI, ChatGPT Pro mode), effort Pro, 2026-11-17; Claude Fable 5.2 (Anthropic), effort high, 16 November 2026}{the hypotheses, notation and definitions stated here; \cref{push:src:C.3.14}; \cref{push:src:C.3.18}}

For the actual cap $M_\psi$ and meridian class $x$ of \ref{push:src:C.3.14},
\[
 \bigl|\langle p_1(M_\psi)x,[M_\psi]\rangle\bigr|=24.
\]
Thus the framing coordinate $\beta'$ in Wall's convention has
$|\beta'|=6$. The orientation sign and the cubic
$k_\psi=\langle x^3,[M_\psi]\rangle$ are not determined by
this proposition, and no numerical $J_\psi$ or monodromy class
is asserted.
\end{proposition}


\begin{proof}
\begin{enumerate}[label=\arabic*.]
\item We use the cohomology suspension of the actual universal class
$p_1$, denoted $s(p_1)\in H^3(SO;\Z)$, with the clutching
convention fixed once throughout. In particular we do not
replace it by an unspecified primitive generator of
$H^3(SO;\Z)$ or of $H^3(\operatorname{Spin};\Z)$.
All computations below may first be made in real cohomology;
they determine the integral characteristic number, and torsion
cannot contribute to its pairing.

\item Write $N=S^2\times D^4$ and $Y=\partial N=S^2\times S^3$,
so that the cap is $N\cup_{\widetilde\psi}(-N)$. Choose a
stable tangent frame on $N$. Its restriction supplies a stable
frame on the outer copy $P=S^2\times D^3\subset Y$. The
stable frame on $Q=S^5$, restricted to $P$, is homotopic to
this one: the difference is a map $P\to SO$, and
$P\simeq S^2$ while $\pi_2(SO)=0$. Thus the frame comparison
period obtained in \ref{push:src:C.3.18} is also the one computed from the
chosen frame on $N$.

\item Let $\mathcal A=L\times I\subset P$ be the annular support,
where
\[
 L=S^1_\gamma\times\operatorname{Nil}^3(u,w,q).
\]
Use a global affine nilmanifold frame on $\mathcal A$.
The nontrivial $q$-monodromy in $u,w$ is accounted for by
the invariant nilmanifold frame; on any fixed $q$-angle
three-torus it differs from the marked period frame only
by a constant triangular change. Translation in $\gamma$
preserves this global frame. Let
$A:\mathcal A\to SO$ be the stable change-of-frame map
from this frame to the restricted global frame.

\item In the affine frame, the derivative of the annular map
$(r,p)\mapsto(r,\rho_{\chi(r)}p)$ is the identity plus
the single shear $\chi'(r)\,dr\otimes\partial_\gamma$.
Multiplying the shear by a parameter from one to zero
contracts it to the identity through invertible matrices,
fixed on both end collars. Therefore the actual derivative
clutching map is homotopic, relative to the complement of
$\mathcal A$, to
\[
 G(r,p)=A(r,\rho_{\chi(r)}p)^{-1}A(r,p).
\]
At both ends $G=1$, since the translation is respectively
zero or a full period. Replacing $A(r,p)$ by its restriction
to one middle slice gives a homotopy of $G$ relative to
these ends. We may consequently write $A=A(p)$.

\item The Mayer--Vietoris clutching formula gives
\[
 p_1(M_\psi)=\delta\bigl(G^*s(p_1)\bigr),
\]
where $G$ is extended by the identity on the rest of $Y$.
With the boundary orientation associated to this connecting
map, its cup-product pairing is exactly
\[
 \langle p_1(M_\psi)x,[M_\psi]\rangle
 =\langle G^*s(p_1)\smile x|_Y,[Y]\rangle.
\]
Equivalently this is the relative pairing on
$(\mathcal A,\partial\mathcal A)$. Changing a global
clutching orientation reverses both associated signs;
we will retain the absolute value in the numerical conclusion.

\item We calculate this relative class explicitly. Define
\[
 \Gamma:S^1\times L\longrightarrow SO,
 \qquad\Gamma(t,p)=A(\rho_t p)^{-1}A(p).
\]
It is the identity on $\{0\}\times L$. The suspension
$s(p_1)$ is primitive under multiplication. Slanting its
pullback by the oriented $t$-circle therefore gives
\[
 \int_{S^1}\Gamma^*s(p_1)
 =-\int_{S^1}\rho^*\bigl(A^*s(p_1)\bigr)
 \quad\hbox{in }H^2(L;\R).
\]
There is no additional relative lift ambiguity here:
restriction $H^2(S^1\times L)\to H^2(\{0\}\times L)$
is surjective. Thus $H^3(S^1\times L,\{0\}\times L)$
is exactly the degree-two slant summand. Pulling back by
$t=\chi(r)$ identifies the relative annular class with
the suspension of the right-hand side. This is the
relative trace calculation, including the full-period
endpoint; an absolute equality on $\mathcal A$ alone
would not suffice.

\item Put $g=\theta_\gamma$ and $a=\arg q/(2\pi)$. The meridian
Euler class has, up to the already fixed overall circle
orientation, the representative
\[
 e=x|_L=dg\wedge da.
\]
Indeed the positive $q$ period has
$A_0\gamma=\gamma-\delta$ and $A_0u=u+w$.
Using coefficient coordinates on the covering mapping
torus, its deck identification can be written
\[
 (g,u,w,d,a)\longmapsto(g,u,w-u,d+g,a+1),
\]
where $d$ is the $\delta$ coordinate. The normalized
connection $\mathrm d d-a\,\mathrm d g$ is invariant and has curvature
$dg\wedge da$. This also shows that no additional
component of the Euler class has been suppressed.

\item Decompose the real cohomology class
$\alpha=A^*s(p_1)$ using $L=S^1_g\times\operatorname{Nil}^3$:
\[
 \alpha=dg\smile\beta+\alpha_0,
 \qquad\beta\in H^2(\operatorname{Nil}^3;\R),\quad
 \alpha_0\in H^3(\operatorname{Nil}^3;\R).
\]
Translation in $g$ has trace $\beta$. Therefore the preceding
relative calculation and the cup pairing reduce, with a
fixed orientation sign, to
\[
 \pm\int_L dg\wedge da\wedge\beta
 =\pm\int_{T^2_{u,w}}\beta
 =\pm\int_{T^3_{\gamma,u,w}}A^*s(p_1).
\]
The first equality on the right is integration over the
mapping-torus base, with its fibre orientation. Components
of $\beta$ containing $da$, and the component $\alpha_0$,
give zero. This proves that there are no unaccounted
cross terms in the cap pairing.

\item Finally, the last integral is the relative stable-frame
Pontryagin defect of \ref{push:src:C.3.18}, with at most the inverse
change-of-frame and orientation signs. Its absolute value
is $24$. Hence the cap pairing has absolute value $24$.
Wall's exact formula $\langle p_1x,[M_\psi]\rangle=4\beta'$
from \ref{push:src:C.3.14} gives $|\beta'|=6$, as asserted.

\end{enumerate}
\end{proof}


\begin{remark}[Scope]
The argument uses both the global frame defect furnished by
the bounding four-manifold and the actual Euler class of the
$\delta$ bundle on the annulus. A local unipotent derivative
alone would give neither. The cubic pairing and its orientation
relative to this Pontryagin pairing remain necessary to
evaluate Wall's unframed knot coordinate.
\end{remark}

\begin{proposition}[the cubic term as an actual relative square]\label{push:src:C.3.20}
\StatementRecord{push:src:C.3.20}{proved}{GPT-6 Astra (OpenAI, ChatGPT Pro mode), effort Pro, 2026-11-17; Claude Fable 5.2 (Anthropic), effort high, 16 November 2026}{the hypotheses, notation and definitions stated here; \cref{push:src:C.3.18}; \cref{push:src:C.3.14}; \cref{push:src:C.3.17}; \cref{push:src:C.3.16}; \citet[Section 3, Propositions 3.2 and 3.8]{ChernSimons1974E}}

The four-manifold in \ref{push:src:C.3.18} can be chosen disjoint from the fixed
sphere. Write it as $M\subset P$, retain its Lefschetz orientation,
and let $\mathcal L$ be the actual principal $\delta$ circle bundle
on $P$, regarded as a complex line bundle. Its restriction to
$\partial M=T_{\gamma uw}$ is trivial. Put
\[
 s_\delta=\langle c_1(\mathcal L|_M,\mathrm{triv})^2,
                         [M,\partial M]\rangle.
\]
This integer is independent of the choice of boundary trivialization.
For the cap of \ref{push:src:C.3.14}, there is one overall sign $\epsilon_{\rm cap}$,
coming from its comparison with the oriented cut-and-reglue trace,
such that
\[
 k_\psi=\epsilon_{\rm cap}\,3s_\delta.
\]
The same relative square, with the compatible Hopf-meridian
orientation, is the linking coefficient of the marked
$\delta=0$ lift of $T_{\gamma uw}$ in $X\setminus F\simeq S^3$.
No numerical value of $s_\delta$ is assumed in this statement.
\end{proposition}


\begin{proof}
\begin{enumerate}[label=\arabic*.]
\item First we arrange the claimed disjointness. In the marked product
of \ref{push:src:C.3.17}, the section of the second Tate family avoids $F$.
Consequently the product circle coordinate of the compact sphere
$F$ has image in a compact proper arc of $S^1\setminus\{0\}$.
An orientation-preserving reparametrization of this circle can
compress that arc to length less than $1/12$. Perform this
reparametrization over a neighbourhood of the central first
elliptic fibre, interpolate through circle diffeomorphisms to the
identity, and leave the outer collar fixed. This changes only the
interior circle coordinate and the extending action, not the
actual annular endpoint $\psi$ or its boundary marking.

\item Rebuild the map of \ref{push:src:C.3.18} with this coordinate. The twelve
intersection parameters of one Milnor page with the critical
knot are equally spaced in the product circle. Choose the phase
of the page so that all twelve avoid the compressed arc.
Since $F$ lies over the central elliptic fibre, the inverse image
$M$ of this page is now disjoint from $F$. It has the same
boundary, twelve positive nodes, orientation and affine frame
as in \ref{push:src:C.3.18}. Shrink the deleted fixed tube so that $M\subset P$.

\item Let $g,u,w$ denote the marked circle coordinates on a positive
real regular cusp fibre, let $a$ be the positive argument of
$q$, and let $r=|q|$. Coefficient coordinates have deck rule
\[
 (g,u,w,d,a)\longmapsto(g,u,w-u,d+g,a+1).
\]
Thus the invariant normalized circle connection is
$\alpha=dd-a\,dg$, with Euler curvature
\[
 E=dg\wedge da.
\]
Choose a global connection on $\mathcal L$ agreeing with this
one on the annular collar. This is possible because the displayed
connection is on the actual restricted bundle. Also choose its
curvature near $\partial P$ to be the pullback of the Hopf
curvature on the normal two-sphere. Then $E^2=0$ near that
boundary. Since $H^4(P)=H^3(P)=0$, there is a three-form $B$
with $dB=E^2$, unique up to an exact form. Because
$H^3(\partial P)=0$, it can moreover be chosen zero on a
boundary collar.

\item The line bundle on the positive real torus is trivialized by
the marked $d=0$ section. On that torus the displayed connection
is zero in this section. Hence relative Chern--Weil theory and
Stokes' theorem give, with $\partial M=-T_{\gamma uw}$,
\[
 s_\delta=\int_M E^2
          =-\int_{T_{\gamma uw}}B.
\]
Changing the trivialization changes the relative Chern class by
the connecting image of $H^1(\partial M;\Z)$. Such a class is
in the radical of the relative intersection pairing: its
pairing with a relative class is the boundary pairing with the
restriction of that class's absolute image, which is zero here.
Products of two connecting classes are zero as well. Thus the
square is independent of the boundary gauge.

\item We calculate the cut-and-reglue variation of $x^3$. Double $P$
to $B_5=DP=\partial N$, where $N=S^2\times D^4$, and extend
$\psi$ and all variations by the identity on the second half.
Here is an explicit seven-dimensional bordism used for this
comparison. Form the $B_5$-bundle over an annulus with monodromy
$\widetilde\psi$. Divide its outer boundary circle into four
successive arcs. Attach two copies of $N\times[0,1]$ along the
$B_5\times[0,1]$ faces over opposite arcs, reversing the interval
parametrization on the second attachment. Choose the
trivializations over these two arcs so that transport over one
remaining arc is the identity and transport over the other is
$\widetilde\psi$. After rounding corners, those remaining arcs
join the endpoint copies of $N$ into, respectively, the identity
double and the $\widetilde\psi$-twisted double. The inner boundary
is the mapping torus. Thus the three closed boundary components
are the two actual caps and the mapping torus, with the two caps
oppositely oriented. Their comparison has a single orientation
sign fixed by these conventions, independent of the map.

\item The meridian line extends over this actual bordism: use its
specified lift of $\widetilde\psi$ on the annular bundle and
the pullback of its line on $N$ on each attached $N\times[0,1]$.
The attaching identifications are precisely the chosen line
identifications, so there is no additional clutching datum.
The tangent bundle is the tangent bundle of this smooth
seven-manifold; on each boundary it restricts stably to that
boundary's tangent bundle. Consequently the same bordism
compares $x^3$ and $p_1x$, with the same orientation sign.
The identity cap is $S^2\times S^4$, and the identity mapping
torus has a line pulled back from $B_5$; both characteristic
pairings vanish on both identity manifolds. Equivalently one
may subtract the identity mapping torus, which bounds
$B_5\times D^2$ with its pulled-back line. Thus the desired
cubic is the mapping-torus variation, up to this one common
boundary-orientation sign.

\item On the annulus $\psi:g\mapsto g+\chi(r)$, where $\chi$ decreases
from one to zero. Consequently
\[
 \Delta=\psi^*E-E=d\chi\wedge da.
\]
There is a global one-form $\eta$ with $d\eta=\Delta$ whose
annular value is
\[
 \eta=-(1-\chi)\,da-12\,dg.
\]
Indeed the outer value $-(da+12\,dg)$ extends as a closed form
over the exterior by \ref{push:src:C.3.16}, while the inner value $-12\,dg$
extends over the product cusp by \ref{push:src:C.3.17}. Closed representatives
can be made to agree exactly on collars by subtracting exact
forms. An exact correction makes $\eta$ zero near $\partial P$,
since this boundary has $H^1=0$. This correction does not affect
any integral below, as $B$ vanishes there.

\item We use the normalized characteristic-form transgression of
\citet[Section 3, Propositions 3.2 and 3.8]{ChernSimons1974E}.
The interpolating curvature on the mapping cylinder is
$E+t\Delta+dt\wedge\eta$. Its endpoint connections can be
matched to the chosen clutching lift by an exact gauge correction,
because $H^1(B_5;\R)=0$. The transgression computation gives
\[
 \int(E+t\Delta+dt\wedge\eta)^3
 =\int_P\eta\wedge(3E^2+3E\Delta+\Delta^2)
 =3\int_P\eta\wedge E^2.
\]
Here $E\Delta=\Delta^2=0$ on the support of $\Delta$. Since
$d(\eta\wedge B)=\Delta\wedge B-\eta\wedge E^2$,
the last expression is $3\int_P\Delta\wedge B$.
Use the orbit orientation
$dg\wedge du\wedge dw\wedge dr\wedge da$, obtained by
taking the positive $\delta$ direction first in the zero-unit
complex model. Integration over the annulus, with $\chi$ going
from one to zero, now gives
\[
 \int_P d\chi\wedge da\wedge B
   =-\int_{T_{\gamma uw}}B=s_\delta.
\]
This proves the asserted factor of three and fixes it without
using a primitive-generator normalization in $H^3(SO)$.

\item Finally on the total circle bundle over $P$ the form
\[
 \Omega=\alpha\wedge\pi^*E-\pi^*B
\]
is closed. Its integral on the normal Hopf three-sphere is one:
the first term is the normalized Hopf integral and the second
vanishes in the chosen fixed collar. It therefore represents
the integral generator of $H^3(X\setminus F;\Z)$. On the marked
lift of the positive real $T_{\gamma uw}$, the connection and
Euler curvature restrict to zero, and its integral is
$-\int_T B=s_\delta$. This is precisely the claimed linking
coefficient. In particular the square is an invariant of the
actual inclusion, not merely of the abstract genus-one surface.

\end{enumerate}
\end{proof}


\begin{remark}[Scope]
To evaluate this number by a height calculation one must identify
the actual line bundle $\mathcal L|_M$, including its mixed
base--fibre class, with the Poincar\'e line attached to the
original normal-function section. The numerical square of an
auxiliary split or torsor surface does not make that
identification. The signed cancellation with the Pontryagin
term likewise uses the same cut-and-reglue orientation in both
calculations.
\end{remark}

\begin{proposition}[the actual delta-line square]\label{push:src:C.3.21}
\StatementRecord{push:src:C.3.21}{proved}{GPT-6 Astra (OpenAI, ChatGPT Pro mode), effort Pro, 2026-11-17; Claude Fable 5.2 (Anthropic), effort high, 2026-11-16; GPT-6 Astra (OpenAI, ChatGPT Pro mode), S6\_06\_BLIND, effort Pro, 2026-11-17; Claude Fable 5.2 (Anthropic), effort high, 16 November 2026}{the hypotheses, notation and definitions stated here; \cref{push:src:C.3.20}; \cref{push:src:C.3.18}}

For the actual four-manifold and line bundle in \ref{push:src:C.3.20},
\[
 s_\delta=-2.
\]
The comparison is with the actual relative Picard section of
the circle bundle; it does not assume that the compact-core
model supplies the excluded analytic filling of the original six-sphere.
\end{proposition}


\begin{proof}
\begin{enumerate}[label=\arabic*.]
\item We first identify the base precisely. On the Milnor page used in
\ref{push:src:C.3.18}, choose its phase so that
$F_\epsilon=(1-\epsilon)z_2^3-z_1^4$ is positive real.
Positive weighted rescaling identifies its interior with
\[
 C^\circ:\quad b^3=a^4+1,
 \qquad
 a=\frac{z_1}{F_\epsilon^{1/4}},\qquad
 b=\frac{(1-\epsilon)^{1/3}z_2}{F_\epsilon^{1/3}}.
\]
Conversely, every point of this affine curve has a unique
positive weighted rescaling to the unit three-sphere. A different
page phase changes this by constant complex factors and gives
the same marked cover. The compactification $C$ has genus three
and one point at infinity. Its actual map to the original
$t$-line is
\[
 t_\epsilon=\frac{a^4+1}{1+\epsilon a^4}.
\]
It is ramified with orders three and four over $0$ and $1$,
and with order twelve at the added point over $t=1/\epsilon$.
The twelve points over $t=\infty$ satisfy
$a^4=-1/\epsilon$ and are unramified. Thus this is the
actual page cover; the cover ramified twelvefold over the
original cusp is its $\epsilon=0$ comparator, not the same
marked map when $\epsilon>0$.

\item Keep the product modification inside $|q|<r_0<\epsilon$.
The page equation there has twelve disjoint graph discs
$\theta=\theta_j(q)$, since its derivative in $\theta$ is
twelve. Consequently $M$ consists of twelve nodal elliptic
neighbourhoods attached to an entirely untouched raw torus
quotient. Denote the latter part by $M^\circ$ and its base by
$C_0$; it is the page with these twelve discs removed. It has
nonempty boundary and retracts to a graph.

\item We identify the actual circle bundle on this unmodified part.
On an ordinary marked chart let $g$ be a real lift of its
$\gamma$ coefficient and set
\[
 \xi=\zeta_1-6\mu g,\qquad
 \eta=\zeta_2-\beta g.
\]
Then $\xi=\tau u+w$ and $\eta=\mu u+d$. Exponentiating
$\eta$ gives the actual degree-zero line on the first elliptic
torus with multipliers
\[
 (\xi,z)\sim(\xi+\tau,e^{2\pi i\mu}z),\qquad
 (\xi,z)\sim(\xi+1,z).
\]
These are the Poincar\'e multipliers of the original normal
function $P$, up to the single dualization convention for the
associated circle line. The identification is an identification
of the varying fibrewise Picard section, not just its fibre
degree. Under a change of marked elliptic periods the maps are
fibrewise affine holomorphic and the line maps are its
semiabelian lifts. The translations may depend smoothly on
$g$; no global holomorphic-torsor assertion is needed.

\item Over the graph base $C_0$, the affine elliptic torsor has a
smooth section. Indeed its obstruction lies in the twisted
group $H^2(C_0,\Lambda)$, which is zero; equivalently use the
sequence from the lattice local system to its real vector
bundle and then to the torus bundle, whose vector-bundle
section sheaf is fine. Identify the torsor with its Jacobian
using such a section. Translation of a degree-zero Poincar\'e
line changes it by a line pulled back from the base. Two
fibrewise holomorphic lines with the same varying Picard
section likewise differ by a base line: their quotient has a
one-dimensional space of fibrewise holomorphic sections and
evaluation trivializes it fibrewise. Since $H^2(C_0;\Z)=0$,
this base line has zero topological Chern class. The actual
class on $M^\circ$ is therefore the Poincar\'e class of
$P\circ t_\epsilon$, up to sign, including its mixed
base--fibre component.

\item There is no hidden correction supported at the twelve modified
nodes. By excision and Poincar\'e--Lefschetz duality, each
local summand of
\[
 H^2(M,M^\circ;\Z)
 \simeq\bigoplus_{j=1}^{12}H^2(E_{I_1},\partial E_{I_1};\Z)
\]
is $\Z$, generated by the Thom class of a regular elliptic
fibre. The regular fibre is primitive in $H_2(E_{I_1};\Z)$,
as its degree-one specialization generates the nodal fibre's
two-cycle. Its image in $H^2(M;\Z)$ is the pullback of the
point class of the base, and is zero because that base has
boundary. The restriction $H^2(M;\Z)\to H^2(M^\circ;\Z)$
is consequently injective. This addresses actual local
clutching classes, rather than only their holonomy.

\item Trim a smaller boundary collar, which does not change the
relative square because its line is trivial, and identify the
remaining annular end of the page with the punctured disc at
the added point of $C$. Fill it by its regular marked
$T^2\times D^2$. The boundary affine torus bundle is
fibrewise complex-trivialized by a path of translations.
After matching the Picard section $P$, a remaining line
identification is a scalar base-line gauge. Thus one can choose
the product-cap Poincar\'e extension whose section is
$P\circ t_\epsilon$ near the added point; the displayed
multipliers are regular there. The line is topologically
trivial on the cap, and its square contribution is zero.
Boundary gauge changes cannot change the relative square, by
\ref{push:src:C.3.20}. Call the resulting
closed smooth genus-one Lefschetz surface $\widehat M$ and its
line $\widehat{\mathcal L}$. Its square is the same relative
square $s_\delta$; the cap contributes zero.

\item We spell out why the square can now be compared with the
Jacobian without a projectivity or multisection assumption on
$\widehat M$. Let $j:C^*=C\setminus\{12\text{ nodal values}\}
\hookrightarrow C$ and let $\Lambda=R^1\pi_*\Z|_{C^*}$.
For a nodal elliptic fibre, degree-one specialization identifies
$R^2\pi_*\Z$ with the constant sheaf $\Z$, while degree one
is the invariant subgroup of the unipotent monodromy. Thus
\[
 R^1\pi_*\Z=j_*\Lambda,\qquad R^2\pi_*\Z=\Z.
\]
These sheaves and their fibre intersection pairing agree for
the actual surface and its Jacobian. The degree-zero class
$c_1(\widehat{\mathcal L})$ lies in the first Leray filtration
step. Its component in
$H^1(C,j_*\Lambda)$ is the connecting class of the varying
Picard section $\pm P$. This holds on $C^*$ by the preceding
actual multiplier calculation. It holds on all of $C$ because
\[
 H^1(C,j_*\Lambda)\longrightarrow H^1(C^*,\Lambda)
\]
is injective, by the first terms of the Leray sequence for $j$.
Local nodal choices can therefore add only a base-filtration
class; the excision argument gives the same conclusion directly.

\item Multiplicativity of Leray now determines the square. In degree
four its only nonzero graded piece is
$H^2(C,R^2\pi_*\Z)=\Z$. The square of a class in filtration
one is the cup square of its $H^1(C,j_*\Lambda)$ component,
using the fibre intersection pairing. Adding filtration-two
base classes contributes only filtration three or four, both
zero on this base. Hence the actual square is exactly the
square of the Poincar\'e class on the Jacobian; dualizing the
line does not change it.

\item Finally let $J_\epsilon\to C$ be the relatively minimal
Jacobian obtained by pulling back the original elliptic surface
along $t_\epsilon$. It has the explicit affine model
\[
 Y^2=X^3-3\lambda bX+2\lambda^2a^2,
 \qquad P=(0,-\sqrt2\lambda a),\qquad
 \lambda=1-\epsilon.
\]
Writing $h=1+\epsilon a^4$, the substitution from the original
Weierstrass coordinates is
$x=(ab^2/h)^2X$, $y=(ab^2/h)^3Y$.
The discriminant is
\[
 \Delta=1728\lambda^3(1+\epsilon a^4).
\]
It has twelve simple zeroes on $C^\circ$, and $c_4$ is a unit
at all of them. They are precisely the twelve $I_1$ fibres.
At the unique point at infinity, $a,b$ have pole orders
three and four. Scaling $X,Y$ by a local parameter to orders
two and three gives a regular Weierstrass model whose
discriminant is a unit when $\epsilon\ne0$; the point at
infinity is therefore a good fibre. This also shows that the
fundamental line has degree one, so
$\chi(\mathcal O_{J_\epsilon})=1$.

\item The displayed section $P$ is finite on every affine chart and
also in this integral chart at infinity, hence it is disjoint
from the zero section. At a simple discriminant zero it is a
smooth point of the nodal fibre: a node has $Y=0$, whereas
$a\ne0$ there and $Y(P)=-\sqrt2\lambda a\ne0$.
All singular fibres are irreducible, so there is no component
correction. These explicit identities are independently checked
in \texttt{ab\_push\_actual\_line\_comparison.py}.

\item The divisor $P-O$ has fibre degree zero and represents the
required relative Picard section, up to a base twist. Both
sections have square $-\chi=-1$ and $P\cdot O=0$. Therefore
\[
 (P-O)^2=-1-1-2(P\cdot O)=-2.
\]
Base twists do not change this square. The actual Leray
comparison above proves $s_\delta=-2$, as asserted.

\end{enumerate}
\end{proof}




\begin{lemma}[one common signed cap convention]\label{push:src:C.3.22}
\StatementRecord{push:src:C.3.22}{proved}{GPT-6 Astra (OpenAI, ChatGPT Pro mode), effort Pro, 2026-11-17; Claude Fable 5.2 (Anthropic), effort high, 2026-11-16; Claude Fable 5.2 (Anthropic), effort high, 16 November 2026}{the hypotheses, notation and definitions stated here; \cref{push:src:C.3.18}; \cref{push:src:C.3.20}; \cref{push:src:C.3.19}}

Retain the relative Pontryagin number
$p_{\mathrm{rel}}=-24$ of \ref{push:src:C.3.18} and the actual relative
Euler square $s_\delta$ of \ref{push:src:C.3.20}. With the single
cut-and-reglue orientation sign $\epsilon_{\rm cap}$ used
in \ref{push:src:C.3.20}, the two cap pairings satisfy
\[
 \langle p_1(M_\psi)x,[M_\psi]\rangle
       =\epsilon_{\rm cap}\,p_{\mathrm{rel}},
 \qquad
 k_\psi=\epsilon_{\rm cap}\,3s_\delta.
\]
In particular
\[
 J_\psi=\epsilon_{\rm cap}
              \frac{p_{\mathrm{rel}}-12s_\delta}{24}.
\]
This lemma does not assign a value to $s_\delta$.
\end{lemma}


\begin{proof}
\begin{enumerate}[label=\arabic*.]
\item We fix the normalization by connections, so that a boundary
frame and a coordinate clutch are not silently identified.
For a real orthogonal connection with curvature $F$, use
\[
 p_1(F)=-\frac{\operatorname{tr}(F\wedge F)}{8\pi^2},
 \qquad
 \operatorname{CS}_{p_1}(\omega)
 =-\frac{\operatorname{tr}(\omega\wedge d\omega
                          +\tfrac23\omega^3)}{8\pi^2}.
\]
For $K$ a change of coordinates,
\[
 \operatorname{CS}_{p_1}(K^{-1}dK)
 =\frac{\operatorname{tr}(K^{-1}dK)^3}{24\pi^2}
 =:\tau_{p_1}(K).
\]
This is the suspension used for a mapping torus oriented
with its interval parameter first. Indeed interpolating
the connection from zero to $K^{-1}dK$ gives curvature
$dt\wedge\omega+(t^2-t)\omega^2$, whose first Pontryagin
form integrates in $t$ to the displayed three-form.

\item In contrast, let a boundary frame be
$\mathfrak f=\mathfrak t B$ relative to a global stable
frame $\mathfrak t$. The connection for which
$\mathfrak f$ is parallel has matrix $-dB B^{-1}$ in
the $\mathfrak t$ coordinates. Therefore
\[
 p_1(E,\mathfrak f)[M,\partial M]
 =-\int_{\partial M}\tau_{p_1}(B).
\]
Equivalently, capping the boundary frame by a standard
frame uses the inverse coordinate clutch $B^{-1}$.
This minus sign is forced by the connection computation;
it is not an orientation chosen to obtain a cancellation.

\item Now use exactly the change-of-frame convention of \cref{push:src:C.3.19}:
\[
 \mathfrak t=\mathfrak f_{\mathrm{aff}} A,
 \qquad
 G=A(\psi p)^{-1}A(p)
\]
after contracting the affine derivative shear. Set
\[
 a_0=\int_{T_{\gamma,u,w}}\tau_{p_1}(A).
\]
Here $B=A^{-1}$ in the boundary-frame formula. Since
the boundary orientation of $M$ is
$\partial M=-T_{\gamma,u,w}$, that formula gives
\[
 p_{\mathrm{rel}}
 =-\int_{\partial M}\tau_{p_1}(A^{-1})
 =-a_0.
\]

\item For the actual annulus use the orbit orientation
$dg\wedge du\wedge dw\wedge dr\wedge da$ and the
Euler curvature $E=dg\wedge da$, with $\chi$ decreasing
from one to zero as $r$ increases. The relative trace
calculation of \ref{push:src:C.3.19} has representative
\[
 \tau_{p_1}(G)
 =-d\chi\wedge\beta
 \quad\hbox{in relative cohomology},
 \qquad
 \int_{T^2_{u,w}}\beta=a_0.
\]
Consequently
\[
 \int_P\tau_{p_1}(G)\wedge E=-a_0=p_{\mathrm{rel}}.
\]
For example, in the only contributing term the form is
$-a_0\,d\chi\wedge du\wedge dw\wedge dg\wedge da$;
relative to the displayed orbit orientation its coefficient
is $a_0\chi'(r)$, and the integral of $\chi'$ is $-1$.

\item This is also the Pontryagin--Euler pairing of the actual
mapping torus in \ref{push:src:C.3.20}. Its line curvature restricts to
$E+t\Delta$ on the $t$-slices. The possible extra term
has integral zero, because
\[
 \int_{DP}\tau_{p_1}(G)\wedge\Delta
 =\int_{DP}\tau_{p_1}(G)\wedge d\eta=0.
\]
Thus no line-clutch correction changes the preceding
pairing. The same cut-and-reglue bordism of \ref{push:src:C.3.20} is a
smooth bordism of tangent bundles and of the specified
line bundle. Both characteristic polynomials $p_1x$
and $x^3$ therefore acquire the same boundary-orientation
sign $\epsilon_{\rm cap}$. Their identity-cap and
identity-mapping-torus values vanish. Combining the
computed mapping-torus pairing with the cubic calculation
of \ref{push:src:C.3.20} gives the two asserted signed equations and
then their stated combination.

\end{enumerate}
\end{proof}




\begin{proposition}[an additive integer before reduction modulo 28]\label{push:src:C.3.23}
\StatementRecord{push:src:C.3.23}{proved}{GPT-6 Astra (OpenAI, ChatGPT Pro mode), effort Pro, 2026-11-17; Claude Fable 5.2 (Anthropic), effort high, 2026-11-16; Claude Fable 5.2 (Anthropic), effort high, 16 November 2026}{the hypotheses, notation and definitions stated here; \cref{push:src:C.3.14}; \citet[Theorem A, page 621]{Smale1959E}; \citep[Sections 1 and 5.16]{Haefliger1966D}; \citet[Example 7.14 and Proposition 7.15]{Davis1982D}}

Fix the linear semifree circle action of \ref{push:src:C.3.14}, with its
circle direction fixed. On the group of orientation-preserving
equivariant diffeomorphisms of $S^6$ there is an integer
\[
 \mathfrak q:\pi_0\operatorname{Diff}^+_{S^1}(S^6)
                    \longrightarrow\Z
\]
defined by the unframed Haefliger coordinate of the orbit
three-knot of equivariant clutching. It is a group homomorphism.
For a representative normalized to be the identity on a fixed
tubular neighbourhood, it is Wall's unframed knot coordinate,
so $\mathfrak q=-J$ in the convention of \ref{push:src:C.3.14}. Its image in
the ordinary group $\Theta_7$ is the fixed signed reduction
modulo $28$ described there. In particular, if an equivariant
representative has a sixth power equivariantly isotopic to a
representative with $J=0$ as an integer, its ordinary mapping
class is zero.
\end{proposition}


\begin{proof}
\begin{enumerate}[label=\arabic*.]
\item We first make the normalization and its independence explicit.
The fixed set is $F=S^2$, and its equivariant normal bundle is
the trivial scalar complex rank-two bundle after conjugating
the negative weight, by \ref{push:src:C.3.14}. An equivariant real derivative
on that normal representation is complex linear and has
positive real determinant. An orientation-preserving ambient
diffeomorphism therefore restricts to an orientation-preserving
diffeomorphism of $F$.

\item By \citet[Theorem A, page 621]{Smale1959E}, this restriction is
isotopic to the identity. Lift such an isotopy using the product
normal bundle and an invariant tubular cutoff, extending its
velocity field equivariantly. After this adjustment the map
fixes $F$ pointwise. Its normal derivative is a map
$F\to GL_2(\C)$. Polar decomposition and $\pi_2(U(2))=0$
give a homotopy to the identity, again realized in the fixed
tube by an equivariant isotopy. Finally the usual rescaling
of a tubular germ to its derivative, followed by an invariant
cutoff, makes the whole map the identity on a smaller fixed
tube. Equivariance rules out a derivative from the nontrivial
normal representation to the fixed tangent representation;
the rescaled germ is smooth at the zero parameter. This
normalization uses equivariant isotopies throughout.

\item Clutch two standard linear equivariant seven-discs by this
diffeomorphism. Their fixed three-discs clutch to $S^3$.
Their orbit six-discs clutch to a homotopy six-sphere, which
is standard because $\Theta_6=0$. The fixed image is thus a
smooth oriented three-knot in $S^6$. Its Haefliger integer is
the proposed $\mathfrak q$, using the connected-sum group of
oriented three-knots in the six-sphere
\citep[Sections 1 and 5.16]{Haefliger1966D}.

\item This integer does not depend on the chosen normalization.
An equivariant isotopy of boundary clutching maps extends
over a collar to a diffeomorphism of the clutched equivariant
pairs. The resulting orbit pairs are diffeomorphic. One can
choose compatible unitary normal tubes along the entire
isotopy, so that their Hopf orbit charts form a smooth family;
this does not invoke functorial smoothness for arbitrary
equivariant maps in a fixed raw quotient atlas.
There is also no ambiguity from identifying the orbit homotopy
sphere with $S^6$. An orientation-preserving diffeomorphism of
$S^6$ is isotopic to one supported in a ball disjoint from a
given compact three-knot: make its germ the identity near a
point, use stereographic coordinates on the complement, and
conjugate its compact support into a small ball by translations
and positive dilations, all isotopic to the identity. Such a
representative fixes that knot. Consequently the ambient
identification cannot change its knot integer.

\item We prove additivity geometrically. Normalize two equivariant
diffeomorphisms to be the identity near $F$. In particular each
is the identity near a chosen fixed point. Stereographic
projection from that point identifies the complement with
$\R^2\oplus\C^2$, with the circle acting scalarly on the
second summand. The diffeomorphism has compact support there.
Positive dilations and translations along the fixed $\R^2$
commute with the circle and are equivariantly isotopic to the
identity on the one-point compactification. By conjugation,
place the two supports inside disjoint small invariant balls
centered at distinct fixed points. The resulting maps remain
in their original equivariant components.

\item For disjoint supports, the composition clutch is the
equivariant connected sum of the two separate clutches.
This follows directly by cutting the equatorial sphere between
the two invariant balls and inserting a product collar; the
map on the separating collar is the identity. The fixed
three-discs undergo the same boundary connected-sum operation.
On the orbit pairs this is the ordinary connected sum of
the two oriented three-knots. The defining group operation
of their Haefliger integers is connected sum. Hence
\[
 \mathfrak q(\phi_1\phi_2)
 =\mathfrak q(\phi_1)+\mathfrak q(\phi_2).
\]
This proves additivity without passing first to $\Theta_7$.

\item For a fixed-tube-relative representative, the surgery cap and
Wall calculation of \ref{push:src:C.3.14} apply verbatim: the inherited product
normal framing gives $\mathfrak q=-J$. Altering this framing
changes only Wall's framing coordinate, whose equivariant lift
extends over the seven-disc by the explicit quaternionic
formula in \ref{push:src:C.3.14}. Finally the orbit-knot generator comparison
of \citet[Example 7.14 and Proposition 7.15]{Davis1982D}
identifies the underlying ordinary sphere class with the fixed
signed reduction of $\mathfrak q$ modulo $28$.

\item If $\phi^6$ is equivariantly isotopic to a representative with
integer $J=0$, then $6\mathfrak q(\phi)=0$ in $\Z$.
Thus $\mathfrak q(\phi)=0$, and its underlying ordinary
mapping class is zero. This last implication uses a torsion-free
integer before reduction; it is not inferred from the weaker
equation $6[\phi]=0$ in $\Z/28$.

\end{enumerate}
\end{proof}




\begin{theorem}[the actual smooth monodromy is trivial]\label{push:src:C.3.24}
\StatementRecord{push:src:C.3.24}{proved}{GPT-6 Astra (OpenAI, ChatGPT Pro mode), effort Pro, 2026-11-17; Claude Fable 5.2 (Anthropic), effort high, 2026-11-17; Claude Fable 5.2 (Anthropic), effort high, 2026-11-16; GPT-6 Astra (OpenAI, ChatGPT Pro mode), S6\_06\_BLIND, effort Pro, 2026-11-17; Claude Fable 5.2 (Anthropic), effort high, 16 November 2026}{the hypotheses, notation and definitions stated here; \cref{push:src:C.4.1}; \cref{push:src:C.3.10}; \cref{push:src:C.3.18}; \cref{push:src:C.3.21}; \cref{push:src:C.3.22}; \cref{push:src:C.3.14}; \cref{push:src:C.3.23}; \cref{push:src:C.3.6}}

For the actual retained $S^6$ family and a positive generator
of its punctured $w$-parameter disc, the oriented smooth
monodromy satisfies
\[
 [\Phi]=0\quad\hbox{in }\pi_0\operatorname{Diff}^+(S^6)
                       \simeq\Theta_7\simeq\Z/28.
\]
Its mapping torus is consequently oriented-diffeomorphic to
$S^1\times S^6$. Every finite cyclic base change has trivial
smooth monodromy as well. By \ref{push:src:C.4.1}'s marked family comparison,
the same conclusion applies to the admissible sphere twist
data covered by that comparison.
\end{theorem}


\begin{proof}
\begin{enumerate}[label=\arabic*.]
\item Let $\widehat\psi$ be the actual smooth equivariant lift of
the annular orbit map in \ref{push:src:C.3.10}.
The actual cap calculation is an integer calculation.
\ref{push:src:C.3.18} gives $p_{\mathrm{rel}}=-24$ and \ref{push:src:C.3.21} gives
$s_\delta=-2$. The common signed comparison in \ref{push:src:C.3.22} yields
\[
 J_\psi=\epsilon_{\rm cap}
          \frac{-24-12(-2)}{24}=0\quad\hbox{in }\Z.
\]
By \ref{push:src:C.3.14} and \ref{push:src:C.3.23}, the unframed orbit-knot integer of the
actual annular lift is therefore zero. Neither the overall cap
orientation sign nor the Wall--Davis generator sign can change
this conclusion.

\item We verify that this is the sixth power at the equivariant
integer level. The family carries the fixed positive
$\delta$ circle, and all its closing maps commute with that
circle. Choose a circle-invariant Ehresmann connection along
one positive $w$ loop. Concatenating it six times gives the
sixth power of the same equivariant return diffeomorphism.
Changing to the invariant transport used in \ref{push:src:C.3.10} changes
this return by an equivariant isotopy: invariant horizontal
connections form an affine space, and their interpolated
flows give that isotopy.

\item The closing cocycle itself introduces no hidden discrepancy.
In the notation of Paper~I, Proposition~18.5,
\[
 b_{h,c}=-\frac{h}{12}\Pi_c\widehat\gamma
             +\left(h\ell-\frac{h^2}{24}\right)e,
 \qquad \Pi_{c+h}\widehat\gamma
                   =\Pi_c\widehat\gamma+h e.
\]
Direct substitution gives
\[
 b_{h,c}+b_{k,c+h}=b_{h+k,c}.
\]
Thus six $h=2$ closings are the $h=12$ closing used for the
annular reduction. Returning from these coherent closings to
any stipulated Paper~I, Proposition~3.3 and Paper~I, Proposition~3.4 and Paper~I, Theorem~3.8 choices changes them by a scalar in the
connected group $\C^*$, which centralizes $S^1_\delta$;
this is again an equivariant isotopy. Every subsequent unit isotopy, finite-chart
transport and circle-gauge comparison in \ref{push:src:C.3.6} and \ref{push:src:C.3.10} is
equivariant. Consequently the annular lift and $\Phi^6$ have
the same integer $\mathfrak q$, not merely the same underlying
ordinary mapping class.

\item Additivity from \ref{push:src:C.3.23} now gives
\[
 6\mathfrak q(\Phi)=\mathfrak q(\Phi^6)
                   =\mathfrak q(\widehat\psi)=-J_\psi=0.
\]
Since this equality is in $\Z$, it implies
$\mathfrak q(\Phi)=0$. Its signed reduction modulo $28$
is the ordinary class $[\Phi]$, so that class is zero.
This removes the possible order-two ambiguity by an integer
argument; it does not divide the equation $6[\Phi]=0$ inside
$\Z/28$.

\item A smooth isotopy from the return map to the identity gives
the usual oriented product identification of its mapping
torus. A finite cyclic base change replaces the return map
by a positive power, whose class is again zero. The final
twist-data statement follows from the local family isomorphisms
and connected automorphism transitions already proved in \ref{push:src:C.4.1}.

\end{enumerate}
\end{proof}




\begin{proposition}[an explicit eight-dimensional coboundary]\label{push:src:C.3.25}
\StatementRecord{push:src:C.3.25}{proved}{GPT-6 Astra (OpenAI, ChatGPT Pro mode), effort Pro, 2026-11-17; Claude Fable 5.2 (Anthropic), effort high, 16 November 2026}{the hypotheses, notation and definitions stated here; \cref{push:src:C.3.14}; \cref{push:src:C.2.1}; \citet[Section 2.2, equation (2.1)]{CrowleyNordstrom2020}; \citet[Section 6, equation (11)]{EellsKuiper1962}; \cref{push:src:C.1.1}}

Let $\varphi$ be an orientation-preserving equivariant return map,
normalized to be the identity on a fixed tubular neighbourhood.
Let $V_\varphi$ be its actual framed-surgery six-manifold, constructed
as in \ref{push:src:C.3.14}. Choose the sign of $x$ so that the unit circle
bundle of its complex line, with the boundary fibre orientation,
is the actual oriented free circle bundle. Orient its meridian
accordingly, so that $\langle x,S^2\rangle=1$. This fixes the
possible associated-line versus dual-line convention. Write
\[
 p=\langle p_1(V_\varphi)x,[V_\varphi]\rangle,
 \qquad k=\langle x^3,[V_\varphi]\rangle.
\]
There is a specified smooth spin$^c$ coboundary $Z_\varphi$ of
the actual clutching sphere $\Sigma_\varphi$. Its fourth
intersection form and characteristic covectors are
\[
 Q=\begin{pmatrix}k&1\\1&0\end{pmatrix},\qquad
 p_1(Z_\varphi)=(p+k,3),\qquad c^2=(k,1),
\]
where $c$ is the determinant class and the covectors are evaluated
on the ordered integral basis specified below. In particular
\[
 \sigma(Z_\varphi)=0,\qquad
 p_1^2=6p-3k,\qquad p_1c^2=p+k,\qquad c^4=k.
\]
With the boundary-consistent orientations, its spin$^c$ index
formula gives
\[
 \mu(\Sigma_\varphi)\equiv\frac{p-4k}{672}\pmod\Z.
\]
Boundary connected sum with $S^1\times D^7$ gives a specified
smooth spin$^c$ coboundary of the actual mapping torus, with the
same displayed characteristic numbers. This construction makes
no use of Wall's knot coordinate or the Davis generator comparison.
\end{proposition}


\begin{proof}
\begin{enumerate}[label=\arabic*.]
\item We first identify the boundary geometrically. Write $C$ for the
exterior of the actual fixed three-knot in the orbit six-sphere.
The framed-surgery construction is
\[
 V_\varphi=C\cup_{S^3\times S^2}(D^4\times S^2).
\]
This description, including its framing, is the direct cut-and-paste
construction of \ref{push:src:C.3.14}; its later use of knot classification is not
needed here. It is spin and has $H^2=H^4=\Z$, $H^3=0$.
Let $L_x\to V_\varphi$ be the complex line with first Chern class
$x$, and put $Z_0=D(L_x)$, oriented by the base followed by its
complex line fibre. Its boundary $T=S(L_x)$ is the corresponding
principal circle bundle.

\item On the surgery piece $N=S^2\times D^4$, the line is the product
of the Euler-one line $\mathcal O(1)\to S^2$ and $D^4$.
Thus $T$ contains the specified product $S^3_{\rm Hopf}\times D^4$.
Attach one four-handle
\[
 D^4_{\rm Hopf}\times D^4
 \quad\hbox{along}\quad S^3_{\rm Hopf}\times D^4
\]
using this product identification and framing. Denote the resulting
eight-manifold by $Z_\varphi$. On the boundary this replaces
$S^3_{\rm Hopf}\times D^4$ by $D^4_{\rm Hopf}\times S^3$.
Over $C$, the remaining boundary is exactly the actual free circle
bundle of the clutched sphere: principal circle bundles there are
classified by their Euler class, which is the specified restriction
of $x$. On their common boundary $S^3\times S^2$, a discrepancy
between two such bundle identifications is a circle gauge with a
real logarithm, because $H^1(S^3\times S^2;\Z)=0$.
A collar cutoff corrects it to the prescribed identification.
Alexander duality gives $H^2(C;\Z)=\Z$, generated by the
meridian, so specifying this oriented Hopf restriction fixes the
entire line class on $C$. There is no additional line-bundle
parameter on the exterior. The replacement
$D^4_{\rm Hopf}\times S^3$ is now exactly the
actual fixed tubular neighbourhood with its Hopf normal action.
Hence
\[
 \partial Z_\varphi=\Sigma_\varphi
\]
as an oriented smooth manifold, with the orientation fixed by
agreement on the common free exterior. This is an identification
with the actual boundary, not merely with a sphere having a
conjectured invariant.

\item Explicitly, the boundary orientation of $D(L_x)$ is the positive
line-circle direction followed by the orientation of its
six-dimensional base. This is the circle-first orientation of
the actual free exterior by the defining choice of $x$.
The surgery product framing then gives the boundary-consistent
orientation on the new fixed tube. If a prior principal-connection
convention uses the dual line, replace $x$ by $-x$ throughout;
both $p$ and $k$ reverse, as does the corresponding integer
conversion. The actual zero evaluation below is unchanged.

\item The added four-cycle has a particularly concrete description.
Inside the displayed product region, the attachment replaces
\[
 D(\mathcal O(1))\times D^4
 \quad\hbox{by}\quad \mathbb{CP}^2\times D^4.
\]
Indeed a tubular neighbourhood of a projective line in
$\mathbb{CP}^2$ is $D(\mathcal O(1))$, and its complement is a
four-disc, with the same standard Hopf boundary identification.
Let $b$ be this positively oriented $\mathbb{CP}^2\times\{0\}$.
Its normal four-plane bundle is trivial, with the explicit
$D^4$ product framing. In particular $b^2=0$; it is not a
four-sphere with zero Pontryagin evaluation.

\item Let $a$ be the image of the class $\operatorname{PD}_{V_\varphi}x$
in the zero section of $D(L_x)$. A generic section of $L_x$
represents it by a closed oriented four-submanifold. Its normal
bundle in $Z_0$ is the sum of its normal line in the zero section
and the vertical line, both isomorphic to $L_x$ on that
submanifold. Therefore
\[
 a^2=\langle x^3,[V_\varphi]\rangle=k.
\]
Moreover $a\cdot b=1$: in the surgery product this intersection
is the pairing of $x$ with the meridian $S^2$, which is one by
the Euler-one convention. The four-handle exact sequence and
$H_3(V_\varphi)=0$ give
\[
 0\longrightarrow H_4(V_\varphi;\Z)
 \longrightarrow H_4(Z_\varphi;\Z)
 \longrightarrow\Z\longrightarrow0.
\]
The class $b$ uses the handle core once and maps to its generator.
Thus $(a,b)$ is an integral basis, and its intersection matrix is
the displayed $Q$. Its determinant is $-1$, so it has one
positive and one negative eigenvalue and signature zero.

\item The handle does not change second cohomology. Let $c$ be the
unique extension of $\pi^*x$ from $Z_0$ to $Z_\varphi$.
On $Z_0$ the tangent bundle restricts stably to
$\pi^*(TV_\varphi\oplus L_x)$, so
\[
 w_2(TZ_0)=\pi^*x\pmod2.
\]
The same degree-two handle isomorphism with $\Z/2$ coefficients
proves $w_2(TZ_\varphi)=c\pmod2$. Thus $c$ is a determinant
class of a spin$^c$ structure. It is trivial on the boundary,
which is a homotopy sphere and has its unique spin structure.
Notice that this coboundary need not be spin: $c$ is primitive.

\item On $a$ the stable tangent splitting gives
\[
 \langle p_1(TZ_\varphi),a\rangle=p+k,
 \qquad \langle c^2,a\rangle=k.
\]
On $b$ the normal framing gives
$TZ_\varphi|_b=T\mathbb{CP}^2\oplus\R^4$.
The class $c|_b$ is the hyperplane class, since it evaluates as
one on the original meridian projective line. Hence
\[
 \langle p_1(TZ_\varphi),b\rangle=3,
 \qquad \langle c^2,b\rangle=1.
\]
Here $p_1(T\mathbb{CP}^2)=c_1^2-2c_2=3H^2$ follows directly
from its Euler sequence. This proves the stated covectors.

\item Because the boundary is a homotopy sphere, the relevant relative
and absolute fourth cohomology groups agree. Their characteristic
pairing is the inverse of the integral intersection matrix:
\[
 Q^{-1}=\begin{pmatrix}0&1\\1&-k\end{pmatrix}.
\]
Multiplying the two covectors by this inverse gives exactly
\[
 p_1^2=6p-3k,\qquad p_1c^2=p+k,\qquad c^4=k.
\]

\item For clarity, the conversion to the boundary invariant uses only
the spin$^c$ index formula, not a knot-classification formula.
The closed eight-dimensional index density is
\[
 [e^{c/2}\widehat A]_8
 =\frac{7p_1^2-4p_2}{5760}
       -\frac{p_1c^2}{192}+\frac{c^4}{384}.
\]
Eliminating $p_2$ by the signature formula and gluing to any spin
filling of the boundary gives \ref{push:src:C.2.1}'s expression
\[
 \mu(\Sigma_\varphi)\equiv
 \frac{p_1^2-4\sigma}{896}
 -\frac{p_1c^2}{192}+\frac{c^4}{384}\pmod\Z.
\]
The primary normalization is
\citet[Section 2.2, equation (2.1)]{CrowleyNordstrom2020},
with its invariant divided by $28$; the spin normalization is
\citet[Section 6, equation (11)]{EellsKuiper1962}.
Substitution now gives, by exact rational arithmetic,
\[
 \frac{6p-3k}{896}-\frac{p+k}{192}+\frac{k}{384}
       =\frac{p-4k}{672}.
\]

\item Finally \ref{push:src:C.1.1} identifies the mapping torus with
$(S^1\times S^6)\#\Sigma_\varphi$. Boundary connected sum
of $Z_\varphi$ with $S^1\times D^7$, using a trivial determinant
line on the latter, realizes this actual boundary. Either spin
structure on the product can be used. The added summand has no
fourth cohomology and has zero signature; the displayed relative
characteristic numbers do not change.

\end{enumerate}
\end{proof}


\begin{remark}[Scope and actual six-loop evaluation]
For the actual annular six-loop cap the independent geometric
calculations \ref{push:src:C.3.18}--\ref{push:src:C.3.22} give
$p=-24\epsilon_{\rm cap}$ and $k=-6\epsilon_{\rm cap}$.
Here the one common sign includes the possible simultaneous
meridian-sign change just specified.
The present coboundary therefore has
\[
 p_1^2=-126\epsilon_{\rm cap},\qquad
 p_1c^2=-30\epsilon_{\rm cap},\qquad
 c^4=-6\epsilon_{\rm cap},\qquad \sigma=0.
\]
Its spin$^c$ formula is zero directly:
\[
 \epsilon_{\rm cap}
 \left(-\frac{126}{896}+\frac{30}{192}
                         -\frac{6}{384}\right)=0.
\]
This does not use the assertion $J_\psi=0$ or the Wall--Davis
conversion. By itself it evaluates the six-loop sphere. The
independent integer additivity needed for the one-loop conclusion
is supplied separately in \ref{push:src:C.3.26}. The characteristic arithmetic
is checked in \texttt{d\_push\_eight\_coboundary.py}.

This is a smooth topological spin$^c$ coboundary. No complex
structure on it is asserted, and it is not identified with the
specified analytic filling excluded by Paper~I, Theorem~21.58.
\end{remark}

\begin{proposition}[direct characteristic-index additivity]\label{push:src:C.3.26}
\StatementRecord{push:src:C.3.26}{proved}{GPT-6 Astra (OpenAI, ChatGPT Pro mode), effort Pro, 2026-11-17; GPT-6 Astra (OpenAI, ChatGPT Pro mode), S6\_06\_BLIND, effort Pro, 2026-11-17; Claude Fable 5.2 (Anthropic), effort high, 16 November 2026}{the hypotheses, notation and definitions stated here; \cref{push:src:C.3.14}; \cref{push:src:C.3.20}; \citet[Theorem A]{Smale1959E}}

For a fixed-tube-relative orbit diffeomorphism $\varphi$ of the
linear semifree sphere, form its actual surgery cap $M_\varphi$
and meridian line $L_x$ as in the geometric construction of
\ref{push:src:C.3.14}. Define
\[
 \lambda(\varphi)=\frac{1}{24}
      \langle p_1(M_\varphi)x-4x^3,[M_\varphi]\rangle.
\]
This integer is additive under composition. It is unchanged by
the equivariant isotopies and fixed-tube normalizations of the
upstairs diffeomorphism, and therefore defines an integer on its
equivariant mapping class. These conclusions follow directly
from spin index theory and characteristic bordisms; they do not
use the classification or integer of three-knots, Wall's surgery
formula, or the comparison of a knot generator with a Milnor sphere.
\end{proposition}


\begin{proof}
\begin{enumerate}[label=\arabic*.]
\item The geometric cap is a simply connected spin six-manifold with
$H^2=H^4=\Z$ and $H^3=0$. These facts follow from the explicit
two-copy construction $N\cup_{\widetilde\varphi}(-N)$,
$N=S^2\times D^4$, by van Kampen, Mayer--Vietoris and the
unique boundary spin structure. They do not require a knot
classification. The twisted spin Dirac formula gives
\[
 \operatorname{ind}(D\otimes L_x)
 =\langle x^3/6-p_1x/24,[M_\varphi]\rangle
 =-\lambda(\varphi),
\]
so the displayed quantity is an integer.

\item We first prove additivity for maps fixed on the chosen tube.
Put $Y=\partial N=DP=S^2\times S^3$ and extend each map by
the identity on the second half. The meridian line on $Y$ has
a lift of each extended map: the actual upstairs circle bundle
provides it on $P$, and the identity provides it on the other
half. Their boundary values agree. Choose the lift of a
composition to be the composition of these lifts.

\item An oriented pair of pants has free fundamental group on its
two incoming boundary generators. The two diffeomorphisms
therefore specify a smooth $Y$-bundle over that surface, with
outgoing boundary monodromy their composition. The chosen line
bundle lifts specify a line bundle on its total space. This
total space is a compact oriented seven-manifold whose boundary
is the outgoing mapping torus minus the two incoming mapping
tori. Its stable tangent bundle restricts to their tangent
bundles plus the outward normal line. Consequently both
characteristic numbers $x^3$ and $p_1x$ are additive on these
mapping tori, by bordism invariance.

\item Use the explicit seven-manifold of \cref{push:src:C.3.20}: a $Y$-bundle over
an annulus with monodromy the doubled map, with two
$N\times[0,1]$ pieces attached over opposite outer arcs.
The remaining outer arcs produce the twisted and identity
doubles, and the inner circle produces the mapping torus.
The chosen line lift on the annular bundle and the pullback
lines on the attached pieces agree on their attaching faces.
Thus this is an actual bordism carrying both the meridian line
and its tangent bundle, not merely a relation between abstract
intersection matrices. Its boundary orientations compare a cap
minus the identity cap with the mapping torus, with one sign
that is independent of the map. The identity cap
$S^2\times S^4$ and identity mapping torus have zero $x^3$
and zero $p_1x$. Applying the same sign to all three boundary
terms of the pair-of-pants bordism proves that $x^3$ and $p_1x$
are individually additive for fixed-tube-relative maps. Their
displayed integral combination is therefore additive as well.
In particular the proof does not infer additivity from a
three-knot group law. The line class is unambiguous here:
$H^1(Y;\Z)=0$, so changing its specified lift by a gauge does
not add a mixed mapping-torus Chern class.

\item It remains to check normalization choices. An orientation-
preserving equivariant sphere diffeomorphism can be made the
identity on a fixed tube by equivariant isotopies: its fixed
$S^2$ restriction is isotopic to the identity by
\citet[Theorem A]{Smale1959E}; its complex normal derivative
is homotopic to the identity since $\pi_2(U(2))=0$; and
tubular rescaling removes the remaining higher-order germ.
These are the local geometric steps, independent of any
three-knot invariant.

\item Two such normalizations give equivariantly diffeomorphic
clutched sphere pairs. Extend their boundary isotopy over a
collar to obtain the diffeomorphism. Compatible unitary normal
tubes along the isotopy give compatible smooth orbit structures.
Thus their orbit pairs are the same smooth pair up to
diffeomorphism. The fixed image is a three-sphere and its
oriented normal three-plane bundle is trivial. The two inherited
normal framings differ by a map $r:S^3\to SO(3)$.

\item We compute directly the change caused by this framing. If
$E$ is the exterior of the embedded three-sphere, its two
framed surgeries fill $E$ by $S^2\times D^4$ with boundary
identifications differing by
\[
 (a,z)\longmapsto(r(z)a,z),\qquad(a,z)\in S^2\times S^3.
\]
Cutting off the common exterior gives a characteristic bordism
between the difference of these two surgery caps and
\[
 Z_r=(S^2\times D^4)\cup_r-(S^2\times D^4),
\]
which is the oriented sphere bundle $S(\eta_r)\to S^4$ of
the rank-three bundle clutched by $r$. The meridian lines
extend over this bordism. They agree on the common boundary,
and their possible gauge discrepancy has a real logarithm
because $H^1(S^2\times S^3;\Z)=0$.

\item Let $\pi:Z_r\to S^4$ and choose $x$ with positive fibre
evaluation one. The oriented vertical tangent two-plane has
Euler class $2x$, since $H^2(Z_r;\Z)=\Z$ and its Euler
evaluation on the fibre is two. Its stable isomorphism
$T_{\mathrm{vert}}Z_r\oplus\R\simeq\pi^*\eta_r$ gives
\[
 4x^2=\pi^*p_1(\eta_r).
\]
Also $p_1(TS^4)=0$, so
\[
 p_1(TZ_r)=\pi^*p_1(\eta_r)=4x^2.
\]
Therefore $\langle p_1x-4x^3,[Z_r]\rangle=0$ directly.
No generator normalization of $\pi_3(SO(3))$ is needed.
The two framed-surgery indices are equal. This proves
normalization independence, and completes the direct
integer-additivity argument.

\end{enumerate}
\end{proof}




\begin{corollary}[a second determination by the spin$^c$ coboundary]\label{push:src:C.3.27}
\StatementRecord{push:src:C.3.27}{proved}{GPT-6 Astra (OpenAI, ChatGPT Pro mode), effort Pro, 2026-11-17; Claude Fable 5.2 (Anthropic), effort high, 2026-11-17; GPT-6 Astra (OpenAI, ChatGPT Pro mode), S6\_06\_BLIND, effort Pro, 2026-11-17; Claude Fable 5.2 (Anthropic), effort high, 2026-11-16; Claude Fable 5.2 (Anthropic), effort high, 16 November 2026}{the hypotheses, notation and definitions stated here; \cref{push:src:C.3.25}; \cref{push:src:C.3.26}; \cref{push:src:C.3.24}; \cref{push:src:C.3.18}; \cref{push:src:C.3.22}; \cref{push:src:C.3.10}; \cref{push:src:C.3.6}; \cref{push:src:C.1.2}; \cref{push:src:C.1.1}}

The actual positive-loop smooth monodromy satisfies
$[\Phi]=0\in\Theta_7$. This determination uses the explicit
eight-dimensional coboundary of \ref{push:src:C.3.25} and the direct index
additivity of \ref{push:src:C.3.26}. It does not use the three-knot integer,
Wall's formula for it, the Davis generator comparison, or the
conclusion of \ref{push:src:C.3.24}.
\end{corollary}


\begin{proof}
\begin{enumerate}[label=\arabic*.]
\item Apply \ref{push:src:C.3.26} to equivariant representatives of the actual return
maps. For the annular six-loop representative $\psi$, the two
geometric characteristic calculations, with their common
orientation sign, give
\[
 p_\psi=-24\epsilon_{\rm cap},\qquad
 k_\psi=-6\epsilon_{\rm cap}.
\]
These are the actual Pontryagin and cubic pairings proved in
\ref{push:src:C.3.18}--\ref{push:src:C.3.22}. Therefore the directly defined index is
\[
 \lambda(\psi)=\frac{p_\psi-4k_\psi}{24}=0
 \quad\hbox{in }\Z.
\]

\item We check the relation to one positive loop before reducing
modulo $28$. Choose an invariant connection on the descended
$w$-family. Its return over six consecutive positive loops is
the sixth power of its one-loop return. A positive loop changes
$c$ by $-2$. The coherent closing translations of Paper~I, Proposition~18.5 obey
\[
 b_{h,c}+b_{k,c+h}=b_{h+k,c};
\]
in particular the six successive positive-$2$ closings from
$c_0-12$ to $c_0$ equal the positive-$12$ closing. Interpolating
between invariant connections gives equivariant isotopies.
The changes of transport in \ref{push:src:C.3.10}, its unit and real-scalar
isotopies, and the lift comparisons of \ref{push:src:C.3.6} are equivariant
as well. Changing back to the stipulated closing choices
adds only a connected $\C^*$ automorphism commuting with the
circle action. Consequently the actual lift of $\psi$ and
$\Phi^6$ define the same equivariant mapping class.

\item If necessary, a lift of the orbit identity on the fixed tube
is made the identity there by its smooth circle gauge; the
gauge has a real logarithm and can be extended with a collar
cutoff. This is an equivariant isotopy and does not change the
orbit map. \ref{push:src:C.3.26} therefore applies with exactly the cap used
in the calculation above, and its normalization independence
gives
\[
 6\lambda(\Phi)=\lambda(\Phi^6)=\lambda(\psi)=0.
\]
The group here is $\Z$, so $\lambda(\Phi)=0$.

\item Now construct the actual eight-dimensional spin$^c$ coboundary
$Z_\Phi$ of \ref{push:src:C.3.25} for a normalized one-loop representative.
Its explicit intersection form and characteristic covectors
give, by the spin$^c$ index formula alone,
\[
 \mu(\Sigma_\Phi)
   \equiv\frac{p_\Phi-4k_\Phi}{672}
   =\frac{\lambda(\Phi)}{28}=0\pmod\Z.
\]
The positive generator in \ref{push:src:C.1.2} has invariant $1/28$, and
$\Theta_7$ has order $28$. Thus this invariant detects its
elements and $[\Sigma_\Phi]=0$. The sphere-gluing comparison
of \ref{push:src:C.1.1} gives $[\Phi]=0$. It also gives
$M_\Phi\cong S^1\times S^6$ with either spin structure having
zero invariant.

\end{enumerate}
\end{proof}


\begin{remark}[Independence and shared geometric inputs]
The first determination uses the integer of equivariant
orbit knots and its Wall--Davis comparison. The present
determination uses the explicit line-disc bundle and
four-handle coboundary, the spin$^c$ index formula, and an
independent integer-additivity argument from actual
seven-dimensional characteristic bordisms. Neither conversion
nor the integer-level passage from six loops to one in this
second argument invokes the first determination. The two
methods share the proved annular representative and the actual
geometric values $p_\psi=-24\epsilon_{\rm cap}$ and
$k_\psi=-6\epsilon_{\rm cap}$. Their independence is at the
boundary-invariant and one-loop arguments; they are not
calculations from disjoint geometric data. Independent
referees and exact arithmetic checks are validation, not
additional mathematical methods.

The smooth spin$^c$ coboundary here is different from the
specified analytic filling excluded by Paper~I, Theorem~21.58. Its construction
does not supply that filling or its almost-complex Chern data.
\end{remark}



\begin{remark}[Established surgery and index machinery]
The Hopf-handle construction used in \ref{push:src:C.3.25} is established
surgery machinery. In the proof of Lemma~3.5, Case~2,
\citet{Jiang2014E} attaches the same
$S^3_{\rm Hopf}\times D^4$ four-handle to a line-disc bundle,
expresses the resulting piece as $\mathbb{CP}^2\times D^4$,
and uses its trivial normal bundle to compute the spin$^c$
invariant. For his $k=0$ case the resulting value is
$p/(24\cdot28)$; Lemma~3.6 also exhibits the characteristic
integer $(p-4k)/24$ in his oriented decomposition conventions.
\ref{push:src:C.3.25} supplies the direct calculation with general cubic
$k$ and verifies that its boundary is the actual clutched
sphere of this family. No new general surgery construction
or invariant formula is claimed.

The relation between semifree actions, fixed three-sphere
knots, and surgery to homotopy complex projective three-spaces
is classical. The primary introduction of
\citet[pp.~125--126]{MontgomeryYang1968E} describes these
correspondences and the Pontryagin offset divisible by $24$.
The second proof here derives its needed additivity and
framing independence from the explicit bordisms and the
vertical Euler identity; it does not import the knot
classification as a premise. This identifies both the
standard literature and the exact scope of the derivation.
\end{remark}



\paragraph{The projection formula and its missing inputs.}
The current literature supplies another precise test for a proposed
orbit-knot computation. Theorem 4.1 of
\citet{GauniyalTurchin2026Root}, published online on
25 September 2026, expresses the Haefliger integer using both
double-point linking and invariants of the projected immersion.
The latter include a rationally normalized Smale invariant.
Consequently the double circle of \ref{push:src:C.3.17} alone is insufficient.
An application here would require an actual compressed projection
and its tangent-framing comparison. We opened the published
formula and its definitions, rather than importing the linking
term in isolation.

Likewise, \citet[Theorem 2.5]{Sakai2015Root} computes the
\emph{change} under a specified crossing change from oriented
double-preimage linking. To obtain an absolute value by that
route one must also identify the comparison embedding and its
class. Neither formula identifies the present trace with a
standard reference embedding. The characteristic-number cap
retains those geometric comparisons as explicit requirements.
Searches on 15 November 2026 for the six-sphere family,
its degeneration monodromy, and these knot formulas did not
locate a numerical determination for this particular family.
This states the search outcome.

\begin{proposition}[twist independence in the sphere cases]\label{push:src:C.4.1}
\StatementRecord{push:src:C.4.1}{proved}{GPT-6 Astra (OpenAI, ChatGPT Pro mode), effort Pro, 2026-11-17; GPT-6 Astra (OpenAI, ChatGPT Pro mode), S6\_06\_BLIND, effort Pro, 2026-11-17; Claude Fable 5.2 (Anthropic), effort high, 16 November 2026}{the hypotheses, notation and definitions stated here; \cref{push:src:C.3.24}}

For the fixed multiplicities $(3,4)$ and all admissible data in Paper~I, Proposition~4.5 and Paper~I, Proposition~4.6 and Paper~I, Proposition~4.7 and Paper~I, Proposition~4.8
with $|p|=1$, the monodromy class of the positive $w$-circle is the
same class in $\Theta_7$, with the complex orientation on each fibre.
\end{proposition}



\begin{proof}
\begin{enumerate}[label=\arabic*.]
\item The explicit sufficiency maps in E.1 depend holomorphically on a
small $c$-disc: the lattice defects and toric integral exponents are
constant, the periods and unit factors are holomorphic, and the
translation correction is solved in the fixed sheaf $V$ with
$H^1(B,V)=0$. This is the parameter-dependent assertion of Paper~I, Corollary~23.3,
not just the fibrewise classification. For the second twist
coefficient $s_2$, the parameter identification is
$W=e^{-\pi i(c+s_2)}=(-1)^{s_2}w$.
The rotation of the punctured disc preserves its positive
fundamental-group generator. Local family identifications differ
by elements of $\Aut(X)=\C^*$, a connected group, so their
transition changes are smoothly isotopic to the identity. Following
these identifications around the circle conjugates the monodromy
in the oriented mapping-class group. Since $\Theta_7$ is abelian,
the resulting class is unchanged. If the lattice sign in Paper~I, Proposition~4.5 and Paper~I, Proposition~4.6 and Paper~I, Proposition~4.7 and Paper~I, Proposition~4.8 is
$-1$, the comparison map is nevertheless holomorphic and preserves
the complex orientation; no extra sign is introduced.
The common class is zero by \ref{push:src:C.3.24}. This last numerical consequence
uses only the already proved equality of the classes, not an
additional deformation or conjugation argument.

\end{enumerate}
\end{proof}




\begin{proposition}[conjugation gives no torsion constraint]\label{push:src:C.4.2}
\StatementRecord{push:src:C.4.2}{proved}{GPT-6 Astra (OpenAI, ChatGPT Pro mode), effort Pro, 2026-11-17; Claude Fable 5.2 (Anthropic), effort high, 16 November 2026}{the hypotheses, notation and definitions stated here; the explicit constructions and calculations in the proof}

The antiholomorphic identification in Paper~I, Proposition~4.9 and Paper~I, Proposition~4.10 and Paper~I, Theorem~5.10 covers
$w\mapsto-\bar w$. At a fixed base point on the imaginary
$w$-axis it identifies the positive-loop monodromy with the
inverse-loop monodromy by an orientation-reversing fibre map.
Its action on $\Theta_7$ is negation, so the resulting equation is
the tautology $-[\Phi]=-[\Phi]$, not $2[\Phi]=0$.
\end{proposition}



\begin{proof}
\begin{enumerate}[label=\arabic*.]
\item From $c'=1-\bar c$ we obtain
$e^{-\pi ic'}=-\overline{e^{-\pi ic}}$.
Complex conjugation reverses the parameter circle. An
antiholomorphic map in complex dimension three reverses the real
orientation because its determinant has sign $(-1)^3$.
Conjugating a sphere clutching map by a reflection changes the
orientation of its glued homotopy sphere, hence negates its
$\Theta_7$ class. Equivalently the Eells--Kuiper invariant changes
sign. Applying this to the inverse loop proves the assertion.
The antiholomorphic symmetry therefore supplies no order-two restriction.

\end{enumerate}
\end{proof}




\begin{corollary}[Counterexample: the specified analytic filling premise]\label{push:src:C.5.1}
\StatementRecord{push:src:C.5.1}{false as stated; correction proved}{GPT-6 Astra (OpenAI, ChatGPT Pro mode), effort Pro, 2026-11-17; Claude Fable 5.2 (Anthropic), effort high, 16 November 2026}{the hypotheses, notation and definitions stated here; \cref{push:src:C.3.24}; \cref{push:src:C.3.25}; \cref{push:src:C.3.27}; \cref{push:src:C.4.1}; \cref{push:src:C.4.2}}

The descended analytic filling excluded by Paper~I, Theorem~21.58,
and its proposed resolution $W_8$ do not exist. Consequently
there are no values of $c_1,c_2,p_1,p_1^2,p_1c_1^2,c_1^4$,
relative lifts, intersection form or signature for that
specified analytic object. The actual positive-loop
monodromy is nevertheless
\[
                    [\Phi]=0\in\Theta_7\simeq\Z/28.
\]
\end{corollary}


\begin{proof}
\begin{enumerate}[label=\arabic*.]
\item A descended analytic filling of the specified family
would be a proper analytic filling before any resolution,
and its finite pullback would remain proper. Paper~I, Theorem~21.58
excludes both possibilities. Characteristic numbers of
an unrelated manifold cannot supply the requested values.

\ref{push:src:C.3.24} determines the actual monodromy by the equivariant
orbit-knot integer. \ref{push:src:C.3.25}--\ref{push:src:C.3.27} construct an actual
smooth spin$^c$ coboundary of the clutching sphere and
mapping torus and give a second determination by the
index formula. Both use the verified actual pairings
$p_\psi=-24\epsilon_{\rm cap}$ and
$k_\psi=-6\epsilon_{\rm cap}$ and pass from six loops
to one in an integer group before reduction modulo $28$.
Their distinct boundary-invariant arguments and shared
geometric inputs are stated in \ref{push:src:C.3.27}. \ref{push:src:C.4.1} gives the
same zero for all the compared sphere twist data, and
\ref{push:src:C.4.2} gives the antiholomorphic transformation law.

\item The smooth coboundary of \ref{push:src:C.3.25} is fully specified and
has the intersection and spin$^c$ characteristic data
computed there. No complex structure on it is asserted,
and its determinant class is not identified with the first
Chern class of an almost-complex tangent bundle. This
distinction is essential now that the prescribed analytic
construction has been refuted.

\end{enumerate}
\end{proof}


This subsection settles B.6 for the original family and every
sphere-producing higher-multiplicity family of Paper~I, Proposition~23.1 and Paper~I, Theorem~23.2. We use the complex
orientation of $X$, the positive loop in the $w$-plane, and the generator
of $\Theta_7$ with Eells--Kuiper invariant $1/28$, as in the conventions
and the supplied result \cref{push:src:C.1.1} and \cref{push:src:C.1.2} and \cref{push:src:C.2.1}. Write $a=[\Phi]\in\Z/28$. The mapping torus
has the oriented diffeomorphism type $(S^1\times S^6)\#\Sigma_\Phi$.
The reduction from this mapping torus to the glued sphere, and the
independence of its Eells--Kuiper invariant from the two spin structures,
are those of \cref{push:src:C.1.1} and \cref{push:src:C.1.2} and \cref{push:src:C.2.1}; no different convention is introduced here.

\begin{proposition}[The sixfold congruence]\label{maps:src:B.6.1}
\StatementRecord{maps:src:B.6.1}{proved}{GPT-6 Astra (OpenAI, ChatGPT Pro mode), effort Pro, 2026-11-17; Claude Fable 5.2 (Anthropic), effort high, 2026-11-15; GPT-6 Astra (OpenAI, ChatGPT Pro mode), S6\_06\_BLIND, effort Pro, 2026-11-17; Claude Fable 5.2 (Anthropic), effort high, 16 November 2026}{the hypotheses, notation and definitions stated here; \cref{push:src:C.3.4}}


If the monodromy of the sixfold parameter pullback has computed class
$b=6a\in\Z/28$, then $b$ is even and, writing $b=2r$, the original
class satisfies
\[
 a\equiv5r\pmod {14}.
\]
There are exactly two residues modulo $28$ satisfying this condition,
and they differ by $14$. Without a value of $b$, this assertion excludes
none of the $28$ possibilities for $a$.
\end{proposition}


\begin{proof}
\begin{enumerate}[label=\arabic*.]
\item Pulling a mapping torus back by a degree-six map of its oriented base
circle replaces its monodromy by $\Phi^6$. Thus its class is $6a$.
The kernel of multiplication by six on $\Z/28$ is $\{0,14\}$ and
its image consists of the even residues. Dividing $6a=2r$ by two gives
$3a=r$ modulo $14$; multiplication by five inverts three modulo $14$.
Every $a$ has an image under this multiplication, so the conditional
statement does not itself restrict $a$. This is the exact arithmetic
content of \cref{push:src:C.3.4}, not a computation of the sixfold monodromy.

\end{enumerate}
\end{proof}




\begin{theorem}[complete answer to B.6]\label{maps:src:B.6.2}
\StatementRecord{maps:src:B.6.2}{proved}{GPT-6 Astra (OpenAI, ChatGPT Pro mode), effort Pro, 2026-11-17; GPT-6 Astra (OpenAI, ChatGPT Pro mode), S6\_06\_BLIND, effort Pro, 2026-11-17; Claude Fable 5.2 (Anthropic), effort high, 16 November 2026}{Definitions and the displayed constructions}


The original class is $a=0$, and the corresponding class for
every sphere-producing family $X_{3k,4l}$ is also zero for
the positive loop with period $2kl$.  The original computation
is \ref{maps:src:B.6.21}; the actual higher-family replacement and cup
calculation are \ref{maps:src:B.6.22}--\ref{maps:src:B.6.23}, and the complete spin-trace
comparison proving the assertion is \ref{maps:src:B.6.24}.  Corollary \ref{maps:src:B.6.25}
then gives an explicit spin filling of every resulting
mapping torus, with one zero-handle and one one-handle.
Propositions \ref{maps:src:B.6.3}--\ref{maps:src:B.6.6} give a separate analytic conifold
reduction for every pair; they do not by themselves evaluate
or compare the numerical classes.  The following records retain
the earlier unsuccessful parameter-filling approaches, followed
by the completed smooth calculation for all requested pairs.

The strongest explicit closing datum available is Paper~I, Proposition~18.5. With
$\ell=(2\pi i)^{-1}\log(F/G)$, $F^3=t$, $G^4=t-1$, its coherent
even-shift comparison has translation
\[
 b_{h,c}=-\frac h{12}\Pi_c\widehat\gamma+
          \left(h\ell-\frac{h^2}{24}\right)\widehat\delta,
 \qquad h\in2\Z.
\]
For the sixfold closing map, $h=12$, subtracting the integral period
$-\Pi_c\widehat\gamma$ produces the scalar $t^4/(t-1)^3$ on the
ordinary locus. At the cusp the same comparison is a holomorphic
toric unit. These representatives differ by one
$\widehat\gamma$ deck transformation. This proves equality of the
endpoint maps, but scaling the logarithms separately does not give
an isotopy on the overlap: an intermediate comparison would require
a nonintegral multiple of that deck transformation. Such a multiple
is not an identification in the lattice quotient. We have therefore
not replaced this mismatch by a purported global scalar gauge.

The CW model of B.5 gives the actual underlying topological attachment
and the integral maps needed for the homotopy-sphere calculation. Its
use of the positive toric quotient and a relative product
homeomorphism does not supply the missing smooth framing of the
parameter-dependent attachment. Carrying its cellular boundary
matrices around the parameter circle would give homology information,
but would not distinguish the $28$ smooth mapping classes.

We also checked the real representatives and their cap transport in
A.3. For the positive-imaginary parameter, the actual order-three
transport is $C=(A_2^{-1}A_1^{-1}A_2)^2$: the two reflection matrices
have product $B_1=A_2^{-1}A_1^{-1}A_2$, and conjugating one reflection
to the other requires $2k=1$ modulo three. Thus $k=2$ is forced.
All four reflection lifts at the multiplicity-four fibre were checked
independently in \texttt{maps\_real\_audit.py}. These verifications
fix the real attachment markings but do not determine the components
of the resulting orientation-reversing diffeomorphisms in the full
diffeomorphism group. In particular, their fixed-set topology alone
does not evaluate their product in $\Theta_7$.
The precise comparison theorem is
\cite[Theorem 1.3]{mapsOlbermann10}: a fixed three-sphere has even
Haefliger knot class, and each even knot supports four involutions up
to equivariant diffeomorphism relative to that sphere. Thus even the
ambient knot class, if computed, would leave additional involution data.

A further possible comparison is the explicit smooth plumbing
relation of Krylov \cite[Theorem 3]{mapsKrylov03}, used in
\cite[Theorem 4.1]{mapsRandalWilliams25}. On $S^3\times S^3$, the
quaternionic maps $Y(u,v)=(u,uv)$ and $U(u,v)=(v^{-1}u,v)$ satisfy
$(YUY)^4=\Sigma_{\rm Milnor}^{-1}$. Their relation detects a smooth
sphere class that their homology matrices do not detect. To apply
this relation here requires an actual relative smooth identification
of our parameter transport with a word in such plumbing maps; no such
identification has been proved. The finite semistable models are treated in
Paper~I, Proposition~21.1 and Paper~I, Theorem~21.2 and Paper~I, Proposition~21.3 and Paper~I, Theorem~21.6.
Their component monodromy
obstruction prevents gluing them to the plain cusp models and hence
prevents using those particular models as a bounding manifold.
These explain why those earlier approaches did not evaluate $a$.
The completed fixed-circle calculation below does not require those
unproved identifications.


\end{theorem}




\begin{proposition}[the smoothing monodromy is Torelli]\label{maps:src:B.6.3}
\StatementRecord{maps:src:B.6.3}{proved}{GPT-6 Astra (OpenAI, ChatGPT Pro mode), effort Pro, 2026-11-17; Claude Fable 5.2 (Anthropic), effort high, 16 November 2026}{the hypotheses, notation and definitions stated here; \cref{maps:eq:B.6.3}; \cref{maps:src:D.3.2}; \cref{maps:eq:B.6.5}}


Fix the sphere twist data and a sufficiently small nonzero smoothing
parameter $s$ in D.3. The even-shift comparisons extend to the joint
moving-zero family. Its monodromy $\Psi_s$ around the positive $w$
loop acts as the identity on $H_3(Y_{u,s};\Z)\simeq\Z^2$.
Thus the conifold route places this monodromy in the Torelli subgroup
of the mapping class group of $S^3\times S^3$; its homology action
does not determine its remaining smooth class.
\end{proposition}


\begin{proof}
\begin{enumerate}[label=\arabic*.]
\item The moving section multiplies the linearization by $1-st$. On an
additive period cover set
\[
 L_s=(2\pi i)^{-1}\log(1-st),\qquad
 \beta_{c,s}=\beta_c-L_s,\qquad
 \ell_s=\ell-L_s/12.
\]
The sign follows directly from the multiplicative period
$\exp(-2\pi i\beta_c)$ of the $\widehat\gamma$ generator.
The same formula for the even-shift comparison is
\[
 b_{h,c,s}=-\frac h{12}\Pi_{c,s}\widehat\gamma+
             \left(h\ell_s-\frac{h^2}{24}\right)\widehat\delta.
\]
The two added $L_s$ terms cancel, so this additive vector equals
$b_{h,c}$ on the common cover. Near the finite fibres $1-st$ is a
holomorphic unit and the checked finite collar proof is unchanged.
At the new zero $t=1/s$, the logarithmic long period cancels the
logarithm in $\ell_s$, giving an additive translation holomorphic
on the transverse elliptic factor and a holomorphic unit in the
vanishing toric coordinate. At infinity the removed logarithm in
$\beta_{c,s}$ is matched by the removed logarithm in $\ell_s$;
the same assertion holds for the other rank-one degeneration.
These are translations preserving the primitive rank-one fans, so
the maps and their inverses extend on their toric charts. The finite
affine defects and the composition law are the same as before because
the additive vector is unchanged. This constructs the even-shift
comparisons on the moving-zero family, not just on its ordinary locus.
For $h=12$ the scalar representative is now
\[
 \frac{t^4}{(t-1)^3(1-st)},
 \qquad \operatorname{div}=4[0]-3[1]-[1/s].
 \tag{B.6.1}\label{maps:eq:B.6.1}
\]
In particular its pole has moved from infinity to the new zero.

\item Write $E=d\otimes\gamma$. The four fibre monodromies, in order
with product one, are
\[
 A_1,\quad A_2,\quad C_w=I+w\otimes u^*,\quad C_d=I-E.
 \tag{B.6.2}\label{maps:eq:B.6.2}
\]
Indeed the cusp period formerly containing $-\log q/(2\pi i)$
now contains $-\log(q-s)/(2\pi i)$, whereas the other nodal period
still contains $\log q/(2\pi i)$. Thus the two primitive rank-one
monodromies are separated, and $C_wC_d=M_0$. The parameter return
has linear part $S=I-2E$; its affine translation acts trivially on
fibre cohomology. All four matrices commute with $S$.

\item Let $j$ include the complement of these four values in $B$, and let
$V$ be the rational degree-one cohomology local system. For each
$q=0,\ldots,4$ we have
\[
 R^q g_*\Q=j_*\bigwedge^qV.
 \tag{B.6.3}\label{maps:eq:B.6.3}
\]
At a finite fibre this follows by transfer from its finite free torus
cover: the central and nearby invariant groups agree. At a primitive
rank-one fibre, its normalization is a ruled surface over an elliptic
curve, with degree-zero normal line bundle and two sections glued by
a translation. Both the line bundle and that translation are
topologically trivial after choosing a trivialization and a translation
isotopy. Its specialization is therefore the product calculation for
a nodal two-torus and a smooth two-torus. In the nodal factor the
groups in degrees zero, one and two map isomorphically to the
monodromy-invariant groups; the product calculation proves \eqref{maps:eq:B.6.3}
in all degrees. This argument only asserts the rational comparison
needed below, and includes both separated rank-one fibres.

\item Here is an explicit calculation of the middle sheaf cohomology. Put
$G_i=\bigwedge^2(A_i^{-t})$ at the two finite values, and use
$\bigwedge^2(C_w^{-t}),\bigwedge^2(C_d^{-t})$ at the other two.
Parabolic cocycles are tuples $(z_i)$ with
$z_i\in\operatorname{im}(G_i-I)$ and
\[
 z_1+G_1z_2+G_1G_2z_3+G_1G_2G_3z_4=0,
\]
modulo tuples $((G_i-I)v)_i$. This computes
$H^1(B,j_*\bigwedge^2V)$: the local condition is exactly zero
restriction to each punctured-disc cohomology group.
In the exterior basis of covectors
$(\gamma u^*,\gamma w^*,\gamma d^*,u^*w^*,u^*d^*,w^*d^*)$,
where juxtaposition denotes exterior product, two representatives
are the columns of the following four matrices:
\[
 P_1=\begin{pmatrix}1&-8\\2&-7\\0&-6\\0&1\\-2&1\\-1&-1\end{pmatrix},
 \quad
 P_2=\begin{pmatrix}6&-3\\3&0\\6&0\\-1&0\\-1&-1\\1&-1\end{pmatrix},
 \quad
 P_3=\begin{pmatrix}0&0\\0&0\\0&0\\0&0\\0&-1\\0&0\end{pmatrix},
 \quad P_4=0.
 \tag{B.6.4}\label{maps:eq:B.6.4}
\]
The four local image ranks are $4,4,2,2$. The relation matrix
has rank five, and the coboundary matrix has rank five; the two
displayed cocycles are independent modulo its image. Hence the
parabolic group has dimension two and these form a basis. Direct
multiplication gives, for its two basis cocycles $p_1,p_2$,
\[
 (\bigwedge^2S^{-t}-I)p_1=d(-2\gamma u^*),\qquad
 (\bigwedge^2S^{-t}-I)p_2=d(-2\gamma w^*).
 \tag{B.6.5}\label{maps:eq:B.6.5}
\]
Its induced parameter action is therefore the identity. All the
matrices and rank assertions, in all exterior degrees, are checked
exactly in \texttt{maps\_smoothing\_monodromy.py}.

\item The middle parabolic term is $E_2^{1,2}$ in the Leray spectral
sequence and survives to $E_\infty$: the base has real cohomological
dimension two. It already has dimension two, which is the full
$b_3(Y)$ proved in \ref{maps:src:D.3.2}. Consequently the other filtration quotients
of $H^3(Y;\Q)$ vanish and the natural filtration identifies this
group with $E_\infty^{1,2}$. Naturality identifies its parameter
action with \eqref{maps:eq:B.6.5}. Since $H^3(Y;\Z)$ is torsion free, identity
over $\Q$ implies identity over $\Z$, and duality gives the claim
for homology.

\end{enumerate}
\end{proof}




\begin{proposition}[the normal-framing obstruction vanishes]\label{maps:src:B.6.4}
\StatementRecord{maps:src:B.6.4}{proved}{GPT-6 Astra (OpenAI, ChatGPT Pro mode), effort Pro, 2026-11-17; Claude Fable 5.2 (Anthropic), effort high, 16 November 2026}{the hypotheses, notation and definitions stated here; \cref{maps:src:B.6.3}; \cref{maps:eq:B.6.6}; \cref{maps:eq:B.6.7}; \cref{maps:eq:B.6.8}; \cref{maps:eq:B.6.9}; \cref{maps:eq:B.6.10}; \cref{maps:src:D.3.2}; \cite[Section 2.1, equation (2)]{mapsKrylov03}; \cite[proof of Theorem 2]{mapsKrylov03}}


For the monodromy $\Psi_s$ in \ref{maps:src:B.6.3}, let $M_s$ be its oriented mapping
torus. Then $p_1(M_s)=0$ in $H^4(M_s;\Z)$. Its Torelli
normal-framing invariant is zero. Consequently $[\Psi_s]$ belongs
to the injected subgroup $\Theta_7$ of
$\pi_0\operatorname{Diff}^+(S^3\times S^3)$.
This statement does not evaluate that central class or identify it
with $[\Phi]$ through the conifold surgery without further relative
framing information.
\end{proposition}


\begin{proof}
\begin{enumerate}[label=\arabic*.]
\item We first compute the two swept critical-curve classes. Use $k$ for
the parameter generator and $a_j$ for the finite base generators in
the boundary markings of B.5. Their conjugation defects are
\[
 k a_j k^{-1}=t_{d_j}a_j,\qquad
 d_1=(0,1,-1,-1)^t,\quad d_2=(0,0,-1,1)^t.
 \tag{B.6.6}\label{maps:eq:B.6.6}
\]
Indeed the finite return is $x\mapsto Sx+\beta_j$, where the
noninvariant part of $\beta_j$ is $-g/6$. The remaining part is
fixed by $A_j$: it is the real inverse image of a complex multiple
of the invariant vector $d$. Conjugating the affine finite generator
$x\mapsto A_jx+v_j/m_j$ therefore gives
\[
 d_j=(S-I)v_j/m_j-(I-A_j)g/6.
\]
Substitution gives \eqref{maps:eq:B.6.6} and $d_1+A_1d_2=0$. At the two
rank-one values the holomorphic unit representatives have
single-valued additive lifts in the toric and transverse coordinates;
the compatible zero-section boundary marking gives zero $k$-defect.

\item The regular five-dimensional fibre $N$ of $M_s\to B$ has fundamental
group $\Lambda\rtimes_S\Z$. It is the product of a three-dimensional
nilmanifold and the $(u,w)$ two-torus. Its invariant one-forms can be
written
\[
 \gamma,\upsilon,\omega,\vartheta,\kappa,
 \qquad d\vartheta=2\kappa\wedge\gamma,
 \qquad d\gamma=d\upsilon=d\omega=d\kappa=0.
 \tag{B.6.7}\label{maps:eq:B.6.7}
\]
The first three closed forms are dual to $g,u,w$, and $\kappa$ is
the parameter class. These formulas follow either from
$[k,g]=-2d$ or from the connection form of the central circle bundle.
The finite base actions on its Lie lattice are
\[
 \widetilde A_j=
 \begin{pmatrix}A_j&-d_j\\0&1\end{pmatrix};
\]
the two other actions are the corresponding block lifts of $C_w,C_d$.
For the indicated cohomology classes this gives
\[
 (C_w^*-I)[\gamma\omega]=-[\gamma\upsilon],\qquad
 (\widetilde A_2^*-I)[\gamma\upsilon]
       =[\gamma\omega]-[\gamma\upsilon].
 \tag{B.6.8}\label{maps:eq:B.6.8}
\]
Here the star denotes cohomological transport, namely inverse
transpose on the one-forms in the chosen homology marking.
In the second identity the additional form $-\gamma\kappa$ is
exact by \eqref{maps:eq:B.6.7}. Thus both $[\gamma\upsilon]$ and
$[\gamma\omega]$ vanish in the global coinvariants of $H^2(N;\Q)$.
The full rational cohomology and these exact actions are computed in
\texttt{maps\_mapping\_torus\_cohom.py}; only the two identities
\eqref{maps:eq:B.6.8} are needed for the argument.

\item Let $p_s=1/s$. The central surface of the rank-one filling there is
the degree-zero ruled surface over the $(u,w)$ elliptic curve with
its sections identified by translation. The parameter return acts
on its normalization by a translation and a constant fibre scaling.
Both admit paths of such maps commuting with the conductor
identification. Its local mapping torus consequently has the
cohomology of this central surface times a circle. In degree three
it contains two independent classes
\[
 z_u=\kappa\gamma\upsilon,\qquad
 z_w=\kappa\gamma\omega.
 \tag{B.6.9}\label{maps:eq:B.6.9}
\]
Here $\gamma$ is the nodal graph-circle class. A toroidal
neighbourhood retracts to the central surface, by radial collapse
in its node and core charts; choosing the collapse continuously
around the parameter circle identifies these with classes of the
whole local mapping torus. Their restrictions to a regular $N$
are zero by \eqref{maps:eq:B.6.7}. They therefore define independent skyscraper
kernel sections in the third direct image sheaf for $M_s\to B$,
supported at $p_s$.

\item We compute their Leray differential, including its boundary term.
On the boundary of this filling, add the closed base-meridian form
$\alpha$. The two commuting monodromies $C_d$ and $S=C_d^2$ give
\[
 d\vartheta=(\alpha+2\kappa)\wedge\gamma.
\]
Consequently the restrictions of \eqref{maps:eq:B.6.9}, in the Wang summand for
the base boundary circle, are respectively
\[
 -\tfrac12[\gamma\upsilon],\qquad
 -\tfrac12[\gamma\omega].
 \tag{B.6.10}\label{maps:eq:B.6.10}
\]
For example
$\kappa\gamma\upsilon=-\alpha\gamma\upsilon/2+
d(\vartheta\upsilon)/2$ on that boundary. All statements here are
over $\Q$, so the division by two is legitimate.

\item For clarity, the comparison of this Wang term with the global
$d_2$ is the following cochain construction. Cover $B$ by a small
disc at $p_s$ and its complement, and then refine the complement
by the other three discs and a graph. A local degree-three class
with zero regular-fibre restriction restricts on the annulus to
its degree-two Wang class. The connecting homomorphism of the
base cover sends this to $H^2(B,R^2)$, which is exactly its Leray
$d_2$. The natural generization $R^2\to j_*H^2(N;\Q)$ has kernel
and cokernel supported at finitely many points. Their positive
cohomology vanishes, so
\[
 H^2(B,R^2)=H^2(B,j_*H^2(N;\Q))
          =H^2(N;\Q)_{\pi_1(B^*)}.
\]
In the graph-and-discs cochain model the connecting map is the
natural map from the local coinvariants to these global
coinvariants, with the sign fixed by the chosen boundary
orientation. Equations \eqref{maps:eq:B.6.8}--\eqref{maps:eq:B.6.10} prove $d_2z_u=d_2z_w=0$.
There is no later differential because the base has cohomological
dimension two.

\item By \ref{maps:src:B.6.3} and the Wang sequence, $b_3(M_s)=2$. The two surviving
classes just constructed therefore exhaust its rational third
cohomology. Their local restrictions are zero at infinity. At
$p_s$ they are \eqref{maps:eq:B.6.9}, whose restrictions to the swept critical
elliptic curve vanish: the graph-circle class is zero on each
elliptic fibre, so its restriction on the swept three-manifold is
a multiple of $\kappa$, and wedging with $\kappa$ kills it.
Thus both swept critical three-manifolds pair to zero with all of
$H^3(M_s;\Q)$. Their classes in $H_3(M_s;\Q)$ are zero.

\item It remains to relate these classes to $p_1$. Let $\mathcal E$ be
the complex rank-three tangent bundle along the fibres of the
parameter family, restricted to $M_s$. Factor the differential
of the map to $B$ by its divisorial zeros at the two multiple
fibres; its target is then a line bundle $\mathcal L$. The
saturated vertical tangent sheaf $\mathcal V$ is a rank-two
bundle. At either rank-one critical curve the local equation is
$xy=q$ times an elliptic coordinate, and its kernel is freely
generated by
$x\partial_x-y\partial_y$ and the elliptic tangent direction.
The invariant vector field $d$ is nowhere zero as a section of
$\mathcal V$: it is the first generator at $p_s$ and the second
at infinity. Hence $\mathcal V$ is an extension of a line bundle
by the trivial line, and $c_2(\mathcal V)=0$.

\item Write $\mathcal C$ for the disjoint union of the two swept critical
curves. The local ideal of each is $(x,y)$, giving the global
exact sequence, or equivalently its restricted topological $K$-class,
\[
 0\longrightarrow\mathcal V\longrightarrow\mathcal E
 \longrightarrow\mathcal L\otimes I_{\mathcal C}\longrightarrow0,
 \qquad
 [\mathcal E]=[\mathcal V]+[\mathcal L]-[\mathcal L|_{\mathcal C}].
\]
The Wang sequence and \ref{maps:src:D.3.2} give $H^2(M_s;\Z)=0$. All the line
first Chern classes in this identity vanish. The degree-four
Chern-character term of the final sheaf is the oriented dual of
$\mathcal C$, as follows from its codimension-two Koszul resolution.
It follows that
\[
 c_2(\mathcal E)=\operatorname{PD}[\mathcal C],\qquad
 p_1(M_s)=-2\operatorname{PD}[\mathcal C]=0
       \quad\hbox{in }H^4(M_s;\Q).
\]
Again by the Wang sequence $H^4(M_s;\Z)\simeq\Z^2$, so this
vanishing is integral.

\item Finally apply the Torelli exact sequence quoted from Kreck in
\cite[Section 2.1, equation (2)]{mapsKrylov03}. For $S^3\times S^3$
its left term is the injected $\Theta_7$, because this manifold
bounds the framed manifold $S^3\times D^4$. The right term is
the homomorphism from $H_3$ to the image of
$\pi_3SO(3)\to\pi_3SO(4)$. Its mapping-torus characteristic class
is detected by $p_1$: the stable clutching bundle of a generator
of $\pi_3SO(3)$ has first Pontryagin number of absolute value four,
as computed in \cite[proof of Theorem 2]{mapsKrylov03}. Thus the
zero $p_1$ class forces the entire integral normal-framing
homomorphism to vanish. Exactness gives the assertion.

\end{enumerate}
\end{proof}




\begin{proposition}[the conifold comparison preserves the sphere class]\label{maps:src:B.6.5}
\StatementRecord{maps:src:B.6.5}{proved}{GPT-6 Astra (OpenAI, ChatGPT Pro mode), effort Pro, 2026-11-17; Claude Fable 5.2 (Anthropic), effort high, 16 November 2026}{the hypotheses, notation and definitions stated here; \cref{maps:src:B.6.4}; \cref{maps:src:B.6.3}; \cref{maps:eq:B.6.12}; \cite[Section 1.4, equations (2)--(3)]{mapsCrowleyNordstrom19}; \cref{maps:eq:B.6.14}}


Let $a\in\Theta_7$ be the original six-sphere monodromy class,
and let $b\in\Theta_7$ be the central class of the smoothing
monodromy supplied by \ref{maps:src:B.6.4}. With the same positive parameter
loop and the complex orientations, $a=b$. Equivalently, if
$M_\Phi$ and $M_s$ are their seven-dimensional mapping tori, then
\[
 \mu(M_\Phi)=\mu(M_s)=\frac{a}{28}\quad\hbox{in }\Q/\Z.
 \tag{B.6.11}\label{maps:eq:B.6.11}
\]
Thus no additional unknown framing correction is introduced by
the reverse conifold surgery. This identifies the remaining
invariant; it does not evaluate it.

\end{proposition}


\begin{proof}
\begin{enumerate}[label=\arabic*.]
\item We first construct the surgery trace around the whole parameter
circle, including its gluing. The even-shift comparisons in \ref{maps:src:B.6.3}
preserve the actual smoothing parameter $s$, not merely its zero
set. In the node chart, the single-valued substitution $x'=x/u$
for $u\ne0$ makes the equation
\[
 x'y-zv=-s.
\]
The relative holomorphic Morse coordinates consequently have
transition germs preserving this fixed quadratic function.
There is no conformal scalar multiplying the quadratic form.
Such a germ $F$ retracts to its derivative by
$F_r(z)=r^{-1}F(rz)$, $0<r\leq1$, with $F_0=dF_0$;
each germ in this deformation still preserves the quadratic
function. The chosen small resolution is part of the family,
so these derivatives preserve its ruling. They therefore belong
to $SO(4,\mathbb C)$ rather than the component which exchanges
the two rulings. Reduce this group to its maximal compact
$SO(4)$ and choose small invariant collars. The deformation of
germs gives the same boundary bundles up to isotopy, and collar
isotopies extend to their adjacent pieces.

\item Here is an explicit equivariant handle trace for this reduction.
Write $D^3\subset\operatorname{Im}\mathbb H$ and
$D^4\subset\mathbb H$. The action
\[
 (g,h)(v,q)=(gvg^{-1},gqh^{-1}),\qquad
 (g,h)\in SU(2)\times SU(2),
 \tag{B.6.12}\label{maps:eq:B.6.12}
\]
on $D^3\times D^4$ factors through $SO(4)$, since $(-1,-1)$
acts trivially. Its boundary faces are the two local surgery
models. On $S^2\times D^4$, choose the identification of the
twistor two-sphere with the exceptional $\mathbb{CP}^1$ for
which its stabilizer line is the tautological line of degree
minus one. With this convention the fibrewise complex structure
$J_v(q)=vq$ identifies the real normal bundle with
$\mathcal O(-1)\oplus\mathcal O(-1)$ over $\mathbb{CP}^1$:
at $v=i$ the stabilizer $U(1)$ acts by scalar multiplication
on the complex two-dimensional fibre, and its associated
tautological line has degree minus one. On $D^3\times S^3$,
the map $(v,q)\mapsto vq\in T_qS^3$ identifies the same action
with that on the cotangent-disc model of the Milnor fibre.
Their common face is $S^2\times S^3$, with the common action
induced on the quadric link. Identify both faces with the
given complex-oriented local models along that link; reversing
the twistor convention would require reversing its sphere
identification too. The outward normal conventions
are those of replacing $S^2\times D^4$ by $D^3\times S^3$.
Thus \eqref{maps:eq:B.6.12}, together with the unchanged complement, builds
a bundle of oriented seven-dimensional surgery traces over
the parameter circle. Denote its total space by $Z$. Its
boundary is $M_s\sqcup(-M_\Phi)$.

\item For one trace $W$, the handle relative to $S^6\times I$ has
index three. Its cellular description gives
\[
 H^3(W;\Z)=\Z,\qquad H^4(W;\Z)=0,
 \qquad H^6(W;\Z)=\Z.
\]
The restriction $H^3(W;\Z)\to H^3(S^3\times S^3;\Z)$
is a primitive injection. Indeed the reverse trace has one
index-four handle, so its homology sequence gives a surjection
$H_3(S^3\times S^3)\to H_3(W)=\Z$ with primitive rank-one
kernel; dualize this sequence. \ref{maps:src:B.6.3} then implies that trace
monodromy acts as the identity on $H^3(W;\Z)$. The Wang
sequence, naturally with respect to its smoothing boundary,
now gives a commutative square
\[
 \begin{array}{ccc}
 H^4(Z;\Z)&\longrightarrow&H^4(M_s;\Z)\\
 \big\downarrow\scriptstyle\simeq&&
       \big\downarrow\scriptstyle\simeq\\
 H^3(W;\Z)&\longrightarrow&H^3(S^3\times S^3;\Z).
 \end{array}
 \tag{B.6.13}\label{maps:eq:B.6.13}
\]
The top restriction is therefore an injection of a rank-one
free group. Its value on $p_1(Z)$ is $p_1(M_s)=0$, by \ref{maps:src:B.6.4},
so $p_1(Z)=0$. Moreover the exact cohomology sequence of
$(Z,\partial Z)$ shows that
$H^4(Z,\partial Z;\Q)\to H^4(Z;\Q)$ has zero image, since
restriction even to the single component $M_s$ is injective.
The intersection form defining the signature has consequently
zero rank, and $\sigma(Z)=0$.

\item The trace $W$ is simply connected and spin: the index-three
attachment extends the unique spin structure, as its attaching
two-sphere has no first cohomology. Trace monodromy preserves
this unique structure. Either of its lifts makes $Z$ spin,
with compatible induced boundary spin structures. For a spin
seven-manifold with rationally zero first Pontryagin class,
the Eells--Kuiper invariant is
\[
 \mu(M)=\frac{p_1(V)^2-4\sigma(V)}{896}\pmod{\Z},
 \tag{B.6.14}\label{maps:eq:B.6.14}
\]
using a spin coboundary $V$ and a rational relative lift of
$p_1(V)$. The square is independent of the lift. This definition
does not require $H^4(M;\Q)=0$; see
\cite[Section 1.4, equations (2)--(3)]{mapsCrowleyNordstrom19},
whose normalization is $28$ times \eqref{maps:eq:B.6.14}. Applying its gluing
formula to $Z$, with the zero lift of $p_1(Z)$, proves
$\mu(M_s)-\mu(M_\Phi)=0$.

\item Finally, an oriented diffeomorphism supported in a six-disc
changes its mapping torus from the product by connected sum
with its corresponding homotopy seven-sphere. This follows
directly by cutting out that disc times a small parameter
interval and regluing its two faces. Every class of
$\pi_0\operatorname{Diff}^+(S^6)$ has such a representative,
and \ref{maps:src:B.6.4} provides such a representative for $\Psi_s$ as well.
The product mapping tori have zero $\mu$: they bound respectively
$D^7\times S^1$ and $S^3\times D^4\times S^1$, with zero
Pontryagin square and signature. These bounds work for either
circle spin structure. Additivity of \eqref{maps:eq:B.6.14}, and its
identification $\Theta_7\simeq\frac1{28}\Z/\Z$, prove
\eqref{maps:eq:B.6.11} and $a=b$.

\end{enumerate}
\end{proof}




\begin{proposition}[the reduction for all coprime multiplicities]\label{maps:src:B.6.6}
\StatementRecord{maps:src:B.6.6}{proved}{GPT-6 Astra (OpenAI, ChatGPT Pro mode), effort Pro, 2026-11-17; Claude Fable 5.2 (Anthropic), effort high, 16 November 2026}{the hypotheses, notation and definitions stated here; \cref{maps:src:D.3.2}; \cref{maps:eq:B.6.2}; \cref{maps:eq:B.6.3}; \cref{maps:eq:B.6.5}; \cref{maps:eq:B.6.8}; \cref{maps:src:B.6.4}; \cref{maps:src:B.6.5}}


For every sphere member $X_{3k,4l}$ of Paper~I, Proposition~23.1 and Paper~I, Theorem~23.2, the monodromy associated
to the positive generator $2kl$ has the same class in $\Theta_7$
as a central monodromy on its conifold smoothing $S^3\times S^3$.
The smoothing monodromy acts trivially on integral $H_3$, and
its seven-dimensional mapping torus has $p_1=0$. The numerical
values of these central classes, and any relation between
different pairs $(k,l)$, remain undetermined here.

\end{proposition}


\begin{proof}
\begin{enumerate}[label=\arabic*.]
\item Put $h=2kl$. The D.3 construction applies with the finite
replacements of orders $3k,4l$: their affine actions are still
free, and multiplying the linearization by $1-st$ is a unit
change on each of their fixed root discs. The node and its
small resolution are unchanged. The one-surgery proof of \ref{maps:src:D.3.2}
therefore identifies the smoothing with $S^3\times S^3$.

\item We spell out why the comparison from Paper~I, Proposition~23.1 and Paper~I, Theorem~23.2 survives the separated
cusps. Use the fixed-data case of its proof, so the first
coordinates of its rational finite translation vectors $b_j$,
integral defects $\ell_j$, and connector vectors $d_j$ are all
zero. Changing $\beta$ to $\beta-L_s$ changes only the
$\widehat\gamma$ column of the period matrix. It consequently
leaves these finite translation formulas unchanged. In Paper~I, Proposition~23.1 and Paper~I, Theorem~23.2 the
remaining cusp defect has already been adjusted to belong to
$\mathbb Z w\oplus\mathbb Z d$; write it as $aw+bd$.
Assign $aw$ to the old elliptic node and $bd$ to the new
$\delta$ node. Their product is the old defect, since
$C_wd=d$. Each is an allowed integral winding in its primitive
rank-one toric direction. Realize it locally by the appropriate
integer power of the base coordinate, times a holomorphic unit;
the transverse elliptic translation is holomorphic. These
are precisely the extending rank-one translations and they
agree with the unsplit two-direction prescription when the
two zeros coincide.

\item The differences of these local prescriptions have the same
coarse vertical-vector sheaf as in Paper~I, Proposition~23.1 and Paper~I, Theorem~23.2:
\[
 0\longrightarrow\mathcal O_B\longrightarrow\mathcal V_B
 \longrightarrow\mathcal O_B(-p_0)\longrightarrow0.
\]
To check this assertion locally, quotient by the invariant
$\delta$ vector. The quotient is the Lie line of the unchanged
elliptic surface; its finite-root exponents and its simple
pole at the old elliptic node are unchanged. The new
$\delta$-direction node has a nonvanishing logarithmic
$\delta$ generator in the abstract vertical bundle and a
regular transverse elliptic generator, so it adds no zero
or pole to that quotient. The extra subgroups in the enlarged
finite roots act trivially on vector coefficients, just as
in Paper~I, Proposition~23.1 and Paper~I, Theorem~23.2. Both end terms have zero $H^1$, hence the Cech
correction glues the new local prescriptions. On a small
parameter disc $\Delta_s$, including the coalescence, this is
the relative exact sequence
\[
 0\longrightarrow\operatorname{pr}_B^*\mathcal O_B
 \longrightarrow\mathcal V_{B\times\Delta_s}
 \longrightarrow\operatorname{pr}_B^*\mathcal O_B(-p_0)
 \longrightarrow0.
\]
The local logarithmic generators described above also define
this sequence at $s=0$. The two-chart Cech complex for either
end term has zero first cohomology, so
$R^1\operatorname{pr}_{\Delta_s*}\mathcal V_{B\times\Delta_s}=0$.
The same holds after adjoining a small Stein $c$-disc. Solving
that relative complex supplies a correction holomorphic in
$c,s$, through the coalescence. The integral windings are constant
and the local units vary holomorphically; at the coalescence
their two toric directions are the original ones. The
comparison thus covers the identity on $s$ and on $B$, with
the chosen small resolution. Different corrections differ
by the connected vertical translation group, so they have
the same smooth monodromy class.

\item The four homological matrices are still \eqref{maps:eq:B.6.2}. Finite transfer
still gives \eqref{maps:eq:B.6.3}, since the root actions are free even when
their multiplicity exceeds the linear order. Replacing
$S=I-2E$ by $I-hE$ multiplies the two coboundaries in \eqref{maps:eq:B.6.5}
by $h/2$. Thus the same rank-two parabolic group proves the
integral Torelli assertion.

\item For the characteristic-class calculation the regular nilpotent
fibre now has
\[
 d\vartheta=h\kappa\gamma.
\]
Changing the finite twist data changes the affine parameter
defects but not their first coordinate, which remains zero.
The base actions therefore still fix $\gamma$; on wedging
their action on $\upsilon,\omega$ with $\gamma$, the extra
terms are multiples of $\gamma\kappa$, which are exact.
The two coinvariant vanishings in \eqref{maps:eq:B.6.8} follow unchanged.
Choose the local parameter lift at the new rank-one point
by subtracting its integral toric winding, so its long
$\gamma$ translation has zero representative. Its boundary
equation is then
$d\vartheta=(\alpha+h\kappa)\gamma$.
The two local skyscraper classes $\kappa\gamma\upsilon$ and
$\kappa\gamma\omega$ consequently have Wang components
$-[\gamma\upsilon]/h$ and $-[\gamma\omega]/h$.
They survive and exhaust the mapping-torus $H^3$ exactly as
in \ref{maps:src:B.6.4}. Their restrictions to both swept critical curves
are zero. The same saturated-vertical tangent sequence proves
$p_1=0$ integrally, and the Kreck sequence places the smoothing
class in the central $\Theta_7$. Finally the fixed-$s$ trace
argument of \ref{maps:src:B.6.5} applies unchanged and identifies this class
with that of the original sphere family.

\end{enumerate}
\end{proof}




\begin{proposition}[Counterexample: the scalar-shear claim already fails on the ordinary locus]\label{maps:src:B.6.7}
\StatementRecord{maps:src:B.6.7}{false as stated; correction proved}{GPT-6 Astra (OpenAI, ChatGPT Pro mode), effort Pro, 2026-11-17; Claude Fable 5.2 (Anthropic), effort high, 2026-11-16; Claude Fable 5.2 (Anthropic), effort high, 16 November 2026}{the hypotheses, notation and definitions stated here; \cref{maps:eq:B.6.6}; \cref{maps:eq:B.6.16}}


The proposed isotopy over the ordinary locus from the marked return to the fibrewise scalar twist $x\mapsto\rho(-2\theta(x))x$ is false for the original family. On $X^\circ=f^{-1}(B^\circ)$,
that return is not homotopic over $B^\circ$ to
\[
 x\longmapsto\rho(-2\theta(x))x
\]
for a globally defined monodromy-invariant $\gamma$ coordinate
$\theta:X^\circ\to\mathbb R/\mathbb Z$.
In fact, after allowing arbitrary additional circle-valued
scalar gauges on the base, a quotient of the affine defect
has exact order six. Its vanishing under the sixth iterate
does not compute the sixth iterate in $\Theta_7$.

\end{proposition}


\begin{proof}
\begin{enumerate}[label=\arabic*.]
\item Take a compact pair-of-pants core $P\subset B^\circ$ and the
marked zero section on its ordinary torus bundle. A connection
for the compact family may be chosen to equal the flat real
period connection over this core and then extended over the
three collars. Consequently its induced return on the core
fundamental group is the one calculated from the closing map
in Paper~I, Proposition~18.5. With base generators $a_1,a_2$ and fibre translations
$t_\lambda$, that group and its return are
\[
 \pi_1(f^{-1}(P))=\Lambda\rtimes F(a_1,a_2),\qquad
 \Phi_*t_\lambda=t_{S\lambda},\qquad
 \Phi_*a_j=t_{d_j}a_j,
\]
where $S=I-2E$ and $d_1,d_2$ are exactly \eqref{maps:eq:B.6.6}.
Different choices of the path from a basepoint to its image
conjugate this homomorphism by an integral fibre translation.

\item The pure scalar shear has the same linear part $S$.
For the zero-section coordinate its two affine defects are
zero. Changing the coordinate by a map $P\to S^1$ adds
defects in $\mathbb Z d$; allowing arbitrary such defects
only weakens the proposed comparison. A homotopy over $P$
tracks the basepoint inside its fibre. It would therefore
require, after quotienting by the invariant $\mathbb Z d$,
an integral vector $\bar\lambda\in\bar\Lambda:=\Lambda/\mathbb Z d$
with
\[
 (\bar A_j-I)\bar\lambda=\bar d_j,\qquad j=1,2.
 \tag{B.6.15}\label{maps:eq:B.6.15}
\]
The possible sign from the direction of basepoint conjugation
has been absorbed into $\bar\lambda$.
Writing $\bar\lambda=(x,y,z)$ in the $(g,u,w)$ basis,
the nonzero equations are
\[
 6x-y+z=1,\qquad -6x-y-2z=-1,\qquad
 -y-z=0,\qquad -6x+y-z=-1.
\]
Their unique rational solution is
\[
 (x,y,z)=(1/6,0,0).
 \tag{B.6.16}\label{maps:eq:B.6.16}
\]
It is not integral, proving the stated obstruction.

\item This is independent of the auxiliary ordinary marking.
Changing the basepoint path changes $(\bar d_j)$ by one
simultaneous coboundary. Replacing the lifted base generators
by $t_{r_j}a_j$ changes $d_j$ by $(S-I)r_j$, which lies in
$\mathbb Z d$ and disappears in the quotient. An integral
linear change of marking transports both the local system
and the displayed class. Thus none of these changes removes
the nonzero class in $H^1(P,\bar\Lambda)$.

\item Moreover $6\bar d_j=(\bar A_j-I)g$, and \eqref{maps:eq:B.6.16} shows
that no smaller positive multiple is a coboundary. Its
order is exactly six; the three nonzero Smith factors of
the stacked quotient coboundary matrix are $1,1,6$.
Since $Sd_j=d_j$, the defect of the $r$th iterate is $rd_j$,
so the sixth iterate is the first one for which this
particular obstruction vanishes. In the full lattice,
subtracting that integral coboundary gives
\[
 6d_1-(A_1-I)g=-4d,\qquad
 6d_2-(A_2-I)g=3d.
 \tag{B.6.17}\label{maps:eq:B.6.17}
\]
These are the two scalar windings seen in Paper~I, Proposition~18.5, with the
inverse sign if one reverses its comparison direction.
They still have to be matched to the finite and cusp collars
to compute a diffeomorphism class of the compact sphere.

\item The arithmetic also explains a distinction in terminology.
For a shift $h=2r$, a literal shear $-h\theta$ with a global
$\gamma$ coordinate would require the base windings of
$\theta$ to be $1/3,-1/4$, after solving the quotient
equations. Those are not integers. A scalar gauge whose
$h$th root is not globally defined is a different map; the
sixfold scalar representative in Paper~I, Proposition~18.5 is of that latter kind.
On a simply connected ordinary base chart the translation
has an additive lift and can be interpolated, which proves
the local shear statement there. It does not give the
asserted isotopy over all of $B^\circ$.
Every rank, Smith factor and rational solution used here
is checked in \texttt{maps\_ordinary\_shear\_obstruction.py}.

\end{enumerate}
\end{proof}




\begin{corollary}[every sphere pair has the same order-six affine obstruction]\label{maps:src:B.6.8}
\StatementRecord{maps:src:B.6.8}{proved}{GPT-6 Astra (OpenAI, ChatGPT Pro mode), effort Pro, 2026-11-17; Claude Fable 5.2 (Anthropic), effort high, 16 November 2026}{the hypotheses, notation and definitions stated here; \cref{maps:src:B.6.7}; \cref{maps:eq:B.6.6}}


For the positive-generator return of every $X_{3k,4l}$, the
quotient affine class in $H^1(P,\Lambda/\mathbb Z d)$ has exact
order six. It obstructs an ordinary-locus isotopy to a pure
scalar shear, even after arbitrary base circle gauges.
This assertion does not identify the different pairs'
classes in $\Theta_7$ with one another.

\end{corollary}


\begin{proof}
\begin{enumerate}[label=\arabic*.]
\item Use the Paper~I, Proposition~23.1 and Paper~I, Theorem~23.2 coordinates with $b=4lN_0+3kN_1=\pm1$ and
$h=2kl$. The fixed-data repeated-norm solutions in its
comparison proof are
\[
 \ell_1=(0,0,-2lN_0,0)^t,\qquad
 \ell_2=(0,-kN_1,0,0)^t.
\]
Their norms are respectively $-hN_0d$ and $-hN_1d$, as
required by the shift of the affine vectors. The remaining
cusp defect is
\[
 -M_0(\ell_1+A_1\ell_2)
       =(0,0,2lN_0-kN_1,kN_1)^t,
\]
which is already in $\mathbb Zw\oplus\mathbb Zd$; no
additional adjustment of its long coordinates is needed.
Thus these are the integral defects of the allowed fixed-data
closing comparison. Its quotient cocycle has the unique
rational coboundary
\[
 \bar\lambda=
 \left(\frac{4lN_0-3kN_1}{6},\,
       2lN_0-kN_1,\,2kN_1-2lN_0\right)^t.
\]
The last two coordinates are integral and the first is
$b/6-kN_1$. Since $b=\pm1$, the least positive integer
making all three coordinates integral is exactly six.
The same basepoint and marking argument as in \ref{maps:src:B.6.7} proves
both the order assertion and the obstruction to a scalar
shear. For $(k,l,N_0,N_1)=(1,1,1,-1)$ the quotient defects
here and those of \eqref{maps:eq:B.6.6} differ by the integral coboundary
of $(-1,-3,4)^t$, checking their comparison convention.
The symbolic identities are verified in the same exact
program as \ref{maps:src:B.6.7}.

\end{enumerate}
\end{proof}


The complete credited statement and re-derivation are in \cref{maps:src:B.6.9}.

\begin{proposition}[the scalar-return cover and its elliptic surface]\label{maps:src:B.6.10}
\StatementRecord{maps:src:B.6.10}{proved}{GPT-6 Astra (OpenAI, ChatGPT Pro mode), effort Pro, 2026-11-17; Claude Fable 5.2 (Anthropic), effort high, 2026-11-16; GPT-6 Astra (OpenAI, ChatGPT Pro mode), S6\_06\_BLIND, effort Pro, 2026-11-17; Claude Fable 5.2 (Anthropic), effort high, 16 November 2026}{the hypotheses, notation and definitions stated here; \cref{maps:src:B.6.3}; \cite[Proposition 8.12]{mapsShioda90}; \cite[Theorem 8.6]{mapsShioda90}}


Fix $p\in\mathbb P^1\setminus\{0,1,\infty\}$ and let $C_p$ be
the smooth projective normalization of
\[
 h^{12}=\frac{t^4}{(t-1)^3(t-p)}.
 \tag{B.C1}\label{maps:eq:B.C1}
\]
Multiplying the right side by a nonzero constant does not affect
the statements below.  Write $\eta(h)=\zeta_{12}h$ for its deck
generator.  This is a connected genus-three cover, with
ramification indices $3,4,12$ at $0,1,p$, respectively.
Let $S_p\to C_p$ be the relatively minimal elliptic surface obtained
by pulling back Engel's marked surface $S$ of type $IV^*,III,I_1$.
Then $S_p$ has precisely twelve singular fibres, all of type $I_1$,
and
\[
 \chi(\mathcal O_{S_p})=1,\quad q(S_p)=p_g(S_p)=3,\quad
 (b_2^+,b_2^-)=(7,15),\quad \sigma(S_p)=-8.
 \tag{B.C2}\label{maps:eq:B.C2}
\]
The section induced by $P$ has height two, meets the zero section
with intersection number zero, and satisfies $(P-O)^2=-2$.
For $k=0,\ldots,11$, the dimensions of the $\zeta_{12}^k$
eigenspaces of $\eta^*$ on $H^2(S_p;\mathbb C)$ and on the
complexification of its positive subspace are, respectively,
\[
 \begin{array}{c|rrrrrrrrrrrr}
 k&0&1&2&3&4&5&6&7&8&9&10&11\\\hline
 \dim H^2_k&3&1&2&2&2&1&3&1&2&2&2&1\\
 \dim H^+_k&1&0&1&1&0&0&2&0&0&1&1&0 .
 \end{array}
 \tag{B.C3}\label{maps:eq:B.C3}
\]
In particular the invariant intersection form has rank three
and signature minus one.  This proposition computes the cover
of the scalar return gauge in \ref{maps:src:B.6.3}.  It does not identify that
cover with a complete semistable component normalization.

\end{proposition}


\begin{proof}
\begin{enumerate}[label=\arabic*.]
\item The three valuations of the function in \eqref{maps:eq:B.C1} are $4,-3,-1$.
Their greatest common divisor with twelve is one, which proves
connectedness.  At these values there are respectively $4,3,1$
points above the base, with the displayed ramification indices.
Riemann--Hurwitz gives
\[
 2g(C_p)-2=-24+(12-4)+(12-3)+(12-1)=4.
\]
There are twelve unramified points over infinity.

\item The order-three and order-four base changes at $0$ and $1$
are exactly the good-reduction charts used in D.3.3:
$t=z^3,x=z^4X,y=z^6Y$ at zero and
$t-1=z^4,x=z^2X,y=z^3Y$ at one.  Thus the relatively
minimal pullback is smooth over those points.  The fibre at
$p$ was already smooth, so ramified base change there remains
a smooth elliptic fibration.  Over infinity are twelve
unramified copies of $I_1$.  Euler localization gives
$c_2(S_p)=12$.  The canonical bundle of an elliptic surface
with a section and no multiple fibres is pulled back from
its base, so $c_1^2=0$.  Noether's formula yields
$\chi(\mathcal O)=1$.

\item For completeness the usual direct-image calculation is enough
to determine the irregularity: $\pi_*\mathcal O=\mathcal O_{C_p}$
and $R^1\pi_*\mathcal O=\mathcal L^{-1}$, where
$\deg\mathcal L=\chi(\mathcal O)=1$.  The negative-degree
line bundle has no section, so $q=g(C_p)=3$; hence $p_g=3$.
The Euler number gives $b_2=22$, and the K\"ahler decomposition
gives $b_2^+=1+2p_g=7$.  This proves \eqref{maps:eq:B.C2}.
Shioda's base-change identity
\cite[Proposition 8.12]{mapsShioda90} multiplies the original
height $1/6$ by twelve.  Since all fibres are irreducible,
his height formula \cite[Theorem 8.6]{mapsShioda90} becomes
$\langle P,P\rangle=2+2(P\cdot O)$; it gives $P\cdot O=0$.
Both section self-intersections are $-\chi=-1$, proving the
last intersection assertion.

\item We compute the Hodge characters, including their signs.
Choose the rational two-form $\Omega_S$ whose divisor on the
original rational elliptic surface is minus the fibre at infinity.
Its pullback to the minimal surface has orders $0$ at each of the
four points above zero, $2$ at each of the three points above one,
$11$ at the point above $p$, and $-1$ at the twelve points above
infinity.  The first two values follow by substituting the two
good-reduction charts in $dt\wedge dx/y$; the third is the
ramification factor of $dt$, and the last is unramified pullback.
At every other fibre the form is regular and nonvanishing as
a canonical generator.  Its divisor on the base therefore has
degree $6+11-12=5$, as also follows from $K_{C_p}\otimes\mathcal L$.

\item Every eigenform is a rational function of $t$ times
$h^j\Omega_S$, for some $0\leq j\leq11$.  Its least allowed
factors $t^a(t-1)^b(t-p)^c$ satisfy
\[
 a=-\lfloor j/3\rfloor,\qquad
 b=\lceil(j-2)/4\rceil,\qquad
 c=\lceil(j-11)/12\rceil=0.
\]
Regularity at infinity allows an additional polynomial of
degree at most $-a-b-c-1$.  This degree is nonnegative exactly
for $j=6,9,10$, and is zero in each case.  A basis is consequently
\[
 \frac{h^6(t-1)}{t^2}\Omega_S,\qquad
 \frac{h^9(t-1)^2}{t^3}\Omega_S,\qquad
 \frac{h^{10}(t-1)^2}{t^3}\Omega_S .
 \tag{B.C4}\label{maps:eq:B.C4}
\]
The invariant K\"ahler class and these forms together with their
conjugates give the positive row of \eqref{maps:eq:B.C3}.

\item We determine the full second-cohomology characters by fixed points.
For $1\leq j\leq11$, the number of fixed points of $\eta^j$ on
$C_p$ is
$1+4\mathbf 1_{4\mid j}+3\mathbf 1_{3\mid j}$.
Thus its first-cohomology trace is two minus this number.
Pullback identifies $H^1(C_p)$ with $H^1(S_p)$: it is injective
by the zero section and both groups have dimension six.
The traces on $H^1$ and $H^3$ agree by duality and conjugation.
Above $p$ the whole smooth elliptic fibre is fixed and has Euler
number zero.  At a point above zero, a nontrivial order-three
elliptic automorphism fixes three points.  At a point above one,
an order-four automorphism fixes two points, and its square fixes
four.  There are no other fixed points.  The local Lefschetz signs
are positive, since the normal actions are complex rotations.
Subtracting the contributions of $H^0,H^1,H^3,H^4$ gives
\[
 \operatorname{Tr}(\eta^{j*}\mid H^2)=
 \begin{cases}
 22&j=0,\\
 4&j=4,8,\\
 6&j=6,\\
 0&\text{otherwise}.
 \end{cases}
\]
The finite Fourier transform gives the first row of \eqref{maps:eq:B.C3}.

\item The program \texttt{fourfolds\_B6\_cover.py} performs these twelve
character calculations in $\mathbb Q[x]/(x^4-x^2+1)$ and checks
every displayed dimension exactly; no numerical eigenvalues are used.

\end{enumerate}
\end{proof}




\begin{proposition}[the untwisted plumbing has signature two]\label{maps:src:B.6.11}
\StatementRecord{maps:src:B.6.11}{proved}{GPT-6 Astra (OpenAI, ChatGPT Pro mode), effort Pro, 2026-11-17; Claude Fable 5.2 (Anthropic), effort high, 2026-11-16; Claude Fable 5.2 (Anthropic), effort high, 16 November 2026}{the hypotheses, notation and definitions stated here; \cref{maps:eq:B.P2}; \cite[Lemma 2.1]{mapsEngel}; \cref{maps:eq:B.P3}; \cref{maps:src:B.6.9}}


Let $W^\pm$ be a sufficiently small closed parameter-disc
neighbourhood in either smooth semistable total fourfold of
Paper~I, Theorem~20.4, with its complex orientation.  Its boundary is the
mapping torus of the \emph{untwisted} family $X(u)$ under one
positive $u$ loop, and this boundary is not spin.  One has
\[
 \sigma(W^\pm)=2,\qquad c_1(W^\pm)^2=0
       \quad\text{in }H^4(W^\pm;\mathbb Z).
 \tag{B.P1}\label{maps:eq:B.P1}
\]
The signature form modulo its radical is rationally represented by
$\left(\begin{smallmatrix}2&1\\1&6\end{smallmatrix}\right)$.
In particular the $E_8$ summand of the rational elliptic surface
does not give an $E_8$ signature for this parameter filling.
These are not the finite logarithmic replacements used for the
sphere; no smooth monodromy value for those replacements follows.

\end{proposition}


\begin{proof}
\begin{enumerate}[label=\arabic*.]
\item We first give the local-to-global intersection calculation.
Write $Q$ for the irreducible central normal-crossings fibre and
$\nu:R\to Q\subset W$ for its smooth normalization; here $R$ means
$R^+$ or $R^-$.  Its conductor sections $i_0,i_\infty:S\to R$
are glued by $T=t_{6P}$.  A sufficiently small semistable tube
retracts to its central fibre: in the product charts $xy=u$ shrink
the normal coordinates toward their zero loci, fixing that locus,
and use positive radial normal fields on the single-branch charts.
Collars make these contractions agree, and compactness gives a
finite choice of charts.  The normalization mapping cylinder
therefore computes the homology of $W$.  In particular
\[
 H_4(S;\mathbb Q)\longrightarrow H_4(R;\mathbb Q)
 \longrightarrow H_4(W;\mathbb Q)\longrightarrow H_3(S;\mathbb Q)=0.
 \tag{B.P2}\label{maps:eq:B.P2}
\]
Thus all middle classes are pushed forward from the normalization.

\item The normalization is a complex codimension-one immersion with
transverse double locus.  Let $D_0,D_\infty$ be its two conductor
divisors.  Its normal line has
\[
 c_1(N_\nu)=-[D_0]-[D_\infty].
\]
Indeed the parameter differential in $u=xy$, restricted to either
branch, is the normal differential multiplied by the coordinate
vanishing on the conductor.  Since the parameter divisor is
principal, these local identities identify the normal line with
$\mathcal O_R(-D_0-D_\infty)$.

\item For $a,b\in H^2(R;\mathbb Q)$, put $a_i=i_i^*a$ and $b_i=i_i^*b$.
Push the four-cycles Poincar\'e dual to $a,b$ into $W$, and perturb
one transversely.  Their same-branch intersections are counted by
the Euler class of $N_\nu$; their two other intersections lie on
the two branches of the conductor.  The complex orientations give
positive signs to both cross terms.  Thus their intersection is
\[
 -a_0b_0-a_\infty b_\infty
       +a_0T^*b_\infty+b_0T^*a_\infty
   =-\bigl(a_0-T^*a_\infty\bigr)
         \bigl(b_0-T^*b_\infty\bigr),
 \tag{B.P3}\label{maps:eq:B.P3}
\]
where products on the right are evaluated on $[S]$.
This also follows by doing the perturbation in each $xy=0$ chart;
the two unequal branch intersections are the two last terms.
Define $U(a)=a_0-T^*a_\infty$.
Equation \eqref{maps:eq:B.P2} shows that the signature is precisely the negative
of the surface intersection form restricted to $\operatorname{im}U$.
The relation coming from $H_4(S)$ causes no discrepancy:
its dual divisor is $D_0-D_\infty$, and its image under $U$ is zero
because the two conductor normal bundles have trivial product.

\item We compute this image using the actual surface classes
\[
 F=[S_t],\qquad m=[P]-[O],\qquad
 v=[6P]-[O]-3F.
\]
They satisfy
\[
 F^2=Fm=Fv=0,
 \qquad m^2=-2,\quad mv=-1,\quad v^2=-6.
 \tag{B.P4}\label{maps:eq:B.P4}
\]
These equalities follow from the section intersections
$P\cdot O=0$, $[6P]\cdot O=P\cdot[6P]=2$ and the
self-intersection $-1$ of every section.  The two latter
intersection numbers follow directly from height $1/6$ and
the narrow section $6P$, as in \cite[Lemma 2.1]{mapsEngel}.
Translation acts by the integral isometry
\[
 T_*x=x+(xF)v-(xv-3xF)F.
\]
It is checked on $F,O,v$ and the seven component-root classes,
which $t_{6P}$ fixes.  In particular
$\operatorname{im}(I-T^*)_{\mathbb Q}=\mathbb QF+\mathbb Qv$.
On $\mathbb P(\mathcal O\oplus L)$, the relative hyperplane
class has two section restrictions differing by $m$; hence its
image under $U$ is $m$ modulo $\mathbb QF+\mathbb Qv$.
The blow-up in Paper~I, Proposition~20.1 and Paper~I, Theorem~20.2 and Paper~I, Proposition~20.3 and Paper~I, Theorem~20.4 and Paper~I, Proposition~20.11 and Paper~I, Proposition~20.12 is along the fibre $S_\infty$ in one section.
Its exceptional divisor restricts to $F$ on that section and zero
on the other.  The small resolution introduces no further divisor
class: the normalization homology, or the two displayed conifold
charts, gives one additional second-cohomology generator in total.
Consequently
\[
 \operatorname{im}U=\mathbb QF+\mathbb Qv+\mathbb Qm.
\]
Its radical is $\mathbb QF$.  The matrix on $(m,v)$ is negative
definite, with leading entry $-2$ and determinant eleven.
Negating it in \eqref{maps:eq:B.P3} gives the positive matrix in the statement,
and proves $\sigma(W)=2$.  The signs and determinant are checked
exactly in \texttt{fourfolds\_B6\_plumbing.py}.

\item Finally on the initial projective bundle,
$c_1(TR)=\pi_R^*c_1(S)+D_0+D_\infty$.
Blowing up the conductor curve subtracts the exceptional divisor
both from $c_1(TR)$ and from $D_0+D_\infty$; the subsequent small
conifold resolution is crepant.  The tangent-normal sequence
therefore gives
\[
 \nu^*c_1(W)=\pi_R^*c_1(S)=\pi_R^*F.
\]
Its square is zero.  The cohomological normalization sequence,
with $H^3(S;\mathbb Z)=0$, makes
$H^4(W;\mathbb Z)\to H^4(R;\mathbb Z)$ injective.
This proves the integral assertion in \eqref{maps:eq:B.P1}.
We check the boundary assertion integrally as well.
The normalization sequence gives $H_6(W;\mathbb Z)=\mathbb Z$,
generated by the reduced irreducible central fibre, and
$H_5(W;\mathbb Z)=0$.  For the latter, $H_5(R)=0$ and the next
map $H_4(S)\to H_4(R)$ is nonzero: after pushing through the
normalization modification its dual is $m\ne0$.
Poincar\'e--Lefschetz duality identifies
$H^2(W,\partial W)\to H^2(W)$ with the class of that central
divisor.  It is zero because the divisor is the zero set of the
global parameter function.  The next relative group is
$H^3(W,\partial W)=H_5(W)=0$.  The exact sequence thus gives
\[
 H^2(W;\mathbb Z)\xrightarrow{\ \simeq\ }
 H^2(\partial W;\mathbb Z).
\]
The class $c_1(W)$ is primitive: its pullback $\pi_R^*F$ pairs
to one with the zero-section curve in either conductor copy of
$S$.  Its restriction to the boundary is therefore not divisible
by two.  For a complex tangent bundle $w_2=c_1\bmod2$, and the
stable tangent restriction to the boundary has the same $w_2$.
Consequently $w_2(\partial W)\ne0$, proving that this is not the
spin mapping torus required in \ref{maps:src:B.6.9}.

\item When relative characteristic numbers for a compatible boundary
are defined, the absolute vanishing also kills the terms
$c_1^4$ and $p_1c_1^2$ by pairing a relative lift with the zero
absolute square.  It does not supply the missing relative
$p_1^2$ or replace the finite logarithmic caps.

\end{enumerate}
\end{proof}




\begin{proposition}[a generic cyclic-norm obstruction]\label{maps:src:B.6.12}
\StatementRecord{maps:src:B.6.12}{proved}{GPT-6 Astra (OpenAI, ChatGPT Pro mode), effort Pro, 2026-11-17; Claude Fable 5.2 (Anthropic), effort high, 16 November 2026}{the hypotheses, notation and definitions stated here; \cref{maps:src:B.6.10}; \cite[Theorem (1.1)]{mapsGriffiths68II}; \cite[Lemmas 8.1--8.3 and Theorem 8.9]{mapsShioda90}; \cite[Theorems 8.6--8.7]{mapsShioda90}; \cref{maps:src:B.6.13}; \cref{maps:src:B.6.14}}


Use the sectioned elliptic surface $S_p\to C_p$ and deck action
$\eta$ of \ref{maps:src:B.6.10}. Outside a countable subset of
$\mathbb P^1\setminus\{0,1,\infty\}$, its Mordell--Weil group
$L_p$ satisfies
\[
 L_p^\eta=\mathbb ZP,\qquad
 (1+\eta+\cdots+\eta^{11})L_p=12\mathbb ZP.
 \tag{B.6.19}\label{maps:eq:B.6.19}
\]
In particular $6P$ represents a nonzero class of order two
in the invariant group modulo the cyclic norm image.
There is no elliptic-surface automorphism covering $\eta$
whose twelfth power is translation by $6P$ or by $-6P$.
This obstructs a conductor-step construction on this
\emph{sectioned Jacobian pullback}. It does not exclude a
normalization which is an elliptic torsor without a section,
an additional division cover, or a different semistable model.
It does not compute an element of $\Theta_7$.

\end{proposition}


\begin{proof}
\begin{enumerate}[label=\arabic*.]
\item We first determine the two Hodge summands relevant to the norm.
Let $\Phi_d$ denote the $d$th cyclotomic polynomial. In the
notation of \ref{maps:src:B.6.10}, the rational $\Phi_2$ summand of $H^2$ has
rank three and signature $(2,1)$; the rational $\Phi_4$ summand
has rank four and signature $(2,2)$. On either of the complex
$\eta$-eigenspaces of order four the Hermitian form has
signature $(1,1)$. We will prove that each associated
one-dimensional period map is nonconstant.

\item Here is an exact calculation of their monodromy as $p$ moves.
The homological matrices of the original elliptic surface are
\[
 R_1=\begin{pmatrix}0&1\\-1&-1\end{pmatrix},\qquad
 R_2=\begin{pmatrix}0&-1\\1&0\end{pmatrix},\qquad
 R_0=\begin{pmatrix}1&0\\1&1\end{pmatrix},\qquad
 R_1R_2R_0=I.
\]
For the character $k=3$ of the twelve-fold cover, use the
cohomological local system with four matrices
\[
 (G_1,G_2,G_3,G_4)
       =(R_1^{-t},\,iR_2^{-t},\,R_0^{-t},\,-iI).
 \tag{B.6.20}\label{maps:eq:B.6.20}
\]
For $k=6$ replace the character factors $(i,-i)$ by $(-1,-1)$,
giving
$(R_1^{-t},-R_2^{-t},R_0^{-t},-I)$.
Changing the deck-character convention conjugates the
order-four computation and does not change its conclusion.

\item These systems compute the nontrivial-character pieces of
$H^2(S_p)$ by parabolic cohomology. To check this identification,
push the elliptic fibration through the finite map $C_p\to\mathbb
P^1$. The $R^0$ and $R^2$ terms are the permutation system of
the cover; their global contributions have trivial character.
For $R^1$, good reduction at the ramification points identifies
the stalk with the invariant subspace of the displayed local
matrix. At an $I_1$ point the nodal specialization gives the
same invariant subspace. Thus its character summand is exactly
the middle extension $j_*G$, and its $H^1$ is the stated
parabolic group. The possible Leray differentials have trivial
character targets or sources, so none changes these summands.

\item Explicitly, a parabolic cocycle is a tuple
$z_j\in\operatorname{im}(G_j-I)$ satisfying
\[
 z_1+G_1z_2+G_1G_2z_3+G_1G_2G_3z_4=0,
\]
modulo the tuples $((G_j-I)v)_j$.
The four local image ranks are $(2,1,1,2)$ for $k=3$ and
$(2,2,1,2)$ for $k=6$. In each case the relation has rank two
and the coboundary space has rank two. The quotient dimensions
are therefore two and three. Representatives for bases are
the columns of the following matrices, grouped in two-row
blocks for $j=1,2,3,4$:
\[
 P_3=\begin{pmatrix}
 -1-i&i\\-1&0\\-1&0\\i&0\\0&-1\\0&0\\0&0\\0&0
 \end{pmatrix},\qquad
 P_6=\begin{pmatrix}
 0&2&-1\\-1&1&0\\-1&1&0\\-1&-1&0\\
 0&0&-1\\0&0&0\\0&0&0\\0&0&0
 \end{pmatrix}.
\]
They satisfy the image and relation conditions, and are
independent modulo the displayed coboundaries.

\item Let the Artin generator $\sigma_j$ send
$a_j\mapsto a_ja_{j+1}a_j^{-1}$ and
$a_{j+1}\mapsto a_j$. Moving the fourth puncture around the
first two gives the pure braids
\[
 \sigma_3\sigma_2\sigma_1^2\sigma_2^{-1}\sigma_3^{-1},
 \qquad \sigma_3\sigma_2^2\sigma_3^{-1}.
\]
The fourth matrix is scalar, so these braids preserve the
local representation. Evaluate them on the above cocycles
using $z(ab)=z(a)+G(a)z(b)$ and
$z(a^{-1})=-G(a)^{-1}z(a)$, then reduce by coboundaries.
For $k=3$ the resulting two matrices are
\[
 T_1=\begin{pmatrix}-1+i&1\\i&1\end{pmatrix},\qquad
 T_2=\begin{pmatrix}2i&1-2i\\i&1-i\end{pmatrix};
\]
their commutator has characteristic polynomial
$x^2-5x+1$. For $k=6$ they are
\[
 T_1=\begin{pmatrix}1&-2&1\\1&0&0\\1&-1&1\end{pmatrix},
 \qquad
 T_2=\begin{pmatrix}-1&-2&2\\1&0&-1\\-1&-1&2\end{pmatrix},
\]
whose commutator has polynomial $(x-1)(x^2-7x+1)$.
Both have an eigenvalue of absolute value greater than one.
All cocycle, quotient and word identities are checked exactly
in \texttt{maps\_component\_periods.py}; the displayed matrices
also give direct small certificates of noncompact monodromy.
A possible scalar deck factor in identifying a moving-root
basepoint cancels in the commutator.

\item The minimal elliptic pullbacks form a smooth projective family
as $p$ varies away from $0,1,\infty$: the two good-reduction
coordinate changes in \ref{maps:src:B.6.10} work relatively, and the twelve
nodal fibres remain disjoint. Their period maps are holomorphic
by \cite[Theorem (1.1)]{mapsGriffiths68II}.
If the period in the order-four summand were constant,
monodromy would preserve its positive line and its negative
orthogonal line, hence lie in a compact unitary subgroup.
If the order-two period were constant, monodromy would
preserve the positive two-plane and negative line, again a
compact subgroup. The computed hyperbolic commutators exclude
both possibilities. Each period domain has complex dimension
one. For any fixed nonzero rational vector, being of type
$(1,1)$ is a proper period hyperplane condition, so its
preimage in this nonconstant one-dimensional family is
discrete. On the universal cover take the countable union
over rational vectors and a countable set of marking charts.
Its projection to the parameter curve is countable.
Outside that set, neither rational $\Phi_2$ nor rational
$\Phi_4$ contains a nonzero $(1,1)$ class.

\item Now let $p$ be outside this countable set.
The Mordell--Weil lattice embeds equivariantly into rational
$(1,1)$ cohomology by the orthogonal projection away from the
fibre and zero-section classes; see
\cite[Lemmas 8.1--8.3 and Theorem 8.9]{mapsShioda90}.
All fibres are irreducible, so its height lattice is integral,
even and torsion-free, with every nonzero vector having norm
at least two, by \cite[Theorems 8.6--8.7]{mapsShioda90}.
Since $P$ has norm two, it is primitive. The invariant
second cohomology has rank three by \ref{maps:src:B.6.10}; the fibre and
zero section account for two dimensions and $P$ for the
third. Hence $L_p^\eta=\mathbb ZP$.

\item The only eigenvalues of an operator of order twelve whose
fourth power is one have orders $1,2,4$. The last two are
absent from $L_p\otimes\mathbb Q$, by the period argument.
Consequently $L_p^{\eta^4}=\mathbb ZP$ as well. For any
$Q\in L_p$, write
\[
 (1+\eta^4+\eta^8)Q=mP.
\]
Pairing with $P$ gives $3\langle P,Q\rangle=2m$.
The integral height pairing forces $m=3r$ for an integer
$r$. Applying $1+\eta+\eta^2+\eta^3$ now gives
$NQ=12rP$. Conversely $N(P)=12P$, proving \eqref{maps:eq:B.6.19}.
In particular $6P$ is not a norm but twice it is.

\item Finally any elliptic-surface automorphism covering the base
deck generator is $t_Q\eta$ or $t_Q[-1]\eta$: its value on
the zero section gives $Q$, and the remaining automorphism
fixes zero and is $\pm1$ because the elliptic $j$-function
is nonconstant. The twelfth power in the first case is
$t_{NQ}$, which cannot be $t_{\pm6P}$. In the second case
the alternating norm has image in the $\eta=-1$ eigenspace
of $L_p\otimes\mathbb Q$, which is zero. Torsion-freeness
makes that norm zero integrally. This proves the assertion
about the sectioned conductor step. Replacing the elliptic
surface by a torsor removes the zero-section argument.
\ref{maps:src:B.6.13} constructs such torsor data, while \ref{maps:src:B.6.14} isolates
a further obstruction to using them as the requested filling.

\end{enumerate}
\end{proof}




\begin{proposition}[torsor and conductor-line data exist]\label{maps:src:B.6.13}
\StatementRecord{maps:src:B.6.13}{proved}{GPT-6 Astra (OpenAI, ChatGPT Pro mode), effort Pro, 2026-11-17; Claude Fable 5.2 (Anthropic), effort high, 16 November 2026}{the hypotheses, notation and definitions stated here; \cref{maps:src:B.6.12}; \cref{maps:src:B.6.14}; \cref{maps:eq:B.6.23}; \cref{maps:src:B.6.10}; \cite[Theorem 8.6 and Proposition 8.12]{mapsShioda90}}


For every $p\notin\{0,1,\infty\}$ there is a compact smooth
holomorphic genus-one fibration $\pi:T_p\to C_p$, locally
isomorphic over $C_p$ to $S_p$, and an automorphism $F$ covering
$\eta$, with
\[
 F^{12}=t_{6P}.
 \tag{B.6.21}\label{maps:eq:B.6.21}
\]
Its torsor class can be chosen to be annihilated by four.
For $p$ outside the countable exceptional set in \ref{maps:src:B.6.12} it
has no global section. There is also a holomorphic line bundle
$M$ of relative Picard class $P$ such that, for the unique
point $c_p\in C_p$ over $p$,
\[
 M\otimes F^*M^{-1}\simeq\pi^*\mathcal O_{C_p}(c_p).
 \tag{B.6.22}\label{maps:eq:B.6.22}
\]
Thus a linearization with a simple zero at that fibre exists.
These are conductor data, not an identification of a parameter
filling with the prescribed mapping torus. The obstruction in
\ref{maps:src:B.6.14} still applies to the direct one-component plumbing.

\end{proposition}


\begin{proof}
\begin{enumerate}[label=\arabic*.]
\item Let $G=\langle\eta\rangle\simeq\mathbb Z/12$, and let
$\mathcal E$ be the abelian sheaf of holomorphic sections of
the identity Neron group of $S_p\to C_p$ in the analytic
topology. Since all singular fibres are $I_1$, this is the
smooth locus of the elliptic fibration. The analytic
exponential sequence is
\[
 0\longrightarrow\mathcal A\longrightarrow\mathcal V
   \longrightarrow\mathcal E\longrightarrow0,
 \tag{B.6.23}\label{maps:eq:B.6.23}
\]
where $\mathcal V$ is the coherent Lie line and
$\mathcal A=j_*\Lambda$ is the constructible period lattice.
This assertion also holds at an $I_1$ fibre. In the Tate chart
the exponential is $z\mapsto e^{2\pi iz}$; at the central
$\mathbb C^*$ fibre its kernel is the invariant rank-one
lattice. A local nonvanishing section has a holomorphic
logarithm near each value, proving local surjectivity.

\item We first compute the obstruction group on the quotient
orbifold $\mathcal B=[C_p/G]$, whose coarse curve is
$B=\mathbb P^1$. Write $q:\mathcal B\to B$.
The higher direct images $R^jq_*\mathcal E$, $j>0$, are
skyscraper sheaves at the three branch values and are the
cohomology of the corresponding stabilizers on section germs.
At such a point, choose a small invariant root disc. The
formal subgroup at zero is uniquely divisible and its
logarithm is a complex vector space, so its higher finite
group cohomology vanishes by averaging. Only the action on
the central elliptic curve contributes.

\item Here are the local calculations, with actual sections.
Engel's generator can be taken as
$P=(0,\sqrt{2}\,t^2(t-1))$. At zero the root chart
$t=z^3,x=z^4X,y=z^6Y$ has limiting curve
$Y^2=X^3+2$, section $(0,-\sqrt2)$, and stabilizer action
$(X,Y)\mapsto(\zeta_3^2X,Y)$. Its fixed group consists of
$O,(0,\sqrt2),(0,-\sqrt2)$ and has order three.
At one, use $t=1+z^4,x=z^2X,y=z^3Y$. The limit is
$Y^2=X^3-3X$, with section $(0,0)$ and action
$(X,Y)\mapsto(-X,iY)$. Its fixed group is
$\{O,(0,0)\}$ of order two. The norms on either of these
nontrivial complex-multiplication elliptic curves are zero;
hence their cyclic $H^2$ groups are respectively
$\mathbb Z/3$ and $\mathbb Z/2$, generated by the
specialization of $P$. At $c_p$, the stabilizer has order
twelve but acts trivially on the elliptic fibre, and
multiplication by twelve is onto, so its $H^2$ is zero.
The nodal fibres are unramified and have trivial stabilizers.

\item The restrictions of $6P$ vanish in all these local groups.
For an explicit order-three witness use $2P$.
At the order-four point $6P$ specializes to zero; the unique
local inverse of multiplication by four in the formal group
is invariant, so its norm is $6P$. At $c_p$ choose a local
holomorphic half of $P$ on the ordinary downstairs disc;
its pullback has norm $6P$. The root equations, actions and
local witnesses are checked in
\texttt{fourfolds\_B6\_torsor\_local.py}.

\item There is no additional degree-two obstruction on the coarse
curve. To see this without omitting a component sheaf, take
invariants in \eqref{maps:eq:B.6.23}. The Lie invariants are coherent and
higher finite-group cohomology of them is zero. The resulting
sequence is
\[
 0\longrightarrow q_*\mathcal A\longrightarrow q_*\mathcal V
  \longrightarrow q_*\mathcal E
  \longrightarrow R^1q_*\mathcal A\longrightarrow0.
\]
The last sheaf is a finite skyscraper sheaf, and the first is
constructible on a real two-dimensional curve. Splitting this
sequence at its middle image, coherent acyclicity in degrees
at least two and $H^j(B,q_*\mathcal A)=0$ for $j>2$ give
\[
 H^2(B,q_*\mathcal E)=H^3(B,q_*\mathcal E)=0.
\]
Skyscraper sheaves have no higher cohomology. The coarse-map
spectral sequence therefore identifies restriction to the
stabilizers with an isomorphism
\[
 H^2(\mathcal B,\mathcal E)\simeq
       \mathbb Z/3\oplus\mathbb Z/2.
 \tag{B.6.24}\label{maps:eq:B.6.24}
\]
In particular the inflation of the cyclic class $[6P]$
vanishes in this group.

\item The group--Cech double complex for $G$ acting on
$\mathcal E$ gives the five-term exact sequence, whose
relevant part is
\[
 H^1(C_p,\mathcal E)^G\longrightarrow
 H^2(G,H^0(C_p,\mathcal E))\longrightarrow
 H^2(\mathcal B,\mathcal E).
 \tag{B.6.25}\label{maps:eq:B.6.25}
\]
Here $H^0(C_p,\mathcal E)=L_p$, and cyclic group cohomology
identifies the middle group with $L_p^\eta/NL_p$.
Thus $[6P]$ is the transgression of an invariant analytic
torsor class $\alpha$. In Cech terms, invariance supplies
an isomorphism $F$ of the torsor covering $\eta$; its
twelfth power is translation by an invariant section $R$,
and the transgression is $[R]$ modulo norms. Its sign is
immaterial here because $[6P]$ has order at most two.
Replacing $F$ by $t_QF$ adds $NQ$ to $R$ and therefore
arranges \eqref{maps:eq:B.6.21} exactly.

\item We justify the finite-order choice of $\alpha$.
The elliptic monodromy on $C_p$ minus its twelve nodal
points contains both $R_1$ and $R_2$: in the four-punctured
downstairs group, multiply either old loop by a power of
the new-point loop to kill its covering character. The
latter has identity elliptic monodromy and covering
character generating $\mathbb Z/12$. Filling the ramified
points kills only loops with identity elliptic monodromy.
The columns $(1,-2)^t$ of $R_1-I$ and $(-1,1)^t$ of
$R_2-I$ form a unimodular pair. Thus the global elliptic
lattice coinvariant group is zero. The exact sequence
$j_!\Lambda\to j_*\Lambda$ has skyscraper cokernel,
so compact-support duality on the punctured curve gives
\[
 H^2(C_p,\mathcal A)=H^2_c(C_p^\circ,\Lambda)=0.
\]
Consequently $H^1(C_p,\mathcal V)\to H^1(C_p,\mathcal E)$
is onto. By relative Serre duality and \ref{maps:src:B.6.10} the characters
of $H^1(C_p,\mathcal V)$ are $6,3,2$ modulo twelve.
There is no invariant vector, so the group norm on this
space, and hence on $H^1(C_p,\mathcal E)$, is zero.
Every invariant class is therefore annihilated by twelve.
Replace $\alpha$ by $9\alpha$: its transgression is still
$[6P]$, while $4(9\alpha)=0$.

\item Represent this class by holomorphic section-valued Cech
translations. Gluing the corresponding copies of $S_p$
gives a compact smooth genus-one surface locally isomorphic
to $S_p$, with no multiple fibres. Such translations extend
over $I_1$ because they act by multiplication on the nodal
curve's normalization. Invariance and the preceding norm
adjustment give the required holomorphic $F$. If this torsor
had a section at a generic $p$, it would be the sectioned
model excluded by \ref{maps:src:B.6.12}, proving its nontriviality.

\item It remains to construct and normalize $M$. On a local
trivialization of the torsor use the line bundle
$\mathcal O(P-O)$. On overlaps these lines have the same
relative degree-zero Picard class; their difference is a
line from the base. The scalar obstruction on triple
overlaps lies in $H^2(C_p,\mathcal O^*)=0$, by the ordinary
exponential sequence and the real dimension of the curve.
Thus they glue to a line bundle $M$ of relative class $P$.
Since $F$ fixes this relative class, there is a line $J$ on
$C_p$ with
\[
 M\otimes F^*M^{-1}=\pi^*J,\qquad
 \bigotimes_{i=0}^{11}\eta^{i*}J
       \simeq M\otimes t_{6P}^*M^{-1}
       \quad\hbox{after pullback to }T_p.
\]
We verify the degree on the right for this torsor, rather
than assuming the sectioned height formula applies to it.
A torsor annihilated by four has a finite splitting cover:
lift its class through the local multiplication-by-four
sequence and trivialize the resulting finite torsion torsor
on the ordinary curve. At an $I_1$ point the added division
branches have Tate coordinates equal to a unit times
$q^{j/4}$; they extend after finite base change and minimal
resolution. Thus compactification gives a finite curve
$D\to C_p$, of degree $d$, on whose regular pullback model
the torsor has a section.

\item On that model, the pulled-back $M$ has degree zero on every
fibre component and differs from $\mathcal O(P_D-O_D)$ by
a base line. All pulled-back sections $P_D,6P_D$ are narrow:
the original $I_1$ sections lie in its smooth group, and
their pullbacks meet the identity components of the
resolved $I_n$ fibres. Restriction to $O_D$ therefore gives
the degree
\[
 \chi+(P_D\cdot O_D)+(6P_D\cdot O_D)-(P_D\cdot6P_D)
       =\langle P_D,6P_D\rangle=12d.
\]
The first equality is the narrow height formula and the
second is base-change functoriality
\cite[Theorem 8.6 and Proposition 8.12]{mapsShioda90}.
Dividing by $d$ shows that the line on the right has degree
twelve. Hence $\deg J=1$.

\item Finally the deck characters on $H^1(C_p;\mathbb C)$ are
$1,2,5,7,10,11$, each once. They follow either from the
base fixed-point traces in \ref{maps:src:B.6.10} or from the three
holomorphic forms of characters $7,10,11$.
The characteristic polynomial is $\Phi_6\Phi_{12}$, so
\[
 \det(1-\eta^*\mid H^1(C_p;\mathbb Z))
      =\Phi_6(1)\Phi_{12}(1)=1.
\]
Thus $1-\eta^*$ is an automorphism of $\operatorname{Pic}^0(C_p)$.
There is a unique degree-zero line $D_0$ with
\[
 D_0\otimes\eta^*D_0^{-1}
       =\mathcal O(c_p)\otimes J^{-1}.
\]
Replace $M$ by $M\otimes\pi^*D_0$ to obtain \eqref{maps:eq:B.6.22}.
The canonical section of $\mathcal O(c_p)$ supplies the
claimed simple zero. The lattice coinvariants, base
characters and norm arithmetic are also recorded in
\texttt{maps\_component\_periods.py}.

\end{enumerate}
\end{proof}




\begin{proposition}[the one-component torsor plumbing cannot fill the required mapping torus]\label{maps:src:B.6.14}
\StatementRecord{maps:src:B.6.14}{proved}{GPT-6 Astra (OpenAI, ChatGPT Pro mode), effort Pro, 2026-11-17; Claude Fable 5.2 (Anthropic), effort high, 16 November 2026}{the hypotheses, notation and definitions stated here; \cref{maps:src:B.6.10}; \cref{maps:src:B.6.11}; \cref{maps:eq:B.T1}}


Let $W\to\Delta_r$ be a proper semistable family with smooth
complex four-dimensional total space, restricted to a small closed
disc.  Suppose that its reduced central fibre is irreducible and
that its normalization is a projective-line bundle, or the stated
conductor blow-up of such a bundle, over a genus-one torsor
$T\to C_p$ with Jacobian $S_p$ from \ref{maps:src:B.6.10}.  Suppose the map to
$B=\mathbb P^1$ extends over $W$ and restricts on that normalization
to the composite through $T\to C_p\to B$.
Then $H^2(\partial W;\mathbb Q)\ne0$.
Consequently this model does not fill the mapping torus of any
finite iterate of either the sphere monodromy or its
$S^3\times S^3$ smoothing monodromy.

More concretely, the proposed plumbing with conductor line
$\mathcal O_{C_p}(c_p)$ and its simple zero at $c_p$ produces
a rank-one fibre of multiplicity twelve at $p$.  A free quotient
of this chart retains that multiplicity.  Thus the analytic
torsor existence assertion does not by itself produce the
required primitive rank-one filling.

\end{proposition}


\begin{proof}
\begin{enumerate}[label=\arabic*.]
\item The small semistable total neighbourhood retracts to its central
fibre, as in \ref{maps:src:B.6.11}.  Because that fibre is reduced and irreducible,
its top homology is $\mathbb Z$, generated by its fundamental class;
the normal-crossing locus has real codimension two and introduces
no further top-dimensional class.  Therefore $H_6(W;\mathbb Q)$
is generated by this central divisor.
Poincar\'e--Lefschetz duality identifies the map
\[
 H^2(W,\partial W;\mathbb Q)\longrightarrow H^2(W;\mathbb Q)
\]
with its divisor class.  That class is zero: the central divisor
is the reduced zero divisor of the global parameter $r$.
The exact sequence of the pair hence makes
$H^2(W;\mathbb Q)\to H^2(\partial W;\mathbb Q)$ injective.

\item We show that the extended base class is nonzero on $W$.
The elliptic degree-one local system on $C_p$ has no global
invariants.  Indeed the cover is unramified above the original
$I_1$ point and is a normal cyclic cover.  Its monodromy group
therefore contains the original nodal transvection and all its
conjugates, in particular its conjugate by the order-three
elliptic generator.  The fixed line of the transvection is real
and one-dimensional.  An order-three elliptic generator has no
real invariant line, since its characteristic polynomial is
$T^2+T+1$.  The two fixed lines are distinct, proving the claim.

\item A genus-one torsor has the same rational degree-one local system
as its Jacobian: local translation changes act trivially on
cohomology.  The local smooth and nodal charts identify its first
direct image with the invariant extension of this system.
In the Leray spectral sequence of $T\to C_p$, the only possible
differential killing $H^2(C_p;\mathbb Q)$ consequently has zero
source $H^0(C_p,R^1)=0$.  It follows that
\[
 H^2(C_p;\mathbb Q)\longrightarrow H^2(T;\mathbb Q)
 \quad\hbox{is injective}.
\]
Pullback from $B$ to $C_p$ multiplies the degree by twelve and is
nonzero.  Pullback to a projective-line bundle is injective by
the bundle formula.  The conductor blow-up and its small
resolution preserve this injectivity: a degree-one proper
modification has a cohomological left inverse given by its
oriented pushforward, or here directly by the unchanged
conductor section.  Thus the class of $\mathcal O_B(1)$ pulls
back nontrivially to the normalization, and hence nontrivially
to $W$.  Its nonzero restriction to $\partial W$ proves the
first assertion.

\item For any diffeomorphism of $S^6$ or $S^3\times S^3$, the Wang
sequence gives zero second cohomology for its mapping torus,
since both the first and second cohomology of the fibre vanish.
This remains true for every finite iterate.  It contradicts
the preceding necessary property of the proposed filling.

\item The local multiplicity makes the mismatch visible without
cohomology.  At the unique point $c_p$ choose a coordinate $z$
with $t-p=z^{12}$.  The proposed conductor section has a simple
zero in $z$, so after multiplying coordinates by units its node
chart is
\[
 xy=rz.
\]
For $r\ne0$ this is smooth and the base map is
\[
 t-p=(xy/r)^{12}.
 \tag{B.T1}\label{maps:eq:B.T1}
\]
The divisor over $p$ is therefore twelve times each reduced local
branch.  An etale quotient preserves the order of vanishing at
a generic divisor point, so a free component-permuting quotient
does not turn \eqref{maps:eq:B.T1} into a primitive node.
The required D.3 filling instead has base function a unit times
$xy$ at a generic rank-one chart.  This proves the precise local
failure.  It does not exclude a model with additional central
components, a suitably ramified total-space quotient, or a
different normalization.  Those additional attachments would
require their own boundary comparison and characteristic-number
calculation; blowing up only the central fibre of the incorrect
model leaves its boundary unchanged and is not such a comparison.
The exact root-chart and multiplicity checks are
saved in \texttt{fourfolds\_B6\_torsor\_local.py}.

\end{enumerate}
\end{proof}




\begin{proposition}[the free-circle filling formula]\label{maps:src:B.6.15}
\StatementRecord{maps:src:B.6.15}{proved}{GPT-6 Astra (OpenAI, ChatGPT Pro mode), effort Pro, 2026-11-17; Claude Fable 5.2 (Anthropic), effort high, 16 November 2026}{the hypotheses, notation and definitions stated here; \cref{maps:src:B.6.9}; \cref{maps:src:B.6.5}}


Let $Q$ be the oriented mapping torus of a diffeomorphism of
$S^2\times S^3$ preserving orientation and the positive generator
of $H^2$.  Choose generators $x\in H^2(Q;\mathbb Z)$ and
$y\in H^4(Q;\mathbb Z)$ so that $\langle xy,[Q]\rangle=1$, and write
\[
 x^2=ky,\qquad p_1(Q)=\ell y.
\]
Let $M\to Q$ be the principal circle bundle of Euler class $x$,
oriented as the boundary of the associated complex disc bundle
$W=D(L_x)$.  For $k\ne0$ its rational first Pontryagin class vanishes
and its normalized Eells--Kuiper invariant is
\[
 \mu(M)=\frac{3\ell^2/k-8\ell-4k-12\operatorname{sign}(k)}{2688}
       \pmod {\mathbb Z}.
 \tag{B.6.26}\label{maps:eq:B.6.26}
\]
Moreover $\ell=4k+24r$ for an integer $r$, and equivalently
\[
 \mu(M)=\frac{k-\operatorname{sign}(k)+32r+144r^2/k}{224}
       \pmod {\mathbb Z}.
\]
For $k=0$, the condition $p_1(M)=0$ is equivalent to $\ell=0$,
and under this condition $\mu(M)=0$.
These formulas do not specify the integers for the surgery on the
fixed curve of the given family, nor the characteristic correction
of that surgery trace.
\end{proposition}


\begin{proof}
\begin{enumerate}[label=\arabic*.]
\item The fibre diffeomorphism acts trivially on integral cohomology, so
its Wang sequence gives the stated infinite cyclic groups in
degrees two and four.  Their product pairing is unimodular:
$y$ is the degree-three fibre generator swept around the mapping
circle.  The fibre has a unique spin structure, hence its mapping
torus is spin.  The index of its spin Dirac operator twisted by
$L_x$ is the degree-six characteristic number
\[
 \left\langle \widehat A(Q)e^x,[Q]\right\rangle
    =\frac{k}{6}-\frac{\ell}{24}\in\mathbb Z.
\]
Thus $\ell\equiv4k\pmod {24}$.

\item The tangent bundle of the disc bundle is stably the pullback of
$TQ\oplus L_x$, so $W$ has a spin-c structure of determinant
$L_x$ and
\[
 z=x,\qquad p_1(W)=(\ell+k)y.
\]
On its boundary $L_x$ is canonically trivialized by the unit vector,
so this structure induces a spin structure on $M$.  The Gysin
sequence gives, when $k\ne0$,
\[
 H^4(M;\mathbb Z)\cong\mathbb Z\oplus\mathbb Z/|k|,
 \qquad p_1(M)=\ell\,\pi^*y
\]
with $\pi^*y$ generating the displayed torsion subgroup.  Hence
$p_1(M)$ vanishes rationally.  When $k=0$, the pullback of $y$
is injective and primitive, which proves the asserted equivalence.

\item Let $U$ be the integral Thom class of the oriented complex line
bundle.  The map $H^4(W,M)\to H^4(W)$ sends $Ux$ to $x^2=ky$.
Its intersection form therefore has matrix $(k)$ and signature
$\operatorname{sign}(k)$ when $k\ne0$.  The rational relative lift
of $p_1(W)$ is $(\ell+k)Ux/k$.  Thom integration gives
\[
 \langle p_1(W)^2,[W,M]\rangle=\frac{(\ell+k)^2}{k},\qquad
 \langle p_1(W)z^2,[W,M]\rangle=\ell+k,\qquad
 \langle z^4,[W,M]\rangle=k.
\]
For the last two products one may use the relative lift $U$ of
$z$, so their definitions do not require $z^2$ to be nonzero.
Insert these values into the spin-c formula in \ref{maps:src:B.6.9} to obtain
\eqref{maps:eq:B.6.26}; substitution of $\ell=4k+24r$ gives its second form.
If $k=\ell=0$, the middle intersection map is zero, $p_1(W)=0$,
and all the displayed relative products vanish, giving $\mu(M)=0$.
The coefficient identities are checked exactly in
\texttt{code/maps\_free\_circle.py}.

\item For the proposed fixed-curve surgery, the disc bundle has the actual
free-circle mapping torus as boundary once its quotient bundle is
identified.  Comparison with the original sphere mapping torus
still requires its own parametrized spin surgery trace.  The
analytic conifold surgery of \ref{maps:src:B.6.5} is along a different curve and
supplies no such comparison for this model.  In particular neither
vanishing of the ordinary affine obstruction after six loops nor
homotopy of the quotient return fixes the integer $\ell$.

\end{enumerate}
\end{proof}




\begin{proposition}[a smooth fixed-circle replacement]\label{maps:src:B.6.16}
\StatementRecord{maps:src:B.6.16}{proved}{GPT-6 Astra (OpenAI, ChatGPT Pro mode), effort Pro, 2026-11-17; Claude Fable 5.2 (Anthropic), effort high, 2026-11-16}{the hypotheses, notation and definitions stated here; \cref{maps:src:B.6.4}; \cref{maps:src:D.2.14}; \cref{maps:src:B.6.3}; \cref{maps:eq:B.6.28}; \cref{maps:eq:B.6.29}; \cref{maps:eq:B.6.31}; \cref{maps:eq:B.6.30}}


The actual family and its even-shift closing maps admit a smooth
equivariant conifold replacement along the fixed curve of the
$\delta$ circle action.  Its fibres are $S^3\times S^3$, the circle
action on them is free, and the finite logarithmic fillings are
unchanged.  In the original ordered lattice $(g,u,w,d)$ the four
local monodromies of the resulting map to $B$ are
\[
 A_1,\quad A_2,\quad
 C_1=I+w\otimes(u^*-g^*),\quad
 C_2=I+(w-d)\otimes g^*.
 \tag{B.6.27}\label{maps:eq:B.6.27}
\]
Their product is one.  The marked parameter return has linear part
$S=I-2d\otimes g^*$ and the same finite affine defects as in
\ref{maps:src:B.6.4}; the two new cusp defects are zero.  On rational $H^3$ it
preserves a two-step filtration with one-dimensional quotients,
acting as the identity on each quotient.  The resulting integral
unipotent shear is not evaluated by these assertions.
\end{proposition}


\begin{proof}
\begin{enumerate}[label=\arabic*.]
\item Use the following short and long bases at the cusp:
\[
 a_1=w,\qquad a_2=w-d,\qquad
 b_1=u,\qquad b_2=g+u.
\]
They form an integral basis of the fibre lattice, and the cusp
nilpotent sends $b_i$ to $a_i$ and kills both $a_i$.  In these
coordinates the normalized period matrix is
\[
 \Omega_{c,0}(q)=\frac{\log q}{2\pi i}I+C_c(q),
 \tag{B.6.28}\label{maps:eq:B.6.28}
\]
where $C_c$ is holomorphic at zero.  The original fan has edge
directions $a_1,a_2,a_1-a_2$; thus it is the square grid with
each square subdivided along the $d=a_1-a_2$ diagonal.
This is precisely the small resolution of the product of two
infinite nodal-chain models with the same degeneration parameter.
There is one orbit of the resolved diagonal curves under the two
long-period translations.  Its quotient is the actual fixed curve
of \ref{maps:src:D.2.14}, not the different diagonal curve used in \ref{maps:src:B.6.3}.

\item Here is an explicit construction including its boundary matching.
Fix a small cusp disc and restrict $c$ to a compact interval
representing the parameter loop.  On a smaller disc replace
\eqref{maps:eq:B.6.28} by
\[
 \Omega_{c,s}(q)=
  \frac1{2\pi i}
   \begin{pmatrix}\log q&0\\0&\log(q-s)\end{pmatrix}
      +C_c(q),\qquad 0<|s|\ll1.
 \tag{B.6.29}\label{maps:eq:B.6.29}
\]
For clarity this inner filling can be constructed before taking
any lattice quotient.  Take the product of the two infinite
nodal-chain families, with parameters $q$ and $q-s$.  Its node
charts are
\[
 xy=q,\qquad zv=q-s,
 \quad\hbox{or equivalently}\quad xy-zv=s.
 \tag{B.6.30}\label{maps:eq:B.6.30}
\]
The two long generators shift the respective chain indices by one
and multiply the two torus coordinates by the columns of
$\exp(2\pi iC_c(q))$.  These multiplications extend on all the
chain charts, since their factors are nonvanishing.  They commute.
The usual bounded-unit lattice estimate proves that the action is
free and properly discontinuous and that its quotient is proper
over the disc: the two diagonal logarithmic lengths tend to
infinity at their respective nodes, their minimum is uniformly
large for small $q,s$, and the off-diagonal logarithms remain
bounded.  A bounded rectangle in the two chain indices is a
compact fundamental set.  This is the same estimate as the
square-grid special case of the supplied manuscript's
Theorem 4.5, Steps 1--4, with the two diagonal lengths allowed
to differ.  In particular no positivity or symmetry of $C_c$
is required.  At $s=0$ the square-grid quotient has ordinary
double points, and its indicated small resolution recovers the
given cusp, including its unit factors.

\item On an annulus separated from both $0$ and $s$ the difference
between \eqref{maps:eq:B.6.29} and \eqref{maps:eq:B.6.28} is the single-valued, small matrix
\[
 \frac{\log(1-s/q)}{2\pi i}E_{22}.
\]
Multiply it by a smooth radial cutoff equal to one on the inner
boundary and zero near the outer boundary.  Real nonsingularity
of the period matrix is open; after reducing $|s|$ this interpolation
therefore defines a smooth marked torus bundle on the annulus.
It agrees exactly with the inner chain quotient and with the
original cusp on their respective collars.  Glue the original
exterior, and retain both finite logarithmic fillings verbatim.
This constructs a smooth global replacement.  Holomorphicity of
the cutoff, or a global holomorphic smoothing, is not asserted.
All changes occur inside the cusp disc, and the outer peripheral
identification is the original marked one.

\item The construction also retains the actual parameter return.  The
period perturbation just made is independent of $c$.  The change
$c\mapsto c+h$ changes the long periods only by the integral
short vector $h d\otimes g^*$; its exponentiated lattice
multipliers are consequently unchanged, for integral $h$.
Use the existing cusp representative of the closing map, namely
the holomorphic torus unit obtained by exponentiating
\[
 \overline b_{h,c}(q)=
 \left(-\frac h2\mu(q),
 -\frac h{12}C_{c,21}(q)+h\kappa_0(q)-\frac{h^2}{24}\right)
 \tag{B.6.31}\label{maps:eq:B.6.31}
\]
in the original short coordinates $(w,d)$.  Formula \eqref{maps:eq:B.6.31} is
the checked cusp formula in Paper~I, Proposition~18.5, with $\ell=-s_0/12+\kappa_0(q)$ in the complete companion proof; here $s_0$
denotes the additive logarithm and is distinct from the new
smoothing parameter.  Change short coordinates integrally to
$(a_1,a_2)$ and use this same unit throughout the new filling.
On \eqref{maps:eq:B.6.30} a torus unit acts by
\[
 (x,y,z,v)\longmapsto
   (\lambda_1x,\lambda_1^{-1}y,
    \lambda_2z,\lambda_2^{-1}v).
\]
It preserves $q$ and the actual parameter $s$, extends over both
nodes and commutes with the two chain shifts.  It also defines the
same affine translation on the interpolated annulus.  Near its
outer boundary it is the original closing map, so the original
finite comparisons and their composition law remain unchanged.
In particular the map has not been replaced by a fractional
long-period translation.  Its additive cusp lift is single-valued
on both new punctures, giving zero new affine defects.

\item This replacement is the usual local small-resolution-to-smoothing
surgery at the one orbit of nodes.  It is therefore surgery on
the actual null-homotopic fixed $S^2$, with its unique real
framing; \ref{maps:src:D.2.14} identifies its underlying smooth fibre with
$S^3\times S^3$.  In \eqref{maps:eq:B.6.30} the $d=a_1-a_2$ circle has weights
$(1,-1,-1,1)$ on $(x,y,z,v)$.  Its only possible fixed point is
the origin, which is absent when $s\ne0$.  At either rank-one
central fibre it has a primitive component in the surviving
elliptic factor, so it is free there as well.  Together with the
unchanged exterior and finite fillings, where it was free, this
proves global freeness.

\item The primitive monodromy at $q=0$ sends $b_1$ to $b_1+a_1$;
at $q=s$ it sends $b_2$ to $b_2+a_2$.  Expressing these two maps
in the original basis gives \eqref{maps:eq:B.6.27}, and direct multiplication
gives $C_1C_2=M_0$.  They commute with $S$.  The finite affine
defects remain $(0,1,-1,-1)$ and $(0,0,-1,1)$, as already
computed in \ref{maps:src:B.6.4}.

\item Finally the invariant-specialization calculation of \ref{maps:src:B.6.3} applies
to each of these two primitive rank-one fibres and to the unchanged
finite fibres.  For $j$ the inclusion of the four-value complement
and $V=R^1g_*\mathbb Q$ there is again
$R^qg_*\mathbb Q=j_*\bigwedge^qV$.  The parabolic cocycle
calculation with the new matrices gives, in degrees $q=0,\ldots,4$,
\[
 \dim H^0(B,j_*\bigwedge^qV)=(1,1,0,1,1),\qquad
 \dim H^1(B,j_*\bigwedge^qV)=(0,0,0,0,0).
\]
The degree-one global coinvariant space is one-dimensional, so
$\dim H^2(B,j_*V)=1$.  The exact matrix calculation is recorded
in \texttt{maps\_fixed\_curve\_monodromy.py}.  Thus the only
possible Leray pieces of degree three are the one-dimensional
groups $E_2^{0,3}$ and $E_2^{2,1}$.  Since the total space is
$S^3\times S^3$, both survive and exhaust $H^3$.  The return acts
as the identity on each: on the first by direct invariant-form
calculation, and on the second because $S-I$ has image in the
vanishing coboundaries.  Saturating the rank-one integral
filtration gives an integral upper triangular return with
diagonal entries one.  The off-diagonal integer is an extension
in this filtration.  The earlier analytic smoothing had all its
$H^3$ in $E_2^{1,2}$, so its Torelli conclusion cannot be
transferred to this different surgery.

\end{enumerate}
\end{proof}




\begin{proposition}[the actual quotient Pontryagin class]\label{maps:src:B.6.17}
\StatementRecord{maps:src:B.6.17}{proved}{GPT-6 Astra (OpenAI, ChatGPT Pro mode), effort Pro, 2026-11-17; GPT-6 Astra (OpenAI, ChatGPT Pro mode), S6\_06\_BLIND, effort Pro, 2026-11-17; Claude Fable 5.2 (Anthropic), effort high, 2026-11-16}{the hypotheses, notation and definitions stated here; \cref{maps:src:B.6.16}; \cref{maps:src:B.6.1}; \cref{maps:src:B.6.15}; \cref{maps:src:D.2.14}; \cref{maps:eq:B.Q3}; \cref{maps:eq:B.Q9}; \cref{maps:eq:B.Q6}; \cref{maps:eq:B.Q7}; \cref{maps:eq:B.Q8}; \cref{maps:eq:B.Q4}}


Use the fixed-circle replacement of \ref{maps:src:B.6.16}, the positive parameter
loop, and the mapping-torus convention
$(z,1)\sim(\Phi z,0)$ of \ref{maps:src:B.6.1}.  Let $M_s$ be its
seven-dimensional mapping torus and let
\[
 \pi:M_s\longrightarrow Q=M_s/S^1_\delta
\]
be the principal circle bundle, with Euler class $x$.
Orient $Q$ by putting the positive $\delta$ direction first
in the orientation of $M_s$.  The bundle over the parameter
circle has fibre $S^2\times S^3$.  The class $x$ is primitive,
and there is a unique $y\in H^4(Q;\mathbb Z)$ with
$\langle xy,[Q]\rangle=1$.  If $Z_1,Z_2\subset Q$ are the
quotients of the two swept critical elliptic curves, oriented
by their complex normal bundles, then
\[
 p_1(Q)=-2\operatorname{PD}_Q(Z_1+Z_2),\qquad
 \langle x,[Z_1]\rangle=2,\quad
 \langle x,[Z_2]\rangle=0.
 \tag{B.Q1}\label{maps:eq:B.Q1}
\]
Consequently
\[
                       p_1(Q)=-4y.
 \tag{B.Q2}\label{maps:eq:B.Q2}
\]
This computes the integer $\ell$ in \ref{maps:src:B.6.15}, with these
orientations.  It does not compute $k$ in $x^2=ky$ or the
relative Pontryagin number of the parametrized surgery trace.

\end{proposition}


\begin{proof}
\begin{enumerate}[label=\arabic*.]
\item The quotient of each six-dimensional fibre is
$B_5=S^2\times S^3$, by the free-circle description in
\ref{maps:src:D.2.14} applied to the nullsphere surgery.  Its circle Euler
class is primitive.  One can see its generator geometrically:
the vanishing $S^3$ of the local conifold smoothing is
invariant under the free circle, and its quotient is the new
$S^2$ in $B_5$.  The restricted circle bundle is a Hopf bundle,
so its Euler number has absolute value one.  In particular
this sphere generates $H_2(B_5;\mathbb Z)$.
The return preserves the circle orientation and hence its
Euler class.  The integral Wang sequence, with
$H^1(B_5)=0$, gives
\[
 H^2(Q;\mathbb Z)\xrightarrow{\;\cong\;}H^2(B_5;\mathbb Z).
 \tag{B.Q3}\label{maps:eq:B.Q3}
\]
It also gives $H^4(Q;\mathbb Z)=\mathbb Z$ and no torsion
in these degrees.  Integral duality proves the assertion
about $y$.

\item We keep the complex structures on all torus fibres and on
the local finite and nodal models.  On the interpolating
annulus in \ref{maps:src:B.6.16}, the period matrix is smooth rather than
holomorphic in the base.  It nevertheless defines a smooth
vertical complex structure.  An invariant horizontal
distribution extends it to an almost complex structure on
the six-dimensional fibres, making the map to $B$
pseudoholomorphic wherever it is a submersion and agreeing
with the local complex structures near every critical
fibre.  Such horizontal choices form an affine space off
the critical collars, so they can be patched there and
averaged over $S^1_\delta$.  They are made on the parameter
mapping torus itself.  Thus they respect the actual closing
map, not just a chosen trivialization over an interval.

\item Let $E$ be the resulting complex rank-three bundle tangent
to the six-dimensional fibres of $M_s\to S^1$.  The circle
action on $E$ is complex linear, and it is free on the base.
It descends to a complex bundle $\overline E$ on $Q$.
The vector field $\delta$ and the parameter direction are
invariant trivial real lines, so the real bundles satisfy
\[
 (\overline E)_{\mathbb R}\oplus\mathbb R
                         \cong TQ\oplus\mathbb R.
 \tag{B.Q4}\label{maps:eq:B.Q4}
\]
In particular $p_1(Q)=p_1(\overline E)$.

\item Write $F:Q\to B$ for the descended base map.  At a multiple
fibre of order $m_i$, divide the derivative by its known
factor of order $m_i-1$.  Its resulting target line on $Q$
is
\[
 \overline L=F^*TB\otimes
        \mathcal O_Q\!\left(-\sum_i(m_i-1)\overline S_i\right),
 \tag{B.Q5}\label{maps:eq:B.Q5}
\]
where $\overline S_i$ is the quotient of the reduced
multiple-fibre divisor in $M_s$.  The notation
$\mathcal O_Q(\overline S_i)$ means the complex line defined
by its oriented complex normal bundle; no complex structure
on all of $Q$ is required.
The generating sphere just described lies inside the
conifold patch, whose base image is a disc.  Therefore
$F^*[\mathrm{point}]$ evaluates to zero on that sphere.
Equation \eqref{maps:eq:B.Q3} proves that it is zero in $H^2(Q;\mathbb Z)$.
The local root-disc model also gives, as an equality of
integral divisor classes on $Q$,
\[
 m_i[\overline S_i]=F^*[\mathrm{point}]=0.
\]
Since $H^2(Q;\mathbb Z)$ is torsion free, every
$[\overline S_i]$ is zero.  Thus
\[
                         c_1(\overline L)=0.
 \tag{B.Q6}\label{maps:eq:B.Q6}
\]

\item The determinant of $\overline E$ is also trivial in
cohomology.  It suffices by \eqref{maps:eq:B.Q3} to test it on the new
sphere.  Its inverse image is the local vanishing $S^3$ in
\[
                         xy-zv=s\subset\mathbb C^4.
\]
The circle weights are $(1,-1,-1,1)$.  The holomorphic
residue volume form of this hypersurface is nowhere zero
and invariant: the numerator has total weight zero and
the defining equation has weight zero.  It trivializes the
equivariant determinant on this $S^3$, hence trivializes
the descended determinant on its Hopf quotient.  Therefore
\[
                         c_1(\overline E)=0.
 \tag{B.Q7}\label{maps:eq:B.Q7}
\]
This checks the possible determinant contribution on the
new sphere; it is not inferred from a determinant computation
only on the ordinary torus locus.

\item Let $V$ denote the rank-two logarithmic vertical tangent
bundle.  It is the ordinary vertical tangent bundle away
from the critical curves.  Near a rank-one node with
$q=xy$ and elliptic coordinate $z$, it is the locally free
bundle with evaluation generators
$x\partial_x-y\partial_y$ and $\partial_z$.
These local extensions glue to the ordinary kernel off
the critical loci.  They descend to $\overline V$ on $Q$.
The complex vector field generated by $\delta$ is an
invariant nowhere-zero section of $V$: the circle is free,
and at either node it has a primitive component in the
surviving elliptic direction.  Hence
\[
                         c_2(\overline V)=0.
 \tag{B.Q8}\label{maps:eq:B.Q8}
\]

\item After division at the multiple fibres, the derivative
has no remaining failure of surjectivity except the two
rank-one critical loci.  Near each, the complex Koszul
resolution of $(x,y)$ identifies its cokernel with the
target line restricted to the critical curve.  These local
resolutions and the actual kernel off the critical curves
give the following identity in complex topological
$K$-theory of $Q$:
\[
 [\overline E]=[\overline V]+[\overline L]
       -[\overline L|_{Z_1}]-[\overline L|_{Z_2}].
 \tag{B.Q9}\label{maps:eq:B.Q9}
\]
Here the last two terms mean the supported $K$-classes,
extended by their complex-normal Thom classes.  Their
degree-four Chern characters are
$\operatorname{PD}_Q(Z_1)$ and
$\operatorname{PD}_Q(Z_2)$, respectively.  The identity
is consequently valid for the smooth interpolation as
well: it is exact off the critical curves, and on their
neighbourhoods it is the displayed holomorphic Koszul
calculation.  No integrability on the annulus is used.

\item The degree-two part of \eqref{maps:eq:B.Q9}, together with \eqref{maps:eq:B.Q6}--\eqref{maps:eq:B.Q7},
gives $c_1(\overline V)=0$.  Its degree-four part and
\eqref{maps:eq:B.Q8} then give
\[
 c_2(\overline E)=\operatorname{PD}_Q(Z_1+Z_2).
\]
Equations \eqref{maps:eq:B.Q4} and \eqref{maps:eq:B.Q7} prove the first formula of
\eqref{maps:eq:B.Q1} integrally.  Equality first in rational cohomology
suffices here because $H^4(Q;\mathbb Z)$ is torsion free.

\item It remains to fix the two Euler signs.  Use the short and
long bases of \ref{maps:src:B.6.16} and let $\kappa$ denote the positive
parameter direction.  At the first node the surviving
elliptic basis is $(a_2,b_2)$, its complex orientation is
that order, and
\[
 \delta=-a_2,\qquad S b_2=b_2+2a_2.
\]
With circle generator $a_2$, the critical mapping torus
has unit-line clutching
\[
 (\beta,1,z)\sim(\beta,0,e^{2\pi i\,2\beta}z).
\]
Its Chern number is $+2$ on the torus oriented by
$(b_2,\kappa)$.  For example, a unitary connection with
local form $2\pi i\,2\kappa\,d\beta$ has curvature
$-2\pi i\,2\,d\beta\wedge d\kappa$, and the convention
$c_1=iF/(2\pi)$ gives that number.  The ambient quotient
orientation with circle $\delta=-a_2$ orients $Z_1$ as
$-(b_2,\kappa)$.  Reversing both the circle and the base
orientation leaves the Euler pairing $+2$ unchanged.
At the second node the surviving basis is $(a_1,b_1)$,
$\delta=a_1$, and the return is the identity on that
lattice.  Its Euler pairing is zero.  These are the last
two assertions of \eqref{maps:eq:B.Q1}.  Pairing its first assertion
with $x$ proves \eqref{maps:eq:B.Q2}.

\end{enumerate}
\end{proof}




\begin{proposition}[the full quotient class has one integer left]\label{maps:src:B.6.18}
\StatementRecord{maps:src:B.6.18}{proved}{GPT-6 Astra (OpenAI, ChatGPT Pro mode), effort Pro, 2026-11-17; Claude Fable 5.2 (Anthropic), effort high, 2026-11-16}{the hypotheses, notation and definitions stated here; \cref{maps:src:B.6.16}; \cref{maps:src:B.6.17}; \cite[Corollary 1.3 and Theorem 4.9]{mapsJin26}}


Let $\overline\Psi$ be the quotient return on $S^2\times S^3$
in \ref{maps:src:B.6.16}--\ref{maps:src:B.6.17}.  Let $a$ be the Dehn twist about the primitive
$S^2$, using the generator of $\pi_3SO(3)$ whose associated
three-plane bundle over $S^4$ has Pontryagin number $4$.
Let $b$ be the same twist about an embedded sphere representing
twice that primitive class.  Then there is a unique integer $B_0$
such that
\[
 [\overline\Psi]=a^{-1-2B_0}b^{B_0},\qquad
 k=-1+6B_0,\qquad B_0=\frac{k+1}{6}.
 \tag{B.6.32}\label{maps:eq:B.6.32}
\]
This is an equality of smooth isotopy classes.  In particular,
there is no additional torsion ambiguity in the quotient return.
The integer $B_0$ has not yet been evaluated for the given family,
and this statement does not identify its reverse-surgery image
in $\Theta_7$.
\end{proposition}


\begin{proof}
\begin{enumerate}[label=\arabic*.]
\item The quotient return preserves its primitive circle Euler class
and its orientation, so it acts trivially on the integral
homology of $S^2\times S^3$.  Jin's mapping-torus construction
identifies this Torelli group with
$\Omega^{\operatorname{Spin}}_6(\mathbb CP^\infty)$;
his Proposition 3.1 and Theorem 4.9 identify the latter group
with $\mathbb Z^2$ by the characteristic numbers
\[
 \left(\frac{4k-\ell}{24},\ k\right).
\]
The injectivity is explicitly proved in his Lemma 2.1 using
modified surgery and the pseudo-isotopy theorem, so this is
an isotopy classification, not just a detection of rational
invariants.  His Lemmas 4.4 and 4.8, and the table on
pp.~23--24 of version 2, give
\[
 (k(a),\ell(a))=(1,4),\qquad
 (k(b),\ell(b))=(8,8).
\]
These two twists freely generate the Torelli group
\cite[Corollary 1.3 and Theorem 4.9]{mapsJin26}.
The framing is unique because $\pi_2SO(3)=0$; the indicated
sphere embeddings are determined by their homotopy classes,
as stated immediately before his Theorem 1.2.

\item Write the unique quotient class as $a^A b^{B_0}$.  Additivity
of the mapping-torus bordism invariant gives
\[
                  k=A+8B_0,\qquad \ell=4A+8B_0.
\]
The actual value $\ell=-4$ from \ref{maps:src:B.6.17} therefore gives
$A+2B_0=-1$, which is exactly \eqref{maps:eq:B.6.32}.  Conversely these
equations recover $A,B_0$ from the two characteristic numbers,
so no further smooth quotient class has the same pair.
This last arithmetic is checked in
\texttt{maps\_quotient\_classification.py}.

\end{enumerate}
\end{proof}




\begin{proposition}[the fixed-curve trace correction]\label{maps:src:B.6.19}
\StatementRecord{maps:src:B.6.19}{proved}{GPT-6 Astra (OpenAI, ChatGPT Pro mode), effort Pro, 2026-11-17; Claude Fable 5.2 (Anthropic), effort high, 2026-11-16}{the hypotheses, notation and definitions stated here; \cref{maps:src:B.6.16}; \cref{maps:eq:B.6.30}; \cref{maps:src:B.6.5}; \cref{push:src:C.1.1}; \cref{push:src:C.1.2}; \cref{push:src:C.2.1}; \cref{maps:src:B.6.17}; \cref{maps:eq:B.6.37}; \cref{maps:eq:B.6.36}; \cref{maps:src:B.6.15}}


For the actual equivariant surgery in \ref{maps:src:B.6.16}, let $Z$ be the
eight-dimensional parametrized surgery trace, with oriented
boundary $M_s\sqcup(-M_\Phi)$.  It is spin.  If $k\ne0$,
its middle relative intersection form has rank one and absolute
determinant $|k|$.  There is a primitive generator
$\alpha\in H^4(Z;\mathbb Z)$ with
\[
                 p_1(Z)=(\ell-2k)\alpha.
 \tag{B.6.33}\label{maps:eq:B.6.33}
\]
In particular, if the actual integer $k$ is $-1$, then
\[
 \langle p_1(Z)^2,[Z,\partial Z]\rangle
                  -4\operatorname{sign}(Z)=0,
 \qquad [\Phi]=0\in\Theta_7.
 \tag{B.6.34}\label{maps:eq:B.6.34}
\]
The last conclusion is conditional here on that still required
integer calculation; it does not follow from $\ell=-4$ alone.
\end{proposition}


\begin{proof}
\begin{enumerate}[label=\arabic*.]
\item First the trace exists with the actual parameter closing.
The local equation \eqref{maps:eq:B.6.30} is a fixed quadratic smoothing,
and the closing maps preserve both $q$ and $s$.  Their torus
unit actions preserve its small-resolution ruling.  Retraction
of this ruling-preserving quadratic group to its compact form
gives the usual equivariant $D^3\times D^4$ surgery handle,
as in the explicit $SO(4)$ trace in \ref{maps:src:B.6.5}.  In the present
case the $\delta$ circle fixes the $D^3$ factor and acts
freely by Hopf rotations on the punctured $D^4$ factor.
The fixed set of the parametrized trace is therefore a
$D^3$ bundle over the parameter circle.  In particular it
has the homotopy type of $S^1$.  The attaching $S^2$ has
its unique spin structure and its unique stable framing
class; the spin surgery extends the unique spin structure
of each $S^6$.  Since the seven-dimensional trace fibre
is simply connected, its spin structure is unique and
is preserved by the closing map.  It gives a spin
structure on $Z$.  Either resulting spin choice on
$M_\Phi$ has the same invariant by \cref{push:src:C.1.1} and \cref{push:src:C.1.2} and \cref{push:src:C.2.1}.

\item Let $W$ be the seven-dimensional trace fibre.  It is obtained
by adding one index-three handle to $S^6\times I$ along
a null-homotopic $S^2$.  Its nonzero reduced integral
homology groups are $H_3(W)=\mathbb Z$ and
$H_6(W)=\mathbb Z$.  The parameter return preserves
their generators: the handle orientation is preserved
by the fixed quadratic model.  The Wang sequence gives
\[
 H^1(Z)=\mathbb Z\kappa,\quad H^2(Z)=0,\quad
 H^3(Z)=\mathbb Z\beta,\quad
 H^4(Z)=\mathbb Z\alpha,\qquad \alpha=\kappa\beta.
 \tag{B.6.35}\label{maps:eq:B.6.35}
\]
All groups here have integral coefficients.  The product
assertion follows as well from the two columns of the
spectral sequence of $W\to Z\to S^1$.

\item We use integral Borel equivariant cohomology for the
$\delta$ action.  Put $u=c_1(\mathbb C_1)$ in
$H^2(BS^1;\mathbb Z)$, with its sign chosen so that
its restriction on the free boundary is the Euler
class $x$ of \ref{maps:src:B.6.17}.  The fixed set provides a section
over $BS^1$ and over $BS^1\times S^1$ after choosing
a section of its contractible disc bundle.  Thus
$u^2$ and $u^2\kappa$ inject into equivariant cohomology.
In the Borel spectral sequence the only possible
differential from $\beta$ is $d_4\beta$ to the cyclic
group generated by $u^2$; the preceding injection forces
it to vanish.  The $d_2$ target vanishes because $H^2(Z)=0$.
Choose an equivariant lift of $\beta$ and multiply it by
$\kappa$ to lift $\alpha$.  On the fixed set the lift
of $\beta$ can only restrict to a multiple of $u\kappa$,
so this lift of $\alpha$ restricts to zero.  It follows that
\[
 H^4_{S^1}(Z;\mathbb Z)=\mathbb Zu^2\oplus\mathbb Z\alpha,
 \qquad \alpha|_{Z^{S^1}}=0.
 \tag{B.6.36}\label{maps:eq:B.6.36}
\]
Here and below the same letter denotes the lift of the
ordinary primitive class.  No torsion extension is
possible between the two free displayed summands.

\item At a fixed point the normal real four-plane has the two
unit circle weights, so its equivariant first Pontryagin
class is $2u^2$.  The fixed set has no ordinary positive
even cohomology, hence supplies no additional normal
or tangent characteristic term.  Consequently
\[
                   p_1^{S^1}(Z)=2u^2+c\alpha
 \tag{B.6.37}\label{maps:eq:B.6.37}
\]
for an integer $c$.  Forgetting the group action gives
$p_1(Z)=c\alpha$.

\item The relative equivariant handle calculation makes the
outgoing restriction integral and primitive.  Excision
and the equivariant Thom isomorphism for its normal
$D^4$ give
\[
 H^4_{S^1}(Z,M_s)=\mathbb ZU,\qquad
 H^5_{S^1}(Z,M_s)=\mathbb ZU\kappa.
\]
The fixed restriction of the Thom class is its normal
Euler class, namely $\varepsilon u^2$, where
$\varepsilon=\pm1$ allows the chosen real normal
orientation.  Thus $U\kappa$ maps to a class whose
fixed restriction is $\varepsilon u^2\kappa\ne0$.
The map from relative to absolute degree-five
cohomology is injective.  The exact sequence of the
pair therefore gives a surjection
\[
 H^4_{S^1}(Z)\longrightarrow H^4_{S^1}(M_s)
                                  =H^4(Q)=\mathbb Z.
\]
Its kernel is generated by $U$, whose $u^2$ coefficient
is $\varepsilon$.  Hence the restriction of $\alpha$
is primitive.  Choose its sign so that
$\alpha|_Q=y$.  Restricting \eqref{maps:eq:B.6.37} then gives
\[
                    \ell y=2x^2+cy=(2k+c)y,
\]
which proves \eqref{maps:eq:B.6.33}.

\item Write the image of the relative Thom class in \eqref{maps:eq:B.6.36}
as $\varepsilon u^2+m\alpha$.  Its restriction to $Q$
is zero, so $m=-\varepsilon k$.  Forgetting equivariance
shows that the ordinary relative intersection map sends
the primitive Thom generator to $m\alpha$.  Passing
from the pair $(Z,M_s)$ to $(Z,\partial Z)$ changes
neither group in degree four, since
$H^3(M_\Phi)=H^4(M_\Phi)=0$.  Integral
Poincare--Lefschetz duality identifies the pairing of
the two primitive generators with a unit.  Thus the
middle form is a rank-one form of absolute determinant
$|m|=|k|$.  In particular, when $k=-1$ its matrix is
$(\epsilon)$ for some $\epsilon=\pm1$.

\item In this case \ref{maps:src:B.6.17} gives $\ell=-4$ and hence $c=-2$.
The rational relative lift of $p_1(Z)$ has square
$\epsilon c^2=4\epsilon$, while the signature is
$\epsilon$.  Their difference is zero for either
choice of the common orientation sign.  The spin
cobordism formula therefore gives
$\mu(M_\Phi)=\mu(M_s)$.  Substitution of
$k=-1,\ell=-4$ into \ref{maps:src:B.6.15} gives $\mu(M_s)=0$.
Finally \cref{push:src:C.1.1} and \cref{push:src:C.1.2} and \cref{push:src:C.2.1} identifies $28\mu(M_\Phi)$ with the
specified class $[\Phi]\in\mathbb Z/28$, proving
the conditional conclusion \eqref{maps:eq:B.6.34}.

\end{enumerate}
\end{proof}




\begin{theorem}\label{maps:src:B.6.20}
\StatementRecord{maps:src:B.6.20}{proved}{GPT-6 Astra (OpenAI, ChatGPT Pro mode), effort Pro, 2026-11-17; Claude Fable 5.2 (Anthropic), effort high, 2026-11-17; GPT-6 Astra (OpenAI, ChatGPT Pro mode), S6\_06\_BLIND, effort Pro, 2026-11-17; Claude Fable 5.2 (Anthropic), effort high, 2026-11-16}{Definitions and the displayed constructions}


For the marked fixed-curve surgery of the minimal $(3,4)$ family in
\ref{maps:src:B.6.16}, let $Q$ be the six-dimensional free-circle quotient of the
parameter mapping torus. Its primitive Euler class $e$ satisfies
$|\langle e^3,[Q]\rangle|=1$. In the oriented convention of \ref{maps:src:B.6.17},
$\ell=-4$, hence the spin congruence of \ref{maps:src:B.6.15} gives $k=-1$.
Consequently the spin-trace calculation of \ref{maps:src:B.6.19} gives the zero class
in $\Theta_7$ for the original minimal-family monodromy.

\emph{Proof.}
Let $g,u,w,k$ be the integral fibre lattice of the free-circle quotient,
with dual forms $G,U,W,K$.  In this lattice notation the last
letter is the parameter direction, distinct from the integer cup
coefficient $k$ in the statement.  The four homology matrices are
\[
 A_1=\begin{pmatrix}1&0&0&0\\6&0&1&-1\\-6&-1&-1&1\\0&0&0&1\end{pmatrix},\quad
 A_2=\begin{pmatrix}1&0&0&0\\0&0&-1&0\\-6&1&0&1\\0&0&0&1\end{pmatrix},
\]
\[
 C_1=I+w(u^*-g^*),\qquad C_2=I+wg^*,\qquad A_1A_2C_1C_2=I.
\]
The finite caps have orders $3,4$ and power vectors
$v_1=(1,2,-4,0)$, $v_2=(-1,-3,3,0)$.
The other two caps are the ordinary primitive nodal products after
quotienting by the free transverse circle. The exact normal-torus
Mayer--Vietoris calculation in \texttt{manifolds\_mapping\_quotient.py} gives
$H_2(Q)=\mathbb Z$, and the regular $gk$-torus maps to twice its generator.
The primitive degree-two class $x$ therefore has fibre restriction
$E_0=2G\wedge K$ after choosing the sign of $x$.

\paragraph{Closed forms on the exterior.}
Remove four discs from the oriented base sphere and choose three disjoint
arcs from the fourth boundary to the first three.  A flat torus bundle over
this surface retracts over the resulting graph. Fibrewise constant forms
and forms supported in disjoint narrow strips about these arcs give a
multiplicative de Rham model: the fibre part is the exterior algebra of
$G,U,W,K$; the first three meridian parts are one-forms with values in this
algebra, with differential $A_i^T-I$.  This assertion can also be checked
by cutting the base to a disc, using the actual linear torus gluings, and
applying de Rham Mayer--Vietoris. In particular, representatives whose
mixed base terms have disjoint support have no product containing two
such base terms.

The integral normal-torus calculation gives the following representative
of $x$ on the exterior:
\[
 E_0=2GK,\qquad
 \eta_1=-4G-W,\quad\eta_2=-3G+U+K,\quad\eta_3=0.
\]
Here a boundary path with parameter $t$ carries the two-form
$E_0+h(t)dt\wedge\eta_i$, where $h$ is a smooth bump of integral one
vanishing near the chosen seam. The notation with $h=1$ below is its
piecewise polynomial path-model shorthand. Its cube is zero, as $E_0^2=0$ and the mixed
terms can be supported in disjoint strips. Independently, if $\alpha_i$
are the extra $\delta$-rows of the actual lifted monodromy, then
\[
 \alpha_1=-2G+U+K,\quad \alpha_2=3G+W-K,\quad\alpha_3=0,
 \qquad\eta_i=-\alpha_i+(A_i^T-I)(-U).
\]
Thus this representative is the actual principal-circle curvature class
up to the stated convention for its overall sign, not a class inferred
from its vertical part alone.

\paragraph{Geometric cap restrictions.}
On a finite cap the constant one-forms $G,K$ descend rationally to the
finite affine torus quotient. Untwisting the affine path by
$x\mapsto x-tv_i/m_i$ pulls a constant form $\zeta$ back to
$\zeta-dt\wedge\iota_{v_i/m_i}\zeta$. This is an actual pullback of forms,
so it preserves products. Consequently a decomposable representative
of $x$ on each finite cap restricts to
\[
 E_0+dt\wedge\beta_i,
 \qquad \beta_1=-\tfrac23K,\quad\beta_2=\tfrac12K.
\]
This convention uses the linear boundary gluing $(x,1)\sim(A_ix,0)$;
its affine cap map is $[x,t]\mapsto[x-tv_i/m_i]$. Thus its normal
meridian is $m_i\xi+v_i$. Reversing the angular convention changes
the final orientation sign, not the absolute cube. The sign also follows
from evaluating on the exact normal tori
$(-\omega(k,v_i),j_i(N_i k))$ in the companion integral calculation.
On each nodal cap, $G$ and $K$ extend as closed one-forms: in the marking
$a_1=w,a_2=w-\delta,b_1=u,b_2=g+u$, the quotient uses the transverse
primitive circle, and $G=b_2^*$: this is the transverse circle on the first cap and
the surviving longitude of the pinched factor on the second cap. The longitude form can be made constant
near the node by an exact, monodromy-invariant change. Its products with
$K$ give a cap representative restricting to $E_0$, so $\beta_3=\beta_4=0$.
All four cap representatives have square zero.

\paragraph{Boundary primitives.}
Put $b_i=\eta_i-\beta_i$.  In the piecewise polynomial notation, a one-form
$\Theta_i(t)=\lambda_i+t b_i$ is globally defined on the boundary mapping
torus whenever $(A_i^T-I)\lambda_i=b_i$; then
$d\Theta_i=dt\wedge b_i$ is precisely exterior minus cap curvature.
For smooth representatives replace $t\eta_i$ by $H(t)\eta_i$,
where $H'=h$ and $H(0)=0,H(1)=1$, retaining the term $-t\beta_i$.
Since $\beta_i\wedge E_0=0$, this change leaves every cube integral below
unchanged. The endpoint comparison is then smooth as well. One can take
\[
\begin{array}{c|c|c}
i& b_i&\lambda_i\\\hline
1&-4G-W+\tfrac23K&(-U+W)/3\\
2&-3G+U+\tfrac12K&(-U+W)/2\\
3&0&0
\end{array}
\]
On the last boundary the based inverse word gives three successive
transported segments
\[
 -C_2^TC_1^TA_2^T\eta_1=-2G+U+K,\quad
 -C_2^TC_1^T\eta_2=3G-U-K,\quad
 -C_2^T\eta_3=0.
\]
Their sum is $G=(C_2^T-I)W$, so start with $\lambda_4=W$ and add the
successive segment integrals. The only $W$ coefficient of $\Theta_4$
remains $1$, while the two nonzero $U$ increments are $+1,-1$.
The last-boundary contribution below therefore cancels exactly;
it is unchanged if the two nonzero segments are reversed.

\paragraph{Multiplicative collar calculation.}
Across a collar with coordinate $s$ from cap to exterior use the closed
form
\[
 \Omega_i=E_{\rm cap}+d(s\Theta_i)
          =E_{\rm cap}+s\,d\Theta_i+ds\wedge\Theta_i.
\]
At its endpoints it equals the chosen cap and exterior representatives.
Both $E_{\rm cap}^2$ and $(d\Theta_i)^2$ vanish. Hence
\[
 \Omega_i^3=6s\,ds\wedge\Theta_i\wedge E_0\wedge dt\wedge b_i.
\]
Integrating $s$ gives the factor $3$. With the algebraic boundary
orientation which records the coefficient of
$\Theta_i\wedge E_0\wedge b_i$ in $G\wedge U\wedge W\wedge K$,
the companion code uses fibre orientation $G\wedge U\wedge W\wedge K$
and obtains
\[
 3[\lambda_1E_0b_1]=2,\qquad
 3[\lambda_2E_0b_2]=-3,\qquad
 3[\Theta_3E_0b_3]=3[\Theta_4E_0b_4]=0.
\]
Thus the sum has absolute value $1$. For completeness, the common
overall orientation sign can be left unassigned: moving $dt$ past $\Theta_iE_0$ contributes a minus sign if
the collar orientation is $ds\wedge dt\wedge G\wedge U\wedge W\wedge K$.
The raw algebraic sum in the code is $-1$.

\paragraph{Explicit change to the inverse mapping-torus convention.}
The calculation above uses $(x,1)\sim(Ax,0)$. The convention
$(x,1)\sim(A^{-1}x,0)$ is obtained by reversing the peripheral
parameter and changing fibre coordinates by $A$ so that the marked fibre
at the basepoint is unchanged. Explicitly,
$[x,t]\mapsto[Ax,1-t]$ is the comparison and is the identity on the marked
$t=0$ fibre. The inverse-convention cap map
$[y,\theta]\mapsto[y+\theta v/m]$ becomes
$[Ax+(1-t)v/m]=[x-tv/m]$ after applying the inverse affine deck
transformation. Thus it is exactly the negative cap term used above.
The ribbon order and all fourth-boundary segments are reversed together.
The data transform as
\[
 A'=A^{-1},\quad \eta'=-A^{-T}\eta,\quad
 \beta'=-A^{-T}\beta=-\beta,\quad\lambda'=\lambda,
 \qquad b'=-A^{-T}b=(A^{-T}-I)\lambda.
\]
These are simultaneous changes of the actual forms, not merely changes
of the displayed matrices. Since $E_0$ is invariant and $\det A=1$,
\[
 [\lambda E_0b']=-[(A^T\lambda)E_0b]
                =-[(\lambda+b)E_0b]=-[\lambda E_0b].
\]
For the multi-segment last boundary use the full reversed primitive
$\Theta'(t)=A^{-T}\Theta(1-t)$; its initial value is again $\lambda$,
and the whole collar integral changes sign. Thus reversing the common
boundary convention multiplies all four contributions by the same sign.
The absolute value $1$ is independent of this choice.

Finally, $H^2(Q;\mathbb Q)$ is one-dimensional and restriction sends its
primitive generator to $2GK$, so the global closed form constructed
above represents the actual Euler class up to its independently fixed
sign. This proves the asserted absolute cube. The independently oriented
value $\ell=-4$ and the spin congruence $\ell\equiv4k\pmod{24}$ select
$k=-1$ from the two possibilities $\{1,-1\}$. The conclusion about
$\Theta_7$ is now \ref{maps:src:B.6.19}.\hfill$\square$



\end{theorem}




\begin{theorem}[the original smooth monodromy is zero]\label{maps:src:B.6.21}
\StatementRecord{maps:src:B.6.21}{proved}{GPT-6 Astra (OpenAI, ChatGPT Pro mode), effort Pro, 2026-11-17; Claude Fable 5.2 (Anthropic), effort high, 2026-11-17; Claude Fable 5.2 (Anthropic), effort high, 2026-11-16; GPT-6 Astra (OpenAI, ChatGPT Pro mode), S6\_06\_BLIND, effort Pro, 2026-11-17; Claude Fable 5.2 (Anthropic), effort high, 16 November 2026}{the hypotheses, notation and definitions stated here; \cref{push:src:C.3.4}; \cref{maps:src:B.6.16}; \cref{maps:src:B.6.20}; \cref{maps:src:B.6.17}; \cref{maps:src:B.6.15}; \cref{maps:src:B.6.19}; \cref{push:src:C.1.1}; \cref{push:src:C.1.2}; \cref{push:src:C.2.1}}


For the original $(3,4)$ family, with its complex orientation,
the positive $w$ loop and the generator of $\Theta_7$ fixed by
$\mu=1/28$, the monodromy satisfies
\[
                         [\Phi]=0\in\mathbb Z/28.
 \tag{B.6.38}\label{maps:eq:B.6.38}
\]
It is therefore smoothly isotopic to the identity.  Its sixfold
class is zero, and among the two resulting \cref{push:src:C.3.4} residues $0,14$
the original class is $0$.  This statement does not yet assert
the same value for the other multiplicity pairs.
\end{theorem}


\begin{proof}
\begin{enumerate}[label=\arabic*.]
\item \ref{maps:src:B.6.16} constructs the free-circle replacement with the original
peripheral identifications and closing map.  \ref{maps:src:B.6.20} gives
$|k|=1$ by an integral-normalized multiplicative calculation on
those actual finite and nodal caps.  Independently \ref{maps:src:B.6.17} gives
$\ell=-4$.  The spin congruence $\ell\equiv4k\pmod {24}$
in \ref{maps:src:B.6.15} then forces $k=-1$.

\item The circle disc bundle consequently has invariant zero by
\ref{maps:src:B.6.15}.  The actual parametrized spin surgery trace has
$p_1$ coefficient $-4-2(-1)=-2$, and its rank-one middle form
is unimodular by \ref{maps:src:B.6.19}.  Thus its relative Pontryagin square
is four times its signature, and its normalized spin
cobordism correction is zero.  It follows that
\[
 \mu(M_\Phi)=\mu(M_s)=0.
\]
\cref{push:src:C.1.1} and \cref{push:src:C.1.2} and \cref{push:src:C.2.1} compares the oriented mapping torus with the glued
homotopy sphere and identifies this invariant with
$[\Phi]/28$; this proves \eqref{maps:eq:B.6.38}.  No framed local bubble,
choice of a scalar interpolation, or unknown reverse-surgery
correction is used in this conclusion.  The integral sign
of the intermediate cup calculation is fixed by the
independent spin congruence, so its preliminary collar
orientation convention cannot change the answer.

\end{enumerate}
\end{proof}




\begin{proposition}[the actual higher-multiplicity replacement]\label{maps:src:B.6.22}
\StatementRecord{maps:src:B.6.22}{proved}{GPT-6 Astra (OpenAI, ChatGPT Pro mode), effort Pro, 2026-11-17; Claude Fable 5.2 (Anthropic), effort high, 2026-11-16}{the hypotheses, notation and definitions stated here; \cref{maps:src:B.6.16}; \cref{maps:src:B.6.19}; \cref{maps:src:B.6.8}; \cref{maps:eq:B.6.29}; \cref{maps:eq:B.6.30}; \cref{maps:src:D.2.14}; \cref{maps:src:B.6.17}}


Let $m=3k,n=4l$, with $k,l>0$ and $\gcd(m,n)=1$, and use
any of the admissible fixed-data sphere families of Paper~I, Proposition~23.1 and Paper~I, Theorem~23.2.  Its
positive-generator return, with period $h=2kl$, has an actual
smooth fixed-circle replacement as in \ref{maps:src:B.6.16}.  With
\[
 A=lN_0,\qquad B=kN_1,\qquad P=4A+3B=\pm1,
\]
its four marked affine defects can be taken to be
\[
 e_1=(0,0,-2A,0),\quad e_2=(0,-B,0,0),\quad
 t_1=2A a_1,\quad t_2=-B a_2,
 \qquad a_1=w,\ a_2=w-d.
 \tag{B.6.39}\label{maps:eq:B.6.39}
\]
Its linear return is $S_h=I-hd\otimes g^*$.  The free-circle
quotient mapping torus has the oriented invariant
\[
                           \ell=-2h=-4kl.
 \tag{B.6.40}\label{maps:eq:B.6.40}
\]
The spin trace and its equivariant characteristic calculation
in \ref{maps:src:B.6.19} apply to this return as well.
\end{proposition}


\begin{proof}
\begin{enumerate}[label=\arabic*.]
\item The invariant power vectors in the question's marking are,
by the checked dictionary Paper~I, Proposition~23.1 and Paper~I, Theorem~23.2,
\[
 v_1=N_0\epsilon-\rho_0d,\qquad
 v_2=N_1\epsilon'-(\rho_1+N_1)d,
\quad \epsilon=g+2u-4w,\quad\epsilon'=g+3u-3w.
\]
The sphere condition is $4lN_0+3kN_1=\pm1$.  In particular
$N_0$ is coprime to $3k$ and $N_1$ is coprime to $4l$.
Applying the invariant covector $g^*$ to a possible finite
stabilizer equation therefore proves that the $d$ circle
is free in both logarithmic caps.  Its only fixed locus
is the same cusp $S^2$, with normal bundle and weights
$\mathcal O(-1)_{+1}\oplus\mathcal O(-1)_{-1}$.

\item For the fixed-data shift $h=2kl$, the explicit repeated-norm
solutions in Paper~I, Theorem~23.2, give the two defects $e_1,e_2$
in \eqref{maps:eq:B.6.39}, as already checked in \ref{maps:src:B.6.8}.  Their combined
cusp defect is
\[
 e_0=-M_0(e_1+A_1e_2)=(2A-B)w+Bd.
\]
Make the following integral change of lifted parameter
generator, leaving the underlying parameter circle fixed:
\[
 \varkappa' =\varkappa+r,
 \qquad r=(-B,\ 2A-B,\ 2B-2A,\ 0).
 \tag{B.6.41}\label{maps:eq:B.6.41}
\]
Since $(M_0-I)r=e_0$, the new cusp defect is zero.
This is a change of lift by an actual integral fibre
period, not a fractional-period interpolation.  In the
toric cusp cover the Paper~I, Proposition~23.1 and Paper~I, Theorem~23.2 map is a holomorphic torus unit
times an integral toric power.  Composing its lift with
this integral long-period deck transformation removes
that power.  Hence the same original quotient closing
map now has a holomorphic unit representative on the
entire cusp.  Both its nonzero limit and its parameter
dependence are retained.

\item Use this representative in exactly the product nodal-chain
construction \eqref{maps:eq:B.6.29}--\eqref{maps:eq:B.6.30}, with the same period
matrix perturbation independent of $c$.  The unit extends
on the two node charts and preserves the actual smoothing
parameter.  The cutoff on the outer regular annulus
matches it to the unchanged original closing map.
The finite maps of Paper~I, Proposition~23.1 and Paper~I, Theorem~23.2 are also unchanged: their new
defects merely record the same change of lifted
parameter generator.  Explicitly they are
\[
 d_1=e_1+(I-A_1)r=(0,P,-P,-2A-B),\qquad
 d_2=e_2+(I-A_2)r=(0,0,-P,2A+B),
 \tag{B.6.42}\label{maps:eq:B.6.42}
\]
and the two nodal defects are zero in this marking.
Returning to the old lift changes these last two
defects to
\[
 (C_1-I)r=2Aa_1,\qquad (C_2-I)r=-Ba_2,
\]
which proves \eqref{maps:eq:B.6.39} with the actual marked caps.
This also checks the peripheral product relation,
because $C_1$ fixes $a_2$ and $t_1+C_1t_2=e_0$.

\item Each original fibre is $S^6$ by the stipulated sphere
theorem.  The same nullsphere surgery and spin-lift
argument as in \ref{maps:src:D.2.14} gives a free circle action on
$S^3\times S^3$ with quotient $S^2\times S^3$.
Thus the geometric generator and torsion-free groups
used in \ref{maps:src:B.6.17} are unchanged.  Its determinant proof
uses the same invariant residue form on the new Hopf
sphere.  Every reduced multiple-fibre divisor still
has zero class, since its positive multiple is the
pullback of a base point and $H^2$ of the quotient
is torsion free.  Consequently its descended Koszul
calculation gives the identical formula
$p_1(Q)=-2\operatorname{PD}(Z_1+Z_2)$.

\item At the first node the surviving elliptic return is now
$b_2\mapsto b_2+h a_2$ modulo the vanishing $a_1$.
At the second it is the identity on $(a_1,b_1)$.
The preceding integral changes of lift only change
affine translations on these tori; they do not change
these linear returns.  The oriented unit-line clutch
calculation of \ref{maps:src:B.6.17} therefore gives Euler pairings
$h,0$, proving \eqref{maps:eq:B.6.40}.

\item Finally the replacement uses the same fixed quadratic
germ and ruling-preserving closing maps, so its
parametrized spin handle is the one in \ref{maps:src:B.6.19}.
Its ordinary cohomology and its fixed-set Borel Thom
calculation depend only on that handle, not on the
finite multiplicities.  Thus the trace coefficient
is again $\ell-2k_{\rm cup}$ and its middle form has
absolute determinant $|k_{\rm cup}|$.

\end{enumerate}
\end{proof}




\begin{theorem}\label{maps:src:B.6.23}
\StatementRecord{maps:src:B.6.23}{proved}{GPT-6 Astra (OpenAI, ChatGPT Pro mode), effort Pro, 2026-11-17; Claude Fable 5.2 (Anthropic), effort high, 16 November 2026}{Definitions and the displayed constructions}


For every higher fixed-data family with $m=3k,n=4l$,
$\gcd(m,n)=1$, and $4lN_0+3kN_1=\pm1$, the oriented cup/shear
invariant of the fixed-curve-surgery mapping torus is
\[
                         k_{\rm cup}=-kl.
\]
The independently oriented tangent invariant of \ref{maps:src:B.6.22} is
$\ell=-4kl$.

\emph{Proof.} Write $A=lN_0$, $B=kN_1$, $P=4A+3B\in\{1,-1\}$,
and $h=2kl$. Distinguish the parameter generator $\varkappa$ from the
positive integer $k$. The ordinary finite matrices are the same $A_1,A_2$
as in \ref{maps:src:B.6.20}. The actual four affine defects from \ref{maps:src:B.6.22} are
\[
 e_1=(0,0,-2A,0),\quad e_2=(0,-B,0,0),\quad
 t_1=2A a_1,\quad t_2=-B a_2,
 \qquad a_1=w,\ a_2=w-\delta.
\]
Use the integral change of parameter generator
\[
 \varkappa'=\varkappa+r,\qquad
 r=(-B,\,2A-B,\,2B-2A,\,0)\in\mathbb Z\langle g,u,w,\delta\rangle.
\]
This has determinant one. A defect transforms by
$d_i=e_i+(I-M_i)r$, giving exactly
\[
 d_1=(0,P,-P,-2A-B),\quad d_2=(0,0,-P,2A+B),\qquad
 d_3=d_4=0.
\]
The statement concerns the actual marked cap models from \ref{maps:src:B.6.22};
it does not infer a filling merely from these matrices.
After quotienting by $\delta$, the two finite matrices are those in
\ref{maps:src:B.6.20} with their last columns multiplied by $P$. The nodal matrices
are unchanged, and their full ordered product is the identity.

Let $G,U,W,K$ be the new dual fibre forms. Use the same common
curvature-sign convention as in \ref{maps:src:B.6.20}, so the Euler representative
has restriction $E_0=hGK$. Its possible overall sign relative to
the oriented principal-circle Euler class is independent of $h$:
the regular $g$-parameter torus has the same unit-line clutching,
with exponent $-h$, for every positive $h$. The actual $\delta$ rows, with the same exterior gauge $-U$
as in \ref{maps:src:B.6.20}, give
\[
 \eta_1=-4G-W+(P-2A-B)K,\qquad
 \eta_2=-3G+U+(2A+B)K,\qquad\eta_3=0.
\]
The power vectors are
$v_1=N_0\epsilon-\rho_0\delta$ and
$v_2=N_1\epsilon'-(\rho_1+N_1)\delta$, with
$\epsilon=g+2u-4w$, $\epsilon'=g+3u-3w$.
Thus the geometric finite cap pullbacks have mixed terms
\[
 \beta_1=-\frac{hN_0}{3k}K=-\frac{2A}{3}K,
 \qquad
 \beta_2=-\frac{hN_1}{4l}K=-\frac B2K.
\]
The central coefficients $\rho_i$ do not occur in these quotient forms.
Using $P=4A+3B$, their differences simplify to
\[
 b_1=\eta_1-\beta_1=-4G-W+\frac{2P}{3}K,
 \qquad
 b_2=\eta_2-\beta_2=-3G+U+\frac P2K.
\]
Exactly the same boundary primitives work:
\[
 \lambda_1=(-U+W)/3,\qquad\lambda_2=(-U+W)/2.
\]
The last boundary has transported segments
\[
 -2G+U+(2A+B)K,\qquad
 3G-U-(2A+B)K,\qquad0.
\]
Their sum is again $G=(C_2^T-I)W$, and the opposite $U$ increments
cancel its cube contribution. All smooth seam and inverse-peripheral
comparisons are exactly those proved in \ref{maps:src:B.6.20}. Consequently the raw
four collar contributions are
\[
                         h,\quad-\frac{3h}{2},\quad0,\quad0,
 \qquad \sum=-\frac h2=-kl.
\]
This is a symbolic identity, verified without any multiplicity bound.
The program is \texttt{manifolds\_mapping\_general.py} in the code directory.

For $P=1$, the parameter shear has determinant one and the oriented base,
ordinary torus fibre, free circle and positive parameter retain the same
conventions as in \ref{maps:src:B.6.20}. Every collar expression is a positive $kl$
multiple of the minimal expression. The calibration $k_{\rm cup}=-1$
from \ref{maps:src:B.6.20} therefore gives $k_{\rm cup}=-kl$ for every such datum.
For $P=-1$, global elliptic fibre inversion sends every $v_i$ to $-v_i$
and changes $\rho_i$ to $-\rho_i$, reducing to $P=1$ while preserving
the complex orientation of the original total space and the positive
parameter. It reverses the free-circle orientation and the quotient
orientation together. Accordingly $e\mapsto-e$ and $[Q]\mapsto-[Q]$,
so both $\langle e^3,[Q]\rangle$ and
$\langle ep_1(Q),[Q]\rangle$ remain unchanged. This proves the stated
signed formula also for $P=-1$.\hfill$\square$



\end{theorem}




\begin{theorem}[all requested monodromy classes vanish]\label{maps:src:B.6.24}
\StatementRecord{maps:src:B.6.24}{proved}{GPT-6 Astra (OpenAI, ChatGPT Pro mode), effort Pro, 2026-11-17; GPT-6 Astra (OpenAI, ChatGPT Pro mode), S6\_06\_BLIND, effort Pro, 2026-11-17; Claude Fable 5.2 (Anthropic), effort high, 16 November 2026}{the hypotheses, notation and definitions stated here; \cref{maps:src:B.6.22}; \cref{maps:src:B.6.23}; \cref{maps:src:B.6.19}; \cref{push:src:C.1.1}; \cref{push:src:C.1.2}; \cref{push:src:C.2.1}; \cref{maps:src:B.6.15}; \cref{maps:eq:B.6.44}}


For every admissible coprime sphere pair $(3k,4l)$ and every
fixed twist datum of Paper~I, Proposition~23.1 and Paper~I, Theorem~23.2, the monodromy for the positive
generator $2kl$ of its parameter-period lattice satisfies
\[
                    [\Phi_{3k,4l}]=0\in\Theta_7.
 \tag{B.6.43}\label{maps:eq:B.6.43}
\]
The identifications and orientation are those specified in
the question.  In particular this completes the numerical
part of B.6 for both the original and higher families.
\end{theorem}


\begin{proof}
\begin{enumerate}[label=\arabic*.]
\item Write $d=kl>0$.  \ref{maps:src:B.6.22}--\ref{maps:src:B.6.23} compute the actual
free-circle quotient invariants as
\[
              k_{\rm cup}=-d,\qquad \ell=-4d=4k_{\rm cup}.
\]
The trace has $p_1$ coefficient $2k_{\rm cup}$, by
\ref{maps:src:B.6.19} and its extension in \ref{maps:src:B.6.22}.  We determine the
sign of its rank-one intersection form where that sign
could affect the answer.
The two choices of spin structure on the trace restrict
to the two choices on $M_s$: restriction on $H^1(-;\mathbb Z/2)$
is an isomorphism, generated by the parameter circle.
Choose the one agreeing with the spin structure induced
by the disc-bundle spin-c filling.  The induced choice
on $M_\Phi$ does not change its invariant, by \cref{push:src:C.1.1} and \cref{push:src:C.1.2} and \cref{push:src:C.2.1}.

\item Put $k_0=k_{\rm cup}$ temporarily.  The disc bundle
$D(L_x)$ has intersection form $(k_0)$ and boundary
$M_s$, by \ref{maps:src:B.6.15}.  Its primitive absolute degree-four
generator $y$ restricts to the torsion class $\pi^*y$
of order $|k_0|$ on $M_s$.  The trace has form $(j)$
with $j=\pm k_0$, and its primitive absolute generator
$\alpha$ restricts to the \emph{same} $\pi^*y$,
by the Borel restriction proof in \ref{maps:src:B.6.19}.  Its other
boundary component $M_\Phi$ has no degree-four
cohomology and contributes no torsion linking term.

\item The boundary torsion linking pairing evaluated on this
fixed generator is the signed reciprocal of the
rank-one intersection coefficient, with one common
boundary-orientation sign for the two fillings.
One can see this directly by lifting the generator
rationally through the relative intersection map:
the relative generator divided by its intersection
coefficient is that lift, and evaluation modulo
$\mathbb Z$ gives the reciprocal.  The outgoing
orientation of $M_s$ is the same for the disc bundle
and the trace.  Hence $j=-k_0$ would imply
\[
                         \frac{2}{k_0}=0\quad\hbox{in }\mathbb Q/\mathbb Z.
\]
For $|k_0|>2$ this is impossible, so $j=k_0$.

\item In that case the normalized spin trace correction is
\[
 \mu(M_s)-\mu(M_\Phi)
 =\frac{(2k_0)^2/k_0-4\operatorname{sign}(k_0)}{896}
 =\frac{k_0-\operatorname{sign}(k_0)}{224}.
 \tag{B.6.44}\label{maps:eq:B.6.44}
\]
The free-circle formula \ref{maps:src:B.6.15}, with $\ell=4k_0$, gives
exactly the same value for $\mu(M_s)$.  Their difference
is therefore $\mu(M_\Phi)=0$.

\item For $|k_0|=1$ both the free invariant and the trace
correction vanish for either possible sign, as already
proved in \ref{maps:src:B.6.19}.  For $|k_0|=2$, suppose instead that
$j=-k_0$.  The trace correction would be the negative
of \eqref{maps:eq:B.6.44}, and subtraction from the free invariant
would give
\[
 \mu(M_\Phi)=\frac{k_0-\operatorname{sign}(k_0)}{112}
                         =\pm\frac1{112}\pmod{\mathbb Z}.
\]
\cref{push:src:C.1.1} and \cref{push:src:C.1.2} and \cref{push:src:C.2.1} requires $\mu(M_\Phi)$ to lie in
$(1/28)\mathbb Z/\mathbb Z$, so this is impossible.
The trace sign is thus $j=k_0$ also in this case,
and the same cancellation proves zero.

\item All positive integers $d$ have now been covered.  \cref{push:src:C.1.1} and \cref{push:src:C.1.2} and \cref{push:src:C.2.1}
identifies the last invariant with the specified smooth
homotopy-sphere class divided by $28$, proving
\eqref{maps:eq:B.6.43}.  The proof used the actual lifted parameter
generator and all four cap maps, and does not equate
the different multiplicity families before computing
their invariants.  The symbolic characteristic identities
and the exceptional $|k_0|=2$ denominator check are
recorded in \texttt{maps\_trace\_linking.py}.

\end{enumerate}
\end{proof}




\begin{corollary}[an explicit spin filling after the computation]\label{maps:src:B.6.25}
\StatementRecord{maps:src:B.6.25}{proved}{GPT-6 Astra (OpenAI, ChatGPT Pro mode), effort Pro, 2026-11-17; GPT-6 Astra (OpenAI, ChatGPT Pro mode), S6\_06\_BLIND, effort Pro, 2026-11-17; Claude Fable 5.2 (Anthropic), effort high, 16 November 2026}{the hypotheses, notation and definitions stated here; \cref{maps:src:B.6.24}}


Every mapping torus in \ref{maps:src:B.6.24} has a spin eight-manifold
filling diffeomorphic to $D^7\times S^1$, with one
zero-handle and one one-handle, and with
$p_1^2=\operatorname{sign}=0$.
\end{corollary}


\begin{proof}
\begin{enumerate}[label=\arabic*.]
\item The computed zero mapping class supplies a smooth
isotopy $h_t$ from the identity to $\Phi$, chosen
constant near its endpoints.  For the convention
$(z,1)\sim(\Phi z,0)$, the map
\[
 [z,t]\longmapsto(h_t(z),e^{2\pi it})
\]
is a well-defined oriented diffeomorphism from
$M_\Phi$ to $S^6\times S^1$.  Give $D^7\times S^1$
its product orientation; its boundary orientation is
the stipulated positive-parameter orientation of
$S^6\times S^1$.  This product is parallelizable and
has zero middle cohomology, proving the two
characteristic-number assertions.  The usual product
handle presentation has one handle in each of
indices zero and one.  Either spin structure on its
boundary extends across the product.  This filling
is a consequence of the preceding actual cap and
trace computation; its existence was not assumed
to establish the zero monodromy class.

\end{enumerate}
\end{proof}


\section{The marked sources and finite caps}

\par\noindent\textbf{Current physical-line convention and correction.}
P3K.68 supplies the actual Hermitian dual, both cap maps and the
first-to-second Hopf-gluing arrows. The physical nonzero-vector line
is \(M\); the outward-meridian-one line is \(U=M^{-1}\).
The actual oriented sphere class is \(A_M=-A_U\pmod{28}\).
P3K.69 proves the corrected reference and source-fan consequences.
Earlier false physical-sign assertions are explicitly labelled
Refuted and retained as historical claims with precise correction
notices. Their source constructions and the separately identified
local Cox, topology and derivative maps retain their stated scope.
The full sign ledger is \path{records/physical_line_correction_scope.json}.
\par\medskip

% Generated from the current-run complete proofs by code/part3/render_part3_tex.py.
\providecommand{\tightlist}{\setlength{\itemsep}{0pt}\setlength{\parskip}{0pt}}
\noindent The following proofs retain their explicit construction maps and scope. Uncompleted auxiliary comparisons are identified at their exact remaining step.

\begin{theorem}\label{night:P3F.1}
\StatementRecord{night:P3F.1}{proved}{S6\_06\_BLIND: GPT-6 Astra (OpenAI, ChatGPT Pro mode, effort Pro), recorded re-derivation in the checkpoint dated 18 November 2026}{Definitions and the displayed constructions}
The torsion-free rank-one Mordell--Weil rows are exactly eight, and their Kodaira refinements give exactly twenty-one realized configurations.
\end{theorem}
\begin{proof}
\begin{enumerate}[label=\arabic*.]
\item The rank-one, torsion-free Mordell--Weil cases in the Oguiso--Shioda Main Theorem are precisely the following. Here \(P\) denotes a generator, \(h(P)=\langle P,P\rangle\), and \(I_1^r\) denotes \(r\) distinct nodal fibres. Every configuration in the last column occurs over the complex numbers with Mordell--Weil group isomorphic to \(\mathbb Z\).

\begin{longtable}[]{@{}
  >{\raggedright\arraybackslash}p{(\columnwidth - 6\tabcolsep) * \real{0.1000}}
  >{\raggedright\arraybackslash}p{(\columnwidth - 6\tabcolsep) * \real{0.2100}}
  >{\raggedright\arraybackslash}p{(\columnwidth - 6\tabcolsep) * \real{0.1000}}
  >{\raggedright\arraybackslash}p{(\columnwidth - 6\tabcolsep) * \real{0.5900}}@{}}
\toprule\noalign{}
\begin{minipage}[b]{\linewidth}\raggedright
OS91 row
\end{minipage} & \begin{minipage}[b]{\linewidth}\raggedright
Root lattice
\end{minipage} & \begin{minipage}[b]{\linewidth}\raggedright
\(h(P)\)
\end{minipage} & \begin{minipage}[b]{\linewidth}\raggedright
Configurations
\end{minipage} \\
\midrule\noalign{}
\endhead
\bottomrule\noalign{}
\endlastfoot
43 & \(E_7\) & \(1/2\) & \(III^*+I_1^3\); \(III^*+II+I_1\) \\
45 & \(A_7\) & \(1/8\) & \(I_8+I_1^4\); \(I_8+II+I_1^2\) \\
46 & \(D_7\) & \(1/4\) & \(I_3^*+I_1^3\); \(I_3^*+II+I_1\) \\
47 & \(A_6+A_1\) & \(1/14\) & \(I_7+I_2+I_1^3\); \(I_7+I_2+II+I_1\); \(I_7+III+I_1^2\) \\
49 & \(E_6+A_1\) & \(1/6\) & \(IV^*+I_2+I_1^2\); \(IV^*+I_2+II\); \(IV^*+III+I_1\) \\
50 & \(D_5+A_2\) & \(1/12\) & \(I_1^*+I_3+I_1^2\); \(I_1^*+I_3+II\); \(I_1^*+IV+I_1\) \\
55 & \(A_4+A_3\) & \(1/20\) & \(I_5+I_4+I_1^3\); \(I_5+I_4+II+I_1\) \\
56 & \(A_4+A_2+A_1\) & \(1/30\) & \(I_5+I_3+I_2+I_1^2\); \(I_5+I_3+I_2+II\); \(I_5+I_3+III+I_1\); \(I_5+IV+I_2+I_1\) \\
\end{longtable}

\item The primary table was read in the scanned original, printed pp.85--86. The program \path{code/part3/enumerate_mw_one.py} transcribes every rank-one row, 43--61, including all eleven rows with torsion. It then selects torsion order one. Thus the eight-entry count is obtained by a reproducible selection from the whole relevant table.

\item We now justify the refinement from root lattices to Kodaira fibres. The Kodaira/root correspondences are \(A_1=I_2\) or \(III\), \(A_2=I_3\) or \(IV\), \(A_n=I_{n+1}\) for \(n\geq3\), \(D_n=I^*_{n-4}\), \(E_6=IV^*\), and \(E_7=III^*\). Any further fibre is irreducible, hence \(I_1\) or \(II\). Their Euler numbers must sum to 12. The program takes the Cartesian product of these choices, then enumerates every nonnegative solution \(r+2s=12-\sum e(F_{\rm red})\). It obtains 24 candidates.

\item For a candidate let \(d=\deg j=\sum n(i_n+i_n^*)\), and use \(j/1728\), so that the elliptic values are \(0,1,\infty\). Put

\[
 a_0=ii+iv^*,\qquad b_0=iv,
 \qquad a_1=iii+iii^*.
\]

\item The minimal ramification profiles over 0 and 1 are respectively

\[
 1^{a_0}2^{b_0}3^{(d-a_0-2b_0)/3},\qquad
 1^{a_1}2^{(d-a_1)/2}.
\]

\item Over infinity they are the multiplicative fibre indices, including the indices of \(I_n^*\). Hence the excess ramification

\[
 x=2d-2-\sum_{u\in\{0,1,\infty\}}\sum_{e\in\rho_u}(e-1)
\]

is nonnegative. This is Miranda90 \(1.11\), and follows directly from Riemann--Hurwitz: additional ramification can only increase the displayed sum. Precisely three of the 24 candidates have \(x=-1\):

\[
 I_8+II+II,\qquad I_7+III+II,\qquad I_5+IV+III.
\]

\item They cannot occur. The other 21 have \(x=0\) or 1. For each of them the program constructs four explicit permutations

\[
 \sigma_0\sigma_1\sigma_*\sigma_\infty=1
\]

on \(d\) letters. The first, second, and fourth have the three profiles above. The third is the identity for \(x=0\), a transposition for \(x=1\). The generated group is transitive, and the program checks \(\sum_\nu(d-\#\mathrm{cycles}(\sigma_\nu))=2d-2\). The permutations are printed in full in \texttt{enumerate\_mw\_one.out} and stored in \texttt{mw\_one\_configurations.json}.

\item To explain exactly what these certificates establish, take positively oriented meridians about \(0,1,z_*,\infty\) on the Riemann sphere, with their product relation in the displayed order. The permutations define a connected finite unramified cover of the complement. At a cycle of length \(e\), attach a disc by the chart \(z\mapsto z^e\). This gives a compact connected Riemann surface and a holomorphic branched map. The verified Euler equation gives genus zero, so the source is a Riemann sphere. This is a \(j\)-map with the specified profiles. Miranda90 Proposition3.4 constructs the corresponding elliptic surface by pullback of its displayed universal \(j=t\) surface, followed by the specified quadratic twists and minimalizations. Its Euler number is 12 and it has a section, and the proposition proves it is rational. Thus each certificate is an existence certificate, not merely a necessary numerical test.

\item The root system \(A_7\) requires an extra check: OS91 rows44 and45 have the same root system but different torsion. Here are explicit torsion-free models for both configurations. Write

\[
 R=t^2+s,\qquad A=R^2+bt,\qquad
 B=R^3+\frac32btR+\frac38b^2,
 \qquad y^2=x^3-3Ax+2B.
\]

\item For \((s,b)=(1,1)\),

\[
 4(-3A)^3+27(2B)^2
 =\frac{27}{16}(48t^4+8t^3+96t^2+72t+57).
\]

\item The program verifies that this quartic is square-free and coprime to \(A,B\), so there are four \(I_1\)'s. At infinity the minimal Weierstrass coefficients are units and the discriminant has order8, giving \(I_8\). For \((s,b)=(3,48)\), the same discriminant expression is

\[
 186624(t+3)^2(3t^2-10t+51).
\]

\item Here \(A,B\) each have a simple zero at \(-3\); hence that fibre is \(II\). The quadratic factor is square-free, coprime to \(t+3\), and gives two \(I_1\)'s. Again the infinity fibre is \(I_8\).

\item Neither model has a 2-torsion point. After \(x=X+R\), its cubic is

\[
 X^3+3(t^2+s)X^2-3btX+\frac34b^2.
\]

\item A root in \(\mathbb C(t)\) would be a polynomial because the cubic is monic over the integrally closed ring \(\mathbb C[t]\). Its degree is at most two. For degree at most one, comparing the coefficient of \(t^4\) first eliminates the linear term, then the coefficient of \(t^2\) eliminates the constant term. The nonzero constant \(3b^2/4\) remains. For a quadratic root, its leading coefficient must be \(-3\); the next two coefficients force \(X=-3(t^2+s)\), which leaves \(9bt(t^2+s)+3b^2/4\ne0\). Thus no root exists. OS91's only two rank-one (A\_7) cases have torsion 0 and \(\mathbb Z/2\); both displayed models therefore belong to row45. The remaining seven root systems each have a unique torsion-free rank-one row, so no such ambiguity occurs.

\item This proves the exact count \(2+2+2+3+3+3+2+4=21\).

\end{enumerate}
\end{proof}




\begin{theorem}\label{night:P3F.2}
\StatementRecord{night:P3F.2}{proved}{S6\_06\_BLIND: GPT-6 Astra (OpenAI, ChatGPT Pro mode, effort Pro), recorded re-derivation in the checkpoint dated 18 November 2026}{\cref{night:P3F.1}, \cref{night:P3F.31}}
The section-intersection formula has the positive local-correction term, and the primitive sections in the eight rows have intersection zero with the zero section.
\end{theorem}
\begin{proof}
\begin{enumerate}[label=\arabic*.]
\item The primary source is Shioda90 Theorem8.6, equation\(8.12\), with local contributions in table\(8.16\), printed pp.228--229. Proposition8.12, equation\(8.25\), printed p.231, gives multiplication of heights by the degree of a field extension. The scanned statements and their proofs were read; OCR is retained only as a discovery aid.

\item For a section \(R\ne O\) of a relatively minimal elliptic surface with section,

\[
 h(R)=2\chi(\mathcal O)+2(R\cdot O)
       -\sum_v\operatorname{contr}_v(R).
\]

\item Consequently the exact rearrangement is

\[
 (R\cdot O)=\frac{h(R)-2\chi(\mathcal O)
                    +\sum_v\operatorname{contr}_v(R)}2. \tag{P3F.2.1}
\]

\item The local contributions in Shioda's convention are nonnegative. The general formula stated in the prompt with a minus sign before these contributions therefore has the wrong sign in this convention. This is an algebraic discrepancy independent of the unfinished identification of the expression with a monodromy invariant.

\item For a primitive generator of each of the eight rational-surface cases, the local contributions and component orders are as follows. Either choice of a pair of components exchanged by inversion is allowed.

\begin{longtable}[]{@{}
  >{\raggedright\arraybackslash}p{(\columnwidth - 6\tabcolsep) * \real{0.2500}}
  >{\raggedright\arraybackslash}p{(\columnwidth - 6\tabcolsep) * \real{0.2500}}
  >{\raggedright\arraybackslash}p{(\columnwidth - 6\tabcolsep) * \real{0.2500}}
  >{\raggedright\arraybackslash}p{(\columnwidth - 6\tabcolsep) * \real{0.2500}}@{}}
\toprule\noalign{}
\begin{minipage}[b]{\linewidth}\raggedright
Root system
\end{minipage} & \begin{minipage}[b]{\linewidth}\raggedright
Component data for \(P\)
\end{minipage} & \begin{minipage}[b]{\linewidth}\raggedright
Sum of contributions
\end{minipage} & \begin{minipage}[b]{\linewidth}\raggedright
Orders
\end{minipage} \\
\midrule\noalign{}
\endhead
\bottomrule\noalign{}
\endlastfoot
\(E_7\) & nonidentity component & \(3/2\) & 2 \\
\(A_7\) & component3 or5 of (I\_8) & \(15/8\) & 8 \\
\(D_7\) & a spinor component & \(7/4\) & 4 \\
\(A_6+A_1\) & component2 or5 of (I\_7), nonidentity (A\_1) & \(10/7+1/2\) & 7,2 \\
\(E_6+A_1\) & nonidentity in both & \(4/3+1/2\) & 3,2 \\
\(D_5+A_2\) & spinor, nonidentity & \(5/4+2/3\) & 4,3 \\
\(A_4+A_3\) & component2 or3 of (I\_5), component1 or3 of (I\_4) & \(6/5+3/4\) & 5,4 \\
\(A_4+A_2+A_1\) & component1 or4 of (I\_5), nonidentity (A\_2,A\_1) & \(4/5+2/3+1/2\) & 5,3,2 \\
\end{longtable}

\item Here \(A_1\) and \(A_2\) have the same nonzero local contribution for either of their two Kodaira realizations. To derive, rather than assume, the table, use (h(P)) from P3F.1 and enumerate all multiplicity-one components in Shioda's local table. The sum of the largest permitted contributions is always smaller than (4-h(P)). Since \(P\cdot O\geq0\) is an integer, the height formula excludes \(P\cdot O\geq1\). It follows that \(P\cdot O=0\). The equation \(\sum\operatorname{contr}_v(P)=2-h(P)\) then gives exactly the component choices displayed. Thus for every one of the 21 configurations the primitive generator is disjoint from the zero section.

\item For the OS49 configuration \(IV^*+II+I_2\), place \(IV^*\) at0 and \(I_2\) at infinity, and pull back by the degree3 map \(t=z^3\). The relatively minimal pullback has three \(II\) fibres and one \(I_6\), so \(\chi=1\). The (IV\^{}*) fibre has good reduction on the cubic root chart. The generator has height \(3/6=1/2\). Its component1 at \(I_2\) pulls back to component3 at \(I_6\), contributing \(3\cdot3/6=3/2\). Equation\(P3F.2.1\) gives \(P'\cdot O'=0\). Substituting the negative sign from the prompt instead would give \(-3/2\), not the intersection number. This calculation concerns the explicit compact elliptic pullback. The actual regular-binding page and its boundary comparison are constructed in P3F.31; its map has three separate I2 fibres.

\end{enumerate}
\end{proof}




\begin{theorem}\label{night:P3F.3}
\StatementRecord{night:P3F.3}{proved}{S6\_06\_BLIND: GPT-6 Astra (OpenAI, ChatGPT Pro mode, effort Pro), recorded re-derivation in the checkpoint dated 18 November 2026}{Definitions and the displayed constructions}
The exact Torelli presentation is $I(S^2\times S^3)=\Z a\oplus\Z b$, with $(k,\ell)(a)=(1,4)$ and $(k,\ell)(b)=(8,8)$.
\end{theorem}
\begin{proof}
\begin{enumerate}[label=\arabic*.]
\item The primary source is Huize Jin, arXiv:2608.26698v2 (12September2026), Theorem1.2, Corollary1.3, Proposition3.1, Lemmas4.4 and4.8, and Theorem4.9. The downloaded PDF is retained in \path{records/primary_sources}. The ambient manifold is the oriented product \(M=S^2\times S^3\). Its orientation-preserving Torelli group is

\[
 I(M)=\mathbb Za\oplus\mathbb Zb.
\]

\item Here \(a=t_{S_1,\alpha}\) twists a standard \(S^2\times D^3\), and \(b=t_{S_1\#S_1,\alpha}\) twists the connected sum of two oriented copies of that core \(S^2\). The normalization of \(\alpha\in\pi_3SO(3)\) is that the rank3 bundle over (S\^{}4) clutched by \(\alpha\) has first Pontryagin number4. The twist is explicitly

\[
 (x,y)\longmapsto(\alpha(y)x,y),
\]

where \(\alpha:D^3\to SO(3)\) equals the identity near its boundary. Extending by the identity gives the same formula with \(\alpha:S^3\to SO(3)\). In particular the core sphere has dimension2.

\item For \(f\in I(M)\), Jin uses

\[
 M_f=M\times[0,1]/((x,0)\sim(f(x),1)),
\]

with the preferred spin structure specified by a fixed disc and a class \(x\in H^2(M_f;\mathbb Z)\) restricting to the positive generator on \(S^2\). Put

\[
 k(f)=\langle x^3,[M_f]\rangle,\qquad
 \ell(f)=\langle p_1(TM_f)x,[M_f]\rangle,
 \qquad A(f)=\frac{4k(f)-\ell(f)}{24}.
\]

\item The theorem identifies \(f\mapsto(A(f),k(f))\) with an isomorphism \(I(M)\to\mathbb Z^2\). Its normalization gives

\[
 (k(a),\ell(a))=(1,4),\qquad
 (k(b),\ell(b))=(8,8),
\]

so \(b-8a\) and \(a\) map to the standard ordered basis of \(\mathbb Z^2\). In particular

\[
 [f]=(k(f)-8A(f))a+A(f)b. \tag{P3F.3.1}
\]

\item These assertions follow from the cited theorem by direct substitution, and verify that the congruence is \(\ell\equiv4k\pmod{24}\). The pair ((-1,-4)) gives ({[}f{]}=-a). It is a nonzero element of \(I(S^2\times S^3)\); its effect on the associated glued 7-manifold is computed explicitly next.

\end{enumerate}
\end{proof}




\begin{theorem}\label{night:P3F.4}
\StatementRecord{night:P3F.4}{proved}{S6\_06\_BLIND: GPT-6 Astra (OpenAI, ChatGPT Pro mode, effort Pro), recorded re-derivation in the checkpoint dated 18 November 2026}{\cref{night:P3F.3}, \cref{night:P3F.8}, \cref{night:P3K.6}}
The Hopf lift of every $a$-power extends over $D^4\times S^3$; its specified sphere gluing is standard.
\end{theorem}
\begin{proof}
\begin{enumerate}[label=\arabic*.]
\item Every Hopf lift of \(a^n\), for \(n\in\mathbb Z\), gives the standard oriented 7-sphere by the following gluing construction. This specifies the morphisms required to pass from these mapping classes to homotopy spheres.

\item Use unit quaternions \(S^3\subset\mathbb H\) and \(S^2\subset\operatorname{Im}\mathbb H\), with Hopf map \(h(u)=uiu^{-1}\). The principal circle acts on the right on (u). On

\[
 p:S^3\times S^3\longrightarrow S^2\times S^3,
 \qquad p(u,y)=(h(u),y),
\]

fix its Chern class to be the chosen positive generator; reversing the orientation of the Hopf fibre changes the sign and does not change the extension argument. Lift \(\alpha:S^3\to SO(3)\) to \(\widetilde\alpha:S^3\to SU(2)\), using the double cover given by quaternionic conjugation. The lift exists because \(S^3\) is simply connected. The map

\[
 A_n:S^3\times S^3\to S^3\times S^3,
 \qquad (u,y)\longmapsto(\widetilde\alpha(y)^n u,y)
\]

is a smooth principal-circle-bundle automorphism covering \(a^n\). The equality \(h(vu)=v h(u)v^{-1}\) verifies this statement, and left quaternionic multiplication commutes with the right Hopf action. Crucially it is the boundary restriction of the explicit diffeomorphism

\[
 E_n:D^4\times S^3\longrightarrow D^4\times S^3,
 \qquad (q,y)\longmapsto(\widetilde\alpha(y)^n q,y). \tag{P3F.4.1}
\]

\item Its inverse uses the inverse unit quaternion, and it preserves orientation.

\item Fix the orientations in the decomposition

\[
 S^7=\partial(D^4\times D^4)
     =(D^4\times S^3)\cup_{\mathrm{id}}(S^3\times D^4)
\]

by the outward-normal-first convention and rounding the corner. Gluing the first piece to the second by \(A_n\) gives a manifold carried to this standard boundary by \(E_n\) on the first piece and the identity on the second. This is an explicit orientation-preserving diffeomorphism.

\item Any two principal-circle lifts of the same base map differ by a gauge map \(g:S^2\times S^3\to S^1\). Because \(H^1(S^2\times S^3;\mathbb Z)=0\), it lifts to a smooth real function and is isotopic to the constant map through gauge transformations. Base isotopies lift to bundle isotopies, by the homotopy lifting property for bundle isomorphisms (or by horizontal lifting for a chosen connection). Boundary-isotopic gluing maps give diffeomorphic glued manifolds: extend the isotopy over a collar. Thus the conclusion depends only on the base mapping class, and holds for every lift and every representative of \(a^n\).

\item For completeness, the same Hopf-lift construction assigns a homotopy 7-sphere to any \(f\in I(S^2\times S^3)\). Its lift acts on \(H_3(S^3\times S^3)\) by a matrix

\[
 \begin{pmatrix}1&c\\0&1\end{pmatrix}
\]

for an integer \(c\): the lower row follows by projection to the base \(S^3\), and the first diagonal entry follows from the positive Hopf bundle map over the positive \(S^2\) generator. In Mayer--Vietoris the map to the two piecewise \(H_3\)'s is \(\left(\begin{smallmatrix}0&1\\1&c\end{smallmatrix}\right)\), of determinant \(-1\). All intermediate reduced homology vanishes; van Kampen gives simple connectivity. Hence the glued manifold is a homotopy sphere. No assertion that this assignment is a group homomorphism is needed for the explicit computation on \(\mathbb Za\).

\item Combining this construction with \(P3F.3.1\) computes the Hopf-glued sphere for the numerical pair ((-1,-4)): its oriented class is zero. P3F.8 supplies the actual geometric identification with the family, and P3K.6--9 derive the characteristic numbers from its charts and its specified descended rotating framing. The classification theorem is used only after these geometric comparisons.

\end{enumerate}
\end{proof}




\begin{theorem}\label{night:P3F.5}
\StatementRecord{night:P3F.5}{proved}{S6\_06\_BLIND: GPT-6 Astra (OpenAI, ChatGPT Pro mode, effort Pro), recorded re-derivation in the checkpoint dated 18 November 2026}{\cref{night:P3F.25}, \cref{night:P3F.27}, \cref{night:P3F.31}, \cref{night:P3F.7}, \cref{night:P3K.6}, \cref{night:P3K.65}, \cref{night:P3K.8}, \cref{night:P3K.9}}
The Haefliger formula applies to an actual Seifert surface. The compact elliptic page has characteristic class $K=5F$, so the determinant correction must be retained.
\end{theorem}
\begin{proof}
\begin{enumerate}[label=\arabic*.]
\item Takase's primary correction paper, \emph{The Hopf invariant of a Haefliger knot}, \nolinkurl{arXiv:math/0506356}, Theorem2.3 and Corollary2.6, gives the formula

\[
 \Omega(F)=-\frac{\sigma(V)-e(\nu V)^2}{8} \tag{P3F.5.1}
\]

for an actual Seifert embedding \(V^4\hookrightarrow S^6\) with \(\partial V=S^3\). The orientation conventions are stated on p.2 of that paper. The square is computed using the boundary identification and the cap, not an unspecified square on an arbitrary open surface.

\item The algebraic expression in the prompt would match this formula after the explicit comparison \(e(\nu V)^2=4s_\delta\) and the compatible relative Pontryagin/trivialization comparison. The following test shows that the compact elliptic pullback cannot simply be punctured and used with this normal Euler class. The degree12 pullback over the compact genus3 curve has \(\chi(\mathcal O)=1\), no multiple fibres, and

\[
 K=(\chi+2g-2)F=5F,\qquad w_2=F\bmod2.
\]

\item Its section intersects \(F\) once, so this class is nonzero, including after removing a small ball disjoint from the section. An oriented embedding of this punctured surface in \(S^6\) would satisfy \(w_2(TV)=w_2(\nu V)=e(\nu V)\bmod2\), by the Whitney formula and the spin structure on (S\^{}6). The particular normal Euler class \(2c_1(M)\) has zero reduction. Thus that specific puncturing with that specific normal bundle cannot furnish the required Seifert embedding. The prompt itself requires only its square: the characteristic class \(K+2c_1(M)\) has the same square \(4c_1(M)^2\), because \(K^2=K\cdot M=0\). Therefore this parity test does not refute the square required in the prompt.

\item This is a proved obstruction to that comparison. It does not rule out the additional removals, handle attachments, and compatible framings that an actual open-book page may possess. Those modifications have to be constructed and included when computing its characteristic numbers.

\subsection{State of the family-monodromy application}\label{state-of-the-family-monodromy-application}

\item The 21 configurations above are now completely enumerated and realized. The statement in the supplied CSV about a bounded slice of 50 inputs does not define its bounds or which section pairs, linearizations, and local fillings constitute an input. The number50 is therefore not used as a computed result here. The further supplied rows assert additional larger searches, but provide no proof data and do not define the missing bounded slice.

\item The initial main-case gap has been resolved by the actual orbit and Hopf comparison P3F.7--13, the real trace P3K.6--7, the independent analytic index calculation P3K.8, and the rotating-frame computation P3K.9. These prove zero positive-loop class for the original family and give \((k,\ell)=(-1,-4)\) in the stated descended framing. P3F.25 applies the same actual trace to the further verified families and computes the nonzero weighted OS47 example. P3F.31--32 construct the actual positive-diagonal OS49 page and prove its zero return index. With the full marking of P3F.27, P3K.65 proves that the negative-diagonal filling has class \(-3k\) for the actual positive \(r\)-loop \(\lambda=r^{3k}\), and class \(-k\) for the shorter \(w\)-loop \(\lambda=w^k\); the positive fan has class zero on both parameters. The exact earlier status passages are preserved, with their superseding proof identifiers, in \path{records/part3_superseded_status.md}.

\end{enumerate}
\end{proof}




\begin{theorem}\label{night:P3F.6}
\StatementRecord{night:P3F.6}{proved}{S6\_06\_BLIND: GPT-6 Astra (OpenAI, ChatGPT Pro mode, effort Pro), recorded re-derivation in the checkpoint dated 18 November 2026}{\cref{night:P3F.5}, \cref{night:P3K.6}}
The two explicit degree-twelve maps of the same genus-three curve give respectively one $I_{12}$ fibre and twelve $I_1$ fibres. Both give $\chi(\mathcal O)=1$, signature $-8$, height $2$, and $(mP)\cdot O=m^2-1$.
\end{theorem}
\begin{proof}
\begin{enumerate}[label=\arabic*.]
\item Let C be the smooth compactification of b\textsuperscript{3}=a\textsuperscript{4}+1. The function x=b\textsuperscript{3} defines a degree12 map to P1 with ramification profiles

\begin{longtable}[]{@{}ll@{}}
\toprule\noalign{}
x-value & profile \\
\midrule\noalign{}
\endhead
\bottomrule\noalign{}
\endlastfoot
0 & 3,3,3,3 \\
1 & 4,4,4 \\
infinity & 12 \\
\end{longtable}

\item Indeed near b=0 the four roots of a\textsuperscript{4}=-1 are distinct, so b is a local parameter and x=b\textsuperscript{3}. Near a=0 there are three points and x-1=a\textsuperscript{4}. The coprimality of 3 and4 gives one point at infinity, with poles of orders3 and4 for a,b respectively, and hence pole order12 for x. The total ramification is8+9+11=28, and Riemann--Hurwitz gives 2g-2=-24+28=4.

\item For the E26 base with IV\emph{,III,I1 over0,1,infinity, the pullback by t=x has a single I12 fibre after minimalization, by local semistable base change of the nodal model. There is a different map on the very same curve that gives twelve I1 fibres. Choose a regular finite value t} different from0,1 and set

\begin{lstlisting}
t=mu(x)=t* x/(x+t*-1).
\end{lstlisting}

\item This sends the three branch values0,1,infinity to0,1,t\emph{. Its inverse image of the source cusp infinity is the regular x-value1-t}, and so has12 distinct points, all unramified. The pullback has twelve I1 fibres. Its order3 and order4 branches still remove IV* and III, while the order12 branch lies over a smooth elliptic fibre. The curve alone therefore does not specify which of these fibrations is being used. The Milnor page of a regular \(3,4\) Seifert fibre uses a regular binding value t* and leads to the second map. The actual four-dimensional page, its boundary surgeries and its real trace are constructed in P3K.6--7.

\item In either compact minimal pullback, the canonical formula gives K=5F, chi\(O\)=1, and sigma=-8. The height pulls back to2. All multiples of P meet the identity component in the I12 case, because P already meets the smooth locus of the original I1 fibre; there are no reducible fibres in the twelve-I1 case. Hence \(mP\).O=m\textsuperscript{2}-1 in both cases.

\item The parity test in P3F.5 only excludes the normal Euler class2M for a ball-puncture of this nonspin compact surface. It does not exclude the square required in the prompt: e=K+2M is characteristic, and e\textsuperscript{2}=4M\textsuperscript{2} because K\textsuperscript{2}=K.M=0. No conclusion about the actual relative normal class follows until its cap and framing have been identified.

\end{enumerate}
\end{proof}




\begin{theorem}\label{night:P3F.7}
\StatementRecord{night:P3F.7}{proved}{S6\_06\_BLIND: GPT-6 Astra (OpenAI, ChatGPT Pro mode, effort Pro), recorded re-derivation in the checkpoint dated 18 November 2026}{Definitions and the displayed constructions}
The actual circle orbit pair is $(S^5,S^2)$, with fixed normal bundle $\mathcal O(-1)\oplus\mathcal O(-1)$ of weights $(-1,+1)$. The specified conifold replacement is free with quotient $S^2\times S^3$.
\end{theorem}
\begin{proof}
\begin{enumerate}[label=\arabic*.]
\item For the AC26 sphere member, its delta circle is semifree, has one fixed curve F=P1, and its orbit pair is diffeomorphic to the standard pair \(S5,S2\). Its fixed normal bundle is O\(-1\) plus O\(-1\) with respective weights\(-1,+1\). Replacing that fixed neighbourhood by the conifold smoothing gives a free circle manifold S3 times S3, with quotient S2 times S3.

\item On an ordinary fibre the delta circle is primitive and acts freely. At a finite cap, a nontrivial deck power has nonintegral gamma character, since gamma\(v_j\) is prime to m\_j. It therefore cannot agree with a delta circle rotation. This proves freeness there. In the cusp fan the directions are ±e1,±e2,±\(e1-e2\). A finite circle stabilizer at an edge of direction d has order dividing det\(d,e2\), which is±1 off the e2 edges; on an e2 edge the entire circle fixes it. At the two-dimensional torus orbit corresponding to a vertex of the height-one triangulation there is no stabilizer. The zero-dimensional torus orbits corresponding to triangles are fixed, and each lies on one of the fixed edges. The period lattice acts transitively on the e2 edges, with no stabilizer, so their quotient is one fixed P1.

\item For the vertical wall use the four rays

\begin{lstlisting}
r0=(0,0,1), r1=(0,1,1), r+=(1,0,1), r-=(-1,1,1).
\end{lstlisting}

\item They satisfy r++r-=r0+r1. The two adjacent unimodular charts have coordinates\(x0,x1,z\) and(x0',x1',z'), with exact transition

\begin{lstlisting}
x0'=z x0,  x1'=z x1,  z'=1/z,
q=x0 x1 z=x0' x1' z'.
\end{lstlisting}

\item The fixed curve is x0=x1=0. The transitions give O\(-1\)\textsuperscript{2}, and the cocharacter e2 has weights\(-1,+1,0\). These identities are independently checked from the ray matrices in code/part3/orbit\_mv.py. The map

\begin{lstlisting}
A=x0, B=z x0, C=x1, D=z x1
\end{lstlisting}

contracts the curve to AD=BC. It intertwines the circle with weights \(-1,-1,+1,+1\) on\(A,B,C,D\). On AD-BC=epsilon nonzero the action is free. For epsilon=1, polar decomposition of SL2\(C\) as U H gives the explicit equivariant diffeomorphism SU2 times R3, with the circle acting on U by left multiplication and acting trivially on H. A compact Milnor neighbourhood is thus S3 times D3 and its quotient is S2 times D3.

\item For completeness, the orbit is identified integrally. Let A=X minus F and B=A/S1. General position gives pi1\(A\)=0. Alexander duality gives H*\(A;Z\)=Z in degrees0,3 and zero otherwise. The oriented normal link S3 is a primitive generator of H3\(A\). Its quotient is the normal link S2, and the circle bundle on these links is the Hopf bundle. In the Gysin sequence, e generates H2\(B\), and integration H3\(A\)-\textgreater H2\(B\) is multiplication by1: evaluate on this linked S2 and its Hopf preimage. Every other positive-degree group of B vanishes. Also pi1\(B\)=0 by the homotopy sequence of the principal circle bundle.

\item Conjugating the negative-weight line converts the fixed normal bundle to a positive-weight complex bundle O\(1\) plus O\(-1\). Its clutching loop diag\(z,z^-1\) in SU2 bounds the explicit disc with SU2 coordinates \(a,b\)=(r exp(i theta),sqrt(1-r\textsuperscript{2})); joining its centre to the identity gives a nullhomotopy. Thus it is a trivial unitary rank2 bundle. Its orbit normal bundle is a trivial oriented rank3 bundle. With this trivialization the boundary is F times S2, and the Euler class is just the meridian generator, with no F term.

\item Now Y=X/S1 is the union of B (truncated at that boundary) and F times D3. The maps on H0 and H2 in Mayer--Vietoris, with domain H2 ordered (F,meridian) and target(complement,tube), are respectively

\begin{lstlisting}
[1;-1],       [[0,1],[-1,0]].
\end{lstlisting}

\item The latter has determinant1. H4(boundary)=Z maps to zero and gives H5\(Y\)=Z; every other positive intermediate group is zero. All matrices and their Smith forms are printed in orbit\_mv.out. Van Kampen gives pi1\(Y\)=0. Thus Y is a homotopy5-sphere. Kervaire--Milnor1963, the table on p505 and the Smale comparison stated immediately above it, gives Theta5=0, so Y is diffeomorphic to S5. Haefliger1966 Corollary6.6 says C\_n\^{}q=0 for n\textless2q-3; with n=2,q=3, and Theorem1.2 converting concordance to isotopy, every embedding S2 in S5 is unknotted. This identifies the orbit pair.

\item One may now construct the equivariant identification with the standard circle action on the sphere boundary of D3 times D4. On the free complement the principal circle bundle over S2 times D3 has Euler generator and is the product Hopf bundle S3 times D3. On the fixed tube it is the trivial positive-weight rank2 bundle just constructed. Their two boundary identifications differ by a gauge map from S2 times S2 to S1. Such a map has a global logarithm because H1(S2 times S2;Z)=0, so it extends over a collar and removes the discrepancy. This gives an equivariant diffeomorphism. It is also the specialization of Levine's classification as stated in Harvey--Kerin--Shankar2020 Theorem4.1, whose hypotheses are checked here: smooth codimension4 fixed set and H2\(F;Z\) without2-torsion. Replacing the fixed tube S2 times D4 by the free piece D3 times S3 doubles the free complement, giving S3 timesS3; the quotient doubles S2 timesD3, giving S2 timesS3. The normal framing is unique up to homotopy because pi2\(SO4\)=0. This is the asserted conifold surgery.

\end{enumerate}
\end{proof}




\begin{theorem}\label{night:P3F.8}
\StatementRecord{night:P3F.8}{proved}{S6\_06\_BLIND: GPT-6 Astra (OpenAI, ChatGPT Pro mode, effort Pro), recorded re-derivation in the checkpoint dated 18 November 2026}{\cref{night:P2.1}, \cref{night:P3F.4}, \cref{night:P3F.7}, \cref{night:P3K.6}, \cref{night:P3K.7}, \cref{night:P3K.9}}
The actual equivariant family return has a typed Hopf-lift image in the Torelli group, and its clutched sphere is the displayed Hopf gluing.
\end{theorem}
\begin{proof}
\begin{enumerate}[label=\arabic*.]
\item Let Phi be the actual positive parameter return map of the sphere family, with the complex orientation, represented by an equivariant Ehresmann trivialization. There is a Torelli diffeomorphism f of S2 timesS3, supported in S2 timesD3, for which the oriented sphere clutched by Phi is exactly the Hopf-lift gluing lambda\(f\) of P3F.4. This statement identifies the morphism; it does not assume values of the two integer invariants of f.

\item The parameter family carries the delta circle action. Its fixed curves form a P1 bundle around the parameter circle; the two normal weight lines have degree-1. The fixed-curve bundle is smoothly trivial because PGL2\(C\) is connected. Each weight line over P1 timesS1 is isomorphic as a smooth complex line bundle to the pullback of O\(-1\), since H2(P1 timesS1;Z)=Z. Thus its equivariant normal disc bundle can be trivialized around the parameter circle. Choose an Ehresmann connection agreeing with that trivialization near the fixed set; averaging over the compact circle makes it equivariant. Its return map is the identity near the fixed tube.

\item Under P3F.7 write the sphere as

\begin{lstlisting}
(S2_y timesD4_u) union (D3_y timesS3_u),
\end{lstlisting}

with the circle the Hopf action on the u coordinate. Conjugating the negative normal line reverses its real two-plane orientation; choose the orientation of the y factor so that this identification preserves the complex orientation of X. The restriction of Phi to the free part is an equivariant map tilde f of D3\_y timesS3\_u equal to the identity near its boundary. Its quotient is a map of D3\_y timesS2\_u relative boundary. Extend it by the identity over the other hemisphere of S3\_y. This defines f:S3\_y timesS2\_u -\textgreater{} S3\_y timesS2\_u. It fixes the H2 generator, since that generator is represented in the identity hemisphere. It preserves orientation, so Poincare duality makes it fix H3 as well. Consequently f is in the Torelli group.

\item Its Hopf lift A on S3\_y timesS3\_u equals tilde f on the first hemisphere and the identity on the second. Any other lift differs by a gauge map S2 timesS3 -\textgreater{} S1. Its H1 is zero, so that map has a global logarithm and yields an isotopy of lifts.

\item Here is the cut-and-reglue comparison. In the standard sphere boundary(D4\_u timesD4\_y), put

\begin{lstlisting}
H=D4_u timesS3_y,  K=S3_u timesD4_y,
N=S3_u timesD3_y in boundary H=boundary K.
\end{lstlisting}

\item Our gluing convention identifies x in boundary H with A\(x\) in boundary K, after the standard product-boundary identification. The ball B7=D4\_u timesD3\_y has boundary equal to N union (D4\_u timesS2\_y). Its complement in the rounded product boundary is the other standard7-ball: this follows by taking the two hemispheres of the y boundary and rounding their product collars. Regluing H to K by A changes the identification only on N. Equivalently cut out B7 and reglue its boundary by the map equal to tilde f on N and the identity on D4\_u timesS2\_y. That boundary map is precisely Phi in the displayed decomposition of S6. The identity on the common collar makes the two quotient constructions coincide pointwise outside N, and their collar identifications coincide on N. Hence

\begin{lstlisting}
D7 union_Phi D7  =  (D4_u timesS3_y) union_A (S3_u timesD4_y)
\end{lstlisting}

as oriented smooth manifolds, with outward boundary orientation on the first piece and the complex sphere orientation used to choose the initial identification. This is lambda\(f\).

\item The actual sixfold descent of P2.1 is

\begin{lstlisting}
H(t)delta, H(t)=-(4l1 log t+3l2 log(t-1))/(2pi i),
\end{lstlisting}

with cusp coefficient -p after carrying the collar integer l0. P3K.6--7 constructs its real trace with the complete line clutch. P3K.9 evaluates its characteristic pair as\(-6,-24\), including the normal-frame comparison, and descends it to\(-1,-4\) for one parameter loop. The forward geometric mapping-torus orientation and Jin's reverse convention are both included in P3K.7.

\end{enumerate}
\end{proof}




\begin{theorem}\label{night:P3F.9}
\StatementRecord{night:P3F.9}{proved}{S6\_06\_BLIND: GPT-6 Astra (OpenAI, ChatGPT Pro mode, effort Pro), recorded re-derivation in the checkpoint dated 18 November 2026}{Definitions and the displayed constructions}
The resolved cusp orbit has a relative product comparison with a nodal elliptic neighbourhood times the marked gamma circle, retaining the stable outer framing.
\end{theorem}
\begin{proof}
\begin{enumerate}[label=\arabic*.]
\item The local calculation giving the product is as follows. Start with the ordinary conifold xy=uv, with the circle acting on u,v with weights+1,-1 and fixing x,y. A small resolution has charts

\begin{lstlisting}
(x,v,z),   u=xz, y=vz;
(u,y,w),   x=uw, v=yw;      w=1/z.
\end{lstlisting}

\item The orbit charts, with the Hopf smoothing defined using these normal coordinates, are

\begin{lstlisting}
(x,y,b), b=(|z|^2-|v|^2)/2;
(x,y,b'), b'=(|u|^2-|w|^2)/2.
\end{lstlisting}

\item On the overlap r=\textbar z\textbar\textsuperscript{2}\textgreater0 they satisfy

\begin{lstlisting}
2b=r-|y|^2/r,       2b'=|x|^2r-1/r,
db'/db=(|x|^2+r^(-2))/(1+|y|^2r^(-2))>0.       (P3F.9.1)
\end{lstlisting}

\item Thus projection to\(x,y\) is a smooth submersion with oriented one-dimensional fibres. Every fibre is a line: away from\(x,y\)=\(0,0\) the two charts join its negative and positive ends in the displayed increasing order; over\(0,0\) the exceptional interval joins the two half-lines. A smooth section can be chosen in the free part. For epsilon\textgreater0 choose the unique r\textgreater0 solving

\begin{lstlisting}
|x|^2r-|y|^2/r+2epsilon(r-1)/(r+1)=0.
\end{lstlisting}

\item The left side increases strictly from a negative to a positive value; its derivative is \textbar x\textbar\textsuperscript{2}+\textbar y\textbar\textsuperscript{2}/r\textsuperscript{2}+4epsilon/\(r+1\)\textsuperscript{2}\textgreater0. Consequently r is a smooth positive function of\(x,y\), and the representatives z=sqrt\(r\),v=y/sqrt\(r\),u=x sqrt\(r\) give the section in the orbit. A positive vertical vector field, patched by a partition of unity and rescaled to be complete, identifies each line with R by flow from this section. It therefore gives a product diffeomorphism over C\textsuperscript{2}.

\item This comparison can be made relative to a prescribed product identification outside the node neighbourhood. The quotient smoothing uses a choice of invariant metric in the fixed normal bundle; choose it to agree with the original product normal metric outside that neighbourhood. There projection to\(x,y\) is the original smooth submersion. On overlapping local line trivialisations the transitions are increasing by\(P3F.9.1\). Taking a partition of unity on the base and a convex combination of increasing, proper fibre coordinates produces an increasing proper coordinate agreeing with the prescribed one where the outer partition has weight1. This is a relative product trivialization. Its derivative over the fixed \(x,y\) base is an upper triangular real matrix with identity base block and a positive fibre derivative. Contracting the off-diagonal block and then the positive scalar gives a homotopy to the identity relative to the prescribed outer product. Thus this comparison retains the stable boundary framing, not just the orbit diffeomorphism type. In particular one must not equate all quadratic invariants in different normal metrics at a fixed point: their transition can involve a radial square root. The metric choice and the positive-coordinate construction avoid that error.

\item Apply this calculation to the square Mumford fan before triangulation. The square model is the fibre product of the two nodal elliptic models; the A2 triangulation is its small resolution, one per lattice cell. The circle acts in the second nodal factor. Its orbit is a smooth circle bundle over the base disc, since its local quotient at the node is\(q=uv,(|u|\textsuperscript{2}-|v|\textsuperscript{2})/2\). It is a product over the disc. Before resolving the conifold, the total orbit is therefore E\_I1 timesS1, where E\_I1 is a compact neighbourhood of one nodal elliptic fibre in a smooth elliptic surface. The relative comparison above replaces the node orbit without changing this product diffeomorphism type or its outer boundary marking.

\item The bounded toric units of the actual construction do not change this statement. Choose their logarithms on the cusp disc. Their interpolation to constant units gives the same unimodular charts and a proper smooth family after shrinking the disc uniformly; the logarithmic terms dominate every bounded unit. Averaging an Ehresmann field gives an equivariant comparison. In the source elliptic factor, the gamma deck translates by the section6P. Its Tate coordinate is a unit a\(q\), and the explicit path Z-\textgreater exp(t log a(q)) Z extends across the nodal fibre. Its suspension is consequently a product E\_I1 timesS1\_gamma. The source Picard line restricted to an elliptic neighbourhood selected disjoint from the fixed sphere is topologically trivial: its degree on the nodal fibre is zero, and that neighbourhood retracts onto the nodal fibre. An extending Hermitian metric removes its radial unit on this slice. This statement does not trivialize the principal-circle line on the complement of the fixed sphere in the entire orbit neighbourhood. That line has Chern number one, up to the fixed orientation convention, on a meridian two-sphere; its meridian class is retained in every subsequent clutch or index comparison. These comparisons preserve the integral gamma,u,w and q markings on the outer boundary. Changes in the affine stable frames consist of unit scalings and triangular shears and are homotopic to the indicated diagonal factors; no three-dimensional Pontryagin framing class is introduced by a logarithm of a unit.

\item The fixed sphere projects in E\_I1 to the normalization of the nodal fibre. The small-resolution interval separates its two preimages of the node in the gamma direction. Its projected normal complex line has transition a'=z\textsuperscript{2}a, namely O\(-2\). Its Euler number, with the complex orientation of that line and of the normalization, is-2. The gamma coordinate of this sphere can be confined to an arbitrarily short arc by taking the relative node comparison sufficiently small. Rotation of the S1\_gamma factor is a free action on the cusp orbit body and extends the marked boundary gamma circle action.

\item This proves the claimed product comparison with the explicit boundary marking, rather than identifying the singular central fibre itself with a product. The map to the elliptic surface changes the transverse leaves in the node neighbourhood, as its displayed small-resolution charts show.

\end{enumerate}
\end{proof}




\begin{theorem}\label{night:P3F.10}
\StatementRecord{night:P3F.10}{proved}{S6\_06\_BLIND: GPT-6 Astra (OpenAI, ChatGPT Pro mode, effort Pro), recorded re-derivation in the checkpoint dated 18 November 2026}{\cref{night:P3F.6}, \cref{night:P3F.9}, \cref{night:P3K.6}, \cref{night:P3K.9}}
The actual regular-binding page has twelve nodal fibres and boundary $T^3$. Its circle-line square is $s_\delta=-2$ and its relative Pontryagin pairing is $p_{\rm rel}=-24$.
\end{theorem}
\begin{proof}
\begin{enumerate}[label=\arabic*.]
\item Here the data are the original sphere twists\(0,1,-1\). On the ordinary and finite-cap orbit pieces, use the invariant gamma character in smooth flat coordinates. The resulting three-dimensional quotient is formed from S1\_gamma times the punctured base and the two solid tori

\begin{lstlisting}
(S1_gamma timesDelta_s)/(g,s)->(g+l_j/m_j,zeta_j s).
\end{lstlisting}

\item The two meridian relations are3x=gamma and4y=-gamma; after the cusp solid torus the third isx+y=0. The relevant integer matrix is {[}{[}3,0,-1{]},{[}0,4,1{]},{[}1,1,0{]}{]}, of determinant-1. The union of the two Seifert solid tori is explicitly the standard weighted circle fibration on S3 with weights3,4. Its meridian/longitude basis change is unimodular, and the two stabilizer orders are3,4. This supplies the Seifert identification, including the labelled exceptional fibres.

\item To compare the smooth flat coordinates with the holomorphic logarithmic gluing, subtract the angular translation from\(log s\)/\(2pi i\)Pi v\_j. The remaining radial translation is single valued on the annulus and has a real scalar homotopy to zero. Translation by this homotopy on a collar gives the required equivariant identification. At the cusp use P3F.9, relative to the outer gamma marking. This constructs a smooth map of the actual orbit Y=S5 to the weighted S3. Its fibres are the source elliptic curves; the critical curve is the nodal critical point in E\_I1 timesS1\_gamma. Near that curve its formula is the ordinary positive node map(x,y,gamma)-\textgreater(xy,gamma). Thus it is an actual one-parameter family of positive genus-one Lefschetz maps. The cusp modification is included in the construction of this map.

\item Choose as binding a regular weighted fibre at a value t* different from0,1,infinity. Its Milnor open-book page is explicitly the affine curve b\textsuperscript{3}-a\textsuperscript{4}=1, embedded in S3 by

\begin{lstlisting}
(a,b)->(rho^3 a,rho^4 b),
rho^6|a|^2+rho^8|b|^2=1, rho>0.
\end{lstlisting}

\item The inverse divides by the positive twelfth root of b\textsuperscript{3}-a\textsuperscript{4} on the chosen argument page. Its projection to the weighted base is the degree12 map of P3F.6, followed by the unique Möbius map sending the binding value to t*. After adding the boundary circle the page is a genus3 surface C° with one boundary component. The regular cusp fibre in the weighted S3 intersects this page in12 points. Their signs are positive in the complex Milnor orientation, equivalently the map to the base is holomorphic and unramified there. Its inverse image V0 in Y is therefore a genus-one Lefschetz fibration over C° with twelve positive nodal fibres. The boundary is the torus mapping torus of a translation, hence T3. It need not be S3.

\item The page can be chosen disjoint from the fixed sphere: by P3F.9 its gamma image lies in an arbitrarily short arc, and the phase of the Milnor page can be chosen so that its twelve intersections with the critical gamma circle avoid that arc. The delta circle bundle therefore restricts to an honest principal circle bundle over V0.

\item We justify the comparison with its sectioned Jacobian without requiring a section of a compact torsor. The affine elliptic bundle over C° is locally the stated Jacobian, including the irreducible nodal fibres. Its torsor obstruction lies in H\textsuperscript{2}(C°,j\_\emph{Lambda\_E). This group is zero: cut the surface along its handle generators and disjoint paths to the nodal values. In the punctured surface fundamental group, every nodal meridian is an independent free generator because the boundary word is not imposed as a relation. The restriction map from cocycles to the direct sum of their cusp coinvariants is consequently surjective. The exact sequence for j\_}Lambda\_E gives H\textsuperscript{2}=0. Equivalently this is the integer Cech obstruction to matching the local origins. The sheaf of smooth Lie-algebra sections is fine, so this vanishing constructs a smooth global section and an affine identification with the Jacobian.

\item The circle line has the fibrewise Picard class of P, not6P. This follows both from the source definition M=O\(P-O\) and from its two elliptic period multipliers: the u column contains mu while the gamma column contains6mu. The latter records translation in the suspension. Under an affine identification of a torus torsor with its Jacobian, a degree zero line has the same Picard point. The quotient of two lines with the same normalized fibrewise Picard point is pulled back from the base: normalize a trivializing holomorphic fibre section at the chosen origin, and its overlap ratios are constant on the compact fibres. This argument remains valid at each irreducible node, where the multiplicative gluing unit is the one supplied by M and a local ball has no line-bundle obstruction. Thus on V0 its Chern class is that of the Jacobian line O(P'-O'), up to a base pullback. Since the base has boundary, that pullback is topologically trivial.

\item On the boundary regular torus the degree-zero line has unitary flat fibre multipliers and the translating point has a flat lift. Thus its restriction to T3 has zero Chern class. A relative lift of its Chern class exists. If u and u+delta g are two lifts, with g in H1(boundary V0;Z), the absolute class of delta g is zero and restriction of the absolute class of u to the boundary is zero. The relative cup-product boundary formula gives

\begin{lstlisting}
u cup delta g=delta(i*u cup g)=0,
(delta g)^2=0.
\end{lstlisting}

\item The relative square is therefore independent of every boundary gauge. Choose the affine boundary identification furnished by the section above and cap the Jacobian by its regular fibre timesD2. This gives the compact elliptic pullback under the regular-binding cover of P3F.6, with twelve I1 fibres. Its P'-O' line extends the actual page line with some boundary gauge; gauge independence now makes

\begin{lstlisting}
s_delta=(P'-O')^2=-2chi-2(P'.O')=-2.
\end{lstlisting}

\item This is the square of the section line P, not the square of6P. The proof uses a section only on the surface with boundary; it does not make any claim about a section on a compact genus-one torsor.

\item The same affine comparison identifies the relative stable tangent framing with the regular product cap framing up to a map with zero Pontryagin winding. Indeed its derivative is triangular with identity linear derivative on the elliptic translation directions. The remaining diagonal factors are complex line scalings; their maps factor through a product of U1 groups. The universal degree-three Pontryagin transgression vanishes on that torus (equivalently each line factor has zero pi3), even when those scalings depend on all boundary variables. The triangular part is contracted by scaling its off-diagonal block. The regular product cap has zero relative contribution. Therefore

\begin{lstlisting}
chi(O of the capped Jacobian)=1,
signature(V0)=signature(capped Jacobian)=-8,
p_rel=integral p1(capped Jacobian)=3signature=-24.
\end{lstlisting}

\item The signature equality uses the regular fibre neighbourhood T2 timesD2, whose intersection form has zero signature; gluing along the closed T3 boundary gives additivity. These statements give the page and its actual line and framing data. The required actual return comparison is the real spin bordism P3K.6--7, and its normal-frame term is computed in P3K.9. Thus the later finite-group conclusion uses both the page calculation and that geometric comparison.

\end{enumerate}
\end{proof}




\begin{theorem}\label{night:P3F.11}
\StatementRecord{night:P3F.11}{proved}{S6\_06\_BLIND: GPT-6 Astra (OpenAI, ChatGPT Pro mode, effort Pro), recorded re-derivation in the checkpoint dated 18 November 2026}{\cref{night:P3F.3}, \cref{night:P3F.4}, \cref{night:P3F.7}, \cref{night:P3F.8}, \cref{night:P3K.7}}
The integer $A=(4k-\ell)/24$ is independent of the equivariant normal trivialization and is additive. Six parameter loops give exactly six times the integer of one loop.
\end{theorem}
\begin{proof}
\begin{enumerate}[label=\arabic*.]
\item Fix the orbit, Hopf and gluing conventions in P3F.7--8. For an orientation-preserving equivariant return preserving the two complex normal weight lines, its conifold replacement determines an integer

\begin{lstlisting}
A(Phi)=(4k(f)-ell(f))/24.
\end{lstlisting}

\item This integer is independent of the choice of equivariant tubular trivialization used for the replacement. It is additive under composition. In particular the actual return of the sixth parameter cover has integer exactly6A\(Phi\). This assertion concerns the integral Torelli invariant; it does not infer a class in Z/28 from multiplication by6 in Z/28.

\item Here is the complete framing comparison. After conjugating the negative normal line, the positive-weight bundle is E+=O\(1\) direct-sum O\(-1\), trivial as an SU2 bundle as in P3F.7. A loop of its unitary bundle automorphisms projectivizes to a loop of SO3 gauge transformations of the normal quotient. To compare two replacement fillings, take the product of this loop with the collar of the boundary S2 and collapse its two ends. Its adjoint is a map

\begin{lstlisting}
alpha:S3 -> SO3.
\end{lstlisting}

\item On the replacement S2-meridian timesD3 this comparison is exactly \(x,y\)-\textgreater(alpha(y)x,y). This follows in each normal trivialization by applying the projectivized unitary matrix to the Hopf line {[}u{]}; on the common boundary the two original trivializations are those same matrices. Hence the comparison is the primitive-sphere twist a to an integral power, with no second Torelli component.

\item For precision about all loops, evaluation at one point gives the fibration Map\_*\(S2,SO3\)-\textgreater Map\(S2,SO3\)-\textgreater SO3 with its constant-map section. Its fundamental group splits as pi3\(SO3\) plus pi1\(SO3\). The latter factor consists of matrices independent of the S2 variable and extends over the D3; it contributes no boundary-clutching map. The first factor is exactly the displayed alpha. A loop of orientation-preserving parametrizations of the fixed S2 also extends by its radial SO3 action on D3, after Smale1959 Theorem A (the strong deformation retraction of Diff+(S2) to SO3); equivalently its only fundamental-group ambiguity is that same extending constant rotation. Thus it adds no second Torelli component. These statements describe the actual morphism from a normal-trivialization change to the return-map difference.

\item For the holomorphic phase loops the integer can be computed without any homotopy shorthand. Let n\_-,n\_+ be their winding numbers on the original weight-minus and weight-plus lines. Over S2\_F timesS2\_theta, write h,t for the positive generators, and form the bundle obtained by clutching those phase loops. Its two positive-weight line classes are

\begin{lstlisting}
h-n_- t,      -h+n_+ t.
\end{lstlisting}

\item Consequently c1=\(n_+-n_-\)t and c2=\(n_-+n_+\)h t. The adjoint oriented three-plane bundle has

\begin{lstlisting}
p1=c1^2-4c2=-4(n_-+n_+)h t.
\end{lstlisting}

\item With Jin's generator alpha normalized by p1=+4, the comparison is -\(n_-+n_+\)a. In particular it can alter k and ell by an arbitrary common multiple of\(1,4\), but it leaves A unchanged. This fixes the sign of this phase-loop convention as well as its divisibility.

\item Jin Theorem4.9, with g=1, identifies the mapping-torus map with the homomorphism to Spin bordism over CP-infinity; its characteristic coordinates are\(A,k\). The representatives a,b have coordinates \(0,1\) and\(1,8\), respectively, as verified in P3F.3. For two returns choose tubular connections that are constant on a common collar and concatenate them. Their quotient replacements concatenate as the actual diffeomorphisms f and g, so the replacement of the composite is fg. Additivity of the cited bordism homomorphism gives A\(fg\)=A\(f\)+A\(g\). The framing calculation proves that this remains true for every separate choice of the tubular connections.

\item There is also a direct zero-class implication requiring no division in a finite group: the kernel of A in the exact presentation I(S2 timesS3)=Za direct-sum Zb is Za. P3F.4 constructs the Hopf lift of every a-power and extends it over D4 timesS3. Together with the pointwise cut-and-reglue identification P3F.8, it proves that the actual clutched sphere of a return in this kernel is standard. P3K.7--8 proves that the AC26 return belongs to this kernel, by two computations through the actual trace. The integral framework here then gives the reduction from six parameter loops to one.

\end{enumerate}
\end{proof}




\begin{theorem}\label{night:P3F.12}
\StatementRecord{night:P3F.12}{proved}{S6\_06\_BLIND: GPT-6 Astra (OpenAI, ChatGPT Pro mode, effort Pro), recorded re-derivation in the checkpoint dated 18 November 2026}{\cref{night:P2.2}}
The compact fixed-sheet torsor is smoothly translation-trivial through the specified $I_{12}$ model; the transported circle line differs from $\mathcal O(P-O)$ only by an integral fibre class.
\end{theorem}
\begin{proof}
\begin{enumerate}[label=\arabic*.]
\item The central-node component cover of the actual partial parameter filling is y\textsuperscript{12=t}4/\(t-1\)\^{}3. Its compact normalization C has genus3, with one point above the cusp. The Jacobian after minimalization has one I12 fibre there. The affine defect of the actual fixed sheet around that point is translation by6P: this is the first-coordinate part of Delta\_Q F=gamma-6delta in P2.2, with q=Q\^{}12. We retain this defect in the following argument; the component cover alone did not remove it.

\item Let Lambda\_E be the source elliptic lattice on the complement of the cusp of C. The original source matrices and cyclic-cover shifts are

\begin{lstlisting}
A=[[0,1],[-1,-1]], shift4;
B=[[0,-1],[1,0]], shift-3.
\end{lstlisting}

\item The commutator belongs to the cover subgroup, and

\begin{lstlisting}
G=[A,B]=[[1,-1],[-1,2]],
G-I=[[0,-1],[-1,1]], det(G-I)=-1.
\end{lstlisting}

\item Thus H0(C minus cusp;Lambda\_E)=0 integrally. The program compact\_torsor.py prints this matrix, its determinant and Smith form; root's independent full Schreier enumeration is in torsor\_schreier.py. For C0=C minus an open cusp disc, Poincare--Lefschetz duality with the oriented local system identifies H\textsuperscript{2}(C0,boundary;Lambda\_E) with this coinvariant group. The restriction H\textsuperscript{1}(C0;Lambda\_E) to H\textsuperscript{1}(boundary; Lambda\_E) is therefore surjective, including its torsion. Gluing the disc with the invariant-lattice extension j\_\emph{Lambda\_E gives H\textsuperscript{2}(C,j\_}Lambda\_E)=0. Equivalently, the top compact-support cohomology of the punctured curve is its integral coinvariant group; adding the skyscraper of local invariant sections changes no degree-two group.

\item The local sheaf here has a concrete meaning. On an I\_n Tate chart with parameter Q, the multiplicative parameter is Q\^{}n times a unit. A translation in the identity component is multiplication of the Tate coordinate by a nonzero smooth unit a\(Q\). It acts in every resolved node chart by opposite unit multiplications; these formulas extend across all components of the I\_n fibre. Locally a\(Q\) has a smooth logarithm. The kernel of this exponential consists of the invariant integral period: the other period contains log Q and cannot be a smooth single-valued section through Q=0. Hence, for the sheaf A\^{}0 of these extending identity-component translations, there is an exact sequence

\begin{lstlisting}
0 -> j_*Lambda_E -> L_infty -> A^0 -> 0,
\end{lstlisting}

where L\_infty is the sheaf of smooth sections of the extended complex Lie line bundle, and is fine. This verifies the sheaf at the singular fibre, rather than using the group law of the smooth fibres only. The local defect6P is in this identity component: P has a nonzero unit Tate coordinate at the original I1 fibre, and its pullback and all its multiples have component zero in I12. Its local logarithm therefore gives an extending translation.

\item It follows from the exact sequence and the vanishing just proved that H\textsuperscript{1}\(C,A^0\)=0. The actual affine fixed sheet, after the chosen local I12 extension, has a smooth translation section. Its translations extend through the fixed I12 resolution by the displayed unit formulas. It is therefore smoothly isomorphic, over C and by translations, to the sectioned Jacobian. This does not supply a holomorphic section or a holomorphic parameter filling.

\item The circle line also needs comparison under this isomorphism. Its fibrewise Picard point is P, and its multidegree on the components of I12 is zero. The latter follows from the actual source line P-O: both sections meet the identity component once, and neither meets any other component. The unit translation used above preserves the components and the fibrewise Picard point. Divide the transported line by O(P'-O') on the Jacobian. On every smooth elliptic fibre this quotient is holomorphically trivial, and at I12 its unit gluing has product1 and multidegree zero. Thus it is trivial on that fibre too.

\item Here is the local see-saw construction in this situation. Normalize a nonzero fibrewise holomorphic trivializing section at the zero section. Away from the nodal value it varies smoothly: it is the normalized one-dimensional kernel of the fibrewise Dolbeault operator. At I12 the construction uses the actual multiplicative unit charts: choose a frame on the first rational component and continue across the successive node units eta\_i. The continued frame on the last component agrees with the first because product eta\_i=1. The units and their logarithms extend smoothly in Q, so this normalized frame extends over the nodal charts. Two such normalized local frames differ by a function of the base alone, because a holomorphic function on a connected compact fibre is constant. Their transition functions define a line bundle on C. This proves, for the actual transported line, the equality

\begin{lstlisting}
c1(L)=c1(O(P'-O'))+n F       for an integer n.
\end{lstlisting}

\item It proves the required filtration statement via explicit normalized frames. In particular it does not infer that equality merely from triviality of the action on the associated graded of Leray. Since F\textsuperscript{2}=F.(P'-O')=0, its square is exactly-2. The smooth affine comparison preserves orientation and signature, so p1 of the compact sheet pairs to-24. The unspecified integer n records the permitted boundary gauge and has no effect on either of these pairings.

\end{enumerate}
\end{proof}




\begin{theorem}\label{night:P3F.13}
\StatementRecord{night:P3F.13}{proved}{S6\_06\_BLIND: GPT-6 Astra (OpenAI, ChatGPT Pro mode, effort Pro), recorded re-derivation in the checkpoint dated 18 November 2026}{\cref{night:P2.2}, \cref{night:P3K.6}, \cref{night:P3K.8}, \cref{night:P3K.9}}
The angular origin comparison is $\theta\gamma+6\theta(1-\theta)\delta$ and its full $Q$-deck line clutch has winding $-12$. The earlier complex corner attempt remains uncompleted; the actual real trace P3K.6--9 now proves the monodromy application.
\end{theorem}
\begin{proof}
\begin{enumerate}[label=\arabic*.]
\item For the default twists let F\(s,v\) be the exact origin difference in P2.2, and keep gamma\(s,v\)=\(6mu,C-s-12v\). Direct polynomial subtraction gives, for real theta,

\begin{lstlisting}
F(s,v+theta)-F(s,v)
   = theta gamma(s,v)+6theta(1-theta)delta.       (P3F.13.1)
\end{lstlisting}

\item On a fixed-radius parameter circle, v=ib+theta. The final summand is real in the delta direction. It is a unitary gauge, equal to1 after exponentiation at both endpoints. It can be removed by a homotopy of periodic gauges. The circle quotient of the origin comparison is therefore exactly the gamma-translation loop, with all radial terms independent of theta.

\item The base deck must still be retained. Since gamma\(s+12,v\)= gamma\(s,v\)-12delta,

\begin{lstlisting}
[F(s+12,v+theta)-F(s+12,v)]
   -[F(s,v+theta)-F(s,v)]=-12theta delta.        (P3F.13.2)
\end{lstlisting}

\item Thus the two angular lifts differ on the Q deck by the unitary gauge exp(-24pi i theta) in the exp(2pi i zeta2) convention, of winding-12. This gauge is the integral clutch of the circle line. It does not disappear on passing to the orbit: its degree is retained by the principal circle bundle. compact\_torsor.py checks both exact identities and the endpoint values.

\item The normal-crossing chart suggested by the second Tate multiplier is XY=Qr after q=Q\^{}12, because q r\textsuperscript{12=\(Qr\)}12. Its elementary small resolution has charts

\begin{lstlisting}
X=Qz, r=zY;          Q=Xw, Y=wr.
\end{lstlisting}

\item At a node Q=uv of the resolved source I12 surface the remaining chart uv=Xw is resolved by

\begin{lstlisting}
X=uk, v=kw;          u=Xk, w=kv.
\end{lstlisting}

\item These four-variable charts are smooth and separate the central fixed surface from the fixed cusp trace. They give the local candidate for an orbit bordism after deleting fixed tubes. The central fixed-sheet normals have opposite weights and product trivial, so their positive Hopf projectivization is a product S2 bundle with Euler class u+c1\(L\). This inner boundary is nonspin when the compact sheet has K=5F. The canonical quotient structure is spin\^{}c with determinant -K on that boundary; no claim that this compact product is spin is made.

\item Historical scope of this complex corner attempt: its cyclic descent and outer identification after resolving the source I12 charts were not completed. P3F.13 proves the exact local clutch, including its line gauge. The actual global application is now proved by the real trace P3K.6--7, the analytic symbol comparison P3K.8, and the rotating normal framing P3K.9. Those constructions retain the same clutch and compute the original return without asserting this complex corner completion.

\end{enumerate}
\end{proof}




\begin{theorem}\label{night:P3F.14}
\StatementRecord{night:P3F.14}{proved}{S6\_06\_BLIND: GPT-6 Astra (OpenAI, ChatGPT Pro mode, effort Pro), recorded re-derivation in the checkpoint dated 18 November 2026}{\cref{night:P3F.11}, \cref{night:P3F.12}, \cref{night:P3K.6}}
The inherited spin$^c$ index of a fixed-sheet Hopf boundary is precisely $\chi(S,L_+)$, independently of the negative-weight normal line.
\end{theorem}
\begin{proof}
\begin{enumerate}[label=\arabic*.]
\item This is the general calculation for the explicitly specified Hopf boundary. The earlier complex corner attempt did not supply a global boundary identification; the actual real trace P3K.6--9 now supplies the application and its complete orientation and line comparison. Let S be a compact complex surface, and let L\_+,L\_- be its two complex normal lines of circle weights+1,-1 in a stable complex eight-dimensional chart. Write a=c1\(L_+\), b=c1\(L_-\), d=c1\(TS\). Its free normal-sphere quotient is B=P(L\_+ direct-sum conjugate(L\_-)). Orient its sphere fibre so that the circle Euler class x has integral1. Then

\begin{lstlisting}
pi_*(x)=1, pi_*(x^2)=a-b,
pi_*(x^3)=a^2-ab+b^2,
p1(TB)=pi^*(p1(TS)+(a+b)^2).
\end{lstlisting}

\item These follow from the projective-bundle relation with roots a,-b. The stable complex structure inherited from the original two opposite weights has spin\^{}c determinant D=pi\^{}*\(d+a+b\). Indeed the associated normal lines have opposite circle characters, so those characters cancel in their determinant. The degree-six spin\^{}c index polynomial for this structure twisted by the circle line is

\begin{lstlisting}
integral_B exp(D/2) exp(x) Ahat(TB).
\end{lstlisting}

\item Substitution of the four displayed pushforwards gives exactly

\begin{lstlisting}
(3d^2-p1(TS))/24 + (a^2+ad)/2
   = chi(O_S)+(a^2-a.K_S)/2
   = chi(S,L_+).
\end{lstlisting}

\item The equality on the right is the surface Riemann--Roch polynomial; the calculation also proves that the negative-weight normal line b cancels entirely. No spin structure on S or on B is required. For an actual spin-six boundary with D=0 the same index polynomial is (4x\textsuperscript{3}-p1 x)/24, the integer A used in P3F.11.

\item For the compact fixed sheet of P3F.12, a=M+nF, chi\(O\)=1, M\textsuperscript{2}=-2, M.K=F\textsuperscript{2}=M.F=0. Its index contribution is therefore0 for every such fibre-supported shift. The fixed cusp curve has positive normal line O\(-1\); its Dolbeault cohomology on P1 is zero in both degrees. These are exact local index inputs. P3K.6--8 subsequently constructs the actual real bordism and identifies its spin-c symbol. The resulting return index therefore follows through that comparison; no analytic corner completion is asserted.

\end{enumerate}
\end{proof}




\begin{theorem}\label{night:P3F.15}
\StatementRecord{night:P3F.15}{proved}{S6\_06\_BLIND: GPT-6 Astra (OpenAI, ChatGPT Pro mode, effort Pro), recorded re-derivation in the checkpoint dated 18 November 2026}{\cref{night:P3F.13}, \cref{night:P3K.6}}
The explicit cyclic action on the twelve labelled node charts extends by units and has the required twelfth-iterate gamma period and circle-line winding. Its free labelled central quotient has multiplicity twelve and is not asserted to be the outer cusp chart.
\end{theorem}
\begin{proof}
\begin{enumerate}[label=\arabic*.]
\item Here is the explicit descent on the rank-one node charts. Bounded nonzero period units are harmless in these identities: choose their logarithms on the cusp disc and split their twelfth roots in the following multipliers. Put rho=Qr, zeta=exp\(2pi i/12\), eta=exp\(pi i/12\), and index the node charts by i in Z, modulo the gamma period of length12:

\begin{lstlisting}
X_i=W/rho^i, Y_i=rho^(i+1)/W, X_i Y_i=rho.
\end{lstlisting}

\item The deck map on the dense torus is

\begin{lstlisting}
D:(Q,i,W)->(zeta Q,i-1,eta rho^(-1) W).
\end{lstlisting}

\item Consequently its formulas in the node charts are

\begin{lstlisting}
X'_(i-1)=eta zeta^(1-i) X_i,
Y'_(i-1)=eta^(-1) zeta^i Y_i.
\end{lstlisting}

\item Their product is zeta rho, as required for the new parameter. Both multipliers are units, so this map extends across each node and preserves the opposite circle weights. It is free on the twelve node labels even at Q=0. After twelve iterations the W multiplier is eta\^{}12 zeta\^{}\(-66\) rho\textsuperscript{\(-12\)=rho}\(-12\), since the two first factors both equal-1. Accompany each step by the locally defined elliptic translation-P/2; the twelfth step is-6P. Thus the twelfth iterate is the inverse gamma period, including its source-translation part. This proves descent on the labelled rank-one charts.

\item On a fixed-radius r-circle, each factor rho\^{}\(-1\) has angular winding-1. The twelve-step gauge has winding-12. It is exactly the winding in \(P3F.13.2\), with the same W/exp(zeta2) sign conversion. This verifies the integral circle-line clutch, not just the orbit action. It leaves the outer contraction in the earlier complex corner attempt unproved. The actual monodromy application is instead completed by the real trace P3K.6--9.

\item Multiplicity check on the outstanding outer identification: a free action on the twelve labels, with reduced divisor Q=0 on those charts, would descend q=Q\^{}12 to a divisor of multiplicity twelve. The actual outer cusp for r nonzero has multiplicity one. Thus the labelled fixed-sheet descent above cannot itself be its outer chart. A valid comparison must explicitly account for the ramification or fixed strata, or for a changed divisor multiplicity. The missing morphism is therefore stronger than merely contracting the additional source I12 exceptional curves.

\end{enumerate}
\end{proof}




\begin{theorem}\label{night:P3F.16}
\StatementRecord{night:P3F.16}{proved}{S6\_06\_BLIND: GPT-6 Astra (OpenAI, ChatGPT Pro mode, effort Pro), recorded re-derivation in the checkpoint dated 18 November 2026}{Definitions and the displayed constructions}
The actual labelled comb-smoothing parameter loop has identity smooth monodromy, by a single-valued radial trivialization of the entire glued family.
\end{theorem}
\begin{proof}
\begin{enumerate}[label=\arabic*.]
\item Let \(s\) denote the smoothing parameter of the labelled comb in Viaclovsky, arXiv:2610.10451, in the germ used to construct the compact threefolds. The smooth family of these threefolds over a sufficiently small punctured \(s\)-disc is a product. Its return diffeomorphism is isotopic to the identity. For a sphere member its class in \(\Theta _7\) is therefore zero.

\item Here the loop is the parameter of the comb smoothing appearing in the source. No identification of that parameter with the scalar parameter \(w=e^{-\pi i c}\) of another construction is used.

\subsubsection{Construction and proof}\label{construction-and-proof}

\item Write \(B'=B\setminus\Sigma\) for the base with its critical values removed, and \(\Lambda\) for the fixed integral homology local system in the source's Section ``Global extension,'' equations \texttt{Jac} and \texttt{fixedreal}. Every admissible Hodge plane \(\mathcal F_s\) gives the real-linear evaluation isomorphism \[
 \Lambda_{\mathbb R,b}\longrightarrow\mathcal F_s(b)^*,
 \qquad v\longmapsto(\omega\longmapsto\omega(v)).
\] Injectivity follows because a real vector annihilating \(\mathcal F_s\) also annihilates its conjugate, whose direct sum is the whole cohomology space. Both sides have dimension four. Thus evaluation identifies the smooth family on \(B'\) with the fixed bundle \(\mathscr M=\Lambda_{\mathbb R}/\Lambda\). This identification is smooth jointly in \((b,s)\) for every \(s\ne0\), even near a comb tooth. No limiting Hodge plane at that tooth is required.

\item Choose pairwise disjoint closed discs about the critical values and put \(B^\circ=B\setminus\bigcup\operatorname{int}D_i\). All their collars avoid the comb teeth. The local finite and toric fillings, their periods and their gluing translations are holomorphic on these collars for the entire parameter disc, including \(s=0\). This is the construction in the source's Proposition \texttt{topology}: its gluing translations are extended from \(s=0\) by the exact sequence \(0\to\Lambda\to\mathcal L\to\mathcal T\to0\), the vanishing of \(H^1(A_i\times\Delta,\mathcal L)\), and extension of a vector-bundle section from the divisor \(s=0\). Consequently the actual attachments, not merely their isotopy classes, extend smoothly across \(s=0\).

\item For precision, one can make all the trivializations single valued in \(s\), as follows. Set \(W=\mathscr M|_{B^\circ}\). On each compact filled piece \(\mathcal V_i\to\Delta\), choose smooth horizontal lifts of the two real coordinate vector fields on the parameter disc. The local defining equations have parameter-independent strata; the lifts can be chosen to preserve the base coordinate and be tangent to every boundary collar. On a finite quotient average the lifts over its finite group. On a toric filling patch the lifts on the regular charts; the units vary smoothly with the parameter. Integrate their combination along the radial path \(\tau\mapsto\tau s\), \(0\leq\tau\leq1\). Compactness and a smaller common disc give diffeomorphisms \[
 a_{i,s}:V_{i,0}\longrightarrow V_{i,s},\qquad a_{i,0}=\mathrm{id},
\] smooth jointly in \(s\). In the fixed source of these trivializations, the actual attachments are smooth maps \[
 \varphi_{i,s}:\partial V_{i,0}\longrightarrow\partial_iW,
 \qquad |s|<\epsilon.
\] Their central value is the original attachment, by the prescribed translation extension just described.

\item Put \(g_{i,\tau,s}=\varphi_{i,\tau s}\varphi_{i,0}^{-1}\). This is a smooth family of isotopies beginning at the identity. Its generating vector field is \[
 Z_{i,\tau,s}=(\partial_\tau g_{i,\tau,s})\circ g_{i,\tau,s}^{-1}.
\] Extend it constantly in the radial coordinate of a fixed collar of \(\partial_iW\), multiply by a cutoff which is one at the boundary and zero at the inner edge, and extend by zero on the rest of \(W\). The collars are disjoint. Integrating their sum gives a smooth family \(K_s\in\operatorname{Diff}(W)\), with \[
 K_s|_{\partial_iW}=g_{i,1,s},\qquad K_0=\mathrm{id}.
\] These maps glue to the \(a_{i,s}\), because \(K_s\varphi_{i,0}=\varphi_{i,s}\) on each common boundary. We have therefore constructed an actual smooth isomorphism of bundles \[
 X_0\times\Delta^*\longrightarrow\mathcal X|_{\Delta^*}.
\] Every step depends on the complex number \(s\), rather than a chosen argument of \(s\). The isomorphism is single valued around the positive parameter circle, so its return is the identity. Any other connection on the same bundle has isotopic return: interpolating connections produces a smooth isotopy of their holonomies. This proves the claim with its actual parameter and without using the annular index formula.

\end{enumerate}
\end{proof}




\begin{theorem}\label{night:P3F.17}
\StatementRecord{night:P3F.17}{proved}{S6\_06\_BLIND: GPT-6 Astra (OpenAI, ChatGPT Pro mode, effort Pro), recorded re-derivation in the checkpoint dated 18 November 2026}{\cref{night:P3F.10}, \cref{night:P3F.25}, \cref{night:P3F.9}, \cref{night:P3K.11}, \cref{night:P3K.6}}
For coprime multiplicities $(3k,4l)$ the regular-binding page has degree $12kl$, Euler characteristic $12kl$, signature $-8kl$, section-line square $-2kl$ and relative Pontryagin pairing $-24kl$.
\end{theorem}
\begin{proof}
\begin{enumerate}[label=\arabic*.]
\item Put \(m=3k\), \(n=4l\), with \(\gcd(m,n)=1\), and use any of the logarithmic data in Viaclovsky, arXiv:2610.39102, giving a sphere. The regular-fibre page of the circle quotient has base genus \((mn-m-n+1)/2\), degree \(d=mn=12kl\) over the original elliptic base, and its sectioned elliptic completion has precisely \(d\) fibres of type \(I_1\). Its invariants are \[
 \chi(\mathcal O)=kl,\quad \sigma=-8kl,\quad
 h(P)=2kl,\quad (aP)\cdot O=kl(a^2-1)
\] for every nonzero integer \(a\). For the actual circle line on a page disjoint from the fixed sphere, with the boundary marking of P3F.10, \[
 s_\delta=-2kl,\qquad p_{\rm rel}=-24kl.
\] These equations compute the page data. The actual real trace P3K.6--9 supplies the original return comparison and P3F.25 records its higher-family application. The fully signed higher-parameter comparison is developed in P3K.11 and is under audit at this checkpoint.

\subsubsection{Proof}\label{proof}

\item The source retains the same rational elliptic surface, section \(P\), line \(\mathcal O(P-O)\), and cusp as E26. Its finite caps are the further root pullbacks of the degree-three and degree-four good reductions. The generators act on the invariant real circle coordinate \(\psi\) by \(n_i/m_i\), with \(\gcd(n_i,m_i)=1\); see the source's equations \texttt{logdat}, \texttt{cyclic}, \texttt{glue}, and Proposition \texttt{logprop}. Since \(\psi(\delta)=0\), no nonidentity finite deck element can be cancelled by the \(\delta\)-circle. Thus the circle has no new stabilizers in these caps. The cusp orbit is the same relative product constructed in P3F.9.

\item After quotienting by this circle, the two base solid tori have Seifert multiplicities \(m,n\). The source's sphere condition is the integral equation \(m_1n_0+m_0n_1=\pm1\); its meridian relation matrix is unimodular. The two solid tori therefore glue to \(S^3\), with the weighted circle fibration of weights \(m,n\), with the sign of its orientation fixed by that equation. A regular binding has a page described by \(b^m-a^n=1\), using positive real radial rescaling with weights \(m,n\), exactly as in P3F.10. Its map to the weighted base has degree \(mn\), ramification indices \(m,n,mn\) at the two exceptional values and the chosen regular binding value. Riemann--Hurwitz gives \(2-2g=m+n-mn+1\).

\item Choose the binding away from all three critical values. The root indices at the original \(IV^*\) and \(III\) fibres are multiples of their good-reduction indices, so their relatively minimal pullbacks are smooth. The original nodal fibre is unramified and has \(d\) preimages. Thus the elliptic completion has \(d\) positive nodes and no reducible fibres. Its Euler number is \(d\), its holomorphic Euler characteristic is \(d/12=kl\), and its signature is \(-2d/3=-8kl\). Height pulls back by \(d\), giving \(h(P)=d/6=2kl\). Shioda's formula, with zero reducible-fibre corrections, gives \(2kl a^2=2kl+2(aP)\cdot O\). In particular \(P\cdot O=0\). The section line square is \((P-O)^2=-2\chi-2P\cdot O=-2kl\).

\item The actual page may be chosen disjoint from the fixed sphere by the same phase choice as P3F.10. Its affine elliptic torsor is smoothly translation-trivial over the base with boundary. The normalized fibrewise Picard point of the circle line is \(P\), and a change of this trivialization changes its line by a base pullback; such a class vanishes on this base. The boundary first Chern class is zero, so the relative square is independent of its relative lift, by the displayed boundary cup-product calculation in P3F.10. The relative tangent framing comparison there, and signature additivity on capping the regular elliptic boundary, give \(p_{\rm rel}=3\sigma=-24kl\).

\end{enumerate}
\end{proof}




\begin{theorem}\label{night:P3F.18}
\StatementRecord{night:P3F.18}{proved}{S6\_06\_BLIND: GPT-6 Astra (OpenAI, ChatGPT Pro mode, effort Pro), recorded re-derivation in the checkpoint dated 18 November 2026}{\cref{night:P3F.2}}
The two supplied one-logarithmic families have degree-four regular-binding page covers; their precise source and page fibre types are those listed below.
\end{theorem}
\begin{proof}
\begin{enumerate}[label=\arabic*.]
\item The following source data have been read directly in the supplied construction sources, rather than inferred from their arXiv titles.

\begin{longtable}[]{@{}
  >{\raggedright\arraybackslash}p{(\columnwidth - 8\tabcolsep) * \real{0.2000}}
  >{\raggedright\arraybackslash}p{(\columnwidth - 8\tabcolsep) * \real{0.2000}}
  >{\raggedright\arraybackslash}p{(\columnwidth - 8\tabcolsep) * \real{0.2000}}
  >{\raggedright\arraybackslash}p{(\columnwidth - 8\tabcolsep) * \real{0.2000}}
  >{\raggedright\arraybackslash}p{(\columnwidth - 8\tabcolsep) * \real{0.2000}}@{}}
\toprule\noalign{}
\begin{minipage}[b]{\linewidth}\raggedright
Construction
\end{minipage} & \begin{minipage}[b]{\linewidth}\raggedright
Source fibres
\end{minipage} & \begin{minipage}[b]{\linewidth}\raggedright
Defining section and translation
\end{minipage} & \begin{minipage}[b]{\linewidth}\raggedright
Finite cap
\end{minipage} & \begin{minipage}[b]{\linewidth}\raggedright
Regular-binding elliptic page
\end{minipage} \\
\midrule\noalign{}
\endhead
\bottomrule\noalign{}
\endlastfoot
Viaclovsky 2609.33785 & \(III^*,I_1,I_1,I_1\) & \(P=(0,1)\) on \(y^2=x^3+tx+1\), \(M=P-O\), translation \(2P\) & order four & degree four, base genus zero, twelve \(I_1\) \\
Liu--Rollenske 2609.33826 & \(III^*,II,I_1\) & \(P=(-t^2,it^3)\) on their displayed Weierstrass model, \(M=P-O\), translation \(2P\) & order four & degree four, base genus zero, four \(II\) and four \(I_1\) \\
\end{longtable}

\item Both primitive sections have height \(1/2\) and are disjoint from \(O\), as also follows from the E7 entry in P3F.2. The first source chooses its scalar linearization to vanish at one nodal fibre; the other two nodal fibres have rank-one torus fillings. The second source chooses the linearization to vanish at its nodal fibre and leaves the \(II\) fibre as the product cusp with an elliptic factor; its sole logarithmic cap is the order-four cap at \(III^*\).

\item For either source, quotienting its finite cap by the \(\delta\)-circle gives a Seifert solid torus of multiplicity four. The meridian of the other solid torus has multiplicity one: the source twist has primitive invariant circle coordinate equal to one. The resulting Seifert \(S^3\) has weights \((1,4)\). Its regular-binding page is a disc. Its compactification maps with degree four to the source base, branched only at the order-four value and the regular binding. Thus \(III^*\) has good reduction and the remaining source critical values have four unramified preimages each. This proves the last column, including the four cuspidal fibres in the Liu--Rollenske case.

\item The Euler number is twelve in both cases, since \(e(II)=2\) and \(e(I_1)=1\). Consequently their sectioned page completions have \(\chi(\mathcal O)=1\), \(\sigma=-8\), \(h(P)=2\), and no reducible fibres. Therefore \(P\cdot O=0\), \((P-O)^2=-2\), and \((aP)\cdot O=a^2-1\). These are exact page computations. The scalar loop's comparison with its global equivariant return is not supplied by these numbers.

\end{enumerate}
\end{proof}




\begin{theorem}\label{night:P3F.19}
\StatementRecord{night:P3F.19}{proved}{S6\_06\_BLIND: GPT-6 Astra (OpenAI, ChatGPT Pro mode, effort Pro), recorded re-derivation in the checkpoint dated 18 November 2026}{Definitions and the displayed constructions}
Yang Type B and Type C have explicit nontrivial cyclic quotients of their fundamental groups, excluding homeomorphisms to the six-sphere.
\end{theorem}
\begin{proof}
\begin{enumerate}[label=\arabic*.]
\item Every Type B and Type C construction in the supplied Yang source has a nontrivial cyclic quotient of its fundamental group: \(\mathbb Z/2\) in Type B and \(\mathbb Z/3\) in Type C. This applies to every index and every admissible twist in those two constructions. In particular none is diffeomorphic, or homeomorphic, to \(S^6\).

\item Here is a direct quotient of the attaching presentation, sufficient for the obstruction and independent of any assertion of the full group. On the regular part, kill every fibre loop. Send the two finite base meridians to \(1,-1\) in \(\mathbb Z/g\), where \(g=\gcd(m_1,m_2)\), and the cusp meridian to zero. All semidirect product relations hold because the fibre loops have been killed. In the local finite quotient chart, the only relation after killing the fibre loops is the \(m_i\)-th power of the corresponding base meridian: the affine translation in that power is a fibre loop. This relation holds in \(\mathbb Z/g\). The local cusp attachment kills its base meridian up to a fibre loop and kills some fibre loops; these relations also hold. The sphere relation holds because \(1+(-1)+0=0\). The map thus extends over all three attachments and is surjective, since a finite meridian maps to one.

\item The construction has \((m_1,m_2)=(4,6)\) in Type B and \((3,6)\) in Type C, giving the displayed nontrivial quotients. The quotient maps are explicit morphisms of the Van Kampen presentation; a hypothetical homeomorphism to \(S^6\) would induce an isomorphism on fundamental groups, whereas the latter sphere has trivial fundamental group.

\item For the canonical twists \((\ell_0,\ell_1,\ell_2)=(0,1,-1)\), the source's invariant-closure calculation kills exactly the other three fibre generators: its cusp kills \(w,\delta\), and its first finite matrix sends \(w\) to \(u-\delta\) in Type B and to \(u-w-\delta\) in Type C. The resulting presentation is abelian on \((z,x,y)\), with relation matrix \[
 \begin{pmatrix}-1&m_1&0\\1&0&m_2\\0&1&1\end{pmatrix}.
\] Its Smith diagonal is \((1,1,2)\) or \((1,1,3)\), respectively. The program \texttt{other\_family\_checks.py} recomputes these matrices.

\end{enumerate}
\end{proof}




\begin{theorem}\label{night:P3F.20}
\StatementRecord{night:P3F.20}{proved}{S6\_06\_BLIND: GPT-6 Astra (OpenAI, ChatGPT Pro mode, effort Pro), recorded re-derivation in the checkpoint dated 18 November 2026}{\cref{night:P3F.1}, \cref{night:P3F.21}, \cref{night:P3F.31}, \cref{night:P3K.65}}
Wang's supplied threefold is simply connected with $H_2=H_3=\Z/2$, as computed by the full integral exterior and clutch complex.
\end{theorem}
\begin{proof}
\begin{enumerate}[label=\arabic*.]
\item The supplied Wang construction is simply connected and has \(H_2=H_3=\mathbb Z/2\), so it cannot be homeomorphic to \(S^6\). The computation below retains the primitive \((2,1)\) rank-one exponent at \(I_2\); the full integer matrices have been executed in \path{code/part3/other_family_checks.py} and are saved in its JSON output.

\item The source configuration is \(IV^*,I_2,I_1,I_1\), with line \(P-O\), translation \(-6P\), and one logarithmic cap of order three. At \(I_2\) the Tate periods are \[
 h(z,v)=(q^2c_1z,qc_2v),\qquad
 g(z,v)=(c_3z,\lambda c_4v),
\] where \(g\) is obtained from the original lift by multiplying it by \(h^3\). The monomial coordinates \(x=v,y=v^2/z\) send the exponent \((2,1)\) to \((1,0)\); this is a unimodular coordinate change. The zero-exponent period contracts the \(y\) coordinate by \(\lambda^2\) times a unit. Thus the central rank-one fibre has one component and the primitive pinching-circle specialization. The ordinary rank-one \(I_1\) has the same pinching map. At the remaining \(I_1\) the scalar zero gives the unimodular two-period toric model, with its standard hexagonal component and integral specialization. These are the displayed local charts in the supplied source's Lemma \texttt{fill:tate} and Propositions \texttt{fill:rankone}, \texttt{fill:a2}; the exponent and primitive collapse calculations above verify their use here.

\item The cohomological monodromies in basis \((\psi,u,w,\delta)\) are \[
 T_* =\begin{pmatrix}1&0&-6&-2\\0&-1&-1&1\\0&1&0&-1\\0&0&0&1\end{pmatrix},\quad
 T_2=\begin{pmatrix}1&0&6&3\\0&1&2&1\\0&0&1&0\\0&0&0&1\end{pmatrix},
\] \[
 T_+=\begin{pmatrix}1&0&0&0\\0&2&1&0\\0&-1&0&0\\0&0&0&1\end{pmatrix},\quad
 T_-=\begin{pmatrix}1&0&0&-1\\0&1&1&0\\0&0&1&0\\0&0&0&1\end{pmatrix}.
\] Their ordered product is the identity, \(T_*^3=I\), and the invariant homology vector of the logarithmic transformation is \(\ell=(1,-2,-4,0)^t\). Its invariant circle coefficient is one, so the order-three affine action is free. The three Mumford sections extend, and their meridians are killed. Their homological monodromy images kill \(\widehat w,\widehat\delta,\widehat u+\widehat w\), leaving only \(\widehat\psi\). The remaining meridian is killed by the base relation, and its cube relation kills \(\ell\), hence the remaining generator. Every boundary map is onto its filled fundamental group by the displayed primitive collapse or affine quotient presentation. This proves simple connectivity.

\item For every exterior degree, use the same integral parabolic complex as in P3F.21, with this ordered tuple. The ranks of its four local images are \((2,1,1,2)\), \((4,2,2,2)\), and \((2,1,1,2)\) in degrees one, two and three. Both differentials have ranks three, five and three. The stacked first differential has only unit nonzero Smith factors. Thus its middle cohomology is zero integrally. The invariant generators are \[
 1,\quad\psi,\quad\xi=uw+6\psi\delta,\quad\phi=\psi uw,\quad\nu=\psi uw\delta,
\] and the free coinvariant generators are \[
 1,\quad[\delta],\quad[\psi\delta],\quad[uw\delta],\quad[\nu],
 \qquad[uw]=6[\psi\delta].
\] The complete local image bases and both integer differentials are recorded in the program output, not just their ranks.

\item At the order-three fibre, the reduced quotient has torsion-free cohomology: its affine abelianization is free of rank two, and Poincaré duality applies. The invariant degree-one basis is \(\psi,h=-2u+w-3\delta\), with \(h(\ell)=0\). Its pullback image is \(\mathbb Z(3\psi)\oplus\mathbb Zh\). The invariant degree-two basis \((\psi h,\xi)\) has intersection matrix \[
 \begin{pmatrix}0&-3\\-3&12\end{pmatrix}.
\] Its determinant has absolute value nine, exactly the determinant of the degree-three pullback of the downstairs unimodular lattice, so the degree-two image is the full invariant lattice. The perfect pairing with degree one then gives degree-three index one. Top degree has index three. The global specialization indices are therefore \((3,1,1,3)\). The two nontrivial cokernels are reached by the global generators \(\psi,\nu\), so all first-cohomology terms of the direct images remain zero.

\item The Leray page consists of columns zero and two; its first column is generated by \(1,3\psi,\xi,\phi,3\nu\). Its second column is generated by the five displayed coinvariants times the positive base class \(H\). The only affine clutch is the logarithmic translation, so its cellular prism gives \[
 d_2(a)=\frac H3[\iota_\ell a].
\] The four contractions in coinvariants, before the indicated source multiplicities, are \(1,6[\delta],6[\psi\delta],[uw\delta]\). Hence the four integral differentials are multiplication by \[
 (1,2,2,1).
\] No higher differential is possible. Thus \(H^3=H^4=\mathbb Z/2\), \(H^0=H^6=\mathbb Z\), and the other cohomology groups vanish; Poincaré duality gives the stated homology. A homeomorphism to \(S^6\) would induce an isomorphism on integral homology, contradicting this explicit nonzero two-torsion.

\item The original line circle also illustrates a local feature relevant to further comparisons: at the \(I_2\) pinching locus its direction \((0,1)\) intersects the collapsing torus circle \((2,1)\) in a subgroup of order two. Therefore its isotropy there is exactly \(\mathbb Z/2\). This follows by solving \(2\theta\in\mathbb Z\) in \(\mathbb R/\mathbb Z\), and must be retained in any circle-orbit calculation for this construction.

\subsubsection{Further component-supported cases}\label{further-component-supported-cases}

\item For the component-supported OS49 target in the prompt with \(Q=3P\), the regular-binding degree-three page has three \(II\) and three \(I_2\) fibres, hence \(\chi(\mathcal O)=1\), \(\sigma=-8\), \(h(P)=1/2\), and total reducible-fibre correction \(3/2\). The height formula gives \(P\cdot O=0\). Its component-permuting translation is carried through the actual positive-fan page in P3F.31--32, giving positive-fan index zero. P3K.65 retains the full native C2 closing and Cox ruling arrow and gives negative-fan class \(-3\) for \(\lambda=r^3\), or \(-1\) for the actual shorter scalar loop \(\lambda=w\). No negative-fan integer index is assigned here.

\item The eight Mordell--Weil entries and twenty-one configurations are enumerated in P3F.1. Enumerating those source configurations does not yet enumerate every compact sphere filling produced from all choices of translation, vertical linearization divisor and logarithmic data. The bounded slice mentioned in the prompt has no specified numerical bounds in the supplied claim. We have not silently supplied bounds or converted its asserted count into a theorem.

\end{enumerate}
\end{proof}




\begin{theorem}\label{night:P3F.21}
\StatementRecord{night:P3F.21}{proved}{S6\_06\_BLIND: GPT-6 Astra (OpenAI, ChatGPT Pro mode, effort Pro), recorded re-derivation in the checkpoint dated 18 November 2026}{\cref{night:P3F.10}, \cref{night:P3F.22}, \cref{night:P3F.23}, \cref{night:P3F.25}, \cref{night:P3F.7}, \cref{night:P3K.6}}
The weighted OS47 data $M=\mathcal O(7G-O)$ and $Q=2G$ give an actual holomorphic sphere family over $\lambda=r^4$. Its actual page has $s_\delta=-14$ and $p_{\rm rel}=-72$; its circle action is semifree through the entire $I_7$ filling.
\end{theorem}
\begin{proof}
\begin{enumerate}[label=\arabic*.]
\item There is a proper holomorphic family of complex structures on the standard six-sphere over a punctured parameter disc, obtained from the rational elliptic configuration \(I_7+III+I_1+I_1\) by taking \[
 M=\mathcal O_S(7G-O),\qquad Q=2G,
\] where \(G\) generates its Mordell--Weil group, choosing the linearization zero at one of the nodal fibres, and making the specified order-four logarithmic transformation at \(III\). An explicit parameter is \(r\), with scalar linearization parameter \(\lambda=r^4\).

\item The actual regular-binding page of its circle quotient has four \(I_7\) fibres and eight \(I_1\) fibres. Its elliptic completion has \[
 \chi(\mathcal O)=3,\quad \sigma=-24,\quad
 h(7G)=14,\quad (7G)\cdot O=4.
\] The circle-line square on the page is \(-14\), and its relative Pontryagin pairing is \(-72\). Thus the \textbf{page expression} is \[
 \frac{p_{\rm rel}-12s_\delta}{24}
   =\frac{-72+168}{24}=4.
\] The actual signed comparison is proved in P3F.25 by the real trace P3K.6--7: the return invariant is the negative of this displayed page expression, so \(A=-4\), with oriented class \(24\) in \(\Theta_7\). The original construction and page calculation are retained below as the inputs to that completed comparison. Its line section is \(7G\); replacing it by the primitive section \(G\) changes the line and the calculation.

\subsubsection{1. An explicit rational elliptic surface and primitive section}\label{an-explicit-rational-elliptic-surface-and-primitive-section}

\item Put \[
 f=1+8t+40t^2+48t^3,\qquad
 g=1+12t+84t^2+280t^3+504t^4,
\] and take the minimal regular model of \[
 y^2=x^3-3f(t)x+2g(t).
\] The exact identity \[
 g^2-f^3=-1728t^7(64t^2+13t+2)
\] gives \(I_7\) at zero, simple nodal fibres at \(t_\pm=(-13\pm7i\sqrt7)/128\), and \(III\) at infinity. At infinity the minimal Weierstrass vanishing orders of \((a_4,a_6,\Delta)\) are \((1,2,3)\). At zero \(a_4,a_6\) are units and the discriminant has order seven. The other two discriminant roots are distinct and have unit \(a_4\). These checks give exactly the stated Kodaira types and no additional singular fibres.

\item The point \[
 G=(1+4t,12\sqrt3\,t^2)
\] satisfies the equation. Its polynomial coordinates show that it is disjoint from the zero section. At zero, write \(\sqrt f=1+4t+12t^2+O(t^3)\). In the two local branches of the node, the point has orders two and five: one of \(y\pm\sqrt3(x-\sqrt f)\) has leading term \(24\sqrt3t^2\), and the product of suitable analytic branch coordinates has order seven. Therefore \(G\) meets component two, up to reversing the cyclic numbering, of \(I_7\). At infinity its transformed coordinates tend to the singular point, so it meets the nonidentity component of \(III\).

\item Shioda's intersection formula now gives \[
 h(G)=2-\frac{2(7-2)}7-\frac12=\frac1{14}.
\] The root rank is \(6+1=7\), so Shioda--Tate gives Mordell--Weil rank one. There is no torsion: the possible local contributions, multiplied by fourteen, are \(0,12,20,24\) from \(I_7\) and \(0,7\) from \(III\). Their sums never equal twenty-eight, and their maximum thirty-one is below fifty-six. Thus none can equal \(28+28R\cdot O\), as a nonzero torsion section would require. The unimodularity of the rational elliptic surface's Néron--Severi lattice then gives primitive height \(1/14\). Thus \(G\) is a generator.

\item For later use, exact addition gives \[
 x(7G)=-\frac{140t^3-609t^2-420t+92}{(t+10)^2}.
\] Its numerator is nonzero at \(t=-10\). The height formula gives \((7G)\cdot O=1\): its \(I_7\) contribution is zero and its \(III\) contribution is \(1/2\). The displayed pole therefore locates the whole intersection at the ordinary fibre \(t=-10\), with intersection multiplicity one. It is not at a reducible fibre.

\subsubsection{2. The linearization and every local filling}\label{the-linearization-and-every-local-filling}

\item Write \(R=7G\) and \(Q=2G\). Translation acts trivially on the Picard group of degree zero of a smooth elliptic fibre, so \[
 L=\operatorname{Hom}(t_Q^*M,M)
\] is vertical. At \(I_7\), \(R\) is in the identity component, hence \(M\) has degree zero on every component; translating by \(Q\) preserves this multidegree. At \(III\), \(Q\) is in the identity component and preserves every component. Thus \(L\) has zero degree on every fibre component. The kernel of the component intersection matrix is generated by the full fibre, so \(L=\pi^*\mathcal O_{\mathbb P^1}(d)\).

\item Intersecting its divisor \(R-O-(R-Q)+(-Q)\) with \(O\) gives \[
 d=1+1-1+0=1.
\] Here \((5G)\cdot O=1\) and \((2G)\cdot O=0\), obtained directly from the same height formula and the component classes \(2n\pmod7\), \(n\pmod2\). Consequently \(L=\mathcal O(F)\).

\item Choose a section of \(L\) whose sole zero is the full fibre at \(t_c=t_-\). It gives a lift of translation by \(-Q\) on \(M^\times\) away from that fibre. Multiply the lift by \(\lambda\), with \(0<|\lambda|\) uniformly small. This defines the smooth two-torus family on the regular base. The contraction/proper-discontinuity argument uses an invariant Hermitian norm: over a compact base set the lift is bounded, and its only degeneration is a zero, so scaling makes the norm strictly decrease. A compact annulus in the norm direction is a fundamental domain.

\item At the zero fibre \(I_1\), the two torus periods have valuation vectors \((1,0)\) and \((0,1)\). The unit-square triangulation along a diagonal has unimodular cones. The affine toric charts are \(q=z_0z_1z_2\), and the two periods identify the opposite sides of one hexagonal component. This gives a smooth total space and a proper one-component rank-two Mumford filling. Its reduced fibre normalizes to \(dP_6\), has three double curves and two triple points, and has Euler number two.

\item At the other \(I_1\), the exponent vectors are \((1,0)\) and \((0,0)\), with the zero-exponent period contracting by a scalar multiple of \(\lambda\). Its quotient first makes the second factor an elliptic curve. The rays \((j,1)\), \(j\in\mathbb Z\), give the standard smooth rank-one filling with equation \(xy=q\); one period shifts \(j\) by one. It has one component, an elliptic ruled normalization, and Euler number zero.

\item The \(I_7\) filling must retain both periods. The section \(R\) is in the identity component, whereas \(Q\) meets component four. In a Tate coordinate and a frame of \(M\) which has degree zero on all components, the two periods have the form \[
 h:(z,v)\longmapsto(q^7u_1z,u_2v),\qquad
 g:(z,v)\longmapsto(q^{-4}u_3z,\lambda u_4v),
\] with holomorphic nonzero units \(u_i\). The integral change of period generators \[
 H=h^{-1}g^{-2},\qquad K=h^4g^7
\] has determinant one. Their exponent vectors are \((1,0)\) and \((0,0)\), respectively. The second multiplier of \(K\) is \(\lambda^7\) times a unit, so it is uniformly contracting. The same ray fan therefore gives a \textbf{one-component primitive rank-one} filling. This explicit change of generators is essential; replacing the source \(I_7\) by seven independent components would give the wrong local attachment. Its charts, properness, Hausdorff property and Euler number are exactly those just described for the primitive rank-one filling. There is also a direct identification retaining the source elliptic map: at this fibre the chosen linearization is invertible on the whole original elliptic neighbourhood, so the filling is the quotient of its actual smooth total line complement \(M^\times\) by the contracting lift of \(-Q\). Locally at a source node its charts are \(xy=q\) times \(\mathbb C^*\), and translation permutes the seven components and their nodes cyclically. The displayed unimodular change of period generators identifies this quotient with the primitive toric model. Thus no elliptic component map has been discarded.

\item At \(III\), take the fourth-root base change. The relatively minimal elliptic surface has good reduction; the section \(R\) specializes to the nonzero two-torsion point and \(Q\) to zero. The pulled-back degree-zero line and its nonvanishing linearization therefore give a smooth family of two-tori across the root origin. The order-four linear deck action is the homology matrix specified below. Its affine translation is \(\ell/4\), where \[
 \ell=(1,-1,4,0)^t.
\] The invariant covector \(\psi\) evaluates to one on \(\ell\). Hence the \(j\)-th power of the action changes the invariant real circle coordinate by \(j/4\); it has no fixed point for \(1\leq j\leq3\). The quotient is smooth and has a reduced bielliptic central surface of multiplicity four.

\item Its punctured identification is the explicit logarithmic translation \[
 (x,s)\longmapsto
 \left(x+\frac{\log s}{2\pi i}\ell(s),s\right),\qquad q=s^4.
\] Changing \(\log s\) by \(2\pi i\) adds the integral period \(\ell\). Under the root deck action it changes by \(\ell/4\), proving descent and giving the inverse by the negative translation. These are exactly the charts used in the global gluing, with no extra affine clutch.

\item Use nested discs and annular overlaps to glue these four local models to the regular torus family. All changes of charts are the displayed holomorphic translations or monomial toric identifications. A finite cover by a compact regular-base set and four smaller closed discs proves compactness. The compact fundamental domains in the toric models prove Hausdorffness. The unimodular charts and the free finite action prove smoothness.

\subsubsection{3. Actual integral monodromy, not an abstract substitute tuple}\label{actual-integral-monodromy-not-an-abstract-substitute-tuple}

\item The elliptic monodromy can be marked as \[
 A=\begin{pmatrix}-5&26\\-1&5\end{pmatrix},\quad
 B=\begin{pmatrix}1&7\\0&1\end{pmatrix},\quad
 C=\begin{pmatrix}1&0\\-1&1\end{pmatrix},\quad
 D=\begin{pmatrix}-2&9\\-1&4\end{pmatrix},
 \qquad ABCD=I.
\] Here \(A^2=-I\); \(B\) is the seventh power of the positive transvection along \((1,0)\); \(C,D\) are positive transvections along \((0,1),(3,1)\). The positive base meridians are chosen so that this is the ordered cohomological product. Equivalently, the corresponding based path product is read in the reverse order. The choice is used consistently in every complex below.

\item We verify that this is a marking of the elliptic surface, rather than assuming it from the fibre labels. In \(PSL_2(\mathbb Z)=\langle S,R\mid S^2=R^3=1\rangle\), put \(T=SR\). The subgroup generated by \[
 T^5ST^{-5},\quad T^7,\quad STS^{-1},\quad T^3STS^{-1}T^{-3}
\] has the following exact coset permutations: \[
 S=(12)(34)(57)(6)(89),\qquad
 R=(123)(456)(789).
\] The product corresponding to \(T\) has cycles \((1,3,5,8,7,6,4),(2),(9)\). The complete nine-row coset table and its exact verification are in \texttt{os47\_topology.json}. It gives a degree-nine genus-zero modular cover, with ramification \(3^3\) over zero, \(2^4 1\) over one, and cusp widths \(7,1,1\). Miranda's Proposition 3.4 constructs its rational elliptic surface with fibres \(III,I_7,I_1,I_1\); all its numerical conditions hold, and the displayed permutations supply its required map. The local positive lifts give the matrices above. Fixing these four local fibre types removes quadratic-twist ambiguity: an odd valuation of a twist changes the local type, while a rational function on \(\mathbb P^1\) with every valuation even is a square over \(\mathbb C\).

\item For completeness this surface is the explicit Weierstrass model of Section 1. Place \(I_7\) at zero and \(III\) at infinity. After scaling, write \(f=1+at+bt^2+ct^3\); the degree-four polynomial \(g\) is the Taylor polynomial of \(f^{3/2}\). The coefficients of degrees five and six must vanish. Solving these exact two equations gives either \(b=a^2/4,c=0\), for which the discriminant vanishes identically, or \(b=5a^2/8,c=3a^3/32\). The nondegenerate case has \(a\ne0\), and rescaling \(t\) makes \(a=8\), giving precisely Section 1. Thus the modular construction and the displayed algebraic source are related by an actual Weierstrass change of coordinates.

\item We next determine the section extension in this marking. Let \(k_i\) be the integral Abel--Jacobi cocycle of \(G\). Its class at \(I_7\) is two modulo seven; reverse \(G\) if necessary. A change of its integral Abel--Jacobi lift first sets \(k_B=(2,0)^t\), then sets \(k_D=0\). Indeed the first changes by \((7x_2,0)^t\), while the coefficient of \(k_D\) along \((3,1)^t\) changes by \(-x_1+3x_2\). These two choices can therefore be made successively with integral \(x_1,x_2\). Write \(k_C=(0,j)^t\). The sphere relation then forces \[
 k_A=-A(k_B+Bk_C)=(9j+10,2j+2)^t.
\]

\item In the cohomology basis \((\psi,u,w,\delta)\), a section cocycle \(k=(k_1,k_2)^t\), elliptic matrix \(E=\left(\begin{smallmatrix}a&b\\c&d\end{smallmatrix}\right)\), line weight seven and translation weight two give the integral lift \[
 T(E,k,z)=
 \begin{pmatrix}
 1&2(ck_1-ak_2)&2(dk_1-bk_2)&z\\
 0&a&b&7k_1\\0&c&d&7k_2\\0&0&0&1
 \end{pmatrix}.
\] This formula follows by changing the Abel--Jacobi representative in the two period columns. It also preserves \(\omega=7u\wedge w+2\psi\wedge\delta\), as direct multiplication checks. At \(B\) the rank-one condition gives \(z_B=8\); at \(C\) it gives \(z_C=14j^2\). At \(A\), the fourth-power identity gives \(z_A=7(5j^2+8j+4)\). At the scalar-zero \(I_1\), \(k_D=0\) and its period multiplier \(\lambda q\) gives \(z_D=-1\). These signs use the positive circle in the line as the dual homology class \(\widehat\delta\).

\item The only possibly nonzero entry of the product minus the identity is now the upper-right entry, equal to \[
 7(j+1)(7j+5).
\] Its only integral zero is \(j=-1\). This determines the actual geometric lift uniquely in the normalized marking, without selecting it on the basis of a desired homology answer. The four matrices are \[
 T_A=\begin{pmatrix}1&-2&10&7\\0&-5&26&7\\0&-1&5&0\\0&0&0&1\end{pmatrix},\quad
 T_B=\begin{pmatrix}1&0&4&8\\0&1&7&14\\0&0&1&0\\0&0&0&1\end{pmatrix},
\] \[
 T_C=\begin{pmatrix}1&2&0&14\\0&1&0&0\\0&-1&1&-7\\0&0&0&1\end{pmatrix},\quad
 T_D=\begin{pmatrix}1&0&0&-1\\0&-2&9&0\\0&-1&4&0\\0&0&0&1\end{pmatrix}.
\] Their homology transports are \(\widehat T_i=(T_i^{-1})^t\). The vector \(\ell=(1,-1,4,0)^t\) is fixed by \(\widehat T_A\).

\subsubsection{4. Fundamental group and every integral cohomology differential}\label{fundamental-group-and-every-integral-cohomology-differential}

\item The regular family has a section. One can choose it from a meromorphic frame of \(M|_O\simeq\mathcal O(2)\) whose divisor is twice the selected scalar-zero fibre. At the two rank-one models its valuation is zero. At the rank-two model its valuation is an integral period vector and is moved to zero by a deck period, so it extends. The good-reduction comparison at \(III\) is chosen with this section as punctured origin before the displayed logarithmic translation. This is the allowed translation identification in the finite-cap construction.

\item The three Mumford fillings kill their section meridians and their primitive vanishing lattices. These lattices are exactly the images of \(\widehat T_i-I\), \(i=B,C,D\): the two rank-one matrices have primitive rank-one image and the rank-two matrix has Smith factors \((1,1)\). The primitive fans in Section 2 give the corresponding collapsing circles, with no index omitted. The sum of these three images has Smith factors \((1,1,1)\) and quotient generated by \(\widehat\psi\). All boundary maps are onto their filled fundamental groups: in a primitive rank-one chart the pinched circle is the only new kernel; in the unimodular rank-two chart the two angular circles are the only new kernels. The finite quotient has generators from the torus and its deck meridian.

\item The base sphere relation kills the remaining finite meridian. Its fourth-power relation then kills \(\ell\), whose image in the remaining fibre group is \(\widehat\psi\). Hence \(\pi_1(X)=1\). The exact integer matrices consisting of these twelve vanishing-image columns and then the additional column \(\ell\) are saved in the program output. Their Smith forms are respectively \((1,1,1)\) and \((1,1,1,1)\); this also checks the integral abelian calculation.

\item We give the full cohomology calculation. In degree \(q\), put \(L_q=\bigwedge^q\mathbb Z^4\), \(E_i=\bigwedge^qT_i\), using the lexicographic exterior basis. The parabolic cellular complex is \[
 L_q\xrightarrow{d^0}\bigoplus_i(E_i-I)L_q
       \xrightarrow{d^1}L_q,
\] \[
 d^0x=((E_i-I)x)_i,\qquad
 d^1(y_A,y_B,y_C,y_D)=y_A+E_Ay_B+E_AE_By_C+E_AE_BE_Cy_D.
\] Cutting the four-punctured sphere along its marked meridians and attaching each disc with its invariant stalk gives this complex integrally. Local image bases are computed by integer Hermite normal form. The program saves every entry of \(d^0,d^1\), checks their product is zero, and computes their Smith forms.

\item The ranks of the four local images in degrees one, two and three are \((2,1,1,2)\), \((4,2,2,2)\), and \((2,1,1,2)\). In those degrees both differentials have ranks three, five and three, respectively. The stacked ambient matrix \(((E_i-I)^t)_i^t\) has only unit nonzero Smith factors. Therefore the middle cohomology is both torsion-free and rationally zero, hence zero. The last cokernels are free of rank one. The invariant generators in degrees zero through four are \[
 1,\quad\psi,\quad\omega=7uw+2\psi\delta,\quad
 \phi=\psi uw,\quad\nu=\psi uw\delta.
\] The coinvariant generators are \[
 1,\quad[\delta],\quad\eta,\quad[uw\delta],\quad[\nu],
 \qquad[\psi\delta]=7\eta,\quad[uw]=2\eta.
\] All omitted exterior basis classes map to zero in these quotients.

\item The primitive rank-one collapse is the cellular map \(E\times T^2\to E\times(S^2\vee S^1)\), killing one primitive one-cell, carrying the other one-cell with degree one and carrying the two-cell with degree one. It therefore identifies integral specialization with the full invariant lattice in every degree. The rank-two local fan is the unimodular one in E26, Lemma 7.4; its normalization-and-double-curve cellular calculation applies verbatim because both period vectors form an integral basis. Its hypotheses are exactly the checked valuation vectors \((1,0),(0,1)\), the same triangulation and unit changes of coordinates. Its top specialization has degree one off the double curves. Thus all three Mumford specializations are integral isomorphisms onto the invariant lattice.

\item Here is a separate integral calculation at the finite cap. Let \(B_0\) denote its reduced bielliptic surface and \(p:T^4\to B_0\) the degree-four covering. Abelianizing its deck presentation gives the five-generator relation matrix \[
 \begin{pmatrix}\widehat T_A-I&-\ell\\0&4\end{pmatrix},
\] whose Smith factors are \((1,1,1)\). Thus \(H_1(B_0)=\mathbb Z^2\). Poincaré duality and the universal coefficient theorem imply that its cohomology is torsion-free. Its Euler number is zero, giving \(b_2(B_0)=2\). Its canonical bundle is torsion (the fourth power of the fibre volume form descends); since \(H^2(B_0)\) is torsion-free, \(c_1(K_{B_0})=0\). Its unimodular intersection lattice is therefore even.

\item The invariant first-cohomology basis is \[
 \psi,\qquad h=-28u-7w+2\delta.
\] Because \(\psi(\ell)=1\), \(h(\ell)=0\), its pullback image is \(\mathbb Z(4\psi)\oplus\mathbb Zh\). In degree two, the full invariant lattice has basis \[
 h_2=\psi h,\qquad v_2=uw+4\psi u+\psi w,
\] intersection matrix \(\left(\begin{smallmatrix}0&2\\2&0\end{smallmatrix}\right)\), and \(\omega=h_2+7v_2\). Since the covering has degree four, the pullback of the downstairs unimodular lattice has determinant \(4^2=16\), so its index in this invariant lattice is two. An index-two sublattice whose divided-by-four intersection is even cannot contain a vector with both coordinates odd: that vector has square \(4ab\), with \(a,b\) odd, and would have odd square downstairs. Hence \(\omega\) does not descend and \(2\omega\) does.

\item An invariant third-cohomology basis is \(\phi,\,4\psi u\delta+\psi w\delta+uw\delta\). Its pairing with \((\psi,h)\) has matrix \(\left(\begin{smallmatrix}0&1\\-2&0\end{smallmatrix}\right)\). The first-degree pullback index is four, so perfect-pairing comparison gives third-degree index two. The class \(\phi\) cannot descend, because its pairing with the descended class \(h\) is \(-2\), which is not divisible by four; hence the global generator is \(2\phi\). In top degree the index is four. Thus the four global specialization indices are precisely \((4,2,2,4)\). Each global invariant generator surjects onto the corresponding local cokernel. There are no extra first-cohomology terms from these skyscraper cokernels.

\item The resulting Leray page has only columns zero and two. Its column zero in degrees zero through four is generated by \[
 1,\quad4\psi,\quad2\omega,\quad2\phi,\quad4\nu;
\] column two has the coinvariant generators listed above, multiplied by the positive base class \(H\). The only clutch is the order-four logarithmic translation. We justify the absence of an intrinsic differential as follows. Replace that single affine cap temporarily by the coarse quotient of the same root torus family under its linear order-four action. This is a specified compact topological filling; smoothness of this auxiliary quotient is not needed. Its rational fibre cohomology is the invariant lattice tensored with \(\mathbb Q\). Indeed an equivariant barycentric subdivision identifies the quotient cochains with invariant cochains, and averaging is exact over \(\mathbb Q\). Inversion about the global origin commutes with the linear cap and extends over the centrally symmetric toric fans. It acts by \((-1)^q\) in degree \(q\), so naturality forces its rational \(d_2\) to vanish. On the boundary, the linear and affine mapping tori are identified by the explicit translation path \(z\mapsto z+t\ell/4\). Four lifted copies of its prism give the integral translation prism for \(\ell\); rational transfer therefore divides that contraction by four. This computes the rational differential of the actual affine cap. Since the source and target groups already computed for the actual family are free, that rational computation determines the integral differential on the stated specialization lattice uniquely.

\item A cellular translation prism through a cycle \(z\) induces contraction \(\iota_z\) on exterior cochains. Going once around the base changes the root logarithm by one quarter of \(\ell\); consequently \[
 d_2(a)=\frac H4[\iota_\ell a].
\] This formula is integral on the stated specialization lattice. The exact contractions are \[
 [\iota_\ell\psi]=1,\quad
 [\iota_\ell\omega]=2[\delta],\quad
 [\iota_\ell\phi]=2\eta,\quad
 [\iota_\ell\nu]=[uw\delta].
\] Thus all four differentials are multiplication by one. There are no higher differentials because the base has dimension two. The sole surviving groups are \(H^0(X)=H^6(X)=\mathbb Z\), with no extension ambiguities. Poincaré duality gives the same integral sphere homology. Together with simple connectivity, Hurewicz and Whitehead identify \(X\) as a homotopy six-sphere. The group \(\Theta _6\) is zero (Kervaire--Milnor, the table on printed p.505), so it is diffeomorphic to the standard \(S^6\).

\subsubsection{5. The punctured holomorphic family and its closing map}\label{the-punctured-holomorphic-family-and-its-closing-map}

\item The local constructions depend holomorphically on \(\lambda\) over the universal logarithm half-plane, with a common bound on its absolute value. Put \(c_{47}=(\log\lambda)/(2\pi i)\), so \(c_{47}\mapsto c_{47}+4\) adds \(4\delta\) to the period \(\gamma\) and to \(\ell\). At the finite cap \(q=1/t=s^4\), the logarithmic gluing translation changes by \((\log q)/(2\pi i)\delta\).

\item The ordinary-fibre multiplication by \[
 h(t)=t-t_c
\] is the closing map from \(c_{47}+4\) back to \(c_{47}\). Equivalently the positive deck from \(c_{47}\) to \(c_{47}+4\) multiplies by \(h^{-1}\). The direction follows directly from the displayed logarithmic gluing: its change is \(\log(q)/(2\pi i)\delta\), so the positive ordinary deck has that same logarithmic coefficient at the finite cap. There \(h^{-1}\) equals \(q\) times the inverse of the unit \(1-t_cq\). Thus it cancels the change when passing back to the finite-chart coordinate; the remaining unit acts holomorphically on the filled torus quotient. At the chosen scalar-zero cusp, it equals the cusp parameter \(q_c\). Its period generator has source translation \(-Q\) and line multiplier \(\lambda q_c\) times a unit. Modulo that period, multiplication by \(q_c\) is source translation by \(+Q\) and multiplication by \(\lambda^{-1}\) times a unit. These extend on the resolved cusp, since \(Q\) has an invertible Tate representative at this \(I_1\) fibre. At the other two fillings \(h\) is a unit and extends directly. The inverse has the reverse descriptions, so these are genuine biholomorphisms. They agree on every annular overlap by the explicit logarithmic equation.

\item Therefore the family closes holomorphically under \(c_{47}\mapsto c_{47}+4\). Its quotient parameter is \(r=e^{2\pi ic_{47}/4}\), with \(\lambda=r^4\). The action is properly discontinuous on the parameter half-plane. The same compact local fundamental domains prove that the quotient family is proper over every closed annulus in the punctured disc. All its local total-space charts are smooth over \(r\), and the finite actions are free, so this is a proper holomorphic submersion with the sphere fibres just proved. Its positive parameter loop is the positive \(r\)-loop. This explicit closing map fixes which family monodromy is being investigated.

\subsubsection{6. The actual page and its section line}\label{the-actual-page-and-its-section-line}

\item The line circle is free on both primitive rank-one fillings, including all the nodes of the original \(I_7\). Here is a verification on the actual line complement, not just its open torus. Choose a Hermitian norm on \(M\) over a compact elliptic neighbourhood where the linearization is invertible. For sufficiently small \(|\lambda|\), the map \(g=\lambda\widetilde t_{-Q}\) has norm ratio strictly below one everywhere. If \(e^{i\theta}v=g^n(v)\), the equality of norms forces \(n=0\), also for negative \(n\) by applying \(g^{-n}\). Then \(e^{i\theta}=1\). This proves freeness through every source node and every component permutation in the contracting quotient. The order-four cap has no extra stabilizers, because its invariant circle displacement is \(j/4\) and the line circle has zero displacement in that coordinate. The sole fixed locus is the standard toric \(\mathbb P^1\) at the rank-two cusp, with normal \(\mathcal O(-1)\oplus\mathcal O(-1)\) and weights \((-1,+1)\). Thus the quotient and conifold construction of P3F.7--8 applies to this actual sphere family.

\item The quotient base is the Seifert \(S^3\) of weights \((1,4)\), because there is one order-four exceptional fibre with primitive invariant translation. A regular binding away from all special values has a disc page. Its compactification covers the original base with degree four, branched at \(III\) and at that regular binding. The \(III\) fibre has good reduction; the other three values each have four unramified preimages. The source \(I_7\) fibres are retained in the elliptic page: the rank-one quotient was made with the actual component-permuting translation \(Q\), and fixing its Seifert circle coordinate recovers the source elliptic fibre. More explicitly, quotienting the actual line complement by its circle gives a positive real line bundle over the source elliptic neighbourhood. It is trivial there because its Chern class has degree zero on every component. The remaining generator acts as \(-Q\) with a strictly positive (or strictly negative) real displacement, uniformly after scaling. Rescaling a smooth collar coordinate in its fundamental domain gives the mapping torus of \(-Q\), over the elliptic base. Its seven critical points form a curve covering the suspension circle seven times. A fixed suspension coordinate therefore retains the entire \(I_7\) elliptic fibre. Consequently the page has four \(I_7\) fibres and eight \(I_1\) fibres, with Euler number \(4\cdot7+8=36\). Its elliptic completion has \(\chi(\mathcal O)=3\) and \(\sigma=-24\).

\item The page line's Picard point is \(R=7G\), since the original line is \(M\); the translation point \(Q=2G\) is the distinct suspension period. Height multiplies by the cover degree, giving \(h(R)=14\). It lies in the identity component of every \(I_7\), so all local height corrections are zero. The formula \(h(R)=2\chi+2R\cdot O\) gives \(R\cdot O=4\). Equivalently these are the four unramified pullbacks of the already located ordinary intersection at \(t=-10\). Therefore \[
 c_1(M)^2=-2\chi-2R\cdot O=-14.
\] Choose the page phase disjoint from the fixed sphere. The argument of P3F.10 identifies the actual circle line with this Picard point on the affine torsor, with any difference pulled back from the base with boundary. Its boundary Chern class is zero; the relative cup-product identity there makes its relative square independent of the chosen boundary trivialization. The same relative tangent comparison gives \(p_{\rm rel}=3\sigma=-72\). This proves the page expression four.

\item The actual return-map comparison is now supplied by the real trace P3K.6--7, with the positive parameter direction fixed in P3F.23. The two orientation reversals are written explicitly in P3F.25. Consequently this page expression gives the signed return integer A=-4, and P3F.22 gives its oriented class24 modulo28. This conclusion uses that actual morphism in addition to the topology and page arithmetic above.

\subsubsection{Executed verification}\label{executed-verification}

\path{code/part3/os47_topology.py} has executed successfully. Its JSON file contains the modular coset table, the complete integral Jacobi tuple, the symbolic uniqueness equation for the geometric section lift, the fundamental-group relation matrices, every exterior-power local image basis, every entry of the parabolic differentials, the finite specialization indices and all four integral clutch contractions. The printed output records \(d_2=(1,1,1,1)\), sphere homology and the page expression four. No prior-run code or proof is imported.

\end{enumerate}
\end{proof}




\begin{theorem}\label{night:P3F.22}
\StatementRecord{night:P3F.22}{proved}{S6\_06\_BLIND: GPT-6 Astra (OpenAI, ChatGPT Pro mode, effort Pro), recorded re-derivation in the checkpoint dated 18 November 2026}{\cref{night:P3F.4}}
The specified Hopf gluing has Eells--Kuiper invariant $\mu(\lambda(f))=A(f)/28$. It is a homomorphism to $\Theta_7$, and $\lambda(b)$ is the generator with $\mu=1/28$.
\end{theorem}
\begin{proof}
\begin{enumerate}[label=\arabic*.]
\item Use the oriented Hopf bundle with Euler class the positive generator of the oriented two-sphere. Orient its total three-sphere by the circle followed by the base, and orient the four-ball which it bounds accordingly. Use the product orientation on each four-ball and the outward-normal-first convention on \(\partial(D^4_u\times D^4_y)\). For the gluing convention \[
 \lambda(f)=(D^4_u\times S^3_y)\cup_{\widetilde f}
                  (S^3_u\times D^4_y),
 \qquad x\in\partial(D^4_u\times S^3_y)
       \sim\widetilde f(x),
\] where \(\widetilde f\) is a principal-circle lift of \(f\in I(S^2\times S^3)\), the oriented Eells--Kuiper invariant is \[
 \mu(\lambda(f))=\frac{A(f)}{28}\pmod{\mathbb Z},
 \qquad A(f)=\frac{4k(f)-\ell(f)}{24}.                 \tag{P3F.22.1}
\] Consequently \(\lambda\) is a homomorphism to \(\Theta _7\), \(\lambda(a)=0\), and \(\lambda(b)\) is the generator with \(\mu=1/28\). Its kernel is \(\mathbb Za\oplus28\mathbb Zb\).

\item The assertion is a transfer computation for the specified gluing map. An assignment of its value to a parameter loop also requires that loop's actual oriented return map. In particular, this proposition by itself does not select the sign of the OS47 parameter loop.

\subsubsection{1. The marked double and its bordism}\label{the-marked-double-and-its-bordism}

\item Put \(K=S^2\times D^4\), with boundary \(M=S^2\times S^3\), and form \[
 N_f=K_+\cup_f(-K_-),\qquad x\in\partial K_+\sim f(x)\in\partial K_-.
\] Let \(x\in H^2(N_f;\mathbb Z)\) restrict to the positive generator on both copies of \(S^2\). Mayer--Vietoris and van Kampen give \[
 \pi_1(N_f)=0,\qquad H^j(N_f;\mathbb Z)=
 \begin{cases}\mathbb Z&j=0,2,4,6,\\0&\text{otherwise}.
 \end{cases}
\] The two spin structures on the pieces agree on the simply connected boundary, so \(N_f\) is spin. There is a unique \(y\in H^4(N_f;\mathbb Z)\) with \(\langle xy,[N_f]\rangle=1\).

\item For clarity the required bordism is an actual collar construction, the one in Jin26v2 Lemma6.2. Write an inward collar of \(K\) as \(M\times[0,1]_s\), and identify the two collars at the ends of \(K\times[0,1]_t\) by \((z,s,0)\sim(f(z),s,1)\). Rounding corners gives a seven-manifold. Its \(s=0\) boundary is the mapping torus in the convention \((z,0)\sim(f(z),1)\); its other boundary is \(-N_f\). Indeed, near that collar the orientation is \(\operatorname{or}(M)\,ds\,dt\); contraction with the outward normal \(-\partial_s\) gives \(\operatorname{or}(M)\,dt\). At \(t=0,1\) the other boundary pieces are \(-K\) and \(+K\), respectively, with exactly the displayed identification. The class \(x\) extends, as does the preferred spin structure (choose the appropriate one of the two spin structures of the bordism). This verifies the direction and the orientation of the comparison, not only equality of the unoriented manifolds.

\item It follows that \[
 x^2=k(f)y,\qquad p_1(TN_f)=\ell(f)y.                 \tag{P3F.22.2}
\]

\subsubsection{2. Removal of the cubic coordinate by an extending map}\label{removal-of-the-cubic-coordinate-by-an-extending-map}

\item Set \(f_0=f\circ a^{-k(f)}\). Jin's additive integral coordinates give \[
 k(f_0)=0,\qquad \ell(f_0)=\ell(f)-4k(f)=-24A(f).
\] The explicit quaternionic map in P3F.4 extends every lift of \(a^n\) over \(D^4_u\times S^3_y\). Apply that extension on the first gluing piece and the identity on the second. With the convention above it identifies \(\lambda(f_0)\) and \(\lambda(f)\) by an orientation-preserving diffeomorphism. This step does not assume in advance that \(\lambda\) is a group homomorphism.

\item We can therefore suppose \(k=0\). Write \(p=\ell(f_0)\), so that \(x^2=0\) and \(p_1(TN_{f_0})=py\).

\subsubsection{3. The surgery and its bounding eight-manifold}\label{the-surgery-and-its-bounding-eight-manifold}

\item Let \(L\to N_{f_0}\) have first Chern class \(x\), let \(E=S(L)\), and let \(D(L)\) be its oriented disc bundle. The preimage of the first \(K=S^2\times D^4\) is \(S^3\times D^4\). Perform the framed surgery which replaces that piece by \(D^4\times S^3\). Its framing is the product framing of the \(D^4\) factor, so the attaching map is completely specified. The result is a homotopy sphere, either by its explicit identification below or by the elementary Mayer--Vietoris computation in P3F.4.

\item The orientation is relevant. The oriented circle bundle decomposes as \[
 E=(S^3\times D^4)_+\cup_{\widetilde f_0}
                                      -(S^3\times D^4)_-.
\] The boundary of \(S^3\times D^4\) has orientation \(-\operatorname{or}(S^3\times S^3)\), whereas the boundary of \(D^4\times S^3\) has the positive product orientation. Thus the replacement has orientation \(-(D^4\times S^3)\). The oriented surgery sphere is exactly \(-\lambda(f_0)\).

\item An eight-manifold bounding that surgery sphere is \[
 W=D(L)\cup (D^4\times D^4),
\] where the four-handle attaches along the specified \(S^3\times D^4\). Equivalently it is obtained by replacing the Hopf-disc neighbourhood of \(S^2\) by \(\mathbb CP^2\times D^4\). This is the explicit construction in Jiang14, proof of Lemma3.5, Case2. It has an embedded \(\mathbb CP^2\) with trivial normal four-plane and a second four-cycle coming from the generator of \(H_4(N_{f_0})\). Their intersection matrix is \[
 \begin{pmatrix}0&1\\1&0\end{pmatrix}.              \tag{P3F.22.3}
\] The first self-intersection is zero by the trivial normal bundle; the second is zero because the disc-bundle Thom map \(H^2(N_{f_0})\to H^4(N_{f_0})\) is multiplication by \(x\), and \(x^2=0\). The mixed intersection is one because the core \(S^2\) is dual to \(y\). These facts also give a direct check of the intersection matrix in the cited proof.

\item Let \(z\in H^2(W;\mathbb Z)\) be the class extending \(x\), positive on the projective line in that \(\mathbb CP^2\). It lifts \(w_2(TW)\): on \(D(L)\) the latter is \(x\bmod2\), and on \(\mathbb CP^2\times D^4\) it is the hyperplane class modulo two. Let \(u,v\) be the relative degree-four classes dual to the two cycles in \(P3F.22.3\). The restrictions to those two cycles give \[
 z^2=jv,\qquad p_1(TW)=p\,ju+3jv,
\] where \(j:H^4(W,\partial W)\to H^4(W)\) is the natural map. Indeed the restrictions of \(z^2\) are \((1,0)\), and those of \(p_1\) are \((3,p)\). Hence \[
 \sigma(W)=0,\quad
 \langle p_1^2,[W,\partial W]\rangle=6p,\quad
 \langle z^2p_1,[W,\partial W]\rangle=p,\quad
 \langle z^4,[W,\partial W]\rangle=0.              \tag{P3F.22.4}
\]

\subsubsection{4. The invariant and the resulting homomorphism}\label{the-invariant-and-the-resulting-homomorphism}

\item For a spin boundary homotopy seven-sphere and this characteristic extension \(z\), the spin-c Eells--Kuiper formula is \[
 \mu(\partial W)=
 \frac{p_1^2-4\sigma(W)}{896}
 -\frac{z^2p_1}{192}+\frac{z^4}{384}\pmod{\mathbb Z}.
\] The formula is Jiang14 Section3.1\(v\), with its stated hypotheses; here \(\partial W\) is a homotopy sphere, so the required rational relative-to-absolute degree-four map is an isomorphism. Its normalization also follows by subtracting the spin-c index \(\int e^{z/2}\widehat A\) from the signature polynomial. Substitution of \(P3F.22.4\) gives \[
 \mu(-\lambda(f_0))=
 \frac{6p}{896}-\frac{p}{192}=\frac{p}{672}.
\] Since \(p=-24A(f)\), reversal of the boundary orientation proves \(P3F.22.1\).

\item The Eells--Kuiper invariant is additive under connected sum and identifies \(\Theta _7\) with \(\frac1{28}\mathbb Z/\mathbb Z\). For the positive normalization one can also check the Milnor sphere \(S_{2,-1}(\mathbb H)\): its disc bundle has Euler number one, Pontryagin number six and signature one, giving \((36-4)/896=1/28\). Davis82 Section3 computes the same value. Thus (P3F.22.1), together with Jin's additivity of \(A\), proves that \(\lambda\) is a homomorphism and gives its exact image and kernel. No classification statement about circle actions is needed for this last deduction.

\end{enumerate}
\end{proof}




\begin{theorem}\label{night:P3F.23}
\StatementRecord{night:P3F.23}{proved}{S6\_06\_BLIND: GPT-6 Astra (OpenAI, ChatGPT Pro mode, effort Pro), recorded re-derivation in the checkpoint dated 18 November 2026}{\cref{night:P3F.21}}
The positive OS47 $r$-loop has angular clutch $\theta\gamma+2\theta(\theta-1)\delta$, hence one positive gamma turn; its fourth-root cusp deck gauge has winding positive four.
\end{theorem}
\begin{proof}
\begin{enumerate}[label=\arabic*.]
\item For the explicitly closed OS47 family of P3F.21, put \(\lambda=r^4\) and \(v=\log r/(2\pi i)\). In the fixed line character \(V=\exp(2\pi i\zeta _2)\), its positive \(r\)-loop has angular clutch \[
 F(s,v+\theta)-F(s,v)
    =\theta\gamma(s,v)+2\theta(\theta-1)\delta,
 \qquad 0\leq\theta\leq1.                         \tag{P3F.23.1}
\] Thus the rotation in the circle quotient is one positive gamma turn. The remaining line gauge is periodic and has value one at both endpoints. On the fourth-root cusp deck its winding is positive four.

\item \textbf{Proof.} Use the cusp parameter \(q=t-t_c\), and set \(s=\log q/(2\pi i)\). Both the line section \(R=7G\) and the translation point \(Q=2G\) meet the smooth identity component at this \(I_1\) fibre. In a Tate coordinate and an invertible local line frame the suspension period is consequently \[
 \gamma(s,v)=\binom{\Gamma(q)}{C(q)+s+4v},\qquad
 \delta=\binom01,
\] where \(\Gamma\) is a holomorphic logarithmic representative of \(-Q\), and \(C\) is holomorphic. The sign in the second coordinate comes from the actual multiplier \(\lambda q\): exponentiating gives \(e^{2\pi iC}q r^4\).

\item The positive parameter deck must be distinguished from its inverse closing isomorphism. At the fourth-root cap, the logarithmic gluing is \(z\mapsto z+\log(s_\infty)/(2\pi i)\ell(c_{47})\), and \(\ell(c_{47}+4)-\ell(c_{47})=4\delta\). Therefore the ordinary positive deck translates by \(-\log(t-t_c)/(2\pi i)\delta\), since its finite-cap expansion is \(+\log q_\infty/(2\pi i)\delta\) plus a holomorphic unit logarithm. In the chosen cusp coordinate this is exactly \[
 (v,\zeta)\longmapsto(v+1,\zeta-s\delta).          \tag{P3F.23.2}
\] Multiplication by \(h(t)=t-t_c\) is the inverse map, from \(c_{47}+4\) to \(c_{47}\). Both directions extend by the local argument in P3F.21. Parameter-independent changes of the ordinary frame add the same term to the two origins below and cancel in their difference.

\item On \(q=Q_0^4\), take the two explicit origins \[
 o(s,v)=\binom{-\Gamma s/4}{-Cs/4-s^2/8-sv},\qquad
 \sigma(s,v)=\binom{\Gamma v}{Cv+2v^2-2v}.          \tag{P3F.23.3}
\] Their deck differences are \[
 \begin{aligned}
 o(s,v+1)-o(s,v)&=-s\delta,\\
 \sigma(s,v+1)-\sigma(s,v)&=-s\delta+\gamma(s,v),\\
 o(s+4,v)-o(s,v)&=-\gamma(s,v)-2\delta,\\
 \sigma(s+4,v)-\sigma(s,v)&=0.
 \end{aligned}                                   \tag{P3F.23.4}
\] Each discrepancy beyond \(P3F.23.2\) is an actual integral lattice period. Hence these origins descend to sections of the specified torus torsor, with the cyclic cusp cover retained. The cusp origin \(\sigma\) is holomorphic and bounded in \(q\) for each \(r\ne0\); its characters act by units in the cusp toric charts, so it extends there. Subtracting \(o\) removes the parameter deck. Over any fixed nonzero \(q\), the remaining period multiplier is \(e^{2\pi iC}q r^4\), and the standard four-node rank-one chain \(X_i=V/r^i,\ Y_i=r^{i+1}/V,\ X_iY_i=r\) supplies the partial parameter filling. Its zero-coordinate origin extends. Thus these are the extending origins for the two actual collars, not arbitrary polynomials satisfying only a finite-difference identity.

\item Their difference is \[
 F=\sigma-o=
 \binom{\Gamma(s+4v)/4}
       {C(s+4v)/4+(s+4v)^2/8-2v}.
\] Direct subtraction proves \(P3F.23.1\). The additional real delta coefficient vanishes at \(\theta=0,1\), so its exponentiation is a periodic unitary gauge homotopic to the constant gauge through periodic gauges. It vanishes in the circle quotient. Furthermore \[
 \begin{split}
 &[F(s+4,v+\theta)-F(s+4,v)]
       -[F(s,v+\theta)-F(s,v)]\\
 &\hspace{35mm}=4\theta\delta.
 \end{split}
\] The corresponding transition is \(e^{8\pi i\theta}\), with winding four. This is the retained line clutch for the actual positive parameter loop. All equalities, including both endpoint values, are independently executed in \path{code/part3/os47_origins.py}.

\end{enumerate}
\end{proof}




\begin{theorem}\label{night:P3F.24}
\StatementRecord{night:P3F.24}{proved}{S6\_06\_BLIND: GPT-6 Astra (OpenAI, ChatGPT Pro mode, effort Pro), recorded re-derivation in the checkpoint dated 18 November 2026}{Definitions and the displayed constructions}
The forward sphere mapping torus is $(S^1\times S^6)\#\Sigma_\Phi$; both spin structures have Eells--Kuiper invariant $\mu(\Sigma_\Phi)$. A nonzero value excludes even an unoriented diffeomorphism with the product.
\end{theorem}
\begin{proof}
\begin{enumerate}[label=\arabic*.]
\item Let \(\Phi\) be an orientation-preserving diffeomorphism of \(S^6\). Write \(\Sigma_\Phi=D^7\cup_\Phi(-D^7)\), with the direction of the boundary identification fixed from the first disc to the second. Its forward mapping torus satisfies \[
 T^+_\Phi=(S^6\times[0,1])/((z,1)\sim(\Phi(z),0))
       \cong(S^1\times S^6)\mathbin{\#}\Sigma_\Phi.   \tag{P3F.24.1}
\] Its reverse mapping torus, the convention used for the Jin six-dimensional mapping torus, satisfies \[
 T^-_\Phi=(S^6\times[0,1])/((z,0)\sim(\Phi(z),1))
       \cong(S^1\times S^6)\mathbin{\#}(-\Sigma_\Phi).\tag{P3F.24.2}
\] These are orientation-preserving diffeomorphisms for the indicated product and gluing orientations. For both spin structures on each mapping torus, its Eells--Kuiper invariant is respectively \(+\mu(\Sigma_\Phi)\) and \(-\mu(\Sigma_\Phi)\). In particular a nonzero such invariant excludes even an unoriented diffeomorphism with \(S^1\times S^6\).

\item \textbf{Proof.} First isotope \(\Phi\) to be the identity on a small oriented disc. The needed isotopy of disc embeddings is explicit: shrink the parametrized image disc toward its centre, move its centre along a path, move its oriented frame along the oriented frame bundle, and expand to the chosen target disc. Compactness gives a common sufficiently small radius along the path. Isotopy extension carries this isotopy to the ambient sphere. The result is a diffeomorphism equal to the identity on the target disc, and so supported in the complementary disc \(D^6\). Mapping tori of isotopic maps are diffeomorphic, by applying the isotopy along the interval; the same is true for their boundary-clutched spheres by extension over a collar. Thus this modification loses no data.

\item In the standard product \(S^1\times S^6\), choose a seam of the circle and the seven-ball \(B=D^6\times[-\epsilon,\epsilon]\) crossing that seam; round its corners. Its upper face has the positive orientation of \(D^6\), because the latter has even dimension. Define a boundary diffeomorphism of \(B\) to equal \(\Phi|_{D^6}\) on this face and the identity on the other face and the cylindrical collar. It is smooth since \(\Phi\) is the identity near \(\partial D^6\). Regluing this ball into its complement changes precisely the upper-to-lower seam identification to \(z\mapsto\Phi(z)\), and hence produces \(T^+_\Phi\). The boundary diffeomorphism, as an oriented map of \(\partial B\), is the original disc-supported \(\Phi\): both are its extension by the identity to the complementary six-disc.

\item The same regluing is connected sum with \(\Sigma_\Phi\). To check the map, remove an interior seven-ball from the second disc in \(D^7\cup_\Phi(-D^7)\). That remaining second piece is a collar \(S^6\times[0,1]\). Inserting this collar in the ball removed from \(S^1\times S^6\) gives exactly the specified boundary identification \(\Phi\). This proves (P3F.24.1) with its gluing direction. Replacing \(\Phi\) by \(\Phi^{-1}\) gives the reverse mapping torus and the reversed clutched sphere, proving \(P3F.24.2\).

\item There are exactly two spin structures because \(H^1(S^1\times S^6;\mathbb Z/2)=\mathbb Z/2\), and connected sum with a homotopy sphere does not change this group. Both product spin structures extend over \(S^1\times D^7\). That explicit eight-manifold has zero signature and zero first Pontryagin square. A homotopy seven-sphere bounds a spin eight-manifold; for instance the usual parallelizable representative supplied by \(\Theta _7=bP_8\) suffices. The boundary connected sum of this bounding manifold with \(S^1\times D^7\) bounds (P3F.24.1) for either spin structure. Its middle intersection form and its Pontryagin square are direct sums, since the joining handle has index one. Hence \[
 \mu(T^+_\Phi)=\mu(\Sigma_\Phi)
\] for each spin structure. The rational relative-to-absolute degree-four map is an isomorphism here because the mapping torus has \(H^3(-;\mathbb Q)=H^4(-;\mathbb Q)=0\). Thus the usual spin Eells--Kuiper expression \((p_1^2-4\sigma)/896\) applies without a choice of rational lift. Its invariance follows by gluing two such bounds: on the resulting closed spin eight-manifold this expression is the integral \(\widehat A\)-genus. Orientation reversal changes its sign. Both values for the product are zero, so neither preserving nor reversing orientation can identify a nonzero example with that product.

\end{enumerate}
\end{proof}




\begin{theorem}\label{night:P3F.30}
\StatementRecord{night:P3F.30}{proved}{S6\_06\_BLIND: GPT-6 Astra (OpenAI, ChatGPT Pro mode, effort Pro), recorded re-derivation in the checkpoint dated 18 November 2026}{\cref{night:P3F.26}, \cref{night:P3F.27}, \cref{night:P3F.28}, \cref{night:P3F.31}, \cref{night:P3K.65}}
The actual quadratic norm morphism has target line $N=M\otimes T^*M=M^2(-C_0)$ and scalar $\lambda^2$. It completes to $X/\mu_2$, with transverse $A_1$ equation along the fixed curve. The two fixed sections of its resolution have opposite Hopf-line base shifts $+1,-1$, which must be retained in this auxiliary norm route. The positive-fan return index is zero by P3F.31; P3K.65 gives the original negative-fan \(r\)-loop class \(-3\) by the full native Cox comparison.
\end{theorem}
\begin{proof}
\begin{enumerate}[label=\arabic*.]
\item The OS49 component line admits the norm modification \[
 N=M\otimes T^*M=M^{\otimes2}(-C_0),
 \qquad T^*N=N(-F).                                  \tag{P3F.30.1}
\] The actual morphism \(\phi:T^*M\to M\) induces \(\phi\otimes T^*\phi:T^*N\to N\), with divisor \(F\). On the ordinary locus the quadratic map to this norm construction is the quotient by the order-two subgroup of the delta circle. Its completion is the actual analytic quotient \(X/\mu_2\). This quotient has a transverse \(A_1\) singularity along the fixed curve \(F_{\mathrm{fix}}\). Its minimal resolution has two fixed sections over that curve. Their normal weight decompositions are \[
 \bigl(\mathcal O(-2)_{-1},\mathcal O(0)_{+1}\bigr),
 \qquad
 \bigl(\mathcal O(0)_{-1},\mathcal O(-2)_{+1}\bigr).
                                                               \tag{P3F.30.2}
\] After conjugating the negative-weight line, the two determinant degrees are respectively \(+2\) and \(-2\). Thus the two Hopf-line base shifts are \(+1\) and \(-1\), in the same circle convention. These charges are part of this exact quotient and resolution; no return-index equality is inferred by discarding them.

\item \textbf{Proof: the morphism on the ordinary locus.} Tensor the exact line identities of P3F.28. Since \(T\) interchanges \(C_0,C_1\), the two successive zero divisors add to \(F=C_0+C_1\), giving (P3F.30.1). In local frames write \(\phi_p\) as multiplication by \(f(p)\). Away from its zero divisor the precise quadratic map is \[
 M_p\longrightarrow N_p=M_p\otimes M_{Tp},
 \qquad v\longmapsto v\otimes\phi_p^{-1}(v).
                                                               \tag{P3F.30.3}
\] It has kernel \(\{1,-1\}\) on the nonzero fibre. A scalar \(\lambda\phi\) in the original suspension gives scalar \(\lambda^2(\phi\otimes T^*\phi)\) in the norm suspension. Compatibility is the explicit identity \[
 \frac{(\lambda f(p)v)^2}{f(p)}
 =\lambda^2 f(p)f(Tp)\frac{v^2}{f(Tp)}.
\] Consequently \(P3F.30.3\) descends to the specified quotient of the ordinary torus family, with its actual scalar parameter.

\item There is also an exact Tate-chart description. In the source period marking of P3F.26 the change \[
                        W=V^2/Z
\] puts the norm line into its degree-zero extension on both \(I_2\) components. Multiplication of this frame by a nonzero base monomial merely changes the chosen integral toric origin. The period valuations change by the integral matrix \((\begin{smallmatrix}1&0\\-1&2\end{smallmatrix})\): \[
 (2,1)\longmapsto(2,0),\qquad
 (-3,-1)\longmapsto(-3,1).
\] The scalar becomes \(\lambda^2\). These are the periods of the whole-fibre linearization of \(N\). The matrix has determinant two, exactly the degree of \(P3F.30.3\). It maps the unimodular source fan to the quotient fan in the larger target lattice. The affine invariant rings of those cones are precisely the order-two invariant rings, so the toric map extends to the analytic quotient. At the other source fibres the original delta action is free, and the quotient is its usual unramified local quotient. Thus these chart maps glue to \(X\to X/\mu_2\), not merely to a comparison of valuation matrices.

\item \textbf{Proof: resolution and the retained charges.} At the actual fixed curve the normal coordinates \((a,b)\) lie in \(L_-\oplus L_+\), where both line bundles are \(\mathcal O_{\mathbb P^1}(-1)\), and delta acts with weights \((-1,+1)\). The order-two invariant coordinates are \[
 x=a^2,\qquad z=ab,\qquad y=b^2,
 \qquad xy=z^2.
\] They give exactly the transverse \(A_1\) equation, with the fixed curve as its singular locus. The two charts of its minimal resolution are \[
 (x,\zeta):\quad z=x\zeta,\quad y=x\zeta^2;
 \qquad
 (y,\eta):\quad z=y\eta,\quad x=y\eta^2,
 \quad\eta=\zeta^{-1}.
\] They glue over the original fixed curve by its line-bundle transitions. In the first chart, \(x\) is a coordinate in \(L_-^{\otimes2}\) and \(\zeta\) in \(L_+\otimes L_-^{-1}\). In the second, \(y\) lies in \(L_+^{\otimes2}\) and \(\eta\) in \(L_-\otimes L_+^{-1}\). The residual circle coordinate is \(\tau=\delta^2\), and its two chart actions are \[
 (x,\zeta)\longmapsto(\tau^{-1}x,\tau\zeta),
 \qquad
 (y,\eta)\longmapsto(\tau y,\tau^{-1}\eta).
\] Their two fixed sections are \(x=\zeta=0\) and \(y=\eta=0\). The stated bundle transitions and weights prove \(P3F.30.2\). They also show that the exceptional ruled surface is \(\mathbb P(L_-\oplus L_+)\cong\mathbb P^1\times\mathbb P^1\). All maps are defined over the parameter and commute with its actual closing identifications, so the quotient and resolution are family morphisms.

\item Conjugating the negative-weight line makes the circle representation positive on both summands. The first resulting bundle is \(\mathcal O(2)\oplus\mathcal O\), and the second is \(\mathcal O\oplus\mathcal O(-2)\). Factoring out respectively \(\mathcal O(1)\) and \(\mathcal O(-1)\) leaves \(\mathcal O(1)\oplus\mathcal O(-1)\), a topologically trivial \(SU(2)\) bundle on the sphere. The factored line is precisely the additional base factor in the associated Hopf line. Its degrees are \(+1\) and \(-1\). In particular an index comparison through this actual norm quotient must retain both opposite base factors. This norm-quotient route does not yet identify its total correction with the mixed-circle trace of P3F.27. The positive-fan zero index is proved separately by P3F.31--32. P3K.65 retains the missing outer normal-line marking and computes the negative-fan sphere class by the native Cox ruling arrow; no negative-fan integer index is assigned by this norm route.

\end{enumerate}
\end{proof}




\begin{theorem}\label{night:P3F.35}
\StatementRecord{night:P3F.35}{proved}{S6\_06\_BLIND: GPT-6 Astra (OpenAI, ChatGPT Pro mode, effort Pro), recorded re-derivation in the checkpoint dated 18 November 2026}{\cref{night:P3F.26}, \cref{night:P3K.15}}
Every one-vertex unimodular cusp fan is an integral-basis triangulation, and explicit periodic diagonal flops connect it to the vertical-edge fan while preserving sphere topology. For every prime, its cusp fixed set has total mod-prime Betti number two; Smith--Floyd therefore forces every separate cap in a sphere member to have free delta action, without a semifree hypothesis at the cusp.
\end{theorem}
\begin{proof}
\begin{enumerate}[label=\arabic*.]
\item The full argument has passed independent audit.

\item For the whole-zero \(I_1,e=1\) cusp, let the allowed triangulation have vertex set \(\mathbb Z^2\), period group \(\Gamma=\mathbb Z^2\), and unimodular triangles. Every such triangulation is the periodic triangulation \[
 T(v,w)=\{(0,v,v+w),(0,w,v+w)\}+\mathbb Z^2                 \tag{P3F.35.1}
\] for an integral basis \((v,w)\). It is connected to the standard vertical-edge triangulation by a finite sequence of the explicitly described periodic diagonal flops. Each step preserves the integral homology and fundamental group of the compact threefold, and therefore preserves its status as the standard smooth six-sphere.

\item For every prime \(p\), the \(\mu_p\)-fixed subset inside this cusp has total \(\mathbb F_p\)-Betti number exactly two, whether or not the physical delta action is semifree. Consequently every separate finite or logarithmic cap in a sphere member has free delta action. This last assertion supplies the cap-freeness input to P3K.15 without assuming semifreeness at the chosen cusp. The present statement does not yet assign a return class to every fan.

\item \textbf{Proof: the exact form of the triangulation.}

\item The quotient torus has area one. Every triangle has area one half, so there are exactly two triangle orbits. Translation fixes no finite edge or triangle. Counting incidences gives three edge orbits. Choose one triangle \((0,a,b)\); unimodularity says \(\det(a,b)=\pm1\). Its three unoriented edge classes are \(a,b,b-a\), which are distinct and exhaust all edge classes. After the integral linear change \((a,b)\mapsto(e_1,e_2)\), the only possible third vertices of a triangle incident to \([e_1,e_2]\), using those edge classes, are \(0\) and \(e_1+e_2\). The opposite-side condition selects the latter. Likewise, across \([0,e_1]\) it selects \(e_1-e_2\), and across \([0,e_2]\) it selects \(e_2-e_1\). These are exactly the adjacent triangles of the periodic unit-square triangulation. Propagating across edges determines the entire plane triangulation. Taking \(v=a\), \(w=b-a\) gives (P3F.35.1). This is an equality of the actual periodic fans after the stated integral coordinate map, not a replacement by an unmarked abstract triangulation.

\item The three edge classes of \(T(v,w)\) are \(\{\pm v,\pm w,\pm(v+w)\}\). In the parallelogram \((0,v,w,v+w)\), replacing the diagonal \(v+w\) by \(v-w\) gives \(T(v-w,w)\). Conversely replacing the edge class \(v\) by \(v+2w\) gives \(T(v+w,w)\). Thus adding or subtracting one column of a basis from the other is exactly one periodic diagonal flip. Interchanging columns does not change the triangulation. Reversing one column changes the sole diagonal class \(v+w\) to \(v-w\), again one flip.

\item For completeness, an explicit terminating reduction is available. Write the two first coordinates as \((x,y)\). Repeatedly replace the first column by the first column minus \(\lfloor x/y\rfloor\) times the second, and interchange columns. Each integer multiple is executed as its finite sequence of unit column moves. Euclid's algorithm reaches \((\pm1,0)\), because the determinant is \(\pm1\). Reverse the first column if needed. The second column is now \((0,\pm1)\); reverse it if needed and subtract its appropriate multiple from the first column. The result is \((e_1,e_2)\). Every intermediate basis and every flip have been specified, and the standard second edge is the physical vertical delta direction.

\item \textbf{Proof: the actual local flop and its topology.}

\item For one selected edge orbit, each triangle contains exactly one edge from that orbit. Distinct corresponding toric curves therefore share no three-cone, hence no torus fixed point and no intersection. Their lattice translates are disjoint. No nonzero lattice translation fixes a finite edge, so in the quotient the center is one embedded rational curve. The four height-one rays of either adjacent parallelogram satisfy the actual integral relation \(r_0+r_2=r_1+r_3\). Their determinant-one adjacent cones give the normal bundle \(\mathcal O(-1)\oplus\mathcal O(-1)\). Blowing up its zero section gives exceptional divisor \(\mathbb P^1\times\mathbb P^1\), with normal line \(\mathcal O(-1,-1)\). The two ruling contractions are the two fan subdivisions just described. All maps commute with lattice translation and agree with the identity off the compact center; they therefore descend to the actual compact models.

\item The integral blowup decomposition adds precisely one free summand in degrees two and four. Primitivity can also be read directly: the exceptional divisor has degree minus one on either contracted ruling. Thus the common blowup gives \[
 H^j(\widetilde X;\mathbb Z)
   =H^j(X_-;\mathbb Z)\oplus H^{j-2}(\mathbb P^1;\mathbb Z)
   =H^j(X_+;\mathbb Z)\oplus H^{j-2}(\mathbb P^1;\mathbb Z).
\] There is no torsion quotient hidden in the additional summands. Removing a real codimension-four center preserves fundamental groups by general position; its normal sphere bundle and disk bundle are simply connected, so van Kampen also proves invariance under the blowup and either contraction. Starting with a simply connected integral homology six-sphere therefore gives the same groups after every step. The smooth sphere identification is the same Smale and \(\Theta_6=0\) argument used in P3F.26.

\item \textbf{Proof: the exact prime fixed sets.}

\item Let \(\delta\in\mathbb Z^2\) be the primitive physical circle cocharacter, and write \(\delta=a v+b w\). On the three open edge strata the finite isotropy orders are \[
 |\det(\delta,v)|=|b|,\qquad
 |\det(\delta,w)|=|a|,\qquad
 |\det(\delta,v+w)|=|a-b|.                              \tag{P3F.35.2}
\] Here order zero denotes the full circle stabilizer. The formula comes from the exact quotient lattice \(\mathbb Z^2/\mathbb Z e\) for an edge direction \(e\): the delta cocharacter has image \(\det(e,\delta)\) in that rank-one lattice. Since \(\gcd(a,b)=1\), any prime divides at most one of the three numbers in \(P3F.35.2\), including zero.

\item In a determinant-one triangle chart with coordinates \((z_0,z_1,z_2)\), the circle has integral weights \((c_0,c_1,c_2)\) with \(c_0+c_1+c_2=0\). Their greatest common divisor is one: the triangle rays are an integral basis of the three-dimensional toric lattice and delta is primitive. Thus modulo a prime no two weights can vanish; if two did, the third would also vanish, contradicting primitivity. The \(\mu_p\)-fixed locus in this chart is therefore either the isolated origin or one coordinate axis. Those axes are exactly the edge strata selected by \(P3F.35.2\). If a prime divides one edge order, the two local axes glue through that edge to its whole rational curve; both triangle fixed points lie on it and there are no further fixed points. If it divides none, the fixed set consists of the two triangle points. Open toric surface strata have no finite stabilizer, because their two-dimensional torus lattice is the original \(\mathbb Z^2\) and delta is primitive.

\item On every regular fibre the delta action is primitive translation, so no \(\mu_p\)-fixed component can leave the central cusp. Consequently the compact cusp fixed set is exactly \(\mathbb P^1\) or exactly two points. Both have total \(\mathbb F_p\)-Betti number two. A nonfree delta action in a separate cap has a point fixed by some prime subgroup. Its nonempty compact fixed subset is disjoint from the cusp fixed set and contributes at least one to the total Betti number. Floyd's Theorem 4.4, applied to the order-\(p\) diffeomorphism of the compact six-sphere, bounds that total by \(\sum_i b_i(S^6;\mathbb F_p)=2\), a contradiction. This proves cap freeness for every allowed fan.

\item The primary theorem was checked in E. E. Floyd, \emph{On periodic maps and the Euler characteristics of associated spaces}, Trans. AMS 72 \(1952\), Theorem 4.4, p.146; the coefficient convention is stated on p.139. The exact elementary fan reduction is reproduced in \path{code/part3/arbitrary_fan.py}.

\end{enumerate}
\end{proof}




\begin{theorem}\label{night:P3F.36}
\StatementRecord{night:P3F.36}{proved}{S6\_06\_BLIND: GPT-6 Astra (OpenAI, ChatGPT Pro mode, effort Pro), recorded re-derivation in the checkpoint dated 18 November 2026}{\cref{night:P3K.15}, \cref{night:P3K.21}, \cref{night:P3K.22}}
Let the marked rational elliptic source have either profile
\(I_1^*+I_3+2I_1\) or \(I_1^*+IV+I_1\), and take the
OS50 data \((a,b)=(6,2)\). Put
\(M=\mathcal O(6G-O-L_0)\), with the actual component \(L_0\)
of P3K.21, and place the simple whole-fibre zero at an \(I_1\)
value. At the \(I_1^*\) value the actual primitive quadratic
cap exists for every odd norm numerator \(n\), with its integral
delta phase retained. Leave the other source values ordinary.
The circle orbit of this actual construction satisfies
\[
 \pi_1(Y)=\mathbb Z/|n|,\qquad
 H_j(Y;\mathbb Z)=
 \begin{cases}
 \mathbb Z,&j=0,5,\\
 \mathbb Z/|n|,&j=1,3,\\
 0,&j=2,4.
 \end{cases}
\]
The total six-manifold is a smooth sphere exactly when \(n=1\)
or \(n=-1\). This statement supplies the actual cap and its
integral topology; it assigns no parameter-return class.
\end{theorem}

\begin{proof}
\begin{enumerate}[label=\arabic*.]
\item \par\medskip\noindent\textbf{The actual root and line correspondence.}
We first give the local construction for each surviving starred
line point \(a=2,6\) and translation \(b\equiv2\pmod4\).
The minimal quadratic root \(S'\to\Delta_z\) has fibre
\(I_2\). Its canonical deck \(D\) combines \(z\mapsto-z\)
with inversion on the Tate fibre. It fixes each of the two central
components and has two fixed smooth points on each. Blow up these
four points by \(\beta:\widehat S\to S'\). The lifted deck
fixes each exceptional curve pointwise; in a chart whose coordinates
before blowing up have signs \((-1,-1)\), the blowup coordinates
\((w/z,z)\) have signs \((1,-1)\), and the quotient coordinates
are \((w/z,z^2)\). The quotient \(\pi:\widehat S\to S\)
is therefore the smooth starred source of P3K.21.

\item Let \(E_0,E_1\) lie over the leaves \(L_0,L_1\) met by
\(O\) and \(P=aG\). Write \(\widehat O,\widehat P\)
for the strict transforms of the root sections. The actual divisor
pullbacks are
\[
 \pi^*L_0=2E_0,\qquad
 \pi^*(P-O-L_0)=\widehat P-\widehat O-2E_0,
\]
\[
 \beta^*(P'-O')=\widehat P-\widehat O+E_1-E_0.
\]
Consequently, with \(N=\mathcal O(P'-O')\), we have the exact
line comparison
\[
                  \beta^*N=\pi^*M+E_0+E_1.                 \tag{P3F.36.1}
\]
The two root sections meet the same \(I_2\) component, so \(N\)
has degree zero on both components. Translation by \(bG\) swaps
\(E_0,E_1\), hence preserves their sum. Thus (P3F.36.1) transports
the specified suspension Hom line and its lift, including the
vertical correction.

\item More explicitly, write the Tate period and generator as
\[
 q_T=z^2u(z^2),\qquad Z_G=zv(z),\qquad
                  v(z)v(-z)=-u(z^2).
\]
The degree-zero root point is represented by the nonvanishing unit
\[
 Z_P=v(z)^a/u(z^2)^{a/2},\qquad
 Z_P(-z)=Z_P(z)^{-1},\qquad Z_P(0)=-1.
\]
Its line has Tate transitions
\((Z,\lambda)\sim(q_TZ,Z_P\lambda)\). These unit transitions
extend over the two-component model and realize precisely the
line in (P3F.36.1). They also retain its nonzero order-two central
Picard character.

\par\medskip\noindent\textbf{The actual norm lift and its sign.}
The positive global Hom section gives the contracting map covering
\(t_{-bG}\). Let \(F\) be a line lift on \(N\) of
\(t_G\), provided directly by the Tate unit transitions, and
let \(D\) denote the canonical order-two line lift as well.
Every other lift of this translation differs by a base unit.
Write
\(DFD^{-1}=c(z)F\). Since \(D^2=1\),
\(c(z)c(-z)=1\). Dividing \(F\) by a base unit obtained from
the odd logarithm of \(c(z)/c(0)\) makes
\[
                  DFD^{-1}=\chi F,\qquad \chi\in\{1,-1\}.
\]
The actual nonzero-linearization lift is
\(B_0=\beta_0(z)F^{-b}\). Its descent and even \(b\) imply
\(\beta_0(-z)=\beta_0(z)\). Put \(k=-nb/2\) and choose
analytic square roots of the even nonvanishing base units. Then
\[
 A_0=\chi^{-k/2}\beta_0(z)^{n/2}F^kD
 \quad\hbox{satisfies}\quad
 A_0B_0=B_0A_0,\qquad A_0^2=B_0^n.                         \tag{P3F.36.2}
\]
Indeed \((F^kD)^2=\chi^kF^{2k}\); this character factor
is cancelled by the displayed scalar, not omitted. The integral
choice of its phase is retained in the outer line collar.

\item For the family parameter \(r\), take
\(\lambda=r^2\), \(B_\lambda=\lambda B_0\), and
\(A_r=r^nA_0\). These are actual maps and
\(A_r^2=B_\lambda^n\). Choose Bezout integers \(r_0,s_0\) with \(r_0n+2s_0=1\). The generator
\(H=A_r^{r_0}B_\lambda^{s_0}\) satisfies
\(A_r=H^n\), \(B_\lambda=H^2\). Its source map is exactly
\[
                         h=t_{-bG/2}D.
\]
Its multiplier contains the positive first power of the small
parameter \(r\), times a bounded line map on a compact smaller
root neighbourhood. It is uniformly contracting for small
\(r\ne0\). Thus no nonzero iterate fixes a nonzero line vector,
and the norm-shell fundamental region proves properness and
compactness as in P3K.15. The quotient is smooth and its line-circle
action is free. The source action exchanges the two central
components.

\item We fix the outer signs explicitly. Normalize the real line
coordinate so that \(B\) sends \(t\) to \(t-1\), and put
\(\theta=\arg z/(2\pi)\). The actual root action is
\[
 A:(\theta,t)\longmapsto(\theta+1/2,t-n/2).
\]
The original outer angular coordinates are
\[
                       x=2\theta,\qquad y=t+n\theta.
                                                               \tag{P3F.36.3}
\]
The positive base meridian \(\alpha\) changes \(x\) by one
at fixed \(y\), while the positive outward suspension circle
\(\gamma\) changes \(y\) by one. The positive cap generator
\(h\) changes \(t\) by \(1/2\). Hence the free cap map is
\[
                         \alpha\longmapsto-nh,
                  \qquad \gamma\longmapsto2h,               \tag{P3F.36.4}
\]
and its killed meridian is \(2\alpha+n\gamma\).
The exact norm-shell coordinate of P3K.15 straightens the bounded
norm units; the remaining source and delta periods stay in the
integral collar coefficients below.

\par\medskip\noindent\textbf{The source complement and its actual chain maps.}
Now take \(b=2\) and either stated source profile. Let \(P\)
be the complement of a small \(I_1^*\) tube. Its elliptic
fibration over the complementary disc has a section, so van Kampen
presents \(\pi_1(P)\) as a quotient of the elliptic lattice by
the singular-fibre relations. It is abelian. The primitive component
lattice and the unimodular source intersection form give
\(H_1(P)=0\); therefore \(\pi_1(P)=0\). The same sequence
gives \(H_2(P)=\mathbb Z^4\), with no other positive reduced
homology.

\item Use the actual source basis
\((O,F,C_0,C_1,L_1,R_0,R_1,A_{2,1},A_{2,2},G)\).
Its Gram determinant is \(-1\). The dual columns indexed by
\((O,G,A_{2,1},A_{2,2})\) give the complement pairing
\[
 J=\begin{pmatrix}
 -12&12&8&4\\12&-12&-8&-4\\8&-8&-6&-3\\4&-4&-3&-2
 \end{pmatrix}.
\]
Modulo the starred components, translation by \(2G\) has matrix
\[
 T=\begin{pmatrix}
 -1&-2&0&0\\2&3&0&0\\1&2&-1&1\\0&1&-1&0
 \end{pmatrix},\qquad U=T^{-t}.
                                                               \tag{P3F.36.5}
\]
For example \([2G]=-[O]+2[G]+[A_{2,1}]\) and
\([3G]=-2[O]+3[G]+2[A_{2,1}]+[A_{2,2}]\); translation sends
\(A_{2,1}\) to \(-A_{2,1}-A_{2,2}\) and \(A_{2,2}\)
to \(A_{2,1}\). These are the actual section and component
classes for both profiles. Thus no identification of their complex
sources is needed. The certificate checks \(JU=TJ\).
The actual contracting translation is the inverse, with matrices
\(T^{-1},U^{-1}\).

\item For clarity we give the exact chain comparison to the auxiliary
positive-translation calculation. The diffeomorphism
\[
 T_{t_{-2G}}P\longrightarrow T_{t_{2G}}P,
                     \qquad[p,t]\longmapsto[p,-t]             \tag{P3F.36.6}
\]
fixes source cycles and reverses the suspension circle. The boundary
complex has groups \(\mathbb Z,\mathbb Z^4,\mathbb Z^4,
\mathbb Z\) in degrees zero through three and sole differential
\(d_2=J\). It is the Poincare--Lefschetz cone for \(P\).
Both absolute and relative split complexes have free homology in
even degrees, so the naturality homotopy has no degree-one term.
Taking the mapping-torus cone therefore gives the actual boundary
and inclusion maps. In degrees one, two, three, (P3F.36.6) has
chain matrices
\[
 \Phi_1=\operatorname{diag}(I_4,-1),\quad
 \Phi_2=\operatorname{diag}(I_4,-T^{-1}),\quad
 \Phi_3=\operatorname{diag}(1,-U^{-1}).
\]
The certificate checks both chain identities before computing
homology. For example its positive-model degree-two cycle columns
are
\[
 R_2=\begin{pmatrix}
 -1&0&1\\0&0&0\\-1&0&0\\-1&0&2\\
 -2&3&-2\\2&-3&0\\2&-2&0\\1&-1&0
 \end{pmatrix};
\]
the actual columns are exactly \(\Phi_2^{-1}R_2\).

\item Write \((\epsilon,r_2,s_2)\) for these boundary classes,
with orders \((2,0,0)\), and \((\tau,g)\) for the exterior
classes, of orders \((2,0)\). Their actual inclusion and fibre
classes are
\[
 i_2=\begin{pmatrix}1&0&-2\\0&0&1\end{pmatrix},\qquad
 F_{\partial}=\epsilon-2r_2,\qquad F_W=\tau.                  \tag{P3F.36.7}
\]
In particular the fibre is the nonzero order-two exterior class.
The cap retracts onto the mapping torus of the \(I_2\) reflection.
Its \(H_2\) is generated by one component-orbit class \(C\),
and the actual fibre specializes to \(2C\). Thus, for the actual
integral collar coefficient \(\kappa\), its boundary map is
\[
             (\epsilon,r_2,s_2)\longmapsto(0,-1,\kappa)C.
\]
The entire degree-two Mayer--Vietoris map is
\[
 D_2=\begin{pmatrix}
 1&0&-2\\0&0&1\\0&1&-\kappa
 \end{pmatrix}.                                               \tag{P3F.36.8}
\]
The first source and target generators have order two and map
identically; the free block is unimodular. This is an integral
isomorphism for the actual \(\kappa\). Equivalently the explicit
basis change \(s_{2,\mathrm{new}}=s_2+\kappa r_2\) makes its
cap column zero and preserves its exterior column. We do not infer
\(\kappa=0\) from an inversion: the exterior inversion can
change the marked cap by a gamma shear.

\par\medskip\noindent\textbf{The remaining maps and fundamental group.}
The cap central graph has generator \(\eta\), reversed by the
affine reflection. Thus
\[
 \pi_1(\mathrm{cap})=
       \langle h,\eta\mid h\eta h^{-1}=\eta^{-1}\rangle,
\]
\(H_1=\mathbb Zh\oplus\mathbb Z/2\eta\), and
\(H_2=H_3=\mathbb Z\). Let \(u,w\) be the vanishing and
graph periods on the source boundary. In the positive outward
\(\gamma\) convention its affine relation is
\(\alpha\gamma\alpha^{-1}=\gamma w^b\). The actual map,
retaining the graph-path integer \(c\), is
\[
 u\mapsto1,\qquad w\mapsto\eta,\qquad
 \alpha\mapsto h^{-n}\eta^c,\qquad
 \gamma\mapsto h^2\eta^{-b/2}.                               \tag{P3F.36.9}
\]
The commutator identity verifies the graph correction. In degree
one, with the target cap signs absorbed into its chosen generators,
the complete matrix for \(b=2\) is
\[
 D_1=\begin{pmatrix}1&c&-1\\0&-n&2\\0&0&1\end{pmatrix}.
\]
Its first source and target entries have order two. Their map is
the identity and the free determinant is \(-n\).

\item In degree three use the oriented source boundary and the positively
oriented fibre--gamma trace. The exterior row is \((0,1)\).
The cap generator is the trace of the sum of the two components
through one positive \(h\) turn. The norm powers and
(P3F.36.4) give cap row \((-n,2)\). A graph correction is a
path inside the two-dimensional central source and adds no
three-homology class. Thus
\[
                       D_3=\begin{pmatrix}0&1\\n&-2\end{pmatrix}.
\]
In degree zero the map is \((1,-1)^t\); in degree four the
domain is \(\mathbb Z\) and the target is zero. These specify
every nonzero Mayer--Vietoris map and give the stated integral
homology. All matrices, their chain representatives and the
symbolic collar coefficients are executed in
\path{code/part3/i1star_boundary.py} and its JSON output.

\item For van Kampen, \(\pi_1(W)=\mathbb Z\gamma\). Gluing kills
\(u,w,\alpha\), hence \(\eta=1\), \(h^n=1\), and
\(\gamma=h^2\). Since \(n\) is odd, the resulting group is
exactly \(\mathbb Z/|n|\).

\par\medskip\noindent\textbf{The actual six-dimensional Hopf filling.}
For \(n=\pm1\), \(Y\) is simply connected and has the
integral homology of the five-sphere, hence is the standard smooth
five-sphere by the smooth classification already used in P3K.22.
The simple whole zero still has the same actual cusp chart and
fixed sphere \(F\cong S^2\), with normal circle weights
\((-1,+1)\). On \(Y\setminus\operatorname{int}\nu F\),
Alexander duality gives the cohomology of \(S^2\). The normal
Hopf meridian has Euler value one, so the principal-circle Gysin
sequence gives the homology of \(S^3\), and its homotopy sequence
gives trivial fundamental group. Gluing \(F\times D^4\) along
\(F\times S^3\), using the actual normal \(SU(2)\) framing,
has maps \((1,-1)^t\) in degree zero, \(-1\) in degree two,
\(1\) in degree three, and zero from \(\mathbb Z\) in degree
five. Thus the resulting \(X\) is simply connected with the
integral homology of \(S^6\). The same six-dimensional smooth
classification as P3K.22 gives the standard smooth sphere.

\item For general odd \(n\), removing a codimension-three fixed sphere
preserves \(\pi_1(Y)\). Its local Hopf meridian still makes the
circle-boundary map onto \(\pi_1(S^1)\), so the free-circle
piece has the same fundamental group as its orbit. The simply
connected fixed tube and its simply connected boundary do not
change that group. Hence \(\pi_1(X)=\mathbb Z/|n|\), which
excludes every other odd numerator from a sphere. All line phases
and collars were retained throughout; no return invariant follows
from this topological calculation alone.

\end{enumerate}
\end{proof}




\begin{theorem}\label{night:P3F.37}
\StatementRecord{night:P3F.37}{proved}{S6\_06\_BLIND: GPT-6 Astra (OpenAI, ChatGPT Pro mode, effort Pro), recorded re-derivation in the checkpoint dated 18 November 2026}{\cref{night:P3F.2}, \cref{night:P3F.9}, \cref{night:P3K.22}, \cref{night:P3K.24}}
Let $S$ be the minimal regular surface of
\[
 y^2=x^3+tx+1,\qquad G=(0,1),\qquad
 M=\mathcal O_S(2G-O),\qquad Q=G.
\]
Choose any one of its three nodal fibres as the simple whole-fibre
zero of the linearization. For every $m=4l+3$, $l\ge0$, the
explicit root and affine norm below give a reduced type III cap
at the type III* fibre. Its marked meridian is
$m q+\gamma$, where $q$ is the positive base meridian and
$\gamma$ is the outward suspension loop.
The resulting construction is a proper holomorphic family over a
punctured $r$-disc, with scalar $\lambda=r^m$. Every fibre is
orientation-preservingly diffeomorphic to the standard $S^6$.
The construction and topology asserted here do not assign the
oriented parameter-return class.
\end{theorem}
\begin{proof}
\begin{enumerate}[label=\arabic*.]
\item \par\medskip\noindent\textbf{The source and its actual root.}
The three finite discriminant roots are the distinct roots of
$4t^3+27$ and give $I_1$ fibres. At infinity put
$q=1/t$, $x_0=q^2x$, $y_0=q^3y$. The exact equation is
\[
 y_0^2=x_0^3+q^3x_0+q^6,\qquad G=(0,q^3).
\]
It has type III*. The source is rational, as witnessed on a
nonempty open set by $t=(y^2-x^3-1)/x$. The marked section data
of P3F.2 give $h(G)=1/2$, $G\cdot O=0$, and the nonidentity
component of the $E_7$ fibre. The height discriminant identity
then proves that $G$ is primitive and that the Mordell--Weil group
has no torsion. The section $2G$ has identity component and height
two, hence $2G\cdot O=0$ on this rational surface. Therefore
$M$ has multidegree zero on every component. Restricting the
actual Hom line to $O$ gives its degree
\[
 (2G)\cdot O+1-G\cdot O+(-G)\cdot O=1.
\]
Its vertical divisor has zero degree on each component, so it is
a whole-fibre pullback. Its section can consequently have its
sole simple zero at the selected $I_1$. Denote its line map by
$B_0$; it covers $t_{-G}$, and the contracting lift is
$B_\lambda=\lambda B_0$.

\item For $m=4l+3$ and $k=3l+2$ make the actual substitution
\[
 q=z^m,\qquad x_0=z^{2k}X,\qquad y_0=z^{3k}Y.
\]
The minimal equation and its marked section become
\[
 Y^2=X^3+zX+z^{6(l+1)},\qquad G'=(0,z^{3(l+1)}).       \tag{F37.1}
\]
Its only surface singularity over zero is the ordinary double
point with quadratic part $Y^2-zX$. One point resolution gives
a smooth source with a reduced III fibre. For example, at $m=3$
the chart $Y=Xv$, $z=Xw$ is
\[
 v^2=X+w+X^4w^6.
\]
The components $X=0$ and $w=0$ meet with tangency order two.
For general $l$ the same chart has the corresponding higher
monomial and the same smoothness and component configuration.
The exact deck is
\[
 D(X,Y,z)=(\zeta_m^{-2k}X,\zeta_m^{-3k}Y,\zeta_m z).   \tag{F37.2}
\]
It extends on the resolution and fixes both component labels.
At $m=3$ its chart formula is
$(X,v,w)\mapsto(\zeta_3^2X,\zeta_3v,\zeta_3^2w)$.
The section $G'$ meets the nonidentity component, while $2G'$
meets the identity component. This also follows from composing
with the good fourth root: its marked limit is the nonzero fixed
two-torsion point, as proved by the coordinate comparison in the
reduced-additive root calculation.

\par\medskip\noindent\textbf{The source norm and the actual line.}
Multiplication by the odd integer $m$ is a bijection on the
central smooth Neron group. On its additive identity component
its derivative is $m$, and on its component group $\mathbb Z/2$
it is invertible. The analytic inverse-function theorem gives the
unique local section $R$ with $mR=G'$ and the prescribed central
point. Uniqueness gives $D(R)=R$. The source translation extends
over the whole minimal regular surface: resolve its graph, or
apply the resolved addition and canonical-form comparison used
in P3K.22. A remaining exceptional curve would have to be a
vertical $(-1)$-curve on a relatively minimal source. Thus
\[
                   a_0=t_{-R}D,\qquad a_0^m=t_{-G'}.
\]
It exchanges the two III components.

\item Put $N=\mathcal O(2G'-O)$. Pull $M$ to the resolved fibre
product first. Its degree on a curve mapping to a point is zero,
and its degree on the transform of every source-fibre component
is zero as well. It therefore descends through each actual
$(-1)$-contraction of the relative minimalization, by the blowup
Picard decomposition. Its comparison with $N$ has a vertical
divisor of degree zero on the two III components. The matrix
\[
                    \begin{pmatrix}-2&2\\2&-2\end{pmatrix}
\]
has integral kernel $\mathbb Z(1,1)$, so this difference is an
integer whole fibre, principal on the disc. This identifies the
line, including its descended deck lift; that lift has $D^m=1$
and commutes with the pullback $B_0$.

\item An invertible line lift $L$ covering $t_{-R}$ can be constructed
without discarding its vertical divisor. For $P=2G'$, the chord
through $P,-R$, divided by the vertical function through $P-R$,
has fibre divisor
\[
                        P+(-R)-(P-R)-O.
\]
This is the divisor of $\operatorname{Hom}(t_R^*N,N)$ in the
specified direction. Its remaining vertical divisor has degree
zero on both components. The displayed matrix makes it an
integer whole fibre, removed by the corresponding base power.
The tangent formulas at coincident points give the same
meromorphic section and its extension.

\item Deck conjugation acts on the rank-one free module of such Hom
sections. Let $\chi$ be its central character. The finite average
\[
 L_{\rm eig}=\frac1m\sum_{j=0}^{m-1}
                 \chi^{-j}D^jLD^{-j}
\]
is nonzero at the central fibre and invertible after shrinking.
It satisfies $DL_{\rm eig}D^{-1}=\chi L_{\rm eig}$ and
$\chi^m=1$. Both $L_{\rm eig}^m$ and $B_0$ cover $t_{-G'}$.
Their ratio is a nonvanishing base function: scalar ratios are
constant on compact connected fibres. Deck invariance makes it
a function $u(q)$. Consequently
\[
 (L_{\rm eig}D)^m=\chi^{m(m-1)/2}u(q)B_0.
\]
Divide by an analytic $m$-th root of this base unit. The resulting
line map $A_0$ satisfies
\[
                   A_0^m=B_0,\qquad A_0B_0=B_0A_0.     \tag{F37.3}
\]
The character factor is one for this odd $m$, but was retained
before normalization. With $\lambda=r^m$ and $A_r=rA_0$ the
identity is exactly $A_r^m=B_\lambda$.

\item On a smaller compact source tube the norm ratios of $A_0$ and
its inverse are bounded. For small $|r|$, $A_r$ is strictly
contracting. Thus $N^\times/\langle A_r\rangle$ is a smooth
proper cap, with its compact norm shell as a fundamental region.
The line circle is free everywhere on this cap: equality of a
unit rotation and an iterate $A_r^j$ forces $j=0$ by strict
contraction, and then forces the rotation to be the identity.
This includes both components and their tangency point.

\par\medskip\noindent\textbf{The annular map and marked meridian.}
Use an outward increasing compatible log-norm coordinate $u$.
The actual cap identification is
\[
       (p,z,u)\sim(a_0p,\zeta_mz,u-1).
\]
The positive suspension generator is $h=A_r^{-1}$, and the
ordinary outward generator is $\gamma=B_\lambda^{-1}=h^m$.
Choose the positive root cut crossing to be the single map
$A_r$, without an additional $B_\lambda$ iterate. In real
angular coordinates its marking is
\[
 \frac{\arg z}{2\pi}=\frac{\arg q}{2\pi m},\qquad
 u_{\rm cap}=m u_{\rm ordinary}-\frac{\arg q}{2\pi}.
\]
Thus $q\mapsto-h$, $\gamma\mapsto mh$, and the killed
meridian is exactly $mq+\gamma$.

\item Here is a holomorphic realization of that marking. On a simply
connected sector a degree-zero source line has the constant-factor
presentation
\[
 (a,V)\sim(a+1,V),\qquad
 (a,V)\sim(a+\tau,e^{2\pi i\beta}V).
\]
The multiplication $(a,V)(a',V')=(a+a',VV')$ is its actual
commutative extension group. A line map covering a translation
lifts with a doubly periodic holomorphic multiplier, hence a
constant multiplier. The annular cap is a torsor for this same
group, with its cocycle relation precisely (F37.3). Lift the
cocycle to the holomorphic Lie bundle. Its integral period
two-cocycle is a coboundary because the annulus retracts onto a
circle. Remove it, and solve the resulting Lie cocycle using
$H^1=0$ for a holomorphic vector bundle on the Stein annulus.
Translation by the resulting zero-cochains gives mutually inverse
holomorphic gluing maps. The same exact sequence makes the
component map onto $H^1$ of the period local system surjective.
Its gamma class is primitive: the invariant integral quotient
$\psi$ takes gamma to one and the other periods to zero.
Choose the component with cut crossing $A_r$ itself; an additional
$B_\lambda^j$ is removed by the section of gamma winding $j$.
All source and line-circle phases remain in these translations.

\item Fix the arrow $F_c:\text{ordinary}\to\text{cap}$ and use the
physical character $V=e^{2\pi i\zeta_2}$, with
$c=\log\lambda/(2\pi i)$. Write $s=\log q/(2\pi i)$.
Its Lie translation obeys $g_c(s+1)=Dg_c(s)+a_c$.
Solving once at $c_0$ and adding $(c-c_0)s\delta/m$ gives its
exact parameter dependence. On a real angular cut its multiplier
is $r^s$, which is the norm marking displayed above. Since the
cap chart is identical at $c$ and $c+m$, the inverse closing map
is $F_c^{-1}F_{c+m}$, with multiplier $q$ near infinity.
Globally its multiplier is
\[
                    (t-t_c)^{-1}=\frac{q}{1-t_cq}.
\]
The positive parameter deck is its inverse. At the chosen cusp
multiplication by $t-t_c$ is, modulo $B_\lambda$, translation
by $+G$ and an invertible scalar unit. Both extend in the actual
nodal and Mumford charts, as does the inverse. At the other
fillings it is a unit; the displayed unit at infinity extends
over the III cap. The maps agree on their specified annuli.
They descend the construction to $r=e^{2\pi ic/m}$, and compact
norm shells give properness on each compact parameter annulus.

\par\medskip\noindent\textbf{The source-complement lattice.}
Use the integral divisor basis $O,F,C_0,\ldots,C_6,G$, where the
$E_7$ Cartan graph has edges
\[
 (0,1),(1,2),(2,3),(3,4),(4,5),(2,6),
\]
and $G$ meets $C_5$. The inverse Cartan column is
\[
 R_1=(1,2,3,5/2,2,3/2,3/2)^t.
\]
The full source Gram matrix has determinant $-1$, with
$O^2=G^2=-1$, $OF=GF=1$ and $GC_j=\delta_{j5}$.
These are therefore an integral basis. The full III* component
lattice, generated by $F$ and the seven $C_j$, is primitive.
For $V=S\setminus\operatorname{int}N_{III^*}$, excision and
the unimodular intersection pairing give
\[
 H_1(V)=H_3(V)=0,\qquad H_2(V)=\mathbb ZF\oplus\mathbb ZM,
 \quad M=P_{2G}-O=2G-2O-F+2R_1.                       \tag{F37.4}
\]
Orthogonality to $F$ and the root block forces the $G$ coefficient
to be even, proving this integral kernel description. The actual
source suspension $T=t_{-G}$ acts by
\[
                         TF=F,\qquad TM=M-F.           \tag{F37.5}
\]
For example the exact section formula
$P_{-G}=2O+2F-G-2R_1$ verifies the second equality directly.

\item The three nodal vanishing cycles generate the full elliptic
lattice. Rank one would make their monodromy product have trace
two. If their rank-two span had index $d>1$, every pairwise
determinant would be divisible by $d$. The trace of the product
of the three transvections is then $2$ modulo $d^2$: its quadratic
terms are determinant squares and its cubic term is their product.
The product is inverse to the III* monodromy and has trace zero.
Thus $d^2\mid2$, so $d=1$. Van Kampen with the section over
the disc base now gives $\pi_1(V)=1$.

\item The actual circle orbit exterior is the mapping torus of $T$ on
$V$. Its inclusion of the chosen cusp uses the relative nodal
orbit-product map of P3F.9, so no cusp piece is omitted. Denote
it by $W_V$. Wang's sequence and (F37.5) give
\[
 \pi_1(W_V)=\mathbb Z\gamma,\quad H_1(W_V)=\mathbb Z,
 \quad H_2(W_V)=\mathbb Z[M],\quad H_3(W_V)=\mathbb Z[F],
 \quad H_4(W_V)=0.
\]
The compact middle pairing is zero: the $H_3$ generator is the
trace of $F$, and $F\cdot M=0$ in a source slice.

\par\medskip\noindent\textbf{Every needed primitive gluing map.}
For the source boundary $L=\partial N_{III^*}$,
\[
 H_1(L)=\mathbb Zq\oplus\mathbb Z/2\,\eta,
 \qquad H_2(L)=\mathbb ZF.
\]
The inverse Cartan column $R_1$ modulo the root lattice is the
nonzero class $\eta$. Thus $Tq=q-\eta$ and $T\eta=\eta$.
An equivalent marked three-torus bundle for $E=\partial W_V$
has matrix
\[
 J=\begin{pmatrix}0&1&-1\\-1&0&0\\0&0&1\end{pmatrix}
\]
over $q$, in the elliptic and gamma basis. Both $J-I$ and
$\bigwedge^2J-I$ have Smith factors $(1,1,0)$. Their primitive
invariant vector and degree-two coinvariant functional are
$(-1,1,2)^t$ and $(2,-1,-1)$, respectively. Hence
\[
 H_1(E)=\mathbb Zq\oplus\mathbb Z\gamma,
 \quad H_2(E)=\mathbb Zz_1\oplus\mathbb Zz_2,
 \qquad F=2z_1.                                      \tag{F37.6}
\]
Poincare--Lefschetz duality and the zero compact middle pairing
make $H_2(E)\to H_2(W_V)$ onto with primitive kernel. Since $F$
maps to zero, that kernel is $\mathbb Zz_1$.

\item The signed second column has an actual relative-Wang representative.
Let $D_O$ be the punctured zero section, with boundary $q$, and
choose a boundary two-chain $C$ with $\partial C=2(Tq-q)$.
Then $2(TD_O-D_O)-C$ is a cycle. Changing $C$ adds a multiple
of $F$, zero in the target. Closing the section discs gives
\[
                  2(P_{-G}-O)=-M+3F-2R_1.
\]
The last two terms are in the removed fibre lattice, so the
second column is exactly $-M$.

\item The cap orbit $W_U$ is the mapping torus of $a_0$ on the III
tube. That tube retracts onto two spheres meeting at one point,
and $a_0$ exchanges them. Thus
\[
 \pi_1(W_U)=\mathbb Zh,\qquad H_2(W_U)=\mathbb ZC,
 \qquad H_3(W_U)=\mathbb Z(C_0+C_1).
\]
The incoming elliptic fibre specializes to $C_0+C_1$, giving
$F\mapsto2C$. Equation (F37.6) gives $z_1\mapsto C$.
Write $z_2\mapsto\kappa C$. The actual degree-two
Mayer--Vietoris matrix is
\[
 D_2=\begin{pmatrix}0&-1\\-1&-\kappa\end{pmatrix},
 \qquad \det D_2=-1.                                 \tag{F37.7}
\]
The determinant-one boundary operation $z_2\mapsto z_2-\kappa z_1$
makes its lower-right entry zero and leaves the exterior column
$-M$ unchanged. It does not alter the degree-one meridian.

\item The cap norm already gave $q\mapsto-h$, $\gamma\mapsto mh$;
the exterior sends $q$ to zero and gamma to its generator.
Consequently the full degree-one matrix is
\[
              D_1=\begin{pmatrix}0&1\\1&-m\end{pmatrix},
              \qquad\det D_1=-1.                     \tag{F37.8}
\]
We specify the higher maps as well. A primitive basis of $H_3(E)$
is the oriented source boundary $L$ followed by the source-fibre
gamma trace $F\gamma$. In the three-torus bundle description,
the former projects to the primitive invariant $F$ over $q$,
and the latter is the fundamental three-torus in a base slice.
The exterior row is $(0,1)$: $L$ bounds in the source complement
and $F\gamma$ is its Wang generator. For the cap, use the actual
projection to its norm-circle coordinate. On $L$ it has degree
$-1$, and on $F\gamma$ it has degree $m$, by the marked formula
$u=m\gamma-\arg q/(2\pi)$. At a regular norm value the corresponding
source cycle is the generic elliptic fibre. Its actual specialization
is $C_0+C_1$, the primitive invariant generator of the cap's source
$H_2$. Thus the cap row is exactly $(-1,m)$; there is no further
three-class from the source tube, which retracts onto its
two-dimensional central fibre. The full degree-three matrix is
\[
 D_3=\begin{pmatrix}0&1\\1&-m\end{pmatrix},
 \qquad \det D_3=-1.                                  \tag{F37.9}
\]
In degree four the map is $D_4:\mathbb Z\to0$, since both
pieces have $H_4=0$. In degree zero it is the incidence column
$D_0=(1,-1)^t$. These specify every nonzero gluing degree.

\item The elliptic boundary loops die in both cyclic fundamental groups.
These exact inclusions give $\pi_1(Y)=1$ by van Kampen.
Equations (F37.7)--(F37.9), $D_4$, and $D_0$ give all integral
sphere homology groups: only $H_0(Y)=H_5(Y)=\mathbb Z$ remains.
The smooth five-dimensional classification already used in
P3K.22 identifies this orbit with the standard $S^5$.

\item The only fixed set of the line circle is the single index-one
cusp sphere, with its actual $\mathcal O(-1)^{\oplus2}$ chart
and primitive Hopf meridian. The integral Gysin and fixed-tube
Mayer--Vietoris calculation of P3K.22 therefore gives
$\pi_1(X)=0$ and $H^*(X;\mathbb Z)=H^*(S^6;\mathbb Z)$.
The equality $\Theta_6=0$ identifies its oriented smooth type
with the standard sphere. The exact source, Smith and gluing
matrices are independently certified in
\path{code/part3_root/reduced_iii_audit.py} and its saved output.
The complete marked gluing degrees are also exported in
\path{code/part3/reduced_iii_cap_maps.py}. The physical-character
cochain and origin arithmetic are certified in
\path{code/part3/reduced_iii_return.py}. The separate actual
parameter-return comparison, with its signed Thom map, is P3K.24.

\end{enumerate}
\end{proof}




\begin{theorem}\label{night:P3F.38}
\StatementRecord{night:P3F.38}{proved}{S6\_06\_BLIND: GPT-6 Astra (OpenAI, ChatGPT Pro mode, effort Pro), recorded re-derivation in the checkpoint dated 18 November 2026}{\cref{night:P3F.35}, \cref{night:P3F.36}, \cref{night:P3K.15}}
For the OS50 source \(I_1^*+IV+I_1\), use
\(M=\mathcal O(aG-O-L_0)\) and the simple whole zero at its
\(I_1\) value. The following actual two-cap constructions exist:
\[
\begin{array}{c|c|c}
(a,b)&\text{star cap}&\text{IV cap}\\\hline
(6,2)&(2,n_*),\ n_*\text{ odd}&(m,n_4),\ m\equiv1\pmod3\\
(2,6)&(2,n_*),\ n_*\text{ odd}&(m,n_4),\ 3\mid m
\end{array}
\]
Here \(m>0\), \(\gcd(m,n_4)=1\), and the actual elliptic,
gamma and delta collars are retained. Put
\(D=mn_*+2n_4\). Their circle orbits have
\[
 \pi_1(Y)=\mathbb Z/|D|,\qquad
 H_j(Y;\mathbb Z)=
 \begin{cases}\mathbb Z,&j=0,5,\\
 \mathbb Z/|D|,&j=1,3,\\0,&j=2,4.
 \end{cases}
\]
The actual six-manifold is a smooth sphere exactly when
\(D=1\) or \(D=-1\). Equivalently, in the first row
\(m\equiv1\pmod6\), and in the second row \(m=3k\) with
\(k\) odd; in both rows
\[
 n_*\text{ is odd},\qquad
             n_4=(\varepsilon-mn_*)/2,
                    \quad\varepsilon\in\{1,-1\}.
\]
The case \(m=1,n_4=0\) in the first row agrees with P3F.36.
No parameter-return class is assigned by this statement.
\end{theorem}

\begin{proof}
\begin{enumerate}[label=\arabic*.]
\item \par\medskip\noindent\textbf{The two actual IV roots.}
We use the marked reduced-IV equation
\[
 y^2=x^3+q^2a(q)x+q^2b(q),\qquad b(0)\ne0,
\]
where these coefficient functions are distinct from the source
integers \((a,b)\). For \((a,b)=(6,2)\), put
\(m=3\ell+1\), \(q=z^m\), \(x=z^{2\ell}X\),
\(y=z^{3\ell}Y\). The exact equation becomes
\[
 Y^2=X^3+z^{2\ell+2}a(z^m)X+z^2b(z^m).
\]
Its minimal regular central fibre is IV. The deck is
\[
 (X,Y,z)\longmapsto
 (\zeta_m^{-2\ell}X,\zeta_m^{-3\ell}Y,\zeta_mz).
\]
In the blowup chart \(X=zX_1,Y=zY_1\), the two nonidentity
components are labelled by \(Y_1=\pm\sqrt{b(0)}\).
The character on \(Y_1\) is \(\zeta_m^{-(3\ell+1)}=1\),
so the deck fixes each component label. The root line
\(N=\mathcal O(6G'-O')\) has degree zero on each component,
and \(Q'=2G'\) has component two. Multiplication by \(m\)
is an analytic local isomorphism of the smooth Neron group:
it has derivative \(m\) on the additive identity component and
is invertible on the component group \(\mathbb Z/3\).
The selected local section \(R\) with \(mR=Q'\) is therefore
unique and deck-invariant.

\item The root line is compared with the actual source line on the
resolved fibre product. Every curve contracted in relative
minimalization has degree zero for the pulled-back source line,
since its source component degree is zero. Descending through the
actual point contractions gives the root line. Its ratio with
\(\mathcal O(6G'-O')\) is generically trivial and has degree
zero on every component, so the proper-direct-image construction
of P3K.15 gives an integral whole-fibre divisor for that ratio.
Multiplication by the corresponding base power supplies the
actual line isomorphism. This retains its deck character.
The source correction \(-L_0\) is supported at the other,
starred value and has no local term here.

\item For \((a,b)=(2,6)\), put \(m=3\ell\),
\(q=z^m\), \(x=z^{2\ell}X\), \(y=z^{3\ell}Y\).
The exact root equation is
\[
 Y^2=X^3+z^{2\ell}a(z^m)X+b(z^m).
\]
Its central curve is the smooth elliptic curve
\(E_0:Y^2=X^3+b(0)\). The deck on it is
\((X,Y)\mapsto(\zeta_3X,Y)\). Its fixed points are
\(P_0=O,P_1,P_2\), where \(P_1,P_2\) are the two nonzero
three-torsion points. The original marked nonidentity component
of \(G\) identifies \(G'_0=P_1\), and \(2G'_0=P_2\).
Thus \(Q'=6G'\) has central value zero. The unique analytic
\(m\)-division \(R\) of \(Q'\) with \(R_0=0\) is
deck-invariant.

\par\medskip\noindent\textbf{The good-root line and its exact exceptional correction.}
First take the minimal cubic root \(q=s^3\). At each of the
three fixed points of
\((X,Y,s)\mapsto(\zeta X,Y,\zeta s)\), the two tangent
weights are \((1,1)\). Blow up these points by
\(\beta:\widehat S\to S_{\rm root}\), with exceptional
curves \(E_i\). In coordinates \((w/s,s)\), the quotient
coordinates are \((w/s,s^3)\). Hence the quotient
\(\pi:\widehat S\to\widetilde S\) is smooth and
\(\pi^*L_i=3E_i\). Since \(E_i^2=-1\), its image has
\(L_i^2=-3\).

\item The proper central elliptic curve has square \(-3\) after the
three blowups. Its quotient image \(C\) has square \(-1\),
because its pullback is that proper transform and the quotient
has degree three. It meets the three \(L_i\) at distinct points.
The actual smooth point contraction
\(\alpha:\widetilde S\to S\) contracts \(C\). The three
images of \(L_i\) then have square \(-2\) and meet at one
point with distinct tangent directions, giving the minimal regular
IV fibre. The invariant functions
\(q=s^3,x=s^2X,y=s^3Y\) identify the displayed Weierstrass
model off the exceptional curves. The resolved graph and relative
minimality identify this model with the original source over the
disc.

\item The section curves avoid \(C\): on the root they are transverse
to the base and their proper transforms meet \(E_i\) away from
the proper central curve. Pullback of the actual two section
divisors therefore gives
\[
 \boxed{\ \beta^*\mathcal O(2G'-O')
       =\pi^*\alpha^*\mathcal O(2G-O)+E_2-E_0\ }.             \tag{P3F.38.1}
\]
The deck weights of this line at \((P_0,P_1,P_2)\) are
\((\zeta^{-1},1,\zeta)\), up to one common character. Indeed
for \(\mathcal O(P_2)\) the canonical section is a local
parameter of weight \(\zeta\) times its frame, so the frame
has weight \(\zeta\); the inverse divisor at \(P_0\) gives
\(\zeta^{-1}\). These are exactly the weights of the
correction in (P3F.38.1).

\item Translation by \(Q'=6G'\) fixes all three central points and
preserves \(E_2-E_0\). Its canonical divisor lift combined with
the pulled-back source lift therefore acts on the right side of
(P3F.38.1), and descends through \(\beta\) to the actual map
on \(N\). This agrees with the original punctured lift and
commutes with the deck. For \(m=3\ell\), pull this entire
correspondence through \(s=z^\ell\) and resolve the fibre
product. The line identity pulls back exactly; all extra
exceptional terms are its actual pullbacks.

\item The nonzero fixed central Picard point gives the generator of the
order-three Euler torsion of the cap. Explicitly, for the Picard
torus \(V/\Lambda\), the connecting map is
\[
 (V/\Lambda)^D\longrightarrow\Lambda/(D-1)\Lambda,
                  \qquad[v]\longmapsto[(D-1)v].
\]
It is an isomorphism because \(D-1\) is invertible on \(V\),
and its target has order three. The line point \(P_2\) is
nonzero, so its clutch is the nonzero torsion class. A common
deck character has zero absolute first Chern class but remains
in the actual line lift.

\par\medskip\noindent\textbf{The commuting norm lifts and free quotients.}
For either IV root, first lift \(t_{-R}\) on \(N\), before
taking its \(n_4\)-th power. The actual meromorphic chord section
of its Hom line has an integral vertical correction. The
component-degree equations make that correction a whole fibre;
multiplying by its base power gives an invertible holomorphic
line map \(F\). Average in the deck-conjugation action on the
rank-one Hom module, choosing its nonzero central character.
The resulting lift is still invertible after shrinking and satisfies
\[
                   DFD^{-1}=\chi F,\qquad\chi^m=1.
\]
The maps \(F^m\) and the specified \(B_0\) cover the same
translation \(t_{-Q'}\), so
\(B_0=\beta(z)F^m\) for a base unit. Its actual descent gives
\(\beta(\zeta_mz)=\beta(z)\). Thus \(F,D\) both commute
with \(B_0\), and
\[
 (F^{n_4}D)^m/B_0^{n_4}
              =\chi^{n_4m(m-1)/2}\beta(q)^{-n_4}.              \tag{P3F.38.2}
\]
This is an invariant nonvanishing unit. Divide by its analytic
\(m\)-th root to obtain \(A_0\), retaining the displayed
character factor, so
\(A_0^m=B_0^{n_4}\) and \([A_0,B_0]=1\).
The source map is \(t_{-n_4R}D\).

\item Combine this with the actual star construction of P3F.36. For
\(L=\operatorname{lcm}(2,m)\), use the parameter map
\(\lambda=r^L\) and multiply the star and IV lifts respectively
by \(r^{Ln_*/2}\) and \(r^{Ln_4/m}\). Their exact norm
identities hold for \(B_\lambda\). The primitive generators
from Bezout have positive scalar powers \(r^{L/2}\) and
\(r^{L/m}\); their remaining maps and inverses have bounded
norm ratios on compact smaller root tubes. They are therefore
uniformly contracting for small \(r\ne0\). Their quotients
are smooth, proper and compact over the original discs by the
actual norm shells. A unit line rotation cannot equal a nonzero
iterate of such a contraction, so the delta action is free there.
The line correspondences identify their punctured collars with
the actual exterior. This constructs the stated global families
under the displayed parameter map, rather than inferring them
from their homology.

\item The positive angular convention is the same at both caps. On a
root of order \(m_i\), write \(A_i:(\theta,t)\mapsto
(\theta+1/m_i,t-n_i/m_i)\), with \(B:t\mapsto t-1\).
The original coordinates are \(x=m_i\theta,y=t+n_i\theta\).
Thus \(\alpha_i\mapsto-n_ih_i\),
\(\gamma\mapsto m_ih_i\), and the killed meridian is exactly
\(m_i\alpha_i+n_i\gamma\).

\par\medskip\noindent\textbf{The actual common source complement.}
Remove both additive source tubes and call the remaining surface
\(P\). Its fundamental group is exactly the base annulus group.
To see this without losing an elliptic kernel, let the starred
monodromy be \(-T_u\), and let \(T_v\) be the monodromy of
the remaining \(I_1\), for primitive cycles \(u,v\). The IV
monodromy has trace \(-1\). The product relation and transvection
formula give
\[
 -1=-\operatorname{tr}(T_u^{-1}T_v^{-1})
                         =-2+(u\cdot v)^2.
\]
Hence \(|u\cdot v|=1\). The nodal vanishing disc kills \(v\),
and conjugation by the retained starred meridian kills
\((-T_u)v=-v\pm u\). Thus the elliptic lattice dies, leaving
\(\pi_1(P)=\mathbb Z\alpha\), with
\(\alpha_4=\alpha_*^{-1}\). Translation fixes this generator,
so \(\pi_1(W)=\mathbb Z\alpha\oplus\mathbb Z\gamma\).

\item We compute the homology from actual component inclusions. In the
NS basis of P3F.36, the form \(B\) is fixed by
\(O^2=G^2=-1\), \(O\cdot F=G\cdot F=1\), the negative
\(D_5+A_2\) root matrix, and the intersections of \(G\) with
\(R_0,A_{2,1}\); the other pairings are zero. Its determinant
is \(-1\). Put
\[
 L_*=(L_0,L_1,R_0,R_1,C_0,C_1),\quad
 L_4=(F-A_{2,1}-A_{2,2},A_{2,1},A_{2,2}),\quad L=(L_*,L_4),
\]
where \(L_0=F-2C_0-2C_1-L_1-R_0-R_1\).
Both local fibre neighbourhoods have free homology only in degrees
zero and two. The cone of \(C(S)\to C(S,P)\), shifted once,
therefore gives the actual complement complex
\[
 C_4=\mathbb Z,\ C_3=\mathbb Z^2,\ C_2=\mathbb Z^{10},
 \ C_1=\mathbb Z^9,\ C_0=\mathbb Z,
 \qquad d_4=(1,1)^t,\quad d_2=L^tB.                         \tag{P3F.38.3}
\]
All other differentials are zero. The two complexes before taking
this cone are even and free, so their degree-one homotopy group
is zero. This makes the construction functorial for the actual
translation, with no omitted annular term.

\item For \(b=6\), the actual section intersections are
\(O\cdot6G=G\cdot6G=1\), \(O\cdot7G=2\),
\(G\cdot7G=1\). The sections meet respectively \(L_1\)
with the identity IV component, and \(R_1\) with \(A_{2,1}\).
Solving against \(B\) determines their integral vectors and
hence the translation matrix. For \(b=2\) use the actual
\(2G,3G\) classes of P3F.36. In each case invert the positive
translation matrix to obtain the contracting map \(t_{-bG}\).
The certificates check the isometry of \(B\), all component
permutations, and the equivariance of (P3F.38.3).

\item For each boundary use the complex with sole differential
\(L_i^tBL_i:\mathbb Z^{k_i}\to\mathbb Z^{k_i}\).
Its map into (P3F.38.3) is \(L_i\) in degree two, the coordinate
inclusion into \(\mathbb Z^9\) in degree one, and the appropriate
coordinate of \(\mathbb Z^2\) in degree three. Taking the actual
mapping-torus cones gives every boundary map. We obtain
\[
 H_1(W)=\mathbb Z^2,\quad H_3(W)=\mathbb Z^2,\quad H_4(W)=\mathbb Z,
 \qquad
 H_2(W)=\begin{cases}\mathbb Z^2\oplus\mathbb Z/2,&b=2,\\
                         \mathbb Z^2\oplus\mathbb Z/6,&b=6.
          \end{cases}
\]

\par\medskip\noindent\textbf{Every integral degree-two map.}
For \(b=6\), the boundary bases are
\((\epsilon_*,r_*,s_*)\), of orders \((2,0,0)\), and
\((\epsilon_4,r_4,s_4)\), of orders \((3,0,0)\).
In the exterior basis \((\tau,g_1,g_2)\), of orders
\((6,0,0)\), the actual matrix is
\[
 \begin{pmatrix}
 -15&-5&83&-4&5&-4\\0&0&-1&0&0&0\\0&0&-2&0&0&1
 \end{pmatrix}.
\]
The fibre classes are \(F_*=-\epsilon_*+2r_*\) and
\(F_4=r_4\), both mapping to \(5\tau\). The star fibre
maps to \(2C_*\), and the good-root fibre maps to \(C_4\)
with degree one. Thus the cap rows are
\((0,1,\kappa_*)\) and \((0,1,\kappa_4)\), for their actual
integral collar coefficients. The torsion map sends
\(\epsilon_*\) to \(3\tau\) and \(\epsilon_4\) to
\(2\tau\). The exact Chinese-remainder map
\(\mathbb Z/6\to\mathbb Z/2\oplus\mathbb Z/3\) is
\(1\mapsto(1,-1)\). In this single torsion generator followed
by the four free generators, the complete map is
\[
 D_2^{\rm good}=\begin{pmatrix}
 1&1&5&5&2\\0&0&-1&0&0\\0&0&-2&0&1\\
 0&-1&-\kappa_*&0&0\\0&0&0&-1&-\kappa_4
 \end{pmatrix}.                                               \tag{P3F.38.4}
\]
Its first row is read modulo six; these are the exact preceding
entries reduced by \(6\tau=0\). Its determinant is one.

\item For \(b=2\), the star boundary basis is the same, while the IV
boundary has two free generators \((r_4,s_4)\). The exterior
basis has orders \((2,0,0)\), and its matrix is
\[
 \begin{pmatrix}-5&2&33&3&9\\0&0&1&0&-2\\0&0&-2&0&3\end{pmatrix}.
\]
Here \(F_*=-\epsilon_*+2r_*\), \(F_4=3r_4\). The reduced
IV root has three rational components permuted transitively, so
its fibre maps to \(3C_4\). Its cap row is again
\((1,\kappa_4)\). Thus the complete matrix, with its first row
read modulo two, is
\[
 D_2^{\rm red}=\begin{pmatrix}
 1&0&1&1&1\\0&0&1&0&-2\\0&0&-2&0&3\\
 0&-1&-\kappa_*&0&0\\0&0&0&-1&-\kappa_4
 \end{pmatrix}.                                               \tag{P3F.38.5}
\]
Its free core block has determinant
\(\det\left(\begin{smallmatrix}1&-2\\-2&3\end{smallmatrix}\right)=-1\),
and the full map is unimodular. Both maps are therefore integral
isomorphisms for their actual collars. No source inversion has
been extended across a marked cap.

\par\medskip\noindent\textbf{The other maps and the specialization indices.}
The good-root generic elliptic fibre transports through its smooth
root family to \(E_0\) with degree one. On \(H_1\), that
same transport is the identity lattice \(\Lambda\). If
\(u n_4+vm=1\), the cap monodromy is \(D^u\). The two
boundary and cap torsion sublattices are equal integrally:
\[
 (D^u-1)\Lambda=(D-1)\Lambda,
\]
because \(u\) is a unit modulo three and, for \(u\equiv2\),
\(D^2-1=(D-1)(D+1)\) with \(D+1=-D^2\) unimodular.
Thus the actual boundary \(\mathbb Z/3\) maps isomorphically
to cap torsion. The starred \(\mathbb Z/2\) map is the
primitive graph specialization of P3F.36. For the reduced IV cap
both groups have no degree-one torsion.

\item Retain all path coefficients in the free columns. The preceding
torsion isomorphisms eliminate them without changing any free
relation. In the positive meridian basis
\((\alpha_*,\gamma_*,\alpha_4,\gamma_4)\), the free
Mayer--Vietoris matrix is
\[
 D_1=\begin{pmatrix}
 1&0&-1&0\\0&1&0&1\\n_*&-2&0&0\\0&0&n_4&-m
 \end{pmatrix}.                                               \tag{P3F.38.6}
\]
Its determinant is \(-D\). In degree three use the oriented
source boundaries and their positive fibre--gamma traces. The
exterior rows are \((0,1,0,1)\) and \((-1,0,1,0)\); in
the reduced case reverse the certificate's second target generator
to obtain this same convention. The cap rows are
\((-n_*,2)\) and \((-n_4,m)\), by the actual norm powers.
Graph corrections lie inside the two-dimensional central source
and add no three-class. Hence
\[
 D_3=\begin{pmatrix}
 0&1&0&1\\-1&0&1&0\\n_*&-2&0&0\\0&0&n_4&-m
 \end{pmatrix},\qquad\det D_3=-D.                           \tag{P3F.38.7}
\]
The degree-four map is \((1,-1)\), and the degree-zero map is
\(\left(\begin{smallmatrix}1&1\\-1&0\\0&-1\end{smallmatrix}\right)\).
These specify every nonzero map. The chain matrices, representatives,
source meridians and oriented fibre traces are executed in
\path{code/part3/i1star_iv_boundary.py} and
\path{code/part3/i1star_iv_reduced_boundary.py}, with their JSON
outputs.

\item For van Kampen, all elliptic periods die in \(W\). Their cap
images generate the star graph group and, in the good case, the
whole elliptic lattice. The reduced-IV central curve is simply
connected. After killing these normal subgroups, the exact remaining
presentation is
\[
 \langle h_*,h_4\mid h_*^2=h_4^m,
                         \ h_*^{n_*}=h_4^{-n_4}\rangle.
\]
Choose \(r_0n_*+2s_0=1\). Both right-hand powers are powers of
\(h_4\), so \(h_*=h_4^{-r_0n_4+s_0m}\). Thus the group is
cyclic, and the two exact relations give \(\mathbb Z/|D|\).
The matrices give the asserted integral homology. Here \(D\ne0\):
otherwise primitivity would imply \(m\mid2\), while the allowed
\(m\) and odd \(n_*\) exclude zero determinant also at \(m=1\).

\item For an arbitrary allowed one-vertex cusp fan, use the finite sequence
of actual periodic diagonal flops of P3F.35. Each step is the common
blowup of its embedded normal \(\mathcal O(-1)\oplus\mathcal O(-1)\)
curve followed by the other ruling contraction. These maps are the
identity outside the cusp, so they preserve the present marked caps.
Their integral blowup decompositions and van Kampen maps preserve
homology and fundamental group. Thus the sphere question is exactly
the same as for the vertical fan used in the orbit calculation.

\item Finally, the unchanged whole-zero \(I_1\) has the same actual
fixed sphere and primitive Hopf meridian as P3F.36. Its Gysin and
fixed-tube gluing gives the smooth six-sphere when \(|D|=1\),
and gives \(\pi_1(X)=\pi_1(Y)\) for every tuple, excluding
all other determinants. The stated arithmetic form follows because
\(|D|=1\) forces odd \(m\), and then
\(2n_4\equiv\varepsilon\pmod m\) automatically gives
\(\gcd(m,n_4)=1\). The actual line and norm lifts, rather than
only this arithmetic, were constructed above.

\end{enumerate}
\end{proof}




\begin{theorem}\label{night:P3F.39}
\StatementRecord{night:P3F.39}{proved}{S6\_06\_BLIND: GPT-6 Astra (OpenAI, ChatGPT Pro mode, effort Pro), recorded re-derivation in the checkpoint dated 18 November 2026}{\cref{night:P3F.35}, \cref{night:P3F.36}, \cref{night:P3F.38}, \cref{night:P3F.8}, \cref{night:P3K.15}, \cref{night:P3K.16}, \cref{night:P3K.19}}
For any allowed one-vertex cusp fan, take the star
core of P3F.36 or either star/IV core of P3F.38, retaining the other
singular source values. Add any finite collection of allowed
logarithmic caps at smooth source values, with their full elliptic,
gamma and delta collars. For a sphere member, every additional
smooth-site cap has effective order one after the exact
source-and-line finite-kernel quotient of P3K.19. The integral
order-one clutches are retained by this assertion.
\end{theorem}

\begin{proof}
\begin{enumerate}[label=\arabic*.]
\item We first use the actual finite flop chain of P3F.35 to reach its
vertical-edge fan. Each step blows up one embedded curve with normal
line sum \(\mathcal O(-1)\oplus\mathcal O(-1)\) and contracts
the other ruling of the exceptional \(\mathbb P^1\times\mathbb P^1\).
The maps agree with the identity outside the cusp and retain every
marked cap. Their integral blowup decompositions and van Kampen
maps preserve homology and fundamental group. A sphere for the
original fan therefore gives a sphere for this vertical fan, whose
orbit is the standard fixed-sphere orbit of P3F.8. We now work with
its actual source exterior before any cap is filled. No inversion
is assumed to extend across a marked cap.

\par\medskip\noindent\textbf{The prime forced by a nontrivial smooth order.}
The actual projection of P3K.16 has a closed Seifert three-manifold
as base. With its retained integral gamma clutch \(b_0\), the
first-homology determinant is
\[
 D_{\rm base}=b_0\prod_i m_i+
                         \sum_i n_i\prod_{j\ne i}m_j.        \tag{P3F.39.1}
\]
The integral clutch can equivalently be represented by an order-one
site; it has not been discarded. Loops in this base lift away from
its finitely many discriminant and core circles and close in the
connected elliptic fibre, as explicitly proved in P3K.16. A sphere
member has the five-sphere orbit of P3F.8, so its base has trivial
fundamental group. Hence \(|D_{\rm base}|=1\). Two cap orders
cannot have a common prime divisor, since that prime would divide
every summand of (P3F.39.1).

\item At a smooth root we first use the actual quotient of P3K.19. If
\(g\) is the finite zero-gamma order, its invariant source-and-line
coordinates give the map \((e,z)\mapsto(e,z^g)\), and the line
is pulled back by this same map. The resulting norm pair is
primitive. We use that proved morphism here, including its marked
meridians, rather than forgetting a finite isogeny.

\item Suppose an additional smooth-site order \(m_i>1\) has prime
divisor \(p\). Its primitive norm numerator is nonzero modulo
\(p\). The existing starred order two makes \(p\ne2\).
For the good-IV core the existing order is divisible by three,
so also \(p\ne3\). Thus \(p\) does not divide the source
translation coefficient: \(b=2\) for the star and reduced-IV
cores, and \(b=6\) for the good-IV core. Put \(k=\mathbb F_p\).

\par\medskip\noindent\textbf{Source parity and the exact rim sequence.}
Let \(W_0\) be the mapping torus of \(t_{-bG}\) on the source
complement of its additive fibre neighbourhoods. The computed
integral groups and the actual source inversion followed by
suspension reversal are
\[
\begin{array}{c|c|c}
\text{core}&H_2(W_0;\mathbb Z)&I_*\\\hline
\text{star}&\mathbb Z/2\oplus\mathbb Z&
 \left(\begin{smallmatrix}1&2\\0&-1\end{smallmatrix}\right)\\
\text{reduced IV}&\mathbb Z/2\oplus\mathbb Z^2&
 \left(\begin{smallmatrix}-1&216&138\\0&-1&0\\0&0&-1\end{smallmatrix}\right)\\
\text{good IV}&\mathbb Z/6\oplus\mathbb Z^2&
 \left(\begin{smallmatrix}-11&120&2\\0&-1&0\\0&0&-1\end{smallmatrix}\right).
\end{array}                                                   \tag{P3F.39.2}
\]
Here \(I[p,t]=[-p,-t]\). The matrices are obtained from the
actual source and cone maps in P3F.36 and P3F.38, and are exported
by their certificates. Their first generators are the displayed
torsion generators. Since \(p\nmid b\), torsion disappears
over \(k\), and \(I_*=-1\) on all \(H_2(W_0;k)\).

\item Remove the \(r\) smooth source-fibre tubes
\(N_i=T^3\times D^2\), and call the result \(W\). On a
regular source disc choose a real lift of the actual translation
section. The mapping-torus coordinate change
\(p'=p-tQ\), for the contracting source map \(t_{-Q}\),
identifies the tube with \(T^3\times D^2\). In this exact
coordinate, \(I\) sends each of \(a_i,b_i,\gamma\) to its
negative and preserves the positive normal meridian \(\alpha_i\).
It reverses the orientation of the five-dimensional exterior and
preserves the orientation of the normal disc.

\item Every core two-class is therefore \(I\)-even. Its image in
\(H_2(W_0;k)\), which is \(I\)-odd, is zero because \(p\)
is odd. The Thom map
\(H_3(W_0;k)\to\bigoplus_iH_1(T^3;k)\) is dual to
intersection with these core two-classes, and is consequently
zero. The Thom map
\(H_2(W_0;k)\to\bigoplus_iH_0(T^3;k)\) is also zero.
Indeed each core fundamental three-class is \(I\)-odd; a two-class
and such a three-class both change sign, whereas their ambient
orientation changes sign once. Naturality of the intersection
pairing gives \(x\cdot z=-x\cdot z\), so it vanishes over
\(k\).

\item The actual Thom sequence now gives
\[
 0\longrightarrow R\longrightarrow H_2(W;k)
          \longrightarrow H_2(W_0;k)\longrightarrow0,
\]
\[
 R=\bigoplus_i\langle
 A_i=a_i\alpha_i,\ B_i=b_i\alpha_i,\ C_i=\alpha_i\gamma
 \rangle_k.                                                   \tag{P3F.39.3}
\]
All three rim generators are \(I\)-odd. The idempotents
\((1\pm I)/2\) split the two eigenspaces over \(k\); the
end terms of (P3F.39.3) are odd, so the middle term is entirely
odd as well. Hence the actual smooth-boundary core columns
\[
                      a_ib_i,\qquad a_i\gamma,\qquad b_i\gamma
\]
map to zero in \(H_2(W;k)\). This uses \(I\) only on the
source exterior.

\par\medskip\noindent\textbf{Elimination of the additive fillings.}
The additive degree-two Mayer--Vietoris map to
\(H_2(W_0)\) plus the additive-cap groups is an integral
isomorphism. For the star it is (P3F.36.8), retaining its actual
\(\kappa\). For the joint cores it is (P3F.38.4) or
(P3F.38.5), retaining both actual collar coefficients and their
respective torsion maps. Over \(k\), the projection of the
additive boundary map from \(H_2(W;k)\) onto this quotient is
still that isomorphism. Its image has zero intersection with
\(R\), and the quotient by that image is exactly \(R\).
This is the elimination of a graph over the quotient in
(P3F.39.3); it asserts no extension of \(I\) over a cap.
The smooth-boundary core columns remain zero, while their three
rim columns remain the independent \(A_i,B_i,C_i\).

\par\medskip\noindent\textbf{The actual smooth-cap kernel.}
The primitive smooth cap has the full marked meridian
\[
 \mu_i=m_i\alpha_i+n_i\gamma+
                         \lambda_{ai}a_i+\lambda_{bi}b_i.
\]
The norm-shell map of P3K.15 identifies its orbit with
\(T^3\times D^2\). The actual first-homology map is the
quotient by \(\mu_i\); its degree-two map is the exterior
square of that quotient. The delta coordinate maps to zero under
the actual line-circle projection, while its character remains
in the principal-circle line. No elliptic period is omitted.

\item Over \(k\), the kernel of the degree-two map is precisely
\(\mu_i\wedge H_1(T^4;k)\). Eliminating the cap target
identifies the remaining cokernel with \(R\) modulo the images
of these kernel vectors. Since the even core columns vanish,
these images have, up to the chosen wedge orientation signs, matrix
\[
 \begin{pmatrix}
 m_i&0&0&-\lambda_{ai}\\
 0&m_i&0&-\lambda_{bi}\\
 0&0&m_i&n_i
 \end{pmatrix}                                                \tag{P3F.39.4}
\]
in the rim basis \((A_i,B_i,C_i)\). For the selected
\(p\mid m_i\), its first three columns vanish and its final
column is nonzero because \(p\nmid n_i\). Its rank is exactly
one. Two independent classes therefore survive in this actual
rim summand. Different smooth sites have disjoint summands, so
their kernel images cannot remove these classes.

\item The cokernel of the full degree-two Mayer--Vietoris map injects
into \(H_2(Y;k)\). It is nonzero, contradicting the five-sphere
orbit of a sphere member. Thus every additional effective smooth
order is one. The source inversion matrices and the exact rank-one
kernel calculation are executed in
\path{code/part3/i1star_smooth_caps.py} and its JSON output.
The proof leaves all integral order-one clutches and their
elliptic and delta coordinates in place.

\end{enumerate}
\end{proof}




\begin{theorem}\label{night:P3F.40}
\StatementRecord{night:P3F.40}{proved}{S6\_06\_BLIND: GPT-6 Astra (OpenAI, ChatGPT Pro mode, effort Pro), recorded re-derivation in the checkpoint dated 18 November 2026}{\cref{night:P3F.35}, \cref{night:P3K.22}}
Take any of the ten primitive-translation sources of P3K.22,
with its specified line, simple whole zero and ordinary remaining
source values. Replace one nonzero-linearization multiplicative
fibre by an actual transitive root-line cap. After the exact
finite-kernel cancellation of P3K.19, write its primitive norm
pair as $(m,n)$, with $m>0$, and retain its elliptic and delta
collars. The degree-two Mayer--Vietoris cokernel of its circle
orbit is exactly
\[
                         (\mathbb Z/m)^2.              \tag{F40.1}
\]
For $n\ne0$ this is $H_2(Y;\mathbb Z)$, and
$\pi_1(Y)=\pi_1(X)=\mathbb Z/|n|$. The sphere members among
these actual one-site constructions have exactly $m=1$ and
$n=\pm1$. The assertion concerns one multiplicative replacement;
it does not identify simultaneous replacements or their returns.
It holds for every allowed one-vertex cusp fan.
\end{theorem}
\begin{proof}
\begin{enumerate}[label=\arabic*.]
\item \par\medskip\noindent\textbf{The marked boundary and its ordinary comparison.}
Let the fibre be $I_N$. Its marked generator has Tate valuation
$c_G$ prime to $N$. The physical outward suspension is the
inverse of the contracting lift and covers $+G$. Set $c=-c_G$.
In the actual boundary generators $a,b,\alpha,\gamma$, where
$a$ is the vanishing circle, $b$ the graph period and $\alpha$
the positive base meridian, the relations are
\[
 [\alpha,b]=a^N,\qquad[\alpha,\gamma]=a^{-c},\qquad
 [a,b]=[a,\gamma]=[b,\gamma]=1.                       \tag{F40.2}
\]
Indeed the Tate periods have first coordinates $Ns$ and
$c_Gs$ plus single-valued units, with $s=\log q/(2\pi i)$.
The indicated commutators are their actual changes under
$s\mapsto s+1$. Thus the boundary $E$ is the $a$-circle bundle
over $T^3_{b,\alpha,\gamma}$ with Euler form
$d\alpha\wedge(N\,db-c\,d\gamma)$.

\item Equivalently it is the three-torus mapping torus with matrix
\[
 J=\begin{pmatrix}1&N&-c\\0&1&0\\0&0&1\end{pmatrix}
\]
over $\alpha$. In its degree-two exterior-square quotient the
sole relation is
$c(a\wedge b)+N(a\wedge\gamma)=0$.
Choose integers $x,y$ with $Nx-cy=1$. The following are
actual integral degree-two generators:
\[
 \begin{split}
 z&=x(a\wedge b)+y(a\wedge\gamma),\qquad
 R=a\wedge\alpha,\\
 B&=b\wedge\gamma,\qquad
 L=(cb+N\gamma)\wedge\alpha.
 \end{split}                                                \tag{F40.3}
\]
The two suspended loops in (F40.3) are fixed by $J$, so their
tori are actual cycles. The change from the relation and $z$
to $(a\wedge b,a\wedge\gamma)$ has determinant $-1$.
Consequently $H_2(E)=\mathbb Z^4$, with
\[
 F=Nz,\qquad a\wedge\gamma=-cz.
\]
For boundary orientation $(a,b,\gamma,\alpha)$ its intersection
matrix in the displayed order is
\[
 Q_E=\begin{pmatrix}0&0&0&1\\0&0&1&0\\0&1&0&0\\1&0&0&0
 \end{pmatrix}.                                           \tag{F40.4}
\]
These signs follow by wedging the actual fibre and suspended
cycles: $zL=Nx-cy=1$ and $RB=1$.

\item Let $P$ be the complement of this fibre tube and
$W=T_{t_G}P$, using the physical outward suspension. The
component lattice is primitive. In fact it is a primitive
direct sublattice of the full trivial lattice, and the latter
has torsion-free quotient in $\operatorname{NS}(S)$ because
$\operatorname{MW}(S)=\mathbb Z$. Excision and the unimodular
source form give $H_1(P)=0$. The elliptic fibration with section
over the complementary disc presents its fundamental group as
a quotient of the abelian elliptic lattice; hence $\pi_1(P)=0$
and $\pi_1(W)=\mathbb Z\gamma$.

\item Denote the unchanged ordinary cap by $U$. Translation cycles
all $N$ components of its central fibre and preserves its
oriented graph circle. Its integral Wang sequence gives
$H_1(U)=\mathbb Z^2$, $H_2(U)=\mathbb Z^2$ and $H_3(U)=\mathbb Z$.
Write $C$ for the component-orbit class and $T$ for a lift of
graph times suspension. Then $F\mapsto NC$. The invariant sum
of the components maps to $F$ in the source Wang sequence, so
\[
 H_3(U)\longrightarrow H_3(T_{t_G}S)=\mathbb ZF
\]
is onto with degree one. The section generator of
$H_2(T_{t_G}S)=\mathbb Z[O]$ has boundary the primitive
meridian $\alpha$, since it meets the removed fibre once.
In the ordinary cap $\alpha$ bounds its section disc, while
$b,\gamma$ give the two primitive generators. These two
facts in the full-source Mayer--Vietoris sequence prove that
\[
 H_2(E)\longrightarrow H_2(W)\oplus H_2(U)              \tag{F40.5}
\]
is an integral isomorphism. In particular $H_2(W)=\mathbb Z^2$.
The ordinary cap map is $z\mapsto C$, $R,L\mapsto0$ and
$B\mapsto T$, after choosing the lift $T$. Thus the exterior
images of $R,L$ form an integral basis.

\par\medskip\noindent\textbf{The actual line frame and the zero fibre class.}
Write $M=\mathcal O(dG-O)$ as in P3K.22. Its degree on each
component of this retained fibre is zero. Its actual Tate Picard
representative $P(q)$ is a nonvanishing holomorphic unit, and
its line transition is
\[
                    (Z,V)\sim(q_TZ,P(q)V).
\]
Choose an analytic logarithm of $P(q)$. The expression
\[
 V\exp\left(-\log P(q)\frac{\log|Z|}{\log|q_T|}\right)
                                                               \tag{F40.6}
\]
is unchanged by this transition and gives a nowhere-zero smooth
frame of $M$ on the source boundary. A translation lift
$(Z,V)\mapsto(G(q)Z,\beta(q)V)$ acts in that frame by
\[
 \beta(q)\exp\left(-\log P(q)
                         \frac{\log|G(q)|}{\log|q_T|}\right).
                                                               \tag{F40.7}
\]
The degree-zero Tate charts have unit line transitions on every
component, so $\beta$ is the actual nonzero holomorphic unit
at this retained value. The exponential in (F40.7) has a
single-valued logarithm. Its graph and vanishing windings are
zero, and its base winding is zero as well.

\item Over the complementary disc, divide the specified degree-one
Hom section by an equation $g$ for its sole simple zero. This
gives an invertible lift on $P$. In the frame (F40.6) its ratio
differs from (F40.7) only by $g^{-1}$, whose boundary winding
is one signed base meridian. It has no graph winding. Form
the relative class $c_1(M,\text{frame})\in H^2(P,\partial P)$
and let $K\in H_2(P)$ be its Poincare--Lefschetz dual. The
relative Chern connecting map therefore gives
\[
                        (t_G)_*K-K=\pm F.              \tag{F40.8}
\]
Here the dual of $d\alpha$ on the source boundary is exactly
its fibre $F$; a graph winding would instead have given a rim
term, and (F40.7) proves that winding is zero. Equation (F40.8)
is a compact relative-cycle identity, not an assertion just in
the closed source. Hence $F=0$ in $H_2(W)$. Its torsion-freeness
and $F=Nz$ give $i(z)=0$ in the exterior.

\item Use $i(R),i(L)$ as the exterior basis and write
$i(B)=e\,i(R)+f\,i(L)$. The kernel of a boundary inclusion
into an oriented compact five-manifold is isotropic for the
boundary intersection form: bounding relative three-chains
give this immediately by the boundary intersection identity.
Pair $z$ with $B-eR-fL$ in (F40.4); this gives $f=0$.
The square of $B-eR$ is then $-2e$, so $e=0$. The complete
actual exterior map is thus
\[
                 i_2=\begin{pmatrix}0&1&0&0\\0&0&0&1\end{pmatrix}.
                                                               \tag{F40.9}
\]

\par\medskip\noindent\textbf{The marked cap and every degree-two column.}
Retain its primitive meridian
\[
                     \mu=m\alpha+jb+n\gamma.          \tag{F40.10}
\]
An additional vanishing-circle coefficient has zero class in
$H_1(E)$ because $\gcd(N,c)=1$ in (F40.2). The actual delta
phase stays in the line lift and has no orbit column.
Put $A=Nj-cn$. In the Tate root the affine deck has component
displacement $-A$: its norm differs from $B^n$ on the Tate
cover by the period $-j$, exactly as (F40.10) records. The
suspension component displacement is $mc$. Their component
orbits in $\mathbb Z/(Nm)$ therefore have number
\[
           \gcd(Nm,A,mc)=\gcd(m,A).
\]
The assumed actual transitive cap gives $\gcd(m,A)=1$.

\item Its central-root Wang sequence again gives
$H_2(\mathrm{cap})=\mathbb ZC\oplus\mathbb ZT$.
The generic fibre specializes to all $Nm$ components, so
$z\mapsto mC$. Choose integers $p,q$ with $mp+nq=-j$.
The actual degree-one cap map, retaining these graph-path
choices, is
\[
 \alpha\mapsto p\eta-nh,\qquad b\mapsto\eta,
                    \gamma\mapsto q\eta+mh.           \tag{F40.11}
\]
It is the quotient by (F40.10); the primitive graph and outward
suspension mark its two target generators. Its exterior square
computes projection to the $T$ summand. Thus $B\mapsto mT$
and $L\mapsto AT$, while $z,R$ have zero $T$ projection.
Write the still undetermined $C$ coefficients of $R,B,L$
as $r,s,t$. The isotropic-kernel identity for this actual cap
map, using (F40.4) and clearing denominators in its kernel
cycles, gives
\[
                         r=-A,\qquad mt=As.            \tag{F40.12}
\]
Equivalently multiply its matrix by $Q_E^{-1}$ and its transpose;
the off-diagonal entry is $m(r+A)$ and the first diagonal
entry is $2(mt-As)$. Both vanish. Since $m,A$ are coprime,
$s=mk$ and $t=Ak$ for an actual integer $k$. The determinant-one
target change $T_{\rm new}=T+kC$ retains this coefficient and
makes the cap matrix
\[
                 j_2=\begin{pmatrix}m&-A&0&0\\0&0&m&A\end{pmatrix}.
                                                               \tag{F40.13}
\]
No exterior inversion has been imposed on a marked cap.

\item The complete degree-two map is $(i_2,-j_2)^t$.
Eliminating its exterior basis columns $R,L$ leaves exactly
$\operatorname{diag}(m,m)$ on the kernel basis $z,B$.
This proves (F40.1), including both torsion summands and all
integer collar coefficients.

\item In the degree-one boundary order $(b,\alpha,\gamma)$ the
complete matrix is
\[
 D_1=\begin{pmatrix}0&0&1\\-1&-p&-q\\0&n&-m\end{pmatrix},
                         \qquad\det D_1=-n.            \tag{F40.14}
\]
For $n\ne0$ it is injective, so Mayer--Vietoris identifies
$H_2(Y)$ with (F40.1). The cap central mapping torus has
fundamental group $\mathbb Z^2$, since its source graph is
a circle preserved in orientation. In van Kampen the exterior
kills $b,\alpha$, hence $\eta=1$, $h^n=1$ and
$\gamma=h^m$. Coprimality of $m,n$ gives precisely
$\pi_1(Y)=\mathbb Z/|n|$. For $n=0$, primitivity already
forces $m=1$, and the same presentation is infinite cyclic.

\item A sphere therefore requires $m=1,n=\pm1$: (F40.1) injects
into orbit $H_2$ for every numerator, and the fundamental group
gives the other restriction. For these two numerators all
degree-one and degree-two groups vanish, and the closed oriented
five-manifold duality gives the integral sphere groups. The
same smooth classification and actual primitive Hopf-meridian
Gysin/fixed-tube maps as P3K.22 identify the orbit with $S^5$
and the total space with the oriented standard $S^6$.
The Hopf meridian also identifies $\pi_1(X)$ with $\pi_1(Y)$.
P3F.35's actual blowup/flop chain, equal to the identity outside
the cusp, preserves these marked caps and the sphere criterion
for every allowed fan.

\item The executable certificate
\path{code/part3_root/multiplicative_cap_boundary.py}
constructs both mapping-torus differentials, the integral cycle
bases, their intersection forms, all cap and exterior matrices,
and the retained target change. It checks every primitive-source
multiplicative marking and the exact symbolic determinant and
Smith form. The relative Chern framing above has an independent
audit recorded with the statement.

\end{enumerate}
\end{proof}




\begin{theorem}\label{night:P3F.41}
\StatementRecord{night:P3F.41}{proved}{S6\_06\_BLIND: GPT-6 Astra (OpenAI, ChatGPT Pro mode, effort Pro), recorded re-derivation in the checkpoint dated 18 November 2026}{\cref{night:P3F.37}, \cref{night:P3K.24}, \cref{night:P3K.26}}
For each of the four marked sources and the eight source/root sectors in the proof, let $M=\mathcal O(aG-O)$ and $Q=bG$. For every indicated root order $m$ and every nonzero integer $n$ with $\gcd(m,n)=1$, the displayed maps construct a proper holomorphic family with scalar $\lambda=r^m$, one simple whole-fibre $I_1$ zero, and marked cap meridian $mq+n\gamma_{\rm out}$. Its orbit has fundamental group $\mathbb Z/|n|$ and the integral homology displayed in \textup{(SC.4)}. Its total space is the standard oriented $S^6$ exactly for $n=\pm1$. The actual line maps and every gluing degree are retained below.
\end{theorem}
\begin{proof}
\begin{enumerate}[label=\arabic*.]
\item \par

\medskip\noindent\textbf{SC.1. Explicit sources and the parameter domains}

\item The following source models retain their actual marked section:

\par\medskip{\small\noindent
\begin{tabularx}{\linewidth}{@{}XXX@{}}
\toprule
Source & Equation and G & Fibres \\
\midrule
E7-A & y\textsuperscript{2}=x\textsuperscript{3}+tx+1, G=(0,1) & III* at infinity, three I1 \\
E7-B & y\textsuperscript{2}=x\textsuperscript{3}+tx-t, G=(1,1) & III* at infinity, II at0, I1 at-27/4 \\
E6-A & y\textsuperscript{2}=x\textsuperscript{3}-3x+t\textsuperscript{2}+2, G=(1,t) & IV* at infinity, I2 at0, I1 at(\pm)2i \\
E6-B & y\textsuperscript{2}=x\textsuperscript{3}+tx+t\textsuperscript{2}, G=(0,t) & IV* at infinity, III at0, I1 at-4/27 \\
\bottomrule
\end{tabularx}
}\par\medskip

\item The first two are verified in P3F.37/P3K.24. For E6-A the discriminant
is 27t\textsuperscript{2}(t\textsuperscript{2}+4). The factorization
(y-t)(y+t)=(x-1)\textsuperscript{2}(x+2) gives a rational chart
u=(y-t)/(x-1), v=(y+t)/(x-1), with inverse
x=uv-2, t=(x-1)(v-u)/2, y=(x-1)(u+v)/2.
For E6-B its discriminant is t\textsuperscript{3}(4+27t). Put
u=(y-t)/x, v=(y+t)/x; the inverse is
x=uv-(v-u)/2, t=x(v-u)/2, y=x(u+v)/2.
Thus these are rational elliptic surfaces, with the stated minimal
coefficient orders. At infinity their cubic good roots send G to a
nonzero fixed three-torsion point. At the I2 in E6-A, the two local
branch orders of G are one and one; at the III in E6-B the fourth
good root sends G to the nonzero fixed two-torsion point. Thus G has
nonidentity components in E6 and A1, height 1/6, and G.O=0.
The root discriminant and height identity show that G is primitive
and the Mordell--Weil group is torsion free.

\item Set M=O(aG-O), Q=bG. Use the following actual data, with a single
cap at infinity, no other logarithmic sites, and one chosen I1 as
the simple whole-fibre zero:

\par\medskip{\small\noindent
\begin{tabularx}{\linewidth}{@{}XXXX@{}}
\toprule
Sources & (a,b) & Root order m & Root type \\
\midrule
E7-A,B & (2,1) & m=4l+3, l(\ge)0 & III \\
E7-A,B & (1,2) & m=4l, l(\ge)1 & smooth \\
E6-A,B & (6,1) & m=3l+2, l(\ge)0 & IV \\
E6-A,B & (2,3) & m=3l, l(\ge)1 & smooth \\
\bottomrule
\end{tabularx}
}\par\medskip

\item For each entry choose an integer n with gcd(m,n)=1. The actual cap
norm will be A\_r\textsuperscript{m=B\_lambda}n, lambda=r\^{}m, with marked meridian
m q+n gamma\_out. In these domains n is nonzero. Every retained
reducible fibre has translation transitive on its components:
at the retained A1 it is odd b, while a is even. The source Hom
line has whole-fibre degree ab/d=1 (d=2 or6). Its degree follows
also directly from the section intersections: for E7 it is1;
for E6 (6,1), (6G).O=(5G).O=2 and G.O=0; for (2,3),
(2G).O=(-G).O=(-3G).O=0. At each component either M has zero
multidegree or Q fixes the component, so the vertical Hom divisor
is a whole fibre. The selected section of that base line supplies
the prescribed one simple I1 zero.

\par

\medskip\noindent\textbf{SC.2. Root equations and exact affine source norms}

\item For E7 reduced roots use (F37.1) and the equation in K24 for E7-B.
For E6-A at infinity,
\[
 y_0^2=x_0^3-3q^4x_0+q^4+2q^6,\qquad G=(q^2,q^2).
\]
For E6-B,
\[
 y_0^2=x_0^3+q^3x_0+q^4,\qquad G=(0,q^2).
\]
For m=3l+2, k=2l+1, substitute q=z\^{}m,
x0=z\^{}\{2k\}X, y0=z\^{}\{3k\}Y. The respective equations are
\[
\begin{array}{ll}
Y^2=X^3-3z^{4l+4}X+z^2+2z^{6l+6},&
G'=(z^{2l+2},z),\\
Y^2=X^3+z^{l+2}X+z^2,&G'=(0,z).
\end{array}
\]
They have the reduced IV minimal regular fibre. In the chart
X=zX1,Y=zY1 its two nonidentity component labels are distinguished
by Y1=(\pm)1. The actual deck
\[
 D(X,Y,z)=(\zeta_m^{-2k}X,\zeta_m^{-3k}Y,\zeta_mz)
\]
has character zeta\_m\textsuperscript{\{-(3k+1)\}=zeta\_m}\{-2m\}=1 on Y1.
It fixes the three component labels. G' is nonidentity and 6G'
has identity component. Resolving the displayed A2 surface point
retains these labels and the source translation.

\item For good roots, take q=z\^{}m with k=3m/4 for E7 and k=2m/3 for E6.
The central smooth curve and its deck have order four or three
respectively. In a coordinate w at any marked fixed section,
the minimal good-root chart has weights (r-1,1), r=4 or3.
The good-root limit of Q is O, since b is the component-group
order (two or three). The limit of aG is a nonzero fixed point
of that order. These are the exact marked limits from the source
coordinates, not unmarked smooth elliptic replacements.

\item There is a unique analytic m-division R of Q' with the selected
central point. At a reduced root multiplication by m is invertible
on its additive identity group and on its component group, since
m is prime to two or three. At a good root use the local inverse
of multiplication by m at O. Deck invariance follows from
uniqueness. The actual source map
\[
                         a_0=t_{-nR}D
\]
satisfies a0\^{}m=t\_\{-nQ'\}. Translation extends to the actual minimal
regular source by the graph/minimality argument used in F37.

\par

\medskip\noindent\textbf{SC.3. The good-root line map as a toric correspondence}

\item Let r=3 or4 be the minimal good root. Near a marked root section
choose its holomorphic coordinate w=0. Centering at the actual
equivariant section preserves the diagonal weights (r-1,1).
The quotient lattice is
\[
 N'=\mathbb Z^2+\mathbb Z(r-1,1)/r.
\]
The minimal quotient resolution has rays
\[
 v_j=((r-j)/r,j/r),\qquad 0\le j\le r.
\]
Consecutive cones have determinant 1/r in the ordinary lattice,
hence are regular in N'. The following actual common refinements
are regular in the original root lattice:
\[
\begin{array}{ll}
r=3:&(1,0),(2,1),(1,1),(1,2),(0,1),\\
r=4:&(1,0),(3,1),(2,1),(1,1),(1,2),(1,3),(0,1).
\end{array}
\]
Every adjacent determinant is one. Let beta be the resulting root
blowup and pi its toric map to the quotient minimal resolution.
The underlying lattice inclusion gives every chart map by its
dual monomials. Thus pi is an actual proper generically finite
map; no finiteness of its exceptional fibres is assumed.

\item The resolved strict section is the divisor of the w-axis ray.
Its integral Cartier support function is
\[
 h_w(v)=\max(v_w-(r-1)v_z,0).
\]
Indeed it is w-(r-1)z on the first minimal cone and zero on the
remaining cones. The character w-(r-1)z belongs to the quotient
dual lattice. Pulling that Cartier function to the common
refinement, while beta pulls the original section with valuation
v\_w, gives the exact exceptional correction
\[
 \beta^*(\hbox{root section})-\pi^*(\hbox{resolved section})
 =\sum_{\text{exceptional }v}\min(v_w,(r-1)v_z)E_v.       \tag{SC.1}
\]
For r=3 the coefficients along the listed exceptional rays are
(2,1,1); for r=4 they are (3,2,1,1,1). Subtract the identical
O-chart contribution from the aG-chart contribution. This gives
an exact isomorphism
\[
 \beta^*N=\pi^*M+D_{aG}-D_O,\qquad N=O(aG'-O),           \tag{SC.2}
\]
including each exceptional coefficient.

\item At r=3 these are the three A2 quotient points of the IV* model.
At r=4 the two marked fixed points have A3 quotients; the other
two points form an order-two orbit and give the remaining A1
arm. No section divisor in (SC.2) passes through that latter
orbit. The minimal resolved quotients are the original IV* and
III* source models: the invariant Weierstrass substitutions give
their birational maps, and their crepant resolution graphs are
the stated affine E6 and E7 fibres, with no vertical (-1)-curve.
The relative minimal-model graph comparison gives the isomorphism.

\item The correction is equivariant for the actual suspension. At a
fixed root point, Q'\_0=O and deck equivariance give
\[
 w\longmapsto w\,u(w,z)+z^{r-1}h(z^r),\qquad z\longmapsto z,
 \qquad u(0,0)=1.
\]
This preserves each ideal (w,z\^{}j) and (z,w\^{}j), 1\(\le\)j\(\le\)r-1.
The normalized blowup of their invariant product gives exactly
the displayed simultaneous fan refinement. Therefore the actual
source translation lifts, preserves every exceptional correction
in (SC.2), and gives an extending line map on N commuting with D.
For m=r l, pull this whole correspondence through the additional
base change and resolve the fibre product. The same invariant
ideals, now (w,z\^{}\{jl\}), retain the lifted source map and the
pullback divisor identity. This supplies the actual line and B0,
not only their generic-fibre Picard classes.

\item The natural deck weights of N at the marked nonzero point and O
are respectively zeta\_r\^{}\{r-1\} and zeta\_r\^{}\{-(r-1)\}, up to a common
character. They are the weights of (SC.1) at the section. The
central Picard class is nonzero of order p=3 for E6 and p=2 for
E7. The exact connecting map
\[
 (\mathbb C/\Lambda)^D\longrightarrow\Lambda/(D-1)\Lambda,
 \qquad[v]\longmapsto[(D-1)v]
\]
identifies it with the primitive p-torsion class of the cap circle
line. A common deck character tensors by a flat line from the
free suspension circle and does not remove this torsion.

\item At a reduced root the source line is already zero on all original
components. Pullback to the resolved fibre product therefore has
degree zero on every contracted curve and descends through the
actual contractions. Its comparison with N has a whole-fibre
divisor, exactly as in F37, using the reduced III or IV
intersection matrix. This is the complementary actual line map.

\par

\medskip\noindent\textbf{SC.4. Commuting line norms, freeness and holomorphic parameter closing}

\item The full lift argument is the one verified in IV.3. Lift t\_\{-R\}
first by F, correcting its meromorphic chord section's vertical
divisor by a whole-fibre power. Average by the finite deck character
to arrange DFD\^{}\{-1\}=chi F, chi\^{}m=1. Since mR=Q', write the actual
descended map as B0=beta(z)F\^{}m. It commutes with D by (SC.2);
therefore beta is a nonzero function of q. In particular F and D
both commute with B0. The map F\^{}nD has norm ratio
\[
 (F^nD)^m/B_0^n=\chi^{nm(m-1)/2}\beta(q)^{-n}.
\]
An invariant analytic m-th root of this unit gives A0 with
A0\textsuperscript{m=B0}n and {[}A0,B0{]}=1. The character factor is retained also
for the even good-root orders.

\item Put A\_r=r\^{}nA0, B\_lambda=r\^{}mB0. Choose u,v with nu+mv=1.
Then
\[
 H=A_r^uB_\lambda^v=rH_0,\qquad H^m=B_\lambda,\quad H^n=A_r.
\]
On a smaller compact source tube H0 has bounded norm ratios, so
H is strictly contracting for small \textbar r\textbar. The cap is exactly
N\^{}times/. Its norm shell proves smoothness and properness,
and its strict contraction proves freeness of the delta circle.
For outward h=H\^{}\{-1\}, the single-A cut gives
\[
                         q\mapsto-nh,\qquad\gamma\mapsto mh,
\]
hence precisely m q+n gamma. All source translations and line
characters are retained in this cap and its annular comparison.

\item The holomorphic ordinary-to-cap translation cochain is solved
once at c0 and then changed by n(c-c0)s\_infty delta/m.
Its positive cut is the actual A\_r. At c-\textgreater c+m its change is
n s\_infty delta, so the inverse closing is multiplication by
(t-tc)\^{}\{-n\}. The positive deck is its inverse. Near the selected
cusp this multiplier and its inverse are B\_lambda iterates
composed with bounded source translations and invertible units.
They extend through its specified Mumford charts. At the other
fillings they are units. This constructs the proper holomorphic
family over the punctured r-disc, with lambda=r\^{}m, including
negative n.~No source root alone is used to choose a parameter
return class.

\par

\medskip\noindent\textbf{SC.5. Integral source and boundary maps}

\item All maps below are executed with full Smith transformations in
\path{code/part3/iii_star_good_complement.py} and
\path{code/part3/iv_star_complement.py}. Their helpers are retained in
integer\_chain.py. The E7 integral basis is O,F,E7,G as in F37.
For E6 use O,F,C0,\ldots,C5,A1,G. The E6 graph has edges
(0,1),(1,2),(2,3),(3,4),(2,5), highest-root marks
(1,2,3,2,1,2), and G meets C0 and the A1 component.
The full Gram determinant is -1. Translation by G cyclically
permutes the three affine E6 arms and swaps the two A1 components.
Its section columns come from the exact height and component
formula. The resulting matrix is checked to preserve the Gram
form, and the actual contracting source map is its negative b-th
power. No separate unmarked lattice is substituted.

\item Let V be the source minus its starred fibre tube. That tube has
k=8 or7 component classes and no odd homology. Its inclusion
matrix L is the affine component matrix, retaining the multiplicity
vector. The exact complement chain model has
\[
 C_4=\mathbb Z,\ C_3=\mathbb Z,\ C_2=\mathbb Z^{10},\
 C_1=\mathbb Z^k,\ C_0=\mathbb Z,\qquad
 d_4=1,\ d_2=L^tB.
\]
The source and relative complexes before taking this cone are
free and even, so their degree-one homotopy group is zero.
The displayed translation and component permutation therefore
give the actual functorial mapping-torus cone. The source boundary
has C3=Z,C2=Z\textsuperscript{k,C1=Z}k,C0=Z,d2=L\^{}tBL; its maps to the complement
are 1,L,1 in the corresponding degrees. These are the chain maps
used in the exported calculation, so all boundary inclusions are
retained.

\item The source complement is simply connected. In E7-A this is the
three-cycle trace argument of F37, and E7-B has the two vanishing
cycles of II meeting once. E6-A has I2+I1+I1. If its primitive
vanishing cycles spanned an index d sublattice, the trace of their
product would be 2 modulo d\textsuperscript{2} (expand the product of rank-one
transvections, including the repeated I2 cycle). The inverse IV*
monodromy has trace-1, so d\textsuperscript{2} divides3, forcing d=1; rank one is
excluded by the same trace. E6-B has a III fibre whose vanishing
cycles already meet once. The source section then kills all base
meridians. Thus the exterior orbit W\_V has pi1=Z gamma.

\item For the reduced partners, H2(W\_V)=Z, H3(W\_V)=ZF and H4=0.
The boundary H1 and H2 are Z\textsuperscript{2}, with F=p z1, p=2 or3.
Its H2 inclusion is (0,-1) in a primitive basis. The second
column is the actual relative-Wang cycle
p(P\_\{-G\}-O) plus its integral removed-root correction, and is
a unit. The original line need not be that primitive generator:
for E6 its free coefficient is -2. The cap has H2=ZC, with
F specializing to pC, so z1 maps to a unit C. The other column
is an arbitrary retained integer kappa C. The full degree-two
map is unimodular for every kappa.

\item For the good partners, use the exported boundary basis
(eta,r,s), with orders (p,0,0), and exterior basis (tau,z), with
orders (p,0). The exact H2 inclusions and fibre classes are:
\[
\begin{array}{c|c|c|c}
 & H_2(E)\longrightarrow H_2(W_V)&F\text{ in }H_2(E)&F\text{ in }H_2(W_V)\\
E7,\ b=2&
\begin{pmatrix}-7&-1&3\\0&0&1\end{pmatrix}&(0,-1,0)^t&(1,0)^t\\
E6,\ b=3&
\begin{pmatrix}-2&2&-2\\0&0&-1\end{pmatrix}&(0,1,0)^t&(2,0)^t .
\end{array}                                                     \tag{SC.3}
\]
The torsion coefficient in each first column is a unit modulo p.
The cap has H2=ZC; F maps with degree one to C by actual parallel
transport through the smooth good-root source. Thus its H2 row
is (0,-1,kappa) for E7 and (0,1,kappa) for E6. The free two-by-two
Mayer--Vietoris block is unimodular, and the torsion block is an
isomorphism Z/p-\textgreater Z/p.~This proves the complete H2 isomorphism
without choosing or deleting kappa.

\item In degree one the good boundary torsion maps primitively to the
cap torsion. The generic elliptic lattice transports identically
to the central good fibre. The two quotient lattices are
Lambda/(D-1)Lambda and Lambda/(D\^{}u-1)Lambda, with u inverse
to n modulo m. Their sublattices agree integrally: for order3
use D+I=-D\textsuperscript{2}, and for order4 an odd power is D or D\^{}\{-1\}.
The multiplying matrices are unimodular. Eliminate these actual
torsion isomorphisms; the remaining free degree-one matrix is
\[
                         D_1=\begin{pmatrix}0&1\\ n&-m\end{pmatrix}.
\]
The reduced partners have that same free matrix without torsion.

\item In degree three the primitive boundary basis is the oriented
source boundary followed by F gamma. The exterior row is (0,1).
The norm-circle projection has degrees -n and m on those two
classes. Their transverse source cycle specializes to the primitive
invariant full central fibre, both for the smooth root and for
the cyclically permuted reduced components. Hence the cap row
is (-n,m), and D3=D1. There is no extra source H3 class. Degree
four is Z-\textgreater0; degree zero is the incidence column (1,-1)\^{}t.
These specify every gluing degree.

\item Van Kampen kills the cap elliptic lattice by its actual incoming
generic elliptic loops, since the exterior source complement is
simply connected. In the reduced cases there is no cap elliptic
fundamental group. The remaining relations identify gamma=m h
and impose n h=0. Thus pi1(Y)=Z/\textbar n\textbar. The displayed integral
Mayer--Vietoris maps give
\[
 H_0(Y)=H_5(Y)=\mathbb Z,\quad
 H_1(Y)=H_3(Y)=\mathbb Z/|n|,\quad H_2(Y)=H_4(Y)=0.       \tag{SC.4}
\]
For n=(\pm)1 the orbit is the standard S5. The sole fixed sphere has
the primitive O(-1)\^{}2 Hopf model, and all other delta actions are
free by the actual norm inequality. The Gysin and fixed-tube maps
of K22 give the standard S6 total space. For nonunit n the
same primitive local Hopf meridian kills the principal-circle
kernel and gives pi1(X)=pi1(Y), so these actual total spaces
have nontrivial fundamental group.

\item This proves the one-cap sphere sectors of the table with their
full local line maps and all integral orbit gluing degrees.
The complete source, line and integral-map proof has passed the
independent audits recorded above. P3K.26 proves the actual return
through its separate signed page and line comparison.

\end{enumerate}
\end{proof}




\begin{theorem}\label{night:P3F.42}
\StatementRecord{night:P3F.42}{proved}{S6\_06\_BLIND: GPT-6 Astra (OpenAI, ChatGPT Pro mode, effort Pro), recorded re-derivation in the checkpoint dated 18 November 2026}{\cref{night:P3K.19}, \cref{night:P3K.22}}
The six source/root sectors in the proof admit the displayed actual line and affine cap maps for every $\gcd(m,n)=1$. Their parameter is $\lambda=r^m$ and their marked cap meridian is $mq+n\gamma_{\rm out}$. For nonzero $n$ their orbit has $H_0=H_5=\mathbb Z$, $H_1=H_3=\mathbb Z/|n|$ and $H_2=H_4=0$, with fundamental group $\mathbb Z/|n|$. They give the standard oriented $S^6$ exactly for $n=\pm1$. The complementary good-root norm for a primitive $Q=G$ is excluded by its exact central elliptic equation.
\end{theorem}
\begin{proof}
\begin{enumerate}[label=\arabic*.]
\item \par

\medskip\noindent\textbf{OA.1. Sources, root domains and the exact excluded norm}

\item Use the marked rational elliptic sources and their actual translations
from P3K.22. The six data are as follows. A chosen I1 is the sole
simple whole-fibre zero of the suspension linearization, and the
displayed additive site is the only logarithmic cap.

\par\medskip{\small\noindent
\begin{tabularx}{\linewidth}{@{}XXXX@{}}
\toprule
Source profile & Capped fibre & (a,b) in M=O(aG-O), Q=bG & Allowed root \\
\midrule
I7+III+2I1 & III & (14,1) & m=4l+1, l\textgreater=0, reduced III \\
I7+III+2I1 & III & (7,2) & m=4l, l\textgreater=1, good \\
I5+I3+III+I1 & III & (30,1) & m=4l+1, l\textgreater=0, reduced III \\
I5+I3+III+I1 & III & (15,2) & m=4l, l\textgreater=1, good \\
I5+IV+I2+I1 & IV & (30,1) & m=3l+1, l\textgreater=0, reduced IV \\
I5+IV+I2+I1 & IV & (10,3) & m=3l, l\textgreater=1, good \\
\bottomrule
\end{tabularx}
}\par\medskip

\item Write p=2 at III and p=3 at IV. At every retained reducible fibre,
Q permutes components transitively and aG meets the identity
component. At the capped fibre, a is divisible by p in the reduced
case; in the good case b=p and a is prime to p.~These assertions
follow in the marked basis from (N,k)=(7,2),(2,1) and
(5,1),(3,1),(2,1), respectively. The exact integral matrices below
also retain these actions. The whole-fibre Hom line has degree
ab*h(G)=1. Thus its chosen scalar section has precisely the stated
simple zero; no component-supported zero is introduced.

\item The norm obstruction for the complementary disallowed good roots
is an equation on the actual central elliptic curve. The canonical
good-root deck has linear part E of order four or three. An affine
source lift has the form t\_R E. Since r divides m and E has no
eigenvalue one,

\[
(t_R E)^m=t_{(1+E+\cdots+E^{m-1})R}=\mathrm{id}.
\]

\item The endomorphism sum is zero: its product with E-1 is E\^{}m-1=0,
and E-1 is an isogeny; the connected image of the sum cannot be
a nontrivial subgroup of that finite kernel. The actual norm
A\textsuperscript{m=B}n therefore requires n Q\_0=0. For Q=G its central limit is
the nonzero fixed point of exact order p.~Since p divides m and
gcd(m,n)=1, this equation is impossible. This proves the excluded
good-root norm for all three primitive cases directly. It does
not exclude the weighted partner Q=pG, whose marked limit is O.

\item For a reduced III or IV root, the deck fixes all component labels,
and m is prime to p.~The affine component norm is
m{[}R{]}=n{[}Q{]} in Z/p.~Its generated component subgroup is exactly
that generated by Q. Thus the weighted partner Q=pG cannot be
transitive on this reduced root. The other residues give the
nonreduced I0\emph{, III}, or IV* fibres; their actual multiplicities
and the intersection obstruction P3K.19 exclude them from this
one-cap sphere sector. No root has been replaced by an unmarked
smooth fibre in these exclusions.

\par

\medskip\noindent\textbf{OA.2. Resolved roots, lines and commuting affine norms}

\item For III use its actual minimal equation

\[
y^2=x^3+q a(q)x+q^2b(q),\qquad a(0)\ne0.
\]

\item For m=4l+1 substitute q=z\^{}m, x=z\^{}(2l)X, y=z\^{}(3l)Y:

\[
Y^2=X^3+z a(z^m)X+z^{2l+2}b(z^m).
\]

\item The A1 resolution gives the reduced III fibre. Its two component
labels are fixed by the deck
D(X,Y,z)=(zeta\_m\textsuperscript{(-2l)X,zeta\_m}(-3l)Y,zeta\_m z).
For m=4l the same substitution gives the smooth family

\[
Y^2=X^3+a(z^m)X+z^{2l}b(z^m).
\]

whose central deck is (X,Y)-\textgreater(-X,iY).

\item For IV the actual equation is

\[
y^2=x^3+q^2a(q)x+q^2b(q),\qquad b(0)\ne0.
\]

\item For m=3l+1 the substitution with k=l gives

\[
Y^2=X^3+z^{2l+2}a(z^m)X+z^2b(z^m).
\]

\item Resolve its A2 point. In X=zX1,Y=zY1 the nonidentity components
are Y1=+/-sqrt(b(0)), and the deck character on Y1 is
zeta\_m\^{}(-(3l+1))=1. Thus all three labels are fixed. For m=3l
the same substitution is good, with equation

\[
Y^2=X^3+z^{2l}a(z^m)X+b(z^m).
\]

and central deck (X,Y)-\textgreater(zeta\_3 X,Y).

\item The marked source G meets the nonidentity component of III or IV.
Its limit on the good root is the corresponding nonzero fixed
point of order p.~This follows by the actual coordinate
w=-X/Y at O: deck-invariant sections with limit O have
w(z)=z\^{}k c(z\^{}r), and substitution in the original x,y coordinates
lands on the original identity component. The converse follows
from the same inverse chart. Thus the component marking used
above is the one carried by the given Weierstrass map.

\item At a reduced root, N=O(aG'-O) has degree zero on every component.
Pull back M to the resolved fibre product and carry it through
the actual relative contractions: each contracted curve has
degree zero. The descended line and N have the same generic
Picard point; their vertical comparison is an integral multiple
of the whole fibre, because the reduced III/IV intersection
matrix has precisely that primitive kernel. The base monomial
with that exponent retains the deck character. This supplies
the extending actual suspension map B0 on N.

\item For good IV roots the full correspondence and its line divisor
are IV.2 in \path{notes/part3_iv_caps_working.md}. Its three fixed points
are blown up, the quotient is taken, and its central (-1)-curve
is contracted. In that actual diagram,

\[
\beta^*N=\pi^*\alpha^*M+E_{aG}-E_O.
\]

\item The following gives the good III correspondence explicitly.
On the minimal fourth root the central elliptic curve has two
order-four fixed points O,P and one orbit \{R1,R2\} with stabilizer
of order two. The local weights at O,P are (1,1) for mu4, and
those at R1,R2 are (1,1) for mu2. Blow up these four points by
beta. In the quotient pi, the images D0,D1 of the first two
exceptional curves have square -4, the image D2 of the latter
pair has square -2, and the central image C has square -1.
Indeed pi\^{}\emph{D0=4E0, pi\^{}}D1=4E1,
pi\^{}\emph{D2=2(E\_R1+E\_R2), and pi\^{}}C is the proper central elliptic
curve of square -4. The quotient fibre is

\[
4C+D_0+D_1+2D_2.
\]

\item Contract C. The three remaining squares become -3,-3,-1,
and their pairwise intersections are one. Contract D2 next.
The remaining two curves have square -2 and intersection two;
the displayed invariant Weierstrass map identifies their
tangency as the original III fibre. These are the actual
relative-minimal contractions. The strict marked sections pass
through D0,D1 away from both contracted curves. Therefore

\[
\beta^*N=\pi^*\alpha^*M+E_P-E_O.\tag{OA.1}
\]

\item This identity is a divisor identity, including its exceptional
coefficients. The root limit of Q=pG is O. Its source translation
is the identity on the central elliptic curve, fixes all four
centres, lifts to their blowup, and preserves the correction in
(OA.1). The quotient and contractions consequently carry the
actual extending B0 and its deck commutation to N. For m=rl,
pull back the whole diagram through the further base change
and resolve that fibre product. This retains the divisor identity
and the maps; it does not discard the ramification characters.

\item The central class of N in a good root is nonzero p-torsion.
The connecting map from fixed elliptic points to
Lambda/(D-1)Lambda sends {[}v{]} to {[}(D-1)v{]}; it identifies this
class with the primitive p-torsion Euler class of the circle
line on the cap. A common deck character contributes the flat
line from the free suspension circle and does not remove it.

\item Now take gcd(m,n)=1. At a reduced root multiplication by m
is invertible on the additive identity group and on Z/p; at
a good root use its inverse branch at O. This gives the unique
marked analytic R with mR=Q', and D-invariance follows from
uniqueness. The source map t\_(-nR)D has mth power t\_(-nQ').
Lift t\_(-R) first to F on N. The exact chord/vertical comparison
and deck eigenspace averaging give DFD\^{}(-1)=chi F, chi\^{}m=1.
The actual descended lift is B0=beta(q)F\^{}m, with beta a base
unit: deck commutation forces beta to be invariant. Thus F,D
commute with B0 and

\[
(F^nD)^m/B_0^n=\chi^{nm(m-1)/2}\beta(q)^{-n}.
\]

\item Divide by an invariant analytic mth root of this unit. The
result A0 satisfies A0\textsuperscript{m=B0}n and {[}A0,B0{]}=1, including the
displayed character factor for the even good-root orders.

\item Put A\_r=r\^{}n A0 and B\_lambda=r\^{}m B0. If nu+mv=1, then
H=A\_r\^{}u B\_lambda\^{}v=r H0 satisfies H\^{}m=B\_lambda and H\^{}n=A\_r.
A smaller compact source tube makes H strictly contracting for
small \textbar r\textbar. Its norm shell gives the actual smooth proper cap
N\^{}times/ and proves that delta is free there. For outward
h=H\^{}(-1), the one-A cut is

\[
q\longmapsto-nh,\qquad\gamma_{\rm out}\longmapsto mh.
\]

so the marked killed meridian is m q+n gamma\_out. The ordinary
to cap cochain has scalar variation n(c-c0)s delta/m. At
c-\textgreater c+m the inverse closing is (t-t\_c)\^{}(-n); the positive
deck is its inverse. As in SC.4, the cusp comparison is by
actual B iterates, bounded source translations and invertible
toric units. Thus these give holomorphic compact families on
the punctured r-disc, not only isolated smooth caps.

\par

\medskip\noindent\textbf{OA.3. Every integral source and boundary map}

\item The full certificates are
\path{code/part3/ordinary_additive_boundary.py/.json/.out}, executed PASS.
They use only the actual NS Gram and negative-translation matrices
of P3K.22. Fourfold independently supplied the primitive and
weighted complement calculations in
\path{code/part3_fourfold/primitive_additive_complements}.* and
weighted\_additive\_complements.*. The code retains all chain maps,
Smith representatives, component columns, and marked fibre traces.

\item Let L be the ten-by-p matrix of all components of the removed
reduced III/IV fibre, and B the unimodular source Gram matrix.
The actual source-complement chain model is

\[
\begin{gathered}C_4=\mathbb Z\longrightarrow C_3=\mathbb Z\longrightarrow C_2=\operatorname{NS}(S)\longrightarrow C_1=\mathbb Z^p\longrightarrow C_0=\mathbb Z,\\ d_4=1,\quad d_3=0,\quad d_2=L^tB,\quad d_1=0.\end{gathered}
\]

\item The boundary model has C3=Z, C2=C1=Z\^{}p, C0=Z and sole
nonzero differential d2=L\^{}t B L. Its inclusion is identity
in degrees zero, one and three, and L in degree two.
The source translation T=t\_(-bG) acts by T on NS and by
the actual permutation P on components, TL=LP. On dual
meridians it acts by P\^{}(-t). The code forms both integral
mapping cones and their inclusion. The even source and relative
complexes have free groups, so these displayed maps compute
the actual relative Wang maps without an omitted torsion choice.

\item The complement V is simply connected. It has a section, so its
fundamental group is a quotient of the abelian elliptic lattice.
The component intersection map L\^{}tB is surjective: replace its
rows by F and the p-1 root-component rows, which are selected
rows of the unimodular B. Therefore H1(V)=0, proving pi1(V)=0.
Equivalently the remaining nodal vanishing cycles span the full
lattice: their product has trace zero at III or -1 at IV;
trace modulo the square of the spanning index gives d\^{}2\textbar2 or
d\^{}2\textbar3, hence d=1. The exact matrices supply the integral proof
without requiring that alternative trace argument.

\item In the primitive case b=1, the exterior W and its boundary E have

\[
\begin{gathered}H_1(W)=H_2(W)=H_3(W)=\mathbb Z,\quad H_4(W)=0,\\ H_1(E)=H_2(E)=H_3(E)=\mathbb Z^2,\quad H_4(E)=\mathbb Z.\end{gathered}
\]

\item In the exported H2 basis z1,z2 the actual fibre is p z1 and
the exterior inclusion is (0,1). The second column is primitive:
its actual relative-Wang cycle is p(T O-O), corrected integrally
in the removed A\_(p-1) root lattice. It has free cokernel value
plus or minus one. The original line M has values 7,15,10 in
the respective cokernels; it has not been substituted for that
primitive cycle.

\item In the good-root weighted cases b=p,

\[
\begin{gathered}H_1(W)=\mathbb Z,\quad H_2(W)=\mathbb Z/p\oplus\mathbb Z,\quad H_3(W)=\mathbb Z,\quad H_4(W)=0,\\ H_1(E)=H_2(E)=\mathbb Z/p\oplus\mathbb Z^2,\quad H_3(E)=\mathbb Z^2,\quad H_4(E)=\mathbb Z.\end{gathered}
\]

\item Use the exported H2 coordinates (eta,z1,z2), with eta of order
p and actual F=z1. The inclusions into (tau,z) of W are exactly

\[
\begin{array}{c|c}\text{OS47, III}&\begin{pmatrix}-1&-7&-1\\0&0&1\end{pmatrix}\\[3pt]\text{OS56, III}&\begin{pmatrix}-1&-15&-1\\0&0&1\end{pmatrix}\\[3pt]\text{OS56, IV}&\begin{pmatrix}-2&-10&-2\\0&0&1\end{pmatrix}\end{array}\tag{OA.2}
\]

\item The torsion row is taken modulo p.~In particular the torsion
boundary maps are units, and F generates exterior torsion. The
line M has nonzero removed-component degrees in these cases;
it is retained as a cohomology functional on V, not as a compact
cycle in its kernel lattice.

\item For both cases the exported H3 basis is precisely the oriented
source boundary L and the actual fibre-gamma trace F gamma.
Its exterior map is (0,1). These marked chain representatives
are stored in the JSON. This also fixes the D3 sign below.

\par

\medskip\noindent\textbf{OA.4. Cap specialization and all Mayer--Vietoris degrees}

\item In the reduced case the cap source is a tree of p rational
components, cyclically permuted by the actual H. Its orbit cap
has H1=Zh, H2=Zc, H3=Z(sum C\_i), H4=0. The actual generic
fibre specializes to sum C\_i, hence its image in the coinvariants
is p c.~As the boundary fibre is p z1 and the target is free,
z1 maps to c.~Write the retained second cap column as kappa c.
The complete degree-two MV matrix, with source rows first, is

\[
D_2=\begin{pmatrix}0&1\\-1&-\kappa\end{pmatrix},\qquad\det D_2=1.\tag{OA.3}
\]

\item No value or vanishing of the collar coefficient kappa is assumed.

\item In the good case the central elliptic monodromy is D\^{}u, where
nu+mv=1. The cap has H1=Zh + Z/p, H2=ZE0, H3=Z(E0 h),
H4=0. Generic elliptic transport in the actual smooth root maps
F to E0 with degree one. The source lattice transport maps the
boundary torsion isomorphically to the cap torsion: for order
four u is odd and D\^{}u-I differs from D-I by a unimodular
factor; for order three D+I=-D\^{}2 supplies that factor. Thus
the torsion and source marking have been compared by actual
root transport. The H2 cap row is (0,1,kappa). Together with
(OA.2), the torsion row is an isomorphism and the remaining
free block is {[}{[}0,1{]},{[}-1,-kappa{]}{]}, hence the full D2 is an
integral isomorphism for every retained kappa.

\item The H1 boundary torsion in the good case also maps isomorphically
to the cap's Z/p and to zero in W. After eliminating precisely
this actual isomorphism, the free degree-one matrix in the
marked basis (q,gamma\_out) is

\[
D_1=\begin{pmatrix}0&1\\n&-m\end{pmatrix}.\tag{OA.4}
\]

\item In degree three, the cap source boundary maps to -n times its
primitive invariant fibre trace and F gamma maps to m times
that trace. This follows by its actual norm projection
u\_cap=m u\_ordinary-n arg(q)/(2pi), followed by the degree-one
specialization of the invariant sum C\_i (reduced case) or E0
(good case). The cap source retracts to a two-dimensional
fibre, so there is no additional H3 source term. Therefore

\[
D_3=\begin{pmatrix}0&1\\n&-m\end{pmatrix}.\tag{OA.5}
\]

in the stored (L,F gamma) boundary basis. Finally D4 is
Z-\textgreater0 and D0 is the incidence column (1,-1)\^{}t. This lists
every MV degree, including the pre-elimination finite groups.

\item Van Kampen kills all elliptic generators in V. In the good case
the cap fundamental group contains the full normal elliptic
lattice, not merely its finite abelianization. The actual
degree-one regular-fibre transport sends the two boundary
elliptic generators onto that full lattice and kills it in
the union. The remaining presentation has q=1 in W,
q=h\^{}(-n), gamma=h\^{}m, so pi1(Y)=Z/\textbar n\textbar{} (Z for n=0).
The integral MV sequence gives H0=H5=Z,
H1=H3=Z/\textbar n\textbar{} and H2=H4=0 for n nonzero. For n=0 (possible
only at the reduced order-one cap) it gives the corresponding
free kernel and cokernel groups directly from (OA.4)--(OA.5).
Thus the sphere members in these six one-cap sectors are exactly
n=+1 or -1. For those values, the simply connected integral
homology sphere Y is the standard S5. Its actual delta line
and the single selected I1 fixed sphere satisfy the Gysin and
Hopf-filling calculation of P3K.22, giving the standard S6
total space. For \textbar n\textbar\textgreater1 the same primitive Hopf meridian kills
only the circle kernel, leaving pi1(X)=Z/\textbar n\textbar{} and excluding a
sphere member.

\item All statements here retain the one-cap domain in OA.1. Multiple
additive caps and their signed Seifert determinant require their
joint boundary presentation. Return indices require the actual
parameter/page line comparison and are not inferred from these
topology matrices.

\end{enumerate}
\end{proof}




\begin{theorem}\label{night:P3F.43}
\StatementRecord{night:P3F.43}{proved}{S6\_06\_BLIND: GPT-6 Astra (OpenAI, ChatGPT Pro mode, effort Pro), recorded re-derivation in the checkpoint dated 18 November 2026}{\cref{night:P3K.19}}
For the type II profiles in rows 45, 47 and 55, with the marked primitive data $M=\mathcal O(dG-O)$, $Q=G$, the root orders $m=6l+1$ and $m=6l$ construct the actual reduced or good cap families below. The good range has $l\ge1$, and the reduced range has $l\ge0$. For $\gcd(m,n)=1$ their cap meridian is $mq+n\gamma_{\rm out}$ and the scalar parameter is $\lambda=r^m$. The standard oriented six-sphere members are exactly $n=\pm1$. The intermediate reduced IV root is excluded by its actual affine component reflection.
\end{theorem}
\begin{proof}
\begin{enumerate}[label=\arabic*.]
\item \par

\medskip\noindent\textbf{II.1. Actual roots and the section marking}

\item Write the actual minimal II equation as

\[
y^2=x^3+q a(q)x+q b(q),\qquad b(0)\ne0.
\]

\item For q=z\textsuperscript{m,x=z}(2l)X,y=z\^{}(3l)Y, the two admissible residues are

\[
\begin{array}{ll}m=6l+1:&Y^2=X^3+z^{2l+1}a(z^m)X+z b(z^m),\\m=6l:&Y^2=X^3+z^{2l}a(z^m)X+b(z^m).\end{array}
\]

\item The first is a reduced II fibre in a smooth total surface. The
second is good, with central curve Y\textsuperscript{2=X}3+b(0) and deck
(X,Y)-\textgreater(zeta\_3\^{}(-1)X,-Y). Its order-six automorphism fixes only O:
a finite fixed point would have both X=0 and Y=0, contrary to
b(0) nonzero. Consequently every original marked section,
including Q=G and dG, has root limit O. In the exact root
coordinate w=-X/Y, its expansion is w=z\^{}l c(z\^{}m), since the
minimal sixth root has weight one. Thus the central good-root
norm requires nQ0=0 and is satisfied here. It is not the
nonzero-torsion norm occurring at primitive III/IV sites.

\item For the intermediate residue m=6l+2, k=l gives a reduced IV
root. In X=zX1,Y=zY1 the deck character on Y1 is
zeta\_m\^{}(-(3l+1))=-1. It fixes the identity component and
interchanges the other two. Q remains on the identity
component: composing this root with its cubic good root gives
the same order-six fixed point O. The affine component norm
therefore cannot generate all three labels. The actual
reflection has two component orbits, and the compact
intersection pairing of P3K.19 excludes a sphere cap.
The residues three, four and five give the nonreduced I0\emph{,
IV} and II* fibres, excluded by their unequal multiplicities
and that same proved intersection obstruction.

\par

\medskip\noindent\textbf{II.2. The good-root line correspondence}

\item On the minimal sixth root, the central curve has one order-six
fixed point O, one orbit of two points with stabilizer order
three, and one orbit of three points with stabilizer order two.
At each point its stabilizer acts on the two surface coordinates
with weights (1,1). Blow up all six points. The quotient by
mu6 has central curve C of square -1 and three exceptional
curves D6,D3,D2 of squares -6,-3,-2. In the actual quotient,

\[
\begin{gathered}\pi^*C=\widehat E,\qquad\pi^*D_6=6E_O,\\ \pi^*D_3=3(E_{P_1}+E_{P_2}),\qquad\pi^*D_2=2(E_{R_1}+E_{R_2}+E_{R_3}).\end{gathered}
\]

\item Here E\_hat has square -6 and maps with degree six. The central
fibre is 6C+D6+2D3+3D2. Contract C, then D2, then D3. Their
successive self-intersections are -1 at the moment of contraction.
After the first step the squares are -5,-2,-1 with pairwise
intersection one; after the second the remaining squares are
-4,-1 with intersection two; after the third D6 has square zero
and the cuspidal singularity specified by the invariant equation
y\textsuperscript{2=x}3+q a(q)x+q b(q). Thus the displayed invariant map and
these contractions identify this quotient with the original
minimal II source, including its smooth total surface.

\item Both marked sections dG' and O pass through the root point O.
In the blowup chart w=z u they meet E\_O at the finite marked
coordinates u=c\_dG(q) and u=0. The proper central elliptic
curve meets E\_O at u=infinity, so both sections avoid every
contracted curve. Each section has exceptional coefficient one
at this single blowup. Subtracting their divisors therefore
cancels that coefficient exactly and gives

\[
\beta^*\mathcal O(dG'-O)=\pi^*\alpha^*\mathcal O(dG-O).\tag{II.1}
\]

\item This identity also holds if the two sections have further
contact: both sides are differences of their actual Cartier
divisors in the displayed (u,q) chart. No other exceptional
point lies on either section. For m=6l pull back the complete
diagram and resolve its fibre product. The actual root source
translation Q' has limit O, fixes the six centres, and lifts
through the whole diagram. Thus (II.1) carries its actual
suspension line map B0 and its deck commutation to the line
N=O(dG'-O). On the good central fibre N is the trivial Picard
class, as required by its actual marked section limit.

\item At a reduced II root the pullback/relative-minimal line comparison
has only one irreducible component. Its vertical difference is
therefore a whole-fibre divisor. The corresponding base monomial
retains the complete deck character and supplies the extending
B0. This is the same actual divisor comparison as at a reduced
III/IV root, with a one-component intersection kernel.

\par

\medskip\noindent\textbf{II.3. Source norm, actual cap and the family closing}

\item Choose gcd(m,n)=1. Multiplication by m has a unique inverse
branch at the marked central point: it acts on the additive
group at a II root, and on the elliptic group at O at a good
root. This gives a D-invariant analytic R with mR=Q'.
The source map a0=t\_(-nR)D satisfies a0\^{}m=t\_(-nQ').
Lift t\_(-R) to F on N using the actual chord divisor and its
whole-fibre correction. Normalize DFD\^{}(-1)=chi F, chi\^{}m=1.
Write B0=beta(q)F\^{}m. Its actual deck descent makes beta
invariant. Then

\[
(F^nD)^m/B_0^n=\chi^{nm(m-1)/2}\beta(q)^{-n}.
\]

\item Dividing by an invariant analytic mth root gives the commuting
line norm A0\textsuperscript{m=B0}n.~This retains the possible even-order
character factor. Set A\_r=r\^{}n A0,B\_lambda=r\^{}m B0. Bezout
gives H=A\_r\^{}u B\_lambda\^{}v=rH0 with H\textsuperscript{m=B\_lambda,H}n=A\_r.
Strict norm contraction on a smaller compact tube constructs
the actual smooth proper cap N\^{}times/ and proves delta
free. Its outward meridian maps are q-\textgreater-nh,gamma-\textgreater mh.
The ordinary-to-cap cochain and closing are precisely the
ones in SC.4/OA.2: scalar variation n(c-c0)s delta/m,
inverse closing (t-t\_c)\^{}(-n), lambda=r\^{}m. They extend by
the actual scalar-zero cusp charts and by units elsewhere.

\par

\medskip\noindent\textbf{II.4. Every integral map for one II cap}

\item The full source/boundary cones are executed in
\path{code/part3/ii_cap_boundary.py/.json/.out}. Here the removed
component column is L=F and its intersection matrix is zero.
The source-complement complex is C4=Z -\textgreater{} C3=Z -\textgreater{} C2=NS
-\textgreater{} C1=Z -\textgreater{} C0=Z with d4=1 and d2=F\^{}t B. The boundary
complex has one Z in each degree zero through three and
zero differentials. Its actual translation and inclusion
are obtained by the same uncancelled cones as OA.3.

\item The source complement is simply connected. Its fundamental
group is a quotient of the elliptic lattice because the section
kills the base loops; F\^{}tB is surjective, so its H1 vanishes.
Its translation coinvariants are free Z, generated by the
actual relative-Wang cycle T O-O, and its invariant lattice
is ZF. One can also read this directly from the exact sequence
0 -\textgreater{} F-perp -\textgreater{} NS -\textgreater{} Z -\textgreater{} 0: in the K22 full Wang calculation
the invariant F maps to zero and the cokernel generator O
maps to one, forcing the relative-Wang column to be primitive.

\item Thus H1(W)=H2(W)=H3(W)=Z,H4(W)=0. The boundary E has
H1=H2=H3=Z\^{}2,H4=Z. In the stored marked H2 basis the actual
fibre is primitive and its exterior image is zero; the other
boundary column is a unit. In both the reduced II and good
cap, H1=Zh,H2=ZC,H3=Z(C h),H4=0. For good reduction
the determinant of D-I on the elliptic lattice is one, so
there is no hidden finite cap group. The actual generic
fibre transport/specialization has degree one in either cap.

\item Consequently, retaining the integer collar coefficient kappa,
the full middle map is D2={[}{[}0,1{]},{[}-1,-kappa{]}{]} after the
explicit oriented boundary choice. The other degree maps in
the actual bases (q,gamma) and (source boundary,F gamma) are
D1=D3={[}{[}0,1{]},{[}n,-m{]}{]}, while D4 is Z-\textgreater0 and D0=(1,-1)\^{}t.
The cap rows follow from u\_cap=m u\_ordinary-n arg(q)/(2pi);
there is no H3 source term because each cap source retracts
to its two-dimensional fibre. In the good case the actual
degree-one regular-fibre transport carries the two boundary
elliptic generators onto the cap's full normal elliptic Z\^{}2,
which is killed in the simply connected source complement.
The reduced II cap source is already simply connected. The
remaining marked relations give pi1(Y)=Z/\textbar n\textbar{} by Van Kampen.
For n nonzero the integral homology is H0=H5=Z,
H1=H3=Z/\textbar n\textbar,H2=H4=0. Hence exactly n=+/-1 are spheres
in these one-cap domains. Their actual Hopf filling is the
standard S6 by the proved K22 Gysin calculation. No return
class has been assigned from this topology computation.

\end{enumerate}
\end{proof}




\begin{theorem}\label{night:P3F.44}
\StatementRecord{night:P3F.44}{proved}{S6\_06\_BLIND: GPT-6 Astra (OpenAI, ChatGPT Pro mode, effort Pro), recorded re-derivation in the checkpoint dated 18 November 2026}{Definitions and the displayed constructions}
Every actual affine line lift on one of the constructed additive root sources is carried by the explicit source translation and line maps below to $\zeta_m^d A_0$, with $d\bmod m$ retained. The comparison preserves the actual suspension $B_0$ and transports the outer annular map. It does not discard the transported source translation or integral line gauge.
\end{theorem}
\begin{proof}
\begin{enumerate}[label=\arabic*.]
\item \par

\medskip\noindent\textbf{AE.1. The section group and its exact cyclic action}

\item Let \(D\) be the canonical order-\(m\) root deck on the
constructed minimal regular source \(S'\to\Delta_z\). Its
action on an analytic section is
\[
 (D\cdot T)(z)=D\bigl(T(\zeta_m^{-1}z)\bigr).
\]
All sums below are sums in this actual section group. The marked
section \(Q'\) is \(D\)-invariant. The constructed branch
\(R\), with \(mR=Q'\), is also \(D\)-invariant. At a
reduced II, III or IV root it is the unique division branch
because the central identity group is additive and multiplication
by \(m\) is invertible on the component group. At a good root
it is the branch with central value \(O\).

\item Every affine source map with the specified linear deck is
\(a=t_S D\) for an actual analytic section \(S\). Its exact
norm equation is
\[
 a^m=t_{-nQ'},\qquad
 \sum_{j=0}^{m-1}D^jS=-nQ'.                         \tag{AE.1}
\]
Set \(Z=S+nR\). Then \(\sum D^jZ=0\). Conjugation of the
constructed source map \(a_0=t_{-nR}D\) by translation is
\[
 t_T a_0t_T^{-1}=t_{T-DT-nR}D.                    \tag{AE.2}
\]
Thus the precise required equation is \((1-D)T=Z\).

\par

\medskip\noindent\textbf{AE.2. An explicit solution on every constructed root}

\item At a reduced root multiplication by \(m\) is a bijection on
germs of analytic sections. On the additive central identity
group its inverse is multiplication by \(1/m\), and on the
finite component group it is a bijection because \(m\) is
prime to its order. Each central inverse extends uniquely by
the analytic inverse function theorem. An analytic section in
its kernel has central value zero, so that same uniqueness
shows that it is the zero section. Define
\[
 T=-[m]^{-1}\left(\sum_{j=0}^{m-1}jD^jZ\right).
                                                               \tag{AE.3}
\]
Since \(D^m=1\), direct subtraction gives
\[
 (1-D)\sum_{j=0}^{m-1}jD^jZ
   =\sum_{j=0}^{m-1}D^jZ-mZ=-mZ.
\]
Equation (AE.3) therefore solves (AE.2) exactly. The translation
extends over the resolved minimal regular source by its actual
source group law and graph comparison, already used to extend
the norm map in SC/OA/II.

\item At a good root the central linear automorphism \(D_0\) has
order three, four or six and has no eigenvalue one. Hence
\(1-D_0\) is a surjective elliptic isogeny. Choose a central
point \(T_0\) with
\((1-D_0)T_0=Z_0\), and lift it to an analytic section
\(\widetilde T\) of the smooth elliptic family. Put
\[
 Z_1=Z-(1-D)\widetilde T.
\]
It has central value zero and still has zero cyclic norm.
The subgroup of analytic sections with central value zero has
unique \(m\)-division with central value zero, by the analytic
inverse function theorem at \(O\). Apply (AE.3) to \(Z_1\)
in this subgroup and call the result \(T_1\). Then
\(T=\widetilde T+T_1\) solves \((1-D)T=Z\). Substitution
in (AE.2) is the required actual source isomorphism. No choice
of a rational homology marking was used to obtain it.

\item This also verifies the good-root norm restriction without losing
the central point. On the good central fibre
\(1+D_0+\cdots+D_0^{m-1}=0\), so (AE.1) forces
\(nQ'_0=0\). For III/IV, \(Q'_0\) has order two/three
when its component is nonzero; coprimality with \(m\) excludes
that case. For II its actual order-six fixed-point calculation
gives \(Q'_0=O\), and the constructed norm is admissible.

\par

\medskip\noindent\textbf{AE.3. The actual translation lift on the specified line}

\item Let \(N\) be the resolved root line already constructed. On
every reduced central component it has degree zero, and on a
good central fibre it has degree zero. Translation by \(T\)
preserves these degree vectors. On a smooth generic fibre,
\(t_T^*N\otimes N^{-1}\) is trivial: in a degree-zero
uniformization its factors of automorphy are constant in the
elliptic variable, so a lift of the translation leaves them
unchanged.

\item There is an extending holomorphic line isomorphism on the whole
source, not only a generic Picard-class equality. Apply proper
direct-image coherence to
\(L=t_T^*N\otimes N^{-1}\). Its direct image is torsion-free
of generic rank one on the disc, hence free after shrinking.
A nonzero generator evaluates to a section of \(L\) whose
divisor is vertical: on each generic compact elliptic fibre it
is a nonzero section of a trivial degree-zero line. Its degree
on each central component is zero. The exact component
intersection matrix consequently puts that divisor in its
primitive whole-fibre kernel. At these reduced or good roots
the fibre has multiplicity one, so dividing by its integral
base monomial removes the divisor. The resulting section is
nowhere zero and gives the required line map \(\widetilde t_T\).

\item This lift commutes with the actual suspension \(B_0\). On a
generic elliptic fibre write the degree-zero line as the quotient
of \(\mathbb C\times\mathbb C^*\) by constant factors of
automorphy. Source translations lift as
\((z,v)\mapsto(z+t,cv)\); these maps commute. The commutator
of the extending maps is a holomorphic unit on the root source.
It is one on every generic fibre and hence one everywhere.
Thus conjugating by \(\widetilde t_T\) transports the actual
line root while preserving \(B_0\), not merely its source
translation. The base monomial used above retains its deck
character in this transported map.

\par

\medskip\noindent\textbf{AE.4. The remaining finite delta character and the boundary map}

\item After this source-and-line conjugation, compare the given line
root \(A\) with the constructed \(A_0\), which have the same
source map. Define the ratio by \(AA_0^{-1}=M_c\), where
\(M_c\) multiplies the line at base coordinate \(z\) by
\(c(z)\). This is a nowhere-zero holomorphic function
on \(S'\), hence a base unit \(c(z)\), because the source
fibres are compact and connected. The equality of their exact
norms is
\[
 \prod_{j=0}^{m-1}c(\zeta_m^jz)=1.                \tag{AE.4}
\]
Its central value is \(c(0)=\zeta_m^d\) for a retained
integer \(d\bmod m\). Divide this constant and take the
analytic logarithm \(h(z)\) with \(h(0)=0\). Equation
(AE.4) implies \(\sum D^jh=0\); the possible integral
logarithm constant is zero at the centre. The additive formula
\[
 b(z)=-\frac1m\sum_{j=0}^{m-1}jD^jh(z)
\]
gives \((1-D)b=h\), with \((Db)(z)=b(\zeta_m^{-1}z)\).
The complete conjugation identity, with its arrow fixed, is
\[
 A=M_{e^b}\,(\zeta_m^dA_0)\,M_{e^{-b}},
\]
because its scalar multiplier is
\(\zeta_m^d\exp(b(z)-b(\zeta_m^{-1}z))=c(z)\).
It leaves \(B_0\) unchanged
because that map fixes the base.

\item Consequently every actual affine line lift on each constructed
root is carried to \(\zeta_m^dA_0\) by the displayed maps,
with its finite delta character retained. Their restrictions to
the outer annulus are the boundary comparison maps. They may
contain a source translation and an integral line gauge; these
are transported, not set to zero. The one-cap topology proofs
allow the resulting integer collar coefficient in full, and
their degree-two matrix remains unimodular for every value.
Global parameter-return comparisons continue to use the actual
closing map and these transported characters, as in K24.

\end{enumerate}
\end{proof}




\begin{theorem}\label{night:P3F.45}
\StatementRecord{night:P3F.45}{proved}{S6\_06\_BLIND: GPT-6 Astra (OpenAI, ChatGPT Pro mode, effort Pro), recorded re-derivation in the checkpoint dated 18 November 2026}{\cref{night:P3F.35}, \cref{night:P3F.40}, \cref{night:P3F.42}, \cref{night:P3F.7}, \cref{night:P3F.9}, \cref{night:P3K.22}}
Use any of the ten actual primitive-translation sources of P3K.22,
with its marked line, translation $G$, and selected simple whole
$I_1$ zero. Install finitely many actual root-line caps away from
that zero. Retain all elliptic, line and annular coefficients.
After the exact finite-kernel comparison of P3K.19, write their
effective norm pairs as $(m_i,n_i)$, with $m_i>0$ and
$\gcd(m_i,n_i)=1$.
At a multiplicative site require the actual component-transitive
root; at an additive site use the complete local domains of
P3F.42--43. Thus the additive possibilities for these primitive
sources are reduced III of order $4l+1$, reduced IV of order
$3l+1$, and type II with reduced order $6l+1$ or good order
$6l$, the latter with $l\geq1$. Every source has at most one
such additive site.

For the circle orbit $Y$, the complete degree-two gluing map has
cokernel
\[
 \operatorname{coker}D_2
   =\bigoplus_{i\text{ smooth or multiplicative}}
                         (\mathbb Z/m_i)^2.             \tag{F45.1}
\]
Its kernel is zero. This cokernel injects into $H_2(Y;\mathbb Z)$.
Consequently every smooth or multiplicative effective order of a
sphere member is one. Among the actual constructions in this
statement the sphere criterion is exactly this order restriction
together with
\[
 \begin{cases}
 n_*+m_*\displaystyle\sum_{i\ne *}n_i=\pm1,
                         &\text{one additive site }*,\\[2mm]
 \displaystyle\sum_i n_i=\pm1,
                         &\text{no additive site}.
 \end{cases}                                           \tag{F45.2}
\]
In that sector the full fundamental group is cyclic of the
displayed order, with $\mathbb Z$ for order zero. The criterion
holds for every permitted one-vertex cusp fan. It retains the
actual order-one source and line clutches; no equality of their
parameter returns is asserted by this topology calculation.
\end{theorem}
\begin{proof}
\begin{enumerate}[label=\arabic*.]
\item \par\medskip\noindent\textbf{The ordinary comparison and integral exterior bases.}\enspace
The relative cusp product of P3F.9 identifies the orbit before
these replacements with the actual mapping torus
$W_0=T_{t_{-G}}S$, with $(p,1)\sim(p-G,0)$ and the physical
outward gamma marking. The integral calculation in P3K.22 gives
$H_2(W_0)=\mathbb Z[O]$ and $H_3(W_0)=\mathbb ZF$.
Let $q\geq1$ be the number of replaced tubes, $W$ their exterior,
$U_i$ the ordinary tubes, and $E_i=\partial U_i$.
All their central source components have multiplicity one, and
$G$ acts transitively on them. The ordinary tube's degree-three
component trace maps to the primitive class $F$ of $W_0$.
At a smooth fibre this trace is the ordinary three-torus.
The sum of these maps is therefore onto $H_3(W_0)$.
The actual section $O$ has boundary $\sum_i\alpha_i$ in
$\bigoplus H_1(E_i)$, where each $\alpha_i$ is its primitive
positive base meridian. This is a nonzero primitive vector.
The full ordinary Mayer--Vietoris sequence now gives the integral
isomorphism
\[
 \bigoplus_iH_2(E_i)
       \ \xrightarrow{\ (i_2,-j_2^{\rm ord})\ }\
 H_2(W)\oplus\bigoplus_iH_2(U_i).                    \tag{F45.3}
\]
In particular $H_2(W)$ is free, and the images of the ordinary
cap kernels form an integral basis of it.

\item For clarity, all the other exterior groups used below have the
following actual source-pair description. Let $L_i$ be the columns
of the source component classes of the $i$th fibre, and let $B$
be the integral source Gram matrix already computed in P3K.22.
Writing $L=(L_1,\ldots,L_q)$, the relative source complex has
\[
 C_4=\mathbb Z,\quad C_3=\mathbb Z^q,\quad
 C_2=\operatorname{NS}(S),\quad C_1=\mathbb Z^{\sum r_i},
 \quad C_0=\mathbb Z,
 \qquad d_4=\mathbf1,\quad d_2=L^tB.                 \tag{F45.4}
\]
The other differentials are zero. The boundary inclusions are
$L_i$ in degree two, their component-coordinate embeddings in
degree one, and the $i$th coordinate in degree three. These are
the actual pair maps: the source and the tube retractions have
free even homology, so their split even complexes have no odd
derived homotopy parameter. This argument includes the tangencies
of III and IV. The primitive trivial lattice of the verified
torsion-free Mordell--Weil source makes the cokernel of $d_2$
free of rank $q-1$. Its generators are the base meridians with
their single sum relation. Translation acts identically on them
and on $H_3$ of the source exterior. The mapping-torus cone of
(F45.4) consequently gives
\[
 H_1(W)=\mathbb Z^q,\qquad H_4(W)=\mathbb Z^{q-1},
 \qquad H_3(W)\text{ is free}.                       \tag{F45.5}
\]
The stated primitivity can also be read directly from the
integral bases in P3K.22: adjoining $O$ and $G$ to the fibre and
the nonidentity components gives the full source lattice, with
the actual Mordell--Weil quotient $\mathbb Z$.

\par\medskip\noindent\textbf{The actual inversion and every degree-two column.}\enspace
On $W_0$ the map $I[p,t]=[-p,-t]$ is an involution. It preserves
the base and every chosen tube, reverses the five-dimensional
orientation, and preserves the normal base-disc orientation.
On a smooth tube choose the actual Lie lift $G(q)$ and the map
\[
 (a,b,t)\longmapsto[a+b\tau(q)+tG(q),t].
\]
In these coordinates $I$ is exactly $(-a,-b,-t)$ and fixes the
base meridian $\alpha$. Thus $ab,a\gamma,b\gamma$ are even,
whereas $a\alpha,b\alpha,\alpha\gamma$ are odd.
The same formula on a Tate sector descends through
$a\mapsto a+Nb+c_Gt$. In the exact basis $(z,R,B,L)$ of
P3F.40, $z,B$ are even and $R,L$ are odd. The ordinary cap kernel
is precisely $\mathbb ZR\oplus\mathbb ZL$.

\item At a reduced additive ordinary site, $H_2(E_i)$ is free of
rank two. Write $F=r z$, with $r=1,2,3$ for II, III, IV.
The primitive class $z$ maps to the primitive component-orbit
class $C$ in the ordinary cap. The boundary is a three-torus
mapping torus and is parallelizable: a path from its oriented
linear monodromy to the identity trivializes its vertical
tangent bundle. Its rank-two unimodular intersection form is
therefore even. Choose $w$ with $z\cdot w=1$ and put
$s=w-(w^2/2)z$. The boundary kernel is isotropic, hence equals
$\mathbb Zs$. Inversion fixes $z$, preserves this kernel, and
reverses the intersection form. The equality
$Iz\cdot Is=-z\cdot s$ gives $Is=-s$. No $z$ correction is
possible in the preserved cap kernel.

\item Every ordinary kernel vector in (F45.3) is thus odd. Hence
$I=-1$ on the entire free group $H_2(W)$. Every even boundary
core column has zero exterior image, because twice that image
is zero. This proves their simultaneous integral vanishing.
The map $I$ is used on the original exterior throughout.

\par\medskip\noindent\textbf{The full local replacement maps.}\enspace
At a smooth site retain
\[
 \mu=m\alpha+n\gamma+\lambda_a a+\lambda_b b,
 \qquad mu+nv=1.
\]
In cap basis $(a,b,h)$ the actual first-homology quotient is
\[
 Q=\begin{pmatrix}
 1&0&-u\lambda_a&-v\lambda_a\\
 0&1&-u\lambda_b&-v\lambda_b\\
 0&0&-n&m
 \end{pmatrix}                                      \tag{F45.6}
\]
on $(a,b,\alpha,\gamma)$. Its kernel is exactly
$\mathbb Z\mu$. The higher maps are its exterior powers,
because the actual smooth root and extending translation give
the marked product $T^3\times D^2$ as in P3K.22.
In degree-two order
$(ab,a\gamma,b\gamma,a\alpha,b\alpha,\alpha\gamma)$,
the complete local block is
\[
 D_2^{\rm sm}=
 \begin{pmatrix}
 0&0&0&1&0&0\\0&0&0&0&1&0\\0&0&0&0&0&1\\
 -1&v\lambda_b&-v\lambda_a&u\lambda_b&-u\lambda_a&0\\
 0&-m&0&n&0&\lambda_a\\
 0&0&-m&0&n&\lambda_b
 \end{pmatrix}.                                    \tag{F45.7}
\]
Eliminating the first three exterior rows and the unit core
column leaves the two diagonal entries $m,m$.
In degree-three order
$(N=ab\alpha,F\gamma=ab\gamma,a\alpha\gamma,b\alpha\gamma)$,
the cap row is
\[
                    (-n,m,-\lambda_b,\lambda_a).     \tag{F45.8}
\]

\item At a multiplicative site use exactly the marked data of P3F.40,
including $c=-c_G$, $A=Nj-cn$, and $\gcd(m,A)=1$.
Choose $p,r$ with $mp+nr=-j$. In degree one, on
$(\alpha,b,\gamma)$, the actual cap map is
\[
       \begin{pmatrix}p&1&r\\-n&0&m\end{pmatrix}.
\]
The complete degree-two block on $(z,R,B,L)$ is
\[
 D_2^{\rm mult}=
 \begin{pmatrix}
 0&1&0&0\\0&0&0&1\\-m&A&0&0\\0&0&-m&-A
 \end{pmatrix}.                                    \tag{F45.9}
\]
This retains the actual integral cap basis change proved by
boundary isotropy in P3F.40. Its cokernel is $(\mathbb Z/m)^2$.
For degree three use $(N=ab\alpha,F\gamma=ab\gamma,
V=a\gamma\alpha)$. The intersection pairing with
$(\alpha,b,\gamma)$ is
\[
 \begin{pmatrix}0&-1&0\\0&0&-1\\1&0&0\end{pmatrix}.
\]
The primitive meridian bounds in the cap. Poincare--Lefschetz
duality identifies the cap degree-three row with its negative
dual, with the sign fixed by the actual fibre trace
$F\gamma\mapsto m$. Thus this row is exactly
\[
                             (-n,m,j).                \tag{F45.10}
\]
Primitivity follows also from the cap's component-trace
generator; there is no further index in this duality map.

\item For the actual additive roots of P3F.42--43, the cap has
$H_1=H_2=H_3=\mathbb Z$. Reduced III/IV specialize the fibre
to $rC$ and have boundary $F=r z$, hence $z\mapsto C$.
Reduced II has the same assertion with $r=1$.
For good II the actual smooth-root transport has degree one,
and $D-1$ is unimodular on the elliptic lattice; again
$z\mapsto C$. With its remaining actual collar coefficient
$\kappa$, the block on $(z,s)$ is
\[
 D_2^{\rm add}=\begin{pmatrix}0&1\\-1&-\kappa\end{pmatrix},
 \qquad j_1=j_3=(-n,m)                             \tag{F45.11}
\]
in the corresponding $(\alpha,\gamma)$ and $(N,F\gamma)$
bases. These specialization statements use the actual resolved
root and line maps, not the numerical Kodaira symbols alone.
The simultaneous integral exterior basis makes (F45.7),
(F45.9), and (F45.11) independent diagonal blocks of $D_2$.
They are injective and prove (F45.1).

\par\medskip\noindent\textbf{Every remaining gluing degree.}\enspace
Put $c_i=3,2,1$ at smooth, multiplicative, additive sites.
In $H_1(W)$ use the $q-1$ independent base meridians and the
common gamma generator. Elliptic boundary loops map to zero;
$\alpha_q=-\sum_{i<q}\alpha_i$, and all $\gamma_i$ have the
same image. These statements follow from (F45.4) and from the
actual source-slice paths, so they retain the original markings.
Together with (F45.6), (F45.10)'s degree-one precursor, and
(F45.11), they specify every entry of $D_1$.

\item The following integral description likewise specifies $D_3$.
For an ordinary cap its kernel in degree three is generated by
$N_i$ and $c_i-1$ local elliptic-rim classes. All these classes
are inversion-even, by the coordinates just used. The source
fundamental class of $W_0$ has boundary $\sum_iN_i$.
The ordinary degree-three exact sequence, whose cap traces
surject onto the primitive $F$, shows that the images of these
kernel classes have this one primitive relation and no other
relation. Adjoin one image of $F\gamma$ to obtain an integral
basis of $H_3(W)$. Every other $F\gamma$ image equals it:
their differences lie in the even kernel and are also odd.
The freeness in (F45.5) makes those differences zero.
Consequently the exterior columns are $N_q=-\sum_{i<q}N_i$,
one common $F\gamma$, and the independent local rim columns.
The complete cap rows are (F45.8), (F45.10), and (F45.11).

\item Eliminating the local elliptic columns in $D_1$ and the local
rim rows in $D_3$ gives identity blocks and the same $2q$ by
$2q$ matrix $S_q$. Its columns are
$(\alpha_1,\gamma_1,\ldots,\alpha_q,\gamma_q)$, or
$(N_1,F\gamma_1,\ldots,N_q,F\gamma_q)$, respectively.
Its first $q-1$ rows have $1$ in column $\alpha_i$ and $-1$
in column $\alpha_q$. Its next row has $1$ in every gamma
column. Its last $q$ rows have $(n_i,-m_i)$ in the $i$th
pair and zero elsewhere. For example,
\[
 S_2=\begin{pmatrix}
 1&0&-1&0\\0&1&0&1\\n_1&-m_1&0&0\\0&0&n_2&-m_2
 \end{pmatrix}.                                    \tag{F45.12}
\]
Expanding after eliminating the $q-1$ meridian rows gives
\[
 \det S_q=(-1)^{q(q+1)/2}
       \sum_i n_i\prod_{j\ne i}m_j.                 \tag{F45.13}
\]
Equivalently, over the rational function field eliminate each
cap row to set $\gamma_i=(n_i/m_i)\alpha_i$; the remaining
meridian relation has coefficient $\sum_i n_i/m_i$.
Multiplication by $\prod_i m_i$ and the displayed column
permutation gives (F45.13) as a polynomial identity over
$\mathbb Z$, with no restriction on the signed $n_i$.
Finally
\[
 D_4=(I_{q-1}\mid-\mathbf1),\qquad
 D_0=\begin{pmatrix}\mathbf1^t\\-I_q\end{pmatrix}. \tag{F45.14}
\]
Thus every nonzero degree of the full Mayer--Vietoris map is
specified as an integer matrix, including every local collar.
The program \path{code/part3_root/primitive_multicap_complete.py}
exports them, checks the exterior-power and duality identities,
and executes 121 full integral cases with positive, negative and
zero determinant. The symbolic determinants and the integral
eliminations prove the assertions for all orders.

\item When $D=\sum_i n_i\prod_{j\ne i}m_j\ne0$, the complete
sequence in particular gives
\[
 \begin{gathered}
 H_0(Y)=H_5(Y)=\mathbb Z,\qquad H_4(Y)=0,\\
 H_1(Y)=H_3(Y)=\operatorname{coker}S_q,\qquad
 H_2(Y)=\bigoplus_{i\text{ smooth or multiplicative}}
                                  (\mathbb Z/m_i)^2.
 \end{gathered}                                    \tag{F45.15}
\]
For a sphere all the latter orders are therefore one. Eliminating
their cap generators in $S_q$ leaves $D=n_*+m_*\sum_{i\ne *}n_i$,
or $D=\sum_i n_i$ with no additive site. Its Smith factors are
units followed by $|D|$; for $D=0$ the final factor is zero.

\par\medskip\noindent\textbf{The full fundamental group and the sphere implication.}\enspace
On the complement of the source critical values the section
gives the actual group extension by the elliptic lattice.
Its normal elliptic subgroup is abelian. Every multiplicative
replacement kills its original primitive vanishing circle:
transport to the actual root central polygon collapses precisely
that circle. This statement includes its graph translation.
A reduced additive replacement kills the elliptic subgroup by
its simply connected reduced central fibre. At a good II cap,
the boundary elliptic lattice maps onto the full elliptic
normal subgroup of the cap. The other source fibres already
kill that lattice. Indeed the complement of the sole II fibre
has base a disc; its fundamental group is the quotient of the
elliptic lattice by the retained vanishing cycles. Its integral
first homology is zero from (F45.4) with $L=F$, so that abelian
group is zero. With no additive replacement, the original
vanishing-cycle relations of the simply connected rational
source are all retained in the unchanged fibres or in the
multiplicative caps. This again kills the entire lattice.
These are relations in the actual group, and every source
translation acts trivially on that elliptic normal subgroup.

\item After eliminating it, van Kampen therefore gives precisely
\[
 \left\langle\gamma,\alpha_i,h_i\ \left|\
 \prod_i\alpha_i=1,\quad [\alpha_i,\gamma]=1,\quad
 h_i^{m_i}=\gamma,\quad h_i^{-n_i}=\alpha_i
                         \right.\right\rangle .       \tag{F45.16}
\]
All elliptic path coefficients have disappeared here because
their full subgroup has just been killed, not because a
different boundary marking was substituted. With every smooth
and multiplicative $m_i=1$, each corresponding $h_i$ equals
$\gamma$. There is at most one other generator $h_*$, and
$\gamma=h_*^{m_*}$. The sole remaining relation is
$h_*^{n_*+m_*\sum_{i\ne *}n_i}=1$. This proves the cyclic
claim and the exact signs in (F45.2). Without an additive site
the relation is $\gamma^{\sum_i n_i}=1$.

\item For $D=\pm1$, (F45.7)--(F45.14) give the integral homology
of $S^5$ and (F45.16) gives simple connectivity. The resulting
oriented orbit is the standard $S^5$. The actual delta circle
is free in every added cap, and its single fixed sphere at the
selected cusp has the primitive positive Hopf meridian. The
Gysin and fixed-tube maps of P3K.22 then give simple connectivity
and the integral homology of $S^6$ for $X$. The vanishing of
$\Theta_6$ gives its standard oriented diffeomorphism type.
Conversely the orbit calculation of P3F.7 applied to a sphere
member forces these same orbit groups to vanish, giving
(F45.1)--(F45.2). The actual common blowup and finite flop chain
of P3F.35 transports this topology criterion to every allowed
cusp fan. All local finite characters, base powers and
order-one gluing maps remain those of the constructed family.

\end{enumerate}
\end{proof}




\begin{theorem}\label{night:P3F.46}
\StatementRecord{night:P3F.46}{proved}{S6\_06\_BLIND: GPT-6 Astra (OpenAI, ChatGPT Pro mode, effort Pro), recorded re-derivation in the checkpoint dated 18 November 2026}{\cref{night:P3F.35}, \cref{night:P3K.22}}
Let $S$ be the minimal regular surface of
\[
 y^2=x^3+tx+t^2,\qquad G=(0,t),\qquad
 M=\mathcal O_S(aG-O),\qquad Q=bG.
\]
Use its $I_1$ fibre at $t_c=-4/27$ as the sole simple whole-fibre
zero. The following four pairs admit the actual two-cap families
constructed below:
\[
\begin{array}{c|c|c}
(a,b)&m_6\text{ at }IV^*&m_2\text{ at }III\\ \hline
(1,6)&3l,\ l\ge1\text{ (good)}&4k,\ k\ge1\text{ (good)}\\
(2,3)&3l,\ l\ge1\text{ (good)}&4k+1,\ k\ge0\text{ (reduced)}\\
(3,2)&3l+2,\ l\ge0\text{ (reduced)}&4k,\ k\ge1\text{ (good)}\\
(6,1)&3l+2,\ l\ge0\text{ (reduced)}&4k+1,\ k\ge0\text{ (reduced)}.
\end{array}                                                     \tag{F46.1}
\]
For each pair choose signed integers $n_i$ with
$\gcd(m_i,n_i)=1$. The marked meridians are
$m_i\alpha_i+n_i\gamma$, where $\gamma$ is the outward
suspension generator. With $L=\operatorname{lcm}(m_6,m_2)$,
the construction is a proper holomorphic family over a punctured
$r$-disc with scalar $\lambda=r^L$. Its fibres are the standard
oriented $S^6$ exactly when
\[
                  D=m_2n_6+m_6n_2\in\{1,-1\}.                 \tag{F46.2}
\]
This criterion holds for every allowed fan. The statement assigns
the integral topology and retains the actual parameter closing;
it assigns no parameter-return class.
\end{theorem}
\begin{proof}
\begin{enumerate}[label=\arabic*.]
\item \par\medskip\noindent\textbf{The source and the exact root maps.}
The discriminant is $t^3(4+27t)$, so the fibres are $III$ at zero,
$I_1$ at $t_c$, and $IV^*$ at infinity. On $x\ne0$ put
$u=(y-t)/x$ and $v=(y+t)/x$. The inverse formulas are
\[
 x=uv-(v-u)/2,\quad t=x(v-u)/2,\quad y=x(u+v)/2.
\]
Thus $S$ is rational. The marked section meets nonidentity
components in $E_6$ and $A_1$, has $G\cdot O=0$, and has
height $1/6$. The integral source Gram matrix below has determinant
$-1$; the discriminant and height identity give primitive $G$
and no Mordell--Weil torsion.

\item The source map is $h=t_{-bG}$. Write $T$ for its pushforward on
divisor cycles. The exact section formula below gives
$TM-M=-F$. Its pullback is $T^{-1}$, so $h^*M-M=F$.
Consequently the forward line map $M_p\to M_{h(p)}$ is a section
of $h^*M\otimes M^{-1}=\mathcal O(F)$. Choose its section with
the sole simple zero at $t_c$, denote it by $B_0$, and put
$B_\lambda=\lambda B_0$. This specifies the line direction.

\item At infinity set $q=1/t$, $x_0=q^2x$, $y_0=q^3y$. Then
\[
 y_0^2=x_0^3+q^3x_0+q^4,\qquad G=(0,q^2).
\]
For $m_6=3l+2$, substitute
$q=z^{m_6}$, $x_0=z^{4l+2}X$, $y_0=z^{6l+3}Y$. We obtain
\[
 Y^2=X^3+z^{l+2}X+z^2,\qquad G'=(0,z).                         \tag{F46.3}
\]
Its $A_2$ resolution has a reduced $IV$ fibre. In
$X=zX_1,Y=zY_1$, the nonidentity labels are $Y_1=1,-1$.
For $k=2l+1$ the actual deck is
\[
 D(X,Y,z)=(\zeta_m^{-2k}X,\zeta_m^{-3k}Y,\zeta_m z).
\]
Its character on $Y_1$ is
$\zeta_m^{-(3k+1)}=\zeta_m^{-2m}=1$, so it fixes all
three labels. For $m_6=3l$, the substitution $k=2l$ instead gives
\[
 Y^2=X^3+z^lX+1,\qquad G'=(0,1),                               \tag{F46.4}
\]
with central deck $(X,Y)\mapsto(\zeta_3^2X,Y)$.
The point $(0,1)$ has exact order three: its horizontal tangent
has triple intersection with $Y^2=X^3+1$.

\item At the $III$ site take $q=t$. For $m_2=4k+1$, put
$q=z^{m_2}$, $x=z^{2k}X$, $y=z^{3k}Y$. The equation is
\[
 Y^2=X^3+zX+z^{2k+2},\qquad G'=(0,z^{k+1}).                    \tag{F46.5}
\]
The $A_1$ resolution gives the reduced $III$ fibre, whose two
labels are fixed by the deck. For $m_2=4k$ the same substitution
gives
\[
 Y^2=X^3+X+z^{2k},\qquad G'=(0,z^k).                           \tag{F46.6}
\]
Its central deck is $(X,Y)\mapsto(-X,iY)$ and its marked
limit $(0,0)$ is the nonzero fixed two-torsion point.
These substitutions and the stated resolutions are the actual
maps to the source models.

\par\medskip\noindent\textbf{The two line comparisons.}
At a reduced root, $aG'$ has identity component and
$N=\mathcal O(aG'-O)$ has degree zero on every component.
The original $M$ has the same zero multidegrees at that site.
Pull it to the resolved fibre product first. Every curve
contracted in relative minimalization has degree zero, so the
blowup Picard decomposition descends the line through each
actual contraction. Its comparison with $N$ is a vertical
divisor with zero component degrees. The reduced $III$ or $IV$
intersection matrix has integral kernel the whole-fibre vector.
The corresponding base monomial therefore gives the line
isomorphism and retains its deck character and the specified $B_0$.

\item For a good $IV^*$ root, the minimal cubic deck has tangent
weights $(2,1)$ at its three fixed points. Its quotient lattice
is $\mathbb Z^2+\mathbb Z(2,1)/3$, and the minimal quotient
rays are $((3-j)/3,j/3)$. A common smooth root refinement has rays
\[
                  (1,0),(2,1),(1,1),(1,2),(0,1).
\]
Let $\beta$ be this root blowup and $\pi$ its toric map to the
minimal quotient resolution. The section $w=0$ has Cartier
support $h(v)=\max(v_w-2v_z,0)$ there. Pulling back this
function gives the exact section-divisor difference
\[
 \beta^*(\text{root section})-\pi^*(\text{resolved section})
       =\sum_v\min(v_w,2v_z)E_v,
\]
with coefficients $(2,1,1)$. Subtract the $O$ contribution from
the $aG'$ contribution to obtain
\[
                    \beta^*N=\pi^*M+D_{aG}-D_O.                \tag{F46.7}
\]
The invariant Weierstrass substitution is the source map, and
its crepant quotient resolution is the original $IV^*$ model.
Since $Q'_0=O$ for $b=3,6$, its translation is locally
$w\mapsto w\,u+z^2h(z^3)$. It preserves
$(w,z),(w,z^2),(z,w),(z,w^2)$, hence the common normalized
blowup and every correction in (F46.7). For $m_6=3l$, pull this
entire correspondence through the extra degree-$l$ root and
resolve the fibre product. The divisor equality and maps pull
back exactly.

\item For a good $III$ root, the minimal fourth deck has two
order-four fixed points $O,P$ and one orbit $\{R_1,R_2\}$
with stabilizer two. The tangent weights are $(1,1)$.
Blow up these four points by $\beta$ and take the quotient
$\pi$. Its exceptional images $D_0,D_1,D_2$ have squares
$-4,-4,-2$, and the central image $C$ has square $-1$.
The actual fibre is $4C+D_0+D_1+2D_2$. Contract $C$, then the
resulting $(-1)$-curve $D_2$. The remaining curves have squares
$-2$ and intersection two; the invariant Weierstrass map gives
their original $III$ tangency. The marked sections avoid both
contracted curves. Thus, for these contractions $\alpha$,
\[
                 \beta^*N=\pi^*\alpha^*M+E_{aG}-E_O.           \tag{F46.8}
\]
For $b=2,6$ we have $Q'_0=O$, so its translation fixes all four
centres and preserves the correction. This carries the actual
$B_0$ and its deck commutation to $N$. Higher fourth roots
pull back this whole diagram and its characters.

\item The good central line has nonzero torsion Picard class of
order three at $IV^*$ and order two at $III$. The actual map
\[
 (\mathbb C/\Lambda)^D\longrightarrow\Lambda/(D-1)\Lambda,
                \quad[v]\longmapsto[(D-1)v]
\]
carries it to the primitive torsion Euler class of the cap
circle line. A common deck character is retained and does not
change that torsion order.

\par\medskip\noindent\textbf{The exact norms and the joint family.}
At either site write $m,n,N,D$ for its data. Multiplication by
$m$ is invertible on the reduced additive identity group and
component group. At a good root use its analytic inverse branch
at $O$. Thus the unique marked analytic division $mR=Q'$ is
$D$-invariant. The source map $t_{-nR}D$ has $m$th power
$t_{-nQ'}$. It extends on the minimal regular root: resolve the
addition graph and compare the invariant holomorphic two-form;
the common exceptional support gives the holomorphic map and
its inverse.

\item Lift $t_{-R}$ to $F$ on $N$. The fibre chord divisor gives its
Hom section, and its residual vertical divisor has zero
component degrees. Remove the resulting whole-fibre base power.
Finite deck averaging then gives
$DFD^{-1}=\chi F$, with $\chi^m=1$.
The specified descended lift is $B_0=\beta(q)F^m$.
The ratio is a base unit because the fibres are compact and
connected, and deck commutation makes it invariant. Hence
\[
       (F^nD)^m/B_0^n
          =\chi^{nm(m-1)/2}\beta(q)^{-n}.
\]
Divide by an invariant analytic $m$th root of this unit.
The resulting $A_0$ satisfies $A_0^m=B_0^n$ and
$[A_0,B_0]=1$, retaining the character factor for even $m$.

\item Set $\lambda=r^L$ and
$A_{i,r}=r^{Ln_i/m_i}A_{i,0}$. For integers $u_i,v_i$ with
$n_i u_i+m_i v_i=1$, put
\[
 H_i=A_{i,r}^{u_i}B_\lambda^{v_i}
       =r^{L/m_i}H_{i,0},\qquad
 H_i^{m_i}=B_\lambda,\quad H_i^{n_i}=A_{i,r}.                  \tag{F46.9}
\]
On smaller compact source tubes the norm ratios of $H_{i,0}$
and its inverse are bounded. Thus $H_i$ is strictly contracting
for small $r$, and its norm shell gives a smooth proper cap
$N_i^\times/\langle H_i\rangle$. A unit rotation equal to
$H_i^j$ forces $j=0$ by contraction, then forces the unit to
be one. The delta circle is therefore free on both caps.
For outward $h_i=H_i^{-1}$, the single-$A$ cut gives exactly
\[
 \alpha_i\longmapsto-n_i h_i,\qquad
 \gamma\longmapsto m_i h_i,\qquad
 \mu_i=m_i\alpha_i+n_i\gamma.                                 \tag{F46.10}
\]

\item The holomorphic annular map is constructed in the actual
constant-period extension $(\mathbb C\times\mathbb C^*)/\Lambda$.
Remove its integral period two-cocycle on the annulus, whose
$H^2$ is zero, then solve the holomorphic Lie-bundle cocycle
on the Stein annulus. Choose the gamma component whose positive
cut is $A_{i,r}$ itself. For $c=\log\lambda/(2\pi i)$ its scalar
variation is $n_i(c-c_0)s_i\delta/m_i$. This is (F46.10) with
every elliptic and delta phase retained.

\item For a positive $r$-loop put $k_i=Ln_i/m_i$. The inverse closing
near cap $i$ is $q_i^{k_i}$ times its extending unit. On the
ordinary region the exact global multiplier is
\[
                  f(t)=t^{k_2}(t-t_c)^{-(k_6+k_2)}.              \tag{F46.11}
\]
Its orders are $k_2$ at zero, $k_6$ at infinity, and
$-(k_6+k_2)$ at the selected cusp. Its ratios to the two local
multipliers are units. The cap charts close by the identity,
and their annular maps supply these multipliers. At the cusp,
(F46.11) is, modulo the specified $B_\lambda$ iterates, a bounded
source translation and an invertible scalar unit. It and its
inverse extend in the actual nodal Mumford charts. This proves
holomorphic parameter closing. Compact norm shells give
properness on each compact parameter annulus.

\par\medskip\noindent\textbf{The full source and boundary complexes.}
Use the integral basis $O,F,C_0,\ldots,C_5,A_1,G$.
The $E_6$ graph has edges
$(0,1),(1,2),(2,3),(3,4),(2,5)$, Cartan matrix $A$, and
highest-root marks $v=(1,2,3,2,1,2)^t$.
The source Gram matrix $B$ has $O^2=G^2=-1$, $OF=GF=1$,
root blocks $-A$ and $(-2)$, and $GC_0=GA_1=1$;
all unspecified pairings vanish. Its determinant is $-1$,
so these are the entire integral source lattice.

\item For an integer $j$, let $R_j$ be the inverse-Cartan column at
$C_0$ or $C_4$ when $j=1$ or $2$ modulo three, together
with $A_1/2$ when $j$ is odd. Let $c_j$ be the sum of the
corresponding diagonal entries $4/3$ and $1/2$.
The actual section formula is
\[
 \phi=G-O-F+R_1,\quad \phi^2=-1/6,\qquad
 P_j=O+\frac{j^2/6+c_j}{2}F+j\phi-R_j.                          \tag{F46.12}
\]
It is integral, has square $-1$, and has the prescribed component
intersections. Translation by $G$ sends $O$ to $G$, $G$ to
$P_2$, cyclically permutes the three affine $E_6$ arms, and
swaps $A_1$ with $F-A_1$. These actual columns define the
integral isometry $T_+$, and $T=T_+^{-b}$.

\item Put
\[
 L_6=(F-\textstyle\sum v_jC_j,C_0,\ldots,C_5),\quad
 L_2=(F-A_1,A_1),\quad L=(L_6,L_2).
\]
Their multiplicity vectors are $(1,v)^t$ and $(1,1)^t$;
both products with $L_i$ equal $F$. If $P$ is the actual
component permutation, then $TL=LP$.
For the complement $V$ of the two fibre tubes the integral
source/relative complex is
\[
 C_4=\mathbb Z,\ C_3=\mathbb Z^2,\ C_2=\operatorname{NS}(S),
 \ C_1=\mathbb Z^9,\ C_0=\mathbb Z,\qquad
 d_4=(1,1)^t,\quad d_2=L^tB,\quad d_3=d_1=0.                  \tag{F46.13}
\]
Boundary $i$ has $C_3=\mathbb Z$, $C_2=C_1=\mathbb Z^{r_i}$,
$C_0=\mathbb Z$, and $d_2=L_i^tBL_i$.
Its inclusion is $L_i$ in degree two, the component-coordinate
embedding in degree one, and the $i$th coordinate in degree
three. These are the complexes of the actual pair. The source
and the tube retracts have free even homology, so the degree-one
derived homotopy group of their even split complexes is zero.
The displayed maps therefore have no omitted odd homotopy
parameter. This argument includes the actual tangencies.

\item The source action is $T$ in degree two, $P^{-t}$ in degree one,
and identity otherwise. The boundary action is $P_i$ in degree
two, $P_i^{-t}$ in degree one, and identity otherwise.
The program forms both actual mapping-torus cones, with
\[
       d_{{\rm cone},j}=
       \begin{pmatrix}d_j&T_{j-1}-I\\0&-d_{j-1}\end{pmatrix},
\]
and their block-diagonal inclusions. All chain matrices,
integral Smith representatives, projections, and marked
$\alpha,\gamma,N_3,F\gamma$ cycles are exported. They give
\[
 H_1(W)=\mathbb Z^2,\quad H_3(W)=\mathbb Z^2,\quad
 H_4(W)=\mathbb Z,\quad
 H_2(W)=\mathbb Z/b\oplus\mathbb Z^2,                           \tag{F46.14}
\]
with the trivial summand omitted when $b=1$.

\item We also compute the full fundamental group. The remaining
$I_1$ kills its primitive vanishing cycle $v$. If $A$ is the
$III$ monodromy, then $A^2=-I$ and $\operatorname{tr}A=0$.
The product relation with $IV^*$ and this $I_1$ gives
$|\langle v,Av\rangle|=1$: expand
$\operatorname{tr}(A^{-1}T_v^{-1})$, whose absolute value is
that pairing and is one from the inverse $IV^*$ trace.
Thus $v,Av$ form an integral elliptic basis. Conjugating the
vanishing relation by the retained base loop kills both.
The section gives the remaining base-annulus generator, so
$\pi_1(V)=\mathbb Z$. Translation fixes it and
\[
             \pi_1(W)=\mathbb Z\alpha\oplus\mathbb Z\gamma.
\]
The two boundary base meridians map to $\alpha,-\alpha$;
both outward suspension loops map to $\gamma$.

\par\medskip\noindent\textbf{The cap specializations.}
A good cap has
$H_1=\mathbb Zh\oplus\mathbb Z/p$, $H_2=\mathbb ZC$,
$H_3=\mathbb Z(Ch)$, $H_4=0$, with $p=3$ or two.
Actual smooth-root transport sends the generic fibre to its
central elliptic fibre with degree one. The same transport is
the identity on the elliptic lattice. If $nu+mv=1$, then $u$
is prime to the deck order and
\[
                    (D^u-I)\Lambda=(D-I)\Lambda.
\]
For order three this follows from $D+I=-D^2$; for order four
the relevant odd powers are $D$ or $D^{-1}$.
The multiplying matrices are unimodular. Thus the incoming
finite torsion maps isomorphically to cap torsion, and the
incoming full elliptic group maps onto the cap elliptic lattice.

\item A reduced cap has $H_1=\mathbb Zh$, $H_2=\mathbb ZC$,
$H_3=\mathbb Z(\sum C_j)$, $H_4=0$.
Its primitive $H_i$ cyclically permutes all central components:
its component translation is $-Q/m_i$, independent of $n_i$
by (F46.9), and is nonzero modulo three or two.
The actual fibre specializes to $\sum C_j$. Its image in
coinvariants is $r_iC$, while its invariant trace is primitive.

\item Use boundary $H_2$ coordinates $(\eta_i,r_i,s_i)$ in the good
case, with $\eta_i$ of order $p_i$ and $F=r_i$.
In the reduced case use $(r_i,s_i)$, with $F=3r_i$ at $IV^*$
and $F=2r_i$ at $III$. The actual cap maps $r_i\mapsto C_i$.
Write its remaining column as $s_i\mapsto\kappa_iC_i$ and
retain that integer.
The full exterior inclusions, with torsion row first, are
\[
\begin{array}{c|c|c}
(a,b)&IV^*&III\\ \hline
(1,6)&\begin{pmatrix}-2&1&6\\0&0&-1\\0&0&-1\end{pmatrix}
 &\begin{pmatrix}3&1&3\\0&0&0\\0&0&1\end{pmatrix}\\[5pt]
(2,3)&\begin{pmatrix}1&2&7\\0&0&0\\0&0&-1\end{pmatrix}
 &\begin{pmatrix}1&0\\0&-1\\0&2\end{pmatrix}\\[5pt]
(3,2)&\begin{pmatrix}-1&18\\0&1\\0&3\end{pmatrix}
 &\begin{pmatrix}-1&-3&-1\\0&0&0\\0&0&1\end{pmatrix}\\[5pt]
(6,1)&\begin{pmatrix}0&-1\\0&3\end{pmatrix}
 &\begin{pmatrix}0&-1\\0&2\end{pmatrix}.
\end{array}                                                     \tag{F46.15}
\]
The first row is taken modulo $b$ when $b>1$.
For $b=6$ the exact CRT map is
$\mathbb Z/6\to\mathbb Z/3\oplus\mathbb Z/2$,
$1\mapsto(\eta_6,\eta_2)$, with exterior coefficient
$-2+3=1$. For $b=3$ the torsion generator already maps to one.
For $b=2$ use minus the incoming $III$ torsion generator.

\item In domain order $(\eta,r_6,s_6,r_2,s_2)$ and target order
$(\tau,g_1,g_2,C_6,C_2)$, the complete degree-two matrices are
\[
 D_{2,6}=\begin{pmatrix}
1&1&0&1&3\\0&0&-1&0&0\\0&0&-1&0&1\\
0&-1&-\kappa_6&0&0\\0&0&0&-1&-\kappa_2
\end{pmatrix},\quad
 D_{2,3}=\begin{pmatrix}
1&2&1&1&0\\0&0&0&0&-1\\0&0&-1&0&2\\
0&-1&-\kappa_6&0&0\\0&0&0&-1&-\kappa_2
\end{pmatrix}.                                                  \tag{F46.16}
\]
\[
 D_{2,2}=\begin{pmatrix}
1&1&0&1&1\\0&0&1&0&0\\0&0&3&0&1\\
0&-1&-\kappa_6&0&0\\0&0&0&-1&-\kappa_2
\end{pmatrix},\qquad
 D_{2,1}=\begin{pmatrix}
0&-1&0&-1\\0&3&0&2\\-1&-\kappa_6&0&0\\
0&0&-1&-\kappa_2
\end{pmatrix}.                                                  \tag{F46.17}
\]
Their determinants are $1,1,-1,-1$. For $b>1$, the first
coordinate and row have order $b$; the first column has no
free entry. Thus these are integral group isomorphisms for
every value of both actual collar coefficients.

\par\medskip\noindent\textbf{All remaining gluing degrees.}
In degree one every good boundary torsion maps isomorphically
to cap torsion and to zero in $W$. The exported full matrices
retain arbitrary torsion path coefficients $c_i,d_i$ in the
images of $\alpha_i,\gamma$. Eliminating precisely those finite
isomorphisms gives, in physical order
$(\alpha_6,\gamma_6,\alpha_2,\gamma_2)$,
\[
 D_1=\begin{pmatrix}
1&0&-1&0\\0&1&0&1\\n_6&-m_6&0&0\\0&0&n_2&-m_2
\end{pmatrix}.                                                  \tag{F46.18}
\]
In degree three use the actual oriented source boundary $N_i$
followed by $F\gamma_i$. The cap map is $(-n_i,m_i)$:
the norm-circle degrees in (F46.10) are followed by specialization
to the primitive invariant central fibre. The cap source
retracts to its two-dimensional central fibre, so it contributes
no additional three-class. The exterior has $N_6=-N_2$ and
the same $F\gamma$ at both ends. In the exported Smith bases
the $IV^*$ boundary change is $\operatorname{diag}(1,-1)$,
the $III$ change is identity, and the exterior target change
for $b=1,2$ is $\left(\begin{smallmatrix}0&-1\\1&0\end{smallmatrix}\right)$;
for $b=3,6$ it is identity. Consequently the exact physical
$D_3$ equals (F46.18), and
\[
                   \det D_1=\det D_3=-D.                        \tag{F46.19}
\]
In degree four the map is $(1,-1):\mathbb Z^2\to\mathbb Z$,
up to the recorded common target orientation. In degree zero
it is the actual incidence matrix
$\left(\begin{smallmatrix}1&1\\-1&0\\0&-1\end{smallmatrix}\right)$.
This specifies every nonzero Mayer--Vietoris degree.

\par\medskip\noindent\textbf{The exact fundamental group and sphere criterion.}
The full incoming elliptic lattice dies in $V$. In each good
cap it is the full normal elliptic subgroup, so van Kampen
kills it. A reduced cap has no remaining elliptic subgroup.
The exact residual presentation is
\[
 \left\langle h_6,h_2\ \middle|\
 h_6^{m_6}=h_2^{m_2},\quad h_6^{n_6}=h_2^{-n_2}\right\rangle.
\]
Choose $u,v$ with $n_6u+m_6v=1$. Then
\[
 h_6=(h_6^{n_6})^u(h_6^{m_6})^v
       =h_2^{-n_2u+m_2v},
\]
so the group is cyclic. Its abelian relation matrix has a
primitive relation and determinant $D$, giving
$\pi_1(Y)=\mathbb Z/|D|$, with $\mathbb Z$ when $D=0$.
The full integral Mayer--Vietoris sequence gives, for $D\ne0$,
\[
 H_0(Y)=H_5(Y)=\mathbb Z,\quad
 H_1(Y)=H_3(Y)=\mathbb Z/|D|,\quad H_2(Y)=H_4(Y)=0.             \tag{F46.20}
\]
For $D=0$, the first three Smith factors of (F46.18) are one
and the last is zero; all $H_j(Y)$, $0\le j\le5$, are
$\mathbb Z$. Thus $Y$ is the standard $S^5$ exactly when
$D=1$ or $-1$.

\item Both added caps have free delta action. The only fixed locus
is the chosen nodal-fibre sphere, with the primitive
$\mathcal O(-1)^{\oplus2}$ Hopf chart. The exact Gysin and
fixed-tube maps of P3K.22 give $\pi_1(X)=\pi_1(Y)$ and, for
$D=\pm1$, the integral $S^6$ homology. Since $\Theta_6=0$,
the oriented total space is the standard sphere. For all other
$D$ its displayed nontrivial fundamental group excludes a
sphere. The actual finite flop chain of P3F.35 is the identity
outside the cusp and transports this criterion to every allowed
fan. Equation (F46.2) also forces coprimality of the two orders
and of each local norm pair.

\item The complete certificate is
\path{code/part3/iv_star_iii_joint.py}, with its saved JSON and
output. It contains every source and boundary chain matrix,
Smith representative, marked cycle, finite-group degree-one
map, symbolic collar coefficient, and degree-three orientation
change. The four degree-two determinants and the common
degree-one/three determinant are symbolic identities. A further
$1056$ exact Smith cases, including $D=0$, supplement these
all-order calculations.

\end{enumerate}
\end{proof}




\begin{theorem}\label{night:P3F.47}
\StatementRecord{night:P3F.47}{proved}{S6\_06\_BLIND: GPT-6 Astra (OpenAI, ChatGPT Pro mode, effort Pro), recorded re-derivation in the checkpoint dated 18 November 2026}{\cref{night:P3F.35}, \cref{night:P3F.37}, \cref{night:P3F.41}, \cref{night:P3F.43}, \cref{night:P3F.46}, \cref{night:P3K.22}}
On the minimal regular rational elliptic surface
\[
 y^2=x^3+tx-t,\qquad G=(1,1),
\]
use $M=\mathcal O(aG-O)$ and the actual suspension translation
$-bG$, for $(a,b)=(2,1)$ or $(1,2)$. The fibres are
$III^*$ at infinity, II at zero, and $I_1$ at $t_c=-27/4$.
Choose the simple whole zero of the line map at that $I_1$.
The following are actual simultaneous root-line cap domains:
\[
\begin{array}{c|c|c}
(a,b)&m_*\text{ at }III^*&m_0\text{ at II}\\ \hline
(2,1)&4l+3, l\geq0&6k+1, k\geq0,\quad\text{or }6k, k\geq1\\
(1,2)&4l, l\geq1&6k+1, k\geq0,\quad\text{or }6k, k\geq1
\end{array}                                            \tag{F47.1}
\]
Retain signed $n_*,n_0$ with $\gcd(m_i,n_i)=1$, including
$n_0=0$ at an order-one ordinary II cap. Put
\[
                         D=m_0n_*+m_*n_0.              \tag{F47.2}
\]
For the constructed joint family,
$\pi_1(Y)=\pi_1(X)=\mathbb Z/|D|$, with $\mathbb Z$ for $D=0$.
For $D\ne0$ the orbit groups are
\[
 H_0(Y)=H_5(Y)=\mathbb Z,\qquad
 H_1(Y)=H_3(Y)=\mathbb Z/|D|,\qquad H_2(Y)=H_4(Y)=0.
\]
For $D=0$ every $H_j(Y)$, $0\leq j\leq5$, is $\mathbb Z$.
The actual sphere tuples are exactly $D=\pm1$, for every
permitted cusp fan. The condition forces $\gcd(m_*,m_0)=1$;
in particular the good $III^*$ and good II roots cannot occur
together in a sphere tuple. This construction and computation
retain all line characters and collar coefficients.
\end{theorem}
\begin{proof}
\begin{enumerate}[label=\arabic*.]
\item \par\medskip\noindent\textbf{The source, its actual roots, and the line maps.}\enspace
The source and its primitive height-$1/2$ section were derived
in P3F.37 and P3F.41. Its displayed discriminant has precisely
the three stated fibres. The integral source calculation below
also verifies the marked height and the full lattice.
For $h=t_{-bG}$, the exact section formula gives
$h^*M\otimes M^{-1}=\mathcal O(F)$ in both cases $ab=2$.
Its section with simple zero at $t_c$ supplies the actual forward
line map $B_0$, and $B_\lambda=\lambda B_0$ fixes its scalar
parameter and direction.

\item At infinity put $q=1/t$, $x_0=q^2x$, $y_0=q^3y$.
Then
\[
 y_0^2=x_0^3+q^3x_0-q^5,\qquad G=(q^2,q^3).
\]
For $m_*=4l+3$, set $q=z^{m_*}$ and divide $x_0,y_0$
by $z^{2(3l+2)},z^{3(3l+2)}$. The exact result is
\[
 Y^2=X^3+zX-z^{2l+3},\qquad
 G'=(z^{2l+2},z^{3l+3}).                            \tag{F47.3}
\]
Resolving its ordinary double point gives the reduced III
root of P3F.37. The marked $G'$ meets its nonidentity component,
whereas $2G'$ meets the identity component. For $m_*=4l$,
divide instead by $z^{6l},z^{9l}$ to obtain
\[
 Y^2=X^3+X-z^{2l},\qquad G'=(z^{2l},z^{3l}).         \tag{F47.4}
\]
Its good central fibre has deck $(X,Y)\mapsto(-X,-iY)$,
and $G'_0=(0,0)$ is the nonzero fixed two-torsion point.
Thus $bG'_0=O$ for the good pair, while the central line
$\mathcal O(G'_0-O)$ retains exact order two.

\item At zero use $q=t$. For $m_0=6k+1$ and $m_0=6k$, respectively,
put $q=z^{m_0}$, $x=z^{2k}X$, $y=z^{3k}Y$.
The actual equations are
\[
 \begin{aligned}
 Y^2&=X^3+z^{2k+1}X-z &&(m_0=6k+1),\\
 Y^2&=X^3+z^{2k}X-1   &&(m_0=6k).
 \end{aligned}                                         \tag{F47.5}
\]
The first has a reduced II fibre; the second has good reduction
with order-six central deck. The marked section is
$(X,Y)=(z^{-2k},z^{-3k})$ on the affine overlap. In the actual
smooth chart at $O$, its parameter is $-X/Y=-z^k$.
Hence its good central limit, and the limit of $bG'$, is $O$.
The formal group law gives $w(aG')=-az^k+O(z^{2k})$.
This proves the central section data used in the line comparison,
including the further contact of nonminimal root orders.

\item Write $N_i=\mathcal O(aG'_i-O_i)$ on either root. At a
reduced root the actual resolved fibre-product comparison
descends the pulled line through each contracted exceptional
curve, since its degree there is zero. Its difference from
$N_i$ is a vertical divisor with zero component degrees.
The exact III or II intersection matrix makes it an integral
whole fibre. Dividing its base monomial supplies the line
isomorphism and retains the corresponding deck character.

\item At the good starred root the minimal fourth deck has weights
$(3,1)$ at each fixed section. The common smooth toric refinement
has rays
\[
 (1,0),(3,1),(2,1),(1,1),(1,2),(1,3),(0,1).
\]
For its root blowup $\beta$ and resolved quotient map $\pi$,
the Cartier support function of the section is
$\max(v_w-3v_z,0)$. Consequently
\[
 \beta^*N_* =\pi^*M+D_{aG'}-D_O,\qquad
 D_P=\sum_v\min(v_w,3v_z)E_{P,v}.                    \tag{F47.6}
\]
The five exceptional coefficients are $(3,2,1,1,1)$.
The actual translation by $bG'$ at the minimal fourth root
has form $w\mapsto w\,u(w,z)+z^3h(z^4)$, with $u$ a unit.
It preserves the ideals $(w,z^j),(z,w^j)$ for $j\leq3$,
hence the normalized common blowup and (F47.6).
The further degree-$l$ root pulls this diagram back, including
its exceptional coefficients and the specified line map.

\item At the good II root, the minimal sixth deck has one point of
stabilizer six, one orbit of two points of stabilizer three,
and one orbit of three points of stabilizer two. Blow them up.
Their tangent weights are $(1,1)$. The quotient central curve
and exceptional images have squares $-1,-6,-3,-2$ and fibre
\[
                          6C+D_6+2D_3+3D_2.
\]
Contract $C$, then $D_2$, then $D_3$. The remaining curve
has square zero and the actual invariant Weierstrass map has
the cusp parametrization $(u^2,u^3)$ of the original II fibre.
Both $aG'$ and $O$ meet the exceptional divisor over $O$ at
finite $w/z$; all three contraction centres are at its infinity.
Each section has exceptional coefficient one. They cancel in
the exact divisor identity
\[
                  \beta^*N_0=\pi^*\alpha^*M.           \tag{F47.7}
\]
Here $\alpha$ denotes the stated contractions. Translation by
$bG'$ fixes the original blowup centres and preserves the
diagram and this equality. Pulling the diagram through the
additional degree-$k$ root proves the same assertion for
every good order in (F47.1). These are the actual line
correspondences of P3F.41 and P3F.43, now with the marked
$aG'$ and $bG'$ of this source explicitly verified.

\par\medskip\noindent\textbf{Commuting norms and the simultaneous family.}\enspace
At either root, choose the actual $D_i$-invariant division
section $R_i$, $m_iR_i=bG'_i$. At a reduced root multiplication
by $m_i$ is invertible on both the additive central identity
group and its component group. At a good root take the unique
analytic inverse branch with central value $O$.
The extending source action $t_{-n_iR_i}D_i$ has $m_i$th
power $t_{-n_i bG'_i}$. Resolve its addition graph as in
P3F.41--43; the source maps and their inverses extend on the
minimal regular root.
Lift $t_{-R_i}$ first to the specified line $N_i$.
Its exact deck commutator is a retained character $\chi_i$.
Writing $B_0=\beta_i(q_i)F_i^{m_i}$ gives the invariant
base unit $\beta_i$ and the norm factor
\[
 (F_i^{n_i}D_i)^{m_i}/B_0^{n_i}
    =\chi_i^{n_im_i(m_i-1)/2}\beta_i(q_i)^{-n_i}.
\]
Its invariant analytic $m_i$th root gives the actual commuting
line maps $A_{i,0}$ with $A_{i,0}^{m_i}=B_0^{n_i}$.
This retains the order-two central line in (F47.6) and every
scalar character in (F47.7).

\item Put $L=\operatorname{lcm}(m_*,m_0)$, $\lambda=r^L$, and
$A_{i,r}=r^{Ln_i/m_i}A_{i,0}$. For $n_iu_i+m_iv_i=1$,
the literal primitive generator
\[
 H_i=A_{i,r}^{u_i}B_\lambda^{v_i},\qquad
 H_i^{m_i}=B_\lambda,\quad H_i^{n_i}=A_{i,r}
\]
contains the positive scalar power $r^{L/m_i}$. Its norm shell
is uniformly contracting on smaller compact source tubes for
small $r$. Thus $N_i^\times/\langle H_i\rangle$ is a smooth
proper holomorphic cap. A unit line rotation equal to $H_i^j$
forces $j=0$ by contraction, so the delta circle is free there.
For outward $h_i=H_i^{-1}$ the actual cut marking is
\[
                 \alpha_i\longmapsto-n_i h_i,\qquad
                 \gamma\longmapsto m_i h_i.          \tag{F47.8}
\]
The holomorphic annular comparison is obtained by the actual
constant-factor degree-zero line uniformization, removing its
period two-cocycle on the annulus and then its holomorphic
Lie-bundle cocycle. Its positive cut is $A_{i,r}$ itself.
For $c=\log\lambda/(2\pi i)$ its line variation is exactly
$n_i(c-c_0)s_i\delta/m_i$; elliptic and delta periods are
retained in this equation.

\item For one positive $r$-loop set $k_i=Ln_i/m_i$. The ordinary
inverse closing has the global multiplier
\[
                 t^{k_0}(t-t_c)^{-(k_*+k_0)}.         \tag{F47.9}
\]
Its orders at zero and infinity are $k_0,k_*$, and its ratios
to the prescribed local multipliers are extending units.
The cap charts close by the identity in these root coordinates.
At the selected cusp, use the specified $B_\lambda$ iterates
to replace the remaining integral power by its bounded source
translation and invertible line unit. They and their inverses
extend in the actual unimodular Mumford charts. This proves
the simultaneous proper punctured-disc family and fixes its
parameter closing. For $D=\pm1$, $L=m_*m_0$, so its cusp
gamma coefficient is exactly the signed integer $D$.

\par\medskip\noindent\textbf{The source cone and all integral boundary maps.}\enspace
Use the basis $(O,F,C_0,\ldots,C_6,G)$.
The $E_7$ graph has edges
$(0,1),(1,2),(2,3),(3,4),(4,5),(2,6)$ and Cartan matrix $A$.
Its highest-root vector is $v=(2,3,4,3,2,1,2)^t$.
The Gram matrix has $O^2=G^2=-1$, $OF=GF=1$,
root block $-A$, and $G C_5=1$, with every other unspecified
pairing zero. Its determinant is $-1$.
Put $R$ equal to the inverse-Cartan column at $C_5$,
$\varphi=G-O-F+R$, and $R_j=R$ for odd $j$, zero for even $j$.
Then the exact section is
\[
 P_j=O+\left(1+\frac{j^2/2-2+c_j}{2}\right)F
                 +j\varphi-R_j,\quad
 c_j=\begin{cases}3/2&j\text{ odd},\\0&j\text{ even}.
 \end{cases}                                         \tag{F47.10}
\]
Its columns are integral, its square is $-1$, and
$\varphi^2=-1/2$. Translation sends $O$ to $G$, $G$ to $P_2$,
and reverses the two long affine $E_7$ arms. These data define
the full integral isometry $T_+$ in the certificate;
the actual source map is $T=T_+^{-b}$.
Set
\[
 L_*=(F-\sum v_jC_j,C_0,\ldots,C_6),\qquad L_0=F,
 \qquad L=(L_*,L_0).
\]
The source exterior complex is
\[
 C_4=\mathbb Z, C_3=\mathbb Z^2, C_2=\mathbb Z^{10},
 \ C_1=\mathbb Z^9, C_0=\mathbb Z,\qquad
 d_4=(1,1)^t,\quad d_2=L^tB.                         \tag{F47.11}
\]
For each boundary the complex is
$C_3=\mathbb Z$, $C_2=C_1=\mathbb Z^{r_i}$, $C_0=\mathbb Z$,
with $d_2=L_i^tBL_i$. Its inclusions are $L_i$, the component
embedding, and the $i$th coordinate in degrees two, one, three.
The actual mapping-torus differential is
\[
 d_{{\rm cone},j}
       =\begin{pmatrix}d_j&T_{j-1}-1\\0&-d_{j-1}\end{pmatrix},
\]
where the action on component coordinates is the inverse
transpose of their permutation, and the degree-three and
degree-four actions are identity. The program exports every
one of these integer matrices, every inclusion, and every
Smith representative and projection. The even free source-pair
argument in P3F.46 verifies that these are the actual derived
maps without an omitted odd homotopy term.

\item The exterior $W$ has $H_1=H_3=\mathbb Z^2$, $H_4=\mathbb Z$,
and $H_2=\mathbb Z^2$ for $b=1$, or $\mathbb Z/2\oplus\mathbb Z^2$
for $b=2$. In the physical degree-one marking,
$\alpha_*\mapsto\alpha$, $\alpha_0\mapsto-\alpha$,
and both gammas map to $\gamma$.
In degree three the corresponding markings are
$N_*\mapsto N$, $N_0\mapsto-N$, and both $F\gamma$ map to $Z$.
The certificate records the exact unimodular changes from its
raw Smith bases to these marked bases.

\item At the starred boundary choose $r_*$ with $F=2r_*$ for the
reduced root, or $F=r_*$ for the good root. At the II boundary
choose $F=r_0$. Actual specialization gives $r_i\mapsto C_i$.
Keep $s_i\mapsto\kappa_i C_i$ without restricting either
integer. In domain order $(r_*,s_*,r_0,s_0)$ the complete
degree-two matrix for $(a,b)=(2,1)$ is
\[
 \begin{pmatrix}
 0&-1&0&0\\0&-2&0&1\\
 -1&-\kappa_*&0&0\\0&0&-1&-\kappa_0
 \end{pmatrix}.                                      \tag{F47.12}
\]
For $(a,b)=(1,2)$ prepend the boundary and exterior torsion
coordinates, both of order two. The full mixed-group map is
\[
 \begin{pmatrix}
 1&1&1&1&0\\0&0&1&0&0\\0&0&1&0&1\\
 0&-1&-\kappa_*&0&0\\0&0&0&-1&-\kappa_0
 \end{pmatrix}.                                      \tag{F47.13}
\]
The first row is read modulo two. The determinants of these
free representatives are $1,-1$, respectively, independently
of the two collar coefficients. Each respects the specified
order-two relation and is an isomorphism of the actual groups.
The original unreduced integer representatives of the torsion
row are also retained in the certificate.

\item At a good starred cap the full boundary elliptic torsion maps
isomorphically to its order-two cap torsion. This is the
degree-one transport through the smooth root. The elliptic
lattice quotients by $D-1$ and $D^u-1$ agree since $u$ is odd.
Keep arbitrary torsion path coefficients $c,d$ in the images
of $\alpha_*,\gamma$. With column order
$(\eta,\alpha_*,\gamma_*,\alpha_0,\gamma_0)$, the full map is
\[
 D_1=\begin{pmatrix}
 0&1&0&-1&0\\0&0&1&0&1\\-1&-c&-d&0&0\\
 0&n_*&-m_*&0&0\\0&0&0&n_0&-m_0
 \end{pmatrix},                                      \tag{F47.14}
\]
whose third target coordinate has order two. Eliminating that
exact torsion isomorphism leaves, for both source pairs,
\[
 D_1^{\rm free}=D_3=
 \begin{pmatrix}
 1&0&-1&0\\0&1&0&1\\n_*&-m_*&0&0\\0&0&n_0&-m_0
 \end{pmatrix}.                                      \tag{F47.15}
\]
For degree three the cap row is $(-n_i,m_i)$ on $(N_i,F\gamma_i)$.
The actual invariant central component trace is primitive for
the reduced roots; smooth-root fibre transport has degree one
for the good roots. These facts prove both entries, with no
source three-class left over in a cap that retracts to its
two-dimensional central source and suspension circle.
The two determinants in (F47.15) are $-D$.
Finally $D_4=(1,-1)$ and
$D_0=\left(\begin{smallmatrix}1&1\\-1&0\\0&-1\end{smallmatrix}\right)$.
These give every nonzero gluing degree.

\par\medskip\noindent\textbf{The full group and the sphere conclusion.}\enspace
Let $v$ be the retained $I_1$ vanishing cycle and let $A$ be
the $III^*$ monodromy on the elliptic lattice. Then
$A^2=-I$, $\operatorname{tr}A=0$. The product relation with
the II monodromy, whose trace is one, gives
$|\langle v,Av\rangle|=1$: substitute the transvection
$T_v=I+v\otimes\langle v,\cdot\rangle$ into
$\operatorname{tr}(A^{-1}T_v^{-1})$.
Thus the vanishing relation and its conjugate by the actual
starred base loop kill an integral elliptic basis in the source
exterior. The source section then identifies its remaining
fundamental group with the base annulus group $\mathbb Z$.
The suspension adds the commuting gamma generator, giving
$\pi_1(W)=\mathbb Z\alpha\oplus\mathbb Z\gamma$.

\item At each good cap the same boundary lattice maps onto the full
elliptic normal subgroup, so it dies in the gluing. Reduced
caps have simply connected central source support. The exact
remaining presentation is
\[
 \langle h_*,h_0\mid h_*^{m_*}=h_0^{m_0},\quad
                         h_*^{n_*}=h_0^{-n_0}\rangle .
\]
Choose $n_*u+m_*v=1$. Substitution gives
$h_*=h_0^{-n_0u+m_0v}$, so the group is cyclic.
Its primitive relation and determinant give precisely
$\mathbb Z/|D|$. The same substitution proves that the first
three Smith factors of (F47.15) are one, also when $D=0$.
Together with the isomorphisms (F47.12)--(F47.13) and the
degree-zero and degree-four maps, the full integral sequence
gives the groups in the statement.
The symbolic identities and 1681 further exact Smith cases
are executed in \path{code/part3_root/iii_star_ii_joint.py}.

\item For $D=\pm1$ the orbit is simply connected with the integral
homology of $S^5$. Its standard diffeomorphism type, the
primitive fixed-sphere Hopf link, and the Gysin calculation of
P3K.22 give the actual standard oriented $S^6$ total space.
For other $D$ the displayed fundamental group excludes a
sphere. The finite common flop chain of P3F.35 carries this
topology criterion through every allowed cusp fan. All line
and annular maps used to obtain it remain the maps in
(F47.6)--(F47.9).

\end{enumerate}
\end{proof}




\begin{theorem}\label{night:P3F.48}
\StatementRecord{night:P3F.48}{proved}{S6\_06\_BLIND: GPT-6 Astra (OpenAI, ChatGPT Pro mode, effort Pro), recorded re-derivation in the checkpoint dated 18 November 2026}{\cref{night:P3F.35}, \cref{night:P3F.40}, \cref{night:P3K.19}}
For the actual additive cores constructed in
P3F.36, P3F.38, and P3F.41--43, P3F.46--47,
add any finite collection of logarithmic caps at smooth or
multiplicative source values, retaining their marked elliptic,
gamma, and delta collar data. Every such added cap in a sphere
member has effective order one. This holds for every allowed
cusp fan. The statement retains all order-one clutches and
assigns no parameter-return identity to them.
\end{theorem}
\begin{proof}
\begin{enumerate}[label=\arabic*.]
\item \par\medskip\noindent\textbf{The actual source exterior.}
Let $W_0$ be the original negative-translation orbit source with
only its prescribed additive-core tubes removed. It is the
actual mapping torus of $t_{-bG}$ on that source complement,
where $b\in\{1,2,3,6\}$. Write $q=1$ or two for the number
of removed additive tubes. The complete source-cone calculations
in the cited statements give
\[
 H_2(W_0;\mathbb Z)=\mathbb Z/b\oplus\mathbb Z^q,
 \qquad H_3(W_0;\mathbb Z)=\mathbb Z^q.                         \tag{F48.1}
\]
The trivial summand is omitted for $b=1$.
An integral $H_3$ basis consists of the actual common generic-fibre
trace $F\gamma$ and $q-1$ source-boundary classes $N_i$ at
the old removed tubes. For $q=2$ their relation is
$N_1+N_2=0$. These are the marked degree-three representatives
in the actual source-cone certificates.

\item The source involution
\[
                         I[p,t]=[-p,-t]
\]
reverses the five-dimensional orientation, fixes the base, and
preserves every chosen tube. We verify its action on degree two
using the actual height decomposition. The rational source
lattice has the marked classes $O,F$, the root components, and
the Shioda class $\phi$. Inversion fixes $O,F$, sends
$\phi$ to $-\phi$, and takes every component to its inverse
component label. All retained multiplicative source fibres have
transitive $t_{-bG}$ component action. Their component
coinvariants therefore span only the whole-fibre direction.
The removed additive components have already been imposed by
the complement intersection map. Translation takes $\phi$ to
$\phi$ plus a nonzero rational multiple of $F$, so $F$ dies
in the rational degree-two coinvariant. Its remaining direction
is $\phi$, with negative inversion. The other $q-1$ directions
are traces of the base-annulus $H_1$ classes; source inversion
fixes those classes and time reversal negates their traces.
Thus $I=-1$ on $H_2(W_0;\mathbb Q)$.

\item For a prime $p\nmid b$, put $R=\mathbb Z_{(p)}$.
Equation (F48.1) makes $H_2(W_0;R)$ free. It injects into
its rationalization, so the actual action on it is exactly
$-1$. This argument does not divide by two. The full source
inversions are also retained in the $I_1^*$ certificates and
in the four cases of
\path{code/part3/iv_star_iii_joint.py}.

\par\medskip\noindent\textbf{The exact ordinary comparison.}
Remove all the added ordinary smooth and multiplicative tubes
from $W_0$, and call the resulting exterior $W$.
Let $U_j$ be these actual ordinary tubes and $E_j=\partial U_j$.
We first use their original ordinary fillings. Every
$H_2(U_j;R)$ and $H_2(E_j;R)$ is free. At a multiplicative
site its source Euler pair $(N,c)$ is primitive, since the
actual translation is transitive on components.
The ordinary cap maps on $H_2$ are surjective. Their primitive
kernels are the three smooth rims
\[
                    a\alpha,\quad b\alpha,\quad\alpha\gamma
\]
or the two actual multiplicative classes $R_j,L_j$ of P3F.40.
Denote the direct sum of these ordinary kernels by $K$.

\item The map $H_3(W;R)\to H_3(W_0;R)$ is onto by actual
representatives. Each old $N_i$ lies in its old additive
boundary collar, disjoint from every new tube. The common
$F\gamma$ trace can be placed over a regular source value
outside all new tubes. These representatives are the entire
integral basis in (F48.1). Hence the ordinary Mayer--Vietoris
connecting map from $H_3(W_0)$ to the new boundary $H_2$
is zero.

\item The next boundary map from $H_2(W_0;R)$ to
$\bigoplus H_1(E_j;R)$ is also zero. Its image lies in the
ordinary cap kernels, the primitive meridian groups
$\bigoplus R\alpha_j$. They are $I$-even, whereas its source
is $I$-odd. Equivariance makes twice each image zero.
The meridian group is torsion-free, also for $p=2$, so the
image itself is zero. The exact ordinary sequence is therefore
\[
 0\longrightarrow\bigoplus_jH_2(E_j;R)
 \longrightarrow H_2(W;R)\oplus\bigoplus_jH_2(U_j;R)
 \longrightarrow H_2(W_0;R)\longrightarrow0.
\]
Using the actual surjective cap projections gives exactly
\[
             0\longrightarrow K\longrightarrow H_2(W;R)
                  \longrightarrow H_2(W_0;R)\longrightarrow0.  \tag{F48.2}
\]
Indeed, subtract boundary lifts of all cap coordinates from an
ordinary preimage to obtain surjectivity of the last map.
A boundary preimage with zero cap coordinates lies in the
ordinary kernels, proving the middle kernel assertion.
The vanishing $H_3$ connecting map proves their injectivity.
This retains every rim and every integral index.

\item Let $g_-(q)$ be the Lie lift of the actual source translation
$-bG(q)$. With convention
$(p,1)\sim(t_{g_-}p,0)$, the exact local product is
\[
        (a,b,t)\longmapsto[a+b\tau(q)-t g_-(q),t].              \tag{F48.3}
\]
At $t=1$ the equality
$t_{g_-}(a+b\tau-g_-)=a+b\tau$ proves its closing.
In these coordinates $I$ is precisely $(-a,-b,-t)$ and
$\alpha$ is fixed. Thus the ordinary smooth kernels are
$I$-odd, while $ab,a\gamma,b\gamma$ are $I$-even.
On a Tate sector, write the continuation of $g_-$ as
$g_-+c_-$. The continuation is
$a\mapsto a+Nb-c_-t$, and commutes with the same minus
identity. The physical coefficient of P3F.40 is consequently
$c=c_-$, namely minus the component integer of $+bG$.
Its exact classes $z,B$ are even and $R,L$ are odd.
No involution has been extended over any replacement cap.

\item Both ends of (F48.2) are free and odd. Choose an $R$-linear
splitting. The involution has block matrix
\[
                           \begin{pmatrix}-I&C\\0&-I\end{pmatrix}.
\]
Its square is one, so $2C=0$. The coefficient ring has no
two-torsion, hence $C=0$. Thus all $H_2(W;R)$ is free and
odd. Every even new boundary core column maps to zero in it.
This proves the needed vanishing at $p=2$ without an
eigenprojection or an inverse of two.

\par\medskip\noindent\textbf{Eliminating the actual additive-core blocks.}
The prescribed additive-core degree-two Mayer--Vietoris map
is an integral isomorphism in each cited construction, for
all its actual collar coefficients. Project the new target
$H_2(W;R)$ plus the old cap groups onto $H_2(W_0;R)$ plus
those cap groups. On the old boundary summand its projection
is exactly that old isomorphism, by the ordinary inclusion
maps. Its image is therefore a graph over the whole projected
module and meets $K$ trivially. Eliminating that graph leaves
$K$ and the new replacement-cap targets. This is an exact
module quotient. It requires neither an involution on the old
caps nor a nonzero old Seifert determinant.

\item At a smooth replacement retain
\[
 \mu=m\alpha+n\gamma+\lambda_a a+\lambda_b b,
                    \qquad\gcd(m,n)=1.
\]
The cap map is the exterior square of the exact primitive
quotient by $\mu$, whose kernel is $\mu\wedge H_1(E)$.
Since the core columns vanish in $W$, its ordinary rim block is
\[
 \begin{pmatrix}
 m&0&0&-\lambda_a\\0&m&0&-\lambda_b\\0&0&m&n
 \end{pmatrix}.                                                   \tag{F48.4}
\]
Its columns are the rim projections of
$-\mu\wedge a,-\mu\wedge b,\mu\wedge\gamma,-\mu\wedge\alpha$.
For $p\mid m$, the integer $n$ is a unit in $R$.
The vector $(-\lambda_a,-\lambda_b,n)$ can therefore be
sent to $(0,0,1)$ by the actual invertible row map
\[
 \begin{pmatrix}1&0&\lambda_a/n\\0&1&\lambda_b/n\\0&0&1/n\end{pmatrix}.
\]
It preserves $mR^3$ as a submodule. The quotient of (F48.4)
is exactly $(R/mR)^2$, with both elliptic coefficients retained.

\item At a multiplicative replacement use the exact basis
$(z,R,B,L)$ with $z\cdot L=R\cdot B=1$ and zero
self-pairings. The exterior $z,B$ columns vanish and $R,L$
are its independent ordinary kernel. With the retained Tate
integer $j$, set $A=Nj-cn$. The actual cap rows are
\[
                 \begin{pmatrix}m&-A&s&t\\0&0&m&A\end{pmatrix}.
\]
Its isotropic-kernel equation is $mt=As$ and its actual
component transitivity gives $\gcd(m,A)=1$. Hence $s=mk$
and $t=Ak$. The actual target change $T_{\rm new}=T+kC$
gives the cap rows
\[
                 \begin{pmatrix}m&-A&0&0\\0&0&m&A\end{pmatrix}. \tag{F48.5}
\]
These are precisely the local maps of P3F.40. Their local
hypotheses are the primitive source Euler pair and
transitivity; the original global application is replaced
here by the exact zero core columns just proved.
Together with the exterior $R,L$ rows, permuting columns
to $(R,L,z,B)$ and subtracting their multiples from the last
two rows gives
$\operatorname{diag}(1,1,-m,-m)$.
The quotient is again $(R/mR)^2$.

\item The site kernels in (F48.2) form a direct sum, so all these
blocks occupy independent summands after the old additive
graphs are eliminated. A prime divisor of any new effective
order therefore contributes nonzero degree-two cokernel.
The complete Mayer--Vietoris sequence injects that cokernel
into $H_2$ of the closed orbit, excluding a sphere orbit.

\par\medskip\noindent\textbf{The localization applies to every sphere member.}
Keep the actual marked base meridians and gamma generator,
and quotient the orbit group by its normal elliptic generators.
Its abelianized marked Seifert relations include
\[
 m_i\alpha_i+n_i\gamma=0,\qquad
                  \sum_i\alpha_i+b_0\gamma=0,                    \tag{F48.6}
\]
where $b_0$ is the retained integer base framing.
The elliptic and delta path terms disappear under this group
quotient, not under a change of gamma marking.
If a prime divides two effective orders, their primitive norm
relations force $\gamma=0$ modulo that prime. Other
meridians are then determined, while the two indicated
meridians have only their one sum relation. The quotient has
nonzero $H_1$ over the residue field. Thus sphere members have
pairwise coprime effective orders. The $b_0$ term remains in
the argument and vanishes only after $\gamma=0$.

\item Every prime divisor of $b$ divides a required old good or
starred order in the actual weighted cores. For $b=2$ there
is a quadratic $I_1^*$ cap or a good $III/III^*$ order
divisible by four. For $b=3$ there is a good $IV/IV^*$ order
divisible by three. For $b=6$ both primes occur: in
$I_1^*$ plus good $IV$, or good $IV^*$ plus good $III$.
There is no restriction when $b=1$.
Consequently a prime divisor of a new smooth or multiplicative
order must be prime to $b$ in a sphere member. It is a prime
to which the preceding localized proof applies, giving a
contradiction. Every added effective order is therefore one.

\item The reduction to effective primitive pairs is the actual
finite-kernel quotient map of P3K.19. No ineffective factor
has merely been discarded. The finite common flop chain of
P3F.35 is the identity off the selected cusp and preserves
all additional sites, so the exclusion holds for every fan.

\item The certificate \path{code/part3/weighted_extra_sites.py}
exports the exact localized row and column maps, the
involution-extension identity, seven actual source-inversion
records, and the required-prime implications. Its saved
output also verifies $924$ supplementary prime cases,
including $p=2$. The ordinary extension (F48.2) is the
geometric proof supplying those algebraic blocks.

\end{enumerate}
\end{proof}




\begin{theorem}\label{night:P3F.49}
\StatementRecord{night:P3F.49}{proved}{S6\_06\_BLIND: GPT-6 Astra (OpenAI, ChatGPT Pro mode, effort Pro), recorded re-derivation in the checkpoint dated 18 November 2026}{\cref{night:P3F.36}, \cref{night:P3F.38}, \cref{night:P3K.15}}
At the quadratic $I_1^*$ root and the good or reduced IV roots of
P3F.36 and P3F.38, every extending affine source norm solution
with the prescribed canonical deck and the permitted transitive
central component action is conjugate to the constructed solution
by an extending translation commuting with the suspension.
For the actual multidegree-zero root line, this conjugacy lifts
to an invertible line map. After that lift, the remaining norm
solution differs by a retained finite delta character.
The actual outer collar coefficients remain in the comparison.
\end{theorem}
\begin{proof}
\begin{enumerate}[label=\arabic*.]
\item \par\medskip\noindent\textbf{The translation equation.}
Let $D$ be the canonical order-$m$ root deck. Write the constructed
and the given source maps as
\[
 A_0=t_{R_0}D,\qquad A=t_R D,\qquad
 A^m=A_0^m=t_{-nQ}.
\]
Here $R,R_0$ are holomorphic sections of the root Neron smooth
group. Its group law is written additively, except in the explicit
multiplicative central calculation below. The action of $D$ on
sections is the one defined by $Dt_TD^{-1}=t_{DT}$.
Then $S=R-R_0$ satisfies
\[
 \operatorname{Nm}_D(S)=\sum_{j=0}^{m-1}D^jS=0.               \tag{F49.1}
\]
Conjugation by the actual translation $t_T$ changes $R_0$ to
$R_0+(1-D)T$. We solve $(1-D)T=S$ on the whole root disc.

\item For a good IV root, the central group is the smooth elliptic
curve $E_0$, and $D$ acts with order three. The actual isogeny
$1-D:E_0\to E_0$ has degree three and is onto. Choose a central
solution $T_0$ and extend it to a local holomorphic section by
the smooth submersion coordinate of the root family.

\item For a reduced IV root, $D$ fixes the component group
$\mathbb Z/3$ and $m\equiv1\pmod3$. Equation (F49.1) therefore
puts $S_0$ in the additive identity component. If $\lambda$ is
the scalar action of $D$ on that component and $\lambda\ne1$,
take $T_0=S_0/(1-\lambda)$. If $\lambda=1$, the central norm
equation is $mS_0=0$, so $S_0=0$ and we take $T_0=0$.
Again the smooth group submersion extends this central section.

\item For the quadratic starred root, $m=2$. Both allowed source
maps exchange the two $I_2$ components, so both translation
components are odd. Their difference lies in the multiplicative
identity component $\mathbb C^*$. On it $D$ is inversion, so
the actual homomorphism $1-D$ is
\[
                     T\longmapsto T^2.
\]
Choose $T_0$ with $T_0^2=S_0$ and extend it in the identity
component. This calculation retains the transitive component
class required by the cap construction.

\item Subtracting $(1-D)T_0$ leaves a norm-zero section whose central
value is the identity. Shrink the root disc so that it lies in
the analytic formal-group neighbourhood of the identity. Its
fibrewise logarithm is a section $s$ of the Lie line. The action
of $D$ on this line is linear over its finite base rotation.
The norm of $s$ is lattice valued and vanishes at the centre;
after the same shrinking it is zero. The explicit finite sum
\[
 t=-\frac1m\sum_{j=0}^{m-1}jD^js,
 \qquad
 (1-D)t=s-\frac1m\sum_{j=0}^{m-1}D^js=s                    \tag{F49.2}
\]
solves the residual equation. Exponentiate and add this
translation to the first one. We have obtained the actual maps
\[
             t_TA_0t_T^{-1}=A,
 \qquad t_Tt_{-Q}=t_{-Q}t_T.                               \tag{F49.3}
\]
The translation extends on the minimal root surface. On a good
root it is the smooth elliptic group action. On the $I_2$ root
it is multiplication by the corresponding Neron unit point.
On the reduced IV root it is the extending section translation
in the resolved charts used in P3F.38. Thus (F49.3) is an
identity of the actual extending source maps.

\item The quadratic case also has a literal Tate-unit formula. With
the prescribed boundary marking, write
\[
 A(Z,z)=(R(z)/Z,-z),\quad R(z)=z w(z),\qquad
 R_0(z)=z w_0(z).
\]
Equality of the two actual norms makes $S=w/w_0$ an even
nowhere-zero unit. Choose its even analytic square root $U$.
Then $C_U(Z,z)=(U(z)Z,z)$ gives
\[
 C_UA_0C_U^{-1}=A,
 \qquad U(z)U(-z)R_0(z)=R(z).                              \tag{F49.3a}
\]
This map can move the four canonical inversion centres.
We transport the Cartier comparison by the literal blowup
isomorphism
$\operatorname{Bl}_{Z_0}S'\to
\operatorname{Bl}_{C_U(Z_0)}S'$ induced by $C_U$; the new
quotient map is the old quotient map composed with its inverse.
Thus the exceptional divisor comparison is transported, without
requiring $C_U$ to fix its original four centres.

\par\medskip\noindent\textbf{The line lift and its finite phase.}
Let $N$ be the actual root line, including the Cartier corrections
proved in P3F.36 and P3F.38. It has multidegree zero. Translation
by $T$ preserves its generic fibrewise Picard class. Apply the
proper-direct-image and vertical-degree calculation of P3K.15
to the ratio of $N$ and its translate. Evaluation gives a base
line, and the remaining divisor is an integral whole fibre.
Its local base equation supplies an actual invertible lift
$L_T$; we retain that base factor and its outer phase.

\item These degree-zero translation lifts commute with the specified
suspension $B$. On every smooth fibre, write the punctured line
as $(\mathbb C\times\mathbb C^*)/\Lambda$ with constant
multiplicative period characters. A translation lift acts by
a constant translation in the first factor and a constant scalar
in the second. The two such maps commute. Every other lift
differs by a base scalar, so this proves $[L_T,B]=1$ on the
generic fibres and, by extension, on the root neighbourhood.

\item The conjugated canonical lift $L_TA_0L_T^{-1}$ and the given
lift of $A$ now cover the same source map. Their ratio is a
holomorphic base unit. More precisely, if
$C=L_TA_0L_T^{-1}$, define $f$ by $A=m_fC$, where $m_f$
multiplies the vector in its \emph{target} fibre by $f$.
Thus conjugation by a base scalar $b$ has target ratio
$b(z)/b(\zeta_m^{-1}z)$, in agreement with the action of $D$
used in (F49.2). Their common $m$th power $B^n$
gives the exact norm equation
\[
               \prod_{j=0}^{m-1}f(\zeta_m^jz)=1.             \tag{F49.4}
\]
Put $\chi=f(0)$, so $\chi^m=1$. The unit $f/\chi$ has an
analytic logarithm $g$ with $g(0)=0$, and its norm is zero.
Formula (F49.2), applied to this scalar logarithm with the same
conjugation action, removes $\exp(g)$ by a nowhere-zero base
scalar $b$, explicitly
$A=\chi\,m_b C m_b^{-1}$. For $m=2$, $g$ is odd and
$b=\exp(g/2)$ has the required target ratio
$b(z)/b(-z)=\exp(g)$. The remaining constant $\chi$ is retained as the finite
delta character of the cap.

\item Consequently the line comparison carries the exact norm maps,
the finite character, and all outer source and line markings.
The homology calculations of P3F.36 and P3F.38 already retain
their actual coefficients $c,\kappa$; no choice of a numerical
outer coefficient is needed here. The displayed conjugacy is
local at each cap and makes no change between distinct source
sites. The exact averaging and central squaring identities are
recorded in \path{code/part3/os50_affine_conjugacy.py} and its
executed JSON output.

\end{enumerate}
\end{proof}




\begin{theorem}\label{night:P3F.50}
\StatementRecord{night:P3F.50}{proved}{S6\_06\_BLIND: GPT-6 Astra (OpenAI, ChatGPT Pro mode, effort Pro), recorded re-derivation in the checkpoint dated 18 November 2026}{\cref{night:P3K.23}}

Every sphere family obtained from a rational elliptic source with
Mordell--Weil group \(\mathbb Z\) in the uniformly expanding case
of E26 Remark~8.5 has the following exact inverse-action comparison.
The parameter map \(\mu=\lambda^{-1}\), the identity on the original
line complement, the unchanged root-cap groups and the identity on
the same physical Mumford charts give a biholomorphic contracting
family. Its cap norms and base framing are
\(n_i^\vee=-n_i\) and \(b_0^\vee=-b_0\).
For a counterclockwise large-\(\lambda\) scalar circle, its return is
the inverse of the positive small-\(\mu\) return:
\[
 [\Phi_{\lambda,+}]=-[\Phi_{\mu,+}]\quad\hbox{in }\Theta_7.
\]
On the K23 reference products the transported individual characteristic
pair and relative integers are negated as well.
The physical fan and line circle are retained.
The same statement holds on each specified root-parameter cover;
no additional shorter-loop descent is asserted.
\end{theorem}
\begin{proof}
\begin{enumerate}[label=\arabic*.]
\item The source data are exactly the fixed marked data of E26 Remark 8.5
on a rational elliptic surface with Mordell--Weil group \(\mathbb Z\).
This statement supplies the comparison for the large-scalar expanding
case as well as the small-scalar contracting case. Every parameter
circle below is specified in its own coordinate.

\item Let \(\pi:S\to\mathbb P^1\) be the marked source, let \(M\) be its
actual line, and put \(S^\circ=S\setminus\pi^{-1}E\). Write
\(B_\lambda=\lambda B_0\) for the expanding automorphism of
\(M^{\times}|_{S^\circ}\), where \(B_0\) covers a translation \(t\).
The zero/pole assertion in the remark means that its actual
meromorphic line map has poles along the whole inverse image of
the effective base divisor \(E\). Define
\[
 \mu=\lambda^{-1},\qquad
 B^\vee_\mu=B_\lambda^{-1}=\mu B_0^{-1}.
 \tag{EX.1}
\]
The inverse covers \(t^{-1}\). Since \(t\) is over the source base,
the divisor of \(B_0^{-1}\) is the whole zero divisor \(\pi^{-1}E\).
It is the contracting case of the same construction. The inverse
formula is an identity of holomorphic automorphisms on the indicated
line complement, and a meromorphic line-map identity on the source.

\par

\medskip\noindent\textbf{1. The family quotient and every finite cap}

\par\nobreak

\item For every permitted scalar and every point of the line complement,
\[
 \langle B_\lambda\rangle=\langle B^\vee_{\lambda^{-1}}\rangle.
 \tag{EX.2}
\]
Consequently
\[
 (\lambda,[v])\longmapsto(\lambda^{-1},[v])
 \tag{EX.3}
\]
is a biholomorphism from the expanding quotient family over
\(\{|\lambda|>R\}\) to the pullback of the contracting quotient
family over \(\{0<|\mu|<R^{-1}\}\). Its inverse has the identical
coordinate formula. It preserves the actual associated line
circle and every marked source base point.

\item Retain an actual cap on the source root \(q=z^m\), with commuting
line maps \(A_\lambda,B_\lambda\) satisfying
\[
 A_\lambda^m=B_\lambda^n,\qquad \gcd(m,n)=1.
\]
Set \(A^\vee_\mu=A_{\mu^{-1}}\), \(n^\vee=-n\). Then, as literal
maps of the same root line,
\[
 (A^\vee_\mu)^m=(B^\vee_\mu)^{n^\vee},\qquad
 \langle A^\vee_\mu,B^\vee_\mu\rangle
       =\langle A_{\mu^{-1}},B_{\mu^{-1}}\rangle.
 \tag{EX.4}
\]
The root source, its resolved components, its line, all finite
characters, and its divisor comparison diagrams are unchanged by
this identity. It gives inverse holomorphic maps on the cap
quotients and the same map on their outer collars. In particular
it retains their full nonreduced fibre structure and their
integral boundary maps.

\item In the actual outward-period convention
\(\gamma_{\rm out}=B_\lambda^{-1}\), the new outward period is
\[
 \gamma_{\rm out}^\vee=(B^\vee_\mu)^{-1}
             =B_\lambda=-\gamma_{\rm out}
\]
in the additive integral period notation. The source meridian
\(\alpha\) and the positive line circle \(\delta\) are retained.
Thus the marked killed meridian has the exact expression
\[
 m\alpha+n\gamma_{\rm out}
       =m\alpha+n^\vee\gamma_{\rm out}^\vee.
 \tag{EX.5}
\]
The identity also retains the signed norm at an order-one site.
For a fixed base framing
\(\sum\alpha_i+b_0\gamma_{\rm out}=0\), its new coefficient is
\(b_0^\vee=-b_0\). Hence each effective signed determinant changes
sign. Its absolute value, every effective order, and every
primitivity condition are preserved by the displayed maps.

\par

\medskip\noindent\textbf{2. The exact extension across the Mumford charts}

\par\nobreak

\item Retain the actual root Tate coordinates and their physical
toric cocharacter lattice \(N\), including the line cocharacter
\(\delta\). Let \(H\) be the source Tate deck map. The period
action on the punctured toric chart is generated by the actual
commuting maps \(H\) and \(B_\lambda\). Its lattice of deck
maps is denoted \(\Gamma\). Inverting the second generator gives
the exact equality of subgroups of holomorphic automorphisms
\[
 \langle H,B_\lambda\rangle
       =\langle H,B_\lambda^{-1}\rangle.
 \tag{EX.6}
\]
This equality retains their bounded Tate and line units.

\item The relation to the fan lattice is the actual valuation
homomorphism \(\nu:\Gamma\to N\): take the integral powers
of the source root coordinate in each of the two torus
coordinates of a deck map. In particular
\(\nu(B_\lambda^{-1})=-\nu(B_\lambda)\). If \(V\) is its
integer matrix in the ordered deck generators \((H,B_\lambda)\),
then the matrix in \((H,B_\lambda^{-1})\) is
\[
 V^\vee=V
       \begin{pmatrix}1&0\\0&-1\end{pmatrix},
 \qquad {\rm im}(V^\vee)={\rm im}(V)\subset N.
 \tag{EX.7}
\]
This is the exact period-to-cocharacter map being used.
The physical cocharacter lattice and the period image are
retained.

\item Use precisely the same periodic fan \(\Sigma\subset N_{\mathbb R}\)
and the same cone-height lattice \(\widetilde N=N\oplus\mathbb Z\).
Periodicity under \(\nu(B_\lambda)\) includes periodicity under
its negative. The extending action of \(B_\lambda^{-1}\) on the
toric chart system is the inverse of the extending action of
\(B_\lambda\), with all its original units. On every cone
\(\sigma\), the extending chart map is the identity
\[
 \mathbb C[\sigma^\vee\cap\widetilde N^\vee]
 \longrightarrow
 \mathbb C[\sigma^\vee\cap\widetilde N^\vee],
 \qquad\chi^a\longmapsto\chi^a .
 \tag{EX.8}
\]
It fixes the source monomial, commutes with all face
localizations, and is equivariant for the equal action groups
in (EX.6). It therefore descends on the Mumford quotient to
the extension of (EX.3) across every stratum. All cone
primitivity, local smoothness, properness and period
identifications are the original ones under this map.
The same identity on the common toric charts carries every
finite flop diagram.

\item The resulting map is a biholomorphism preserving the complex
orientation. It retains the physical line circle and every
fan edge. In particular, in an actual fan \(T(v,w)\) the
identity \(\delta=a v+b w\) is the same on both sides.
The integer \(a^2-ab+b^2\) used in the common-exterior
calculation is unchanged. This follows from the displayed
identity of charts and valuation images and does not use
an unproved coordinate identification for the suspension
period.

\par

\medskip\noindent\textbf{3. The specified parameter arrows and return classes}

\par\nobreak

\item Fix \(R\) in the expanding range. Its counterclockwise scalar
circle is
\[
 \lambda(\theta)=R e^{2\pi i\theta},\qquad 0\le\theta\le1.
 \tag{EX.9}
\]
Under (EX.3) it is exactly
\(\mu(\theta)=R^{-1}e^{-2\pi i\theta}\), the reverse of the
positive small-\(\mu\) circle. The quotient and cap identities
give a diffeomorphism of the corresponding mapping-torus
families over this actual base-circle map. Hence their endpoint
returns, with initial fibres identified by (EX.3), are inverse
mapping classes. Under the fixed oriented identification
\(\Theta_7=\mathbb Z/28\), this gives
\[
 [\Phi_{\lambda,+}]=-[\Phi_{\mu,+}].
 \tag{EX.10}
\]
The sign follows from the displayed loop map. A positive
meridian of the end \(\lambda=\infty\), defined in its local
coordinate \(\mu=1/\lambda\), instead gives
\([\Phi_{\mu,+}]\) itself.

\item The individual characteristic numbers and relative group
coordinates use the product framing transported by (EX.3)
and (EX.8). Reversing the parameter cylinder reverses the
mapping-torus circle orientation and the exact relative
return. The characteristic integral and group morphisms
already proved in K23 therefore give
\[
 (k,\ell)_{\lambda,+}=-(k,\ell)_{\mu,+},\qquad
 (m_{\rm rel},n_{\rm rel},s)_{\lambda,+}
       =-(m_{\rm rel},n_{\rm rel},s)_{\mu,+}.
 \tag{EX.11}
\]
These are integral equalities for \(k,\ell,m_{\rm rel},n_{\rm rel}\)
and the stated equality modulo \(28\) for \(s\).
They retain the normal reference frame; no division of a
covered-loop class occurs.

\item For a root parameter \(\lambda=r^L\), use the exact inverse
root parameter \(r^\vee=r^{-1}\), so
\(\mu=(r^\vee)^L\). Equations (EX.4)--(EX.11) hold on this
same cover. For any shorter descended loop whose original
endpoint map has been proved, (EX.3) carries that endpoint
map as well. This last comparison applies to the proved
endpoint map itself and supplies no unproved descent.

\par

\medskip\noindent\textbf{4. Generator signs}

\par\nobreak

\item The generator relabelling has the exact lattice map
\(\mathbb Z\to\mathbb Z,\ j\mapsto-j\). Put
\(G^\vee=-G,\ a^\vee=-a,\ b^\vee=-b\). Then
\[
 a^\vee G^\vee=aG,\qquad b^\vee G^\vee=bG.
 \tag{EX.12}
\]
The actual generic line point and translation are literally
unchanged. Keeping every vertical divisor as well, the
comparison of these relabelled source data is the identity
of the original line and source.

\item There is also an exact map when the target convention fixes
the name \(G\). Let \(\iota:S\to S\) be holomorphic elliptic
inversion, and put \(M^\iota=\iota^*M\). The canonical
holomorphic map of line complements is
\[
 (M^\iota)^\times\longrightarrow M^\times,\qquad
 (p,v)\longmapsto(\iota(p),v),
 \quad v\in M_{\iota(p)}.
 \tag{EX.13}
\]
It is an isomorphism, with the same formula for its inverse.
Pullback sends the generic line point \(aG\) to \(-aG\),
sends every vertical divisor to its actual inverse image,
and carries \(t_Q\) to \(t_{-Q}\). Carry the root sources,
their resolution centres, the line maps and the norm maps
by this same isomorphism.

\item At a Tate node choose the pulled-back line frame. The
actual dense-torus expression of (EX.13) is
\((Z,V)\mapsto(Z^{-1},V)\). Thus its cocharacter map is
\(J=\operatorname{diag}(-1,1)\) in these \emph{physical Tate}
coordinates. Carry the fan by \(J\), and let
\(\widetilde J=J\oplus1\) fix the source height. On each
cone the full contravariant chart map is
\[
 \mathbb C[(\widetilde J\sigma)^\vee\cap\widetilde N^\vee]
 \longrightarrow
 \mathbb C[\sigma^\vee\cap\widetilde N^\vee],
 \qquad\chi^a\longmapsto\chi^{\widetilde J^*a}.
 \tag{EX.14}
\]
These are the extending
toric chart isomorphisms. They carry the period group
and every flop diagram and fix the line cocharacter
\(\delta\). In \(\delta=a_0v+b_0w\), all three vectors
are carried by \(J\), so the coefficients and the fan
integer are retained. This cocharacter map comes from
the displayed actual torus map; (EX.7) separately records
the basis change for the inverse suspension generator.

\item Both comparisons preserve the scalar parameter and the
complex orientation. They introduce no parameter reversal
beyond (EX.9), and every normal frame is carried by the
derivative of the stated isomorphism.

\end{enumerate}
\end{proof}




\begin{theorem}\label{night:P3F.51}
\StatementRecord{night:P3F.51}{proved}{S6\_06\_BLIND: GPT-6 Astra (OpenAI, ChatGPT Pro mode, effort Pro), recorded re-derivation in the checkpoint dated 18 November 2026}{\cref{night:P3F.36}, \cref{night:P3F.37}, \cref{night:P3F.38}, \cref{night:P3F.39}, \cref{night:P3F.40}, \cref{night:P3F.41}, \cref{night:P3F.42}, \cref{night:P3F.43}, \cref{night:P3F.44}, \cref{night:P3F.45}, \cref{night:P3F.46}, \cref{night:P3F.47}, \cref{night:P3F.48}, \cref{night:P3F.49}, \cref{night:P3F.50}, \cref{night:P3K.26}, \cref{night:P3K.27}, \cref{night:P3K.28}, \cref{night:P3K.29}, \cref{night:P3K.30}}

Every fixed marked rational elliptic source with Mordell--Weil group
\(\mathbb Z\) in the surviving source profiles is covered by the
source, line, cap, integral topology and actual scalar-return
comparisons of P3F.36--49 and P3K.26--30. In the two starred
Mordell--Weil entries, the full marked forms are
\[
\begin{aligned}
 E_7:\quad &y^2=x^3+3r x^2+t x+1,\qquad G=(0,1),\\
 E_6+A_1:\quad &y^2=x^3+3r x^2+A t x+t^2,\qquad G=(0,t).
\end{aligned}
\]
The explicit coordinate, finite-root, resolved-quotient and
Cartier-line maps below retain all allowed coefficients.
The ordinary III, IV, II and \(I_1^*\) cap proofs already retain
the full analytic coefficients, and their integral source-pair
maps are the actual maps on each source's marked divisor basis.
Thus the stated sphere criteria, page-line squares and signed
scalar-return integers hold for every such fixed source, with
the specified parameter \(\lambda=r_{\rm par}^{\,M}\).
No transport around a loop varying the source in its moduli
space is inferred from this fixed-source assertion.

\end{theorem}
\begin{proof}
\begin{enumerate}[label=\arabic*.]
\item \par

\medskip\noindent\textbf{S1. Exact marked normal forms for the two starred source entries}

\par\nobreak

\item Let \(S\to\mathbb P^1\) be a relatively minimal rational elliptic
surface with zero section \(O\), and choose a primitive generator
\(G\) of its torsion-free Mordell--Weil group. A minimal short
Weierstrass equation has coefficients in
\(H^0(\mathcal O(4))\), \(H^0(\mathcal O(6))\), and its section
coordinates, when disjoint from \(O\), are polynomials of degrees
at most two and three. These coordinates are the restrictions
of the actual relative Weierstrass functions \(x,y\), whose pole
orders along \(O\) are two and three.

\item For the E7 entry the source height is \(1/2\). The only nonzero
component correction is \(3/2\) at \(III^*\), and the height
identity
\[
 \frac12=2+2(G\cdot O)-\operatorname{cor}_{III^*}(G)
\]
forces \(G\cdot O=0\) and the nonidentity component. Put this
fibre at infinity. The actual short equation is
\[
                 y^2=x^3+(a t+b)x+(c t+d),\qquad a\ne0.
 \tag{S51.1}
\]
Indeed the local \(III^*\) valuations are
\(v_\infty(f)=3\), \(v_\infty(g)\ge5\), so these degree bounds
are exact.

\item Here is the component arrow used to control the marked section,
including its minimal-model map. Put \(q=1/t\),
\(x_0=q^2x,y_0=q^3y\), then make the fourth root
\(q=z^4,x_0=z^6X,y_0=z^9Y\). Its good central fibre is
\(Y^2=X^3+aX\), and its deck is
\((X,Y,z)\mapsto(-X,-iY,iz)\). The two order-four fixed points
are \(O\) and \((0,0)\). At either, the actual tangent weights
are \((3,1)\). The quotient's minimal toric resolution is the
A3 chain in the lattice
\(\mathbb Z^2+\mathbb Z(3,1)/4\). The remaining order-two
orbit contributes its A1 chain. The resolved quotient maps
properly to the original Weierstrass model by the displayed
invariant monomials \(q,x_0,y_0\). Its invariant two-form
\(dX\wedge dz/Y\) descends crepantly; each compact rational
vertical curve has self-intersection minus two by adjunction.
Thus no exceptional minus-one curve remains, and this proper
birational model is the minimal regular elliptic model of
(S51.1). In each A3 chart the successive primitive rays are
\((3,1)/4,(2,2)/4,(1,3)/4\); the divisor of the actual base
function \(q=z^4\) has multiplicities \(1,2,3\) on them.
The order-two chart contributes multiplicity two, and the
central quotient curve has multiplicity four. Thus the two
multiplicity-one leaves are precisely the leaves over the two
order-four fixed points. An equivariant section through either
point has local equation \(w=z^3h(z^4)\), so its strict transform
meets that first ray's leaf, including \(h\equiv0\). The zero
section meets the leaf over \(O\). Hence a section on the
nonidentity component lifts to the finite fixed point \((0,0)\),
not \(O\). The displayed invariant monomials generate the
invariant function field. Consequently the map to the original
Weierstrass model is birational, and the minimal-resolution
comparison is a map of the original source, not a replacement
by a surface with merely the same fibre graph.

\item Consequently the actual root coordinates of \(G\) are
\[
              X_G=z^2x_G(z^{-4}),\qquad
              Y_G=z^3y_G(z^{-4}),
\]
and are holomorphic at zero. Since \(x_G,y_G\) were polynomials
of degrees at most two and three, both must be constant,
say \(r_0,s_0\). The nonconstant coefficient in the section
equation gives \(a r_0+c=0\). Also \(s_0\ne0\), since a
section with \(y\equiv0\) is a nonzero two-torsion section.
Translate \(x\) by \(r_0\), choose a constant scaling making
\(s_0=1\), and use the remaining nonconstant linear coefficient
of \(x\) as the affine base coordinate. These are explicit
holomorphic coordinate maps with nonzero constant Jacobian.
They give
\[
                  y^2=x^3+3r x^2+t x+1,\qquad G=(0,1).
 \tag{S51.2}
\]
Thus (S51.2) captures every marked E7 source, rather than only
the displayed \(r=0\) example.

\item Up to the conventional factor16 the discriminant is
\[
 -4t^3+9r^2t^2+54rt-108r^3-27.
 \tag{S51.3}
\]
Its roots are distinct for \(r^3\ne1\). At \(r^3=1\) there
is a type II fibre at \(t=3r^2\) and a separate \(I_1\).
For example, translating to short form gives
\(f=t-3r^2\), \(g=2r^3-rt+1\), so at that value both vanish
and \(g\) has a simple zero. No additional double nodal
discriminant root is possible: at a double root of the cubic
in \(x\), the derivative with respect to \(t\) is that root
\(x\), which cannot be zero because the cubic's constant
term is one. The certificate also verifies the discriminant
of (S51.3). When \(r^3=1\), constant cube-unit scaling makes
\(r=1\), then \(x_{\rm old}=x+1\),
\(t_{\rm old}=t-3\) gives exactly
\[
             y^2=x_{\rm old}^3+t_{\rm old}x_{\rm old}
                       -t_{\rm old},\qquad G=(1,1).
 \tag{S51.4}
\]
This identifies all three cusp parameters with the E7-B
source by actual coordinate maps.

\item For the E6+A1 entry, \(h(G)=1/6\). The maximum possible
component corrections are \(4/3\) and \(1/2\). Hence
\[
                 \frac16=2+2(G\cdot O)-4/3-1/2
\]
forces \(G\cdot O=0\) and nonidentity components at both
reducible sites. Put \(IV^*\) at infinity. The short
Weierstrass coefficients have degrees at most one and two.
On the minimal cubic good root
\(q=z^3,x_0=z^4X,y_0=z^6Y\), the central deck is
\((X,Y)\mapsto(\zeta_3^2X,Y)\). Its three fixed points are
\(O,(0,\pm\sqrt{g_4})\), with tangent weights \((2,1)\).
Resolving their actual A2 quotient charts and using the
same invariant two-form argument gives the minimal \(IV^*\)
model. The actual rays \((2,1)/3,(1,2)/3\) have base
multiplicities one and two under \(q=z^3\), and the central
quotient curve has multiplicity three. An equivariant section
through a fixed point has equation \(w=z^2h(z^3)\) and meets
its first ray. Hence the three multiplicity-one leaves
correspond exactly to these three fixed points; \(O\) determines
the identity leaf. Thus \(G\) has a finite nonzero three-torsion
limit. The root coordinates are
\[
                 X_G=z^2x_G(z^{-3}),\qquad
                 Y_G=z^3y_G(z^{-3}).
\]
It follows that \(x_G\) is constant and \(y_G\) is linear
with nonzero leading coefficient. Set the new affine base
coordinate equal to \(y_G\), and translate \(x_G\) to zero.
The marked section is now \(G=(0,t)\), and the equation has
the form
\(y^2=x^3+3r x^2+(A t+B)x+t^2\).
The nonidentity A1 section maps to the singular point of
its finite Weierstrass fibre. There \(y_G=0\), so its base
value is \(t=0\); the derivative with respect to \(x\)
then forces \(B=0\). We have proved the full marked form
\[
                   y^2=x^3+3r x^2+A t x+t^2,\qquad G=(0,t).
 \tag{S51.5}
\]
Its discriminant, up to factor16, is
\[
 t^2\{9r^2(A^2-12r)+(54rA-4A^3)t-27t^2\}.
 \tag{S51.6}
\]
The discriminant of the bracket is \(16(A^2-9r)^3\).
Accordingly the retained-\(I_1\) E6+A1 sectors are exactly
\[
 \begin{array}{ll}
 r\ne0,\ A^2\ne9r,12r &: IV^*+I_2+2I_1,\\
 r=0,\ A\ne0 &: IV^*+III+I_1 .
 \end{array}
 \tag{S51.7}
\]
The case \(A^2=9r\ne0\) is \(IV^*+I_2+II\), with no
\(I_1\) whole-zero site. The case \(A^2=12r\ne0\) has
\(I_3\) in place of \(I_2\), so its root rank is eight
and it is outside the Mordell--Weil-\(\mathbb Z\) entry.
The case \(r=A=0\) is also outside that entry.
For the III sector the literal scaling
\(x=s^2X,y=s^3Y,t=s^3T\), with \(s=A\), makes \(A=1\).
Thus every E6-B source is actually the displayed model.
The I2 sector retains its coefficient \(A\); it is not
identified with the \(A=0\) representative by assertion.

\par

\medskip\noindent\textbf{S2. Every starred root map with all coefficients retained}

\par\nobreak

\item For (S51.2), at infinity the exact equation and section are
\[
 y_0^2=x_0^3+3r q^2x_0^2+q^3x_0+q^6,\qquad
                         G=(0,q^3).
\]
The actual substitutions of F41 give
\[
 \begin{array}{ll}
 m=4l,\ k=3l:
 &Y^2=X^3+3r z^{2l}X^2+X+z^{6l},
       \quad G'=(0,z^{3l}),\\
 m=4l+3,\ k=3l+2:
 &Y^2=X^3+3r z^{2l+2}X^2+zX+z^{6l+6},
       \quad G'=(0,z^{3l+3}).
 \end{array}
 \tag{S51.8}
\]
Here \(q=z^m,x_0=z^{2k}X,y_0=z^{3k}Y\).
The good range has \(l\ge1\), and the reduced range \(l\ge0\).
The good central deck is exactly \((-X,-iY)\), and the
nonidentity fixed point is still \((0,0)\). In the reduced
root the A1 tangent cone \(Y^2-zX\) is unchanged; its actual
blowup gives the same two reduced III component labels and
the same marked nonidentity section.

\item For the full E6 form (S51.5), the infinity equation is
\[
 y_0^2=x_0^3+3r q^2x_0^2+Aq^3x_0+q^4,\qquad G=(0,q^2).
\]
The actual root maps give
\[
 \begin{array}{ll}
 m=3l,\ k=2l:
 &Y^2=X^3+3r z^{2l}X^2+A z^lX+1,\quad G'=(0,1),\\
 m=3l+2,\ k=2l+1:
 &Y^2=X^3+3r z^{2l+2}X^2+A z^{l+2}X+z^2,
       \quad G'=(0,z).
 \end{array}
 \tag{S51.9}
\]
The good central deck is \((X,Y)\mapsto(\zeta_3^2X,Y)\).
At a reduced root the blowup \(X=zX_1,Y=zY_1\) gives
the two nonidentity labels \(Y_1=\pm1\), and its deck
fixes them because its \(Y_1\) character is
\(\zeta_m^{-(3k+1)}=\zeta_m^{-2m}=1\).
The additional coefficient terms vanish to the displayed
orders and do not change these resolved charts.

\item The source maps on the good central fibres retain their
precise section limits. For the E7 good root, direct doubling
in (S51.8) gives
\[
 \begin{split}
 X(2G')&=\frac1{4z^{6l}}-3r z^{2l},\\
 Y(2G')&=-\frac1{8z^{9l}}+\frac{3r}{2}z^{-l}-z^{3l},\\
 w(2G')&=-X(2G')/Y(2G')
                  =2z^{3l}+\text{higher terms}.
 \end{split}
 \tag{S51.10}
\]
Thus at the minimal fourth root its germ is exactly
\(z^3 h(z^4)\), with \(h(0)=2\).
For the E6 good root, put \(c=A^2/4-3r\).
Then \(X(2G')=c z^{2l}\) and
\(Y(2G')=-1-(Ac/2)z^{3l}\). Addition to \(G'\) gives
\[
 w(3G')=\frac{12r-A^2}{8}z^{2l}
                         +\text{higher terms}.
 \tag{S51.11}
\]
The leading coefficient is nonzero in the surviving
Mordell--Weil-\(\mathbb Z\) sector. At the minimal cubic
root this is exactly \(z^2 h(z^3)\).
These are actual section calculations, not only assertions
about central component groups.

\par

\medskip\noindent\textbf{S3. The generic coefficient terms preserve the actual line maps}

\par\nobreak

\item For a good starred root the local finite deck may be
linearized while retaining its marked section as follows.
Start with a holomorphic transverse coordinate vanishing
on that section and average it against the required finite
character. Its first derivative is nonzero. Together with
the unchanged root coordinate \(z\), this gives a local
holomorphic coordinate map with invertible Jacobian.
The weights are therefore literally \((r_0-1,1)\), with
\(r_0=4\) or \(3\), in the charts used in F41.
The invariant quotient monomials and the actual root
substitutions in (S51.8)--(S51.9) identify their resolved
quotients with the corresponding marked minimal source.

\item The section Cartier function is consequently still
\[
             h(v_w,v_z)=\max(v_w-(r_0-1)v_z,0).
\]
The common regular root refinements, their maps and their
exceptional pullback coefficients are exactly those of F41:
\[
 \begin{array}{c|c|c}
 r_0&\text{interior root rays}&
       \min(v_w,(r_0-1)v_z)\\ \hline
 4&(3,1),(2,1),(1,1),(1,2),(1,3)&(3,2,1,1,1)\\
 3&(2,1),(1,1),(1,2)&(2,1,1).
 \end{array}
\]
They give the actual equality
\[
                 \beta^*N=\pi^*\mathcal M+D_P-D_O.
 \tag{S51.12}
\]
The translation \(Q=pG'\), with \(p=2\) or \(3\), has
the formal germ just computed. In each fixed-section chart
it is
\[
          w\longmapsto w\,u(w,z)+z^{r_0-1}h_0(z^{r_0}),
 \tag{S51.13}
\]
where \(u\) is a nonzero holomorphic unit. The inverse
translation has the same form. Hence it preserves every
ideal \((w,z^j)\) and \((z,w^j)\), \(1\le j\le r_0-1\).
The normalized blowup of their invariant product is the
specified common refinement, so these literal ideal
equalities lift the actual translation and preserve each
Cartier correction in (S51.12).
For higher good root orders, pull back the entire diagram
through the additional root and resolve the fibre product.
The maps and equality pull back, with their ramification
coefficients retained.

\item At each reduced root the line point in the allowed pair
meets the identity component. The actual resolved
fibre-product map therefore pulls the source line to a
line of degree zero on every exceptional curve that is
contracted to reach the minimal root. The blowup Picard
decomposition descends it through each contraction.
Its difference from \(N=\mathcal O(aG'-O)\) is a vertical
divisor of zero component degrees; the actual reduced
III or IV intersection matrix has integral kernel the
whole fibre. Division by that exact base monomial gives
the required line isomorphism and retains its deck
character. The extra terms in (S51.8)--(S51.9) do not change
any of these component degrees or contraction maps.

\item The source division sections, actual commuting line norms,
strict contraction estimate and annular gluing of F41,
F44, F46 and F47 apply to these full-coefficient roots.
For clarity, this is a direct application to each source,
not a claim that its complex surface is isomorphic to the
special displayed representative. The source action is
lifted on the actual line just compared; its remaining
ratio is a base unit, finite deck averaging retains its
character, and the invariant analytic root gives
\(A_i^{m_i}=B^{n_i}\) and \([A_i,B]=1\) exactly.
The same norm-shell construction gives the smooth proper
caps and the physical delta action. Thus every coefficient
in the marked source forms has been carried through
the actual source, line and quotient maps.

\par

\medskip\noindent\textbf{S4. The other source sectors are already coefficient-generic}

\par\nobreak

\item The ordinary III, IV and II root proofs F42--43 were
written with the full holomorphic unit coefficients:
\[
 \begin{array}{ll}
 III:&y^2=x^3+q\,a(q)x+q^2b(q),\quad a(0)\ne0,\\
 IV:&y^2=x^3+q^2a(q)x+q^2b(q),\quad b(0)\ne0,\\
 II:&y^2=x^3+q\,a(q)x+q\,b(q),\quad b(0)\ne0 .
 \end{array}
 \tag{S51.14}
\]
The displayed root substitutions, finite-deck characters,
blowup centres, quotient maps and contractions in those
proofs use precisely these equations. The marked generator's
central point is fixed by the original component map:
nonidentity for III/IV and identity for II. Finite
linearization retaining the marked sections is the
explicit averaging map in S3. The good-III and good-IV
Cartier differences \(E_P-E_O\), and the cancellation
at the six-point good-II correspondence, are therefore
unchanged as actual divisor identities. Higher section
contact changes the finite point on the exceptional
curve, but every contracted centre is the recorded
infinity point, so the identities remain valid.

\item For the \(I_1^*\) sources the F36/F38/F49 proof is also
coefficient-generic. On the actual minimal quadratic
Tate root it retains
\[
 q_T=z^2u(z^2),\qquad Z_G=zv(z),\qquad
                  v(z)v(-z)=-u(z^2).
 \tag{S51.15}
\]
Here \(u,v\) are arbitrary nonvanishing holomorphic units.
The first equality follows from the analytic inverse
of the Tate \(j\)-series and the even source base; the
second is the marked odd component; the last is the
literal deck equation \(Z_G(-z)=q_T/Z_G(z)\).
Thus the source norm, root Picard character, four-point
blowup quotient and the equality
\(\beta^*N=\pi^*\mathcal M+E_0+E_1\) all retain arbitrary
source coefficients. The translation exchanging the
two indicated exceptional curves is the actual \(Q\)
translation, not a representative-model identification.
F49 then gives the exact affine and line conjugacies,
with the annular map and any moved blowup centres
transported as specified there.

\item At the selected whole-zero \(I_1\), the actual Tate
period and Picard characters are arbitrary nonzero
holomorphic units. The only vanishing scalar is its
specified simple base factor. K26--30 already retain
these units, their logarithms, the exact cusp-origin
difference and its bounded Cox extension. Thus their
cusp, page and normal-frame comparisons do not impose
a hidden special-source coefficient.

\par

\medskip\noindent\textbf{S5. Why the integral topology computations apply to every source}

\par\nobreak

\item No identification of whole source surfaces is required
for this step. The exact Oguiso--Shioda component
markings and primitive height determine the actual
intersection numbers of \(O,F,G\) and the labelled
nonidentity components. The same displayed Gram matrix
in each source proof has determinant minus one. Since
the source is rational, its rank-ten integral
Neron--Severi lattice is unimodular; those actual
classes therefore form a basis on every marked source.

\item The actual translation matrix is determined by its
action on \(O\), \(G\), the fibre and the labelled
components. The section \(2G\), or any required \(jG\),
has its integral vector given by the proved height
formula and root corrections. Thus the programs'
translation matrices are the matrices of the actual
maps on those actual classes, independently of all
Weierstrass coefficients. Component relabelling is
the corresponding explicit permutation of the rows
and columns; no biholomorphism of the whole source
is asserted from that relabelling.

\item The source-pair complexes used in F36--48 are obtained
from these actual component inclusions and the
Poincare--Lefschetz pairing. Their chain maps and
translation cones are functorial for the just
identified actual maps. Thus every differential,
Smith representative, finite torsion column and
retained collar coefficient in those proofs remains
the actual marked map for each source. The geometric
source-complement fundamental-group proofs likewise
use the actual section and the retained singular
fibres: the one-cap argument uses the abelian
elliptic-lattice quotient and its computed zero
abelianization, and the two-cap argument uses the
trace-unit pair of conjugate nodal vanishing cycles.
Neither argument assumes a special monodromy
factorization of a single displayed model.

\item This proves the same sphere criteria for every fixed
source in the surviving profiles. The actual page
maps of K26/K28 use these cap charts and retain their
line correspondences. Their integral relative
period obstruction vanishes by the same source
torsion or conjugate-nodal argument, and their
rigidified Picard comparison gives the same page
line up to the specified base pullback. Consequently
the exact height, Euler and signed-return formulas,
and the matched individual pair of K30, hold on
every such fixed source, with the actual parameter
\(\lambda=r^M\) retained.

\par

\medskip\noindent\textbf{S6. Precise scalar and source-moduli loop scope}

\par\nobreak

\item Remark8.5 specifies an elliptic source, a line bundle
and a translation, then scales the resulting
linearization. The theorem proved here covers every
choice of that fixed source and all the proved cap,
affine, character and annular data. It does not
replace a loop in source coefficients by a scalar
loop. A source-moduli loop can braid the nodal values
or change the marking of \(G\), the cap sites and
the fan lattice. Its return requires transport of
that additional actual marking before it can be
compared with the scalar return. No such comparison
is inferred merely from the coefficient-generic
calculation above.

\item The positive scalar parameter is exactly the one
specified in K26/K28, and no division by the root
degree is made. The shorter simultaneous-Cox scalar
descent remains the separately assigned problem.
Expansion is treated by the exact inverse-action
and parameter-inversion maps of F50, including its
positive-loop reversal; it is not absorbed into a
choice of a positive norm numerator.

\item The executed exact certificate is \path{code/part3/generic_source_models.py}; its JSON and text outputs record all 43 identities.

\end{enumerate}
\end{proof}




\begin{theorem}\label{night:P3F.52}
\StatementRecord{night:P3F.52}{proved}{S6\_06\_BLIND: GPT-6 Astra (OpenAI, ChatGPT Pro mode, effort Pro), recorded re-derivation in the checkpoint dated 18 November 2026}{\cref{night:P3F.36}, \cref{night:P3F.37}, \cref{night:P3F.38}, \cref{night:P3F.39}, \cref{night:P3F.40}, \cref{night:P3F.41}, \cref{night:P3F.42}, \cref{night:P3F.43}, \cref{night:P3F.44}, \cref{night:P3F.45}, \cref{night:P3F.46}, \cref{night:P3F.47}, \cref{night:P3F.48}, \cref{night:P3F.49}, \cref{night:P3F.50}, \cref{night:P3F.51}, \cref{night:P3K.13}, \cref{night:P3K.14}, \cref{night:P3K.15}, \cref{night:P3K.19}, \cref{night:P3K.20}, \cref{night:P3K.21}, \cref{night:P3K.26}, \cref{night:P3K.27}, \cref{night:P3K.28}, \cref{night:P3K.30}, \cref{night:P3K.32}, \cref{night:P3K.5}, \cref{night:P3K.73}, \cref{night:P3K.74}}
\par\noindent\textbf{Completed comparison, 17 November 2026.}
The actual invariant-divisor, root-evaluation and collar morphisms
below identify the constructed families and preserve their actual
returns. P3K.74 supplies the precise return-map conjugacy and
P3K.73 the corrected complete numerical algorithm. The former
appeal to the false P3K.31 recurrence is replaced; its exact prior
text is preserved in the correction archive.
\medskip


Fix one of the whole-fibre sphere source tuples of K20--21 and
F36--51, its marked translating point \(Q=bG\), its generic line
point \(aG\), a selected simple \(I_1\) zero, and its effective
cap data. Every source line satisfying the literal Hom condition
is related to the displayed representative by the actual
meromorphic line map \(\mathcal M=\mathcal M_0\otimes\mathcal O(D)\),
where \(t_Q^*D=D\) as divisors.
Its canonical section, the actual root-line evaluation, and the
explicit cone-height map below identify the resulting family
with the constructed family carrying the retained base factors,
finite characters and integral collar windings.
For the canonical scalar loop, these maps preserve the individual
characteristic pair, the relative index and the actual return
computed by P3K.73 in the transported native and ordinary products.
\end{theorem}
\begin{proof}
\begin{enumerate}[label=\arabic*.]
\item The source \(S\to\mathbb P^1\), generator \(G\), translating point
\(Q=bG\), chosen simple whole-zero \(I_1\) value, and primitive
effective cap data are fixed as in K20--21 and F36--51.
Use the contracting data given by the explicit inverse-action map
of F50. Thus both source Hom lines below have the same positive
simple whole-fibre zero; the expanding parameter is recovered
through that actual inverse map.
Let \(\mathcal M_0\) be the displayed representative source line:
\(\mathcal O(aG-O)\), with the actual correction \(-L_0\) at
the \(I_1^*\) source when specified in F36/F38. Every line
considered has the same generic Picard point \(aG\) and satisfies
the literal whole-base Hom condition. We retain its actual
linearization and all cap line lifts.

\par

\medskip\noindent\textbf{V1. The vertical ratio is invariant as a divisor}

\par\nobreak

\item The restriction of \(\mathcal M\otimes\mathcal M_0^{-1}\) to the
generic elliptic fibre is the trivial line. A nonzero section of
that trivial line over \(\mathbb C(t)\), viewed as a rational
section on \(S\), therefore has an integral vertical divisor \(D\).
It supplies the actual meromorphic line isomorphism
\[
 \mathcal M=\mathcal M_0\otimes\mathcal O_S(D).
 \tag{VL.1}
\]
The chosen source zero and K14's exact height computation give
the same Hom line for both source lines. Their ratio is consequently
\[
 \mathcal O_S(D-t_Q^*D)\simeq\mathcal O_S.
 \tag{VL.2}
\]
Choose a rational function whose divisor is \(D-t_Q^*D\).
Its restriction to every smooth fibre in a nonempty open
base set has no zeros or poles. A meromorphic function on
a compact connected elliptic curve with that property is
constant. Thus this rational function is the pullback of
a rational function on the base, and
\[
 D-t_Q^*D=\pi^*\operatorname{div}(f),\qquad
                       f\in\mathbb C(t)^* .
 \tag{VL.3}
\]
The fibrewise constancy gives this factorization in function
fields, so no assertion about a merely numerical divisor
class is used.

\item Choose a positive integer \(r\) for which \(t_{rQ}\) fixes
every component in the finite support of \(D\). Such an
integer exists because translation is over the base and
permutes the finite component set of each fibre. Pull back
(VL.3) by \(1,t_Q,\ldots,t_{(r-1)Q}\) and add. The left
side telescopes to zero, while the base divisor is fixed:
\[
 0=D-t_{rQ}^*D=r\,\pi^*\operatorname{div}(f).
\]
The group of divisors on the smooth source is free abelian.
The pullback of a nonzero base divisor has a nonzero fibre
coefficient. Therefore \(\operatorname{div}(f)=0\), and
\[
                             t_Q^*D=D
 \tag{VL.4}
\]
as actual divisors.

\item Use the canonical \(t_Q\)-equivariant structure on
\(\mathcal O(D)\) that sends its canonical meromorphic
section \(s_D\) to itself. On the complement of the support
of \(D\), multiplication by \(s_D\) gives an honest
holomorphic line isomorphism
\[
 \mathcal M_0^\times\longrightarrow\mathcal M^\times,
                       \qquad v\longmapsto v\,s_D .
 \tag{VL.5}
\]
The two Hom sections, with the same specified simple zero,
then differ by a nonzero constant: their ratio is a
meromorphic function on the compact source with zero
divisor. Multiplying the scalar parameter by that constant
is the exact parameter map identifying their suspension
lifts in (VL.5). It preserves the direction and degree
of every specified scalar circle.

\par

\medskip\noindent\textbf{V2. The retained fibres and the selected zero}

\par\nobreak

\item At every retained multiplicative fibre the actual source
translation is transitive on its components by K13/K19.
Equation (VL.4) makes the coefficients of \(D\) equal on
that orbit. Hence \(D\) is an integral multiple of the
whole fibre there. At a retained reduced additive fibre,
the same transitivity proof K15 gives the same conclusion:
the reduced components form one orbit, so the equal
coefficients again give the whole fibre. At a smooth or
irreducible fibre this conclusion is immediate.

\item Every nonreduced additive fibre has been placed inside
one of the required cap discs in the complete sphere
tuples. Thus, outside those cap discs, all of \(D\) is a
sum of whole fibres. Their defining base monomials make
(VL.5) a base line-frame change, with its exact integral
winding retained in the collar. At the selected \(I_1\)
zero the same statement applies, since its full reduced
fibre has one component. A local coefficient \(kF\)
gives the actual torus-frame map
\((Z,V,q)\mapsto(Z,q^kV,q)\). On the cone-height lattice
it is
\[
 \widetilde L_k(v,h)=(v+k h\,\delta,h).
 \tag{VL.6}
\]
On the height-one fan it translates every vertex by
\(k\delta\). The whole-\(I_1\), simple-zero period lattice
is the full index-one lattice by K20, so this is an
actual period reindexing. More explicitly, it carries
every cone by \(\widetilde L_k\), with chart pullback
\(\chi^a\mapsto\chi^{\widetilde L_k^*a}\). These maps
commute with all face localizations and with the
actual deck maps. They extend (VL.5) holomorphically
on the carried charts. The differences of fan vertices
and the line cocharacter \(\delta\) are unchanged.
Thus its coefficient pair in the actual edge basis
and the fan integer are retained.

\item This statement does not set \(D\) equal to zero. Its
canonical section gives (VL.5), and its full-fibre
coefficients remain in the base-line transition functions.

\par

\medskip\noindent\textbf{V3. The capped fibres and their actual root lines}

\par\nobreak

\item Use K19's actual quotient to remove a zero-gamma normal
root kernel. This changes the source coordinate and the
line frame by the displayed quotient maps and leaves a
primitive effective cap. It does not leave a hidden
elliptic isogeny or finite line character. On its minimal
root source, K15 proves that its actual line \(N\) has
degree zero on every component of the reduced central
fibre. Its generic Picard point is the pullback \(aG'\).
The corresponding constructed root line \(N_0\) has
the same point and the same zero component degrees.

\item The proper direct-image evaluation for
\(N\otimes N_0^{-1}\), rigidified on a regular fibre,
gives a base line and an effective vertical evaluation
divisor \(E\). The intersection of \(E\) with every
component is zero. The actual connected Kodaira
intersection matrix has integral kernel the whole
fibre multiplicity vector. Its components here are
reduced and their multiplicities have greatest common
divisor one. Consequently the evaluation divisor is
an integral whole-fibre divisor. Division by its
actual base monomial yields the line isomorphism
\[
             N\simeq N_0\otimes f_{\rm root}^*L_{\rm root}.
 \tag{VL.7}
\]
This is the direct-image map of K15, applied to these
two actual root lines; it retains their resolutions
and their marked generic Picard point.

\item Transport the actual affine source action to the
constructed root source by F44 or F49, including the
transported blowup centres in the quadratic case.
Those maps retain the entire outer annulus. Apply
(VL.7) before comparing the line lifts. On each
proper punctured root fibre, the ratio of two line
isomorphisms over the same source map is a
holomorphic nonzero function on that fibre, hence
constant there. It is therefore a base unit.
The conjugated suspension map agrees with the
representative suspension on the punctured root by
V1, hence on the full root by holomorphic extension.
The remaining ratio for the finite-deck lift
therefore has norm one.
The exact scalar averaging in F44/F49 removes its
norm-one base-unit coboundary and retains its finite
deck character.

\item Changing a root frame by \(z^k\) changes the deck
character by \(\zeta_m^k\). Keep both of these
quantities. The canonical source map (VL.5), the
actual root evaluation (VL.7), and their base-unit
ratio therefore give precisely the integer delta
collar and finite-character data already retained
in F44/F49 and K32. In particular the source
vertical coefficients supported inside a removed
cap are processed through its actual root divisor
diagram; they are not discarded from the line.

\par

\medskip\noindent\textbf{V4. Page and return comparison}

\par\nobreak

\item The preceding maps identify the actual ordinary
source pieces and every actual root-cap line with
the constructed data, together with the displayed
base line and retained finite characters. They
commute with the suspension lift after the constant
scalar reparametrization in V1. All maps are fixed
while that scalar varies. Their outer annular
maps are exactly the ones used to glue the family.

\item Pull them to the actual one- or two-cap page of
K26/K28/K32. The remaining line ratio is the base
pullback in K32's evaluated line comparison. On
the compact page its class is \(kF\), and
\[
 c_1(N_{\rm page})F=F^2=0,\qquad
 (c_1(N_{\rm page})+kF)^2=c_1(N_{\rm page})^2 .
 \tag{VL.8}
\]
On the regular binding, the remaining base-line
transition is the one from an extending ordinary
disc, and the fibre degree is zero. The exact
relative boundary-gauge formula of K5/K32 gives
zero square ambiguity, including odd gauges.
The corresponding source isomorphism carries
the stable tangent pairing.

\item At the selected toric zero the monomial map of V2
and its extending base unit carry the entire Cox
product and the three complex curve labels. K27/K30's
explicit normal-coordinate comparison applies to
this transported product. On its normal three-plane
the residual bounded character contributes the
extending constant rotation loop already calculated
there. Thus the same individual pair and relative
coordinates are obtained in the transported fan
and framing. The transported family connection has literally conjugate holonomy
and specified inverse endpoint. The typed conjugacy in P3K.74
therefore preserves its physical sphere class. Each full period
unit, native product and Cox ruling arrow of P3K.73 is carried by
this family isomorphism, so its exact algorithm gives that same
actual value on the transported fan.

\item This supplies the actual comparison for every
allowed vertical source line in the fixed-source
scalar construction. Its source, root line,
parameter, collar and fan maps are (VL.1)--(VL.7);
the characteristic equality follows through the
actual page and Cox comparisons just specified.

\end{enumerate}
\end{proof}




\begin{theorem}\label{night:P3F.54}
\StatementRecord{night:P3F.54}{proved}{S6\_06\_BLIND: GPT-6 Astra (OpenAI, ChatGPT Pro mode, effort Pro), recorded re-derivation in the checkpoint dated 18 November 2026}{\cref{night:P3F.41}, \cref{night:P3F.51}}

For the marked E7 source
\[
 y^2=x^3+3r x^2+t x+1,\qquad G=(0,1),
\]
the source together with its selected finite nodal zero has the exact
parameter
\[
 U=\mathbb C^*\setminus\{u:u^3=-1\},\qquad
 r=\frac{1-2u^3}{3u^2},\quad t_0=u^2-\frac2u .
\]
Fix an actual E7 sphere datum of P3F.41 whose only logarithmic cap
is at \(III^*\), a sufficiently small nonzero scalar, and a marked
index-one unimodular periodic fan. In particular no additional
logarithmic disc meets either member of the coalescing finite-node pair.
Over every sufficiently small disc in \(U\) the full construction
extends as a proper holomorphic submersion. The return on its boundary
is therefore isotopic to the identity, including a collision of two
unselected nodes to a type II fibre.

At a selected-cusp puncture \(u_*^3=-1\), one positive source loop has
the exact Tate and source-Cartier transition
\[
 (Z,V,q_T)\longmapsto(q_T/Z,V,q_T),\qquad
 (v_Z,v_V,h)\longmapsto(-v_Z+h,v_V,h).
\]
Its reflection on edge directions preserves no index-one unimodular
three-edge periodic fan. Thus this specified marked source, line and
annular transport does not close the chosen completion after one turn.
Its twofold cover \(a(u)=u+u^{-2}=\rho^2\) has literally closed source,
line, toric and finite-cap charts and gives the marked holomorphic
family there. This statement assigns no mapping-class return to the
closed \(\rho\)-loop.

\end{theorem}
\begin{proof}
\begin{enumerate}[label=\arabic*.]
\item \par

\medskip\noindent\textbf{SM.1. The complete selected-node cover for the E7 marked form}

\par\nobreak

\item Use the full F51 source
\[
 y^2=P(x,t,r)=x^3+3rx^2+tx+1,\qquad G=(0,1).
\]
A singular finite fibre has a double root \(u\ne0\). The equations
\(P(u,t,r)=P_x(u,t,r)=0\) give the exact maps
\[
 r(u)=\frac{1-2u^3}{3u^2},\qquad
 t_0(u)=u^2-\frac2u,\qquad
 P(x,t_0(u),r(u))=(x-u)^2(x+u^{-2}).
 \tag{SM.1}
\]
Conversely these displayed functions give every marked finite
singular fibre of the normalized E7 source. A double root
determines \(u\), so this is its normalization parameter, with
no unmarked choice of a discriminant branch.

\item The selected fibre is nodal exactly when \(u^3\ne-1\). At
\(u^3=-1\) the displayed polynomial has a triple root, and
the selected fibre is the type II cusp. Thus the explicit
parameter for the marked source together with a selected
whole-zero \(I_1\) is
\[
                         U=\mathbb C^*\setminus\{u:u^3=-1\}.
 \tag{SM.2}
\]
The other finite fibres are allowed to coalesce to the II
fibre of F51 while the selected one stays nodal.
For example \(r=1,u=1/2\) selects the separate nodal
fibre \(t_0=-15/4\), whereas \(u=-1\) selects the cusp.
The maps (SM.1), not an equality of fibre profiles, distinguish
these two situations.

\item Their derivatives are
\[
 r'(u)=-\frac{2(1+u^3)}{3u^3},\qquad
 t_0'(u)=\frac{2(1+u^3)}{u^2}.
 \tag{SM.3}
\]
In particular the selected-node normalization is ramified
over a cusp parameter of the \(r\)-line.

\item The remaining constant marked coordinate equivalence is also
explicit. For \(s^3=1\), the map
\[
 x=s^2X,\quad y=s^3Y=Y,\quad t=s^4T
\]
sends \(r\) to \(r/s^2\) and \(u\) to \(u/s^2\), while
retaining \(G=(0,1)\). Hence its finite quotient on (SM.2)
is described by \(u^3\in\mathbb C^*\setminus\{-1\}\).
No monodromy of the total threefold is inferred from this
finite source-coordinate quotient.

\par

\medskip\noindent\textbf{SM.2. Actual local extension in the permitted source parameters}

\par\nobreak

\item Fix a point of \(U\), its selected zero and a proved E7
datum whose only cap is at \(III^*\). Over a sufficiently small disc in
\(U\), retain the displayed functions \(r(u),t_0(u)\) and
the actual section \(G\). The finite Weierstrass total source
is smooth at a nodal or cuspidal fibre: at a singular point
of the cubic, its derivative with respect to \(t\) is the
nonzero root \(x\). Thus the source family is a submersion
to this small parameter disc. At infinity the full root
equations of F51, their relative quotient charts and their
actual marked section resolutions give the simultaneous
minimal-source map.

\item The cap division section is holomorphic in this parameter
disc. Indeed it is the inverse branch of multiplication by
the cap order at the recorded central point; on a reduced
root both the identity group and component norm are the
ones already checked in F41. Every finite-character
averaging is the same explicit finite sum, and its remaining
unit has a holomorphic logarithm after shrinking the disc.
Its analytic root therefore gives the literal family of
commuting norm lifts. The Cartier correspondence of F51
uses the same invariant ideals in these relative charts.

\item At the selected node the exact holomorphic Morse chart is
the one in SM.3 below; at a point of \(U\) its square-root
branches are fixed on the small disc. The chosen fan and
the physical line character therefore have fixed marked
charts. Their monomial face maps and all cap annular maps
remain holomorphic in the source parameter. A fixed
sufficiently small scalar works uniformly on the compact
closure of this disc, by the strict norm estimate on the
common ordinary complement. The resulting glued family
is proper and is a holomorphic submersion on that disc.

\item Consequently the return around the boundary of such a
disc is isotopic to the identity by the actual smooth
bundle over the disc: choose a smooth horizontal complement
to the vertical tangent bundle and lift radial paths from
its centre. Properness makes these lifts complete on the
closed disc, giving a smooth fibre trivialization there.
Its restriction contracts the boundary return. This
includes a collision of two
unselected nodal fibres to a type II fibre: the source
surface remains smooth, the selected nodal chart persists,
and the starred cap maps remain the same full-coefficient
maps. It does not include a puncture \(u^3=-1\) of (SM.2),
where the selected zero itself ceases to be nodal.
This local extension keeps the finite source fibres unmodified.
The joint \(III^*+II\) cap belongs to the locus \(r^3=1\);
no extension of its II cap across a deformation to two
separate nodal fibres is asserted here.

\par

\medskip\noindent\textbf{SM.3. The exact nodal-coordinate transport at a selected cusp}

\par\nobreak

\item Put \(X=x-u\), \(q=t-t_0(u)\), and
\[
                         a(u)=u+u^{-2}.
\]
In the actual smooth source near its selected singular
point the equation becomes
\[
 y^2=X^2(X+a(u))+q(X+u),\qquad
 q=\frac{y^2-X^2(X+a(u))}{X+u}.
 \tag{SM.4}
\]
For \(a(u)\ne0\), choose a small chart
\(|X|<c|a(u)|\), with \(c<1\), and choose the displayed
square roots. Define
\[
 \xi=\frac{y-X\sqrt{X+a(u)}}{\sqrt{X+u}},\qquad
 \eta=\frac{y+X\sqrt{X+a(u)}}{\sqrt{X+u}}.
 \tag{SM.5}
\]
These are actual holomorphic coordinates near \(X=y=0\):
their Jacobian in \((X,y)\) is
\(-2\sqrt{a(u)}/u\), which is nonzero. Their product is
exactly \(q\), so (SM.5) is a base-preserving node chart.

\item At \(u_*\) with \(u_*^3=-1\), \(a'(u_*)=3\). On a small
positive loop about \(u_*\), the square root of \(X+a(u)\)
changes sign, while that of \(X+u\) can be retained on a
single analytic branch. The exact final chart map is
\[
                              (\xi,\eta)\longmapsto(\eta,\xi).
 \tag{SM.6}
\]
Thus the two normalization branches of the selected nodal
fibre are exchanged. Their ordered Picard gluing map is
reversed: an identification \(L_{p_+}\to L_{p_-}\) is
replaced by its inverse \(L_{p_-}\to L_{p_+}\). This is an
explicit source and line-branch transport, not yet the
clutch of the completed threefold.

\par

\medskip\noindent\textbf{SM.5. Exact Tate, Cartier-line and fan transition}

\par\nobreak

\item The source-line transition in SM.3 can be identified on the
actual Tate model, retaining the line frame. In the selected
node chart let \(\Delta(t,r)\) denote the discriminant
polynomial (S51.3), without its factor 16. The exact identities
are
\[
 f(t_0)=-a(u)^2/3,\qquad
 \Delta_t(t_0)=-4u\,a(u)^3,\qquad
 j=-\frac{64a(u)^3}{u\,q}+O(1).
 \tag{SM.7}
\]
The actual Tate parameter \(q_T\), obtained from the inverse
of its \(j\)-series, is consequently
\[
                q_T=-\frac{u}{64a(u)^3}q+O(q^2).
 \tag{SM.8}
\]
For a sufficiently small fixed loop about \(u_*\), choose the
base collar with \(|q|\) uniformly small relative to \(|a|^3\).
Then \(q_T/q\) is an invertible holomorphic base function
throughout that collar.

\item The Tate Weierstrass model has \(c_4\) a unit with value one.
The source has \(c_4=16a^2+O(q)\). Its actual marked
Weierstrass scaling is therefore \(h=2\sqrt a\) times a
holomorphic unit with a fixed fourth-root branch. On one
positive loop about \(u_*\), that unit has no monodromy and
\(h\) changes sign. The source coordinate isomorphism
\(x=h^2x_T+\text{translation},y=h^3y_T\), in short form,
therefore changes by the marked elliptic inversion.
Since the zero section is fixed, there is no translation
factor. In the multiplicative Tate coordinate this is
literally \(Z\mapsto Z^{-1}\), with \(q_T\) unchanged.
The normalized period has the retained change \(-3\),
as follows also from (SM.8).

\item Here is its line map. Write \(P\) for the actual multiplicative
point of \(aG\). The line \(\mathcal O(aG-O)\) is represented by
\[
       (Z,V)\sim(q_T Z,P V).
 \tag{SM.9}
\]
The convergent Tate theta product has the exact identities
\[
 \Theta(q_TZ)=-Z^{-1}\Theta(Z),\qquad
 \Theta(Z^{-1})=-Z^{-1}\Theta(Z).
 \tag{SM.10}
\]
Its ratio \(f_P(Z)=\Theta(Z/P)/\Theta(Z)\) has divisor
\((P)-(O)\) and multiplier \(P\). It identifies (SM.9)
with the actual source Cartier line and its canonical
meromorphic section. On the continued chart \(P'=P^{-1}\),
and
\[
 f_{P^{-1}}(Z^{-1})=P f_P(Z).
 \tag{SM.11}
\]
Thus the divisor-compatible line transition is exactly
\((Z,V)\mapsto(Z^{-1},PV)\). Reindex once by the new
Tate deck \((Z',V')\mapsto(q_TZ',P^{-1}V')\).
The two line factors cancel, giving the literal
node-annulus transition
\[
                (Z,V,q_T)\longmapsto(q_T/Z,V,q_T).
 \tag{SM.12}
\]
Its square is the identity. The root of \(a\) exchanges
the two local branches; every finite multiple of \(G\)
avoids the critical point at the limiting cusp because a
section cannot meet a critical point of the elliptic
map. Hence this Cartier frame extends across that local
source neighbourhood. No component-supported line shear
has been inserted.

\item The source translation and its Hom lift are also transported:
their Tate form is
\((Z,V)\mapsto(g^{-1}Z,\lambda q C\,V)\), with actual
nonzero base units \(g,C\). Equation (SM.12) changes
\(g\) to \(g^{-1}\) and leaves its scalar factor fixed.
It intertwines the old and continued suspension maps
literally. The \(III^*\) cap data have holomorphic
coefficients on the full small \(u\)-disc, including
\(u_*\); their division section, unit root and finite
character can be chosen there by the exact constructions
of F41/F51. Their annular comparison remains part of
this transport, rather than being replaced by a different
cap.

\item On the cone-height lattice (SM.12) is
\[
             (v_Z,v_V,h)\longmapsto(-v_Z+h,v_V,h).
 \tag{SM.13}
\]
Its action on edge directions is
\(J=\operatorname{diag}(-1,1)\). It admits no invariant
index-one unimodular periodic three-edge triangulation,
even allowing an affine translation. Indeed the three
unoriented primitive edge directions have the three
distinct nonzero reductions in \(\mathbb F_2^2\).
Since \(J\) is the identity modulo two, an invariant
triangulation would require it to preserve each edge
line separately. A reflection has only two eigenspaces,
whereas these are three distinct lines.

\item Thus one loop about this selected-cusp puncture does not
close the chosen one-vertex toroidal completion under this
specified marked source, line and annular transport. Its
second iterate has the identity lattice and Cartier-line
transition (SM.12), and the finite cap coefficient
charts return as above. Using these literal transition maps
in the construction gives the marked holomorphic family
on the double cover \(a=\rho^2\). The actual return of that
closed double source loop remains the smooth relative
transport problem described next. An identity endpoint
coordinate transition alone does not compute a
mapping-class return.

\item The exact algebraic certificate is \path{code/part3/source_moduli_nodes.py} with its saved outputs.

\end{enumerate}
\end{proof}




\begin{theorem}\label{night:P3F.55}
\StatementRecord{night:P3F.55}{proved}{S6\_06\_BLIND: GPT-6 Astra (OpenAI, ChatGPT Pro mode, effort Pro), recorded re-derivation in the checkpoint dated 18 November 2026}{\cref{night:P3F.56}}
Let
\[
 f(s,\rho)=(X,Y,R)
       =(s^2-\rho^2,\ s(s^2-\rho^2),\ \rho):
                    \mathbb C^2\longrightarrow\mathbb C^3.       \tag{CL.1}
\]
Give both complex spaces their complex orientations. Choose a small positive
\(\epsilon\), put \(L_\epsilon=f^{-1}(S^5_\epsilon)\), and give this
three-manifold the outward boundary orientation of
\(f^{-1}(B^6_\epsilon)\). Choose a point of \(S^5_\epsilon\) outside its
image and an orientation-preserving stereographic chart
\(\psi:S^5_\epsilon\setminus\{p\}\to\mathbb R^5\).
For any positive constant \(\eta\), define
\[
 h(s,\rho)=\eta\,
       \frac{\operatorname{Re}(s\overline\rho)}
            {\sqrt{|s|^2+|\rho|^2}},\qquad
 e=(\psi f,h):L_\epsilon\longrightarrow\mathbb R^5\times\mathbb R.
                                                                  \tag{CL.2}
\]
Then \(L_\epsilon\) is an oriented three-sphere, \(e\) is a smooth
embedding, and its Haefliger invariant in the convention of
Gauniyal--Turchin, Theorem 4.1, is \(+1\).

This statement computes the specified local lifted link. Its comparison with a completed source-parameter family requires the separate ambient and circle-line morphisms.
\end{theorem}
\begin{proof}
\begin{enumerate}[label=\arabic*.]
\item \par

\medskip\noindent\textbf{1. The sphere, the immersion, and its complete double locus}

\par\nobreak

\item The function
\[
 \|f(s,\rho)\|^2
   =|s^2-\rho^2|^2+|s(s^2-\rho^2)|^2+|\rho|^2
\]
is proper and has strictly positive derivative along every positive radial
ray away from the origin: its three summands scale by \(r^4,r^6,r^2\),
respectively, and cannot vanish simultaneously away from zero. Each ray
therefore meets \(L_\epsilon\) once and transversely. Radial projection
is a smooth orientation-preserving diffeomorphism
\(L_\epsilon\to S^3\), with the boundary orientations just specified.

\item The actual derivative is
\[
 df=\begin{pmatrix}
       2s&-2\rho\\
       3s^2-\rho^2&-2s\rho\\
       0&1
     \end{pmatrix}.                                               \tag{CL.3}
\]
Its minor ideal is \((s,\rho^2)\); hence its only rank drop is at the
origin. The radial transversality above implies that
\(f|_{L_\epsilon}\) is an immersion into \(S^5_\epsilon\).

\item Two points with the same image have the same \(\rho\) and the same
\(s^2\). The unequal pair is \(s,-s\), and the equality of the second
coordinates is \(2s(s^2-\rho^2)=0\). Thus the complete double locus is
\[
 L_+=\{s=\rho,\ |\rho|=\epsilon\},\qquad
 L_-=\{s=-\rho,\ |\rho|=\epsilon\}.                                \tag{CL.4}
\]
Its image is \(D=\{(0,0,\rho):|\rho|=\epsilon\}\).
There are no triple points. At a double point the two transverse complex
directions are \((2\rho,2\rho^2,0)\) and
\((-2\rho,2\rho^2,0)\). They are independent for \(\rho\ne0\);
the sheets meet transversely along their common complex \(\rho\)-line.
Their restrictions to the links therefore meet transversely along \(D\).
The projected immersion is generic in the precise sense of the theorem.

\item On the two double branches \(h=\pm\eta|\rho|/\sqrt2\).
Consequently (CL.2) separates every pair not separated by \(f\).
It is an injective immersion of a compact manifold, hence an embedding.
This proves the actual lift, including its ordering.

\par

\medskip\noindent\textbf{2. The double-link orientation and the exact pushed circle}

\par\nobreak

\item We use the outward-normal-first boundary convention and the transverse
intersection orientation specified in Section 1 of Gauniyal--Turchin.
For the two complex sheets in \(\mathbb C^3\), their common complex
line has its complex orientation: the complementary complex directions
give a positive real determinant. On the link the outward common radial
vector precedes its phase vector. Thus the Ekholm orientation of \(D\)
is increasing \(\arg\rho\), and the induced orientations on both
\(L_+\) and \(L_-\) are also increasing \(\arg\rho\).
The two transverse complex lines \(s=\rho\) and \(s=-\rho\) meet at
the origin with intersection number \(+1\). Equivalently their radial
images form the positively oriented Hopf link. Hence
\[
                         \operatorname{lk}(L_+,L_-)=1.            \tag{CL.5}
\]

\item The theorem's invariant \(E\) uses a normal field on the entire double
preimage whose pushed union has zero class in the complement. We give
that field explicitly. On either branch set
\[
                         v_\pm=\overline\rho\,\partial_s.         \tag{CL.6}
\]
These vectors are tangent to \(L_\epsilon\): at \(X=Y=0\), the
differential of \(\|f\|^2\) depends only on the real radial
\(\rho\)-coordinate, which (CL.6) leaves fixed. They are transverse
to the double circle. The constant \(\partial_s\) framing has
self-linking zero. Indeed the holomorphic coordinates
\((s\mp\rho,\rho)\) identify its small push with a parallel complex
disc disjoint from the core disc. Multiplication by \(\overline\rho\)
has winding \(-1\), so the self-linking of each pushed branch is
\(-1\). Together with (CL.5), each total linking is
\(-1+1=0\), exactly the zero-class requirement of their equations
(4.1)--(4.2).

\item No normal field is now replaced by a homotopic one without its derivative.
The actual two images of (CL.6) satisfy
\[
 df_{\rho,\rho}(v_+)+df_{-\rho,\rho}(v_-)
                    =(0,4|\rho|^2\rho,0).                       \tag{CL.7}
\]
The pushed image circle is therefore, after its positive radial
readjustment, isotopic in the complement to the link of the complex line
\[
                         (X,Y,R)=(0,c\rho,\rho),\qquad c>0.       \tag{CL.8}
\]
The image of (CL.1) is the reduced hypersurface
\[
                         Y^2-X^2(X+R^2)=0.                      \tag{CL.9}
\]
Its normalization is exactly (CL.1): \(s\) is integral, satisfying
\(s^2=X+R^2\), and \(s=Y/X\) on the dense set \(X\ne0\).
On the disc bounded by (CL.8), the left side of (CL.9) is
\(c^2\rho^2\). Its winding is \(+2\). A small transverse analytic
perturbation gives the same two positive intersections, so the oriented
linking number of the pushed circle and the immersed link is
\[
                                  E=2.                           \tag{CL.10}
\]
This computes the specific tangent push in the cited formula; it does
not substitute a differently framed normal push.

\par

\medskip\noindent\textbf{3. The actual Stiefel and quaternion comparison}

\par\nobreak

\item Némethi--Pintér, Theorem 1.2.2, Proposition 4.2.3 and Proposition 6.1.1,
apply to the holomorphic germ (CL.1), whose isolated rank drop was proved
above. The quotient of its minor ideal has basis \(1,\rho\). Their
complex Smale class is therefore \(2[u]\), and its realification is
\[
                  \Omega(f)=-2[L]\in\pi_3(SO(5)).                \tag{CL.11}
\]
Here we retain their literal convention
\(q=z+w j\), with \((z,w)\) carrying the complex orientation,
\(L_q(p)=qp\), and \(u_q(p)=pq\) in its complex two-dimensional
matrix representation. Their realification arrow sends \([u]\) to
\(-[L]\). The local stereographic chart is the same kind of
oriented link chart used in Proposition 6.1.1, so (CL.11) concerns the
actual projected link, with no independent source framing substituted.

\item We next compute the morphism needed by the coefficient
\(\Omega_{S^2}\) in Gauniyal--Turchin. The fibre inclusion
\(S^2\to V_{5,3}\) is the orbit of the last vector under a fixed
three-plane rotation. Its quaternion lift is
\(\operatorname{Ad}_q(p)=qpq^{-1}\). Proposition 4.2.2(c) of
Némethi--Pintér gives \([\operatorname{Ad}]=2[L]\) stably.
The projection of this lift is
\[
                            q\longmapsto qi q^{-1}.              \tag{CL.12}
\]
Its Hopf invariant is \(-1\) in the specified complex source
orientation. To check the sign directly, its circle fibres are
\[
          (z,w)\longmapsto(ze^{i\theta},we^{-i\theta}).
\]
At \(q=1\) the derivative on the complementary \(w\)-plane is
\((u,v)\mapsto(2v,-2u)\) in the ordered \(j,k\)-plane, with
positive determinant. Thus the oriented fibre over \(i\) is the
positive \(z\)-circle. The oriented fibre over \(-i\) is the
negative \(w\)-circle, by the same right-action vector \(qi\).
Their linking number in the complex-oriented three-sphere is \(-1\),
which proves the assertion.

\item The Whitehead square \([\operatorname{id}_{S^2},
\operatorname{id}_{S^2}]\) has Hopf invariant \(+2\): the attaching
map of the oriented top cell of \(S^2\times S^2\), followed by
the fold of its two two-spheres, has cup square
\((a+b)^2=2ab\). The Hopf fibration exact sequence identifies
\(\pi_3(S^2)\) with \(\mathbb Z\) by this linking invariant.
Consequently this precise Whitehead square is \(-2\) times
(CL.12), and its image in the Stiefel group is
\[
       [\operatorname{id},\operatorname{id}]
                \longmapsto-2[\operatorname{Ad}]=-4[L].         \tag{CL.13}
\]
Combining (CL.11) and (CL.13), with both arrows typed, gives
\[
                             \Omega_{S^2}(f)=\frac12.             \tag{CL.14}
\]

\par

\medskip\noindent\textbf{4. The embedding invariant and its exact scope}

\par\nobreak

\item All hypotheses of Gauniyal--Turchin, Theorem 4.1, have been verified:
the lift is an actual embedding; its projection is a generic immersion;
both double-preimage components are explicitly oriented; their required
total-link-zero framing and its push have been computed; and the
Stiefel coefficient uses their Whitehead generator. The spherical
version stated directly before that theorem applies after the chosen
stereographic chart. Its formula gives
\[
 \mathcal H(e)
   =\frac12\operatorname{lk}(L_+,L_-)
      +\frac16E(f)+\frac13\Omega_{S^2}(f)
   =\frac12+\frac13+\frac16=1.                                  \tag{CL.15}
\]
Changing \(\eta\) through positive values is a literal isotopy of
embeddings, since the projected double pairs remain separated.
The same holds for any interpolation of the height that preserves
the displayed two-branch order.

\item This proves the invariant of (CL.2). The ambient source-to-orbit map,
the outside cap, the circle Euler line, and the boundary-normal
comparison of a completed threefold are additional morphisms.
Their proof belongs to P3F.56; (CL.15) alone does not identify that
threefold's return.

\par

\medskip\noindent\textbf{5. Exact comparison with the corrected Seifert-surface formula}

\par\nobreak

\item We identify the normalization of \(\mathcal H\) just computed with
the corrected Takase invariant, rather than infer a sign from their
common name. Both use the outward-normal-first boundary convention.
Takase's correction paper, Corollary 2.6, gives the group isomorphism
\[
 \mathcal H_{\rm T}(e)
      =\frac{e(\nu V)^2-\sigma(V)}8                 \tag{CL.16}
\]
for an actual oriented Seifert embedding \(V^4\subset S^6\).
Gauniyal--Turchin, Theorem 4.1, gives the group isomorphism used in
(CL.15). We compare them on an explicitly nonzero embedded projection.

\item Takase, Remark 4.2, gives a punctured Kummer surface
\(V\subset S^5\), with \(\sigma(V)=-16\); its existence is
also the example of Haefliger (1966), Theorem 5.17. Its normal line
in \(S^5\) is trivial because it is oriented. Composing with the
equatorial inclusion in \(S^6\) therefore gives a trivial normal
two-plane, so \(e(\nu V)=0\). Choose a point of \(S^5\) outside
this compact surface and an orientation-preserving stereographic
chart. The boundary becomes an actual embedding
\(g:S^3\hookrightarrow\mathbb R^5\), with empty double locus,
and the same Seifert surface in \(\mathbb R^5\).

\item Use the oriented boundary identification of Némethi--Pintér, Section
2.4, with the same outward-normal-first convention. Their sign-refined
Theorem 9.1.6, applied to the embedded Seifert surface, gives
\[
                 \Omega_a(g)=\frac32\sigma(V)=-24
                         \quad\hbox{in the }[L]\hbox{ convention}.
                                                               \tag{CL.17}
\]
The literal fibre-inclusion morphism (CL.13) sends the Whitehead
square to \(-4[L]\). It follows that
\(\Omega_{S^2}(g)=6\). Since the double locus is empty, the two
linking terms in Gauniyal--Turchin's formula vanish, and hence
\[
                  \mathcal H(g)=6/3=2
                    =\frac{0-(-16)}8=\mathcal H_{\rm T}(g).
                                                               \tag{CL.18}
\]
In particular this is a nonzero element of the embedding group.
The two displayed invariants are isomorphisms of that same group
\(C_3^3\cong\mathbb Z\). Their composite automorphism of
\(\mathbb Z\) fixes the nonzero integer two; the only such
automorphism is the identity. Thus
\[
                           \mathcal H=\mathcal H_{\rm T}
                                                               \tag{CL.19}
\]
on the entire group, with its specified orientation convention.
This identifies the exact local value \(+1\) in (CL.15) with
the corrected Seifert-surface formula. It still makes no inference
about a source-family return without the separate morphisms of
P3F.56. As an additional direct source check, Sakai (2015),
Proposition 6.8 and its proof, state precisely the same embedded-
projection comparison \(\mathcal H(g)=-\Omega_a(g)/12
=-\sigma(V)/8\).

\item The exact certificates are \path{code/part3_root/source_cusp_link.py} and
\path{code/part3_root/haefliger_normalization_bridge.py}, with their JSON and
captured output. The second verifies every integer and rational arrow
in the nonzero normalization test (CL.17)--(CL.19). The local-link
certificate also checks the perturbation
\((s^2-\rho^2,s(s^2-\rho^2+\varepsilon),\rho)\): it has two ordinary
crosscaps at \(s=0,\rho^2=\varepsilon\) and no triple point.
The first structural-factor comparison failure is preserved in
\path{code/part3_root/attempts/}; its correction uses an exact polynomial
zero-difference check.

\end{enumerate}
\end{proof}




\begin{theorem}\label{night:P3F.56}
\StatementRecord{night:P3F.56}{proved}{S6\_06\_BLIND: GPT-6 Astra (OpenAI, ChatGPT Pro mode, effort Pro), recorded re-derivation in the checkpoint dated 18 November 2026}{\cref{night:P3F.22}, \cref{night:P3F.24}, \cref{night:P3F.5}, \cref{night:P3F.51}, \cref{night:P3F.54}, \cref{night:P3F.55}, \cref{night:P3F.9}, \cref{night:P3K.23}, \cref{night:P3K.30}, \cref{night:P3K.68}, \cref{night:P3K.7}}

Use the exact marked E7 source of P3F.54,
\[
 y^2=x^3+3r(u)x^2+t x+1,\qquad G=(0,1),\qquad
 r(u)=\frac{1-2u^3}{3u^2},\quad t_0(u)=u^2-\frac2u.
\]
Fix either actual one-cap sphere datum of P3F.41:
\[
 \begin{array}{c|c|c}
 (M,Q)&\text{order at }III^*&\text{root fibre}\\ \hline
 (\mathcal O(2G-O),G)&m=4l+3,\ l\ge0&III\\
 (\mathcal O(G-O),2G)&m=4l,\ l\ge1&\text{smooth}.
 \end{array}
\]
Keep its norm numerator \(n=\pm1\), its line lift and finite character,
and its sufficiently small nonzero scalar fixed. There is no additional
logarithmic disc on either member of the coalescing finite-node pair.
Use the standard vertical-edge toric fan of this construction.

Let \(u_*^3=-1\), and take the actual closed twofold source family
\[
 u+u^{-2}=\rho^2,\qquad u(0)=u_*.
\]
For one positive \(\rho\)-loop its forward physical return \(\Phi\)
has
\[
 A_{M,\rm Jin}(f)=-1,\qquad [\Phi]=-1\in\Theta_7\cong\mathbb Z/28.
\]
Here \(f\) is the quotient of that physical return with the orientations
and physical-line Hopf transfer of P3F.22 as explicitly identified in P3K.68. This is two positive turns about \(u_*\)
in the marked source parameter; the scalar has zero winding.
The actual sphere mapping torus is
\((S^1\times S^6)\mathbin{\#}\Sigma_{-1}\). Both its spin structures
have Eells--Kuiper invariant \(-1/28\), so it is not diffeomorphic to
\(S^1\times S^6\), including without an orientation requirement.

This particular full \(\rho\)-family has no descent as an oriented
smooth sphere bundle through the actual degree-two cover
\(p:D_\rho^*\to D_a^*\), \(a=\rho^2\), with \(u\) the analytic
inverse of \(a(u)=u+u^{-2}\) near \(u_*\). Indeed its clutch class
minus one cannot be twice a class in \(\Theta_7\).

The selected-node separation has braid \(\sigma^{12}\), and its
retained elliptic boundary map is the once-punctured-torus relation
\((gh)^6=T_{\partial}\). The proof below gives the actual ambient,
line, normal and cap morphisms carrying the local lifted link to this
return. It assigns no return to another fan or another source puncture.

\end{theorem}
\begin{proof}
\begin{enumerate}[label=\arabic*.]
\item \par

\medskip\noindent\textbf{SC.1. The exact normalization and the ordered critical-value braid}

\par\nobreak

\item Retain the F54 functions and put
\[
 a(u)=u+u^{-2}=\rho^2,\qquad u(0)=u_*,\quad u_*^3=-1.
\]
The selected fibre is at \(q=t-t_0(u)=0\). In the actual source-parameter
threefold, the coordinates \((X,y,\rho)\) are smooth coordinates and
\[
 q=\frac{y^2-X^2(X+\rho^2)}{X+u(\rho)}.                 \tag{SC.1}
\]
The entire selected rational fibre is normalized by
\[
 f(s,\rho)=\bigl(s^2-\rho^2,
                    s(s^2-\rho^2),\rho\bigr).         \tag{SC.2}
\]
At \(s=\infty\) this is the zero-section point, so its domain is
\(N=\mathbf P^1_s\times D_\rho\). Its two preimages of the node are
the actual sections
\[
 D_+=\{s=\rho\},\qquad D_-=\{s=-\rho\}.              \tag{SC.3}
\]
They meet only at \((s,\rho)=(0,0)\).

\item The derivative minors of (SC.2) are
\[
 2s,\qquad 3s^2-\rho^2,\qquad
                      2\rho(s^2-\rho^2).
\]
Their ideal is \((s,\rho^2)\). Thus \(f\) is a complex immersion off
the origin. Its only distinct double pairs there are
\((\rho,\rho)\) and \((-\rho,\rho)\): equality of the last and first
coordinates first gives the same \(\rho\) and \(s'=\pm s\), and equality
of the second then gives \(s(s^2-\rho^2)=0\).

\item The two nearby critical values of (SC.1) can also be read without a
change of the base coordinate. Besides \(X=0,y=0\), the second critical
point is \(X=aR(a),y=0\), where
\[
 2aR^2+(a+3u)R+2u=0,\qquad R(0)=-2/3.
\]
The implicit derivative in \(R\) is \(3u_*\ne0\). Its actual critical
value is
\[
 q_1(\rho)
   =-\rho^6\frac{R(\rho^2)^2(R(\rho^2)+1)}
                    {\rho^2R(\rho^2)+u(\rho)},\qquad
 \frac{q_1(\rho)}{\rho^6}\longrightarrow-\frac4{27u_*}.
                                                               \tag{SC.4}
\]
The remaining critical point stays outside this small source chart.
For one positive \(\rho\)-loop the ordered pair \((0,q_1)\) therefore
has six positive full turns of its difference. In the convention that
the positive generator \(\sigma\) of the two-point braid group is a
positive half turn, this is exactly \(\sigma^{12}\).

\item Here is the retained Hurwitz convention. For an ordered pair of positive
Picard--Lefschetz twists \((g,h)\), the right Hurwitz action is
\[
 (g,h)\cdot\sigma=(h,h^{-1}gh).
\]
Writing \(C=gh\), direct multiplication gives
\[
 (g,h)\cdot\sigma^{2j}
       =(C^{-j}gC^j,C^{-j}hC^j).                    \tag{SC.5}
\]
The two vanishing arcs are adjacent arcs between the three roots of the
cubic in (SC.1). Their lifts to its double cover meet once. The local
regular fibre is the double cover of a disc branched at these three
points, hence a once-punctured torus. The positive twists are the lifts
of the two positive adjacent half twists. The braid
\((\sigma_1\sigma_2)^3\) is one full turn of the three branch points.
Its lift rotates the inner boundary coordinate by half a turn, because
the boundary covering has degree two. Squaring it gives the actual
boundary twist:
\[
                        (gh)^6=T_{\partial}.        \tag{SC.6}
\]
For example, on a boundary collar with covering angle
\(\vartheta=2\widetilde\vartheta\), the base twist
\(\vartheta\mapsto\vartheta+2\pi\chi(r)\) lifts to
\(\widetilde\vartheta\mapsto\widetilde\vartheta+\pi\chi(r)\);
its square is the positive boundary twist of the lifted collar.
This is the boundary map, not merely the order-six homology matrix.
Thus (SC.5) retains conjugation by \(T_{\partial}^{-1}\) after the
positive \(\rho\)-loop. The ordered vanishing curves return because the
boundary twist is central, while the boundary twist itself is retained.
The weighted chart rotates the three cubic roots twice and has the
opposite, inner-to-outer boundary arrow \(T_{\partial}\). These two
arrows are the same gluing relation read in opposite directions.

\par

\medskip\noindent\textbf{SC.2. The explicit compact pencil and its plumbing}

\par\nobreak

\item The substitution
\[
 X=\rho^2\widetilde X,\qquad
 y=\rho^3\widetilde Y,\qquad q=\rho^6Q
\]
is the \(\rho\)-chart of the weighted blowup of the smooth coordinates
\((X,y,\rho)\) with weights \((2,3,1)\). The complete equation is
\[
 \widetilde Y^2=\widetilde X^2(\widetilde X+1)
                    +Q\bigl(u(\rho)+\rho^2\widetilde X\bigr).
                                                               \tag{SC.7}
\]
On the exceptional weighted plane \(\mathbf P(2,3,1)\), with coordinates
\([X:Y:R]\), the actual pencil is
\[
 Q=[Y^2-X^2(X+R^2):u_*R^6].                          \tag{SC.8}
\]
The constant base unit \(u_*\) is retained.

\item Its only base point is the smooth point \([1:1:0]\). In local coordinates
\(R,z=Y^2-X^2(X+R^2)\), the ideal is \((z,R^6)\). The six successive
blowups with charts \(z=R^jz_j\), \(j=1,\ldots,6\), resolve the
pencil. The first five exceptional curves are in the infinite fibre,
with multiplicities \(5,4,3,2,1\); the sixth is its section. The proper
curve \(R=0\) has multiplicity six. The quotient point of order two on
this curve has the scalar \((1,1)\) chart; its resolution contributes a
single curve with multiplicity three. The order-three quotient has
weights \((2,1)\); the two regular rays have \(R\)-valuations
\(2/3,1/3\), and contribute multiplicities four and two, in that order
away from the central curve.

\item Every vertical curve has self-intersection \(-2\). For the central
curve this can also be checked before the point blowups: its weighted
self-intersection \(1/6\), minus the two resolution corrections
\(1/2\) and \(2/3\), is \(-1\); the first ordinary base-point blowup
makes it \(-2\). Each of the first five ordinary exceptional curves
is blown up once more, while the last section remains a \((-1)\)-curve.
Thus the infinite fibre is the literal affine \(E_8\) plumbing with
central multiplicity six and three arms
\[
 (5,4,3,2,1),\qquad(3),\qquad(4,2).                  \tag{SC.9}
\]
The finite discriminant of (SC.7) at \(\rho=0\) is
\(-u_*Q(4+27u_*Q)\), giving the two displayed nodal fibres. The compact
source pencil therefore has \(II^*+I_1+I_1\). Its integral intersection
matrix, including the last section, is unimodular of signature \(-8\);
the executed certificate supplies the integral basis consisting of the
fibre, section and the eight finite \(E_8\) roots.

\item The pullback of the original source Cartier line to this blowup and all
six further blowups is trivial near the exceptional pencil: both original
divisor sections avoid the source cusp point. This assertion is about
that literal pullback. It is not a substitution of the compact pencil
for the actual principal-circle line. The latter comparison is supplied
by the real trace below.

\par

\medskip\noindent\textbf{SC.3. The actual ambient orbit product and the embedded trace}

\par\nobreak

\item Choose a small base disc containing just the two coalescing finite
critical values, and a small source-parameter disc. On its boundary the
elliptic map is a submersion for every parameter. The source total family
is smooth, including the cuspidal member, and is proper with this boundary.
A horizontal field tangent to the boundary therefore gives a smooth
source trivialization over the parameter disc. The III-star cap and the
ordinary complement extend there by the actual full-coefficient maps of
F51/F54. These choices retain their original annular boundary marking.

\item The suspension translation also has a logarithm on this small base and
parameter polydisc. At the central cuspidal fibre its smooth group is
additive. Divide its central section value by a sufficiently large
integer, continue that division section by the inverse function theorem,
and take the small-group logarithm of this section. Multiplication by the
same integer gives a logarithm of the actual section. Exponentiating
\(t\) times it, \(0\leq t\leq1\), supplies an actual extending
translation isotopy. The invertible line lift obtained after dividing
the chosen Hom section by its base factor \(q\) has the same construction
in the line-extension group: its scalar unit has an analytic logarithm
on the polydisc. Thus this isotopy retains the source-line lift as well
as the source section.

\item Use this isotopy to untwist the source suspension. At the selected node
use the exact relative orbit product of P3F.9. Its local formulas
\[
 2b=r-|y_0|^2/r,\qquad
 2b'=|x_0|^2r-r^{-1},\qquad db'/db>0
\]
give a proper increasing gamma coordinate, fixed on the prescribed
outside collar. They identify the actual orbit body with the source
neighbourhood times its marked gamma circle. They are periodic on the
\(\rho\)-cover, since the ordered node coordinates return literally
by F54. The target source product exists also at \(\rho=0\). Gluing it
to the already extending exterior gives a smooth ambient orbit-body
family over the disc. We use its trivialization as \(Y\times D_\rho\).
No extension of the complex threefold at \(\rho=0\) is made here.

\item There is a global real gamma lift of the fixed-sphere graph on the whole
closed parameter annulus. The ordinary gamma coordinate is formed from
real logarithmic norms, divided by the positive logarithmic norm of the
contracting suspension. The source and invertible-line logarithms just
constructed extend across the full parameter disc. At the node the F9
coordinate is a proper real coordinate with positive derivative, fixed
relative to that ordinary coordinate. Its two endpoint values therefore
have a real lift, not only a circle-valued difference. The full \(\rho\)
closing has no gamma reindex: by F54 the two literal node transitions
\((Z,V)\mapsto(q_T/Z,V)\) compose to \((Z,V)\mapsto(Z,V)\).
A nonzero suspension deck power would change \(|V|\) by a nontrivial
power of its strictly contracting norm; hence it is not present in this
identity. The chosen gamma cell and its real coordinate return to
themselves. An increasing real change of the F9 coordinate, fixed on the
outer collar, puts the entire compact fixed graph in one common short
arc. This proves its zero parameter winding and supplies the real lift
used in the convex interpolation below. A scalar Cox-line gauge is a
different map and is retained in SC.4--SC.5.

\item The source projection of the fixed sphere is the actual normalization
(SC.2). Fractional source translations used in the suspension product
do not change this local normalization germ. Here is their exact
normalization map. If \(Z=(s-\rho)/(s+\rho)\) and a translation has
Tate multiplier \(g(\rho)\), with \(g(0)=1\) and
\((1-g)/\rho\) analytic, then
\[
 s\longmapsto
 \frac{(1+g)s+(1-g)\rho}
      {((1-g)/\rho)s+(1+g)}.                        \tag{SC.10}
\]
It is analytic at the origin, fixes \(s=\pm\rho\), and has invertible
derivative there. The same formula for the fractional multiplier
\(\exp(t\log g)\) provides the actual isotopy. The positive product
coordinate may be chosen in an arbitrarily short gamma arc. For an exact
quantitative check, write \(g_t=\exp(t\rho A(\rho))=1+\rho h_t\).
Formula (SC.10) gives
\[
 T_t(s,\rho)-s
       =\frac{h_t(s^2-\rho^2)}{2+h_t(\rho-s)}.       \tag{SC.10a}
\]
For bounded \(t\) with bounded first derivative, \(h_t\) and its
first derivative are bounded. Uniformly for \(0\leq\alpha\leq1\),
substitution of \(\alpha t\) gives
\[
 T_{\alpha t}=s+O((|s|+|\rho|)^2),\qquad
 D(T_{\alpha t},\rho)=I+O(|s|+|\rho|).
\]
On one sufficiently small common ball the error in the derivative has
norm less than one half. The maps are therefore injective diffeomorphism
germs, with positive real determinant, and the displayed isotopy fixes
both branches exactly. The height (SC.11) has bounded first derivative
off the origin. Inside the removed inner ball it can be cut off to the
smooth polynomial \(\operatorname{Re}(s\overline\rho)\), without
changing \(W\) or its inner boundary and with a bounded first derivative.
Thus the same estimate gives a smooth isotopy of actual germs.

\item To avoid circular use of this estimate, first lift the fixed graph to the
universal real gamma cover furnished above. In that cover its source
projection is the normalization itself. Perform the height interpolation
there, in the common arc of length less than one. Quotienting cannot
create a new identification by a nonzero gamma deck, since two heights
in that arc do not differ by a nonzero integer. Now untwist the
suspension by the fractional translation. Equation (SC.10a) is the
actual resulting normalization reparametrization and removes that
source shear by the explicit isotopy. This proves the use of (SC.2)
as the projected germ, with its marking retained.

\item The order of the two gamma heights is fixed by the charts, rather than
chosen afterward. The source branch \(y_0=0\) is \(s=-\rho\) and
occurs at the lower endpoint \(b=0\). The branch \(x_0=0\) is
\(s=+\rho\) and occurs at the upper endpoint \(b'=0\). The increasing
overlap identifies them in precisely this order.

\item On the local normalization put
\[
 H(s,\rho)=\frac{\operatorname{Re}(s\overline\rho)}
                  {\sqrt{|s|^2+|\rho|^2}}.          \tag{SC.11}
\]
It has the same strict ordering on the two branches. Multiplying by a
fixed small positive constant puts it in the chosen gamma arc. Extend
it with a cutoff over the rest of \(\mathbf P^1\). On a collar of the
outer parameter circle retain the actual gamma graph. Interpolate by
a convex combination on the intervening annulus. The difference of
the heights on every double pair remains positive. Away from those
pairs the source map is already injective. Hence this is an isotopy
through embedded graphs, with the actual outer trace unchanged.

\item Remove the small complex-oriented star-shaped four-ball specified below from \(N\), and call
the result \(W\). The map \((f,H)\), with the just specified outer
interpolation, is an actual smooth embedding of \(W\) in
\(Y\times D_\rho\), and its outer boundary is the actual fixed-sphere
trace with orientation \(F\times S^1_\rho\). Its inner boundary is the
lifted link of (SC.2).

\item Here is the literal inner boundary used in F55. Choose a small positive
\(\epsilon\), and remove from the normalization the star-shaped four-ball
\(B_f=f^{-1}(B^6_\epsilon)\). Its boundary is exactly
\(L_\epsilon=f^{-1}(S^5_\epsilon)\). The norm square of \(f\) is
strictly increasing on every positive domain ray, so this is a smooth
three-sphere with its stated complex boundary orientation. A smooth
radial interpolation, fixed near a smaller Euclidean ball, identifies
\(B_f\) with a four-ball without changing its boundary orientation.

\item Choose the ambient inner seven-ball to be the rounded product
\(B^6_\epsilon\times[-2\eta\epsilon,2\eta\epsilon]\), with
\(\eta H\) as height. Indeed \(|H|\leq|\rho|\leq\epsilon\) whenever
\(\|f\|\leq\epsilon\); thus the embedded graph misses the two flat
caps and every rounded corner. Its intersection with the boundary is
literally
\[
 (f,\eta H):L_\epsilon\longrightarrow
                 S^5_\epsilon\times(-2\eta\epsilon,2\eta\epsilon).
                                                               \tag{SC.11a}
\]
This boundary is the smooth six-sphere obtained by rounding the product
of a six-ball and an interval. In its cylindrical chart, use precisely
the orientation-preserving stereographic map \(\psi\) of F55. Equation
(SC.11a) is then exactly \((\psi f,\eta H)\), with its prescribed
ordering. The chart is an actual open chart of this six-sphere; shrinking
a compact subchart by its radial coordinate gives the usual chart
comparison, through embeddings. No comparison of two different radial
projections is required. The weighted radial construction of the first
draft is retained in the attempts record; the present product-ball map
supplies the missing literal boundary identification.

\par

\medskip\noindent\textbf{SC.4. The extending normal spinor and actual Hopf circle line}

\par\nobreak

\item The finite-base source has a nowhere-zero holomorphic two-form. In the
coordinates (SC.1), its residue expression is a nonzero holomorphic
unit times \(dX\wedge dy\). Together with \(d\rho\) it trivializes the
ambient complex determinant along the normalization. The complex normal
line of the immersion (SC.2) on \(W\) is consequently
\[
 L=N_f\cong\mathcal O(-D_+-D_-).                     \tag{SC.12}
\]
This is a determinant isomorphism, not merely a degree equality. In the
finite coordinate chart its normal covector is the explicit cross product
\[
 (3s^2-\rho^2)\,dX-2s\,dy
                    +2\rho(s^2-\rho^2)\,d\rho.     \tag{SC.13}
\]
It is nonzero everywhere on \(W\) in this chart. Changing to the
infinite normalization coordinate gives the \(\mathcal O(-2)\)
normal transition. The divisor identity
\(\operatorname{div}(s^2-\rho^2)=D_++D_--2\{s=\infty\}\)
gives the displayed actual isomorphism (SC.12).

\item The normal bundle of the real graph is isomorphic to
\(L_{\mathbb R}\oplus\mathbb R_\gamma\): quotient by its tangent
graph gives the identity on the source normal two-plane and the gamma
coordinate, with only a triangular shear. Its linear contraction is
the precise graph-normal comparison. Retain the same constant orientation
change as in the positive-Hopf convention of K7.

\item Define the rank-two complex spinor bundle
\[
 E=(\mathcal O\oplus L)\otimes\mathcal O(D_+)
       =\mathcal O(D_+)\oplus\mathcal O(-D_-).        \tag{SC.14}
\]
Its determinant is \(\mathcal O(D_+-D_-)\), trivialized by the literal
Cartier ratio \((s-\rho)/(s+\rho)\). Thus \(E\) is an actual
\(SU(2)\) bundle, and
\[
 \operatorname{ad}E\cong
              \mathbb R\oplus L_{\mathbb R}.        \tag{SC.15}
\]
The off-diagonal summand is exactly
\(\operatorname{Hom}(\mathcal O(D_+),\mathcal O(-D_-))=L\).
Its Hopf sphere bundle carries the meridian-one Thom line U with
determinant class b=0. The physical nonzero-vector line M is its dual:
in the literal normal chart (v,z), the positive spinor is (bar v,z),
whose physical circle is the unit circle of the tautological line.
The actual outward normal meridian has its projective complex orientation,
so that physical line has degree -1. The Hermitian dual map supplies
the degree-one Thom line and inverts its circle character. P3K.68
gives this complete coordinate and orientation computation.

\item This identifies the actual outer physical circle line and its meridian-one dual; their characters are retained separately. In the normalization coordinate
\(Z=(s-\rho)/(s+\rho)\), the two Cox normal coordinates have transition
\(c_i'=Zc_i\) and physical weights \((-1,+1)\). Conjugating the
negative-weight coordinate gives the unitary transition
\[
       \operatorname{diag}(\overline Z/|Z|,Z/|Z|),
\]
which is precisely that of (SC.14), after the explicit Hermitian
identification of the conjugate tautological line with its dual. It
intertwines the physical circle actions. The quaternion formula
\[
 q_{\rm H}=(1+jZ)(c_0+j\overline{c_1})
\]
is the same formula as K30; its second-chart expression gives the
identical vector. The ratio \([s-\rho:s+\rho]\) is a genuine map from
\(W\) to \(\mathbf P^1\), since its only common zero was removed.
Consequently this Cox/Hopf identification extends over all of \(W\),
including the inner link, with the original ordered branch marking.

\item The pullback of the source Picard line to
\(\mathbf P^1\times D_\rho\) has degree zero on every normalization
fibre and hence is a line from the parameter disc. Rigidifying it along
the zero section gives a literal trivialization. The source-line unit
in the suspension isotopy above is retained in this trivialization.
It contributes the same extending scalar gauge on both sides of
(SC.14). No degree-zero line or scalar winding is discarded.

\item The circle line also extends on the whole excision complement; this
does not follow just from constructing its restriction to the normal
sphere bundle. Put \(Z'= (Y\times D_\rho)\setminus\operatorname{int}B^7\).
It has \(H^2(Z';\mathbb Z)=H^3(Z';\mathbb Z)=0\). The codimension-three
Thom sequence of the properly embedded connected \(W\) therefore gives
\[
 H^2(Z'\setminus W;\mathbb Z)\xrightarrow{\ \cong\ }
             H^0(W;\mathbb Z)=\mathbb Z.             \tag{SC.15a}
\]
The isomorphism is evaluation on the oriented normal meridian. Choose
the line with meridian value one. On the outer complement the actual
meridian-one dual U has this value, and \(H^2\) of that complement is also
\(\mathbb Z\): each fibre is the complement of the unknotted \(S^2\)
in \(S^5\), with \(H^1=0,H^2=\mathbb Z\), and the oriented return
preserves its generator. The Wang sequence has no other degree-two
term. The ratio of the two lines is therefore topologically trivial;
multiplication by a chosen nonzero section of that ratio is the actual
boundary bundle isomorphism. It retains the stated dual circle character.

\item Furthermore restriction
\(H^2(S(\nu W))\to H^2(S(\nu W)|_{F\times S^1})\) is an isomorphism:
it is an isomorphism both on the base \(H^2\) and on the meridian
summand of the sphere-bundle sequence. Thus the restriction of the line
in (SC.15a) is exactly the line (SC.14), whose outer identification was
just given in Cox coordinates. In particular its determinant class on
\(W\) is zero. On the inner link it is the same oriented meridian line.
Possible boundary isomorphisms differ by circle-valued gauges; their
outer parameter winding is extended over the specified
\(D^3\times S^1\) cap below. This retains that winding rather than
assuming it vanishes in the excision complement.

\par

\medskip\noindent\textbf{SC.5. Excision, its zero-index correction, and the transfer sign}

\par\nobreak

\item Take a small seven-ball around the local lifted germ inside
\(Z=Y\times D_\rho\); its boundary contains the inner embedded
three-sphere \(e\). Remove the normal disc bundle of \(W\) from the
remaining ambient cobordism. Attach the product surgery collar
\((D^3\times S^1)\times S^2\) on the outer side, as in K7, and the
\(D^4\times S^2\) collar on the inner side. The outer closed boundary
is the actual physical mapping torus \(P_f\) after fixed-sphere
replacement. The inner closed boundary is the surgery manifold
\(N_e\) of the actual embedded link, with the same meridian line.
The further boundary is the normal sphere bundle over
\[
 \overline W=W\cup(D^3\times S^1)\cup D^4
       =(S^2\times D^2)\cup(D^3\times S^1)
       =\partial(D^3\times D^2)\cong S^4.            \tag{SC.16}
\]
These are the literal boundary identifications from the construction;
the two caps are not exchanged by an unspecified change of marking.

\item The varying outer normalization coordinate also has an explicit cap
extension. On the outer parameter circle put
\[
 A_\rho=\begin{pmatrix}1&-\rho\\1&\rho\end{pmatrix},
 \qquad Z=A_\rho\cdot s=\frac{s-\rho}{s+\rho}.
\]
Represent the closed three-ball by positive semidefinite Hermitian
two-by-two matrices \(P\) of trace one. Its boundary consists of the
rank-one matrices and is \(\mathbf P^1\). The exact ball map is
\[
 \mathcal B_\rho(P)
       =\frac{A_\rho P A_\rho^*}
                    {\operatorname{tr}(A_\rho P A_\rho^*)}.
                                                               \tag{SC.16a}
\]
Its inverse uses \(A_\rho^{-1}\). The denominator is positive on the
closed ball, so both maps are smooth there. On its boundary it is
exactly the displayed Mobius map. The matrix \(A_\rho\) is periodic
on the actual parameter circle, so this ball map closes literally even
when the associated projective loop has nontrivial fundamental-group
class. The source reparametrization isotopies of (SC.10a) extend as
well: extend their time-dependent sphere vector fields over a boundary
collar of the ball using a radial cutoff, and integrate them. These
extensions are periodic in the parameter. Composing them with
(SC.16a) identifies the actual outer cap marking with the product cap
in (SC.16). Its normal clutch is carried along, not set to zero.

\item The meridian-one dual U extends over both collars; the physical M extends by the inverse character. On the outer boundary
its class is the positive meridian class, by (SC.14); every possible
boundary gauge extends over \(D^3\times S^1\), since restriction is
an isomorphism on \(H^1\). On the inner \(S^3\) there is no
\(H^1\) gauge obstruction. Both normal bundles extend after choosing
their actual boundary framings; any resulting normal Pontryagin class
is retained. The extra sphere-bundle boundary has index zero, since its
base is (SC.16), \(H^2(S^4)=0\), and \(\sigma(S^4)=0\). The sphere-
bundle formula cancels the normal Pontryagin class and gives this zero
without a normal-framing assumption.

\item The orientations are the same two reversals as K7. The complex
orientation on \(W\) has outer boundary \(F\times S^1_\rho\) and
inner boundary the negative of the complex-domain link. The outer
boundary of \(Y\times D_\rho\) is \(-P_f\), and the excised inner
seven-ball supplies \(-N_e\). The remaining normal-sphere boundary
has zero index. Hence
\[
                         A(P_f)=-A(N_e).
\]
The actual positive family is Jin's mapping torus for \(f^{-1}\), so
\[
                         A_{U,\rm Jin}(f)=A_U(N_e).      \tag{SC.17}
\]

\item For an actual Seifert embedding of the lifted link in its six-sphere,
repeat the identical normal-tube excision, cap its boundary by \(D^4\),
and use the complex normal two-plane plus the radial real line. The
meridian line on the surgery boundary is the one in (SC.14); on the
three-sphere a line or its gauge has no additional class. The sphere-
bundle formula gives
\[
 A_U(N_e)=\frac{e(\nu V)^2-\sigma(V)}8
                      =H(e),                        \tag{SC.18}
\]
where the last equality is Takase's corrected Haefliger formula,
recorded in P3F.5. The normalization comparison with the invariant
computed directly in P3F.55 is its equations (CL.16)--(CL.19): an
actual punctured Kummer surface embedded in the oriented five-sphere
has signature \(-16\), trivial normal Euler class and value two
for both invariants. The exact Whitehead-to-Stiefel morphism and the
sign-refined immersed Seifert formula give that equality on this
nonzero element. The two group isomorphisms to \(\mathbb Z\)
therefore agree on every embedding. This supplies the sign bridge
used here.

\item The displayed local Seifert computation uses the meridian-one U
line. Its raw complex-plus-height orientation and the physical orbit
orientation are opposite; the outward-meridian-one line reverses at
the same time, so its U index is the displayed H(e)=1. The actual
nonzero-vector physical line M is U inverse. Therefore its odd
characteristic index in the same oriented Jin manifold is the negative.
The exact spinor arrow is anti-linear, not the previously written
linear swap: Cartier coefficients (a,b) map to g(x,y) overline(a,b)
in the right-complex Cox spinor. The globally defined quaternionic
anti-linear map commutes with all SU2 clutches and extends over both
cap products; its normal antipodal map and the dual circle character
are retained. P3K.68 gives every unit and both cap extensions.

\item The physical Hopf transfer has positive-line base -S2\_z and y'=-y,
so its first and second pieces are -R,-L in the K23 plumbing order.
Its first-to-second clutch r maps to F(r); rewriting with L first
gives minus Sigma\_(F inverse). The removed ball has the original
complex six-sphere as its outward boundary, and its free-face map is
the actual physical return, with the identity on the fixed tube.
Thus this is the actual clutched sphere. The meridian-one U transfer
instead has pieces +R,+L and is its orientation reverse. There is
no extra inverse after this explicit piece exchange.

\item Consequently the direct local calculation H(e)=1 gives
\[
 A_{U,\rm Jin}(f)=1,\qquad A_{M,\rm Jin}(f)=-1,
 \qquad[\Phi]=-1\in\Theta_7\cong\mathbb Z/28.          \tag{SC.19}
\]
The last equality is the physical-line oriented Hopf transfer of P3F.22, with the exact dual and piece maps of P3K.68. For the
forward physical return \(\Phi\), P3F.24 gives the literal cut-and-paste
identification
\[
 M_\Phi\cong(S^1\times S^6)\mathbin{\#}\Sigma_{-1}.
                                                               \tag{SC.20}
\]
Both spin structures have Eells--Kuiper invariant \(-1/28\), whereas
both spin structures of the product have value zero. Orientation
reversal changes the sign and still does not give zero. Consequently
this mapping torus is not diffeomorphic to the product, even without
an orientation requirement.

\item The loop in (SC.19) is one positive \(\rho\)-loop. Since
\(a'(u_*)=1-2u_*^{-3}=3\), the equation \(a(u)=\rho^2\) makes it
exactly two positive turns of the original marked source coordinate
\(u-u_*\). The scalar and all finite-cap choices stay fixed.

\item There is also an obstruction to smooth descent of this entire closed
family. Let
\[
 p:D_\rho^*\longrightarrow D_a^*,\qquad a=\rho^2,
                                                               \tag{SC.21}
\]
with \(u\) the analytic inverse of \(a(u)\) near \(u_*\). An oriented
smooth sphere bundle on \(D_a^*\) is obtained by cutting an annulus
along a radial interval and gluing its ends by an orientation-preserving
diffeomorphism \(\Phi_a\). Pulling this exact clutch construction back
along (SC.21) concatenates two copies of the interval; its closing map
is \(\Phi_a^2\). A change of fibre marking conjugates this map.
In \(\pi_0\operatorname{Diff}^+(S^6)\cong\Theta_7\) conjugation is
trivial because the group is abelian. Therefore an oriented smooth
bundle isomorphism from the actual \(\rho\)-family to such a pullback
would imply
\[
                -1=[\Phi]=2[\Phi_a]\quad\hbox{in }\mathbb Z/28.
                                                               \tag{SC.22}
\]
Reduction modulo two contradicts (SC.22). Thus this particular full
\(\rho\)-family does not descend as an oriented smooth bundle through
(SC.21). This does not exclude all conceivable one-turn endpoint maps:
it excludes exactly those that would give a descent of this specified
closed family.

\item No value for a different toric fan or another source-moduli puncture
is assigned by this proof.

\item Exact symbolic identities and the integral plumbing certificate are
saved separately in \path{code/part3/source_cusp_trace.py} and its outputs.

\end{enumerate}
\end{proof}


\section{The free characteristic calculation}



\begin{theorem}\label{night:P3R.1}
\StatementRecord{night:P3R.1}{proved}{S6\_06\_BLIND: GPT-6 Astra (OpenAI, ChatGPT Pro mode, effort Pro), recorded re-derivation in the checkpoint dated 18 November 2026}{Definitions and the displayed constructions}
Let $W$ be a compact almost complex eight-manifold whose
outward-oriented boundary is the homotopy sphere $\Sigma$.
In the presentation $\Theta_7=\Z/28$ fixed on the first page,
\[
 [\Sigma]=
 \frac{-c_1^4+4c_1^2c_2+3c_2^2-3\sigma(W)}{24}
 \pmod {28}.
\]
All characteristic products are evaluated on $[W,\partial W]$
using their unique relative rational lifts. A spin assumption
on $W$ is unnecessary.
\end{theorem}
\begin{proof}
\begin{enumerate}[label=\arabic*.]
\item The almost complex structure induces a spin-c structure
with determinant class $c_1=c_1(TW)$ and the complex
orientation. The homotopy sphere has $H^1=H^2=H^3=H^4=0$
over $\Q$ and a unique spin structure. Thus the boundary
determinant is trivial, the induced boundary spin structure
is the required one, and the exact sequences of the pair
give unique relative rational lifts in degrees two and four.
The spin-c Eells--Kuiper formula in \cite[Section~3.1(v)]{Jiang2014},
with these hypotheses, is
\[
 \mu(\Sigma)=
 \frac{p_1^2-4\sigma(W)}{896}
 -\frac{c_1^2p_1}{192}+\frac{c_1^4}{384}
 \pmod\Z.
\]
The formula is for the outward boundary, as used here.
For the underlying real tangent bundle of a complex bundle,
$p_1=c_1^2-2c_2$. Multiplication by $28$ and substitution
give
\[
 \begin{split}
 28\mu(\Sigma)
 &=\frac{3p_1^2-14c_1^2p_1+7c_1^4-12\sigma(W)}{96}\\
 &=\frac{-c_1^4+4c_1^2c_2+3c_2^2-3\sigma(W)}{24}
 \pmod {28}.
 \end{split}
\]
This is precisely the specified presentation of $\Theta_7$.
The substitution is independently executed in
\texttt{code/root\_comparisons/characteristics.py}.

\end{enumerate}
\end{proof}




\begin{theorem}\label{night:P3R.2}
\StatementRecord{night:P3R.2}{proved}{S6\_06\_BLIND: GPT-6 Astra (OpenAI, ChatGPT Pro mode, effort Pro), recorded re-derivation in the checkpoint dated 18 November 2026}{\cref{night:P3K.7}}
The equation $6a=b$ in $\Z/28$ has no solution for odd $b$.
For $b=2r$, its solutions are exactly
$a=5r$ and $a=5r+14$ modulo $28$.
Thus the sixfold class alone leaves precisely these two
possibilities for the original class; absent a value of $b$,
this calculation excludes no original class.
\end{theorem}
\begin{proof}
\begin{enumerate}[label=\arabic*.]
\item Reduction modulo two gives the first assertion. For even
$b$, division by two identifies the equation with
$3a=r$ in $\Z/14$. The inverse of $3$ in that ring is
$5$, since $3\cdot5=1+14$. Its unique solution modulo
$14$ has the two displayed lifts modulo $28$.
Conversely each lift satisfies the original equation.
Every original class occurs for its own value $b=6a$.
The attached program verifies the complete table of
twenty-eight target residues. The integral index used
in P3K.7--8 supplies the additional information that
settles the actual original family.

\end{enumerate}
\end{proof}




\begin{theorem}\label{night:P3K.1}
\StatementRecord{night:P3K.1}{proved}{S6\_06\_BLIND: GPT-6 Astra (OpenAI, ChatGPT Pro mode, effort Pro), recorded re-derivation in the checkpoint dated 18 November 2026}{Definitions and the displayed constructions}
Let

\[
 R=\mathbf C[u,w,X,v]/(uv-Xw),\qquad Z=\operatorname{Spec}R.
\]

Use the small resolution with charts

\[
 U_0=(u,w,z),\quad X=uz,\ v=wz;
 \qquad U_1=(X,v,t),\quad u=Xt,\ w=vt,\quad t=z^{-1}.
\]

Its exceptional curve is the zero section of
\(\mathcal O_{\mathbf P^1}(-1)^{\oplus2}\).  Write
\(L=\pi_{\mathbf P^1}^*\mathcal O(-1)\), with frames
\(e_1=z^{-1}e_0\).  Then the resolution map \(\rho\) satisfies

\[
 \rho_*L\cong I=(u,w),\qquad R^j\rho_*L=0\quad(j>0).
\]

Nevertheless \(I\) is a nonperfect rank-one module at the conifold
point.  The restriction of \(L\) to the link has primitive first Chern
class; it does not extend as a line bundle over the conifold smoothing.
\end{theorem}

\begin{proof}
\begin{enumerate}[label=\arabic*.]
\item The two affine charts and their affine intersection compute
coherent cohomology.  A Laurent monomial on the intersection has the
form \(u^a w^b z^c e_0\), with \(a,b\geq0\).  It extends over \(U_0\)
when \(c\geq0\).  A monomial from \(U_1\) is

\[
 X^a v^b t^k e_1=u^a w^b z^{a+b-k-1}e_0,
\]

so it extends over that chart precisely for \(c\leq a+b-1\).
Those two integer intervals cover \(\mathbf Z\), including the case
\(a+b=0\).  Thus the first Cech cohomology vanishes.  Higher cohomology
vanishes because the cover has two affine opens and affine intersection.
The intersection of the two spaces of sections is spanned by the
monomials with \(0\leq c\leq a+b-1\).  Each has a factor \(u\) or \(w\)
without a factor \(z\); its remaining factor belongs to the subring
generated by \(u,w,uz,wz\), namely \(R\).  Hence the global sections are
generated by \(u e_0,w e_0\), which on the other chart are \(X e_1,v e_1\).
Sending these generators to \(u,w\) identifies the module with \(I\).
Since the target is affine, these cohomology calculations give all the
stated direct images.

\item The exact presentation and its matrix factorization are

\[
 R^2\xrightarrow{\Phi}R^2\xrightarrow{[u\ w]}I\longrightarrow0,
 \qquad
 \Phi=\begin{pmatrix}w&v\\-u&-X\end{pmatrix},\quad
 \Psi=\begin{pmatrix}-X&-v\\u&w\end{pmatrix},
\]

\[
 \Phi\Psi=\Psi\Phi=(uv-Xw)\,1_2.
\]

\item To check the presentation without appealing to a classification,
lift a relation \(au+bw=0\) to the polynomial ring.  Write
\(au+bw=c(uv-Xw)\).  Then
\(u(a-cv)+w(b+cX)=0\).  The coprimality of \(u,w\) gives
\((a,b)=d(w,-u)+c(v,-X)\), exactly the two displayed columns.
For exactness of the periodic tail, lift an element in
\(\ker\Phi\) to a polynomial column \(z\), so that
\(\Phi z=(uv-Xw)y\).  Multiplication by \(\Psi\) and cancellation in
the polynomial domain gives \(z=\Psi y\).  The other kernel is treated
by the same calculation with the matrices reversed.  Consequently

\[
 \cdots\xrightarrow{\Phi}R^2\xrightarrow{\Psi}R^2
 \xrightarrow{\Phi}R^2\longrightarrow I\longrightarrow0
\]

is an infinite free resolution.  At the maximal ideal
\(\mathfrak m=(u,w,X,v)\), every matrix entry lies in \(\mathfrak m\).
After tensoring with the residue field all differentials vanish and
all positive Tor groups have dimension two.  Thus \(I\) has no finite
locally free resolution.  In particular it has not become a zero
neighborhood class merely because \(H^*(\mathbf P^1,\mathcal O(-1))=0\).
It has generic rank one and \(\dim I/\mathfrak m I=2\).

\item The link of the resolution is the unit sphere bundle of the underlying
real rank-four bundle \(\mathcal O(-1)^{\oplus2}\) on \(S^2\).  Oriented
rank-four bundles on \(S^2\) are classified by
\(\pi_1(SO(4))=\mathbf Z/2\); the class here is
\(c_1\bmod2=(-2)\bmod2=0\).  The link is therefore \(S^2\times S^3\),
and projection to the exceptional \(S^2\) induces an isomorphism in
degree-two cohomology.  The restricted line has \(c_1=-h\), so its
rank-reduced K-class is the primitive generator of
\(\widetilde K^0(S^2\times S^3)=\mathbf Z\).  The smoothing is
\(S^3\times D^3\), whose second integral cohomology vanishes.  The
restriction of a line from that smoothing therefore has zero first
Chern class.  This proves the claimed obstruction to extending this
particular line.  Reversing the chosen Euler-line convention replaces
\(L\) by its dual and changes the generator's sign, without removing
the obstruction.  This establishes only the stated local obstruction;
it does not decide the actual return index.

\end{enumerate}
\end{proof}


All polynomial identities are independently checked in
\path{code/part3_fourfold/conifold_module.py} and exported in the adjacent
JSON file.  The monomial bounds in the proof are symbolic; the finite
loop in the code is an additional sanity check, not its justification.

\begin{theorem}\label{night:P3K.2}
\StatementRecord{night:P3K.2}{proved}{S6\_06\_BLIND: GPT-6 Astra (OpenAI, ChatGPT Pro mode, effort Pro), recorded re-derivation in the checkpoint dated 18 November 2026}{Definitions and the displayed constructions}
Put \(B=\mathbf C[u,v]\),
\(R'=B[x,y]/(xy-uv)\), and \(J=(x,u)\subset R'\).
The module \(J\) is flat over \(B\).  Its fibre at \((u,v)=(0,0)\) is
the normalization module

\[
 J\otimes_B\mathbf C\cong
 \mathbf C[x]\oplus\mathbf C[y]
 \quad\text{over}\quad\mathbf C[x,y]/(xy).
\]

For an integral projective nodal genus-one curve \(C\) with
normalization \(\nu:\mathbf P^1\to C\), the rank-one torsion-free
degree-zero sheaf \(\nu_*\mathcal O_{\mathbf P^1}(-1)\) has no
cohomology.  Thus the local magnetic module has an exact realization
inside a flat family of torsion-free nodal sheaves.  This supplies a
candidate coefficient object for index excision.  The local lemma
does not by itself compare the actual return symbol; P3K.6–7 below
provide that comparison by a real bordism.
\end{theorem}

\begin{proof}
\begin{enumerate}[label=\arabic*.]
\item As a \(B\)-module, \(R'\) has the free monomial basis
\(1,x,x^2,\ldots,y,y^2,\ldots\), by reducing each mixed product using
\(xy=uv\).  The submodule \(J\) has the \(B\)-basis

\[
 \{x^n:n\geq1\}\ \cup\ \{u y^n:n\geq0\}.
\]

\item Spanning follows from \(xy^n=uvy^{n-1}\); independence follows from
the basis of \(R'\) and the fact that \(u\) is a nonzero divisor of
\(B\).  Hence \(J\) is free, in particular flat, over \(B\).
It is presented by the matrix
\(\left(\begin{smallmatrix}u&y\\-x&-v\end{smallmatrix}\right)\).
At \(u=v=0\) the relations become \(x e_2=0\), \(y e_1=0\), proving
the displayed normalization module.  The normalization map is finite,
so the cohomology of its direct image equals that of \(\mathcal O(-1)\),
which vanishes in both degrees.  Since \(\chi(\mathcal O_C)=0\), the
degree \(\chi(\nu_*\mathcal O(-1))-\chi(\mathcal O_C)\) is zero.
Its generic rank is one and it has no subsheaf supported at finitely
many points, so it is rank-one torsion-free.

\item There is also an exact connection to the local compactified-Picard
kernel.  On the fibre product of two nodal families, written
\(xy=UV\), the diagonal ideal is \((x-U,y-V)\).  Set
\(a=x-U\), \(b=y-V\).  The equation becomes
\(a(V+b)+Ub=0\).  With coordinates
\((x',u',y',v')=(a,b,V+b,-U)\), it is
\(x'y'=u'v'\), and the diagonal ideal is exactly \((x',u')\).
The kernel
\(I_\Delta\otimes\pi_1^*\mathcal O(O)\otimes\pi_2^*\mathcal O(O)\)
therefore has precisely this local matrix module, since the section
lines are trivial near the two nodes.

\end{enumerate}
\end{proof}


For the global algebraic kernel construction, Burban–Kreussler,
*On a relative Fourier–Mukai transform on genus one fibrations*,
arXiv:math/0410349v2, equation (1), pp. 5–7, constructs this sheaf flat
over both factors for a flat projective family of integral Gorenstein
genus-one curves with a section in its smooth locus.  Theorem 2.12
proves the resulting derived equivalence.  The theorem concerns that
specified algebraic family and kernel.  The required identification
with the present circle quotient's symbol still needs its own map.

\begin{theorem}\label{night:P3K.3}
\StatementRecord{night:P3K.3}{proved}{S6\_06\_BLIND: GPT-6 Astra (OpenAI, ChatGPT Pro mode, effort Pro), recorded re-derivation in the checkpoint dated 18 November 2026}{Definitions and the displayed constructions}
Let \(B=S^3\times S^1\), let \(K\subset S^3\) be the cusp circle,
and set \(T=K\times S^1\).  Give the trivial real normal two-plane
bundle of \(T\) its oriented complex-line structure.  Its K-theory
Thom map and extension by zero give

\[
 i_!:K^0(T)\cong K^0_c(N_T)\longrightarrow K^0(B).
\]

In the rank/degree basis of \(K^0(T)=\mathbf Z^2\), this map kills
the rank generator and maps the degree-one reduced line class to the
primitive top generator of \(\widetilde K^0(B)=\mathbf Z\).  Every
integer is consequently realized by a K-class supported on the cusp
torus and vanishing on its complement.  The spin index on \(B\)
reads exactly this integer.
\end{theorem}

\begin{proof}
\begin{enumerate}[label=\arabic*.]
\item The oriented normal bundle of a knot in an oriented
three-manifold is trivial: its Euler class lies in \(H^2(S^1)=0\).
The same trivialization pulls back to \(T\).  The K-theory Kunneth
formula, or the cellular filtration and Bott periodicity, gives

\[
 K^0(T)=\mathbf Z\{1\}\oplus\mathbf Z\{\beta\},\quad
 \operatorname{ch}(\beta)=[T]^*,
\]

where \(\beta=[H]-1\) for a line \(H\) of degree one, and

\[
 K^0(B)=\mathbf Z\{1\}\oplus\mathbf Z\{\omega\},\quad
 \operatorname{ch}(\omega)=[B]^*.
\]

\item The normal Todd class is one, so the Chern character of the Thom
pushforward is the ordinary cohomological Gysin map.  For the rank
generator, \(\operatorname{ch}(i_!1)=\operatorname{PD}(T)=0\), because
\(H^2(B;\mathbf Z)=0\).  Both K-groups are torsion-free, so \(i_!1=0\).
For \(\beta\), the Gysin image integrates to one in degree four;
hence \(i_!\beta=\omega\).  Extension by zero makes its restriction
to \(B\setminus T\) zero.  Finally \(B\) is parallelizable, so its
spin index is \(\int_B\operatorname{ch}_2\), which sends \(\omega\)
to one.

\end{enumerate}
\end{proof}


This proves an explicit limitation of the inference from acyclicity
on regular elliptic fibres: even a known pushforward on the complement
does not choose the extension's top K-class.  It is an actual
nonzero image under a specified Thom/Gysin map, rather than an
unsupported assertion that the local and global calculations are
unrelated.



The following records the precise target of the initial algebraic
index-excision attack.  P3K.6–7 subsequently construct the required
actual comparison by a real trace and a sphere-bundle bordism.

For the actual conifold replacement mapping torus \(Q^6\), write
\(L_\delta\) for its circle Euler line.  The required integer is

\[
 A=\operatorname{ind}(D_{Q,\mathrm{spin}}\otimes L_\delta).
\]

On the ordinary elliptic pieces, the coefficient line has Picard point
\(P\), and the corresponding degree-zero Dolbeault complexes are
acyclic.  On the compact fixed sheet, P3F.12 supplies
\(R f_*(M+nF)=0\).  P3F.14 supplies the exact local transfer to that
complex.  P3K.1 shows why the conifold neighborhood cannot be discarded
merely from the fibre \(\mathcal O(-1)\)'s cohomology.  P3K.2 gives an
explicit flat torsion-free alternative retaining its link module.

The mathematical object sought at this stage was a relative symbol comparison
between the actual spin symbol twisted by \(L_\delta\) and this
nodal Dolbeault complex, with its two outer trivializations.  It must
identify the primitive meridian class and the degree \(-12\) angular
gauge in P3F.13–15, and must account for the added source-resolution
curves on the outer boundary.  None of P3K.1–3 constructs that relative
symbol map.  The universal Fourier–Mukai theorem supplies the coherent
kernel once a family and pullback map are given; its hypotheses do
not themselves provide the map from this orbit family.

The unsuccessful algebraic attack used the candidate outer morphism
given by the minimal source
resolution of \(UV=Q^{12}\), followed by
\((U,V,Q)\mapsto(U,V,q=Q^{12})\), crossed with the nondegenerate second
direction \(r\ne0\), and then the labelled cyclic descent of P3F.15.
At that stage the descended outer quotient had not been identified
with a proper equivariant modification of the original smooth
family.  P3K.4 supplies the precise multiplicity obstruction to
that proposed map.  The real trace P3K.6–8 supplies the actual
return comparison without using that failed algebraic extension.

\begin{theorem}\label{night:P3K.4}
\StatementRecord{night:P3K.4}{proved}{S6\_06\_BLIND: GPT-6 Astra (OpenAI, ChatGPT Pro mode, effort Pro), recorded re-derivation in the checkpoint dated 18 November 2026}{\cref{night:P2.7}, \cref{night:P3F.13}}
The free twelve-label construction of P3F.15, when applied to a
reduced \(Q\)-central divisor with no stabilizers at its generic
points, has multiplicity twelve over \(q=Q^{12}\) on every image
component.  A normal space with all those multiplicities equal to
twelve admits no proper bimeromorphic map over the \(q\)-disc to the
actual reduced cusp model.  Thus a common proper toric modification
of these two outer extensions cannot be inferred from their
punctured-family comparison.
\end{theorem}

\begin{proof}
\begin{enumerate}[label=\arabic*.]
\item A finite quotient by an action with trivial stabilizer is
etale near the generic point under consideration.  For a reduced
component \(D\) of \(Q=0\), its divisorial valuation gives
\(\operatorname{ord}_D(q)=\operatorname{ord}_D(Q^{12})=12\).
The ramification index of the quotient at \(D\) is one, so the order
on its image is also twelve.  A proper bimeromorphic morphism between
normal spaces is an isomorphism at the generic point of the strict
transform of any target prime divisor.  It therefore preserves that
divisor's valuation of \(q\).  A target component of order one would
give a source component of order one, contradicting the displayed
orders.  This proves the assertion about proper modifications.

\item The exact origin itself detects why the two extensions differ.
For the original default twists, P2.7 gives

\[
 F(s,v)=\left(\frac{\mu(s+12v)}2,
 \frac{C(s+12v)}{12}-\frac{(s+12v)^2}{24}+6v\right),
 \qquad \gamma(s,v)=(6\mu,C-s-12v).
\]

\item Here \(\mu,C\) are unchanged by the integer cusp deck.  Direct
subtraction, with no limiting argument, yields

\[
 F(s+1,v)-F(s,v)=\frac{\gamma(s,v)}{12}-\frac{\delta}{24},
 \qquad
 F(s+12,v)-F(s,v)=\gamma(s,v)-6\delta.
\]

\item Thus the partial-origin one-step deck is the fractional gamma
translation that shifts the twelve-cycle.  It is not the reduced
cusp extension furnished by the extending cusp origin.  The exact
identities are exported by \path{code/part3_fourfold/outer_descent.py}.
They do not contradict the angular clutch in P3F.13, nor the partial
filling.  They prevent identifying the freely descended outer model
with the original reduced cusp by the proposed proper-map shortcut.

\end{enumerate}
\end{proof}




\begin{theorem}\label{night:P3K.5}
\StatementRecord{night:P3K.5}{proved}{S6\_06\_BLIND: GPT-6 Astra (OpenAI, ChatGPT Pro mode, effort Pro), recorded re-derivation in the checkpoint dated 18 November 2026}{Definitions and the displayed constructions}
The integral connecting class from the marked boundary \(T^3\) has square zero and pairs trivially with every extension of the page line.  In particular an odd cap correction does not change the normal Euler square.
\end{theorem}

\begin{proof}
\begin{enumerate}[label=\arabic*.]
\item Let \(V=V_0\) be the actual regular-binding page, with its marked
boundary \(T^3\), and let \(m_V=c_1(M)\).  Its boundary restriction is
zero.  For the inverse zero-framed Borromean surgery trace \(H\),
\(H^2(H;\mathbf Z)\to H^2(T^3;\mathbf Z)\) is an isomorphism: the
three-handle intersection matrix is zero, and the relevant
Poincare–Lefschetz exact sequence has last term \(H_1(H)=0\).
Consequently any class on \(V\cup H\) that restricts to \(2m_V\) on
\(V\) and to zero on the common boundary restricts to zero on \(H\).
Choose an integral extension \(m\) of \(m_V\), using a boundary
trivialization and the trivial line on \(H\).  The difference between
that class and \(2m\) lies in the image of
\(\delta:H^1(T^3;\mathbf Z)\to H^2(V\cup H;\mathbf Z)\).
For \(k=\delta g\), the relative cup-product formula gives

\[
 k^2=\delta(g\smile k|_{T^3})=0,\qquad
 km=\delta(g\smile m|_{T^3})=0.
\]

\item Thus its square is always \(4m_{V,\mathrm{rel}}^2\), even when the
gluing correction is odd.  The odd correction is relevant here:
\(V\subset S^5\) is spin and has signature \(-8\), so its boundary
spin structure has Rokhlin invariant eight.  The zero-framed
Borromean trace is spin of signature zero and induces Rokhlin
invariant zero.  Those spin structures cannot be identified.  This
does not change the square calculation, which retains the integral
connecting class rather than setting it to zero.

\item The following real trace construction uses a different cap and does
not require an embedding of this Borromean trace in the six-sphere.

\end{enumerate}
\end{proof}




\begin{theorem}\label{night:P3K.7}
\StatementRecord{night:P3K.7}{proved}{S6\_06\_BLIND: GPT-6 Astra (OpenAI, ChatGPT Pro mode, effort Pro), recorded re-derivation in the checkpoint dated 18 November 2026}{\cref{night:P3F.10}, \cref{night:P3F.11}, \cref{night:P3F.4}, \cref{night:P3F.7}, \cref{night:P3F.8}, \cref{night:P3K.5}, \cref{night:P3K.6}}
The real trace gives a spin bordism carrying the Euler line from the actual covered conifold return to a normal sphere bundle.  Its index is zero.  The original return also has index zero, and its actual Hopf-lift gluing is the standard sphere.
\end{theorem}

\begin{proof}
\begin{enumerate}[label=\arabic*.]
\item Cap \(X\) along \(F\times S^1\) with \(C=D^3\times S^1\), using
the fixed normal framing chosen for the actual conifold replacement.
Denote the resulting closed four-manifold by \(\bar X\).  Its cap
relative to \(V\) is

\[
 H_a=(W\cup_FD^3)\times S^1
      =D^2_a\times S^1_b\times S^1_\gamma.
\]

\item The subscripts retain the marked boundary cycles.  This cap kills
\(a\).  The usual elliptic cap kills \(\gamma\); their boundary
markings are not silently equated.  Both have signature zero.

\item Delete the normal disc bundle of \(X\) from \(Z\), and glue
\((C\times S^2)\times I\) along the resulting boundary collar.
The resulting two closed boundary components are the physical
positive return family \(P_f\) and the normal sphere bundle
\(B=S(\bar\nu)\to\bar X\).  On the outer side this is precisely
the replacement of the fixed \(F\times D^3\) by \(D^3\times S^2\),
performed in every parameter fibre.  On the inner side it completes
the normal sphere bundle over the cap.  We now specify its two
boundary orientations and the return convention explicitly.

\item Let the circle-first quotient convention be
\(\Omega_{X_6}=d\delta\wedge\Omega_Y\).  In the fixed normal
complex two-plane the weights are \((-1,+1)\).  Conjugating the
negative-weight complex coordinate reverses its real orientation;
contracting the positive Hopf volume by the circle vector reverses
the outward cone convention once more.  The two signs give
\(\Omega_Y=\Omega_F\wedge\Omega_{\nu,\mathrm{Hopf}}\), where
the meridian sphere has the positive Hopf orientation.  The graph
part of \(X\) has the complex page orientation, and the oriented
compression has \(\partial X=F\times S^1_\theta\), in that order.
Give \(Z=Y\times D^2\) its product orientation.  Since \(\dim Y=5\),
the outward boundary orientation is

\[
 \Omega_{\partial Z}=-\Omega_Y\wedge d\theta.
\]

\item The inner boundary of the deleted normal disc bundle is oriented
\(-\Omega_{\bar X}\wedge\Omega_{S^2}\).  The product cap preserves
these signs.  Thus the actual oriented bordism has boundary
\(-P_f\sqcup(-B)\), with \(P_f\) oriented by the fibre followed
by positive parameter, and \(B\) by the complex page followed
by the positive meridian sphere.

\item Jin's mapping torus \(M_f\) identifies \((z,0)\sim(f(z),1)\).
The physical positive family of the return \(f\), which identifies
\((z,1)\sim(f(z),0)\), is therefore \(P_f=M_{f^{-1}}\).
Characteristic-number additivity under inversion, together with
the two boundary signs just computed, gives

\[
 (k,\ell,A)(P_f)=-(k,\ell,A)(B),\qquad
 (k,\ell,A)_{\mathrm{Jin}}(f)=(k,\ell,A)(B).
\]

\item This records both reversals.  The following integrations on \(B\)
therefore compute Jin's invariants of the specified physical
positive return \(f\), with the stated positive meridian class.

\item The Euler line extends over the bordism.  On
\(F\times S^1\times S^2\) its Chern class is the meridian generator,
by the positive-weight SU(2) trivialization of P3F.7.  It therefore
extends over \(D^3\times S^1\times S^2\).  An ambiguity in the
boundary bundle isomorphism is a map to \(S^1\); its homotopy class
lies in \(H^1(S^2\times S^1\times S^2)=\mathbf Z\), and the
restriction from \(H^1(D^3\times S^1\times S^2)\) is an isomorphism.
It extends as well.  The bordism is spin: \(Y\times D^2\) is spin,
and the two possible spin structures on the gluing boundary differ
only on the same \(S^1\) factor, where either extends over the cap.
This identifies the actual spin return symbol and Euler line with
the inner boundary symbol by spin bordism, with no analytic corner
or contraction of exceptional curves.

\item Write \(e_v\) for the vertical Euler class of \(S(\bar\nu)\), and
define \(b\in H^2(\bar X;\mathbf Z)\) by

\[
 e_v=2x-\pi^*b.
\]

\item Existence and uniqueness follow from the sphere-bundle Gysin
sequence, since both classes have fibre integral two.  Since the
normal bundle of the graph page is framed and its chosen section
has \(s^*e_v=0\), \cref{night:P3K.68} gives \(b|_V=-2c_1(M)\).  On the product
cap \(C\), \(b=0\).  Its restriction on the intermediate product
compression is zero as well: \(H^2(W\times S^1)\) injects into the
cohomology of its two boundary components.  This is seen from
\(H_1(W)=\mathbf Z\), \(H_2(W)=\mathbf Z\), and the boundary
torus and sphere, whose products with the circle generate both
classes of \(H_2(W\times S^1)\).  On the two boundaries the relevant
restriction is zero.  Equivalently,
\(H^2(H_a)\to H^2(T^3)\) is injective, so the vanishing boundary
restriction of \(b\) gives its vanishing on this whole cap.

\item Extend \(M\) trivially across \(H_a\), with class \(m\).  The
Mayer–Vietoris difference \(b+2m\) belongs to
\(\delta H^1(T^3;\mathbf Z)\).  As in P3K.5, its square and its
product with \(m\) vanish.  Therefore

\[
 \langle b^2,[\bar X]\rangle=4m_{V,\mathrm{rel}}^2=-8,
 \qquad
 \sigma(\bar X)=\sigma(V)=-8,
 \qquad
 \langle p_1(T\bar X),[\bar X]\rangle=-24.
\]

\item These use the actual page invariants P3F.10.  No identification of
the two different marked caps is needed.

\item Finally \(T_{\mathrm{vert}}B\oplus\mathbf R\cong\pi^*\bar\nu\),
so \(e_v^2=\pi^*p_1(\bar\nu)\).  With \(\pi_*x=1\), this gives

\[
 \pi_*x^2=b,\qquad
 \pi_*x^3=\frac{p_1(\bar\nu)+3b^2}{4},\qquad
 p_1(TB)=\pi^*(p_1(T\bar X)+p_1(\bar\nu)).
\]

\item The actual spin return index is consequently

\[
 A(f)=\frac{4\int_Bx^3-\int_Bp_1(TB)x}{24}
     =\frac{3\int_{\bar X}b^2-\int_{\bar X}p_1(T\bar X)}{24}
     =\frac{3(-8)-(-24)}{24}=0.
\]

\item The normal Pontryagin class cancels.  Thus no unproved normal
framing or spin claim about the elliptic cap enters the result.
For the original parameter return, P3F.11 states that this covered
return has index six times the original integer.  Since that
integer is integral, the original index is also zero.  The actual
Hopf-lift gluing P3F.8 and the extending \(a\)-powers P3F.4 then
identify the clutched sphere as standard.

\end{enumerate}
\end{proof}




\begin{theorem}\label{night:P3K.8}
\StatementRecord{night:P3K.8}{proved}{S6\_06\_BLIND: GPT-6 Astra (OpenAI, ChatGPT Pro mode, effort Pro), recorded re-derivation in the checkpoint dated 18 November 2026}{\cref{night:P3F.11}, \cref{night:P3F.12}, \cref{night:P3K.6}, \cref{night:P3K.7}}
The same actual return index vanishes by a second calculation: vertical sphere quantization, an explicit marked cap switch, and acyclicity of the elliptic Picard line.  This calculation does not evaluate the page signature, its line square, or its first Pontryagin number.
\end{theorem}

\begin{proof}
\begin{enumerate}[label=\arabic*.]
\item The actual bordism of P3K.6–7 also gives an analytic computation of
\(A\) that does not evaluate the signature of the page, the square
of its line, or a Pontryagin number.  Keep the inner sphere bundle
\(\pi:B=S(\bar\nu)\to\bar X\) and the integral class
\(e_v=2x-\pi^*b\).

\par\medskip\noindent\textbf{Vertical quantization.}\enspace The spin Dirac symbol on \(B\),
twisted by the Euler line of class \(x\), is a spin-c symbol of
determinant \(2x\).  The oriented vertical two-plane has its
canonical complex spin-c symbol, of determinant \(e_v\).  Dividing
these two spin-c structures gives the pullback of a base spin-c
structure of determinant \(b\).  This follows directly from the
stable tangent splitting into vertical and horizontal bundles and
the displayed determinant identity.  Fibrewise the vertical symbol
is the Dolbeault symbol of \(\mathcal O_{\mathbf P^1}\), not that of
\(\mathcal O(-1)\).  Its kernel consists of constant functions and
its cokernel is zero.  The constants are fixed by every oriented
sphere transition, so its family index is the trivial line, with
its actual transition functions.  This symbol composition can be
verified directly at the operator level.  Give the fibres their
round metrics and use an \(SO(3)\) connection on \(\bar\nu\).
The vertical Dolbeault Dirac operator \(D_v\) is invariant under
this structure group, so the horizontal connection commutes with
\(D_v\).  Its squared restriction orthogonal to the constants
has a positive uniform lower bound.  If \(\epsilon_v\) is the
vertical grading and \(D_h\) the horizontal base spin-c Dirac
operator on the bundle of vertical spinors, the operator

\[
 D_t=tD_v+\epsilon_vD_h,\qquad t>0,
 \qquad D_t^2=t^2D_v^2+D_h^2
\]

is invertible on that orthogonal complement.  The square identity
uses \(D_v\epsilon_v=-\epsilon_vD_v\) and the commuting invariant
connection.  On the constant vertical kernel it is precisely the
base Dirac operator.  The displayed operator has the same elliptic
symbol as the twisted total-space spin Dirac operator after a
positive rescaling of the vertical metric; zeroth-order terms do
not change its index.  Consequently

\[
 A=\operatorname{ind}(D_{\bar X,b}).
\]

\item This is a factorization of the specified actual symbol.  It keeps
the meridian Dirac index one that is essential at the conifold
replacement.

\par\medskip\noindent\textbf{Switch the marked cap.}\enspace The compression cap is
\(H_a=D^2_a\times S^1_b\times S^1_\gamma\), whereas the regular
elliptic cap is \(H_\gamma=S^1_a\times S^1_b\times D^2_\gamma\).
Use their common marked \(T^3\) boundary.  Their closed difference
is

\[
 H_a\cup_{T^3}(-H_\gamma)
 =S^1_b\times
   \bigl((D^2_a\times S^1_\gamma)
       \cup_{T^2}(S^1_a\times D^2_\gamma)\bigr)
 \cong S^1\times S^3.
\]

\item The indicated meridian-longitude exchange is the usual two-solid-
torus decomposition of \(S^3\); its meridian intersection is one.
Both cap determinant classes vanish and their boundary spin-c
structures agree: on \(T^3\), a spin-c structure with determinant
class zero is unique because \(H^2(T^3;\mathbf Z)\) has no
two-torsion.  Glue with this boundary identification.  The closed
difference has \(H^2=0\), so it has its spin structure as a spin-c
structure and extends over \(S^1\times D^4\).  Its Dirac index is
zero by spin-c bordism.

\item For completeness, additivity here uses the cut-and-paste bordism
obtained by taking the cylinders on the common page and on the
two caps, and joining their common boundary collars by a pair of
pants.  The three boundary components are the old capped page,
the oppositely oriented new capped page, and the oppositely
oriented closed cap difference.  The same spin-c structure on the
common collar extends over this product construction.  Thus the
index of \((\bar X,b)\) is equal to that of the compact sectioned
elliptic page \(S\) with the resulting determinant \(b_\gamma\).

\par\medskip\noindent\textbf{Identify the base spin-c structure.}\enspace The two integral
vanishing cycles needed here can be exhibited directly.  In the
source marking of P3F.12, take

\[
 A_E=\begin{pmatrix}0&1\\-1&-1\end{pmatrix},\qquad
 B_E=\begin{pmatrix}0&-1\\1&0\end{pmatrix}.
\]

\item The nodal monodromy \(A_EB_E\) has image of its nilpotent part
generated by \(a=(0,1)^t\).  The twelve unramified cusp lifts of
the regular-binding page include the transport \(A_Ea=(1,-1)^t\).
The two columns \(a,A_Ea\) have determinant minus one.  Van Kampen
therefore kills the whole fundamental group of a torus fibre.
The section then identifies the fundamental group of the compact
elliptic page with that of its genus-three base, and identifies
the fundamental group of the open page with that of the base
with one boundary component.  In particular \(H_1(S;\mathbf Z)
\cong\mathbf Z^6\), so \(H^2(S;\mathbf Z)\) has no torsion.

\item On the marked boundary \(T^3\), the two source loops die in
\(H_1(V)\), and the base boundary loop is a product of
commutators and also dies.  Hence the restriction
\(H^1(V)\to H^1(T^3)\) is zero.  The image of
\(H^1(H_\gamma)\) is the span of the two source characters.
The Mayer–Vietoris kernel of restriction to \(V\) and
\(H_\gamma\) is consequently the infinite cyclic group generated
by the fibre class \(F\).  It is primitive, since the section
intersects a fibre once.

\item The elliptic canonical class \(K_S\) is pulled back from the base,
so it restricts to zero on both the page with boundary and the
regular cap.  The actual determinant and the canonical spin-c
determinant twisted by the Picard line thus satisfy

\[
 b_\gamma-\bigl(-K_S+2c_1(M)\bigr)=kF
\]

for an integer \(k\).  Both are characteristic classes.  Since
\(F\) is primitive, their common mod-two class makes \(k=2n\)
even.  Therefore

\[
 b_\gamma=-K_S+2c_1(M+nF).
\]

\item There is no two-torsion ambiguity in the spin-c structure on
\(S\).  Its Dirac operator is exactly the Dolbeault spin-c
operator twisted by \(M+nF\).

\par\medskip\noindent\textbf{Acyclicity from explicit sections.}\enspace The source
primitive-generator identification fixes, up to sign,

\[
 P_*=(0,\sqrt2\,t^2(t-1))
 \quad\hbox{on}\quad
 y^2=x^3-3t^3(t-1)x+2t^4(t-1)^2.
\]

\item We use that identification only to specify the line in the
construction.  The following direct calculation of its fibrewise
nontriviality does not use the height, the signature, or the square
of the page line.  In the ordinary affine charts the displayed
point is finite and hence different from the zero section.  At the
\(IV^*\) fibre take \(t=z^3\), \(x=z^4X\), \(y=z^6Y\).
Then

\[
 P'_*=(0,\sqrt2(z^3-1))\longrightarrow(0,-\sqrt2)
 \quad\hbox{on }Y^2=X^3+2.
\]

\item At the \(III\) fibre take \(t=1+z^4\), \(x=z^2X\),
\(y=z^3Y\).  This gives

\[
 P'_*=(0,\sqrt2\,z(1+z^4)^2)\longrightarrow(0,0)
 \quad\hbox{on }Y^2=X^3-3X.
\]

\item Both limiting points are finite smooth points of the good-reduction
elliptic curves and are different from the point at infinity.
At the nodal fibre put \(q=1/t\), \(X=q^2x\), \(Y=q^3y\).
The section is \((0,\sqrt2(1-q))\), with limiting point
\((0,\sqrt2)\) on \(Y^2=(X-1)^2(X+2)\).  It is smooth and
different from both the node \((1,0)\) and the zero section.
The same statements hold for \(-P_*\); fibre inversion gives
the exact comparison between these two choices.  Pullback under
the actual regular-binding cover preserves these explicit
sections and their nonidentity specializations.  Every smooth
elliptic fibre therefore carries a nontrivial degree-zero line,
whose zeroth and first cohomology vanish.  On each of the twelve
nodal fibres the line has degree zero and a nonidentity
multiplicative Picard parameter.  The normalization sequence is
a two-term complex
\(\mathbf C\xrightarrow{1-\lambda}\mathbf C\), where
\(\lambda\ne1\); it is acyclic.  The bundle \(nF\) is pulled
back from the base and does not affect any of these restrictions.
All source identities and root-chart substitutions are checked
exactly in \path{code/part3_fourfold/direct_section.py} and its recorded
outputs.
Proper cohomology and base change give the vanishing below.
More precisely, Stacks Project Lemma 30.22.1, tag 07VK,
applies to this proper flat projective elliptic fibration and
its line bundle: on each Noetherian affine base open it provides
a perfect complex computing the direct image and commuting with
arbitrary derived base change.  Every fibre complex is zero by
the preceding calculation.  A bounded complex of finite projective
modules with this property is zero: at any local ring, cancel
every differential entry that is a unit; the remaining minimal
complex has zero differentials after reduction to the residue
field, so its zero fibre cohomology forces every remaining rank
to be zero.  This proves

\[
 R f_*(M+nF)=0.
\]

\item Consequently the Dolbeault index is zero.  Combining the actual
symbol factorization and cap comparison gives

\[
 A=\operatorname{ind}(D_{\bar X,b})
  =\operatorname{ind}\bar\partial_{S,M+nF}=0.
\]

\item The computation uses fibrewise acyclicity and the explicit cap
switch; it does not reuse the numerical page characteristic
calculation in P3K.7.  The integral cover argument P3F.11 again
gives zero for the original parameter return.

\end{enumerate}
\end{proof}




\begin{theorem}\label{night:P3K.9}
\StatementRecord{night:P3K.9}{proved}{S6\_06\_BLIND: GPT-6 Astra (OpenAI, ChatGPT Pro mode, effort Pro), recorded re-derivation in the checkpoint dated 18 November 2026}{\cref{night:P3F.11}, \cref{night:P3K.6}, \cref{night:P3K.7}}
For the rotating framing specified below, the one-gamma-turn
return has Jin characteristic numbers \((k,\ell)=(-6,-24)\).
After removing an extending constant-rotation loop, this framing
descends through the sixfold parameter cover.  With the descended
framing the original positive return has

\[
 k=-1,\qquad \ell=-4.
\]

Every other tubular framing differs by the explicitly described
changes \((k,\ell)\mapsto(k+n,\ell+4n)\) of P3F.11.  The asserted
individual values include the framing; their index is independent
of that choice.
\end{theorem}

\begin{proof}
\begin{enumerate}[label=\arabic*.]
\item \par\medskip\noindent\textbf{Frames on the two pieces.}\enspace Give the graph copy of the page
its normal frame consisting of its oriented normal in \(Y\),
followed by the two Cartesian parameter directions.  The graph
derivative changes this frame only by a triangular shear, whose
off-diagonal block contracts linearly to zero.  On the compression
use a normal frame of \(W\subset Y\times I\), then transport it
by the actual gamma rotation \(R_\theta\).  This normal rank-three
bundle is trivial.  Indeed \(W\) is a solid torus minus a ball and
has trivial tangent bundle, while \(TY\oplus\mathbf R\) is trivial;
the Whitney formula gives \(w_2(\nu W)=0\).  Over a three-complex,
an oriented rank-three bundle is classified by \(w_2\), since
\(\pi_1(SO(3))=\mathbf Z/2\) and \(\pi_2(SO(3))=0\).
At the sphere end, its comparison with the fixed positive-Hopf
frame is a map \(S^2\to SO(3)\).  This extends over \(D^3\) because
\(\pi_2(SO(3))=0\).  Make this extension at parameter zero and
transport it by \(R_\theta\).  The resulting frame extends over
the whole marked compression cap \(H_a\).

\item We compute the transition between these two actual frames.  Use
radian angular coordinates in this paragraph.  Let \(\phi\) be
the fixed outward-collar phase, namely the Milnor-page phase plus
\(\pi\), and let \(d=12\) be its degree along the gamma circle.
At the binding its outward page ray has angle
\(\alpha=\phi-d\gamma\) in the complex \(q-\epsilon\) plane.
The exact graph boundary from P3K.6 is
\(\gamma=\theta+g_\epsilon(p)\), where
\(p\in T^2\) is the source point and \(g_\epsilon\) is now the
radian-valued lift of that graph.  No source-dependent change of
the binding phase is made.

\item After quotienting by the source tangent directions and the tangent
to this boundary graph, the four remaining coordinates are the
real and imaginary \(q-\epsilon\) directions, the parameter radial
direction, and the relative angular direction.  The chosen page
collar has positive parameter radial derivative, pointing from the
page towards the compression.  We use the outward page collar and
the inward compression tangent with this same direction.  In the
four indicated coordinates its tangent is
\(u=(\cos\alpha,\sin\alpha,\rho,0)\), with \(\rho>0\).
The explicit quotient morphism, from this four-space modulo
\(\mathbf Ru\) to the reference three-space, is

\[
 (v_q,v_r,v_\eta)\longmapsto
 \left(v_q-\frac{v_r}{\rho}(\cos\alpha,\sin\alpha),v_\eta\right).
\]

\item The Cartesian parameter axes have radial-relative coordinates
\((\cos\theta,-\sin\theta)\) and
\((\sin\theta,\cos\theta)\), after a positive scaling of the
relative angular coordinate.  Project them by this formula,
reverse both parameter axes, and apply the positive reference
rescaling \(\operatorname{diag}(\rho,\rho,1)\); normalize the
page-normal column by the positive factor \(\rho^{-1}\).
The graph normal frame in
this reference space, after constant oriented changes of frame and
positive rescalings, is

\[
 G(\alpha,\theta)=
 \begin{pmatrix}
 -\sin\alpha&\cos\theta\cos\alpha&\sin\theta\cos\alpha\\
 \cos\alpha&\cos\theta\sin\alpha&\sin\theta\sin\alpha\\
 0&\sin\theta&-\cos\theta
 \end{pmatrix}.
\]

\item Here the first column is the oriented page normal, and the other
two columns are the Cartesian parameter directions projected along
the collar tangent.  The coefficients of the source derivatives
of \(g_\epsilon\) enter only the source tangent directions already
quotiented out; they give triangular shears before this quotient.
The positive radial collar coefficient can be rescaled to one.
During collar rounding the transverse vector remains in the same
open radial half-space.  This space is contractible, so this
rounding contributes no additional three-dimensional clutching
class.  These operations specify a homotopy of the actual normal
transition, not merely an equality of its characteristic numbers.

\item Direct multiplication gives \(G\in SO(3)\) and
\(G(\alpha,\theta)=R_z(\alpha)G(0,\theta)\), where \(R_z\)
rotates the first two reference axes.  The compression-cap frame
in these same coordinates has the form \(h(p)\), independent of
\(\theta\): both the ambient product coordinates and this frame
are transported by \(R_\theta\).  Consequently the full transition
on the marked \(T^3\) is homotopic to

\[
\begin{aligned}
 h(p)^{-1}G(\phi-d(\theta+g_\epsilon(p)),\theta)
 &=\bigl[h(p)^{-1}R_z(-d g_\epsilon(p))\bigr]\\
 &\quad\cdot\bigl[R_z(\phi-d\theta)G(0,\theta)\bigr].
\end{aligned}
\]

\item The first factor is pulled back from the source \(T^2\), and the
second from the parameter circle.  The cubic Maurer--Cartan form
of a product of two maps is the sum of their cubic forms plus an
exact form; this identity follows by expanding
\((ab)^{-1}d(ab)=b^{-1}(a^{-1}da)b+b^{-1}db\) and using the
Maurer--Cartan equations.  Each of the two cubic forms here is
zero by dimension.  Its integral over \(T^3\) is therefore zero.
The Pontryagin clutching formula can be obtained directly by using
the interpolating connection \(t g^{-1}dg\) on the collar between
two framed pieces: its Chern--Weil integral is a fixed multiple
of that cubic integral.  We conclude

\[
 \langle p_1(\bar\nu),[\bar X]\rangle=0.
\]

\item This calculation allows nontrivial one-dimensional and
two-dimensional frame data.  It proves the required vanishing
without replacing the actual normal bundle by a trivial bundle.
The matrix identities and ensuing exact arithmetic are exported
in \path{code/part3_fourfold/normal_clutch.py} and its JSON and text outputs.

\par\medskip\noindent\textbf{The integers and their descent.}\enspace P3K.7 now gives, with
its full orientation comparison,

\[
 k_{\gamma}=\frac{0+3(-8)}4=-6,
 \qquad \ell_{\gamma}=-24+0=-24.
\]

\item We must compare the rotating framing to framings pulled back from
one original parameter circuit.  The fundamental group of the
normal-frame gauge space is

\[
 \pi_1\operatorname{Map}(S^2,SO(3))
 =\pi_1(SO(3))\oplus\pi_3(SO(3))
 =\mathbf Z/2\oplus\mathbf Z.
\]

\item Evaluation at one point, with its constant-map section, gives the
displayed splitting.  Pullback through the positive sixfold circle
cover sends \((\varepsilon,n)\) to \((0,6n)\).  Choose any
framing over the original parameter circle and write \(k_0\) for
its Jin integer.  Its sixfold pullback has integer \(6k_0\).
By the exact framing-to-return morphism P3F.11, the free part of
the difference between this pullback and the rotating framing is
\(-6-6k_0\), a multiple of six.  The remaining possible
\(\mathbf Z/2\) part is represented by a constant-axis rotation
loop, independent of the fixed-sphere point.  It extends over the
fixed \(D^3\), contributes zero to the cubic clutch integral, and
changes neither characteristic number.  Remove that extending
loop.  The resulting rotating framing is the sixfold pullback,
up to homotopy, of a framing over the original circle.  We take
this as the descended rotating-framing convention.

\item The replacement diffeomorphism for the pullback is the sixth
iterate of that original framed replacement.  Jin's integral
mapping-torus homomorphism, in the convention of P3F.11, gives
\(6k=-6\) and \(6\ell=-24\).  This proves the stated
\((k,\ell)=(-1,-4)\), including the exact descent of the
normalization.

\end{enumerate}
\end{proof}




\begin{theorem}\label{night:P3K.11}
\StatementRecord{night:P3K.11}{proved}{S6\_06\_BLIND: GPT-6 Astra (OpenAI, ChatGPT Pro mode, effort Pro), recorded re-derivation in the checkpoint dated 18 November 2026}{\cref{night:P2.2}, \cref{night:P3F.11}, \cref{night:P3F.17}, \cref{night:P3K.6}, \cref{night:P3K.7}, \cref{night:P3K.9}, \cref{night:P5.6}}
Fix logarithmic orders \(3K,4L\), with
\(\gcd(3K,4L)=1\), and primitive logarithmic coefficients
\(n_0,n_1\), so \(\gcd(n_0,3K)=\gcd(n_1,4L)=1\).
For the sphere members impose
\(a=4L n_0+3K n_1=\pm1\), with zero additional cusp collar.
Use the period coordinate \(c\) and scalar \(u\) of the original
family.  The closed higher parameter is
\[
 w_h=\exp\!\left(-\frac{\pi i c}{KL}\right),
 \qquad u=\kappa w_h^{2KL},\qquad \kappa\ne0.
\]
In particular, the original expression \(w=\exp(-\pi i c)\)
does not give one closed higher loop when \(KL>1\).
The positive sixfold cover \(w_h=r^6\) gives the signed
gamma turn \(a\gamma\).  For the positive-gamma sphere marking
\(a=1\), with the descended rotating framing of P3K.9, the Jin
integers of the positive \(w_h\)-loop are
\[
                  (k,\ell)=(-KL,-4KL).
\]
The integer return index is zero for either sphere sign.
These statements identify
both the parameter and the framing in the supplied higher-order
formulas.
\end{theorem}

\begin{proof}
\begin{enumerate}[label=\arabic*.]
\item \par\medskip\noindent\textbf{The integral orbifold complex.}\enspace We use the marked
lattice and period convention of P5.6.  Its cohomological
monodromy matrices are
\[
 T_1=\begin{pmatrix}
 1&0&-6&2\\0&-1&1&1\\0&-1&0&1\\0&0&0&1
 \end{pmatrix},\qquad
 T_2=\begin{pmatrix}
 1&6&0&-3\\0&0&-1&1\\0&1&0&0\\0&0&0&1
 \end{pmatrix}.
\]
Put \(A_i=T_i^{-t}\), \(A_0=(A_1A_2)^{-1}\),
\(N_1=1+A_1+A_1^2\), and
\(N_2=1+A_2+A_2^2+A_2^3\).  The lattice local system
on the higher root base is the pullback of the original one.
Consequently its finite cyclic norm operators are exactly
\(K N_1\) and \(L N_2\).  This assertion retains the
ineffective stabilizers: the generator of the additional
order-\(K\) or order-\(L\) stabilizer acts as the identity
on the lattice, so the full cyclic sum repeats the original
sum precisely \(K\) or \(L\) times.

\item Here is a complete block specification of the integer complex
used in the computation.  Let
\[
 C^0=\Lambda\oplus\Lambda\oplus\Lambda\oplus\mathbb Z^2,
 \quad C^1=\Lambda^7,\quad C^2=\Lambda^5,
 \quad C^3=\Lambda^2,
 \qquad \Lambda=\mathbb Z^4,
\]
and let \(J:\mathbb Z^2\to\Lambda\) have columns
\(e_3,e_4\).  In the displayed ordered summands, the three
differentials are
\[
\begin{split}
 d^0(a,b,c,z)={}&\bigl((A_1-1)a,(A_2-1)a,
 (A_1-1)b,(A_2-1)c,\,a-b,a-c,a-Jz\bigr),\\
 d^1(p,q,r,s,u,v,w)={}&\bigl(KN_1r,LN_2s,
 p-r-(A_1-1)u,\ q-s-(A_2-1)v,\\
 &\hspace{39mm}-A_0p-A_0A_1q-(A_0-1)w\bigr),\\
 d^2(x,y,z_1,z_2,z_0)={}&\bigl((A_1-1)x,(A_2-1)y\bigr).
\end{split}
\]
The pair-of-pants and the three annuli give the last three
summands of \(C^2\).  The first two are the degree-two terms
of the finite cyclic resolutions; the cusp contributes its
invariant lattice \(J\mathbb Z^2\) in degree zero.  These
are precisely the gluing maps of the lattice sheaf, including
the cyclic degree-two terms.  In particular, no finite norm
relation is replaced by a rational invariant projector.
Direct multiplication gives \(d^1d^0=d^2d^1=0\).

\item Write
\[
 \delta=(0,0,0,1)^t,\quad
 \epsilon_1=(1,2,-4,0)^t,\quad
 \epsilon_2=(1,3,-3,0)^t.
\]
The two invariant lattices have primitive bases
\(B_i=(\delta,\epsilon_i)\), and the norm matrices in
these bases are
\[
 B_1^{-1}N_1=
 \begin{pmatrix}0&2&1&3\\3&0&0&0\end{pmatrix},
 \qquad
 B_2^{-1}N_2=
 \begin{pmatrix}0&2&2&4\\4&0&0&0\end{pmatrix}.
\]
Here \(B_i^{-1}\) denotes the inverse coordinate map on its
image, so these displayed expressions are integral maps
\(\Lambda\to\mathbb Z^2\).  A primitive basis of
\(\ker d^2\) is \(\operatorname{diag}(B_1,B_2,I_{12})\).
The exact computation eliminates twelve unit pivots in the
matrix of \(d^1\) in that basis.  Every operation is an
elementary operation over \(\mathbb Z[K,L]\); the exported
left and right matrices have determinant \(\pm1\).
The remaining matrix has the same image as the three columns
\[
 R_0=
 \begin{pmatrix}
 K&0&0\\0&0&-3K\\0&2L&0\\0&0&4L
 \end{pmatrix}.
\]
This last image equality is also certified integrally: the
export gives matrices \(C,Q\) with
\(R_{\rm tail}=R_0C\), \(CQ=I_3\), and
\(R_{\rm tail}Q=R_0\).  Thus it is a proof for symbolic
\(K,L\), in addition to the executed Smith forms.
Equivalently, adjoining the eliminated generator \(h_c\),
the presentation is
\[
\begin{gathered}
 K\delta_1=0,\qquad 2L\delta_2=0,\\
 3K\epsilon_1=h_c=4L\epsilon_2.
\end{gathered}
\]
Since \(\gcd(3K,4L)=1\), the last relation gives a free
rank-one group and no torsion.  Moreover \(K\) is odd and
\(\gcd(K,L)=1\).  Therefore
\[
 H^2(\widehat B,\mathcal H)
 =\mathbb Z\oplus\mathbb Z/K\oplus\mathbb Z/(2L)
 =\mathbb Z\oplus\mathbb Z/(2KL).
\]
The complete matrices and elementary transformations are in
\path{code/part3_fourfold/higher_parameter.json}.
The same script computes the full integer complex, rather than
only this reduced presentation, for
\((K,L)=(1,1),(1,2),(1,5),(3,1),(5,1),(5,2),(7,5),(13,7)\).
Their torsion orders are respectively
\(2,4,10,6,10,20,70,182\), with free rank one in every case.

\par\medskip\noindent\textbf{The actual torsor period.}\enspace In the fixed period
coordinates, changing \(c\) to \(c'=c-b\) is the integral
shear
\[
 S_b=I+bE_{41},\qquad
 \Pi_c=\Pi_{c'}S_b,\qquad b\in\mathbb Z.
\]
It fixes \(\delta\) and sends each \(\epsilon_i\) to
\(\epsilon_i+b\delta\).  The logarithmic vectors are
\(r_i\delta+n_i\epsilon_i\).  Their actual logarithmic
annulus cocycles therefore change the characteristic class
by
\[
             b(n_0\delta_1+n_1\delta_2).
\]
This follows either by inserting the two vectors in the
displayed cochain complex, or by taking the logarithms in the
gluing maps \(x\mapsto x+\log(s_i)v_i/(2\pi i)\).
The coefficient \(r_i\) is fixed by the shear and contributes
zero to the difference.  Primitivity implies
\(\gcd(n_0,K)=\gcd(n_1,2L)=1\), so the displayed torsion
element has order exactly \(2KL\).  A change of logarithm
branch changes its cochain representative by a coboundary
and leaves this conclusion unchanged.

\item For completeness, the zero of this characteristic obstruction
is realized by actual holomorphic translations.  The vector
bundle of Lie algebras on the higher root base is the pullback
of the original bundle.  The added stabilizers act trivially
on its fibres.  Taking their invariants is exact in
characteristic zero, and the coarse pushforward is unchanged.
Consequently the explicit original calculation
\(H^1(\widehat B,\mathcal V)=H^2(\widehat B,\mathcal V)=0\)
continues to hold.  In the exponential sequence
\(0\to\mathcal H\to\mathcal V\to\mathcal V/\mathcal H\to0\),
the connecting map from torsors to \(H^2(\mathcal H)\) is
therefore an isomorphism.  When \(2KL\mid b\), solve the
integer coboundary equation above and subtract its logarithms
on the overlapping charts.  The remaining vector-valued
cocycle is a coboundary by \(H^1(\mathcal V)=0\); exponentiating
its chartwise primitives gives the required translation
isomorphism.  Reversing these primitives is its inverse.
At the cusp the shear changes the bounded beta term by an
integer; its toric factors \(\exp(2\pi i C\lambda)\), fan,
and monomial transition maps are identical.  Thus these chart
translations extend through the actual filling, as in the
explicit sheaf construction of P5.6.  Conversely a marked
translation isomorphism makes the characteristic difference
zero, so \(2KL\mid b\).  This proves the asserted exact
period in the marked family, with an extending morphism in
both directions.

\par\medskip\noindent\textbf{The loop and its characteristic numbers.}\enspace The scalar
relation is \(u=\kappa\exp(-2\pi i c)\).  Quotienting the
marked parameter by its just-proved period \(2KL\) gives
\(w_h=\exp(-\pi i c/(KL))\), and consequently
\(u=\kappa w_h^{2KL}\).  Put \(d=12KL\),
\(u=\kappa r^d\), and \(v=\log r/(2\pi i)\).
We retain the sign of the actual origin calculation.  The
finite-cap parameter descent is translation by
\[
 H(t)\delta=-\frac{4L n_0\log t+3K n_1\log(t-1)}{2\pi i}\delta.
\]
At the cusp it is \((as+U(q))\delta\), where \(U\) is a
single-valued holomorphic unit logarithm.  Write the period
\(\gamma=(\Gamma,C-s-dv)^t\).  On \(q=Q^d\), the
actual extending origins are
\[
\begin{split}
 o&=\binom{-a\Gamma s/d}{a s^2/(2d)-aCs/d+v(as+U)},\\
 \sigma&=\binom{a\Gamma v}{vU+aCv-ad v^2/2+ad v/2}.
\end{split}
\]
Their parameter differences are respectively \(H\delta\)
and \(H\delta+a\gamma\).  Their \(Q\)-differences are
integral lattice periods because \(d\) is even.  The
bounded cusp expression for \(\sigma\), and the invariant
zero coordinate after subtracting \(o\), give precisely the
extending sections in the node charts used in P2.2--3.
This checks the same morphisms with \(12\) replaced by
\(d\); it retains the signed twist.  Subtracting gives
\[
 F=\sigma-o
 =a\binom{\Gamma(s+dv)/d}{C(s+dv)/d-(s+dv)^2/(2d)+dv/2},
\]
and hence
\[
 F(s,v+\theta)-F(s,v)
 =a\theta\gamma+\frac{ad}{2}\theta(1-\theta)\delta.
\]
The last term is a periodic unitary gauge with both endpoint
values zero.  Thus the quotient clutch is the signed gamma
loop \(a\gamma\).  Equating the scalar expressions gives
the exact sixfold map \(w_h=r^6\), independently of the sign.

\item For \(a=1\), the page calculation P3F.17 gives
\[
 s_\delta=-2KL,\qquad p_{\rm rel}=-24KL.
\]
The normal-frame calculation of P3K.9 applies with page degree
\(12KL\).  Its transition factors into a two-dimensional
factor and a circle factor, exactly as in that proof;
the cubic clutch integral remains zero.  Thus
\(p_1(\bar\nu)=0\), whereas the Euler determinant has
square \(b^2=4s_\delta=-8KL\).  The oriented sphere-bundle
comparison P3K.7 gives
\[
 k_\gamma=\frac{3(-8KL)}4=-6KL,
 \qquad \ell_\gamma=-24KL.
\]
To descend the framing, choose any framing over the
\(w_h\)-circle.  Its sixfold pullback differs from the
rotating framing by a class in
\(\pi_1\operatorname{Map}(S^2,SO(3))=\mathbb Z/2\oplus\mathbb Z\).
The free difference is \(-6KL-6k_0\), which is divisible
by six.  Remove the constant-axis torsion loop, which extends
over the fixed \(D^3\) and changes neither characteristic
number.  The remaining class is the sixfold pullback of an
integral framing class, exactly as proved in P3K.9.  The
resulting descended rotating framing therefore gives
\(6k=-6KL\) and \(6\ell=-24KL\).  This proves the
stated pair and \(A=(4k-\ell)/24=0\).  For \(a=-1\),
the just-computed clutch traverses the gamma loop inversely.
The actual real trace of P3K.6--7 applies to this source
page with the corresponding boundary orientation; the index
of the inverse loop is the negative of the forward index.
That forward index is zero since
\(12s_\delta-p_{\rm rel}=0\).  Consequently the covered
index is zero also for \(a=-1\), and integral additivity
implies that the original \(w_h\)-loop has index zero.
The displayed individual pair is assigned to the stated
positive-gamma marking, so it involves no implicit reversal
of a negative sphere datum.

\item The original expression \(w\) and the higher expression
\(w_h\) satisfy \(w=w_h^{KL}\).  A single positive
\(w\)-loop only advances \(c\) by its original period two;
the higher characteristic obstruction vanishes for that
advance exactly when \(KL=1\).  This explicitly accounts
for the otherwise different-looking cover counts.  Changing
the final normal framing still changes \((k,\ell)\) by
\(n(1,4)\), as in P3F.11; the displayed values use the
descended rotating convention throughout.

\end{enumerate}
\end{proof}




\begin{theorem}\label{night:P3K.12}
\StatementRecord{night:P3K.12}{proved}{S6\_06\_BLIND: GPT-6 Astra (OpenAI, ChatGPT Pro mode, effort Pro), recorded re-derivation in the checkpoint dated 18 November 2026}{Definitions and the displayed constructions}
In the rank-one Mumford model above a smooth elliptic source fibre,
every double curve of the central cycle is elliptic and fixed by the
line circle.  A smooth cyclic descent, with the actual faithful root
character on the base, retains an elliptic fixed curve.  The same
conclusion holds after potential good reduction of the source and
for any number of permuted double curves.  Such a descended local
piece cannot occur in a smooth circle action on \(S^6\).
\end{theorem}

\begin{proof}
\begin{enumerate}[label=\arabic*.]
\item Let \(E\) be one of the double curves.  At a point of
\(E\) the smooth threefold has coordinates \((e,x,y)\), with
\(e\) a local elliptic coordinate and root-base map
\[
                         z=xy.
\]
The line circle acts by \((e,x,y)\mapsto(e,\tau x,\tau^{-1}y)\).
For a higher vanishing order the refined Mumford chain has these
same charts at each double curve.  A finite root deck transformation
commutes with this circle, so a stabilizer at the chosen point
preserves the two normal weight lines separately.  A nonidentity
element of the stabilizer acts faithfully on \(z\), since it is
a subgroup of the cyclic root group acting by multiplication on
the root coordinate.

\item Write a cyclic stabilizer of order \(h\) on its tangent space as
\[
 (e,x,y)\longmapsto
 (\zeta^{a}e,\zeta^{b}x,\zeta^{c}y),
 \qquad b+c\equiv1\pmod h,
\]
after choosing a generator for its faithful base character.  The
equation on the characters follows from the quadratic leading term
of \(z=xy\); it is invariant under changes of local coordinates.
An element which acts nontrivially on the elliptic tangent cannot
be a pseudoreflection.  Indeed, a pseudoreflection with nontrivial
elliptic character would have to fix both normal lines pointwise,
and would therefore fix \(z\).  Faithfulness makes that element
the identity, a contradiction.  Hence every pseudoreflection lies
in the kernel of the elliptic tangent character.

\item We spell out the smooth-quotient criterion used here.  A finite
analytic action fixing the origin is holomorphically linearized by
the germ
\[
 F(v)=\frac1{|H|}\sum_{g\in H}(Dg_0)^{-1}g(v).
\]
It is equivariant, its derivative is the identity, and the inverse
function theorem supplies its inverse.  Averaging over the finite
group commutes with completion of the invariant local ring.  Thus
the analytic quotient is smooth exactly when the algebraic linear
quotient has a regular local ring at the origin.  Satriano,
Corollary 1.8, applies to that linear action; constant finite
groups are stable by Proposition 2.2.  It says that the stabilizer
of a smooth quotient is generated by pseudoreflections
\cite[Corollary 1.8 and Proposition 2.2]{Satriano2012}.
If the elliptic tangent character were nontrivial, the preceding
kernel argument would contradict this criterion.  It is therefore
trivial for every point stabilizer.

\item A finite-order automorphism of an elliptic curve which fixes a
point and has derivative one is the identity: the same averaging
linearization in one dimension proves identity near that point,
and analytic continuation proves identity on the connected curve.
Consequently the effective action of the setwise stabilizer of
\(E\) is free on \(E\).  Its quotient is again an elliptic
curve, by the unramified Riemann--Hurwitz formula.  If the root group
permutes double curves, choose one representative from any orbit;
its image is precisely the quotient by its setwise stabilizer.
It is a smooth elliptic curve fixed by the descended line circle.
Possible stabilizer elements acting identically on \(E\) can
ramify in a normal direction, but they do not change this curve
quotient or its genus.

\item Choose a prime not dividing the two nonzero normal circle weights
along the elliptic fixed curve.  Average a Riemannian metric and
use its equivariant normal exponential map; in that neighbourhood
the fixed set of the prime-order subgroup is exactly the zero
section.  Thus the elliptic curve is a connected component of
that finite-subgroup fixed set.  The Smith--Floyd inequality
\cite[Theorem 4.4, p.~146]{Floyd1952} bounds its total mod-prime
Betti number by that of \(S^6\), which is two.  The theorem's
locally compact, finite-dimensional Hausdorff hypotheses apply to
this compact smooth manifold; its Cech homology agrees with
singular homology here and on the smooth fixed submanifold.
An elliptic fixed component already has total Betti number four,
a contradiction.  This proves the
assertion without assuming that the finite descent is unramified
in the normal directions.

\end{enumerate}
\end{proof}




\begin{theorem}\label{night:P3K.13}
\StatementRecord{night:P3K.13}{proved}{S6\_06\_BLIND: GPT-6 Astra (OpenAI, ChatGPT Pro mode, effort Pro), recorded re-derivation in the checkpoint dated 18 November 2026}{Definitions and the displayed constructions}
In a retained nonzero-linearization \(I_n\) neighbourhood, the
circle quotient is the mapping torus of the actual source
translation.  If its global oriented orbit space has
\(H_2(Y;\mathbb Z)=0\), that translation is transitive on the
\(I_n\) components.  In particular this conclusion applies to
the retained neighbourhoods of the semifree sphere members with
the single standard fixed \(\mathbb P^1\).
\end{theorem}

\begin{proof}
\begin{enumerate}[label=\arabic*.]
\item Let \(S_D\) be the smooth elliptic-surface
neighbourhood of the fibre, and let \(Q:S_D\to S_D\) be
translation by the specified section.  The actual nonzero
linearization gives \(M^\times/(\lambda\phi)\).  Quotient
first by the unit line circle.  A Hermitian norm identifies
\(M^\times/S^1\) with \(S_D\times\mathbb R\), where
the second coordinate is log norm.  The lifted translation becomes
\((p,t)\mapsto(Qp,t+\rho(p))\), after reversing the chosen
generator if necessary, with \(\rho(p)>0\).  This inequality
holds uniformly on a compact smaller neighbourhood by contraction.
Its fundamental region is
\(0\le t\le\rho(p)\), with the upper boundary identified to
the lower boundary by \(Q\).  Reparametrizing each interval
smoothly to \([0,1]\), with compatible collar germs at its
ends, gives an orientation-preserving diffeomorphism to
\[
 T_Q=(S_D\times[0,1])/(p,1)\sim(Qp,0).
\]
This is the actual orbit map; no period-only replacement has been
used.  The global orientation on this neighbourhood is fixed by
the complex orientation of \(S_D\), followed by the positive
suspension coordinate.  Reversing that convention would multiply
all intersection numbers below by the same sign and would not
affect their nonvanishing.

\item Label the fibre components cyclically by \(C_j\),
\(j\in\mathbb Z/n\), and let \(Q\) shift labels by
\(c\).  Put \(g=\gcd(n,c)\) and \(h=n/g\).  Its
component orbits have the integral cycle sums
\[
 Z_r=\sum_{j\equiv r\;({\rm mod}\;g)}C_j,
 \qquad 0\le r<g.
\]
Suppose \(g>1\).  No two components in one orbit are adjacent
in the cycle.  Hence each \(Z_r\) is a disjoint union of smooth
oriented two-spheres in the slice \(S_D\times\{0\}\).
Its translation mapping torus
\[
 D_r=(Z_r\times[0,1])/(p,1)\sim(Qp,0)
\]
is a smooth closed oriented three-manifold in \(T_Q\),
because \(Q\) permutes these disjoint spheres by holomorphic
isomorphisms.  It defines an integral three-cycle.  These cycles
are compact and lie in the interior of the local
neighbourhood, so inclusion defines classes in the global
\(H_2(Y;\mathbb Z)\) and \(H_3(Y;\mathbb Z)\).

\item The local intersection pairing is
\[
                    [Z_r]\cdot[D_s]=Z_r\cdot_S Z_s.
\]
To verify it, represent the two-cycle by a small transverse
perturbation in a source slice and represent the mapping-torus
cycle in a product collar of that slice.  Its extra tangent
direction is the positive suspension direction.  The local
five-dimensional intersection sign is therefore exactly the
complex-surface intersection sign.  The same calculation using
a normal pushoff computes the diagonal self-intersections.
All perturbations are compactly supported in this neighbourhood,
so the calculation also equals the global homological pairing.

\item Let \(B_n\) be the intersection matrix of the \(I_n\)
cycle: diagonal entries are \(-2\), successive distinct
components meet once, and the two \(I_2\) components meet
twice.  Put \(B_1=(0)\).  Summing by the orbits gives
\[
                 (Z_r\cdot Z_s)_{r,s}=h B_g.
\]
This identity follows by counting the \(h\) equal edges
between adjacent orbits.  The quadratic form is
\[
 (t_0,\ldots,t_{g-1})B_g(t_0,\ldots,t_{g-1})^t
       =-\sum_{r\in\mathbb Z/g}(t_r-t_{r+1})^2.
\]
Thus its rank is \(g-1\), and it has a nonzero entry whenever
\(g>1\).  Such an entry is impossible when
\(H_2(Y;\mathbb Z)=0\), because every displayed global
two-cycle would then have zero pairing with every closed
three-cycle.  Hence \(g=1\), proving transitivity.
The program \path{code/part3_fourfold/vertical_scope.py} exports
every orbit matrix for \(2\le n\le12\), all shifts, and
their exact Smith forms.  The displayed cycle and intersection
arguments prove the general formula, independently of that
finite check.

\end{enumerate}
\end{proof}




\begin{theorem}\label{night:P3K.14}
\StatementRecord{night:P3K.14}{proved}{S6\_06\_BLIND: GPT-6 Astra (OpenAI, ChatGPT Pro mode, effort Pro), recorded re-derivation in the checkpoint dated 18 November 2026}{\cref{night:P3F.2}, \cref{night:P3K.13}}
Let \(S\) be one of the rational elliptic surfaces with
\({\rm MW}(S)=\mathbb ZG\).  Write the generic Picard point
of \(M\) as \(aG\), and the translating section as \(bG\).
For a literal Remark 8.5 line with
\(M\otimes t_{bG}^*M^{-1}=\pi^*\mathcal O(k)\),
\[
                         k=ab\,h(G).
\]
At a retained \(I_n\) on which the translation is transitive,
the whole-fibre linearization forces the generic Picard section
to meet the identity component.  At that fibre every vertical
correction in \(M\) is a whole-fibre divisor.  In particular,
for the transitive component markings in the rank-one table,
\(n\mid a\).
\end{theorem}

\begin{proof}
\begin{enumerate}[label=\arabic*.]
\item On the generic fibre,
\(M\otimes\mathcal O(aG-O)^{-1}\) is trivial.  A nonzero
rational trivializing section has divisor supported on fibres,
so
\(M=\mathcal O(aG-O+D)\) for an integral vertical divisor
\(D\).  Choose a positive integer \(N\) such that
translation by \(NbG\) fixes every irreducible fibre
component.  Iterating the actual line isomorphism telescopes to
\[
 M\otimes t_{NbG}^*M^{-1}=\pi^*\mathcal O(Nk).
\]
The pullback of \(D\) by this translation equals \(D\)
as a divisor, so it cancels from the left-hand side.  Intersecting
with the zero section computes its degree as
\[
 (aG)\cdot O+1-(aG-NbG)\cdot O+(NbG)\cdot O.
\]
In the section-height formula of P3F.2, the local corrections
for \(aG\) and \(aG-NbG\) agree, since \(NbG\) meets
the identity component everywhere; those for \(NbG\) are
zero.  Substitution cancels these local terms and the constant
terms, leaving
\[
 \frac{a^2-(a-Nb)^2+(Nb)^2}{2}\,h(G)
                         =Nab\,h(G).
\]
Division by \(N\) proves the asserted degree, with no
restriction on the allowed vertical divisor.

\item For the local statement, intersect the whole-fibre line relation
with every \(C_j\).  It gives
\(M\cdot C_j=M\cdot t_{bG}(C_j)\).  Under a transitive
permutation these degrees are equal.  Their sum is the degree
of \(M\) on the fibre, namely zero, so every degree is zero.
If \(aG\) meets \(C_r\), the component-degree equation is
\[
                   e_r-e_0+B_n d=0,
\]
where \(d\) is the vector of coefficients of \(D\) at this
fibre.  For every integral vector \(d\),
\(\sum_j j(B_nd)_j\) is divisible by \(n\), as follows
by subtracting consecutive cyclic differences.  Applying this
functional to the displayed equation gives \(r\equiv0\pmod n\).
Thus \(aG\) meets the identity component.  The equation
now reduces to \(B_nd=0\).  The displayed quadratic form
in P3K.13 shows that \(d\) is constant, so this divisor is
an integer multiple of the full fibre.  The actual generator
markings derived in P3F.2 are: components \(3,5\) at \(I_8\);
\(2,5\) at \(I_7\); \(2,3\) at the OS55 \(I_5\)
and \(1,4\) at the OS56 \(I_5\); \(1,3\) at \(I_4\);
\(1,2\) at \(I_3\); and \(1\) at \(I_2\).
P3F.2 obtains these labels from the enumerated rank-one source
lattices and the local height contributions, rather than choosing
an arbitrary component marking.  Every displayed label is a unit
in its \(\mathbb Z/n\).  Thus \(r=a[G]=0\) is exactly
\(n\mid a\).  This also proves that the local removal of
the vertical divisor is multiplication by the actual meromorphic
base factor \(t^j\), a holomorphic invertible map on the
punctured annulus.  Its exponent is the integral delta collar
change; no fractional or unrecorded gauge is introduced.

\end{enumerate}
\end{proof}


P3K.15 below transports the component cycles, the line, and its integer
delta character through a supplied affine cap on an actual extending
root line.  Applying it to a proposed new cap still requires constructing
that root source, its line, and the complete cyclic affine norm.  Partial
roots may leave a reduced reducible fibre, and zero-gamma finite kernels
must be retained.  Thus neither the retained-neighbourhood obstruction
nor the cap comparison alone is a complete enumeration of the literal
arbitrary-site scope.

\begin{theorem}\label{night:P3K.15}
\StatementRecord{night:P3K.15}{proved}{S6\_06\_BLIND: GPT-6 Astra (OpenAI, ChatGPT Pro mode, effort Pro), recorded re-derivation in the checkpoint dated 18 November 2026}{\cref{night:P3K.13}}
Every smooth sphere member with the standard fixed sphere and
normal circle weights \((-1,+1)\) has a semifree line-circle
action.  For a cyclic affine log cap at a nonzero-linearization
site, written on a smooth root source as a linearized line quotient, its actual circle orbit is
a mapping torus of a smooth source-surface quotient.  The finite
zero-gamma subgroup and its line character are retained in this
comparison.  If the global orbit has \(H_2(Y;\mathbb Z)=0\),
the resulting monodromy is transitive on the components of its
central fibre.  The descended line has multidegree zero there,
and a vertical change with the same generic Picard point is an
integer multiple of the primitive whole-fibre divisor.
\end{theorem}

\begin{proof}
\begin{enumerate}[label=\arabic*.]
\item \par\medskip\noindent\textbf{Semifreeness from the actual fixed sphere.}\enspace Let
\(F\cong S^2\) be the fixed sphere whose two normal weights
are \(-1,+1\).  For every prime \(p\), an invariant normal
exponential map identifies the \(\mu_p\)-fixed set in a
neighbourhood of \(F\) with \(F\) itself.  Thus \(F\)
is an entire connected component of \(X^{\mu_p}\).
It contributes two to the total mod-\(p\) Betti number.
Floyd's inequality \cite[Theorem 4.4]{Floyd1952} bounds the
total by two.  Any point of \(X^{\mu_p}\setminus F\)
would belong to another component and contribute at least one,
which is impossible.  Every nontrivial finite subgroup of the
circle contains a subgroup of prime order.  Consequently no
finite nontrivial stabilizer occurs outside \(F\), proving
semifreeness without an assumption on the distant log caps.

\par\medskip\noindent\textbf{The lifted group and its finite kernel.}\enspace Write the
actual root source as \(S'\to\Delta_z\), a smooth surface,
and its generic-degree-zero line as \(N\).  Its scalar
suspension lift is a holomorphic line automorphism \(B\) of
\(N^\times\), covering the source translation \(Q\).
For sufficiently small nonzero scalar it is uniformly contracting
on a compact smaller source neighbourhood.  The actual order-\(m\)
affine cap action on \(N^\times/\langle B\rangle\) lifts
to a holomorphic line map \(A\) on the root source, with
\[
 AB=BA,\qquad A^m=B^n,
\]
and with root action \(z\mapsto\zeta_m z\).  Here \(n\)
is the integral gamma coefficient of the actual cyclic norm.
More explicitly, on the vector cover write the affine action
as \((p,R)\).  Its \(m\)-th power translates by
\[
            p+Rp+\cdots+R^{m-1}p
                 =n\gamma+\lambda_0,
\]
where \(\lambda_0\) is an elliptic-and-delta period.  This
is the actual order-\(m\) descent condition.  The latter
period acts trivially on \(N^\times\), whereas the gamma
period is exactly \(B\); hence \(A^m=B^n\).
For the supplied caps with invariant vector \(v\),
\(p=\Pi v/m\) and the norm is exactly \(\Pi v\), so
\(n=\psi(v)\).  For a noninvariant affine translation the
displayed norm, rather than its initial gamma coefficient,
is required.  The root deck preserves this gamma point in
the semiabelian group, so its affine lift commutes with \(B\).
This retains any delta phase of \(A\); it is not discarded
from the line lift.

\item First take \(\gcd(m,n)=1\).  Choose integers \(r,s\)
with \(rn+sm=1\), and set \(H=A^rB^s\).  The displayed
relations give
\[
                  A=H^n,\qquad B=H^m.
\]
These are identities of actual line maps, not merely identities
of their homological actions.  Thus the original cap quotient
\(N^\times/\langle A,B\rangle\) is exactly
\(N^\times/\langle H\rangle\).  In a root choice for the
scalar, \(A\) contains its \(n/m\)-th power, while \(B\)
contains its first power.  Hence \(H\) contains the positive
power \(1/m\), times a line map bounded on the chosen compact
source neighbourhood.  Taking the scalar smaller makes \(H\)
uniformly contracting.  This conclusion also holds when some of
\(n,r,s\) are negative; the exponent is the exact Bezout
sum \(rn/m+s=1/m\).

\item For the nonprimitive case put \(g=\gcd(m,n)\),
\(m'=m/g\), \(n'=n/g\), and define
\[
                  C=A^{m'}B^{-n'}.
\]
Then \(C^g=1\), and the root-base character shows that its
order is exactly \(g\).  It is the complete finite subgroup
of the abelian group generated by \(A,B\).  Choose
\(rn'+sm'=1\).  Modulo \(C\), the map
\(H=A^rB^s\) satisfies
\[
             \bar A=\bar H^{n'},\qquad
             \bar B=\bar H^{m'}.
\]
This describes the actual finite quotient followed by the
residual contracting suspension.  In particular it does not
replace \((m,n)\) by a reduced pair while forgetting its
finite kernel.

\item The marked outer meridians are retained under this change of
generators.  In the ordered generators \((C,H)\), they are
exactly
\[
                   A=C^sH^{n'},\qquad
                   B=C^{-r}H^{m'}.
\]
The matrix with columns \((m',-n')^t\) and \((r,s)^t\)
has determinant one.  It transforms the sole relation
\((m,-n)^t\) to \((g,0)^t\).  Thus even for an order-one
cap the integer annular clutch is still the specified meridian;
it has not been lost in the local mapping-torus description.
\path{code/part3_fourfold/finite_cap_group.py} exports these matrices
and checks them for all \(1\le m\le24\),
\(-2m\le n\le2m\), as well as retaining the general
Bezout proof above.

\item Let \(\bar S=S'/\langle C\rangle\).  A stabilizer of a
point \(p\in S'\) acts on the complex line \(N_p\) by
a root of unity.  If that root were nontrivial, the corresponding
unit line rotation would identify two points of the cap quotient
and produce nontrivial finite line-circle isotropy.  Semifreeness
therefore forces every source-point stabilizer to act trivially
on its line fibre.  The line descends to a line bundle
\(\bar N\) on \(\bar S\).  This descent can be seen
without a categorical shortcut: choose a local frame nonzero at
the point and average it over its stabilizer.  Its value remains
nonzero because the stabilizer acts trivially on that fibre.
The invariant frame identifies the quotient line with a product
over the local source quotient.  Translates of these frames give
the transition functions on the quotient.

\item The finite source quotient is smooth.  Indeed the line-complement
quotient \(N^\times/\langle C\rangle\) is smooth, since
its further quotient by the residual contracting infinite cyclic
group is the given smooth cap and that quotient is locally a
covering.  In the invariant frames just constructed it is
\(\bar S\times\mathbb C^*\).  A local analytic product
with \(\mathbb C^*\) is smooth exactly when its base is
smooth, as follows from the local dimension and cotangent-space
dimension.  This proves smoothness of \(\bar S\).
The source quotient may be branched in a normal direction; the
argument only uses the proved triviality of the line stabilizers.

\par\medskip\noindent\textbf{The real orbit map and its line.}\enspace On \(\bar N\)
choose a Hermitian metric and write log norm as \(t\).
The contracting residual map becomes
\[
(p,t)\longmapsto(hp,t-\rho(p))
\]
on the line-circle quotient, with
\(\rho(p)=-\log(\|H(v)\|/\|v\|)>0\).  It is
bounded away from zero on a compact smaller neighbourhood.
The upper face \((p,\rho(p))\) is identified by this actual
map with \((hp,0)\).  As in P3K.13, this interval fundamental
region is smoothly
reparametrized to the unit interval with compatible end collars.
This gives an explicit orientation-preserving diffeomorphism
from the actual circle orbit of the cap to
\[
           T_h=(\bar S\times[0,1])/(p,1)\sim(hp,0).
\]
The principal circle line is also identified: it is the unit
circle bundle of \(\bar N\) on this mapping torus, glued at
the ends by the unit normalization of the actual line map
\(\bar H\).  Both its scalar phase and any finite-kernel
character have therefore been transported, not reconstructed
from the orbit alone.

\par\medskip\noindent\textbf{The component intersection obstruction through the cap.}\enspace
Let the connected central fibre of \(\bar S\) be
\(\sum_i m_i C_i\), and let \(B=(C_i\cdot C_j)\).
All \(m_i\) are positive.  The relation \(B m=0\) and
nonnegative off-diagonal intersections give the exact identity
\[
 x^tBx=-\sum_{i<j}(C_i\cdot C_j)m_i m_j
       \left(\frac{x_i}{m_i}-\frac{x_j}{m_j}\right)^2.
\]
Connectedness of the fibre implies that the kernel is exactly
the line spanned by \(m\).  The map \(h\) permutes its
components and preserves the multiplicities.  If there are
\(r\) component orbits, the invariant coefficient space has
dimension \(r\).  The restriction of \(B\) to it has
rank \(r-1\).

\item We realize this restricted form by actual cycles in \(T_h\),
including the cases where components within one orbit meet.
Let \(z\) be an integral two-cycle representing an orbit-sum
divisor.  Its class is \(h\)-invariant, so choose a compact
singular three-chain \(c\) in \(\bar S\) with
\(\partial c=h_*z-z\).  The product trace of \(z\)
over the suspension interval, corrected by \(-c\) in the
seam slice, is a closed integral three-cycle \(D_z\) in the
mapping torus.  A test two-cycle in a slice halfway along the
interval misses the correcting chain.  Its intersection with
\(D_z\) is exactly its surface intersection with \(z\),
with the complex-surface orientation followed by the positive
suspension direction.  Thus all entries of the invariant
intersection form are realized by compact two- and three-cycles
in the actual cap orbit.  Inclusion in the global orbit preserves
these intersection numbers.  If \(H_2(Y;\mathbb Z)=0\),
every such number is zero, forcing \(r=1\).

\item Finally \(\bar H\) is a line automorphism, so the degrees
\(\bar N\cdot C_i\) are invariant under the component
permutation.  Transitivity makes them all equal.  Their positive
multiplicity-weighted sum is the generic fibre degree, zero;
hence every one is zero.  For two such line extensions with the
same generic Picard point, let \(L\) be their ratio line.
The proper source map to the disc has coherent direct image
\(\pi_*L\) by proper direct-image coherence.  It has rank
one because the generic restriction is the trivial line on a
connected elliptic curve.  It is torsion-free: multiplication
by a nonzero base function is injective on the line sheaf on
the smooth source, hence also on its direct image.  On the
disc this coherent sheaf is therefore free of rank one after
shrinking.  A nonzero local generator pulls back to a section
of \(L\), nonzero on generic fibres, whose divisor \(D\)
is vertical.  This gives the required meromorphic line
comparison from the proper source itself, without extending
an arbitrary punctured-disc scalar.  The component-degree
equations now say \(BD=0\).
Thus \(D\) is an integer multiple of the primitive integral
fibre vector \(F_{\rm prim}\).  If the actual fibre is
\(\nu F_{\rm prim}\), this need not yet be a base pullback.
On the normalization of the root which removes that multiplicity,
\(\mathcal O(kF_{\rm prim})\) is represented by the
base monomial \(z^k\).  Its deck character is
\(\zeta_\nu^k\), so the comparison records exactly the
integer delta character of the log cap.  This proves the stated
vertical rigidity with its full remaining collar ambiguity.

\end{enumerate}
\end{proof}


The statement concerns caps supplied with their actual extending
linearized root line and affine action at a nonzero-linearization
site.  Scalar-zero Mumford fillings are treated separately.  For any proposed broader
cap, constructing these maps remains necessary before applying
the statement.  The OS50 quadratic-root lead is kept separately
until its source marking and full integer topology are certified.

\begin{theorem}\label{night:P3K.16}
\StatementRecord{night:P3K.16}{proved}{S6\_06\_BLIND: GPT-6 Astra (OpenAI, ChatGPT Pro mode, effort Pro), recorded re-derivation in the checkpoint dated 18 November 2026}{\cref{night:P3F.8}, \cref{night:P3F.9}, \cref{night:P3K.13}, \cref{night:P3K.15}}
For the whole-fibre \(I_1\) cusp construction of P3F.9, together
with the actual nonzero-linearization affine caps of P3K.15, the
delta orbit has a global map to a closed Seifert three-manifold
over the original base \(\mathbb P^1\).  Its regular fibres are
connected elliptic tori.  Every sphere member has at most two
nontrivial effective gamma multiplicities.  For two such
multiplicities, write the marked local pairs as
\((m_1,n_1),(m_2,n_2)\), with each pair primitive, and let
\(b\) be the remaining integral gamma clutch.  The exact restriction is
\[
       m_1m_2b+m_2n_1+m_1n_2=\pm1.                 \tag{P3K.16.1}
\]
In particular \(\gcd(m_1,m_2)=1\).  Order-one clutches and the
finite zero-gamma subgroups are retained in the source and line
data; this restriction does not remove them.
\end{theorem}

\begin{proof}
\begin{enumerate}[label=\arabic*.]
\item \par\medskip\noindent\textbf{Construction of the map.}\enspace On an ordinary source disc,
the delta orbit of the contracting line quotient is the source
mapping torus of P3K.13.  Projection of that source surface to
its disc, together with the suspension coordinate, gives the
map to the disc times the gamma circle.  Straightening the norm
roof is over the source and its base.  On the whole-fibre \(I_1\)
cusp, the actual product of P3F.9 preserves the Seifert gamma
coordinate, so its map is exactly
\[
       E_{I_1}\times S^1_\gamma\longrightarrow
                    \Delta_q\times S^1_\gamma.
\]
This includes the nodal fibre and the circle-fixed sphere.
The specified annular period maps intertwine these projections;
their integral gamma shifts are their recorded Seifert clutches.

\item Here is the map at an arbitrary supplied cap, including its
finite kernel.  Use the actual maps \(A,B,C,H\) of P3K.15,
and write \(g=\gcd(m,n)\), \(m'=m/g\), \(n'=n/g\).
The root source has coordinate \(z\), with original base
coordinate \(q=z^m\).  The finite map \(C\) has faithful
base character of order \(g\).  Hence the smooth source
quotient \(\bar S=S'/\langle C\rangle\) has base coordinate
\(u=z^g\), and \(q=u^{m'}\).  If
\(rn'+sm'=1\), then \(\bar H\) acts on this base by
\[
                 u\longmapsto e^{2\pi i r/m'}u.
\]
Projecting the actual mapping torus of \(\bar S\) to the
mapping torus of this disc rotation gives the cap map.
The target is a solid torus: the explicit map
\((u,t)\mapsto(e^{2\pi i rt/m'}u,e^{2\pi it})\)
trivializes its disc bundle.  The regular Seifert gamma orbit
has multiplicity \(m'\), and its primitive marked meridian
relation is
\[
                     m'a-n'\gamma=0.             \tag{P3K.16.2}
\]
Indeed these two loops are the actual \(\bar A=\bar H^{n'}\)
and \(\bar B=\bar H^{m'}\) of P3K.15, not an inferred
replacement for their affine lifts.  This transports the stated
outer meridians and glues the target solid torus to the ordinary
target charts.  Doing so at all sites constructs the closed
three-manifold \(B\) and the map \(p:Y\to B\).

\item Off the finitely many discriminant and cap-core circles, the
map is a smooth proper submersion with connected elliptic
fibres.  Connectedness at a cap can be checked before any
topological conclusion: the \(g\) roots of a nonzero \(u\)
are permuted transitively by \(C\), so their quotient fibre
is one elliptic curve.  The same fact is preserved by the
suspension.  The finite-root and ordinary maps above agree
with the annular source maps, so there is no disconnected
generic-fibre substitution in this construction.

\par\medskip\noindent\textbf{Fundamental groups and the number of sites.}\enspace Every
loop in the smooth three-manifold \(B\) can be moved off the
finite union of discriminant and core circles: perturb its
finitely many smooth arcs transversely to these curves, whose
dimension sum is two, less than three.  Over the resulting
loop, a choice of horizontal complement to the submersion
kernel lifts the loop from a chosen point of its fibre.  The
lift exists for the whole compact interval because the
submersion is proper over that interval.  Its endpoint lies
in the same connected elliptic fibre as its starting point.
Join them by a path in that fibre.  The resulting closed loop
projects to the original loop followed by a constant path.
Thus \(p_*:\pi_1(Y)\to\pi_1(B)\) is surjective.
For a sphere member, P3F.8 gives \(Y=S^5\), so
\(\pi_1(B)=1\).

\item Remove discs around the nontrivial effective Seifert sites.
The gamma-circle bundle over the remaining punctured sphere
is trivial, since its base retracts onto a graph.  Choose a
section, retain the marked boundary loops \(a_i\), and write
its integer discrepancy from the actual cap framings as \(b\).
Van Kampen applied to the actual meridians (P3K.16.2) gives
\[
 \pi_1(B)=\left\langle a_1,\ldots,a_s,\gamma\ \middle|\
 [a_i,\gamma]=1,\quad a_i^{m_i}=\gamma^{n_i},\quad
                 a_1\cdots a_s\gamma^b=1\right\rangle.
                                                        \tag{P3K.16.3}
\]
Order-one sites can be included in this presentation and
eliminated using their actual relation \(a_i=\gamma^{n_i}\).
Their integers are then added to \(b\); their source and
line lifts are still retained in P3K.15.

\item Quotient (P3K.16.3) by \(\gamma\).  With at least three
nontrivial multiplicities, further killing all but three
generators maps this quotient onto
\[
       \langle x,y,z\mid x^{m_1}=y^{m_2}=z^{m_3}=xyz=1\rangle.
\]
This group is nontrivial for every \(m_i\ge2\), as the
following explicit real matrix representation proves.  On
\(\mathbb R^3\), take the symmetric bilinear form with matrix
\[
 G=\begin{pmatrix}
 1&-\cos(\pi/m_3)&-\cos(\pi/m_2)\\
 -\cos(\pi/m_3)&1&-\cos(\pi/m_1)\\
 -\cos(\pi/m_2)&-\cos(\pi/m_1)&1
 \end{pmatrix},
\]
and set \(R_i(v)=v-2(v,e_i)_G e_i\).  Direct substitution
gives \(R_i^2=I\).  The span of each pair \(e_i,e_j\)
has positive definite Gram matrix, since its determinant is
\(\sin^2(\pi/m_k)>0\).  It therefore has a one-dimensional
orthogonal complement even when \(G\) itself is degenerate.
On the pair span the two reflections have angle \(\pi/m_k\),
and their product has characteristic polynomial
\(z^2-2\cos(2\pi/m_k)z+1\); on the complement it is
the identity.  Thus \(R_2R_3,R_3R_1,R_1R_2\) have
orders exactly \(m_1,m_2,m_3\), and their product is
exactly the identity by cancellation of adjacent reflections.
Their first member is nonidentity.  This gives the promised
nontrivial representation, contradicting
\(\pi_1(B)=1\).  Thus \(s\le2\).

\par\medskip\noindent\textbf{The exact determinant and the marked base.}\enspace For two
sites eliminate \(a_2\) from (P3K.16.3).  The remaining
generators \(a_1,\gamma\) commute, and their relation matrix is
\[
 \begin{pmatrix}
      m_1&-n_1\\
      -m_2&-m_2b-n_2
 \end{pmatrix}.
                                                        \tag{P3K.16.4}
\]
Triviality of the group is therefore exactly (P3K.16.1).
Any common divisor of \(m_1,m_2\) divides its left-hand
side, which proves their coprimality.  With one exceptional
site the same calculation is \(|m_1b+n_1|=1\); with none
it is \(|b|=1\).  Equivalently one may insert multiplicity-one
solid tori and use the displayed two-site matrix in all cases.

\item The resulting target is the actual marked Hopf three-sphere,
with both signs retained.  Here is an explicit Heegaard
comparison, rather than an appeal to a classification by
fundamental group.  In the middle torus with basis
\((a_1,\gamma)\), its two meridians are
\[
 \mu_1=(m_1,-n_1),\qquad
 \mu_2=(-m_2,-m_2b-n_2).
\]
Put \(\delta=\det(\mu_1,\mu_2)=\pm1\).
The integer matrix with columns \(\mu_1,\delta\mu_2\)
has determinant one and sends these meridians to the two
coordinate meridians after inversion.  This torus map
extends across each solid torus, because in meridian-longitude
coordinates every required boundary shear extends by
\[
 (\rho e^{i\theta},e^{i\varphi})\longmapsto
 (\rho e^{i(\varepsilon\theta+k\varphi)},
                         e^{i\varepsilon'\varphi}),
 \qquad \varepsilon,\varepsilon'\in\{\pm1\}.
\]
The two target solid tori are the standard decomposition of
\(S^3\subset\mathbb C^2\).  In these coordinates the
original gamma vector is exactly \((\delta m_2,m_1)\).
Consequently its signed weighted Hopf action and its two
exceptional orders have both been transported explicitly.

\item Finally the replacement \(a_i\mapsto a_i\gamma^{k_i}\)
changes the integer records by
\[
         n_i\mapsto n_i+m_i k_i,\qquad
         b\mapsto b-\sum_i k_i.
\]
This is also an identity of actual cap quotients: replace
\(A_i\) by \(A_iB_i^{k_i}\), which generates the same
group with \(B_i\).  Its cyclic norm changes by
\(m_i k_i\gamma\); its finite kernel \(C_i\) is unchanged.
Thus residue representatives for \(n_i\) can be chosen
without deleting the recorded finite kernel, delta character,
or source-line map.  The determinant is invariant under
exactly these transformations.

\end{enumerate}
\end{proof}


The initial draft and its explicit audit boundary are preserved in
\path{code/part3_fourfold/history/part3_flop_cobordism_audit.md}.
Independent blind-control audit checked the residual semifree weights,
both associated-line signs, the common Cox delta gauge, and the extra
normal-sphere term.  Root audit additionally required and then accepted
the transported truncations and the free period-orbit quotient before
the two-overlap Cech gauge cancellation.  The orientation transport
between the two complex fixed-sphere markings was explicitly retained.
The exact matrix and gauge certificate is \path{code/part3_fourfold/cox_flop.json}.
This proof treats the two stated OS49 diagonal fans; it does not assign
the same quotient model to every periodic unimodular fan.

\begin{theorem}\label{night:P3K.18}
\StatementRecord{night:P3K.18}{proved}{S6\_06\_BLIND: GPT-6 Astra (OpenAI, ChatGPT Pro mode, effort Pro), recorded re-derivation in the checkpoint dated 18 November 2026}{\cref{night:P3K.16}}
In the actual whole-fibre \(I_1\) orbit projection of P3K.16,
every sphere member has primitive source-fibre multiplicities
at every logarithmic cap.  In particular no cap whose smooth
source quotient has a genuinely multiple central fibre occurs
in such a sphere member.
\end{theorem}

\begin{proof}
\begin{enumerate}[label=\arabic*.]
\item The signed meridian matrix in P3K.16 has determinant
\(\pm1\); its explicit genus-one Heegaard map identifies
the target \(B\) with \(S^3\).  Let \(K\subset B\)
be one oriented cap-core knot.  Suppose that the source
fibre multiplicities on that core have a common divisor
\(\nu>1\).  The linking-number class of this single knot is
\[
            \eta\in H^1(S^3\setminus K;\mathbb Z/\nu),
\]
with value one on its positive meridian.  It is represented,
for example, by intersection with an oriented Seifert surface
for \(K\), reduced modulo \(\nu\).  Pull it back to
\(Y\setminus p^{-1}K\).

\item The relevant local normal exponent is that of the actual
three-dimensional target, not the old branched base coordinate.
In a cap chart of P3K.16 it is
\[
        \xi=e^{2\pi i rt/m'}u,
        \qquad u=z^g,
\]
the transverse disc coordinate of its target solid torus.
At a smooth point of a source component of multiplicity
\(\nu a_j\), choose a complex normal coordinate \(w\).
Then \(u=w^{\nu a_j}v\), where \(v\) is a nonvanishing
holomorphic unit.  Consequently a small positive component
meridian has linking number \(\nu a_j\) under \(p\).
The pulled-back class kills every such meridian modulo
\(\nu\).

\item The remaining singular strata of the central source fibre,
crossed with the suspension direction, have real codimension
at least four in the smooth five-manifold \(Y\).  A loop
or a homotopy disc can be perturbed off these strata; after
this perturbation, van Kampen says that filling the smooth
divisorial pieces only kills their normal meridians.  Thus
the pulled-back class extends to a homomorphism
\[
                         \pi_1(Y)\longrightarrow\mathbb Z/\nu.
\]
It is nonzero.  Indeed a meridian loop of \(K\) can first
be perturbed off all the other critical and cap-core circles.
Lift it through the connected regular elliptic fibres by the
horizontal-path construction of P3K.16, and close the lifted
path inside its starting fibre.  The resulting loop lies in
\(Y\setminus p^{-1}K\) and evaluates to one.  This
contradicts \(\pi_1(Y)=0\).  Therefore the common divisor
is one at every cap.

\end{enumerate}
\end{proof}




\begin{theorem}\label{night:P3K.19}
\StatementRecord{night:P3K.19}{proved}{S6\_06\_BLIND: GPT-6 Astra (OpenAI, ChatGPT Pro mode, effort Pro), recorded re-derivation in the checkpoint dated 18 November 2026}{\cref{night:P3K.14}, \cref{night:P3K.15}, \cref{night:P3K.18}}
For the actual affine root-line caps of P3K.15 in the sphere
orbit scope of P3K.16, the root fibre is reduced.  A singular
root fibre has no nontrivial finite zero-gamma kernel.  At a
smooth elliptic root fibre that kernel is an actual normal
root change, and the invariant source coordinates and line
frames cancel it exactly.  Thus the cap has a primitive
effective pair without an unrecorded source or line quotient.

At an original multiplicative fibre \(I_N\), a sphere cap
retains the following necessary restrictions on the original
source data \(M=\mathcal O(aG-O+D)\), \(Q=bG\):
\[
              \gcd(N,b[G])=1,
              \qquad a[G]=0\quad\text{in }\mathbb Z/N.       \tag{P3K.19.1}
\]
For the explicit generator markings in P3F.2 these are
\(\gcd(N,b)=1\) and \(N\mid a\).  These restrictions
hold at logarithmically replaced multiplicative fibres as
well as at retained ones.
\end{theorem}

\begin{proof}
\begin{enumerate}[label=\arabic*.]
\item \par\medskip\noindent\textbf{The actual component maps and their multiplicities.}\enspace
Use \(A,B,C,H\) from P3K.15, with
\(g=\gcd(m,n)\), \(C=A^{m/g}B^{-n/g}\).
The components of the finite source quotient are precisely
the \(C\)-orbits of root components.  The residual
\(H\)-action is transitive on those quotient components
by P3K.15.  Since \(C,H\) generate the same group as
\(A,B\), the latter group is transitive on all root
components.  Both maps preserve the multiplicities of
the root fibre \(\operatorname{div}(z)\).  The root
source retains the pulled-back section of the original
elliptic surface, so these multiplicities have greatest
common divisor one: intersecting the fibre with that
section gives one.  Transitivity makes every multiplicity
equal, and hence all are one.

\item Let \(H_j\subset C\) be the stabilizer of a general
point on root component \(C_j\).  It is the subgroup
acting pointwise on that component.  Since \(C\) is
central and \(A,B\) are transitive, all these subgroups
are the same subgroup \(H_0\), of order \(h\).
At a general point its action on the normal parameter has
order exactly \(h\), because the root-base character is
faithful.  The finite quotient is therefore ramified to
order \(h\) along each root component.  The quotient
base coordinate is \(u=z^g\).  If \(\bar C\) is an
image component, comparing divisors in
\(\pi^*u=z^g\) gives its multiplicity exactly
\[
                              g/h.                         \tag{P3K.19.2}
\]
This computation includes component permutations and the
degree of the component map; it uses the generic normal
ramification index, not a declaration that the quotient is
unramified.  P3K.18 forces \(g/h=1\).  Hence \(C\)
fixes every root component pointwise.

\item At a singular point of a reduced root fibre, a nontrivial
finite-order surface automorphism cannot fix that fibre
germ pointwise.  Average a local coordinate map over the
finite group to linearize the action, with derivative the
identity at the base point of the averaging map.  The fixed
germ of a nonidentity linear map in two complex dimensions
is either a point or a smooth line.  Neither can contain
the reduced singular fibre germ.  The derivative must
therefore be the identity; the linearization then makes
the automorphism the identity in a neighbourhood, and
analytic continuation makes it the identity on the root
source.  This contradicts its faithful root-base character.
Thus \(g=1\) at a singular reduced root fibre.

\item At a smooth elliptic root fibre, the group fixes that curve
pointwise.  Near each of its points average a tangent
coordinate to obtain an invariant coordinate \(e\), and
keep the given normal base coordinate \(z\).  These are
coordinates by the inverse function theorem, and the
generator acts exactly by
\[
                         (e,z)\longmapsto(e,\zeta_g z).
\]
The actual quotient coordinates are \((e,u=z^g)\).
They identify the root source with the pullback of its
smooth quotient along \(u=z^g\).  The invariant local
line frames of P3K.15 identify the root line with the
pullback quotient line as well; no fibre character remains
because every point stabilizer acts trivially on the line.
These local maps agree on overlaps by their quotient
construction and give the required global map after
shrinking the source disc.  On the quotient,
\[
  \bar A=\bar H^{n/g},\qquad
  \bar B=\bar H^{m/g},\qquad
  q=u^{m/g}.
\]
Thus the root order and gamma numerator reduce to the
primitive pair \((m/g,n/g)\) by an exact source-and-line
quotient.  The marked meridians and any remaining integral
delta collar factor are still those transported in P3K.15.

\par\medskip\noindent\textbf{The multiplicative root and its integral component action.}\enspace
On the minimal \(m\)-root of an \(I_N\) neighbourhood,
the Tate period is \(z^{Nm}\), and the components of
the \(I_{Nm}\) fibre are labelled by
\(\operatorname{ord}_z Z\pmod{Nm}\).  The canonical
root deck sends \(z\) to \(\zeta_m z\) and preserves
these labels.  A translation by a Tate point of valuation
\(\alpha\) adds \(\alpha\) to them.  If the original
point \(Q\) has component \(c\in\mathbb Z/N\), its
root lift has component \(mc\).  The actual affine
root map \(A\) therefore acts by addition of some
integer \(\alpha\), and \(B\) by addition of \(mc\).
The norm identity \(A^m=B^n\) gives
\[
       m\alpha\equiv nmc\pmod{Nm},\qquad
                   \alpha\equiv nc\pmod N.          \tag{P3K.19.3}
\]
This is the actual source-component action of the root
deck and its translation.  An additional inversion would
change the marked linear monodromy from trace two to
trace minus two and is not the outer monodromy of this
\(I_N\) source cap.

\item Every component orbit is a coset of the subgroup generated
by \(\alpha,mc\) in \(\mathbb Z/(Nm)\).  Its number
is \(\gcd(Nm,\alpha,mc)\).  By (P3K.19.3) every
common divisor of \(N,c\) divides this number.  The
proved transitivity thus forces \(\gcd(N,c)=1\).

\item The root line is the pullback of the quotient line of
P3K.15 and consequently has degree zero on every root
component.  Its generic Picard point is the pulled-back
point \(aG\), whose component is \(ma[G]\pmod{Nm}\).
Its extension differs from the section line by an integral
vertical divisor, by the coherent-direct-image argument
of P3K.15.  The component-degree equations are therefore
\[
                  e_{ma[G]}-e_0+B_{Nm}d=0.
\]
Multiply by the cyclic component labels and sum.  The
sum of labels times \(B_{Nm}d\) is divisible by \(Nm\),
by the exact cyclic difference calculation of P3K.14.
It follows that \(ma[G]=0\pmod{Nm}\), or equivalently
\(a[G]=0\pmod N\).  This proves (P3K.19.1).

\item The argument applies to the minimal root source used in
the actual Tate cap.  A superfluous equivariant blowup
cannot evade the transitivity statement: blowing up a
node introduces a component of different multiplicity,
and a rational vertical minus-one component in a reduced
transitive fibre would make every component a rational
minus-one curve.  The fibre intersection equations would
then give each vertex weighted valency one, so the
connected fibre would consist of two rational curves
with one transverse intersection.  Its arithmetic genus
would be zero, contradicting the elliptic fibre genus
one.  Thus the transitive root is relatively minimal.

\end{enumerate}
\end{proof}


These conclusions concern the actual root-line caps and the
global projection explicitly constructed in P3K.15–16.  They
do not assert a replacement of an arbitrary periodic cusp fan
by that fixed-sphere orbit model.  At such a fan the requisite
orbit comparison remains a separate part of the exhaustive
scope question.

\begin{theorem}\label{night:P3K.20}
\StatementRecord{night:P3K.20}{proved}{S6\_06\_BLIND: GPT-6 Astra (OpenAI, ChatGPT Pro mode, effort Pro), recorded re-derivation in the checkpoint dated 18 November 2026}{\cref{night:P3K.12}, \cref{night:P3K.14}}
For a semistable multiplicative source fibre \(I_N\), a
whole-fibre zero of order \(e\) has period determinant
\(Ne\) and central Mumford fibre Euler number \(2Ne\),
for every periodic unimodular triangulation.  At an
\(I_s^*\) zero, \(s>0\), the smooth quadratic
semistable Mumford root, followed by its actual monomial
involution quotient and toric quotient resolutions, has
Euler contribution at least \(4se\).  Potentially good
zero sites are excluded from sphere members by P3K.12.

In this actual smooth-root toric construction, a sphere
member therefore has exactly one whole-fibre zero: a
simple \(I_1\) zero.  For a rational rank-one source,
writing its generic line point and translation as
\(aG\) and \(bG\), respectively, the height denominator
\(d=1/h(G)\) satisfies \(ab=d\) after choosing the
contracting generator and the sign of \(G\).
\end{theorem}

\begin{proof}
\begin{enumerate}[label=\arabic*.]
\item \par\medskip\noindent\textbf{The actual line-frame determinant.}\enspace On the infinite
Tate chain let the source period and suspension period have
valuation vectors
\[
                    h=(N,\alpha),\qquad g=(-c,\beta).
\]
Write the fibre-coordinate graph of the actual line extension
as heights \(h_j\), so \(h_{j+N}=h_j+\alpha\).
An arbitrary vertical line divisor changes these heights
integrally; it is not omitted.  The line lift has the same
zero order \(e\) on every component precisely when its
actual fibre-coordinate transition has defect
\[
                      h_j+\beta-h_{j-c}=e.             \tag{P3K.20.1}
\]
Put \(r=\gcd(N,c)\), and iterate this identity \(N/r\)
times.  The period relation on the left gives
\[
      -\frac c r\alpha=\frac N r(\beta-e),
      \qquad
      \det(h,g)=N\beta+c\alpha=Ne.                    \tag{P3K.20.2}
\]
This is an identity of the actual component line frames.
It proves cancellation of the off-diagonal valuation terms
without changing or normalizing the vertical line divisor.

\item A fundamental polygon of this period lattice has area
\(Ne\).  Every unimodular triangle has area \(1/2\),
so it contains exactly \(2Ne\) triangles, counted with
their periodic boundary identifications.  In the central
Mumford fibre each triangle gives one zero-dimensional
torus orbit.  All other orbit strata are positive-dimensional
complex tori and have Euler number zero.  Additivity of
Euler number over the finite orbit stratification gives
\(2Ne\).  No vertical edge or semifree property of the
chosen delta circle is assumed in this calculation.

\par\medskip\noindent\textbf{The starred multiplicative quotient.}\enspace The quadratic
semistable root of an \(I_s^*\) fibre is \(I_{2s}\),
and the pulled-back whole-fibre zero has order \(2e\).
The determinant just proved is \(4se\).  The smooth
root Mumford fibre therefore has \(8se\) zero-dimensional
torus strata and Euler number \(8se\).

\item The quadratic deck is monomial on the central strata.
Indeed on the source nodal identity torus it is
\(Z\mapsto uZ^{-1}\), with a nonzero central unit
\(u\).  It commutes with the scalar circle on the line,
so in a line frame it sends \(V\mapsto f(Z)V\).
The source, line and their linearization are algebraic,
and this central torus map is regular and invertible.
Thus \(f\) is a unit of \(\mathbb C[Z,Z^{-1}]\),
so \(f(Z)=vZ^{\ell}\).  The induced maps on preserved
boundary torus strata are the restrictions of this same
integral monomial map.  Nonpreserved strata are permuted.
An invariant refinement can be used for descent; every
subdivision increases the number of maximal cones and
does not lower the preceding root Euler count.

\item There is also an explicit invariant unimodular height-one
choice, without a restriction to the rank-one generator
marking.  Write the actual monomial deck on valuations as
\(R(x,y)=(-x,y-jx)\), with \(j=-\ell\in\mathbb Z\).
It sends the source period \(h=(2s,\alpha)\) to
\(-h\); comparison of the second coordinate gives
\[
                         j=\alpha/s.
\]
Thus the integrality of this shear follows from the
actual deck morphism.  The whole-zero identity is now
\(2\beta+cj=4e\).  For odd \(j\), the integral
shear \(y'=y-(j-1)x/2\) gives
\(R(x,y')=(-x,y'-x)\).  The positive unit-square
triangulation is invariant, since its edge vectors
\(u,v,u+v\) go to \(-(u+v),v,-u\).
For even \(j\), use \(y'=y-jx/2\).  The reflection
is then \((-x,y')\) and the periods are
\((2s,0),(-c,2e)\).  Give the square with lower-left
corner \((v,k)\) its diagonal according to the parity
\[
                       v+c\left\lfloor k/(2e)\right\rfloor.
\]
Both periods preserve this parity.  Reflection sends the
square to \((-v-1,k)\), changes the parity, and exchanges
the two diagonals.  This is the required invariant
unimodular periodic triangulation.

\item In the marked rank-one case the generic point \(aG\)
has root Tate lift \(Z_G^a\), hence \(\alpha=as\),
\(c=bs\), and \(j=a\).  Changing that lift by a
Tate-period power \(Q_T^k\) is the actual line-frame
map \(V\mapsto Z^kV\).  Vertical component divisors
change the graph heights periodically and do not change
the quasiperiod.  These explicit maps recover the
rank-one spelling while the preceding deck argument
proves the starred-zero bound before that specialization.

\item The canonical-form normalization also identifies the
quotient-resolution comparison.  In Weierstrass coordinates
\(q=z^2,x=z^2X,y=z^3Y\), one has exactly
\[
                  (dx/y)\wedge dq=2(dX/Y)\wedge dz.
\]
After wedging the logarithmic line form, this is the
nowhere-zero Kulikov form
\(\Omega=d\log Z\wedge d\log V\wedge dz\).
The source inversion and \(z\mapsto-z\) give two
minus signs, while the additional \(Z^{\ell}\) in the line
map drops out of the wedge.  Thus the deck preserves
\(\Omega\).  Its nontrivial stabilizers on a smooth
root have eigenvalues \((-1,-1,+1)\), and their quotient
is a transverse \(A_1\) surface singularity along a fixed
curve.  In a resolution chart \(v=ut,w=ut^2\) for
\(v^2=uw\), the invariant two-form is
\(\tfrac12du\wedge dt\); it is nonvanishing.  This
is the actual crepant quotient resolution.

\item A divisor in this resolved model cannot be contracted to
a smooth model carrying the same nonvanishing form: the
Jacobian of a birational morphism vanishes along every
divisor mapped into codimension at least two, contradicting
the equality of those forms on their common open set.
Changes of the height-one Mumford triangulation are already
covered by the period-area count for each such fan.  In
particular a pure base-root contraction, whose pullback
canonical form has a ramification zero, is not being
silently used as a comparison of these normalized forms.

\item On an invariant torus stratum a fixed point equation is
\(x^{M-I}=a^{-1}\).  Its solution is empty or a coset
of the kernel of an integral torus homomorphism.  That
kernel is a finite union of complex tori.  Its Euler
number is nonnegative: positive-dimensional components
give zero and points give one.  Euler additivity therefore
gives \(\chi(X'^{\iota})\ge0\) on the root central fibre.
The finite-action orbit count, applied stratumwise, is
\[
 \chi(X'/\langle\iota\rangle)
       =\frac{\chi(X')+\chi(X'^{\iota})}{2}
       \ge4se.                                        \tag{P3K.20.3}
\]

\item For completeness, the quotient-resolution correction has
a definite sign in the stated toric construction.  Along
each fixed torus stratum the normal cyclic quotient has
constant diagonal characters.  Resolve it by a toric
subdivision.  A compact exceptional fibre has finitely
many torus orbits; its positive-dimensional orbits have
Euler number zero and it has at least one zero-dimensional
orbit.  Its Euler number is consequently at least one.
Euler additivity over the base fixed torus strata shows
that replacing a point-normal fibre by this exceptional
fibre adds a nonnegative quantity.  This proves the same
lower bound after the actual toric quotient resolutions.
Flips between unimodular height-one triangulations keep
the triangle count unchanged.

\par\medskip\noindent\textbf{The actual invariant-twist map at a semistable zero.}\enspace
The logarithmic modification in the supplied construction
\texttt{con-log} uses an extending root automorphism and
an invariant lattice translation, with punctured comparison
\(\sigma=(\log z)/(2\pi i)\,v\).  A whole-zero line
lift itself is not invertible on the zero fibre.  If the
same invariant-translation formula is interpreted on its
semistable Mumford torus, its quotient can nevertheless
be described exactly.  Full-rank unipotent monodromy has
invariant lattice equal to the primitive vanishing-torus
lattice.  Thus its twist has no gamma component.
On \(q=z^m\) write its deck as
\[
 (Z,V,z)\longmapsto
       (\zeta_m^r Z,\zeta_m^t V,\zeta_m z).
\]
The invariant coordinates are
\[
               Z'=Z/z^r,\qquad V'=V/z^t,\qquad q=z^m.
\]
They preserve every original period \(c(q)q^\Gamma\).
On cocharacters the quotient lattice and its exact
isomorphism are
\[
 \mathbb Z^3+\mathbb Z(r,t,1)/m
      \longrightarrow\mathbb Z^3,
 \qquad (x,y,s)\longmapsto(x-rs,y-ts,ms).
\]
The two horizontal basis vectors and \((r,t,1)/m\)
map to the three standard basis vectors.  Conjugating
the root period shear with translation \(m\Gamma\)
by this map gives the original shear \(\Gamma\).
Hence a smooth Kulikov target is again a Mumford fan
for the original period lattice and has the same Euler
count.  This gives the quotient directly, including
cases where resolving its root first would add curves.
No noninvariant affine operation at a zero germ has
been assumed beyond these actual published maps.

\par\medskip\noindent\textbf{The global zero count and height equation.}\enspace Away from
zero fibres, the line-complement construction has no fixed
point for the scalar circle, including its finite affine
quotients.  Its compact fibre Euler number is zero.
This also follows directly from the corresponding real
mapping-torus description.  On the smooth part of the
base the fibre is a complex torus and also has Euler
number zero.  Stratifying the base by its finitely many
special values therefore expresses the total Euler number
as the sum of the zero-site contributions above.

\item A smooth or potentially good zero site has the elliptic
fixed component excluded by P3K.12.  Every semistable
multiplicative zero contributes the positive even integer
\(2Ne\).  Every starred multiplicative zero in the
stated quotient-resolution construction contributes at
least four.  Since a six-sphere has Euler number two,
there is exactly one contribution, with \(N=e=1\).

\item The literal whole-fibre Hom relation and the telescoping
height calculation of P3K.14 give
\(\deg E=ab\,h(G)\).  Hence \(ab=d\).
For the pole version of the construction replace the
actual suspension generator \(B\) by \(B^{-1}\).
Their cyclic groups are identical, and the corresponding
period lattices are identical with the gamma generator
reversed.  Thus their toric quotient is identified by the
identity on its common charts; the inverse lift has zeros
in place of poles.  Finally replacing the name of the
generator \(G\) by \(-G\) makes the two coefficients
positive without changing the represented line point or
translation.  This gives the stated positive factorization.

\end{enumerate}
\end{proof}


The starred-fibre bound retains the normalized canonical form
and its actual monomial quotient/resolution maps.  The
invariant-twist paragraph supplies the direct quotient map
for the stated logarithmic formula; it does not assume an
additional noninvariant zero-germ operation.  The later
Seifert restrictions still require the orbit map of K16;
the present Euler count itself applies to every periodic
unimodular Mumford triangulation and does not impose a
vertical-edge or semifree delta condition.

\begin{theorem}\label{night:P3K.21}
\StatementRecord{night:P3K.21}{proved}{S6\_06\_BLIND: GPT-6 Astra (OpenAI, ChatGPT Pro mode, effort Pro), recorded re-derivation in the checkpoint dated 18 November 2026}{\cref{night:P3F.1}, \cref{night:P3F.2}, \cref{night:P3K.14}, \cref{night:P3K.19}, \cref{night:P3K.20}}
In the established sphere orbit and root-line cap scope, a
nonzero-linearization \(I_s^*\) source fibre can occur only
with \(s=1\), a primitive quadratic cap, and
\[
                         a\equiv0\pmod2,
                         \qquad b\equiv2\pmod4.       \tag{P3K.21.1}
\]
The two \(I_3^*\) rank-one source configurations are excluded.
Combined with P3K.19–20 and the actual source table, the
remaining generic source coefficients are the 26 pairs in
the 16 profiles displayed below.  This is a proved finite
restriction on the generic source data; it does not declare
all remaining affine cap tuples to be sphere tuples without
their integral gluing calculation.
\end{theorem}

\begin{proof}
\begin{enumerate}[label=\arabic*.]
\item \par\medskip\noindent\textbf{The marked quadratic root.}\enspace The quadratic root of
\(I_s^*\) is the Tate \(I_{2s}\) fibre.  Its deck
acts on component labels by \(j\mapsto-j\), since
it is elliptic inversion and fixes the zero component.
The two invariant components are \(j=0,s\).  On each,
inversion has two fixed smooth points.  Near such a point
the two local coordinates have signs \((-1,-1)\), so
the quotient has equation \(v^2=uq\).  Resolving it
adds a multiplicity-one outer component.  The quotient
of the cycle is a path of \(s+1\) multiplicity-two
components, with two such outer leaves at each end.
This is the actual affine \(D_{s+4}\) graph.

\item The end \(j=0\) contains the zero section.  Remove
its corresponding outer leaf to obtain the finite
\(D_{s+4}\) intersection matrix.  Its inverse has
diagonal entry one at the other leaf of that end, and
\((s+4)/4\) at the two far-end leaves.  These entries
follow directly from the graph matrix; they are also the
local-height entries used in P3F.2.  The primitive
rank-one generator meets a spinor leaf, so its quadratic
root component is exactly \(s\), not zero.
The further degree-\(m/2\) root multiplies component
labels by \(m/2\).  Thus on an even \(m\)-root
the source fibre is \(I_{sm}\), the generator has
component \(sm/2\), and the canonical deck still
acts by reflection \(j\mapsto-j\).

\item An odd root leaves a starred fibre with unequal
multiplicities, excluded by the reduced-root conclusion
of P3K.19.  For an even root write the source component
action of the affine cap map as
\(A(j)=-j+\alpha\).  The suspension map acts as
\(B(j)=j+bsm/2\).  Its root fibre is singular, so
P3K.19 makes the pair \((m,n)\) primitive.  In
particular \(n\) is odd.  Since \(A^m\) is the
identity on components, the actual norm relation
\(A^m=B^n\) forces \(b\) even.  Then \(B\)
has trivial component action.  A reflection has orbits
of size at most two; transitivity forces \(sm\le2\).
Hence \(s=1,m=2\), and \(\alpha\) is odd.
The multidegree-zero root line has generic component
\(asm/2\); its component-Laplacian equation forces
\(a\) even, as in P3K.19.

\item There is one further unit constraint, which the component
labels alone would miss.  Write the quadratic Tate period
as \(Q_T=z^2u(z^2)\) and the generator as
\(Z_G=zv(z)\).  Its actual descent under inversion says
\(Z_G(-z)=Q_T/Z_G(z)\).  Thus
\(v(z)v(-z)=-u(z^2)\), and the identity-component
limit of \(2G\) is exactly \(-1\).  For the
affine reflection write
\(A(Z,z)=(R(z)/Z,-z)\), where
\(R=z^\alpha w(z)\).  Its square multiplies
\(Z\) by \(R(-z)/R(z)\), whose central value is
\((-1)^\alpha=-1\).  The central value of the
right-hand side \(B^n\) is
\((-1)^{bn/2}\).  Since \(n\) is odd, equality
forces \(b\equiv2\pmod4\).  This proves (P3K.21.1)
with both the component and unit parts of the norm retained.

\par\medskip\noindent\textbf{The finite generic factor list.}\enspace A simple whole
zero requires an \(I_1\) source value and \(ab=d\)
by P3K.20.  The three profiles with no \(I_1\) are
therefore removed.  P3K.19 imposes \(N\mid a\)
and \(\gcd(N,b)=1\) at every multiplicative fibre.
The preceding argument removes both \(I_3^*\) profiles
and imposes (P3K.21.1) at \(I_1^*\).  Enumerating the
positive divisors of each of the eight height denominators
now gives exactly the following list.

\begin{center}\small\begin{tabular}{c l l}\hline
OS row & Source fibre profile & Remaining \((a,b)\) \\
\hline
43 & \(III^*+3I_1\) & \((1,2),(2,1)\) \\
43 & \(III^*+II+I_1\) & \((1,2),(2,1)\) \\
45 & \(I_8+4I_1\) & \((8,1)\) \\
45 & \(I_8+II+2I_1\) & \((8,1)\) \\
47 & \(I_7+I_2+3I_1\) & \((14,1)\) \\
47 & \(I_7+I_2+II+I_1\) & \((14,1)\) \\
47 & \(I_7+III+2I_1\) & \((7,2),(14,1)\) \\
49 & \(IV^*+I_2+2I_1\) & \((2,3),(6,1)\) \\
49 & \(IV^*+III+I_1\) & \((1,6),(2,3),(3,2),(6,1)\) \\
50 & \(I_1^*+I_3+2I_1\) & \((6,2)\) \\
50 & \(I_1^*+IV+I_1\) & \((2,6),(6,2)\) \\
55 & \(I_5+I_4+3I_1\) & \((20,1)\) \\
55 & \(I_5+I_4+II+I_1\) & \((20,1)\) \\
56 & \(I_5+I_3+I_2+2I_1\) & \((30,1)\) \\
56 & \(I_5+I_3+III+I_1\) & \((15,2),(30,1)\) \\
56 & \(I_5+IV+I_2+I_1\) & \((10,3),(30,1)\) \\
\hline\end{tabular}\end{center}

\item The script \path{code/part3_fourfold/source_reduction.py} starts
from all 21 source profiles of P3F.1, retains their marked
height denominators, and executes these exact tests.  It
also exports the affine \(D\) matrices and their inverse
entries, so the starred marking is independently checkable.

\par\medskip\noindent\textbf{The indispensable vertical correction in row 50.}\enspace
The list allows arbitrary vertical line divisors, as the
literal construction does.  For its two row-50 factors,
the necessary correction is explicit.  Label the affine
\(D_5\) fibre by central vertices \(C_0,C_1\),
identity outer leaf \(L_0\), its order-two mate \(L_1\),
and spinor leaves \(R_0,R_1\).  The two left leaves
meet \(C_0\), the two right leaves meet \(C_1\),
and \(C_0,C_1\) meet.  Translation by \(2G\) or
\(6G\) swaps \(L_0,L_1\) and \(R_0,R_1\),
while fixing the two central components.

\item For \(a=2\) or \(6\), the section \(aG\) meets
\(L_1\), so the bare section-line degree vector is
\(e_{L_1}-e_{L_0}\).  If \(B_D\) is the actual
affine intersection matrix, then
\[
 e_{L_1}-e_{L_0}-B_D e_{L_0}
                  =e_{L_0}+e_{L_1}-e_{C_0}.          \tag{P3K.21.2}
\]
The right-hand side is invariant under the stated
translation.  Thus the actual line
\(\mathcal O(aG-O-L_0)\), with \(L_0\) here a
vertical fibre component, has the required invariant
component degrees.  At the other reducible fibre the
displayed factor pairs likewise have invariant component
degrees: at \(I_3\) the line point is in the identity
component; at \(IV\), either the line point is in the
identity component or the translation preserves all
components.  The Hom ratio is generically trivial and
has degree zero on every fibre component.  Its vertical
divisor is therefore a sum of whole fibres, and P3K.14
gives its base degree one.  Choosing its section places
the simple zero at the selected \(I_1\).

\item This proves that these corrected source line data satisfy
the literal whole-fibre Hom condition.  It does not
replace the remaining integral cap/gluing verification
for their total threefolds.

\end{enumerate}
\end{proof}




\begin{theorem}\label{night:P3K.26}
\StatementRecord{night:P3K.26}{proved}{S6\_06\_BLIND: GPT-6 Astra (OpenAI, ChatGPT Pro mode, effort Pro), recorded re-derivation in the checkpoint dated 18 November 2026}{\cref{night:P3F.24}, \cref{night:P3F.31}, \cref{night:P3F.9}, \cref{night:P3K.5}, \cref{night:P3K.6}, \cref{night:P3K.68}, \cref{night:P3K.7}}
For every one-cap sphere family of P3F.37 and P3F.41--43, with the actual parameter $\lambda=r^m$, the positive single $r$-loop has the same signed return integer for the two norms $n=1$ and $n=-1$. With source height $h(G)=1/d$, line point $aG$, original capped Euler number $e$ and resolved root Euler number $e'$, it is
\[
A^{\rm phys}=\frac{ma^2}{2d}-\frac{m(12-e)+e'}{12}.
\]
The exact values are listed below. By the actual dual-line and inverse-closing map P3K.68, the physical Hopf-gluing class is $A^{\rm phys}\bmod28$ with $\mu=A^{\rm phys}/28\bmod\mathbb Z$. Every nonzero residue gives the nonproduct mapping torus described in P3F.24.
\end{theorem}
\begin{proof}
\begin{enumerate}[label=\arabic*.]
\item \par

\medskip\noindent\textbf{AR.1. Physical characters and the complete parameter square}

\item Move the capped source value to infinity and denote the selected
scalar-zero nodal value by \(t_c\). Near that cusp write
\(q=t-t_c\), \(s=\log q/(2\pi i)\), and
\(v=\log r/(2\pi i)\). In physical characters the source
Abel coordinate is \(a=\zeta_1\) and the scalar line coordinate
is \(V=e^{2\pi i\zeta_2}\). Let \(g\) be the marked Abel
logarithm of \(Q=bG\). The actual suspension map is
\[
 B_\lambda(a,V)=(a-g,e^{2\pi iC}qr^mV),\qquad
 \gamma_{\rm out}=(g,-C-s-mv),\qquad \delta=(0,1).
                                                        \tag{AR.1}
\]
Thus \(B_\lambda\) is translation by \(-\gamma_{\rm out}\),
and \(\gamma_{\rm out}=B_\lambda^{-1}\) is exactly the
outward loop used in the cap meridian. Delta is positive physical
line rotation. The functions \(g,C\) have bounded holomorphic
characters in the specified nodal chart; their constants and
any fixed line character are retained.

\item The annular arrow is \(F_c:\text{ordinary}\to\text{cap}\),
where \(c=\log\lambda/(2\pi i)\). Its scalar variation is
\(n(c-c_0)s_\infty\delta/m\), as follows directly from
the positive \(A_r\) root cut. At \(c\mapsto c+m\) the
cap is identical. The inverse closing \(c+m\to c\) is
\(F_c^{-1}F_{c+m}\), whose global multiplier is
\((t-t_c)^{-n}\). Its positive deck is consequently
\[
              (v,\zeta)\longmapsto(v+1,\zeta+n s\delta).
                                                        \tag{AR.2}
\]
This is fixed by the actual ordinary-to-cap arrow, including its
base unit; it is not deduced from a later orientation convention.

\item On the auxiliary cusp cover \(q=Q_0^{2m}\), use the two
extending origins
\[
\begin{split}
 o_n(s,v)&=n\bigl(-gs/m,\ Cs/m+s^2/(2m)+sv\bigr),\\
 \sigma_n(s,v)&=n\bigl(gv,\ -Cv-m(v^2-v)/2\bigr).
\end{split}                                             \tag{AR.3}
\]
They satisfy the exact identities
\[
\begin{split}
 o_n(s,v+1)-o_n(s,v)&=n s\delta,\\
 \sigma_n(s,v+1)-\sigma_n(s,v)&=n s\delta+n\gamma_{\rm out},\\
 o_n(s+2m,v)-o_n(s,v)&=-2n\gamma_{\rm out}+2mn\delta,\\
 \sigma_n(s+2m,v)-\sigma_n(s,v)&=0.
\end{split}                                             \tag{AR.4}
\]
The characters of \(\sigma_n\) are bounded units for either
sign, so they extend in the actual Mumford charts. The ordinary
origin removes (AR.2). On their common punctured collar their
difference is the explicit translation
\[
 F_n=\sigma_n-o_n
 =n\left(g(s/m+v),-C(s/m+v)-\frac{(s+mv)^2}{2m}+\frac{mv}{2}\right).
\]
Its angular change is
\[
 F_n(s,v+\theta)-F_n(s,v)
   =n\theta\gamma_{\rm out}
                   -\frac{nm}{2}\theta(\theta-1)\delta.  \tag{AR.5}
\]
The quadratic line gauge vanishes at both endpoints. The whole
displayed difference changes under \(s\mapsto s+2m\) by
\(-2mn\theta\delta\), an integral retained line winding.
On the orbit the actual positive \(r\)-loop is therefore
\(n\gamma_{\rm out}\). The auxiliary root is in the cusp
base and does not replace this loop by a parameter cover.

\item Changing the root of the norm unit by \(\zeta_m^d\) adds a
fixed delta character to the cap. Its actual graph line is still
the norm sphere of \(N\); that character is retained in the
sheet identifications. It is independent of \(v\), so it cancels
in (AR.2)--(AR.5). The affine and line conjugations of F44 are
also independent of \(r\), preserve \(B_0\), and transport
the annular map. Thus the same identities apply after those
exact local comparisons, without dropping their boundary gauges.

\par

\medskip\noindent\textbf{AR.2. The signed Seifert map and the actual complex page}

\item Use the actual root coordinate \(q_\infty=z^m\) and the
cap norm coordinate
\[
 u=m\gamma_{\rm out}-n\frac{\arg q_\infty}{2\pi}.
\]
Its quotient relation is
\((z,u)\sim(\zeta_mz,u-n)\), with \(u\) of period \(m\).
For both signs let \(R>0\) solve
\(R^2+|z|^2R^{2/m}=1\). The actual map is
\[
 (z,u)\longmapsto(z_1,z_2)=
 \left(zR^{1/m}e^{2\pi i n u/m},R e^{2\pi i n u}\right).
                                                        \tag{AR.6}
\]
It is invariant under the quotient relation because \(n^2=1\).
The positive derivative of the radial equation and the inverse
\(R=|z_2|\), followed by a lift of its argument and
\(z=z_1R^{-1/m}e^{-2\pi inu/m}\), prove that it is a smooth
marked solid-torus diffeomorphism. Changing the lift is exactly
a power of the original root action. Its base projection is
\[
 t=z_2/z_1^m=q_\infty^{-1}.
\]
On the ordinary base disc the same map is
\(z_1=\rho e^{2\pi in\gamma_{\rm out}}\),
\(z_2=t z_1^m\), with the unique positive \(\rho\) satisfying
\(\rho^2+|t|^2\rho^{2m}=1\). The two displayed maps agree
on the overlap by the exact equation for \(u\). They identify
the actual three-dimensional orbit-base gluing with \(S^3\),
not merely its abstract Seifert invariants.

\item The positive weighted Hopf coordinate is
\[
                         \beta=n\gamma_{\rm out}.       \tag{AR.7}
\]
Choose a regular binding value \(\tau\). The affine complex
page \(\eta-\tau\xi^m=1\) embeds in this sphere by
\((\xi,\eta)\mapsto(\rho\xi,\rho^m\eta)\), where \(\rho>0\)
puts it on the sphere. Its projection is
\[
                         t=\tau+\xi^{-m}.                \tag{AR.8}
\]
The compactified page base is \(\mathbb P^1_\xi\). It is
ramified to order \(m\) at the additive cap and at the
regular binding, and it is unramified at the selected nodal
cusp and at every other source singular fibre. In particular
the selected cusp contributes \(m\) separate \(I_1\) fibres.

\item The root-chart page equation is
\[
 n u+\frac{\arg(1-\tau z^m)}{2\pi}=\frac\phi{2\pi}
                                  \pmod{\mathbb Z}.     \tag{AR.9}
\]
Its \(m\) sheets are cyclically permuted by the actual
\(A_r\), which decreases the sheet index by one for either
norm sign. Choosing sheet zero gives the exact source-root
slice. For phase zero the page coordinates satisfy
\[
 \xi^m=\frac{z^m}{1-\tau z^m},\qquad
 \eta=\frac1{1-\tau z^m}.                                 \tag{AR.10}
\]
These are holomorphic, independent of \(n\). Thus the page
inherits the actual complex source orientation for both signs.

\par

\medskip\noindent\textbf{AR.3. The actual root line, finite-base torsor, and square}

\item On the cap, the sheet-zero map is
\(p\mapsto[p,u_0(z(p))]\) on the complete resolved root
surface, including its exceptional components and tangencies.
Every root sheet has a unique representative in this one. Indeed
a relation taking the same norm level to itself has
\(nk+mh=0\), hence \(k=mj,h=-nj\); the actual identity
\(A_r^{mj}B_\lambda^{-nj}=1\) then proves uniqueness.
No residual finite stabilizer quotient remains. On the circle line,
normalize a nonzero vector in \(N_p\) to have the indicated
norm coordinate. This is an explicit equivariant map from the
norm sphere of \(N\) to the actual delta bundle on the page.
It uses the resolved line identities of F41--43, including the
good-starred toric corrections, the good III/IV exceptional
differences, and the cancelling II correction. It therefore
identifies the actual line, not only its generic Picard point.

\item On the ordinary finite-base charts the page is the unramified
source pullback. At the selected nodal pieces the actual
circle-line/orbit charts of P3F.9 give the same source Picard
line, with the conifold Hopf charge retained in the later
compression. The page phase and the regular-binding bands in
P3K.6 and P3F.31 are chosen using (AR.7); since its positive
parameter angle is \(\beta\) for both signs, their exact
avoidance and line-carrying band maps apply unchanged. The
finite-base page is consequently a source translation torsor
with the specified line and the marked binding collar.

\item Here is its comparison with the sectioned compact page, relative
to that collar. The integral monodromy coinvariants of the source
lattice \(\Lambda\) are zero, for the following exact reason.
Any nonzero finitely generated coinvariant group has a quotient
\(\mathbb Z/p\) for some prime \(p\). The unimodular
symplectic pairing would give a nonzero common invariant vector
in \(\Lambda/p\). It is a nonzero monodromy-invariant
\(p\)-torsion point of the smooth elliptic family. The finite
algebraic division cover then has a degree-one component: its
single-valued branch is an algebraic section over the regular
base. Properness extends it over the missing points of the
minimal regular source. This contradicts its proved
Mordell--Weil group \(\mathbb Z\).

\item The capped meridian is the inverse product of the retained
singular-fibre meridians, so those retained meridians generate
the same elliptic monodromy group. Each of them lifts as a
closed based loop in one cut sector of (AR.8), since that
cover is unramified there with trivial sheet permutation.
Their actual lattice transports are the original ones.
The page monodromy coinvariants therefore vanish integrally
as well. This argument retains the actual images of
\(T-I\), including their multiplicities; it does not replace
them by primitive vanishing cycles. For the constructible
lattice of periods,
\[
 H^2(C,\partial C;j_*\Lambda)=0.                       \tag{AR.11}
\]
Indeed compact-support duality on the punctured oriented base
identifies the corresponding \(j_!\)-group with the coinvariant
lattice. The quotient \(j_*\Lambda/j_!\Lambda\) is a finite
skyscraper and has no first cohomology. This proves (AR.11)
with the actual boundary marking. The component quotient of the
extending translation sheaf is also a skyscraper and has no
relative first cohomology. After removing it, lift the remaining
translation cocycle to the smooth Lie sheaf. Its period
two-cocycle vanishes by (AR.11), and a partition of unity that
vanishes on the prescribed collar kills the remaining Lie
cocycle. Exponentiation gives the actual smooth translation
isomorphism with the sectioned page, fixed on that collar.

\item The line comparison under these smooth base translations is
made through the rigidified fibrewise Picard map, rather than
an unproved holomorphic divisor for a smooth total-space map.
On a local holomorphic translation chart form
\(\operatorname{Hom}(N,T^*N)\). Its restriction is trivial on
every fibre. At a smooth fibre this is the constant-factor
degree-zero uniformization. At a reduced additive fibre the
component degrees are zero; the proper direct-image generator
has vertical divisor in the whole-fibre kernel, so division by
that base monomial gives a nowhere-zero section. At a node it
is the normalization gluing unit, preserved by translation.
Thus the fibre \(h^0\) is one everywhere, and proper base
change plus evaluation gives a holomorphic line isomorphism
after pulling out a base line. Rigidify it along the zero
section, then pull back by the actual smooth translation
section. This is the desired smooth line isomorphism.

\item On the compactification its remaining base ambiguity is an
integral fibre class \(kF\). The actual line is therefore
compared, for its square, with
\[
                         M_{\rm page}=P_{aG'}-O.
\]
Here \(M_{\rm page}F=F^2=0\), so adding \(kF\) changes
neither its square nor the characteristic pairings. On the
regular binding collar the line extends over the actual
\(T^2\)-bundle over a disc and has fibre degree zero; its
boundary first Chern class is zero. The relative boundary-gauge
ambiguity consequently has zero pairing with the line and
with itself, by the exact relative product calculation P3K.5.
This retains even an odd boundary gauge.

\par

\medskip\noindent\textbf{AR.4. The Thom sign for each physical return}

\item By (AR.5) and (AR.7), the positive physical parameter coordinate
\(\theta\) is the positive weighted coordinate \(\beta\)
for both signs of \(n\). The complex boundary of the affine
page has positive \(\arg\xi\), hence
\[
                         \partial V=+E\times S^1_\theta.
\]
Orient the actual compression by \(\partial W=F-E\).
Attaching its rotated product gives the trace with
\(\partial X=+F\times S^1_\theta\). In the actual
fixed-sphere Hopf chart the ambient orientation is
\[
 {\rm or}(F)\,{\rm or}(\nu_{\rm Hopf})\,dr\,d\theta,
\]
whereas the trace tangent is \({\rm or}(F)\,dr\,d\theta\).
Moving the three normal directions past the two last directions
has positive sign. Thus the induced normal three-plane is the
positive Hopf model and its positive-meridian line U satisfies
\(\pi_*x=+1\) for each norm sign. The possible reversal of
the ambient source-product orientation in the signed \(S^3\)
identification is already included in this normal comparison;
it supplies no additional sign. This is the full signed
calculation of \path{notes/part3_signed_cap_page.md}, independently
executed in its sign certificate.

\item The two outer-bordism/Jin reversals are exactly those in P3K.7.
The exact physical-line comparison of P3K.68 identifies the actual
circle line with U*, and the actual clutch, including its inverse
closing arrow, with lambda\_M=-lambda\_U. Thus the page computation
on U changes sign when assigned to the physical return. The square
of the page line and all source/cap maps above are unchanged.
Consequently the actual positive-return invariant is
\[
 A^{\rm phys}=\frac{p_1(TS_{\rm page})-12M_{\rm page}^2}{24}.
                                                        \tag{AR.12}
\]
Both the circle line and its orientation have now been compared
to the actual page, so this equation applies to the constructed
return rather than to a formal characteristic-number model.

\par

\medskip\noindent\textbf{AR.5. Exact arithmetic and the oriented classes}

\item Let \(e\) be the original capped Kodaira Euler number and
\(e'\) the resolved root Euler number. The compact Jacobian
page in (AR.8) has
\[
 \chi=\frac{m(12-e)+e'}{12},\quad
 \sigma=-8\chi,\quad p_1=-24\chi.                    \tag{AR.13}
\]
All other source fibres are unramified, and the binding is
regular, so the displayed Euler count includes every fibre.
The actual line point meets the identity component at every
retained reducible fibre and at each reduced additive root.
Good roots are irreducible. The height base-change formula
therefore gives, with \(h(G)=1/d\),
\[
 h(aG')=\frac{ma^2}{d}=2\chi+2(P_{aG'}\cdot O),
 \qquad M_{\rm page}^2=-\frac{ma^2}{d}.
\]
Combining this actual line comparison with (AR.12) yields
\[
 A^{\rm phys}=\frac{ma^2}{2d}-\chi=P_{aG'}\cdot O.
                                                        \tag{AR.14}
\]

\par\medskip{\small\noindent
\begin{tabularx}{\linewidth}{@{}XXX@{}}
\toprule
Source/capped site & Root order & Actual A, for both n=1 and n=-1 \\
\midrule
E7-A or E7-B, III*, (a,b)=(2,1) & m=4l+3 & \(3 l + 2\) \\
E7-A or E7-B, III*, (1,2) & m=4l, l\textgreater=1 & \(0\) \\
E6-A or E6-B, IV*, (6,1) & m=3l+2 & \(8 l + 5\) \\
E6-A or E6-B, IV*, (2,3) & m=3l, l\textgreater=1 & \(0\) \\
OS47, III, (14,1) & m=4l+1 & \(25 l + 6\) \\
OS47, III, (7,2) & m=4l, l\textgreater=1 & \(4 l\) \\
OS56, III, (30,1) & m=4l+1 & \(57 l + 14\) \\
OS56, III, (15,2) & m=4l, l\textgreater=1 & \(12 l\) \\
OS56, IV, (30,1) & m=3l+1 & \(43 l + 14\) \\
OS56, IV, (10,3) & m=3l, l\textgreater=1 & \(3 l\) \\
OS45, II, (8,1) & m=6l+1 & \(19 l + 3\) \\
OS45, II, (8,1) & m=6l, l\textgreater=1 & \(19 l\) \\
OS47, II, (14,1) & m=6l+1 & \(37 l + 6\) \\
OS47, II, (14,1) & m=6l, l\textgreater=1 & \(37 l\) \\
OS55, II, (20,1) & m=6l+1 & \(55 l + 9\) \\
OS55, II, (20,1) & m=6l, l\textgreater=1 & \(55 l\) \\
\bottomrule
\end{tabularx}
}\par\medskip

\item Here l\textgreater=0 in every reduced row. The E7 and E6 paired source
equations are those retained in F41. The signed physical return
class is \([\Sigma]=A^{\rm phys}\pmod{28}\) with
\(\mu(\Sigma)=A^{\rm phys}/28\pmod{\mathbb Z}\), by the
actual physical Hopf-gluing transfer P3K.68, with the explicit inverse-closing convention. A nonzero residue gives the
nonproduct mapping torus by P3F.24, with both spin structures
retained. The table is for the actual single positive \(r\)-loop;
no unrecorded multiple of that loop is used. Its origin, root,
Euler, height, square and residue identities are executed in
\path{code/part3/additive_returns.py/.json/.out}.

\end{enumerate}
\end{proof}




\begin{theorem}\label{night:P3K.28}
\StatementRecord{night:P3K.28}{proved}{S6\_06\_BLIND: GPT-6 Astra (OpenAI, ChatGPT Pro mode, effort Pro), recorded re-derivation in the checkpoint dated 18 November 2026}{\cref{night:P3F.24}, \cref{night:P3F.31}, \cref{night:P3F.9}, \cref{night:P3K.26}, \cref{night:P3K.5}, \cref{night:P3K.6}, \cref{night:P3K.68}, \cref{night:P3K.7}}

For the actual two-cap sphere families of P3F.38, P3F.46 and P3F.47,
retain their signed norms and write
\[
 D=m_0n_\infty+m_\infty n_0=\varepsilon\in\{1,-1\},
 \qquad M=m_\infty m_0,\qquad \lambda=r^M .
\]
One positive loop in the actual parameter \(r\) has the signed return
\[
 A^{\rm phys}
 =\frac{Ma^2}{2d}
       -\frac{M(12-e_\infty-e_0)+m_0e_\infty'+m_\infty e_0'}{12},
\]
where \(h(G)=1/d\), the marked line point is \(aG\), and
\(e_i,e_i'\) are the original and actual resolved-root Euler numbers.
The exact values in all nine nonempty sectors appear below and hold
for both determinant signs. The single quadratic \(I_1^*\) cap of
P3F.36, on either source profile and with \(\lambda=r^2\),
has \(A^{\rm phys}=2\) for both \(n=1\) and \(n=-1\).
By the actual dual-line and inverse-closing map P3K.68, the physical Hopf-gluing class is \(A^{\rm phys}\bmod28\), with
\(\mu=A^{\rm phys}/28\bmod\mathbb Z\).

\end{theorem}
\begin{proof}
\begin{enumerate}[label=\arabic*.]
\item \par

\medskip\noindent\textbf{K28.1. The marked norm data and the actual parameter}

\par\nobreak

\item Place the two capped source fibres at infinity and zero. Write their root
orders as \(m_\infty,m_0\), their signed norms as \(n_\infty,n_0\), and put
\[
 D=m_0n_\infty+m_\infty n_0.
 \tag{K28.1}
\]
The complete integral topology proofs F38, F46 and F47 show that their sphere
sector is exactly \(D=\varepsilon\), \(\varepsilon\in\{1,-1\}\). In particular
\(\gcd(m_\infty,m_0)=1\), and each local pair \((m_i,n_i)\) is primitive.
Throughout this proof put \(M=m_\infty m_0\).

\item The actual source line is denoted \(\mathcal M\), to distinguish it from the
integer \(M\). Its contracting lift \(B_\lambda\) covers translation by
\(-Q=-bG\). At each resolved root there are actual commuting maps
\[
 A_{i,0}^{m_i}=B_0^{n_i},\qquad [A_{i,0},B_0]=1 .
\]
The family in the topology proofs uses
\[
 \lambda=r^M,\qquad
 A_{i,r}=r^{M n_i/m_i}A_{i,0},\qquad
 B_\lambda=r^M B_0 .
 \tag{K28.2}
\]
This is the parameter used here. The positive return means one positive
loop in \(r\).

\item Let \(\gamma\) be the outward suspension coordinate, so a positive unit
of \(\gamma\) is the inverse contracting map \(B_\lambda^{-1}\). Write
\(\alpha=\arg(t)/(2\pi)\) on the ordinary annulus. The actual cap norm
coordinates, of periods \(m_i\), are
\[
 u_0=m_0\gamma-n_0\alpha,\qquad
 u_\infty=m_\infty\gamma+n_\infty\alpha.
 \tag{K28.3}
\]
Indeed, on the root \(q_i=z^{m_i}\), the defining cut is
\((z,u_i)\mapsto(\zeta_{m_i}z,u_i-n_i)\). Thus a positive source
meridian maps to \(-n_i h_i\), and \(\gamma\) maps to \(m_i h_i\),
where \(h_i\) is the positive primitive cap circle. This is the actual
marked meridian \(m_i\alpha_i+n_i\gamma\) of F38, F46 and F47.
The source translations and the integral line characters remain in the
other coordinates of these maps.

\item The norm-shell coordinate can be used with the prescribed boundary metric.
To justify this without supposing that a fibrewise lift has a constant
pointwise norm, use the actual strictly contracting primitive generator
\(H_i\), with \(H_i^{m_i}=B_\lambda\) and \(H_i^{n_i}=A_{i,r}\).
An equivariant radial coordinate is chosen on one norm shell and transported
by \(H_i\). On the overlap with the flat boundary norm, the increasing lift
of the circle-coordinate comparison has \(F(s+1)=F(s)+1\). The radial map
\(v\mapsto c^{F(s)-s}v\), in the notation of P3F.31, commutes with the
actual contraction and with the physical delta circle. It fixes the source
and base coordinates and carries the actual circle line. This gives the
same marked norm coordinates (K28.3) through the entire cap.

\par

\medskip\noindent\textbf{K28.2. An exact map of the orbit base to the weighted three-sphere}

\par\nobreak

\item On the ordinary annulus set
\[
 \phi_1=\varepsilon(m_0\gamma-n_0\alpha),\qquad
 \phi_2=\varepsilon(m_\infty\gamma+n_\infty\alpha).
 \tag{K28.4}
\]
The phase matrix, with columns \((\gamma,\alpha)\), has determinant
\(D=\varepsilon\). Its entries are integers, so it is an actual
automorphism of the marked two-torus. Moreover
\[
 m_0\phi_2-m_\infty\phi_1=\alpha.
 \tag{K28.5}
\]
Choose the unique positive radii satisfying
\[
 r_1^2+r_2^2=1,\qquad
 \frac{r_2^{m_0}}{r_1^{m_\infty}}=|t|.
\]
The ratio is strictly increasing as \(r_2\) increases from zero to one.
The map \(z_j=r_j e^{2\pi i\phi_j}\) therefore identifies the actual
ordinary orbit-base annulus with the corresponding annulus in
\(S^3\subset\mathbb C^2\), with
\[
                         t=\frac{z_2^{m_0}}{z_1^{m_\infty}}.
 \tag{K28.6}
\]

\item Here are its extensions across the two actual cap solid tori. At infinity,
\(q_\infty=1/t=z^{m_\infty}\), let \(R>0\) be defined by
\[
 R^2+|z|^2R^{2m_0/m_\infty}=1.
\]
The extending map is
\[
 (z,u_\infty)\longmapsto
 \left(
 zR^{m_0/m_\infty}e^{2\pi i\varepsilon m_0u_\infty/m_\infty},
 Re^{2\pi i\varepsilon u_\infty}
 \right).
 \tag{K28.7}
\]
Under the actual \(A_\infty\) cut the first phase changes by
\[
 \frac{1-\varepsilon m_0n_\infty}{m_\infty}
       =\varepsilon n_0,
\]
and the second by \(-\varepsilon n_\infty\). Both are integers.
Thus (K28.7) descends through the actual marked quotient. Its inverse first
recovers \(R=|z_2|\), then a lift of \(\arg(z_2)/(2\pi)\), and finally
\(z\). The possible lifts are identified by precisely the powers of
the displayed \(A_\infty\) cut. The derivative of the radial equation
with respect to \(R\) is positive, so the map and inverse are smooth.

\item At zero, \(q_0=t=z^{m_0}\), use
\[
 R^2+|z|^2R^{2m_\infty/m_0}=1
\]
and the map
\[
 (z,u_0)\longmapsto
 \left(
 Re^{2\pi i\varepsilon u_0},
 zR^{m_\infty/m_0}e^{2\pi i\varepsilon m_\infty u_0/m_0}
 \right).
 \tag{K28.8}
\]
The two phase changes under \(A_0\) are
\(-\varepsilon n_0\) and
\((1-\varepsilon m_\infty n_0)/m_0=\varepsilon n_\infty\).
The same inverse argument applies. Substitution of (K28.3) shows that
(K28.7) and (K28.8) agree exactly with (K28.4) on both ordinary overlaps.
These formulas, rather than the abstract Seifert invariants alone,
give the required global marked orbit-base identification.

\item The positive weighted Hopf angle, of weights
\((m_0,m_\infty)\), is
\[
                              \beta=\varepsilon\gamma.
 \tag{K28.9}
\]

\par

\medskip\noindent\textbf{K28.3. The physical parameter cochain and the complex page}

\par\nobreak

\item The annular arrow is ordinary-to-cap, as fixed in the construction proofs.
If \(c=\log\lambda/(2\pi i)\) and \(s_i=\log q_i/(2\pi i)\), its scalar
variation is \(n_i(c-c_0)s_i\delta/m_i\). For a positive \(r\)-loop put
\[
 k_\infty=m_0n_\infty,\qquad k_0=m_\infty n_0.
\]
The actual inverse closing multiplier is
\[
                      f(t)=t^{k_0}(t-t_c)^{-(k_\infty+k_0)}.
 \tag{K28.10}
\]
Its orders at zero and infinity are \(k_0,k_\infty\), respectively.
In the sphere sector \(k_\infty+k_0=\varepsilon\).
At the selected cusp \(q=t-t_c\), the positive deck is therefore the
inverse of (K28.10), and its logarithmic part is
\(+\varepsilon\log q\,\delta/(2\pi i)\).
The remaining factor is a specified nonvanishing unit. Its logarithm is
added to both compared origins, so it cancels in their difference.

\item For clarity the full cusp origin computation is retained. Put
\[
 s=\frac{\log q}{2\pi i},\qquad
 v=\frac{\log r}{2\pi i},\qquad
 \gamma=(g,-C-s-Mv).
\]
The physical line coordinate is \(V=e^{2\pi i\zeta_2}\), the source
Abel coordinate is \(\zeta_1\), and
\[
 B_\lambda(a,V)=(a-g,e^{2\pi iC}q r^M V).
\]
Thus the displayed \(\gamma\) is exactly \(B_\lambda^{-1}\).
Use the two origins
\[
 \begin{split}
 o_\varepsilon&=\varepsilon
       \left(-\frac{gs}{M},\frac{Cs}{M}+\frac{s^2}{2M}+sv\right),\\
 \sigma_\varepsilon&=\varepsilon
       \left(gv,-Cv-\frac{M(v^2-v)}2\right).
 \end{split}
 \tag{K28.11}
\]
They have the actual parameter deck just found, and their difference
\(F_\varepsilon=\sigma_\varepsilon-o_\varepsilon\) satisfies
\[
 F_\varepsilon(v+\theta)-F_\varepsilon(v)
   =\varepsilon\theta\gamma
      -\frac{\varepsilon M\theta(\theta-1)}2\,\delta .
 \tag{K28.12}
\]
The quadratic delta gauge is zero at both endpoints. On the auxiliary
cusp root \(q=Q^{2M}\), the whole angular difference changes by
\(-2\varepsilon M\theta\delta\) under \(s\mapsto s+2M\); this is the
retained degree-\(2M\) boundary gauge. That auxiliary source root does
not change the parameter. One positive \(r\)-loop moves by
\(\varepsilon\gamma=\beta\), with the same positive sign for either
sphere determinant.

\item Choose a regular value \(\tau\), distinct from the capped sites and
the selected cusp. The actual complex page is
\[
 \mathcal C^\circ:\quad \eta^{m_0}-\tau\xi^{m_\infty}=1.
 \tag{K28.13}
\]
Its weighted radial embedding into the above three-sphere is
\[
 (\xi,\eta)\longmapsto
       (\rho^{m_0}\xi,\rho^{m_\infty}\eta),
 \qquad
 \rho^{2m_0}|\xi|^2+\rho^{2m_\infty}|\eta|^2=1 .
 \tag{K28.14}
\]
It is the positive-argument page of the weighted homogeneous polynomial
\(z_2^{m_0}-\tau z_1^{m_\infty}\), and its map to the source base is
\[
                      t=\tau+\xi^{-m_\infty}.
 \tag{K28.15}
\]
The affine curve is smooth: its two partial derivatives cannot both
vanish on (K28.13). Its map has degree \(M\). Over infinity there are
\(m_0\) points, each of ramification order \(m_\infty\); over zero there
are \(m_\infty\) points, each of order \(m_0\); over the regular binding
\(\tau\) there is one point of order \(M\). The last assertion follows
also from \(\gcd(m_\infty,m_0)=1\): the expansion
\(\xi\sim c w^{-m_0}\), \(\eta\sim d w^{-m_\infty}\)
has one place at infinity. Riemann--Hurwitz gives
\[
                  g(\overline{\mathcal C})
                    =\frac{(m_\infty-1)(m_0-1)}2.
 \tag{K28.16}
\]
The selected source \(I_1\) is unramified and has \(M\) separate
nodal inverse images.

\item The page is also the exact graph in each actual root cap. At infinity
its phase-\(\phi\) graph is
\[
 \varepsilon m_0 u_\infty+
     \frac{\arg(1-\tau z^{m_\infty})}{2\pi}
           =\frac{\phi}{2\pi}\pmod{\mathbb Z},
 \tag{K28.17}
\]
and at zero it is
\[
 \varepsilon m_\infty u_0+
     \frac{\arg(z^{m_0}-\tau)}{2\pi}
           =\frac{\phi}{2\pi}\pmod{\mathbb Z}.
 \tag{K28.18}
\]
There are \(M\) sheets before quotienting. The actual \(A_\infty\)
changes the sheet index by \(-\varepsilon m_0n_\infty\);
there are exactly \(m_0\) orbits of length \(m_\infty\).
The \(A_0\) sheet step is \(-\varepsilon m_\infty n_0\),
giving \(m_\infty\) orbits of length \(m_0\). For phase zero the
holomorphic coordinates are, respectively,
\[
 \begin{array}{ll}
 \displaystyle \xi^{m_\infty}
       =\frac{z^{m_\infty}}{1-\tau z^{m_\infty}},
 &\displaystyle \eta^{m_0}
       =\frac1{1-\tau z^{m_\infty}},\\[6pt]
 \displaystyle \xi^{m_\infty}
       =\frac1{z^{m_0}-\tau},
 &\displaystyle \eta^{m_0}
       =\frac{z^{m_0}}{z^{m_0}-\tau}.
 \end{array}
 \tag{K28.19}
\]
Thus these local complex maps do not depend on the signs of the norms.

\par

\medskip\noindent\textbf{K28.4. The exact line on the page, including the quadratic starred cap}

\par\nobreak

\item On each orbit of the sheet permutation choose one representative graph
\(u_i=u_{i,j}(z)\). The map from the actual resolved root source is
\[
                        p\longmapsto [p,u_{i,j}(q_i(p))].
 \tag{K28.20}
\]
It extends across every resolved component, since the phase in
(K28.17) or (K28.18) is smooth there. A relation preserving this
representative norm level has \(n_i k+m_i h=0\). Primitivity gives
\(k=m_i j\), \(h=-n_i j\), and the actual identity
\(A_i^{m_i j}B_\lambda^{-n_i j}=1\) proves uniqueness. Thus no residual
stabilizer or further root quotient remains in (K28.20).
For a nonzero vector of the actual root line
\[
                         N_i=\mathcal O(aG_i'-O_i),
\]
normalize its norm to \(u_{i,j}(z)\). This is an explicit
delta-equivariant map to the principal circle over (K28.20).
Consequently its associated complex line is \(N_i\), including its
specified lift and all finite characters.

\item For F46 and F47 this uses the precise line maps of those propositions:
the reduced-root contractions with degree-zero exceptional restriction;
the good \(IV^*\) and \(III^*\) toric Cartier corrections; the good
III exceptional difference; and the cancelling good-II exceptional
coefficients. Each correction was proved invariant under the actual
translation \(Q\), so the comparison carries \(B_0\) and \(A_{i,0}\).
The sheet map does not replace these corrections by numerical degrees.

\item Here is the additional line map required for F36 and F38. On the
quadratic \(I_1^*\) root let \(S'\) be the minimal \(I_2\) model,
let \(\beta:\widehat S\to S'\) blow up the four inversion fixed
points, and let \(\pi:\widehat S\to S\) be the smooth resolved
quotient. The actual source line is
\(\mathcal M=\mathcal O(aG-O-L_0)\), with \(a=2\) or \(6\).
If \(E_0,E_1\) are the two exceptional curves at the section leaves,
the literal Cartier equalities are
\[
 \begin{split}
 \pi^*(aG-O-L_0)&=\widehat P-\widehat O-2E_0,\\
 \beta^*(P'-O')&=\widehat P-\widehat O+E_1-E_0,\\
 \beta^*N_*&=\pi^*\mathcal M+E_0+E_1 .
 \end{split}
 \tag{K28.21}
\]
The actual \(Q=bG\), \(b=2\) or \(6\), exchanges \(E_0,E_1\)
and exchanges the other two exceptional curves. Hence (K28.21)
is equivariant for the specified suspension and its affine norm lift.
The Tate character of \(N_*\) is
\[
             Z_P=v(z)^a/u(z^2)^{a/2},\qquad Z_P(0)=-1 .
 \tag{K28.22}
\]
Its transitions are units on both \(I_2\) components, and its degree is
zero on each. The full F36 norm
\(A_*^2=B_\lambda^{n_*}\) includes the retained line eigencharacter.
Applying (K28.20) with \(m_*=2\) therefore gives the actual line
\(N_*=\mathcal O(aG'-O')\), not the uncorrected pullback of
\(\mathcal M\). This proves the needed line identification even at
the two section points and all four fixed-point blowups.

\item On the ordinary part of the page the projection is the unramified
source pullback. At the selected nodes use the actual source-line/orbit
charts of P3F.9. The page comparison concerns that source line; the
Hopf charge around the fixed sphere is retained in the compression
and is not extended through the fixed sphere as a trivial line.
The regular-binding band construction of P3K.6 and P3F.31 applies
in the coordinate \(\beta\) of (K28.9); it keeps the base coordinate
and carries the line by the radial maps in K28.1.

\item We now trivialize the remaining finite-base source torsor relative
to this actual binding collar. The integral monodromy coinvariants
vanish for a direct reason in each two-cap case. For F46 and F47 let
\(v\) be the retained nodal vanishing cycle and let \(A\) be the
order-four source monodromy at III or \(III^*\). The trace calculation
in those topology proofs gives
\[
                          |\langle v,Av\rangle|=1.
 \tag{K28.23}
\]
For F38 the starred monodromy is \(-T_u\), and its product with
the retained nodal transvection \(T_v\) and the IV monodromy gives
\[
       -1=-2+\langle u,v\rangle^2,\qquad
       |\langle v,(-T_u)v\rangle|=1 .
 \tag{K28.24}
\]
The nodal meridian has trivial sheet permutation in the actual page
cover. Its conjugate by the starred or order-four meridian also
has trivial sheet permutation, even though that conjugating path
itself need not close on the page. Thus both nodal loops lift as
closed page loops. Their monodromies are \(T_v,T_{Av}\), with
images \(\mathbb Z v,\mathbb Z Av\) for \(T-1\). Equations
(K28.23)--(K28.24) make these images the entire integral elliptic
lattice. This proves zero page coinvariants without assuming that
an arbitrary branched cover preserves the original Mordell--Weil
torsion group.

\item Let \(C\) be the page base with its marked binding boundary,
\(j:C_{\rm reg}\to C\), and \(\Lambda\) its elliptic lattice.
The quotient \(j_*\Lambda/j_!\Lambda\) is supported at finitely
many singular points and has no first cohomology. Compact-support
Poincare duality and the preceding coinvariants give
\[
                  H^2(C,\partial C;j_*\Lambda)=0 .
 \tag{K28.25}
\]
The smooth identity-component Lie sheaf is fine. Its lattice
exponential sequence, together with the finite component sheaf
at the singular fibres, consequently gives a translation
trivialization of the actual torsor relative to the specified
binding collar. Every translation used is a section of the sheaf
of translations extending on the actual reduced resolved fibres.
The \(I_2\) root is included: its identity group is multiplicative,
its lattice at the node is the invariant \(j_*\Lambda\), and its
component quotient is the finite skyscraper just retained.

\item The line comparison through this smooth translation is the
rigidified fibrewise Picard comparison of P3K.26, rather than
a claimed holomorphic divisor for a merely smooth map. Over the
universal translating-section parameter, the actual Hom line
between \(N\) and its translate has trivial restriction to every
fibre. The constant-factor line uniformization proves this on the
smooth fibres. The reduced component degrees and the retained
unit transitions prove it on III, IV and \(I_2\) fibres, and
proper base change and evaluation give the relative isomorphism.
Rigidify at \(O\), then pull it back by the constructed smooth
translation section. The remaining ambiguity on the compact
base is a line pulled back from that base, with Chern class
\(kF\). Therefore
\[
                    (c_1(N)+kF)^2=c_1(N)^2,
 \qquad c_1(N)F=F^2=0 .
 \tag{K28.26}
\]
The binding restriction has zero first Chern class and extends over
the regular \(T^2\)-cap. The relative boundary gauges, including
the retained odd ones, have zero square ambiguity by P3K.5.
This completes the actual page-line comparison.

\par

\medskip\noindent\textbf{K28.5. The signed trace and the characteristic calculation}

\par\nobreak

\item Near the single point at infinity of (K28.13), write
\(\xi\sim c w^{-m_0}\), \(\eta\sim d w^{-m_\infty}\).
The complex boundary of the punctured page has parameter
\(-\arg w=\beta\). By (K28.12), the positive physical parameter
is exactly \(\theta=\beta\). Hence
\(\partial V=+E\times S^1_\theta\).
Orient the actual nodal compression by \(\partial W=F-E\).
Its rotated product attaches to give a trace with
\(\partial X=+F\times S^1_\theta\).

\item In the actual Hopf chart the ambient orientation is
\(\operatorname{or}(F)\operatorname{or}(\nu_{\rm Hopf})\,dr\,d\theta\),
whereas the trace orientation is
\(\operatorname{or}(F)\,dr\,d\theta\).
Moving the three normal directions past the last two is even.
The induced normal sphere is therefore the positive Hopf sphere
and its positive-meridian line U satisfies \(\pi_*x=+1\).
This is the full signed Thom comparison, for both
\(\varepsilon=1\) and \(\varepsilon=-1\). Any source-product
orientation reversal in the three-sphere map is already included
in this normal comparison and introduces no additional sign.

\item The two outer-boundary/Jin reversals are the ones proved in
P3K.7. P3K.68 retains the actual positive-meridian-to-physical
line map M=U*, together with the inverse closing arrow in the Hopf
gluing. Its physical clutch is lambda\_M=-lambda\_U. Consequently
the physical return is the negative of the U-line page index.
The actual trace bordism, with the page-line map just constructed,
therefore gives
\[
            A^{\rm phys}
              =\frac{p_1(TS_{\rm page})-12\,c_1(N)^2}{24}.
 \tag{K28.27}
\]
The compression cap and the regular elliptic cap have the exact
boundary line comparison of P3K.5--7, zero signature, and zero
line-square ambiguity. Thus the right side may be computed on
the compact sectioned elliptic page. This uses the displayed
relative maps, not a holomorphic compactification of the original
parameter family.

\item Let the original capped Euler numbers be \(e_\infty,e_0\),
and the actual minimal root Euler numbers be
\(e_\infty',e_0'\). There are \(m_0\) copies of the first root
and \(m_\infty\) of the second. Every other fibre is unramified.
Consequently
\[
 \chi(\mathcal O_{S_{\rm page}})
   =\frac{M(12-e_\infty-e_0)
            +m_0e_\infty'+m_\infty e_0'}{12}.
 \tag{K28.28}
\]
Its signature is \(-8\chi\) and its tangent Pontryagin number is
\(-24\chi\). The actual marked line section \(aG'\) meets the
identity component at every remaining reducible fibre and at
each reduced root in the listed cases; good roots are irreducible.
The height base-change formula, with \(h(G)=1/d\), gives
\[
 h(aG')=\frac{Ma^2}{d}
       =2\chi+2(P_{aG'}\cdot O),\qquad
 c_1(N)^2=-\frac{Ma^2}{d}.
 \tag{K28.29}
\]
For the quadratic starred line this formula uses \(N\) from
(K28.21), so no root correction from \(-L_0\) is omitted.
Combining the actual comparison with (K28.27) yields
\[
                    A^{\rm phys}
                         =\frac{Ma^2}{2d}-\chi
                         =P_{aG'}\cdot O .
 \tag{K28.30}
\]

\par

\medskip\noindent\textbf{K28.6. All joint return values and the one-star specialization}

\par\nobreak

\item The following tables are evaluated only on the sphere tuples
\(D=\pm1\), with the exact signed parameter (K28.2). For each
fixed row and fixed root orders, both determinant signs give
the displayed integer. The actual proof of this equality of
signs is (K28.4)--(K28.27).

\item For the E6 source of F46, \(d=6\), the other source fibre is
the selected \(I_1\), and the values are:

\par\medskip{\small\noindent
\begin{tabularx}{\linewidth}{@{}>{\hsize=.65\hsize\linewidth=\hsize}X>{\hsize=1.05\hsize\linewidth=\hsize}X>{\hsize=1.05\hsize\linewidth=\hsize}X>{\hsize=1.25\hsize\linewidth=\hsize}X@{}}
\toprule
\((a,b)\) & \(m_\infty\) at \(IV^*\) & \(m_0\) at III & Actual \(A^{\rm phys}\) \\
\midrule
\((1,6)\) & \(3l\), good & \(4k\), good & \(0\) \\
\((2,3)\) & \(3l\), good & \(4k+1\), reduced III & \(3 k l\) \\
\((3,2)\) & \(3l+2\), reduced IV & \(4k\), good & \(8 k l + 4 k\) \\
\((6,1)\) & \(3l+2\), reduced IV & \(4k+1\), reduced III & \(35 k l + 22 k + 8 l + 5\) \\
\bottomrule
\end{tabularx}
}\par\medskip

\item The good-root indices are positive and the reduced-root indices
are nonnegative. The unit determinant also retains every
coprimality restriction; for example the third row requires
\(3l+2\) odd. These restrictions are not removed from the table.

\item For the E7 source of F47, \(d=2\), the other source fibre is again
the selected \(I_1\):

\par\medskip{\small\noindent
\begin{tabularx}{\linewidth}{@{}>{\hsize=.65\hsize\linewidth=\hsize}X>{\hsize=1.05\hsize\linewidth=\hsize}X>{\hsize=1.05\hsize\linewidth=\hsize}X>{\hsize=1.25\hsize\linewidth=\hsize}X@{}}
\toprule
\((a,b)\) & \(m_\infty\) at \(III^*\) & \(m_0\) at II & Actual \(A^{\rm phys}\) \\
\midrule
\((2,1)\) & \(4l+3\), reduced III & \(6k+1\), reduced II & \(22 k l + 15 k + 3 l + 2\) \\
\((2,1)\) & \(4l+3\), reduced III & \(6k\), good & \(22 k l + 15 k\) \\
\((1,2)\) & \(4l\), good & \(6k+1\), reduced II & \(4 k l\) \\
\bottomrule
\end{tabularx}
}\par\medskip

\item The fourth formal root combination, good \(III^*\) and good II,
has both orders even. Its determinant is even and it contains
no sphere tuple. No sphere return class is assigned to it.

\item For F38, the source is \(I_1^*+IV+I_1\), with \(d=12\).
At infinity \(m_\infty=2\), the actual root is \(I_2\),
and \(e_\infty'=2\). Its page consists of \(m_0\) copies of this
root and two copies of the IV root. The two sectors are:

\par\medskip{\small\noindent
\begin{tabularx}{\linewidth}{@{}>{\hsize=.65\hsize\linewidth=\hsize}X>{\hsize=1.05\hsize\linewidth=\hsize}X>{\hsize=1.05\hsize\linewidth=\hsize}X>{\hsize=1.25\hsize\linewidth=\hsize}X@{}}
\toprule
\((a,b)\) & \(m_\infty\) at \(I_1^*\) & \(m_0\) at IV & Actual \(A^{\rm phys}\) \\
\midrule
\((6,2)\) & \(2\), root \(I_2\) & \(6l+1\), reduced IV & \(16 l + 2\) \\
\((2,6)\) & \(2\), root \(I_2\) & \(3k\), good, \(k\) odd & \(0\) \\
\bottomrule
\end{tabularx}
}\par\medskip

\item Here the sphere equation is \(m_0n_\infty+2n_0=\pm1\).
It forces \(m_0\) odd; together with the reduced root domain
\(m_0\equiv1\pmod3\), this gives the stated \(6l+1\).
Equation (K28.28) is \(\chi=(2m_0+e_0')/6\).
In the reduced row (K28.29) gives
\(c_1(N)^2=-6m_0\), and hence
\(A=(8m_0-2)/3=16l+2\). In the good row it gives
\(h(aG')=2m_0/3\) and \(\chi=m_0/3\), hence zero.

\item Finally consider the single quadratic \(I_1^*\) cap of F36
on either \(I_1^*+I_3+2I_1\) or \(I_1^*+IV+I_1\).
Here \((a,b)=(6,2)\), \(\lambda=r^2\), and \(n=\pm1\).
The exact one-cap map of K26, with \(m=2\), gives the complex
page \(\eta-\tau\xi^2=1\) and the positive Hopf angle \(n\gamma\).
The root sheet and its actual line are precisely (K28.20)--(K28.22).
The integral page coinvariants vanish by the one-cap argument
of K26: the retained meridians recover the whole source
monodromy, and a nonzero prime quotient of its coinvariants
would give a nonzero rational torsion section, contradicting
the proved source Mordell--Weil group \(\mathbb Z\).
Thus the same relative line and signed Thom maps apply.
The Euler count is
\[
              \chi=\frac{2(12-7)+2}{12}=1,\qquad
              h(6G')=\frac{2\cdot36}{12}=6 .
\]
The remaining \(I_3\) or IV component correction is zero for
\(6G'\). It follows that the actual return is
\[
                       A^{\rm phys}=2
 \tag{K28.31}
\]
for either norm sign and for both source profiles.

\item By the exact physical transfer and inverse-closing map of P3K.68, the physical Hopf-gluing class of
every row is \(A^{\rm phys}\bmod28\), with
\(\mu=A^{\rm phys}/28\bmod\mathbb Z\). A nonzero residue has the
nonproduct mapping torus proved in P3F.24. The present proof
uses the actual norm, page, line and trace maps before applying
that transfer.

\par

\medskip\noindent\textbf{Certificate and scope}

\par\nobreak

\item The executed certificate \path{code/part3/joint_additive_returns.py} records
the exact phase matrix and determinant, both cap-cut identities,
the origin difference, the sheet-orbit gcd calculation, the
Riemann--Hurwitz genus, and every arithmetic polynomial in the
tables. Its JSON links the actual construction and integral topology
proofs. The proof of the global maps and line identification is
the text above; finite numerical samples do not replace it.

\par

\medskip\noindent\textbf{Independent audit record}

\par\nobreak

\item The parent passed the complete K28.1--6 proof, all actual maps and every arithmetic row. Fourfold independently passed K28.1--3 and K28.5, including both solid-torus quotient inverses, the actual positive parameter, the degree-2M boundary gauge, the sheet orbits, and the signed Hopf Thom comparison; no extra sign or factor M occurs. Checks independently passed K28.4 and the complete table, including the literal I1-star Cartier correction, its retained Tate character, the I2 local period sheaf, the relative rigidified line comparison, and all single/joint I1-star values. These audits were received after the standalone export; the exported proof substance is unchanged.

\end{enumerate}
\end{proof}


\section{The ordinary product and its complete transport maps}



\begin{theorem}\label{night:P3K.38}
\StatementRecord{night:P3K.38}{proved}{S6\_06\_BLIND: GPT-6 Astra (OpenAI, ChatGPT Pro mode, effort Pro), recorded re-derivation in the checkpoint dated 18 November 2026}{\cref{night:P3F.26}, \cref{night:P3F.35}, \cref{night:P3F.36}, \cref{night:P3F.41}, \cref{night:P3F.46}, \cref{night:P3F.51}, \cref{night:P3F.52}, \cref{night:P3F.53}, \cref{night:P3K.32}, \cref{night:P3K.34}, \cref{night:P3K.35}, \cref{night:P3K.36}, \cref{night:P3K.37}}
Fix any of the thirty nonempty branches of P3F.53, including its actual
source coefficients, vertical line divisor, logarithmic norms, affine
cap maps, finite line characters, order-one clutches, and allowed
index-one fan. We write its line and source translation as
\[
 \mathcal M=\mathcal O(aG-O+D),\qquad Q=-bG,
 \qquad ab=d,
\]
with the notation and actual divisor comparisons of P3F.51--53. Let
\(m_i>1\) be its effective cap orders and put
\(M=\prod_i m_i\); an absent or order-one slot contributes one.
All norms \(n_i\) and the ordinary framing term are retained in
P3K.32's signed Seifert equation. In particular,
\(\gcd(m_i,n_i)=1\), and the nontrivial \(m_i\) are pairwise
coprime.

The actual scalar is the multiplier in
\[
                         B_\lambda=\lambda B_0.       \tag{IM.1}
\]
We continue the fixed construction on the logarithm cover of a small
punctured \(\lambda\)-disc. Its data at two logarithm values differing
by \(2\pi iN\) have a biholomorphic endpoint identification exactly
when
\[
                              N\in (M/b)\mathbb Z.     \tag{IM.2}
\]
This assertion allows every biholomorphic map between the resulting
threefolds. It imposes no preservation of a chosen elliptic marking,
line-circle action, source section, or toric chart on that map.
The positive integer \(M/b\) is therefore the least positive scalar
winding. Sufficiency is the actual family of P3K.35,
\[
                     \lambda=w^{M/b},\qquad w=r^b,
                     \lambda=r^M.                     \tag{IM.3}
\]
We prove necessity by intrinsic invariants of its multiple fibres.

For the component-supported OS49 family of P3K.34 and P3K.37, with
\(m=3k\), logarithmic vector
\(n(e_\gamma-e_u)+d_\delta\delta\), and \(n=\pm1\), the
corresponding least positive scalar winding is \(k\). Its actual
closed parameter is \(\lambda=w^k\), \(w=r^3\).
The argument for that statement uses its actual central Picard point
\(G'_0\); it uses no whole-fibre height formula.
\end{theorem}
\begin{proof}
\begin{enumerate}[label=\arabic*.]
\item \par

\medskip\noindent\textbf{2. Every endpoint map preserves each nontrivial multiple fibre}

\par\nobreak

\item We first make the fibration intrinsic. Every threefold in question is
diffeomorphic to \(S^6\), by the integral topology of P3F.53 and
P3K.32, with the actual fan changes of P3F.35. Thus
\(H^2(Y;\mathbb R)=0\). A general fibre \(F\) of the displayed
map \(f:Y\to\mathbb P^1\) is a compact complex torus, hence is
Kähler. A horizontal effective divisor \(Z\) on \(Y\) would
restrict to a nonzero effective divisor \(Z|_F\) with zero real
first Chern class. For a Kähler form \(\omega\) on \(F\),
\[
 0=\int_F c_1(\mathcal O(Z)|_F)\wedge\omega
   =\int_{Z|_F}\omega>0,
\]
a contradiction. Every effective divisor is consequently vertical.

\item Let \(g\) be a meromorphic function on \(Y\). Its polar divisor
is vertical, so on a general compact torus fibre its restriction
has no polar divisor. Hartogs extension across its possible
codimension-two indeterminacy set makes that restriction holomorphic,
and hence constant. To see the descent at each base value, choose a
local coordinate \(q\). A sufficiently large power of \(q\circ f\)
cancels all divisorial poles on a small proper inverse image. The
resulting function is holomorphic by Hartogs. The fibres are connected;
the Stein factorization of this proper map to a normal disc therefore
gives \(f_*\mathcal O_Y=\mathcal O\). The holomorphic function
descends, and division by the same power of \(q\) gives a meromorphic
function on the disc. We have proved the literal field equality
\[
                    \mathbb C(Y)=f^*\mathbb C(t).       \tag{IM.4}
\]

\item A biholomorphism between endpoint threefolds induces an automorphism
of this field and hence a Möbius transformation of the base. The
equality of meromorphic maps extends over the entire threefold.
Divisorial multiplicities of its fibres are preserved. In the actual
cap chart the base function is \(q=z^{m_i}\), the root central
fibre is reduced, and the quotient is free with transitive component
action; its fibre is exactly \(m_i S_i\). The ordinary models
have divisorial gcd one. At the toric filling the same fact is the
primitivity of its physical valuation vector: the component orders
have gcd one, and changing the index-one fan retains that vector.
Thus the displayed effective caps are precisely the fibres of
multiplicity greater than one.

\item There are at most two such caps. Their coprime orders are distinct.
Every endpoint biholomorphism consequently carries each \(S_i\)
to the continuation of the same \(S_i\). This conclusion allows
permutations of ordinary fibres, including ones with the same Kodaira
type, and arbitrary source automorphisms. A biholomorphism of the
reduced supports lifts uniquely to their normalizations, by the
universal property of normalization. We may therefore use the
Albanese curve of a smooth cap support, or of its normalization,
without requiring any additional equivariance of the endpoint map.

\par

\medskip\noindent\textbf{3. The exact period lattice at a good-reduction cap}

\par\nobreak

\item Fix one good cap, and write its order and norm as \(m,n\).
The actual root line has central restriction
\[
                          N_0=\mathcal O_E(aG'_0-O).
                                                               \tag{IM.5}
\]
This is the line supplied by the root/quotient divisor diagrams in
P3F.41--42 and P3F.46--47. For every compatible vertical divisor
and fixed clutch, P3F.52's direct-image evaluation compares it with
(IM.5) by an actual base line and a fixed deck character. A base
line has trivial restriction to \(E\), and a character changes
the lift, not the underlying central line. Let
\[
                         t=\operatorname{ord}(aG'_0).
                                                               \tag{IM.6}
\]
This order includes the value one for the zero Picard point.

\item The literal minimal-root substitutions of P3F.51 give the following
points and automorphisms. At a good III or III* root, \(G'_0\)
is a nonzero point of order two and the deck derivative has order
four. At a good IV or IV* root it is a nonzero point of order three
and the derivative has order three. At a good II root the point
is \(O\), and the derivative has order six. In particular, with
\(o=\operatorname{ord}(G'_0)\),
\[
  t=\frac{o}{\gcd(o,a)},\qquad
  o=2,3,1\quad\hbox{in those three cases}.             \tag{IM.7}
\]
The program retains this multiplication by \(a\). For example the
good E7 line with \(a=2\) has order one, whereas that with \(a=1\)
has order two.

\item Write the central source cap map as \(z\mapsto Dz+R\) on
\(E=\mathbb C/\Lambda_E\). Here \(D\ne1\), its order divides
\(m\), and \(1+D+\cdots+D^{m-1}=0\). The source part of
\(A^m=B^n\) gives \(nbG'_0=0\). We also have \(mG'_0=0\).
Since \(\gcd(m,n)=1\), the Bezout identity gives \(bG'_0=0\).
The actual \(B_0\) therefore covers the identity on \(E\), and
acts on \(N_0\) by a single nonzero constant \(c\). Put
\[
                  \tau=\frac{\log(c\lambda)}{2\pi i}.
                                                               \tag{IM.8}
\]

\item We now exhibit every period rather than just an invariant subspace.
The torsion line (IM.5) has its finite unitary character
\(\chi:\Lambda_E\to\mu_t\), of exact image order \(t\).
In a logarithmic line coordinate \(\eta\), its universal cover is
\(\mathbb C^2\), with generators
\[
 (z,\eta)\longmapsto(z+\omega,\eta+\theta_\omega),
 \quad e^{2\pi i\theta_\omega}=\chi(\omega),
 \qquad (z,\eta)\longmapsto(z,\eta+1).
                                                               \tag{IM.9}
\]
Here every \(\theta_\omega\) is real, and their classes together
with \(\mathbb Z\) generate exactly \((1/t)\mathbb Z\).
The quotient by \(B_\lambda\) adds the period \((0,\tau)\)
and is the actual compact central abelian surface.

\item The equality \(D^*N_0=N_0\) makes the finite unitary character
\(D\)-invariant. Thus an affine source map \(Dz+R\) has a
constant line lift in (IM.9). Every other holomorphic lift differs
by an invertible holomorphic function on the compact elliptic curve,
which is constant. The actual cap generator can consequently be
written, with its fixed character retained, as
\[
 A_\lambda(z,\eta)
       =(Dz+R,\eta+n\tau/m+\kappa).                  \tag{IM.10}
\]
The exact norm gives \(m\kappa\in\mathbb Z\): the source lift
has zero translation after \(m\) iterations because the geometric
sum of the powers of \(D\) is zero, and a deck period with zero
source coordinate has integral line coordinate. We retain \(\kappa\) in every
formula below; no finite character is discarded.

\item All holomorphic one-forms on the abelian cover are constant.
The affine linear part \(\operatorname{diag}(D,1)\) fixes just
the line generated by \(d\eta\). Integration over the full
affine deck group (IM.9)--(IM.10) gives its complete period image
\[
 \mathcal P=(1/t)\mathbb Z+\tau\mathbb Z
                         +(n\tau/m+\kappa)\mathbb Z. \tag{IM.11}
\]
Choose integers \(x,y\) with \(mx+ny=1\). Then the actual
period combination
\(x\tau+y(n\tau/m+\kappa)\) is
\(\beta=\tau/m+y\kappa\), and conversely
\[
  \tau=m\beta-my\kappa,\qquad
  n\tau/m+\kappa=n\beta+mx\kappa.
                                                               \tag{IM.12}
\]
The remainders belong to \((1/t)\mathbb Z\). These equations
prove the lattice equality, with its arrows in both directions,
\[
                 \mathcal P=(1/t)\mathbb Z
                             +(\tau/m+y\kappa)\mathbb Z.
                                                               \tag{IM.13}
\]
Replacing \((x,y)\) by \((x+ns,y-ms)\) changes the second
generator by the first lattice, so the result is independent of
the Bezout representative.

\item The map induced by \(\eta\) to \(\mathbb C/\mathcal P\)
is the Albanese map. This follows directly from the invariant
one-form and its complete integration lattice; equivalently its
fibre is the elliptic cover with period kernel \(\ker\chi\).
Multiplication by \(t\) gives its Tate coordinate
\[
           Q_i=e^{2\pi i t y\kappa}(c\lambda)^{t/m}.
                                                               \tag{IM.14}
\]
The norm \(n\), source affine term, and fixed line character occur
in the displayed constant. The exponent is determined by the actual
Picard point (IM.6).

\par

\medskip\noindent\textbf{4. The normalization and its period at a reduced-root cap}

\par\nobreak

\item We now use an actual reduced root of type II, III, IV, or I2. Its
central curve has, respectively, \(r=1,2,3,2\) rational components
in its normalization. We retain the precise root line, whose degree
on each component is zero by the line diagrams in P3F.36--49 and
P3F.52. Its pullback to each \(\mathbb P^1\) is therefore
trivial. This assertion does not trivialize the line on the singular
curve, nor discard its nodal or infinitesimal gluing data.

\item On the actual root line the two maps satisfy
\[
 B_\lambda=\lambda B_0,\qquad
 A_\lambda=\lambda^{n/m}A_0,\qquad
 A_\lambda^m=B_\lambda^n,
 \qquad[A_\lambda,B_\lambda]=1.                       \tag{IM.15}
\]
The fractional power here is on the chosen logarithm cover of the
scalar disc. For integers \(x,y\) with \(mx+ny=1\), form the
literal holomorphic line map
\[
                         H=B_\lambda^xA_\lambda^y.
                                                               \tag{IM.16}
\]
Negative powers are allowed because both maps are invertible in
this cap. Their exact norm and commutativity prove
\[
 H^m=B_\lambda^{mx+ny}=B_\lambda,\qquad
 H^n A_\lambda^{-1}
        =(B_\lambda^nA_\lambda^{-m})^x=1.
                                                               \tag{IM.17}
\]
Consequently their actual action group is precisely \(\langle H\rangle\).

\item The sphere cap's component action is transitive, as established by
the integral cap constructions. Since it is now generated by \(H\),
that generator cycles all \(r\) normalization components. Choose
fixed trivializations of their degree-zero lines and label the
components in this cyclic order. On them the literal maps have form
\[
 (z,V,j)\longmapsto
        (h_j(z),c_j\lambda^{1/m}V,j+1).               \tag{IM.18}
\]
The source maps \(h_j\) are Möbius transformations, independent
of \(\lambda\). The coefficients \(c_j\) are constant in
\(z\): an invertible map between the two trivial lines on compact
\(\mathbb P^1\) is a nowhere-zero holomorphic function on that
curve. The exponent is exactly
\(x+yn/m=1/m\). The constants retain every finite character and
the fixed singular-line gluing.

\item The stabilizer of one component is \(\langle H^r\rangle\).
Its quotient is therefore the explicit surface
\[
  \widetilde S_i=
  (\mathbb P^1\times\mathbb C^*)/\langle
       (z,V)\longmapsto(h(z),c_i\lambda^{r/m}V)
                                      \rangle .       \tag{IM.19}
\]
This is the normalization of the reduced cap support. To check the
map, normalize the root central curve and pull back its line. The
resulting finite map is equivariant for (IM.16). For sufficiently
small \(|\lambda|\), \(H^r\) strictly contracts the line norm,
so its action is free and properly discontinuous. Quotienting the
finite equivariant map, using compact fundamental annuli, gives a
finite map from (IM.19) to the actual central support. It is
birational because it is an isomorphism over the smooth locus of
the root central curve. The source is normal, so the finite map is
its normalization. This establishes the quotient arrow, including
the cuspidal type II and the multiple-intersection type III/IV
central curves.

\item The actual projection \([(z,V)]\mapsto[V]\) makes (IM.19) a
holomorphic \(\mathbb P^1\)-bundle over
\[
            E_{Q_i}=\mathbb C^*/Q_i^{\mathbb Z},
            \qquad Q_i=c_i\lambda^{r/m}.             \tag{IM.20}
\]
All holomorphic one-forms come from this elliptic curve, because
they vanish along the \(\mathbb P^1\) fibres. The homotopy exact
sequence of the displayed locally trivial bundle gives an
isomorphism on fundamental groups. Thus its full integration
lattice, and its Albanese curve, are exactly those of (IM.20).
In particular the source Möbius monodromy \(h\) introduces no
unrecorded isogeny into the Albanese period.

\item For the quadratic I1* cap, \(r=m=2\). Here \(B\) preserves
the two components while \(A\) exchanges them; (IM.17) still
cycles them and gives \(Q_i=c_i\lambda\). This cap imposes
no scalar-cover factor on its own. The reduced III and IV root
progressions have \(\gcd(m,r)=1\); they retain the full factor
\(m\). The type II normalization has \(r=1\).

\par

\medskip\noindent\textbf{5. Intrinsic obstruction to every smaller winding}

\par\nobreak

\item For a small nonzero Tate coordinate we have
\(j(Q)=Q^{-1}+744+O(Q)\), with the same normalization as in
P3K.36. For any root of unity \(\zeta\ne1\),
\[
      Q\bigl(j(Q)-j(\zeta Q)\bigr)
                         \longrightarrow1-\zeta^{-1}\ne0.
                                                               \tag{IM.21}
\]
There are only finitely many phase ratios for each fixed cap and
finitely many caps. We shrink the scalar disc once so that all the
nontrivial differences in (IM.21) are nonzero there.

\item After scalar winding \(N\), formula (IM.14) multiplies \(Q_i\)
by \(e^{2\pi iNt_i/m_i}\) at a good cap, and formula (IM.20)
multiplies it by \(e^{2\pi iNr_i/m_i}\) at a reduced cap. An
endpoint biholomorphism would give an isomorphism of the intrinsic
Albanese curves in Section 2, so it would make their \(j\)-invariants
equal. Equations (IM.21) therefore force, respectively,
\[
 m_i\mid Nt_i,
 \qquad m_i\mid Nr_i.                                \tag{IM.22}
\]
Thus the exact local necessary denominators are
\[
 d_i=\frac{m_i}{\gcd(m_i,t_i)}\quad\hbox{for a good root},
 \qquad
 d_i=\frac{m_i}{\gcd(m_i,r_i)}\quad\hbox{for a reduced root}.
                                                               \tag{IM.23}
\]
Each \(d_i\) divides \(m_i\), so distinct \(d_i\) are coprime.
The combined necessary winding is their product.

\item The actual root cases have the following removed factors:

\begin{longtable}[]{@{}llr@{}}
\toprule\noalign{}
Root type & Actual central datum & Factor \(m_i/d_i\) \\
\midrule\noalign{}
\endhead
\bottomrule\noalign{}
\endlastfoot
Good III or III* & \(t_i=2/\gcd(2,a)\) & \(t_i\) \\
Good IV or IV* & \(t_i=3/\gcd(3,a)\) & \(t_i\) \\
Good II & \(t_i=1\) & 1 \\
Reduced III, orders \(4j+1\) or \(4j+3\) & \(r_i=2\) & 1 \\
Reduced IV, orders \(3j+1\) or \(3j+2\) & \(r_i=3\) & 1 \\
Reduced II & \(r_i=1\) & 1 \\
Quadratic I1* root I2 & \(m_i=r_i=2\) & 2 \\
Absent or order-one slot & ordinary marked cap & 1 \\
\end{longtable}

\item In every nonempty F53 branch the product of these factors is exactly
the specified \(b\). The complete branch list and the exact
residue-class verification are recorded below and in the certificate.
Consequently
\[
                       \prod_i d_i=M/b.                \tag{IM.24}
\]
This proves the necessity in (IM.2). The quotient constructed in
P3K.35 proves its sufficiency, including every retained order-one
clutch and all allowed fans. Pulling that actual family back by
\(w\mapsto w^q\), or reversing the continuation direction,
gives the endpoint identifications for every integer multiple.
No statement about the mapping class of an additional endpoint
automorphism is required or deduced.

\par

\medskip\noindent\textbf{6. The component-supported OS49 cap}

\par\nobreak

\item The good IV* root of P3F.26 has the actual central line
\(\mathcal O(G'_0-O)\), where \(G'_0\) is a nonzero
three-torsion point. Its order is \(m=3k\). The component-supported
line at the I2 fibre changes neither this Picard point nor the
central IV* root diagram. The cap calculation of Section 3 therefore
has \(t=3\). Its intrinsic Albanese Tate parameter has exponent
\(3/m=1/k\).

\item For completeness we retain the mixed logarithmic marking directly.
With the actual covector marking of P3F.26, invariant covectors are
\(\psi\) and \(h=-w+3\delta\). In the coherent central period
matrix
\[
 \Pi=[(-3\mu,\beta),(-\tau_E,-\mu),(1,0),(0,1)],
 \qquad \tau_E^2-3\tau_E+3=0,\quad \mu=\tau_E/3,
\]
the invariant holomorphic form is \(-da+3db\). Its period row
is \(h+T\psi\), where
\[
 T=3\beta+\tau_E,
 \qquad \beta=\frac{\log\lambda}{2\pi i}+c_0.
\]
The actual affine vector
\(n(e_\gamma-e_u)+d_\delta\delta\) adds
\((nT+3d_\delta)/m\). For \(n=\pm1\), the complete period
lattice is
\[
 \mathbb Z+\mathbb Z\frac{T+3nd_\delta}{m}.
                                                               \tag{IM.25}
\]
Indeed the added period is \(n\) times the displayed second
generator, and \(T\) differs from \(m\) times that generator
by the integer \(3nd_\delta\). The generator \(h\) gives the
primitive real period one. This proves directly that the fixed
delta term and either norm sign affect the constant, while
\[
                        Q_i=c_i\lambda^{1/k}.          \tag{IM.26}
\]
The algebraic-reduction argument of Section 2 applies to these
actual spheres as well. Their unique nontrivial multiple fibre has
order \(3k\), so every endpoint map preserves it. Equation (IM.21)
forces \(k\mid N\). P3K.37 supplies the exact quotient
\(\lambda=w^k\), proving sharpness. This conclusion uses the
actual cap periods; the component line square remains the separate
value \(-2k\) established in P3K.34.

\end{enumerate}
\end{proof}




\begin{theorem}\label{night:P3K.42}
\StatementRecord{night:P3K.42}{proved}{S6\_06\_BLIND: GPT-6 Astra (OpenAI, ChatGPT Pro mode, effort Pro), recorded re-derivation in the checkpoint dated 18 November 2026}{\cref{night:P3F.22}, \cref{night:P3F.55}, \cref{night:P3F.56}, \cref{night:P3F.8}, \cref{night:P3F.9}, \cref{night:P3K.40}, \cref{night:P3K.68}}
Retain the actual positive selected-cusp source loop of P3F.54 and
P3F.56, with fixed sufficiently small scalar, the III-star cap, no
logarithmic modification on the coalescing finite pair, and the
initial physical-delta-edge fan. Its parameter is one positive
\(\rho\)-loop, where \(a(u)=\rho^2\), with the literal double
source closure and the actual Cox normal product. In Jin's endpoint
and physical nonzero-vector-line conventions its characteristic pair and
first-resolution relative coordinates are
\[
       (k_M,\ell_M)=(-6,0),\qquad(k_U,\ell_U)=(6,0),\qquad
       (m_{\rm rel},n_{\rm rel},s)=(0,-6,-1).
\]
The framed local surgery has pair \((5,-4)\). The actual capped
normal sphere bundle in its source-family bordism has pair
\((1,4)\). Their oriented sum gives the meridian-one U pair; the actual physical line M is its Hermitian dual and gives the displayed negative pair.
This statement concerns only the specified initial fan. The complete physical line and Hopf-transfer correction is P3K.68.
\end{theorem}
\begin{proof}
\begin{enumerate}[label=\arabic*.]
\item \par

\medskip\noindent\textbf{SP.1. The framed local surgery and its actual four-cycle}

\par\nobreak

\item Let \(e:S^3\hookrightarrow S^6\) be the actual lifted link of P3F.55.
Its projection is the immersion of complex Smale length two, and its
third normal field is the actual increasing height direction. Fix the
first normal field in the oriented complex normal plane of the
projection, and complete the frame \((\nu _1,\nu _2,\nu _3)\) with
that height field. Normalizing and applying Gram--Schmidt uses a
positive triangular homotopy. It therefore retains the actual
projection-normal plus height framing, rather than adding an element
of \(\pi _3(SO(3))\).

\item Takase, Corollary 2.15, supplies a Seifert embedding
\(V\hookrightarrow S^6\) whose outward boundary normal is this
specified \(\nu _1\). The proof of his Theorem 3.2 gives, for this
projection and this boundary normal,
\[
              2\omega=3\sigma(V)+e(\nu V)^2.             \tag{SP.2}
\]
The invariant and generator comparison in P3F.55, equations
(CL.16)--(CL.19), fixes the convention in (SP.2): the punctured
Kummer test has signature \(-16\) and Smale invariant \(-24\)
in both conventions. Hence the actual local projection has
\(\omega=-2\). The same comparison gives
\[
             8H(e)=e(\nu V)^2-\sigma(V),\qquad H(e)=1.
                                                               \tag{SP.3}
\]
Solving these two displayed equalities gives
\[
                    e(\nu V)^2=5,\qquad \sigma(V)=-3.    \tag{SP.4}
\]

\item We now identify these numbers with the characteristic numbers of the
\emph{framed} surgery manifold, including its meridian line. Let
\[
 N_e=(S^6\setminus\operatorname{int}(S^3\times D^3))
                  \cup(D^4\times S^2)
\]
use precisely the preceding normal frame. Along the Seifert collar,
the vector from its boundary into \(V\) is \(-\nu _1\).
Consequently the truncated Seifert surface has boundary
\(S^3\times\{-\nu _1\}\), and caps by the literal constant section
\(D^4\times\{-\nu _1\}\). Denote this closed four-manifold by
\(\widehat V\).

\item Here are the orientations of this cap. Write the ambient orientation
near the knot as
\(\Omega_{S^3}\wedge dn_1\wedge dn_2\wedge dn_3\).
The Seifert orientation is
\(dn_1\wedge\Omega_{S^3}\), so its oriented normal two-plane is
\(-dn_2\wedge dn_3\). At \(-\nu _1\), the outward-first
orientation of the normal meridian sphere is also
\(-dn_2\wedge dn_3\). The boundary of the knot complement is
\(+\Omega_{S^3}\wedge\Omega_{S^2}\); the surgery cap therefore
has orientation \(-\Omega_{D^4}\wedge\Omega_{S^2}\).
The cap of \(\widehat V\) has orientation \(-\Omega_{D^4}\).
Thus its oriented normal plane is exactly the positive meridian
plane, with no sign change in its Euler class.

\item Let \(x=x_U\) be the meridian class with integral one on this sphere.
It is the Thom-line convention; the physical nonzero-vector line is its
dual, as proved explicitly in SP.2 and SP.5 below.
The Mayer--Vietoris map in degree two identifies the boundary
\(H_2(S^3\times S^2)=\mathbb Z\) with the meridian summands in
both pieces; it gives \(H^2(N_e)=\mathbb Zx\). In degree four,
the connecting map identifies \([\widehat V]\) with the oriented
boundary \([S^3]\). The normal-plane calculation above therefore
gives the actual duality identity
\[
       \operatorname{PD}_{N_e}(x)=[\widehat V],\qquad
       x|_{\widehat V}=e(\nu\widehat V).                 \tag{SP.5}
\]
The cap frame is constant. Its Euler square is the relative square
in (SP.4), and capping by a four-ball leaves the signature unchanged.
Restriction of the ambient tangent bundle to this embedded cycle is
\(T\widehat V\oplus\nu\widehat V\). It follows directly that
\[
 \begin{aligned}
 k(N_e)&=\int_{N_e}x^3=\int_{\widehat V}e(\nu\widehat V)^2=5,\\
 \ell(N_e)&=\int_{N_e}p_1(TN_e)x
       =3\sigma(\widehat V)+e(\nu\widehat V)^2=-4.
 \end{aligned}                                             \tag{SP.6}
\]
This constructs the four-cycle and its Euler-line restriction; it
does not recover the pair from the Haefliger number alone.

\par

\medskip\noindent\textbf{SP.2. The actual Cartier-to-Cox coordinate map}

\par\nobreak

\item Use the normalization and branch notation of F56:
\[
 N=\mathbf P^1_s\times D_\rho,\quad
 x=s-\rho,\quad y=s+\rho,\quad Z=x/y,\quad
 D_+=(x=0),\quad D_-=(y=0).
\]
Remove the inner four-ball to obtain \(W\). The actual determinant
normal isomorphism of F56 is
\[
 L=\mathcal O(-D_+-D_-),\qquad
 E=\mathcal O(D_+)\oplus\mathcal O(-D_-),\qquad
 \operatorname{ad}E=\mathbb R_\gamma\oplus L_{\mathbb R}.
                                                               \tag{SP.7}
\]
The determinant trivialization is multiplication by \(x/y\).
In the finite \((s,\rho)\) ball, the Cartier frames of \(E\) are
\((1/x,y)\); their determinant is \(y/x\), so the stated
trivialization sends their wedge to one. These are nonvanishing
frames of the indicated line bundles even on the corresponding
divisors.

\item The morphism \(r=[x:y]:W\to\mathbf P^1\) identifies \(E\)
with \(r^*(\mathcal O(1)\oplus\mathcal O(-1))\). On the chart
\(y\ne0\), the coefficient arrow from the Cartier frames to the
standard dual-tautological and tautological frames is exactly
\[
                         \operatorname{diag}(1/y,y).      \tag{SP.8}
\]
Indeed the first pulled-back divisor frame is \(y/x\), and the
second divisor is absent from this chart. Equation (SP.8) follows by
expressing \(1/x\) and \(y\) in those frames. The corresponding
equation on \(x\ne0\) agrees on the overlap.

\item The actual positive-weight spinor coefficients are
\((\overline{c_0},c_1)\), while the right-complex Cox coefficients are
\((c_0,\overline{c_1})\). Put \(Y=y/|y|\). The full coefficient map
from (SP.8) is therefore
\[
 (\overline{c_0},c_1)=\operatorname{diag}(\overline Y,Y)(a,b),
 \qquad (c_0,\overline{c_1})
      =\operatorname{diag}(Y,\overline Y)(\overline a,\overline b).
\]
Multiplication by the normalized quaternion \(1+jZ\) gives the actual
anti-linear arrow
\[
 (a,b)\longmapsto g(x,y)\overline{(a,b)},\qquad
 g(x,y)=\frac1{\sqrt{|x|^2+|y|^2}}
 \begin{pmatrix}y&-\overline x\\x&\overline y\end{pmatrix}.
                                                        \tag{SP.9}
\]
It intertwines the positive scalar action on the spinor with the
physical negative right scalar on the Cox quaternion. The previous
linear summand swap omitted this conjugation.

\item For \(J=\left(\begin{smallmatrix}0&-1\\1&0\end{smallmatrix}\right)\)
and every SU(2) transition, \(\overline g=JgJ^{-1}\). Thus the
global anti-linear map \(Q=J^{-1}C\),
\[
 Q(a,b)=(\overline b,-\overline a),\qquad
 Q(e^{i\theta}v)=e^{-i\theta}Q(v),                       \tag{SP.10}
\]
commutes with all clutches. It has real four-dimensional determinant
one, while its adjoint is the antipodal map of the real normal
three-plane, of determinant minus one. It extends over both cap
normal discs by minus the identity. The Hermitian dual associated-line
map is retained on both cap collars; every extending scalar gauge is
inverted pointwise. The conjugated complex clutch has the same second
Chern number, since conjugation by J is an actual constant bundle map.

\item The physical meridian sign follows in the literal F9 chart
\((x,v,z)\), with \(y=vz\), weights \((0,-1,+1)\), and
\(b=(|z|^2-|v|^2)/2\). Its positive spinor is
\((\overline v,z)\). At \(v=1,z=0\), the projective coordinate
\(z/\overline v\) has the complex orientation of \(y\), while
the outward raw normal sphere at \(b<0\) has its negative. The
physical circle-first normal orientation is the negative raw
orientation. Thus the physical spinor projective sphere has that
physical normal orientation. Its physical nonzero-vector line is
tautological, with Chern number minus one. The meridian-one Thom
line U is its Hermitian dual; on the source page it is M inverse.
This is a coefficient and circle map, not merely a degree comparison.

\item The inner projection-normal frame is the Cartier frame used here,
up to an extending gauge. To verify that assertion rather than
postulate it, use F56's explicit nonzero normal covector
\[
       (3s^2-\rho^2)\,dX-2s\,dy
                          +2\rho(s^2-\rho^2)\,d\rho .     \tag{SP.11}
\]
The determinant exact sequence identifies its normal quotient with
the canonical line of the normalization. The ratio of this line
identification and the Cartier identification is a holomorphic unit
on the punctured finite ball. Hartogs extension extends it to the
ball, and a holomorphic function in two variables cannot have its
only zero at the origin. Its extension is therefore a unit. The
resulting real two-plane frame comparison extends over the inner
four-ball. Adding the height line and the explicit triangular
graph-normal shear gives the actual frame of SP.1, with no free
\(SO(3)\) clutch.

\item The normal-coordinate unit can also be checked directly in both
actual fixed-edge charts. We use the physical source coordinate
\(q\) of F56, not its Tate multiple. With \(X=xy\), restriction
of the differential of that displayed source equation to the
normalization gives
\[
 dq=-\frac{xy}{xy+u(\rho)}\,n,
 \qquad n\text{ is the covector in (SP.11)}.              \tag{SP.11a}
\]
The first Cox chart has \(q=Zc_0c_1\), and the second has
\(q=Z^{-1}c'_0c'_1\), with \(c'_i=Zc_i\). Their quotient
complex normal coordinates consequently satisfy, on the fixed curve,
\[
 \begin{aligned}
 d(c_0c_1)&=-\frac{y^2}{xy+u(\rho)}\,n,\\
 d(c'_0c'_1)&=-\frac{x^2}{xy+u(\rho)}\,n.
 \end{aligned}                                             \tag{SP.11b}
\]
These are exactly the two Hom arrows of (SP.8), multiplied by the
same explicit unit \(-1/(xy+u(\rho))\). Since \(u(0)=u_*\ne0\),
it has an analytic logarithm and square root on the finite inner
ball. Its spinor adjustment therefore extends over that ball. This
verifies the actual normal-unit comparison, as well as its line
transition. It supplies no additional relative Chern number.

\item There is also an exact check on the separate normal factors. Once
their product has been matched by (SP.11b), any remaining
weight-preserving linear normal gauge has the form
\((c_0,c_1)\mapsto(g c_0,g^{-1}c_1)\). Remove its positive
modulus by the actual equivariant positive rescaling. For its unitary
phase, both right-complex spinor coordinates are multiplied by
\(g/|g|\); hence
\[
                  q_H\longmapsto q_H\,g/|g|.             \tag{SP.11c}
\]
This is the physical circle action, and its adjoint Hopf normal
coordinate is exactly unchanged. It contributes no additional
\(SO(3)\) clutch. A possible parameter winding of this circle-line
gauge is retained and extends over \(D^3\times S^1\), whose
restriction map on \(H^1\) is an isomorphism. We therefore do not
assume that each separate normal gauge has a logarithm across the
parameter origin.

\par

\medskip\noindent\textbf{SP.3. The outer cap: the actual product arrow}

\par\nobreak

\item The outer coordinate change is the literal map of F56
\[
 A_\rho=\begin{pmatrix}1&-\rho\\1&\rho\end{pmatrix},
 \qquad
 \mathcal B_\rho(P)=
       \frac{A_\rho P A_\rho^*}{\operatorname{tr}(A_\rho P A_\rho^*)}.
                                                               \tag{SP.12}
\]
Here the closed three-ball is the positive semidefinite Hermitian
matrices of trace one. Its boundary consists of rank-one matrices.
This is a periodic diffeomorphism of the ball, including its centre,
and its boundary map is exactly \(r\). We use its actual product
extension on the surgery collar, rather than replacing the normal
frame by its derivative.

\item The collar in P3K.40 and F56 is the literal first-resolution product
\((h,q_H)\), where \(h\) is the base sphere coordinate and
\(q_H=(1+jZ)(c_0+j\overline{c_1})\) is the normalized actual
Cox quaternion. The physical circle acts by
\(q_H\mapsto q_H e^{-i\theta}\). Thus on the normal sphere
cover the exact outer cap map is
\[
 (P,q_H,\rho)\longmapsto
                  (\mathcal B_\rho(P),q_H,\rho)
 \quad\text{on }D^3\times S^3\times S^1.                 \tag{SP.13}
\]
It is a periodic diffeomorphism, is equivariant for that same circle,
and restricts to the actual base-coordinate change on the boundary.
Its quotient is
\((\mathcal B_\rho,\operatorname{Id}_{S^2})\). It therefore
extends both the actual Cox normal frame and the meridian circle
line over the product cap. This is precisely the product surgery
collar used in F56: its second factor is the actual Hopf normal
sphere, and (SP.13) holds that factor, with its orientation, fixed.

\item In particular (SP.13) keeps the full determinant winding. For
clarity the exact determinant discrepancy can be written down.
The Hom arrow induced by (SP.8) takes the Cartier frame \(xy\)
of \(L\) to \(y^2\) times the standard \(\mathcal O(-2)\)
frame. Under the standard canonical identification this is
\[
                  y^2\,dZ=2\rho\,ds.                    \tag{SP.14}
\]
Thus the actual Cox arrow has the factor
\(\det A_\rho=2\rho\) relative to the cotangent inverse.
Its winding is one. We neither remove it nor identify (SP.13) with
the tangent derivative of \(\mathcal B_\rho\). It remains in
the Cartier-to-Cox coordinate matrix (SP.9) and hence in the clutch
computed below. The earlier attempted derivative comparison that
omitted this factor is preserved as a superseded unpromoted attempt.

\item All coordinates in (SP.13) are the actual coordinates already
matched in F56 and P3K.40. In particular we make no additional
replacement of \(q_H\) by a Tate-scaled quaternion. The winding
\(-6\) of \(q_T/q\) is retained by those matched charts.
Equations (SP.11a)--(SP.11b) use the physical \(q\) in those
charts; the nonzero factor \(q_T/q\) remains in the period
identification of the adjacent cell. Replacing that period by one
with factor one would change the marking and is not done here.
The original source and line unit gauges used in F56 extend over
the full source disc. Their displayed logarithms give extending
isotopies on the same normal collar. The induced isotopies of its
boundary extend over the product cap by extending their invariant
vector fields with a collar cutoff; averaging preserves the
physical circle action. This carries these actual gauges in
(SP.13), rather than changing the cap framing.

\par

\medskip\noindent\textbf{SP.4. The relative second Chern number and the normal sphere bundle}

\par\nobreak

\item Cap \(W\) by the inner four-ball in its projection-normal frame
and the outer \(D^3\times S^1\) in its actual Cox frame. The
resulting base and spinor bundle will be denoted by
\(\overline W\) and \(\overline E\). The base orientation is
\[
 \overline W=(S^2\times D^2)\cup(D^3\times S^1)
                      =\partial(D^3\times D^2).
                                                               \tag{SP.15}
\]
The first piece has its complex orientation, and the second has
minus its displayed product orientation. Thus the boundary of the
outer cap is the negative of
\(K=\mathbf P^1\times S^1\), with \(K\) oriented as the
boundary of \(N\).

\item Use the conjugate spinor on the inner ball, with its conjugate Cartier
frame, and the actual right-complex Cox frame on the outer cap.
The coefficient map between these is exactly the complex-linear matrix
g in (SP.9), after retaining the global anti-linear map (SP.10).
Its Chern number equals that of E by the explicit J conjugacy.
The resulting bundle extends across the inner ball and is a trivial
\(SU(2)\) bundle over \(N\simeq S^2\). Choose a global frame
agreeing there with the Cartier frame. Such agreement can be made
on the ball because the difference of two extending frames is
itself an extending gauge. Compare this frame with the Cox frame on
\(W\). On the inner boundary its complex-linear coordinate matrix after
that conjugation is exactly g in (SP.9).
Since this comparison extends over \(W\), its degree on the outer
boundary \(K\) is the same as its degree on the inner sphere in
its positive complex-domain orientation.

\item We calculate that degree with its precise trace convention. The
first column of (SP.9) is \((y,x)\); the map
\((s,\rho)\mapsto(y,x)\) is complex linear with nonzero
determinant \(-2\). Hence its normalized boundary map has degree
one for the complex first-column orientation. At the identity the
positive tangent matrices in that orientation are
\[
 I=\begin{pmatrix}i&0\\0&-i\end{pmatrix},\quad
 J=\begin{pmatrix}0&-1\\1&0\end{pmatrix},\quad
 K_0=\begin{pmatrix}0&i\\i&0\end{pmatrix}.
                                                               \tag{SP.16}
\]
They satisfy \(IJ=-K_0\) and \(\operatorname{tr}(IJK_0)=2\).
Consequently the right Maurer--Cartan form has
\[
           \int_K\operatorname{tr}(dg\,g^{-1})^3=24\pi^2.
                                                               \tag{SP.17}
\]
This explicitly fixes the sign; the quaternion convention with the
opposite group orientation cannot be substituted in (SP.17).

\item Equation (SP.9) is a coordinate arrow from the frame on \(N\) to
the cap frame. A connection that is zero in the first frame is
therefore \(-dg\,g^{-1}\) in the latter frame near their boundary.
Writing \(\omega=dg\,g^{-1}\), we have
\(d\omega=\omega^2\) and
\[
 \operatorname{tr}\bigl((-\omega)d(-\omega)
                    +\tfrac23(-\omega)^3\bigr)
                    =\tfrac13\operatorname{tr}(\omega^3).
\]
Extend this connection into the cap with a collar cutoff. Stokes'
formula and its negative boundary orientation give the actual
clutch number
\[
 \begin{aligned}
 c_2(\overline E)[\overline W]
    &=\frac1{8\pi^2}\int_{\text{cap}}\operatorname{tr}(F^2)\\
    &=-\frac1{24\pi^2}
                       \int_K\operatorname{tr}(dg\,g^{-1})^3=-1.
 \end{aligned}                                             \tag{SP.18}
\]
For the actual adjoint bundle,
\(p_1(\operatorname{ad}\overline E)=-4c_2(\overline E)\).
This identity follows directly on the diagonal circle: the spinor
has Chern roots \(a,-a\), and the nontrivial real adjoint plane has
Euler class \(2a\). Therefore
\[
             p_1(\operatorname{ad}\overline E)[S^4]=4.    \tag{SP.19}
\]

\item Let \(B=S(\operatorname{ad}\overline E)\to S^4\), with the
actual positive Hopf meridian class \(x\) from F56. Its determinant
class is zero, so the oriented vertical tangent Euler class is
\(e_v=2x\). The exact splitting
\(T_vB\oplus\mathbb R\cong\pi^*(\operatorname{ad}\overline E)\)
gives \(4x^2=\pi^*p_1(\operatorname{ad}\overline E)\).
Since \(\pi_*x=1\) and \(p_1(TS^4)=0\), we obtain
\[
                         (k(B),\ell(B))=(1,4).             \tag{SP.20}
\]
Its index is zero, as already used in F56. The present calculation
retains the individual numbers of this actual normal bundle.

\par

\medskip\noindent\textbf{SP.5. Both line conventions and the actual physical return}

\par\nobreak

\item The base orientation is the complex orientation on W. In the actual
F9 normal coordinates the ambient orbit orientation is
\(-\Omega_{\rm source}\wedge d\gamma\), and the physical normal
orientation is \(-d\gamma\wedge\Omega_L\). The anti-linear map
(SP.10) supplies the actual normal reflection between the coefficient
frames. Its extensions over the inner D4 and outer D3 times S1 caps
carry the dual circle character and every retained boundary gauge.
In particular the meridian-one line U extends through both product
surgery collars with exactly the line used in (SP.5) and (SP.20).

\item The outward boundary of the seven-dimensional quotient bordism is
\(-P_f\sqcup-N_e\sqcup-B\): the first sign is the five-dimensional
orbit followed by the outward parameter radial coordinate, and each
inner sign is the outward boundary of a removed normal disc. Thus
Stokes for U gives
\[
 (k_U,\ell_U)(P_f)=-(5,-4)-(1,4).
\]
The physical family is Jin's mapping torus for the inverse return;
reversing that parameter gives
\[
                (k_U,\ell_U)_{\rm Jin}=(6,0).
                                                        \tag{SP.21}
\]
The raw local complex-plus-height six-dimensional orientation and the
physical orbit orientation are opposite. The outward-meridian-one
Thom line changes sign with that normal orientation; hence the raw
Seifert values (SP.6) enter the U calculation with the signs displayed.
The physical nonzero-vector line M instead satisfies M=U inverse,
so in the same oriented Jin manifold
\[
           (k_M,\ell_M)=-(k_U,\ell_U)=(-6,0),
           \qquad A_M=-1.                              \tag{SP.22}
\]

\item For completeness the Hopf transfer uses this physical line. The
actual complex exterior is D3\_y times S3\_q, and the sphere is
minus the boundary of D3\_y times D4\_q. Physical M is positive on
the oppositely oriented Hopf base; reverse the y factor at the same
time, so the quotient orientation is unchanged. The F22 positive
Hopf factor is then u=+q and y'=-y. If L=D4\_y times S3\_q and
R=S3\_y times D4\_q, its first and second pieces are -R and -L.
The first-to-second clutch r maps to F(r); rewriting with L first
sets x=F(r), so the second point is F inverse of x. Consequently
lambda\_M(F) is minus Sigma\_(F inverse). The removed ball is
D4\_u times D3\_y', whose boundary is the original complex sphere;
the F8 cut-and-reglue maps coincide on its actual free face and are
the identity on its fixed-tube face. This proves that lambda\_M is
the physical clutched sphere. For U the pieces are +R,+L and the
same clutch gives lambda\_U(F)=Sigma\_(F inverse), its orientation
reverse. All circle gauges are carried by the Hermitian dual map.
These exact arrows are written in mathematical notation in P3K.68.

\item Therefore the physical source class is -1 modulo 28. The exact
relative-coordinate identities in the physical convention are
\(m=k_M-6A_M\), \(n=6A_M\), giving
\[
                    (m,n,s)=(0,-6,-1).                  \tag{SP.23}
\]
The native central generator still has sign plus one: exchanging
the two Hopf pieces inverts the clutch, and the outer orientation
reversal cancels that inverse on a disc-supported central class.
No value on another fan is assigned here.

\end{enumerate}
\end{proof}




\begin{theorem}\label{night:P3K.47}
\StatementRecord{night:P3K.47}{proved}{S6\_06\_BLIND: GPT-6 Astra (OpenAI, ChatGPT Pro mode, effort Pro), recorded re-derivation in the checkpoint dated 18 November 2026}{\cref{night:P3F.54}, \cref{night:P3F.56}, \cref{night:P3K.45}, \cref{night:P3K.46}}
Retain the actual selected-cusp positive \(\rho\)-loop of
P3F.54/P3F.56 and its initial fan. Its three-curve plumbing has no
periodic smooth product which intertwines the physical delta
action and carries all three labelled complex curves, their two
labelled vertices and their positive complex orientations to
constant copies. The obstruction is the actual relative normal
determinant winding
\[
 \operatorname{wind}\kappa+2\operatorname{wind}\alpha
                  -\operatorname{wind}\eta=-6.
\]
This concerns that simultaneous marked product. It does not obstruct
the separate single-star products or the actual reflected fixed-edge
cap map of P3K.46, and it does not assign a nonfixed relative integer
or an arbitrary-fan return formula.
\end{theorem}
\begin{proof}
\begin{enumerate}[label=\arabic*.]
\item \par

\medskip\noindent\textbf{PO.1. The actual family and the derivative marks}

\par\nobreak

\item Use the two actual vertex charts \((a,b,c)\) and \((u,v,w)\), with
\(q=abc=uvw\). The fixed curve \(C_0\) has tangent coordinates
\(b,w\), the other two labelled curves \(C_1,C_2\) have tangent
coordinates \((c,v)\) and \((a,u)\), respectively. The actual
transition arrows from OE.4 are
\[
 \begin{array}{c|ccc}
 &u&v&w\\ \hline
 C_0&\kappa\alpha bc&\alpha^{-1}ab&\kappa^{-1}b^{-1}\\
 C_1&g\eta bc&\eta^{-1}c^{-1}&g^{-1}ac\\
 C_2&a^{-1}&ab&ac.
 \end{array}                                                   \tag{PO.1}
\]
Only their restrictions and normal derivatives along the curves are
used below. Thus all units in the proof mean their restrictions at
\(q=0\), as functions of rho. The exact F54/F56 units satisfy
\[
 \operatorname{wind}\kappa=-6,\qquad
 \operatorname{wind}\alpha=\operatorname{wind}\eta=0. \tag{PO.2}
\]
The nonzero scalar is fixed, and \(\eta=\lambda\beta\). No physical
height character is changed to \(q_T\).

\item Suppose the stated periodic marked product existed. Read it as a map
from the actual plumbing over rho to one fixed reference plumbing.
At the two vertices its derivatives on the tangent lines of
\(C_2,C_1\) have nonzero complex factors
\[
               A,C\quad\hbox{and}\quad U,V,
                                                               \tag{PO.3}
\]
respectively. These factors are complex linear even though the product
is only smooth. Indeed these real two-planes have physical circle
weights \(-1,+1\) at the first vertex and \(+1,-1\) at the second;
the map intertwines that circle, retains each labelled axis, and
retains its complex orientation. The commutant on an individual such
real rotation plane is the complex scalar group. Put
\[
 a_* =\operatorname{wind}A,\quad
 c_* =\operatorname{wind}C,\quad
 u_* =\operatorname{wind}U,\quad
 v_* =\operatorname{wind}V.                            \tag{PO.4}
\]
All are integers because the product is periodic. We do not assume
the derivative stays diagonal along the fixed curve away from the
vertices.

\par

\medskip\noindent\textbf{PO.2. The two curve products force their endpoint winding equations}

\par\nobreak

\item Here is the elementary pointed-curve argument, with its actual unit
retained. For a curve having tangent coordinates \(z_+,z_-\) at its
two labelled vertices and transition
\(z_-=\xi(\rho)^{-1}/z_+\), write the reference transition with the
fixed nonzero constant \(\xi_*\). A marked equivariant product on
the open orbit cylinder gives an orientation-preserving equivariant
map \(F_\rho:\mathbb C^*\to\mathbb C^*\). The ratio
\(F_\rho(z_+)/z_+\) is invariant under the rotation of \(z_+\), so
its rho winding is independent of the radius and of the angle of
\(z_+\). This remains true on a rho-dependent nonzero path by
radially interpolating its modulus to a fixed modulus.

\item Near the first endpoint that ratio tends to its nonzero derivative
factor \(P(\rho)\). At the second endpoint the derivative factor
\(Q(\rho)\), in the two displayed endpoint coordinates, gives
\[
 \frac{F_\rho(z_+)}{z_+}
          \longrightarrow\frac{\xi(\rho)}{\xi_*Q(\rho)}.
\]
Taking winding numbers proves
\[
                   \operatorname{wind}P+
                   \operatorname{wind}Q
                         =\operatorname{wind}\xi.       \tag{PO.5}
\]
The assertion concerns the actual smooth map, not a replacement by
a complex automorphism of the curve. Equivariance is enough for the
ratio argument; the positive complex orientation fixes its degree
on the open orbit cylinder.

\item For \(C_2\) the unit \(\xi\) is one, while for \(C_1\) it is
\(\eta\), by (PO.1). Consequently (PO.5) gives
\[
                 a_*+u_*=0,\qquad
                 c_*+v_* =\operatorname{wind}\eta.     \tag{PO.6}
\]

\par

\medskip\noindent\textbf{PO.3. The full U(2) determinant cannot match these endpoints}

\par\nobreak

\item On the normal real four-plane to \(C_0\), the infinitesimal physical
delta action defines a canonical complex structure: conjugate its
negative-weight complex line and keep its positive-weight line. At
the first vertex use the ordered complex normal coordinates
\((\overline a,c)\); at the second use
\((\overline v,u)\). Along the fixed curve any smooth equivariant
normal derivative is a complex-linear invertible two-by-two map in
this structure. After a positive Hermitian homotopy it can be viewed
as U(2), but the argument below uses its nonzero complex determinant
directly, without assuming diagonality.

\item The actual fixed-curve normal transition obtained from (PO.1) is
\[
 \begin{pmatrix}\overline v\\u\end{pmatrix}
 =\begin{pmatrix}
       \overline\alpha^{-1}\overline b&0\\
       0&\kappa\alpha b
   \end{pmatrix}
   \begin{pmatrix}\overline a\\c\end{pmatrix}.
                                                               \tag{PO.7}
\]
Its determinant is
\[
              \kappa\,\frac\alpha{\overline\alpha}\,|b|^2.
                                                               \tag{PO.8}
\]
In the reference plumbing the corresponding determinant has constant
phase: its units are their fixed initial values and the source-coordinate
factor is the positive real norm square. In particular its phase has
zero rho winding, regardless of the smooth source-coordinate map in
the assumed product.

\item Let \(G_+,G_-\) be the normal-derivative matrices of that product,
expressed in the two actual normal charts and a nonzero global
determinant frame on the reference normal bundle. Such a reference
determinant frame exists explicitly: (PO.8) at the fixed reference
parameter has constant phase and positive variable factor, so its
two chart frames glue after positive rescaling and a constant phase.
The winding of \(\det G_+\) on any overlap parameter loop equals
its winding at the first vertex. This follows by contracting the
point in the first source-coordinate disc to that vertex while
keeping the determinant nonzero. The same argument in the second
disc applies to \(G_-\), even when its overlap coordinate depends
on rho. At the vertices (PO.3) gives, in the stated positive-weight
complex structure,
\[
 \operatorname{wind}\det G_+=-a_*+c_*,\qquad
 \operatorname{wind}\det G_-= -v_*+u_* .             \tag{PO.9}
\]

\item The derivative is one actual bundle map. Therefore its two matrices
and (PO.7) must commute with the reference normal transition. Taking
determinants and winding numbers in this exact square gives
\[
 (-v_*+u_*)+\operatorname{wind}\kappa
             +2\operatorname{wind}\alpha
                     -(-a_*+c_*)=0.                   \tag{PO.10}
\]
This equality holds for arbitrary invertible complex normal matrices
between the vertices. Off-diagonal U(2) terms cannot change it.
Substituting the two actual endpoint equations (PO.6) leaves
\[
 \operatorname{wind}\kappa+2\operatorname{wind}\alpha
                     -\operatorname{wind}\eta=0.       \tag{PO.11}
\]
Its left-hand side is \(-6\) by (PO.2). This is the required
contradiction. Thus the stated simultaneous marked equivariant
product does not exist.

\par

\medskip\noindent\textbf{PO.4. Exact scope of the obstruction}

\par\nobreak

\item The obstruction is the relative determinant winding of the actual
fixed-curve normal family, with its endpoint frames supplied by the
two other labelled pointed curve products. It is already detected
by the determinant; a contraction in SU(2) does not alter (PO.10).
The earlier even-winding factorization in OE.4 transfers the seam
between vertices and retains a physical-circle loop. It does not
satisfy both endpoint equations together with the determinant square.

\item In particular a periodic equivariant product connection on this
actual labelled plumbing, with identity return on that entire
marked neighbourhood, would give the forbidden product by its
parallel transport. This identifies the precise missing hypothesis
in an attempted direct application of the scalar-loop simultaneous
product argument. It makes no assertion about another global
representative whose isotopy moves these labelled curves, or about
an all-fan computation retaining the nontrivial boundary clutch.

\item The separate fixed-edge product and the nonfixed-star products
remain valid. The actual half-source reflected-family and cap maps
of P3K.45/P3K.46 remain valid as well: they transport one fixed-edge
product and do not assert a simultaneous product of the three
labelled curve stars.

\item The exact determinant and integer-consistency certificate is
\path{code/part3/source_cusp_product_obstruction.py}. No numerical
nonfixed relative return is assigned by this proof.

\end{enumerate}
\end{proof}




\begin{theorem}\label{night:P3K.48}
\StatementRecord{night:P3K.48}{proved}{S6\_06\_BLIND: GPT-6 Astra (OpenAI, ChatGPT Pro mode, effort Pro), recorded re-derivation in the checkpoint dated 18 November 2026}{\cref{night:P3K.23}, \cref{night:P3K.42}, \cref{night:P3K.47}}
The actual selected-cusp family of P3F.54/P3F.56 admits a periodic
orientation-preserving ordinary smooth product of its labelled
three-curve plumbing, based at its initial fibre and retaining its
two labelled vertices and complex curve orientations.
In the explicitly constructed product, after the actual based
unit-logarithm contraction, the physical circle on the fixed-curve
normal coordinate has axis
\[
 \nu(s,\theta)=H(s,\theta)iH(s,\theta)^{-1},\qquad
 \deg\nu=+3,
\]
where the domain is the cylinder with each end circle collapsed,
oriented by \(ds\wedge d\theta\), and the target sphere has its
outward orientation. The closed fixed-curve normal-frame comparison has
factor degrees \(6,0\) in complex-curve-times-parameter orientation,
and \((-6,0)\) in the Hopf-sphere orientation of P3K.23.
The full-unit circle action is displayed in
the proof. This product is not delta-equivariant. No circle-trace
correction, common-exterior variation matrix, nonfixed relative
integer or additional fan return is assigned here.
\end{theorem}
\begin{proof}
\begin{enumerate}[label=\arabic*.]
\item \par

\medskip\noindent\textbf{OP.1. One actual vertex product and all three normal transition loops}

\par\nobreak

\item Use PO.1 of P3K.47 and its ordered charts \((a,b,c)\), \((u,v,w)\).
Keep the first vertex coordinates and use at the second vertex the
actual unit map
\[
                  (u,v,w)\longmapsto(u,\eta v,\kappa w).
                                                               \tag{OP.1}
\]
Here the units are first restricted to the curves, hence to \(q=0\).
The higher terms of their original chart maps are restored below.
The reference vertex at \(\rho_*\) uses the same map with its initial
constant units. This makes (OP.1) an actual based vertex comparison,
rather than a replacement of the initial coordinate system.

\item The tangent coordinates on each pair of curve charts now have literal
reciprocal transition. With the following ordered complex normal
coordinates the three transitions are
\[
 \begin{array}{c|c|c|c}
  &z&\text{first normal pair}&\text{second normal pair}\\ \hline
 C_0&b&(a,c)&(\eta v,u)\\
 C_1&c&(a,b)&(\kappa w,u)\\
 C_2&a&(b,c)&(\eta v,\kappa w).
 \end{array}
\]
Their normal transition matrices are \(zD_i(\rho)\), where
\[
 D_0=\operatorname{diag}(\eta/\alpha,\kappa\alpha),\quad
 D_1=\operatorname{diag}(\kappa/g,g\eta),\quad
 D_2=\operatorname{diag}(\eta,\kappa).                 \tag{OP.2}
\]
These are direct substitutions in the three rows of PO.1. Put
\(K_i=D_i(\rho_*)^{-1}D_i(\rho)\). In particular
\[
 \begin{aligned}
 K_0&=\operatorname{diag}
   ((\eta/\eta_*)(\alpha_*/\alpha),
             (\kappa/\kappa_*)(\alpha/\alpha_*)),\\
 K_1&=\operatorname{diag}
   ((\kappa/\kappa_*)(g_*/g),(g/g_*)(\eta/\eta_*)),\\
 K_2&=\operatorname{diag}(\eta/\eta_*,\kappa/\kappa_*).
 \end{aligned}                                               \tag{OP.3}
\]
Each determinant has winding \(-6\), and each \(K_i\) is the identity
at the initial parameter. Thus their real oriented normal loops have
even parity. We will not rely on parity alone: OP.2 gives their
explicit contractions.

\item Write the rho circle as \(\rho=\rho_*e^{i\theta}\). Each unit of
winding zero in (OP.3) has its actual periodic logarithm, chosen zero
at \(\theta=0\). For \(\kappa/\kappa_*\), retain the factor
\(e^{-6i\theta}\), and take the periodic logarithm of the remaining
unit. Multiplication by the exponentials of scalar multiples of these
logarithms supplies exact based homotopies of the unit factors. After
these carried unit maps, the remaining real normal loops are
\[
 K_0^{\rm ang}=K_2^{\rm ang}
        =\operatorname{diag}(1,e^{-6i\theta}),\qquad
 K_1^{\rm ang}=\operatorname{diag}(e^{-6i\theta},1).
                                                               \tag{OP.4}
\]
The same logarithms can be inserted in the final maps, as explained
in OP.3; the original units are not discarded.

\par

\medskip\noindent\textbf{OP.2. Explicit based quaternion contractions in the actual order}

\par\nobreak

\item For an ordered pair of complex normal coordinates \((n_0,n_1)\),
we use exactly the K23 convention
\[
                         Q=n_0+j\overline{n_1}.       \tag{OP.5}
\]
This real coordinate has the original complex normal orientation.
Choose a smooth function \(\vartheta:[0,1]\to[0,\pi/2]\), constant
near both ends and strictly increasing between them, and put
\[
 \begin{aligned}
 p(s,\theta)&=\sin\vartheta(s)\,e^{-3i\theta}
                                      +j\cos\vartheta(s),\\
 H(s,\theta)&=p(s,\theta)p(s,0)^{-1}.
 \end{aligned}                                               \tag{OP.6}
\]
Both factors are unit quaternions. The maps are smooth and periodic,
and satisfy
\[
 H(0,\theta)=1,\qquad H(1,\theta)=e^{-3i\theta},\qquad
 H(s,0)=H(s,2\pi)=1.                                  \tag{OP.7}
\]
Define two actual oriented real normal maps
\[
    T_-(s,\theta)(Q)=H QH^{-1},\qquad
    T_+(s,\theta)(Q)=H QH.                            \tag{OP.8}
\]
Their inverses are displayed by replacing the two factors with their
inverses in reverse order. They are orthogonal and orientation
preserving. At \(s=1\), the first sends the quaternion coefficients
\((n_0,\overline{n_1})\) to
\((n_0,e^{6i\theta}\overline{n_1})\). Thus in the original complex
coordinates it is
\((n_0,n_1)\mapsto(n_0,e^{-6i\theta}n_1)\).
The second is \((n_0,n_1)\mapsto(e^{-6i\theta}n_0,n_1)\).
These are precisely the two loops in (OP.4). In particular no
unconjugated quaternion coordinate is substituted for (OP.5).

\item We use \(T_-\) for \(C_0,C_2\) and \(T_+\) for \(C_1\).
Multiplying these normal maps by the diagonal exponentials of the
actual periodic unit logarithms, with a cutoff in \(s\), gives a
contraction \(\mathcal K_i(s,\theta)\) of each full loop (OP.3).
The cutoff is zero near \(s=0\) and one near \(s=1\), so
\[
 \mathcal K_i(0,\theta)=I,\quad
 \mathcal K_i(1,\theta)=K_i(\theta),\quad
 \mathcal K_i(s,0)=\mathcal K_i(s,2\pi)=I.             \tag{OP.9}
\]
Write these chosen contractions as
\(\mathcal K_i=\mathcal D_iT_i\), where \(\mathcal D_i(s,\theta)\)
is the actual diagonal exponential of the periodic unit logarithms,
and \(T_i=T_-\) for \(i=0,2\), \(T_i=T_+\) for \(i=1\).
All maps are invertible throughout. The diagonal unit homotopies
commute with the physical circle; the explicitly displayed maps
\(T_\pm\) need not do so.

\par

\medskip\noindent\textbf{OP.3. Gluing the actual ordinary plumbing product}

\par\nobreak

\item Form the marked regular neighbourhood from two small actual coordinate
balls at the vertices and three disjoint normal tubes along the curve
interiors. On each curve choose a radial coordinate \(s=s(|z|)\)
which is zero near its first vertex and one near its second. In the
first normal frame send
\[
                    (z,n)\longmapsto
                       (z,\mathcal K_i(s(|z|),\theta)n).
                                                               \tag{OP.10}
\]
It is the identity near the first vertex. Near the second it is
\(K_i\), and the equality
\[
               zD_i(\rho_*)K_i=zD_i(\rho)            \tag{OP.11}
\]
shows that it is exactly the second vertex comparison (OP.1).
Consequently these three maps and the two actual vertex maps agree
on their common collars. There is no unresolved corner: the open
curve tubes are disjoint away from the vertex balls, and (OP.11)
holds in the ordered normal directions of each incident axis.

\item The maps on the actual coordinate germs have the same first normal
derivatives as these displayed bundle maps. The remaining higher
terms are removed on smaller tube collars by normal rescaling.
For a map \(F\) preserving the zero section, put
\(R_t(z,n)=(z,tn)\) and use the literal conjugation
\(R_t^{-1}FR_t\) for \(t>0\). Its limit consists of the given
zero-section map and its invertible normal derivative. Only the
normal vector is scaled in this construction. On the compact parameter circle
one common sufficiently small tube makes this a smooth isotopy of
invertible germs through \(t=0\). On each vertex collar the original
unit maps have nonzero diagonal derivative and are treated by the
same rescaling, with the vertex coordinate balls protected.

\item For precision about gluing this isotopy, take its velocity
\((\partial_tF_t)F_t^{-1}\), multiply it by a smooth collar cutoff,
and integrate. This supplies the actual collar diffeomorphism
between the original germ and the displayed normal map. The cutoffs
are chosen on the three disjoint tube interiors and agree with the
already fixed vertex maps on their protected collars. Thus they
glue without changing any other incident axis. All isotopies are
periodic and are the identity at the reference parameter.

\item Finally transport the rounded boundary and use the resulting collar
isotopy to fix one regular-neighbourhood radius. This constructs a
periodic smooth product of the actual labelled plumbing with the
reference plumbing. The vertex maps are complex-unit maps, the
curve maps preserve their complex orientations, and the normal maps
are in the positive real group. Hence the product preserves the
six-manifold orientation. It is an ordinary smooth product of the
labelled neighbourhood; no source-height-preserving or
delta-equivariant product is asserted.

\par

\medskip\noindent\textbf{OP.4. The physical circle and the exact degree-three axis}

\par\nobreak

\item We now retain the physical circle in this product. After the actual
extending unit maps of OP.1, on the fixed-curve normal quaternion
(OP.5) it is \(Q\mapsto Qe^{-i\phi}\). Conjugating by the actual
map \(T_-\) gives
\[
       Q\longmapsto Q\,H(s,\theta)e^{-i\phi}H(s,\theta)^{-1}.
                                                               \tag{OP.12}
\]
The same right-hand axis occurs after multiplication by the fixed
source Hopf representative used to form \(q_H\). Thus the circle
axis in the ordinary product is
\[
                         \nu(s,\theta)=H iH^{-1}.     \tag{OP.13}
\]
It equals \(i\) on both ends of the cylinder, but it is not the
constant axis on its interior. The zero-winding diagonal unit maps
are the previously displayed transported maps; contracting their
logarithms carries (OP.12) to the formula written here and leaves
the following integer degree unchanged.

\item Orient the cylinder by \(ds\wedge d\theta\), collapse each boundary
circle to its endpoint, and orient the target sphere of imaginary
unit quaternions by its outward normal. We compute the degree rather
than infer it from parity. Put \(A=\sin\vartheta\),
\(B=\cos\vartheta\), \(C=\cos3\theta\), \(S=\sin3\theta\).
Then (OP.6) is literally
\[
 H=(B^2+A^2C)-A^2S\,i+AB(1-C)\,j+AB S\,k.            \tag{OP.14}
\]
The equation \(HiH^{-1}=-i\) is equivalent to zero real and
\(i\)-coefficients of \(H\). Its only solutions are
\[
 \vartheta=\pi/4,\qquad
 \theta=(\pi+2\pi j)/3,
                   \quad j=0,1,2.                    \tag{OP.15}
\]
At each one \(H=j\), and direct differentiation gives
\[
                    \partial_s\nu=4\vartheta'(s)k,
                    \partial_\theta\nu=3j.           \tag{OP.16}
\]
The ordered basis \((k,j)\) is positive at the target point \(-i\),
since \((-i,k,j)\) has the ambient orientation \((i,j,k)\).
Every preimage in (OP.15) is therefore regular and positive. Hence
\[
                              \deg\nu=+3.             \tag{OP.17}
\]

\item For completeness the two nonfixed normal actions are also explicit.
In their ordered normal coordinates of OP.1, the initial actions are
\[
 R_1(\phi)Q=e^{-i\phi/2}Qe^{-i\phi/2},\qquad
 R_2(\phi)Q=e^{i\phi/2}Qe^{-i\phi/2}.
\]
Their base coordinates have weights \(+1,-1\), respectively.
The ordinary product carries their normal actions to
\[
 \begin{aligned}
 C_1:\quad Q&\longmapsto
 (He^{-i\phi/2}H^{-1})Q(H^{-1}e^{-i\phi/2}H),\\
 C_2:\quad Q&\longmapsto
 (He^{i\phi/2}H^{-1})Q(He^{-i\phi/2}H^{-1}).
 \end{aligned}                                               \tag{OP.18}
\]
The paired central signs make both formulas close at \(\phi=2\pi\).
For the full original units, the physical normal action on each tube is
exactly
\[
 \mathcal D_i(s,\theta)
      [T_i(s,\theta)R_i(\phi)T_i(s,\theta)^{-1}]
                          \mathcal D_i(s,\theta)^{-1},
                                                               \tag{OP.19}
\]
with \(R_0(\phi)Q=Qe^{-i\phi}\). Thus the written \(H\)-axis and
(OP.18) are reached from the full-unit action by the displayed based
logarithm contraction, with the real normal metric transported along
that contraction. Together with their base rotations, (OP.19) gives
the actual transported circle actions on all three tubes.
They agree with the vertex action on the protected collars.

\item The product in OP.3 thus provides a possible ordinary common boundary
marking, with the nonconstant circle action (OP.12)--(OP.18). The
circle-trace and stable-derivative arguments for a constant equivariant
product require a further correction in this marking. That correction
is not evaluated here. In particular neither a rank-one initial
variation matrix nor vanishing nonfixed relative derivatives follows
from (OP.17).

\par

\medskip\noindent\textbf{OP.5. The closed fixed-curve frame comparison}

\par\nobreak

\item We also compute the actual normal-frame comparison on the fixed
curve. Write \(z=\tan\chi\,e^{i\psi}\), with
\(0\leq\chi\leq\pi/2\), and put
\[
 U(z)=\frac{1+jz}{\sqrt{1+|z|^2}},\qquad q_H=U(z)Q.
\]
The ordinary product of OP.2 and the physical Cox product of P3K.42
are therefore related by the literal arrow
\[
 q_H\longmapsto L(z,\theta)q_H R(z,\theta),\qquad
 L=UHU^{-1},\quad R=H^{-1}.                         \tag{OP.20}
\]
Here the radial parameter in \(H\) is the previously chosen smooth
function of \(|z|\). Both maps close on \(\mathbb P^1\times S^1\).
At the first pole \(H=1\). Near the second pole \(H=e^{-3i\theta}\),
so \(L=\cos3\theta-(UiU^{-1})\sin3\theta\) extends in the smooth
Hopf axis; its limiting value is \(e^{3i\theta}\). The right factor
extends independently of the source angle. The actual zero-winding
diagonal unit maps in OP.19 are carried in this comparison and then
contracted by their displayed based logarithms. This is a homotopy
of the actual closed frame arrow, rather than a discarded unit.

\item We orient the domain by its complex curve orientation followed by
positive \(\theta\), hence by \(d\chi\wedge d\psi\wedge d\theta\)
on the cylinder. In the quaternion orientation \((1,i,j,k)\), use
the matrix trace with \(\operatorname{tr}(ijk)=-2\), and the positive
normalized form
\[
 \omega(g)=-\frac{\operatorname{tr}(g^{-1}dg)^3}{24\pi^2}.
\]
The Maurer--Cartan equations give the exact product identity
\[
 \omega(gh)=\omega(g)+\omega(h)+
 \frac{1}{8\pi^2}d\operatorname{tr}
       (g^{-1}dg\wedge dh\,h^{-1}).                  \tag{OP.21}
\]
Applying it twice to \(UHU^{-1}\), the two \(U\) cubic terms cancel
and the \(H\) term vanishes because \(H\) depends on only two
variables. Thus its integral is the boundary integral of
\[
 \frac{1}{8\pi^2}\left\{
 \operatorname{tr}(U^{-1}dU\wedge dH\,H^{-1})+
 \operatorname{tr}((UH)^{-1}d(UH)\wedge dU^{-1}U)\right\}.
                                                               \tag{OP.22}
\]
There is no first-pole contribution. At the second pole,
\(U=je^{i\psi}\), \(H=e^{-3i\theta}\), and
\[
 U^{-1}dU=i\,d\psi,\qquad dH\,H^{-1}=-3i\,d\theta.
\]
Each trace in braces in (OP.22) is
\(6\,d\psi\wedge d\theta\). The outward boundary orientation there
is \(d\psi\wedge d\theta\), and consequently
\[
 \deg L=\frac{12(2\pi)^2}{8\pi^2}=6,\qquad
 \deg R=0.                                            \tag{OP.23}
\]
The second equality follows directly from its two-variable
factorization. The Hopf coordinate \(h=UiU^{-1}\) reverses the
complex orientation of the source curve: near \(z=0\), its two
derivatives in \((\Re z,\Im z)\) are \((-2k,-2j)\), negatively
oriented at \(i\). Therefore the same closed comparison has
factor degrees \((-6,0)\) in the \(h\)-sphere-times-parameter
orientation of P3K.23, and \((6,0)\) in complex-curve orientation.
We have computed these closed factor maps and their degrees. The
connecting-homomorphism sign and the resulting relative-return
correction remain unevaluated; no additional fan return follows
without those arrows.

\item The exact normal-matrix, quaternion and degree checks are in
\path{code/part3/source_cusp_ordinary_product.py}.

\end{enumerate}
\end{proof}




\begin{theorem}\label{night:P3K.49}
\StatementRecord{night:P3K.49}{proved}{S6\_06\_BLIND: GPT-6 Astra (OpenAI, ChatGPT Pro mode, effort Pro), recorded re-derivation in the checkpoint dated 18 November 2026}{\cref{night:P3F.37}, \cref{night:P3F.41}}

Retain the marked E7 source, Cartier line
\(M_a=\mathcal O(aG-O)\), suspension by \(-bG\), and one
III-star cap, for \((a,b)=(1,2)\) with the actual good root
\(m=4p\), or \((a,b)=(2,1)\) with the actual reduced root
\(m=4p-1\). Keep the signed norm \(n=\pm1\), every prescribed
finite character and annular \(d\delta\), the chosen whole-zero
\(I_1\) cusp fan and the literal scalar. Use the source arrow
\[
 I_u:(x,y,t,u)\longmapsto(u^2x,u^3y,u^4t,u^3),\qquad v=u^3.
\]
Complete the previously unselected new reference annular map
\(F_0\) by continuation in its actual Hom/source-period torsor
of the full old overlap at one nonzero \(u_*\), including all
of its integral and finite data.

Over a smaller punctured \(u\)-annulus carrying the chosen
positive loop, there is an actual holomorphic,
physical-circle-equivariant family isomorphism from the old
family to the \(v=u^3\) pullback of this newly marked
\(v\)-family. It retains the root endpoint \(D^{-4}\), the
reduced doubled period, all root and Cartier line maps, and
the full annular cochains. For the actual scalar rescaling
\(\lambda=v^N\lambda_0\), \(N=km\), it compares exactly
\(\lambda=u^{3N}\lambda_0\) on the old side. Their physical
returns and oriented mapping tori therefore obey
\[
 [\Phi_{\rm old,u}]=3[\Phi_{\rm new,v}]\quad\text{in }\Theta_7,
 \qquad M_{\Phi_{\rm old,u}}\cong M_{\Phi_{\rm new,v}^{\,3}}.
\]
This statement completes the punctured-family comparison. It
does not identify the raw collision E-chain with the finite-cap
chain across \(v=0\), assert the still-unassembled global
central attachment, or assign a numerical source-zero return.

\end{theorem}
\begin{proof}
\begin{enumerate}[label=\arabic*.]
\item \par

\medskip\noindent\textbf{RCZ.1. The actual common source, line and suspension}

\par\nobreak

\item The coordinate arrow is
\[
 I_u:(x,y,t,u)\longmapsto
 (X,Y,T,v)=(u^2x,u^3y,u^4t,u^3).
\]
It carries G and O to their marked counterparts and carries the
canonical Cartier line \(M_a=\mathcal O(aG-O)\), for
\((a,b)=(1,2),(2,1)\), by its actual divisor map. The primitive
translation lift and selected-base coordinate satisfy
\[
 I_uU_{a,\mathrm{old}}I_u^{-1}=u^{2a}U_{a,\mathrm{new}},
 \qquad q_{\mathrm{old}}=u^{-4}q_{\mathrm{new}}.
\]
Since \(ab=2\), these factors cancel in the actual suspension
\(2(-1)^{b+1}\lambda qU_a^b\). The scalar is therefore the same
literal scalar, not a scalar with an omitted u-power.

\item The cusp Tate period and marked section point agree under \(I_u\).
Its physical-height unit is retained as
\(q_{\mathrm{old}}/q_{\mathrm{new}}=u^{-4}\). On the u-annulus
this is an invertible integer-power toric multiplier. It preserves
the chosen fan and physical delta character. Its parameter winding
is retained; no global logarithm of u is chosen here.

\item For a cap of order m, the root arrow on the logarithmic u cover is
\[
 z_{\mathrm{old}}=u^{4/m}z_{\mathrm{new}}.
\]
For the good root \(m=4p\), the minimal coordinates satisfy
\(X'_{\mathrm{new}}=x'_{\mathrm{old}}\) and
\(Y'_{\mathrm{new}}=y'_{\mathrm{old}}\). For the reduced root
\(m=4p-1\), they satisfy
\[
 X'_{\mathrm{new}}=u^{-2/m}x'_{\mathrm{old}},\qquad
 Y'_{\mathrm{new}}=u^{-3/m}y'_{\mathrm{old}}.
\]
These maps carry the actual section to
\((0,vz_{\mathrm{new}}^{3p})\), preserve the specified resolution
ideals, and commute with the root deck D. Their full positive
u-loop endpoint is \(D_{\mathrm{new}}^{-4}\). In the reduced
coordinates this is exactly
\[
 (z,X',Y')\longmapsto
 (\zeta_m^{-4}z,\zeta_m^{-2}X',\zeta_m^{-3}Y'),
\]
because \(4p=m+1\). This is the root endpoint; it is not silently
replaced by the identity on the affine cap.

\par

\medskip\noindent\textbf{RCZ.2. The extending fourth translation and its retained period}

\par\nobreak

\item Let R be the actual division section \(mR=Q\), with the branch
specified in F41 or F37. The signed affine cap is represented by
\(A=L D\), where L covers \(-nR\), and its scalar normalization
is the one giving \(A^m=B^n\). We retain the finite eigenspace
construction used for L.

\item Its central commutator character is one in both actual branches.
In the good branch R specializes to O, so L on the central compact
elliptic fibre is a scalar line automorphism. Such a scalar commutes
with the actual D lift. In the reduced branch the line is
\(N=\mathcal O(2G'-O)\). The doubling formula puts 2G' at O on
the reduced central fibre. Both meet its smooth locus, so the
restriction of this Cartier line is literally the trivial line;
its component degrees are also zero. In a trivialization of that
line, the central line multipliers of D and L are holomorphic
units on a compact connected reduced curve, hence constants.
Their source maps commute because D fixes R. Their line maps
therefore commute as well. Thus the finite eigencharacter is one.

\item The actual finite eigenspace lift can consequently be chosen with
\(DLD^{-1}=L\). Its norm normalization is a base-invariant unit
and preserves this equality. This is the exact normalized lift
used in the cap; any initial unit discrepancy is removed by its
finite averaging conjugation, with that conjugation retained on
the annulus. In these charts
\[
                         A^4D^{-4}=L^4.
\]
The map on the right is an actual extending line map. In the good
branch its central source translation is zero. In the reduced
branch R specializes to the nonidentity two-torsion point, so
\(-4nR\) also specializes to zero. Hence \(L^4\) restricts to a
nonzero scalar on the central fibre. On a simply connected
parameter chart divide by that scalar, take its scalar logarithm,
and use the inverse of the exponential near the identity of the
extending translation group. This produces a small extending Lie
logarithm of \(L^4\). It does not identify that logarithm with four
times an annular logarithm of L.

\item In particular, on the marked good fourth root of the reduced
branch, write its square lattice as \(\mathbb Z+i\mathbb Z\),
with \(R(0)=(1+i)/2\). A logarithm tending to this half-period
satisfies
\[
             4\log R-\log(4R)=2(1+i),
\]
where the second logarithm is the small one. With signed n the
source part is multiplied by n.~The full line-period vector is
the actual difference of the two lifted line logarithms; it
includes the Cartier factor of automorphy. It is retained as an
integral annular cochain. No parameter-dependent multiple of
this source period is asserted to extend across the reduced cap.

\item The exact finite-unit conjugation has a useful elementary form.
In the input-scalar convention, let an actual cap map be
\(A'=\beta A\), where \(\prod_{j=0}^{m-1}\beta(D^jz)=1\)
and its central character has already been matched. Put
\(g=\log\beta\), with central value zero, and
\[
 f(z)=-\frac1m\sum_{j=0}^{m-1}j\,g(D^jz).
\]
Then \(f(Dz)-f(z)=-g(z)\), and multiplication by \(e^f\)
conjugates A' to A. The maps and logarithms depend holomorphically
on the parameter on each chosen parameter chart. The finite
character is matched before this zero-central logarithm is taken.

\par

\medskip\noindent\textbf{RCZ.3. The relative Lie solve and the full groupoid}

\par\nobreak

\item Over a simply connected small parameter chart in the u-annulus,
the actual extending translation Lie bundle has
\[
 0\longrightarrow\mathcal O\delta\longrightarrow V
       \longrightarrow\mathcal O(-1)\longrightarrow0
\]
on the coarse source base. This is the relative version of CP.1.
Indeed the root substitution sends the source dualizing vector
to \(z^k\) times the root vector, with its exact deck character
\(\zeta_m^k\), and \(0\le k<m\). Its invariant extension is
therefore the original source Lie line. The Weierstrass invariant
differential and \(I_u\) supply the actual relative line frames;
their extra u-powers are units on this parameter chart.

\item The extension splits by the ordinary Cauchy/Laurent solution for
\(\mathcal O(1)\). The same fixed-circle integrals solve every
V-valued base one-cocycle for \(\mathcal O\) and
\(\mathcal O(-1)\). These integrals preserve holomorphic
parameter dependence. Finite averaging applies to these complex
vector cochains. Thus the relative first coherent cohomology of
V over the source base is zero, with an explicit parameter-dependent
solution operator.

\item The exponential sequence is used on the entire root atlas and
its groupoid,
\[
 0\longrightarrow H\longrightarrow V\longrightarrow A
       \longrightarrow0.
\]
At a reduced root, the actual suspension-power removal of its
component translation is made before taking the identity-component
logarithm, exactly as in CP.3. Lie exponentiation alone is not used
to erase a disconnected component.
Atlas logarithms are not presumed invariant. Their ordinary
overlap discrepancies and all finite cyclic-group discrepancies
form a total integral two-cocycle. Its finite terms retain the
literal equation \(A^m=B^n\), the central line character and
the actual period lift. The doubled period in RCZ.2 is carried
by that total cochain whenever parameter charts are changed.
Only after the two actual total two-cocycles have been compared
is the remaining V-valued one-cocycle solved by the preceding
operator. No assertion about surjectivity of exponentiation on
coarse invariant sections is used.

\item All resulting local exponentials are the extending source/line
maps of CP.3. At the cusp their two characters become bounded
holomorphic units after the actual integral component reindexing;
at every finite root they are the line lifts on the actual
resolved source, with the stated finite character. These local
maps therefore transport the annular and Cox normal charts as
actual maps.

\par

\medskip\noindent\textbf{RCZ.4. The parameter scalar group and explicit analytic factorization}

\par\nobreak

\item On the punctured source-parameter annulus every source has the
same marked Mordell--Weil group \(\mathbb ZG\). The endpoint
calculation CP.3--4 applies relative to this annulus. In particular
an actual global extending translation lifts off the selected cusp,
its source part is \(jG\) with j locally constant, and its line
part is a rational base function times the primitive lift \(U_G^j\).
The cap orders and root characters give its exact valuations.
The selected cusp uses its actual integer reindexing. Taking the
degree of that rational divisor gives the same equality
\(j=bk\) as CP.4, and the norm ratio identifies the full map
with a physical scalar. Evaluation at one ordinary marked origin
shows that this scalar is holomorphic in the source parameter.
Conversely every such holomorphic nonzero scalar extends through
all the actual charts. The relative global translation group is
therefore the physical \(\mathcal O^*\), with no omitted finite
kernel.

\item Here is the explicit parameter factorization used for locally
constructed family comparisons. Cover a slightly larger round
u-annulus by two angular charts U\_+,U\_-, with simply connected
overlap components. Write the actual scalar transition as c and
choose a holomorphic logarithm \(\ell\) on each overlap
component. Choose a smooth angular partition \(\chi_++\chi_-=1\)
subordinate to these charts, and put
\[
 b_+=\chi_-\ell,\qquad b_-=-\chi_+\ell.
\]
These are extended by zero away from the respective overlap
supports. They satisfy \(b_+-b_-=\ell\), and
\(\bar\partial b_+=\bar\partial b_-\) is a globally defined
smooth (0,1)-form \(\eta=f(u)d\bar u\). Choose a radial cutoff
equal to one on a smaller closed annulus carrying the loop and
compactly supported in the larger annulus. Extend the cut-off
coefficient by zero to the plane and set
\[
 v(u)=\frac1\pi\int_{\mathbb C}
             \frac{\widetilde f(w)}{u-w}\,dA(w).
\]
The distribution identity
\(\bar\partial_u(1/(\pi(u-w)))=\delta_w\) gives
\(\bar\partial v=\eta\) on the smaller annulus. Thus
\[
 h_+=e^{b_+-v},\qquad h_-=e^{b_--v}
\]
are nonvanishing holomorphic functions with \(h_+/h_-=c\).
When the local comparison maps have the actual arrow
\(I_+=cI_-\), the maps \(h_+^{-1}I_+\) and
\(h_-^{-1}I_-\) glue literally. All multipliers are physical
scalars; their action on the line and normal coordinates is the
specified one. This gives the analytic factorization, including
both overlap components and any winding of c.

\item For the actual scalar rescaling \(\lambda=v^N\lambda_0\),
write \(N=km\), \(v=u^3\), and
\(t=\log u/(2\pi i)\). The old and new ordinary forward decks
are \(q_{\mathrm{old}}^{3kn}\) and
\(q_{\mathrm{new}}^{3kn}\); their ratio is
\(u^{-12kn}\). The literal compensating scalar cochain is
\[
 c_0(t)=\exp\bigl(2\pi i\,6kn\,t(t-1)\bigr),
 \qquad \frac{c_0(t+1)}{c_0(t)}=u^{12kn}.
\]
This is retained in the forward endpoint square. It is not
declared to be a single-valued multiplier by itself, or to extend
over \(u=0\) in the old coordinates. Any remaining scalar transitions
are treated by the factorization just given.

\par

\medskip\noindent\textbf{RCZ.5. Completing the previously unselected annular marking}

\par\nobreak

\item Section 21 of \texttt{part3\_source\_zero\_completion\_working.md} left \(F_0\)
specified as a holomorphic actual reference-scalar overlap, without
choosing its absolute integral annular marking. We now make that
choice. Fix a nonzero \(u_*\) and put \(v_*=u_*^3\). Transport the \textbf{entire}
old ordinary-to-cap map at \(u_*\) through the canonical source/Cartier
arrow \(I_{u_*}\), including its chosen root cut, signed norm, prescribed
\(d\delta\), and every period in its lifted logarithm. Call this
actual map \(F^*\). The new \(F_0\) is constructed below to restrict exactly
to \(F^*\). This is completion of the previously unselected data; it is
not a proof that an unrecorded different choice had the same integer.
The new analytic norm root uses the transported old value at this
reference point and the same prescribed finite character. Its
nonvanishing norm unit extends on the root polydisc, so this branch
continues analytically there. The finite unit conjugation in RCZ.2
is included before \(F^*\) is regarded as a map into that new cap.

\item Choose a fixed cap-overlap annulus C in the T-base sufficiently far
from the colliding finite fibres. All source fibres above C remain
smooth elliptic curves for v in a full sufficiently small disc D,
including \(v=0\). The root source, its actual corrected Cartier line,
its division branch and its normalized \(A_0\) extend on this product by
the displayed good/reduced equations and resolved line diagrams.
At a fixed sufficiently small reference scalar, the ordinary and
cap norm quotients are therefore smooth proper families of compact
complex tori over \(C\times D\).

\item There are literal local comparison maps on this whole product.
Cover C by two angular charts \(C_+,C_-\). Choose a root branch \(z_i(q)\)
on each. Restriction to that chosen sheet gives a map from the
ordinary pullback line quotient to the cap quotient. It has no
residual quotient: an element \(A^jB^h\) fixing that sheet has
\(j=mk\), and is exactly \(B^{nk+h}\). Thus the fibre on the
chosen sheet is precisely \(N^\times/\langle B\rangle\).
The source/Cartier identification supplies the claimed local map.
On overlap, its map between sheets is the literal indicated power
of \(A_0\), including its line character and the relation \(A_0^m=B_0^n\).
This description continues through \(v=0\) on C; no singular source
fibre occurs in this overlap annulus.

\item The two local comparisons differ by translations in the actual
Jacobian family J over \(C\times D\). Its Lie bundle V is holomorphically
trivializable there by the source dualizing vector, physical delta,
and the same additive Cauchy splitting of their extension. One can
also use the literal constant-factor-of-automorphy presentation of
the degree-zero elliptic line on each simply connected angular
chart. Its periods and marked Picard point are holomorphic in
(q,v), including at \(v=0\). These are frames of the actual line
translation group, not an identification with an unrelated torus.

\item Correct each local comparison on \(C_i\times\{v_*\}\) to equal \(F^*\). The required
translation has a holomorphic logarithm on the simply connected
\(C_i\). Extend its coefficients in the just specified Lie frame
holomorphically in v, and exponentiate. This makes local actual
maps \(E_i\) on \(C_i\times D\) satisfying \(E_i|_{v_*}=F^*\).

\item Their overlap translation is the identity at \(v_*\). It therefore has
a unique holomorphic Lie logarithm \(\ell_{ij}(q,v)\) whose value
at \(v_*\) is zero. To construct it, lift along the simply connected
v-disc through the exponential covering of the smooth torus
family, starting at zero; uniqueness gives holomorphic dependence
on q as well. The overlap components are simply connected. Any
ordinary or finite-group discrepancy of these logarithms is an
integral period and is zero at \(v_*\). It remains zero throughout D.
The finite root groupoid has not been forgotten: its free-sheet
description above carries every \(A_0\) power and its norm to these
ordinary overlap maps. In particular its bar-resolution terms are
carried by those exact maps before the logarithms are chosen.

\item Thus \(\ell_{ij}\) is a genuine V-valued one-cocycle, vanishing
at \(v_*\). Apply the q-variable Cauchy splitting with contours and
partition independent of v. The integrals preserve holomorphic
v-dependence and vanish at \(v_*\). Exponentiating the resulting
cochains glues the \(E_i\) to an actual holomorphic comparison
\[
             F_0:\mathcal Y|_{C\times D}
                         \longrightarrow\mathcal C|_{C\times D},
                   \qquad F_0|_{v_*}=F^*.
\]
This constructs continuation in the actual Hom/source-period
torsor. It does not merely transport an abstract lattice vector.
The integral T-annulus cocycle of this map is locally constant in
v in its transported period lattice. At \(v_*\) it is exactly that of
\(F^*\), so it retains the prescribed \(d\delta\) and full reduced doubled
period throughout D. Inserting an extra \(q^j\delta\) would
change \(F^*\); no such insertion occurs in this construction.
The physical delta part has the literal coordinate check
\(q_{\rm old}^d=u^{-4d}q_{\rm new}^d\). Its q-order is d
on both sides. The remaining u-unit is retained in the physical
scalar comparison; it is not a new annular integer.

\par

\medskip\noindent\textbf{RCZ.6. The chain extension and the common base-direction cocycle}

\par\nobreak

\item On a cut angular chart, the constructed \(F_0\) lifts to an actual
nonzero-line map. Lift the value of the marked origin through the
regular covering \(N^\times\to N^\times/\langle B_0\rangle\).
The cut annulus times D is simply connected, so this gives a
holomorphic source point and a nonzero line vector. The
degree-zero Hom-line argument of CP.3 then gives an invertible
line lift over the source fibre; its multiplier is holomorphic
and nonzero through \(v=0\). It commutes with \(B_0\) by the literal
constant-factor-of-automorphy calculation. On two such charts
the maps differ by the actual integral deck power, not by a
newly chosen pointwise identification.
Every such line map also commutes with physical scalar
multiplication. Hence its \(B_0\) commutation is the exact commutation
with the rescaled suspension \(v^N\lambda_0B_0\) as well;
no change of the source lift is made at that step.

\item Consequently these maps act in the finite-cap chain charts by
the actual source map and
\[
 (X_j,Y_j)\longmapsto
           (a(q,v)X_j,a(q,v)^{-1}Y_j),
\]
with a a nonzero holomorphic unit through \(v=0\). Their integral
deck ambiguities are retained as the corresponding chain-index
translations. In particular, for \(N=km\) the root A shifts
the index by kn, B by km, and the actual sheet relation remains
\(A^m=B^n\). The prescribed \(d\delta\) changes only the retained
single-valued line multiplier on the overlap; q is nonzero
there, so that multiplier also extends through \(v=0\).
Here a(q,v) denotes the multiplier in the chosen rigidified
constant-factor line coordinates of the actual degree-zero
line. The Cartier-to-rigidified coordinate maps are carried
with it; the formula does not identify every original Cartier
chart's multiplier with the same scalar function.

\item For the rescaled scalar, the exact Section 21 cochain is still
\(F=F_0+nk\,t_v s_\infty\delta\), where
\(t_v=\log v/(2\pi i)\). Its ordinary and cap origins obey
the already checked literal equality
\[
 F+o_{\rm ord}-o_{\rm cap}=F_0.
\]
A positive root-cut continuation multiplies the new line
coordinate by \(v^{nk}\), exactly its chain shift kn, with
the same \(A_0\) source and line map. The parameter deck is the
retained pair \(q_1^{nk},h^{nk}\), and
\(F_{t_v+1}R_{\rm ord}=R_{\rm cap}F_{t_v}\) is unchanged.
Thus the new choice of \(F_0\) extends in the existing chain atlas
and commutes with its chosen root, Cartier and norm maps. No
fractional power of v is inserted in a single central chart.

\item The ordinary origin itself is not a single-valued central
coordinate. With the exact convention
\(\eta=\zeta-o_{\rm ord}\) and
\(o_{\rm ord}=nk\,t_v s_1\delta\), positive continuation
of q1 at fixed old coordinate changes
\[
 \eta\longmapsto\eta-nk\,t_v\delta,
 \qquad V_\eta\longmapsto v^{-nk}V_\eta.
\]
Its inverse gluing arrow multiplies by \(v^{nk}\). Positive
q-infinity continuation has \(s_1\mapsto s_1-1\) and hence
the opposite sign, giving the \(v^{nk}\) cap-cut factor above.
These are the actual two arrows. Their compatibility with the
literal h,g collision-fan atlas is a separate remaining assembly
calculation. This note does not use the ordinary origin to
identify central E-chain components in that fan.

\item Now compare the old u-family with the pullback of this marked
v-family on the u-annulus. At \(u_*\) the source, line, cap and cusp
maps agree, with the reference annular arrow just constructed.
All their integral base-direction total two-cocycles therefore
agree at that point. Under the actual \(I_u\) identification of
period lattices, their difference is a locally constant class:
choose atlas logarithms on a small simply connected parameter
chart and retain each ordinary, finite-group and cusp period
discrepancy. These discrepancies are integral and continuous.
The explicit root maps, line-unit conjugations and toric unit
maps in RCZ.1--2 identify the local models in every such chart.
The zero difference at \(u_*\) consequently stays zero under their
actual parameter continuation around the connected annulus.
In particular the doubled period is carried as an integral
one-cochain on changing parameter charts, not discarded as an
extending Lie logarithm.

\item This establishes the base-direction discrete equality required
by RCZ.3 for the newly completed marking. The relative coherent
solution in that section therefore gives actual local family
isomorphisms on each simply connected parameter chart, with all
root/cusp and annular maps retained.

\par

\medskip\noindent\textbf{RCZ.7. The actual global family arrow}

\par\nobreak

\item The ratio of two such local comparisons is a global extending
translation of the same marked compactified family. RCZ.4
identifies it with a holomorphic physical scalar, including
every finite-kernel endpoint. Apply its explicit two-angular-chart
factorization on a smaller u-annulus carrying the chosen loop.
The corrected local maps glue to a literal holomorphic
delta-equivariant family isomorphism from the old u-family to
the pullback under \(v=u^3\) of the constructed v-family, with \(F_0\)
marked as in RCZ.5. For the rescaled scalar the quadratic
cochain of RCZ.4 is included in its forward endpoint square;
its one-step ratio is \(u^{12kn}\), the exact compensating ratio.

\item Every map used here is the actual source/Cartier coordinate
arrow, an extending translation on its stated root or toric
chart, or physical scalar multiplication. Their inverses have
the same extensions. Thus the map carries the normal circle
and the prescribed fan as well as the complex fibres. It is a
family morphism over the annulus. This argument assigns no
numerical return or extension of the old coordinates across
\(u=0\); those remain the separate collision/fixed-stratum tasks.

\item The construction deliberately fixes the previously unselected
\(F_0\) by the full old reference map. It makes no assertion that an
independently chosen different integral annular marking would
be equal to that one.

\par

\medskip\noindent\textbf{RCZ.8. The oriented mapping-torus consequence}

\par\nobreak

\item Use the actual rescaled new family
\(\lambda=v^N\lambda_0\), with \(N=km\), and its pullback
\(\lambda=u^{3N}\lambda_0\). The isomorphism just constructed
is a map of these actual families, including their closing
arrows. A positive u-circle maps to three positive v-circles.
Pull back a connection from the v-family and transport it by
the isomorphism. The resulting actual parallel transport around
u is the third iterate of that around v, conjugated by the
reference-fibre map. Consequently their oriented mapping tori
satisfy
\[
 M_{\Phi_{\rm old,u}}\cong M_{\Phi_{\rm new,v}^{\,3}},
 \qquad
 [\Phi_{\rm old,u}]=3[\Phi_{\rm new,v}]\quad\text{in }\Theta_7.
\]
Equivalently the left mapping torus is the degree-three
parameter pullback of the actual new mapping torus. The
physical return uses the inverse closing arrow on both sides;
the family isomorphism intertwines those arrows, so this
argument introduces no inverse-loop sign. It retains the
specified scalar rescaling and assigns no value to either
unknown class. No numerical class division is used.

\end{enumerate}
\end{proof}




\begin{theorem}\label{night:P3K.51}
\StatementRecord{night:P3K.51}{proved}{S6\_06\_BLIND: GPT-6 Astra (OpenAI, ChatGPT Pro mode, effort Pro), recorded re-derivation in the checkpoint dated 18 November 2026}{\cref{night:P3K.29}}
Retain the actual ordinary plumbing product of P3K.48 and the
positive selected-cusp source loop. In the positive complex normal
meridians \((m_0,m_1,m_2)\) and their integral relative dual cycles
\((D_0,D_1,D_2)\), its common-exterior variation is exactly
\[
 V=\begin{pmatrix}0&6&-6\\6&0&0\\-6&0&0\end{pmatrix}.
\]
The matrix is computed by the actual orbit-trace maps, generic
normal sections, vertex-link paths and signed relative chain
homotopy. No nonfixed stable derivative or further fan return is
assigned by this statement.
\end{theorem}
\begin{proof}
\begin{enumerate}[label=\arabic*.]
\item \par

\medskip\noindent\textbf{CT.1. Three literal trace maps}

\par\nobreak

\item On each curve use \(z=\tan\chi>0\), \(0\leq\chi\leq\pi/2\),
and set \(U=\cos\chi+j\sin\chi\). Use the actual normal sphere
coordinate \(q_i=UQ\) and the section \(q_i=1\). The based ordinary
normal product is \(T_-=HQH^{-1}\) for \(C_0,C_2\), and
\(T_+=HQH\) for \(C_1\), with \(H\) from OP.6. We first display
the angular maps; the full diagonal units are carried by OP.19 and
its based homotopy. The physical circle angle is \(\phi\), and
\(E=e^{i\phi/2}\).

\item On \(C_0\) the source point is fixed. On \(C_1,C_2\) it moves
by \(z\mapsto e^{i\phi}z,e^{-i\phi}z\), respectively. The exact
Hopf representatives satisfy
\[
 U(e^{i\phi}z)=E^{-1}UE,\qquad
 U(e^{-i\phi}z)=EUE^{-1}.                         \tag{CT.1}
\]
Substitute \(Q=U^{-1}\), undo \(T_i\), apply the original
physical normal action OP.18, and then redo \(T_i\) and the
source Hopf representative at the moved point. This gives
\[
 \begin{aligned}
 F_0&=H E^{-2}H^{-1},\\
 F_1&=E^{-1}UEH E^{-1}H^{-1}U^{-1}H^{-1}E^{-1}H,\\
 F_2&=EUE^{-1}H E H^{-1}U^{-1}H E^{-1}H^{-1}.
 \end{aligned}                                           \tag{CT.2}
\]
These are actual maps of the parameter/physical-orbit cylinders into
the corresponding normal \(S^3\). Each closes at \(\phi=2\pi\)
and \(\theta=2\pi\). At each radial endpoint it is independent
of \(\theta\). Thus, subtracting its initial-parameter trace gives
a closed relative degree-three cycle with a well-defined integer
normal-sphere degree. We use domain orientation
\(d\chi\wedge d\theta\wedge d\phi\).

\par

\medskip\noindent\textbf{CT.2. The three degrees}

\par\nobreak

\item Write
\[
 h=UiU^{-1},\qquad a=UHiH^{-1}U^{-1},\qquad
 v=HiH^{-1},\qquad w=H^{-1}iH.
\]
Each is constant in \(\theta\) at both radial ends. As maps from
the cylinder with its ends collapsed, their outward-sphere degrees
are
\[
             \deg h=0,\qquad \deg a=3,
             \deg v=3,\qquad \deg w=-3.             \tag{CT.3}
\]
The first is independent of \(\theta\), and the third is OP.17.
For the second, the exact homotopy
\(U_t=\cos(t\chi)+j\sin(t\chi)\) gives
\(U_tvU_t^{-1}\); it keeps each end independent of \(\theta\),
so it is a homotopy on the collapsed cylinder. For the fourth, the
same three inverse images of \(-i\) as OP.15 have \(H=j\).
Direct differentiation gives
\[
             \partial_s w=-4\vartheta'(s)k,
             \partial_\theta w=3j.
\]
They are all negative regular inverse images in the same orientation
as OP.16. This proves the last integer without assigning a sign by
analogy.

\item For an axis map \(n\) of degree \(d\), the map
\(\exp(\phi n/2)\) on this cylinder times \([0,2\pi]\) has
normalized third integral \(d\), while its inverse has integral
\(-d\). Indeed the positive quaternion volume is
\(\sin^2 t\,dt\wedge\Omega_{S^2}\); substituting
\(t=\phi/2\) gives
\(4\pi d\int_0^\pi\sin^2t\,dt/(2\pi^2)=d\).
A full-angle inverse \(\exp(-\phi n)\) gives \(-2d\).

\item The latter formula with axis \(v=HiH^{-1}\) immediately gives the degree of \(F_0\).
For the others, (CT.2) factors literally as
\[
 \begin{aligned}
 F_1&=E^{-1}\exp(\phi h/2)
                  \exp(-\phi a/2)\exp(-\phi w/2),\\
 F_2&=E\exp(-\phi h/2)
                  \exp(\phi a/2)\exp(-\phi v/2).
 \end{aligned}                                           \tag{CT.4}
\]
The Maurer--Cartan product formula OP.21 makes the integrals additive
here. At the radial ends every factor is independent of \(\theta\),
so the boundary two-forms vanish; at \(\phi=0,2\pi\) every factor
is a constant central sign, so they also vanish. The two parameter
faces cancel because the factors are periodic. Consequently
\[
               \deg F_0=-6,\qquad
               \deg F_1=0-3+3=0,\qquad
               \deg F_2=0+3-3=0.                    \tag{CT.5}
\]
The zero-winding unit factors have their carried based homotopy;
through it the full-unit physical action is exactly OP.19. This
computes the degrees of these specified local orbit trace maps.
At this stage of the proof, (CT.5) computes local traces only.
Section CT.3 supplies the actual relative two-chains, their graph-loop
boundaries, and the signed boundary-to-variation map used below.

\par

\medskip\noindent\textbf{CT.3. Closing the graph loops and the full variation map}

\par\nobreak

\item The following exact chain calculation has passed both independent
audits. We retain its integral signs and the actual vertex paths.

\item Choose the actual invariant coordinate balls and Hermitian normal
tubes first, with equivariant rounding. Restrict the ordinary product
to their boundaries and use positive radial normalization on each
normal sphere; the transported metric is retained as in OP.19. Thus
the physical action preserves every actual boundary before applying
the ordinary coordinates. Choose the two graph generators of P3K.29 as
\(\gamma_j=a_j-a_0\), each \(a_i\) running from the first vertex
to the second. Lift their curve-interior paths to the plumbing
boundary by a constant generic normal-sphere section \(q_i=q_*\).
Generic means that at both ends both actual normal coordinates are
nonzero, avoiding the other two curve axes. Such constants form an
open dense set. Join the endpoints by fixed paths in the reference
vertex link, the five-sphere with the three axis-circle tubes
removed. The connecting paths can be explicit. Arrange the two
endpoints in the open set with all three coordinates nonzero, and
write their coordinates as \(r_ke^{it_k}\), \(r'_ke^{it'_k}\), with
\(\sum r_k^2=\sum(r'_k)^2=1\). Choose one real argument for each
endpoint coordinate and use
\[
 z_k(s)=\sqrt{(1-s)r_k^2+s(r'_k)^2}\,
                 e^{i((1-s)t_k+st'_k)}.
\]
All coordinates remain nonzero. The three axis-tube inequalities
are linear inequalities in the squared moduli, so these paths stay
outside the same sufficiently small axis tubes. The physical
circle changes only phases, hence its traces avoid them as well.
The connecting paths are fixed in the reference vertex and do not
depend on the parameter.

\item The ordinary vertex maps OP.1 commute with the physical circle, so
the transported action on these vertex connectors is independent
of the parameter. The edge-end traces have the same property. The
local calculation (CT.5) remains unchanged for a constant generic
\(q_*\): varying the constant normal point gives a homotopy of
normal-sphere trace maps with parameter-independent radial ends.
The integers are therefore the ones computed at \(q_*=1\).

\item Here is the exact meridian projection, including the endpoints.
Each edge block is an annulus times its actual normal three-sphere.
The trace at either end is independent of the parameter circle.
Fill its two-torus by \(D^2_\theta\times S^1_\phi\), with the map
constant in \(D^2_\theta\). This closes the local three-chain
inside the edge block by chains of rank at most one. Its third
homology is the normal \(\mathbb Z\), measured by (CT.5), and its
inclusion in the common exterior takes that generator to the
positive complex normal meridian \(m_i\). The same endpoint
fillings on the vertex connectors leave cycles factoring through
a two-dimensional space, since the entire vertex trace is
parameter-independent. Their third homology classes vanish.
Thus the closed graph-loop trace, in the orientation
parameter--circle--path, is exactly
\[
 [K\gamma_j]=[F_j]-[F_0]=6m_0,\qquad j=1,2.       \tag{CT.6}
\]
The cyclic permutation from path--parameter--circle in CT.2 to this
orientation has positive sign. No meridian contribution from a
vertex link has been omitted: the displayed fillings and the
parameter-independent factorization remove it explicitly.

\item Choose relative two-chains \(x_j\) with \(\partial x_j=\gamma_j\).
They exist because \(H_1(E_C)=0\). For the positive physical circle
trace \(\Delta c=a_*([S^1_\phi]\times c)\), the boundary identity
is \(\partial\Delta=-\Delta\partial\). The map from the plumbing
boundary to its regular-neighbourhood retract identifies its second
homology with \(H_2(C)\), with the positive complex sphere basis.
Indeed the relative groups in degrees two and three of the
neighbourhood and its boundary vanish by duality and
\(H^4(C)=H^3(C)=0\). On \(C_1\), the source rotation is
\(z\mapsto e^{i\phi}z\); circle--path orientation is opposite to
its complex orientation. On \(C_2\) the rotation has the opposite
sign and gives its positive complex sphere. The trace of the
\(C_0\) path projects to a one-dimensional path. Hence
\[
             \Delta x_1=D_1,\qquad
             \Delta x_2=-D_2.                         \tag{CT.7}
\]
These equalities are integral: P3K.29 identifies the relative
three-group by the boundary sphere basis, with no torsion.

\item Let \(P_t\) be ordinary-product parallel transport on the common
exterior, and set \(a_t=P_t^{-1}\delta_tP_t\). If \(F=P_1\) is
the actual return, then \(a_1=F^{-1}a_0F\). Define the prism
\(Kc=a_*([I_t]\times[S^1_\phi]\times c)\). The literal boundary
formula is
\[
 \partial Kc=(\Delta_1-\Delta_0)c+K\partial c,
          \qquad F\Delta_1=\Delta_0F.                 \tag{CT.8}
\]
Consequently
\[
 (F-1)\Delta_0x
 =-F\partial Kx+F K\partial x+\Delta_0(Fx-x).
\]
Here \(F\) is the identity on the common boundary. It acts trivially
on absolute third homology because the boundary meridians generate
that group integrally. Moreover,
\(H_2(E_C)=0\) makes the last term a boundary. Thus
\[
             \operatorname{Var}_3(\Delta_0x)
                           =[K\partial x].            \tag{CT.9}
\]
Combining (CT.6)--(CT.9) gives
\(\operatorname{Var}(D_1)=6m_0\) and
\(\operatorname{Var}(D_2)=-6m_0\). The integral symmetry proof
of P3K.29 applies to this ordinary relative return without any
circle-equivariance hypothesis. OR.3 gives \(V_{00}=0\).
The resulting exact matrix is therefore
\[
 V=\begin{pmatrix}0&6&-6\\6&0&0\\-6&0&0\end{pmatrix}. \tag{CT.10}
\]
No stable derivative on the two nonfixed cycles, and no further fan
return, is determined by this chain calculation alone.

\item The exact certificate is \path{code/part3/source_cusp_circle_trace.py}, with its JSON and captured output.

\end{enumerate}
\end{proof}




\begin{theorem}\label{night:P3K.53}
\StatementRecord{night:P3K.53}{proved}{S6\_06\_BLIND: GPT-6 Astra (OpenAI, ChatGPT Pro mode, effort Pro), recorded re-derivation in the checkpoint dated 18 November 2026}{\cref{night:P3F.51}, \cref{night:P3K.41}}

Use the marked source
\[
 Y^2=X^3+(1-2v)X^2+TX+v^2,\qquad G=(0,v),
 \qquad T_0(v)=v^2-2v.
\]
Retain one III-star cap, the selected whole-zero \(I_1\) cusp,
its chosen periodic fan, and all source, Cartier line, finite
character and annular data of P3K.49. The cap is the good
\(m=4p\) branch for \((a,b)=(1,2)\), or the reduced
\(m=4p-1\) branch for \((a,b)=(2,1)\), with the signed norm
\(n=\pm1\). There is no added logarithmic disc on the two
unselected fibres that collide at \(v_*=1/8\).

The canonical continuation with the literal scalar
\(\lambda(v)=v^N\lambda_0\), \(N=km\), extends to an actual
smooth proper holomorphic family over a disc around \(v_*\).
The root, line and annular maps are the continuation of their
full reference data; no additional cap-data loop is inserted.
The scalar is in the common contracting range, with its
actual radial comparison retained. Consequently the positive
local return has oriented class zero in \(\Theta_7\), and
its mapping torus is oriented-diffeomorphic to \(S^1\times S^6\).
An averaged connection gives the corresponding smooth
trivialization carrying the physical circle.

This result supplies the unselected-collision disc in the
global marked-source problem. It assigns neither a
source-zero return nor an infinity endpoint.

\end{theorem}
\begin{proof}
\begin{enumerate}[label=\arabic*.]
\item \par

\medskip\noindent\textbf{UI.1. The source and selected cusp extend as actual maps}

\par\nobreak

\item Use the literal Weierstrass model
\[
 Y^2=X^3+(1-2v)X^2+TX+v^2,\qquad G=(0,v),
 \qquad T_0=v^2-2v.
\]
The other two critical values are
\[
 T_\pm=\frac{1+4v\pm(1-8v)^{3/2}}8.
\]
At \(v_*=1/8\) they coalesce at \(T_*=3/16\). Their singular
point is \((X,Y)=(-1/4,0)\). Write the equation as \(F=0\), with
\[
 F=Y^2-X^3-(1-2v)X^2-TX-v^2.
\]
Then \(F_T=1/4\) at that point. The total source near the
collision is literally the graph
\[
 T=\frac{Y^2-X^3-(1-2v)X^2-v^2}{X},
\]
in the coordinates \((X,Y,v)\), with \(X\ne0\). Its map to the
\(v\)-disc is therefore a submersion, including the type-II fibre.
No separate labels or separate discs for \(T_+\) and \(T_-\) are
used in this neighbourhood.

\item The selected value is \(T_0(v_*)=-15/64\), which is distinct
from \(T_*\). Its nodal equation is exactly
\[
 Y^2=(X-v)^2(X+1).
\]
Choose the holomorphic branch \(\rho(v)=\sqrt{1+v}\) on a small
disc \(D\) around \(v_*\), with its reference value fixed. The node
preimages are \(s=\pm\rho(v)\) in the actual normalization
\(X=s^2-1\), \(Y=s(s^2-1-v)\). The branch does not exchange on
\(D\). The physical selected height is
\(q_1=T-T_0(v)\); the parameter Morse map and the Tate coordinate
retain the holomorphic nonzero unit \(q_T/q_1\). The marked point
\[
 P_G=\frac{1+\rho(v)}{1-\rho(v)}
\]
is finite and nonzero on this disc. The actual theta/Cartier map,
including its integral reindexing, consequently extends on the
selected cusp charts. A fixed periodic fan gives the same toric
coordinate charts over \(D\), with holomorphic nonzero unit
transitions. There is no reflection or additional height winding.

\item The sections \(G\) and \(O\) remain smooth and avoid the colliding
point. Thus their actual Cartier line
\(M_a=\mathcal O(aG-O)\) extends over the entire source family.

\par

\medskip\noindent\textbf{UI.2. Uniform minimal source and cap roots}

\par\nobreak

\item On \(D\), choose the actual nonzero holomorphic cube root
\(u(v)^3=v\). The map
\[
 (X,Y,T)\longmapsto(x,y,t)=(u^{-2}X,u^{-3}Y,u^{-4}T)
\]
gives
\[
 y^2=x^3+3r(v)x^2+tx+1,\qquad G=(0,1),\qquad
 3r(v)=\frac{1-2v}{u(v)^2}.
\]
At \(v_*\), \(u=1/2\), \(r=1\), and \(r'(v_*)=-8\) in this
reference branch. These are the full source-coordinate maps of
F51, specialized to a parameter disc through its E7-B member.
They do not discard the \(x^2\) term. The selected node is
\(t=-15/4\), while the type-II collision is \(t=3\).

\item An alternate fixed cube-root reference is also retained. Replacing
\(u\) by \(\zeta u\), \(\zeta^3=1\), gives the literal constant
coordinate map
\[
 (x,y,t)\longmapsto(\zeta x,y,\zeta^2t),\qquad
 3r\longmapsto\zeta\,3r.
\]
It fixes the marked zero and generator sections. The root coordinate
is carried by its chosen constant lift \(z\mapsto c z\), where
\(c^m=\zeta\). Since \(q_{\rm new}=\zeta q_{\rm old}\),
the two infinity coordinates \(x_0=q^2x\), \(y_0=q^3y\) are
unchanged. Thus the full minimal-root map is
\[
 (z,x',y')\longmapsto(cz,c^{-2k_0}x',c^{-3k_0}y'),
 \qquad k_0=3p\ \hbox{or}\ 3p-1
\]
in the good or reduced case. The actual Cartier map transports
the divisors \(aG-O\), and the root-line map retains this chosen
constant character.
The primitive source lift contributes \(\zeta^a\), and the old
selected height expressed in the new coordinate contributes
\(\zeta\). Their factor in the suspension is
\(\zeta^{1+ab}=1\). Thus this actual constant map transports an
alternate reference to the one just displayed without creating a
parameter winding or changing the literal scalar.

\item The III-star source resolution is uniform here. In the actual
infinity coordinates \(q=1/T\), \(x_0=q^2X\), \(y_0=q^3Y\), its
equation is
\[
 y_0^2=x_0^3+(1-2v)q^2x_0^2+q^3x_0+v^2q^6.
\]
For the good cap \(m=4p\), put \(q=z^m\) and divide \(x_0,y_0\)
by \(z^{6p},z^{9p}\). The exact root equation is
\[
 y'^2=x'^3+(1-2v)z^{2p}x'^2+x'+v^2z^{6p}.
\]
For the reduced cap \(m=4p-1\), divide by
\(z^{6p-2},z^{9p-3}\). Its equation is
\[
 y'^2=x'^3+(1-2v)z^{2p}x'^2+zx'+v^2z^{6p}.
\]
In both cases the actual section is \(G'=(0,vz^{3p})\), and
the root deck is
\[
 D:(z,x',y')\longmapsto
       (\zeta_m z,\zeta_m^{2p}x',\zeta_m^{3p}y').
\]
The good central elliptic curve, reduced central III model,
resolution ideals and section-leading coefficients are exactly
those of F51. All coefficients above are holomorphic on \(D\),
and \(v\) is a unit there. The actual equivariant blowup ideals
and their Cartier support functions therefore give the same
source/root/line correspondence over this disc. The section does
not cross an additional centre. The minimal source is smooth
over \(D\): near the finite collision this was the graph in
UI.1, and at infinity these are the uniform resolved charts.

\item The cap division branch, norm lift and finite character also
extend. On the good-root atlas, multiplication by \(m\) is
locally invertible at the specified division point. In the
reduced case use the actual good-root atlas and its retained
finite descent. The inverse branch selected at the reference
point continues holomorphically across \(D\), with its finite
kernel label unchanged. The norm-ratio is a nonzero base unit;
its root and the finite averaging cochain are taken on the
polydisc with that reference value. Thus the literal
\(A^m=B^n\), the finite character and every root-line Cartier
correction are preserved as maps, not only as fibrewise classes.

\par

\medskip\noindent\textbf{UI.3. Translation, line lift and annular continuation}

\par\nobreak

\item The source translation extends through the unselected type-II
point. A direct local verification uses the holomorphic vertical
vector field
\[
 W=2Y\,\partial_X+
       (3X^2+2(1-2v)X+T)\,\partial_Y.
\]
It is tangent to the displayed equation at fixed \(T,v\). Its
local expressions near \(O\) are the invariant elliptic vector
field. On the smooth part of the central cuspidal fibre the
group law can be checked directly. Its normalization is
\[
 X=s^2-\tfrac14,\qquad Y=s^3,
\]
with \(O\) at \(s=\infty\), the cusp at \(s=0\), and
\(G\) at \(s=1/2\). The invariant additive coordinate is
\(\tau=1/s\), with \(\tau(O)=0\) and \(\tau(G)=2\).
Hence \(Q=bG\) has \(\tau=2b\), or \(s=1/(2b)\), for
\(b=1,2\). In particular it is a smooth point and
\(2Y(Q)=1/(4b^3)\ne0\), so \(W(Q)\ne0\). Since
\(dX/(2Y)=-d\tau\), the displayed vector field takes \(O\)
to \(Q\) at the exact central time \(\beta_*=-2b\).
The time parameter \(\beta(T,v)\) taking \(O\) to \(Q\)
is holomorphic near the collision: at the central value the
time derivative is \(W(Q)\ne0\), so the analytic inverse
function theorem gives that branch. Its flow is an actual
holomorphic translation on nearby smooth fibres and extends
through the cusp, with inverse the flow of \(-\beta\). The
flow stays on compact fibres; the finite time interval is
covered by finitely many ordinary analytic flow charts. This
also proves the parameter-holomorphic extension.

\item The canonical Cartier/Hom map has its already specified whole
fibre divisor at \(q_1=0\), which is disjoint from this
collision. It has no new parameter divisor at \(v_*\): evaluate
its rational source formula at one ordinary point where all
Cartier frames are units. Its value is holomorphic and nonzero
there, and \(v\) is a unit. The total minimal source is smooth,
so a zero or pole supported only at the type-II curve would
have codimension two and cannot be a Cartier divisor. The map
therefore extends invertibly across this curve in the actual
line frames. This is the same source/line map used in the
suspension, not a new linearization.

\item The normalizing coordinate arrow in UI.2 preserves its literal
scalar. In the direction from \(E_v\) to the \(G=(0,1)\) model,
the primitive lift changes by \(u^{-2a}\), while the selected
height changes by \(q_{1,E}=u^4q_{1,\mathrm{norm}}\). Since
\(ab=2\), the factors in \(q_1U_a^b\) cancel exactly. The
remaining powers of the nonzero function \(u(v)\) are retained
in the root/line coordinate maps; their logarithms extend on
\(D\).

\item For the annular arrow, use a fixed cap-overlap annulus in the
\(T\)-plane outside all these finite critical values. This is
the same type of fixed annulus used to define the full \(F_0\)
in RCZ.5. Its source fibres are smooth over annulus times \(D\).
On its fixed angular charts, the actual source/root/Cartier
maps and their finite deck powers just constructed are
holomorphic over \(D\). Continue their Lie logarithms from
the full reference overlap, with the integer \(d\delta\),
finite characters and all ordinary and cyclic lattice
cochains included. Each discrepancy is integral and locally
constant in \(v\); it keeps its full reference value. The
fixed-contour Cauchy splitting in RCZ.5 preserves parameter
holomorphy and this reference. This gives the actual annular
map over \(D\).

\item This continuation uses the same fixed contours and source/line
root expressions as the continuation along the preceding
regular parameter path. Their inputs have no singularity at
\(v_*\). Thus it extends that canonical continuation itself;
no independent \(q^j\delta\) or parameter-scalar winding is
inserted at the endpoint. Extra prescribed cap-data loops
would be the separate loops of K41 and are outside this
statement.

\item Finally \(\lambda(v)=v^N\lambda_0\) is nonzero on \(D\) and
has the logarithm fixed by its reference value. Its actual
root multipliers and the scalar term in the ordinary-to-cap
cochain extend with that logarithm. This does not replace
\(\lambda(v)\) by a constant.

\par

\medskip\noindent\textbf{UI.4. The compact family and its actual positive return}

\par\nobreak

\item Choose the scalar in the common contracting range on the compact
parameter disc and the finitely many retained source/base charts.
The line maps in UI.3 have bounded operator norms in a fixed
Hermitian metric on these compact pieces. The usual inner and
outer norm shells therefore give a uniform compact fundamental
region for the suspension. At the finite cap use its actual free
root quotient and the root-line norm shell; at the selected cusp
use the fixed toric charts of UI.1. All their transitions are
the maps just displayed, holomorphic on \(D\). Each total chart
is smooth and its projection to \(D\) is a submersion.

\item The resulting family is proper over \(D\). Indeed each local
model is proper over its fixed base chart times \(D\), the
overlaps are the actual fibre isomorphisms above, and a finite
cover of the compact base by relatively compact such charts
gives compact inverse images of compact subsets of \(D\).
This is the actual compact suspension/cap/cusp family, including
the prescribed physical circle and fan.

\item A smooth connection for this proper submersion, averaged over
the physical circle, gives an actual smooth trivialization
over the contractible disc by transport along its radial paths.
Its restriction to the positive boundary loop has an endpoint
isotopy to the identity. Hence this unselected-collision loop
has oriented class zero and its mapping torus is
\(S^1\times S^6\). The comparison uses the actual inverse
closing arrow, which is intertwined by this disc trivialization;
there is no loop reversal.

\item For a smaller reference scalar within the same contracting
region, the radial scalar path gives the same smooth family
construction over the path times \(D\); its transport conjugates
the two boundary returns. This records the actual comparison
used when choosing a uniform small scalar for the global
source-loop construction.

\par

\medskip\noindent\textbf{UI.5. Scope in the global loop relation}

\par\nobreak

\item The proof supplies the disc-extending unselected-collision
branch of the marked global family. It does not identify the
two source-zero lifts on \(\rho^2=1+v\), discard their J-fan
transition, or compute the infinity endpoint. Those exact
maps remain the separate global-relation inputs.

\end{enumerate}
\end{proof}




\begin{theorem}\label{night:P3K.55}
\StatementRecord{night:P3K.55}{proved}{S6\_06\_BLIND: GPT-6 Astra (OpenAI, ChatGPT Pro mode, effort Pro), recorded re-derivation in the checkpoint dated 18 November 2026}{\cref{night:P3F.56}, \cref{night:P3K.29}, \cref{night:P3K.51}, \cref{night:P3K.52}}
Retain the actual positive selected-cusp source loop of P3F.54/P3F.56,
its physical height, fixed scalar, full line and finite-cap maps. For
every integral unimodular matrix \(P\), the two-triangle fan obtained
from OE.1 by \((x,h)\mapsto(Px,h)\) admits the periodic oriented
ordinary plumbing product constructed below. Its complete period
winding matrix and normal first exponents are
\[
 B_P=(b_{ij})=-6P^{-1}e_1e_1^TP,\qquad
 (l_0,l_1,l_2)=(b_{22}-b_{21},b_{11}-b_{12},b_{22}).
\]
Writing \((\xi,\upsilon)^T=P^{-1}e_2\), define
\[
 w=(\xi,\upsilon,-\xi-\upsilon),\qquad
 d=(-\xi-2\upsilon,-2\xi-\upsilon,\xi-\upsilon).
\]
Its actual local orbit and stable-derivative trace degrees are
\[
 I_i=-3w_i+d_i(l_i+3),\qquad W_i=-3w_i-d_i(l_i+3).
\]
In the positive complex meridians and relative dual cycles, its
actual symmetric variation matrix \(V\) and ambient derivative
\(\delta\) satisfy the integral identities
\[
 V(w_je_j-w_0e_0)=I_je_j-I_0e_0,\qquad
 w_j\delta_j-w_0\delta_0=W_j-W_0\quad(j=1,2).
\]
On each native star the ordinary-product comparison is the explicit
pure-left endpoint \(L^6\). The proof supplies every character,
period, unit, normal and chain map used in these formulas. It leaves
the transverse invariants and the general central return unassigned.
\end{theorem}
\begin{proof}
\begin{enumerate}[label=\arabic*.]
\item \par

\medskip\noindent\textbf{GF.1. The actual character and period changes}

\par\nobreak

\item Let \(P\) be an integral two-by-two matrix with determinant \(+1\) or \(-1\), and
write \(R=P^{-1}\). The fan here is the image under \(P\) of OE.1: its two
triangles at each lattice point \(n\) have vertices

\[
 n,\ n+Pe_1,\ n+Pe_2,
 \qquad n+Pe_1,\ n+Pe_2,\ n+P(e_1+e_2).
\]

\item These are the specified two triangles in each period parallelogram.
The homogeneous lattice map is \((x,h)\) to \((Px,h)\); its determinant is
\(\det P\), and its inverse is integral. On the dense torus use the literal
holomorphic character change

\[
 Z'=Z^{R_{11}}V^{R_{12}},\qquad
 V'=Z^{R_{21}}V^{R_{22}},\qquad q'=q.                 \tag{GF.1}
\]

\item Its inverse uses the rows of \(P\). Thus it extends to an isomorphism
between the toric charts with the two matched fans; on each matched
unimodular triangle it is the identity in the corresponding three
character coordinates. This is a coordinate description of the actual
chosen fan filling. The original source and line coordinates on its
overlap with the fixed exterior are still obtained by the inverse of
(GF.1), so the source and all finite caps are retained.

\item The original commuting periods are exactly OE.2:

\[
 h:(Z,V,q)\longmapsto(\kappa qZ,\alpha V,q),\qquad
 B:(Z,V,q)\longmapsto(gZ,\eta qV,q),\qquad\eta=\lambda\beta.
\]

\item Replace their basis by \(h'_j=h^{P_{1j}}B^{P_{2j}}\), for \(j=1,2\).
This has the inverse integral period change \(R\) and therefore gives
the same actual quotient. In the characters (GF.1) the two periods
are

\[
 h'_1:(Z',V')\longmapsto(qA_{11}Z',A_{21}V'),\qquad
 h'_2:(Z',V')\longmapsto(A_{12}Z',qA_{22}V'),           \tag{GF.2}
\]

where every complete multiplier is retained:

\[
 A_{ij}=(\kappa^{P_{1j}}g^{P_{2j}})^{R_{i1}}
                (\alpha^{P_{1j}}\eta^{P_{2j}})^{R_{i2}}.
                                                               \tag{GF.3}
\]

\item The power of \(q\) in the \(i\)-th character of \(h'_j\) is \((RP)_{ij}\).
This proves the height factors in (GF.2), including their signs.
The units are nonzero on the same protected collars as the original
ones. At \(q=0\) their parameter winding matrix is exactly

\[
 B_P=(b_{ij})=-6P^{-1}e_1e_1^TP,
                    \qquad\operatorname{tr}B_P=-6.   \tag{GF.4}
\]

\item There is no omitted unit homotopy in this identity. Write the based
ratio of \(\kappa\) as \(\exp(-6i\theta)\) times \(\exp(f_\kappa)\), and write the
based ratios of \(\alpha,g,\eta\) as \(\exp(f_\alpha)\), \(\exp(f_g)\), \(\exp(f_\eta)\).
The four \(f\)'s are their actual based periodic logarithms from F56.
The periodic logarithm of the ratio \(A_{ij}/A_{ij,*}\) after its displayed
angular factor is

\[
 R_{i1}(P_{1j}f_\kappa+P_{2j}f_g)
       +R_{i2}(P_{1j}f_\alpha+P_{2j}f_\eta).          \tag{GF.5}
\]

\item This formula supplies every subsequent unit homotopy.

\par

\medskip\noindent\textbf{GF.2. The two vertices and the three full normal loops}

\par\nobreak

\item In the new character coordinates put \((a,b,c)=T^+_0\) and
\((u,v,w)=T^-_0\). Direct substitution in (GF.2) gives

\[
\begin{array}{c|ccc}
 &u&v&w\\ \hline
 C_0&A_{11}A_{21}bc&A_{21}^{-1}ab&A_{11}^{-1}b^{-1}\\
 C_1&A_{12}A_{22}bc&A_{22}^{-1}c^{-1}&A_{12}^{-1}ac\\
 C_2&a^{-1}&ab&ac.
\end{array}                                                   \tag{GF.6}
\]

\item Each triple product is \(q\). Keep the first vertex coordinates and
use \((u,A_{22}v,A_{11}w)\) at the second, with the corresponding
initial constant units in the reference fibre. The tangent
coordinates on each curve are then literal reciprocals. The
ordered normal pairs are the same as OP.1, and their normal
transition matrices are \(zD_i\) with

\[
 D_0=\operatorname{diag}(A_{22}/A_{21},A_{11}A_{21}),\quad
 D_1=\operatorname{diag}(A_{11}/A_{12},A_{12}A_{22}),\quad
 D_2=\operatorname{diag}(A_{22},A_{11}).                \tag{GF.7}
\]

\item For the based loops \(K_i=D_{i,*}^{-1}D_i\), the first diagonal
exponents and the sum of the two exponents are therefore

\[
 l_0=b_{22}-b_{21},\qquad l_1=b_{11}-b_{12},\qquad
 l_2=b_{22},\qquad l_i+\widetilde l_i=-6.              \tag{GF.8}
\]

\item Every determinant has the same even winding \(-6\), including matrices
\(P\) whose entries have arbitrary signs. We now construct the actual
normal contractions, rather than infer them from that parity.

\item For an integer \(d\) define

\[
 H_d(s,\theta)=
 (\sin\vartheta(s)e^{id\theta}+j\cos\vartheta(s))
                    (\sin\vartheta(s)+j\cos\vartheta(s))^{-1},
                                                               \tag{GF.9}
\]

where \(\vartheta\) is OP.6's fixed smooth radial function. These maps
are smooth, based, and periodic, and satisfy \(H_d(0,\theta)=1\),
\(H_d(1,\theta)=\exp(id\theta)\). Put

\[
 A=H_{-3},\qquad B_i=H_{l_i+3},\qquad
                  T_i(Q)=A Q B_i,
 \quad Q=n_0+j\overline{n_1}.                         \tag{GF.10}
\]

\item At \(s=1\) its two complex normal factors are \(\exp(il_i\theta)\)
and \(\exp(i(-6-l_i)\theta)\), respectively: left and right quaternion
phases with exponents \(x,y\) give complex exponents \(x+y,x-y\).
Thus (GF.10) is the required contraction of each angular loop.
Multiplying by the diagonal exponentials from (GF.5), with OP.9's
radial cutoff, gives a contraction of the full \(K_i\). In particular
the original units, norms, and reference values remain present.

\item At \(P=I\) the right factor on \(C_0\) and \(C_2\) is \(H_3\), while OP.8 used
\(H_{-3}^{-1}\). We retain their exact comparison. Set \(a=\sin\vartheta\),
\(b=\cos\vartheta\) and \(n=ai-bk\), a unit imaginary quaternion. The
expansion of (GF.9) is

\[
 H_d=b^2+a^2\cos(d\theta)
       +ab(1-\cos(d\theta))j+a\sin(d\theta)n.
\]

\item Conjugation by \(n\) fixes \(n\) and sends \(j\) to \(-j\). Consequently
\(nH_d^{-1}n^{-1}=H_{-d}\). The explicit path

\[
 C_t=\cos(\pi t/2)+n\sin(\pi t/2),\qquad
 C_tH_{-3}^{-1}C_t^{-1},\qquad 0\leq t\leq1,
                                                               \tag{GF.10a}
\]

joins OP's right factor to \(H_3\). It is fixed at \(s=0\), where \(H=1\),
and at \(s=1\), where \(n=i\) commutes with the diagonal endpoint. It is
also fixed at both parameter endpoints. Thus it glues with the
protected vertex collars and the carried unit maps to an actual
based homotopy of the two entire plumbing products. This verifies
the comparison to the initial matrices of K51/K52.

\item Send \((z,n)\) to \((z,K_i(s(|z|),\theta)n)\) on each tube. The exact
collar identity \(zD_{i,*}K_i=zD_i\) glues this to the two vertex
maps. The normal rescaling \(R_t(z,n)=(z,tn)\), followed by the
cutoff flow of its velocity, carries the original higher normal
terms and their protected vertex collars as in OP.3. The curves'
interiors are disjoint, so these three flows have the stated
disjoint supports. This constructs a periodic oriented ordinary
product of the entire labelled plumbing on this actual fan.

\item Each separate native star is also explicit. Its second coordinates
are, respectively,

\[
\begin{array}{c|ccc}
 C_0&z'=A_{11}w&n'_0=A_{21}v&n'_1=(A_{11}A_{21})^{-1}u\\
 C_1&z'=A_{22}v&n'_0=A_{12}w&n'_1=(A_{12}A_{22})^{-1}u\\
 C_2&z'=u&n'_0=v&n'_1=w.
\end{array}                                                   \tag{GF.11}
\]

\item The first normal pairs are \((a,c),(a,b),(b,c)\). Each row gives
exactly \(z'=z^{-1}\), \(n'=zn\). These are actual periodic holomorphic
star products with the complete units shown. The ordinary product
comparison on its normal Hopf boundary is

\[
              q_H\longmapsto (U(z)AU(z)^{-1})q_HB_i.
                                                               \tag{GF.12}
\]

\item The left factor is precisely OP.20 and has degree \(-6\) on the
Hopf-oriented base sphere times the parameter. The right factor
has only radial and parameter variables and has degree zero.
The explicit cylinder extension and orientation of OR.2 give
the pure-left endpoint \(L^6\). Thus on every separate star the
ordinary relative return is the native relative return composed
with \(L^6\); its right integer is unchanged. This is an actual
boundary-product comparison on the given fan.

\par

\medskip\noindent\textbf{GF.3. Physical weights and exact local orbit maps}

\par\nobreak

\item The physical delta character in the new coordinates has integral
weights \((\xi,\upsilon)^T=Re_2\). At the first vertex its weights on
\((a,b,c)\) are \((-\xi-\upsilon,\xi,\upsilon)\); at the second they are
\((\xi+\upsilon,-\upsilon,-\xi)\). Define the tangent and normal-difference
weights in the above ordered pairs by

\[
\begin{array}{c|rr}
 &w_i&d_i\\ \hline
 C_0&\xi&-\xi-2\upsilon\\
 C_1&\upsilon&-2\xi-\upsilon\\
 C_2&-\xi-\upsilon&\xi-\upsilon.
\end{array}                                                   \tag{GF.13}
\]

\item The original physical normal action on \(Q\) is left multiplication
by \(\exp(-iw_i\phi/2)\) and right multiplication by \(\exp(id_i\phi/2)\).
In the ordinary product (GF.10) it becomes exactly

\[
 Q\longmapsto
 (A e^{-iw_i\phi/2}A^{-1})Q
                 (B_i^{-1}e^{id_i\phi/2}B_i).         \tag{GF.14}
\]

\item The full unit version is obtained by the actual conjugations in
(GF.5), with the same based homotopy. Its source point moves by
\(z\) to \(\exp(iw_i\phi)z\). On the positive radial path \(z=\tan\chi\),
write \(U=\cos\chi+j\sin\chi\) and \(E_i=\exp(iw_i\phi/2)\). The moved
Hopf representative is exactly \(E_i^{-1}UE_i\). At \(q_H=1\) the
normal-sphere orbit trace is consequently

\[
 F_i=E_i^{-1}UE_i A E_i^{-1}A^{-1}U^{-1}
                       B_i^{-1}e^{id_i\phi/2}B_i.     \tag{GF.15}
\]

\item This includes \(w_i=0\): its left \(A\) factors cancel and the result is
the right-axis map in (GF.14). Every expression closes at \(\phi=2\pi\)
because \(w_i\) and \(d_i\) have the same parity. It is periodic in \(\theta\),
and its radial-end values are \(\theta\)-independent.

\item For completeness the axis degree in (GF.9) is computed directly.
For \(d=0\) the map is constant. For \(d\) nonzero, the inverse images
of \(-i\) under \(H_diH_d^{-1}\) have \(\vartheta=\pi/4\) and \(\exp(id\theta)=-1\).
There are \(|d|\) of them. At each, \(H_d=j\), and differentiation gives

\[
 \partial_s(H_diH_d^{-1})=4\vartheta'(s)k,\qquad
 \partial_\theta(H_diH_d^{-1})=-d j.
\]

\item Since \((k,j)\) is the positive outward tangent basis at \(-i\), the
degree is \(-d\). For \(H_d^{-1}iH_d\) the first derivative changes
sign and the second remains \(-dj\), giving degree \(d\).
Thus the axes \(AiA^{-1}\) and \(B_i^{-1}iB_i\) have degrees
\(3\) and \(l_i+3\). Conjugating the first by \(U\) preserves its degree
through \(U_t=\cos(t\chi)+j\sin(t\chi)\); both collapsed ends remain
\(\theta\)-independent. The axis \(UiU^{-1}\) has degree zero.

\item Factor (GF.15) into \(E_i^{-1}\), \(\exp(iw_i\phi\,\operatorname{Ad}_U(i)/2)\),\newline
\(\exp(-iw_i\phi\,\operatorname{Ad}_{UA}(i)/2)\), and
\(\exp(id_i\phi\,\operatorname{Ad}_{B_i^{-1}}(i)/2)\). Integration of the positive
quaternion volume gives degree \(ad\) for an axis of degree \(d\) and
angle \(a\phi/2\), for every integer \(a\). The Maurer product boundary
forms vanish at the radial ends and at the constant central-sign
angle ends; the parameter faces cancel. Therefore its exact local
normal-meridian degree, in path--parameter--circle orientation, is

\[
                         I_i=-3w_i+d_i(l_i+3).         \tag{GF.16}
\]

\item For the full tangent derivative, use K52's covering group of
\(GL^+\) and its normalized third form. The local normal block in
(GF.14) has stable degree \(\deg(\mathrm{left})-\deg(\mathrm{right})\). The base rotation
has only the \(\phi\) variable. At \(\phi=2\pi\) its Spin endpoint is
\((-1)^{w_i}\), while the normal Spin endpoint has the same sign;
their product is \(+1\). Thus the full tangent lift is closed even
for odd weights. The same exact boundary-form calculation gives

\[
                         W_i=-3w_i-d_i(l_i+3).         \tag{GF.17}
\]

\item The ambient/local frame comparison extends through the annulus
times \(D^4\) and has zero third pairing with each capped edge trace.
The original nonzero unit factors and higher derivatives are
carried by the actual homotopies above, so these are degrees of
the full source family, with its original physical circle.

\par

\medskip\noindent\textbf{GF.4. The resulting actual integral relations}

\par\nobreak

\item Choose generic normal sections at the two ends of each edge and
the fixed vertex-link paths of CT.3. Their three coordinate
moduli are nonzero, and the interpolation of squared moduli
preserves all axis-tube inequalities. The physical rotations,
now with the integer weights in (GF.13), preserve these moduli.
The vertex actions are \(\theta\)-independent, because the vertex
unit maps commute with the diagonal physical circle.

\item Cap each \(\theta\)-independent endpoint by \(D^2_\theta\) times \(S^1_\phi\),
constant in its \(D^2\) variable. The capped edge trace has class
\(I_im_i\) in the common exterior; the vertex trace has zero third
class. For the graph loops \(\gamma_j=a_j-a_0\) and the relative
two-chains \(x_j\) with those boundaries, the actual positive circle
trace and the complex sphere orientation give

\[
 \Delta x_j=w_jD_j-w_0D_0,\qquad j=1,2.               \tag{GF.18}
\]

\item This follows on the regular-neighbourhood retract: circle--path
orientation on a curve with weight \(w_i\) has degree \(-w_i\), and
the boundary identity is \(\partial\Delta=-\Delta\partial\). The
integral boundary map identifies the relative three-group with
the complex curve sphere basis, exactly as in CT.3. No claim
that the two vectors in (GF.18) form a primitive sublattice is
needed here.

\item Let \(V\) be the actual ordinary-product variation matrix on this
fan and \(\delta\) its actual ambient stable-derivative vector. The
two exact prism identities proved in CT.3 and ST.1 now give

\[
 V(w_je_j-w_0e_0)=I_je_j-I_0e_0,\qquad
 w_j\delta_j-w_0\delta_0=W_j-W_0,\qquad j=1,2.        \tag{GF.19}
\]

\item The matrix \(V\) is symmetric by the integral doubled-exterior
intersection proof of K29, since the boundary meridians generate
\(H_3\) and the return fixes its boundary. Equations (GF.19) are
complete integral identities for the actual constructed product.
They leave the transverse variation/derivative entries to be
computed by the remaining global or flop comparison. In particular
this proof does not assign the missing central return or replace
the needed product-reset map by a scalar-family recurrence.

\item At \(P=I\), \((l_i)=(0,-6,0)\), \((w_i)=(0,1,-1)\),
\((d_i)=(-2,-1,-1)\); the formulas give \(I=(-6,0,0)\) and \(W=(6,-6,6)\).
Substitution of the completely audited initial matrices of K51/K52
verifies both relations, with every sign retained. This equality
is also executed in the accompanying exact certificate.

\item The exact certificate is \path{code/part3_root/source_cusp_general_fan_product.py}. It includes symbolic character, period, contraction and orientation identities, and all 1384 integral bases with entries in \([-8,8]\). The written proof establishes the statements for every integral unimodular basis.

\end{enumerate}
\end{proof}




\begin{theorem}\label{night:P3K.61}
\StatementRecord{night:P3K.61}{proved}{S6\_06\_BLIND: GPT-6 Astra (OpenAI, ChatGPT Pro mode, effort Pro), recorded re-derivation in the checkpoint dated 18 November 2026}{\cref{night:P3K.22}, \cref{night:P3K.25}, \cref{night:P3K.29}}
For an actual primitive scalar family of P3K.22, including
the OS45 family of P3K.10, retain its positive physical circle, full
source and line units, actual cap-to-ordinary gluing and forward
parameter arrow. On every nonzero signed degree-\(M\) pullback of
its primitive scalar loop, the three labelled cusp curves, their two
labelled vertices and their positive complex curve orientations admit
no periodic smooth product intertwining the physical delta action and
carrying all these marked submanifolds to constant copies.

The literal comparison at the second outer fixed branch between the
fixed-curve native product and the horizontal-curve native product is
\[
 z_1=(g/\kappa)z_0,\qquad
 (N_{1,-},N_{1,+})=
 \bigl((\eta/\alpha)N_{0,-},\,
       (\kappa\alpha/(g\eta))N_{0,+}\bigr).
\]
After conjugating the negative-weight first line, its determinant is
\[
 \det_+\Gamma
 =\frac\kappa g\frac\alpha{\overline\alpha}
                 \frac{\overline\eta}{\eta},
 \qquad\operatorname{wind}(\det_+\Gamma)=-2M.
\]
The other outer branch comparison is the identity. This identifies
the retained outer line clutch missing from P3K.25's claimed
parameter-free determinant, and disproves the marked simultaneous
equivariant-product premise used in P3K.29, Section3. It assigns no
corrected global characteristic number. The separately proved
cellular topology, reciprocal-section basis maps, variation symmetry,
and stable-frame change identities do not require that premise.
\end{theorem}
\begin{proof}
\begin{enumerate}[label=\arabic*.]
\item \par\medskip\noindent\textbf{SOC.1. Literal periods and the closed scalar charts.}\par\nopagebreak[4]

\item The physical line coordinate is \(V\). The actual cusp periods are
\[
 h:(Z,V,q)\longmapsto(\kappa qZ,\alpha V,q),\qquad
 B:(Z,V,q)\longmapsto(gZ,q\eta V,q),\qquad
 \eta=\lambda\beta.
\]
Here \(\kappa,\alpha,g,\beta\) are independent of the scalar parameter.
For the primitive K22 cap-to-ordinary convention the forward arrow is
\(R_{\rm ord}(Z,V)=(Z,q^{-1}V)\). Reindexing by the actual period
\(B\) leaves \(E_\lambda(Z,V)=(gZ,\eta V)\). On a cut scalar
interval the maps
\[
 C_\tau(Z,V)=(A_\tau Z,D_\tau V),\qquad
 A_\tau=e^{-\chi\log g},\quad D_\tau=e^{-\chi\log\eta}
\]
have \(C_{\rm end}E_\lambda=1\) on full collars. The same actual
logarithms and cutoff as SCN.1 retain every spatial and parameter
derivative. They commute with the full periods and the physical delta
circle, preserve \(q\), and act on the two vertex charts by
\[
 (a,b,c)\mapsto((AD)^{-1}a,Ab,Dc),\qquad
 (u,v,w)\mapsto(ADu,D^{-1}v,A^{-1}w).
\]
Substitution in all three curve transition rows leaves their complete
coefficients unchanged. Thus in these genuinely closed charts the
primitive scalar unit windings remain
\[
 \operatorname{wind}(\kappa,\alpha,g,\eta)=(0,0,0,1).
 \tag{P3K.61.SOC.1}
\]
For a positive degree-\(M\) pullback they are \((0,0,0,M)\).
The pullback atlas is explicit. On its \(j\)-th cut sheet, write
\(u=j+\tau\), \(0\le j<M\), and put
\[
 C^{(M)}_{j+\tau}=C^{(1)}_\tau E_{\lambda(\tau)}^{-j}.
\]
At the next seam both collar expressions equal
\(E_\lambda^{-(j+1)}\), with all derivatives. The last cancels the
forward arrow \(q^{-M}\) after the actual reindexing \(B^M\).
This is the pullback of the original family, not a separately chosen
normal product. For negative \(M\), reverse the parameter of the
positive \(|M|\)-fold pullback; all the displayed windings acquire
the same sign. No change from an endpoint-to-initial closing map to a
forward map is implicit in this calculation.

\par\medskip\noindent\textbf{SOC.2. The outer fixed-normal map is not the inner common cap gauge.}\par\nopagebreak[4]

\item At the second vertex the actual native fixed-curve product uses
\[
 z_0=\kappa w,\qquad N_{0,-}=\alpha v,\qquad
 N_{0,+}=(\kappa\alpha)^{-1}u.
 \tag{P3K.61.SOC.2}
\]
These are exactly the \(C_0\) row of GF.11. On that same physical
fixed branch, the native \(C_1\) star has coordinates
\[
 z_1=g w,\qquad N_{1,-}=\eta v,\qquad
 N_{1,+}=(g\eta)^{-1}u.
 \tag{P3K.61.SOC.3}
\]
Indeed its native star is
\((\eta v,g w,(g\eta)^{-1}u)\); the middle coordinate is the
zero-weight tangent to this fixed branch, while the first and last
have physical weights \(-1,+1\). Therefore the literal map is
\[
 z_1=(g/\kappa)z_0,\qquad
 \begin{pmatrix}N_{1,-}\\N_{1,+}\end{pmatrix}
 =\begin{pmatrix}
 \eta/\alpha&0\\0&\kappa\alpha/(g\eta)
 \end{pmatrix}
 \begin{pmatrix}N_{0,-}\\N_{0,+}\end{pmatrix}.
 \tag{P3K.61.SOC.4}
\]
All factors here are the original functions. Relative to their fixed
reference values, the scalar phase part of the normal matrix is
\(\operatorname{diag}(\lambda/\lambda_*,\lambda_*/\lambda)\).
The canonical positive-weight complex normal structure conjugates
the negative first line. Its determinant transition is consequently
\[
 \det_+\Gamma
 =\frac{\kappa}{g}\frac{\alpha}{\overline\alpha}
                     \frac{\overline\eta}{\eta}.
 \tag{P3K.61.SOC.5}
\]
For the fixed-source scalar loop its winding is exactly \(-2\), or
\(-2M\) on the degree-\(M\) pullback. At the other outer vertex the
two products use the same first-vertex coordinates: the fixed branch
has tangent \(b\), negative normal \(a\), and positive normal \(c\).
The corresponding comparison is the identity. Thus the two outer
branch determinant comparisons have windings \(0\) and \(-2M\).

\item This is distinct data on explicitly different overlaps from the equal
inner charged-cap gauges in K25. Those inner gauges can be removed by
their one common cap frame change; that operation has not provided a
frame change on the two outer fixed branches matching (P3K.61.SOC.4). A
nonvanishing determinant homotopy on the fixed annulus, preserving the
specified endpoint products, would have the same parameter winding at
both ends. Equations (P3K.61.SOC.4)--\allowbreak (P3K.61.SOC.5) prove that such a homotopy does
not exist for \(M\ne0\).

\item In particular the K25 sentence which makes the determinant class of
the additional fixed-locus sphere boundary a pullback from its
three-dimensional parameter-free base has not carried this outer
clutch. Its asserted \(b^2=0\) does not follow from the inner common
gauge calculation. The needed correction is the actual line obtained
by adjoining (P3K.61.SOC.4) to the already written charged-cap maps, with
their oriented base cycle retained.

\par\medskip\noindent\textbf{SOC.3. Direct obstruction to K29's stated marked equivariant product.}\par\nopagebreak[4]

\item The full proof of PO.5--PO.11 applies to the closed charts in SOC.1;
its argument uses the actual three rows and arbitrary smooth
equivariant normal derivatives, not a diagonal-normal assumption.
Here is its scalar specialization in full.

\item At the first vertex write the tangent derivative factors on the
labelled \(C_2,C_1\) axes as \(A,C\), and at the second as \(U,V\).
Their parameter windings are integers \(a_*,c_*,u_*,v_*\).
For a pointed nonfixed curve whose actual coordinate transition is
\(z_-=\xi^{-1}/z_+\), an orientation-preserving equivariant product
has \(F(z_+)/z_+\) rotation-invariant. Its winding is constant along
the orbit cylinder. Comparing its two endpoint derivative limits
gives \(\operatorname{wind}P+\operatorname{wind}Q
=\operatorname{wind}\xi\). The actual \(C_2,C_1\) transitions
therefore force
\[
 a_*+u_*=0,\qquad c_*+v_*=M.
 \tag{P3K.61.SOC.6}
\]

\item On \(C_0\), the canonical positive-weight normal coordinates at the
first and second vertices are \((\overline a,c)\) and
\((\overline v,u)\). The actual normal transition is
\[
 \begin{pmatrix}\overline v\\u\end{pmatrix}
 =\begin{pmatrix}
 \overline\alpha^{-1}\overline b&0\\0&\kappa\alpha b
 \end{pmatrix}
 \begin{pmatrix}\overline a\\c\end{pmatrix},
\]
with determinant \(\kappa(\alpha/\overline\alpha)|b|^2\).
In the fixed reference product this determinant has constant phase.
An arbitrary equivariant normal derivative is an invertible complex
two-by-two map in this structure. Its determinant winding is constant
as its source point contracts to either vertex. Consequently the
single bundle-map square gives
\[
 (-v_*+u_*)+\operatorname{wind}\kappa
   +2\operatorname{wind}\alpha-(-a_*+c_*)=0.
 \tag{P3K.61.SOC.7}
\]
Substituting (P3K.61.SOC.6) and the literal scalar windings reduces the
left-hand side to \(-M\). For the positive primitive loop it is
\(-1\), a contradiction. Thus the actual labelled three-curve
plumbing admits no periodic smooth delta-equivariant product carrying
the labelled curves, vertices and complex curve orientations to
constant copies. A product connection with identity return on the
entire marked neighbourhood would supply precisely that forbidden
product.

\item This locates the failure of the premise asserted at K29 Section3.
The subsequent use of that premise to put the original scalar
variation in the form \(\operatorname{diag}(v,0,0)\), and to make
both nonfixed relative right integers zero, is unsupported by that
construction. The independently established cellular groups,
reciprocal-section basis maps, variation symmetry, and frame-change
formula remain their separate exact statements. The scalar fixed-star
product and its individual pair are also separate maps and are not
contradicted by the marked three-curve obstruction.

\end{enumerate}
\end{proof}




\begin{theorem}\label{night:P3K.62}
\StatementRecord{night:P3K.62}{proved}{S6\_06\_BLIND: GPT-6 Astra (OpenAI, ChatGPT Pro mode, effort Pro), recorded re-derivation in the checkpoint dated 18 November 2026}{\cref{night:P3K.10}, \cref{night:P3K.17}, \cref{night:P3K.22}, \cref{night:P3K.23}, \cref{night:P3K.25}, \cref{night:P3K.26}, \cref{night:P3K.28}, \cref{night:P3K.29}, \cref{night:P3K.31}, \cref{night:P3K.32}, \cref{night:P3K.47}, \cref{night:P3K.58}}

Retain the actual fixed-source whole-zero \(I_1\) scalar families,
their source and Cartier line, finite-root and annular cap maps,
physical delta circle, and specified forward parameter closing.
On their protected cusp disc the full marked periods are
\[
 h:(Z,V)\mapsto(\kappa(q)qZ,\alpha(q)V),\qquad
 B:(Z,V)\mapsto(g(q)Z,q\lambda\beta(q)V),
\]
with the displayed units independent of the scalar parameter.
Use the initial physical-vertical-edge fan and its labelled native
stars \(C_1,C_2\), with the coordinates constructed below.
For an integral scalar winding \(J\), including the actual
\(\lambda=r^M\) covers and either norm sign, the single native
relative free pairs are
\[
 (m_{C_1},n_{C_1})=(J,-J),\qquad
 (m_{C_2},n_{C_2})=(0,0).
 \tag{P3K.62.0}
\]
The actual first \(C_1\) flop changes the oriented central return
by \(-J\) in \(\Theta_7\cong\mathbb Z/28\). The actual first
\(C_2\) flop preserves that return and its integer characteristic
index. In particular the initial and reflected fan scalar
indices used in P3K.58 are equal by this explicit \(C_2\) map.

For \(C_1\), the two outer positive-Hopf normal coordinate maps
have determinant windings \(0\) and \(-2J\). This is the missing
outer line marking in P3K.25's earlier zero-correction argument.
For \(J\ne0\), the actual labelled delta-equivariant simultaneous
three-curve plumbing product invoked in P3K.29 is obstructed by
its full normal determinant square. The local topology and
single-star frame maps remain the ones used below. This statement
does not assign a universal all-fan formula.
\end{theorem}
\begin{proof}
\begin{enumerate}[label=\arabic*.]
\item \par

\medskip\noindent\textbf{SCN.1. The actual primitive scalar seam}

\par\nobreak

\item Keep the physical line character \(V\) and the original positive
\(\lambda\)-loop of K22/K25. The full periods are
\[
 h:(Z,V,q)\mapsto(\kappa(q)qZ,\alpha(q)V,q),\qquad
 B:(Z,V,q)\mapsto(g(q)Z,q\lambda\beta(q)V,q),
\]
where all displayed units are independent of the scalar parameter.
Put \(\eta=\lambda\beta\). The primitive cap-to-ordinary gluing of
K10/K22 has the literal forward parameter arrow
\[
                  R_{\rm ord}:(Z,V,q)\mapsto(Z,q^{-1}V,q).
                                                        \tag{SCN.1}
\]
Its cap arrow is the holomorphic unit recorded in K10; the exact
overlap square there is retained unchanged. Reindexing the shifted
cone by the actual period \(B\) gives
\[
                    E_\lambda:(Z,V)\mapsto(gZ,\eta V).
                                                        \tag{SCN.2}
\]
This is the K25.6 endpoint in global torus coordinates; in its
local square coordinates \((Z_{\rm loc},V_{\rm loc})=(Z,ZV)\)
it is \((gZ_{\rm loc},g\eta V_{\rm loc})\).

\item On a cut scalar interval \(\tau\in[0,1]\), let
\(\lambda=\lambda_*e^{2\pi i\tau}\). Retain the actual logarithms
of \(g,\eta\) on the \(q\)-disc times this interval, and take a
cutoff \(\chi=0\) near the first end, \(\chi=1\) near the last.
Use
\[
 C^{(1)}_\tau:(Z,V)\mapsto(AZ,DV),\qquad
 A=e^{-\chi\log g},\quad D=e^{-\chi\log\eta}.         \tag{SCN.3}
\]
On the terminal collar \(C^{(1)}=E_\lambda^{-1}\), so the
entire seam identity \(C_{\rm end}^{(1)}E_\lambda=1\) holds
with every derivative. These maps preserve \(q\), commute with
both periods and the physical circle, and have the two actual
chart expressions
\[
 (a,b,c)\mapsto((AD)^{-1}a,Ab,Dc),\qquad
 (u,v,w)\mapsto(ADu,D^{-1}v,A^{-1}w).
                                                        \tag{SCN.4}
\]
Substitution in each of DZ.8's three transition rows preserves
every coefficient. Thus these are actual closed vertex charts,
with the original units and forward \(q^{-1}\) arrow.

\item The doubled seam is specified as the literal pullback
\[
 C^{(2)}_{j+\tau}=C^{(1)}_\tau E_{\lambda(\tau)}^{-j},
                       \qquad j=0,1.                \tag{SCN.5}
\]
Its middle and terminal collar identities are respectively
\(E^{-1}\) and \(E^{-2}\). It has not been independently chosen
from the doubled endpoint.

\par

\medskip\noindent\textbf{SCN.2. The complete native C1 product and Cox lift}

\par\nobreak

\item In these closed charts the horizontal native star is given by
\[
 (z,n_0,n_1)=(c,a,b),\qquad
 (z',n'_0,n'_1)=(\eta v,g w,(g\eta)^{-1}u).
                                                        \tag{SCN.6}
\]
The full \(C_1\) transition row is
\[
              (u,v,w)=(g\eta bc,\eta^{-1}c^{-1},g^{-1}ac).
\]
Hence (SCN.6) satisfies exactly
\[
                         z'=z^{-1},\qquad n'_i=zn_i.
                                                        \tag{SCN.7}
\]
This is a single native star product with all parameter units
included. On its first chart (SCN.4) acts by
\[
                 (z,n_0,n_1)\mapsto
                    (Dz,(AD)^{-1}n_0,An_1).
                                                        \tag{SCN.8}
\]
On the second it is
\((z',n'_0,n'_1)\mapsto(D^{-1}z',A^{-1}n'_0,ADn'_1)\);
the equality (SCN.7) is retained literally.

\item For Cox coordinates with auxiliary weights \((1,1,-1,-1)\),
put \(z=a_1/a_0,n_i=a_0c_i\). The lift of (SCN.8) is
\[
 (a_0,a_1,c_0,c_1)\mapsto
       (a_0,D a_1,(AD)^{-1}c_0,Ac_1).                \tag{SCN.9}
\]
It preserves \(q=a_0a_1c_0c_1\), so its inverse has the same
\(q\) and inverse multipliers even for the full \(q\)-dependent
units. It extends through both small resolutions. In this order
the physical circle weights are \((0,1,-1,0)\). The angular
factor of (SCN.9) is the corresponding physical-circle factor;
the inverse direction from K25.7 is retained by (SCN.3).
The native product pulled back under the double parameter cover
is exactly the one using (SCN.5), including (SCN.6).

\par

\medskip\noindent\textbf{SCN.3. The doubled ordinary normal maps and trace degrees}

\par\nobreak

\item The raw period-unit winding is
\(\left(\begin{smallmatrix}0&0\\0&1\end{smallmatrix}\right)\)
for one positive primitive loop. For two turns, with
\(\lambda=\lambda_*e^{2i\theta}\), the actual based normal
ratios in the second-vertex product
\((u,\eta v,\kappa w)\) are exactly
\[
 K_0=K_2=\operatorname{diag}(e^{2i\theta},1),
 \qquad K_1=\operatorname{diag}(1,e^{2i\theta}).       \tag{SCN.10}
\]
All other units cancel against their identical reference values.
The original higher normal terms are carried by the same
normal-rescaling and cutoff-flow construction as DZ.2. Thus the
closed ordinary product uses
\[
 t=1,\qquad l=(2,0,2),\qquad
 T_i(Q)=H_1 QH_{l_i-1}.
\]
The physical weights remain \(w=(0,1,-1)\),
\(d=(-2,-1,-1)\). Substitution in the literal orbit map DZ.16
and its closed stable derivative gives
\[
                  I=(-2,2,-2),\qquad W=(2,0,0).
                                                        \tag{SCN.11}
\]
The actual capped graph-loop prism and the ambient-frame
comparison then give on \(C_1\)
\[
 V_{11}=2,\qquad\delta_1=-2,\qquad
       2n_1=\delta_1-V_{11}=-4,\qquad m_1=0.
                                                        \tag{SCN.12}
\]
The native-to-ordinary boundary arrow in this star is
\[
                    q_H\mapsto
                   (U H_1U^{-1})q_H H_{-1}.
                                                        \tag{SCN.13}
\]
Its Hopf-oriented left degree is \(+2\), and its right degree
is zero. Extending its inverse by OR.2's actual cylinder map
gives the pure-left endpoint \(L^{-2}\). The doubled native
pair obtained from (SCN.12) is therefore \((2,-2)\).

\item The product (SCN.5)--(SCN.7) is the literal pullback of the
single native product, so the same fixed-boundary connection
argument of DZ.4 identifies its relative return with \(f^2\).
Its two free coordinates consequently give the single native
pair
\[
                               (m(f),n(f))=(1,-1).
                                                        \tag{SCN.14}
\]
This uses the free coordinates of the actual relative square;
it does not divide a seven-sphere class.

\par

\medskip\noindent\textbf{SCN.5. The actual outer normal marking omitted by K25}

\par\nobreak

\item The following calculation identifies the outer normal-line mismatch.

\item Keep every unit of SCN.1. The native fixed-curve \(C_0\) charts
have first base \(b\), normal pair \((a,c)\), and second base
\(\kappa w\), normal pair
\[
             \bigl(\alpha v,(\kappa\alpha)^{-1}u\bigr).
                                                        \tag{SCN.15}
\]
The negative physical weight is the first normal coordinate
in both ordered pairs. In the actual \(C_1\) native chart
\((z,n_0,n_1)=(c,a,b)\), the fixed branch is
\(z=n_0=0\), with base \(n_1=b\) and normal pair \((a,c)\).
Thus its first endpoint agrees literally with (SCN.15).

\item In its second chart
\[
              (z',n'_0,n'_1)
                   =(\eta v,g w,(g\eta)^{-1}u),
\]
the fixed branch is \(z'=n'_1=0\), with base \(n'_0=g w\)
and ordered negative/positive normal pair
\[
                    \bigl(\eta v,(g\eta)^{-1}u\bigr).
                                                        \tag{SCN.16}
\]
Consequently the actual coordinate map from the native
\(C_0\) normal pair to this \(C_1\) pair at the second branch
has base factor \(g/\kappa\) and normal matrix
\[
       R=\operatorname{diag}
          \left(\frac{\eta}{\alpha},
                       \frac{\kappa\alpha}{g\eta}\right).
                                                        \tag{SCN.17}
\]
This is an equality of the displayed coordinates, with no
change of period or omission of a bounded multiplier.
The first branch has the identity normal matrix.

\item Divide (SCN.17) by its reference value at \(\lambda_*\).
The units \(\kappa,\alpha,g,\beta\) are independent of the
scalar parameter. Hence on the positive fixed-radius loop,
\[
        R R_*^{-1}
           =\operatorname{diag}(e^{i\theta},e^{-i\theta}).
                                                        \tag{SCN.18}
\]
The physical positive-weight complex normal structure
conjugates the first coordinate. Thus the actual Hopf
spinor coordinate map in (SCN.18) is
\[
                         e^{-i\theta}I_2.             \tag{SCN.19}
\]
At the other endpoint it is \(I_2\). Its determinant winds
\(-2\) at the second endpoint and \(0\) at the first.
For \(\lambda=r^m\) these windings are \(-2m\) and \(0\).
They have been calculated from the actual arrows, not
selected from SCN.14 or from a desired flop sign.

\item K25.7 and K25.8 compare the two \emph{inner} charged-cap maps.
Their identical physical-circle gauge can be removed
there by one cap frame. Equations (SCN.15)--(SCN.19)
compare those native master charts to the two \emph{outer}
fixed-curve caps. They give different endpoint gauges.
Multiplying both by a common parameter scalar leaves
their ratio \(e^{-i\theta}\) unchanged. Therefore the
inner common-gauge calculation alone does not establish
the claimed pullback of the outer normal-sphere line.
The precise unsupported sentence is K25's assertion,
after identifying the four fixed-boundary charges,
that its parameter circle is a separate factor and
consequently \(b^2=0\).

\par

\medskip\noindent\textbf{SCN.6. The integral Cech class with an explicit orientation}

\par\nobreak

\item Here is the line class of the actual mismatch, separately
from the orientation of the full characteristic bordism.
Let a spatial circle be the union of a master interval
from endpoint \(0\) to endpoint \(1\) and a cap interval
traversed from endpoint \(1\) back to endpoint \(0\).
Orient its parameter circle by increasing \(\theta\).
Use the coordinate arrows from the master frame to the
cap frame
\[
        y_{\rm cap}=g_j(\theta)y_{\rm master},\qquad
                 g_0=1,\quad g_1=e^{-im\theta}.
                                                        \tag{SCN.20}
\]
These are (SCN.19), with the indicated arrow and ordering.
They determine a line \(L_m\) on the resulting two-torus.
In this explicitly prescribed spatial orientation,
\[
       \langle c_1(L_m),[S^1_{\rm spatial}
                              \times S^1_\theta]\rangle=-m.
                                                        \tag{SCN.21}
\]
One direct verification keeps the connection arrow.
Set \(A_{\rm master}=0\). For the convention
\(\nabla=d+A\), equation (SCN.20) forces
\(A_{\rm cap}=-dg_j\,g_j^{-1}\), hence \(0\) at endpoint
\(0\) and \(im\,d\theta\) at endpoint \(1\).
Parametrize the cap in the temporary direction \(0\to1\)
and interpolate with \(A_{\rm cap}=im\rho(s)d\theta\),
where \(\rho\) is constant \(0\) and \(1\) on the
respective collars. Its curvature integral in that
temporary direction is \(m\). The oriented spatial
circle traverses the cap in the opposite direction,
giving (SCN.21). This also gives the integral Cech
calculation without differentiating a multivalued log.
The positive Hopf normal two-plane has the scalar
transition (SCN.20) on both summands; its determinant
crossing class is therefore \(2c_1(L_m)\), with
evaluation \(-2m\).

\item This calculation identifies a nonzero relative line
class and fixes its sign in the displayed two-torus
orientation. We have not yet identified that oriented
two-torus with the oriented four-dimensional fixed
base in K25's full bordism, nor computed its sphere
class or the resulting value of \(b^2\). No numerical
characteristic correction is assigned by SCN.21.

\par

\medskip\noindent\textbf{SCN.7. The simultaneous-product premise also fails directly}

\par\nobreak

\item Apply the complete determinant square PO.5--PO.11
of K47 to the actual scalar charts (SCN.4), retaining
their labelled vertices and all three labelled curves.
The pointed \(C_2\) and \(C_1\) products force
\[
                 a_*+u_*=0,\qquad c_*+v_*=m.
\]
The arbitrary invertible positive-weight complex
normal derivative along \(C_0\) forces
\[
      (-v_*+u_*)+\operatorname{wind}\kappa
            +2\operatorname{wind}\alpha-(-a_*+c_*)=0.
\]
Substitution yields
\(\operatorname{wind}\kappa+2\operatorname{wind}\alpha
 -\operatorname{wind}\eta=-m=0\), which is impossible
for \(m\ne0\). The argument allows full off-diagonal
\(GL_2(\mathbb C)\) normal matrices; it does not assume
diagonality away from the vertices. Thus the actual
marked delta-equivariant simultaneous plumbing product
asserted in K29 Section 3 does not exist for this
positive scalar loop. The separate native \(C_1\)
product of SCN.6--SCN.9 and the ordinary quaternion
product are not excluded by this determinant square.

\item The exact local maps and integral mismatch are recorded
by \path{code/part3/scalar_c1_outer_marking.py} and its output.
The dependency notice is kept separately in
\path{notes/part3_scalar_consistency_dependencies.md}.

\par

\medskip\noindent\textbf{SCN.8. The complete native C2 comparison}

\par\nobreak

\item The diagonal \(C_2\) star in these same seam-closed
charts has the actual coordinates
\[
        (z,n_0,n_1)=(a,b,c),\qquad
        (z',n'_0,n'_1)=(u,v,w).
\]
Its full transition is
\[
       (u,v,w)=(a^{-1},ab,ac),\qquad
                 z'=z^{-1},\quad n'_i=zn_i.           \tag{SCN.22}
\]
There is no omitted unit or parameter factor in this
transition. The seam (SCN.4) acts on its first chart by
\[
       (z,n_0,n_1)\longmapsto
                    ((AD)^{-1}z,An_0,Dn_1)
\]
and on the second by
\((z',n'_0,n'_1)\mapsto(ADz',D^{-1}n'_0,A^{-1}n'_1)\).
With \(z=a_1/a_0,\ n_i=a_0c_i\), its full Cox lift is
\[
   (a_0,a_1,c_0,c_1)\longmapsto
                (a_0,(AD)^{-1}a_1,Ac_0,Dc_1).        \tag{SCN.23}
\]
It preserves \(q=a_0a_1c_0c_1\), the auxiliary action,
the physical circle and every actual period translate.
It extends over both small resolutions and has the
displayed inverse at the same \(q\).

\item For two primitive scalar turns, SCN.11 and the actual
prisms DZ.19, with \(w_2=-1,w_0=0\), give
\[
                    V_{22}=2,\qquad \delta_2=2.
\]
The local ambient-to-product frame comparison therefore
gives the ordinary relative pair \((-2,0)\).
The native-to-ordinary boundary map is literally
\[
                q_H\longmapsto
                 (UH_1U^{-1})q_HH_1.                 \tag{SCN.24}
\]
Its left degree is \(+2\) and its right degree is zero.
The inverse extends by the explicit cylinder retraction
of OR.2, with endpoint \(L^{-2}\); its coefficient maps
extend over \(D^3\times D^4\). Thus the native doubled
pair is \((0,0)\).

\item Equations (SCN.5), (SCN.22) and (SCN.23) are the literal
pullback of one native \(C_2\) product. Consequently
the doubled relative return is its actual square.
The free coordinates of its single relative return
are \((0,0)\), since multiplication by two is injective
on these two integral coordinates. The actual
single-star Cox conjugation changes the central
seven-sphere class by its right integer, which is
zero. This establishes equality of the initial and
reflected scalar return classes without division in
\(\Theta_7\).

\item The same maps retain every integral scalar winding
\(J\). On the double pullback use
\[
       t=J,\quad l=(2J,0,2J),\quad
       I=(-2J,2J,-2J),\quad W=(2J,0,0).
                                                        \tag{SCN.25}
\]
The ordinary \(C_2\) pair is \((-2J,0)\), its native
pair is \((0,0)\), and the single native pair is
again \((0,0)\). This includes negative winding and
zero winding, with the same physical circle.

\item For the canonical cap families the ordinary closing
may be \(q^e\) with either sign of \(e\), rather than
the primitive cap-to-ordinary \(q^{-1}\) of SCN.1.
The exact reindexing is then \(B^{-e}\), and its
bounded map is
\((Z,V)\mapsto(g^{-e}Z,\eta^{-e}V)\).
Use the cut seam
\[
       (Z,V)\mapsto
         \bigl(e^{e\chi\log g}Z,e^{e\chi\log\eta}V\bigr).
                                                        \tag{SCN.26}
\]
It obeys \(C_{\rm end}E=1\) with all derivatives.
Its double is still the literal pullback (SCN.5).
Thus the full source and cap closing squares already
proved in K26/K28/K32 are carried unchanged; no norm
sign is reversed. Parameter-independent bounded
closing units are carried through their fixed
logarithms. The unit winding in the actual period
matrix is still \((0,0;0,J)\), so (SCN.25) applies.
The scalar equality needed by K58 for its initial
and reflected fans follows from these actual maps.

\par

\medskip\noindent\textbf{SCN.9. The C2 characteristic bordism has no outer mismatch}

\par\nobreak

\item There is also a direct narrow repair of the integer
characteristic comparison. At the first fixed branch
of the \(C_2\) native star, its fixed base and normal
pair are exactly \(b\) and \((a,c)\). At the second,
they are \(w\) and \((v,u)\). Relative to the native
\(C_0\) coordinates (SCN.15), the second arrow is
\[
 \text{base factor }\kappa^{-1},\qquad
 \text{normal matrix }\operatorname{diag}
                           (\alpha^{-1},\kappa\alpha).
                                                        \tag{SCN.27}
\]
Both this map and its first-branch identity are
independent of the scalar parameter. After carrying
their fixed logarithms, the outer normal maps are
constant in that parameter.

\item In the \(C_2\) Cox coordinates of (SCN.23), the
parameter factor has physical weights
\((0,-1,0,+1)\). The two inner charged boundary maps
therefore have one identical scalar gauge, as in
the actual master and joint-cap maps of K17.
Use that same scalar on every Cox chart and on the
joint cap. It commutes with their actual transition
functions. The outer identifications (SCN.27) are
constant in the parameter, so this gauge matches
the two outer conifold products as well. Unlike
SCN.17, there is no endpoint-dependent parameter
transition left on the fixed-cobordism closure.

\item The line on the closed five-dimensional orbit
correction is consequently the pullback of its
parameter-zero line. The same is true of its
tangent bundle under the literal moment-chart
maps and constant bounded-unit isotopies. Its
six-dimensional numbers \(x^3,p_1x\) vanish.
On the additional normal-sphere boundary the
closed four-dimensional fixed base is the
oriented \(S^2\) mapping torus crossed with the
parameter circle, with all its normal and line
transitions pulled back from the three-dimensional
factor. Thus its actual class \(b\) has \(b^2=0\),
and the base tangent Pontryagin number is zero.
The sphere-bundle identity
\[
       \frac{4\int x^3-\int p_1x}{24}
                 =\frac{3\int b^2-\int p_1(T\bar X)}{24}
\]
gives zero for this retained boundary. This proves
equality of the two integer characteristic indices
for this particular \(C_2\) comparison. It proves
no equality for the \(C_1\) comparison, whose
actual outer map is (SCN.17).

\par

\medskip\noindent\textbf{SCN.10. The full actual single-star flop and its central coordinate}

\par\nobreak

\item For each of SCN.9 and SCN.23, transport the Hermitian
metric by the displayed invertible Cox map and then
apply its positive radial normalization on the common
link. Its inverse is literal at the unchanged height
\(q\), so this construction retains all bounded
moduli and all parameter collars. The unrestricted
four-weight identity FF.1/SF.8 intertwines the two
resolution links with
\[
  \tau(h,q_H)=(q_H^{-1}h q_H,q_H),\qquad
  T(y,q_H)=(q_H^{-1}yq_H,q_H).                         \tag{SCN.28}
\]
Both are independent of the parameter in the native
products. The normal-rescaling homotopy carries the
full \(q\)-dependent unit germs, rather than only
their central linear terms. The coordinate maps
commute with every actual period translate, so one
free representative extends to the full periodic
orbit of disjoint flop curves. The source and cap
maps off these matched charts stay the original ones.

\item The first and second small resolutions in this arrow
use the same \(q_H\). The second complex-normal
convention is restored by the fixed anti-quaternion
map of K23/K31; its oriented central coordinate is
unchanged. In these actual conventions the relative
map is the literal conjugation
\[
 T(L^\mu R^\nu\mathcal C^s)T^{-1}
          =L^{-\nu}R^{\mu+2\nu}\mathcal C^{s+\nu}.
                                                        \tag{SCN.29}
\]
The unrestricted four-weight calculation is needed
here because each seam acts on all four Cox
coordinates, not only one pair of normal weights.

\item For scalar winding \(J\), SCN.25 gives ordinary
\(C_1\) pair \((0,-2J)\) on the double cover. Its
native boundary change has left degree \(2J\) and
right degree zero, yielding native double pair
\((2J,-2J)\). The connection pulled back from the
single native product gives the actual relative
square, exactly as in SCN.8. Its single free pair
is therefore \((J,-J)\). Equation (SCN.29) changes
the actual central return by \(-J\) for this first
\(C_1\) flop. For the first \(C_2\) flop the single
free pair is \((0,0)\), and the central change is
zero. No central class is divided by two.

\end{enumerate}
\end{proof}




\begin{theorem}\label{night:P3K.63}
\StatementRecord{night:P3K.63}{proved}{S6\_06\_BLIND: GPT-6 Astra (OpenAI, ChatGPT Pro mode, effort Pro), recorded re-derivation in the checkpoint dated 18 November 2026}{\cref{night:P3F.54}, \cref{night:P3F.9}, \cref{night:P3K.42}, \cref{night:P3K.66}, \cref{night:P3K.68}}
Retain the actual selected-cusp source loop and physical vertical-edge
fan \(T_n=T(e_1+ne_2,e_2)\). Its Pontryagin pairing is zero for
both the physical circle line \(M\) and the meridian-one dual \(U=M^*\).
With \(A_U=(4k_U-\ell_U)/24\), the exact native and ordinary identities are
\[
(m_0,n_0)=(0,-6A_U),\qquad
V_{00}=6A_U-6,\qquad \delta_0=-6A_U-6.
\]
The full integer line/cochain proof P3K.66 gives
\[
(m_0,n_0)=(0,-6-18n(n-1)),\quad
V_{00}=18n(n-1),\quad \delta_0=-12-18n(n-1).
\]
The proof retains the actual quotient-piece, cap, normal and line
maps, and computes the entire boundary tangent transgression.
\end{theorem}
\begin{proof}
\begin{enumerate}[label=\arabic*.]
\item \par

\medskip\noindent\textbf{SQ.1. The two actual quotient pieces}

\par\nobreak

\item Write the original periods on the protected cusp disc as

\[
 h:(Z,V)\mapsto(\kappa qZ,\alpha V),\qquad
 B:(Z,V)\mapsto(gZ,\eta qV),\qquad Q=\kappa q.
                                                        \tag{SQ.1}
\]

\item The functions are the full ones in F54 and GF.1. On the positive \(\rho\)
circle their windings are respectively -6,0,0,0. In particular g,
\(\alpha\) and \(\eta\) have their actual periodic logarithms. The new fan is
described by the exact character and period arrows

\[
 Z_n=Z,\quad V_n=Z^{-n}V,\qquad h_n=hB^n,\quad B_n=B.
\]

\item Their complete multipliers, including the ones with nonzero winding,
are

\[
 Q_n=Qg^n,\quad
 \alpha_n=\alpha\eta^n\kappa^{-n}g^{-n^2},\quad
 g_n=g,\quad\eta_n=\eta g^{-n}.                       \tag{SQ.2}
\]

\item Here \(h_{n}\) acts by (\(Q_{n}\) \(Z_{n}\),\(\alpha_{n}\) \(V_{n}\)) and \(B_{n}\) by
(\(g_{n}\) \(Z_{n}\),q \(\eta_{n}\) \(V_{n}\)). The coordinate q is unchanged. This is the
literal GF.1--3 arrow, rather than a change of the height character.

\item In this note \(\gamma_{n}\) is the real coefficient of the contracting B.
The increasing F9 moment coordinate has the outward orientation
-\(\gamma_{n}\). We carry this explicit minus on both compared pieces;
we do not reverse the physical Hopf normal orientation.

\item Take the \(\delta\) quotient of this cusp piece. The proof of F9 applies
in these new vertical characters. Its node charts are the actual
conifold quotient charts (x,y,b),(x,y,b'), with

\[
 2b=r-|y|^2/r,\quad 2b'=|x|^2r-r^{-1},\qquad
 \frac{db'}{db}=\frac{|x|^2+r^{-2}}{1+|y|^2r^{-2}}>0.    \tag{SQ.3}
\]

\item Its proper increasing outward coordinate -\(\gamma_{n}\) is constructed relative to the
specified exterior coordinate, by the section and positive-vector-field
construction in F9. It is therefore a morphism over the actual source
node, with positive upper-triangular stable derivative. The same
construction with the complete positive norms of (SQ.2) is periodic
in \(\rho\): logarithms of positive norms are single-valued even when the
complex unit has nonzero phase winding. A phase of \(\alpha_{n}\) is a \(\delta\)
rotation in the line factor, so it is absent from this orbit coordinate.
It remains in the principal-circle line, which we compute below.

\item The source part of B has the periodic logarithm log g. Its translation
isotopy exp(t log g), together with its positive line-norm isotopy,
untwists the suspension exactly as in F9. Thus the actual orbit piece,
with its marked outer boundary, is

\[
                     W_n=E_n\times S^1_{\gamma_n},       \tag{SQ.4}
\]

where \(E_{n}\) is a compact nodal elliptic neighbourhood with Tate
parameter \(Q_{n}\). This statement also describes its \(\rho\) family. We use
the actual radial coordinate supplied by (SQ.3), and transport its
metric; we do not set the line phase \(\alpha_{n}\) to one.

\item We compute the endpoint lift in this prescribed boundary marking.
Put \(L_Q=\log|Q_n|\), \(L_a=\log|\alpha_n|\), \(L_g=\log|g|\) and
\(L_S=\log|q\eta_n|\). On the regular annulus the actual real period
coefficients are obtained by inverting the absolute-value period matrix.
Its \(\gamma\) coefficient is exactly

\[
 \gamma_n=
 \frac{L_Q\log|V_n|-L_a\log|Z_n|}{L_QL_S-L_aL_g}.
                                                        \tag{SQ.4a}
\]

\item For a sufficiently small common q disc the denominator is nonzero:
its leading term is \((\log|q|)^2\) and the remaining unit terms are
uniformly bounded on the compact parameter circle. Substitution of
\(h_{n}\) adds (\(L_{Q}\),\(L_{a}\)) to the two logarithms, hence leaves (SQ.4a)
unchanged. Substitution of B adds (\(L_{g}\),\(L_{S}\)), hence increases it by
exactly one. This identifies the actual lifts of both deck arrows.
In particular the continuation along the source normalization uses
zero \(\gamma\) deck index under \(h_{n}\). The n powers of B present in the
old expression \(hB^n\) have been canceled by the actual character
\(V_n=Z^{-n}V\); they cannot be added again as an endpoint integer.

\item In the square model whose two period directions are \(h_{n}\) and B,
cut the second nodal orbit circle at a free point. The F9 conifold
above its node is one line piece between the two adjacent free
ends of this cut. Over the source node its exceptional interval has
the exact representatives

\[
 z=\sqrt r,\quad v=0,\quad u=0,\quad w=1/\sqrt r,
                         \qquad0<r<\infty .             \tag{SQ.4b}
\]

\item Its fixed endpoints are r=0 in the first chart and r=infinity in
the second. They have b=0 and \(b_{prime}\)=0 respectively; the section
of F9 is r=1. On the line-factor nodal model the positive-weight coordinate is
u=\(Z_{n}\) \(V_{n}\). The actual B multiplies it by q g \(\eta_{n}\), whose
norm is less than one after the chosen shrinking. Hence B runs from
the u branch toward the v branch, opposite to the increasing F9
moment direction. On this interval choose the correctly oriented
contracting-period coordinate

\[
 \gamma_n=\tfrac12-\epsilon\frac{r-1}{r+1},
                       \qquad0<\epsilon<\tfrac14.       \tag{SQ.4c}
\]

\item Thus its actual endpoint lifts are 1/2+epsilon and 1/2-epsilon.
In the first chart its derivative with respect to b at r=0 is
-4epsilon; in the second its derivative with respect to \(b_{prime}\) at
r=infinity is also -4epsilon. Thus the specified coordinate extends
smoothly with negative derivative at both fixed endpoints; its
outward negative has positive derivative, exactly as in F9.
They lie between the same two successive B translates. To verify
that cell statement, use the literal initial square in the new
characters: its source sides are identified by \(h_{n}\), whose \(\gamma\)
lift is zero by (SQ.4a), and its suspension sides by B, whose lift
is one. More explicitly, the two conifold endpoint charts are \(T^+_0\) and
\(T^-_0\). The fixed \(C_{0}\) sphere uses \(T^+_0\) and
\(T^-_{-e_1}\); the latter is the \(h_{n}\) inverse translate of the
second conifold endpoint. Its \(\gamma\) lift is unchanged by (SQ.4a).
Thus these are exactly the two endpoint lifts of the same fixed
normalization sphere. The small-resolution interval (SQ.4b) lies in the interior
of that one square and joins its two fixed branches; there is no
further period arrow on this interval. The two infinite half-lines
in the F9 quotient continue to the chosen free cuts in the
F9 increasing outward order. The full nonzero units change their positive
coordinates but not these deck indices or this order. In particular
no integer has been suppressed in (SQ.4c).

\item Extend (SQ.4c) to a proper decreasing coordinate on those half-lines,
and repeat it by \(\gamma_{n}\) to \(\gamma_{n}\)+1 on B translates. It agrees
with the prescribed real-period coordinate after a degree-one change in the marked circle orientation
on the overlap. Interpolate the lifted decreasing F9-coordinate
maps by their convex combination, with a partition equal to one
on the protected outside collar. All interpolating derivatives
remain negative in the F9 moment direction, and their lifts respect translation by one.
This is the explicit relative coordinate construction of F9, now
with its interval and its two endpoint indices specified. The
complete positive unit norms are periodic in \(\rho\); (SQ.4a) has no
\(\rho\) reindex, so this relative construction is periodic as well.
On the remaining normalization cylinder choose its graph height
between these two values and match it by the same relative
increasing-coordinate construction. The source normalization is
one-to-one there, so these choices impose no further identification.
The \(h_{n}\) end identification has the zero lift already computed.

\item Consequently the complete fixed sphere graph has a real lift in the closed short arc
{[}1/2-epsilon,1/2+epsilon{]}, with the two displayed endpoint values.
For two such graphs, their convex interpolation has endpoint
difference -2epsilon throughout; it never reaches any integer.
This is the endpoint verification needed in SQ.2, rather than an
inference from the existence of a real lift on a sphere alone.

\par

\medskip\noindent\textbf{SQ.2. The actual source map, fixed tube and cap}

\par\nobreak

\item Put \(L_g=\log g\), with the periodic branch just specified. Consider the
source family with

\[
                      Q_s=Q\exp(snL_g),\qquad0\le s\le1.
                                                        \tag{SQ.5}
\]

\item For this fixed integer n, shrink the protected q disc uniformly in s
and the compact \(\rho\) circle so that all \(Q_{s}\) lie in the same Tate
annular range and no second critical value enters the disc.
Its nodal charts are \(xy=Q_s\). They are smooth total source charts,
since q times a unit is \(Q_{s}\). Near the node the vector field

\[
                 \partial_s+nL_g\,y\partial_y           \tag{SQ.6}
\]

is tangent to this total space and keeps q and \(\rho\) fixed. In this
formula the y derivative is taken in the displayed complex chart;
the real vector field is its underlying real derivative. It preserves
the node locus and the two branches. On the outside source torus
write a point as \(u+\tau_s v\) with u,v real modulo integers and
\(\tau_s=\log Q_s/(2\pi i)\). Prescribe the actual real marked-period map

\[
              u+\tau_s v\longmapsto u+\tau_0v.           \tag{SQ.7}
\]

\item It gives a smooth horizontal field on an outside collar. It respects
both monodromies: a q turn sends (u,v) to (u-v,v), and a \(\rho\) turn
sends it to (u+6v,v); \(L_{g}\) is periodic. On the overlap, both fields
project to the same vector \(partial_{s}\), keep q and \(\rho\) fixed and are
tangent to the fibres. Their difference is vertical. A partition of
unity on the total source family, equal to one on a smaller node
neighbourhood and on the protected outside collar respectively,
therefore interpolates them in this affine space of horizontal
fields. No singular fibre is crossed by an untyped vector field.
The source is compact with its boundary, and the field is tangent to
that boundary. Its flow is consequently a diffeomorphism \(E_{n}\) to \(E_{0}\),
periodic in \(\rho\) and with boundary exactly (SQ.7). This flow, rather
than the generally nonsmooth real period formula at the node, is the
map used in the interior.

\item On the central normalization this flow is an orientation-preserving
sphere isotopy which preserves its two marked branch points. Transport
\(F_{n}\) by it. Its height and the height of \(F_{0}\) are real functions in the
same short arc, and their two node-preimage values have the same strict
order. Their convex interpolation therefore remains an embedding:
the source normalization is injective away from that one pair, and
the two heights at the pair remain different. Extending the velocity
from this compact sphere to its tube and then by a cutoff gives an
ambient isotopy, relative the outside boundary. All data are periodic
in \(\rho\), so its vector field and flow are periodic as well.

\item We retain the normal arrow, not just the isotopy of the sphere. In
the branch notation x=s-\(\rho\),y=s+\(\rho\),Z=x/y, the shear spinor and its
Cartier morphism are precisely SFIX.5--8:

\[
 E_n^{\rm spin}=E\otimes\mathcal O(n(D_+-D_-)),\qquad
       E_n^{\rm spin}\longrightarrow E,\quad v\longmapsto Z^n v.
                                                        \tag{SQ.8}
\]

\item The source generators \(x^{-n-1}y^n,x^{-n}y^{n+1}\) go to \(1/x,y\).
Their determinant trivializations also match. The Hom of the two
summands and the adjoint normal bundle are thus carried by the
literal typed bundle isomorphism. The physical q equations
\(q=Zc_0c_1=Z^{-1}c'_0c'_1\) are retained; their differential is the
SP.11b normal arrow. The source flow contributes its actual derivative
and the graph interpolation contributes the positive triangular
graph-normal shear. They are families starting at the identity, so
their normal comparisons extend through the corresponding isotopies.
No phase of \(\kappa\) is removed from a normal coordinate.

\item Use these arrows on the normal tube and extend the sphere isotopy to
the three-ball cap by its time-dependent boundary vector field with
a collar cutoff. On the Hopf sphere factor keep the adjoint arrow
of (SQ.8). This is exactly the product-cap construction SP.13 with
the transported base sphere, not a substitution of its tangent
derivative for the Cox frame. It supplies an oriented diffeomorphism
between the completed quotient pieces

\[
 C_n=\bigl(W_n\setminus\operatorname{int}(F_n\times D^3)\bigr)
                      \cup(D^3\times S^2)
                 \longrightarrow C_0 .                 \tag{SQ.9}
\]

\item It extends over the \(\rho\) family. The circle line is carried separately
below; (SQ.8) retains its tensor factor rather than identifying two
unmarked circle lines.

\par

\medskip\noindent\textbf{SQ.3. The exact boundary shear and the Euler line}

\par\nobreak

\item On the boundary use real coefficients (u,v,d,\(\gamma\)) of the marked
source vanishing period, source period h, physical \(\delta\) period and
suspension period B. The source point before suspension untwisting is
u+\(\tau\) v+G \(\gamma\), where exp(2pi i G)=g. The same point in the new
basis is

\[
 u+(\tau+nG)v+G(\gamma-nv).
\]

\item The full integral coefficient arrow, including the line character,
is

\[
 u_n=u,\quad v_n=v,\quad
 \gamma_n=\gamma-nv,\quad d_n=d-nu.                   \tag{SQ.10}
\]

\item It follows directly by substituting \(Z_n=Z,V_n=Z^{-n}V\) into the
period matrix. It is the boundary gluing between (SQ.9) and the
unchanged exterior. In particular its \(\gamma\) shear is not asserted to
extend over the old exterior.

\item Let \(\alpha_{q}\) and \(\theta\) have period one. Retain the full Tate phase

\[
 \alpha_Q=\alpha_q-6\theta+\phi_\kappa,
 \qquad \phi_\kappa=(2\pi)^{-1}\operatorname{Im}f_\kappa,
                                                        \tag{SQ.11}
\]

where \(f_{kappa}\) is the actual periodic logarithm after its winding has
been removed. Then du+\(\alpha_{Q}\) dv is a global source one-form on the
boundary. The q monodromy u to u-v and \(\rho\) monodromy u to u+6v
are both retained by this formula.

\item First display the integral part of the circle connection. In the
old coefficients it is

\[
 A_0=dd+\alpha_q\,d\gamma.
\]

\item In the new coefficients the phase winding of \(\alpha_{n}\) is 6n and
that of \(\eta_{n}\) is zero. Its connection is therefore

\[
 A_n=dd_n+\alpha_q\,d\gamma_n+6n\theta\,dv_n.
                                                        \tag{SQ.12}
\]

\item For a \(\rho\) turn \(d_{n}\) changes to \(d_{n}\)-6nv, and for a q turn it changes
to \(d_{n}\)-\(\gamma_{n}\). Thus these are connections on the actual integral
circle bundles, with the displayed transition functions. Substitution
of (SQ.10) gives

\[
 A_n-A_0=-n\{du+(\alpha_q-6\theta)dv\}.             \tag{SQ.13}
\]

\item The full unit factors have not disappeared. Their connection terms
are obtained from their actual periodic logarithms

\[
 f_{11}=f_\kappa+nf_g,\quad
 f_{21}=f_\alpha+nf_\eta-nf_\kappa-n^2f_g,\quad
 f_{12}=f_g,\quad f_{22}=f_\eta-nf_g.                \tag{SQ.14}
\]

\item In the actual flat real-period coordinates these are globally defined
one-form corrections, including d(arg g) and d(arg \(\eta\)). Accordingly
(SQ.13) becomes -n(du+\(\alpha_{Q}\) dv) plus that explicitly retained
global one-form correction. Its curvature difference is the exterior
derivative of that expression. We use these full connections below.

\item Here is the interior class, with no assumption that b is a relative
class. The source E retracts onto its nodal fibre \(S^{2}\) wedge \(S^{1}\).
Let b denote the graph generator of \(H^{1}(E)\), with boundary value the
actual coefficient v. One concrete extension is given by the Lefschetz
handle presentation of E: start with the circle map v on the regular
torus product. The attaching vanishing circle is the a circle, on
which v is constant. Its real lift on the attaching handle collar
extends over the two-handle, with a cutoff agreeing on the collar.
This produces the required source map, with its actual outer value.
The parameter monodromy (u,v) to (u+6v,v) preserves that value and
the same handle extension; the marked branches return literally,
so the construction is periodic. On its normalization cut,
the line with constant factor \(\alpha_{n}\) is glued by \(\alpha_{n}\) between the
two branches. Formula (SQ.2) makes the parameter winding of that
actual gluing 6n. The line has no vanishing-cycle factor and the
suspension factor \(\eta_{n}\) has zero parameter winding. Thus its
parameter-dependent Chern class differs from the old one by

\[
                         6n\,[d\theta]\smile b.         \tag{SQ.15}
\]

\item One can realize this class directly by cutting in \(\theta\): its circle
coordinate is glued by d to d-6n b at the two ends. On the graph cut
the equivalent gluing is exp(2pi i 6n \(\theta\)). The equality of these
two clutch descriptions is the square with change of coordinate
d to d+6n \(\theta\) b on its doubly cut rectangle. All zero-winding
terms are the actual logarithms in (SQ.14).

\item The same calculation holds after fixed-sphere replacement. To see the
extension, the graph map b restricted to the normalization sphere has
a real lift. Its \(\rho\) winding is zero, since the marked normalization
and branches return literally. Modify that circle-valued map in a
small tube by a compactly supported real homotopy so it is constant
on a smaller fixed tube. It now extends constantly over \(D^{3}\) times \(S^{2}\).
For the line arrow, if b and its modified map \(b_{prime}\) differ by the
supported real function h, change the fibre coordinate on the \(\theta\)
cut by exp(2pi i 6n \(\theta\) h). Its two endpoints give precisely the
change of the displayed \(\theta\) clutch. The modified clutch is constant
on the cap, where this line extends. Thus this supplies an actual
representative of (SQ.15) on \(C_{0}\), including the cap identification. The
meridian Chern number of the new and old dual lines \(U\) is
the same positive one, by (SQ.8); their difference there is zero.
The retained ordinary b times \(\gamma\) torus has Chern number zero:
on a regular source fibre the actual total torus is the quotient of
its real period lattice and its physical \(\delta\) sublattice is primitive,
so the chosen integral period basis gives a section of that torus
bundle. The \(\gamma\) times \(\theta\) difference is also zero by \(\eta_{n}\)'s
zero winding. These maps identify the dual-line difference with
(SQ.15), up to the exact curvature terms already retained in (SQ.14)
and its prescribed outer gauge (SQ.13).

\item For an integral group check, \(H_{1}\)(W)=Zb plus Zgamma and
\(H_{2}\)(W)=ZF plus Z(b times \(\gamma\)). The Thom exact sequence for removal
of F adds its normal meridian in degree two: \(H_{3}\)(W) to \(H_{0}\)(F) is
zero, since a regular fibre times \(\gamma\) is disjoint from F. The
Mayer--Vietoris map for the cap identifies its meridian with that
one and kills the primitive F. Consequently

\[
 H_1(C)=\mathbb Z^2,\qquad
 H_2(C)=\mathbb Z\langle\text{meridian},b\times\gamma\rangle.
                                                        \tag{SQ.16}
\]

\item Both are torsion free. The parameter monodromy fixes b,\(\gamma\) and
the oriented meridian. It cannot add a meridian to b times \(\gamma\):
the actual global Euler line has evaluation zero on that regular
torus and one on the meridian. This gives the two parameter classes
b times \(\theta\) and \(\gamma\) times \(\theta\) in the Wang sequence and
confirms the preceding class computation without dropping torsion.

\par

\medskip\noindent\textbf{SQ.4. The boundary tangent connection and its exact transgression}

\par\nobreak

\item Use the invariant boundary coframe

\[
 du+\alpha_Qdv,\quad dv,\quad d\gamma,\quad
 d\alpha_q,\quad d\theta,
                                                        \tag{SQ.17}
\]

together with the outward normal. This is a frame on the actual
nilmanifold times the two retained circles. It retains the \(\rho\)
monodromy in (SQ.11). Use the corresponding flat tangent connection
on a whole boundary collar; a tangent connection need not be the
Levi--Civita connection to represent p1.

\item We prescribe these actual real period coordinates on the entire
outside collar before performing the F9 construction. They are smooth
coordinates on the regular real torus quotient, formed from the full
real period matrix of (SQ.1)--(SQ.2). On that collar the marked line coordinate is precisely its
contracting \(\gamma_{B}\) coefficient, including the
actual norms of \(\alpha\) and \(\eta\). In the convex combination of proper decreasing F9-coordinate
maps, equivalently proper increasing maps in -\(\gamma_{B}\) on both sides,
take the outside partition weight identically
one on this collar. The actual comparison is orientation-preserving between those two
marked coordinates. The resulting orbit product is therefore exactly
this prescribed coordinate there. It is not merely isotopic to it.
The source horizontal field in SQ.2 was likewise prescribed there by
(SQ.7), so it is the identity on the marked real coefficients u,v.
The variables q and \(\rho\), including the outward radial variable r,
are fixed.

\item Order the stable coframes as
\((du+\alpha_Q\,dv,dv,d\gamma,d\alpha_q,d\theta,dr)\), with \(dr\) outward.
Put \(\phi_g=\operatorname{Im}(L_g)/(2\pi)\). The pullback of the target coframe by the
actual gluing arrow is exactly

\[
 \begin{pmatrix}
 1&n\phi_g&0&0&0&0\\
 0&1&0&0&0&0\\
 0&-n&1&0&0&0\\
 0&0&0&1&0&0\\
 0&0&0&0&1&0\\
 0&0&0&0&0&1
 \end{pmatrix}
 \begin{pmatrix}
 du+\alpha_Qdv\\dv\\d\gamma\\d\alpha_q\\d\theta\\dr
 \end{pmatrix}.
                                                        \tag{SQ.18}
\]

\item Indeed \(\alpha_{Q_n}=\alpha_Q+n\phi_g\), while \(\gamma_n=\gamma-nv\).
The matrix has determinant one. The coordinate matrix on the dual
frames is its inverse transpose; using that matrix instead of the
coframe matrix changes none of the following calculation. There is
no unspecified tangent block. The differentials of arg(\(\alpha\)) and
arg(\(\eta\)) belong to the full circle connection (SQ.14); their norms
have already entered the prescribed \(\gamma\) coordinate. The term
with \(\kappa\) winding remains in \(\alpha_{Q}\) and in (SQ.12). The derivative
d(arg g)/(2pi) is retained explicitly as d(\(\phi_{g}\)) below.

\item Denote the displayed matrix by G. Its Maurer form is literally
\(\omega=G^{-1}dG=n\,d\phi_g\,E_{12}\). The constant \(\gamma\) shear contributes
no derivative, and its product with \(E_{12}\) is zero. Thus \(\omega\) squared
is zero even before restricting to the two boundary base variables.
The dual-frame matrix has the corresponding transposed square-zero
Maurer form. Interpolate the actual flat connections with \(\chi\)(r) \(\omega\)
on the gluing collar. Direct multiplication gives

\[
 F=\chi' dr\wedge\omega+(\chi^2-\chi)\omega^2,
 \qquad\operatorname{tr}(F^2)=0.                       \tag{SQ.18a}
\]

\item This verifies the full collar transgression with the actual matrix.
Flatness at its two ends alone would not suffice. The positive-norm
and graph comparisons are confined to the interior by the prescribed
collar equality, and therefore introduce no additional boundary
matrix. We can compare the two glued tangent connections with p1
identically zero on the entire intervening collar.

\par

\medskip\noindent\textbf{SQ.5. The remaining relative pairing is zero}

\par\nobreak

\item Choose nested fixed tubes \(U_{0}\) inside \(U_{1}\) inside \(U_{2}\) inside \(U_{3}\), all away
from the outer boundary. Confine the surgery and every connection change
to \(U_{1}\). Choose the tangent connection on \(C_{0}\) to agree outside \(U_{1}\)
with a connection pulled from TE in the actual product E times
\(S^1_\gamma\); add its trivial \(\gamma\) line. On the outside collar use
the flat connection in (SQ.17). The source connection can be
chosen with precisely its source part there. Extend the connection
arbitrarily across the surgery region using the actual normal and
cap arrows of SQ.2. In the \(\rho\) family choose it smoothly and
periodically. Its p1 form vanishes near the outer boundary.

\item We evaluate its relative p1 class on (SQ.15). Fix a parameter \(\theta\)
and start with a smooth graph map b:E to \(S^{1}\). At a regular value,
its four-dimensional product level in W is \(N^3\) times \(S^1_\gamma\).
The source-pulled p1 four-form restricts to zero there, since its
source component lies in a three-manifold.

\item The restriction of b to F has a real lift. Extend that lift over a
tube \(U_{3}\). A cutoff homotopy supported there, fixed on the outside
boundary, changes b to a constant on \(U_{2}\). Choose a regular value
different from that constant. The new level misses the entire
surgery region and is literally the same submanifold before and
after attaching \(D^{3}\) times \(S^{2}\). It need not be a product; we do not
make that assertion. The regular-value cobordism of the displayed
homotopy, relative the unchanged outer map, identifies its relative
cycle with the original product level. The source-pulled p1 form is
closed and its outer trace is zero. Its integral on the new level
is therefore also zero. The two normal directions to that level
in the \(\rho\) family are the b direction and the \(\theta\) direction;
both are globally framed. Consequently this is exactly the
relative pairing, with its actual tangent bundle,

\[
              \big\langle p_1(TC_0),[d\theta]\smile b
                         \big\rangle=0.                 \tag{SQ.19}
\]

\item Here p1 is relative to the specified flat boundary connection,
whereas b is absolute. There is no assertion that b vanishes on
the outer torus. This distinction is precisely why the retained
boundary connection and its transgression in SQ.4 were required.

\par

\medskip\noindent\textbf{SQ.6. Gluing the actual characteristic calculation}

\par\nobreak

\item The rest of the completed quotient is the same actual exterior for
every \(T_{n}\). On the cut piece use the explicit diffeomorphism (SQ.9)
to transport its tangent connection; on the common exterior keep
the original one. SQ.4 joins them over the actual boundary shear
with identically zero p1 form. Thus the difference of the two
p1-times-Euler integrals has only the interior term (SQ.15) and
the exact curvature terms in (SQ.13)--(SQ.14). The interior term
is zero by (SQ.19). The exact terms contribute a boundary integral
of their actual one-form primitive times p1, which is zero on the
whole collar. This explicitly includes the gauge
-n(du+\(\alpha_{Q}\) dv); it is not discarded or assumed to extend.

\item The two physical completions consequently have the same ell.
The common parameter reversal retains the Euler line, so their
Jin completions have the same ell as well. P3K.42 gives ell(\(T_{0}\))=0.
The complete comparison therefore proves

\[
                              \ell(T_n)=0.              \tag{SQ.20}
\]

\item Exact algebraic checks are executed in \path{code/part3/source_cusp_shear_pontryagin.py}, with saved JSON, stdout and stderr. The 29 checks retain the full unit winding, both parameter monodromies, the literal source point, the nodal vector field, the integral Thom/Mayer--Vietoris quotient, the six-by-six derivative and its whole-collar transgression. The geometric proof is assessed separately from this certificate.

\par

\medskip\noindent\textbf{SQ.7. Exact relative dictionary with both lines}

\par\nobreak

\item The real curvature forms in SQ.12--15 represent the meridian-one
line \(U=M^*\), where \(M\) is the physical nonzero-vector circle
line. Their exact Hermitian dual arrow reverses the circle character;
\(c_1(U)=dA\) and \(c_1(M)=-dA\). The zero pairing proved above
holds for both lines. K68 retains the actual cap and Hopf-transfer
map, and K66 computes its full relative Euler cube. Write
\(A_U=(4k_U-\ell_U)/24\) and \(A_M=-A_U\). The native product
identities are
\[
 k_U=V_{00}+6=-(m_0+n_0),\qquad
 \ell_U=24+2V_{00}+2\delta_0=-4m_0.
\]
Thus \(m_0=0\), \(n_0=-6A_U=6A_M\), with
\(V_{00}=6A_U-6\), \(\delta_0=-6A_U-6\).
The full actual line/cochain and oriented collar proof of K66 gives
\[
 A_U=1+3n(n-1),\quad A_M=-1-3n(n-1),\quad
 (m_0,n_0)=(0,-6-18n(n-1)),
\]
\[
 V_{00}=18n(n-1),\qquad \delta_0=-12-18n(n-1).             \tag{SQ.21}
\]
The old opposite-sign residue inference is preserved in the
correction history. These exact entries now follow from the complete
K66 integer proof, not from choosing an integer lift of a residue.

\end{enumerate}
\end{proof}




\begin{theorem}\label{night:P3K.64}
\StatementRecord{night:P3K.64}{proved}{S6\_06\_BLIND: GPT-6 Astra (OpenAI, ChatGPT Pro mode, effort Pro), recorded re-derivation in the checkpoint dated 18 November 2026}{\cref{night:P3K.10}, \cref{night:P3K.22}, \cref{night:P3K.25}, \cref{night:P3K.26}, \cref{night:P3K.27}, \cref{night:P3K.28}, \cref{night:P3K.31}, \cref{night:P3K.32}, \cref{night:P3K.47}, \cref{night:P3K.62}, \cref{night:P3K.68}, \cref{night:P3K.69}}

Retain the actual canonical fixed-source scalar loops of the whole-zero
\(I_1\) constructions K22, K26, K28 and K32, with their source,
Cartier line, finite-root cap characters and complete annular maps.
The ordinary forward closing is its specified \(q^e\) times the
bounded source-independent unit, with the original cap arrow.
Write \(J\in\mathbb Z\) for the actual scalar winding, including
integer powers and inverse traversals of these constructed loops.
On the protected cusp disc keep
\[
 h:(Z,V)\mapsto(\kappa qZ,\alpha V),\qquad
 B:(Z,V)\mapsto(gZ,q\lambda\beta V),
\]
where \(\kappa,\alpha,g,\beta\) are independent of the scalar.
For every \(n\in\mathbb Z\), use the actual periodic fan \(T_n\)
whose edge directions are
\[
                    e_2,\qquad (1,n),\qquad(1,n-1).
\]
The actual positive return on this fan has class
\[
        A_n=A_0-J\,\frac{n(n-1)}2
                  \quad\text{in }\Theta_7\cong\mathbb Z/28,
 \tag{P3K.64.0}
\]
where \(A_0\) is the independently constructed initial return.
Its native relative free pairs on the two nonfixed stars are
\[
 \begin{array}{c|cc}
  &m_{\rm rel}&n_{\rm rel}\\ \hline
 C_1&J(1-n)&J(n-1)\\
 C_2&nJ&-nJ.
 \end{array}
\]
The actual \(C_1\) flip goes from \(T_n\) to \(T_{n-1}\);
the actual \(C_2\) flip goes from \(T_n\) to \(T_{n+1}\).
All periods, bounded-unit seams and cap maps are retained by
the explicit arrows in the proof. This statement covers this
vertical line of fans and these canonical closing maps; it
does not supply an arbitrary-fan formula or a shorter-loop
comparison with an additional source-translation endpoint.
\end{theorem}
\begin{proof}
\begin{enumerate}[label=\arabic*.]
\item \par

\medskip\noindent\textbf{SV.1. The actual shear, periods and closing}

\par\nobreak

\item Retain one of K62's actual fixed-source scalar families,
with its protected whole-zero nodal cusp and full periods
\[
 h:(Z,V)\mapsto(\kappa qZ,\alpha V),\qquad
 B:(Z,V)\mapsto(gZ,q\eta V),\qquad \eta=\lambda\beta.
\]
All of \(\kappa,\alpha,g,\beta\) are independent of the
scalar parameter. Write \(J\in\mathbb Z\) for its actual
scalar winding. The forward ordinary closing, with its
already constructed cap arrow, is \(q^e\) times its
specified bounded source-independent unit. The canonical
scalar loops of K26/K28/K32 retain exactly these arrows.

\item For each integer \(n\), use the physical-vertical-edge
fan \(T_n\) with directions
\[
             e_2,\quad (1,n),\quad(1,n-1).
\]
On its torus the literal coordinate and period maps are
\[
 Z_n=Z,\quad V_n=Z^{-n}V,\qquad h_n=hB^n,\quad B_n=B.
                                                        \tag{SV.1}
\]
They preserve the physical circle and the original height
\(q\). In these actual coordinates the full periods are
\[
 \begin{aligned}
 h_n:(Z_n,V_n)&\mapsto(\kappa_n qZ_n,\alpha_n V_n),\\
 B_n:(Z_n,V_n)&\mapsto(g_nZ_n,q\eta_nV_n),\\
 \kappa_n&=\kappa g^n,\qquad
 \alpha_n=\alpha\eta^n\kappa^{-n}g^{-n^2},\\
 g_n&=g,\qquad \eta_n=\eta g^{-n}.
 \end{aligned}                                          \tag{SV.2}
\]
These identities follow by substituting (SV.1) in each
original period, including its line coordinate. The
integer labels of \(h_n,B_n\) are the standard two basis
vectors in the new characters, so their period orbit
still has one vertex and no stabilizer of a chosen
star representative.

\item The actual winding vector is consequently
\[
 \bigl(\operatorname{wind}\kappa_n,
       \operatorname{wind}\alpha_n,
       \operatorname{wind}g_n,
       \operatorname{wind}\eta_n\bigr)
                    =(0,nJ,0,J).                    \tag{SV.3}
\]
There is no assertion that the new source line is an
unmarked replacement for the old one. Its actual
character map is \(V_n=Z^{-n}V\).

\item The ordinary closing remains \(q^e\) in these characters.
Reindexing by \(B_n^{-e}\) gives the full bounded
endpoint \(E_n=(g^{-e},\eta_n^{-e})\), with the retained
bounded closing unit included in its second factor.
On the cut parameter interval take its actual
logarithms and use
\[
 C_\tau:(Z_n,V_n)\mapsto(AZ_n,DV_n),\quad
 A=\exp(e\chi\log g),\quad
 D=\exp(e\chi\log\eta_n)                              \tag{SV.4}
\]
and the inverse logarithm of that bounded closing
unit in \(D\). The cutoff is zero on the first collar
and one on the last. Thus \(C_{\rm end}E_n=1\)
with every derivative. It commutes with all actual
periods and the physical circle. The two vertex
chart expressions are exactly
\[
 (a,b,c)\mapsto((AD)^{-1}a,Ab,Dc),\qquad
 (u,v,w)\mapsto(ADu,D^{-1}v,A^{-1}w).
                                                        \tag{SV.5}
\]
They preserve every transition row with the full units
\(\kappa_n,\alpha_n,g_n,\eta_n\).
Use on the degree-two parameter cover the literal
pullback
\[
          C^{(2)}_{j+\tau}=C_\tau E_n^{-j},\quad j=0,1.
                                                        \tag{SV.6}
\]
This retains the original ordinary and cap endpoint
squares and gives \(E_n^{-1},E_n^{-2}\) at its two
successive collars. No new seam is selected for the
doubled return.

\par

\medskip\noindent\textbf{SV.2. The closed ordinary product and exact trace}

\par\nobreak

\item The three native normal transitions are the ones of
K47/K62 with the full units (SV.2). In the second
vertex product \((u,\eta_n v,\kappa_n w)\), their
based double-loop normal matrices are
\[
 \begin{aligned}
 K_0&=\operatorname{diag}
       (e^{2iJ(1-n)\theta},e^{2inJ\theta}),\\
 K_1&=\operatorname{diag}(1,e^{2iJ\theta}),\\
 K_2&=\operatorname{diag}(e^{2iJ\theta},1).
 \end{aligned}                                          \tag{SV.7}
\]
For example their first coordinate factors before
reference division are \(\eta_n/\alpha_n\),
\(\kappa_n/g_n\), and \(\eta_n\), respectively.
Their second coordinate factors are
\(\kappa_n\alpha_n\), \(g_n\eta_n\), and
\(\kappa_n\). This proves (SV.7) from the actual
native coordinate arrows. The determinant is
\(e^{2iJ\theta}\) in every row.

\item Use the explicit quaternion contractions of DZ.12
with
\[
 t=J,\qquad
 l=\bigl(2J(1-n),0,2J\bigr),\qquad
                 Q\mapsto H_t QH_{l_i-t}.            \tag{SV.8}
\]
The equality of the determinant loops makes the
first vertex product common to all three stars.
Each remaining determinant-one normal loop is
contracted by that displayed \(H\) map. The actual
normal-rescaling and cutoff-flow construction of
DZ.2 carries the higher derivatives and retains
the full nonzero unit germs. The scalar factors
with nonzero winding are the exact powers in
(SV.7); every other based unit has its fixed
logarithm. Thus this is a closed ordinary product
of the doubled marked plumbing. It is not claimed
to intertwine the physical circle.

\item The physical tangent and normal weights in these
characters remain
\[
                  w=(0,1,-1),\qquad d=(-2,-1,-1).
\]
Use the literal orbit map DZ.16 and its derivative
in the covering group of stable \(GL^+\). Their
degrees are
\[
 \begin{aligned}
 I_i&=t w_i+d_i(l_i-t),&
 I&=(-2J+4nJ,\,2J,\,-2J),\\
 W_i&=t w_i-d_i(l_i-t),&
 W&=(2J-4nJ,\,0,\,0).
 \end{aligned}                                          \tag{SV.9}
\]
The same generic vertex paths and rank-one endpoint
caps as in DZ.3 close these exact traces. The
parameter-dependent vertices are (SV.5), so they
commute with the diagonal physical action and
their endpoint traces remain independent of the
parameter. No extra vertex meridian is inserted.
The coordinate-frame changes extend over the
annulus times its normal four-disc, exactly as
in the proved stable-trace comparison.

\item The integral prism identities now give
\[
 \begin{array}{c|cc|cc}
 &V_{ii}&\delta_i&m_i^{\rm ord}&n_i^{\rm ord}\\ \hline
 C_1&2J&-2J+4nJ&-2nJ&2J(n-1)\\
 C_2&2J& 2J-4nJ&2J(n-1)&-2nJ .
 \end{array}                                          \tag{SV.10}
\]
Here \(2n_i^{\rm ord}=\delta_i-V_{ii}\) uses only
the actual local ambient-to-single-star frame map,
not the earlier scalar initial-variation assertion.

\par

\medskip\noindent\textbf{SV.3. The native pairs, exact flops and recurrence}

\par\nobreak

\item The native \(C_1\) coordinates are
\[
 (c,a,b),\qquad
       (\eta_n v,g_nw,(g_n\eta_n)^{-1}u),
\]
and the native \(C_2\) coordinates are
\((a,b,c),(u,v,w)\). Their exact transitions are
\(z'=z^{-1},n'_i=zn_i\). On each star the change
to (SV.8) is
\[
                q_H\mapsto
                   (UH_JU^{-1})q_HH_{l_i-J}.         \tag{SV.11}
\]
Its left degree is \(2J\), and its right degree
is zero. The inverse cylinder extension has
endpoint \(L^{-2J}\), with coefficient maps
extending over \(D^3\times D^4\). Thus the native
doubled pairs are obtained from (SV.10) by adding
\(2J\) to the left coordinate:
\[
 C_1:\ (2J(1-n),2J(n-1)),\qquad
 C_2:\ (2nJ,-2nJ).                                   \tag{SV.12}
\]

\item The boundary products and seams (SV.6) are literal
pullbacks of the single native products. Pull back
a connection with those exact boundary products.
Its return is the actual relative square. Any
connection used in (SV.8)--(SV.11) with the same
boundary product interpolates to that one relative
to the boundary. The free-coordinate map to
\(\mathbb Z^2\) therefore gives the actual single
pairs
\[
 C_1:\ (J(1-n),J(n-1)),\qquad
 C_2:\ (nJ,-nJ).                                     \tag{SV.13}
\]
This is injectivity of multiplication by two in
the free relative group, not division of a
seven-sphere class.

\item The full Cox lifts of (SV.5) are
\[
 \begin{aligned}
 C_1:\ &(a_0,a_1,c_0,c_1)
       \mapsto(a_0,Da_1,(AD)^{-1}c_0,Ac_1),\\
 C_2:\ &(a_0,a_1,c_0,c_1)
       \mapsto(a_0,(AD)^{-1}a_1,Ac_0,Dc_1).
 \end{aligned}                                          \tag{SV.14}
\]
Each preserves the actual height monomial and
extends through both resolutions. Transport the
metric by the full unit map and use its positive
radial normalization on the link. The unrestricted
four-weight common-link identity then gives the
literal \(T(y,q_H)=(q_H^{-1}yq_H,q_H)\) extension,
with the same fixed complex-normal reset as K62.
Equivariance under \(h_n,B_n\) carries one chosen
free star representative to the whole periodic
orbit. The outside source, line and finite-cap
maps are unchanged.

\item In the actual lattice marking, flipping \(C_2\)
replaces direction \((1,n-1)\) by \((1,n+1)\);
the new fan is \(T_{n+1}\). Flipping \(C_1\)
replaces \((1,n)\) by \((1,n-2)\), giving
\(T_{n-1}\). These are the literal two diagonal
replacements in the periodic squares, with the
same period group and positive height.
The actual Cox conjugation changes the central
coordinate by the right native integer. Hence
\[
 A_{n+1}-A_n=-nJ,\qquad
 A_{n-1}-A_n=(n-1)J.                                \tag{SV.15}
\]
The two are inverse statements for the same
oriented elementary comparison. Their exact
telescoping, with the independently computed
initial fixed-star value \(A_0\), gives
\[
                   A_n=A_0-J\,\frac{n(n-1)}2
                          \quad\text{in }\mathbb Z/28.
                                                        \tag{SV.16}
\]
The integer \(n(n-1)/2\) is integral for every
positive, zero or negative \(n\). At \(n=0\)
the two first comparisons agree with K62.
For a concrete primitive counterexample to the earlier all-fan
scalar equality, retain K10/K22's actual source
\[
 y^2=x^3-3f(t)x+2g(t),\quad
 f=(t^2+1)^2+t,\quad
 g=(t^2+1)^3+\tfrac32t(t^2+1)+\tfrac38.
\]
Take its actual section \(G=(t^2+1,\sqrt3/2)\), line
\(\mathcal M=\mathcal O(8G-O)\), and translation section \(Q=G\).
Choose any one of the four simple roots of
\(48t^4+8t^3+96t^2+72t+57\) as the linearization zero,
and retain the multiplicity-one ordinary cap at \(t=0\)
with integral clutch \(\widehat\gamma\). K10 constructs
these charts. Denote the actual initial class by \(A_0\), retaining
the physical Hopf-line identification of K68. K69 gives \(A_0=3\). K27.2's actual positive
scalar loop is \(\lambda=e^{2\pi i v}\), hence \(J=1\).
The literal fan \(T_2\), with directions
\(e_2,(1,2),(1,1)\), therefore has
\[
                  A_2=A_0-1\quad\hbox{in }\mathbb Z/28.
\]
The earlier K25 equality gives the same value on both fans.
The earlier K31 formula likewise gives the same value on both,
because their physical-circle coordinates are \((0,1)\) and their
norm is one. The actual classes differ by one, independently
of its absolute reference. With K68--69, the actual classes are three and two.
The displayed family, loop, cap and fan thus supply the
specified counterexample to both universal assertions.
This computation covers the displayed vertical
line of fans only. It makes no assertion about
the arbitrary-fan recurrence quarantined after
the scalar determinant obstruction.

\item The exact certificate is \path{code/part3/scalar_vertical_fans.py}, with its
JSON and captured output. It verifies 14 symbolic identities and 425
signed instances. The actual maps and relative comparison were independently
audited by root and checks before promotion. The scope retains canonical
ordinary closing by \(q^e\) times the specified source-independent bounded unit;
shorter endpoints with an additional source translation are not included.

\end{enumerate}
\end{proof}




\begin{theorem}\label{night:P3K.65}
\StatementRecord{night:P3K.65}{proved}{S6\_06\_BLIND: GPT-6 Astra (OpenAI, ChatGPT Pro mode, effort Pro), recorded re-derivation in the checkpoint dated 18 November 2026}{\cref{night:P3F.22}, \cref{night:P3F.26}, \cref{night:P3F.29}, \cref{night:P3K.17}, \cref{night:P3K.23}, \cref{night:P3K.6}, \cref{night:P3K.7}}
Let the source and marked section be
\[
 y^2=x^3-3t^3(t-1)x+2t^4(t-1)(t-\tfrac12),\qquad
 G=(t^2,t^2),\qquad M=\mathcal O(G-O).
\]
Use the actual suspension over \(-3G\), with target-point
linearization divisor the identity component \(C_0\) of the
\(I_2\) fibre at \(q=1/t=0\). At the \(IV^*\) fibre use the
cap of order \(m=3k\), \(k>0\), and invariant vector
\(\ell=n(e_\gamma-e_u)+d\delta\), where \(n=\pm1\) and
\(d\in\mathbb Z\). Retain the original type-II line-complement
quotient. Denote by \(\mathcal T_-\) the negative-diagonal fan
of P3F.26 and by \(\mathcal T_+\) the positive-diagonal fan of
P3F.31. For the actual positive parameter loops the oriented
clutched-sphere classes are
\[
\begin{array}{c|c|c|c}
 \text{parameter}&\text{scalar}&\mathcal T_-&\mathcal T_+\\ \hline
 r&\lambda=r^{3k}&-3k&0\\
 w&\lambda=w^k,\quad w=r^3&-k&0
\end{array}\qquad\text{in }\Theta_7=\mathbb Z/28.       \tag{P3K.65.OR.1}
\]
The positive fan has Jin characteristic integer zero in both rows.
The actual native C2 free pairs on the negative fan are respectively
\((-3k,3k)\) and \((-k,k)\). In particular the positive scalar
loop \(\lambda=w\), with \(k=1\), on \(\mathcal T_-\) has class
\(27\), a generator. As \(k\) varies, the shorter negative-fan
families realize every residue modulo 28.

The proof assigns the negative fan's sphere class by the actual
relative diffeomorphism and its oriented central generator. It does
not identify an integer characteristic number with a chosen lift of
that residue.
\end{theorem}
\begin{proof}
\begin{enumerate}[label=\arabic*.]
\item \par\medskip\noindent\textbf{OR65.1. Source, component line, cap and forward parameter arrow.}\par\nopagebreak[4]

\item The source discriminant factor and section substitution are
\[
 [t^3(t-1)]^3-[t^4(t-1)(t-\tfrac12)]^2
       =-\tfrac14t^8(t-1)^2,
\]
so the marked fibres are \(IV^*,II,I_2\) at \(0,1,\infty\).
On the cubic root \(t=z^3,x=z^4X,y=z^6Y\), the good central
curve is \(Y^2=X^3+1\), and \(G\) specializes to \((0,1)\),
of exact order three. At infinity the minimal coordinates of G
are \((1,q)\), meeting component one after resolving the node.
The section is disjoint from O. The line of target-point
linearizations of translation by \(-3G\) has divisor class
\[
 G-O-(-2G)+(-3G)\sim C_0.
\]
Indeed its component degrees at \(I_2\) are \((-2,2)\), those of
\(C_0\); all degrees cancel at \(IV^*\); subtracting \(C_0\)
therefore leaves an integral fibre class, and intersection with O
is zero. We retain that actual section and its scalar multiplier,
not a whole-fibre substitute.

\item On the nonzero line over the Tate chain the source and suspension
periods have literal valuation vectors \((2,1)\), \((-3,-1)\):
\[
 u=(2s,s)^t+U(q),\qquad
 \gamma=(-3s,-s+mv)^t+C(q),
 \quad s=\frac{\log q}{2\pi i},\quad
 v=\frac{\log r}{2\pi i}.                            \tag{P3K.65.OR.2}
\]
All components of U and C are the retained bounded holomorphic
units' logarithms. Their units, rather than their logarithms
alone, are the ordinary torus maps. Put
\[
 H=-u-\gamma=(s,-mv)^t-U-C,\qquad
 K=3u+2\gamma=(0,s+2mv)^t+D,\quad D=3U+2C.
\]
The integral change has determinant one. In the physical characters
\((Z,V)\) these primitive periods are
\[
 h:(Z,V,q)\mapsto(\kappa q Z,\alpha V,q),\qquad
 B:(Z,V,q)\mapsto(gZ,\eta qV,q).                     \tag{P3K.65.OR.3}
\]
For one positive r-loop their full unit windings are
\[
  (\operatorname{wind}\kappa,\operatorname{wind}\alpha,
    \operatorname{wind}g,\operatorname{wind}\eta)
                         =(0,-m,0,2m).              \tag{P3K.65.OR.4}
\]
Here the physical circle is the positive phase of V.

\item The actual good-root cap has linear part \(A^3=1\), with
\(A\ell=\ell\), and affine action
\(z\mapsto Az+\ell/m\). Its j-th power has gamma displacement
\(jn/m\). For \(0<j<m\) this is nonintegral, so the action is
free and no physical delta rotation compensates it. On its source
root \(t=s_0^m\) the actual ordinary-to-cap map is
\[
 \Phi_r(z,s_0)=z+\frac{\log s_0}{2\pi i}\Pi_r(s_0)\ell.
                                                               \tag{P3K.65.OR.5}
\]
Changing the logarithm branch adds an integral period; the root
deck adds exactly the affine displacement. The fixed character
\(\exp(2\pi i d/m)\) is retained. The cap action in the physical
line character is \(A_r=r^nA_0\).

\item The one-r-turn endpoint-to-initial ordinary closing is multiplication
by \(t^n=q^{-n}\), with the cap comparison dictated by (P3K.65.OR.5).
The forward parameter arrow is consequently \(q^n\). Reindex by
the actual K-period \(B^{-n}\) to return a chosen toric cell to
its initial representative. The resulting bounded forward arrow is
\[
 E_r(Z,V,q)=(g(q,r)^{-n}Z,\eta(q,r)^{-n}V,q).          \tag{P3K.65.OR.6}
\]
This equality retains the period and the forward direction for
each sign of n. The type-II and finite-cap maps are untouched
outside the cusp neighbourhood. The independent integral topology
of P3F.26/P3F.29 gives \(\pi_1=0\) and integral sphere homology
for these signs: the four actual transgression maps are
multiplication by n after the specialization indices
\((m,k,k,m)\). Both small resolutions have the same topology,
because their common blowup adds one primitive integral generator
in degrees two and four and preserves the fundamental group.
Thus these are actual smooth six-sphere families.

\par\medskip\noindent\textbf{OR65.2. The independent positive-fan index.}\par\nopagebreak[4]

\item We record why the positive-fan input does not use the claimed
equality in P3K.17. On the source Tate chain the actual line graph
is \(h_j=\lceil j/2\rceil\), with \(h_{j+2}=h_j+1\).
The fibre-coordinate transition gives component degrees
\(-h_{j-1}-h_{j+1}+2h_j=(-1,+1)\). Its alternating edges
\((1,1)\), \((1,0)\) lie in the positive-diagonal unit-square fan.
Projection to the source chain and quotient by u therefore give a
literal toric open embedding of \(M^\times\), through both nodes.
The suspension's height defects are
\[
 h_j-1-h_{j-3}=0\ (j\text{ even}),\quad1\ (j\text{ odd}),
\]
and its N-th positive iterate has defect
\(\lfloor N/2\rfloor\) or \(\lceil N/2\rceil\), respectively.
A point leaving this line-complement open set at a zero cannot
re-enter through a later positive translate.

\item Choose a smooth Hermitian metric, fibrewise flat on the regular
cusp annulus, and a scalar small enough that the actual lift has
operator norm at most \(c<1\). Choose m shell radii with
\(R_1<\cdots<R_m\) and \(cR_m<R_1\). The preceding iterate
formula and this strict norm inequality prove that the shells
embed disjointly in the completed quotient. Their physical-circle
quotients are embedded copies of the full source \(I_2\)
neighbourhood, carrying the actual line M and missing the fixed
sphere.

\item On the regular annulus the lift's norm multiplier c(q) is constant
on each elliptic fibre: the logarithm of the ratio of its flat
metrics is harmonic on that compact fibre. The exact lift of a
circle reparametrization \(F_q(t+1)=F_q(t)+1\), with
\(t=\log\|V\|/\log c(q)\), is
\[
                  V\mapsto c(q)^{F_q(t)-t}V.
\]
It preserves the source and base coordinates, commutes with the
contracting lift, and carries the circle line. This glues the m
shells to the m sheets of the actual regular-binding page of the
weighted Hopf base of weights \((1,m)\). At the type-II fibres
use the original line-complement map; at the order-m cap use the
stated good-root map. The fixed d-character supplies only the
explicit scalar transition \(\exp(2\pi i d/m)\) on each sheet;
its product is \(\exp(2\pi i d)=1\). It therefore adds a base
line, whose restriction to the base page is topologically trivial.
This constructs a page V carrying M through all source components.

\item The source cap of this page has m fibres of type II and m of type
\(I_2\). Its two sections remain disjoint: the degree-3 root
already separates them, and the further degree-k cover branches
only at its smooth good-reduction fibre and the regular binding.
Thus
\[
 \chi(\mathcal O)=k,\quad \sigma=-8k,\quad
 M_{\rm rel}^2=-2k,\quad p_{\rm rel}=-24k.           \tag{P3K.65.OR.7}
\]

\item For completeness the meridian-line comparison uses the actual
regular bands. The fixed-edge charts have normals
\((A,B)=(q/V,V/Z)\), with second chart
\((A',B')=(ZA,ZB)\); their weights are \((-1,+1)\).
The real compression of P3K.6 lies over \(q\in[0,\epsilon]\).
The page misses its fixed end, so all remaining intersections
project to a compact regular interval. The normal pushoff of its
unchanged binding collar is
\[
 q-\epsilon=(\rho+i\eta_0)e^{i\alpha(\gamma)},\quad
 \rho\ge0,\quad \eta_0>0.
\]
Its intersection with the real compression requires
\(\rho=-\eta_0\cot\alpha\) and
\(q-\epsilon=-\eta_0/\sin\alpha\), hence occurs only on m
proper boundary arcs. Enlarge their regular lifts to disjoint
half-disc bands. Each pair is
\((T^2\times D_i,T^2\times A_i)\), where the inner semicircle
\(A_i\hookrightarrow D_i\) is a homotopy equivalence. Its
relative \(H^2\) is zero. Outside these bands the exact radial
graph homotopy reaches the original page pushoff without meeting
the trace. Excision and restriction injectivity therefore identify
the graph-normal meridian line with M on the entire page.

\item The incoming cycle map is exactly
\[
 (u,\gamma)=(2H+K,-3H-K),\qquad
 \det\begin{pmatrix}2&-3\\1&-1\end{pmatrix}=1.
\]
The compression cap kills the primitive vanishing cycle \(w_0\);
the regular elliptic cap kills gamma. Their closed difference is
\(S^1_u\times S^3\), with zero signature. The cap's line class
is zero because its restriction to the boundary is zero and the
restriction map on \(H^2\) is injective. A boundary gauge changes
the glued class by \(d\chi\wedge\alpha_0\); its square and its
product with M vanish. The actual sphere-bundle bordism of P3K.7
therefore has \(b^2=-8k\), \(p_1(T\bar X)=-24k\), and
\[
              A_+=\frac{3(-8k)-(-24k)}{24}=0.        \tag{P3K.65.OR.8}
\]
Both norm signs use their actual positive parameter arrow. The
extending origins on \(q=Q^{2m}\) are
\[
 \sigma=nvD+nmv(v-1)\delta,\qquad
 o=-\frac{nsD}{2m}-\frac{ns^2}{4m}\delta-nsv\delta.
\]
Their difference changes by
\(n\theta K+nm\theta(\theta-1)\delta\); the second term has
zero endpoint values. The signed page/cycle map retains the two
norm signs, and the computed integer is zero in either orientation.
P3F.22's actual Hopf-lift transfer gives positive-fan sphere class
zero. No negative-fan assertion enters this computation.

\par\medskip\noindent\textbf{OR65.3. The actual closed native C2 product.}\par\nopagebreak[4]

\item On a cut r-circle use holomorphic-in-q, smooth-in-parameter
logarithms of g and eta. Let \(\chi\) be zero and one on the
initial and terminal collars. Put
\[
 C_\tau(Z,V,q)=(A_\tau Z,D_\tau V,q),\qquad
 A_\tau=e^{n\chi\log g},\quad D_\tau=e^{n\chi\log\eta}.
                                                               \tag{P3K.65.OR.9}
\]
It has the inverse obtained by negating the exponents, commutes
with both periods and delta, and satisfies
\(C_{\rm end}E_r=C_{\rm start}=1\) on whole collars, with all
derivatives. It acts on the two vertices as
\[
 (a,b,c)\mapsto((AD)^{-1}a,Ab,Dc),\qquad
 (u_0,v_0,w_0)\mapsto(ADu_0,D^{-1}v_0,A^{-1}w_0).
\]
The complete six-chart transitions are
\[
\begin{array}{c|ccc}
 &u_0&v_0&w_0\\ \hline
 C_0&\kappa\alpha bc&\alpha^{-1}ab&\kappa^{-1}b^{-1}\\
 C_1&g\eta bc&\eta^{-1}c^{-1}&g^{-1}ac\\
 C_2&a^{-1}&ab&ac.
\end{array}                                                   \tag{P3K.65.OR.10}
\]
Substitution gives the same three factors on every row. Thus all
full unit coefficients, their windings (P3K.65.OR.4), the physical
weights and the finite-cap closing remain the actual ones. In
particular C2 is a literal closed native star
\(z'=z^{-1},\ c'_i=zc_i\), with first coordinates \((a,b,c)\).
This asserts a product on that single star, not a simultaneous
delta-equivariant product of the three-curve plumbing.

\item Define its doubled product as its actual pullback. On the second
sheet use
\[
 C^{(2)}_{j+\tau}=C_\tau E_{r(\tau)}^{-j},\qquad j=0,1.
                                                               \tag{P3K.65.OR.11}
\]
At the middle seam both maps equal \(E^{-1}\); at the final seam
the map cancels the doubled forward arrow after \(B^{-2n}\).
The complete cap overlap (P3K.65.OR.5) is pulled back at the same time.
A connection on the actual compact family with this prescribed
native boundary product is obtained by extending its velocity
from a collar with a cutoff and joining it to an arbitrary global
connection in the affine space of connections. On the double
cover use the literal pullback connection. Its relative return
is the square of the single native return.

\par\medskip\noindent\textbf{OR65.4. An ordinary doubled product and its actual trace degrees.}\par\nopagebreak[4]

\item For the doubled closed charts set
\[
       t_0=2m,\qquad (l_0,l_1,l_2)=(6m,0,4m).       \tag{P3K.65.OR.12}
\]
At the second vertex use \((u_0,\eta v_0,\kappa w_0)\), with
the same initial constants in the reference chart. The three
normal matrices become
\[
 D_0=\operatorname{diag}(\eta/\alpha,\kappa\alpha),\quad
 D_1=\operatorname{diag}(\kappa/g,g\eta),\quad
 D_2=\operatorname{diag}(\eta,\kappa).
\]
Their determinant winding is \(2t_0\); their first diagonal
windings are the \(l_i\) in (P3K.65.OR.12). For each full based diagonal
entry subtract its displayed angular winding. The remaining
unit has a based periodic logarithm, retained below.

\item The explicit based quaternion contraction is
\[
 H_d(s_0,\theta)=
 (\sin\vartheta(s_0)e^{id\theta}+j\cos\vartheta(s_0))
 (\sin\vartheta(s_0)+j\cos\vartheta(s_0))^{-1},
\]
where \(\vartheta\) is constant at 0 and \(\pi/2\) near its
ends. It is periodic and based in theta. For
\(Q=n_0+j\bar n_1\), use
\[
                    Q\mapsto H_{t_0}QH_{l_i-t_0}.   \tag{P3K.65.OR.13}
\]
At the second radial end its two complex factors have windings
\((l_i,2t_0-l_i)\), exactly those of \(D_{i,*}^{-1}D_i\).
Multiply by the exponential of each retained zero-winding unit
logarithm, with radial cutoff zero and one at the two ends.
The resulting map \(\mathcal K_i\) satisfies
\[
 \mathcal K_i(0,\theta)=1,\qquad
 zD_{i,*}\mathcal K_i(1,\theta)=zD_i.
\]
The three curve interiors are disjoint, so these actual tube maps
glue to the two vertex products. For nonlinear germs conjugate by
\((z,n)\mapsto(z,\tau n)\), which extends at \(\tau=0\)
to its invertible normal derivative, and integrate the cutoff
velocity to restore the full germ. Carry the rounded boundary by
the same isotopy. This constructs a based periodic ordinary
plumbing product, with all unit maps and its boundary specified.
It is not required to intertwine the physical circle.

\item In that product the tangent weights and normal differences are
\[
                    w=(0,1,-1),\qquad d=(-2,-1,-1).
\]
The actual circle action on the normal quaternion is
\[
 Q\mapsto(H_{t_0}e^{-iw_i\phi/2}H_{t_0}^{-1})Q
       (H_{l_i-t_0}^{-1}e^{id_i\phi/2}H_{l_i-t_0}).
                                                               \tag{P3K.65.OR.14}
\]
The retained unit exponentials conjugate this formula and give a
based homotopy to it, with the metric transported. Let
\(U=\cos\chi+j\sin\chi\), \(E_i=e^{iw_i\phi/2}\). At the
normal section \(q_H=1\) the circle trace is literally
\[
 E_i^{-1}UE_i H_{t_0}E_i^{-1}H_{t_0}^{-1}U^{-1}
              H_{l_i-t_0}^{-1}e^{id_i\phi/2}H_{l_i-t_0}.
                                                               \tag{P3K.65.OR.15}
\]
Direct differentiation at the inverse images of \(-i\) gives
\(\deg(H_diH_d^{-1})=-d\),
\(\deg(H_d^{-1}iH_d)=d\), in the path-parameter orientation.
The quaternion volume gives degree ce for an axis of degree e
and angle \(c\phi/2\). The Maurer product boundary terms vanish
on the radial ends and constant central-sign angle ends; the
parameter faces cancel. Therefore the normal-meridian and closed
stable-derivative degrees are
\[
 I_i=t_0w_i+d_i(l_i-t_0),\qquad
 W_i=t_0w_i-d_i(l_i-t_0).                              \tag{P3K.65.OR.16}
\]
For the latter formula a left quaternion factor contributes its
degree and a right factor its negative. The base rotation has
only the phi variable; its Spin endpoint and the normal endpoint
are both \((-1)^{w_i}\), so their product is one. The vertex
caps are actual parameter-independent paths in the vertex links,
interpolating squared moduli and phases; their traces factor
through dimension two. They contribute no additional meridian
or third derivative class.

\item We recall the integral groups which make this a computation of
the actual return. The union of the three spheres has two vertices,
six edges and six faces, with
\[
 D_1=\begin{pmatrix}-1&-1&-1&-1&-1&-1\\1&1&1&1&1&1\end{pmatrix},
 \qquad D_2=\bigoplus_{i=0}^2
             \begin{pmatrix}1&-1\\-1&1\end{pmatrix}.
\]
Its primitive bases are the three sums of paired faces in degree
two and \(a_1-a_0,a_2-a_0\) in degree one. The oriented dual-cell
map in S6 gives the common exterior \(H_3=\mathbb Z^3\),
\(H_4=\mathbb Z^2\), \(H_2=0\). The H3 generators are the
positive complex normal meridians. Their pairwise intersection is
zero. Poincare--Lefschetz duality gives the corresponding relative
cycles \(D_i\). Filling the other two curve tubes projects the H3
basis to the selected meridian. These are independent cellular
and dual-cell maps, not consequences of an equivariant product.

\item Write \(\mathcal V\) for the ordinary-product variation and
\(\delta_i\) for the actual ambient stable-derivative pairings.
For \(a_j-a_0=\partial x_j\), the positive-circle identity
\(\partial\Delta=-\Delta\partial\) gives
\(\Delta x_j=w_jD_j-w_0D_0\). Its action prism and the chain-rule
prism yield
\[
 \mathcal V(w_je_j-w_0e_0)=I_je_j-I_0e_0,\qquad
 w_j\delta_j-w_0\delta_0=W_j-W_0.                    \tag{P3K.65.OR.17}
\]
Here the actual return fixes the common boundary by the ordinary
product just constructed. The ambient-frame change extends
through each annulus times its normal four-disc, whose H3 is zero;
its pairing on the capped trace vanishes. Thus no omitted change
of ambient stable frame occurs in (P3K.65.OR.17).

\item Substitution in (P3K.65.OR.16) gives
\[
                    I=(-8m,4m,-4m),\qquad W=(8m,0,0).
\]
Since \(w_0=0,w_2=-1\), equation (P3K.65.OR.17) gives
\[
                    \mathcal V_{22}=4m,\qquad \delta_2=8m.
\]
The actual ambient-to-first-resolution product comparison has
third class \(+1\) on its complex meridian: write the Cox normal
trivialization \(q_H=u(c_0+j\bar c_1)\), so that
\(c=u^{-1}q_H\). Its derivative is triangular; contract its
off-diagonal block to zero. The only nonconstant stable block
is left multiplication by \(q_H\), of degree +1. The positive
complex meridian and the relative D3 product orientation give
the frame correction \(-\mathcal V_{22}\). In the product frame
the relative generators L and R have stable derivatives 0 and
\(-2\), respectively. Hence the ordinary relative pair is
\[
 2\nu_2=\delta_2-\mathcal V_{22},\qquad
 \mu_2=-\mathcal V_{22}-\nu_2,
 \qquad (\mu_2,\nu_2)=(-6m,2m).                    \tag{P3K.65.OR.18}
\]

\item The native-to-ordinary C2 boundary arrow is
\[
 q_H\mapsto (U H_{t_0}U^{-1})q_H H_{l_2-t_0}.        \tag{P3K.65.OR.19}
\]
Its left factor has degree \(2t_0=4m\) on the Hopf-oriented
sphere times parameter; this is the Maurer boundary integral
of the displayed based H contraction. Its right factor depends
only on radial coordinate and parameter, so has third degree
zero. Extend the inverse factors over \(D^3\times I\) by the
literal cylinder retraction. The endpoint has free degrees
\((-2t_0,0)\), and these coefficient maps extend over
\(D^3\times D^4\). Thus there is no disc-supported correction.
The ordinary return is the native return followed by this
pure-left map. Its native doubled free pair is consequently
\[
                         (-2m,2m).                  \tag{P3K.65.OR.20}
\]
By (P3K.65.OR.11) this is the actual square of the single native return.
The relative group has torsion-free free quotient \(\mathbb Z^2\)
and central disc group \(\Theta_7\). Therefore its single native
free pair is \((-m,m)\). This step divides only two known free
integers of an actual square.

\par\medskip\noindent\textbf{OR65.5. The oriented native flop and the r-loop class.}\par\nopagebreak[4]

\item Retain K17's original Cox order \((z_0,z_4,z_1,z_2)\), auxiliary
weights \((1,1,-1,-1)\), and physical weights \((-1,0,0,1)\).
The positive auxiliary moment has \((z_0,z_4)\ne0\), with common
fan edge \(u_1u_2\), so it is \(\mathcal T_-\). In the chart
\(z_4=1\) take \((a,b,c)=(z_0,z_1,z_2)\); the chart \(z_0=1\)
is exactly \((a^{-1},ab,ac)\). This identifies the first ruling
with the native C2 product used above. The negative auxiliary
moment has \((z_1,z_2)\ne0\), common edge \(u_0u_4\), and is
\(\mathcal T_+\). The actual lifted seam is
\[
 (z_0,z_4,z_1,z_2)\mapsto
       ((AD)^{-1}z_0,z_4,Az_1,Dz_2).                 \tag{P3K.65.OR.21}
\]
It preserves q, commutes with auxiliary and physical circles,
and extends to both quotient rulings. Its inverse uses the
reciprocal units at the same q. Transport the Hermitian metric
by this map before radial normalization; thus its nonunit moduli
are retained on the common link. Full q-dependent germs are
carried by the same normal-rescaling isotopies on both rulings.
Every period representative has trivial stabilizer, and extension
by the original periods gives the map on the whole orbit of curves.
The cap and ordinary exterior are unchanged.

\item On the common link use the actual quaternion products
\[
 u=a_0+ja_1,\quad v=c_0+j\bar c_1,\quad
 h=uiu^{-1},\quad k=v^{-1}iv,\quad q_H=uv.
\]
The old and new boundary coordinates are \((h,q_H)\) and
\((k,q_H)\). Hence their exact arrow and oriented exterior
extension are
\[
 \tau(h,q_H)=(q_H^{-1}h q_H,q_H),\qquad
 T(y,q_H)=(q_H^{-1}yq_H,q_H).                         \tag{P3K.65.OR.22}
\]
The inverse uses \(q_Hyq_H^{-1}\). It is smooth at y=0 and
orientation preserving. The same four-coordinate multipliers
(P3K.65.OR.21) intertwine this map with the actual parameter action,
including any odd loop; thus no additional parameter transition
is inserted between the two native products. In particular the
direction is \(\mathcal T_-\) to \(\mathcal T_+\).

\item Use the oriented relative generators of P3K.23,
\[
 L(y,q_H)=(y,\zeta(y)q_H),\qquad
 R(y,q_H)=(y,q_H\zeta(y)),
\]
where \(\zeta:(D^3,S^2)\to(S^3,1)\) has degree one. They
commute and extend over \(D^3\times D^4\), so their closed
sphere classes are zero. The central disc generator \(\mathcal C\)
has the positive oriented class 1; its sign is fixed by the signed
plumbing proof in P3K.23, with the same complex exterior orientation.
The exact relative conjugation is
\[
 T(L^\mu R^\nu\mathcal C^s)T^{-1}
       =L^{-\nu}R^{\mu+2\nu}\mathcal C^{s+\nu}.     \tag{P3K.65.OR.23}
\]
This group identity imposes no circle-quotient restriction on the
selected curve. The actual native pair \((-m,m)\) therefore
gives
\([\Sigma_+]=[\Sigma_-]+m\). By the independent computation
(P3K.65.OR.8), \([\Sigma_+]=0\). Thus
\([\Sigma_-]=-m=-3k\), proving the first row of (P3K.65.OR.1).

\item The exact outer obstruction which the old charged-cap argument
omitted is visible in the same charts. On the fixed branch at
the second vertex, the C0 native coordinates have base
\(\kappa w_0\), negative normal \(\alpha v_0\), and positive
normal \((\kappa\alpha)^{-1}u_0\). The C2 native coordinates
have base \(w_0\) and normals \((v_0,u_0)\). Their map has
base multiplier \(\kappa^{-1}\), normal matrix
\(\operatorname{diag}(\alpha^{-1},\kappa\alpha)\), and
positive-weight determinant
\[
                  \kappa\alpha/\bar\alpha,
                  \qquad\operatorname{wind}=-2m.    \tag{P3K.65.OR.24}
\]
The other outer branch is the identity. A nonvanishing determinant
homotopy between them would have equal winding at both ends.
Thus equal inner charged-cap gauges do not make the full outer
determinant parameter-independent. Equation (P3K.65.OR.23) computes
the sphere correction without discarding that outer map.

\par\medskip\noindent\textbf{OR65.6. The actual shorter source/line/cap endpoint.}\par\nopagebreak[4]

\item We reproduce the endpoint map needed for the second row, retaining
the component divisor. Let B0 be the actual unscaled contracting
lift over \(-3G\), with zero on source C1, so its inverse over
\(+3G\) has pole on source C0. An actual rational lift U over
\(+G\) is supplied in the section frame by the divisor identity
\[
 \operatorname{div}\frac1{x-x(G)}=2O-G-(-G)
                 =t_G^*(G-O)-(G-O),
\]
taking the inverse of the resulting line isomorphism. Translation
extends on the relatively minimal regular source; its graph
preserves the canonical two-form, and cannot have an exceptional
curve over a smooth centre. In the actual line graph its forward
source shift is j to j+1. Normalize by an integral base monomial
so its line valuation is zero. Its two defects are
\(h_j-h_{j+1}=-1,0\). Thus its vertical pole is C0, through both
node charts. The third translated norm has divisor
\(-2C_0-C_1\), and
\[
                \rho=B_0^{-1}/U^3,\qquad
                \operatorname{ord}_{q=0}\rho=1.     \tag{P3K.65.OR.25}
\]
The ratio is an actual rational base function because its maps
cover the same source translation and their ratio is constant
on every compact generic elliptic fibre.

\item On the actual order-m good root, the central Cartier line has
deck weights \(\varrho,\varrho^{-1},1\) on the cyclic orbit
of G', O and the third order-three point, with
\(\varrho=\zeta_3^2\). An invertible lift L of translation
\(+G'\) has successive weight ratios \(\varrho^{-1}=\zeta_3\).
Finite cyclic averaging removes the nonconstant commutator unit;
an invariant cube root fixes its exact norm. Hence
\[
                    DLD^{-1}=\zeta_3 L,\qquad
                    L^3=B_0^{-1}.
\]
Write the pulled-back U as \(z^e a(z)L\). The target-scalar
conjugation rule \(D m_fD^{-1}=m_{f(\zeta_m^{-1}z)}\) gives
\[
 \zeta_m^e=\zeta_3,\quad a(\zeta_mz)=a(z),\quad
 \operatorname{ord}_{t=0}\rho=-3e/m\equiv-1\pmod3.  \tag{P3K.65.OR.26}
\]
At the type-II and other ordinary fibres, the vertical divisor
is a whole-fibre divisor, so the order of rho is divisible by
three. These statements retain the original d-character, which
is scalar and commutes with L.

\item Put \(h_0=t^n=q^{-n}\). The rational divisor of
\(h_0\rho^{-2n}\) is divisible by three at every fibre: its
cusp order is \(-3n\) modulo three, its IV* order is \(3n\)
modulo three, and the remaining orders are zero modulo three.
Choose the actual rational root
\[
                     g_0^3=h_0\rho^{-2n}.
\]
The endpoint-to-initial ordinary map is
\[
 C_r=r^{2nk}g_0(t)U^{-2n}:X_{\zeta_3r}\longrightarrow X_r.
                                                               \tag{P3K.65.OR.27}
\]
Its whole-fibre powers cancel at every ordinary fibre. On the
IV* root insertion of (P3K.65.OR.26) leaves precisely the base power
\(z^{mn/3}\), an invariant unit and \(L^{-2n}\).
The ordinary-to-cap map (P3K.65.OR.5) changes under
\(r\mapsto\zeta_3r\) by the logarithm of that positive base
power. Consequently its exact square is
\[
 \Phi_r C_r=C_{\rm cap}\Phi_{\zeta_3r},\qquad
 C_{\rm cap}=(\text{invariant unit})L^{-2n}.
\]
The bounded map extends, and
\[
 C_{\rm cap}A_{\zeta_3r}C_{\rm cap}^{-1}
       =\zeta_3^n\zeta_3^{2n}A_r=A_r.               \tag{P3K.65.OR.28}
\]
This proves the cap descent for both signs, with its original
character and the specified closing direction.

\item The cusp vector of (P3K.65.OR.27) is \((-2n,-n)\). Write
\(h_{\rm src}\) for the actual source deck with logarithmic period u;
it is distinct from the primitive H-deck h in (P3K.65.OR.3).
Reindex by \(h_{\rm src}^n\), of vector \((2n,n)\), to make it
bounded on every toric chart. Retain all logarithms by putting
\(R=\log\rho/(2\pi i)\), \(L_U=\log U/(2\pi i)\),
\(g_*=\log g_0/(2\pi i)\), coherently with the norm. Their
literal identities are
\[
 3L_U=-\gamma+mv\delta-R\delta,\qquad
 3g_*=-ns-2nR.
\]
The bounded logarithm of the reindexed closing is therefore
\[
 nu-2nL_U+(2nkv+g_*)\delta
       =\frac{nK}{3}-\frac{ns}{3}\delta
       =\tau(v),\qquad
 \tau(v)=\frac n3(D+2mv\delta).                    \tag{P3K.65.OR.29}
\]
This is an equality of the full units, including their constants.
Change the constant cube root by
\(\exp(-4\pi i nk/3)\), an actual cube root of unity, and write
\(a_*=-2nk/3\). The closing becomes the bounded translation
\(\tau(v)+a_*\delta\). The same extending origin as above obeys
\[
 \sigma(v+1/3)-\sigma(v)=\tau(v)+a_*\delta,\qquad
 \sum_{j=0}^2\{\tau(v+j/3)+a_*\delta\}=nD+2nmv\delta.
                                                               \tag{P3K.65.OR.30}
\]
On the ordinary chart the three closing maps compose exactly to
\(h_0B_\lambda^{2n}\). The second factor is a literal period;
the first is the prescribed one-r-turn closing. Work first over
the universal v-plane, quotient by that one-turn gluing, then
by the resulting order-three parameter deck. Its action is free
because r is nonzero. This constructs the proper holomorphic
family over \(w=r^3\), with \(\lambda=w^k\), and retains both
fans, all source and cap maps, and both vertices at once.

\par\medskip\noindent\textbf{OR65.7. The three native seams and the shorter sphere class.}\par\nopagebreak[4]

\item Use the literal plus translation \(\xi=z+\sigma(v)\) in the
additive torus coordinates of (P3K.65.OR.30); in characters it is
multiplication by \(\exp(2\pi i\sigma(v))\). The closing sends
\(z_{\rm end}\) to \(z_{\rm end}+c_{\rm add}(v)\), where
\(c_{\rm add}(v)=\tau(v)+a_*\delta\). Since
\(\sigma(v+1/3)-\sigma(v)=c_{\rm add}(v)\), it gives exactly
\(\xi_{\rm initial}=\xi_{\rm end}\). The forward map is the
inverse of this closing. The C2 row of
(P3K.65.OR.10) remains \((a^{-1},ab,ac)\), so these coordinates give
an actual closed native C2 product. They do not supply the
quarantined whole-plumbing equivariant product.

\item Here are all three intermediate seams of its pullback. Write
\(c(v)=\exp(2\pi i\{\tau(v)+a_*\delta\})\) for the bounded
closing, and \(e(v)=c(v)^{-1}\) for its forward map, after the
actual source-period reindexing. On the j-th cut interval of the
triple pullback use the same single product and the transition
\[
            P_j(v)=\prod_{h=0}^{j-1}e(v+h/3)^{-1},
            \qquad j=0,1,2,3.                       \tag{P3K.65.OR.31}
\]
The factors commute because they are actual base-dependent torus
translations. At the seam from j to j+1 the outgoing product
gains exactly \(e(v+j/3)^{-1}\); hence its outgoing chart is
the incoming chart multiplied by the prescribed closing. This
is the literal pullback identification, with every collar
derivative retained. Equation (P3K.65.OR.30) gives
the complete endpoint list
\[
\begin{array}{c|c}
 j&\log P_j/(2\pi i)\\ \hline
 0&0\\
 1&nD/3+(2nkv-2nk/3)\delta\\
 2&2nD/3+(4nkv-2nk/3)\delta\\
 3&nD+6nkv\delta .
\end{array}
\]
In particular
\[
 P_3(v)=\exp(2\pi i(nD+2nmv\delta)),
\]
which is precisely the bounded one-r-turn closing after the
actual period reindexing. The ordinary identity
\(h_0B_\lambda^{2n}\) and the cap square (P3K.65.OR.28) identify the
full family, not just these bounded charts.

\item Define the covered r native product to be this triple pullback.
Each coordinate change commutes with H, K and delta and leaves
the complete transition coefficients in (P3K.65.OR.10) unchanged.
Thus its doubled ordinary construction (P3K.65.OR.12)--\allowbreak (P3K.65.OR.20)
computes the same covered free pair \((-3k,3k)\) directly in
these charts. No identification with a presumed whole-plumbing
product is used. A connection with the single w native product,
pulled back to the cover, has relative return \(f_w^3\).
The free quotient of the relative group is \(\mathbb Z^2\),
so its exact free coordinates satisfy
\[
                     3(\mu_w,\nu_w)=(-3k,3k),
                     \qquad(\mu_w,\nu_w)=(-k,k).    \tag{P3K.65.OR.32}
\]

\item For the positive fan choose the actual closed fixed-curve native
product; the vertex maps and their standard fixed-curve unit
normalization close by the same endpoint (P3K.65.OR.30). Its pullback
gives the covered physical circle Euler line and tangent bundle.
The degree-three covering of the completed conifold mapping tori
therefore gives \(3A_{+,w}=A_{+,r}=0\) as an equality of integer
characteristic integrals. Thus \(A_{+,w}=0\) and its actual
Hopf lift has sphere class zero. This is integer naturality on
constructed mapping tori, not division of an element of Theta7.

\item The same Cox arrow (P3K.65.OR.21), now with the shorter endpoint,
carries the native flop and both cap maps. Applying the oriented
identity (P3K.65.OR.23) to (P3K.65.OR.32) gives
\([\Sigma_{+,w}]=[\Sigma_{-,w}]+k\). Hence
\([\Sigma_{-,w}]=-k\), completing (P3K.65.OR.1). For k=1 this is
the actual positive loop of the scalar itself, with class 27.
The integers k=1 through 28 supply every residue; each is an
actual admitted cap order and completed family.

\item The symbolic certificate records the source identity, period
change, all three full-unit seam rows, both norm-sign cap squares,
the doubled trace degrees, actual free-coordinate equations and
three-step norm. Its computations supplement the explicit maps
above. No claim for a third fan or an arbitrary fan is made here.

\end{enumerate}
\end{proof}




\begin{theorem}\label{night:P3K.66}
\StatementRecord{night:P3K.66}{proved}{S6\_06\_BLIND: GPT-6 Astra (OpenAI, ChatGPT Pro mode, effort Pro), recorded re-derivation in the checkpoint dated 18 November 2026}{\cref{night:P3F.22}, \cref{night:P3F.54}, \cref{night:P3F.8}, \cref{night:P3K.23}, \cref{night:P3K.42}, \cref{night:P3K.46}, \cref{night:P3K.6}, \cref{night:P3K.63}, \cref{night:P3K.68}, \cref{night:P3K.7}}
Retain the actual positive selected-cusp source loop, its fixed scalar,
finite caps and physical vertical-edge fan \(T_n=T(e_1+ne_2,e_2)\).
Let \(M\) be the physical circle line and \(U=M^*\) its meridian-one
dual. For every integer \(n\), the actual Jin pairs and indices are
\[
\begin{aligned}
(k_U,\ell_U)&=(6+18n(n-1),0),&A_U&=1+3n(n-1),\\
(k_M,\ell_M)&=(-6-18n(n-1),0),&A_M&=-1-3n(n-1).
\end{aligned}
\]
The physical oriented sphere class is \(A_M\pmod {28}\).
The native fixed-star free pair is \((0,-6-18n(n-1))\), and the
actual ordinary plumbing entries are
\(V_{00}=18n(n-1)\), \(\delta_0=-12-18n(n-1)\).
The proof extends the prescribed boundary line arrow through the
actual completed piece and cap, retaining the anti-equivariant
physical/dual circle map and both orientations.
\end{theorem}
\begin{proof}
\begin{enumerate}[label=\arabic*.]
\item \par

\medskip\noindent\textbf{EC.1. The physical line, its meridian dual, and both transfer arrows}

\par\nobreak

\item Let \(M\) be the nonzero-vector line of the physical principal \(\delta\)
circle, with coefficient \(V=|V|e^{2\pi i d}\). Its real connection
\(A\), normalized by \(A(\partial_d)=1\), gives the unitary connection
\(d+2\pi i A\). Hence the two actual line classes are
\[
 x_M=c_1(M)=-dA,\qquad U=M^*,\qquad x_U=c_1(U)=dA.       \tag{EC.1}
\]
The Hermitian dual arrow sends a nonzero vector \(v\) to the
functional \(w\mapsto\langle v,w\rangle/\langle v,v\rangle\),
with the inner product linear in its second argument. It reverses
the physical circle character and extends over both product caps
by the same formula, carrying every boundary scalar to its inverse.

\item In the actual Cox coordinate
\(q_H=(1+jZ)(c_0+j\overline{c_1})\), the physical action is
\(q_H\mapsto q_He^{-2\pi i t}\). At \(q_H=1\) its vector is
\(-i\), and the Hopf quotient \(z=q_Hiq_H^{-1}\) has derivatives
\(j\mapsto-2k\), \(k\mapsto2j\). Therefore
\[
 [S^3_q]=-[S^1_\delta][S^2_z],\qquad
 \Omega_{\mathbb C^2}=d\delta\wedge dr\wedge\Omega_{S^2_z}.
                                                               \tag{EC.2}
\]
The actual outward meridian is \(S^2_z\). Its connection curvature
has integral \(+1\), so the meridian-one line is \(U\), whereas
\(M\) has value \(-1\). On a regular source page the physical
line restricts to the original source line, and the meridian-one
line restricts to its dual. In particular the corrected K6 page
arrow is \(s^*x_U=-c_1(M_{\rm source})\). The squared page terms
of K7 are unchanged by that explicit duality.

\item We retain the gluing inverse when comparing the two Hopf transfers.
The complex exterior is \(D^3_y\times S^3_q\), and the physical
sphere is \(-\partial(D^3_y\times D^4_q)\). In F22's physical-line
marking choose \(S^2_M=-S^2_z\), \(y_M=-y\), and \(u_M=q\).
Both base-factor reversals preserve the quotient orientation, and
\(\partial(D^4_{u_M}\times D^3_{y_M})\) is the physical sphere.
For the dual line choose \(S^2_U=S^2_z\), the reversed circle,
\(y_U=y\), and \(u_U=q\); this product boundary has opposite
physical orientation. These markings are related by the displayed
anti-equivariant dual map and the reflection of the \(y\) factor,
which extends over its four-ball.

\item Write
\(\Sigma_F=(D^4_y\times S^3_q)\cup_F(S^3_y\times D^4_q)\),
with first boundary point identified with its image under \(F\)
on the second piece. The physical F22 pieces
\(D^4_{u_M}\times S^3_{y_M}\) and
\(S^3_{u_M}\times D^4_{y_M}\) identify respectively with the
negative of the second and first \(\Sigma\) pieces. The dual
pieces identify with the same exchanged pieces with positive
signs. Rewriting the same gluing relation in first-piece order
inverts \(F\). Thus
\[
 \lambda_M(F)=-\Sigma_{F^{-1}}=\Sigma_F,\qquad
 \lambda_U(F)=+\Sigma_{F^{-1}}=-\Sigma_F.                 \tag{EC.3}
\]
For the last equalities the actual relative \(L,R\) factors extend
the seven-ball; for the remaining disc-supported factor the two
opposite product collars cancel to give
\(\Sigma_{F^{-1}}=-\Sigma_F\). F8's cut-out ball lies in its
first piece and uses the literal relation
\(x_{\rm ball}\sim F(x_{\rm ball})_{\rm exterior}\), agreeing
with the physical-line marking above. Therefore the physical
sphere class is \(A_M=-A_U\) modulo 28, where
\(24A_M=4k_M-\ell_M\) and similarly for \(U\).

\item Keep the contracting suspension coefficient \(\gamma_B\).
The complex source-line orientation is
\(\Omega_S\wedge d\log|V|\wedge dd\). Taking \(\delta\) first
and using the negative radial derivative of the contracting
coefficient gives orbit orientation \(\Omega_S\wedge d\gamma_B\).
The actual physical positive-parameter orientation is consequently
\[
 \Omega_+=dr_q\wedge d\alpha\wedge du\wedge dv
                   \wedge d\gamma_B\wedge d\theta.       \tag{EC.4}
\]
The source fibre is \(u+\tau v\), with \(\operatorname{Im}\tau>0\),
so the source orientation in this formula is positive. Jin's endpoint
convention reverses the parameter. Hence
\[
 k_U=-\int_{\Omega_+}(dA)^3,\qquad
 k_M=+\int_{\Omega_+}(dA)^3=-k_U.
\]
The analogous relation holds for \(\ell\). We compute the curvature
of \(U\) in the auxiliary orientation \(\Omega_+\), retaining
this exact orientation dictionary throughout.

\par

\medskip\noindent\textbf{EC.2. Full periods and exact boundary line coordinates}

\par\nobreak

\item Retain SQ.1--2 in full:
\[
 h:(Z,V)\mapsto(\kappa qZ,\alpha_L V),\qquad
 B:(Z,V)\mapsto(gZ,\eta qV),
\]
\[
 Z_n=Z,\quad V_n=Z^{-n}V,\quad h_n=hB^n,\quad B_n=B,
\]
\[
 Q_n=\kappa qg^n,\quad
 \alpha_{L,n}=\alpha_L\eta^n\kappa^{-n}g^{-n^2},
 \qquad\eta_n=\eta g^{-n}.
                                                        \tag{EC.5}
\]
The symbol \(\alpha_L\) here denotes the line multiplier; it is
not the angle \(\alpha\). The \(\rho\) windings of
\((\kappa,\alpha_L,g,\eta)\) are \((-6,0,0,0)\).
The actual periodic logarithms of every zero-winding factor are
retained.

\item Here is a precise connection choice that retains them. On the regular
boundary, use the actual real period coefficients
\[
 (z,\ell)=u(1,0)+v(\tau,\mu)+d(0,1)+\gamma(G,\beta),
 \qquad z=\frac{\log Z}{2\pi i},\quad
        \ell=\frac{\log V}{2\pi i},
                                                        \tag{EC.6}
\]
where
\(\tau=\log(\kappa q)/(2\pi i)\),
\(\mu=\log\alpha_L/(2\pi i)\),
\(G=\log g/(2\pi i)\) and
\(\beta=\log(\eta q)/(2\pi i)\).
The imaginary parts give exactly the invertible matrix in SQ.4a;
the real parts then determine \(u,d\). In particular
\[
 d=\operatorname{Re}\ell-\operatorname{Re}\mu\,v
                           -\operatorname{Re}\beta\,\gamma.
                                                        \tag{EC.7}
\]
Differentiating (EC.7) retains the differentials and values of both
unit logarithms. Thus using \(d\) as a coefficient coordinate does
not set their phases or their derivatives to zero.

\item The q and \(\rho\) coordinate changes are respectively
\[
 (u,d)\mapsto(u-v,d-\gamma),\qquad
 (u,d)\mapsto(u+6v,d).
                                                        \tag{EC.8}
\]
The complete zero-winding logarithms are single-valued and introduce
no further integer change. Consequently
\[
 A_0=dd+\alpha\,d\gamma
                                                        \tag{EC.9}
\]
is an actual global real principal connection on this boundary.
In physical \(V\) coordinates it contains all terms obtained by
differentiating (EC.7). It is a permissible connection choice on the
actual bundle, not an omission of its units.

\item Substitution in (EC.6), including both complex coordinates, gives
exactly
\[
 u_n=u,\quad v_n=v,\quad\gamma_n=\gamma-nv,\quad d_n=d-nu.
                                                        \tag{EC.10}
\]
In the new coordinates the corresponding real connection is
\[
 A_n=dd_n+\alpha\,d\gamma_n+6n\theta\,dv.
                                                        \tag{EC.11}
\]
Its \(\rho\) change is \(d_n\mapsto d_n-6nv\), and its q change is
\(d_n\mapsto d_n-\gamma_n\), so (EC.11) is invariant under the
literal two deck arrows. Equations (EC.9)--(EC.11) give the exact
global difference
\[
 \beta_n=-n\{du+(\alpha-6\theta)dv\}.
                                                        \tag{EC.12}
\]
It is invariant under both changes (EC.8). If one uses the source
coframe \(du+(\alpha-6\theta+\phi_\kappa)dv\), then the additional
one-form \(n\phi_\kappa dv\) is present in (EC.12). It has not been
deleted. Formula (EC.7) explains the same fact for every other unit.

\par

\medskip\noindent\textbf{EC.3. The relative Cech line isomorphism}

\par\nobreak

\item Let \(\mathcal C\) be the six-dimensional family of completed cusp
quotient pieces for \(T_0\), over the positive \(\theta\) circle. Use
SQ.1--2 to identify each actual \(\mathcal C_n\) with this fixed
smooth piece. Those arrows are periodic, preserve the outer real
period marking, and retain the typed Cartier spinor map and the
physical Hopf cap. Its boundary is denoted \(\mathcal B\).

\item Choose the graph map \(b:\mathcal C\to\mathbb R/\mathbb Z\) of
SQ.3, with boundary value \(v\). The same proof modifies it by a
supported real function to be constant near the fixed tube and
extends it over the cap. It is periodic in \(\theta\). Define an actual
complex line \(\mathcal T\) on \(\mathcal C\) by cutting \(\theta\)
and gluing its coefficient by
\[
 (\theta,p,w)\longmapsto
       (\theta+1,p,e^{2\pi i b(p)}w).
                                                        \tag{EC.13}
\]
In these cut unit frames its connection is
\(\nabla_{\mathcal T}=d-2\pi i\theta\,db\). The frame change at
\(\theta\) one is the inverse of the displayed coefficient change, so
this connection descends exactly. Its Chern curvature is
\[
                         \Xi=d\theta\wedge db,\qquad \Xi^2=0.
                                                        \tag{EC.14}
\]

\item The actual normalization constant factor
\(\alpha_{L,n}/\alpha_L\) has \(\rho\) winding \(6n\). In the inverse
line \(\mathcal U\), its Chern class is consequently
\(+6n[d\theta]\smile b\). The other classes of the line difference
are zero: the typed SQ.8 Hopf arrow preserves the meridian value;
the regular \(b\times\gamma\) torus has the integral period section
from SQ.3; and the \(\gamma\) times \(\theta\) value is zero because
\(\eta_n/\eta=g^{-n}\) has its periodic logarithm. The integral
Thom, Mayer--Vietoris and Wang calculation SQ.16 has no torsion
and shows that these are all the evaluations. Thus
\[
 c_1(\mathcal U_n)-c_1(\mathcal U_0)
                   =6n\,c_1(\mathcal T)
             \quad\hbox{on }\mathcal C.               \tag{EC.15}
\]
This is also the direct \(\theta\)-cut/graph-cut clutch square: the
physical \(\delta\) coefficient changes by \(-6n b\), so its dual
coefficient changes by \(+6n b\), exactly (EC.13).
The zero-winding logarithms left in (EC.5) are
\[
 nf_\eta-nf_\kappa-n^2f_g,\qquad -nf_g,
                                                        \tag{EC.16}
\]
together with the actual source-flow pullbacks of their extending
units. They are retained in the bundle whose class is (EC.15).

\item We next prove the stronger relative assertion with its prescribed
outer line arrow:
\[
 J_n:\mathcal U_n\longrightarrow
                    \mathcal U_0\otimes\mathcal T^{6n}.
                                                        \tag{EC.17}
\]
On \(\mathcal B\), this arrow is specified in the doubly cut real
period frames by (EC.10)--(EC.11) and the clutch (EC.13).
Their deck changes agree exactly, so it is a genuine boundary
bundle isomorphism. The assertion is that this particular boundary
arrow extends, including its integral gauge; an absolute equality
of Chern classes alone would not prove that assertion.

\item Here are the exact groups and the only possible relative obstruction.
On \(\mathcal B\), the q monodromy has commutator \(u\) with the
source graph loop \(v\); its coefficient is one. The \(\rho\) monodromy
has commutator \(6u\), and all other retained circle generators
commute after abelianization. Hence
\[
 H_1(\mathcal B;\mathbb Z)
       =\mathbb Z\langle\alpha,\theta,v,\gamma\rangle.
                                                        \tag{EC.18}
\]
The \(u\) generator has been killed primitively, so no finite
summand has been suppressed. SQ.16 and its periodic graph and
\(\gamma\) maps give
\[
 H^1(\mathcal C;\mathbb Z)
       =\mathbb Z\langle d\theta,db,d\gamma\rangle,
\quad
 \operatorname{im}H^1(\mathcal C)
       =\mathbb Z\langle d\theta,dv,d\gamma\rangle
                   \subset H^1(\mathcal B).           \tag{EC.19}
\]
The \(\gamma\) map extends through the cap by its real lift near the
fixed sphere, exactly as for \(b\). A section of the \(\theta\)
family outside the fixed tube splits the Wang sequence, so there
is no extra generator or torsion in (EC.19). Therefore the only
boundary winding that does not extend is the q-angle winding.

\item We test it using an actual disc, rather than a rational cohomology
argument. Choose a constant Tate section \(Z=Z_*\), with
\(|Z_*|=1\), at a smooth normalization point avoiding the
marked zero section, the retained Picard-section divisors and
the two nodal branches. This can be done uniformly: the explicit
Tate points of F54 remain near \(1\) for the chosen small \(\rho\)
circle, whereas a point near \(-1\) stays disjoint from them;
shrink the common q disc once more if necessary. Take the marked \(\gamma\)
value zero. The fixed graph's short arc in SQ.4c is centred at
one half, so this value lies outside it. The section in the
source neighbourhood extends across q zero through a smooth
point of the nodal fibre. Together with this \(\gamma\) value it is
a disc in the orbit piece, disjoint from the fixed tube and its
surgery; it therefore lies in the completed piece as well.
The SQ.2 source horizontal field can be taken to preserve this
section. On the outside collar its real marked-period field is
already zero on \(v=0\); near the node the section is disjoint.
Adding this prescribed zero field to the stated partition of
unity leaves all the node, boundary and cap arrows unchanged.

\item More generally, at a fixed \(\gamma\) value in this free interval the
suspension-untwisted coordinates give the physical source point
\[
             Z_{\rm phys}=Z_*\exp(\gamma\log g).
                                                        \tag{EC.20}
\]
This is a nonzero extending function on the q disc; its logarithm
uses the actual periodic \(\log g\). The line character arrow
on it is literally multiplication by \(Z_{\rm phys}^{-n}\).
The new untwisting factor
\(\eta_n^\gamma=\eta^\gamma\exp(-n\gamma\log g)\)
cancels the displayed \(g\) factor on passing to the common
untwisted line frame. Its remaining factor is \(Z_*^{-n}\).
Equivalently one can retain both factors in (EC.20); both extend
and are nowhere zero. No fractional power of \(q\) is used in
this ratio: the common \(q^\gamma\) factors cancel before the
ratio is formed. The disc lies in one free suspension interval,
so no B reindex is introduced. The tensor line \(\mathcal T\)
is trivial on this disc at fixed \(\theta\) and \(b=v=0\).
For the chosen value \(\gamma\) zero, the untwisting is literally the
identity: the disc is the ordinary source section with \(V\) in
one fixed nonzero line chart. We can take \(V=1\). In the two
actual first vertex charts this disc and its image have coordinates
\[
 (q/Z_*,Z_*,1)\longmapsto
                 (qZ_*^{\,n-1},Z_*,Z_*^{-n}).
                                                        \tag{EC.20a}
\]
Both charts have physical height equal to the product of their
three coordinates. Their second and third coordinates are nonzero
at q zero, so this is the free locus and no fixed-tube or cap
coordinate is involved. The inverse is given by the opposite
nonzero powers of \(Z_*\). The dual line arrow is multiplication
by \(Z_*^{-n}\). This gives an extending isomorphism
with the unchanged physical \(\delta\) frame. The prescribed boundary
arrow in (EC.17) thus extends over this disc by this explicit line
map. Its q-angle winding is zero.

\item For completeness, an explicit Cech construction turns these facts
into the full isomorphism. On a good cover of \(\mathcal C\), use
the actual unitary transition functions of
\(\mathcal R_n=\mathcal U_n\otimes\mathcal U_0^{-1}
\otimes\mathcal T^{-6n}\), formed from all the retained maps.
Choose real logarithms on the overlaps. Their Cech coboundary
is an integral two-cocycle. Equation (EC.15) says that cocycle
is the coboundary of an integral one-cochain; subtract that
one-cochain from the chosen logarithms. Our local frame convention
is \(e_j=\exp(2\pi i r_{ij})e_i\). The resulting real
one-cocycle \(r_{ij}\) is solved by
\[
                  a_i=\sum_k\rho_k r_{ki},\qquad
                  a_j-a_i=r_{ij},                    \tag{EC.21}
\]
for a partition of unity subordinate to the cover. Multiplying
the local unit frames by \(\exp(-2\pi i a_i)\) gives a global
unit frame. Since
\(\mathcal R_n=\operatorname{Hom}
(\mathcal U_0\otimes\mathcal T^{6n},\mathcal U_n)\),
this frame is a nowhere-zero map in that direction. Its inverse
is the absolute isomorphism \(J_n\) in (EC.17).
Thus the direction of the required map is fixed by the actual
transitions, including (EC.16).

\item Compare its boundary value with the prescribed boundary arrow.
Their ratio is a circle-valued function. By the section-disc
calculation its \(\alpha\) coefficient is zero. By (EC.18)--(EC.19)
its other three winding coefficients are integral multiples
of the restrictions of \(\theta\), \(b\), and \(\gamma\). Multiply the
absolute isomorphism by the corresponding extending functions
\(\exp(2\pi i(a\theta+cb+d\gamma))\) to remove these coefficients.
The remaining boundary ratio has a single real logarithm.
Extend that logarithm to the collar, multiply it by a cutoff,
and set it to zero in the remaining interior. Multiplication
by its exponential makes the boundary map exactly the
prescribed one. This proves (EC.17) as a relative bundle map.
It retains every possible boundary integer rather than declaring
the zero-winding unit terms harmless without an extending arrow.

\item The cap is included in this construction. Its actual line is
the one carried by the SQ.8 Cartier/Hopf map. The supported
change \(b-b'\) is implemented on the \(\theta\) cut by its explicit
exponential in SQ.3. A residual circle gauge on the inner cap
boundary has only its \(\theta\) winding, and this extends over
\(D^3\times S^2\times S^1_\theta\). Thus no additional cap
line class or unspecified framing is inserted in (EC.17).

\par

\medskip\noindent\textbf{EC.4. The affine curvature family and the oriented collar integral}

\par\nobreak

\item Choose a connection on \(\mathcal U_0\) whose boundary value is
the one just fixed, and write its real Chern curvature as \(F_0\).
On the fixed completed piece use the relative isomorphism
(EC.17) and the tensor connection of (EC.13). Their curvature is
exactly
\[
                         F_n=F_0+6n\Xi .
                                                        \tag{EC.22}
\]
All connections here are on the actual lines. They need not be
Chern connections of the original complex structure; Chern--Weil
evaluation is independent of that choice. The relative map
(EC.17) ensures their specified boundary values agree with
the actual outer gluing. In particular no omitted Mayer--Vietoris
gauge can add a further quadratic or other term.

\item Keep the original connection on the common exterior, where the
fan change is absent. In piece coordinates its pullback on
\(\mathcal B\) is displayed below. From this point onward, after
pullback by the actual piece map, we write \(\gamma=\gamma_n\)
for the common piece coefficient. The exterior coefficient is
\(\gamma_{\rm ext}=\gamma+nv\), and its \(\delta\) coefficient is
\(d_{\rm ext}=d_n+nu\). Thus the \(\gamma\) in the curvature
formula is the new piece coefficient, not the exterior coefficient.
With this dictionary the literal pullback is
\[
 A_{\mathrm{ext},n}
   =dd_n+\alpha\,d\gamma_n+n(du+\alpha\,dv),
\quad
 F_{\mathrm{ext},n}
   =d\alpha\wedge d\gamma+n\,d\alpha\wedge dv.
                                                        \tag{EC.23}
\]
The difference between the piece and exterior connections is
the global form (EC.12). It has derivative
\[
 d\beta_n=-n\,d\alpha\wedge dv+6n\,d\theta\wedge dv .
                                                        \tag{EC.24}
\]
Directly,
\[
 F_{\mathrm{ext},n}^2=0,\quad (d\beta_n)^2=0,
\quad F_{\mathrm{ext},n}+d\beta_n
            =d\alpha\wedge d\gamma+6n\,d\theta\wedge dv .
                                                        \tag{EC.25}
\]

\item Orient \(\mathcal B\) as the outward boundary of the piece in
(EC.4):
\[
 \Omega_{\mathcal B}
       =d\alpha\wedge du\wedge dv\wedge d\gamma\wedge d\theta .
                                                        \tag{EC.26}
\]
Each marked angular coordinate has period one. Its integral
over the actual boundary is one: the q/\(\rho\) monodromies shear
\(u\) by integral multiples of \(v\), and hence preserve this
volume form and its unit fundamental cell. The full unit
periods have already entered the coordinate construction
(EC.6); they do not alter this cell or its primitive lattice.

\item Let \(t\) increase across a gluing collar from the piece toward
the exterior. Its orientation is \(dt\wedge\Omega_{\mathcal B}\).
Interpolate by
\[
 A(t)=A_{\mathrm{ext},n}+\chi(t)\beta_n,\qquad
                  \chi(0)=1,\quad\chi(1)=0.
                                                        \tag{EC.27}
\]
It is constant near each endpoint. The curvature and its cube
are
\[
 F=F_{\mathrm{ext},n}+\chi d\beta_n
                           +\chi' dt\wedge\beta_n,
\]
\[
 F^3=3\chi'\,dt\wedge
          \beta_n\wedge(F_{\mathrm{ext},n}+\chi d\beta_n)^2 .
                                                        \tag{EC.28}
\]
The term without \(dt\) vanishes on the five-dimensional boundary.
The exact exterior-algebra identity is
\[
 3\int_0^1
   \beta_n\wedge(F_{\mathrm{ext},n}+s\,d\beta_n)^2\,ds
                    =18n^2\,\Omega_{\mathcal B}.        \tag{EC.29}
\]
The actual interpolation in (EC.27) runs from \(s=1\) to \(s=0\).
Its contribution is therefore \(-18n^2\), with the sign fixed
by (EC.26)--(EC.28).

\item The interior contribution is linear because \(\Xi^2=0\):
\[
 F_n^3-F_0^3=18n\,F_0^2\wedge\Xi .
                                                        \tag{EC.30}
\]
The common exterior contributes no change. Its boundary map
has already been accounted for by (EC.23)--(EC.29).
There is consequently one real constant \(L\), independent
of \(n\), such that the physical-line Jin Euler cubes satisfy
\[
                         k_{M,n}=k_{M,0}+Ln-18n^2 .
                                                        \tag{EC.31}
\]
This conclusion follows from the global relative line map
and the actual oriented gluing. It is not a polynomial assigned
from the local fifth form alone.

\par

\medskip\noindent\textbf{EC.5. Actual endpoint evaluations and exact integers}

\par\nobreak

\item The corrected Cartier-to-Cox arrow in K42 is anti-linear:
\(q_H=g(x,y)\overline{(a,b)}\), with the actual retained matrix
\[
 g=\frac1{\sqrt{|x|^2+|y|^2}}
       \begin{pmatrix}y&-\overline x\\x&\overline y\end{pmatrix}.
\]
The quaternionic anti-linear map
\((a,b)\mapsto(\overline b,-\overline a)\) commutes with every
\(SU(2)\) clutch and extends both caps. It reverses both the
projective meridian and the associated line. Their two signs cancel
in the characteristic integrals. The conjugate matrix is a constant
\(SU(2)\) conjugate of \(g\), so its relative second Chern number
remains \(-1\), and its adjoint Pontryagin number is \(+4\).
Thus the normal sphere-bundle contribution for \(U\) is \((1,4)\),
while the actual inner Seifert cap contributes \((5,-4)\).
The normal-tube boundary and physical-to-Jin reversal of K42 give
\((k_{U,0},\ell_{U,0})=(6,0)\), hence
\((k_{M,0},\ell_{M,0})=(-6,0)\). These are the exact corrected
cap computations of K68, not values selected from a residue.

\item The fan \(T_1\) is literally K46's reflected fan: its first
triangle is \((0,e_1+e_2,e_2)\), and translating its second
triangle by \(-e_2\) gives \((0,e_1,e_1+e_2)\). This is the
actual \(B\) period including its full unit map. K46's oriented
cap and line map preserves each of the two line conventions.
Consequently \(k_{M,1}=-6\) and \(k_{U,1}=6\). Equation
(EC.31) now gives \(L=18\), independently of any central class.

\item The full result is
\[
 \begin{aligned}
 k_{U,n}&=6+18n(n-1),&\ell_{U,n}&=0,
      &A_{U,n}&=1+3n(n-1),\\
 k_{M,n}&=-6-18n(n-1),&\ell_{M,n}&=0,
      &A_{M,n}&=-1-3n(n-1).
 \end{aligned}                                                   \tag{EC.32}
\]
Here K63 proves the zero Pontryagin pairing for either dual line.
The actual physical sphere class is \(A_{M,n}\) modulo 28 by
(EC.3). K23's actual native free generators satisfy
\(k_U=-(m+n_0)\), \(\ell_U=-4m\), equivalently
\(k_M=m+n_0\), \(\ell_M=4m\). The native-to-ordinary product
map adds the already proved left integer six. Therefore
\[
 (m_0^{\rm native},n_0^{\rm native})
          =(0,-6-18n(n-1)),\qquad
 V_{00}=18n(n-1),\qquad \delta_0=-12-18n(n-1).            \tag{EC.33}
\]
These follow from the completed integer line and collar calculation.
They do not use the old source-loop residue, whose initial physical
line was misidentified. The original draft, all earlier audits,
and its output are preserved with the correction history.

\item The certificate \path{code/part3/source_shear_euler_cube.py} checks
the exact periods, integral relative line groups, both Chern-line
signs, oriented collar coefficient, actual reflected fan arrows,
and these corrected polynomial identities. The geometric Cech,
section-disc, normal and cap maps are proved in full above.

\end{enumerate}
\end{proof}




\begin{theorem}\label{night:P3K.67}
\StatementRecord{night:P3K.67}{proved}{S6\_06\_BLIND: GPT-6 Astra (OpenAI, ChatGPT Pro mode, effort Pro), recorded re-derivation in the checkpoint dated 18 November 2026}{\cref{night:P3K.32}, \cref{night:P3K.33}, \cref{night:P3K.35}, \cref{night:P3K.62}, \cref{night:P3K.64}}

Retain the actual whole-zero \(I_1\) source, Cartier line, affine
root caps, finite characters and annular maps of the weighted shorter
construction in P3K.35.2--5. Write the contracting source translation
as \(-bG\), the canonical scalar as \(\lambda=r^M\), and the actual
shorter parameter as \(w=r^b\), with \(b\mid M\). Keep its specified
forward endpoint, including the rational source lift, bounded source
factor and constant physical-delta normalization. Put \(J=M/b\).
For the actual periodic fan \(T_n\) with edge directions
\[
                 e_2,\qquad (1,n),\qquad(1,n-1),
                         \qquad n\in\mathbb Z,
\]
the positive shorter-loop return satisfies
\[
 A_n=A_0-J\frac{n(n-1)}2\pmod{28},
 \tag{P3K.67.0}
\]
where \(A_0\) is the actual reference-fan sphere class. No signed
nonzero reference page-index formula is assigned here.
Its actual native relative free pairs are
\[
 C_1:\bigl(J(1-n),J(n-1)\bigr),\qquad
 C_2:\bigl(nJ,-nJ\bigr).
\]
The two oriented Cox flips go respectively to \(T_{n-1}\) and
\(T_{n+1}\). All source, cap and line factors are retained in their
explicit morphisms. No simultaneous equivariant product of the full
three-curve plumbing is used.

For the original P3K.33 endpoint, with root orders \(3K,4L\),
\(\gcd(3K,4L)=1\), either sphere sign, and its prescribed delta
characters, retain exactly
\[
 w_h=e^{-\pi ic/(KL)},\qquad u=\kappa w_h^{2KL}.
\]
Here \(J=2KL\), \(A_0=0\), and
\[
                       A_n=-KL\,n(n-1)\pmod{28}.
\]
In particular the original \(K=L=1\) member on \(T_2\) has class
\(-2=26\), while \(T_0,T_1\) have class zero. The proof supplies
the actual endpoint and cover maps, not a division in \(\Theta_7\).
This statement assigns neither an arbitrary-fan formula nor separate
characteristic integers on the changed fans.
\end{theorem}
\begin{proof}
\begin{enumerate}[label=\arabic*.]
\item \par

\medskip\noindent\textbf{P3K.67.1. The actual shorter family and its finite-cap square}

\par\nobreak

\item Fix a whole-zero I1 source and marked line from the canonical catalogue.
Write the contracting source translation as \(-bG\), the line as
\(\mathcal M=\mathcal O(aG-O+D)\), and the scalar family as
\(\lambda=r^M\). Here \(b\mid M\). The nontrivial root orders and the
actual order-one charts have exponents
\[
 P_i=Mn_i/m_i,\qquad \sum_iP_i=\varepsilon\in\{1,-1\},
 \quad A_{i,r}=r^{P_i}A_{i,0},\quad B_\lambda=r^MB_0.
 \tag{P3K.67.1}
\]
The sum includes the actual framing-chart exponent \(-Mb_0\) in K32.
The finite cyclic equations are \(A_{i,0}^{m_i}=B_0^{n_i}\) and
\([A_{i,0},B_0]=1\), with every prescribed affine and line character.

\item Use the actual rational lift \(U\) of \(+G\), obtained from the chord
divisor in K35.6. Its norm is the rational base function
\[
 \rho=B_0^{-1}/U^b.
 \tag{P3K.67.2}
\]
On root chart \(i\) its exact comparison to an extending lift is
\[
 U=z^{e_i}h_i(z)L_i,\quad L_i^b=B_0^{-1},\quad
 D_iL_iD_i^{-1}=\chi_iL_i,\quad
 \chi_i=\zeta_{m_i}^{e_i}=\zeta_b^{z_i},\quad
 h_i(\zeta_{m_i}z)=h_i(z).
 \tag{P3K.67.3}
\]
These are maps of the actual corrected root line. In particular the
quadratic I1* line retains the two node scalars \((1,-1)\) and
\(\chi_i=-1\). The good III/III* roots have \(\chi_i=-1\), the good
IV/IV* roots have \(\chi_i=\rho_i^{-a}\), and the reduced partners
whose component class is killed by \(a\) have \(\chi_i=1\). The finite
averaging of a norm-one unit in K35.3 is performed before (P3K.67.3); no
nonconstant unit is discarded as a character.

\item Let \(h_0=e^{2\pi iH}\) be the actual one-gamma \textbf{forward} ordinary
closing, including its bounded constant and all order-one charts.
The orders of \(h_0\rho^\varepsilon\) are divisible by \(b\).
At a root chart the required integer congruence is
\(P_i+\varepsilon z_i=0\pmod b\); at the selected I1 the simple
zero cancels the pole of \(B_0^{-1}\). Retain the specified rational
root \(g_0^b=h_0\rho^\varepsilon\). The actual forward arrow is
\[
 D_r=r^{-\varepsilon M/b}g_0U^\varepsilon:
 X_r\longrightarrow X_{\zeta_b r}.
 \tag{P3K.67.4}
\]
At cap \(i\), the ordinary-to-cap cochain changes by
\(+Mn_i\log z\,\delta/(2\pi ib)\), cancelling the literal ordinary
power \(z^{-Mn_i/b}\). Its remaining map is an invariant unit times
\(L_i^\varepsilon\), and
\[
 L_i^\varepsilon A_{i,r}L_i^{-\varepsilon}
   =\chi_i^{-\varepsilon}A_{i,r}
   =\zeta_b^{P_i}A_{i,r}=A_{i,\zeta_b r}.
 \tag{P3K.67.5}
\]
It commutes with \(B_\lambda\). Thus this is the actual cap square,
not a fibrewise identification without an endpoint. The integral
annular/source-period cochain is retained as
\[
 \frac{P_i\delta-\varepsilon(A_i-I)e_\gamma}{b}
 =(0,-\varepsilon k_{i,2},\varepsilon k_{i,1},
                    (P_i+\varepsilon z_i)/b)^t.
 \tag{P3K.67.6}
\]
This supplies every finite group arrow as well as the ordinary overlap.

\par

\medskip\noindent\textbf{P3K.67.2. The bounded forward endpoint and a literal pullback product}

\par\nobreak

\item Use the actual logarithmic cusp characters, with exponential
\((z_1,z_2)\mapsto(Z,V)=(e^{2\pi iz_1},e^{2\pi iz_2})\).
The marked outward gamma and forward ordinary logarithm are
\[
 \gamma=(\Gamma,C-s-Mv),\quad
 H=\varepsilon s+U_0(q),\quad v=\log r/(2\pi i).
\]
Here \(\Gamma=b\mu_G\) is the fixed bounded Tate lift of \(+bG\);
\(U_0\) denotes a bounded base logarithm and is distinct from the
rational source lift \(U\). The bounded displacement of (P3K.67.4) is
\[
 \tau(v)=\frac{H\delta+\varepsilon\gamma}{b}
   =\left(\frac{\varepsilon\Gamma}{b},
           \frac{U_0+\varepsilon C-\varepsilon Mv}{b}\right).
 \tag{P3K.67.7}
\]
Every entry is holomorphic on the whole cusp disc. Put
\[
 \alpha=\frac{\varepsilon M(b-1)}{2b^2},\qquad
 d(v)=\tau(v)+\alpha\delta,
 \quad \widetilde D_r=e^{2\pi i\alpha}D_r.
 \tag{P3K.67.8}
\]
The added factor is the actual global physical-circle action. It
changes neither (P3K.67.5) nor any finite affine vector. The exact norm is
\[
 \sum_{j=0}^{b-1}d(v+j/b)=b\tau(v)
                         =H\delta+\varepsilon\gamma.
 \tag{P3K.67.9}
\]
On the ordinary charts the same equality follows from (P3K.67.2), including
the scalar factor of (P3K.67.4). Explicitly, the uncorrected b-fold product
is
\[
 r^{-\varepsilon M}
 e^{-2\pi i\varepsilon M(b-1)/(2b)}
       g_0^bU^{\varepsilon b}
 =e^{-2\pi i\varepsilon M(b-1)/(2b)}h_0B_\lambda^{-\varepsilon}.
\]
The b copies of \(e^{2\pi i\alpha}\) cancel precisely the displayed
phase. The bounded period \(B_\lambda^{-\varepsilon}\) is retained in
the cusp representative; the ordinary quotient arrow is the specified
\(h_0\). Consequently the b-th iterate is the
specified one-gamma r-closing, and the quotient is the actual proper
family over \(w=r^b\). Its scalar winding is
\[
                         J=M/b.
 \tag{P3K.67.10}
\]

\item The bounded vector
\[
 \sigma(v)=\left(\varepsilon\Gamma v,
   v(U_0+\varepsilon C)-\frac{\varepsilon M}{2}(v^2-v)\right)
 \tag{P3K.67.11}
\]
satisfies \(\sigma(v+1/b)-\sigma(v)=d(v)\) exactly, including
the constant in (P3K.67.8). Since (P3K.67.4) is \textbf{forward}, the closed native
coordinates are \(\xi=z-\sigma(v)\). This arrow differs from the
endpoint-to-initial OS49 closing, whose closed coordinates use a plus
sign. All formulas here use the forward convention (P3K.67.4).

\item For a cut shorter parameter interval beginning at \(v\), set
\[
 C_s(v)=\exp\bigl(\sigma(v)-\sigma(v+s)\bigr),
      \qquad 0\le s\le1/b.
 \tag{P3K.67.12}
\]
The exponential acts on the two torus characters. It is a holomorphic
nonzero unit on the full cusp disc, smooth in the path variable. Its
inverse has the negative exponent. It commutes with the actual torus
periods and with physical delta. Its terminal value is precisely
\(\widetilde D_v^{-1}\). The functional identity in (P3K.67.11) holds
on a whole parameter neighborhood, so the first and last collars
match with all derivatives, rather than only at their endpoints.
Since \(b\mid M\), the endpoint rotation leaves \(\lambda\) and
all actual cusp period multipliers unchanged. The finite-cap action
still changes by its explicitly retained character in (P3K.67.5).

\item Retain the b successive factors
\[
 \begin{aligned}
 P_j(v)&=\exp\bigl(\sigma(v)-\sigma(v+j/b)\bigr),\\
 \log P_j/(2\pi i)
   &=-j\tau(v)-j\alpha\delta
           +\frac{\varepsilon Mj(j-1)}{2b^2}\delta.
 \end{aligned}
 \tag{P3K.67.13}
\]
The logarithm is the displayed vector on the fixed logarithmic cover.
For every sheet, including both collars,
\[
 C_s(v+j/b)P_j(v)=C_{j/b+s}(v),\qquad
 P_b(v)=\exp(-H\delta-\varepsilon\gamma).
 \tag{P3K.67.14}
\]
Thus the r-product is the \textbf{literal b-fold pullback} of the shorter
native product. Every cap chart is pulled back by (P3K.67.5), and the
full lattice cochain (P3K.67.6) accompanies it. No new covered boundary
product is chosen.

\par

\medskip\noindent\textbf{P3K.67.3. Each actual native star and its full unit maps}

\par\nobreak

\item For the fan \(T_n\), use the same literal shear as K64,
\[
 Z_n=Z,\quad V_n=Z^{-n}V,\quad h_n=hB^n,\quad B_n=B.
\]
Its full unit coefficients are
\[
 \kappa_n=\kappa g^n,\quad
 \alpha_n=\alpha_0\eta^n\kappa^{-n}g^{-n^2},\quad
 g_n=g,\quad\eta_n=\eta g^{-n},\qquad \eta=\lambda\beta.
 \tag{P3K.67.15}
\]
Here \(\alpha_0\) is the original period multiplier, not the constant
in (P3K.67.8). Its shorter-loop winding vector is exactly
\((0,nJ,0,J)\). The source factor \(U^\varepsilon\) in the endpoint
has not been removed: it is the first component of (P3K.67.7).

\item Apply the same shear to \(\sigma\):
\(\sigma_n=(\sigma_1,\sigma_2-n\sigma_1)\). If (P3K.67.12) has
character multipliers \((A,D)\) in these coordinates, its two vertex
maps are literally
\[
 (a,b_1,c)\mapsto((AD)^{-1}a,Ab_1,Dc),\qquad
 (u,v_1,w_1)\mapsto(ADu,D^{-1}v_1,A^{-1}w_1).
 \tag{P3K.67.16}
\]
The complete transitions, before suppressing any unit, are
\[
\begin{array}{c|ccc}
 &u&v_1&w_1\\ \hline
 C_0&\kappa_n\alpha_n b_1c&\alpha_n^{-1}ab_1&\kappa_n^{-1}b_1^{-1}\\
 C_1&g_n\eta_n b_1c&\eta_n^{-1}c^{-1}&g_n^{-1}ac\\
 C_2&a^{-1}&ab_1&ac.
\end{array}
 \tag{P3K.67.17}
\]
Substitution of (P3K.67.16) gives the same factors on every row. The
native C1 coordinates are
\((c,a,b_1)\) and \((\eta_nv_1,g_nw_1,(g_n\eta_n)^{-1}u)\);
the native C2 coordinates are \((a,b_1,c)\) and \((u,v_1,w_1)\).
Each has the literal transition
\(z'=z^{-1},\ n'_i=zn_i\). On C1 the factors are
\((D,(AD)^{-1},A)\) and \((D^{-1},A^{-1},AD)\); on C2 they are
\(((AD)^{-1},A,D)\) and \((AD,D^{-1},A^{-1})\).
These are exact automorphisms of the individual
\(\mathcal O(-1)\oplus\mathcal O(-1)\) models. They prove closed
native products on those stars. They do not supply a simultaneous
physical-circle-equivariant product of the three-curve plumbing.
The physical height is \(q=ab_1c\). Along any open curve axis it
vanishes to second order in the two normal coordinates, and at a
vertex it has order three. Thus the normal derivatives of the full
source factors \(A(q,v),D(q,v)\) are their values at \(q=0\), constant
along that curve; their remaining terms are precisely higher normal
jets carried by the normal-rescaling isotopy below. This retains the
q-dependence instead of treating the source factor as a constant map.

\item The Cox lifts, in the native orders, are
\[
 \begin{aligned}
 C_1:&\ (a_0,a_1,c_0,c_1)
       \mapsto(a_0,Da_1,(AD)^{-1}c_0,Ac_1),\\
 C_2:&\ (a_0,a_1,c_0,c_1)
       \mapsto(a_0,(AD)^{-1}a_1,Ac_0,Dc_1).
 \end{aligned}
 \tag{P3K.67.18}
\]
They preserve q and extend to both resolutions. In particular the
extra bounded source factor of the shorter endpoint is present in
both Cox rulings. The free period action transports a representative
to its orbit with the same maps. Outside those neighborhoods the
source, line, and every finite-cap arrow remain (P3K.67.4)--(P3K.67.6).

\par

\medskip\noindent\textbf{P3K.67.4. The ordinary double product and the actual relative square}

\par\nobreak

\item Take the literal double pullback of the single native products just
constructed. Equations (P3K.67.13) and (P3K.67.14) apply to its intermediate
seams as well. In the second vertex frame \((u,\eta_nv_1,\kappa_nw_1)\),
the actual based normal matrices of (P3K.67.17) have double windings
\[
 K_0=\operatorname{diag}(e^{2iJ(1-n)\theta},e^{2inJ\theta}),\quad
 K_1=\operatorname{diag}(1,e^{2iJ\theta}),\quad
 K_2=\operatorname{diag}(e^{2iJ\theta},1).
 \tag{P3K.67.19}
\]
Before reference division the respective factors are
\((\eta_n/\alpha_n,\kappa_n\alpha_n)\),
\((\kappa_n/g_n,g_n\eta_n)\), and \((\eta_n,\kappa_n)\).
All zero-winding based units have their logarithms on the full q-disc.
The seams (P3K.67.16) preserve these rows. Their parameter-dependent vertex
actions commute with physical delta, so the vertex orbit maps have
no additional parameter dependence. Their stable coordinate-frame
changes extend over the vertex four-disc and the curve annulus.

\item Use the actual ordinary quaternion contractions of K64 with
\[
 t=J,\qquad(l_0,l_1,l_2)=(2J(1-n),0,2J),\quad
 Q\mapsto H_tQH_{l_i-t}.
 \tag{P3K.67.20}
\]
The common determinant loop in (P3K.67.19) makes the first vertex product
common. The remaining determinant-one loops are filled by these
literal H maps. Multiply by the exponentials of the retained bounded
unit logarithms, with radial cutoffs, and restore the full germs by
normal rescaling and its cutoff velocity. This gives the same actual
ordinary double product construction as K64 on the supplied seams.

\item The physical weights are \(w=(0,1,-1)\) and \(d=(-2,-1,-1)\).
The meridian and stable-derivative trace degrees of the displayed map
are therefore
\[
 I=(-2J+4nJ,2J,-2J),\qquad W=(2J-4nJ,0,0).
 \tag{P3K.67.21}
\]
The exact vertex paths and rank-one endpoint caps in K64 close the
traces. They remain independent of the parameter under (P3K.67.16),
including its bounded source factor. The two ordinary relative pairs
are \((-2nJ,2J(n-1))\) and \((2J(n-1),-2nJ)\).
The native comparison is
\[
 q_H\mapsto(UH_JU^{-1})q_HH_{l_i-J}.
 \tag{P3K.67.22}
\]
It has left degree 2J and right degree zero, with the inverse cylinder
endpoint \(L^{-2J}\). Hence the actual native doubled pairs are
\[
 C_1:(2J(1-n),2J(n-1)),\qquad C_2:(2nJ,-2nJ).
 \tag{P3K.67.23}
\]
Choose a connection on the actual compact shorter family equal to its
native product velocity on a collar; extend it by a cutoff and the
affine space of connections. Its double pullback has exactly the
boundary products in (P3K.67.14). Interpolation relative to that boundary
compares it to (P3K.67.20). Thus (P3K.67.23) belongs to the actual relative
square. Multiplication by two on the free relative coordinates is
injective, giving the single pairs
\[
 C_1:(J(1-n),J(n-1)),\qquad C_2:(nJ,-nJ).
 \tag{P3K.67.24}
\]
Equivalently its literal b-fold pullback has the K64 pairs with
\(bJ=M\); (P3K.67.14) identifies the products in this check. There is no
division in the seven-sphere group in either calculation.

\item For the flop use the unrestricted four-weight common-link map with
the actual Cox lifts (P3K.67.18). Transport the metric by the full units
and use positive radial normalization. The resulting common-link
map is \(T(y,q_H)=(q_H^{-1}yq_H,q_H)\), followed by the same fixed
complex-normal reset as K62/K64. Its action on the relative group
adds the right native coordinate to the central class. C2 replaces
\((1,n-1)\) by \((1,n+1)\), and C1 replaces \((1,n)\) by
\((1,n-2)\). Therefore the actual shorter returns satisfy
\[
 A_{n+1}-A_n=-nJ,\qquad A_{n-1}-A_n=(n-1)J,
 \quad A_n=A_0-Jn(n-1)/2\pmod {28}.
 \tag{P3K.67.25}
\]

\item The shorter endpoint commutes with physical delta and preserves the
complex orientation. Its b-fold parameter cover is the prescribed
r-family by (P3K.67.9), so the actual Euler line and tangent bundle pull
back. Writing I for their integer characteristic index gives
\[
                             bI_w=I_r.
 \tag{P3K.67.26}
\]
This is naturality on an explicit oriented cover. It uses no
simultaneous normal framing. The signed identification of a nonzero
page expression with the physical index is not needed in (P3K.67.25):
A0 there is the actual reference sphere class. No numerical nonzero
reference value or separate k or ell is assigned by this note.

\par

\medskip\noindent\textbf{P3K.67.5. The original K33 endpoint with its retained constant}

\par\nobreak

\item For the original construction use root orders \(3K,4L\),
\(\gcd(3K,4L)=1\), and \(4Ln_0+3Kn_1=\varepsilon\).
Its actual parameter is
\[
 w_h=e^{-\pi ic/(KL)},\quad u=\kappa w_h^{2KL},\quad
 w_h=r^6,\quad M=12KL.
 \tag{P3K.67.27}
\]
The exact integral cochain and finite norms are K33.4--7:
\[
 p=(0,-\varepsilon,\varepsilon,(4Ln_0+2\varepsilon)/6)^t,
 \quad q=(0,0,\varepsilon,(3Kn_1-3\varepsilon)/6)^t.
\]
The full ordinary, root, annular, and cusp maps obey
\(x_{\rm cov}-6x=d^0(\varepsilon e_1,\varepsilon e_1,
\varepsilon e_1,0)\). Their actual one-step forward cusp vector is
\(\tau=(H\delta+\varepsilon\gamma)/6\), with no scalar correction
silently added. The canonical sigma above has
\(\Delta_{1/6}\sigma=\tau+5\varepsilon KL\delta/6\).
Retain instead the specified original endpoint product
\[
 \sigma_{\rm desc}(v)=\sigma(v)+\varepsilon KLv\delta,
 \qquad \Delta_{1/6}\sigma_{\rm desc}
                    =\tau+\varepsilon KL\delta.
 \tag{P3K.67.28}
\]
The last term is an integral physical period, so its exponential is
one in every actual toric character and on every finite cap. Thus
\(z-\sigma_{\rm desc}(v)\) is a closed native-star coordinate for
the original forward endpoint. Its j-fold product has additive
discrepancy \(j\varepsilon KL\delta\), still an actual integral
period. Formula (P3K.67.14) holds with sigma\_desc and the same exact
original source/root cochain.

\item On the r-cover, subtracting the initial origin makes the difference
from the sigma product the literal global map
\[
 G_\theta:X_{r(\theta)}\longrightarrow X_{r(\theta)},
 \qquad G_\theta=\delta_{-\varepsilon KL\theta},
                         \quad0\le\theta\le1.
 \tag{P3K.67.29}
\]
It is the physical scalar multiplication on the line. It commutes
with every period and finite affine map, and its integral toric
character action extends over all the indicated charts. Since
\(G_0=G_1=1\), it is a based periodic global gauge, not merely a
map on one normal boundary. Transporting a connection by this map
preserves its actual based relative return under the transported
native boundary products. Its intermediate right quaternion phase
is retained by (P3K.67.18). This proves the covered product comparison
without identifying the original endpoint with another sixth root.

\item Alternatively (P3K.67.19)--(P3K.67.24) applies directly to the original
native seam (P3K.67.28), whose actual scalar winding is \(J=2KL\).
The exact base/root arrow \(T=(t-1)/t\) in K33.14--15 identifies
the reference page with the E6 \((a,b)=(1,6)\) row and retains both
norm signs. Its cover reference index is zero. Integer index
naturality gives \(6A_0=0\) in \(\mathbb Z\), hence \(A_0=0\).
The vertical-fan result is therefore
\[
             A_n=-KL\,n(n-1)\pmod {28}.
 \tag{P3K.67.30}
\]
For \(K=L=1\), T2 has class \(-2\), while T0 and T1 have class
zero. No arbitrary-fan formula, source-moduli return, or separate
characteristic pair is assigned by this note.

\item The accompanying exact certificate is
\path{code/part3/shorter_scalar_vertical_fans.py}. It checks the cyclic norm,
all intermediate product logarithms, the native and Cox maps, the
full-unit shear and winding trace, the original integral constant,
and both signed norm choices. These algebraic identities accompany
the actual cap and relative-product maps above.

\end{enumerate}
\end{proof}


\section{The physical transfer and the corrected scalar fan algorithm}



\begin{theorem}\label{night:P3K.68}
\StatementRecord{night:P3K.68}{proved}{S6\_06\_BLIND: GPT-6 Astra (OpenAI, ChatGPT Pro mode, effort Pro), recorded re-derivation in the checkpoint dated 18 November 2026; also checked in this writing by \CurrentChecker}{\cref{night:P3F.22}, \cref{night:P3F.56}, \cref{night:P3F.8}, \cref{night:P3F.9}, \cref{night:P3K.23}, \cref{night:P3K.42}, \cref{night:P3K.57}, \cref{night:P3K.6}, \cref{night:P3K.65}, \cref{night:P3K.7}}
Let \(M\) denote the actual complex line whose nonzero vectors
carry the physical positive delta circle, and let \(U=M^{-1}\) be
the outward-meridian-one Thom line. For the actual source/page/conifold
return constructions specified below, the two oriented Jin pairs satisfy
\[
 (k_M,\ell_M,A_M)=-(k_U,\ell_U,A_U),\qquad
 [\Phi]=A_M=-A_U\pmod {28}.
\]
The exact physical Hopf transfer is
\(\lambda_M(F)=-\Sigma_{F^{-1}}\); the dual transfer is
\(\lambda_U(F)=\Sigma_{F^{-1}}\). In the selected-cusp source
family of F56/K42 the corrected physical pair is \((-6,0)\),
the class is \(-1\), and the initial native relative tuple is
\((0,-6,-1)\). The previously calculated pair \((6,0)\) is the
\(U\) pair. The signed native Cox commutator and its flop increment
are unchanged. This statement supplies the exact line, cap, orientation
and inverse-clutch maps for these corrections.
\end{theorem}
\begin{proof}
\begin{enumerate}[label=\arabic*.]
\item \par

\medskip\noindent\textbf{PL.1. The two oriented quotient spheres}

\par\nopagebreak[4]

\item At the first Cox-chart origin write
\[
 q_H=c_0+j\overline{c_1}.
\]
In the original complex coordinate order \((c_0,c_1)\), this is the
orientation-preserving real map to \((1,i,j,k)\). The physical circle is
\(q_H\mapsto q_He^{-i\theta}\). At \(q_H=1\) its tangent is \(-i\).
For \(z=q_H i q_H^{-1}\), the two horizontal derivatives are
\(j\mapsto-2k\) and \(k\mapsto2j\). Thus, with the outward orientation
on the unit sphere in \(\operatorname{Im}\mathbb H\),
\[
 [S^3_H]=-[S^1_\delta][S^2_z].
\]
The right-complex coordinates are \((c_0,d)\), \(d=\overline{c_1}\).
Their real complex basis \((1,i,j,ji)=(1,i,j,-k)\) reverses quaternion
orientation. In the projective chart \(d/c_0=x+iy\), the Hopf derivative
is \((\partial_x,\partial_y)\mapsto(-2k,-2j)\). Consequently
\[
 [\mathbf P^1_{\mathrm{right}}]=-[S^2_z].
\]

\item Let \(\alpha\) denote the Maurer/contact form normalized to value one
on the positive right phase. Its pullback to the unit section
\((1+j(x+iy))/\sqrt{1+x^2+y^2}\) is
\((x\,dy-y\,dx)/(1+x^2+y^2)\). Hence its curvature has integral one
on the right-complex projective orientation, and minus one on \(S^2_z\).
The physical principal connection is \(A=-\alpha\). The nonzero-vector
line with the physical positive scalar representation has
\[
 c_1(\mathcal L_\delta)=-dA,\qquad
 \int_{S^2_z}dA=1,\qquad
 \int_{S^2_z}c_1(\mathcal L_\delta)=-1,
 \qquad
 \int_{\mathbf P^1_{\mathrm{right}}}c_1(\mathcal L_\delta)=1.
\]
Here the Chern sign follows directly from a local unitary connection
\(\nabla=d+2\pi iA\): \(c_1=(i/2\pi)F_\nabla=-dA\).

\item The radial factor must also be retained. Since
\(\Omega_{\mathbb C^2}=dr\wedge\Omega_{S^3_H}\), moving the circle
factor to the front gives
\[
 \Omega_{\mathbb C^2}=d\delta\wedge dr\wedge\Omega_{S^2_z}.
\]
Therefore the outward circle-first normal quotient meridian, for the
complex orientation on the fixed curve, is \(S^2_z\). The right-complex
projective sphere is its opposite. The names ``positive meridian'' and
``physical line'' cannot identify these two orientations without the
actual accompanying normal/base orientation arrow.

\par

\medskip\noindent\textbf{PL.2. The full Cartier-to-Cox coefficient arrow}

\par\nopagebreak[4]

\item F56 constructs its positive-weight spinor as
\[
 E=\mathcal O(D_+)\oplus\mathcal O(-D_-),
 \qquad (p,q)=(\overline{c_0},c_1).
\]
K42 uses \(x=s-\rho\), \(y=s+\rho\), \(Z=x/y\), with the actual
Cartier vector frames \((1/x,y)\). Its coefficient map to the
dual-tautological/tautological frames is \(\operatorname{diag}(1/y,y)\).
After positive Hermitian normalization, writing \(Y=y/|y|\), this is
\[
 (\overline{c_0},c_1)
       =\operatorname{diag}(\overline Y,Y)(a,b).
\]
It follows literally that
\[
 (c_0,\overline{c_1})
       =\operatorname{diag}(Y,\overline Y)(\overline a,\overline b).
\]
Left multiplication by the first normalized quaternion gives
\[
 q_H^{\mathrm{right\ coefficients}}
   =g(x,y)\overline{(a,b)},\qquad
 g(x,y)=\frac1{\sqrt{|x|^2+|y|^2}}
     \begin{pmatrix}y&-\overline x\\x&\overline y\end{pmatrix}.
\]
Thus the same full-unit matrix \(g\) appearing in K42 is retained,
but its coefficient arrow is anti-linear. The currently written K42
arrow \(gS(a,b)\), with \(S\) the linear summand swap, differs by the
constant anti-linear swap \(S\circ\text{conjugation}\).

\item For adjoint coordinates
\(\gamma=|a|^2-|b|^2\) and \(\ell=\overline a b\), the linear swap
has action \((\gamma,\ell)\mapsto(-\gamma,\overline\ell)\), of real
determinant \(+1\). The anti-linear swap has action
\((\gamma,\ell)\mapsto(-\gamma,\ell)\), of determinant \(-1\).
The latter reverses the projective meridian. This is also forced by
the circle action: \(g\overline{e^{i\theta}(a,b)}
=e^{-i\theta}g\overline{(a,b)}\), whereas the linear swap retains
the positive scalar. A dual/Hermitian interpretation must include
this arrow; the displayed Cartier frames are vector frames.

\par

\medskip\noindent\textbf{PL.2a. The literal F9 source-normal and physical-line map}

\par\nopagebreak[4]

\item P3F.9 removes the remaining ambiguity in the local height convention.
Its first small-resolution chart is \((x,v,z)\), with
\(u=xz\), \(y=vz\), physical weights \(0,-1,+1\), and increasing
quotient height \(b=(|z|^2-|v|^2)/2\). The fixed curve is
\(v=z=0\). Thus its ordered Cox normal is \((c_0,c_1)=(v,z)\), its
source complex normal is \(y=vz\), and its positive-weight spinor is
\((\overline v,z)\). The physical circle acts by a common positive
scalar on that spinor. The positive projective complex orientation of
this spinor equals the actual circle-first normal meridian orientation.
Indeed, in coordinates \((b,\operatorname{Re}y,\operatorname{Im}y)\)
the outward projective sphere has the opposite of the raw orientation,
and F9/F56's physical normal orientation is exactly that opposite raw
orientation. To check the sign at one point, take \(v=1,z=0\). Then
\(b<0\), the projective coordinate is \(z/\overline v\), and
\(y=|v|^2(z/\overline v)\) has the same complex orientation. The
outward sphere orientation in the raw ordered three-plane at its
negative \(b\)-axis is minus this complex orientation. This proves the
stated opposite sign. Equivalently this is the explicit quaternion computation
in PL.1.

\item The physical principal circle over this projective sphere is the unit
circle in the tautological line. Therefore its nonzero-vector line is
\(\mathcal O(-1)\), and its Chern class has integral \(-1\) on the
actual outward meridian. The meridian-one Thom line is its dual
\(\mathcal O(1)\). The duality arrow is explicit: for the invariant
Hermitian metric it sends a nonzero vector \(v\) to the functional
\(w\mapsto\langle v,w\rangle/\langle v,v\rangle\), with the inner
product linear in its second argument. It conjugates the physical
circle character and retains the base point. On a regular source page
this gives the literal line map to \(M^{-1}\), since the original
nonzero-vector line there is \(M\).

\item Thus the source-page comparison and the meridian-one comparison are
related by an actual Hermitian dual map. With the latter convention the
page restriction is \(-c_1(M)\), and the real principal curvature
represents \(+dA\). The assertion in EC.1 that the same convention has
both meridian value one and page restriction \(+c_1(M)\) omits that
dual arrow.

\par

\medskip\noindent\textbf{PL.3. The independent collar sign}

\par\nopagebreak[4]

\item In the boundary orientation
\(\Omega_B=d\alpha\wedge du\wedge dv\wedge d\gamma\wedge d\theta\),
retain the exact forms from EC.4:
\[
 \beta=-n\{du+(\alpha-6\theta)dv\},\quad
 F=d\alpha\wedge d\gamma+n\,d\alpha\wedge dv,
 \quad d\beta=-n\,d\alpha\wedge dv+6n\,d\theta\wedge dv.
\]
They obey \(F^2=(d\beta)^2=0\) and
\(\beta Fd\beta=6n^2\Omega_B\). Hence
\[
 3\int_0^1\beta(F+s\,d\beta)^2ds=18n^2\Omega_B.
\]
The actual outward collar has its cutoff decreasing from one to zero,
so its contribution in the declared source-first auxiliary orientation
is \(-18n^2\). This computation does not identify that orientation
with any unstated Hopf-meridian convention.

\par

\medskip\noindent\textbf{PL.4. A global anti-linear cap map}

\par\nopagebreak[4]

\item The missing arrow extends over the actual two cap products. This can
be seen without choosing a reflection separately in each chart. For
every \(SU(2)\) transition matrix write
\[
 g=\begin{pmatrix}a&-\overline b\\b&\overline a\end{pmatrix},
 \qquad J=\begin{pmatrix}0&-1\\1&0\end{pmatrix}.
\]
Then \(\overline g=JgJ^{-1}\). Complex conjugation is the canonical
anti-linear bundle map \(C:E\longrightarrow\overline E\), and \(J\)
is a complex-linear bundle isomorphism from \(E\) to \(\overline E\).
Their quotient is the globally defined quaternionic anti-linear map
\[
 Q=J^{-1}C:E\longrightarrow E,
 \qquad Q(a,b)=(\overline b,-\overline a).
\]
It commutes with every \(SU(2)\) transition, has \(Q^2=-1\), preserves
the real four-dimensional orientation, and satisfies
\(Q(e^{i\theta}v)=e^{-i\theta}Q(v)\).

\item For a unit vector \(v\), the vectors \(v,Qv\) are an orthonormal
basis. Their rank-one Hermitian projections sum to the identity.
Consequently the projectivization of \(Q\) acts by minus the identity
on the traceless Hermitian three-plane. This is the actual fibrewise
antipodal map of \(S(\operatorname{ad}E)\); its normal three-dimensional
determinant and its sphere degree are both \(-1\). It extends to the
normal three-disc by \(w\mapsto-w\).

\item In the inner cap \(D^4\), use the actual Cartier frame; in the outer
cap \(D^3\times S^1\), use the actual Cox frame. The same formula for
\(Q\) is used in both. Since it commutes with the transition on the
overlap, these extensions glue through \(W\), retaining the complete
matrix \(g(x,y)\) and its parameter dependence. The orientation-
preserving factor \(J\) gives the corresponding map to the conjugate
spinor. The literal two cap maps on their sphere boundaries are thus
identity on the cap base times this projective map, with the specified
clutch conjugation. No homotopy removing the degree of \(g\) is used.

\item The determinant trivialization from the Cartier ratio is retained.
Any extending scalar unit used in the original coefficient maps is
sent to its conjugate by \(C\); the inverse circle character is retained
at the same time. Its normal adjoint action is unchanged. Thus this
construction also carries the actual scalar boundary gauges, rather
than discarding a winding. The associated complex line is dualized by
the exact anti-equivariance of \(Q\).

\par

\medskip\noindent\textbf{PL.5. The oriented cap pair under that exact map}

\par\nopagebreak[4]

\item Let \(B\) be the closed normal sphere bundle obtained from the original
two cap products, and let \(B'\) be the conjugated cap model just
constructed, with its conjugate projective orientation. The above
map \(\Psi:B\longrightarrow B'\) is fibrewise degree \(-1\) over the
unchanged oriented four-dimensional base. Its associated-line map is
the Hermitian dual map, so
\[
 \deg\Psi=-1,\qquad \Psi^*x'=-x,
 \qquad \Psi^*p_1(TB')=p_1(TB).
\]
The last equality is naturality under the displayed diffeomorphism;
it does not depend on orienting the normal bundle. Direct integration
therefore gives
\[
 \int_{B'}(x')^3=\int_Bx^3,
 \qquad
 \int_{B'}p_1(TB')x'=\int_Bp_1(TB)x.
\]
The two minus signs in each equality are respectively the degree of
\(\Psi\) and the dual-line factor. Moreover \(\overline g=JgJ^{-1}\)
shows that the relative second Chern number of the spinor remains
\(-1\) and its adjoint Pontryagin number remains \(4\), with the
base orientation of SP.15 unchanged.

\item These are direct statements about the exact conjugated cap models.
Their application to the physical mapping-torus class uses
the actual normal/meridian and associated-circle dictionary in PL.1:
the underlying nonzero-vector line of the physical circle has value
\(-1\) on the outward circle-first meridian for the complex fixed-base
orientation. This sign cannot be supplied by a declaration that the
same line has meridian value one. The full Jin and physical-clutch comparison is proved in PL.6 below.

\item For the meridian-one Thom line specifically, PL.2a supplies the missing
dual map on both caps. A retained scalar boundary gauge is inverted
pointwise; its extending gauge on the outer \(D^3\times S^1\) is
inverted by the same formula, and the inner \(S^3\) has no additional
gauge class. The normal clutch remains the conjugate of the same
\(g(x,y)\), so its relative Pontryagin number remains \(4\). On the
normal sphere bundle with that line and its physical outward meridian,
\(\pi_*x=1\), \(4x^2=\pi^*p_1(\nu)\), and the resulting pair is
\((1,4)\). On the inner surgery cap the same meridian-one line is
exactly SP.5's oriented normal Euler class. Its previously constructed
Seifert four-cycle therefore retains pair \((5,-4)\). These two
computations give \((6,0)\) after the stated physical-to-Jin parameter
reversal for this meridian-one line. This conclusion uses the explicit
duality and normal arrows, not the prior value of that pair. The physical Hopf-lift convention is then identified in PL.6 below.

\par

\medskip\noindent\textbf{PL.6. The physical product factors and the literal inverse clutch}

\par\nopagebreak[4]

\item The exact exterior of the fixed curve is
\(E=D^3_y\times S^3_{q_H}\), with its complex orientation. Here the
orientation of the \(h\)-sphere bounding \(D^3_y\) is opposite the
complex orientation of the fixed curve. Therefore the actual complex
six-sphere is
\[
 X=-\partial(D^3_y\times D^4_{q_H}).
\]
This is an identity of the two oriented product pieces: on the exterior
the boundary sign of \(D^3_y\times D^4_{q_H}\) is negative, while on
the fixed tube the displayed outer minus agrees with the complex
fixed-curve orientation.

\item The physical Hopf circle and quotient \(S^2_z\) give the circle-first
orientation \(S^3_u=-S^3_{q_H}\). Consequently the same actual
six-sphere is \(\partial(D^4_u\times D^3_y)\), including both pieces.
The associated physical nonzero-vector line is \(M\), with Chern
number \(-1\) on this \(S^2_z\). To use F22's convention that its
generator is positive, orient the factor \(S^2_z\) oppositely and
orient the doubled \(S^3_y\) oppositely as well. The orientation of
the six-dimensional quotient product is unchanged. Its circle-first
Hopf \(S^3_u\) also reverses. Thus both four-ball factors in
\[
 (D^4_u\times S^3_y)\cup_{\widetilde f}
                      (S^3_u\times D^4_y)
\]
reverse orientation, and their eight-dimensional product orientation
is unchanged. The underlying maps and circle action are the same;
this is the full factor-orientation relabelling, rather than a choice
of a sign from a previously calculated invariant.

\item Under this physical-line dictionary, the class entering the characteristic
integrals is \(x_{\rm phys}=c_1(M)=-x_{\rm Thom}\), where
\(x_{\rm Thom}\) is SP.5's outward-meridian-one class. The six-dimensional
mapping-torus orientation is unchanged. Thus the actual oriented odd
pair changes by minus:
\[
 (k,\ell)_{\rm phys}=-(k,\ell)_{\rm Thom}.
\]
This is separate from PL.5's \emph{conjugated cap-model} comparison, which
reverses both the normal projective orientation and the line. The
physical quotient orientation in the present paragraph is fixed by
the original complex threefold and the physical circle.

\item Here is the complete first-to-second clutch comparison, including the
piece exchange. In the plumbing coordinates \((y,q_H)\), put
\[
 L=D^4_y\times S^3_{q_H},\qquad
 R=S^3_y\times D^4_{q_H},
\]
and define \(\Sigma_F=L\cup_F R\) by identifying a point \(x\) in
the first boundary with \(F(x)\) in the second. For the physical line
\(M\), the positive-generator convention of F22 uses
\(u=q_H\) and the oppositely oriented \(y\)-ball.
More explicitly let \(\rho(y_0,y_1,y_2,y_3)=(-y_0,y_1,y_2,y_3)\)
and use \(y'=\rho(y)\) on that ball and its boundary.
The determinant of \(D\rho\) is \(-1\), so this is a literal
orientation-reversing extension across the ball. In this paragraph
\(-L\) and \(-R\) denote these opposite product orientations;
no antipodal map of \(\mathbb R^4\) is being used. Thus its first
and second pieces are
\[
 H_M=D^4_u\times S^3_{y'}=-R,\qquad
 K_M=S^3_u\times D^4_{y'}=-L.
\]
Its prescribed first-to-second clutch sends \(r\in R\) to
\(F(r)\in L\). Rewriting with \(L\) first means setting
\(x=F(r)\), so the second point is \(r=F^{-1}(x)\). Therefore
\[
                    \lambda_M(F)=-\Sigma_{F^{-1}}.
                                                        \tag{PL.1}
\]
The ball removed in the actual F8 cut-and-reglue construction is
\(D^4_u\times D^3_{y'}=-D^4_{q_H}\times D^3_y\); its outward
boundary is exactly the original complex six-sphere. On its free
face the boundary map is the actual relative physical return, and
on its fixed-tube face it is the identity. Both quotients identify
the first boundary point with its image in the second. They agree
pointwise off that free face and on its full collar. Hence
\(\lambda_M\) is the actual oriented clutched sphere.

\item For the dual line \(U=M^{-1}\), its complex positive circle is the
inverse physical circle. Its positive base is \(S^2_z\), so
\(u_U=+q_H\), \(y_U=y\). Now \(H_U=+R\) and \(K_U=+L\), with
the same first-to-second arrow. Consequently
\[
                  \lambda_U(F)=\Sigma_{F^{-1}}
                              =-\lambda_M(F).           \tag{PL.2}
\]
The comparison of the lifts is literal: the Hermitian dual map
\(D\) gives \(\widetilde f_U=D\widetilde f_M D^{-1}\).
It inverts each retained circle gauge. Its real four-dimensional
extension is the anti-linear spinor map of PL.4; the other four-ball
factor changes its orientation. Thus the maps on both seven-dimensional
pieces glue across the displayed clutch and reverse the eight-dimensional
product orientation. This is the geometric arrow in (PL.2).

\item For a disc-supported relative map, cancelling the opposite product
collars gives \(\Sigma_{F^{-1}}=-\Sigma_F\). In particular the
signed commutator generator of K23 still has
\(\lambda_M(\mathcal C)=\Sigma_{\mathcal C}=+1\). Its oriented
plumbing calculation and the native Cox flop increment are unchanged.
There is no extra inverse left to insert after (PL.1).

\item It follows directly that, in the same oriented Jin six-manifold,
\[
 (k_M,\ell_M,A_M)=-(k_U,\ell_U,A_U),\qquad
 [\Phi]=A_M\pmod {28}=-A_U\pmod {28}.                 \tag{PL.3}
\]
For the selected-cusp source loop the two cap computations above give
\[
 (k_U,\ell_U)=(6,0),\qquad(k_M,\ell_M)=(-6,0),
 \quad A_M=-1.
\]
The exact relative free-coordinate formula in this physical convention
is \(m=k_M-6A_M\), \(n=6A_M\); thus the initial native tuple is
\((m,n,s)=(0,-6,-1)\). The values follow from the line, orientation
and clutch maps, not from selecting a lift of a known residue.

\item The same comparison applies to the K6/K7 reference page calculation:
its outward-meridian-one line is \(U=M^{-1}\), and its determinant
restricts to \(-2c_1(M)\). Its square is unchanged. The proved page
expression \((3b^2-p_1(T\overline V))/24\) is therefore \(A_U\),
and the actual physical sphere class is its negative by (PL.3).
The inverse circle character and the exact physical Hopf lift, rather
than an equality of numerical polynomials, supply this conclusion.

\item For K57 specifically, IE.4 uses a real connection with
\(\alpha(\delta)=1\) and takes its curvature as its class \(x\).
It is exactly this \(U\) class. Its finite-chart Stokes calculation,
all normal residues and the scalar-order bordism remain proved.
Thus its good infinity value is still zero, while its reduced
\(A_U=1\) gives physical value \(-1\). This is restricted to the
two actual fans and completed families proved there.

\par

\medskip\noindent\textbf{PL.7. Exact correction scope and preserved records}

\par\nopagebreak[4]

\item The source, Cartier divisor, actual Cox and cap charts, relative
matrix \(g\), its degree, \(c_2=-1\), and normal \(p_1=4\) are
retained. The written linear SP.9/10 arrow is replaced by PL.2/4's
actual anti-linear arrow. K42's pair \((6,0)\) is retained as the
meridian-one \(U\) pair, while its physical pair is \((-6,0)\).
F56's actual selected-cusp class is \(-1\), with its local link and
zero-index normal-sphere correction unchanged. Its nontriviality
and degree-two non-descent still follow because \(-1\) is odd.

\item F22's abstract specified complex-line transfer remains its proved
construction. Equations (PL.1)--(PL.3) give its actual physical
instantiation. In K6/K7 a meridian-one line cannot also be the
un-dualized source page line; the explicit Hermitian dual map is
retained. Every nonzero reference index using that convention must
be converted by (PL.3). An independent native Cox correction is
not negated by that conversion: its signed central generator is
unchanged, as shown above. In particular K65's two-fan correction
uses an independently zero reference and the same native generator,
so its proved classes are unchanged.

\item The pre-correction K42/F56 proof, certificates, outputs, PDF and QA
are preserved verbatim in \path{records/k42_pre_antilinear_arrow/}, with
SHA256 manifest. The exact local orientation, clutch and collar
certificate is \path{code/part3_fourfold/euler_cube_peer_sign.py} and its
saved JSON/output files. Root and checks independently audited the
physical product-factor and inverse-clutch comparison before this
corrigendum was promoted.

\end{enumerate}
\end{proof}




\begin{theorem}\label{night:P3K.69}
\StatementRecord{night:P3K.69}{proved}{S6\_06\_BLIND: GPT-6 Astra (OpenAI, ChatGPT Pro mode, effort Pro), recorded re-derivation in the checkpoint dated 18 November 2026; also checked in this writing by \CurrentChecker}{\cref{night:P3F.22}, \cref{night:P3F.25}, \cref{night:P3F.56}, \cref{night:P3F.9}, \cref{night:P3K.10}, \cref{night:P3K.23}, \cref{night:P3K.24}, \cref{night:P3K.27}, \cref{night:P3K.30}, \cref{night:P3K.42}, \cref{night:P3K.45}, \cref{night:P3K.50}, \cref{night:P3K.51}, \cref{night:P3K.52}, \cref{night:P3K.56}, \cref{night:P3K.57}, \cref{night:P3K.58}, \cref{night:P3K.59}, \cref{night:P3K.6}, \cref{night:P3K.60}, \cref{night:P3K.64}, \cref{night:P3K.65}, \cref{night:P3K.68}, \cref{night:P3K.7}}
Use the physical line \(M\), its meridian-one Hermitian dual
\(U=M^{-1}\), and the complete circle, cap and sphere-gluing arrows
of P3K.68. Write
\[
 A_M=(4k_M-\ell_M)/24,\qquad A_U=(4k_U-\ell_U)/24.
\]
The actual oriented sphere return is \(A_M=-A_U\) modulo 28.
The following are corrected consequences for the already constructed
families and their specified positive parameter loops.

\begin{enumerate}
\item For the canonical whole-zero reference pages of K26, K28 and
K32, with \(H=Ma^2/d\) and page value \(\chi\), the actual physical
integer is \(H/2-\chi\). The retained characteristic and native free
coordinates are
\[
 (k_U,\ell_U)=(-3H,-24\chi),\quad
 (k_M,\ell_M)=(3H,24\chi),\quad
 (m,n)=(6\chi,3H-6\chi).
\]
The primitive K22 groups with \(d=8,14,20,30\) have physical integers
\(3,6,9,14\). The explicit primitive OS45 family of K10 has class
\(3\); its actual \(T_2\) fan of K64 has class \(2\).
The weighted degree-four OS47 example of F25 has class \(4\).
The reduced III families of K24 have physical integer \(3l+2\).
\item An actual fixed-edge flop in K23 has physical class
\(7A_M\). Its four primitive reference examples therefore have
flopped classes \(21,14,7,14\), respectively. The weighted OS47
example has fixed-edge-flopped class zero.
\item The selected-cusp source loop of F56 has initial and reflected
physical class \(-1=27\). Its initial native fixed-star coordinates
are \((0,-6,-1)\), and its ordinary coordinates there are
\((6,-6,-1)\). The first flops at \(C_0,C_1,C_2\) have classes
\(21,21,27\). On every actual physical-vertical-edge fan
\(T_n=T(e_1+ne_2,e_2)\), its physical class is
\[
                  s_n=-1-3n(n-1)\pmod{28},
\]
with exact image \(\{5,7,9,13,19,21,23,27\}\).
\item On the actual two mutation families of K59 the physical classes
are
\[
 [\Phi_{P_j}]=-1-3j(j+1),\qquad
 [\Phi_{R_j}]=-7-24j(j+1)\pmod{28}\quad(j\geq0).
\]
All its variation matrices, stable derivatives and native free pairs
are retained. The second sequence has least positive period seven
and exact image \(\{1,13,17,21\}\).
\item In the two initial/reflected E7 source objects of K57--60, put
\(A_{\rm scalar}=0\) for \((a,b,m)=(1,2,4p)\), and
\(A_{\rm scalar}=3p-1\) for \((a,b,m)=(2,1,4p-1)\).
The fixed-scalar infinity physical returns are respectively \(0,-1\).
For the literal global scalar \(\lambda=v^{km}\lambda_0\),
\[
\begin{array}{c|cc}
 (a,b,m)&I_T&Z_T+Z_{JT}\\\hline
 (1,2,4p)&0&1\\
 (2,1,4p-1)&-1-2k(3p-1)&2+2k(3p-1).
\end{array}
\]
Together with the actual single-loop difference \(Z_{JT}-Z_T=a\),
this leaves exactly the two possibilities
\[
 (Z_T,Z_{JT})=(kA_{\rm scalar}+14e,
                 a+kA_{\rm scalar}+14e),\qquad e\in\{0,1\}.
\]
No value of \(e\) is assigned. The smallest good case has pairs
\((0,1)\) or \((14,15)\); the reduced case \(p=k=1\) has pairs
\((2,4)\) or \((16,18)\).
\end{enumerate}
The reference zeros of the original family and the positive OS49
fan remain zero. The actual OS49 negative-fan classes \(-3k\) and
\(-k\) in K65 are unchanged, since their directed flop comparison
starts from that zero reference.
\end{theorem}
\begin{proof}
\begin{enumerate}[label=\arabic*.]
\item \par\medskip\noindent\textbf{The exact line map and the scope of the correction.}\par\nobreak
K68 constructs the Hermitian functional from each nonzero physical
vector, on both the ordinary source and each normal cap. It reverses
the circle character and gives \(c_1(U)=-c_1(M)\). Its local F9
coordinate calculation gives meridian value one for \(U\).
On the regular page the restriction of the K6 Thom class is therefore
\(-c_1(M)\). In K7 the graph-normal Euler class restricts as
\(-2c_1(M)\), so its square is exactly the previously computed
\(4c_1(M)^2\). No source, divisor, graph or cap is replaced.
The original real trace computes \(A_U\).

\item The physical-to-dual comparison in K68 retains the oriented quotient
and changes its two-sphere marking by the actual antipodal map. On
both Hopf-gluing pieces it is the identity on the quaternionic factor
and the reflection on the other four-ball. The map has degree minus
one in dimension eight. In physical positive-line coordinates the
pieces are \(-R,-L\), whereas in dual positive-line coordinates they
are \(R,L\), where
\(L=D^4_y\times S^3_q\) and \(R=S^3_y\times D^4_q\).
The first-to-second identification becomes \(F^{-1}\) when these
pieces are put back in \(L,R\) order. Consequently the exact arrows
are
\[
 \lambda_M(F)=-\Sigma_{F^{-1}}=\Sigma_F,\qquad
 \lambda_U(F)=\Sigma_{F^{-1}}=-\Sigma_F.
\]
The inverse equality follows by the explicit opposite-collar
cancellation after removing the extending \(L,R\) generators.
This accounts for the orientation and the inverse separately.
Applying F22 to its specified positive Chern line gives
\(\mu(\Sigma_F)=A_M/28=-A_U/28\).

\par\medskip\noindent\textbf{Reference pages and a fixed-edge comparison.}\par\nobreak
The retained page proofs give
\(k_U=-3H\), \(\ell_U=-24\chi\). Duality changes each odd
characteristic pairing by minus one, yielding the two pairs in the
statement and \(A_M=H/2-\chi\). The actual K27/K30 Cox framing
comparison is unchanged: its relative integers were computed from
these same \(U\)-pairings by
\[
 k_U=-m-n,\qquad \ell_U=-4m.
\]
Equivalently \(k_M=m+n\), \(\ell_M=4m\). Solving the two displayed
integral equalities gives \(m=6\chi\), \(n=3H-6\chi=6A_M\).
This proves the claimed characteristic and relative dictionary without
altering a native product.

\item For a primitive rational source \(H=d\) and \(\chi=1\), giving
\(d/2-1\). K10 supplies the complete OS45 equation, section,
linearization zero and ordinary cap, and K27.2 fixes its primitive
scalar winding as one. K64's actual directed \(C_2\) comparison at
\(T_1\) therefore subtracts one from its reference value three;
its \(T_2\) value is two. For F25 the retained actual degree-four
page has \(H=14\), \(\chi=3\), giving four. For K24 use its literal previously proved dual index \(A_U=-(3l+2)\), obtained in
its full page computation. The exact dual map just proved gives
\(A_M=3l+2\). No root Euler number is inferred from a residue.

\item For a fixed curve, K23's full common-link map remains
\(T(y,q)=(q^{-1}yq,q)\). Its proved relative conjugation is
\[
 T(L^mR^nC^s)T^{-1}=L^{-n}R^{m+2n}C^{s+n}.
\]
Its central generator has \(\mu(C)=1/28\) by the oriented
three-ball plumbing. The physical initial central class is
\(s=A_M\), and the exact Gysin/derivative dictionary above gives
\(n=6A_M\). Thus the actual flopped class is \(7A_M\).
Substituting \(A_M=3,6,9,14,4\) gives the stated five examples.

\par\medskip\noindent\textbf{The selected source loop and its vertical fans.}\par\nobreak
K42/K68 retain the actual anti-linear Cartier-to-Cox normal map and
both cap products. They give \((k_U,\ell_U)=(6,0)\), hence
\((k_M,\ell_M)=(-6,0)\), \(A_M=-1\), and native free pair
\((0,-6)\). K50's explicit inverse cylinder still adds \(L^6\),
so its ordinary pair is \((6,-6)\). This operation extends over the
normal four-disc and introduces no central factor. K51 and K52's
literal orbit and stable-derivative prisms are also unchanged; the
three native pairs are \((0,-6),(0,-6),(-6,0)\).
Adding their respective right integers to the physical baseline
\(-1\) gives \((-7,-7,-1)=(21,21,27)\) modulo 28.

\item The K45 half-source arrow carries the full source, line and finite-cap
maps and preserves the physical orientation. Its reflected initial
return therefore has the same class \(-1\). K56's native Cox arrows
carry every bounded unit and preserve the actual height. On \(T_n\)
they have right integers \(6n-6\) and \(-6n\), and their respective
target fans are \(T_{n-1}\) and \(T_{n+1}\). Consequently
\[
 s_{n+1}-s_n=-6n,\qquad s_0=-1.
\]
The telescoping is an equality in the stated group for every signed
integer \(n\), and gives \(s_n=-1-3n(n-1)\).
The actual lattice arrows and enumeration in K56 still identify the
whole stated vertical-edge domain. Its difference under \(n\mapsto
n+28\) is \(28(-6n-81)\). Direct evaluation of the 28 residues,
executed in the certificate below, gives the displayed image.

\par\medskip\noindent\textbf{The two mutation families.}\par\nobreak
K59's actual source/period/normal maps give the unchanged directed
update
\[
 V'_{i'i'}=V_{ii},\qquad
 \delta'_{i'}=2V_{ii}-\delta_i,\qquad
 s'=s+(\delta_i-V_{ii})/2.
\]
Its variation and derivative update equations contain no central
coordinate. Their initial matrices are the independently computed
K51/K52 matrices, unchanged above. Induction through exactly its
specified sequence therefore retains every matrix and free pair.
Only the initial central value changes from the historical one to
\(-1\). Every later central value changes by the same minus two,
because each directed increment uses those retained matrices.
This gives the two formulas in the statement. For the second,
\(-7-24j(j+1)\equiv21+4j(j+1)\pmod{28}\). Its first seven
values are \(21,1,17,13,17,1,21\). The sequence has period seven
and is not constant; since seven is prime, this is its least positive
period. The range follows from those seven values.

\par\medskip\noindent\textbf{The source-infinity comparison and paired source-zero returns.}\par\nobreak
K68 checks the actual connection in K57, IE.4--5: it is normalized
on the positive physical circle, and its curvature represents
\(c_1(U)\). Its complete local residue and collar calculation
therefore gives \(A_U=0,1\). The physical infinity indices are
\(A_M=0,-1\). This uses the same source blowups, finite-root maps,
line clutches and analytic scalar bounds; no new infinity filling
is invoked.

\item K58's actual four-holed parameter surface and its cut maps give
\(Z_T+Z_{JT}+S+I_T=0\), where \(S\) is the selected-cusp return.
The actual selected loop identified in its RG.4 is the same loop as
F56, so the corrected value is \(S=-1\). Its scalar winding at
infinity is exactly \(-2k\) in the independent root scalar.
The same pair-of-pants morphism consequently gives
\(I_T=I_{\rm fixed}-2kA_{\rm scalar}\).
The good reference scalar is zero; the reduced reference scalar is
\(3p-1\), by the reference calculation above. Substitution yields
the table and the common paired sum
\[
                    Z_T+Z_{JT}=a+2kA_{\rm scalar}.
\]
K60's independently constructed single native product and its literal
double pullback give \((m,n)=(1-a,a)\). Its directed Cox map thus
still gives \(Z_{JT}-Z_T=a\). Subtracting this equality from the
paired sum gives
\(2(Z_T-kA_{\rm scalar})=0\) in \(\mathbb Z/28\).
The kernel of multiplication by two is exactly \(\{0,14\}\), so
the displayed two ordered possibilities follow. Both satisfy both
integral congruences, and their remaining order-two choice has not
been supplied by any of these maps.

\item Finally, multiplication of either zero reference integer by minus
one gives zero. K65 starts its OS49 comparison from that independently
constructed positive-fan zero, and uses the actual negative-to-positive
native right exponent, which is unchanged by the line dictionary.
Its negative-fan classes \(-3k\) and \(-k\) are therefore retained.
The exact arithmetic certificate is
\path{code/part3_root/physical_return_consequences.py}.

\end{enumerate}
\end{proof}




\begin{theorem}\label{night:P3K.70}
\StatementRecord{night:P3K.70}{proved}{S6\_06\_BLIND: GPT-6 Astra (OpenAI, ChatGPT Pro mode, effort Pro), recorded re-derivation in the checkpoint dated 18 November 2026}{\cref{night:P3K.10}, \cref{night:P3K.22}, \cref{night:P3K.30}, \cref{night:P3K.31}, \cref{night:P3K.51}, \cref{night:P3K.52}, \cref{night:P3K.54}, \cref{night:P3K.64}, \cref{night:P3K.68}, \cref{night:P3K.69}}
Retain the exact OS45 rational elliptic source, section, line and cap
\[
\begin{gathered}
y^2=x^3-3f(t)x+2g(t),\qquad
f=(t^2+1)^2+t,\\
g=(t^2+1)^3+\tfrac32t(t^2+1)+\tfrac38,\qquad
G=(t^2+1,\sqrt3/2),\\
\mathcal M=\mathcal O(8G-O),\qquad Q=G.
\end{gathered}
\]
Choose one of the four marked simple roots of
\(48t^4+8t^3+96t^2+72t+57\) as its whole-fibre linearization zero,
and retain the ordinary order-one logarithmic cap at \(t=0\) from
K10. Use its primitive scalar parameter \(\lambda=r\), its actual
forward closing \(t/(t-\tau)\), and all the original source, line
and cap maps.

In the original Tate/source and line period lattice, replace the cusp
fan by the periodic triangulation with ordered basis
\[
 v=(2,-1),\qquad w=(-1,1),
\]
and triangles \((n,n+v,n+w)\),
\((n+v,n+w,n+v+w)\) for every \(n\in\mathbb Z^2\).
The resulting actual holomorphic sphere family has positive
single-loop physical return
\[
 [\Sigma_\Phi]=11\quad\hbox{in }\mathbb Z/28,
 \qquad \mu(\Sigma_\Phi)=11/28\pmod{\mathbb Z}.
\]
It is obtained by the two explicit Cox flops below, with initial,
intermediate and final physical classes \(3,21,11\). Its mapping
torus is not diffeomorphic to \(S^1\times S^6\), by the retained
spin/connected-sum argument of F24 and the physical transfer K68.
This statement assigns no universal scalar-fan formula or value at
an unevaluated fixed-star target.
\end{theorem}
\begin{proof}
\begin{enumerate}[label=\arabic*.]
\item \par

\medskip\noindent\textbf{NV.1. The actual family, parameter and two products}

\par\nobreak

\item Fix the constructed OS45 surface
\[
 y^2=x^3-3f(t)x+2g(t),\quad
 f=(t^2+1)^2+t,\quad
 g=(t^2+1)^3+\tfrac32t(t^2+1)+\tfrac38.
\]
Its marked section is \(G=(t^2+1,\sqrt3/2)\), its line is
\(M=\mathcal O(8G-O)\), and the suspension translates by \(-G\).
Retain any one of the four marked simple roots \(\tau\) of
\(48t^4+8t^3+96t^2+72t+57\) as the scalar-zero fibre and retain
K10's actual ordinary multiplicity-one cap at \(t=0\). The parameter
is the primitive \(\lambda=r\), so its scalar winding is one. The
root/cap/source construction, properness and sphere topology are the
ones already proved in K10/K22; the maps below change its cusp fan
and preserve that actual exterior. The forward ordinary closing is
\(t/(t-\tau)\), namely \(q^{-1}\) times the bounded unit \(t\) at
\(q=t-\tau=0\). No cap is added or removed.

\item We keep symbols \(J,H,\chi\) temporarily to record exactly the
source of each contribution. For this concrete family they are
\(J=1,H=8,\chi=1\). On the degree-two parameter cover the full cusp
periods are
\[
 h:(Z,V)\mapsto(\kappa qZ,\alpha V),\qquad
 B:(Z,V)\mapsto(gZ,q\eta V),\qquad \eta=\lambda\beta.
\tag{NV.1}
\]
All four bounded units other than \(\lambda\) are the original ones.
The forward closing has the form \(q^e u(q)\), with its specified
nonzero unit; in the concrete family \(e=-1,u=t\). Reindex by
\(B^{-e}\), retaining its source factor, and use K64's actual cut
map \(C_\tau=(AZ,DV)\). It satisfies \(C_{\rm end}E=1\) on
whole collars, including all derivatives. On the double cover we
use exactly \(C^{(2)}_{j+\tau}=C_\tau E^{-j}\), \(j=0,1\).
The exterior and cap squares are therefore literal pullbacks of the
single family.

\item The initial two vertex charts are \((a,b,c)\), \((u,v,w)\), with
\(q=abc=uvw\). The native fixed star \(C_0\) has charts
\[
 (z,n_0,n_1)=(b,a,c),\qquad
 (z',n'_0,n'_1)=(\kappa w,\alpha v,(\kappa\alpha)^{-1}u).
\tag{NV.2}
\]
Substitution of the full period row gives
\(z'=z^{-1}\), \(n'_i=zn_i\), with every unit retained.
Its Cox quaternion is
\[
 q_H=(1+jz)(n_0+j\overline{n_1}).
\tag{NV.3}
\]
This is literally K30.11 in the same order: \(n_0\) has physical
weight \(-1\), \(n_1\) weight \(+1\), and physical delta is right
multiplication by \(e^{-i\phi}\). Thus the native C0 normal and base
markings are the ones used for K69's retained relative pair.

\item There is one endpoint comparison to retain here. K30 writes its
extending origin as a reference-to-current map. Its inverse gives
closed coordinates. With linear parameter cutoff and \(g=U^{-1}\),
\(\eta=\lambda W^{-1}\), K64's current-to-reference cut map differs
from that inverse by the global physical-circle factor
\[
 \exp\bigl(\pi i eJ(\theta^2+\theta)\bigr)
\tag{NV.4}
\]
on a single \(J\)-winding interval. Both endpoints are one. This is
an actual based circle action on the whole family, commuting with all
periods and finite caps, so carrying the connection by it preserves
the relative return. Replacing the linear cutoff by the protected
cutoff changes the bounded source and line factors through their
explicit interval logarithms with zero endpoint values. Multiplying
these logarithms by a real homotopy parameter gives an actual based
homotopy of the star products. The q-dependent normal jets are carried
by the normal-rescaling flow of K64. No parameter winding is discarded.

\par

\medskip\noindent\textbf{NV.2. The fixed-star native-to-ordinary map}

\par\nobreak

\item For the doubled loop, put \(t=J\) and use exactly K64's common
ordinary plumbing product
\[
 T_i(Q)=H_JQH_{l_i-J},\qquad(l_0,l_1,l_2)=(2J,0,2J).
\tag{NV.5}
\]
At the fixed star, the normal pair in (NV.2) is \((a,c)\), so its
literal map to the common product is
\[
 q_H\longmapsto (U(z)H_JU(z)^{-1})q_HH_J,
 \qquad U(z)=\frac{1+jz}{\sqrt{1+|z|^2}}.
\tag{NV.6}
\]
This is obtained by substituting \(Q=n_0+j\overline{n_1}\) into
(NV.3), with no interchange or inversion of the normal pair.
At the opposite vertex the factors are precisely the full unit
factors \(\eta/\alpha\) and \(\kappa\alpha\) from K64. The
normal homotopy therefore matches both protected collars of (NV.2).

\item The left factor in (NV.6) has degree \(2J\) on the oriented native
Hopf base times parameter, by the same direct quaternion integral as
SV.11; the right factor depends only on radius and parameter and has
zero third degree. The inverse-cylinder extension has endpoint
\(L^{-2J}\). It follows that the ordinary relative pair is the
native pair minus \((2J,0)\). This is the identical arrow to the
already proved nonfixed C1/C2 comparison, now verified in the actual
fixed-star coordinates (NV.2)--(NV.3).

\item The retained K69 native pair on the single family is
\((6\chi,3H-6\chi)\). Its literal double pullback is
\((12\chi,6H-12\chi)\), hence the common-product C0 pair is
\[
 (m_0^{\rm ord},n_0^{\rm ord})
       =(12\chi-2J,6H-12\chi).
\tag{NV.7}
\]
The actual product-frame derivative map K31.7, retained independently
of its false initial-data claim, therefore gives
\[
 V_{00}=2J-6H,\qquad \delta_0=2J+6H-24\chi.
\tag{NV.8}
\]
The physical double return is \(2(H/2-\chi)\), by the actual
parameter pullback and K68, not by a division of a sphere class.

\par

\medskip\noindent\textbf{NV.3. Complete initial ordinary data from the actual prisms}

\par\nobreak

\item The full-unit trace maps of K64 have degrees
\[
 I=(-2J,2J,-2J),\qquad W=(2J,0,0).
\tag{NV.9}
\]
Use the literal generic endpoint paths and rank-one endpoint caps of
K51/K52. Their vertex maps are (SV.5), commute with the physical
circle and are parameter independent after transport. Thus the same
oriented chains give
\[
 \operatorname{Var}(D_1)=2J(m_0+m_1),\qquad
 \operatorname{Var}(D_2)=2J(-m_0+m_2),\qquad
 \delta_1=-2J,\quad\delta_2=2J.
\tag{NV.10}
\]
Here \(\Delta x_1=D_1\), \(\Delta x_2=-D_2\), and the boundary
identity is \(\operatorname{Var}(\Delta x)=[K\partial x]\).
Together with symmetry and (NV.8), the actual data are
\[
 V_0=\begin{pmatrix}
 2J-6H&2J&-2J\\2J&2J&0\\-2J&0&2J
 \end{pmatrix},\qquad
 \delta_0^{\rm vec}=
 \begin{pmatrix}2J+6H-24\chi\\-2J\\2J\end{pmatrix}.
\tag{NV.11}
\]
These replace the withdrawn scalar rank-one seed. Their derivation
uses the actual separate native products and the ordinary double
product; it makes no simultaneous equivariance assertion.

\par

\medskip\noindent\textbf{NV.4. Two literal Cox changes and target products}

\par\nobreak

\item The first flip is C0, with exact new basis and characters
\[
 P_1=\begin{pmatrix}1&-1\\0&1\end{pmatrix},\qquad
 (Z',V')=(ZV,V).
\tag{NV.12}
\]
The native other-ruling map is
\((z,n_0,n_1)\mapsto(n_0/n_1,zn_1,n_1)\); the rank-one matrix
\(\bigl(\begin{smallmatrix}n_0&n_1\\zn_0&zn_1\end{smallmatrix}\bigr)\)
gives both charts. It preserves \(q=zn_0n_1\) exactly. The whole
unit matrix transforms by the actual monomial map
\(A'_{rs}=\prod A_{hj}^{(P_1^{-1})_{rh}(P_1)_{js}}\).
The cut map C transforms by the same characters and has the full Cox
lift. Thus the outside source, finite cap and closing remain the
ones of NV.1.

\item Construct the new ordinary product from these full transformed units,
using the same based logarithms and H maps. On this fan its doubled
winding matrix, tangent/normal weights and first normal exponents are
\[
 B_1=\begin{pmatrix}0&2J\\0&2J\end{pmatrix},\quad
 w=(1,1,-2),\quad d=(-3,-3,0),\quad l=(2J,-2J,2J).
\tag{NV.13}
\]
The trace maps, including their full unit homotopies, have
\[
 I=(-2J,10J,-2J),\qquad W=(4J,-8J,-2J).
\tag{NV.14}
\]
The selected new star is C2. The unrestricted four-weight Cox
common-link and the constant complex reset are precisely K54/CF.5.
On each side the native-to-ordinary map adds the same pure-left
\(-2J\) to the native return. It cancels in the selected-coordinate
comparison. Consequently, if a,b are the old selected ordinary
\(V_{00},\delta_0\), the new selected entries are
\(V'_{22}=a\), \(\delta'_2=2a-b\), and the physical central
increment is the old native right integer. This is a direct comparison
of the two constructed products on the selected star; no common
exterior matrix-reset premise is used.

\item All target weights are nonzero. Its actual prism equations and
symmetry then determine the target uniquely from those two selected
entries: with \(r_i=1/w_i\),
\[
 V'=\operatorname{diag}(I_i/w_i)+c\,rr^t,\quad
 \delta'_i=(W_i+e_*)/w_i,\quad
 c=w_2^2a-w_2I_2,\quad e_*=w_2(2a-b)-W_2.
\tag{NV.15}
\]
This is linear algebra on actual integral cycle identities. Substitution
gives
\[
 V'=\begin{pmatrix}
 2J-24H&4J-24H&-2J+12H\\
 4J-24H&14J-24H&-2J+12H\\
 -2J+12H&-2J+12H&2J-6H
 \end{pmatrix},\quad
 \delta'=\begin{pmatrix}
 2J+36H-48\chi\\-10J+36H-48\chi\\2J-18H+24\chi
 \end{pmatrix}.
\tag{NV.16}
\]
The new C1 ordinary and native right integer is therefore
\[
 n_1^{(2)}=30H-24\chi-12J.
\tag{NV.17}
\]
Its native product is the explicit periodic holomorphic star product
with all transformed units. Its double product is chosen as the
literal pullback of that single product, exactly as in K64. The
ordinary-product comparison just made uses those boundary products.
Thus (NV.17) is the right coordinate of the actual relative square;
injectivity of multiplication by two in the free relative group gives
\[
 n_1=15H-12\chi-6J.
\tag{NV.18}
\]
No division occurs in \(\Theta_7\).

\item The second actual flip has basis
\[
 P_1Q_1=\begin{pmatrix}2&-1\\-1&1\end{pmatrix},\qquad
 Q_1=\begin{pmatrix}1&0\\-1&1\end{pmatrix}.
\tag{NV.19}
\]
Its characters are \((Z'',V'')=(Z',Z'V')\). The same full unit
formula and the two Cox charts give the actual family. The constant
normal interchange is carried by CF.6a's explicit U(2) path. The
unrestricted common-link map, followed by its fixed complex-normal
reset, adds (NV.18) to the physical central class.

\item For the specified OS45 data, the actual initial physical value is 3.
The first fixed-edge flip adds 18 and gives 21. Formula (NV.18) is 102,
so the second specified flip gives
\[
               3+18+102=123\equiv11\pmod{28}.
\tag{NV.20}
\]
This is a concrete nonvertical fan comparison. The root independently audited NV.1--4, including the exact C0 closing,
product comparison and target reconstruction. The complete closing
and repeated-fan certificate below has also passed. This assigns the
return of the stated concrete family and fan only.

\par

\medskip\noindent\textbf{NV.5. The complete closing check and exact finite consistency record}

\par\nobreak

\item For the concrete OS45 cap the ordinary forward map has line multiplier
\(q^{-1}t\). Its reindexing by the actual \(B\) gives
\[
 E:(Z,V)\longmapsto(gZ,t\eta V),\qquad
 C_{\rm end}:(Z,V)\longmapsto(g^{-1}Z,(t\eta)^{-1}V).
\tag{NV.21}
\]
Their composition is the identity on the whole cusp annulus. These
are full nonzero germs in q, so equality retains every derivative.
On the two successive laps the actual products are \(C_\tau\) and
\(C_\tau E^{-1}\). Their adjoining collar values are both
\(E^{-1}\), and their final value is \(E^{-2}\). The finite-cap
map is the original order-one map carried by the same ordinary
closing, hence this change of coordinates preserves its full overlap
square. On either mutated fan the character changes of NV.4 send
(NV.21) to their literal monomial images. All maps fix q, so these
images still give exactly the original source, line and cap family
on the common exterior.

\item For every vertex multiplier \((A,D)\), the three full rows satisfy
\[
 F_i((AD)^{-1}a,Ab,Dc)
      =(AD\,u_i,D^{-1}v_i,A^{-1}w_i),
\tag{NV.22}
\]
where \((u_i,v_i,w_i)\) are the respective period rows. Thus the
bounded source factor, line factor and higher q jets are carried
through every selected star. The Cox lift is precisely the one
obtained by this substitution in the rank-one matrix of NV.4.
It extends to both rulings before taking the free period quotient.
This verifies the actual closing/product square, rather than only
the final winding matrix.

\item The executed certificate is \path{code/part3/scalar_nonvertical.py}, with
its JSON and captured output. It passes 66 exact identities. These
include the OS45 equation, section, simple-cusp polynomial and regular
cap; the full unit rows and quaternion overlap; the based physical
delta comparison; all initial and target trace equations; the actual
native-square integers; and both other-ruling charts.

\item A separate finite consistency calculation in that same file uses the
concrete seed only. Through words of length four it obtains 21 basis
states and 16 literal periodic triangulations. All ten repeated-basis
or repeated-fan checks agree. For a repeated fan the program does
not identify it from a numerical invariant: it compares its actual
two triangle orbits, writes each new triangle as an old triangle plus
an integral vector, and exports the vertex permutation. The vector
\((r,s)\) is the actual period \(h^rB^s\), with its full units; the
induced map on the quotient is therefore the identity of the same
family. For example the reverse of the second flop has basis
\(\bigl(\begin{smallmatrix}2&1\\-1&0\end{smallmatrix}\bigr)\).
Its two triangles are the earlier \(P_1\) triangles translated by
\((1,-1)\) and \((2,-1)\), with exported permutations, and its
physical class returns to 21.

\item The two encountered targets whose zero-weight curve is not the
selected flopped curve are left unevaluated. No fixed-star left
integer is inferred there from a central class. The finite search
is an exact consistency certificate for the written maps, not a
claim to have computed every fan or every word.

\end{enumerate}
\end{proof}




\begin{theorem}\label{night:P3K.71}
\StatementRecord{night:P3K.71}{proved}{S6\_06\_BLIND: GPT-6 Astra (OpenAI, ChatGPT Pro mode, effort Pro), recorded re-derivation in the checkpoint dated 18 November 2026}{\cref{night:P3K.55}, \cref{night:P3K.66}, \cref{night:P3K.68}}
Retain the actual source family of P3F.54 and P3F.56,
\[
 y^2=x^3+3r(u)x^2+tx+1,\quad G=(0,1),\quad
 r(u)=\frac{1-2u^3}{3u^2},\quad t_0(u)=u^2-\frac2u,
 \qquad u+u^{-2}=\rho^2,
\]
where \(u(0)=u_*\), \(u_*^3=-1\), and the parameter is one positive
\(\rho\)-circle. Use either constructed sphere datum
\((M,Q,m)=(\mathcal O(G-O),2G,4h)\) or
\((\mathcal O(2G-O),G,4h-1)\), \(h\ge1\), with its prescribed
norm numerator \(\pm1\), finite character and complete annular
cochain. The scalar is fixed and sufficiently small. There is no
additional logarithmic disc at the coalescing finite sites.

For every \(P\in GL_2(\mathbb Z)\), let \(T(P)\) be the periodic
two-triangle fan with columns \(v,w\) and triangles
\[
 \{z,z+v,z+w\},\quad\{z+v,z+w,z+v+w\},\qquad z\in\mathbb Z^2.
\]
The following terminating exact algorithm computes the actual
oriented return \([\Phi_P]\in\Theta_7=\mathbb Z/28\), with
\(\mu(\Sigma_j)=j/28\).

First use the identical-fan maps in the proof to choose the unique
representative
\[
 T_n=\begin{pmatrix}1&0\\n&1\end{pmatrix},\quad n\in\mathbb Z,
 \qquad\hbox{or}\qquad
 C=\begin{pmatrix}p&-r\\q&s\end{pmatrix},\quad p,r>0,
 \quad ps+rq=1.
\]
The integer subtraction in the proof gives an exact factorization
\[
 C=T_nQ_{i_1}\cdots Q_{i_k},\qquad i_j\in\{0,1\},\qquad
 Q_0=\begin{pmatrix}1&-1\\0&1\end{pmatrix},\quad
 Q_1=\begin{pmatrix}1&0\\-1&1\end{pmatrix}.
\]
Start at \(T_n\) with
\[
 V=\begin{pmatrix}
 18n(n-1)&6-12n&12n-6\\
 6-12n&0&0\\12n-6&0&0
 \end{pmatrix},\quad
 \delta=\begin{pmatrix}-12-18n(n-1)\\12n-12\\-12n\end{pmatrix},
 \quad S=-1-3n(n-1).
 \tag{P3K.71.1}
\]
At a step \(P\mapsto P'=PQ_i\), put \(a=V_{ii}\),
\(b=\delta_i\). Compute the following data at the target \(P'\):
\[
 B=(b_{jk})=-6(P')^{-1}e_1e_1^tP',\qquad
 l=(b_{22}-b_{21},b_{11}-b_{12},b_{22}),
\]
\[
 (x,y)^t=(P')^{-1}e_2,\quad
 w=(x,y,-x-y),\quad d=(-x-2y,-2x-y,x-y),
\]
\[
 I_j=-3w_j+d_j(l_j+3),\qquad
 W_j=-3w_j-d_j(l_j+3),\qquad j=0,1,2.
 \tag{P3K.71.2}
\]
Every target weight \(w_j\) is nonzero. Set
\[
 c=w_2^2a-w_2I_2,\qquad e=w_2(2a-b)-W_2,
\]
\[
 V'_{jk}=\frac{I_j}{w_j}\delta_{jk}+\frac{c}{w_jw_k},\qquad
 \delta'_j=\frac{W_j+e}{w_j},\qquad
 S'=S+\frac{b-a}{2}.                                  \tag{P3K.71.3}
\]
All entries and the increment in (P3K.71.3) are integers. The
output is \([\Phi_P]=S\pmod{28}\). The matrices and vector are
the actual ordinary-product variation and stable derivative in
the canonical fan marking. Thus the algorithm computes each
specified family; its domain has no bound on the entries of \(P\).
For example, \(P=(\begin{smallmatrix}2&-1\\-1&1\end{smallmatrix})\)
has class one.
\end{theorem}
\begin{proof}
\begin{enumerate}[label=\arabic*.]
\item \par\medskip\noindent\textbf{AF.1. The exact identical-fan morphisms.}\par\nopagebreak[4]

\item Put
\[
 R=\begin{pmatrix}-1&-1\\1&0\end{pmatrix},\qquad R^3=1.
\]
The columns of \(PR\) are \(w-v,-v\). Its first triangle is
the first \(P\) triangle translated by \(-v\); its second is
the second \(P\) triangle translated by \(-2v\). The first
triangle of \(-P\) is the second \(P\) triangle translated
by \(-(v+w)\), and its second is the first \(P\) triangle
with the same translation. Swapping the columns just exchanges
the two named vertices. These are equalities of the full
infinite fans, with their primitive rays and heights.

\item Use the identity on the absolute dense-torus characters
\((Z,V,q)\). Each matched absolute cone has the identical
character semigroup, so this map extends to the identity on
its affine toric chart. The overlap maps are the same character
maps. This identity commutes with both full original periods,
including every bounded multiplier. A changed basis of their
generators has an integral inverse and generates exactly the
same group. The map therefore descends to a holomorphic family
isomorphism. On the source/cap overlap it is the original
identity and hence glues to the identity of that entire exterior.
It carries the prescribed closing arrow and the physical
\(\delta\) circle. In particular it preserves the complex
orientation and the physical circle character.

\item For determinant minus one, first swap the columns. The six
remaining candidates are
\[
 P,PR,PR^2,-P,-PR,-PR^2.                              \tag{P3K.71.4}
\]
The star of the lattice origin has rays
\(\pm v,\pm w,\pm(v-w)\), bounding its six actual unimodular
triangles. The physical vector is \(\delta=(0,1)\). When it
is a ray, primitivity makes that ray exactly \(\delta\).
The unique adjacent triangle with a positive-first-coordinate
edge and positive orientation has columns \((1,n),(0,1)\).
Its representative is \(T_n\). Otherwise \(\delta\) lies
in the interior of one of the six cones. Order that cone's
boundary vectors from positive to negative first coordinate.
They are \((p,q),(-r,s)\) with \(p,r>0\) and determinant one.
The coefficients of \(\delta\) are exactly \((r,p)\).
The candidates (P3K.71.4) are the six positively oriented
triangles just identified. This proves the asserted unique
choice and gives its actual family morphism.

\par\medskip\noindent\textbf{AF.2. The finite factorization and its actual Cox maps.}\par\nopagebreak[4]

\item Write \(Q_2=Q_0^{-1}\). Starting with the nonvertical representative
\(C=(\begin{smallmatrix}p&-r\\q&s\end{smallmatrix})\), use
\(C Q_2\) when \(r\ge p\), recording zero. Use \(C Q_1^{-1}\)
when \(p>r\), recording one. Before termination, the new positive
pair is respectively \((p,r-p)\) or \((p-r,r)\), and its sum
strictly decreases. The determinant gives \(\gcd(p,r)=1\).
Equality \(p=r\) is therefore \(p=r=1\), and its next matrix
is \((\begin{smallmatrix}1&0\\q&1\end{smallmatrix})\).
This proves finite termination at a specified \(T_n\).
Reversing the recorded list gives the displayed factorization
by literal matrix multiplication. Every subsequent forward
target has positive first entry and negative second entry;
its tangent weights \((r,p,-r-p)\) are all nonzero.

\item We use the factorization to construct the actual forward families.
For \(C_0\) the new basis is \((v,w-v)\), the dense characters
are \((Z',V')=(ZV,V)\), and its first triangle coordinates are
\((a/c,bc,c)\). For \(C_1\) these are \((v-w,w)\),
\((Z',V')=(Z,ZV)\), and \((a/b,b,bc)\). In the old native
chart \((z,n_0,n_1)\), the actual two other-ruling charts are
\[
 (n_0/n_1,zn_1,n_1),\qquad(n_1/n_0,zn_0,n_0).       \tag{P3K.71.5}
\]
For \(C_1\) the two normal coordinates are interchanged. The
interchange is carried by the explicit constant path
\[
 S_t=\frac12\begin{pmatrix}
 1+e^{\pi it}&1-e^{\pi it}\\
 1-e^{\pi it}&1+e^{\pi it}
 \end{pmatrix},\qquad 0\le t\le1.                    \tag{P3K.71.6}
\]
Its constant eigenspaces have eigenvalues \(1,e^{\pi it}\),
so it is unitary and preserves the real four-plane orientation.
It commutes with the scalar native transition \(z1\).
Both products in (P3K.71.5) are exactly \(q=zn_0n_1\).
Their common rank-one matrix is
\( (\begin{smallmatrix}n_0&n_1\\zn_0&zn_1\end{smallmatrix})\).
Thus these are the two actual rulings of the same local conifold.

\item Let \(A=(A_{hj})\) be the complete old period-multiplier matrix.
For either of these basis changes the new one is exactly
\[
 A'_{ik}=\prod_{h,j}A_{hj}^{(Q^{-1})_{ih}Q_{jk}}.
                                                            \tag{P3K.71.7}
\]
This follows by substitution in the characters and the period
generators, keeping \(q\) fixed. It retains every unit as a
function of \(q,\rho\). The new second native charts give
\(z'=z^{-1}\), \(n'=zn\) with these same factors, as in the
explicit GF.11 rows. The common complement uses the original
source and line coordinates via the inverse character map.
Consequently its cap, finite character and annular closing
are the original maps. The common-link normalization is the
positive radial normalization of these Cox charts, with their
Hermitian metric transported; the normal-rescaling cutoff
flow of GF.2 carries their higher normal terms. We have
therefore constructed the actual periodic family flop at
each step, including its boundary products.

\par\medskip\noindent\textbf{AF.3. The exact seed and full target products.}\par\nopagebreak[4]

\item P3K.66 gives on the specified \(T_n\) family the exact fixed
native pair \((0,-6-18n(n-1))\) and physical integer
\(-1-3n(n-1)\). Its proof includes the actual line map,
full collar cube and cap cochains. GF.12's ordinary product
adds the pure-left endpoint \(L^6\), so its fixed entries
are \(V_{00}=18n(n-1)\), \(\delta_0=-12-18n(n-1)\).
For this basis the full winding matrix is
\[
 B=\begin{pmatrix}-6&0\\6n&0\end{pmatrix},\qquad
 w=(0,1,-1),\quad
 I=(12n-6,0,0),\quad W=(6-12n,-6,6).
\]
Substitution in the integral orbit-prism and derivative-prism
identities of K55 gives all remaining entries in (P3K.71.1).
Thus neither a free coordinate nor an integer lift is inferred
from a residue.

\item We retain the actual products determining (P3K.71.2) on every
intermediate target. The complete periods are
\[
 h:(Z,V)\mapsto(\kappa qZ,\alpha V),\qquad
 B:(Z,V)\mapsto(gZ,q\eta V).
\]
Their parameter windings are \((-6,0,0,0)\). Write the based
ratio of \(\kappa\) as \(e^{-6i\theta}e^{f_\kappa}\), and
the other ratios as \(e^{f_\alpha},e^{f_g},e^{f_\eta}\),
using their actual periodic logarithms. For the basis \(P\)
the logarithm after extracting the integral winding is
\[
 (P^{-1})_{i1}(P_{1j}f_\kappa+P_{2j}f_g)
 +(P^{-1})_{i2}(P_{1j}f_\alpha+P_{2j}f_\eta).
                                                            \tag{P3K.71.8}
\]
Its protected radial cutoff, together with the based smooth
quaternion maps \(H_d\) of GF.9, gives the full normal
contraction \(H_{-3}QH_{l_i+3}\). The complete three-curve
ordinary product, including both vertex collars, is GF.2's
normal-rescaling flow with these exact functions. Its native
Hopf-boundary map is
\[
 q_H\longmapsto (U H_{-3}U^{-1})q_HH_{l_i+3}.
                                                            \tag{P3K.71.9}
\]
The left factor has degree \(-6\) and the right factor
degree zero; the explicit inverse-cylinder extension gives
endpoint \(L^6\). The actual circle and stable derivative
traces of these products are precisely (P3K.71.2), by GF.3's
oriented prism maps. All unit homotopies in (P3K.71.8) are
included in this construction.

\par\medskip\noindent\textbf{AF.4. The selected free integers and unique reconstruction.}\par\nopagebreak[4]

\item At the selected old curve write \(a=V_{ii}\), \(b=\delta_i\).
The product-frame dictionary gives its native right integer
\(n_i=(b-a)/2\). The Cox common-link arrow is literally
\[
 T(y,q_H)=(q_H^{-1}yq_H,q_H),
\]
with the fixed complex-normal reset. Its proved action on
the native relative coordinates is
\[
 (m,n,s)\longmapsto(m+2n,-n,s+n).                    \tag{P3K.71.10}
\]
The commutator generator has \(\mu=+1/28\) with the physical
line and both cap orientations identified in K68. Thus the
last entry in (P3K.71.10) is the actual physical sphere class.
The constant normal interchange (P3K.71.6) conjugates both
Spin factors by orientation-preserving constant lifts and
preserves their third degrees. This verifies the same update
for both selected curves. Adding the same ordinary \(L^6\)
comparison on the new star gives
\[
 V'_{22}=a,\qquad\delta'_2=2a-b,\qquad
 s'=s+(b-a)/2\pmod{28}.                             \tag{P3K.71.11}
\]

\item The actual target family has the symmetric integral identities
\[
 V'(w_je_j-w_0e_0)=I_je_j-I_0e_0,\qquad
 w_j\delta'_j-w_0\delta'_0=W_j-W_0,\qquad j=1,2.
                                                            \tag{P3K.71.12}
\]
All its weights are nonzero by AF.2. The two vectors on the
left span the kernel of the functional with coefficients
\(r=(1/w_0,1/w_1,1/w_2)\). A symmetric difference
annihilating that plane is exactly \(c rr^t\). The diagonal
matrix \(\operatorname{diag}(I_j/w_j)\) is a particular
solution. The derivative equations leave exactly the analogous
single constant \(e\). The selected entries (P3K.71.11)
give the constants in (P3K.71.3), so that formula is the
unique complete solution.

\item Existence here is furnished by the actual target family and
the products just constructed. Its integral invariants must
equal this unique rational solution. The increment is its
actual native right integer. This proves all asserted
integrality, without division in \(\Theta_7\). Induction
through the finite word proves that the algorithm computes
the actual variation, stable derivative and physical class
on its prescribed fan. AF.1 identifies that family with the
original input marking by the explicit identity on the full fan.

\par\medskip\noindent\textbf{AF.5. Exact execution.}\par\nopagebreak[4]

\item The program \path{code/part3_root/source_all_fans.py} uses
integer and exact fraction arithmetic. It checks all 6,504
unimodular inputs with entries of absolute value at most 18,
representing 814 fans. Each execution checks the exact fan
representative, strictly decreasing factorization, matrix
reconstruction, integral prism identities, selected inverse
and native integer. A separate three-edge traversal from
the original seed checks 439 states through depth eight
against this factorization. Its four former transitions
with a separately retained fixed curve are not inputs to
this proof; their targets are computed from the appropriate
K66 seed by AF.2. There are also 123 exact checks against
the already proved vertical, horizontal and adjacent formulas.

\item For \(P=(\begin{smallmatrix}2&-1\\-1&1\end{smallmatrix})\)
the word is \((0,1)\) from \(T_0\). Its actual native right
increments are \(-6,-48\). Consequently
\([\Phi_P]=-1-6-48=1\pmod{28}\). The finite bounds belong
to the accompanying arithmetic certificate; AF.1--4 prove
termination and the claimed value for every unimodular input.

\end{enumerate}
\end{proof}




\begin{theorem}\label{night:P3K.72}
\StatementRecord{night:P3K.72}{proved}{S6\_06\_BLIND: GPT-6 Astra (OpenAI, ChatGPT Pro mode, effort Pro), recorded re-derivation in the checkpoint dated 18 November 2026; also checked in this writing by \CurrentChecker}{\cref{night:P3F.9}, \cref{night:P3K.55}, \cref{night:P3K.62}, \cref{night:P3K.64}, \cref{night:P3K.66}, \cref{night:P3K.68}, \cref{night:P3K.70}}
Retain an actual fixed-source whole-\(I_1\) canonical scalar family
of K26, K28 or K32, its original source, Cartier line, finite caps and
complete ordinary forward closing \(q^e u(q)\). Its scalar winding is
\(J\), with positive canonical parameter \(\lambda=r^J\), and its
initial physical characteristic pair is
\((k_{M,0},\ell_{M,0})=(3H,24\chi)\) in the K68 convention.
Replace only the selected cusp fan by
\(T_n=T(e_1+ne_2,e_2)\), using the literal character and period maps
below. Then its actual integer characteristic data and native fixed-star
free coordinates are
\[
\begin{aligned}
k_{M,n}&=3H-3Jn(n-1),&\ell_{M,n}&=24\chi,\\
k_{U,n}&=-3H+3Jn(n-1),&\ell_{U,n}&=-24\chi,\\
A_{M,n}&=H/2-\chi-Jn(n-1)/2,\\
(m_0^{\rm native},n_0^{\rm native})
  &=(6\chi,3H-6\chi-3Jn(n-1)).
\end{aligned}
\]
Here \(M\) is the physical line and \(U=M^*\) is its meridian-one
dual. Formula (SI.36) gives the full integer variation and derivative
seed in the literal doubled ordinary product. All full-unit maps,
annular integers and finite-cap factors are retained. This statement
concerns the specified canonical closing and its literal covers; it
does not include a new shorter source-translation endpoint or an OS49
component-supported construction.
\end{theorem}
\begin{proof}
\begin{enumerate}[label=\arabic*.]
\item \par

\medskip\noindent\textbf{SI.1. The closed scalar seam and the actual orbit pieces}

\par\nobreak

\item Retain the full cusp periods
\[
 h:(Z,V)\longmapsto(\kappa qZ,\alpha_LV),\qquad
 B:(Z,V)\longmapsto(gZ,q\eta V),\qquad \eta=\lambda\beta.
                                                        \tag{SI.1}
\]
Here \(\kappa,\alpha_L,g,\beta\) are the actual units, independent of the
scalar parameter. Their nonzero holomorphic logarithms on the small
protected \(q\) disc are fixed once. The scalar radius is fixed on the
parameter circle; \(\theta\) has period one. The only parameter winding
in this matrix is \(\operatorname{wind}\eta=J\).

\item Reindex the ordinary forward closing by the actual \(B^{-e}\). Its bounded
arrow is exactly
\[
 E:(Z,V)\longmapsto(g^{-e}Z,u\eta^{-e}V).
                                                        \tag{SI.2}
\]
On a cut parameter interval use the full torus multiplier
\(C_\theta=(A_\theta,D_\theta)\) of K64, with
\[
 C_0=1,\qquad C_1=(g^e,u^{-1}\eta^e),\qquad C_1E=C_0.
                                                        \tag{SI.3}
\]
The equality holds on the full protected collars, with all derivatives.
The functions are the actual logarithmic cut functions of the bounded
units and of \(\eta\); their angular winding is retained. They commute
with both periods and the physical circle. On the two vertex charts
their maps are, literally,
\[
 (a,b,c)\longmapsto((AD)^{-1}a,Ab,Dc),\qquad
 (u,v,w)\longmapsto(ADu,D^{-1}v,A^{-1}w).
                                                        \tag{SI.4}
\]
Both triple products remain \(q\). The complete three period rows
intertwine these maps, so they extend across the toric charts and
retain every \(q\)-normal derivative. The finite-cap closing square is
the original one, carried by the same ordinary arrow. Thus the closed
coordinates used below are coordinates on that actual family.

\item For the \(n\)-th fan take
\[
 Z_n=Z,\quad V_n=Z^{-n}V,\quad h_n=hB^n,\quad B_n=B.
\]
Direct substitution gives
\[
 Q_n=\kappa qg^n,\quad
 \alpha_{L,n}=\alpha_L\eta^n\kappa^{-n}g^{-n^2},\quad
 g_n=g,\quad\eta_n=\eta g^{-n}.                       \tag{SI.5}
\]
The source parameter \(Q_n\) is independent of \(\theta\). The two line
windings are \(nJ,J\). In these characters the seam is exactly
\((A,D A^{-n})\), so (SI.3) remains the actual endpoint equality. On a
multiple parameter cover its successive laps are \(C_\theta E^{-j}\);
this is the literal pullback, not a new seam choice.

\item Take the physical-circle quotient. The F9 node maps are the same
actual maps in these vertical characters:
\[
 2b=r-|y|^2/r,\quad 2b'=|x|^2r-r^{-1},\qquad
 \frac{db'}{db}=\frac{|x|^2+r^{-2}}{1+|y|^2r^{-2}}>0.
                                                        \tag{SI.6}
\]
The contracting \(B\) coefficient is the negative of the increasing
F9 moment coordinate. Its real-period formula is
\[
 \gamma_n=
 \frac{\log|Q_n|\log|V_n|-\log|\alpha_{L,n}|\log|Z_n|}
 {\log|Q_n|\log|q\eta_n|-\log|\alpha_{L,n}|\log|g|}.
                                                        \tag{SI.7}
\]
The denominator has leading term \((\log|q|)^2\). It is nonzero on a
uniform small disc. The complete norms in (SI.7) are independent of
\(\theta\); parameter phases have not been deleted from the circle
bundle. Substitution of \(h_n\) leaves \(\gamma_n\) unchanged and
substitution of \(B\) increases it by one.

\item The F9 conifold interval has representatives \(z=\sqrt r\), \(v=u=0\),
\(w=1/\sqrt r\). Choose on it
\[
 \gamma_n=\tfrac12-\varepsilon\frac{r-1}{r+1},
 \qquad0<\varepsilon<\tfrac14.                       \tag{SI.8}
\]
Its endpoint lifts are \(1/2+\varepsilon,1/2-\varepsilon\) and the
endpoint derivatives in the two F9 charts are both \(-4\varepsilon\).
The fixed sphere's second endpoint is the \(h_n^{-1}\) translate of the
second conifold endpoint, hence has the same gamma lift by (SI.7).
There is no extra integer between them. Extend the two proper
decreasing coordinate maps, using a convex interpolation equal to
the prescribed real-period map on the outer collar. Its derivatives
remain negative and its lifts commute with addition of one. The
normalization cylinder then has a graph in this same short arc.

\item Untwist the source part of \(B\) by its actual \(\exp(\gamma\log g)\)
isotopy, and the line norm by its full positive-norm isotopy. This
gives the orbit piece
\[
 W_n=E_n\times S^1_{\gamma_n},
                                                        \tag{SI.9}
\]
where \(E_n\) is the nodal source neighbourhood with Tate parameter
\(Q_n\). This assertion is relative to the exact outer coordinates
just prescribed. It also holds over the closed parameter: (SI.3)
closes the toric charts, its multipliers commute with all periods,
and the quotient construction uses precisely their unchanged
source and real-norm data. Only the line phase has parameter
winding. That phase is retained in SI.3 below.

\par

\medskip\noindent\textbf{SI.2. The source, fixed tube and marked Cox cap map}

\par\nobreak

\item Write \(L_g=\log g\). Interpolate the actual source parameters by
\(Q_s=\kappa q\exp(snL_g)\), \(0\le s\le1\). Near the node \(xy=Q_s\)
the horizontal vector field is
\[
 \partial_s+nL_g y\partial_y.                         \tag{SI.10}
\]
It keeps \(q,\theta\) fixed and preserves the node and its two branches.
On the outer regular collar prescribe instead the real marked-period
map \(u+\tau_s v\mapsto u+\tau_0v\). Both are horizontal over the same
\(\partial_s\); their difference on the overlap is vertical. A partition
equal to one on the protected node and outside collars interpolates
them in this affine space. On the compact source neighbourhood the
resulting boundary-tangent field has a complete flow. It is periodic
in \(\theta\) and is precisely the real-coefficient map on the entire
outer collar. No real period expression is used as a smooth node map.

\item Transport the fixed normalization sphere by this source flow. Its
gamma graph and the graph at \(n=0\) lie in the short arc (SI.8), with
the same ordered, distinct endpoint values. Their convex interpolation
remains embedded: the source normalization is injective except at its
two node preimages, and those two heights stay distinct. The velocity
extends from the sphere to a tube and then by a cutoff. This gives an
actual ambient isotopy, fixed on the outside collar and periodic in
the parameter.

\item The normal map is also retained. In the native fixed star
\((z,n_0,n_1)=(Z,q/(ZV),V)\) the character map is exactly
\[
 (z,n_0,n_1)\longmapsto(z,z^n n_0,z^{-n}n_1),
 \qquad q=zn_0n_1.                                  \tag{SI.11}
\]
Let \(D_0,D_\infty\) denote the two marked branch sections of the
normalization. In the positive Hopf spinor convention the corresponding
bundle arrow is the Cartier arrow
\[
 E_n^{\rm spin}=E_0^{\rm spin}\otimes
               \mathcal O(n(D_0-D_\infty))
       \longrightarrow E_0^{\rm spin},\qquad v\longmapsto Z^n v.
                                                        \tag{SI.12}
\]
In its Cartier frames the displayed apparent zeros and poles cancel.
In (SI.11) the negative complex normal is conjugated when passing to
the Hopf spinor; transporting the Hermitian norms leaves the common
phase of \(Z^n\), and the remaining real positive factors are the
explicit positive normal rescalings. Thus (SI.12) is the typed
normal-bundle map of (SI.11), with the same physical circle and
physical height. Its adjoint map carries the normal three-plane.
The source flow contributes its actual derivative and the graph
isotopy contributes its positive triangular normal map. These
comparisons extend through their isotopies.

\item Extend the boundary sphere isotopy over its three-ball by a collar
cutoff of its boundary velocity. On the Hopf sphere factor use the
adjoint arrow of (SI.12). This is the same literal product-cap
construction as SQ.8--9, now with the constant marked nodal source;
it does not replace the Cox normal map by a base tangent derivative.
It supplies a diffeomorphism of completed pieces, including their
actual cap and circle-line marking,
\[
 C_n=(W_n\setminus\operatorname{int}(F_n\times D^3))
                      \cup(D^3\times S^2)\longrightarrow C_0.
                                                        \tag{SI.13}
\]
Its circle line is compared next. The finite caps and the rest of the
quotient exterior are the original ones throughout.

\par

\medskip\noindent\textbf{SI.3. The full integral line and its prescribed boundary arrow}

\par\nobreak

\item Use the actual complex period coefficients
\[
 (z,\ell)=u(1,0)+v(\tau,\mu)+d(0,1)+\gamma(G,\zeta),
                                                        \tag{SI.14}
\]
where \(z=\log Z/(2\pi i)\), \(\ell=\log V/(2\pi i)\),
\(\tau=\log(\kappa q)/(2\pi i)\), \(\mu=\log\alpha_L/(2\pi i)\),
\(G=\log g/(2\pi i)\), and \(\zeta=\log(q\eta)/(2\pi i)\).
The imaginary parts invert (SI.7); the real parts then determine
\(u,d\). Thus differentiating \(d=\Re\ell-\Re\mu\,v-\Re\zeta\,\gamma\)
retains every unit and its derivative. The coordinates are the closed
ones in (SI.3). With \(\alpha=\arg q/(2\pi)\), their integer deck maps are
\[
 q\text{-turn}:(u,d)\mapsto(u-v,d-\gamma),\qquad
 \theta\text{-turn}:(u,d)\mapsto(u,d-J\gamma).
                                                        \tag{SI.15}
\]
The bounded source and line factors in the seam give their actual
periodic logarithms, not further unrecorded integral changes.

\item Direct substitution of both characters in (SI.14) gives
\[
 u_n=u,\quad v_n=v,\quad\gamma_n=\gamma-nv,\quad d_n=d-nu.
                                                        \tag{SI.16}
\]
The meridian-one line is \(U=M^*\), with \(M\) the physical line, exactly
as in K68. In these coefficient frames choose the actual boundary
connections
\[
 A_0=dd+(\alpha+J\theta)d\gamma,\qquad
 A_n=dd_n+(\alpha+J\theta)d\gamma_n+nJ\theta\,dv.
                                                        \tag{SI.17}
\]
They are invariant under their full integer deck maps. In particular
the new theta map is \(d_n\mapsto d_n-J\gamma_n-nJv\). In physical
coordinates the connection includes all the differentials in
(SI.14); it does not discard the units. Their exact difference is
\[
 \beta_n=A_n-A_0=-n(du+\alpha\,dv).                  \tag{SI.18}
\]
The coframe with the full Tate phase is
\(du+(\alpha+\phi_\kappa)dv\); consequently (SI.18) retains the
additional \(n\phi_\kappa dv\) when expressed in that coframe.

\item Let \(b:C_0\to\mathbb R/\mathbb Z\) be the source graph map with boundary
value \(v\). It extends over the Lefschetz two-handle because the
attaching vanishing circle is the \(u\) circle and \(v\) is constant on
it. On the fixed sphere its restriction has a real lift. Modify that
lift by a supported real homotopy to make \(b\) constant on a smaller
fixed tube; it then extends constantly over the cap. The marked
source and both branch sections return literally in the closed
coordinates, so this construction is periodic. Its precise theta
clutch change is implemented by \(\exp(2\pi i nJ\theta h)\), where \(h\)
is the supported real difference of the two graph maps.

\item On the completed six-dimensional piece \(\mathcal C\) define the line
\(\mathcal T\) by its theta-cut coefficient gluing
\(w\mapsto e^{2\pi i b}w\). Its connection in those frames is
\(d-2\pi i\theta\,db\), so
\[
 c_1(\mathcal T)=\Xi=d\theta\wedge db,\qquad\Xi^2=0.
                                                        \tag{SI.19}
\]
On the normalization graph cut, the actual constant factor
\(\alpha_{L,n}/\alpha_L\) has winding \(nJ\). The theta-cut description is
\(d_n\mapsto d_n-nJb\) relative to the old line, exactly the dual
coefficient gluing of \(\mathcal T^{nJ}\). On the doubly cut rectangle
the coordinate change \(d_n\mapsto d_n+nJ\theta b\) identifies these
two descriptions. The remaining logarithms are explicitly
\[
 nf_\beta-nf_\kappa-n^2f_g,
 \qquad -nf_g,                                       \tag{SI.20}
\]
with the fixed scalar-radius constant retained. They extend in the
chosen full unit charts.

\item For the integral check, \(H_1(C_0)=\mathbb Z\langle b,\gamma\rangle\) and
\(H_2(C_0)=\mathbb Z\langle\text{Hopf meridian},b\times\gamma\rangle\).
The removal Thom sequence adds the normal meridian, since the
regular fibre times gamma misses \(F\); the cap identifies that
meridian and kills the primitive fibre class. All groups are free.
The parameter maps fix \(b,\gamma\) and the oriented meridian. They
cannot add a meridian to the regular \(b\times\gamma\) class: the actual
Euler line has value zero on that real-period torus and one on the
meridian. Wang therefore adds precisely the two classes
\(b\times\theta,\gamma\times\theta\). The line difference has value
\(nJ\) on the first, zero on the second because \(\eta_n/\eta=g^{-n}\)
has a logarithm, and zero on both old generators by (SI.12) and the
primitive real-period section. Consequently
\[
 c_1(U_n)-c_1(U_0)=nJ\Xi.                            \tag{SI.21}
\]

\item We require the particular outer isomorphism, not only this equality.
The actual boundary \(\mathcal B\) has
\[
 H_1(\mathcal B)=\mathbb Z\langle\alpha,\theta,v,\gamma\rangle,
 \quad \operatorname{im}H^1(\mathcal C)=
       \mathbb Z\langle d\theta,db,d\gamma\rangle.
                                                        \tag{SI.22}
\]
The \(q\) monodromy kills \(u\) primitively. There is no source theta
monodromy in this scalar case. The gamma map extends through the cap
by its real lift on the short fixed graph, and the theta family has
the explicit free section supplied by the closed coordinates.
Thus the only boundary winding that does not extend is the
\(q\)-angle winding.

\item Choose one parameter value \(\theta_*\), gamma zero outside (SI.8),
and a constant Tate section \(Z=Z_*\) with \(|Z_*|=1\), at a generic
smooth normalization point avoiding the finitely many marked source
sections and the two branches at that parameter. On a small parameter
arc the same point and a uniformly smaller \(q\) disc work. The
relations \(\log|Z_*|=\log|V|=0\) give \(v=\gamma=0\) literally in
(SI.7) and (SI.14). In the actual closed first-vertex charts the
section-disc map is
\[
 (q/Z_*,Z_*,1)\longmapsto
              (qZ_*^{n-1},Z_*,Z_*^{-n}).              \tag{SI.23}
\]
The height is \(q\) on both sides. Both nonvanishing coordinates
remain nonzero at \(q=0\), so the disc lies in the free locus and
misses the fixed surgery. The line arrow \(Z_*^{-n}\) extends on
the full disc, and its inverse is an opposite power of the same
nonzero constant. In the original coordinates all seam factors
are retained in the inverse of (SI.4); they are extending units.
The source horizontal flow can fix this section on that parameter
arc, since its prescribed boundary field fixes \(v=0\) and the node
is disjoint; extend this choice periodically by a parameter cutoff.
This disc at \(\theta_*\) tests the generator of the sole alpha
obstruction in (SI.22). That integral winding is constant on the
connected boundary, so its zero value proves the assertion for the
whole prescribed boundary line arrow. No simultaneous avoidance of
moving Cartier sections over the full parameter circle is needed.

\item An explicit Cech construction now supplies
\[
 J_n:U_n\longrightarrow U_0\otimes\mathcal T^{nJ}
                                                        \tag{SI.24}
\]
with exactly the boundary arrow (SI.16)--(SI.18). Form the actual
unit transitions of the ratio line, including (SI.20), on a good
cover of \(\mathcal C\). Their logarithms have an integral two-cocycle
which vanishes by (SI.21). Subtract its integral one-cochain. For
the resulting real one-cocycle \(r_{ij}\) in frame convention
\(e_j=e^{2\pi i r_{ij}}e_i\), set \(a_i=\sum_k\rho_k r_{ki}\); then
\(a_j-a_i=r_{ij}\). Multiplication by \(e^{-2\pi i a_i}\) gives its
global unit frame; invert the corresponding Hom map to obtain
(SI.24). Its boundary ratio with the prescribed map has no alpha
winding by (SI.23). Remove its other integral windings using the
extending functions theta, \(b\), gamma in (SI.22), then extend the
remaining real logarithm with a collar cutoff. This makes the
boundary map literally the prescribed one. On the inner cap any
remaining theta gauge extends over \(D^3\times S^2\times S^1\).
The resulting map retains every annular integer and finite-cap
character.

\par

\medskip\noindent\textbf{SI.4. Relative Pontryagin pairing with the full tangent collar}

\par\nobreak

\item On the outer collar use the stable coframe
\[
 du+(\alpha+\phi_\kappa)dv,\quad dv,\quad d\gamma,
 \quad d\alpha,\quad d\theta,\quad dr_q.
\]
The actual piece map (SI.13), prescribed in the real-period coordinates
on this whole collar, has coframe matrix
\[
 G=I+n\phi_g E_{12}-nE_{32},\qquad
 G^{-1}dG=n\,d\phi_g E_{12}.                         \tag{SI.25}
\]
Here \(\phi_g=\Im\log g/(2\pi)\) is independent of theta. The line
phase \(J\theta\) belongs to (SI.17), not to an unrecorded tangent
block; its positive norm has already entered gamma. The Maurer
form squares to zero. For an interpolation of the actual flat
tangent connections by \(\chi(r)G^{-1}dG\),
\[
 F=\chi' dr\wedge\omega+(\chi^2-\chi)\omega^2,
 \qquad\operatorname{tr}F^2=0.                       \tag{SI.26}
\]
Thus the full tangent \(p_1\) form can be chosen zero throughout
the gluing collar. This uses the exact matrix, not only flatness
at its ends. The source and graph isotopies are fixed there.

\item On \(C_0\) choose nested fixed tubes so that all surgery and
connection modifications lie inside an inner tube, and choose
the tangent connection outside it pulled from the actual source
\(E_0\), with a trivial gamma line. It has the specified flat source
part on the outer collar. At a fixed theta, a regular product level
of \(b:E_0\to S^1\) is \(N^3\times S^1_\gamma\), and the pulled-back
source \(p_1\) four-form restricts to zero on it. Modify \(b\) by its
real lift to be constant on a larger tube. A regular value distinct
from that constant gives a level missing the entire surgery and
all connection changes. The regular-value cobordism of this
supported homotopy is relative to the original outside boundary;
the closed \(p_1\) form and its zero outer trace show its integral
remains zero. The new level need not be a product. Its two normal
directions in \(\mathcal C\) are the framed \(b\) and theta directions.
Consequently, with the retained flat boundary connection,
\[
 \langle p_1(T\mathcal C),\Xi\rangle=0.              \tag{SI.27}
\]
Here \(p_1\) is relative and \(\Xi\) is absolute. Equation (SI.24)
makes the line change exactly \(nJ\Xi\) plus the retained boundary
primitive. The former pairs to zero by (SI.27); the exact terms
pair to a boundary integral which vanishes on the whole collar
by (SI.26). The exterior is unchanged. Therefore
\[
 \ell_{U,n}=\ell_{U,0}=-24\chi,
 \qquad\ell_{M,n}=24\chi.                           \tag{SI.28}
\]

\par

\medskip\noindent\textbf{SI.5. The whole collar cube, its orientation and the integer seed}

\par\nobreak

\item Use the K68 physical/dual dictionary. In the auxiliary positive
orientation
\[
 \Omega_+=dr_q\wedge d\alpha\wedge du\wedge dv
                        \wedge d\gamma_B\wedge d\theta,
\]
the real curvature \(F=dA\) is that of \(U\). The actual physical
line and the physical-to-Jin parameter reversal give
\[
 k_M=\int_{\Omega_+}F^3,
 \qquad k_U=-\int_{\Omega_+}F^3.                     \tag{SI.29}
\]
These are exactly the K66/K68 line and orientation maps, now
applied to the same physical Cox normal and positive scalar
parameter. No generator is selected from its residue.

\item On the common completed piece (SI.24) permits the affine
curvature \(F_n=F_0+nJ\Xi\). Since \(\Xi^2=0\), its cube change is
linear in \(n\). The common exterior connection is kept. In the
piece coefficients, write gamma for \(\gamma_n\); thus
\(\gamma_{\rm ext}=\gamma+nv\), \(d_{\rm ext}=d_n+nu\). Its full
boundary connection and curvature pull back to
\[
 A_{{\rm ext},n}=dd_n+(\alpha+J\theta)d\gamma
                         +n\{du+(\alpha+J\theta)dv\},
 \quad F_{{\rm ext},n}=(d\alpha+Jd\theta)
                                      \wedge(d\gamma+n\,dv).
                                                        \tag{SI.30}
\]
The piece-minus-exterior form is exactly
\(\beta_n=-n(du+\alpha dv)\), and \(d\beta_n=-n d\alpha\wedge dv\).
Both \(F_{{\rm ext},n}^2\) and \((d\beta_n)^2\) vanish. Orient the
outward boundary of the piece by
\[
 \Omega_{\mathcal B}=d\alpha\wedge du\wedge dv
                                    \wedge d\gamma\wedge d\theta.
\]
Its marked unit cell has integral one; the actual nilmanifold
monodromy preserves this cell and volume. Direct exterior algebra
gives
\[
 3\int_0^1\beta_n\wedge
            (F_{{\rm ext},n}+s\,d\beta_n)^2\,ds
                       =3J n^2\Omega_{\mathcal B}.     \tag{SI.31}
\]
The connection on an outward gluing collar is
\(A_{{\rm ext},n}+\chi(t)\beta_n\) with \(\chi(0)=1\), \(\chi(1)=0\).
Its actual cube integral therefore runs from \(s=1\) to \(s=0\),
and is \(-3Jn^2\). Combining the interior linear term with this
complete collar yields one fixed real constant \(L\) such that
\[
 k_{M,n}=k_{M,0}+Ln-3Jn^2.                           \tag{SI.32}
\]
The relative map (SI.24) and its prescribed outer gauge are what
exclude a further Mayer--Vietoris term here.

\item The exact \(n=1\) comparison is K62's C2 Cox bordism. Its first
outer map is the identity and its second is base \(\kappa^{-1}\),
normal \(\operatorname{diag}(\alpha_L^{-1},\kappa\alpha_L)\).
They are scalar-parameter-independent. The two inner charged
boundaries carry one and the same physical-delta gauge, so it
extends through these constant outer maps. On the additional
normal-sphere boundary, both its determinant line and its normal
three-plane are pulled from the three-dimensional fixed factor.
Thus \(b^2=p_1(\nu)=p_1(T\bar X)=0\). The individual sphere-bundle
identities \(4\int x^3=3b^2+p_1(\nu)\) and
\(\int p_1x=p_1(T\bar X)+p_1(\nu)\) show that both characteristic
corrections vanish. This is the narrow integer comparison proved
in SCN.9, with every source, cap and closing factor retained.
The C2 target basis \((e_1,e_1+e_2)\) is exactly \(T_1\):
its second triangle translated by \(-e_1\) is the first \(T_1\)
triangle, and its first translated by \(e_2\) is the second.
These translations are the actual \(h^{-1}\) and \(B\) periods with
all their units, so the induced quotient/cap arrow is the literal
identity in the absolute charts. Thus \(k_{M,1}=k_{M,0}\) in the
same Cox marking used in (SI.13). It follows from (SI.32) that
\(L=3J\).

\item The complete resulting integer pair and index are
\[
 \begin{aligned}
 k_{M,n}&=3H-3J n(n-1),&\ell_{M,n}&=24\chi,\\
 k_{U,n}&=-3H+3J n(n-1),&\ell_{U,n}&=-24\chi,\\
 A_{M,n}&=H/2-\chi-Jn(n-1)/2.
 \end{aligned}                                      \tag{SI.33}
\]
The native fixed-star identities from the actual physical
Hopf transfer are \(k_M=m+n_0\), \(\ell_M=4m\). Hence
\[
 (m_0^{\rm native},n_0^{\rm native})
       =(6\chi,3H-6\chi-3Jn(n-1)).                   \tag{SI.34}
\]
The polynomial has been obtained from the whole relative line
and tangent calculation. The congruence K64 is an independent
consistency check, not an input determining an integer lift.

\item For the literal doubled ordinary plumbing product of K64/K70,
the native-to-ordinary comparison is the actual pure-left
\(-2J\) endpoint. Thus its fixed entries, ready for the actual
nonzero-weight target reconstruction, are
\[
 V_{00}^{(2)}=2J-6H+6Jn(n-1),\qquad
 \delta_0^{(2)}=2J+6H-24\chi-6Jn(n-1).
                                                        \tag{SI.35}
\]
For completeness, the retained K55/K64 trace maps on this literal
doubled product have
\[
 B=P_n^{-1}\operatorname{diag}(0,2J)P_n
       =\begin{pmatrix}0&0\\2Jn&2J\end{pmatrix},\qquad
 w=(0,1,-1),
\]
\[
 I=(2J(2n-1),2J,-2J),\qquad
 W=(2J(1-2n),0,0).
\]
Their exact integral graph/prism equations consequently give the
complete ordinary seed
\[
 V^{(2)}=
 \begin{pmatrix}
 2J-6H+6Jn(n-1)&2J(1-2n)&2J(2n-1)\\
 2J(1-2n)&2J&0\\
 2J(2n-1)&0&2J
 \end{pmatrix},
\]
\[
 \delta^{(2)}=
 \begin{pmatrix}
 2J+6H-24\chi-6Jn(n-1)\\2J(2n-1)\\2J(1-2n)
 \end{pmatrix}.                                      \tag{SI.36}
\]
No division in \(\Theta_7\) occurs in this statement. The exact
scope remains the canonical closing (SI.2), its literal parameter
covers, and the actual physical-vertical-edge fans.

\par\medskip The exact certificate is \path{code/part3/scalar_vertical_integer_seed.py}, with its saved JSON, output and full-unit scope.

\end{enumerate}
\end{proof}




\begin{theorem}\label{night:P3K.73}
\StatementRecord{night:P3K.73}{proved}{S6\_06\_BLIND: GPT-6 Astra (OpenAI, ChatGPT Pro mode, effort Pro), recorded re-derivation in the checkpoint dated 18 November 2026; also checked in this writing by \CurrentChecker}{\cref{night:P3F.24}, \cref{night:P3K.64}, \cref{night:P3K.68}, \cref{night:P3K.69}, \cref{night:P3K.72}}
Retain an actual fixed-source, whole-zero \(I_1\), canonical scalar
sphere family of P3K.26, P3K.28 or P3K.32, or the original coprime
family with the literal AC parameter endpoint of P3K.67, SE.5,
with its marked Cartier
line, source translation, finite characters, integral annular
cochain and forward ordinary closing. Write its parameter as
\(\lambda=r^{M_0}\), and let \(H,\chi\) be its actual reference-page
data in P3K.69. We take either its positive \(r\)-loop, or its
constructed shorter positive \(w=r^b\)-loop with all the endpoint
data of P3K.67. Write \(c=1\) in the first case and \(c=b\) in the
second, and set
\[
 J=M_0/c,\qquad \widehat H=H/c,\qquad
 \widehat\chi=\chi/c.
 \tag{P3K.73.1}
\]
The integers in the ensuing algorithm are obtained from these
actual families; \(\widehat H\) and \(\widehat\chi\) themselves
are permitted to be rational.

For any \(P\in GL_2(\mathbb Z)\), use the identical-fan maps and
terminating subtraction of P3K.71, AF.1--2, to obtain the unique
vertical seed \(T_n\) and the finite word in \(Q_0,Q_1\) ending
at the canonical representative of \(T(P)\). Start with
\[
 V=\begin{pmatrix}
 2J-6\widehat H+6Jn(n-1)&2J(1-2n)&2J(2n-1)\\
 2J(1-2n)&2J&0\\2J(2n-1)&0&2J
 \end{pmatrix},
\]
\[
 \delta=\begin{pmatrix}
 2J+6\widehat H-24\widehat\chi-6Jn(n-1)\\
 2J(2n-1)\\2J(1-2n)
 \end{pmatrix},\qquad
 S=\widehat H/2-\widehat\chi-\frac J2n(n-1).
 \tag{P3K.73.2}
\]
Here \(V,\delta\) are the actual ordinary-product invariants of
the literal double parameter cover, whereas \(S\) records the
single-loop physical class as an integer path sum.

At a step \(P\mapsto P'=PQ_i\), \(i=0,1\), put
\(a=V_{ii}\), \(b_0=\delta_i\). At the target compute
\[
 B=(b_{jk})=2J(P')^{-1}e_2e_2^tP',\qquad
 l=(b_{22}-b_{21},b_{11}-b_{12},b_{22}),
\]
\[
 (x,y)^t=(P')^{-1}e_2,\quad
 w=(x,y,-x-y),\quad d=(-x-2y,-2x-y,x-y),
\]
\[
 I_j=Jw_j+d_j(l_j-J),\qquad
 W_j=Jw_j-d_j(l_j-J),\qquad j=0,1,2.
 \tag{P3K.73.3}
\]
All three target weights are nonzero. Set
\[
 C=w_2^2a-w_2I_2,\qquad E=w_2(2a-b_0)-W_2,
\]
\[
 V'_{jk}=\frac{I_j}{w_j}\delta_{jk}+\frac{C}{w_jw_k},
 \qquad \delta'_j=\frac{W_j+E}{w_j},\qquad
 S'=S+\frac{b_0-a}{4}.
 \tag{P3K.73.4}
\]
Every displayed free invariant and increment is an integer.
The actual return is \([\Phi_P]=S\pmod{28}\), in the physical
complex orientation and with \(\mu(\Sigma_j)=j/28\).
This computes every lattice-basis fan of these specified scalar
families, including their specified shorter endpoints. For the
primitive OS45 datum of P3K.70 on the basis
\(\left(\begin{smallmatrix}2&-1\\-1&1\end{smallmatrix}\right)\)
the value is \(11\).
For the original \((3,4)\) sphere family with its specified positive
\(w=e^{-\pi i c_{\rm AC}}\)-loop, these actual fan choices realize
all twenty-eight elements of \(\Theta_7\). In particular the basis
\(\left(\begin{smallmatrix}1&-2\\0&1\end{smallmatrix}\right)\)
has class \(3\).
\end{theorem}
\begin{proof}
\begin{enumerate}[label=\arabic*.]
\item \par\medskip\noindent\textbf{SA.1. The actual shorter fixed star.}\par\nopagebreak[4]
For the whole \(r\)-loop, P3K.72 constructs the completed orbit
pieces, their physical line and full tangent collar. Its integer
pair at \(T_n\) is
\[
 (k_{M,n},\ell_{M,n})=
 (3H-3M_0n(n-1),24\chi).
 \tag{P3K.73.5}
\]
We spell out the additional arrow for the shorter loop. In SE.7--14
the bounded forward displacement \(d(v)\), the full vector
\(\sigma(v)\), and the cut multipliers \(C_s(v)\) satisfy
\[
 \sigma(v+1/b)-\sigma(v)=d(v),\qquad
 C_s(v+j/b)P_j(v)=C_{j/b+s}(v),
\]
\[
 P_b(v)=\exp(-H_0\delta-\varepsilon\gamma).
 \tag{P3K.73.6}
\]
Here \(H_0\) is the specified ordinary logarithm, including its
bounded unit, and not the page number \(H\). The formula retains
the scalar correction \(\varepsilon M_0(b-1)/(2b^2)\) and every
intermediate cap-character change. Its actual cap squares and
integral lattice cochains are SE.5--6.

\item In the vertical characters the fixed native star has
\[
 (z,n_0,n_1)=(b_1,a,c_1),\quad
 (z',n'_0,n'_1)=
 (\kappa_n w_1,\alpha_n v_1,(\kappa_n\alpha_n)^{-1}u).
\]
For a cut multiplier \((A,D)\), its first-chart factors are
\((A,(AD)^{-1},D)\) and its second-chart factors are
\((A^{-1},D^{-1},AD)\). The Cox lift is literally
\[
 (a_0,a_1,c_0,c_1)\longmapsto
 (a_0,Aa_1,(AD)^{-1}c_0,Dc_1).
 \tag{P3K.73.7}
\]
It preserves \(q=zn_0n_1\) and intertwines the two native charts.
Equation (P3K.73.6) is an equality of both character multipliers,
so it holds on this lift, including every intermediate seam.
Thus the covered native fixed-star product is precisely the
\(b\)-fold pullback of the shorter one. The canonical product
used in (P3K.73.5) is carried to it by NV.4's explicit based
global physical-circle action and the bounded-unit homotopies
with zero endpoint logarithms. These extend over the original
exterior and finite caps. The normal-rescaling flow carries their
full \(q\)-jets.

\item On the actual completed parameter-circle quotients this is an
oriented degree-\(b\) covering. The physical line, its cap
identification and the stable tangent bundle pull back literally.
Integration of their integral characteristic classes gives
\[
 b k_{M,n}^{\rm short}=k_{M,n}^{r},\qquad
 b\ell_{M,n}^{\rm short}=\ell_{M,n}^{r}.
 \tag{P3K.73.8}
\]
The native relative coordinates also multiply by \(b\) on the
specified products. Their values are therefore
\[
 (m_n,n_n)=
 (6\widehat\chi,
 3\widehat H-6\widehat\chi-3Jn(n-1)).
 \tag{P3K.73.9}
\]
These are equalities in free integral groups on the actual
cover, and determine integers uniquely. P3K.68's physical
Hopf transfer gives the single-loop seed in (P3K.73.2).
No sphere-group element has been divided.

\par\medskip\noindent\textbf{SA.2. The exact family on the chosen fan.}\par\nopagebreak[4]
The full periods on the original cusp chart are
\[
 h:(Z,V)\mapsto(\kappa q Z,\alpha V),\qquad
 B_0:(Z,V)\mapsto(gZ,q\eta V),\quad \eta=\lambda\beta.
 \tag{P3K.73.10}
\]
Every bounded factor and its protected logarithm is retained.
On the shorter loop, \(b\mid M_0\), so the periods themselves
return while the exterior finite characters change by their
specified endpoint map. The transformed cut multiplier at a
basis \(P\) is the character-wise image under \(P^{-1}\) of
the multiplier in (P3K.73.6). The multiplication identity in
that equation is preserved by this integral monomial map.

\item More explicitly, the two vertex actions remain
\[
 (a,b_1,c_1)\mapsto((AD)^{-1}a,Ab_1,Dc_1),\qquad
 (u,v_1,w_1)\mapsto(ADu,D^{-1}v_1,A^{-1}w_1)
 \tag{P3K.73.11}
\]
in the new characters. Each preserves its triple product \(q\).
The complete period matrix \(A=(A_{hj})\) changes by
\[
 A'_{ik}=\prod_{h,j}A_{hj}^{(Q^{-1})_{ih}Q_{jk}}.
 \tag{P3K.73.12}
\]
The \(q\)-valuation matrix is the identity, so it remains the
identity under this simultaneous change of characters and
period generators. Thus the displayed factors are full
holomorphic units on the protected disc. They act on every
toric chart and commute with the original full period action.
In each native star their Cox lift is the same diagonal
character lift as (P3K.73.7), with its native order retained.
It extends to both rulings. This proves the required actual
shorter endpoint on all the intermediate fans, rather than
only at the vertical seed.

\item For completeness, the family-identification part of AF.1 uses
the identity on the absolute torus \((Z,V,q)\). The bases
\(P,PR,PR^2,-P,-PR,-PR^2\), with
\(R=\left(\begin{smallmatrix}-1&-1\\1&0\end{smallmatrix}\right)\),
give identical full fans by the explicit triangle translations
there. Swapping columns handles determinant minus one. The
matched cones have identical character semigroups, so the
identity extends on every affine chart, intertwines all full
period multipliers and every closing arrow, and is the original
identity on the source and cap overlap. It preserves the
physical circle. The unique cone containing its primitive ray
gives exactly the canonical representative used in the statement.
AF.2's integer subtraction strictly decreases \(p+r\) and
constructs its finite word from \(T_n\). Every subsequent
target has weights \((r,p,-r-p)\), all nonzero.

\item At each word step the actual ruling maps are the AF.2 Cox maps
\((n_0/n_1,zn_1,n_1)\) or \((n_1/n_0,zn_0,n_0)\), with their
other affine charts and the same rank-one matrix. The constant
unitary interchange displayed there carries the second choice's
normal order. It is connected to the identity by that explicit
path and preserves real orientation. These maps keep \(q\)
exactly, carry (P3K.73.11)--(P3K.73.12), and are the original
source/line coordinates on their common exterior. Thus every
step is an actual holomorphic periodic-family flop with its
original finite caps and specified parameter closing.

\par\medskip\noindent\textbf{SA.3. The full ordinary product and traces.}\par\nopagebreak[4]
Take the literal double cover of the selected single loop.
Its winding matrix on the original characters is
\(2Je_2e_2^t\). The exact monomial map
(P3K.73.12) gives (P3K.73.3) at the target. For each bounded
period unit, subtract the displayed integral winding and retain
its based logarithm. Its transformed logarithm is explicitly
\[
 f'_{ik}=\sum_{h,j}(P^{-1})_{ih}P_{jk}f_{hj}.
 \tag{P3K.73.13}
\]
The full \(q\)- and parameter-dependent functions occur here;
the endpoint identities (P3K.73.6) hold on whole collars.

\item The common determinant winding is \(2J\). The actual ordinary
quaternion products of K64 and GF.2 therefore have normal
contractions
\[
 Q\longmapsto H_J Q H_{l_j-J}.
 \tag{P3K.73.14}
\]
The based maps \(H_d\) are the fixed smooth maps of GF.9.
Multiplication by the exponentials of (P3K.73.13), with the
protected radial cutoffs, carries every bounded unit. GF.2's
normal-rescaling flow carries the remaining higher normal
terms, since the physical height has normal order at least
two along an open curve and three at a vertex. This constructs
one ordinary product on the whole three-curve plumbing. Its
vertex maps, oriented orbit prisms, endpoint caps and stable
derivative maps are the literal GF.3 maps with these complete
functions. The degree computation there gives exactly
\(I_j=Jw_j+d_j(l_j-J)\) and
\(W_j=Jw_j-d_j(l_j-J)\).

\item On a native star the comparison is
\[
 q_H\longmapsto (UH_JU^{-1})q_H H_{l_j-J}.
 \tag{P3K.73.15}
\]
Its left degree is \(2J\), its right degree is zero, and its
inverse-cylinder endpoint is \(L^{-2J}\). On the fixed star
the precise native order was checked in (P3K.73.7) and NV.2;
on the other stars it is the SE.3 order. This proves the same
comparison with each separately retained native product.
The doubled version of (P3K.73.9), followed by this pure-left
comparison and the full vertical prism equations, gives every
entry of (P3K.73.2).

\par\medskip\noindent\textbf{SA.4. The native increment and unique target.}\par\nopagebreak[4]
Let \((m_i,n_i)\) be the old selected star's actual single
native pair. Its doubled ordinary pair is
\((2m_i-2J,2n_i)\). The product-frame variation and derivative
are consequently
\[
 a=2J-2m_i-2n_i,\qquad
 b_0=2J-2m_i+2n_i,\qquad n_i=(b_0-a)/4.
 \tag{P3K.73.16}
\]
The unrestricted four-weight Cox common-link map, with the
fixed complex-normal reset, has the proved physical action
\[
 (m,n,s)\longmapsto(m+2n,-n,s+n).
 \tag{P3K.73.17}
\]
The constant normal interchange conjugates the two Spin
factors by orientation-preserving constants; it preserves
these third degrees and the physical commutator generator.
The new selected star is C2. Adding the same ordinary
comparison on that side gives
\[
 V'_{22}=a,\qquad\delta'_2=2a-b_0,\qquad S'=S+n_i.
 \tag{P3K.73.18}
\]

\item Its full integral target prisms have
\[
 V'(w_je_j-w_0e_0)=I_je_j-I_0e_0,\qquad
 w_j\delta'_j-w_0\delta'_0=W_j-W_0,\quad j=1,2.
 \tag{P3K.73.19}
\]
With every \(w_j\ne0\), symmetry says the difference of two
variation solutions is \(Crr^t\), where \(r_j=1/w_j\).
The diagonal matrix \(\operatorname{diag}(I_j/w_j)\) is a
particular solution. The derivative equations have the analogous
one-constant solution. Equations (P3K.73.18) determine those
constants to be precisely \(C,E\) of (P3K.73.4).

\item Existence comes from the actual family and products just
constructed. Their integral invariants therefore equal this
unique rational solution. Equation (P3K.73.16) proves the
increment's integrality as its actual single native integer.
Induction along the terminating word proves all the asserted
entries and the physical class. The final absolute-fan identity
carries this family to the original input marking with its
unchanged full source, cap and closing maps.

\par\medskip\noindent\textbf{SA.5. Exact execution.}\par\nopagebreak[4]
The program \path{code/part3_root/scalar_all_fans.py} uses
integer and exact fraction arithmetic. Its family inputs are
the actual primitive profiles OS45, OS47, OS55 and OS56 of
discriminants \(8,14,20,30\), respectively, and the
original shorter coprime families \((K,L)=(1,1),(1,5)\).
For the latter the parameter is exactly
\(w_h=e^{-\pi i c_{\rm AC}/(KL)}\), with the original SE.5
endpoint and its degree-six pullback. Its actual data are
\[
 J=2KL,\qquad\widehat H=KL/3,\qquad\widehat\chi=KL/6,
\]
as follows by the explicit degree-six parameter cover and
the K69 physical reference pair \((KL,4KL)\).
The program checks every determinant \(\pm1\) matrix with
entries of absolute value at most eighteen, compares every
repeated fan, verifies all prism identities and selected inverse
pairs, and runs an independent three-edge traversal. It also
checks the vertical formula for \(-30\le n\le30\) and the
complete free-invariant multiplication on the original actual
degree-six cover. The attached output supplies the counts.
For the OS45 degree-eight example the word is \((0,1)\),
the single native increments are \(18,102\), and the actual
class is \(3+18+102=123\equiv11\pmod{28}\).
There are \(39{,}024\) bounded unimodular executions, \(814\)
distinct fans for each of six actual family data, \(366\) checks
of the vertical formula, and \(210\) complete free-invariant
cover checks. The separate three-edge traversals each visit
\(210\) states through depth seven. Their four transitions
retaining a separate fixed curve are not used in the algorithm;
AF.2 supplies the appropriate actual vertical seed for every
such target.

\item For the original \(K=L=1\) shorter loop the following table
gives an actual basis for each of the twenty-eight residues.
The notation \((a,b;c,d)\) means the matrix with rows
\((a,b)\) and \((c,d)\); \(S\) is its computed integer path
sum. Every determinant is one.
\input{notes/part3_scalar_all_fan_witness_table.tex}
The full words and individual native increments are in the
attached JSON output. A separate program,
\path{code/part3/audit_scalar_all_fan_witnesses.py}, solves the
complete target prism systems afresh by exact linear algebra,
without importing the reconstruction routine. Its eighty-six
independent solves verify all these witnesses. Every element
of the stated finite target group therefore occurs. The
class-three example has word \((0,0)\), seed zero and native
increments \(0,3\). Its mapping torus has the nonproduct
conclusion already proved for nonzero physical classes in
P3F.24, with both spin structures.

\item The finite bounds specify the arithmetic audit. SA.1--4 prove
the algorithm for every unimodular input; the complete explicit
list of target-group witnesses proves the asserted surjectivity.

\end{enumerate}
\end{proof}




\begin{theorem}\label{night:P3K.74}
\StatementRecord{night:P3K.74}{proved}{S6\_06\_BLIND: GPT-6 Astra (OpenAI, ChatGPT Pro mode, effort Pro), recorded re-derivation in the checkpoint dated 18 November 2026; also checked in this writing by \CurrentChecker}{\cref{night:P3F.52}, \cref{night:P3K.31}, \cref{night:P3K.33}, \cref{night:P3K.35}, \cref{night:P3K.36}, \cref{night:P3K.41}, \cref{night:P3K.43}, \cref{night:P3K.44}, \cref{night:P3K.65}, \cref{night:P3K.67}, \cref{night:P3K.68}, \cref{night:P3K.69}, \cref{night:P3K.70}, \cref{night:P3K.73}}
The numerical all-fan assertions in the historical compound
statements P3K.33, P3F.53, P3K.35--37, P3K.41, P3K.43 and
P3K.44 are false. In the physical convention proved in P3K.68,
the following actual data give counterexamples; every value
in the last two columns is in \(\mathbb Z/28\).
\begin{center}\small
\begin{tabular}{llrr}
\toprule
Claim & Specified datum and fan & Earlier value & Actual value\\
\midrule
K33 & Original \(K=L=1\), shorter loop, \(T_2\)&0&26\\
F53 & Primitive OS45, \(\lambda=r\), \(T_0\)&25&3\\
K35 & Same OS45 datum, \(b=1\), \(T_0\)&25&3\\
K36 & Corrected Liu--Rollenske \(u=w^2\), \(T_2\)&0&26\\
K37 & OS49, \(k=1\), shorter loop, \(\mathcal T_-\)&0&27\\
K41 & Same OS45 cap data, \(N=1\), \(T_0\)&25&3\\
K43 & Same OS45 datum, constant site path, \(T_0\)&25&3\\
K44 & Same OS49 data, \(k=N=1\), \(\mathcal T_-\)&0&27\\
\bottomrule
\end{tabular}
\end{center}
The actual source, line, finite-cap and endpoint constructions
identified in those proofs retain their stated constructive
scope. For the whole-zero \(I_1\) families their complete
numerical replacement is the algorithm of P3K.73. For the
original shorter family its fan image is all of \(\Theta_7\),
with the twenty-eight explicit witnesses there.

The invariant-divisor family isomorphism in P3F.52 preserves
the actual return computed by this algorithm. For the marked
fixed-source cap-data and moving ordinary-site families of
K41 and K43, a scalar winding \(N M/b\) has return \(N A(P)\),
where \(A(P)\) is the actual shorter return of K73 on the
transported fan. Their fixed-scalar loops have return zero.
For the fixed-chart OS49 cap-data scope of K44, the actual
return is \(N\) times the specified shorter OS49 return;
K65 supplies \(-Nk\) on \(\mathcal T_-\) and zero on
\(\mathcal T_+\). This last statement does not assign values
on the other OS49 fans or on additional ordinary clutches.
\end{theorem}
\begin{proof}
\begin{enumerate}[label=\arabic*.]
\item \par\medskip\noindent\textbf{The complete actual counterexamples.}\par\nopagebreak[4]
Use the exact OS45 source equation, section \(G\), Cartier
line \(\mathcal O(8G-O)\), suspension \(-G\), selected simple
zero and ordinary cap of NV.1 in K70. Its scalar is the
primitive \(\lambda=r\), and \((J,H,\chi)=(1,8,1)\).
K69 proves its physical reference value \(3\) by the complete
physical/dual line and gluing morphism. The old reference
expression is \(\chi-H/2=-3=25\). Taking the actual shorter
degree \(b=1\) gives the identical family in K35. Taking
\(N=1\) gives the same family in K41, and taking each ordinary
site path constant gives the permitted closed family in K43.
Thus all four rows concern the fully specified source and
closing, with no change of parameter orientation.

\item For K33 use the original source, cap orders \(3,4\), either
admissible norm sign, and the literal SE.5 endpoint over
\(w=e^{-\pi i c_{\rm AC}}\). The exact shorter data are
\(J=2,\widehat H=1/3,\widehat\chi=1/6\). The physical
vertical-edge fan \(T_2\) therefore has actual value
\(0-2\cdot2\cdot1/2=-2=26\), by the actual fan maps of
K67 or K73. On every \(T_n\), the exact matrix equality
\(T_n^{-1}e_2=e_2\) gives the old cocharacter norm
\(S=0^2-0\cdot1+1^2=1\). Its displayed polynomial is zero.

\item For K36 retain the precise source and scalar \(u\) of its
LR.1--11 construction. Its good \(III^*\) root has order four,
the source pair is \((a,b)=(1,2)\), and the actual closed
endpoint is \(u=w^2\), with \(w=r^2\), \(u=r^4\).
The degree-four page has \((H,\chi)=(2,1)\), so the actual
shorter data are \(J=2,\widehat H=1,\widehat\chi=1/2\).
The same \(T_2\) computation gives \(26\), whereas its
old \(3(S-1)(S-3)/8\) is zero. The separately proved
one-\(u\)-turn endpoint obstruction and the reference-fan
zero remain exactly their earlier statements.

\item For the two OS49 rows use K65's explicit cubic, section
\((t^2,t^2)\), line \(\mathcal O(G-O)\), translation
\(-3G\), good cap of order three, norm \(\pm1\), and
fixed delta character. Its actual shorter scalar is
\(\lambda=w\). The primitive period change in OR65.2
gives the negative fan physical cocharacter \((0,1)\), so
again \(S=1\). Its actual directed C2 flop has native pair
\((-1,1)\) and goes from this negative fan to the proved
zero positive fan. The actual central increment is therefore
\(+1\), giving negative-fan class \(-1=27\). The
fixed-chart cap-data case with \(N=1\) is this same family.

\item The exact program
\path{code/part3_root/refute_remaining_scalar_polynomials.py}
executes all eight comparisons. The old expressions have a
verified unequal value on an actual family, so the assertions
are refuted as compound assertions. The retained historical
texts identify precisely the erroneous steps: the K31
common-product numerical recurrence was applied to the false
scalar initial variation, and the nonzero physical reference
values also used the dual-line sign corrected by K68.

\par\medskip\noindent\textbf{The invariant-divisor morphism.}\par\nopagebreak[4]
F52, V1--3, proves the actual divisor equality \(t_Q^*D=D\)
and constructs multiplication by its canonical section
\[
 \mathcal M_0^\times\longrightarrow\mathcal M^\times,
 \qquad v\longmapsto v s_D.
\]
It conjugates the suspension after the specified constant
scalar reparametrization. At a whole-fibre coefficient \(kF\)
the actual character map is
\((Z,V,q)\mapsto(Z,q^kV,q)\), with cone-height map
\((v,h)\mapsto(v+kh\delta,h)\). It translates the
height-one fan by the actual integral period \(k\delta\).
Every cone and face localization is carried by the displayed
integral lattice map, including the full unit transitions.
At each cap the proper direct-image evaluation map gives
VL.7; its base monomial, finite character and integral annular
winding are retained. These are exact holomorphic family
isomorphisms, and their source and line maps commute with
the physical circle.

\item Transport a connection by this isomorphism. Its holonomy,
followed by its specified inverse endpoint, is conjugated
literally. The same isomorphism carries the physical line,
the stable tangent bundle and both marked caps, so it
preserves their characteristic pair and the native relative
invariants in the transported products. Hence it preserves
the actual value of K73. This proves the covariance assertion
of F52 without its former appeal to the false K31 formula.

\par\medskip\noindent\textbf{The retained cap-data and site maps.}\par\nopagebreak[4]
Use the actual local root comparisons CP.1--5 of K41 and
the whole-disc moving-site maps of K43. They carry the
source, Cartier line, finite character, cone-height map and
ordinary annular cochain at each parameter. The exact
endpoint calculation CP.6 gives scalar winding \(NM/b\)
and compares its endpoint to the \(N\)-th iterate of the
SE.4--14 shorter endpoint. Their difference is the global
extending physical \(\mathbb C^*\) action; CP.5 supplies
its path to the identity. The moving-site maps retain the
whole cap data under every allowed site permutation, so
the same endpoint comparison holds for their closed braids.

\item Here is the actual return-map equation. Write
\(F_t:X_t\to Y_t\) for those fibre isomorphisms and
\(E_X:X_0\to X_1\), \(E_Y:Y_0\to Y_1\) for the
forward endpoints. Transport the connection by \(F_t\).
For its parallel transports \(\psi_X,\psi_Y\), we have
\(F_1\psi_X=\psi_YF_0\). Define
\[
 A=E_Y^{-1}F_1E_XF_0^{-1}:Y_0\longrightarrow Y_0.
\]
This is the extending physical-circle-complexification
action just identified, with its explicit path to the
identity. The physical returns consequently satisfy
\[
 F_0\Phi_XF_0^{-1}=A^{-1}\Phi_Y,
 \qquad \Phi_X=E_X^{-1}\psi_X,\quad
 \Phi_Y=E_Y^{-1}\psi_Y.
\]
Their isotopy classes agree. The target is the specified
shorter family iterated \(N\) times, so its class is
\(N A(P)\). The scalar-zero-winding target is constant,
giving zero. The exact component-line and cap arrows of
K44 give this same typed equation for its fixed-chart
OS49 scope. Substitution of K65 gives the two asserted
OS49 values. The equation retains the inverse endpoint
and every prescribed circle and orientation marking.

\end{enumerate}
\end{proof}


\section{The completed source and two-cap return sectors}



\begin{theorem}\label{night:P3K.76}
\StatementRecord{night:P3K.76}{proved}{S6\_06\_BLIND: GPT-6 Astra (OpenAI, ChatGPT Pro mode, effort Pro), recorded re-derivation in the checkpoint dated 18 November 2026}{\cref{night:P3F.9}, \cref{night:P3K.62}, \cref{night:P3K.65}, \cref{night:P3K.68}, \cref{night:P3K.70}, \cref{night:P3K.73}}
Retain the actual OS49 source
\[
y^2=x^3-3t^3(t-1)x+2t^4(t-1)(t-\tfrac12),\qquad G=(t^2,t^2),
\]
the line \(M=\mathcal O(G-O)\), suspension over \(-3G\), and
its target-point linearization divisor \(C_0\) on the \(I_2\) fibre.
Retain the \(IV^*\) cap of order \(m=3k>0\), either norm
\(\nu=\pm1\), every fixed character \(d\), and the full root,
type-II and annular arrows of P3K.65. In the physical character
lattice put \(T_n=T(e_1+ne_2,e_2)\). The negative fan is \(T_0\)
and the positive fan is \(T_1\), by the literal cone maps below.

For the positive canonical \(r\)-loop with \(\lambda=r^m\), set
\(j=m\). For the actual shorter positive \(w=r^3\)-loop with
\(\lambda=w^k\), set \(j=k\). On every \(T_n\) the complete
physical integer pair, index and native fixed-star free pair are
\[
\begin{aligned}
k_{M,n}&=2j-6j(n-1)^2,&\ell_{M,n}&=8j,\\
A_{M,n}&=-j(n-1)^2,&
(\mu_n,\rho_n)&=(2j,-6j(n-1)^2).
\end{aligned}
\]
The dual line \(U=M^*\) has the negative characteristic pair.
The actual oriented clutched-sphere class is
\(-j(n-1)^2\pmod {28}\) in the physical convention of P3K.68.
Formula (BW.30) gives the complete variation and stable-derivative
seed on the literal double parameter cover. The proof retains
both norm signs, every cap character, all full unit factors and
the actual shorter endpoint. It establishes these vertical-edge
fans and does not use the withdrawn universal fan polynomial.
\end{theorem}
\begin{proof}
\begin{enumerate}[label=\arabic*.]
\item \par

\medskip\noindent\textbf{BW.1. The actual line tensor, including its outer gauge}

\par\nobreak

\item Use the closed coordinates of the actual torus seam. A source-independent
bounded rescaling \((A,D)\) acts by the two literal vertex maps (SI.4)
and intertwines every period row. The endpoint equality \(C_1E=C_0\)
is retained on full collars. In these coordinates write
\(\alpha=e^{2\pi iL\theta}\alpha_*\) and
\(\eta=e^{2\pi iJ\theta}\eta_*\), with the full zero-winding factors
and their logarithms. The physical-circle quotient retains their
positive norms and the source factors. Its F9 source/graph construction
is therefore the same smooth piece \(C\) for every \(L,J\); the phase
changes remain on its actual circle line. The normal and cap maps
are exactly (SI.11)--(SI.13), with the physical height unchanged.

\item Let \(b,\gamma:C\to\mathbb R/\mathbb Z\) be the two graph maps with
outer values \(v,\gamma\). Each has a real lift on the fixed sphere
and can be made constant on the fixed surgery tube by its supported
real homotopy. The exact theta-cut gauge implementing this change
is its exponential times theta, as in SI.3. They then extend over the
same \(D^3\times S^2\) cap. Define two actual lines on \(\mathcal C\)
by theta-cut coefficient gluings
\[
 T_b:\ w\mapsto e^{2\pi ib}w,\qquad
 T_\gamma:\ w\mapsto e^{2\pi i\gamma}w.
\]
Their real curvatures are
\[
 \Xi=d\theta\wedge db,\quad \Psi=d\theta\wedge d\gamma,
 \qquad \Xi^2=\Psi^2=\Xi\Psi=0.                    \tag{BW.1}
\]
Use the meridian-one line \(U=M^*\), with the K68 physical orientation.
The actual period coefficients have theta deck
\[
 (u,d)\longmapsto(u,d-Lv-J\gamma).
                                                        \tag{BW.2}
\]
Hence the full boundary connection on \(U\) is
\[
 A_{L,J}=dd+(\alpha_q+J\theta)d\gamma+L\theta\,dv.
                                                        \tag{BW.3}
\]
Differentiating the actual complex period expression (SI.14) retains
all unit phases and their derivatives. Equation (BW.3) is written in
those full coefficient coordinates, not in a constant-factor
approximation to the physical characters.

\item The exact coefficient theta clutch is the tensor product of those
for \(T_b^L\) and \(T_\gamma^J\). The meridian and regular \(b\times\gamma\)
evaluations remain respectively one and zero. The integral Thom,
Mayer--Vietoris and Wang computation of SI.3 leaves precisely these
four evaluations. It therefore gives
\[
 U_{L,J}\simeq U_{0,1}\otimes T_b^L\otimes T_\gamma^{J-1}.
                                                        \tag{BW.4}
\]
Here the symbol \(U_{0,1}\) means the actual line obtained from this
piece by the displayed inverse theta-cut tensors. It is defined
before making a comparison to any reference family.

\item The isomorphism in (BW.4) carries its prescribed outer arrow, rather
than only its absolute Chern class. In the doubly cut coefficient
frames that arrow has the literal connection difference
\(L\theta\,dv+(J-1)\theta\,d\gamma\). The only nonextending boundary
winding in \(H^1(\mathcal B)/\operatorname{im}H^1(\mathcal C)\) is the
\(q\) angle. Choose one theta value and the free section disc
\(|Z_*|=1,V=1,\gamma=v=0\) used in (SI.23). On this disc both tensor
lines have their constant frame and the prescribed arrow is the
identity. It has no \(q\)-angle winding. The integral-logarithm and
partition-of-unity Cech construction of (SI.24), followed by the
extending theta, \(b\), gamma gauges and the remaining real boundary
logarithm, makes (BW.4) exactly that outer arrow. Its inner cap gauge
extends over \(D^3\times S^2\times S^1\). Thus (BW.4) is a relative
line morphism with its full boundary and cap marking.

\par

\medskip\noindent\textbf{BW.2. The relative local pairing is calibrated by an actual map}

\par\nobreak

\item We spell out the local reference comparison needed here. All four
zero-winding unit factors are nonzero on a protected \(q\) disc. Choose
their actual logarithms. Between two specified sets of these factors
use
\[
 u_s=u_0\exp\bigl(s\log(u_1/u_0)\bigr),\qquad0\le s\le1.
                                                        \tag{BW.5}
\]
The ratios have the displayed holomorphic logarithms on a common
smaller disc. This is an explicit path of the full commuting torus
periods, not just a path of their leading coefficients. Their
\(q\)-valuation matrix remains the identity. Their logarithmic norms
are uniformly bounded on the compact path, so the same small disc
and the same finite fundamental regions give the toric quotients.
The period translations act freely on the labels of the fixed
unimodular charts; on the regular region the invertible real period
matrix gives the quotient. The resulting compact pieces with
boundary are a smooth family over \(s\).
At the calibration endpoint take \(\kappa,\alpha_*,g,\eta_*\)
to be their fixed-parameter values, hence functions of \(q\) only.
The zero-winding logarithms in (BW.5) give the displayed path to
these values; the endpoint normal matrices below are therefore
literally independent of the circle parameter.

\item The integer closing exponent \(e\) is fixed throughout this path.
For the full bounded closing unit \(u_s\) the actual bounded endpoint
and cut multiplier obey
\[
 E_s=(g_s^{-e},u_s\eta_s^{-e}),\qquad
 (C_s)_1=(g_s^e,u_s^{-1}\eta_s^e),\qquad
 (C_s)_1E_s=(C_s)_0=1.
                                                        \tag{BW.5a}
\]
Use the same endpoint-flat cut function on every \(s\) chart. These
identities then hold on the entire collars with all jets. No
discrete annular exponent, finite root character or exterior cap
datum is varied. Only these compact local pieces and their
prescribed outer arrows occur in the comparison.

\item For the source parameter \(Q_s=q\kappa_s\), the exact node horizontal
field is \(\partial_s+(\partial_s\log\kappa_s)y\partial_y\) in
\(xy=Q_s\). Prescribe the identity on the real source-period
coefficients on the entire outer collar. Interpolate the two
horizontal fields by a partition in their affine space of vertical
differences, as in (SI.10). For the suspension source factor \(g_s\)
use its full \(\exp(\gamma\log g_s)\) untwisting. Positive line norms
use the F9 proper coordinate, with its two fixed endpoint lifts
\(1/2+\varepsilon,1/2-\varepsilon\). Every coordinate comparison is
relative to the specified outer real coefficients. These are actual
flows and hence diffeomorphisms of the compact pieces.

\item The normalization graph isotopy stays in that same short arc with
distinct endpoint heights. Its velocity extends over the fixed tube.
The full unit normal factors in the two toric charts give its
Cartier/Hopf normal arrow; since the unit logarithms in (BW.5) extend,
their parameter-\(s\) homotopies carry the entire normal jets and the
actual physical circle. The boundary sphere velocity extends over
the three-ball, and the adjoint normal arrow extends over the Hopf
sphere cap. This is the exact extension used in (SI.12)--(SI.13).

\item Over this path the principal circle lines, including their Hopf caps,
form an actual line bundle on the smooth family. Choose an invariant
horizontal lift of the just constructed base field. On the outside
collar prescribe the lift which fixes the full real coefficient \(d\).
Near the cap use the carried Cartier/Hopf lift. On their overlap the
difference is vertical, so an invariant partition of unity glues
these lifts. Its flow gives an actual line-bundle isomorphism. It
is the prescribed real-period isomorphism on the whole boundary;
there the connections \(dd+\alpha_qd\gamma+\theta d\gamma\) agree
literally. The \(q\)-disc gauge above also fixes the integral boundary
normalization. This constructs the relative pair morphism, including
its line and cap, between \(U_{0,1}\) and the corresponding local
scalar-winding-one piece with the same integer \(e\) and prescribed
outer coefficients. More directly, inverse theta-cut tensors in
(BW.4) replace the local line multipliers by
\(\alpha_*\) and \(e^{2\pi i\theta}\eta_*\), without changing their
positive norms or source factors. Identity on the circle quotient,
the displayed coefficient tensors, and (BW.5a) are its literal
local arrow. It does not alter or posit a map between unspecified
global exteriors.

\item Let \(F_{0,1}\) be a connection curvature with this fixed boundary
value. The relative pairing
\[
 B_{\rm loc}=\int_{\mathcal C}F_{0,1}^{2}\wedge\Xi
                                                        \tag{BW.6}
\]
is preserved by the constructed morphism. Connections with the same
boundary value give the same number: their curvature difference is
exact with zero boundary primitive, and Stokes applies to the closed
form \(\Xi\).

\item Apply the actual local SCN.9/K62 \(C_2\) bordism to this winding-one
piece, with the same \(e\) and every factor in (BW.5a). Its first
outer normal arrow is the identity and the second is the base
factor \(\kappa^{-1}\) with normal matrix
\(\operatorname{diag}(\alpha_*^{-1},\kappa\alpha_*)\).
Both are parameter-independent. The common inner physical-circle
gauge therefore extends over the retained joint cap and these
outer arrows. The resulting line and normal transitions pull back
from the three-dimensional fixed-cobordism factor; hence
\(b^2=p_1(\nu)=p_1(T\bar X)=0\) there. The two separate sphere-bundle
identities for \(4\int x^3=3b^2+p_1(\nu)\) and
\(\int p_1x=p_1(T\bar X)+p_1(\nu)\) give zero individually.
The common exterior collar is an actual product with its equal
connection. It contributes zero to their difference, without
selecting or changing a global exterior. This proves equality of
the two relative integer Euler cubes for the local pieces.
The SI whole-collar calculation at \(n=1\) is exactly
\(3B_{\rm loc}-3\). Thus
\[
                         B_{\rm loc}=1.               \tag{BW.7}
\]
This is a calibration by one actual integer comparison and the
explicit relative morphism above. It does not select an integer
from a central residue.

\item By (BW.4) we may choose on the original piece the curvature
\(F_{L,J}=F_{0,1}+L\Xi+(J-1)\Psi\), with its exact boundary value.
Since the last two summands both contain \(d\theta\), (BW.1) gives
\[
 \int_{\mathcal C}F_{L,J}^{2}\wedge\Xi
       =\int_{\mathcal C}F_{0,1}^{2}\wedge\Xi=1.
                                                        \tag{BW.8}
\]

\par

\medskip\noindent\textbf{BW.3. The full shear and its two characteristic integrals}

\par\nobreak

\item For the actual shear \(V_n=Z^{-n}V\), \(h_n=hB^n\), the source parameter
remains \(Q_n=\kappa qg^n\), while the line windings are \((L+nJ,J)\).
The exact coefficient map is still \(u_n=u,v_n=v,\gamma_n=\gamma-nv,
d_n=d-nu\). Its new connection is
\[
 A_n=dd_n+(\alpha_q+J\theta)d\gamma_n+(L+nJ)\theta\,dv.
\]
Thus the actual piece-minus-exterior form is still
\[
 \beta_n=-n(du+\alpha_qdv),\qquad d\beta_n=-n\,d\alpha_q\wedge dv.
                                                        \tag{BW.9}
\]
The same section-disc/Cech construction makes the interior line map
\(U_n\to U_0\otimes T_b^{nJ}\) relative to this exact boundary arrow.
The source and graph maps (SI.10)--(SI.13) use the same full units;
their tangent boundary matrix is \(I+n\phi_gE_{12}-nE_{32}\).
Its Maurer form squares to zero, so the entire tangent collar has
\(\operatorname{tr}F^2=0\). The actual regular-level argument (SI.27)
gives \(\langle p_1,\Xi\rangle=0\). Therefore
\[
                         \ell_{M,n}=\ell_{M,0}.        \tag{BW.10}
\]

\item Use the piece coefficient gamma\(=\gamma_n\) and the positive boundary
orientation \(\Omega=d\alpha_q\wedge du\wedge dv\wedge d\gamma\wedge d\theta\).
The retained exterior curvature is now
\[
 F_{{\rm ext},n}=(d\alpha_q+Jd\theta)\wedge(d\gamma+n\,dv)
                                      +L\,d\theta\wedge dv.
                                                        \tag{BW.11}
\]
Its square contributes a linear term; it must not be set to zero.
Exact exterior algebra gives
\[
 3\int_0^1\beta_n\wedge
             (F_{{\rm ext},n}+s\,d\beta_n)^2\,ds
                  =(6Ln+3Jn^2)\Omega.                \tag{BW.12}
\]
The outward collar runs from \(s=1\) to \(s=0\), so its physical Euler
cube contribution is \(-6Ln-3Jn^2\). The exact interior relative line
map and (BW.8) give interior contribution \(3nJ\). The K68 physical
line and Jin orientation are the ones in (SI.29). Combining the
two actual pieces therefore gives
\[
 \begin{aligned}
 k_{M,n}&=k_{M,0}-6Ln-3Jn(n-1),\\
 \ell_{M,n}&=\ell_{M,0},\\
 A_{M,n}&=A_{M,0}-Ln-\tfrac J2n(n-1).
 \end{aligned}                                      \tag{BW.13}
\]
Every boundary primitive has been included. In particular the
linear \(L\) term comes from the actual \(F_{{\rm ext},n}^2\) term and
is fixed as an integer by the relative morphism; no modulo-28
comparison is used to choose it.

\item The domain of this local calculation is the explicitly constructed
primitive nodal quotient piece and its stated closed seam. Its use
for a named global family retains that family's identification of
this piece, physical line, cap and parameter orientation. The
ordinary and finite-cap source data have not been replaced by
an abstract winding matrix in this proof.

\par

\medskip\noindent\textbf{BW.4. The actual OS49 positive reference is the fan T1}

\par\nobreak

\item Now retain the actual source and every datum of OR65.1:
\[
 y^2=x^3-3t^3(t-1)x+2t^4(t-1)(t-\tfrac12),\qquad
 G=(t^2,t^2),\quad M=\mathcal O(G-O),\quad Q=3G.
                                                        \tag{BW.14}
\]
The suspension is over \(-Q\), with its actual target-point
linearization divisor \(C_0\) on the \(I_2\) source fibre. The
\(IV^*\) cap has order \(m=3k>0\), norm sign \(\nu=\pm1\) and
fixed character \(d\in\mathbb Z\). Its exact ordinary-to-cap
arrow is
\[
 \Phi_r(z,s_0)=z+\frac{\log s_0}{2\pi i}\Pi_r(s_0)
                 \{\nu(e_\gamma-e_u)+d\delta\}.
                                                        \tag{BW.15}
\]
The root action is \(A_r=r^\nu A_0\). The ordinary closing is
the endpoint-to-initial map \(q^{-\nu}\); its forward is \(q^\nu\).
On the type-II source neighbourhood retain the original line
complement and contracting suspension. None of these arrows
is replaced by a whole-fibre linearization.

\item The original source and suspension logarithmic periods are
\[
 u=(2s,s)^t+U(q),\quad
 \gamma=(-3s,-s+mv)^t+C(q),\quad
 s=\frac{\log q}{2\pi i},\quad v=\frac{\log r}{2\pi i}.
\]
The actual integral deck-generator change is
\[
 H=-u-\gamma=(s,-mv)^t-U-C,\qquad
 K=3u+2\gamma=(0,s+2mv)^t+D,\quad D=3U+2C.
                                                        \tag{BW.16}
\]
Its determinant is one; these are actual deck transformations,
including their bounded source and line units. Thus in the
original physical characters the periods are exactly those of
BW.1, with
\[
                   (L,J)=(-m,2m).                    \tag{BW.17}
\]
The bounded forward is \(E_r=(g^{-\nu},\eta^{-\nu})\);
the complete seam of BW.5a has \(e=\nu,u=1\). Its actual cap
square is (BW.15). Hence both signs use the same positive
parameter orientation in (BW.17).

\item Write \(T_n=T(e_1+ne_2,e_2)\) in this fixed physical lattice.
The infinite negative-diagonal fan is \(T_0\). The two triangles
of \(T_1\) are
\[
 \{0,e_1+e_2,e_2\},\qquad
 \{e_1+e_2,e_2,e_1+2e_2\}.
\]
The first is a positive-diagonal square triangle; the second
is its other triangle translated by \(e_2\). Translation by
\(e_2\) is the full actual \(K\) deck, with the unit factors in
(BW.16). These equalities hold for every integral translate,
so the two infinite fans and all their localized cone rings
are literally the same. The dense-character identity descends
through the unchanged period action and finite-cap arrows.

\item The fixed-star charts on \(T_n\) use
\(Z_n=Z,V_n=Z^{-n}V\) and
\[
 (Z,n_0,n_1)=(Z,q/(ZV_n),V_n),\qquad
 (Z',n'_0,n'_1)=(Z^{-1},Zn_0,Zn_1).
                                                        \tag{BW.18}
\]
For \(n=1\) they become exactly
\[
 (Z,A,B)=(Z,q/V,V/Z),\qquad
 (Z',A',B')=(Z^{-1},qZ/V,V),
                                                        \tag{BW.19}
\]
the OC.11 positive-fan normal charts. In particular the
actual global quaternion is precisely
\[
 q_H=(1+jZ)(A+j\overline B)
     =(Z^{-1}+j)(A'+j\overline {B'}).
                                                        \tag{BW.20}
\]
The physical circle is right multiplication by \(e^{-i\phi}\).
The decreasing contracting-\(K\) height, the constant signed
normal reference change for \(\nu=-1\), and the full unit
normal arrows are OC.12--19. Thus this identification retains
the actual Cox cap and its framing, not just the common fan.

\item The independent positive page is the actual line-complement
subfan with heights \(\lceil j/2\rceil\) and component degrees
\((-1,+1)\). Its \(m\) norm shells, cap root and regular bands
are OR65.2. The full normal quotient clutch in OC.17--20 has
zero normal Pontryagin integral; OC.21 retains
\(M_{\rm rel}^2=-2k\) and \(p_{\rm rel}=-24k\). K68's exact
physical/dual line arrow converts its pair to
\[
                      (k_{M,1},\ell_{M,1})=(2m,8m).
                                                        \tag{BW.21}
\]
The constant \(d\)-character is carried on every root sheet;
its \(m\)-fold product is one. The relative page line changes
only by its already retained base class, whose square and
pairing with \(M\) vanish. The point-angle determinant-one
map in OC.7--10 proves the same signed pair for each \(\nu\).

\item Apply the actual relative calculation (BW.13) in these charts
and subtract its value at \(n=1\). With (BW.17) and (BW.21)
this gives the complete physical pair and index
\[
 \boxed{\quad k_{M,n}=2m-6m(n-1)^2,\qquad
 \ell_{M,n}=8m,\qquad A_{M,n}=-m(n-1)^2.\quad}
                                                        \tag{BW.22}
\]
The dual-\(U\) pair is the negative of this pair. In particular
the negative fan's integer is now computed by the actual line
and tangent morphisms. Its value was not selected from K65's
residue. The native fixed-star free pair is
\[
               (\mu_n,\rho_n)=(2m,-6m(n-1)^2).       \tag{BW.23}
\]

\par

\medskip\noindent\textbf{BW.5. The actual shorter endpoint and the full vertical seed}

\par\nobreak

\item Retain the exact endpoint of OR65.6. The rational source-line
lift \(U\) satisfies \(\rho=B_0^{-1}/U^3\); its orders at the
cusp and finite cap are respectively \(1\) and \(-1\) modulo
three. The rational root \(g_0^3=q^{-\nu}\rho^{-2\nu}\) gives
the endpoint-to-initial map
\[
 C_r=r^{2\nu k}g_0U^{-2\nu}:X_{\zeta_3r}\longrightarrow X_r.
                                                        \tag{BW.24}
\]
On the cap the full Cartier map is the invariant unit times
\(L^{-2\nu}\), with \(DLD^{-1}=\zeta_3L\) and \(L^3=B_0^{-1}\).
Consequently \(\Phi_rC_r=C_{\rm cap}\Phi_{\zeta_3r}\) and
\(C_{\rm cap}A_{\zeta_3r}C_{\rm cap}^{-1}=A_r\). This retains
the original \(d\)-character and both values of \(\nu\).

\item After the actual source deck reindexing, the full bounded
closing logarithm and extending origin are
\[
 c(v)=\frac{\nu}{3}(D+2mv\delta)-\frac{2\nu k}{3}\delta,
 \qquad \sigma(v)=\nu vD+\nu m v(v-1)\delta,
\]
\[
 \sigma(v+1/3)-\sigma(v)=c(v),\qquad
 \sum_{j=0}^2c(v+j/3)=\nu D+2\nu mv\delta.           \tag{BW.25}
\]
The forward endpoint is the inverse of this closing. On
\(T_n\) apply the literal character map \((z,\ell)\mapsto
(z,\ell-nz)\) to every vector in (BW.25). This is an
integral monomial map and preserves their equalities, all
bounded units and the zero-height endpoint. The sheet product
\(P_j=\exp(2\pi i\sum_{h<j}c(v+h/3))\) is the actual product
of the first \(j\) closing maps. Its values at \(j=0,1,2,3\) are
those of OR65.31, with this same character map applied.

\item For completeness use the fixed native C0 coordinate order
\((b_1,a,c_1)\) from K73. A cut multiplier \((A,D_n)\), with
\(D_n=DA^{-n}\), acts by the factors
\[
 (A,(AD_n)^{-1},D_n),\qquad
 (A^{-1},D_n^{-1},AD_n)
                                                        \tag{BW.26}
\]
on its two normal charts. Its literal Cox lift is
\[
 (a_0,a_1,c_0,c_1)\longmapsto
 (a_0,Aa_1,(AD_n)^{-1}c_0,D_n c_1).                  \tag{BW.27}
\]
It preserves \(q\) and the physical circle. The all-sheet
identity in (BW.25) holds on each of these four character
factors. Thus the covered \(r\) native product is the literal
triple pullback of the shorter \(w=r^3\) native product,
including every collar and the cap square (BW.24).

\item We also compare that product to the canonical seam used in
BW.4. On a linear interior cut the origin product has
logarithm \(\sigma(v+\theta)-\sigma(v)\), whereas the
canonical cut has logarithm
\(\nu\theta D+2\nu m\theta(v+\theta)\delta\).
Their difference is the explicit physical-circle factor
\[
                 -\nu m(\theta^2+\theta)\delta.
                                                        \tag{BW.28}
\]
Its exponential has value one at both endpoints. It is
therefore an actual based global physical-circle action,
including its possible winding. It acts identically on the
orbit normal three-plane and extends on the same line and
Hopf cap. Flatten both cuts on the protected collars using
their common actual endpoint relation. The remaining unit
ratios have zero endpoint logarithm and the explicit
exponential interpolation carries their full jets. This
gives the actual relative cap and line comparison between
the two covered products; it does not discard a closed
circle loop. Hence their completed physical-circle lines
and stable tangent bundles are isomorphic with their
specified orientation and outer frame.

\item The degree-three mapping-torus covering now gives literal
equalities of characteristic integrals, not division in
\(\Theta_7\). Writing \(j=m/3=k\) on the shorter loop gives
\[
 (k_{M,n}^{w},\ell_{M,n}^{w})
       =(2k-6k(n-1)^2,8k),\qquad
 A_{M,n}^{w}=-k(n-1)^2.                              \tag{BW.29}
\]
The sphere classes are the residues of (BW.22) and (BW.29)
by the actual physical Hopf transfer K68.

\item Here is also the full seed for further exact fan comparisons.
Write \(j=m\) for the \(r\) loop and \(j=k\) for the actual shorter
\(w\) loop. Use the literal doubled ordinary product, carrying
the native fixed-star product just specified. Its native
double pair is \((4j,-12j(n-1)^2)\). The K70/NV actual C0
product change is pure left \(-2J=-4j\), so its ordinary
double pair is \((0,-12j(n-1)^2)\). Consequently
\[
 V^{(2)}=\begin{pmatrix}
 12j(n-1)^2&8j(1-n)&8j(n-1)\\
 8j(1-n)&4j&0\\8j(n-1)&0&4j
 \end{pmatrix},\qquad
 \delta^{(2)}=\begin{pmatrix}
 -12j(n-1)^2\\8j(n-1)\\8j(1-n)
 \end{pmatrix}.                                    \tag{BW.30}
\]
The off-diagonal and remaining diagonal entries come from
the full actual trace, not the index. Indeed the doubled
period winding in these characters is
\(B=\left(\begin{smallmatrix}0&0\\2j(2n-1)&4j\end{smallmatrix}\right)\),
\(t_0=2j\), \(w=(0,1,-1)\), \(d=(-2,-1,-1)\). The literal
normal and derivative prisms of OR65.16--17 give
\(I=(8j(n-1),4j,-4j)\) and \(W=(8j(1-n),0,0)\), which
determine every remaining entry of (BW.30).

\item This result computes every vertical-edge fan with the exact
OS49 source, component line, cap and endpoint just retained.
It does not substitute these numbers into a previously
withdrawn universal scalar-fan expression.

\par\medskip The exact certificate is \path{code/part3/bottom_row_winding.py}, with its saved JSON and captured output. It checks the complete period, line, seam, tangent, collar and named OS49 identities.

\end{enumerate}
\end{proof}




\begin{theorem}\label{night:P3K.77}
\StatementRecord{night:P3K.77}{proved}{S6\_06\_BLIND: GPT-6 Astra (OpenAI, ChatGPT Pro mode, effort Pro), recorded re-derivation in the checkpoint dated 18 November 2026}{\cref{night:P3F.35}, \cref{night:P3K.44}, \cref{night:P3K.55}, \cref{night:P3K.62}, \cref{night:P3K.64}, \cref{night:P3K.71}}
Retain the exact OS49 source, marked component line, suspension over
\(-3G\), \(IV^*\) cap, type-II chart, norm sign \(\nu=\pm1\),
character \(d\), and ordinary closing of P3K.76. In particular
\[
 y^2=x^3-3t^3(t-1)x+2t^4(t-1)(t-\tfrac12),\qquad
 G=(t^2,t^2),\quad M=\mathcal O(G-O),\quad m=3k>0.
\]
Set \(j=m\) on its positive canonical \(r\)-loop and \(j=k\)
on its constructed positive \(w=r^3\)-loop with the exact OR65
endpoint. For every \(P\in GL_2(\mathbb Z)\), form the identical-fan
representative and the terminating word of P3K.71, AF.1--2.
Its vertical seed is \(T_n\), and its forward word consists of
\[
 Q_0=\begin{pmatrix}1&-1\\0&1\end{pmatrix},\qquad
 Q_1=\begin{pmatrix}1&0\\-1&1\end{pmatrix}.
\]
Start with the complete doubled ordinary invariants
\[
 V=\begin{pmatrix}
 12j(n-1)^2&8j(1-n)&8j(n-1)\\
 8j(1-n)&4j&0\\8j(n-1)&0&4j
 \end{pmatrix},\quad
 \delta=\begin{pmatrix}-12j(n-1)^2\\8j(n-1)\\8j(1-n)\end{pmatrix},
\]
\[
                         S=-j(n-1)^2.                 \tag{OF.1}
\]
For a step \(P\mapsto P'=PQ_i\), retain \(a=V_{ii}\) and
\(b=\delta_i\), and compute at the target
\[
 B=(b_{hk})=2j(P')^{-1}
           \begin{pmatrix}0&0\\-1&2\end{pmatrix}P',\qquad
 t_0=2j,\qquad
 l=(b_{22}-b_{21},b_{11}-b_{12},b_{22}),
\]
\[
 (x,y)^t=(P')^{-1}e_2,\quad
 w=(x,y,-x-y),\quad d=(-x-2y,-2x-y,x-y),
\]
\[
 I_h=t_0w_h+d_h(l_h-t_0),\qquad
 W_h=t_0w_h-d_h(l_h-t_0).                             \tag{OF.2}
\]
All three target weights are nonzero. The next complete invariants
and single-loop path sum are
\[
 C=w_2^2a-w_2I_2,\qquad E=w_2(2a-b)-W_2,
\]
\[
 V'_{hk}=\frac{I_h}{w_h}\delta_{hk}+\frac{C}{w_hw_k},\qquad
 \delta'_h=\frac{W_h+E}{w_h},\qquad
 S'=S+\frac{b-a}{4}.                                 \tag{OF.3}
\]
Every free invariant and every increment is an integer. The
actual physical return is \([\Phi_P]=S\pmod {28}\), with
\(\mu(\Sigma_s)=s/28\) and the complex orientation of P3K.68.
This gives every lattice-basis fan of these specified OS49
families, retaining both norm signs and every finite character.
For the shorter \(k=1\) loop and the basis
\(\left(\begin{smallmatrix}2&-1\\-1&1\end{smallmatrix}\right)\)
the class is \(7\).

As the fan varies, the set of classes for either the canonical
loop or the shorter loop is exactly
\[
                    \gcd(k,28)\mathbb Z/28\mathbb Z. \tag{OF.4}
\]
In particular the \(k=1\) family realizes all twenty-eight
elements, and every nonzero one has a mapping torus not
diffeomorphic to \(S^1\times S^6\), by P3F.24 with both spin
structures.

For the fixed-source marked affine cap, line and annular
variations of the retained K44 construction, a closed loop
has scalar winding \(\ell=Nk\), and its return on the
same transported fan is \(N S_k(P)\pmod {28}\).
In particular its fixed-scalar return is zero. Here
\(S_k(P)\) is the shorter-loop algorithm above; the
cap orders, sites, source marking and allowed endpoint
class are those specified in K44.
\end{theorem}
\begin{proof}
\begin{enumerate}[label=\arabic*.]
\item \par\medskip\noindent\textbf{OF.1. The actual period and endpoint maps.}\par\nopagebreak[4]
We begin with the integral generator change (BW.16), including
its entire bounded factors:
\[
 H=-u-\gamma=(s,-mv)^t-U-C,\qquad
 K=3u+2\gamma=(0,s+2mv)^t+D,\qquad D=3U+2C.
 \tag{OF.5}
\]
It has determinant one. In the physical characters the full
period maps are
\[
 H:(Z,V)\longmapsto(\kappa qZ,\alpha V),\qquad
 K:(Z,V)\longmapsto(gZ,q\eta V),
 \tag{OF.6}
\]
where the parameter windings of \((\kappa,g,\alpha,\eta)\)
on the selected single loop are \((0,0,-j,2j)\).
The source and line units in (OF.5) remain functions of \(q\).
On the whole loop the bounded forward closing is
\((g^{-\nu},\eta^{-\nu})\), and its cut is the exact same
integer-\(e=\nu\) cut (BW.5a).

\item On the shorter loop the endpoint-to-initial holomorphic map is
the actual map from (BW.24),
\[
 C_r=r^{2\nu k}g_0U^{-2\nu}:X_{\zeta_3r}\longrightarrow X_r,
 \qquad g_0^3=q^{-\nu}\rho^{-2\nu}.
 \tag{OF.7}
\]
Its full cap square uses \(L^{-2\nu}\) and the identities
\(DLD^{-1}=\zeta_3L\), \(L^3=B_0^{-1}\). Its bounded logarithm
and extending origin are exactly
\[
 c(v)=\frac\nu3(D+2mv\delta)-\frac{2\nu k}{3}\delta,
 \qquad \sigma(v)=\nu vD+\nu mv(v-1)\delta,
\]
\[
 \sigma(v+1/3)-\sigma(v)=c(v),\qquad
 \sum_{h=0}^2c(v+h/3)=\nu D+2\nu mv\delta.            \tag{OF.8}
\]
Thus every intermediate root sheet, the character \(d\), and
the full unit in (OF.7) are specified before any fan changes.

\item Here are the comparison maps on an arbitrary basis. For an
integral matrix \(P\), take the character morphism
\((Z,V)\mapsto (Z,V)^{P^{-1}}\), and change the period generators
by the columns of \(P\). These are mutually inverse integral
monomial maps of the dense complex torus. For its full multiplier
matrix \(A=(A_{hk})\) this gives exactly
\[
 A^P_{ik}=\prod_{h,l}A_{hl}^{(P^{-1})_{ih}P_{lk}}.
 \tag{OF.9}
\]
The valuation matrix of (OF.6) is the identity, so (OF.9)
retains that matrix and every remaining factor is a full unit
on the protected \(q\)-disc. Apply the same character map to
\(c,\sigma\) and every partial sheet product in (OF.8).
It preserves those equalities exactly. This also specifies
the parameter endpoint on every new toric chart.

\item For clarity, a cut multiplier \((A_0,D_0)\) acts on the two
vertex charts by
\[
 (a,b_1,c_1)\longmapsto((A_0D_0)^{-1}a,A_0b_1,D_0c_1),
\]
\[
 (u,v_1,w_1)\longmapsto(A_0D_0u,D_0^{-1}v_1,A_0^{-1}w_1).
 \tag{OF.10}
\]
Both triple products are literally \(q\). In the native C0
order \((b_1,a,c_1)\), the Cox lift is
\[
 (a_0,a_1,c_0,c_1)\longmapsto
 (a_0,A_0a_1,(A_0D_0)^{-1}c_0,D_0c_1).             \tag{OF.11}
\]
The other native orders use the corresponding permutation
and the constant normal interchange of AF.2. Equations
(OF.8)--(OF.10) hold on all four factors of each Cox lift,
so the actual endpoint extends to both rulings at every step.
On the common exterior, compose (OF.9) with its inverse
character map. This is the original identity of the source
and marked line. It carries the original ordinary-to-cap
arrow (BW.15), and hence the exact cap square (OF.7), including
the finite character and the type-II chart. No new cap is
inserted in this comparison.

\par\medskip\noindent\textbf{OF.2. The full fan morphisms and products.}\par\nopagebreak[4]
For a determinant-minus-one input, the column swap and then
the six matrices of AF.1 give identical full fans by the
displayed triangle translations. Identity on the absolute
characters extends to identical cone semigroup rings and
commutes with (OF.6), (OF.7) and every finite-cap map.
Thus this is an actual holomorphic family isomorphism with
its physical circle and complex orientation.

\item The unique representative is a vertical \(T_n\), or has
columns \((p,q_0),(-r,s_0)\), where \(p,r>0\) and
\(ps_0+rq_0=1\). The successive right multiplications
\(Q_0^{-1}\) for \(r\ge p\) and \(Q_1^{-1}\) for \(p>r\)
decrease \(p+r\), and terminate at \(T_n\). Reversing the
word constructs its actual forward families. Each target
has weights \((r,p,-r-p)\), which are all nonzero.

\item In native coordinates the two other-ruling charts are
\[
 (n_0/n_1,zn_1,n_1),\qquad(n_1/n_0,zn_0,n_0).
 \tag{OF.12}
\]
They have the same rank-one matrix and preserve
\(q=zn_0n_1\). For C1 the explicit unitary path
(P3K.71.6) carries the normal-coordinate interchange. These
are precisely the Cox morphisms of AF.2, now with the full
periods (OF.6) and endpoints (OF.7)--(OF.11). The common
blowup, its two projections, and the integral blowup
decompositions of P3F.35 transfer the sphere topology at each step.
The unchanged compact exterior and the protected finite
periodic chart quotient give the actual compact holomorphic
sphere family on every fan in the word.

\item Take the literal double cover of the selected parameter
circle. Its winding matrix is
\(2j\left(\begin{smallmatrix}0&0\\-1&2\end{smallmatrix}\right)\).
Extract these integral windings from the full period functions
and retain their based logarithms \(f_{hk}\). Their transformed
zero-winding logarithms are
\[
                    f^P_{ik}=\sum_{h,l}(P^{-1})_{ih}P_{lk}f_{hl}.
 \tag{OF.13}
\]
The full normal contractions of GF.2 and K64 therefore are
\[
                          Q\longmapsto H_{2j}QH_{l_h-2j}.
 \tag{OF.14}
\]
Use the fixed smooth maps \(H_d\) of GF.9, multiply by the
exponentials of (OF.13) with the protected radial cutoffs,
and carry all higher normal terms by the normal-rescaling
flow. This constructs the actual ordinary product on the
whole three-curve plumbing, including its two vertex collars,
source, physical line and stable derivative. Its oriented
orbit-prism and derivative-prism maps are the full GF.3
maps with these functions; their degrees give exactly
\(I_h,W_h\) of (OF.2).

\item On a native star the Hopf-boundary morphism is
\[
             q_H\longmapsto (UH_{2j}U^{-1})q_HH_{l_h-2j}.
 \tag{OF.15}
\]
The left degree is \(4j\), the right degree is zero, and
the actual inverse-cylinder endpoint is \(L^{-4j}\).
The native order at the vertical fixed star is (BW.26)--(BW.27);
(OF.11)--(OF.12) carry it to every selected star. The constant
normal interchange preserves the real orientation and both
third degrees. Thus the pure-left comparison \(-4j\) is
the same on the two retained native products at every step.

\item For the shorter endpoint, (OF.8) makes its three-fold
pullback the actual covered product on each Cox chart.
Its comparison to the canonical seam is the based global
physical-circle action with logarithm
\(-\nu m(\theta^2+\theta)\delta\) of (BW.28), together with
the full bounded-unit logarithms vanishing at the endpoints.
Apply \(P^{-1}\) to this same physical circle; its global
action, cap extension and all normal jets are carried by
(OF.9). The degree-three covering therefore multiplies
the free relative invariants and both characteristic
integrals by three on these actual products. This proves
the use of \(j=k\) in the shorter formulas while retaining
the specified endpoint.

\par\medskip\noindent\textbf{OF.3. The integer seed and one actual step.}\par\nopagebreak[4]
The complete relative line and tangent morphisms of BW.1--5
give the seed (OF.1). In particular its physical pair is
\((2j-6j(n-1)^2,8j)\), and its single native fixed-star
pair is \((2j,-6j(n-1)^2)\). Its doubled native pair is
\((4j,-12j(n-1)^2)\); the comparison \(-4j\) just proved
gives the ordinary pair \((0,-12j(n-1)^2)\).
The full vertical prism equations in (BW.30) give every
off-diagonal and derivative entry in (OF.1).

\item Let \((\mu_i,\rho_i)\) be an old selected star's actual
single native free pair. Its doubled ordinary pair is
\((2\mu_i-4j,2\rho_i)\). Consequently the product-frame
variation and derivative at that star satisfy
\[
 a=4j-2\mu_i-2\rho_i,\qquad
 b=4j-2\mu_i+2\rho_i,\qquad \rho_i=(b-a)/4.
 \tag{OF.16}
\]
The four-weight Cox common-link map, with the physical
complex-normal reset and the transported Hermitian metric,
has the proved relative action
\[
             (\mu,\rho,s)\longmapsto(\mu+2\rho,-\rho,s+\rho).
 \tag{OF.17}
\]
This is the actual local map of K55/K62, with the endpoint
and native products carried in OF.1--2. The new selected
curve is C2. Applying its same ordinary comparison gives
\[
              V'_{22}=a,\qquad \delta'_2=2a-b,\qquad
              S'=S+(b-a)/4.                          \tag{OF.18}
\]
The full target prism equations are
\[
 V'(w_he_h-w_0e_0)=I_he_h-I_0e_0,
 \qquad w_h\delta'_h-w_0\delta'_0=W_h-W_0,
 \quad h=1,2.                                       \tag{OF.19}
\]
For the symmetric variation, the diagonal particular solution
is \(\operatorname{diag}(I_h/w_h)\). Its difference from any
other solution is \(Crr^t\), with \(r_h=1/w_h\), by the
two displayed vector equations. The derivative equations
give \(\delta'_h=(W_h+E)/w_h\). The selected entries
(OF.18) determine \(C,E\) exactly as in (OF.3).
These equations prove uniqueness over \(\mathbb Q\).
The actual compact family, ordinary products and oriented
prisms in OF.1--2 supply an integral solution; hence the
displayed unique solution is integral. Equation (OF.16)
proves that the increment is its actual single native
integer. Induction along the finite word proves the
algorithm and the stated physical class.

\par\medskip\noindent\textbf{OF.4. Exact execution and the complete image.}\par\nopagebreak[4]
The exact program \path{code/part3_root/os49_all_fans.py}
checks every determinant \(\pm1\) matrix with entries
of absolute value at most eighteen. There are \(6504\)
inputs and \(814\) distinct full fans. It runs the actual
shorter and canonical data for \(k=1,2,5,7\), giving
\(52{,}032\) executions, and checks all free invariants
under the actual degree-three cover on \(26{,}016\)
inputs. For each \(k\), a separate three-edge traversal
visits \(439\) states through depth eight and agrees on
the physical classes, including inverse edges. Its four
separate-fixed-target transitions retain an undetermined
vertical integration constant and are excluded from that
traversal; the actual vertical seed (OF.1) supplies those
targets in the main terminating algorithm. There are
\(648\) additional vertical checks for \(-40\le n\le40\)
on both parameter circles.

\item For \(P=\left(\begin{smallmatrix}2&-1\\-1&1\end{smallmatrix}\right)\)
and the shorter \(k=1\) loop, the word is \((0,1)\).
The seed is \(-1\). Its two selected old pairs
\((V_{ii},\delta_i)\) are \((12,-12)\) and \((72,-96)\),
so the single native increments are \(-6,-42\).
Thus \(S=-49\equiv7\pmod {28}\), with complete final
ordinary invariants
\[
 V=\begin{pmatrix}472&300&-200\\300&184&-100\\-200&-100&72\end{pmatrix},
 \qquad\delta=(-592,-392,240)^t.                     \tag{OF.20}
\]

\item The following table supplies one actual \(k=1\) shorter
fan for each target residue. Rows \((a,b;c,d)\) denote
the corresponding determinant-one matrix, and \(S\)
is its computed integer path sum. The attached output
contains every word, selected pair, native increment,
and complete intermediate invariant.
\input{notes/part3_os49_all_fan_witness_table.tex}

\item The separate certificate
\path{code/part3/audit_os49_all_fan_witnesses.py}
imports none of the reconstruction routines. It solves
the full symmetric-variation and derivative equations
afresh for every target in every displayed witness,
checks the selected Cox entries and each increment,
and repeats the target calculations at \(j=3\).
Its \(166\) exact target solves verify all twenty-eight
witnesses and their complete free-invariant cover scaling.

\item Finally, for every fixed input \(P\), (OF.1)--(OF.3)
give the exact identities
\[
 V_j(P)=jV_1(P),\qquad\delta_j(P)=j\delta_1(P),\qquad
                         S_j(P)=jS_1(P).              \tag{OF.21}
\]
Indeed the seed and every term of (OF.2) are linear in
\(j\); the denominators and the finite word depend only
on \(P\). Induction in (OF.3) proves all three equalities
as integers. The twenty-eight displayed witnesses prove
that the image for \(j=1\) equals \(\mathbb Z/28\).
Equation (OF.21) then proves the image \(j\mathbb Z/28\)
for every positive \(j\). On the actual shorter loop
\(j=k\), and on the actual canonical loop \(j=3k\).
Multiplication by three is invertible modulo twenty-eight,
with inverse nineteen. Their images are consequently
exactly (OF.4). This image equality uses complete witnesses
for the finite target together with the unrestricted
geometric algorithm proved above.

\par\medskip\noindent\textbf{OF.5. The retained marked cap-data loops.}\par\nopagebreak[4]
We carry the actual construction in K44, OC.1--5, through
the same maps. At its good root, the finite affine norm
is solved by the explicit section
\[
 T=\widetilde T-[m]^{-1}\sum_{h=0}^{m-1}hD^hZ_1,
 \qquad (1-D)T=Z.
\]
Its source translation has the extending line lift and
retains the finite kernel and character. The full
orbifold Lie sequence has coarse terms \(\mathcal O\)
and \(\mathcal O(-1)\); the Laurent/Cauchy splitting
and finite averaging of OC.1 give the actual interval
comparison. On the cusp the integral period reindexing
of OC.7 leaves the complete unit map
\((Z,V,q)\mapsto(a(q)Z,b(q)V,q)\). Formula (OF.9)
is its exact expression in any chosen fan, and (OF.10)
extends it and its inverse on all affine cone charts.
Thus the same comparison preserves the transported
fan and its full Cox normal maps.

\item The global endpoint calculation OC.14 identifies its
remaining translation with the actual physical
\(\mathbb C^*\) action. The root square and the
mixed-cusp lattice impose the exact integer equation
\[
                         3v_0-j_0=\ell\nu/k.
 \tag{OF.22}
\]
Here \(j_0\) is the source-translation integer and
\(v_0\) the integral scalar order in OC.16--18.
Since \(\nu=\pm1\), (OF.22) proves \(\ell=Nk\).
The actual endpoint (OF.7), iterated \(N\) times,
realizes every such integer, with its original cap
character. Inverse iteration defines negative \(N\).

\item Write the constructed interval isomorphisms as
\(I_t:X_t\to Y_t\), where \(Y\) is that specified
\(N\)-fold shorter family. Write
\(E_X:X_0\to X_1\), \(E_Y:Y_0\to Y_1\) for the
forward endpoints and \(P_Y:Y_0\to Y_1\) for smooth
parallel transport. The endpoint ratio is an extending
physical scalar \(C:Y_1\to Y_1\), and the actual
commutative-square equation is
\[
 I_1E_XI_0^{-1}=CE_Y,
\]
\[
 I_0\bigl(E_X^{-1}I_1^{-1}P_YI_0\bigr)I_0^{-1}
                         =E_Y^{-1}C^{-1}P_Y.            \tag{OF.23}
\]
A chosen logarithm of this nonzero scalar gives
\(C_s=\exp(s\log C)\), a global physical-circle-complexification
isotopy on the same compact family, including its
finite caps and every toric chart. It removes \(C\)
from the right side of (OF.23). The physical return
is therefore the \(N\)-th iterate of the exact return
already calculated, namely \(N S_k(P)\pmod {28}\).
This reuses the proved marked family maps while
replacing the refuted numerical expression OC.22 by
the complete calculation (OF.1)--(OF.3).

\end{enumerate}
\end{proof}




\begin{theorem}\label{night:P3K.78}
\StatementRecord{night:P3K.78}{proved}{S6\_06\_BLIND: GPT-6 Astra (OpenAI, ChatGPT Pro mode, effort Pro), recorded re-derivation in the checkpoint dated 18 November 2026}{\cref{night:P3F.26}, \cref{night:P3F.29}, \cref{night:P3F.33}, \cref{night:P3F.34}, \cref{night:P3K.44}, \cref{night:P3K.76}, \cref{night:P3K.77}}
Retain the fixed OS49 source
\[
y^2=x^3-3t^3(t-1)x+2t^4(t-1)(t-\tfrac12),\qquad G=(t^2,t^2),
\]
the actual Cartier line \(M=\mathcal O(G-O)\), suspension over
\(-3G\), retained type-II fibre at \(1\), selected component cusp
at infinity, and good \(IV^*\) cap at \(0\). Let the finite cap
have order \(m=3k>0\), invariant vector \(n\ell_0+d\delta\), and
an actual free action as in P3F.33, where
\(\ell_0=(1,-1,0,0)^t\). Fix any periodic index-one lattice-basis
fan \(P\), with the original source, line, root and annular markings.

At finitely many fixed, marked, disjoint smooth-site discs, retain
ordinary-to-disc translation clutches and their full line lifts
\[
 \Phi_i(x,z_i)=x+\frac{\log z_i}{2\pi i}\Pi(z_i)\ell_i+u_i(z_i),
 \qquad \ell_i\in\Lambda=\mathbb Z^4.                 \tag{OCS.1}
\]
Here every vector is transported along its marked cut, and every
single-valued Lie term \(u_i\), including its line component and
annular units, is retained. In the basis \((\gamma,u,w,\delta)\), put
\[
 \psi=(1,0,0,0),\qquad
 \sigma=\sum_i\psi(\ell_i),\qquad D=n+m\sigma.        \tag{OCS.2}
\]

The complete integral extending-period groupoid has
\(H^2(\mathcal H)=\mathbb Z\oplus\mathbb Z/k\), with these data
represented by \((D,d\bmod k)\). The explicit integral cochain
and actual holomorphic maps below identify the entire clutched
family, on the same fan and at the same scalar, with the one-cap
datum \((m,D,d)\). They retain every finite group cochain and the
actual parameter-closing maps. Its fundamental group is
\(\mathbb Z/|D|\), with \(D=0\) interpreted as \(\mathbb Z\), and
its sphere members are exactly \(D=\pm1\).

For the canonical positive \(r\)-loop \(\lambda=r^m\) in this
sphere sector, the oriented physical return is \(S_m(P)\) of
P3K.77. Transporting the actual shorter OR65 endpoint by the
constructed family map gives the shorter return \(S_k(P)\),
with its full three-step norm. This last assertion concerns that
specified transported endpoint. The theorem retains fixed smooth
order-one sites, the fixed source and its stated singular models;
it makes no assertion about an additional unspecified short-loop
identification or a change of source.
\end{theorem}
\begin{proof}
\begin{enumerate}[label=\arabic*.]
\item \par

\medskip\noindent\textbf{OCS.1. The actual extending period sheaf}

\par\nobreak

\item Order the homology lattice as \((\gamma,u,w,\delta)\), and set
\(\psi=(1,0,0,0)\), \(\ell_0=(1,-1,0,0)^t\). The actual marked
homology matrices, obtained as inverse transposes of F26.3, are
\[
 A=\begin{pmatrix}1&0&0&0\\0&1&-1&0\\3&3&-2&0\\1&1&-1&1\end{pmatrix},
 \quad B=\begin{pmatrix}1&0&0&0\\0&0&1&0\\0&-1&1&0\\0&0&0&1\end{pmatrix},
 \quad C=\begin{pmatrix}1&0&0&0\\0&1&0&0\\-3&-2&1&0\\-1&-1&0&1\end{pmatrix}.
\tag{OCS.3}
\]
They satisfy \(ABC=1\), \(A^3=1\), and \(\psi A=\psi B=\psi C=\psi\).
The finite root atlas has the constant lattice \(\Lambda=\mathbb Z^4\)
and cyclic group of order \(m\) acting by \(A\), including its
ineffective subgroup on this linear lattice. That subgroup is not
discarded.

\item At the retained II disc the extending period stalk is
\(\ker(B-1)=\mathbb Z\gamma\oplus\mathbb Z\delta\). The source
translation direction is the actual additive direction there; its
elliptic periods have no nonzero invariant combination. The suspension
period and the physical period extend by the actual line-complement
quotient. At the selected cusp the full period valuation matrix has
determinant one. Its two logarithmic radial periods therefore have
no nonzero bounded combination; the two angular periods extend as
the logarithms of the actual source and physical torus actions.
The stalk is exactly
\(\ker(C-1)=\mathbb Zw\oplus\mathbb Z\delta\).
These are the explicit kernel calculations of (OCS.3), and match
the unit charts in K44, OC.2. On each nonfinite disc the period sheaf
is the invariant extension \(j_*\Lambda\).

\item For completeness its higher cohomology on that disc is zero. On a
punctured disc the circle complex is
\(\Lambda\xrightarrow{M-1}\Lambda\). Its degree-one group is
the coinvariant lattice. In the radial extension this is exactly the
local degree-one term removed by passing from the punctured derived
extension to its invariant degree-zero sheaf: the map on that term
is the identity of the same marked circle complex. Equivalently,
using two slit sectors and the centre, the same cancellation gives
only \(\ker(M-1)\) in degree zero. This is an integral cancellation;
no rational averaging occurs. Smooth-site discs give the constant
lattice in degree zero and the same vanishing in higher degrees.

\item Let \(\mathcal A\) be the sheaf of extending translations of the
actual compactified suspension, on the full root atlas and groupoid.
Let \(\mathcal V\) be its holomorphic Lie sheaf and \(\mathcal H\)
its period sheaf just described. The exact sequence is
\[
 0\longrightarrow\mathcal H\longrightarrow\mathcal V
                  \xrightarrow{\exp}\mathcal A\longrightarrow0.
\tag{OCS.4}
\]
At the cusp this follows from the literal integral reindexing and
bounded unit map in K44, OC.7--8; logarithms of those units give
the local Lie lift on every same fan. At II it is the actual source
and line lift. On the good root it is the usual local exponential
cover of the smooth compact torus family, with its entire finite
group cochain retained. It is not an assertion of surjectivity on
invariant coarse sections.

\item The coarse terms of \(\mathcal V\) are \(\mathcal O\) and
\(\mathcal O(-1)\), as proved by the actual IV-star substitution
in K44, OC.3. The original source vector is \(z^{2k}\) times the
root vector and \(0\le2k<3k\), so this coarse extension has not
changed. Its vector extension splits by the Laurent calculation
\(H^1(\mathcal O(1))=0\). Thus
\[
                  H^1(\mathcal V)=H^2(\mathcal V)=0.
\tag{OCS.5}
\]
Finite averaging is used only on these complex vector sheaves.
It is not used on \(\mathcal H\); the latter has the torsion
computed next.

\par

\medskip\noindent\textbf{OCS.2. A complete integral root-groupoid complex}

\par\nobreak

\item Remove the three special discs. On their pair of pants use the
marked fundamental loops with product \(ABC=1\). A one-cochain
is \((p,q)\in\Lambda^2\); its three restrictions are
\[
                    p,\qquad q,\qquad -C(p+Aq).
\tag{OCS.6}
\]
The finite root disc uses the full cyclic resolution, whose
differentials alternate \(A-1\) and
\(N_m=1+A+\cdots+A^{m-1}\). Each of the other discs has only
the invariant degree-zero term described above.

\item A total degree-two cocycle is
\((\eta,b_A,b_B,b_C)\), with \(\eta\in\Lambda^A\) and
\(b_A,b_B,b_C\in\Lambda\). The complete degree-one relations,
for \(p,q,a,r,s,t\in\Lambda\), are
\[
 \begin{split}
 \eta&=N_m a,\\
 b_A&=p-a-(A-1)r,\\
 b_B&=q-(B-1)s,\\
 b_C&=-C(p+Aq)-(C-1)t.
 \end{split}                                           \tag{OCS.7}
\]
These are the Mayer--Vietoris cone maps of the displayed actual
restrictions. In particular the finite group relation appears as
\(N_m a\), and is not replaced by the order-three norm.

\item With extra ordinary boundary loops take the actual presentation
\(a b c\mu_1\cdots\mu_r=1\), where each \(\mu_i\) has identity
linear action in its transported frame. If its pair-of-pants
one-cochain is \(p_i\), the third restriction in (OCS.6) is
\(-C(p+Aq+\sum_i p_i)\), and its own restriction is \(p_i\).
Consequently (OCS.8) gains exactly \(\sum_i b_i\). Eliminating
each \(p_i\) is an integral identity block and gives the addition
of the full transported clutch vector used below. This fixes the
cut order and every vector transport in the calculation.

\item Eliminate \(p,q\) by the unimodular boundary-coordinate change
\[
                 v_*=b_A+A b_B+C^{-1}b_C.
\tag{OCS.8}
\]
Its determinant on the three lattice blocks is \(\det C^{-1}=1\).
The remaining relation in that block is
\[
 v_*=-a-(A-1)r-A(B-1)s-(1-C^{-1})t.                   \tag{OCS.9}
\]
The last three images generate precisely \(\ker\psi\), integrally.
Here is an explicit preimage table, with its three entries ordered
as the arguments of \((A-1),A(B-1),(1-C^{-1})\):
\[
 \begin{array}{c|ccc}
 \text{target}&r&s&t\\\hline
 u&-u&w&0\\
 w&-w&-u-w&0\\
 \delta&2w&2(u+w)&-Cu.
 \end{array}                                           \tag{OCS.10}
\]
Substituting (OCS.3) verifies every column. The row \(\psi\)
kills all three images. Thus the quotient is the primitive lattice
\(\mathbb Z\), without hidden finite index.

\item The invariant lattice of \(A\) is primitively generated by
\(\ell_0,\delta\), and the exact norm is
\[
 N_3=\begin{pmatrix}3&0&0&0\\-3&0&0&0\\0&0&0&0\\0&0&-1&3\end{pmatrix},
 \qquad N_m=kN_3.
\tag{OCS.11}
\]
Writing \(c\) for the remaining ordinary base generator, the
entire reduced relation matrix in the basis
\((\ell_0,\delta,c)\) is consequently
\[
 \begin{pmatrix}m&0&0&0\\0&0&-k&m\\-1&0&0&0\end{pmatrix}.
\tag{OCS.12}
\]
The first column gives \(c=m\ell_0\); the third gives
\(k\delta=0\), and the fourth is its multiple. Hence
\[
 H^2(\mathcal H)=\mathbb Z\ell_0\oplus(\mathbb Z/k)\delta,
 \qquad
 (\eta=n\ell_0+d\delta,v_*)\longmapsto
                   (n+m\psi(v_*),d\bmod k).
\tag{OCS.13}
\]
This is an integral atlas/groupoid computation. In particular its
\(\mathbb Z/k\) summand has not been lost by taking finite
invariants over \(\mathbb C\).

\item Adding a smooth disc is the same cone calculation with identity
linear monodromy and no finite group. Its full clutch vector is
transported along the marked cut, and contributes that vector to
\(v_*\). Thus its class in (OCS.13) is
\(m\psi(\ell_i)\ell_0\). The original cap cocycle is exactly
\(n\ell_0+d\delta\), because its lifted cyclic generator has
\(m\)-th power that period.

\par

\medskip\noindent\textbf{OCS.3. The signed meridians and the exact integer cochain}

\par\nobreak

\item The plus-log arrow (OCS.1) sends the zero section of the filled
disc to the path \(-s\Pi\ell_i\) in the ordinary coordinates.
Its filled positive meridian is therefore
\[
                         \alpha_i-\ell_i.
\tag{OCS.14}
\]
For the finite root the marked filled meridian is
\(m\alpha_A-(n\ell_0+d\delta)\), exactly F34.3. The II and cusp
fillings kill their base meridians and the primitive relation
lattice computed in F26. The surviving fibre quotient is
\(\mathbb Z\gamma\), measured by \(\psi\). The base relation
then gives \(\alpha_A=-\sum_i\psi(\ell_i)\gamma\), and the
finite power relation gives \(D\gamma=0\), with the plus sign
in (OCS.2). These are all generators and relations: each filled
boundary maps surjectively onto its local group, and the two
nonfinite relations reduce the original commuting torus generators
to this one generator. Thus
\[
                         \pi_1(X)=\mathbb Z/|D|,
\tag{OCS.15}
\]
where \(D=0\) means \(\mathbb Z\).

\item We now give the full integer cochain behind redistribution. Put
\(L=\sum_i\ell_i\) after the marked transports and
\(\sigma=\psi(L)\). The difference from the one-cap reference
\((m,D,d)\) has reduced cocycle
\[
                  (\eta,v_*)=(-m\sigma\ell_0,L).
\tag{OCS.16}
\]
Take the finite-root one-cochain \(a=-\sigma\ell_0\) in (OCS.7).
Its contribution is \((-m\sigma\ell_0,\sigma\ell_0)\).
The remaining vector \(L-\sigma\ell_0\) lies in \(\ker\psi\).
Use its actual \((u,w,\delta)\) coefficients in the negative
of the table (OCS.10) to obtain \(r,s,t\) in (OCS.9).
The eliminated \(p,q\) are recovered by the inverse of (OCS.8).
Thus (OCS.16) is the coboundary of a specified integral total
one-cochain, including every component of every original vector.
The slit-sector identifications in (OCS.6) are the actual parallel
transports along the marked cuts, so this is also a cochain on
the refined analytic atlas. No ordinary vector has been replaced
by its gamma projection without this cochain.

\item The finite action itself is retained. Its two displayed affine
lifts differ by
\[
 \frac{D-n}{m}\ell_0=\sigma\ell_0,
\tag{OCS.17}
\]
an invariant integral period of the actual smooth root torus.
They induce exactly the same automorphism of that torus and the
same finite quotient. Their lifted norm and annular logarithm
differ by the recorded one-cochain, not by an omitted choice.
The delta coefficient is unchanged: \(N_m(-\sigma\ell_0)\)
has zero delta component. The fixed finite character \(d\)
therefore remains the original one. Since \(D\equiv n\pmod m\),
the free-action criterion of F33 holds for both displayed data.

\item The same change has a literal source/line-cover expression. Retain
the actual commuting cap and suspension maps and put
\[
 B_r=r^mB_0,\qquad A_{n,d,r}=r^nA_{n,d,0},\qquad
 A_{D,d,r}=B_r^\sigma A_{n,d,r}.
\tag{OCS.17a}
\]
Then \((A_{D,d,r})^m=B_r^D\) and
\(\langle A_{D,d,r},B_r\rangle=\langle A_{n,d,r},B_r\rangle\).
The root coordinate is unchanged. In the root marking the
suspension period is the invariant \(\ell_0\); its full source
and line factors are those of F29, not just its gamma covector.
Thus (OCS.17a) gives exactly the invariant lattice change in
(OCS.17), with central line factor \(B_0^\sigma\) and scalar
factor \(r^{m\sigma}\) retained. Its fixed delta character
commutes with these maps. Identity on this quotient is the
finite-cap map used in (OCS.18).

\par

\medskip\noindent\textbf{OCS.4. Realizing the cochain by actual holomorphic family maps}

\par\nobreak

\item Compare the given clutched family with that one-cap reference
using the unchanged local ordinary, II, cusp, and root models.
Their local overlap differences are actual sections of
\(\mathcal A\). On simply connected atlas sectors choose the
actual holomorphic Lie lifts, including their source and line
coordinates. Their lattice-valued total two-cocycle is precisely
the one computed in OCS.2--3. Subtract the periods of the explicit
integral one-cochain in (OCS.16). The resulting Lie one-cochain
is closed on all ordinary overlaps and on the full cyclic
root groupoid.

\item Solve this closed cochain by (OCS.5). This is an actual analytic
operation: split the negative and nonnegative Laurent parts in
the two coarse summands \(\mathcal O\), \(\mathcal O(-1)\),
use the same Laurent splitting for their extension, and average
only the resulting complex vector cochains at the finite root.
On a refinement of the given discs, fixed Cauchy contours give
the identical solution. All source and line units in (OCS.1)
are inputs to these integrals and are retained.

\item Here is its exact map convention. Write the new overlap as
\(g^{\rm new}_{ij}=t_{\theta_{ij}}g^{\rm ref}_{ij}\), with
linear Lie action \(L_{ij}\). The solved equation is
\(\theta_{ij}=s_i-L_{ij}s_j\). The local maps
\[
 F_i=t_{-s_i},\qquad
                F_i g^{\rm new}_{ij}=g^{\rm ref}_{ij}F_j
\tag{OCS.18}
\]
therefore glue literally. Their inverses are the opposite
translations. On the cusp their exponentials are the full bounded
torus units after the exact period reindexing of K44, OC.8; those
units and their inverses extend on every cone of the same fan.
On the root, (OCS.18) is also the finite-deck square because the
finite group terms were included before averaging. The source
translation lifts to the actual degree-zero root line, and the
generic compact-elliptic commutation argument in K44, OC.1 extends
that equality across the root. At II it uses the original reduced
line-complement quotient. Thus these are actual holomorphic maps
of the filled families and their lines, including their normal
collars, rather than only isomorphisms of cohomology classes.

\item For a varying scalar all of this is analytic on a small parameter
disc: the period matrix, source/line lifts and given overlap units
are analytic, and the Cauchy contours can be fixed. The integral
cochain in OCS.3 is constant, including its finite group terms.
The canonical parameter can be \(\lambda=r^m\); the original
finite lift, the changed lift, and their parameter-closing arrows
remain part of (OCS.18).

\item The local parameter comparisons glue over an annulus as follows.
Their ratios are global extending translations. The fixed-scalar
endpoint proof of K44, OC.3 applies here without a restriction on
the free norm \(n\): the root commutator has \(m\)-fold value
\(B^{ma}\), and \(A^m=B^n\) and commutation with \(B\) still
force \(a=0\). Its rational source/line lift is \(fU^j\);
the actual divisor of \(B_0^{-1}/U^3\) still gives \(v_*=j/3\).
Thus \(j=3a\) and \(fU^{3a}=cB_0^{-a}\), which on the quotient
is the actual physical scalar \(c\lambda^a\). The entire global
translation group is consequently the same physical \(\mathbb C^*\),
with every finite kernel included. This also holds on the new
torsor because translations commute and hence have the same
translation sheaf.

\item The parameter Isom torsor is therefore a holomorphic
\(\mathbb C^*\) torsor. Here is an explicit smooth nonzero section
on a thin annulus. Take local holomorphic comparisons \(F_\alpha\),
with \(F_\alpha=c_{\alpha\beta}F_\beta\), and a smooth partition
of unity \(\rho_\beta\) on their parameter cover. The forms
\[
 \mathfrak a_\alpha=
       \sum_\beta\rho_\beta\,d\log c_{\alpha\beta},
 \qquad
 \mathfrak a_\alpha-\mathfrak a_\gamma
                         =d\log c_{\alpha\gamma}
\tag{OCS.18a}
\]
define a connection on this actual torsor. Angular parallel transport
from one chosen point gives a nonzero comparison \(F^{\rm par}_\theta\),
with holonomy \(c\in\mathbb C^*\). Choose a logarithm of that
one scalar and a smooth lift \(\chi\) satisfying
\(\chi(\theta+1)=\chi(\theta)+1\), constant on the prescribed
parameter collars. Then
\(\exp(-\chi(\theta)\log c)F^{\rm par}_\theta\) is a smooth
periodic section, with all derivatives matched.

\item Extend it to a thin radial annulus by the explicit connection ODE.
In a local comparison frame its scalar coefficient satisfies
\[
 \partial_\varrho z_\alpha
     =-\mathfrak a_\alpha(\partial_\varrho)z_\alpha,
 \qquad
 z_\alpha(\varrho,\theta)=z_\alpha(\varrho_0,\theta)
    \exp\left(-\int_{\varrho_0}^{\varrho}
           \mathfrak a_\alpha(\partial_s)\,ds\right).
\tag{OCS.18b}
\]
The transformation rule in (OCS.18a) makes these solutions agree
under the original transition scalars. The connection and the
initial circle section are periodic, so the extended section is
periodic as well. It is nowhere zero and consists throughout of
actual family isomorphisms.
In the associated line write its induced antiholomorphic derivative
as \(\bar\partial s=\omega s\). Extend the coefficient of
\(\omega\) smoothly with compact support to a plane neighbourhood
and put
\[
 u(r)=\frac1\pi\int\frac{\omega_{\bar w}(w)}{r-w}\,dA(w).
\tag{OCS.19}
\]
On the smaller annulus \(\bar\partial u=\omega\), so
\(e^{-u}s\) is an actual nonzero holomorphic section. Multiplication
by this scalar acts on all original line, toric and finite-cap
charts. It changes no finite character or annular lattice vector.
This proves a closed holomorphic family isomorphism on the parameter
annulus, with the actual endpoint arrows carried. It does not infer
trivial loop transport from a fibrewise conjugacy.

\item The physical return can also be checked directly before this last
scalar adjustment. Retain the actual parameter period matrix:
on the positive canonical \(r\)-loop, with \(\lambda=r^m\),
its cut logarithm obeys
\[
             \Pi(v+1)-\Pi(v)=m\delta\psi,
                \qquad v=\log r/(2\pi i).
\tag{OCS.19a}
\]
Thus the change of the full ordinary clutch (OCS.1) is
\(m\psi(\ell_i)\log(z_i)\delta/(2\pi i)\); its source and
fixed delta components are unchanged. Its inverse forward-closing
factor is exactly \(z_i^{-m\psi(\ell_i)}\). The finite cap
contributes \(t^{-n}\) in the original OS49 base coordinate, with
finite cap at \(t=0\) and selected cusp at infinity. With fixed
clutch units the ordinary forward map is consequently
\[
 f_X(t)=t^{-n}\prod_i(t-c_i)^{-m\psi(\ell_i)},
 \qquad f_Y(t)=t^{-D}.                               \tag{OCS.19b}
\]
At each site all the other factors are retained as its extending
cap unit. At the cusp the first map has valuation \(D\), and its
ratio to \(q^D\) is the full extending unit. Its actual toric
period reindexing is the one in K76. Any given additional
parameter-dependent unit is likewise retained in its prescribed
closing arrow. Formula (OCS.19a), rather than a parameter-constant
period convention, is used in the comparison.

\item Let \(F_t:X_t\to Y_t\) be the constructed interval comparison,
where \(Y\) is the canonical datum \((m,D,d)\) at the same scalar.
Let \(E_X:X_0\to X_1\) and \(E_Y:Y_0\to Y_1\) be the actual
forward endpoints just specified, including their root squares and
all units, and let \(P_Y:Y_0\to Y_1\) be smooth parallel transport.
Both endpoints carry the same marked parameter lattice deck
(OCS.19a); all remaining factors are the actual translation
cochains. Therefore
\[
 C=F_1E_XF_0^{-1}E_Y^{-1}:Y_1\longrightarrow Y_1
\]
is a global extending translation with identity linear lattice
part. The endpoint calculation above identifies it with a physical
scalar, including every finite root kernel. The literal equations
are
\[
 F_1E_XF_0^{-1}=CE_Y,\qquad
 F_0\bigl(E_X^{-1}F_1^{-1}P_YF_0\bigr)F_0^{-1}
                       =E_Y^{-1}C^{-1}P_Y.            \tag{OCS.19c}
\]
The actual scalar isotopy removes \(C\). This is the oriented
physical return conjugacy. In the sphere sector \(D=\pm1\),
\(E_Y\) is precisely the K76 canonical positive \(r\)-endpoint,
so the positive \(r\)-return has class \(S_m(P)\) of K77.

\item There is also a specified shorter endpoint obtained by transporting
the exact OR65 endpoint through this family map. If
\(C^Y_r:Y_{\zeta_3r}\to Y_r\) is its endpoint-to-initial map, set
\[
                 C^X_r=F_r^{-1}C^Y_rF_{\zeta_3r}.
\tag{OCS.19d}
\]
Every cap square and intermediate-sheet product is conjugated by
these same maps, including its full three-step norm. The induced
shorter family is isomorphic to the K77 shorter family and has
return \(S_k(P)\). This concerns this actual transported endpoint;
it assigns no value to an additional unspecified short-loop
identification.

\par

\medskip\noindent\textbf{OCS.5. Topology and the precise completed scope}

\par\nobreak

\item The map (OCS.18)--(OCS.19) identifies the given family with the
one-cap family \((m,D,d)\) at the same scalar and on the same fan.
All original annular integers remain in the explicit comparison.
F34 now computes its entire integral homology with \(n\) replaced
by \(D\); its finite torsion factor is unchanged because
\[
              \gcd(m,D,3d)=\gcd(m,n,3d).
\]
Together with (OCS.15), this proves that its sphere members are
exactly \(D=\pm1\). Their comparison target is precisely the
actual K76 datum of that norm sign, with order \(m\), delta
character \(d\), and the original source/Cartier marking.

\item This concerns finitely many fixed, marked, smooth order-one sites
and their actual translation clutches. The map transports their
full affine and line data. No additional singular-site replacement,
change of source, or arbitrary nontranslation endpoint is included.
The source-zero collision work and its unresolved compact
attachment are separate from this fixed-source calculation.

\item The complete integer matrices, primitive preimages, 14-by-24
total degree-two map, reduced relations and the explicit
redistribution cochain are executed in
\path{code/part3/os49_order_one_groupoid.py} and its saved outputs.

\end{enumerate}
\end{proof}




\begin{theorem}\label{night:P3K.79}
\StatementRecord{night:P3K.79}{proved}{S6\_06\_BLIND: GPT-6 Astra (OpenAI, ChatGPT Pro mode, effort Pro), recorded re-derivation in the checkpoint dated 18 November 2026}{\cref{night:P3F.29}, \cref{night:P3K.37}, \cref{night:P3K.44}, \cref{night:P3K.78}}
Retain an actual fixed-source OS49 sphere model of P3K.76--78,
with the given scalar, component Cartier line, finite cap of order
\(m=3k>0\), all fixed marked smooth order-one clutches and their full
root/line/annular data, and any periodic index-one lattice-basis fan \(P\).
Write \(D=n+m\sum_i\psi(\ell_i)=\pm1\) for its exact sphere datum.
Its full holomorphic automorphism group is the faithful physical
\(\mathbb C^*\) action. Every such automorphism descends to the identity
on the base and has identity linear part on the generic torus. The proof
uses the actual root and same-fan inversion maps to exclude the negative
linear part; no translation or source-descent assumption is imposed.

For a prescribed fixed-source scalar path of integer winding \(L\),
including the marked cap/line variations of P3K.44 and the fixed
smooth-site clutches of P3K.78, every specified holomorphic endpoint has
\(L=Nk\) for an integer \(N\). The explicit canonical endpoint and its
actual powers realize every such winding. All holomorphic endpoints for
that same path have oriented physical return
\[
                              N S_k(P)\in\mathbb Z/28,
\]
where \(S_k(P)\) is the actual all-fan return of P3K.77. The minimum
positive scalar winding in this endpoint scope is \(k\).
The maps retain the full finite group cochain, every line character,
and the original ordinary-clutch integers. This statement concerns the
fixed source and its stated data, with holomorphic endpoints; it does
not identify a different source-moduli path or a smooth endpoint that
has not been given a holomorphic map.
\end{theorem}
\begin{proof}
\begin{enumerate}[label=\arabic*.]
\item \par

\medskip\noindent\textbf{FA.1. Recovering the base and fixing its three marked values}

\par\nobreak

\item Let \(\pi:X\to\mathbb P^1\) be the actual torus fibration. Its generic fibre T is a
compact complex two-torus, hence has a flat Kähler form \(\omega\). Since the
actual X is diffeomorphic to \(S^6\), \(H^2(X;\mathbb Z)=0\). If an effective divisor E
had a horizontal component, its restriction to a generic smooth fibre
would be a nonzero effective divisor \(E_T\). The line \(\mathcal O_X(E)\) has zero first
Chern class, so its restriction has zero first Chern class, whereas
\[
 \int_T c_1\bigl(\mathcal O_T(E_T)\bigr)\wedge\omega
                =\int_{E_T}\omega>0.                    \tag{FA.1}
\]
The equality is the divisor fundamental-class pairing; its last term is
strictly positive on each nonzero effective complex curve with its
multiplicity. This contradiction excludes every horizontal divisor.

\item The polar divisor of a meromorphic function on X is therefore vertical.
Its restriction to a generic compact connected T has no polar divisor;
removable singularities across the remaining codimension-two set make it
holomorphic, and hence constant. The local product over a regular base disc gives its meromorphic
descent there. At any special value choose a generic smooth point of a
reduced fibre component, outside the function's indeterminacy, and a
transverse analytic disc. In its local coordinate the base map is
\(\pi=u^e\) times a holomorphic unit. An extending e-th root of that unit
changes u so that \(\pi=u^e\) exactly. The pulled-back meromorphic function
has a convergent Laurent expansion with a finite principal part on
this disc. On the punctured disc it is constant on the generic fibres,
so this Laurent expansion is invariant under \(u\mapsto\zeta_e u\); every
nonzero exponent is divisible by e. Replacing \(u^e\) by the base
coordinate therefore gives the exact meromorphic extension across
that special value. This proves the descended function meromorphic
on all of \(\mathbb P^1\), and hence rational. Conversely every rational function
of pi pulls back. Thus \(\mathbb C(X)=\mathbb C(t)\), with this equality
implemented by \(\pi^*\). A biholomorphism F of X induces a Möbius transformation
h of \(\mathbb P^1\) and the exact equation \(\pi F=h\pi\). It holds first on regular fibres
and then everywhere by continuation.

\item The three nontrivial linear local monodromies are precisely \(A,B,C\) in
(OCS.3). Their characteristic polynomials are
\[
 \begin{aligned}
 \det(xI-A)&=(x-1)^2(x^2+x+1),\\
 \det(xI-B)&=(x-1)^2(x^2-x+1),\\
 \det(xI-C)&=(x-1)^4,\qquad C\ne I .
 \end{aligned}                                           \tag{FA.2}
\]
The actual finite root has its original effective order-three linear
monodromy A even when \(m=3k\); its ineffective subgroup is still retained in
the affine groupoid. Each added smooth order-one clutch has identity
linear monodromy. A holomorphic base map carries a positive local loop to
a positive local loop, so these conjugacy invariants make h fix 0,1 and
infinity separately. Hence h is the identity. This proves base preservation
for every biholomorphism without a prior source-preservation assumption.

\par

\medskip\noindent\textbf{FA.2. The complete integral linear centralizer}

\par\nobreak

\item A holomorphic map between compact complex tori lifts to their universal
covers. Its derivative is periodic and therefore constant, because its
entries descend to holomorphic functions on a compact torus. Every fibre
map is consequently affine. Its induced integral homology matrix is
locally constant over the regular base. Since F preserves the base, the
matrix X commutes with the actual three monodromies \(A,B,C\).

\item Commutation with B first gives the full block form
\[
 X=\begin{pmatrix}
 a&0&0&b\\0&p&q&0\\0&-q&p+q&0\\c&0&0&d
 \end{pmatrix}.
\]
Substitution in \(XA=AX\) gives \(b=q=0\) and \(a=p=d\). Consequently the complete
rational, and hence integral, centralizer is
\[
                  X=aI+c\,\delta\psi,
       \qquad \delta\psi=E_{41},\qquad \det X=a^4.       \tag{FA.3}
\]
An integral automorphism has \(a=\epsilon\in\{+1,-1\}\) and \(c\in\mathbb Z\). On a fixed
smooth torus its real linear map is complex linear. Subtracting
\(\epsilon I\) still gives a complex-linear map. If c were nonzero, the
result would have real rank one by (FA.3), whereas a nonzero complex-linear
map has even positive real rank. Thus \(c=0\). Every holomorphic automorphism
has linear part exactly \(+I\) or \(-I\). This computation uses the entire
monodromy module; it leaves no exceptional complex-multiplication case.

\par

\medskip\noindent\textbf{FA.3. Actual local inversion and the affine root/cusp square}

\par\nobreak

\item We use the same local origins and rigidified line coordinates as the
OCS groupoid. On every regular source chart, the degree-zero line has
constant factors of automorphy. In the corresponding additive-source,
multiplicative-line cover its group inverse is \((z,V)\mapsto(-z,V^{-1})\).
A source-translation line lift has the form \((z,V)\mapsto(z+s,cV)\): its ratio
with this flat lift is a holomorphic function on the compact elliptic
curve and hence is constant. These are actual translations of the same
commutative extension. Inversion conjugates the complete suspension B,
including its scalar, to \(B^{-1}\). This gives inversion of the local compact
torus quotient with its chosen origin.

\item Its extension at the retained II fibre has a direct Cartier formula.
Write \(i(x,y,t)=(x,-y,t)\), and let e be the canonical meromorphic frame of
\(M=\mathcal O(G-O)\), whose divisor is \(G-O\). On the ordinary source charts,
including the retained II disc and excluding the finite IV-star cap
and the selected cusp,
\[
 f=x-t^2,\qquad \operatorname{div}(f)=G+(-G)-2O,
 \qquad M\otimes i^*M\longrightarrow\mathcal O,
                 \quad e\otimes i^*e\longmapsto f.       \tag{FA.4}
\]
The nonzero-vector inverse is therefore literally
\[
                (p,V e(p))\longmapsto
                    \left(i(p),\frac{i^*e(p)}{f(p)V}\right).
                                                               \tag{FA.5}
\]
It is a map of the actual line bundles: the divisors in (FA.4) cancel the
apparent frame poles and zeros. Its square is the identity since \(f\circ i=f\).
At \(t=1\) the cusp has \(x=y=0\) and \(f=-1\), so the same map and its inverse extend
there. The rigidification arrow is explicit as well. In the source
parameter \(\tau=-x/y\) at O, the regular frame is \(\tau e\); since
\(\tau(i(p))=-\tau(p)\) and \(\tau^2 f\) tends to 1, the pairing (FA.4)
has value -1 on this regular frame at O. If the actual chosen origin
vector has coefficient \(v_i(q)\) in that frame, (FA.5) sends it to
\(-1/v_i(q)\). Compose (FA.5) with the extending physical scalar
\(-v_i(q)^2\). The resulting \(I_i\) fixes that exact origin and is its
group inverse. For the regular-frame origin \(v_i=1\) this factor is
literally -1. These are actual extending units, carried through the
line and all overlap maps; no parameter winding is discarded.

\item On the good root \(t=z^{3k},\ x=z^{4k}X,\ y=z^{6k}Y\), carry first the already proved corrected root-line map of F29/K37,
including its resolved-divisor correction and finite character. Its target
line is \(\mathcal O(G'-O')\), with \(G'=(z^{2k},1)\). The divisor pairing in (FA.4)
is, after removal of the displayed fibre power \(z^{4k}\), exactly
\[
                            f_{\rm root}=X-z^{2k}.        \tag{FA.6}
\]
The same regular-frame calculation at O supplies the identical
explicit factor \(-v_i(z)^2\) for its actual root origin. This defines
the origin-fixing inverse on the smooth root family. All original
finite line characters remain in the affine deck action. In particular
we do not claim that inversion commutes with an affine cap action: its
failure to do so is precisely the finite-group equation below. The
base-root lift of an actual base-preserving F is obtained by the literal
fibre product with \(z^m=t\) and its normalization; it fixes z. Its affine
linear part on the root torus is the same \(\epsilon I\).

\item At the selected cusp use the complete periods (OC.5), including every
coefficient function, and define on the toric covering
\[
                   I_c(Z,V,q)=(Z^{-1},V^{-1},q).           \tag{FA.7}
\]
It conjugates each full h and B period to its inverse. On the fan lattice
its map is \((v,h)\mapsto(-v,h)\). For the basis fan with triangles
\(\{0,u,v\}\) and \(\{u,v,u+v\}\), the negative of the first triangle, translated by
\(u+v\), is the second, and the negative of the second, translated by \(u+v\),
is the first. Thus this is an exact automorphism of the same infinite
periodic fan, compatible with every period and its full unit factors.
It extends on every toric chart, including non-semifree fans. This is an
actual local inverse map; no reflection to a different fan is substituted.
At every additional smooth clutch disc the ordinary local inverse above
is used, with the given local origin.

\item Suppose now the actual F had linear part \(-I\). Composing its local lifts
with these inversions gives identity-linear maps of the local models.
Their values at the chosen origin give extending translations \(s_i\). At
II and on the good root their line maps extend by (FA.4)--(FA.6) and the
actual Hom evaluation argument of OC.1. At the cusp, the origin lies in
the open central torus. Its image is reindexed by the literal h,B periods
using (OC.8); evaluation of the map and its inverse then gives the full
bounded units of (OC.7). Those units extend on every cone. These \(s_i\) are
therefore sections of the actual sheaf A in OCS, including the line and
all normal collars.

\item Write an ordinary overlap in precisely the OCS convention as
\[
                g_{ij}(x)=L_{ij}x+\theta_{ij},
                         \qquad F_i(x)=-x+s_i.
\]
The literal equation \(F_i g_{ij}=g_{ij}F_j\) is
\[
                       s_i-L_{ij}s_j=2\theta_{ij}.        \tag{FA.8}
\]
There is also the full finite-deck square. If the actual root arrow is
\(A_z(x)=L_D(z)x+a_D(z)\), from z to Dz, its square with F is
\[
                  s_{Dz}-L_D(z)s_z=2a_D(z).              \tag{FA.9}
\]
This retains the complete cyclic group of order m and every central line
character, including the ineffective subgroup. Taking all iterates of
(FA.9) retains the original m-th norm; it does not replace m by 3.
Equations (FA.8)--(FA.9) say exactly that twice the actual total translation
cocycle is the coboundary of these extending sections. Passing through
the same atlas exponential sequence (OCS.4), with its full integral
finite-group cochain, therefore gives
\[
                2(D,d\bmod k)=0
                  \quad\hbox{in }\mathbb Z\oplus\mathbb Z/k.     \tag{FA.10}
\]
But \(D=\pm1\) in every sphere datum under discussion. Its free component
\(2D\) is nonzero. Thus a linear \(-I\) automorphism cannot exist.

\par

\medskip\noindent\textbf{FA.4. The remaining group and endpoint consequence}

\par\nobreak

\item Every automorphism is now base preserving and identity linear, so OC.2--3
apply with their actual source/line and toric lifts. K78's explicit
clutch comparison transports that same group for the additional smooth
sites. The full calculation gives a rational lift \(f(t)U^j\), forces \(j=3a\),
and identifies it literally with \(cB_0^{-a}\); on the compact quotient it
is the physical scalar \(c\lambda^a\). Conversely every physical nonzero
scalar extends through every root and every toric chart, and acts
faithfully. Consequently
\[
                      \operatorname{Aut}_{\rm hol}(X)
                                =\mathbb C^*_{\rm physical}.     \tag{FA.11}
\]
A path from 1 to any c in this actual group is an isotopy of the full
complex model, its finite caps and its normal collars.

\item For a fixed closed family with two specified holomorphic endpoint maps,
the ratio of those endpoint maps is an automorphism of the same fixed
sphere member X. Equation (FA.11) supplies a literal isotopy between
them by physical scalar multiplication. Their inverse-endpoint physical
returns are therefore conjugate up to this isotopy and have the same
oriented clutched-sphere class. This conclusion concerns endpoints of the
same fixed data; it does not identify two different norm signs, different
source-moduli paths, or an arbitrary smooth endpoint with a holomorphic
one. No existence of an unspecified short endpoint is inferred.

\item The independent centralizer, characteristic-polynomial, Cartier, root,
fan and finite-cochain identities are recorded in
\path{code/part3/os49_full_automorphisms.py} and its captured output. The complete local inverse and root-groupoid argument have passed the independent audits recorded with this export.

\par

\medskip\noindent\textbf{FA.5. A prescribed scalar path with an arbitrary holomorphic endpoint}

\par\nobreak

\item Retain the actual fixed-source translation/line data of K44 and the fixed
smooth-site data of K78, with \(D=\pm1\). Let a cut scalar path have integer
winding L, and let an actual holomorphic endpoint between its two models
be specified. We now determine that endpoint's linear part without
assuming a translation marking.

\item On the ordinary source complement the endpoints have the same scalar
\(\lambda\) and hence the same literal quotient by \(B_\lambda\). Use this common
source/line-cover marking to compare their smooth compact tori. The
endpoint descends to the identity of the base by FA.1, and its integral
matrix commutes with the same \(A,B,C\). The rank-one argument of FA.2 again
leaves only \(+I\) and \(-I\). The analytic continuation of a chosen period
logarithm along the scalar path is explicitly
\[
                      S_L=I+L\delta\psi.                 \tag{FA.12}
\]
This preserves \(\psi\) and takes \(\ell_0\) to \(\ell_0+L\delta\). Thus every finite
norm vector and every transported ordinary vector retains its free
coordinate in (OCS.13). The actual total groupoid class at both endpoints
has free coordinate D. Its delta coordinate and every finite character
are retained, not equated by this assertion.

\item For an endpoint of linear part \(-I\), the literal affine square between
the two local atlases is now
\[
                    s_i-L_{ij}s_j
                         =\theta^{\rm end}_{ij}
                             +\theta^{\rm start}_{ij},    \tag{FA.13}
\]
with the identical sum in the full finite-deck square. The same local
extensions in FA.3 apply. Its integral free obstruction is therefore
\(D+D=2D\), nonzero. Thus the actual endpoint necessarily has linear part
\(+I\). The interval comparison of K78 consists of translations and hence
carries it to a translation endpoint of the canonical one-cap datum
\((m,D,d)\), with all scalar-origin and finite root terms retained.

\item K44, OC.15--19 now applies to this actual endpoint. Its forward root
square and the rational source-line divisor give
\[
                  3v-j=\frac{L D}{k},\qquad
                         L=N k\quad(N\in\mathbb Z).      \tag{FA.14}
\]
Here j is the actual source-translation coefficient and v its integral
cusp line order. The cap character, including its full norm, is the one
in OC.16--18; the source and line factors are not replaced by winding
numbers alone. The canonical OR65 endpoint and its powers realize every
\(L=Nk\), including negative N through the actual inverse map. K78 transports
those maps through the full family comparison.

\item For a given realized \(L=Nk\), compare any specified endpoint with that
transported canonical endpoint. Their ratio is an automorphism of the
same fixed sphere model; by (FA.11) it is an actual physical scalar.
Its isotopy removes that ratio from the inverse-endpoint physical
return. Thus every such holomorphic endpoint has the return
\[
                                N S_k(P).                \tag{FA.15}
\]
In particular the minimum positive scalar winding for this full
holomorphic endpoint scope is k. This conclusion is obtained from the
actual normalizer, affine and finite-root squares and the explicit
canonical endpoint; it is not a division of a class in \(\Theta_7\). A family
with a changed source or an arbitrary smooth, nonholomorphic endpoint
has not been identified by this argument.

\end{enumerate}
\end{proof}




\begin{theorem}\label{night:P3K.80}
\StatementRecord{night:P3K.80}{proved}{S6\_06\_BLIND: GPT-6 Astra (OpenAI, ChatGPT Pro mode, effort Pro), recorded re-derivation in the checkpoint dated 18 November 2026}{\cref{night:P3F.0}, \cref{night:P3F.22}, \cref{night:P3K.23}, \cref{night:P3K.24}, \cref{night:P3K.42}, \cref{night:P3K.58}, \cref{night:P3K.60}, \cref{night:P3K.62}, \cref{night:P3K.68}, \cref{night:P3K.69}, \cref{night:P3K.72}}
Retain the marked source
\[
 E_v:\quad Y^2=X^3+(1-2v)X^2+TX+v^2,\qquad G=(0,v),
\]
the actual Cartier line \(M_2=O([2G]-[O])\), the suspension
over \(-G\), and the selected whole-zero fibre at
\(T_0(v)=v^2-2v\). Retain the reduced III cap of order
\(m=4p-1\), \(p\ge1\), either norm sign \(n=\pm1\),
and every admitted integral \(d\delta\) character. These
are the sphere families and exact cap markings of P3K.24,
with the physical line convention of P3K.68.

For a fixed sufficiently small nonzero scalar, the positive
\(v\)-loop has physical Jin integer \(0\) in the initial
negative-square fan \(T_0\). Its actual oriented clutched-sphere
class is \(0\in\Theta_7\simeq\mathbb Z/28\). The literal
reflected fan \(JT_0\) has actual class \(2\). For the
actual coupled family \(r=v^k r_0\), \(\lambda=r^m\),
\(k\ge1\), the two classes are
\[
 (Z_{T_0},Z_{JT_0})
   =\bigl(k(3p-1),\ 2+k(3p-1)\bigr)\pmod {28}.
 \tag{P3K.80.RA.1}
\]
The reflected value uses the actual directed relative flop.
This statement treats the reduced \(a=2,b=1\) branch.
\end{theorem}
\begin{proof}
\begin{enumerate}[label=\arabic*.]
\item The proof constructs a finite real bordism carrying the
actual physical circle, its associated line, and the
canonical complex spin-c arrows on every local overlap.
The fixed-scalar loop is separated from the scalar winding
by the actual two-parameter family before the cut is made.
All local source, line, node and finite-root maps needed
for the cut are reproduced below. The local equation
prefixes identify the component calculations within this
single proof. In this proof \(a=2,b=1\), except in the
displayed weight identities which retain \(a\) to show
the actual character. The Tate period is denoted
\(H_T=(Q,P^2)\); the finite-cap unit is
\(h_\infty=1-T_0(v)q_\infty\). They are distinct
explicit maps/functions throughout.

\par\medskip\noindent\textbf{RA.2. Source, Cartier and actual Tate data.}\par\nopagebreak[4]

\item Put \(q=T-(v^2-2v)\).  On the source threefold set
\[
 A=Y-v,\quad B=Y+v,\quad C=X,\quad
 D=T+X^2+(1-2v)X .
\]
The source equation is literally \(AB=CD\).  Its marked small
resolution has the chart
\[
 A=Dz,\qquad C=Bz,\qquad v=(B-Dz)/2.
\]
The chosen section \(G\) has \(z=0\).  The actual closure of
\(2G\) has Cartier equation
\[
 F_2=1-Bz^3-(1-2v)z^2=0.
\]
Indeed the doubling formula gives \(z=-2v/T\) and
\(zT+2v=BF_2\); the factor \(B=0\) is the other section
\(-G\), so the strict closure is \(F_2=0\).
Define the actual target-point primitive on \(O(G-O)\) by
\[
 U:(p,w)\longmapsto\left(p-G,\frac{w}{X(p-G)}\right).
 \tag{P3K.80.SZ.7}
\]
The identity
\((G)-O-t_G^*((G)-O)=\operatorname{div}X\)
makes this the indicated line map; its tensor-square is the
primitive on \(O(2G-2O)\).  The Cartier conjugation below
identifies that tensor-square with the specified actual
line \(M_2\).  In that conjugation the retained rational
function is
\[
 F_{21}=\frac{Y+v+(T/(2v))X}{X}.
\]
Its displayed poles are carried by the rational line map;
the resolved Cartier arrows below are the extensions used
on the finite cut.

\par\smallskip\noindent\textit{The source graph and its Cartier coefficient.}\par\nopagebreak[4]

\item The chamber overlap is the following literal map.
Write \(p'=p-G\), and let \(F_2=1-Bz^3-(1-2v)z^2\) be the actual
Cartier equation just displayed. The addition formula gives
\[
 X'=\frac{F_2}{z^2},\qquad Y'=v-\frac{X'}z.
\]
Set \(A'=Y'-v,\ B'=Y'+v,\ C'=X'\), and
\(D'=T+X'^2+(1-2v)X'\). Direct substitution gives
\[
                   A'=-C'/z,\qquad D'=-B'/z.                \tag{P3K.80.SZ.14}
\]
These are exactly the equations of the other small resolution of
\(A'B'=C'D'\). In particular the map contracts the proper component
\(C_0\), where \(F_2=0\), and the ratio of the other ruling is
\(-1/z\). The marked target small resolution uses the first ruling.
This identifies the graph of translation by \(-G\) with the actual
Atiyah flop in these coordinates.

\item Use its common Cox coordinates
\[
 A'=s_0t_0,\quad B'=s_1t_1,\quad C'=s_0t_1,\quad D'=s_1t_0,
\]
with auxiliary weights \(+1\) on the \(s_i\) and \(-1\) on the
\(t_i\). The target ruling has coordinate \(s_0/s_1\), and its marked
section \(G\) has \(s_0=0\). The domain ruling has coordinate
\(t_0/t_1=-1/z\), and its zero section \(O\) has \(t_1=0\).
Let \(V\) have auxiliary weight \(+1\). On the target chart the
canonical meromorphic coefficient of the line \(O(G-O)\) is
\[
                             w_{\rm target}=V/s_0;
\]
on the domain chart it is
\[
                             w_{\rm source}=Vt_1.
\]
For clarity, on \(s_1=1\) the regular line frame is the divisor
frame \(1/(s_0/s_1)\), whereas on \(t_0=1\) it is the divisor
frame \(t_1/t_0\). These two displayed coefficients are exactly
their expressions in the canonical meromorphic frame. Hence
\[
                 \frac{w_{\rm target}}{w_{\rm source}}
                       =\frac1{s_0t_1}=\frac1{X'},           \tag{P3K.80.SZ.15}
\]
which is the actual multiplier in (P3K.80.SZ.7). Thus the source line map
through the flop is the identity in the common Cox coordinate \(V\),
not merely an abstract equality of line degrees. Fixing the nonzero
\(V\) gives the two moment chambers of the same \(C^4\) model.
For the \(a\)-th tensor power it gives the same two chambers of
\(C^4/\mu_a\), retaining the physical delta action already recorded.

\item The actual degree-two line has a compatible explicit Cartier map
to the square of the degree-one line. Namely
\[
 r(p)=X(p-G)=\frac{2vF_{21}(p)}{X(p)}
              =\frac{F_2}{z^2}.                            \tag{P3K.80.SZ.16}
\]
Its divisor is the difference between the actual \(2G-O\) divisor
and twice \(G-O\), on this collision neighbourhood. Multiplication
by \(r\) is consequently the line isomorphism
\(O(2G-O)\longrightarrow O(G-O)^{\otimes2}\). The expression
\(F_2/z^2\) retains the central component and has no extra factor of
\(v\). At a target point \(p'=p-G\),
\[
 F_{21}(p')=\frac{r(p')r(p)}{2v}.
\]
Transporting the actual suspension through (P3K.80.SZ.16) therefore gives
\[
 \frac{r(p')}{r(p)}\frac{\lambda q}{vF_{21}(p')}
                         =\frac{2\lambda q}{X(p')^2}.       \tag{P3K.80.SZ.17}
\]

\item Let \(U_a\) denote the tensor-power version of (P3K.80.SZ.7), with target
multiplier \(X(p-G)^{-a}\). Both allowed source/line sectors now have
the exact common form
\[
           B_{a,b}=2(-1)^{b+1}\lambda q\,U_a^b,
                       \qquad ab=2.                        \tag{P3K.80.SZ.18}
\]
The source Cox chamber overlap for every factor \(U_a\) is (P3K.80.SZ.15).
The finite construction below uses these actual chambers, the full selected factor, and their finite-root maps.
Its outside boundary is identified with the actual family by the mixed-face maps below.

\item In these same target Cox coordinates the two base functions are
explicit:
\[
 v=\frac{s_1t_1-s_0t_0}{2},\qquad
 q=(s_1-s_0)(t_0+t_1)-(s_0t_1-v)^2.                         \tag{P3K.80.SZ.19}
\]
Both are invariant under the physical delta action. In particular
the \(q\)-factor in (P3K.80.SZ.18) is the actual invariant polynomial
(P3K.80.SZ.19); a product construction replacing it by an independent
disc coordinate would omit its singular central divisor.
The certificate checks this identity as well as the two chamber
line frames. This retained equation is the local input for the
periodic Mumford-boundary comparison.

\par\smallskip\noindent\textit{The exact two-node base coordinates.}\par\nopagebreak[4]

\item There is a direct analytic coordinate map to the two-node conifold,
without dropping any parameter unit. Put \(A_0=1-2v\) and
\[
 \xi(T,v)=\frac{-A_0+\sqrt{A_0^2-3T}}3 ,
\]
where the square root is the unit branch equal to one at the origin.
It satisfies \(3\xi^2+2A_0\xi+T=0\). For \(x=X-\xi\), define
\[
 R=x\sqrt{A_0+3\xi+x},\qquad
 h(T,v)=v^2-A_0\xi^2-2\xi^3 .
\]
Both square roots are actual holomorphic units on a sufficiently
small polydisc. The source equation becomes exactly
\[
                         (Y-R)(Y+R)=h(T,v).                \tag{P3K.80.SZ.20}
\]
The two local critical-value curves are
\[
 T_0=v^2-2v,\qquad
 T_1=\frac{1+4v-(1-8v)^{3/2}}8 .
\]
The latter is the small root of the remaining discriminant factor.
Their tangent slopes are \(-2\) and \(+2\); hence
\[
                 q_1=T-T_0(v),\qquad q_2=T-T_1(v)
\]
are actual analytic coordinates on the base near the origin.
The function \(h\) vanishes simply along both curves away from their
intersection. The explicit formula shows its quadratic term is
\(-\tfrac14(T^2-4v^2)\). Analytic division therefore gives
\[
                       h=\upsilon(q_1,q_2)q_1q_2,
                         \qquad\upsilon(0,0)=-1/4,          \tag{P3K.80.SZ.21}
\]
with \(\upsilon\) holomorphic and nowhere zero after shrinking.
Replacing \(Y-R\) by \((Y-R)/\upsilon\) in (P3K.80.SZ.20) gives the literal
equation \(AB=q_1q_2\). This is an explicit local isomorphism, including
the selected \(q_1\) direction, the other \(q_2\) direction and the
unit \(\upsilon\). The unit has a holomorphic logarithm on the
polydisc, so its restriction to a source-parameter loop has zero
winding.

\par\smallskip\noindent\textit{The full Tate periods and units.}\par\nopagebreak[4]

\item The periods used in the finite construction are obtained from the actual
Tate coordinates on the full two-node base. The source has
\[
 c_4=16((1-2v)^2-3T).
\]
This is a unit at the origin. Its discriminant has the two simple
factors \(q_1q_2\) just defined. Inverting the convergent \(j\)-series
therefore gives
\[
 Q=q_T=q_1q_2u_Q,\qquad u_Q(0)=1/256.                 \tag{P3K.80.SZ.39}
\]
The indicated value follows directly from
\[
 j^{-1}=
 \frac{(1-2v)^2T^2-4T^3-4(1-2v)^3v^2-27v^4+
          18(1-2v)Tv^2}
 {256((1-2v)^2-3T)^3}.
\]
All units in this section are holomorphic and nonvanishing on a
smaller two-node polydisc; their logarithms are fixed there.

\item Here is the marked section map, including its lift rather than just
its class in the Tate group. The convergent Tate Weierstrass sums
extend over the node chart with coordinates \(Z_0,W_0\) and
\(Q=Z_0W_0\). In these coordinates their first terms are
\(x_T=Z_0+W_0+\text{higher terms}\) and
\(y_T=-W_0+\text{higher terms}\). They consequently give the
invertible analytic node map to the Weierstrass node, after the
chosen \(c_4\)-root scaling and the Weierstrass translation. This is
also seen directly by writing each negative-index summand as a
power series in \(W_0\) and \(Q\); the positive-index summands are
power series in \(Z_0\) and \(Q\).

\item The actual section \(G\) lifts through this node chart. Its two
coordinates vanish on opposite discriminant branches. Choose the
branch which, on the selected \(q_1=0\) normalization, represents
the lift \(Q P_G\). The linear terms of the explicit Morse map in
the preceding two-node chart are \(v+T/2\) and \(v-T/2\); hence those two vanishing
orders are one. Analytic division gives
\[
 Z_0(G)=q_1u_G,\qquad W_0(G)=q_2u'_G,\qquad
 P=P_G=\frac{Z_0(G)}Q=q_2^{-1}u_P.                   \tag{P3K.80.SZ.40}
\]
The branch is checked against the exact selected nodal value
\((1+\sqrt{1+v})/(1-\sqrt{1+v})\). Since
\(q_2|_{q_1=0}=-4v+O(v^2)\), it gives \(u_P(0)=16\).
Thus (P3K.80.SZ.40) retains the actual section and the marked period
lift, including its integral reindexing.

\item The line and its primitive translation have an equally direct
formula. Set
\[
 f_P(Z)=\frac{\Theta(Z/P)}{\Theta(Z)}.
\]
The exact divisor and multiplier identify \(f_P\) with the
canonical meromorphic section of the actual Cartier line. The
function \(X\), whose divisor is \(G+(-G)-2O\), consequently has
the literal expression
\[
 X(Z)=c_X f_P(Z)f_{P^{-1}}(Z).                       \tag{P3K.80.SZ.41}
\]
To see its scalar at the collision, evaluate at a fixed ordinary
Tate point, for example \(Z=-1\). The identity
\(\Theta(QZP)=-(ZP)^{-1}\Theta(ZP)\) shows that
\(q_2 f_P(-1)f_{P^{-1}}(-1)\) is a nonzero holomorphic
unit. The actual value \(X(-1)\) is a holomorphic unit there.
Hence
\[
                         c_X=q_2u_X,\qquad u_X(0)=-1/4 .
                                                               \tag{P3K.80.SZ.42}
\]
For the last value, the limiting source is
\(Y^2=X^3+X^2\), so the retained scaling is
\(X=4Z/(1-Z)^2\), and (P3K.80.SZ.41) gives \(c_X\sim-4/P\).
No rational frame is discarded in this evaluation.

\item The target-point primitive lift (P3K.80.SZ.7) is now exactly
\[
 U_a:(Z,V)\longmapsto(P^{-1}Z,c_X^{-a}V).             \tag{P3K.80.SZ.43}
\]
Indeed
\[
 f_P(Z)f_{P^{-1}}(Z/P)=1,
\]
so substituting (P3K.80.SZ.41) into the source and target Cartier
coefficients cancels every \(Z\)-dependent factor. For \(a=2\)
the literal map (P3K.80.SZ.16) is applied first. The resulting identity
retains its central Cartier component.

\item Combining (P3K.80.SZ.18) and (P3K.80.SZ.43), the two actual periods are
\[
 \begin{split}
 h:(Z,V)&\longmapsto
       (QZ,P^aV),\\
 g:(Z,V)&\longmapsto
       (P^{-b}Z,\,
        2(-1)^{b+1}\lambda_0v^Nq_1c_X^{-2}V).
 \end{split}                                                   \tag{P3K.80.SZ.44}
\]
Their valuation matrices along \(q_1=0,q_2=0,v=0\) are precisely
\[
 I,\qquad W=\begin{pmatrix}1&b\\-a&-2\end{pmatrix},
 \qquad E=NE_{22}.
\]
Every remaining factor in (P3K.80.SZ.44) extends as a unit. This proves
the actual period/unit input to the sector fan construction; the
earlier winding calculation alone was not used as that input.

\par\medskip\noindent\textbf{TL.1. The full marked source map.}\par\nopagebreak[4]

\item Use the convergent Tate node map from SZ.39--40, with its chosen
Weierstrass scaling, both inverse branches, and marked lift:
\[
 Z_0W_0=Q=q_1q_2u_Q,\qquad
 P=P_G=q_2^{-1}u_P,\qquad u_G=u_Qu_P,
\tag{P3K.80.TL.1}
\]
where \(u_Q,u_P\) are actual nonvanishing holomorphic functions on
the full two-node base and
\[
 u_Q(0)=1/256,\qquad u_P(0)=16,
\qquad Z_0(G)=q_1u_G,\quad W_0(G)=q_2/u_P.
\]
The source small-resolution chart is given by the literal maps
\[
 S=P W_0,\qquad T_0=Z_0/u_G,
\]
\[
 Z_0=u_GT_0,\qquad W_0=q_2S/u_P,
\qquad q_1=ST_0.                                    \tag{P3K.80.TL.2}
\]
This is the chart of the resolution of \((W_0,q_2)\) in which the
ratio \(W_0/q_2\) is finite. Its other projective chart retains
\(S'=S^{-1}\), \(T'_0=ST_0=q_1\), and \(T'_1=S q_2\), so
\[
 q_2=S'T'_1,\qquad
 Z_0=u_G S'T'_0,\qquad W_0=T'_1/u_P.
\tag{P3K.80.TL.3}
\]
These are inverse coordinate charts of the same proper small
resolution; they agree by the displayed substitutions. In particular
\(S(G)=1\) holds exactly on the whole section, not only to first
order. The two branches in the selected source bidisc are
\(S=0\) and \(T_0=0\). Both are retained in (P3K.80.TL.2).

\item Here is the arrow from the earlier node-adapted Morse chart. In
SZ.20--21 put \(A_M=(Y-R)/\upsilon\), \(B_M=Y+R\), so
\(A_MB_M=q_1q_2\). Write
\(A_M(G)=q_2a_G\), with \(a_G\) an extending unit. The actual
base-unit change
\[
 \alpha_M=-A_M/(2a_G),\qquad
 \beta_M=-2a_GB_M                                  \tag{P3K.80.TL.4}
\]
has product \(q_1q_2\), and sends the marked section to
\(\alpha_M(G)=-q_2/2\), \(\beta_M(G)=-2q_1\). The analogous
Tate functions are
\[
 \alpha_T=-u_PW_0/2,\qquad \beta_T=-2Z_0/u_G,
 \qquad \alpha_T\beta_T=q_1q_2.                     \tag{P3K.80.TL.5}
\]
They are literal holomorphic functions of the Weierstrass point by
the convergent Tate inverse map.

\item The selected branches in SZ.40 choose \(W_0\) to match the
\(\alpha_M\) branch. More explicitly the two maps preserve the
nodal point on each of the smooth punctured discriminant branches;
their nonsingular linear normal maps match \(\alpha_M\) with
\(W_0\), rather than with \(Z_0\). Hence the Cartier zero divisors
of \(\alpha_T\) and \(\alpha_M\) have the same two irreducible
components and multiplicity one. They have no further zero divisor
on a smaller conifold neighbourhood. The exact quotient
\[
                      b_T=\alpha_T/\alpha_M          \tag{P3K.80.TL.6}
\]
is therefore a holomorphic unit on the normal conifold: both the
quotient and its inverse have nonnegative valuation at every
prime divisor. The chosen linear maps give \(b_T(0)=1\), and
the marked evaluations give \(b_T(G)=1\).

\item In the earlier projective chart
\(\alpha_M=-q_2S_M/2\), \(\beta_M=-2T_M\), this supplies the
actual resolution lift
\[
             (S_M,T_M,q_2)\longmapsto
                    (b_TS_M,b_T^{-1}T_M,q_2).        \tag{P3K.80.TL.7}
\]
Its inverse is the inverse of the already specified analytic
Weierstrass/Tate maps; the ideal \((\alpha_M,q_2)\) is carried
to \((\alpha_T,q_2)\) because their ratio is a unit. Thus the
map lifts on both projective charts. On the exceptional projective
line its unit is one. After shrinking, its logarithm is the actual
holomorphic logarithm near one. No parameter winding has been
removed by this lift.

\item For the original Cartesian section equation, the precise additional
Cartier factor is its pullback divided by \((1-S)/2\). Both have
the same marked divisor with multiplicity one under the full map,
so this quotient is an extending unit. It is retained on the source
side, with its inverse on the target side. This defines its full
higher-normal terms; none is replaced by its exceptional value.

\par\medskip\noindent\textbf{TL.2. The exact line map on both source branches.}\par\nopagebreak[4]

\item Use the same theta convention as SZ.41:
\[
 f_P(Z)=\Theta_Q(Z/P)/\Theta_Q(Z),\qquad
 \Theta_Q(z^{-1})=-z^{-1}\Theta_Q(z),\quad
 \Theta_Q(Qz)=-z^{-1}\Theta_Q(z).
\]
The Tate representative on the source chart (P3K.80.TL.2) is
\[
                         Z=P/S=1/W_0.                \tag{P3K.80.TL.8}
\]
It is the original representative obtained from \(Z_0\) by the
literal inverse source period. Theta inversion gives the exact
identity
\[
 f_P(P/S)=P^{-1}\frac{\Theta_Q(S)}{\Theta_Q(S/P)}.
\tag{P3K.80.TL.9}
\]
The tensor Cartier coefficient in the actual marked divisor frame
is
\[
       \xi=\left(\frac{1-S}{2}\right)^a
                      \frac{V}{f_P(P/S)^a}.          \tag{P3K.80.TL.10}
\]
Applying the actual source period
\(h:(Z,V)\mapsto(QZ,P^aV)\) yields
\[
 Z_0=QP/S=u_GT_0,\qquad
 V_0=P^aV=E_+(S,T_0,q_2)^a\xi,
\]
\[
 E_+=\frac{2\Theta_Q(S)}{(1-S)\Theta_Q(S/P)}.        \tag{P3K.80.TL.11}
\]
Every function in this equality is the full function, including
the actual \(P\) and the source-period line factor \(P^a\).

\item The factor in (P3K.80.TL.11) is a holomorphic unit on the entire source
bidisc, including both branches. Indeed the convergent product
for \(\Theta_Q(S)/(1-S)\) has factors in \(Q^rS\) and
\(Q^r/S\), and
\[
 Q^r/S=S^{r-1}T_0^r q_2^r u_Q^r\quad(r\ge1).
\]
For the denominator use
\[
 S/P=W_0=q_2S/u_P,\qquad Q/(S/P)=Z_0=u_GT_0.
\]
Thus its product factors are holomorphic in \(Z_0,W_0,Q\).
On a sufficiently small fixed bidisc all of them are nonzero,
and the products converge normally. On \(q_2=0\), the numerator ratio equals one, whereas the
denominator retains its factor \(1-Z_0\). Thus the exact restrictions
are
\[
 E_+(S,T_0,0)=\frac{2}{1-u_G(ST_0,0)T_0},\qquad
 E_+(S,0,0)=2.                                       \tag{P3K.80.TL.12}
\]
This also handles the removable zero at the marked section
\(S=1\). In particular the source line map and its inverse
extend; a logarithm of \(E_+/2\) extends through the full collision
base. Equation (P3K.80.TL.11), rather than only its value (P3K.80.TL.12), is the
line map used on the collar.

\item On the other projective source chart the coefficient is
\(\xi'=S^{-a}\xi\), because the marked equation becomes
\((S'-1)/2\). Substitution of \(S=1/S'\) in (P3K.80.TL.9)--(P3K.80.TL.11)
gives the same Cartier map wherever the two charts overlap.
This is a statement about the two line charts of the actual
degree-\(a\) line. It does not declare the constant-factor Tate
frame to extend unchanged over that entire projective line.

\par\medskip\noindent\textbf{TL.3. The actual fan label.}\par\nopagebreak[4]

\item For fixed nonzero \(q_2\), the selected node has source coordinates
\((S,T_0)\) with \(ST_0=q_1\). Its associated node and first
small resolution are the actual equations
\[
 U V_{\rm node}=S T_0,\qquad U=S z,\quad T_0=zV_{\rm node},
\]
with \(U=\xi\) after the Cartier arrow just written. In the
characters \((S,U,q_1)\) its two cones have positive-square
triangles
\[
 \{0,e_1,e_1+e_2\},\qquad\{0,e_2,e_1+e_2\}.
\]
The exact map (P3K.80.TL.11) has valuation map
\[
         (x,y,h)\longmapsto(h-x,y,h).                 \tag{P3K.80.TL.13}
\]
The units \(u_G,E_+\), and the preceding Morse/Cartier units
are nonvanishing on the full source bidisc. Hence their pullbacks
are nonvanishing on all of these toric charts. A monomial generator
of each cone semigroup is carried to its stated monomial times an
explicit such unit, and the inverse has the same property. This
proves extension of the character map to the full matched cones.

\item On height one the two triangle images are
\[
 \{e_1,0,e_2\},\qquad\{e_1,e_1+e_2,e_2\}.
\]
They are precisely the two triangles of current Tate \(T_0\),
the negative diagonal. The opposite associated resolution maps
to current Tate \(T_1\). The terminology \(j=-1\) in the
earlier SZ.15 uses its earlier square convention; (P3K.80.TL.13) supplies
the actual character map and does not identify those index symbols.
The physical circle acts on \(\xi\) and \(V_0\) with the same
positive weight. The analytic map preserves the complex orientation.

\par\medskip\noindent\textbf{TL.4. Full period and closing arrows.}\par\nopagebreak[4]

\item The source change above is the actual identity of the marked
Weierstrass source expressed by two specified analytic coordinate
systems. The Cartier ratio is the canonical theta/divisor ratio.
Consequently the primitive translation is transported to precisely
\(U_a:(Z,V)\mapsto(P^{-1}Z,c_X^{-a}V)\) of SZ.43, and the
full suspension is still
\[
 B=2(-1)^{b+1}\lambda q_1 U_a^b,\qquad ab=2.
\]
The equality follows on the common smooth torus from the exact
theta identity used in SZ.43 and extends wherever both charts are
defined. The source period in (P3K.80.TL.11) is an actual quotient deck,
including its line factor, not a replacement of the period by a
valuation vector.

\item On a source-parameter annulus the full functions in (P3K.80.TL.1)--(P3K.80.TL.11)
are single valued. The functions \(P\) and \(P^a\) retain their
nonzero windings before the displayed period reindexing. The
remaining factors have extending logarithms. An ordinary parameter
closing \(q_1^{\nu k}\) commutes with these source/line maps and
is retained literally. At \(a=2\), the same assertion includes
the source and target \(r\) factors. Thus the calculation labels
the selected boundary chart with its actual parameter arrow.

\item The separate cyclic reset of EA.18--23 enters the complementary
charge-\(a\) chart, retains the physical normalization
\(c_\lambda=2\lambda\), and uses an adjacent Tate-node overlap.
The present local comparison does not remove that reset, its
parameter winding, or any transported truncation. In particular
neither a compact collision-to-cap attachment nor an absolute
source-zero return is inferred from the fan label alone.

\par\medskip\noindent\textbf{RT.1. The horizontal trace immersion.}\par\nopagebreak[4]

\item The marked small-resolution coordinates are
\[
 A=Dz,\qquad C=Bz,\qquad v=(B-Dz)/2.
\]
On the normalization of the selected nodal fibre,
\[
 X=s^2-1,\quad Y=s(s^2-1-v),\quad T=v^2-2v,\quad
 z=(s+1)/(s^2+s-v).
\]
Put \(x=s+1,v=xr\). The blowup chart has the exact map
\[
 z=(x-1-r)^{-1},\quad B=x(x-2)(x-1-r),\quad
 D=(B-2xr)/z.
\]
The Jacobian of \((z,B)\) with respect to \((x,r)\) is
\(-2(x-1)/(x-1-r)\), hence is nonzero on the exceptional curve wherever
this chart is finite. At its missing point use
\(w=x-1-r,C=x(x-2)\); the Jacobian is \(2x-2\), again a unit there.
The second blowup chart \(v=y,x=yh\) has
\(z=h/(-1-h+yh^2)\), whose h derivative at y=0 is
\(-1/(-1-h)^2\). These formulas give the actual immersion of the complete
horizontal normalization trace through the exceptional curve.

\item Let \(F=\operatorname{Bl}_{(-1,0)}(\mathbb P^1_s\times D_v)\) and let E
be its exceptional curve. The Cartier ratio
\((s-\sqrt{1+v})/(s+\sqrt{1+v})\) has divisor
\(D_+-D_--E\). The actual source line restricts as \(O(-aE)\).
Consequently the retained spinor determinant on this trace is
\((1-2a)E\); reversing the two selected branches gives
\(-(1+2a)E\).
The complete fixed and cap contributions are calculated in RT.16--17 below.

\par\medskip\noindent\textbf{RT.2. Exact q factorization on the Cox bridge.}\par\nopagebreak[4]

\item Use the actual Cox coordinates
\[
 A'=s_0t_0,\quad B'=s_1t_1,\quad C'=s_0t_1,\quad D'=s_1t_0,
\]
with physical weights \((-1,-1,+1,+1)\) on
\((s_0,s_1,t_0,t_1)\), divided by the retained \(\mu_a\) when appropriate.
The actual base functions, before any quotient completion, are
\[
 v=(s_1t_1-s_0t_0)/2,\qquad
 q=(s_1-s_0)(t_0+t_1)-(s_0t_1-v)^2.
\]
Set \(x=s_0,y=t_1,k=s_1-s_0,l=t_0+t_1\), and choose the analytic unit
\(\sigma=\sqrt{1+xy}\) with value 1 at the origin. Define
\[
 F_+=\frac{(\sigma+1)k-x^2l/(\sigma+1)}2,\qquad
 F_-=\frac{(\sigma+1)l-y^2k/(\sigma+1)}2.
\]
Then
\[
 q=F_+F_-,\qquad
 \det\frac{\partial(F_+,F_-)}{\partial(k,l)}=\sigma.
\]
Thus this is an actual analytic coordinate change over \((x,y)\), with inverse
\[
 k=\frac{(\sigma+1)F_++x^2F_-/(\sigma+1)}{2\sigma},\quad
 l=\frac{y^2F_+/(\sigma+1)+(\sigma+1)F_-}{2\sigma}.
\]
The weights of \(F_+,F_-\) are respectively \(-1,+1\). No physical
character is changed. The source parameter becomes
\[
 v=xy+\frac{yF_+-xF_-}{2\sigma}.
\]
On the chart with invariant coordinates \(u=x/F_+,z=y/F_-\), this reads
\[
 v=q\left(uz+\frac{z-u}{2\sqrt{1+uzq}}\right).
\]
These formulas retain both actual branches of q=0. The finite domain below retains the quotient strata over both branches.

\par\medskip\noindent\textbf{EA.17. Full two-node refinement of the Morse-Cox pairs.}\par\nopagebreak[4]

\item The parameter formula in the exact two-node chart is
\(\widetilde v=q_1-q_2=v u_v(v)\). To retain it together with
the primitive Cartier frame, use the actual analytic source
functions from SZ.20--21 and TL.4. Explicitly, put
\[
 A_0=1-2v,\qquad x_c=\frac{-A_0+\sqrt{A_0^2-3T}}3,
\]
\[
 R_c=(X-x_c)\sqrt{A_0+2x_c+X},\quad
 \upsilon=\frac{v^2-A_0x_c^2-2x_c^3}{q_1q_2},
\]
where both square roots and the extending unit upsilon are the
specified analytic branches. Define
\[
 A_M=(Y-R_c)/\upsilon,\quad B_M=Y+R_c,\quad
 a_G=A_M(G)/q_2,
\]
\[
             \alpha=-A_M/(2a_G),\qquad
             \beta=-2a_GB_M.
 \tag{P3K.80.EA.42}
\]
They satisfy \(\alpha\beta=q_1q_2\),
\(\alpha(G)=-q_2/2\), and \(\beta(G)=-2q_1\).
All functions are pulled back through the actual Cartesian source
map of SZ.19; no linear approximation is used in (P3K.80.EA.42).

\item Keep the exact selected factors F+,F- of RT.2 and define
\[
 S_0=F_+,\qquad T_0=F_-,\qquad
 S_1=-\frac{\beta}{2F_-},\qquad
 T_1=-\frac{2\alpha}{F_+}.
 \tag{P3K.80.EA.43}
\]
The quotients extend holomorphically. For example F+=0 is an
irreducible coordinate hypersurface of the RT.2 chart. The
restriction of beta to it is not identically zero, as its leading
term is \(-2(s_0+s_1)(t_0+t_1)\). Since alpha beta vanishes
there, alpha vanishes identically there and is divisible by F+.
The analogous argument gives divisibility of beta by F-. The
leading maps are exactly
\[
 (S_0,S_1)=(s_1-s_0,s_0+s_1)+O(\|(s,t)\|^3),
\]
\[
 (T_0,T_1)=(t_0+t_1,t_0-t_1)+O(\|(s,t)\|^3).
\]
Their complex Jacobian at the origin is 4. Thus (P3K.80.EA.43) is an
actual auxiliary-equivariant analytic isomorphism on a small
neighbourhood. It has
\[
                 q_1=S_0T_0,\qquad q_2=S_1T_1,
                 \widetilde v=S_0T_0-S_1T_1
 \tag{P3K.80.EA.44}
\]
as literal equations. It preserves the two ruling ideals and
therefore lifts both small resolutions, with their actual source
coordinate maps. This is the full higher-normal comparison behind
the linear display SZ.26.

\item The marked section is retained as well. On the old divisor s0=0,
the exact evaluations in (P3K.80.EA.42) give S1=S0. Hence
\[
                \kappa=\frac{S_1-S_0}{2s_0}
 \tag{P3K.80.EA.45}
\]
is holomorphic; its value at the origin is one, so it is a unit.
The actual line coefficient is carried by
\[
                    \xi_{\rm node}=\kappa^a\xi_{\rm Cartier}.
 \tag{P3K.80.EA.46}
\]
This follows by equating the meromorphic coefficients with
denominators \(s_0^a\) and \(((S_1-S_0)/2)^a\).
The source and target r maps are applied before and after this
formula when a=2. The full reset is consequently conjugated by
the actual kappa factors at its two ends; they are not replaced
by their value on the exceptional curve. The same conjugation
is made in the F0 root square.

\item The unit kappa and the normalized Jacobian of (P3K.80.EA.43) have
extending logarithms. Their exact logarithms carry the line and
canonical spin-c determinant lifts, while the nonunit
\((-1)^a z_s^{-2a}\) reset and its explicit root in EA.27
remain. Thus (P3K.80.EA.43)--(P3K.80.EA.46) connect the full finite-face proof
to the bilinear source map in RT.13--15 and to TL's actual Tate
label. They do not identify the intermediate nonlinear v formula
with the bilinear formula without this morphism.

\par\medskip\noindent\textbf{RT.12. Removing the scalar winding by an actual parameter homotopy.}\par\nopagebreak[4]

\item This step concerns the punctured family only. It supplies the exact
reason a constant-scalar local attachment can be studied separately;
it does not replace v\textasciicircum{}(km) by a unit over the origin.

\item Fix a small source circle |v|=rho and its annular neighbourhood. Choose
the branch of sqrt(1+v) which extends across the source disc and the
same root and Cartier data as RG.1 and RCZ.5--7. On this compact
annulus introduce an independent nonzero root-scalar r. The actual
maps, in the existing full annular marking, are
\[
 B=r^mB_0(v),\qquad A=r^nA_0(v),\qquad
 F=F_0(v)+n\frac{\log r}{2\pi i}s_\infty\delta.
                                                        \tag{P3K.80.RT.17}
\]
The reference F0 here is the actual extending cap-annulus map
constructed in RCZ.5, with its complete integral period, d-delta
cochain and reduced Cartier correction. Its construction is over
the full v-disc on that cap annulus. Source and root equations and
A0 are also holomorphic there, as RG.1 records. The ordinary toric
charts are the same fixed fan over the source annulus; all their
coefficient functions are retained. Compactness of the parameter
annulus permits one sufficiently small scalar annulus on which
these quotients and caps give the proper smooth family.

\item The positive r deck has the literal ordinary and cap arrows
q1\textasciicircum{}n and h\textasciicircum{}n, with h=1-T0(v)q\_infinity. They commute with the
retained period maps and finite groupoid arrows. Their coefficients
are single-valued in v on this annulus. Thus (P3K.80.RT.17) is the same
independent-scalar construction used in RG.5, now on the source-zero
annulus; no scalar or source closing is replaced by an unmarked map.

\item Take r0\textasciicircum{}m=lambda0 and retain the actual base amplitude \(r_{\rm base}=v_*^k r_0\). The original loop is
\[
 v(t)=v_*e^{2\pi it},\qquad r(t)=r_{\rm base}e^{2\pi ikt}.
\]
On the real phase cover its path (t,kt) has the same endpoints as
the broken path
\[
 (\min(2t,1),\ k\max(2t-1,0)).
\]
Their straight interpolation, with the corner smoothed relative
the endpoints, is an explicit based homotopy in the actual parameter
torus of (P3K.80.RT.17). The first part of the broken path is the source
loop at constant \(r_{\rm base}\); the second is k root-scalar loops at the same
source basepoint. Their complete endpoint arrows compose to
q1\textasciicircum{}(kn),h\textasciicircum{}(kn), exactly the original forward arrows. The given root
norm A0\textasciicircum{}m=B0\textasciicircum{}n and the full F0 square remain in this homotopy.

\item Any further comparison to a different fixed scalar reference uses
a separate zero-winding radial path inside the common contracting
range, carrying the actual norm root and F0. No such amplitude change
is hidden in the phase-cover homotopy above.

\item Consequently the actual oriented sphere-return classes satisfy
\[
 Z_{\mathcal T}(k)=Z_{\mathcal T}(0)
                      +k A^{\rm scalar}_{\mathcal T}.   \tag{P3K.80.RT.18}
\]
This follows from the constructed family homotopy and the clutch
composition, before using the abelian identification with Z/28.
For the two reference/reflected vertical fans the current physical
scalar values are 0 in the good branch and 3p-1 in the reduced
branch, by K69/K72. Their equality on those two fans is the literal
n=0,1 comparison in that integer scalar proof.

\item Thus the factor c*lambda in a local constant-scalar study is
constant only after (P3K.80.RT.17)--(P3K.80.RT.18) separate its k actual scalar
turns. In the original family c*v\textasciicircum{}(km)*lambda0 still has its full
winding and zero. No extending logarithm for that factor is asserted,
and the unknown source-zero class Z\_T(0) is not evaluated here.

\par\medskip\noindent\textbf{EA.12. The actual fixed-scalar exterior, finite root and inverse square.}\par\nopagebreak[4]

\item The scalar split uses the actual basepoint from RT.12:
\(r_{\rm base}=v_*^k r_0\), \(r_0^m=\lambda_0\).
The fixed-scalar piece is studied at \(r=r_{\rm base}\), and the
additional k canonical scalar turns are retained by the explicit
based homotopy in RT.12. A subsequent zero-winding radial change of
the scalar is a separate path in the common contracting range.

\item Here is the proper exterior available before the inner attachment.
Choose a source-base disc \(|T|<\epsilon\) containing both fibres
which collide at \((T,v)=(0,0)\), and shrink the source-parameter
disc so that their closures stay inside it. On its compact complement,
the resolved marked source, the section G, and the actual Cartier line
are smooth families through v=0. The selected factor q is nonzero
there. Consequently the actual nonzero-line lift B0 is an invertible
line map on this complement. A Hermitian metric on the compact
source complement bounds its norm by one finite constant C. Taking
\(|r_{\rm base}|^m C<1\) makes B a uniform contraction. Its
nonzero-line quotient has a compact fundamental annulus: the region
between a unit norm hypersurface and its B-image. This construction
retains the source boundary \(|T|=\epsilon\), the physical circle,
and the actual B endpoint; no quotient at the colliding source point
is used.

\item At the finite root cap retain \(A=r_{\rm base}^n A_0\),
\(B=r_{\rm base}^mB_0\), and \(A_0^m=B_0^n\).
Choose integers u,w with \(nu+mw=1\). The literal common generator
\(H=A^uB^w=r_{\rm base}H_0\) satisfies \(H^m=B\) and
\(H^n=A\). Its norm is bounded on the compact resolved root source,
so a further decrease within the same scalar range gives the same
compact fundamental-annulus construction. This applies to both
norm signs. The actual finite character and corrected root line are
carried by H0, not replaced by an unmarked root. The ordinary/cap
overlap is the full F0 of RCZ.5--7. Its lifted source/line maps commute
with B0 and physical scalars, so they identify the two actual quotient
boundaries. This gives a proper smooth exterior over the full v-disc,
with its only new source-base boundary \(|T|=\epsilon\).

\item For the inner finite strip, write \(J_j\) for the actual associated
node chart at chain label j. The following maps and equations are
already fixed:

1. The forward suspension is the actual map
   \(B_j:J_j\to J_{j+1}\), carrying the source by \(-G\) in
   the reduced branch and the line by its Cartier coefficient. In
   the common graph frame its node coordinates satisfy (P3K.80.EA.11).
2. Between consecutive source graphs the reset is (P3K.80.EA.18)--(P3K.80.EA.20),
   including \(\widetilde\xi=\xi/z_s^4\). This enters the
   complementary positive U=1 chart; it does not remain in the
   primary charge-one chart.
3. The adjacent-node arrow \(J_{j+1}\to J_j\) is (P3K.80.EA.23).
   Its physical normalization at fixed scalar is
   \((U,V)\mapsto((2\lambda)^{-1}U,2\lambda V)\).
   In the charged gauges this is (P3K.80.EA.22), with the chambers in EA.11.
   Thus the U=1 and V=1 descriptions are one charge-a boundary point.
4. The full ordinary-to-cap square on each chosen root sheet is
   \[
      F_0 U_a^{\rm Cartier}
         =z_\infty^{ma/2}U_{\rm cap}F_0,
      \qquad
      F_0B=2(-1)^{b+1}\lambda h_\infty\,U_{\rm cap}^{b}F_0.
   \]
   The root character and the good-branch sign cancel as in EA.4;
   the reduced branch commutes. These are equations of actual line
   maps, with \(q=q_\infty^{-1}h_\infty\), not merely source-point equations.
5. The full node-to-Tate map TL.1--4 carries the
   actual h reindex, \(P=q_2^{-1}u_P\), the theta units and the
   degree-two Cartier map. Its positive associated square is the
   current Tate fan T0. It must be used on the \(|T|=\epsilon\)
   collar rather than the earlier normalization-stratum formula EA.8.

\par\medskip\noindent\textbf{EA.5. The finite stable-end metric arrow and its fixed boundary.}\par\nopagebreak[4]

\item This section works with a fixed nonzero scalar, chosen small after the
compact source-origin neighbourhood has been cut out. It does not use the
nonproper infinite central quotient from SZ.29. On the compact exterior
the actual source/line multiplier is bounded, so this scalar gives a
uniform finite contracting cut. The selected divisor is retained by its
node charts, rather than by a limit of a zero scalar in one chart.

\item For the reduced primitive period write the nonzero constant as
\(c_\lambda=2\lambda\). The actual arrow is
\[
 (p,w)\longmapsto(p-G,c_\lambda q\,\rho(p)w),
\]
where \(\rho\) is precisely the Cartier line arrow (P3K.80.EA.6). On the
logarithmic Tate chain put
\[
 U_j=\frac{w}{c_\lambda^j q^j},\qquad
 V_j=\frac{c_\lambda^j q^{j+1}}w,
 \qquad U_jV_j=q.
\]
The period sends \(j\) to \(j+1\) and
\[
 (U_j,V_j)\longmapsto(\rho U_j,\rho^{-1}V_j).
 \tag{P3K.80.EA.11}
\]
In the common Cartier/Cox line frame, \(\rho\) is the identity
arrow between the two source rulings. Thus its node extension is
literally the associated chart \(UV=S_0T_0\), with the weights in
SZ.28. The factor \(q\) has been removed by the stated chain-index
change, not by a limiting argument. The two finite cap arrows in
EA.4 still accompany this chart.

\item Here is the needed change of the finite cut hypersurfaces on the stable
locus. Let \(\xi\) be the common line coordinate of auxiliary weight
\(a\), with \(s\) of weight \(+1\) and \(t\) of weight \(-1\).
The actual Fubini--Study divisor metrics give
\[
 N_t=|\xi|\,\|t\|^a,\qquad
 N_s=\frac{|\xi|}{\|s\|^a},\qquad
 R=\|s\|\|t\|,\qquad N_t=R^aN_s.
 \tag{P3K.80.EA.12}
\]
These are auxiliary-invariant functions. The arrow from the domain end
\(H_t=\{N_t=1\}\) to the target end \(H_s=\{N_s=1\}\) is the
physical positive rescaling
\[
 (\xi,\xi_{\rm dual})\longmapsto
 (R^a\xi,R^{-a}\xi_{\rm dual}).
 \tag{P3K.80.EA.13}
\]
It preserves the product \(\xi\xi_{\rm dual}=q\), the source
coordinates and the physical circle. On \(R\ge\varepsilon>0\)
it is smooth and bounded and extends both node branches. At
\(R=\varepsilon\) its factor is the constant
\(\varepsilon^a\). Only after (P3K.80.EA.13) is the auxiliary positive
normalization used to choose a representative. Auxiliary normalization
alone does not change (P3K.80.EA.12), and was not an end-to-end map.

\item On the reduced moment graph of RT.8,
\(\mu_{\rm aux}+K\mu_\delta=c\), with \(K,c>0\), the derivative
along the positive auxiliary orbit has the five coefficients
\[
             1,\quad1,\quad1,\quad1+K,\quad4+2K.
\]
It is therefore strictly positive away from the origin. The unique
normalizing arrow is smooth on the displayed stable overlaps and
preserves the full functions \(q,v\), because they are auxiliary
invariants. The actual cap metric is transported by \(U_{\rm cap}\);
on the overlap it is joined to (P3K.80.EA.12) by the logarithm of the positive
metric ratio. Averaging that real function over the finite root group
makes the interpolation equivariant and retains its original character.

\item The inner boundary also meets the fixed surface. In the positive node
chart it is
\[
 z=V=0,\qquad (S_0,S_1,T_1)/\mathbb C^*,
 \qquad \text{weights }(1,1,-1).
\]
The invariant radius on this surface is
\(R=\|(S_0,S_1)\||T_1|\). Its level \(R=\varepsilon\) is the
circle bundle of \(O(-1)\), with the following literal \(S^3\)
coordinate. Use the unique positive auxiliary gauge \(\|S\|=1\),
write \(T_1=\varepsilon e^{i\phi}\), and set
\[
                  p=e^{i\phi}(S_0,S_1)\in S^3.
 \tag{P3K.80.EA.14}
\]
The positive-weight spinor is \((z,\overline V)\). In this same
gauge its actual boundary frame is
\[
 \left(e^{i(a-1)\phi}z,e^{ia\phi}\overline V\right).
 \tag{P3K.80.EA.15}
\]
Before positive normalization the first component is
\(T_1^{a-1}z/\varepsilon^{a-1}\), and the second is
\(\|S\|^{2a}T_1^a\overline V/\varepsilon^a\).
The displayed positive factor in the second component is required for
full auxiliary invariance; no phase has been suppressed.

\item Consequently the inner end is not identified here with the bare ruling
mapping torus. Its stable portion has the precise (P3K.80.EA.13) arrow and
the ruling change of EA.4, while its fixed portion has the \(S^3\)
and normal frame (P3K.80.EA.14)--(P3K.80.EA.15). The complete mixed-face construction below joins these portions with their actual boundary orientations.
The certificate \path{code/part3_fourfold/source_zero_cap_end.py}
checks the admitted finite characters and the exact metric/frame weights.

\par\medskip\noindent\textbf{EA.9. The actual cyclic graph reset and its charged-chart continuation.}\par\nopagebreak[4]

\item The common graph identity of EA.4 does not identify consecutive graph
coordinates. Here is the missing reset, with its line map. Use the
marked source-resolution coordinates
\[
 A=Dz_s,\quad C=Bz_s,\quad
 v=(B-Dz_s)/2,\quad X=Bz_s,
\]
and set
\[
 F_2=1-Bz_s^3-(1-2v)z_s^2.
\]
The next negative-G translation has the exact source coordinates
\[
 X'=F_2/z_s^2,\quad A'=-F_2/z_s^3,\quad
 B'=B-Dz_s-F_2/z_s^3,\quad D'=-B'/z_s.
 \tag{P3K.80.EA.18}
\]
They satisfy \(A'B'=X'D'\), preserve \(v\), and preserve the full
source base coordinate
\(T=D-X^2-(1-2v)X\). In the next common Cox gauge take
\[
 (\widetilde t_0,\widetilde t_1)=(1,-z_s),\qquad
 (\widetilde s_0,\widetilde s_1)=(A',D').
 \tag{P3K.80.EA.19}
\]
In the current target gauge \((s_0,s_1)=(z_s,1)\), the meromorphic
line coefficient is \(\xi/z_s^a\). In the next domain gauge it is
\(\widetilde\xi(-z_s)^a\). Therefore the exact coefficient arrow is
\[
            \widetilde\xi=(-1)^a z_s^{-2a}\xi.
 \tag{P3K.80.EA.20}
\]
For the actual reduced line, this reads
\(\widetilde\xi=\xi/z_s^4\), after the retained Cartier map
\(r\) at both ends. Equations (P3K.80.EA.18)--(P3K.80.EA.20) give the reset,
rather than replacing it by the bare projective ruling.

\item Its effect on the associated-node chart is already nontrivial on the
central proper component. Write
\[
 X=s^2-1,\quad Y=s(s^2-1),\quad z_s=1/s,\quad
 B=s(s^2-1),\quad D=s^2(s^2-1).
\]
Near its \(s=-1\) node, the analytic square root in RT.2 is
\(\sigma=-s\). Substitution in the literal functions
\[
 F_+=\frac{(\sigma+1)(s_1-s_0)
                -s_0^2(t_0+t_1)/(\sigma+1)}2,
\]
\[
 F_-=\frac{(\sigma+1)(t_0+t_1)
                -t_1^2(s_1-s_0)/(\sigma+1)}2
\]
gives \(F_+=2,F_-=0\). Thus the finite associated coordinate
\(z_n=\xi/F_+\) is nonzero. At the next target, (P3K.80.EA.18) has
\(\widetilde s_0=\widetilde s_1=0\), whereas
\(\widetilde\xi=\xi s^4\ne0\). Its analytic square root has
value one, and
\[
             \widetilde F_+=0,\qquad
             \widetilde F_-=1-1/s.
\]
Consequently \(z_n\) goes to infinity. The reset enters the
complementary \(U\ne0\) charge-a chart. It does not close inside
the finite nonzero charge-\((a-1)\) chart. This is the exact reason
the Euler sequence \(2M\oplus2M^{-1}\) from that chart alone has
not been descended over the entire cyclic end.

\item There is nevertheless a literal continuation from the resulting
charged chart. Use the complementary small-resolution equations
\[
                  S_0=Uw,\qquad V=wT_0.
\]
In its \(U=1\) gauge the remaining coordinates are
\((w,S_1,T_1,T_0)\), and the full source functions from RT.2 are
\[
 q=wT_0,\qquad
 v=S_1T_1+\frac{T_1w-S_1T_0}{2\sqrt{1+S_1T_1}}.
 \tag{P3K.80.EA.21}
\]
The primary \(V=1\) chart has coordinates
\((S_0,S_1,T_1,z_n)\). Its source functions are exactly (P3K.80.EA.21)
under
\[
                (S_0,S_1,T_1,z_n)=(w,S_1,T_1,T_0).
 \tag{P3K.80.EA.22}
\]
This map has a specific chain meaning. On \(U_j\ne0\), the
actual adjacent node overlap is
\[
 U_{j-1}=c_\lambda qU_j,\qquad
 V_{j-1}=(c_\lambda U_j)^{-1},\qquad c_\lambda=2\lambda.
 \tag{P3K.80.EA.23}
\]
At \((U_j,V_j)=(1,q)\), apply the retained source-fibre scalar
normalization \((U,V)\mapsto(c_\lambda^{-1}U,c_\lambda V)\).
The result is \((q,1)\), precisely (P3K.80.EA.22). Thus this is a node
chart overlap with index \(j\mapsto j-1\), not an additional
cap or an extra charged point. In a forward B cut its index cancels
the preceding \(j\mapsto j+1\) in (P3K.80.EA.11). In particular the
charged origin in these two descriptions is one point, and must
not be counted twice.

\item Both gauges retain the same effective physical weights, in the
ordered coordinates of (P3K.80.EA.22),
\[
                   (-1,-1,+1,+1)/a,
\]
and the same residual \(\mu_a\) action. The displayed coordinate
arrow is complex linear on these charged coordinates and preserves
their real eight-dimensional orientation. Conjugating its two
negative physical-weight coordinates gives the same positive Hopf
projection at both ends. Its physical line is \(O(-a)\), with
the outward-meridian dual \(O(a)\). A charged ball is carried by
(P3K.80.EA.22), with this full normalization retained; its transported
truncation is used if the original radii differ. This fixes the
boundary orientation against the local charged model of RT.10.

\item The source functions \(q,v\) in (P3K.80.EA.21) are unchanged throughout.
On the finite cap, (P3K.80.EA.7)--(P3K.80.EA.10) carry the same primitive map,
including \(q=q_\infty^{-1}h_\infty\), the actual \(z_\infty^m\)
factor in the reduced line, and its trivial commutator with
\(A_0\). The physical normalization in (P3K.80.EA.23) commutes
with this full root/line square. Both norm signs and every retained
finite character remain.

\item These are the cyclic chart transition and local charged orientation used in the complete finite-domain attachment below.
The exact source, line, and charged-chart identities are executed in
\path{code/part3_fourfold/source_zero_associated_reset.py}.

\par\medskip\noindent\textbf{EA.10. The retained scalar winding in the adjacent-node arrow.}\par\nopagebreak[4]

\item The finite-cut construction in EA.5 starts at fixed nonzero
\(\lambda\). For the original rescaled source loop, the exact same
overlap has
\[
 c_\lambda(v)=2\lambda_0v^{km},\qquad
 (U,V)\longmapsto(c_\lambda(v)^{-1}U,c_\lambda(v)V).
 \tag{P3K.80.EA.24}
\]
Thus its two characters have source-parameter windings \(-km\)
and \(+km\). This map is not asserted to extend at \(v=0\) in
those old coordinates. Its full collar derivative is retained: it is
the diagonal character map followed by a rank-one factor
\(I-\eta\otimes d\log c_\lambda\), where \(\eta\) is the
physical complex orbit field and
\(d\log c_\lambda(\eta)=0\). The latter factor is invertible
through the explicit homotopy
\(I-t\eta\otimes d\log c_\lambda\); the diagonal character and
its winding are not removed by this homotopy.

\item The comparison with a fixed-scalar boundary uses the actual two-parameter
family, not an extension of (P3K.80.EA.24). Write \(\lambda=r^m\). Its
cap generator is \(A(v,r)=r^nA_0(v)\), and the ordinary-to-cap
map on the logarithmic r cover is
\[
 F(v,\nu)=F_0(v)+n\nu s_\infty\delta,\qquad
 \nu=\log r/(2\pi i).
\]
Here \(F_0(v)\) is the complete extending annular map of RCZ.5--7,
with its integral cochain fixed at the reference fibre. The positive
r deck has the full arrows \(q_1^n\) and \(h_\infty^n\), since
\(q_1=q_\infty^{-1}h_\infty\), and
\[
          F(v,\nu+1)\,q_1^n=h_\infty^n F(v,\nu).
 \tag{P3K.80.EA.25}
\]
The equation retains both norm signs, the fixed finite character,
and every source/line unit. The same source and line Cartier maps
are used over both parameter directions; the base-only physical
scalars commute with them. The cap division and \(F_0\) have their
stated holomorphic source-parameter extensions, so their source-loop
endpoint is retained literally.

\item On the boundary, \(r=v^k r_0\). In this actual product-parameter
family the diagonal loop is homotopic to the source loop at fixed r
followed by k scalar-root loops at the comparison source. The based
homotopy carries all the endpoint squares (P3K.80.EA.25). Hence its return
class is the fixed-scalar source return plus the actual
\(kA_{\rm scalar}\) term, in the same transported fan. This is
the geometric parameter-loop relation already used in K58; no
class in \(\Theta_7\) is divided. The fixed-scalar attachment is constructed by the complete domains and collar maps below.

\par\medskip\noindent\textbf{EA.11. Chamber ownership of the two charged descriptions.}\par\nopagebreak[4]

\item The coordinate identification (P3K.80.EA.22) is retained together with its two
auxiliary chambers. For the primary resolution
\(U=S_0z,T_0=zV\), the auxiliary and physical weights on
\((S_0,S_1,T_1,z,V)\) are respectively
\[
 (1,1,-1,a-1,-a),\qquad(0,0,0,1,-1).
\]
For the complementary resolution \(S_0=Uw,V=wT_0\), the weights
on \((U,w,S_1,T_1,T_0)\) are
\[
 (a,1-a,1,-1,-1),\qquad(1,-1,0,0,0).
\]
Thus, for a=2, the primary positive chamber is the actual
\(\operatorname{Tot}(O(-1)\oplus O(-2)\to\mathbb P^2)\)
of RT.9. Its z=1 origin has charge one. The primary V=1 origin
instead lies in the negative auxiliary chamber and has charge two.
The complementary U=1 origin lies in its positive chamber and also
has charge two. Equation (P3K.80.EA.22) identifies the last two charts;
it does not place a charge-two origin in the primary positive chamber.

\item In either charged gauge, maintaining U=1 or V=1 under the physical
circle uses the auxiliary factor of exponent \(-1/a\). The ordered
remaining coordinates therefore have the same physical weights
\((-1,-1,+1,+1)/a\). Their residual auxiliary \(\mu_a\)
characters agree modulo a. The map (P3K.80.EA.22) is complex linear in
those coordinates, preserves their real eight-dimensional orientation,
and carries the physical line \(M=O(-a)\) literally.

\item The ungauged continuation, on the complementary chart U nonzero,
is the explicit map
\[
 (U,w,S_1,T_1,T_0)\longmapsto
 (S'_0,S'_1,T'_1,z',V')=(Uw,S_1,T_1,UT_0,U^{-1}).
 \tag{P3K.80.EA.26}
\]
Its inverse is \(U=(V')^{-1},w=S'_0V',T_0=z'V'\), with
S1,T1 unchanged. It preserves the auxiliary action and the physical
circle separately. It also preserves q and the full v in (P3K.80.EA.21).
At U=1 it is exactly (P3K.80.EA.22). It is the adjacent-node arrow with
the full scalar normalization already specified, not the naive
exchange of U and V on the ungauged associated-node space. A
radial defining function on either charged quotient chart can be
transported by (P3K.80.EA.26), with its entire derivative. Choosing the
retained or excised side of that transported function remains part
of the compact-cut construction.

\item Chamber names alone do not determine the boundary normal. On the
normalized charged gauges the two actual moment functions are
\[
 \mu_{\rm comp}=a+(1-a)|w|^2+|S_1|^2-|T_1|^2-|T_0|^2,
\]
\[
 \mu_{\rm prim}=|w|^2+|S_1|^2-|T_1|^2
                         +(a-1)|T_0|^2-a.
\]
Their sum is
\[
 (2-a)|w|^2+2|S_1|^2-2|T_1|^2+(a-2)|T_0|^2,
\]
and is not identically zero. The sign of a charged boundary in a
compact bordism must therefore be taken from the explicitly retained
or excised ball, with its outward radial normal. It is not obtained
by negating one moment function. These weight and moment identities
are checked in \path{code/part3_fourfold/source_zero_chamber_ledger.py}.

\par\medskip\noindent\textbf{EA.13. Canonical complex spin-c and the explicit determinant lifts.}\par\nopagebreak[4]

\item The lower charged arm carries the canonical complex spin-c structure,
with its actual physical associated line. It is not assigned a spin
structure by its rational characteristic classes. In particular, the
two possible lifts over the fundamental group of an RP3 collar are
kept distinct until the actual chart maps have been carried.

\item There are explicit determinant arrows for those maps. In the source
reset (P3K.80.EA.18), use source coordinates \((B,D,z_s)\) on the first
small-resolution chart, and \((A',D',w')\), \(w'=-z_s\), on
the second. Direct differentiation gives
\[
 \det\frac{\partial(A',D',w')}{\partial(B,D,z_s)}=-1.
\]
Including the exact coefficient (P3K.80.EA.20), with the Cartier r map at
both ends for a=2, the four-dimensional complex Jacobians are
\[
  J_1=z_s^{-2},\qquad J_2=-z_s^{-4}.
 \tag{P3K.80.EA.27}
\]
The determinant square roots used to transport the spin reduction
in these coordinate trivializations are the literal functions
\(z_s^{-1}\) and \(i z_s^{-2}\), respectively. Their inverses
are used on the reverse arrow. These functions retain the full
source winding; no choice of a small logarithm of z\_s is made.

\item For (P3K.80.EA.26), in the displayed ordered five coordinates the complex
Jacobian is the constant -1. The map intertwines the auxiliary
Euler vector field, so contraction of the five-form with that
field gives the same constant -1 on the quotient canonical
four-form. The fixed square root i carries its spin reduction.
On the normalized charged covers, the four-coordinate arrow is
the identity; its residual mu-a character is the one in EA.11.
These two descriptions agree by the actual auxiliary gauge and
its determinant. They do not permit an additional sign character
of the RP3 fundamental group.

\item More explicitly, on the lower arm the normal bundle is
\((S^3\times\mathbb C^2)/\{\pm1\}\), with the indicated
minus action on both normal coefficients. Its complex determinant
is trivialized by their ordered wedge. The canonical complex
spin-c lift selects a specific spin lift with respect to that
trivialization. With Clifford convention \(e_i^2=-1\), follow
\(\operatorname{diag}(e^{it},e^{it})\), \(0\le t\le\pi\).
Its canonical spin-c endpoint is
\[
 [e_1e_2e_3e_4,-1]=[-e_1e_2e_3e_4,1].
 \tag{P3K.80.EA.27a}
\]
Thus the normal mu2 lift is the specified \(-\omega_4\), not an
arbitrary choice between \(\omega_4\) and \(-\omega_4\).
The other complex two-plane block has the same canonical lift;
their product is \(+\omega_8\), in agreement with the full
charged-cover computation below. Equation (P3K.80.EA.22) is the identity
on these ordered cover blocks. The reset carries their embedded
collars and normal derivative by (P3K.80.EA.27), so this local lift is
transported together with the physical line.

\item The physical normalization is also carried at the level of its
lift. At a=2, along a chosen path of nonzero scalars d from one
to its actual value, the charged-cover derivative has diagonal
characters
\[
 (d^{-1/2},d^{-1/2},d^{1/2},d^{1/2}).
 \tag{P3K.80.EA.28}
\]
The square roots here are continued along that very path. A full
positive phase turn gives the endpoint -I4. Its canonical complex
spin-c action on \(\Lambda^{0,*}\mathbb C^4\) is
\((-1)^{\deg}\). The determinant-one path in (P3K.80.EA.28), whose
four angles are \((-\pi,-\pi,\pi,\pi)\), has exactly that
spin endpoint: on an occupation vector \((n_1,n_2,n_3,n_4)\)
its sign is
\[
 (-1)^{-n_1-n_2+n_3+n_4}=(-1)^{n_1+n_2+n_3+n_4}.
\]
It therefore closes through the specified canonical residual mu2
action. Replacing the root by its other branch is this same deck
map with this specified lift, not an unrecorded alternative spin
structure. The physical line remains M, not its dual; its charged
Hopf class is \(O(-a)\).

\item For a nonconstant scalar the full derivative also contains the
rank-one factor in EA.10. Its displayed homotopy is based at the
actual diagonal arrow and carries the chosen lift. The theta and
Morse units in TL.1--2 have extending logarithms on the complete
bidisc; their determinant and line maps are carried by those
logarithms. The actual h period and the Cartier r factors precede
this unit homotopy, so their integral characters are not removed.
On the finite cap the canonical complex spin-c maps are those of
the full root/line F0 square in EA.12; its finite character remains
part of the map.

\item The exact determinant computations and the sixteen spinor endpoint
checks are in \path{code/part3_fourfold/source_zero_spin_chart.py}.
These are the actual local lift data carried on the complete finite attachment below.

\par\medskip\noindent\textbf{EA.14. A full charged shape matching the lower arm.}\par\nopagebreak[4]

\item The regular part RT.20 has its two explicit arms; the actual complete retained domain is EA.48, which also retains its upper fibre and charge-one point. In its lower arm use the primary V=1 gauge and write
\(x=(S_0,S_1,z)\),
\[
 H=\|x\|^2,\quad b=H|T_1|^2,\quad c=H^2,
 \quad u=\frac b{b+c},\quad
 \sigma^2=\|x\|^2+|T_1|^2.
\]
The actual disk-bundle defining function satisfies the exact identity
\[
                  \rho=\sigma^2\sqrt{1-u}.
 \tag{P3K.80.EA.29}
\]
Indeed its square is \(b+c\), and
\(\sigma^4(1-u)=H(H+|T_1|^2)=b+c\).
On the lower small-R arm, u is at most the chosen ratio
\(\varepsilon^2/r_0^2<1/4\). Choose a smooth invariant angular
cutoff chi equal to one on a neighbourhood of that closed arm and
supported in \(u<1/2\). Set
\[
             f=1-\chi(1-\sqrt{1-u})>0.
\]
Then the full charged hypersurface \(\sigma^2f=\rho_0\)
agrees literally with \(\rho=\rho_0\) near the arm. It is a
complete smooth star-shaped hypersurface on the charged cover,
invariant under the residual mu2 and the physical circle. The
normal equality on that neighbourhood is an equality of defining
functions, including their derivatives, not only an equality of
zero sets.

\item Use the positive-width compact annulus
\[
 \mathcal A_f=\{\rho_0/2\le\sigma^2\le\rho_0/f\}/\mu_2.
 \tag{P3K.80.EA.30}
\]
It is a smooth interval bundle with disjoint boundary faces. Its
inner boundary is a round charged link; its outer boundary is the
full transported shape. All q and v functions, the physical line,
and the canonical complex spin-c structure are their restrictions
from the actual charged chart. No radial map is asserted to
preserve q or v. The physical circle is free on this annulus:
its nonzero coordinate weights are all plus or minus one after
the retained residual mu2 is accounted for. Product connections
on its collars therefore show that the outer shape and round
charged link have equal physical characteristic integrals.

\item Equation (P3K.80.EA.26) pulls the entire outer defining function to the
complementary U chart. It transports the outward radial normal
and the retained/excised side exactly, and it carries the lifts
in EA.13. The earlier coincident-face proposal with inner radius
rho0 is not used: (P3K.80.EA.30) uses rho0/2 everywhere, including where
f=1.

\item This shape comparison still respects the chamber distinction.
The lower arm itself lies in the common stable overlap with the
primary positive chamber. The charge-two origin filling the full
shape belongs to the primary negative chamber. The complete domain and horn chain EA.48--52 specify where the complement of the lower arm lies among the boundary faces. The charged shape is not
glued as an additional positive GIT chart. Thus (P3K.80.EA.29)--(P3K.80.EA.30)
fix the radial normal and its spin-c/line data for the full mixed-face attachment EA.21--22.

\par\medskip\noindent\textbf{EA.15. The native finite-strip quaternion and the charged seam.}\par\nopagebreak[4]

\item The native ruling coordinates are taken after the actual analytic
Morse-Cox map of RT.2 and the full two-node refinement in EA.17 below.
To make the first arrow explicit, write its old
pairs as s=(s0,s1), t=(t0,t1), and put
\[
       S=(F_+,s_0),\qquad T=(F_-,t_1).
 \tag{P3K.80.EA.31}
\]
This is an auxiliary-equivariant analytic isomorphism with the
inverse already displayed in RT.2. It preserves both ruling
ideals. Indeed, writing \(d=(\sigma+1)/2\),
\[
 F_+=d s_1+s_0\left(-d-
                  \frac{s_0(t_0+t_1)}{2(\sigma+1)}\right),
\]
\[
 F_-=d t_0+t_1\left(d-
                  \frac{t_1(s_1-s_0)}{2(\sigma+1)}\right).
\]
The two coefficient determinants are -d and d, both units.
Consequently (F+,s0)=(s0,s1) and (F-,t1)=(t0,t1) as ideals,
so the map lifts both projective rulings. It keeps the common
line coefficient xi, with the actual r maps in the degree-two
sector. Its full complex Jacobian is -sigma, an extending unit.
The source functions are exactly those of RT.2, expressed in
these coordinates. This is the first source-preserving chart arrow,
not a substitution based only on matching weights. Its displayed
pairs are intermediate selected-node coordinates. EA.17 carries
them to the full two-node pairs; the S,T in the norm and quaternion
formulas below denote those final pairs. The extra source/line
units and their inverses are kept in every reset and cap composition.

\item Use the divisor metrics transported by (P3K.80.EA.31) for the finite
end. In these metrics the exact EA.12 ratio is
\(N_T=R^aN_S\), \(R=\|S\|\|T\|\). On the stable collar
the old and transported metrics have a bounded positive ratio;
its actual real logarithm, and its transported cap average, give
the metric interpolation. No phase or period map is changed.
The positive scaling in EA.13 leaves the unit directions of S
and T unchanged.

\item Set, in these actual pairs,
\[
 u=S_0+jS_1,\qquad c=T_0+j\overline{T_1},\qquad
 \widehat u=u/\|S\|,\quad\widehat c=c/\|T\|,
\]
\[
 q_H=\widehat u\widehat c,\qquad
 h=\widehat u i\widehat u^{-1},\qquad
 k=\widehat c^{-1}i\widehat c.
 \tag{P3K.80.EA.32}
\]
Thus the actual t-ruling to s-ruling map is
\(h=q_H k q_H^{-1}\), with its full normal coordinate qH.

\item In the primary V=1 charged chart, T0=z. Conjugate the two
negative physical-weight coordinates and choose the explicit
constant unitary target coordinates
\[
 p=\overline{S_0}+j\overline{S_1}=j u j^{-1},\qquad
 q_{\rm alt}=z-jT_1=j\overline c j^{-1}.
\]
Both have the same positive right physical-circle action, with
the retained common mu-a quotient. The exact projective ratio is
\[
 w_{\rm alt}=p q_{\rm alt}^{-1}
       =\frac{j(u c)j^{-1}}{\|T\|^2},\qquad
 g_{\rm alt}=\frac{w_{\rm alt}}{|w_{\rm alt}|}
       =j q_H j^{-1}.
 \tag{P3K.80.EA.33}
\]
Its two Hopf base coordinates are
\[
 h_{\rm alt}=-j k j^{-1},\qquad
 h'_{\rm alt}=-j h j^{-1}.
 \tag{P3K.80.EA.34}
\]
Therefore \(h'_{\rm alt}=g_{\rm alt}h_{\rm alt}
g_{\rm alt}^{-1}\) is precisely the native inverse ruling
map with the stated constant arrows. Each minus in (P3K.80.EA.34)
is the antipodal map of S2, of degree -1; neither is discarded.
The exact Hopf lift is
\[
       g_{\rm alt}\frac{q_{\rm alt}}{|q_{\rm alt}|}
                           =\frac p{|p|}.
 \tag{P3K.80.EA.35}
\]

\item A second explicit charged atlas uses instead
\(q_{\rm root}=T_1+jz=j q_{\rm alt}\).
This is the literal constant complex-linear map
\((z,-T_1)\mapsto(T_1,z)\), with complex determinant one.
In that exact coordinate choice the dictionary becomes
\[
 g_{\rm root}=-j q_H,\qquad h_{\rm root}=-k,
 \qquad h'_{\rm root}=-j h j^{-1},
 \qquad
 g_{\rm root}\frac{q_{\rm root}}{|q_{\rm root}|}
                               =\frac p{|p|}.
 \tag{P3K.80.EA.36}
\]
This gives the literal comparison between the two charged atlases. The other quaternion
hemisphere always retains \(w'=w^{-1}\), including its actual
four-dimensional base orientation. The identity reconstructing
the same charged point is
\((p,q)=(rwu,ru)=(r'u',r'w'u')\), with \(r'=|w|r\),
and the positive radius is the actual solution on the transported
shape of EA.14. Hence every q-source and v function is literally
unchanged on the overlap. The inverse of (P3K.80.EA.31) carries that
statement back to the original Weierstrass/Cox chart.

\item For a=2 the associated physical line map is
\([u,z]\mapsto z u^{\otimes2}\) into O(-2); the Hopf arrow
(P3K.80.EA.35) intertwines it and its normal coefficient frame. The
canonical spin-c and determinant lifts are those of EA.13;
the analytic Morse unit -sigma has its extending logarithm
after its fixed constant sign is retained. The full F0 cap
square and the scalar factor in EA.23 accompany this seam.
The exact quaternion and ideal identities are executed in
\path{code/part3_fourfold/source_zero_native_charged_seam.py}.

\item This is the charged/native seam map. Its hemisphere inversion and both sphere antipodal maps are retained in the complete finite-strip boundary calculation EA.47 below.

\par\medskip\noindent\textbf{EA.16. Entire finite faces and the retained line coefficient.}\par\nopagebreak[4]

\item The norm inequalities are first imposed on the complementary
charged chart, after the literal chain-index change and the source
reset. Its common line coefficient is \(\xi=U_j\). Fixing
\(U_j=1\) gives the residual mu-a cover with actual source pairs
\[
                 S=(w,S_1),\qquad T=(T_0,T_1).
\]
The invariant norm functions in this gauge are exactly
\[
                  N_t=\|T\|^a,\qquad
                  N_s=\|S\|^{-a},\qquad R=\|S\|\|T\|.
 \tag{P3K.80.EA.37}
\]
This gauge is taken only after the actual node reindex; it is not
an assertion that the primary line coordinate U equals one.
Indeed (P3K.80.EA.26) sends it to the primary chart V'=1, with the same
ordered pairs \(S'=(w,S_1)\), \(T'=(T_0,T_1)\). Before that
gauge the transported common coefficient is
\[
                         \xi=(V')^{-1},
 \tag{P3K.80.EA.38}
\]
whereas the primary node coefficient is \(U'=q/(V')\).
Consequently the transported norms there are
\(|V'|^{-1}\|T'\|^a\) and
\(|V'|^{-1}/\|S'\|^a\). This retains the q-factor and is the
exact reason the bare inequalities in (P3K.80.EA.37) apply after the
specified arrow. The source-fibre constant c-lambda, its phase,
and the index j to j-1 are those in (P3K.80.EA.23); they precede (P3K.80.EA.38).

\item Choose a sufficiently small charged radius r. The two entire
closed finite faces are
\[
 H_t(r)=\{N_t=r^a,\ R\le r^2\}
       =(D^4_S(r)\times S^3_T(r))/\mu_a,
\]
\[
 H_s(r)=\{N_s=r^{-a},\ R\le r^2\}
       =(S^3_S(r)\times D^4_T(r))/\mu_a.
 \tag{P3K.80.EA.39}
\]
In the positive charged quaternion coordinates of EA.15 these
are respectively \(|p|\le|q|=r\) and
\(|q|\le|p|=r\). Their common face is
\(|p|=|q|=r\), with the exact quaternion and Hopf-line
transition (P3K.80.EA.33)--(P3K.80.EA.36). Thus the comparison specifies the
full two orbit charts, not only their unit-quaternion seam.
Rounding the product-ball corner uses its two actual norm defining
functions. Its charged shape is then compared with the prescribed
lower radial arm by EA.14, by ambient inclusion with q,v retained.

\item On R positive the end normalization from Ht to Hs is
\[
             d=(R/r^2)^a,\qquad
             (S,T)\longmapsto(d^{-1/a}S,d^{1/a}T).
 \tag{P3K.80.EA.40}
\]
It sends \((\|S\|,\|T\|)=(s,r)\) to \((r,s)\),
preserves all source auxiliary invariants and hence q,v, and is
the actual positive physical scaling of the common line before
refixing its gauge. It leaves the unit directions used in qH
unchanged. At a positive inner cut R=epsilon the factor is a
single positive constant. No extension of (P3K.80.EA.40) at R=0 is
claimed: the two cores remain the distinct full-face boundary
pieces in (P3K.80.EA.39).

\item The oriented side follows from the actual contracting direction.
Let eta denote the positive physical real-orbit field on this
charged cover, up to its positive factor 1/a:
\[
                       \eta=-S\partial_S+T\partial_T.
\]
The forward suspension is contracting in its physical line, so
its normalized real direction is -eta. The retained constant
phase acts tangentially and carries the lift in EA.13. On Ht,
-eta differentiates \(\|T\|^2-r^2\) to \(-2r^2\): it
enters the product ball. On Hs it differentiates
\(\|S\|^2-r^2\) to \(+2r^2\): it exits. If each end is
oriented by contraction with the forward direction, the target
minus source end chain therefore equals the positively oriented
boundary of
\[
                       (D^4_S(r)\times D^4_T(r))/\mu_a.
 \tag{P3K.80.EA.41}
\]
This comparison uses the full real eight-dimensional coordinate
arrow. The two negative complex conjugations preserve that
orientation; the S2 antipodal maps and the hemisphere inversion
in EA.15 are retained within it. Thus no side is selected merely
from the name of a ruling or an auxiliary chamber.

\item All maps on these faces are the node-chart maps in EA.9, followed
by the source/Cartier map of TL and the full F0 cap square of
EA.12. The physical associated line and its canonical spin-c
lift are transported by those same maps. The exact finite-face
norm and side computations are in
\path{code/part3_fourfold/source_zero_full_faces.py}.
The complete domain and mixed-face maps in EA.19--22 insert these full faces into the compact period strip and retain the whole fixed cap.

\par\medskip\noindent\textbf{EA.6. The reduced positive moment model as an actual bundle.}\par\nopagebreak[4]

\item For \(a=2\), the positive auxiliary coordinates of the associated node
chart are \((S_0,S_1,z)\), each of weight one. The negative coordinates
\((T_1,V)\) have weights \((-1,-2)\). The strict derivative in EA.5
shows that the level \(\mu_{\rm aux}+K\mu_\delta=c>0\) meets each
positive-stable auxiliary complex orbit exactly once modulo its compact
circle. It therefore has the literal smooth quotient map to
\[
 \operatorname{Tot}\bigl(O(-1)\oplus O(-2)\longrightarrow
                         \mathbb P^2_{[S_0:S_1:z]}\bigr).
 \tag{P3K.80.EA.16}
\]
Its physical action is the one induced by weight one on \(z\), weight
minus one on \(V\), and weight zero on the remaining coordinates. This
keeps the fixed \(O(-1)\) surface over \(z=0\) and the isolated point
\([0:0:1]\), rather than replacing them by a putative free end.

\item In the exact node-adapted source coordinates, \(q_1-q_2=v u_v(v)\),
where \(u_v(0)\ne0\). The parameter change
\(\widetilde v=v u_v(v)\) is an actual biholomorphism of discs.
On (P3K.80.EA.16) its map is
\[
                    \widetilde v=S_0zV-S_1T_1.
 \tag{P3K.80.EA.17}
\]
The two coefficients of the two fibre coordinates vanish together only
at \([1:0:0]\) and \([0:0:1]\). In each corresponding affine base
chart, (P3K.80.EA.17) is exactly a nondegenerate four-variable form
\(xw-yz\). Thus both critical points and their actual circle actions
remain explicit. The physical fixed surface meets the first of these
points; the second is the isolated fixed point.

\item The actual radial and product cuts are compared by the complete domain EA.19, rather than by identifying their defining functions.

\par\medskip\noindent\textbf{RT.13. A compact cut region with its two actual corner arms.}\par\nopagebreak[4]

\item We specify the compact region before making any end identification.
In the positive reduced chart put
\[
 H=|S_0|^2+|S_1|^2+|z|^2,\quad
 \rho^2=H|T_1|^2+H^2|V|^2,\quad
 R^2=(|S_0|^2+|S_1|^2)(|T_1|^2+|zV|^2).
                                                        \tag{P3K.80.RT.19}
\]
All three displayed combinations used below are invariant under
the full auxiliary complex action; rho is the actual Fubini--Study
disk-bundle norm on W. Choose r0>0 and epsilon>0 with
0<epsilon2/r0\textasciicircum{}2<1/4 and take
\[
       \mathcal A_{\varepsilon,r_0}
       =\{\rho\le r_0,\ R\ge\varepsilon\}\subset W.   \tag{P3K.80.RT.20}
\]
This is compact because it is closed in the proper disk bundle over
CP2. It is circle invariant. Its maps q and vtilde are the restrictions
of the actual invariant polynomials; the region makes no change to
those maps or the physical line.

\item Here is a direct regularity and corner calculation. Set
\[
 A=|z|^2/H,\qquad b=H|T_1|^2,\qquad c=H^2|V|^2.
\]
Then rho2=b+c and R2=(1-A)(b+Ac). On the sphere rho=r0 put
t=b/(b+c). Thus
\[
 \frac{R^2}{r_0^2}=(1-A)\{A+(1-A)t\}=:f(A,t).       \tag{P3K.80.RT.21}
\]
The derivative in the simultaneous fibre radial direction is
2R2, so R=epsilon is regular. On the radial sphere, the only
nonzero critical ratios of R2/r0\textasciicircum{}2 are 1/4 and 1. Indeed
partial\_t f=(1-A)\textasciicircum{}2, so an interior critical point cannot occur.
On T1=0 the function is A(1-A), whose interior critical value is
1/4. On V=0 it is 1-A, with only endpoint critical values 0,1.
The remaining coordinate endpoints give 0 or 1. Consequently the
two faces in (P3K.80.RT.20) meet transversely. Their corner has the literal
equations rho=r0,R=epsilon, and may be smoothed by an invariant
rounding in these two regular defining functions.

\item The portion \(\rho=r_0,\ R\le\varepsilon\) has exactly two disjoint components.
Write \(\eta=\varepsilon^2/r_0^2\) and
\[
 A_\pm=(1\pm\sqrt{1-4\eta})/2,
\]
the inequality f(A,t)<=eta permits only A<=A\_- or A>=A\_+.
For A>=A\_+, increasing A to one decreases f; this component is a
tubular neighbourhood of the core over [0:0:1], with T1,V on their
unit three-sphere. For A<=A\_-, decreasing both A and t decreases
f. To check the latter assertion on the allowed sublevel, substitute
t<=(eta-A+A2)/(1-A)\textasciicircum{}2 in partial\_A f; its lower bound is
(1-A-2eta)/(1-A)>0. This component is a tubular neighbourhood of
z=T1=0, with V on the circle bundle O(-2) over P1.

\item The two core total spaces and their normal arrows are explicit.
At the first core, fix z=1. The core is S3 in (T1,V), the normal
coordinates are (S0,S1), and the physical weights are respectively
(+1,+1) and (-1,-1). At the second core, use V=1 with its actual
residual mu2. Its core is S3\_(S0,S1)/mu2=RP3 and its normal disk
bundle is
\[
 (S^3\times D^4)/\{(S,c)\sim(-S,-c)\}.             \tag{P3K.80.RT.22}
\]
The physical weights in that gauge are -1/2 on the core coordinates
and +1/2 on the two normal coordinates (T1,z). Its principal circle
and physical line therefore retain charge two. No free quotient by
mu2 is removed from this description. The first core retains charge
one. Both corner arms are thus part of the compact comparison;
neither is silently replaced by an ordinary sphere cap.

\par\medskip\noindent\textbf{EA.19. Explicit upper-arm filling in the retained positive chamber.}\par\nopagebreak[4]

\item The region RT.20 alone excludes the primary charge-one point.
The finite construction uses the following additional actual
domain in the same positive quotient W. Choose
\(2\varepsilon<r_1<\rho_0\) and put
\[
 \mathcal D=\{\rho\le r_1\}\ \cup\
 \{r_1\le\rho\le\rho_0,\ R\ge\varepsilon\}\ \cup\
 \{r_1\le\rho\le\rho_0,\ A\ge1/2\},
 \qquad A=|z|^2/H.
 \tag{P3K.80.EA.48}
\]
This is a compact subset of the actual disk bundle W. On A=1/2,
RT.21 gives \(R^2/\rho^2\ge1/4\). For rho at least r1 the
artificial A boundary therefore lies strictly in R>epsilon,
already retained by the middle piece. It supplies no additional
boundary face. Below rho=r1 the entire disk bundle is retained.
All intersections are identity intersections in the same smooth
positive quotient, not identifications of different GIT limits.

\item The added upper portion is in the literal z=1 chart. There
\(A\ge1/2\) means \(|S_0|^2+|S_1|^2\le1\), and
\[
 \rho^2=(1+\|S\|^2)|T_1|^2
                 +(1+\|S\|^2)^2|V|^2,
 \qquad R^2=\|S\|^2(|T_1|^2+|V|^2),
\]
\[
             q_1=S_0V,\qquad
             \widetilde v=S_0V-S_1T_1.
 \tag{P3K.80.EA.49}
\]
Thus the full upper fibre S=0 with rho at most rho0 is present,
including its unique charge-one origin. A sufficiently small
round ball around that origin lies entirely in rho<r1 and in
this single chart; it is the one ball deleted in EA.47. Its
physical weights and canonical spin-c maps are exactly those of
RT.15. No second origin is introduced by an overlap.

\item The entire zero section, and hence the exceptional disk of the
fixed surface, also lies in rho<r1. On the fixed surface z=V=0
one has R=rho. Its part above r1 therefore belongs to the
R>=epsilon piece; the whole fixed surface inside rho<=rho0
is retained. It is attached to the rest of F and capped only by
the global C0 construction. There is no separate RT.10 handle.

\item The sole removed horn of (P3K.80.EA.48) is
\(r_1<\rho\le\rho_0,\ A<1/2,\ R<\varepsilon\).
At every rho in this interval it is the lower arm, since
\(\varepsilon^2/\rho^2<1/4\). Its boundary consists of the
two specified rho cuts and the actual R=epsilon side. The
charged defining-function and full-face maps EA.29--41 are
applied to those cuts and this side, with their transported
functions. They retain q,v, the physical line and the stated
spin lift. Invariant rounding uses only these regular defining
functions, with the A-boundary absent as just proved.

\item Finally the outside separation is quantitative in the final
two-node pairs:
\[
 |q_1|+|q_2|=|S_0T_0|+|S_1T_1|
                       \le\|S\|\|T\|=R.
\]
Since \(q_1-q_2=v u_v(v)\) is the actual base-coordinate
change, on the chosen fixed bidisc
\(|v|\le C R\) and \(|T|\le C'R\).
The R<=epsilon modification can consequently be chosen disjoint
from both the prescribed outer source-parameter circle and the
outer source-base boundary. Their actual family maps remain
unchanged. This gives the upper-domain ownership in the complete boundary ledger.

\par\medskip\noindent\textbf{EA.20. The complete lower-horn boundary chain.}\par\nopagebreak[4]

\item Let H be the closure of the lower horn removed in EA.48. Its
three seven-dimensional faces are the outer rim Hr at rho=rho0,
the floor Hf at rho=r1, and the side Hs at R=epsilon. Orient H
by the ambient complex orientation. With positive rho and R
coordinate conventions its boundary chain is
\[
                       \partial H=H_r-H_f+H_s.
 \tag{P3K.80.EA.50}
\]
The retained domain D consequently has
\[
              \partial D=\Sigma_{\rho_0}-H_r+H_f-H_s.
 \tag{P3K.80.EA.51}
\]
The rim is removed from its outer radial face; the floor and
side are actual remaining boundary faces, not extra copies of
the charged end. The artificial A=1/2 face is absent by EA.19.

\item The entire horn has V nonzero, since its lower-arm parameter
u is at most \((\varepsilon/r_1)^2<1/4\). Apply the actual
EA.26 charged map, retaining the c-lambda normalization, to all
of H. Choose the angular cutoff to be one for
\(u\le(\varepsilon/r_1)^2\) and supported in u<1/2.
Then the full outer shape of EA.14 has sigma squared times f
equal to rho on the entire horn, not only on its rim. Take the
inner round charged link with squared radius r1/2. It is
strictly below H because sigma squared is at least rho, hence
at least r1 there.
The resulting full charged annulus A\_f contains the whole H,
and its outer boundary meets H exactly along the prescribed
rim. The same explicit positive shape function applies.
No radial map is used to identify source functions.

\item Put \(C=\mathcal A_f\setminus\operatorname{int}H\).
It is a compact free-circle domain with its inherited source
functions and physical line. Its complete boundary is
\[
           \partial C=Y_f-Y_{\rm round}-H_r+H_f-H_s.
 \tag{P3K.80.EA.52}
\]
Glue D to the oppositely oriented C along their floor and side
faces, using EA.26 and the exact defining functions, line maps
and spin-c lifts already specified. Their identical face
orientations before reversal are precisely why C is reversed.
The two common faces cancel. The remaining outer face is
\[
 (\Sigma_{\rho_0}\setminus\operatorname{int}H_r)
       \cup_{\partial H_r}
       -(Y_f\setminus\operatorname{int}H_r),
\]
and the only additional full face is the positively oriented
\(Y_{\rm round}\). The latter is a single charge-a link.
This is an equality of cut manifolds with their actual collars,
not a subtraction of local characteristic numbers.

\item The charged free portions of this outer face use the full
source and target maps in EA.39, with transition EA.33--36
and transported coefficient V inverse from EA.38. The outer
face also contains its fixed S3 tube and its upper z=1 arm.
Those pieces must be retained in their own actual charts:
EA.39 is not an identification of this entire mixed face.
The exact fixed-tube and upper-chart maps are recorded in
EA.21 below. The ordinary and root face agreements are the actual B and F0 arrows in EA.22.

\item All fixed strata lie in D. C and its glued faces are free.
The exceptional disk remains in the whole F, and the one
primary charge-one point remains in the upper domain of EA.19.
Thus this construction neither adds a second fixed cap nor
counts the two U/V descriptions of the charge-a end twice.

\item The gluing uses product collars obtained from the specified
invariant boundary functions and an averaged connection on
the actual line. Source functions agree on every glued face.
For a smooth real map after collar gluing, use the common face
value on a smaller product collar and interpolate to each
ambient function with an even cutoff flat at the gluing face.
This is the explicit real collar extension; it is not asserted
to be a holomorphic radial normalization. It is supported in
the inner source neighbourhood, disjoint from the actual
outer parameter family and the finite-cap annulus by EA.19.
The physical line and its derivative lift are glued by the
actual face map, constant in this product collar. The opposite
orientation on C uses the opposite spin-c collar obtained by
reflection in its specified outward real normal. That normal
is a global collar vector, so this reflection is a fixed map
and inserts no sign character on the RP3 fundamental group.
The resulting endpoint is therefore the one already carried
in EA.13, not an independently chosen RP3 lift.

\item The chain (P3K.80.EA.50)--(P3K.80.EA.52) supplies the lower-face counting
and the sign of the full residual charged boundary in EA.47.
The complete finite-strip and outside comparison is the one established in EA.21--22.

\par\medskip\noindent\textbf{EA.21. The fixed tube and free parts of the mixed outer face.}\par\nopagebreak[4]

\item The explicit fixed points z=V=0 on the outer radial sphere are retained separately from the charged product faces by the following full chart decomposition.
The following chart decomposition retains the fixed part in
the normal bundle of the whole surface F.

\item On an invariant neighbourhood of its fixed S3, use the unique
positive auxiliary gauge ||S||=1 and write T1=t exp(i phi),
with t positive. Put
\[
 p=e^{i\phi}S,\qquad \zeta=e^{i\phi}z,\qquad
 \nu=e^{2i\phi}\overline V,\qquad \|p\|=1 .
 \tag{P3K.80.EA.53}
\]
These are invariant under the remaining auxiliary circle.
On rho=rho0 the exact inverse is obtained from
\[
 t^2=\frac{\rho_0^2}{1+|\zeta|^2}
                  -(1+|\zeta|^2)|\nu|^2 ,
 \tag{P3K.80.EA.54}
\]
followed by any Hopf representative phi; changing that
representative is exactly the auxiliary identification.
For a sufficiently small normal disk the right side is positive.
Thus (P3K.80.EA.53) is a literal product of the fixed S3 with its
small real normal disk, including the zero vector. Its source
functions are
\[
 q_1=p_0\zeta\overline\nu,\qquad q_2=p_1t,\qquad
 \widetilde v=q_1-q_2 .
 \tag{P3K.80.EA.55}
\]
The compact physical circle multiplies both zeta and nu by
the same phase. They are precisely the positive-weight
spinor coefficients in EA.15; the second coefficient retains
the conjugation of the negative original normal. The
original complex normal and this positive spinor are related
by that displayed real-linear arrow, with the orientation
and physical tautological line of K68.

\item The radius on this tube satisfies
\[
 R^2=\frac{\rho_0^2}{1+|\zeta|^2}-|\nu|^2 .
 \tag{P3K.80.EA.56}
\]
Choose its normal radius so that R>rho0/2>epsilon throughout.
Its punctured part consequently meets the R>0 free norm
strip. On that overlap the physical positive end map is
the same coefficient map
(U,V) -> (d U,d\textasciicircum{}{-1}V), with d=(R/r\textasciicircum{}2)\textasciicircum{}a.
In the primary chart it is
(z,V) -> (d z,d\textasciicircum{}{-1}V); it extends over z=V=0,
fixes the source invariants, and has the stated positive
diagonal derivative there. Its full variable derivative is
the base-preserving rank-one extension from EA.10. This
specifies the overlap of the fixed-tube arrow and the free
norm arrow; no independent normal framing is selected.

\item The local fixed S3 first glues to the rest of the horizontal
normalization surface by the actual RT.1 source immersion,
the two-node map EA.42--46 and the retained Cartier r map.
It receives no local D4 cap. At the global boundary
S2 times S1, the entire resulting normal tube is carried
by the actual Cox quaternion product RT.33a and its ball
extension RT.33b to the single outside cap C0.
The normal clutch on its comparison with C is retained in
RT.29--34. In particular the local S3 and the global
S2 times S1 are connected by this actual trace, not
identified as boundary manifolds.

\item The upper arm has its own chart. In z nonzero take z=1 and
write
\[
 s_i=S_i/z,\qquad t=zT_1,\qquad w=z^2V .
 \tag{P3K.80.EA.57}
\]
The physical weights are (-1,-1,+1,+1) on
(s0,s1,t,w), and
\[
 q_1=s_0w,\qquad q_2=s_1t,\qquad
 \rho^2=(1+\|s\|^2)|t|^2+(1+\|s\|^2)^2|w|^2 .
 \tag{P3K.80.EA.58}
\]
These are the literal upper chart of the same positive
quotient, including s=0. Its origin belongs once to D;
the deleted charge-one ball and its source functions are
those of EA.19.

\item On the overlap w nonzero the complete map to the primary
charged V=1 chart is
\[
 (s_0,s_1,t,w)\longmapsto
 (\sqrt w\,s_0,\sqrt w\,s_1,\ t/\sqrt w,\ \sqrt w)
 \tag{P3K.80.EA.59}
\]
in the ordered coordinates (S0,S1,T1,z). The two roots
differ by the actual residual mu2 action. This map preserves
q1 and q2, is equivariant for the same physical circle,
and carries its canonical complex spin-c map through the
specified auxiliary gauge. Its full determinant and line
characters, including this residual action, are retained.
It is the overlap of the upper chart and the lower charged
chart, not a map between their origins.

\item At w=0,t nonzero the upper chart is still free and has R=0
when s=0. The following actual T1 chart and Tate-period map
retain this endpoint without using an R>0 flow.

\item On T1 nonzero take the auxiliary gauge T1=1, with coordinates
\[
 (A,B,C,D)=(T_1S_0,T_1S_1,T_1z,T_1^{-2}V).
 \tag{P3K.80.EA.60}
\]
Its physical weights are (0,0,+1,-1), and
q1=A C D, q2=B. The upper-chart arrow is
\[
 (s_0,s_1,t,w)\longmapsto
                   (ts_0,ts_1,t,w/t^2).
 \tag{P3K.80.EA.61}
\]
Thus it includes the entire w=0,t nonzero core.
The radial face rho=rho0 is covered by this chart and V
nonzero, since rho>0 implies (T1,V) is not (0,0).
On their overlap choose h with h squared equal to D. The
map to the ordered V=1 cover coordinates is
\[
 (S_0,S_1,T_1,z)=(hA,hB,h^{-1},hC).
 \tag{P3K.80.EA.62}
\]
Both choices of h give the same point of the specified mu2
quotient. The determinant of this four-coordinate map,
in the ordered variables (A,B,C,D), is +1/2. Consequently
the actual charged contraction volume
-2 dS0 wedge dS1 wedge dT1 wedge dz pulls back to
-dA wedge dB wedge dC wedge dD, which is the contraction
volume in the T1=1 chart. This is the full canonical
spin-c overlap, including its root deck, rather than an
orientation assigned from the coordinate names.

\item Here is the exact additional free-line chart. On the dense
overlap used by TL take its source coordinate S=s0/s1
and its tensor-Cartier coefficient \(\xi=s_0/s_1^2\).
The latter denominator is required by auxiliary weight two.
The full TL formulas give
\[
 W_0=s_0t/u_P,\quad Z_0=u_Gs_1w,\quad
 V_0=E_+^2s_0/s_1^2,\quad
 Q=q_1q_2u_Q,\quad P=u_P/q_2,\quad u_G=u_Qu_P .
 \tag{P3K.80.EA.63}
\]
Apply the actual inverse period h inverse to all coordinates,
then the indicated integral toric character map:
\[
 (Z_0,V_0)\longmapsto(Z_0/Q,V_0/P^2)=(1/W_0,V_0/P^2),
 \qquad (Z,V)\longmapsto(Z^{-1},ZV).
 \tag{P3K.80.EA.64}
\]
The resulting coordinates extend as
\[
 Z_t=s_0t/u_P,\qquad T_t=u_Pw/t,\qquad
 L_t=E_+^2t/u_P,\qquad q_2=s_1t .
 \tag{P3K.80.EA.65}
\]
In particular q1=Zt Tt and Lt is nonzero when t is nonzero.
Every theta and Cartier unit in E+ is retained in (P3K.80.EA.65).
The full inverse is
\[
 t=u_PL_t/E_+^2,\qquad w=T_tL_t/E_+^2,\qquad
 s_0=E_+^2Z_t/L_t,\qquad
 s_1=E_+^2q_2/(u_PL_t).
 \tag{P3K.80.EA.66}
\]
Here uP and E+ are evaluated on the recovered full source
chart (Zt,Tt,q2). They are independent of Lt. For example,
S=P Zt and Q/S=Tt q2 squared uQ/uP in the full theta
product, while S/P=Zt and Q/Zt=Tt q2 uQ. All factors
after the extracted 1-S have extending unit logarithms on
the chosen t-dominant collar. Thus (P3K.80.EA.65)--(P3K.80.EA.66) give
a holomorphic free-line chart through w=0,t nonzero, not
only formulas on the dense source torus.

\item The character map in (P3K.80.EA.64) has the literal lattice matrix
\[
 J=\begin{pmatrix}-1&0\\1&1\end{pmatrix},\qquad
 \widehat J(x,y,h)=(-x,x+y,h).
 \tag{P3K.80.EA.67}
\]
For the actual negative-square T0 fan, J sends
{0,e1,e2} to {e1,e2,e1+e2}-e1, and sends the latter
triangle to {0,e1,e2}-e1+e2. The full integral height
translations are retained in these statements. The preceding
inverse h period in (P3K.80.EA.64) is the actual period (Q,P squared),
not just its valuation. All period maps on this new chart are
its literal conjugates through (P3K.80.EA.64).

\item The source marking in this chart is explicit:
\[
 W_0=Z_t,\qquad Z_0=u_Q q_2 T_t,\qquad q_1=Z_tT_t .
 \tag{P3K.80.EA.68}
\]
It is the (Z0,q2) small-resolution chart. TL's original
chart used (W0,q2). The two source projections are the two
projections of the actual common translation graph, and
their comparison must use the source -G map and its
Cartier reset EA.18--20. This resolution arrow, with the
full r map, is retained in the endpoint audit. It is not
replaced by an identification of central Picard classes.

\par\medskip\noindent\textbf{EA.22. The actual upper closing arrow and the indexed ordinary end.}\par\nopagebreak[4]

\item The opposite-resolution arrow in EA.68 is supplied by the
literal inverse suspension. Retain the fixed nonzero scalar
\(c_\lambda=2\lambda\), \(c_X=q_2u_X\), and the full
SZ.43 Cartier-normalized primitive. In the original torus
coordinates the forward map and its inverse are
\[
 B:(Z,V)\longmapsto
       (P^{-1}Z,c_\lambda q_1c_X^{-2}V),\qquad
 B^{-1}:(Z,V)\longmapsto
       (PZ,c_X^2V/(c_\lambda q_1)).
 \tag{P3K.80.EA.69}
\]
This is the reduced a=2,b=1 map. The cX square in the
inverse is the actual source/target r-conjugated line
map proved in SZ.41--43; it is not an inferred action
on generic Picard classes.

\item Write (W,T,L) for (Zt,Tt,Lt). Before the inverse
suspension, the full target torus coefficients recovered
from EA.63--66 are
\[
 Z_0=q_2u_QT,\qquad V_0=P^2WL
                  =\frac{u_P^2}{q_2^2}WL .
 \tag{P3K.80.EA.70}
\]
Applying EA.69 gives the exact expressions
\[
 Z_{\rm pre}=u_GT,\qquad
 W_{\rm pre}=Q/Z_{\rm pre}=q_2W/u_P,\qquad
 V_{\rm pre}=\frac{u_P^2u_X^2}{c_\lambda T}\,L.
 \tag{P3K.80.EA.71}
\]
The source is now precisely the marked TL resolution:
its coordinate S=P Wpre is W and its other node
coordinate is T. Thus the source part of B inverse,
which is the actual +G map, carries the opposite
(Z0,q2) resolution of EA.68 to the prescribed
(Wpre,q2) resolution. This is also the actual inverse
of the graph formulas SZ.14 and EA.18. Both projections
retain the same q1,q2 and the chosen v coordinate.

\item The last expression of EA.71 is read in its actual
toric edge chart. The edge is
\[
 \{(0,0,1),(1,-1,1)\}
       =\{(1,0,1),(0,1,1)\}-(0,1,0).
 \tag{P3K.80.EA.72}
\]
It is the negative diagonal of the square translated
by (0,-1). Its coordinates are
\[
 x=Z_{\rm pre},\qquad y=q_1/Z_{\rm pre},\qquad
 L_{\rm pre}=Z_{\rm pre}V_{\rm pre},
 \qquad xy=q_1,\quad L_{\rm pre}\ne0.
 \tag{P3K.80.EA.73}
\]
The invertible character on this indexed edge is ZV.
On the unshifted edge it would be ZV/q1. The actual
height translate in EA.72 therefore retains precisely
the adjacent-node j to j-1 change in the inverse
suspension; it is not applied a second time.

\item In the marked source coordinates (W,T,q2), the
complete upper closing map is now
\[
 (W,T,q_2,L)\longmapsto
           (W,T,q_2,F(WT,q_2)L),\qquad
 F=\frac{u_Gu_P^2u_X^2}{c_\lambda}.
 \tag{P3K.80.EA.74}
\]
Every factor is the full extending unit from SZ.39--43.
The retained constants uQ(0)=1/256, uP(0)=16 and
uX(0)=-1/4 give F(0)=1/c\_lambda. The quotient
F/F(0) has an extending logarithm on the source
bidisc. The actual fixed scalar and its chosen phase
remain in F(0). Equation EA.74 is invertible, with
complex determinant F. It extends through T=0 and
W=0 and includes the entire t-dominant upper core.
Its inverse is the forward finite-face pairing.
The line is the physical weight-one line in both
charts. In particular no dual circle is introduced
by the source inversion in EA.64.

\item Here is the complete pairwise collar comparison.
The radial face has the open cover V nonzero and
T1 nonzero, since rho is positive. On the T1 chart
the free locus is C nonzero or D nonzero. The part
D nonzero is the V-chart overlap EA.62. The part
C nonzero is the upper z=1 chart with t=C nonzero,
and is carried by EA.65--74 to the actual indexed
next end. The remaining set C=D=0 is exactly the
fixed-tube locus. Thus these named charts cover
the entire mixed face, including its R=0 free
points and its fixed S3.

\item On the fixed-tube chart EA.53 the map to EA.60 is
the literal formula
\[
 (A,B,C,D)=(t p_0,t p_1,t\zeta,\overline\nu/t^2).
 \tag{P3K.80.EA.75}
\]
It gives q1=p0 zeta conjugate(nu), q2=t p1.
Its normal map is the actual positive spinor
with the retained negative-coordinate conjugation.
On z nonzero EA.61 is its inverse change to the
upper chart; on V nonzero EA.62 is its change
to the charged chart. Therefore the fixed tube,
upper free chart and charged chart agree on all
their pairwise and triple overlaps by the displayed
coefficient maps. The source functions agree
literally. The original r and kappa maps precede
these formulas; their full coefficients enter E+
and the cX square in EA.69.

\item The R>0 finite norm strips are attached with their
actual transported norms from EA.37--41. Their
overlap with the upper and fixed charts can be
chosen inside R>=epsilon. On that overlap the
same positive scalar d=(R/r\textasciicircum{}2)\textasciicircum{}a is used on both
maps. To define the cut on the upper chart down
to R=0, take
\[
 \widetilde d(R)=
 \exp\!\left(\chi(R)\,a\log(R/r^2)\right),
 \tag{P3K.80.EA.76}
\]
where chi is one for R>=epsilon and zero on a
neighbourhood of R<=epsilon/2. Set dtilde=1 on
that neighbourhood. This is a positive smooth
coefficient, applied only to transport the cut.
It preserves q1,q2 because R is invariant along
the physical real orbit. It agrees with the
specified finite norm map on its entire attachment
collar and with the exact upper endpoint near its
core. The full charged cores in EA.39 remain their
own source/target faces; EA.76 is not an extension
of the interchange EA.40 across those cores.

\item The finite-cap compatibility follows from the
same actual inverse arrow:
\[
 F_0 B^{-1}=(c_\lambda h_\infty)^{-1}
                         U_{\rm cap}^{-1}F_0,
 \qquad q_1=q_\infty^{-1}h_\infty .
 \tag{P3K.80.EA.77}
\]
This is obtained by inverting the full square in
EA.12, with its root source, r correction, norm sign,
and finite character retained. The full Tate period H\_T in
EA.64 is carried as its full lattice arrow through
F0. The positive cut functions are transported
through F0 on its overlap and averaged over its
actual finite group, as in EA.5. The small-R
modification is disjoint from that annulus by
EA.19, so it introduces no new cap face.

\item For the spin-c comparison, the upper map is the
canonical complex map EA.69 followed by the
actual toric chart EA.73. Its extending unit in
EA.74 has the specified fixed scalar and the
extending logarithm of F/F(0). On the V overlap
the two square roots in EA.62 differ by precisely
the canonical residual mu2 action of EA.13/RT.15.
On the fixed tube EA.75 uses the actual original
complex normal and its written real positive
spinor map. Thus the maps on a triple overlap
are lifts of the same actual composite, with
the determinant factors already fixed in EA.27,
EA.62 and EA.74. There is no additional choice
of an RP3 sign character. The local fixed tube
joins the rest of F before the single global
C0 cap of RT.33a--b is attached.

\item These are the full upper endpoint and mixed-face overlaps used in the finite boundary identity EA.47.

\par\medskip\noindent\textbf{EA.18. Finite cut-and-paste ledger for the reduced branch.}\par\nopagebreak[4]

\item This section assembles the finite reduced-branch attachment.
It uses a=2,b=1. It does not copy the different a=1 moment picture.
All scalar windings have first been separated by the actual based
homotopy RT.12, so the scalar in this construction is the fixed
nonzero value at r\_base. The unscaled infinite central quotient of
SZ.29 is not used.

\item Take the actual compact source translation graph. Its domain and
target maps outside the central curve are the literal addition
formula and the Cartier line map of SZ.7 and SZ.14--17. At that
curve use both projective charts of the common Cox graph: the
domain t-ruling includes its t1=0 chart at the zero section, and
the target s-ruling includes both divisor charts of G. Thus the
whole proper central component and the whole exceptional component
are retained, not only their affine node points. On a small graph
neighbourhood use the analytic node map (P3K.80.EA.42)--(P3K.80.EA.46). Outside
it the actual source/line map is an isomorphism between the two
graph projections; choose the ordinary norm cut and transport its
target face by that map. The finite-root source/line map is the
actual U\_cap of EA.4, with F0 on its annulus. Compactness of the
remaining source and root pieces gives one common contracting
scalar range.

\item Cut before taking the period quotient. On the compact complement
of the graph neighbourhood this is the ordinary fundamental
annulus between a norm hypersurface and its actual B-image. On
the stable graph collar use the finite region between Ht and Hs
from EA.39, with R bounded below by a positive epsilon. The
actual end scalar and its phase are retained before the positive
normalization EA.40. The change of norm cut is an explicit
monotone interpolation of the two logarithmic norm coordinates,
fixed at their transported endpoints; on the node charts its
extension is the positive diagonal map in EA.13. Thus no raw
logarithm of q is extended through q=0. The q factor is absorbed
by the literal j to j+1 node arrow, and the actual cyclic graph
reset and j to j-1 overlap are EA.18--26.

\item The selected node is completed in the primary positive associated
chart. TL.1--4 identifies its boundary, through the full theta,
Cartier and h-period map, with current Tate T0. On the positive
moment piece use the complete retained domain EA.48 and its invariant
corner rounding. It includes the upper fibre and charge-one point. Its defining functions are the actual R and
rho, and q,v are retained by ambient inclusion. The lower arm
is the transported shape EA.29--30. Its full two faces are
EA.39, with their actual source/target line coefficient EA.38;
their full orbit transition is EA.33--36. The upper arm stays
in the primary positive piece. It is not relabelled as a
charge-two arm or silently capped by a second copy of the
lower chart.

\item The following face pairings are the ledger for this finite
construction. They specify maps, not merely a cancellation of
numbers.

\begin{itemize}
\item Ordinary source and target norm faces are paired by the actual
B map, with its Cartier coefficient. Their boundary normals
are opposite. On the graph collar this is the same pairing
after the explicit positive norm interpolation just described.
\item Every finite-root face is paired by the literal H or A map of
EA.12. The F0 square identifies it with the ordinary B pairing;
both norm signs, all finite characters, and the root-component
translations remain. There is no residual cap face from replacing
a finite root by its generic source point.
\item The inner source collar and the ordinary source exterior are
glued by TL's full source/line map. On every node chart its
monomial factors and units are the actual ones, including the
h reindex. The two source-base normals are opposite. The
source functions and their parameter arrow agree literally.
\item In the chosen finite B strip the complementary U and primary
V descriptions are one charged end. They are related by EA.26,
not added as two positive GIT charts. After all other period
faces have been paired, the target-minus-source finite face
cycle is the single positive charged boundary Sigma\_a of
EA.41. The q reindex, scalar phase and determinant root are
retained. In particular the lower charged origin is not
quotiented by its infinite central stabilizer.
\item The fixed surface is the entire
F=Bl(P1 times D), including its exceptional disk. Its annular
part in RT.20 carries the same two normal frames on both
sides. It is capped only by the actual outside cap C0=D3
times S1. No separate local D4 cap from RT.10 is inserted.

\end{itemize}
Every gluing is between disjoint closed collar faces by the
specified diffeomorphism. The construction is finite: it uses a
compact norm interval, both charts of each projective ruling,
the finite root atlas and the two compact corner arms. Thus the
resulting glued space is compact and Hausdorff; it is not obtained
by gluing two open GIT quotients with distinct limits. Invariant
rounding of the stated regular boundary functions gives its smooth
collars. On the outer nonzero v circle the cuts are simply cuts
of the original proper family: choose epsilon so small that the
R-epsilon region lies strictly inside that parameter circle.
The full TL and F0 maps therefore identify the completed outside
boundary with that actual parameter return.

\item The spin-c and physical-line identifications on these face
pairings are the actual coefficient and derivative maps. The
ordinary and finite-cap maps are holomorphic. The norm
interpolations have the explicit positive and rank-one derivative
homotopies in EA.10 and RT.15. The graph reset has the exact
determinant root i*z\_s\textasciicircum{}(-2), conjugated by the full extending
kappa units; the charged overlap has its constant root i.
The mu2 endpoint is the canonical +omega8 in EA.13/RT.15.
These are composed in the same order as the actual face maps.
On an ordinary triple overlap their composition is the
canonical complex spin-c lift of the actual coordinate identity;
on a root overlap it is the specified finite-character action.
No additional sign character is inserted at the RP3 face.

\item The complete face maps give the following oriented boundary equation.
Let N7 be the actual free-circle total boundary obtained by the
outside conifold/Hopf completion, and let N6=N7/S1. Delete the
entire capped fixed tube and the primary positive chart's
charge-one ball. Replace the fixed cap C0 by C=P1 times D2
with the actual normal-map triangle of RT.17. Give C and C0
the same boundary orientation as partial F. Then
\[
 [F-C_0]=[F-C]+[C-C_0],\qquad C-C_0=S^4.
\]
Write B\_F and B\_4 for their actual projectivized normal sphere
bundles, with physical tautological line. The finite free
circle quotient has boundary cycle
\[
       -N_6-\mathbb P^3_a-B_F-B_4+\mathbb P^3_1.
 \tag{P3K.80.EA.47}
\]
Here P3\_a denotes the charged boundary with M=O(-a), not a
discarded residual quotient; P3\_1 is the primary positive
isolated-point boundary. The sign of P3\_a is fixed by the
positive total-space end chain EA.41 and the circle-first
boundary reversal. The deleted fixed tube contributes -B\_F
and -B\_4 by the literal cap identity. The two negative normal
conjugations at the isolated point give +P3\_1, as in RT.13.
The exceptional disk belongs to F once and never to a separate
local fixed-S3 handle.

\item The physical line and stable tangent bundle on this cycle are
the ones carried by the actual free circle action. The S4
normal clutch is retained; no value of its c2 is imposed.
RT.16--17 computes its index and the full fixed-surface term
with both actual parameter windings. The complete outside and cyclic endpoint maps in EA.19--22 establish this boundary identity for the actual family.

\par\medskip\noindent\textbf{RT.16. The whole fixed trace, its actual cap clutches, and the index.}\par\nopagebreak[4]

\item We keep the whole surface F=Bl(P1\_s times D\_v), including its
exceptional disk. No copy of the local fixed-S3 handle of RT.9 is
inserted in this calculation. Let Y=boundary F=S2 times S1, with
its complex boundary orientation. Put c(v)=sqrt(1+v), on the actual
branch equal to one at zero. The marked coordinate

\[
 Z=(s-c(v))/(s+c(v))
\]

is induced by the full matrix ((1,-c),(1,c)); its determinant 2c
is nonzero on the whole parameter disk. Thus this is an extending
Mobius map on the unblown P1 family, not a replacement of its
parameter loop by a constant matrix. It takes D+ to zero and D-
to infinity. The blowup centre lies on D- and not on D+.

\item Cap by a second P1 times disk using this actual boundary coordinate.
Write C for this cap with orientation chosen so that its boundary
is Y. Equivalently, -C has the complex orientation of the second
parameter chart w=constant/v. The closed oriented surface F-C is
the complex blowup of P1 times P1. Let H and B be respectively the
two factor classes, and retain E for the exceptional class. Then

\[
 H^2=B^2=0,\quad H B=1,\quad E^2=-1,\quad
 H E=B E=0,\qquad \sigma(F-C)=-1.                 \tag{P3K.80.RT.27}
\]

\item In the first parameter disk the actual divisors have classes
D+=H and D-=H-E. The positive spinor of the physical normal
model is the same tensor sum carried in RT.1 and RT.9:

\[
 O(D_+)\otimes O(-aE)\ \oplus\
 O(-D_-)\otimes O(-aE).                           \tag{P3K.80.RT.28}
\]

\item The tensor factor is the actual source line. At a=2 this statement
includes the source and target Cartier r maps, not just the tensor
line's divisor class. The full F0 cap maps act on its nonzero
coefficients. The negative normal coefficient is conjugated by the
same physical positive-weight arrow as in K68; no line dual is
silently substituted.

\item Here is how all parameter and cap clutches are retained on the
second disk. Give its two summands the reference divisor frames
of O(H) and O(-H), in the displayed Z coordinate. On Y compare
them with (P3K.80.RT.28), transported by the actual source/Cartier/F0
maps and the actual Cox normal frame. Denote the resulting two
nonzero coefficient ratios by g+ and g-. They are functions on
Y because the common O(1) and O(-1) point transitions cancel.
In a finite s chart these are the ratios of the reference frames
1/(s-c) and s+c to the actual frames of (P3K.80.RT.28), including their
full source-line coefficient and normal units. At their divisor
points the same ratios are taken in the alternate Cartier chart;
the ratios agree, so no zero or pole is introduced at D+ or D-.
The finite character of the cap is part of these ratios. Its
constant value is not replaced by one.

\item Use the frame convention e\_cap=g e\_F. Define the two actual
integers by

\[
 w_\pm=\frac{1}{2\pi i}\int_{\{p_*\}\times S^1}
                              g_\pm^{-1}dg_\pm .   \tag{P3K.80.RT.29}
\]

\item Here p* is a regular point in the marked sphere, away from the
Cartier charts' divisor points. The integer is independent of p*
because g is a nonzero function on the connected product Y.
The direction of (P3K.80.RT.29) is fixed as follows: e\_cap=v e\_F is
the actual transition of O(1) on the doubled parameter sphere,
so a winding w contributes +w B. This also fixes the sign of
the base classes below.

\item There is an exact smooth line map reducing this clutch to its
displayed integer without losing the rest of the map. The ratio
g exp(-i w theta) has zero class in H1(S2 times S1;Z), and
therefore has a single-valued smooth logarithm h. Extend h over
S2 times the cap disk by the parameter-disk Poisson extension,
with a collar interpolation preserving its boundary value. Changing
the cap frame by exp(-h) gives precisely exp(i w theta) on Y.
This is a nonvanishing frame change on that specified cap. It
does not change the original F0, source or normal map on the first
piece. In particular it does not declare either w to be zero.

\item It follows that the actual capped spinor is isomorphic, with these
boundary arrows retained, to the sum of lines with Chern classes

\[
 l_+=H-aE+w_+B,\qquad
 l_-=-H+(1-a)E+w_-B.                               \tag{P3K.80.RT.30}
\]

\item Consequently, for b=c1(E\_spin) and c2=c2(E\_spin),

\[
 b=(1-2a)E+(w_++w_-)B,\qquad b^2=-(1-2a)^2,
\]
\[
 c_2=-a(a-1)+(w_--w_+).                            \tag{P3K.80.RT.31}
\]

\item Thus even the unequal normal-clutch windings have been retained;
their cross term occurs in c2. The determinant has no H term,
which is the precise reason there is no B cross term in b squared.
The local canonical complex spin-c arrows, their determinant
lifts on the charged collars, and the physical line are those
of RT.15 and EA.13. The cap construction carries those actual
bundle maps. It makes no additional choice of a spin character
on the lower RP3 arm.

\item For completeness the characteristic cancellation is obtained from
the physical line itself. On P(E\_spin) put h=c1(O(1)); the actual
physical nonzero-vector line is M=O(-1), hence x=-h. The same
Hopf coefficient map on both caps gives this line arrow. The
projective relation and fibre integrals are

\[
 h^2+b h+c_2=0,\quad \pi_*h=1,\quad
 \pi_*h^2=-b,\quad\pi_*h^3=b^2-c_2.
\]

\item The vertical complex tangent has Chern class 2h+b, so its first
Pontryagin class is b squared minus 4c2, pulled back from the base.
Thus the two physical characteristic numbers, with the stated
base and complex fibre orientations, are

\[
 k=c_2-b^2,\qquad
 \ell=-\langle p_1(F-C),[F-C]\rangle-b^2+4c_2,
\]
\[
 \frac{4k-\ell}{24}
 =\frac{\langle p_1(F-C),[F-C]\rangle-3b^2}{24}
 =\frac{\sigma(F-C)-b^2}{8}=\frac{a(a-1)}{2}.           \tag{P3K.80.RT.32}
\]

\item Neither the normal c2 nor either actual clutch integer was set to
zero in this computation. The full source and finite-cap maps
enter (P3K.80.RT.29); their possible cross terms cancel by (P3K.80.RT.32).

\par\medskip\noindent\textbf{RT.17. The two actual caps and the fixed-boundary correction.}\par\nopagebreak[4]

\item The outside Jin completion uses its actual product cap
C0=D3 times S1, with boundary Y and the retained Cox/Hopf normal
coefficient frame. Its normal bundle and its physical line are
the ones specified by that frame. They are not replaced by the
reference divisor frame used to calculate (P3K.80.RT.30).

\item The actual boundary arrow defining this cap is the following.
Use the full TL.1--3 node/Cartier map to the selected fixed-edge
charts of current Tate T0. In the boundary parameter annulus their
native coordinate is the displayed Z, and their two linear normal
coordinates have overlap c\_i'=Z c\_i and physical weights (-1,+1).
The units in TL.11 and the degree-two r arrow are included in
these coefficient maps. This is the literal two-cone fixed-star
chart, not a statement obtained from its normal degrees alone.
An actual nonzero chart unit on this annulus has the exact form
v\textasciicircum{}d exp(h); its integer d and its full coefficient are retained.
On each of the two physical weight lines, the linear normal
automorphism over the compact fixed P1 is a scalar function of
the parameter, because H0(P1,O)=C. The higher normal terms are
carried to their linear maps by equivariant normal rescaling on
a common smaller tube. The rescaling is the identity on the zero
section and is followed by its collar isotopy, so it carries the
original line and annular maps rather than discarding them.

\item In homogeneous unit Cox factors put

\[
 u=a_0+j a_1,\qquad v_N=c_0+j\overline{c_1},\qquad
 h=u i u^{-1},\qquad q_H=u v_N.                    \tag{P3K.80.RT.33a}
\]

\item The auxiliary circle cancels in q\_H. Choosing a Hopf representative
of h and putting v\_N=u\textasciicircum{}{-1}q\_H is its explicit inverse. On the
finite Z chart it is the normalized expression
(1+jZ)(c0+j conjugate(c1)). These are exactly the first-resolution
products of K23, equations SF.5--6 of
part3\_source\_cusp\_fixed\_edge\_working.md, applied here to the actual
TL fixed-star coefficient maps. Their formulas are independent of
which source supplied those maps. At nonunit normal radius the
same expression is real linear in the normal vector and sends
zero to zero. Thus it identifies the entire small normal tube,
not just its link. The physical action is
q\_H -> q\_H exp(-i phi), with its actual character.

\item All parameter powers act through their actual Cox coefficients
in (P3K.80.RT.33a). They are single-valued integer powers on this annulus.
This yields a periodic normal product over the parameter circle;
it does not claim that its comparison to a different product
extends over the cap. That comparison is the normal clutch on
the S4 in (P3K.80.RT.33), and is retained below. The underlying sphere
coordinate change has the literal extension

\[
 P\longmapsto A_{c(v)}P A_{c(v)}^*/
                    \operatorname{tr}(A_{c(v)}P A_{c(v)}^*),
 \qquad A_c=\begin{pmatrix}1&-c\\1&c\end{pmatrix},  \tag{P3K.80.RT.33b}
\]

on the trace-one positive-Hermitian three-ball. Its denominator
is positive, its inverse uses A\_c inverse, and its boundary map
is exactly Z. Holding the actual q\_H coordinate fixed gives the
product cap arrow (P,q\_H,v)->(B\_c(P),q\_H,v). It is the explicit
K42 SP.13/K68 physical-circle cap arrow, now with the extending
c(v) and the actual TL/Cox coefficient product. It preserves
the physical line and normal orientation. No tangent derivative
is substituted for this product arrow.

\item The positive spinor on this cap is typed explicitly: its two
complex coefficients are the conjugates of the right-complex
coefficients of q\_H. They have positive physical weight because
q\_H has the negative right phase. If e denotes the positive
Cartier spinor on F, K68's actual arrow is q\_H=G\_Z conjugate(e),
including its fixed ordering and constant normal comparison.
It therefore becomes the complex-linear arrow
conjugate(q\_H)=conjugate(G\_Z)e to the cap's positive spinor.
The physical associated line is its O(-1), with the original
circle character. The normal antipodal/meridian reversal in
the first arrow is retained by this coefficient choice. It is
not replaced by a complex-linear swap of e and q\_H. Thus the
bundle used on the cap and on the S4 below is an actual complex
rank-two bundle with the displayed line morphism.

\item The two cap choices have the exact closed base

\[
 S_4=C-C_0
     =(S^2\times D^2)\cup_{S^2\times S^1}
                        -(D^3\times S^1)\cong S^4, \tag{P3K.80.RT.33}
\]

with the orientations chosen above. The last map is the rounded
boundary of D3 times D2; its two product faces are precisely the
two written caps, with their opposite boundary orientations.
Their actual normal coefficient maps along Y define the normal
bundle on this S4 by gluing, including every difference of the
Cartier and Cox frames. The induced physical Hopf line is the
tautological O(-1) on its projectivization. There is no unrecorded
circle gauge: the coefficient ratios on Y are the same g functions
as above, and a circle gauge's S1 degree extends on C0 because
H1(C0;Z) maps isomorphically to H1(Y;Z).

\item There is an explicit bordism implementing this cap replacement.
Take the products of F,C,C0 with intervals and join their boundary
collars by Y times a triangle, one interval product on each edge.
Use the three boundary identifications furnished by the actual
normal frames; the third is the inverse composite of the first
two. They therefore extend over this collar triangle. Rounding
its product corners gives a compact oriented five-dimensional
base bordism with the three closed four-dimensional faces
F-C0, F-C and C-C0. The normal bundles glue on the same products,
and projectivizing gives the required sphere-bundle bordism with
the physical line. This proves the literal cycle identity

\[
 [B_{F-C_0}]=[B_{F-C}]+[B_{S_4}]                 \tag{P3K.80.RT.34}
\]

with the retained normal maps, rather than only an equality of
intersection forms. A nonzero third-class normal clutch may be
present on S4. Its c2 is carried by this construction. Applying
the same formula (P3K.80.RT.32) gives index zero, because H2(S4)=0
and its signature is zero. The c2 term cancels explicitly.

\item Deleting the entire fixed tube in the global orbit construction
therefore supplies -B\_(F-C)-B\_(S4). The exceptional E remains
in the first base. This is the fixed-cap part of the boundary
ledger; it must be combined with the finite-period face chain
and the actual outside/Jin arrow before a global return is assigned.

\item The additional exact certificates are
code/part3/source\_zero\_cut\_spin.py (67 local cut, charged-overlap,
Clifford and spinor identities) and
code/part3/source\_zero\_fixed\_cap.py (462 whole-cap identities,
including arbitrary retained w+ and w-). Their JSON and captured
outputs retain these local scopes. The fixed-cap program's first
JSON serialization error is preserved under
records/attempts/source\_zero\_fixed\_cap\_json\_types; all its
mathematical assertions had completed, and the repair converts
only SymPy output integers to native JSON integers.

\par\medskip\noindent\textbf{RA.3. The physical characteristic equation and actual return.}\par\nopagebreak[4]

\item Let \(N_7\) be the actual completed outside boundary of the
finite construction. The physical circle is free after
the conifold/Hopf replacement on its selected fixed tube.
Let \(N_6=N_7/S^1\), with the Jin orientation of P3K.68.
On every free part of the seven-dimensional quotient
bordism the associated line is the actual nonzero-vector
line \(M\), and \(x=c_1(M)\). The tangent splitting is the
quotient tangent plus the oriented physical circle. Thus
\(p_1\) is the pullback of the quotient \(p_1\).
Choose invariant product connections on the named collars
and patch them on the compact interior. Stokes applied
to the closed Chern--Weil polynomial gives the exact
oriented characteristic identity
\[
 -A(N_6)-A(\mathbb P^3_2)
       -A(B_{F-C})-A(B_{S^4})+A(\mathbb P^3_1)=0,
 \qquad
 A(X)=\int_X\frac{4x^3-p_1(TX)x}{24}.             \tag{P3K.80.RA.3}
\]
The boundary signs are those of EA.47. The full charged
end occurs once by the horn chain, the exceptional disk
belongs to the whole \(F\), and the sole fixed cap is \(C_0\).
The two auxiliary charged descriptions are related by
EA.26 and contribute one boundary.

\item On a positively oriented charged projective boundary the
physical line is \(O(-a)\). Hence, writing \(h\) for the
positive generator, \(x=-a h\) and \(p_1=4h^2\), so
\[
 A(\mathbb P^3_a)=-\frac{a(a^2-1)}6.
\]
For \(a=2\) this is \(-1\); for the isolated charge-one
boundary it is zero. The whole capped surface calculation
RT.27--32 retains both actual normal clutches and gives
\(A(B_{F-C})=1\). The actual comparison cap \(S^4\) has
\(A(B_{S^4})=0\), independently of its retained normal
\(c_2\). Therefore (P3K.80.RA.3) gives
\[
                     A(N_6)=-(-1)-1-0+0=0.      \tag{P3K.80.RA.4}
\]
This is a calculation on the specified compact finite
bordism. The physical line and canonical spin-c overlap
maps, including the residual \(\mu_2\) endpoint, have
already been carried through every face.

\item For the actual return, take a delta-averaged connection on
the original completed family, product on the selected
tube in the exact RT.33a--b coordinates. Its quotient after
the conifold replacement is a diffeomorphism
\(f:S^2\times S^3\to S^2\times S^3\), fixed on that tube.
The preserved Euler class generates \(H^2\), and the
orientation then preserves \(H^3\); thus \(f\) belongs to
the required homology-trivial relative group.
The Hopf lift in P3F.22 is literally the physical return
lift on the complement and uses the same identity
extension on the selected tube. P3K.68 supplies the full
piece and inverse-clutch arrow
\[
       \lambda_M(f)=-\Sigma(f^{-1})=\Sigma(f).
\]
Consequently its class in \(\Theta_7\) is the physical
integer (P3K.80.RA.4), namely zero. This step selects the
actual parameter return by its source, cap, tube and
line maps, rather than only by its characteristic polynomial.

\item For the literal reflected fan, the already proved single
native directed flop of P3K.60 has the actual pair
\((1-a,a)=(-1,2)\). Its signed relative Cox conjugation
adds the right exponent \(a=2\) to the central class.
The complex normal orientation and the physical transfer
are precisely those retained in P3K.68--69. Thus the
reflected fixed-scalar class is \(2\). No division of a
covered \(\Theta_7\) class is made.

\item Finally the actual scalar-root loop has physical index
\(3p-1\) in these two fans: P3K.24 with the physical
correction P3K.68 gives the initial value, and the narrow
actual C2 comparison P3K.62 gives their equality.
The exact based parameter homotopy RT.12 adds \(k\) such
loops, with basepoint \(r_{\rm base}=v_*^k r_0\).
It therefore gives (P3K.80.RA.1), retaining either norm sign
and all prescribed finite characters.

\par\medskip\noindent\textbf{RA.4. Exact certificates and proof scope.}\par\nopagebreak[4]

\item The full source/Cartier, node and shape calculations are
executed in the following retained certificates:
\path{code/part3_fourfold/source_zero_two_node_cox.json},
\path{code/part3/source_zero_tate_label.json},
\path{code/part3_fourfold/source_zero_associated_reset.json},
\path{code/part3_fourfold/source_zero_chamber_ledger.json},
\path{code/part3_fourfold/source_zero_spin_chart.json},
\path{code/part3_fourfold/source_zero_charged_shape.json},
\path{code/part3_fourfold/source_zero_native_charged_seam.json},
\path{code/part3_fourfold/source_zero_full_faces.json},
\path{code/part3_fourfold/source_zero_upper_domain.json},
\path{code/part3_fourfold/source_zero_boundary_ledger.json},
and \path{code/part3_fourfold/source_zero_mixed_outer.json}.
They retain the full units and exact integer/normal
characters. The analytic extension and compact-face proofs
are given above; numerical identities do not replace those
maps. The chronological attempts and withdrawn shortcuts
remain in the working notes and attempt records. This
theorem proves the reduced branch only.

\end{enumerate}
\end{proof}




\begin{theorem}\label{night:P3K.81}
\StatementRecord{night:P3K.81}{proved}{S6\_06\_BLIND: GPT-6 Astra (OpenAI, ChatGPT Pro mode, effort Pro), recorded re-derivation in the checkpoint dated 18 November 2026}{\cref{night:P3F.26}, \cref{night:P3F.35}, \cref{night:P3F.43}}
For the fixed OS49 source and its component Cartier line, every
actual sphere member obtained by the stated finite root caps and smooth
translation caps has only order-one effective caps at smooth source values.
This holds with the full norm vectors, all finite characters and other
specialization groups retained, and for every allowed lattice-basis cusp
fan. At any effective smooth cap of order \(r>1\), the actual degree-two
specialization quotient is \((\mathbb Z/r)^3\). The unique global invariant
\(\omega=uw+3\psi\delta\) leaves a nonzero group in
\(H^1(\mathbb P^1,R^2\pi_*\mathbb Z)\), which survives in \(H^3(X;\mathbb Z)\).

For the explicit retained finite cap \((m,n,d)=(3,-1,0)\) and one smooth
cap of order two with full norm vector \(\gamma=(1,0,0,0)\), the actual
constructed compact manifold has
\[
 \pi_1(X)=1,\qquad H_2(X;\mathbb Z)=H_3(X;\mathbb Z)=(\mathbb Z/2)^2,
\]
\[
 H_0(X;\mathbb Z)=H_6(X;\mathbb Z)=\mathbb Z,
 \qquad H_1(X;\mathbb Z)=H_4(X;\mathbb Z)=H_5(X;\mathbb Z)=0.
\]
The proof supplies the full specialization, integral Mayer--Vietoris,
translation-prism and fundamental-group maps.

At the type-II value, every realized free good root has effective order
divisible by six. Its actual root source and corrected Cartier line give
\(\psi(\ell)=n\) for the full norm vector. Together with the retained
\(IV^*\) cap of order \(3k\), it has the displayed nonzero
\(\mathbb Z/3\) fundamental-group quotient and cannot be a sphere.
The reduced-II orders \(6l+1\) and their return classes are not assigned
by this statement.
\end{theorem}
\begin{proof}
\begin{enumerate}[label=\arabic*.]
\item The following full calculation has passed the recorded independent audit. Retain the actual OS49 source and
component line, the finite IV-star cap \((m,n,d)=(3,-1,0)\), and a fixed smooth
source value c.~Add the ordinary cap of order two whose full norm vector
is \(\gamma=(1,0,0,0)\) in the actual \((\gamma,u,w,\delta)\) homology marking. No
return class is assigned. The completed manifold is simply connected,
but its integral groups are
\[
 H_2(X)=H_3(X)=(\mathbb Z/2)^2,
 \quad H_1(X)=H_4(X)=H_5(X)=0,
 \quad H_0(X)=H_6(X)=\mathbb Z.                     \tag{OT.1}
\]
All maps below are the actual local specialization and integral
Cech/Mayer--Vietoris maps, with the original four fibre directions retained.

\par

\medskip\noindent\textbf{OT.1. The actual local cap and its pullback in every degree}

\par\nobreak

\item On a sufficiently small disc about c, the source and its degree-zero line
are smooth. Retain the actual source translation and its complete line
unit in the suspension B. On this disc choose their holomorphic Lie
logarithms. Half of this actual logarithm gives a commuting lift with
square B, including its line multiplier. On the root \(q=z^2\) the compact
torus action is
\[
                  (x,z)\longmapsto
                 (x+\Pi(q)\gamma/2,-z).                 \tag{OT.2}
\]
Its square adds the actual integral B period. Its nontrivial element is
free because the integral covector \(\psi\) evaluates \(\gamma\) to one. The
ordinary-to-root overlap is the actual plus-log translation with \(\gamma\),
so its filled meridian is \(2\alpha_S-\gamma\). All source and line units used
to form the half logarithm remain in this overlap.

\item The cap retracts onto the compact torus with period lattice
\(\Lambda'=\Lambda+\mathbb Z(\gamma/2)\). In the actual dual basis \((\psi,u,w,\delta)\), its
first-cohomology pullback is exactly
\[
                R_1=\operatorname{diag}(2,1,1,1),
                 \qquad R_q=\bigwedge^q R_1.              \tag{OT.3}
\]
This follows by evaluating an integral covector on \(\gamma\)/2; it is not a
rational rank computation. In lexicographic exterior bases \(R_q\) has
entry two precisely on basis forms containing \(\psi\), and entry one on
all the others. Thus the specialization quotient at this point is
\[
           Q_{S,q}=(\mathbb Z/2)^{\binom3{q-1}}
                 \quad(1\le q\le4).                    \tag{OT.4}
\]
The quotient map is the literal vector of coefficients on the basis
forms containing \(\psi\), reduced modulo two.

\item At the retained order-three cap, the norm is -\(\ell_0\) with
\(\ell_0=(1,-1,0,0)\). In actual root coordinates the map \(x\mapsto-x\) conjugates
the affine translation \(+\ell_0/3\) to \(-\ell_0/3\), keeps its linear action A,
and carries every full fibre period to its inverse. It extends on this
smooth root torus family. Hence the actual specialization lattices are
those computed in F26, with the sign \((-1)^q\) retained on degree q. Their
quotients are \(\mathbb Z/3\) in degrees one and four, and zero in degrees zero,
two and three. More explicitly the invariant degree-one basis is
\(\psi,\ h=-w+3\delta\), and the cap image is generated by \(3\psi,h\). In degree two
the full invariant lattice generated by
\[
                u_2=-\psi w+3\psi\delta,
                 \qquad v_2=\psi w+uw
\]
is its image. Degree three has the full invariant image, and degree four
has image \(3\nu\). These are integral cap maps from the actual finite
quotient, including its norm sign; no cohomology torsion is omitted.
The retained II and selected component cusp have the full integral
invariant specialization images of F26.

\par

\medskip\noindent\textbf{OT.2. The complete sheaf Mayer--Vietoris presentation}

\par\nobreak

\item Write \(L_q\) for the degree-q fibre local system and \(J_q=j_*L_q\) for its
invariant extension across the three original critical values. The
actual cohomological monodromies are
\[
 T_A=\begin{pmatrix}1&-3&-3&-1\\0&-2&-3&-1\\0&1&1&0\\0&0&0&1\end{pmatrix},
 \quad T_B=\begin{pmatrix}1&0&0&0\\0&1&1&0\\0&-1&0&0\\0&0&0&1\end{pmatrix},
 \quad T_C=\begin{pmatrix}1&0&3&1\\0&1&2&1\\0&0&1&0\\0&0&0&1\end{pmatrix}.
                                                               \tag{OT.5}
\]
The smooth new puncture has identity linear action. For \(E_i=\bigwedge^q T_i\),
the integral parabolic complex is
\[
 L_q\xrightarrow{d^0}\bigoplus_i(E_i-I)L_q
       \xrightarrow{d^1}L_q,
 \quad d^0v=((E_i-I)v)_i,
 \quad d^1(y_A,y_B,y_C)=y_A+E_Ay_B+E_AE_By_C.             \tag{OT.6}
\]
The exact Hermite bases, both maps and their full Smith forms are retained
in the certificate. Their middle cohomology is zero integrally, and the
last cokernel is free of rank one. The generators and projections are
\[
\begin{array}{c|c|c}
q&H^0(J_q)&H^2(J_q)/H_B\\\hline
0&1&[1]\\
1&\psi&[\delta]\\
2&\omega=uw+3\psi\delta&\eta,\quad[\psi\delta]=\eta,\ [uw]=3\eta\\
3&\phi=\psi uw&[uw\delta]\\
4&\nu=\psi uw\delta&[\nu].
\end{array}                                                   \tag{OT.7}
\]
Every other exterior-basis coefficient maps to zero in the displayed
coinvariant projection. All these statements use the full rank-four
matrices; the component line has not been replaced by a whole-fibre line.

\item Let \(S_q=R^q\pi_*\mathbb Z\) for the actual two-cap family. Every local restriction
above is injective, including the cap torsion-free cohomology, so the
literal restriction maps give
\[
 0\longrightarrow S_q\longrightarrow J_q
   \longrightarrow Q_{A,q}\oplus Q_{S,q}\longrightarrow0.       \tag{OT.8}
\]
The last map is supported at the two cap points. On a cover by the
ordinary complement and their discs, (OT.8) is precisely the quotient of
the actual integral Cech/Mayer--Vietoris restriction complex by the local
invariant blocks; its finite indices are (OT.3) and the order-three maps
above. It retains all base overlaps. In particular it is not a substitution
of invariant rational cohomology for the finite cap.

\item The global-invariant map to the skyscraper group has the following full
integer presentations, with the last column being that map:
\[
\begin{array}{c|c|c}
q&[\text{local relations}\mid\text{global image}]&H^1(S_q)\\\hline
1&\begin{pmatrix}3&0&1\\0&2&1\end{pmatrix}&0\\[3pt]
2&\begin{pmatrix}2&0&0&0\\0&2&0&0\\0&0&2&3\end{pmatrix}&(\mathbb Z/2)^2\\[3pt]
3&\begin{pmatrix}2&0&0&1\\0&2&0&0\\0&0&2&0\end{pmatrix}&(\mathbb Z/2)^2\\[3pt]
4&\begin{pmatrix}3&0&1\\0&2&1\end{pmatrix}&0.
\end{array}                                                    \tag{OT.9}
\]
For \(q=2\) the three smooth quotient coordinates are \(\psi u,\psi w,\psi\delta\), and the global image is \((0,0,3)\). For \(q=3\) they are \(\psi uw,\psi u\delta,\psi w\delta\), and the image is \((1,0,0)\). Thus these are the
actual complete maps, rather than only their Smith factors.

\item The kernels give the primitive zeroth-column basis
\[
                    1,\quad6\psi,\quad2\omega,
                          \quad2\phi,\quad6\nu.          \tag{OT.10}
\]
The two-cap sheaf still has the same free second column as (OT.7), because
the supported quotient in (OT.8) has no positive cohomology. There are no
other columns. Formula (OT.9) displays the nonzero middle column that a
calculation using only the Seifert determinant would miss.

\par

\medskip\noindent\textbf{OT.3. Actual translation prisms and the four integral differentials}

\par\nobreak

\item The untwisted source origin is the one constructed in F26: a meromorphic
frame of M restricted to O has its divisor at the selected cusp; the
actual integral cusp period is subtracted, so the origin extends there
and at II. Add the new smooth disc with its unchanged zero origin.
For the auxiliary linear root quotients, inversion about these origins
commutes with their linear monodromies and extends on the centrally
symmetric same toric fan. Its action on fibre degree q is \((-1)^q\).
Naturality therefore makes the intrinsic rational \(d_2\) zero. This is the
same actual origin and full line/cusp comparison as F26, with one
additional ordinary zero chart.

\item The actual affine finite-cap replacement is compared on its boundary by
the retained translation prism. Three lifted copies of the first prism
close to the integral translation -\(\ell_0\); two copies of the new prism
close to \(\gamma\). Their full operators on a fibre exterior form are
therefore
\[
             d_2=-\frac13\iota_{\ell_0}
                         +\frac12\iota_\gamma.           \tag{OT.11}
\]
The line and source entries of \(\ell_0\), including its -u entry, remain in
this formula. The source/target of each actual \(d_2\) from column zero to
column two is a free integral rank-one lattice, so its rational value
on the primitive bases (OT.7),(OT.10) determines its integral value.
The four evaluations are all exactly one:
\[
 d_2(6\psi)=H_B,\quad
 d_2(2\omega)=H_B[\delta],\quad
 d_2(2\phi)=H_B\eta,\quad
 d_2(6\nu)=H_B[uw\delta].                              \tag{OT.12}
\]
For example the second raw contraction is \(\delta+(2/3)w\); its w term is
zero under the actual coinvariant projection. The third is
\((uw-2\psi w)/3\), whose projection is \(\eta\). The fourth is
\(uw\delta-2\psi w\delta\), with the displayed target. These checks retain
all coordinates before projection.

\item The middle-column groups in (OT.9) have no possible incoming or outgoing
differential: the base has cohomological dimension two. All free terms
except degrees zero and six cancel by (OT.12). Hence
\[
 H^3(X;\mathbb Z)=H^4(X;\mathbb Z)=(\mathbb Z/2)^2,
 \qquad H^1=H^2=H^5=0.                                  \tag{OT.13}
\]
There are no extension ambiguities because each of these middle degrees
has only the single displayed surviving filtration group.

\par

\medskip\noindent\textbf{OT.4. Full fundamental group and integral homology}

\par\nobreak

\item Use the actual homological matrices \(H_i=(T_i^{-1})^t\). The II and cusp
fillings kill their base meridians and the images of \(H_B-I,H_C-I\). This
rank-three primitive relation module has quotient \(\mathbb Z\gamma\), as in F26.
The remaining base loops are \(\alpha_A,\alpha_S\). The exact finite-cap and
base relations reduce to
\[
 3\alpha_A+\gamma=0,\qquad 2\alpha_S-\gamma=0,
                  \qquad \alpha_A+\alpha_S=0.             \tag{OT.14}
\]
Their matrix on \((\gamma,\alpha_A,\alpha_S)\) is
\[
                  \begin{pmatrix}1&3&0\\-1&0&2\\0&1&1\end{pmatrix},
                     \qquad\det=1.                      \tag{OT.15}
\]
These are the full van Kampen relations, not merely a presentation of an
abelianization: the original torus generators commute, the killed II/cusp
meridians impose the displayed monodromy relations, and the base relation
leaves a single base generator. The surviving invariant \(\gamma\) commutes
with it. Every boundary maps surjectively onto its local filled group.
The full six-generator, eleven-column integer matrix and the actual
preimages of the \(H_A-I\) columns are retained in the certificate. Thus
\(\pi_1(X)=1\). Universal coefficients and integral Poincare duality applied
to (OT.13) give exactly (OT.1).

\item The chosen local cap is smooth, since its nontrivial affine deck is free;
the original root and cusp charts are unchanged. Compactness follows by
the same actual compact fundamental-domain gluing as F26. The calculation
therefore concerns the specified constructed manifold, not only a formal
Seifert presentation.

\par

\medskip\noindent\textbf{OT.5. The general primary obstruction for added smooth caps}

\par\nobreak

\item Retain any actual free OS49 finite cap of order \(m_0=3k\), and add finitely
many effective smooth logarithmic caps of orders \(m_i\) with their complete
norm vectors \(\ell_i\). The invariant covector \(\psi\) measures \(n_i=\psi(\ell_i)\).
Quotient the actual fundamental-group presentation by the normal fibre
subgroup killed by \(\psi\). The remaining marked relations are
\[
 m_i\alpha_i-n_i\gamma=0,\qquad
                  \sum_i\alpha_i+b_0\gamma=0.             \tag{OT.16}
\]
Here \(b_0\) retains all ordinary order-one clutches and the actual base
framing. It is not set to zero.

\item A sphere member must have pairwise coprime effective orders. Indeed if a
prime p divides two orders, all other invertible-order meridians can be
eliminated modulo p, while the indicated two have only their one sum
relation; their own norm equations constrain only \(\gamma\). This quotient
has a nonzero first homology class over \(\mathbb F_p\). When only \(m_i\) is divisible
by p, its coefficient \(n_i\) must be a unit modulo p, for otherwise \(\gamma\)
remains unconstrained after the same elimination. These are exact
consequences of (OT.16), independent of the other entries of \(\ell_i\).

\item Fix a prime p dividing the order of an added smooth cap in a sphere datum,
and use \(R=\mathbb Z_{(p)}\). Then p does not divide \(m_0\), in particular p is not three,
and no other effective order is divisible by p.~All the other finite
cap pullbacks become isomorphisms onto their invariant R lattices by
the actual finite covering and transfer; their degrees are units in R.
The II and cusp maps remain the integral invariant maps. The unchanged
parabolic complex (OT.6) has \(H^1=0\) over R and invariant degree-two
generator \(\omega\).

\item Write the full chosen norm vector as \(\ell=(n,a,b,d)\), with n a unit of R.
An explicit basis of its annihilator in the dual lattice is
\[
 u-(a/n)\psi,\qquad w-(b/n)\psi,
                         \qquad\delta-(d/n)\psi.          \tag{OT.17}
\]
The local degree-two quotient is \((R/m_iR)^3\), by the same actual torus
covering as (OT.3), now with \(\Lambda+R(\ell/m_i)\). Its natural quotient map
is contraction with \(\ell\) modulo \(m_i\). On the retained global invariant,
\[
 \iota_\ell\omega=-b u+a w+3n\delta-3d\psi,
 \qquad\text{coordinates in (OT.17): }(-b,a,3n).           \tag{OT.18}
\]
The third coefficient is a unit in R. The full local/global MV
presentation is therefore
\[
            [\ m_i I_3\mid(-b,a,3n)^t\ ],
 \qquad H^1(R^2\pi_*R)=(R/m_iR)^2\ne0.                 \tag{OT.19}
\]
The source, line and \(\delta\) coefficients are retained in (OT.17)--(OT.18);
this is a calculation in the localized lattice, not a change of the
actual cap or a deletion of its \(\delta\) character.

\item This middle-column group survives to the integral localized cohomology
filtration for the same dimensional reason as in OT.3. Thus \(H^3(X;R)\)
has a nonzero filtration quotient, contradicting the sphere cohomology.
Consequently every added effective smooth cap of an OS49 sphere member
has order one. The actual periodic flop comparisons of F35 are the
identity near every cap and preserve integral homology and \(\pi_1\), so the
obstruction applies to every allowed lattice-basis fan without using an
unproved return formula.

\item The executed certificate is \path{code/part3/os49_order_two_cap.py}, with all
exterior matrices, parabolic maps, local quotient presentations, primitive
indices, prism contractions and the full fundamental relation matrix.
Its first run passes 63 identities. The preceding localized argument is
retained in full. The next section gives a stronger direct obstruction
that does not use the pairwise-coprime reduction or transfer at the other
caps.

\par

\medskip\noindent\textbf{OT.6. Every higher effective smooth cap is excluded directly}

\par\nobreak

\item Retain the fixed OS49 source, its actual component line and linear
period tuple (OT.5). Take any realized free finite caps at the other
sites and any finite collection of actual smooth logarithmic caps. The
word effective records the order of the translation on the central
compact torus. On a smooth root disc the full actual affine map has the
form
\[
 (x,z)\longmapsto(x+\Pi(z^r)\ell/r,\zeta_r z),
 \qquad \ell\in\Lambda,\qquad
       \gcd(r,\ell_1,\ell_2,\ell_3,\ell_4)=1.          \tag{OT.20}
\]
Here is the full map to these coordinates. In the original smooth root
atlas write the actual lifted translation as a(z), with its source and
line entries, and choose its holomorphic universal-cover lift on the
disc. The literal norm gives
\(\sum_{j=0}^{r-1}a(\zeta_r^jz)=\Pi(z^r)\ell\). Put
\[
 f(z)=a(z)-\Pi(z^r)\ell/r,\qquad
 g(z)=\frac1r\sum_{j=0}^{r-1}j f(\zeta_r^j z).
                                                               \tag{OT.20a}
\]
Then \(g(\zeta_rz)-g(z)=f(z)\), so \((x,z)\mapsto(x-g(z),z)\)
intertwines the original deck with (OT.20). This map commutes with the
full period lattice and the physical circle, extends on the root disc,
and is carried through the actual ordinary-to-cap overlap. Its constant
finite character is part of \(\ell\) and is not removed by this averaging.
The gcd condition says precisely that
the translation has order r, so every nonidentity deck element acts
freely at the central fibre. A chosen root with a smaller translation
order has a subgroup acting only on the root coordinate. Its quotient
is the literal map \(z\mapsto z^g\), where
\(g=\gcd(r,\ell_1,\ell_2,\ell_3,\ell_4)\); the residual displayed map has order
\(r/g\) and vector \(\ell/g\). Thus (OT.20) specifies the effective cap without
changing an effective nontrivial translation into an order-one cap.

\item The actual cap retracts to
\(T'=\Lambda_{\mathbb R}/(\Lambda+\mathbb Z\ell/r)\), and the specialization from an ordinary
fibre is the identity-linear torus covering
\[
 c:\Lambda_{\mathbb R}/\Lambda\longrightarrow
       \Lambda_{\mathbb R}/(\Lambda+\mathbb Z\ell/r).
                                                               \tag{OT.21}
\]
In the real flat marking \(x=\Pi(q)y\), the radial retraction
\((y,z)\mapsto(y,tz)\), \(0\le t\le1\), commutes with the deck
\((y,z)\mapsto(y+\ell/r,\zeta_rz)\). It proves both the asserted cap retraction
and the actual specialization covering (OT.21). Its degree is r and
its deck group is generated by \(\ell/r\). In particular
the first-cohomology pullback is the actual sublattice
\[
 c^*H^1(T';\mathbb Z)
 =\{v\in\Lambda^*:v(\ell)\equiv0\pmod r\}
                 \subset\Lambda^*.                         \tag{OT.22}
\]
The evaluation map onto \(\mathbb Z/r\) is surjective by the gcd condition.
Consequently (OT.22) has Smith factors \((1,1,1,r)\). Exterior powers give
the literal degree-two specialization quotient
\[
 Q_{c,2}=\bigwedge^2\Lambda^*/
                 \bigwedge^2c^*H^1(T';\mathbb Z)
                      \simeq(\mathbb Z/r)^3.              \tag{OT.23}
\]
For an entirely integral coordinate construction, let H be the column
Hermite basis of \([rI_4\mid\ell]\). Then \(\Lambda'=r^{-1}H\mathbb Z^4\), the actual
dual pullback matrix is \(rH^{-t}\), and its exterior square is the map
in (OT.23). This retains every entry of \(\ell\) and furnishes the exact
lattice morphism underlying the Smith description.

\item Let \(R=R^2\pi_*\mathbb Z\) and \(J=j_*L_2\) on the fixed base. On the ordinary locus
the specialization map R to J is the identity. Let K be its kernel,
S its image, and Q its cokernel. Their actual stalk maps give
\[
 0\longrightarrow K\longrightarrow R\longrightarrow S
       \longrightarrow0,\qquad
 0\longrightarrow S\longrightarrow J\longrightarrow Q
       \longrightarrow0.                                \tag{OT.24}
\]
K and Q are supported on the finite set of special values. This retains
all kernel and torsion groups of every other cap or cusp completion;
they are not assumed to vanish. Such a finite-support sheaf has no
positive sheaf cohomology. Hence
\[
 H^1(R)\simeq H^1(S),\qquad
 \operatorname{coker}\{H^0(J)\longrightarrow H^0(Q)\}
                         \hookrightarrow H^1(S).         \tag{OT.25}
\]
The fixed tuple gives \(H^0(J)=\mathbb Z\omega\), as checked integrally in OT.2.
Projection of \(H^0(Q)\) to the actual smooth-cap summand (OT.23) induces
a surjection from the cokernel in (OT.25) onto
\[
                  Q_{c,2}/\langle\operatorname{sp}_c(\omega)\rangle.
                                                               \tag{OT.26}
\]
For every prime p dividing r, the quotient in (OT.26), tensored with
\(\mathbb F_p\), has dimension at least two: (OT.23) has three generators and the
displayed image supplies only one relation. It is therefore nonzero.
Equations (OT.25)--(OT.26) prove \(H^1(R^2\pi_*\mathbb Z)\) nonzero without any
condition on specialization at the other caps.

\item The Leray term \(E_2^{1,2}\) has no possible incoming or outgoing
differential. Incoming indices have negative base degree; outgoing
indices exceed the cohomological dimension two of the real base.
Thus it survives as a nonzero associated-graded part of \(H^3(X;\mathbb Z)\).
This excludes sphere cohomology for every effective smooth cap of
order r greater than one.

\item This also proves the asserted all-fan scope directly. On the dense
ordinary locus the actual toric changes are the identity with the same
full period tuple, and (OT.21) is unchanged near every smooth cap.
Changing the selected-cusp fan changes only a special stalk already
retained in K and Q. For the full homology calculation of OT.1, the
actual common blowups and inverse blowdowns of F35 give the integral
homology and fundamental-group transport, relative to all cap charts;
both exceptional rulings pair primitively with the exceptional divisor.
No return formula or identification of the two rulings is used here.

\par

\medskip\noindent\textbf{OT.7. The actual good-II root and its surviving norm}

\par\nobreak

\item At the source value \(t=1\) put \(q=t-1\). The exact OS49 equation is
\[
 y^2=x^3-3(1+q)^3q x+(1+q)^4q(1+2q),
                       \qquad G=((1+q)^2,(1+q)^2).       \tag{OT.27}
\]
For a good root \(m_0=6l\), l at least one, put \(q=z^{6l}\),
\(x=z^{2l}X\), \(y=z^{3l}Y\). Its minimal equation and marked section are
\[
 \begin{split}
 Y^2&=X^3-3(1+z^{6l})^3z^{2l}X
                   +(1+z^{6l})^4(1+2z^{6l}),\\
 X(G')&=z^{-2l}(1+z^{6l})^2,\qquad
 Y(G')=z^{-3l}(1+z^{6l})^2,\qquad
                  -X(G')/Y(G')=-z^l .
 \end{split}                                                \tag{OT.28}
\]
Thus \(G'\text{ and }Q'=3G'\) both tend to O on the central curve
\(Y^2=X^3+1\). The actual deck is
\[
 D:(z,X,Y)\longmapsto
       (\zeta_{m_0}z,\zeta_{m_0}^{-2l}X,
                              \zeta_{m_0}^{-3l}Y).        \tag{OT.29}
\]
Its central order is six. The small formal-group inverse of
\(m_0\) gives the D-invariant section R with \(m_0R=Q'\) and \(R(0)=O\).

\item The local six-point blowup/quotient/contraction diagram of F43 applies
to the two marked sections G' and O individually. Both meet the same
order-six exceptional curve at finite coordinates, whereas all three
contractions occur at its infinity point. Their exceptional coefficients
cancel in their difference. The literal Cartier identity is therefore
\[
 \beta^*\mathcal O(G'-O')
               =\pi^*\alpha^*\mathcal O(G-O).             \tag{OT.30}
\]
This is the local diagram of F43 with \(d=1\); its proof uses the actual two
section divisors and requires no change to the component-supported
line at the distant selected \(I_2\) fibre. Translation Q' tends to O and
fixes all six centres, so (OT.30) carries the actual source/line
suspension and its deck action through that diagram.

\item Lift \(t_{-R}\) on this actual root line to F. In the retained finite
character gauge, \(DFD^{-1}=\chi F\), \(\chi^{m_0}=1\), and
\(B_0=\beta(q)F^{m_0}\). The invariant norm-root correction from F43 gives
the commuting full lift \(A_0\) with \(A_0^{m_0}=B_0^{n_0}\); it retains the
factor \(\chi^{n_0m_0(m_0-1)/2}\) and the prescribed \(\delta\) character.
The coefficient covector \(\psi\) kills the elliptic and physical-\(\delta\)
periods and takes value one on the full suspension period B. Lifting
the literal root relation \(A^{m_0}=B^{n_0}\) to the actual lattice thus
gives
\[
 \psi(\ell_0)=n_0,\qquad
          m_0\alpha_0-\ell_0=0,
          \quad m_0\alpha_0-n_0\gamma=0
                       \text{ after the psi quotient}. \tag{OT.31}
\]
All other source and \(\delta\) entries of \(\ell_0\) remain in the first
relation. This establishes the actual norm before using the base
determinant.

\item Every such effective free good-II cap has order divisible by six, and
the retained IV-star cap has order 3k. There is an explicit surjection
of the full fundamental-group presentation onto \(\mathbb Z/3\): send the
IV-star base meridian to 1, the good-II meridian to -1, every other
base meridian to zero, and every fibre generator to zero. Both
multiple-cap relations vanish because their orders are divisible by
three; the base sum and all affine conjugation relations also vanish.
The local free quotient maps retain exactly these root meridians, so
this is a homomorphism on the filled manifold. Hence a good-II root
replacement cannot be a sphere member. The reduced-II orders
\(6l+1\) are not assigned a topology or return by this calculation.

\end{enumerate}
\end{proof}




\begin{theorem}\label{night:P3K.82}
\StatementRecord{night:P3K.82}{proved}{S6\_06\_BLIND: GPT-6 Astra (OpenAI, ChatGPT Pro mode, effort Pro), recorded re-derivation in the checkpoint dated 18 November 2026}{\cref{night:P3F.35}, \cref{night:P3K.31}, \cref{night:P3K.60}, \cref{night:P3K.68}, \cref{night:P3K.71}}
Retain the actual marked source
\[
 Y^2=X^3+(1-2v)X^2+TX+v^2,\quad G=(0,v),\quad
 q=T-(v^2-2v),\quad\lambda=v^{N}\lambda_0,
\]
with \((a,b,m)=(1,2,4p)\) or \((2,1,4p-1)\), \(p\ge1\),
\(N=km\), both cap norm signs \(\nu=\pm1\), every finite character
and the full annular cochain. The case \(k=0\) is the actual fixed-scalar
loop separated by RT.12 on its compact source annulus. For each
\(j\in\mathbb Z\) use the physical vertical-edge fan
\(T_j=T(e_1+je_2,e_2)\), and write \(Z_j\) for its actual positive
source-zero return in the physical K68 presentation \(\Theta_7=\mathbb Z/28\).
Then
\[
 Z_j-Z_0=2b\binom{j+1}{3}+(3-b)\binom{j+1}{2}
                            -N\binom j2\pmod {28}. \tag{ZV.1}
\]
The binomials are the integer-valued falling-factorial polynomials at
all integers. In the reduced branch \(a=2,b=1,m=4p-1\), P3K.80 gives
the fully evaluated formula
\[
 Z_j=k(3p-1)+2\binom{j+2}{3}
                  -k(4p-1)\binom j2\pmod {28}.       \tag{ZV.2}
\]
The proof constructs the actual single-loop \(C_1\) flop and derives
its free right integer from the literal doubled product. The good
branch retains its independently determined \(Z_0\) in the relative
formula; no absolute value is assigned to it here.
\end{theorem}
\begin{proof}
\begin{enumerate}[label=\arabic*.]
\item \par

\medskip\noindent\textbf{ZV.1. Actual vertical characters, full periods and their seam}

\par\nopagebreak[4]

\item Write the full periods on the protected selected-cusp neighbourhood as

\[
 h(Z,V,q)=(\kappa qZ,\alpha V,q),\qquad
 B(Z,V,q)=(gZ,q\eta V,q).
\]

\item Every coefficient is the full nowhere-zero function of q and the cut
source parameter. The exact SZ.44 calculation gives the single-loop
winding matrix

\[
 W=\begin{pmatrix}1&b\\-a&N-2\end{pmatrix}.
\]

\item Make the literal character map \((Z,V,q)\mapsto(Z_j,V_j,q)=(Z,Z^{-j}V,q)\),
and take the actual period generators \(h_j=hB^j\) and \(B_j=B\). Their
complete coefficients are

\[
 \kappa_j=\kappa g^j,\quad g_j=g,\quad
 \alpha_j=\alpha\eta^j\kappa^{-j}g^{-j^2},\quad
 \eta_j=\eta g^{-j}.                                  \tag{ZV.3}
\]

\item The character matrix and its inverse are integral. On every absolute
matched cone they give inverse maps of its character semigroup with
the stated full units. They therefore identify \(T_j\) with its standard
vertical marking as an actual holomorphic toric neighbourhood. The
inverse maps attach its unchanged common exterior, source line and
finite root charts. Its full multiplier winding matrix is exactly
\(P_j^{-1}WP_j\). No q-height or annular cochain is discarded.

\item Put \(L_0=\nu k\). The actual forward ordinary closing is \((Z,V)\mapsto(Z,q^{L_0}V)\). After the actual \(B^{-L_0}\) reindex, its bounded chart arrow is
\((g_j^{-L_0}Z_j,\eta_j^{-L_0}V_j)\). On the cut parameter interval take
logarithms \(G_j=\log g_j,\ H_j=\log\eta_j\), and a cutoff chi that is zero
on the initial collar and one on the terminal collar. The coordinate
map has multipliers \(A=e^{L_0\chi G_j},\ D=e^{L_0\chi H_j}\).
It has the explicit inverse with negative exponents and cancels the
bounded closing on the entire terminal collar with all derivatives.
All its multipliers depend only on the unchanged q and v, so it
commutes with both full periods and the physical circle. Its vertex
maps are exactly

\[
 (a_0,b_0,c_0)\mapsto((AD)^{-1}a_0,Ab_0,Dc_0),
 \quad(u_0,v_0,w_0)\mapsto(ADu_0,D^{-1}v_0,A^{-1}w_0).
                                                        \tag{ZV.4}
\]

\item The products in both charts remain q. The three period rows are
DZ.8 with the full coefficients (ZV.3). Thus these are closed charts
of the actual family. The doubled seam is its literal pullback:
on lap \(\ell=0,1\) use \(C_\tau E_\tau^{-\ell}\). The two formulas agree on the
middle collar and cancel \(E^2\) at the last collar, just as DZ.8a.
The full cap square and every finite character remain the same.

\par

\medskip\noindent\textbf{ZV.2. The complete doubled product and its integral nonfixed entries}

\par\nopagebreak[4]

\item On the literal double parameter cover, set

\[
 \mathcal B_j=2P_j^{-1}WP_j,
 \quad t=N-1,
 \quad l=(\mathcal B_{22}-\mathcal B_{21},
          \mathcal B_{11}-\mathcal B_{12},\mathcal B_{22}).
                                                        \tag{ZV.5}
\]

\item In particular

\[
 l_0=2\{N+a-2+(3-N-b)j+bj^2\},\quad
 l_1=2\{1+b(j-1)\},\quad l_2=2(N-2-bj).
                                                        \tag{ZV.6}
\]

\item Their normal transition matrices are the actual DZ.10 matrices
\(\operatorname{diag}(\eta_j/\alpha_j,\kappa_j\alpha_j)\),
\(\operatorname{diag}(\kappa_j/g_j,g_j\eta_j)\), and \(\operatorname{diag}(\eta_j,\kappa_j)\).
Their determinant winding is 2t. Each based coefficient equals its
angular factor times the exponential of its actual periodic logarithm.
Use precisely the based quaternion maps \(H_d\) of DZ.12, with integral
d, and the normal map \(H_t Q H_{l_i-t}\), multiplied by those full
logarithmic units with a protected radial cutoff. Its second radial
endpoint is exactly the original diagonal matrix, not its valuation.
The same normal-rescaling isotopy and cutoff velocity as DZ.14
restore all higher terms and keep both vertex collars fixed. This
constructs an ordinary product of the entire doubled three-curve
plumbing and its rounded boundary.

\item The physical weights are \(w=(0,1,-1),\ d=(-2,-1,-1)\). They remain
these actual weights because \(P_j^{-1}e_2=e_2\). The same explicit Hopf
trace and full stable derivative of DZ.16--18 give

\[
 I_i=t w_i+d_i(l_i-t),\qquad
 W_i=t w_i-d_i(l_i-t).                              \tag{ZV.7}
\]

\item The parameter-independent vertex caps, generic normal sections and
ambient-frame homotopy are those in DZ.3, transported by (ZV.3)--(ZV.4).
Consequently the two literal action/derivative prism identities are

\[
 V(w_i e_i-w_0e_0)=I_i e_i-I_0e_0,\qquad
 w_i\delta_i-w_0\delta_0=W_i-W_0.                  \tag{ZV.8}
\]

\item For the nonfixed \(C_1\) curve, \(w_0=0,\ w_1=1\). These determine its complete
selected entries without any fixed-star entry:

\[
 V_{11}=I_1=2t-l_1,\qquad
 \delta_1=W_1-W_0=l_1-2l_0+2t.                     \tag{ZV.9}
\]

\par

\medskip\noindent\textbf{ZV.3. The actual single native \(C_1\) product and oriented flop}

\par\nopagebreak[4]

\item The native chart is explicit. In the first vertex use
\((z,n_0,n_1)=(c_0,a_0,b_0)\); in the second use

\[
 (z',n'_0,n'_1)=(\eta_jv_0,\ g_jw_0,\ u_0/(g_j\eta_j)).
                                                        \tag{ZV.10}
\]

\item The full \(C_1\) row of DZ.8 gives \(z'=z^{-1},\ n'=zn\). Both products
are q. These functions are nonzero and periodic on the closed
single-loop chart; their integer winding has not been replaced by a
logarithm. Thus (ZV.10) is a literal native product of this one star.
The normal-order interchange needed for its \(Q_1\) flop is the constant
unitary path of K71.6. It conjugates both Spin factors by fixed
orientation-preserving lifts and preserves their third degrees.

\item On the doubled cover the native-to-ordinary boundary coefficient map
is \((UH_tU^{-1})q_HH_{l_1-t}\), with the same full unit homotopy as above.
The left degree is 2t; the right factor has third degree zero because
it has only the radial and parameter variables. The actual inverse
cylinder extension has endpoint degrees \((-2t,0)\) and extends over the
normal D4. Therefore the native right free integer equals the ordinary
right free integer. The ambient-to-star frame map of K31, carried by
(ZV.10), gives that doubled integer as \((\delta_1-V_{11})/2\).

\item The doubled native product is the literal pullback of (ZV.10) with
the doubled seam in ZV.1. Choose a complete global connection having
this boundary product and pull it back to the two-sheeted parameter
cover. The relative return is the square of the actual single return.
The right free coordinate is a homomorphism to Z, so the single
native right integer is exactly

\[
 r_j=\frac{\delta_1-V_{11}}4
       =\frac{l_1-l_0}{2}
       =-N+(N+2b-3)j-bj^2.                         \tag{ZV.11}
\]

\item This uses division in a proved integral free coordinate of an actual
square. Every \(l_i\) is even, and the displayed polynomial is integral.
No \(\Theta_7\) class or arbitrary integer lift has been divided.

\item The full seam extends to the common Cox chart of this \(C_1\) star.
With \(z=a_1/a_0,\ n_i=a_0c_i\), the actual lift of (ZV.4) is

\[
 (a_0,a_1,c_0,c_1)\mapsto
       (a_0,Da_1,(AD)^{-1}c_0,Ac_1).                \tag{ZV.12}
\]

\item It preserves \(q=a_0a_1c_0c_1\) and both the auxiliary and physical circle
actions. Since q is unchanged, its inverse uses the reciprocal full
functions at that same q. It therefore extends through both small
resolutions. The other-ruling charts are the actual maps
\((n_0/n_1,zn_1,n_1),\ (n_1/n_0,zn_0,n_0)\), with the stated constant
normal interchange. They retain the common rank-one matrix and all
full period multipliers. Their cone bases are \(P_jQ_1=P_{j-1}\),
\(Q_1=\left(\begin{smallmatrix}1&0\\-1&1\end{smallmatrix}\right)\). On the common link the Hopf arrow is
\((h,q_H)\mapsto(q_H^{-1}hq_H,q_H)\), with K68 physical orientation.
The exact Cox conjugation consequently changes the central sphere
class by the actual right integer \(r_j\). Thus

\[
                         Z_{j-1}-Z_j=r_j.          \tag{ZV.13}
\]

\item All maps are the identity on the matched common exterior after their
inverse character identification. Hence the original source, finite
cap, signed norm and full annular closing are still the same data.
The common blowup and both primitive exceptional rulings preserve
the proved integral sphere topology, as in F35.

\par

\medskip\noindent\textbf{ZV.4. Summation and evaluation}

\par\nopagebreak[4]

\item The integer polynomial

\[
 F(j)=2b\binom{j+1}{3}+(3-b)\binom{j+1}{2}
                                          -N\binom j2
\]

satisfies \(F(0)=0,\ F(j-1)-F(j)=r_j\) as exact identities for every
integer j. The inverse of the actual flop supplies the negative steps
as well. Induction in either direction proves (ZV.1). For \(a=2,\ b=1\),
Pascal's exact polynomial identity gives \(F(j)=2\binom{j+2}{3}-N\binom j2\).
Substitution of the audited reduced attachment and RT.12 gives (ZV.2).
In particular j=1 is the already marked opposite-diagonal fan and
has \(Z_1-Z_0=a\) in both branches, agreeing with the independent K60
single native diagonal map. The new recurrence extends that
calculation to every integer vertical fan without assigning any
unknown fixed-star free coordinate from a finite-group residue.

\item The independent exact program
\path{code/part3_root/source_zero_vertical_fans.py} executes the full
matrix conjugation, primitive prism entries, native-chart products,
Cox seam, recurrence and integer-valued formulas. Its bounded examples
are checks of the symbolic proof, not its arbitrary-integer argument.

\end{enumerate}
\end{proof}




\begin{theorem}\label{night:P3K.83}
\StatementRecord{night:P3K.83}{proved}{S6\_06\_BLIND: GPT-6 Astra (OpenAI, ChatGPT Pro mode, effort Pro), recorded re-derivation in the checkpoint dated 18 November 2026}{\cref{night:P3F.22}, \cref{night:P3K.1}, \cref{night:P3K.60}, \cref{night:P3K.62}, \cref{night:P3K.68}}
Retain the actual source-zero family with
\((a,b)=(1,2)\), good-root order \(m=4p\), either realized norm
sign and every realized finite \(\delta\) character, with all source,
Cartier, root and annular maps prescribed above. In the selected
Tate fan \(T_0\), the positive fixed-scalar source loop has physical
integer index and oriented return class
\[
                         A_M=0,\qquad [\Sigma]=0\in\Theta_7.
\]
Its literal reflected fan \(JT_0\) has oriented return class
\(1\in\Theta_7\). The actual coupled loops
\(\lambda(v)=\lambda_0v^{km}\) have these same two classes,
\(0\) and \(1\), respectively. The reflected conclusion is an
oriented class; no additional integer characteristic pair is
assigned. The proof retains two chronological primitive maps,
their intermediate cap object, both charge-one faces, the entire
fixed trace and the actual physical Hopf completion. The already
proved vertical-fan recurrence is P3K.82.
\end{theorem}
\begin{proof}
\begin{enumerate}[label=\arabic*.]
\item \par

\medskip\noindent\textbf{GC.1. The actual positive quotient and a proper weighted cut}

\par\nobreak

\item Use the full bilinear node coordinates of EA.42--46. The primary
associated resolution has coordinates (\(S_0\),\(S_1\),\(T_1\),z,V), with auxiliary
weights (1,1,-1,0,-1) and physical weights (0,0,0,1,-1). Its positive
auxiliary quotient is literally

\[
 W_1=\operatorname{Tot}\bigl(O(-1)_{T_1}\oplus O(-1)_V
                  \longrightarrow P^1_{[S_0:S_1]}\bigr)\times C_z.
 \tag{GC.1}
\]

\item The transition functions are the actual inverse powers of the two
homogeneous source coordinates. The source functions are
\(q_1\)=\(S_0\) z V and \(q_2\)=\(S_1\) \(T_1\). The original parameter v is the analytic
inverse of the base coordinate \(q_1\)-\(q_2\), with the unit retained as in
EA.17. This is the exact morphism from the original source, not
its homogeneous approximation.

\item Let H=\textbar{}\(S_0\)\textbar{} squared plus \textbar{}\(S_1\)\textbar{} squared, and define invariant functions

\[
 b=H|T_1|^2,\qquad c=H|V|^2,\qquad d=|z|^2,
 \qquad \rho^4=b+c^2+d^2,\qquad R^2=b+cd.          \tag{GC.2}
\]

\item The function \(\rho\) is proper on the vector-bundle neighbourhood of
the compact zero section. Every positive \(\rho\) level is smooth:
the derivative in \(T_1\) is nonzero when \(T_1\) is nonzero; otherwise a
nonzero V or z gives a nonzero derivative of its fourth power.
This is the weighted radial function with weights 2,1,1 on
(\(T_1\),V,z), respectively. R is exactly the source product radius
\textbar\textbar S\textbar\textbar{} \textbar\textbar T\textbar\textbar{} for T=(zV,\(T_1\)).

\item The physical fixed locus is exactly z=V=0. To see this without
discarding auxiliary stabilizers, the nonzero positive auxiliary
coordinates S force the auxiliary compensating phase to be one.
The remaining physical phase therefore fixes a point precisely
when both z and V vanish. This fixed locus is O(-1) over \(\mathbb P^1\), the
same exceptional trace disk as in RT.1. There is no isolated
positive-chamber fixed point in (GC.1). On z nonzero the function
z is an actual equivariant nonzero coordinate of physical weight
one; it supplies the physical line trivialization there.

\item Choose radii satisfying \(\epsilon\) positive and
\(r_1^4>4\epsilon^2\), with \(r_1\)\textless{}\(\rho_0\). Retain the actual compact
domain

\[
 D_1=\{\rho\le r_1\}\ \cup
 \{r_1\le\rho\le\rho_0,\ R\ge\epsilon\}\ \cup
 \{r_1\le\rho\le\rho_0,\ d\ge c\}.              \tag{GC.3}
\]

\item Every union and intersection is in the same quotient (GC.1), by
the identity source, line and tangent map. At d=c one has
R squared=b+c squared and \(\rho\) fourth=b+2c squared, so
R squared is at least \(\rho\) fourth divided by two. Hence the
artificial d=c seam above \(r_1\) lies strictly in R\textgreater{}\(\epsilon\) and
supplies no boundary face. The entire zero section is in the
first piece. On the fixed trace c=d=0, R=\(\rho\) squared, so the
whole trace remains in (GC.3), including its exceptional disk.
The upper z arm is filled inside this same smooth quotient.

\item The removed horn is exactly \(r_1\)\textless{}\(\rho\)\textless=\(\rho_0\), c\textgreater d and R\textless{}\(\epsilon\).
It has V nonzero throughout. The two functions \(\rho\) and R have
regular transverse level faces there. Here is the check in the
actual complex fibre coordinates. At a point with \(T_1\) nonzero,
proportionality of their gradients would force the multiplier
between d(R squared) and d(\(\rho\) fourth) to be one. The nonzero
V coordinate would then give d=2c, contradicting c\textgreater d.~At a
point with \(T_1\) zero and R=\(\epsilon\) positive, both z and V are
nonzero. The ordered two-gradient determinant is
2(c squared-d squared)\textgreater0. If z is zero and \(T_1\) is nonzero,
the corresponding determinant is 2c\textgreater0. In the fully positive
case the \(T_1\) equation fixes the proportionality multiplier
to one, and the two others would force c=d=0. This handles zero
fibre coordinates without treating their squared norms as
smooth coordinates there. Invariant corner rounding retains
the specified outward normals.

\par

\medskip\noindent\textbf{GC.2. The complete charge-one lower shape}

\par\nobreak

\item In the actual V=1 chart there is no residual finite quotient.
The physical weights on (\(S_0\),\(S_1\),\(T_1\),z) are (-1,-1,+1,+1). Put

\[
 \sigma^2=\|S\|^2+|T_1|^2+|z|^2,\quad
 x=\|S\|^2/\sigma^2,\quad \zeta=|z|^2/\sigma^2.
\]

\item Direct substitution, with no radial coordinate map, gives

\[
 \rho^4=\sigma^4\bigl(x(1-\zeta)+\zeta^2\bigr),
 \qquad R^2=\sigma^4 x(1-x).                       \tag{GC.4}
\]

\item On the lower horn c\textgreater d says x\textgreater{}\(\zeta\). Its angular denominator
f0 squared=x(1-\(\zeta\))+\(\zeta\) squared is at most x. Consequently

\[
 1-x\le {R^2\over\rho^4}
       <{\epsilon^2\over r_1^4}<\frac14,
 \qquad 0\le\zeta\le1-x.                          \tag{GC.5}
\]

\item Thus x\textgreater3/4 and \(\zeta\)\textless1/4 on the closed horn. The square root f0 is therefore a positive
smooth angular function on a neighbourhood of that entire horn.
Choose an invariant cutoff equal to one there and supported in
x\textgreater1/2, and extend f=1+\(\chi\)(f0-1) to the whole charged sphere.
On that support \(\zeta\)\textless=1-x=x squared\textgreater1/4.
Together with the upper bound below this gives 1/2\textless f0\textless=1.
The cutoff support avoids every zero of f0; f is positive.
On the horn the exact defining-function equality is

\[
                         \rho^2=\sigma^2 f.         \tag{GC.6}
\]

\item Thus the complete outer charged shape \(\sigma\) squared times f
equal \(\rho_0\) squared agrees literally with the old radial face
on its lower rim. Since x+\(\zeta\)\textless=1, f0 squared is at most
(1-\(\zeta\)) squared plus \(\zeta\) squared, hence at most one. Our
cutoff therefore has f\textless=1. Take the round inner link with
\(\sigma\) squared equal to \(r_1\) squared divided by two. This lies
strictly below the entire horn, where \(\sigma\) squared\textgreater=\(\rho\) squared.
The interval region between that inner sphere and the full
outer shape is a compact free-circle annulus, with all source
functions carried by ambient inclusion. No radial map is
asserted to preserve \(q_1\) or \(q_2\).

\item Removing the horn from that charged annulus and gluing the
opposite resulting piece to D1 along its floor and side gives
the same explicit boundary-chain identity as EA.50--52. The
face maps here are the a=1 versions of EA.26 and EA.39, with
their actual node coefficient and scalar normalization. The
remaining complete charged boundary has physical line O(-1)
over \(\mathbb{CP}^3\): conjugating the two negative complex coordinates
preserves real eight-dimensional orientation, and its actual
coefficient map is {[}u,w{]} to w u in the tautological line.
Its physical pair is (-1,-4), so its characteristic index is
zero. This is the specified lower boundary calculation; it
does not count how the two primitive source steps occur in
the whole B strip.

\item The upper z region does not produce a charge-zero projective
substitution. It remains the actual region (GC.1), with the
equivariant nonzero function z on its free upper arm. Its line
and both boundary maps are retained in the ordinary source
gluing. The source functions satisfy
\textbar{}\(q_1\)\textbar+\textbar{}\(q_2\)\textbar\textless=R, so the same quantitative separation from the
outer source and parameter boundaries used in EA.19 applies.

\par

\medskip\noindent\textbf{GC.3. Two primitive cap steps with the full character and divisor}

\par\nobreak

\item Write m=4p. The actual primitive and suspension are

\[
 U_1:(p,w)\longmapsto(p-G,w/X(p-G)),\qquad
 B=-2\lambda q_1 U_1^2.                            \tag{GC.7}
\]

\item On the root \(q_\infty=z_\infty^m\) the actual normalization from
EA.7 is

\[
 U_{\rm cap}=z_\infty^{-m/2}U_1,\quad
 F_0U_1=z_\infty^{m/2}U_{\rm cap}F_0,\quad
 q_1=z_\infty^{-m}h,\quad
 F_0B=-2\lambda h\,U_{\rm cap}^2F_0.              \tag{GC.8}
\]

\item All source and line factors of \(F_0\) are retained. The root map
Ucap is the literal translation with target coefficient
1/X'(p-G'), not an uncorrected tensor frame. The canonical
deck sends X' to -X', and the exact finite relation is
Ucap \(A_0\) Ucap inverse=-\(A_0\). Hence a single primitive step changes
the finite physical character by m/2; it is not declared a
self-map of the unchanged cap data.

\item Here is the complete intermediate two-step groupoid. For
\(\epsilon\)=0,1,2 set

\[
 A_\epsilon=(-1)^\epsilon A_0,\qquad
 F_\epsilon=z_\infty^{-\epsilon m/2}F_0.
 \tag{GC.9}
\]

\item Then the actual equalities are

\[
 U_{\rm cap} A_\epsilon U_{\rm cap}^{-1}
       =A_{\epsilon+1},\qquad
 F_{\epsilon+1}U_1=U_{\rm cap}F_\epsilon.
 \tag{GC.10}
\]

\item The first equality retains the full root generator, including
its source translation and chosen finite \(\delta\) character.
The scalar \(z_{\infty}\) to the m/2 power has root character -1,
which cancels exactly that of Ucap on the ordinary overlap.
After two steps \(A_2\)=\(A_0\), but \(F_2\)=\(z_{\infty}\) inverse-m times \(F_0\).
This vertical divisor is not removed as a unit. It is
compensated only by the actual \(q_1\) factor in B, giving the last
square in (GC.8), with h still present. Both norm signs n and
every original finite character remain throughout.

\par

\medskip\noindent\textbf{GC.4. The retained derivative and spin maps of the two steps}

\par\nobreak

\item For one graph reset EA.18--20 gives the full complex-four
determinant \(z_{s}\) inverse-squared in this branch. Its chosen
determinant root is the literal function \(z_{s}\) inverse. Thus two
consecutive reset charts carry the product
(\(z_{s}\),0 \(z_{s}\),1) inverse as their determinant root, evaluated at
the actual consecutive source points. Neither source coordinate
is replaced by its value on a central component. On another
projective source chart the canonical volume and this root are
transported by that chart's actual transition; these are bundle
arrows, not an assertion that \(z_{s}\) inverse extends as a function
through zero.

\item The charged overlap EA.26 has complex determinant -1 and the
fixed root i before the normalized charged gauges; it is the
identity on their ordered complex-four coefficients. For a=1
there is no residual mu2 character at this lower link. The
canonical complex spin-c lifts are the exterior-algebra lifts
of these actual derivatives. Their compositions are the lift
of the displayed composed coordinate arrow, with the stated
determinant roots. The positive norm changes have the exact
diagonal and nilpotent derivative homotopies of RT.15, based
at these arrows.

\item The vertical factor \(z_{\infty}\) to the m/2 power in (GC.8) is also
kept as a typed root-line map. In an ordinary nonzero line
coefficient w its derivative has determinant \(z_{\infty}\) to the
m/2 power: the source coordinates are fixed and the additional
base derivative is triangular. Since m/2=2p, the corresponding
determinant root is the single-valued \(z_{\infty}\) to the p power
on this annulus. This statement is in the dw coefficient frame.
In the logarithmic line frame dw/w the same map has determinant
one and the extra derivative is unipotent; the frame arrow
dw=w dlog(w) is retained. In an associated node pair (U,V)
the simultaneous map is \((z_\infty^{m/2}U,z_\infty^{-m/2}V)\),
whose two normal determinants cancel. None of these different
frames is equated without its displayed coefficient arrow.
At the root centre the scalar powers are maps of the
corresponding divisor lines, not nonvanishing scalar units.

\item The canonical complex spin-c lift of the one-line character
-1 has the literal representative {[}e1 e2,i{]}, using the
positive half-rotation in that complex line and \(e_{i}\) squared
equal to -1. Its square is {[}-1,-1{]}, the identity in Spin-c.
On the exterior-algebra spinors the first character acts as
plus one in degree zero and minus one in degree one. Thus the
two character changes in (GC.10) compose to the original
character with its canonical lift, not merely as bare signs
on the underlying line. The determinant frames of the source
reset are still the chronological roots stated above. The full
second-step norm is exactly the B square in (GC.8), so no
additional independent sign or unrecorded ruling reset is
inserted at the cap.

\item Sections GC.5--9 place this entire two-object sequence in the
actual finite source strip, matching the upper z arm, the two
local charged collars and the whole fixed trace. The displayed
local/cap composition is retained in that full assembly.

\par

\medskip\noindent\textbf{GC.5. The entire upper free product, with its coefficient arrow}

\par\nobreak

\item The a=1 upper region z nonzero has a literal product map that
also retains its R=0 exceptional curve. Set W=zV. In the actual
positive auxiliary quotient the map and its inverse are

\[
 (S,T_1,V,z)\longmapsto(S,T_1,W=zV,z),\qquad
 (S,T_1,W,z)\longmapsto(S,T_1,W/z,z).
 \tag{GC.11}
\]

\item Since z has auxiliary weight zero, these maps commute with every
auxiliary chart transition. Thus this region is exactly

\[
 \operatorname{Tot}\bigl(O(-1)_{T_1}\oplus O(-1)_W
          \longrightarrow P^1_S\bigr)\times C^*_z.
 \tag{GC.12}
\]

\item The physical circle acts only on the z factor. The full source
functions are \(q_1\)=\(S_0\) W and \(q_2\)=\(S_1\) \(T_1\), independent of that factor;
the original source parameter is carried by the same analytic
inverse of \(q_1\)-\(q_2\) as in GC.1. In particular \(T_1\)=W=0 is the entire
exceptional \(\mathbb P^1\) times C*, carried by (GC.11), not a point removed
from its domain. Its R=0 locus needs no positive-norm flow to
make this product map defined.

\item The physical line on this piece has a specified trivialization.
Use the convention that a physical positive circle frame p and
its coefficient w represent the vector wp; consequently
(gp,g inverse w) represents that same vector. The actual
weight-one coordinate z gives

\[
                  [p,w]\longmapsto z(p)w.             \tag{GC.13}
\]

\item The formula is invariant and invertible because z is nowhere
zero on this piece. It is the coefficient arrow in this chart,
not an assertion about its extension through z=0 or through the
next source reset. The complex derivative of (GC.11) in the
ordered (V,z) block is

\[
              \begin{pmatrix}z&V\\0&1\end{pmatrix},
 \tag{GC.14}
\]

with all source coordinates retained. This carries the canonical
complex spin-c structure by its actual exterior-algebra map.
Its determinant z is retained; no single-valued square root on
the z circle is asserted. The chart's physical line trivialization
alone does not discard a transition contribution at its boundary.

\item The marked G divisor in the final node coordinates is \(S_1\)-\(S_0\)=0,
by EA.45. In particular the coordinate \(S_0\)=F+ in (GC.11) is not
the old Cartesian coefficient s0. Any next-source transition
uses the actual EA.43--46 unit and the chronological reset,
before one assigns a marked source section to an upper point.
The full two-step outside attachment remains the next map to
be checked; (GC.11)--(GC.14) establish the complete upper chart
without assigning that attachment or an individual return.

\par

\medskip\noindent\textbf{GC.6. The proper upper graph and its actual primitive endpoint}

\par\nobreak

\item Put \(T_0\)=W=zV in (GC.11), and use the four invariant conifold
functions

\[
 q_1=S_0T_0,\quad q_2=S_1T_1,\quad
 C=S_0T_1,\quad D=S_1T_0,\qquad q_1q_2=CD.
 \tag{GC.15}
\]

\item The rational opposite source ruling {[}\(T_0\):\(T_1\){]} is retained by its
actual graph, not made into a coordinate at \(T_0\)=\(T_1\)=0. Its proper
common resolution has homogeneous coordinates {[}\(S_0\):\(S_1\){]}, {[}t0:t1{]},
a coefficient \(\eta\) of O(-1,-1), and the nonzero z. Its two
projections are

\[
 ([S],[t],\eta,z)\longmapsto
       ([S],T_0=\eta t_0,T_1=\eta t_1,z),
\]
\[
 ([S],[t],\eta,z)\longmapsto
       ([t],A_0=\eta S_0,A_1=\eta S_1,z).
 \tag{GC.16}
\]

\item Explicitly the two auxiliary scalings multiply S and t by
\(\alpha\) and \(\beta\) and \(\eta\) by (\(\alpha\) \(\beta\)) inverse. This gives
both inverse ordinary blowup charts wherever t0 or t1 is
nonzero. The first projection is precisely the blowup of the
normal ideal (\(T_0\),\(T_1\)) of the upper exceptional \(\mathbb P^1\); the second
contracts the other projective factor. Both retain the source
functions \(q_1\)=\(\eta\) \(S_0\) t0 and \(q_2\)=\(\eta\) \(S_1\) t1. The physical circle
acts only on z, so the graph and both proper source products
are free on this entire upper piece. Their physical coefficient
arrow is the pullback of (GC.13), including \(\eta\)=0.

\item Here the full \(\theta\) unit has a useful exact extension. Retain
Q=\(q_1\) \(q_2\) uQ, P=uP/\(q_2\) and uG=uQ uP. In the \(S_1\) chart TL uses
S=\(S_0\)/\(S_1\) and

\[
 E=\frac{2\Theta_Q(S)}{(1-S)\Theta_Q(S/P)}.
 \tag{GC.17}
\]

\item This unit descends to the full conifold (GC.15). Indeed its
normally convergent product factors are expressed by

\[
 QS=u_Qq_1C,\quad Q/S=u_Qq_2D,\quad
 S/P=C/u_P,\quad Q/(S/P)=u_GD.
 \tag{GC.18}
\]

\item For every higher power of Q multiply the first two expressions
by that power minus one. The apparent marked-section factor
1-S is already divided out. All numerator and denominator
factors are therefore holomorphic and nonzero on a sufficiently
small full conifold neighbourhood, including both S charts and
both target t charts. Their product converges normally there.
The value at the vertex is exactly two. This supplies an
extending logarithm of E/2 on that neighbourhood, without
replacing E by its vertex value. In particular

\[
                              L=E z
 \tag{GC.19}
\]

is the same actual nonzero physical coefficient over the whole
upper graph. It descends through each projection in (GC.16),
since E is a function of the four common invariant functions.

\item The endpoint is the actual inverse primitive. On the original
Tate torus, with the precise source/Cartier map of TL and SZ.43,
its input coefficients are

\[
              Z_0=u_GD,\qquad V_0=P(C/u_P)L.
 \tag{GC.20}
\]

\item The identity is first read in the \(S_1\) chart, where it is TL.11
with the actual U=\(S_0\) z coefficient. The homogeneous line-frame
arrow on the other S chart is exactly TL's retained Cartier
transition; (GC.18) gives its extension. Applying \(U_1\) inverse,
which has the literal torus map (Z,V) to (PZ,\(c_{X}\) V), gives

\[
 Z_1=Pu_GD,\qquad V_1=u_X C L,
                  \qquad c_X=q_2u_X .
 \tag{GC.21}
\]

\item No \(q_2\) factor or source translation has been discarded. On the
t1 chart of (GC.16), with t=t0/t1 and coordinates (\(A_0\),\(A_1\)),
these formulas read

\[
 q_1=A_0t,\quad q_2=A_1,\quad C=A_0,\quad D=A_1t,
\]
\[
 Z_1=u_Gu_Pt,\qquad V_1=u_XA_0L,\qquad
 L_1=Z_1V_1/q_1=(u_Gu_Pu_X)L .
 \tag{GC.22}
\]

\item The last expression is interpreted in its actual toric edge
chart. The unit character on the unshifted negative diagonal
\{e1,e2\} is ZV/\(q_1\). Thus \(L_1\) is an extending nonzero coefficient
at \(A_0\)=0 and t=0, rather than a division performed only after
restricting to that divisor. The source node coordinate
\(q_1\)/\(Z_1\) is \(A_0\)/(uG uP), with its full unit.

\item On the t0 chart put p=t1/t0, so \(q_1\)=\(A_0\), \(q_2\)=\(A_1\) p,
C=\(A_0\) p and D=\(A_1\). Carry the actual source period
h:(Z,V) to (QZ,PV), including P. Then

\[
 Z_2=QZ_1=q_1u_G^2D,\qquad
 V_2=PV_1=u_Pu_X(A_0/A_1)L,
\]
\[
             L_2=Z_2V_2/q_1^2=u_G L_1.
 \tag{GC.23}
\]

\item These are the other indexed node/edge coordinates. Relative to
the first \(q_1\) fan chart the h arrow translates the edge by e1;
its unit character is ZV/\(q_1\) squared. The full \(q_2\)-dependent
period P is still in (GC.23). The target edge coordinates on this second chart must also retain
the two S charts. With \(S_1\)=t0=1, \(\eta\)=r, \(S_0\)=s and t1=p, they are

\[
 q_1=rs,\quad q_2=rp,\quad
 Z_2/q_1=u_G^2r,\quad q_1^2/Z_2=s/u_G^2.
 \tag{GC.23a}
\]

\item All four quantities extend through r=s=p=0. With \(S_0\)=t0=1,
\(\eta\)=r, \(S_1\)=s and t1=p, one instead has \(q_1\)=r,\(q_2\)=rsp. The
ratio \(q_1\) squared over \(Z_2\) is not a coordinate at s=0. At that
endpoint retain the actual adjacent source Tate-node pair

\[
 w=Z_1^{-1}=p/(u_Gu_P),\qquad
 z_T=QZ_1=u_G^2r^2s,\qquad z_Tw=Q.
 \tag{GC.23b}
\]

\item The corresponding nonzero line coefficient is
Y0=\(Z_1\) \(V_1\)/\(q_1\)=\(L_1\). Its h-side frame is Yh=\(q_1\) uG Y0.
The latter factor is retained as a divisor, not called a unit:
the adjacent associated node coordinates are

\[
           (U_h,V_h)=(q_1u_GY_0,(u_GY_0)^{-1}),
                         \qquad U_hV_h=q_1.
 \tag{GC.23c}
\]

\item These are regular even at r=s=p=0 because Y0 and uG are
nonzero. The proper graph projection in (GC.16), with its
{[}S{]} and {[}t{]} coordinates, is still retained. No inverse to its
exceptional contraction is asserted. In the remaining two
charts t1=1 one has \(Z_1\)=uG uP t0 and
\(q_1\)/\(Z_1\)=\(\eta\) \(S_0\)/(uG uP), with \(q_2\)=\(\eta\) \(S_1\); these are regular for
either \(S_0\)=1 or \(S_1\)=1. Thus every endpoint of \(\mathbb P^1_S\times\mathbb P^1_t\)
has a specified source-node and associated-line-node map, not
an unextended ratio. On the common edge overlap \(L_2\)=uG \(L_1\) is
an extending unit, whose logarithm is carried with the original
Cartier map. This is not a free change of suspension marking.

\item The inverse primitive belongs to the two-object cap sequence
(GC.9)--(GC.10). In the direction used above its root square is

\[
                    F_\epsilon U_1^{-1}
                         =U_{\rm cap}^{-1}F_{\epsilon+1}.
 \tag{GC.24}
\]

\item Both finite characters are retained. The complete B inverse is
\(U_1\) inverse twice followed by the scalar (minus2 \(\lambda\) \(q_1\))
inverse, and its cap square is the inverse of (GC.8).
Equations (GC.16)--(GC.24) give the full upper proper graph and
its one-primitive endpoint through the intermediate object.
They do not replace that composite by its normal-cone limit or
assert that one primitive is an automorphism of the original
cap object. Section GC.9 glues the two chronological primitive
graphs through this intermediate object and the single whole
fixed trace, using the explicit trace of GC.8.

\par

\medskip\noindent\textbf{GC.7. The full two-step upper line character}

\par\nobreak

\item For the complete good suspension, put \(c_{\lambda}\)=minus2 \(\lambda\).
The actual inverse, including its selected \(q_1\) factor, is

\[
 B^{-1}:(Z,V)\longmapsto
       (P^2Z,c_X^2V/(c_\lambda q_1)).
 \tag{GC.25}
\]

\item Substituting (GC.20), with CD=q1q2, gives

\[
 Z_B=P^2u_GD,\qquad
 V_B={q_2u_X^2C\over c_\lambda q_1}L,\qquad
 Z_BV_B={u_Gu_P^2u_X^2\over c_\lambda}L.
 \tag{GC.26}
\]

\item The product in the last expression is the unit character of
the actual height-zero indexed edge. Its coefficient is the
full extending base unit \(F_{B}\)=uG uP squared uX squared divided
by \(c_{\lambda}\). \(F_{B}\)(0)=1/\(c_{\lambda}\), and \(F_{B}\)/\(F_{B}\)(0) has an
extending logarithm. Both primitive source translations and
the \(q_1\) factor were used before this cancellation. In particular
(GC.26) does not discard the p-dependent multiplier in SZ.34.

\item The source graph of the two chronological primitive maps may
be taken with both projective intermediate source coordinates.
Its line factor is the same C* coordinate L, with its target
coefficient changed by the explicit nonzero \(F_{B}\). Thus the
actual upper source/line graph is the product of that proper
source graph with C*\_L under the typed target map
\(Y_{B}\)=\(F_{B}\) L. The physical action is scalar on this factor.
All source functions are pulled back by the graph projections.
This proves the physical-line trivialization on that graph,
including its projective source endpoints; it does not assert
that the source graph itself is smooth or that either of its
projections is an isomorphism at a contracted curve.

\item The full finite-root square is

\[
 F_0B^{-1}=(c_\lambda h)^{-1}U_{\rm cap}^{-2}F_0,
                    \qquad q_1=z_\infty^{-m}h.
 \tag{GC.27}
\]

\item It is the composition of the two different cap-object maps
in (GC.24), with \(F_2\)=\(z_{\infty}\) inverse-m times \(F_0\) still retained.
The scalar character returns only after that composition.
No additional primitive ruling change is inserted at the cap.
The finite smooth cut of this full proper source graph is
constructed in GC.8--9, with its physical line and the whole
fixed trace carried on every collar.

\par

\medskip\noindent\textbf{GC.8. A finite smooth trace for the upper proper graph}

\par\nobreak

\item The upper source graph in (GC.16) has the following explicit
smooth real trace. Use its common source Cox coordinates
(s0,s1,t0,t1), with compact auxiliary circle weights
(+1,+1,-1,-1), and put

\[
 \mu=(\|t\|^2-\|s\|^2)/2,\qquad R=\|s\|\|t\|,
                  \qquad \rho_s^2=\|s\|^2+\|t\|^2.
 \tag{GC.28}
\]

\item Choose positive constants with r squared divided by two smaller
than both \(\delta\) and R0. Take the compact region

\[
 -\delta\le\mu\le\delta,\qquad R\le R_0,
                         \qquad \rho_s^2\ge r^2
 \tag{GC.29}
\]

and divide by that compact auxiliary circle. Compactness follows
from the exact identity \(\rho_{s}\) fourth=4mu squared+4R squared.
The auxiliary action is free off the removed origin because all
four nonzero coordinate weights have absolute value one. The
moment levels have nonzero differential away from the origin;
on the R=R0 side both s and t are nonzero, and the two moment
and product-radius differentials are independent. The round
inner boundary lies strictly between the two moment faces and
strictly inside the product-radius side by the chosen inequalities.
Invariant rounding therefore gives a compact smooth seven-manifold
with precisely its two moment faces, the stable side, and one
full inner quotient of \(S^7\) by this auxiliary circle.

\item The two moment faces are the actual two source small resolutions.
On the s-stable face, positive auxiliary normalization takes a
Cox orbit to that specified moment value and retains {[}s{]} and the
products \(s_{i}\) \(t_{j}\); on the t-stable face it retains {[}t{]} and those
same products. On their R\textgreater0 graph collar both projective factors
are present and the projections are exactly (GC.16). Each moment
face also includes its own small-resolution zero section: t=0
on the s-stable face and s=0 on the t-stable face. The opposite
projective factor is not asserted to be defined at that zero
section. These endpoints are supplied by the actual moment
face, while the proper graph retains both factors only on its
specified comparison collar. The positive auxiliary flow
(s,t) to (a s,a inverse t) has strictly monotone moment on each
nonzero orbit. On the R=R0 side it supplies the unique map to
any prescribed intermediate moment. It preserves every \(s_{i}\) \(t_{j}\),
and hence the full \(q_1\),\(q_2\), the analytic inverse giving v, and the
unit E in (GC.18). This is the typed stable-side identification,
not a radial normalization of q or v.

\item Carry the physical coefficient L of (GC.19) and take its unit
circle with respect to the actual transported metric. The physical
circle acts on this factor alone. The resulting eight-dimensional
trace carries the same physical line with the coefficient arrow
{[}p,w{]} to L(p)w. At a target face use its actual coefficient
\(L_1\) or \(L_2\) from (GC.22)--(GC.23); the metric is transported by
that full nonzero multiplier. The extending logarithms of the
unit factors give the explicit interpolation of coefficient
frames through the trace. Constants and their phases are kept.
For the composed endpoint use the full \(F_{B}\) in (GC.26).
Thus all maps on the stable side agree with the ordinary upper
finite strip, including both source- and line-node frame changes
of (GC.23a)--(GC.23c).

\item The one new full inner face is \((S^7/S^1_{\rm aux})\times S^1_L\). Conjugating
the two negative auxiliary coordinates identifies its first
factor with \(\mathbb{CP}^3\) by the literal positive Hopf map; the physical
circle remains the independent second factor. Denote its oriented
physical quotient by \(Q_{up}\), with the orientation induced by this
actual boundary face. Its physical associated line has the global
nonzero frame just displayed. Consequently its first Chern class
is zero and both x cubed and \(p_1\) times x have zero integral.
No physical charge-one Hopf line is substituted for this trivial
line. In particular no topology of \(Q_{up}\) is inferred from its
vanishing characteristic numbers.

\item The trace's stable tangent and determinant maps are the actual
auxiliary quotient maps of these Cox coordinates. The sum of its
four auxiliary weights is zero, so the determinant frame is the
one descended from their ordered complex volume, with the orbit
and moment directions explicitly removed. The physical L circle
uses its original complex coefficient frame dL and its radial
normal. A passage to dL/L carries the factor L, rather than
silently changing the spin along that circle. Unit coefficient
changes are lifted with their actual extending logarithms; the
finite character maps are the canonical lifts of GC.4. This fixes
the determinant/spin-c maps on every interface of this free trace.
It creates no additional physical fixed stratum.

\par

\medskip\noindent\textbf{GC.9. Chronological two-primitive finite strip and its boundary ledger}

\par\nobreak

\item Fix the scalar by the actual RT.12 based homotopy. Put
\(c_{\lambda}\)=minus2 \(\lambda\) and choose a nonzero constant c0. Factor

\[
 K_0=c_0U_1,\qquad K_1=(c_\lambda/c_0)q_1U_1,
                              \qquad B=K_1K_0.
 \tag{GC.30}
\]

\item The intermediate object is exactly (\(A_1\),\(F_1\)) from (GC.9). The
finite-root arrows on these two pieces are

\[
 F_1K_0=c_0U_{\rm cap}F_0,\qquad
 F_0K_1=(c_\lambda h/c_0)U_{\rm cap}F_1.
 \tag{GC.31}
\]

\item The second equality uses \(F_0\)=\(z_{\infty}\) to the m times \(F_2\) and the
retained \(q_1\)=\(z_{\infty}\) inverse-m times h. Thus the first root
character is shifted by m/2 and the second returns it, with its
canonical spin-c composition. Their product is exactly the
original B square, including h, norm sign and every finite
character. The primitive object is never identified with the
original one merely by its underlying source surface.

\item On the compact source and root complement of the collision graph,
choose c0 small enough to make the first transported norm cut
contracting. Then choose the fixed \(\lambda\) small enough that the
second cut is contracting too. Both maps and their inverses are
bounded there. Changing \(\lambda\) within this same nonzero range
uses the zero-winding radial path already retained in RT.12.
On the graph collars use the actual Fubini--Study norms
Nt=\textbar xi\textbar{} \textbar\textbar t\textbar\textbar{} and Ns=\textbar xi\textbar/\textbar\textbar s\textbar\textbar{} from EA.12. Their ratio R is
retained, and the actual positive norm map is the a=1 instance
of EA.13. Choose the cut between \(K_0\) and \(K_1\) as a distinct marked
face. This gives two chronological compact primitive pieces,
with one intermediate ordinary face and the root object (\(A_1\),\(F_1\)).
Their intermediate faces are glued by the actual coefficient
identity, with opposite outward normals. Their full finite-root
faces are glued by (GC.31), not by discarding a finite character.

\item There is one pure nonzero-line Cox-origin cut for \(K_0\). In its
common xi=1 gauge the source and target finite faces are
\(D^4\)\_s times \(S^3\)\_t and \(S^3\)\_s times \(D^4\)\_t, with their physical
weights (-1,-1,+1,+1). The actual contracting direction enters
the first and exits the second, exactly as in EA.40--41.
Their target-minus-source full face is therefore one positive
\(S^7\) with its charge-one physical circle and line O(-1).
Call its physical quotient \(P_0\). Its physical characteristic pair
is (-1,-4), so its index is zero. Its label is the first primitive
piece; it is not glued to the second primitive origin by a source
period. That period preserves the intermediate root character,
whereas passing from object zero to object one changes it by
m/2. The original B closes only after both pieces.

\item The \(q_1\) factor occurs only in \(K_1\). Use its actual associated node
and the compact domain D1 of (GC.3), with the lower horn and
full charged annulus of GC.2. The complete chain identity
EA.50--52, now with the exact GC.4--6 defining functions, gives
one further positive full charge-one link \(P_1\). Its source/line
node coefficient includes \(c_{\lambda}\)/c0. There is exactly one
adjacent-node height shift from this \(q_1\) factor. It is retained
in the final B endpoint (GC.25)--(GC.27), and is not assigned
a second time to \(K_0\). The physical pair of \(P_1\) is again (-1,-4).
This chart has no isolated positive-chamber physical fixed point:
z has auxiliary weight zero, and the only fixed locus is the
whole surface z=V=0 already identified in GC.1.

\item The remaining upper face of this associated cut is the complete
free product of GC.5, including its exceptional \(\mathbb P^1\). Use the
proper graph and both target charts of GC.6 to carry its actual
inverse primitive to the intermediate object. Insert the finite
source trace of GC.8 on this free upper part. Its stable side
is the exact norm-strip map, because the positive auxiliary flow
preserves the common source matrix and the unit E. Its target
line coefficients are precisely \(L_1\),\(L_2\) and the composed \(F_{B}\) L.
The source Tate-node and associated-line-node arrows (GC.23b)--
(GC.23c) retain the final projective corner. Thus this upper
insertion has one additional full free face \(Q_{up}\), and does not
add another charge-one origin or another fixed surface.

\item The fixed locus is joined to the outside by the actual RT.1
normalization immersion and the full EA.42--46 Cartier map.
It is the whole F=\(\operatorname{Bl}(\mathbb P^1\times D)\), with its one exceptional E.
\(K_0\) is a nonzero-line piece and the inserted upper trace is free,
so neither contributes another fixed component. The local fixed
\(S^3\) is retained inside F and receives no separate \(D^4\) cap. The
one global boundary \(S^2\times S^1\) is capped by the actual C0
product RT.33a--b with its positive-spinor coefficient arrow.
Use the exact RT.34 triangle bordism to compare it with \(C=\mathbb P^1\times D^2\), retaining both g/w clutches and the whole \(S^4\) normal
bundle. This contributes the two deleted-tube faces -\(B_F\) and -\(B_4\).

\item Every remaining ordinary B face is paired by K1K0=B; every
root face is paired through (GC.31); and the source-base collar
is paired through the full TL \(\theta\)/Cartier map. The inequalities
\textbar{}\(q_1\)\textbar+\textbar{}\(q_2\)\textbar\textless=R place all inner changes away from the prescribed
outer parameter circle and source-base boundary. The normal
maps on intersecting collars are their written common Cox maps;
their products on a triple overlap are the actual identity.
The two chronological determinant roots and their finite
characters are those of GC.4. The free upper trace has its
specified determinant frame and physical coefficient. Thus no
unrecorded order-two spin character is used in joining these
faces. The construction uses finitely many compact norm pieces,
the two actual source projective charts at each primitive step,
and the one compact source-moment trace. All identifications are
between its specified closed collars. Rounding yields the finite
smooth oriented orbit bordism.

\item Let N be its actual outside conifold/Hopf completion. Give \(Q_{up}\)
the sign induced by its boundary inclusion in this construction.
The resulting physical orbit boundary has the full face count

\[
                  -N-P_0-P_1-B_F-B_4+Q_{\rm up}.
 \tag{GC.32}
\]

\item The minus signs of \(P_0\) and \(P_1\) follow from their positive total-space
finite-end chains and the circle-first boundary reversal. The
fixed-tube signs are the actual RT.34 cap signs. The free face
\(Q_{up}\) has its independently specified trivial physical line.
There is no charge-zero projective substitution or an additional
isolated positive-chamber point in this ledger.

\item The characteristic expression on \(B_F\) is the full RT.32 expression
with a=1 and both retained clutch integers, hence is zero. On \(B_4\)
it is zero with arbitrary normal \(c_2\) retained. The two charge-one
faces and \(Q_{up}\) have just been computed as zero from their actual
line maps. The outside characteristic expression is therefore
zero by Stokes for the carried physical line and stable tangent.
The outside boundary is the actual fixed-scalar source loop,
because all constructions are identity modifications off the
inner collision collar and retain (GC.31) on the cap. The
physical/Jin inverse-return arrow is the one in K68; the zero
value is unchanged by that explicitly carried reversal. This
completes the good-branch boundary proof. Section GC.10 gives
the retained-clutch evaluation and the full actual-return arrow.

\par

\medskip\noindent\textbf{GC.10. The carried characteristic expression and actual parameter return}

\par\nobreak

\item Write M for the actual physical nonzero-vector line and x=\(c_1\)(M).
On the free quotient bordism choose invariant connections product on
all the specified collars and patch them on its compact interior.
The circle tangent is oriented and trivial, so the actual tangent
splitting identifies its \(p_1\) with the pullback of the quotient \(p_1\).
Stokes applies to the closed polynomial

\[
                      A(Y)=\int_Y(4x^3-p_1(TY)x)/24.
 \tag{GC.33}
\]

\item Every line in this formula is the one carried by the written
coefficient maps. In particular the two chronological charged
faces have x=-h and \(p_1\)=4h squared on their positively oriented
\(\mathbb{CP}^3\) quotients, and hence A(\(P_0\))=A(\(P_1\))=0. The upper free face has
its actual coefficient frame L and x=0, so A(\(Q_{up}\))=0.

\item Here is the whole fixed-cap evaluation with the actual clutches
retained. On F-C=\(\operatorname{Bl}(\mathbb P^1\times\mathbb P^1)\), use H,B,E with H squared and
B squared zero, HB=1, E squared=-1 and HE=BE=0. If the full
normal/cap coefficient ratios have the integers w+ and w- of
RT.29, the two positive spinor line classes for a=1 are

\[
        l_+=H-E+w_+B,\qquad l_-=-H+w_-B.
\]

\item Thus b=\(c_1\)(\(E_{\mathrm{spin}}\))=-E+(w\_++w\_-)B has b squared=-1 and
\(c_2\)=w\_--w\_+. The actual Hopf coefficient arrow makes M=O(-1)
on its projectivization. With h=\(c_1\)(O(1)), the relation
h squared+b h+\(c_2\)=0 and pi\_*h=1 give

\[
 k=c_2-b^2=1+w_--w_+,\qquad
 \ell=-3\sigma-b^2+4c_2=4+4(w_--w_+).
\]

\item Here \(\sigma\)=-1. Consequently A(\(B_F\))=0 with both integers retained.
The exact triangle cap comparison gives a normal bundle on \(S^4\)
with arbitrary retained \(c_2\); its pair is (\(c_2\),4c2), so A(\(B_4\))=0.
There is no use of a vanishing normal clutch. Substitution in
the actual oriented ledger (GC.32) therefore gives A(N)=0.

\item To identify this with the requested parameter return, take a
physical-\(\delta\)-averaged connection on the actual completed
family, product on the selected tube in the exact RT.33a--b
Cox/Hopf coordinates. Its conifold quotient return is a
relative diffeomorphism f of \(S^2\) times \(S^3\), with the stated
identity product on that tube. It preserves the actual Euler
class, which generates \(H^2\) cohomology. Its orientation then
preserves \(H_3\), so f is in the required homology-trivial relative
group. The physical Hopf lift of P3F.22 is literally the return
on this same complement and the same identity extension on
the selected tube. P3K.68 supplies the full oriented pieces
and the inverse-clutch arrow

\[
                  \lambda_M(f)=-\Sigma(f^{-1})=\Sigma(f).
 \tag{GC.34}
\]

\item Hence the fixed-scalar source return has physical integer
A(N)=0 and oriented class zero in \(\Theta_7\). This last step uses
the actual source, cap, tube and line maps; it does not select
a return representative from the zero polynomial alone.

\item For the literal reflected fan the already proved single
native Cox comparison of P3K.60 has free pair (1-a,a)=(0,1).
The same physical relative convention of P3K.68--69 therefore
adds one to the central class. Its reflected fixed-scalar
class is one. This does not assign a new reflected integer
characteristic pair from that residue.

\item Finally the canonical scalar-root loop in this good branch
has physical index zero in these two actual fans, by the
retained scalar page computation and its physical-line
correction, with the narrow native C2 comparison of P3K.62.
The exact RT.12 based phase homotopy, at
\(r_{\rm base}=v_*^k r_0\) and \(r_0^m=\lambda_0\), adds k such actual scalar
turns to the fixed-scalar source loop. Thus every retained
coupled loop \(\lambda=\lambda_0v^{km}\) has the same two oriented
classes zero and one. Both norm signs and all realized finite
\(\delta\) characters were kept in (GC.9)--(GC.10) and (GC.31).
No class in \(\Theta_7\) has been divided.

\item The certificate \path{code/part3/source_zero_good_cut.py} checks the
exact cut, node, \(\theta\)-factor, cap-square, determinant and
moment identities. Its first extended run reached the literal
charge-one index and encountered a Python float-zero versus
exact SymPy zero comparison; the code and output are retained
in \path{records/attempts/source_zero_good_trace_float_zero}. The
correction uses Rational for that literal zero only. All
preceding geometric algebra had already passed. The current
71-check certificate complements the geometric face proof
and its independent reviews.

\end{enumerate}
\end{proof}




\begin{theorem}\label{night:P3K.84}
\StatementRecord{night:P3K.84}{proved}{S6\_06\_BLIND: GPT-6 Astra (OpenAI, ChatGPT Pro mode, effort Pro), recorded re-derivation in the checkpoint dated 18 November 2026}{\cref{night:P3F.26}, \cref{night:P3F.33}, \cref{night:P3F.34}, \cref{night:P3F.35}}
Retain the fixed OS49 source, its actual component Cartier line
\(M=\mathcal O(G-O)\), translation \(-3G\) and selected cusp. In the
specified two-cap construction take the \(IV^*\) root order \(m=3k\)
and reduced-II order \(h=6l+1\), with \(k,l>0\), complete norm vectors
\(n_A\ell_0+d_A\delta\) and \(n_B\gamma+d_B\delta\), and their actual
Cartier/root/annular maps. The sphere sector is exactly
\[
                    h n_A+m n_B=\pm1.
\]
Every such datum is constructed for arbitrary fixed \(d_A,d_B\), with
all source and line factors retained, and yields the standard smooth
\(S^6\) on every allowed cusp fan. These data exist exactly when
\(\gcd(3k,6l+1)=1\). The smallest displayed instance is
\((m,h,n_A,n_B,d_A,d_B)=(3,7,1,-2,0,0)\).
The proof gives every actual specialization, integral Mayer--Vietoris,
translation-prism and fundamental-group map. This statement does not assign
additional ordinary clutches or a parameter return.
\end{theorem}
\begin{proof}
\begin{enumerate}[label=\arabic*.]
\item The full proof has passed the recorded root and independent audits. Keep the exact fixed OS49 source,
\(M=\mathcal O(G-O)\), \(Q=3G\) and the original component-supported selected cusp. Put
\(m=3k\) for the IV-star cap, \(h=6l+1\) for a reduced-II cap, with k,l positive.
The full cap vectors are \(n_A\ell_0+d_A\delta\) at IV-star and
\(n_B\gamma+d_B\delta\) at II. This note constructs the roots and records
all integral specialization and fundamental maps. It assigns no return.

\par

\medskip\noindent\textbf{RI.1. Exact reduced root, marked section and Cartier diagram}

\par\nobreak

\item Write \(q=t-1\). The original equation and its actual marked section are
\[
 y^2=x^3-3(1+q)^3q x+(1+q)^4q(1+2q),\qquad
                         G=((1+q)^2,(1+q)^2).
\]
On \(q=z^h,\ x=z^{2l}X,\ y=z^{3l}Y\), the relatively minimal equation is
\[
 Y^2=X^3-3(1+z^h)^3z^{2l+1}X
                 +z(1+z^h)^4(1+2z^h).                  \tag{RI.1}
\]
It has a reduced II fibre and smooth total surface, since the last
term has derivative one in z at the cusp. Its actual deck and section
satisfy
\[
 D(z,X,Y)=(\zeta_hz,\zeta_h^{-2l}X,\zeta_h^{-3l}Y),
 \quad \tau'(G')=-X(G')/Y(G')=-z^l,\quad \tau'(O)=0.    \tag{RI.2}
\]
On the normalized central cusp \(X=s^2,\ Y=s^3\), D acts by
\(s\mapsto\zeta_h^{-l}s\). Its order is h because \(\gcd(l,6l+1)=1\).
Both \(G'\text{ and }Q'=3G'\) tend to O in the smooth additive group.

\item Here is the actual line map, retaining the marked source coordinates.
At the original smooth O and G the coordinate \(\tau=-x/y\) has values
0 and -1. The source change gives \(\tau'=z^l\tau\) exactly. Resolve its
indeterminacy at the new O by the successive charts
\[
 \tau'=z^j u_j\quad(1\le j\le l),\qquad
                 \tau=u_l,\quad q=z^h.                \tag{RI.3}
\]
These are the consecutive point blowups in the graph of the rational
source map. Before the last blowup both marked sections pass through
the same centre; on the last chart they are \(u_l=0\) and \(u_l=-1\). At every
intermediate exceptional curve the two pulled-back section divisors
have the same coefficient j. Their difference therefore has coefficient
zero. The remaining graph resolution is over the original cusp, away
from G and O, where their divisor line has its original nonzero frame.
Consequently the two actual graph projections \(\beta\) to the root and pi
to the original source give the exact Cartier equality
\[
         \beta^*N=\pi^*M,\qquad N=\mathcal O(G'-O').     \tag{RI.4}
\]
On (RI.3) its local ratio is \((u_j+z^{l-j})/u_j\) and on the last chart
it is \((u_l+1)/u_l\); all vertical powers have literally cancelled.
The formula uses these functions only in their marked-section charts;
the other source points where the rational function \(\tau+1\) vanishes
are not inserted into the divisor. The original Cartier frames are
used on their complement. On the chart (RI.3), D sends \(u_j\) to
\(\zeta_h^{l-j}u_j\), so both numerator and denominator have the same
character. Thus the diagram and the line map commute with D. The
central line \(N|_{II}\) is trivial with the resulting D weight zero.

\item Multiplication by h has a unique formal-group inverse at O. Applied
to Q' it gives the extending D-invariant section R, \(hR=Q'\), \(R(0)=O\).
There is no other extending h-torsion section: the root's elliptic
monodromy is the rank-two elliptic block of F26, since h=1 modulo six:
\[
 B_{\rm ell}=\begin{pmatrix}1&1\\-1&0\end{pmatrix}
 \quad\text{on cohomology},\qquad
 \widehat B_{\rm ell}=(B_{\rm ell}^{-1})^t
       =\begin{pmatrix}0&1\\-1&1\end{pmatrix}
 \quad\text{on homology}.                               \tag{RI.4a}
\]
Both \(B_{\rm ell}-I\) and its displayed homological counterpart have determinant
one. An invariant torsion point in that rank-two elliptic lattice is
therefore zero. No such assertion is made for the full rank-four
matrix, which retains its suspension and \(\delta\) invariant directions.
The source map \(t_{-n_BR}D\) is consequently the actual reduced-root
norm source map; its hth power is \(t_{-n_BQ'}\).

\item Lift \(t_{-R}\) on \(N\) through (RI.4). Its commutator with D is a base unit.
At the central fibre the source lift is the identity and the line is
trivial, so its finite commutator character is one. The exact finite
average of the zero-constant unit logarithm makes the lift F commute
with D; this is a conjugation through extending line units, not an
average of the lattice cochain. Write \(B_0=\beta(q)F^h\). The base unit
\(\beta\) is D-invariant. Choose its actual analytic root with the prescribed
reference value and put
\[
 A_0=\beta(q)^{n_B/h}F^{n_B}D,\qquad A_0^h=B_0^{n_B}.
                                                               \tag{RI.5}
\]
Multiplication by the retained character \(\zeta_h^{d_B}\) and the scalar
norm factor gives the full actual cap. Its ordinary-to-cap arrow
retains these source and line lifts and has normal relation
\(h\alpha_B-n_B\gamma-d_B\delta\). No integer period has been inserted.

\par

\medskip\noindent\textbf{RI.2. Central bundle, freeness and every specialization map}

\par\nobreak

\item Write \(y=(\gamma,\delta)\) in the two actual suspension/physical period
coordinates. The central fibre is the quotient
\[
 (\mathbb P^1_s\times T^2_y)/\langle
       (s,y)\longmapsto(\zeta_h^{-l}s,
                        y+(n_B,d_B)/h)\rangle.          \tag{RI.6}
\]
Normalization of the cuspidal source is a homeomorphism, so this gives
its actual topological central fibre. The two source points \(s=0,\infty\)
are fixed by every deck power. The cap action is free exactly when
\(\gcd(h,n_B,d_B)=1\). Its physical circle is free exactly when \(\gcd(h,n_B)=1\):
a nonintegral \(\gamma\) shift prevents isotropy, while a power with integral
\(\gamma\) shift fixes s=0 and can be cancelled by a physical-\(\delta\) rotation.
These conclusions include the retained line character \(d_B\).

\item The sphere sector below has \(\gcd(h,n_B)=1\). More generally, for every
free (RI.6), choose integers a,b with \(a n_B+b d_B=1\) modulo h. The map
\[
 (s,\gamma,\delta)\longmapsto
 \left(e^{2\pi i l(a\gamma+b\delta)}s,
              [\gamma,\delta]\right)                  \tag{RI.7}
\]
identifies (RI.6) with \(S^2\) times the enlarged torus
\(\mathbb R^2/(\mathbb Z^2+\mathbb Z(n_B,d_B)/h)\). Its equivariance follows by substituting the
literal generator in (RI.6). It fixes both marked source sections.
The physical circle action is carried by the displayed map, including
the induced rotation of its first coordinate. Integral cohomology is
therefore free; this is not only a rational invariant calculation.

\item The source two-class \(\eta\) is normalized to integrate one on the oriented
source sphere and zero on the O torus. The root formal translation flow of Q' gives an actual generic torus
section. In the normalized commuting lift of RI.1 write \(F_t\) for the
small extending flow with \(F_1=F\). It is the exponential of the formal
Lie logarithm of the already normalized commuting source/line lift near
its identity. Its additive Lie coefficient gives \(F_{ht}=F_t^h\) and
commutes with D. Choose the retained full logarithm
of the nonzero scalar unit \(\lambda\beta(q)\). The literal flow
\(B_t=\exp(t\operatorname{Log}(\lambda\beta))F_{ht}\) satisfies \(B_1=B\). At suspension
coordinate \(\gamma\), its value on the rigidified O frame is the source
point \(-\gamma Q'\), with that full line multiplier. The map
\(\sigma(\gamma,\delta)=e^{2\pi i\delta}B_\gamma(O,1)\) is therefore the actual
torus section.
Changing \(\gamma\) by one is exactly B, and changing it by \(n_B/h\) is exactly
the source/line part of A. It extends to the central O torus because
\(Q'(0)=O\). In these literal coordinates its source forms u,w restrict to zero.
This is the original F26 marking: the root formal Abel lift has zero
B-monodromy cocycle, and any difference from that marked lift would be
an integral elliptic vector fixed by the homological matrix \(\widehat B_{\rm ell}\)
in (RI.4a). Its difference from the identity is unimodular in that
rank-two lattice, so the elliptic difference is zero. The retained scalar logarithm may change by its actual integer
\(\delta\) period; it is kept in \(\sigma\) and cannot create a source u,w term. The source collapse has degree one on its oriented two-cell and
kills its two one-cycles. Thus specialization of \(\eta\) is precisely uw,
with no additional \(\psi\delta\) term. This is the marked-section arrow
that fixes the possible integral cross term.

\item At II the original invariant lattices have bases
\[
 \begin{array}{c|c}
 1&(\psi,\delta)\\
 2&(uw,\psi\delta)\\
 3&(\psi uw,uw\delta)\\
 4&\nu.
 \end{array}                                             \tag{RI.8}
\]
The actual specialization images are, respectively,
\[
 \begin{array}{c|c|c}
 q&\text{image}&\text{quotient map}\\\hline
 1&\{a\psi+b\delta:n_Ba+d_Bb=0\bmod h\}&n_Ba+d_Bb\bmod h\\
 2&\mathbb Z(uw)+\mathbb Z(h\psi\delta)&[\psi\delta]\bmod h\\
 3&uw\{a\psi+b\delta:n_Ba+d_Bb=0\bmod h\}&n_Ba+d_Bb\bmod h\\
 4&h\mathbb Z\nu&[\nu]\bmod h .
 \end{array}                                             \tag{RI.9}
\]
These are the torus covering pullbacks in (RI.7), the actual source
class just identified, and their cup products. Every quotient is \(\mathbb Z/h\).
There is no specialization kernel.

\par

\medskip\noindent\textbf{RI.3. The retained IV-star cap maps in the joint sphere sector}

\par\nobreak

\item The actual base/fibre relations already imply that any sphere datum
must satisfy
\[
 D=h n_A+m n_B=\pm1.                                    \tag{RI.10}
\]
In particular \(\gcd(m,h)=\gcd(m,n_A)=\gcd(h,n_B)=1\). Thus every such IV-star
cap satisfies F33's freeness criterion for all \(d_A\), and its flat central
surface has torsion-free cohomology by F34.

\item Retain F34's boundary basis \((u_2,v_2,p,q)\), where
\(u_2=\psi(-w+3\delta)\), \(v_2=(\psi+u)w\), \(p=\eta_\partial\psi\),
\(q=\eta_\partial(-w+3\delta)/3\). The actual normal-circle integration map
in F34 is
\[
 u_2\mapsto3d_A\psi-n_Ah_1,\quad v_2\mapsto0,
 \quad p\mapsto m\psi+n_A\eta_{\partial},\quad
 q\mapsto k h_1+d_A\eta_{\partial},\qquad h_1=-w+3\delta.
                                                               \tag{RI.11}
\]
Its integral kernel is generated by
\(v_2\) and \(ku_2-d_Ap+n_Aq\). Hence the degree-two fibre image is exactly
\(\mathbb Zv_2+\mathbb Z(ku_2)\); its quotient is \(\mathbb Z/k\) and \(\omega=u_2+v_2\) maps to one.
The degree-one fibre image, in basis \((\psi,h_1)\), is
\(n_Aa+3d_Ab=0\) modulo m. Its quotient is \(\mathbb Z/m\).

\item For completeness put \(\phi=\psi uw\) and \(\rho=(\psi+u)w\delta\). The oriented
pairing of the full degree-one and degree-three invariant lattices is
\[
 \begin{pmatrix}0&1\\-3&0\end{pmatrix}
 \quad\text{in }(\psi,h_1)\text{ and }(\phi,\rho).
                                                               \tag{RI.12}
\]
The actual covering T4 to the cap has degree m, so its pullbacks pair
as m times the unimodular downstairs pairing. Therefore the complete
third image is exactly
\[
 \{c\phi+e\rho:n_Ac+d_Ae=0\bmod k\}.                  \tag{RI.13}
\]
Indeed an integer inverse of \(n_A\) modulo m gives the first-image basis
\((m,0),(-3d_An_A^{-1},1)\); substituting in (RI.12) gives (RI.13).
The determinant of that lattice is k, and the perfect downstairs
pairing proves there is no omitted sublattice. The degree-four image
is \(m\mathbb Z\nu\) by the actual covering degree. All \(d_A\) and normal-circle
boundary terms remain in (RI.11)--(RI.13).

\par

\medskip\noindent\textbf{RI.4. Complete global MV and differential calculation}

\par\nobreak

\item The linear tuple is exactly F26/OT.5: the reduced root construction
changes the affine filling, not its marked ordinary monodromy. Its
integral parabolic H1 vanishes and its global invariants/coinvariants
are OT.7. The only finite specialization quotients are the ones in
(RI.9),(RI.11)--(RI.13). In degrees one through four their full global
cokernel presentations are
\[
 \begin{array}{c|c}
 1&\begin{pmatrix}m&0&n_A\\0&h&n_B\end{pmatrix}\\[2pt]
 2&\begin{pmatrix}k&0&1\\0&h&3\end{pmatrix}\\[2pt]
 3&\begin{pmatrix}k&0&n_A\\0&h&n_B\end{pmatrix}\\[2pt]
 4&\begin{pmatrix}m&0&1\\0&h&1\end{pmatrix}.
 \end{array}                                             \tag{RI.14}
\]
Under (RI.10) all four have unit Smith factors. Their primitive global
kernels are generated by
\[
        1,\quad mh\psi,\quad kh\omega,
                        \quad kh\phi,\quad mh\nu.       \tag{RI.15}
\]
The second column remains the free coinvariant column OT.7 and the
middle column is zero.

\item The auxiliary linear-root fillings and the original cusp admit the
actual origin inversion, so their rational intrinsic \(d_2\) is zero by
its \((-1)^q\) action. The actual cap boundary paths, repeated m and h
respectively, close to their full integral norm vectors. Their exact
translation-prism differential is
\[
 d_2=\frac1m\iota_{n_A\ell_0+d_A\delta}
               +\frac1h\iota_{n_B\gamma+d_B\delta}.     \tag{RI.16}
\]
Apply the unchanged coinvariant projections to the four primitive
classes (RI.15). Each value is precisely \(D=h n_A+m n_B\). The \(d_A\),\(d_B\)
contractions are retained before projection and vanish in the specified
coinvariant component; they are not discarded from the cap maps.
Thus for \(D=\pm1\) the four maps are integral units. There is no
higher differential and only \(H^0=H^6=\mathbb Z\) remains.

\par

\medskip\noindent\textbf{RI.5. Full fundamental group and constructed sphere members}

\par\nobreak

\item The reduced-II filling kills the two elliptic one-cycles by its actual
source collapse; its boundary onto the group of (RI.6) retains the two
torus directions and the root meridian. Along with the selected cusp
this gives exactly the primitive fibre relation lattice generated by
\((H_B-I,H_C-I)\) in F26, leaving \(\gamma\). The two surviving base meridians
have the full relations
\[
 m\alpha_A-n_A\gamma=0,\qquad
 h\alpha_B-n_B\gamma=0,\qquad \alpha_A+\alpha_B=0.       \tag{RI.17}
\]
All generators now commute: the remaining fibre generator is invariant
and the base has one generator after the last relation. The full group
is cyclic of order \(|D|\) for D nonzero and infinite cyclic for D zero.
Thus \(D=\pm1\) is necessary and, by RI.4, sufficient for the actual
constructed manifold to be a simply connected integral homology \(S^6\).
The same Hurewicz/Whitehead and \(\Theta_6\) identification used in F26 gives
the standard smooth \(S^6\). The full group and homology are carried to every
allowed fan by the actual relative common blowups of F35.

\item These constructions exist exactly for \(\gcd(3k,6l+1)=1\) in this two-cap
sphere sector. Bezout gives every pair \((n_A,n_B)\) in (RI.10); the root
norm and retained characters above construct each one. The smallest
case is
\[
 (m,h,n_A,n_B,d_A,d_B)=(3,7,1,-2,0,0),\qquad D=1.        \tag{RI.18}
\]
All full source/line and finite-character maps remain in the proof.
Additional ordinary order-one clutches, moving sites, parameter closing
and return classes have not been assigned by this note.

\end{enumerate}
\end{proof}




\begin{theorem}\label{night:P3K.85}
\StatementRecord{night:P3K.85}{proved}{S6\_06\_BLIND: GPT-6 Astra (OpenAI, ChatGPT Pro mode, effort Pro), recorded re-derivation in the checkpoint dated 18 November 2026}{\cref{night:P3F.26}, \cref{night:P3F.29}, \cref{night:P3K.28}, \cref{night:P3K.68}, \cref{night:P3K.76}, \cref{night:P3K.77}, \cref{night:P3K.84}}
Keep the actual OS49 source, Cartier line and two caps of P3K.84, with \(m=3k\), \(h=6l+1\), \(k,l>0\), \(hn_A+mn_B=\epsilon=\pm1\), and arbitrary retained \(d_A,d_B\). For the canonical full positive \(r\)-loop \(\lambda=r^{mh}\), the positive fan \(T_1\) has physical characteristic pair and return
\[(k_M,\ell_M)=(6kh,24k(3l+1)),\qquad A_M=3kl\pmod {28}.\]
The proof retains the actual solid-torus quotients, both root actions, the corrected contracting/outward/Hopf coordinates, the complete page line and the physical dual-line gluing. The smallest tuple \((3,7,1,-2,0,0)\) gives \((42,96)\) and class \(3\).
\end{theorem}
\begin{proof}
\begin{enumerate}[label=\arabic*.]
\item \par

\medskip\noindent\textbf{RP.1. Literal contracting and outward cap phases}

\par\nobreak

\item Let \(y_B\) be the coefficient of the actual contracting B period. On
\(q_i=z_i^{r_i}\), with \(a_i=\arg z_i/(2\pi)\), the plus-log ordinary-to-cap
arrow gives \(y_{\rm cap}=y_B+n_ia_i\). The actual root cut is
\((a_i,y_{\rm cap})\mapsto (a_i+1/r_i,y_{\rm cap}+n_i/r_i)\). Thus \(u_B=r_i\) \(y_{\rm cap}\)
\(=\)\(r_i y_B+n_i\alpha_i\). The filled meridian is
\(r_i\alpha_i-n_i\gamma_B\). Put \(\gamma_{\rm out}=-y_B\) and \(u_{\rm out}=-u_B\);
then \(u_{\rm out}=r_i\gamma_{\rm out}-n_i\alpha_i\) and the root cut shifts
\(u_{\rm out}\) by \(-n_i\). This is the exact coordinate arrow to K28\textquotesingle s cap.
The positive weighted Hopf angle is therefore
\[
                              \beta=-\epsilon y_B.
\tag{RP.4}
\]
No inward cap formula is imported with the outward \(\gamma\) unchanged.

\item Set \(T=t/(t-1)\). The two root sites are \(T=0\) (order \(m,n_A)\) and
T infinity (order \(h,n_B)\). Write \(\alpha=\arg T/(2\pi)\). The exact
ordinary phases are
\(\phi_1=\epsilon(m\gamma_{\rm out}-n_A\alpha)\),
\(\phi_2=\epsilon(h\gamma_{\rm out}+n_B\alpha)\).
The integer phase map has determinant \(\epsilon\), and \(m\phi_2-h\phi_1=\alpha\).
Choose unique positive radii \(r_1^{2}+r_2^{2}=1\) with \(r_2^m/r_1^h=|T|\).
Then \(z_j=r_j e^{2\pi i\phi_j}\) gives \(T=z_2^m/z_1^h\).

\item At infinity \(q_\infty=z^h\) let \(R^{2}+|z|^{2}R^{2m/h}=1\). The actual
extending map is
\[
(z,u_\infty)\longmapsto
 (zR^{m/h}e^{2\pi i\epsilon m u_\infty/h},
                         R e^{2\pi i\epsilon u_\infty}).
\tag{RP.5}
\]
At zero \(q_0=z^m\) use \(R^{2}+|z|^{2}R^{2h/m}=1\) and
\[
(z,u_0)\longmapsto
 (R e^{2\pi i\epsilon u_0},
                  zR^{h/m}e^{2\pi i\epsilon h u_0/m}).
\tag{RP.6}
\]
The root cuts change the phases by the integers
\((\epsilon n_A,-\epsilon n_B)\) and \((-\epsilon n_A,\epsilon n_B)\),
respectively, using \(hn_A+mn_B=\epsilon\). Both maps descend through
the ACTUAL source/line root quotient. Their inverse first recovers R
from the nonzero coordinate, then an argument lift, and then z;
different lifts differ by precisely the root cut. The radial derivative
is positive. Substitution verifies both ordinary overlaps exactly.
The unit changes from \(q_i\) to T or 1/T have their specified analytic
roots and are carried by the same actual cap maps.

\par

\medskip\noindent\textbf{RP.2. Exact page, sheet-zero map and actual line}

\par\nobreak

\item For regular \(\tau\) distinct from 0, 1, infinity the complex page is
\[
\mathcal C:\ \eta^m-\tau\xi^h=1,\qquad
T=\tau+\xi^{-h}.
\tag{RP.7}
\]
Its weighted radial map is \((\xi,\eta)\mapsto (\rho^m \xi,\rho^h \eta)\),
where \(\rho^{2m}|\xi|^{2}+\rho^{2h}|\eta|^{2}=1\). It is the positive-argument
page of \(z_2^m-\tau\) \(z_1^h\), with the complex orientation. Its map has
degree M. There are h points over \(T_0\) of index m, m over infinity
of index h, and one binding point of index M. The last follows
from \(\gcd(m,h)=1\) and the expansion \(\xi\sim c w^{-m},\ \eta\sim d w^{-h}\). Thus
\(g(\overline{\mathcal C})=(m-1)(h-1)/2\). The derivatives of the defining equation cannot
both vanish on its level 1.

\item The sheet-zero map is the actual inverse of RP.5-\/-6 followed by the
specified root source map and Cartier graph. A stabilizer of a lifted
sheet has \(nk+rh=0\), hence \((k,h)=(rj,-nj)\); its action is \(A^{rj}B^{-nj}=1\)
by the literal norm. Thus no residual finite quotient remains. The
full finite characters are the prescribed ones, not replaced by an
abstract sectioned base change. At IV-star use the actual corrected
root line of F26/F29; at II use K84\textquotesingle s exact graph equality \(\beta^*N=\pi^*M\).
The II graph coordinates \(\tau\prime=z^ju_j\) carry the ratio
\((u_j+z^{l-j})/u_j\), and final \((u_l+1)/u_l\); equal exceptional
coefficients cancel. Consequently the sheet map pulls precisely
\(N=\mathcal O(G_{\rm page}-O)\), retaining every \(I_2\) component degree.

\item The relative source torsor is trivialized with its exact binding collar.
On a punctured page, the smooth Lie sheaf is fine and the exponential
obstruction is \(H^2(\mathcal C,\partial\mathcal C;j_*\Lambda)\). The two-sector cancellation
and compact-support duality identify it with monodromy coinvariants.
A lifted II root meridian has elliptic monodromy \(B_{\rm ell}^h=B_{\rm ell}\),
and \(B_{\rm ell}-I=\begin{pmatrix}0&1\\-1&-1\end{pmatrix}\) has determinant 1. Thus these integral
coinvariants vanish. Component terms are finite skyscraper sheaves
with no \(H^1\); they remain in the local diagrams. The actual torsor
translation is chosen relative to the specified binding collar.
Use the rigidified fibrewise Picard/Poincare Hom isomorphism to carry
the line under these translations. Its remaining base line and integral
whole-fibre terms have zero square and zero pairing with N. The
retained local component and finite-character maps fix all other terms.

\item The elliptic page has M unchanged \(I_2\) fibres, m reduced-II fibres,
and h good IV-star roots. Hence its Euler characteristic is
\(2M+2m=12\chi\), giving \(\chi=m(h+1)/6=k(3l+1)\). G and O were disjoint
on the original source and on the good root. At each of the m II
points the exact equality \(\tau'(G')=-z^l\) gives intersection l.
Thus \(G_{\rm page}.O=ml\) and
\[
N^2=-2\chi-2ml=-2kh,\quad
\sigma(S_{page})=-8\chi,\quad p_1(TS_{page})=-24\chi.
\tag{RP.8}
\]
These are the actual section-line numbers; the nonzero \(I_2\) component
degrees have not been replaced by a whole-fibre height.

\par

\medskip\noindent\textbf{RP.3. Full origin square and its actual forward arrow}

\par\nobreak

\item In the original \(q=1/t\) coordinate keep the complete periods
\[
\begin{gathered}u=(2s,s)^t+U(q),\quad\gamma_B=(-3s,-s+Mv)^t+C(q),\\
H=-u-\gamma_B,\quad K=3u+2\gamma_B=\mathcal D+(s+2Mv)\delta.\end{gathered}
\tag{RP.9}
\]
where \(\mathcal D=3U+2C,\ s=\log q/(2\pi i),\ v=\log r/(2\pi i)\). The valuation
matrix \(\begin{pmatrix}2&-3\\1&-1\end{pmatrix}\) has determinant 1. With the actual unit
\((1-q)^{mn_B}\) retained on both ordinary/cusp origins, their basic terms are
\[
o=-\epsilon s\mathcal D/(2M)-\epsilon s^2\delta/(4M)-\epsilon sv\delta,
\qquad\sigma=\epsilon v\mathcal D+\epsilon Mv(v-1)\delta.
\tag{RP.10}
\]
These are PLUS products: \(\xi_{\rm ord}=\zeta+o,\xi_{\rm cusp}=\zeta+\sigma\).
Therefore cusp-to-ordinary is \(o-\sigma\). If \(c=-\epsilon s\delta\) is
the ordinary closing log, then \(o(v+1)-o(v)=c\) and
\(\sigma(v+1)-\sigma(v)=c+\epsilon K\). The actual bounded forward endpoint
is \(K^{-\epsilon}\), and its full angular path is
\[
(Z,V)\mapsto(U_K^{-\epsilon\theta}Z,
 (W_K\lambda_0^2)^{-\epsilon\theta}
                      e^{-2\pi i\epsilon M\theta(\theta-1)}V).
\tag{RP.11}
\]
Every unit uses its prescribed logarithm. The whole angular difference
under \(s\mapsto s+2M\) retains its integral boundary gauge; the auxiliary source
root does not change the positive r parameter. The quadratic physical
factor is carried on the entire source, line and both cap quotients.
The old \(\sigma\)-o arrow with a forward +\(\epsilon\) K is the inverse
direction and is preserved only in the historical attempt.

\par

\medskip\noindent\textbf{RP.4. Mixed boundary map and Hopf orientation}

\par\nobreak

\item Let \((x,y_B)\) be point angles of \((u,\gamma_B)\), and (h0,\(\kappa\)) those of(H,K).
The literal torus maps are
\[
\binom{h_0}{\kappa}=\begin{pmatrix}2&-3\\1&-1\end{pmatrix}
\binom{x}{y_B},\qquad
\binom{x}{y_B}=\begin{pmatrix}-1&3\\-1&2\end{pmatrix}
\binom{h_0}{\kappa}.
\tag{RP.12}
\]
The incoming actual compression graph \(\kappa=g(h_0,w)\) moves under the
positive parameter as \(\kappa=g-\epsilon\) \(\theta\), by RP.11. Combining
\(\beta=-\epsilon\) \(y_B\) with \(\theta=-\epsilon(\kappa-g)\) gives
\[
\binom{h_0}{\theta}=\begin{pmatrix}2&3\epsilon\\-\epsilon&-1\end{pmatrix}
\binom{x}{\beta}+\binom0{\epsilon g},\qquad\det=1.
\tag{RP.13}
\]
The other source cycle w is retained. The source imaginary periods of
u and H have positive leading lengths \(2\operatorname{Im}s\text{ and }\operatorname{Im}s\), so the complex
source orientations agree on the chosen collar. Thus the actual page
boundary maps to \(E_H\times S^1_\theta\) with \(degree+1\) for both signs. With
\(\partial W=F-E_H\), the glued trace has \(\partial X=F\times S^1_\theta\).
In the fixed normal chart the circle-first orientation and the actual
Hopf coordinate below give the positive normal meridian. This is the
line \(U=M^*\), not the physical line M; K68\textquotesingle s exact dual arrow is retained.

\par

\medskip\noindent\textbf{RP.5. Actual Cox frame, full moduli homotopy and normal square}

\par\nobreak

\item The positive-fan charts are
\[
(A,B,Z)=(q/V,V/Z,Z),\quad(A',B',Z')=(ZA,ZB,Z^{-1}),\quad q=ZAB.
\]
They have weights(-1,+1,0). Write \(y=AB,b=(|B|^{2}-|A|^{2})/2\). The
quotient coordinates are (y,b). Positive K decreases \(\log|V|\), so its
height \(\chi_K\) has derivative \((\chi_K)_b<0\). Retain that reflection in
both endpoint references. The global quaternion is
\[
q_H=(1+jZ)(A+j\overline B)=(Z^{-1}+j)(A'+j\overline {B'}).
\tag{RP.14}
\]
Physical \(\delta\) is right multiplication by \(\exp(-i \phi)\). At \(Z=0\) the
Hopf coordinate \(q_Hi\overline{q_H}\) is
\(-2b i+2\operatorname{Im}(y)j-2\operatorname{Re}(y)k\).
Its determinant relative to\((\operatorname{Re}y,\operatorname{Im}y,b)\) is-8. The circle-first
quotient determinant \(at(A,B)=(1,0)\) is-1. Together these give the
actual positive normal Hopf orientation. The raw initial comparison
in \((\operatorname{Re}y,\operatorname{Im}y,-b)\) is triangular with last entry \(-(\chi_K)_b>0\);
its source shears and positive diagonal contract explicitly.
For the signed \(\beta\) coordinate apply \(T_{-\epsilon}=\operatorname{diag}(1,1,-\epsilon)\)
to BOTH initial references. Its residual comparison is the same
positive triangular matrix; conjugation uses \((-\epsilon)T_{-\epsilon}\)
in \(SO(3)\) and adds no free framing term.

\item For a bounded complex torus multiplier \((Z,V)\mapsto (aZ,b_1 V)\), retain
the exact q-preserving chart action \((Z,A,B)\mapsto (aZ,b_1^{-1} A,(b_1/a)B)\).
The quaternion identity is stated for its UNITARY phases: put
\(p=b_1^{-1},q_1=b_1/a\) and choose \(L^{2}=pq1,R^{2}=p/q_1,LR=p\) along the actual
interval. Then \(q_H\mapsto Lq_HR\) and \(\eta\mapsto L\eta\overline L\). The full positive
moduli are carried by the literal homotopy
\[
a_t=(a/|a|)|a|^t,\quad b_t=(b_1/|b_1|)|b_1|^t,\qquad
(Z,A,B)\mapsto(a_tZ,b_t^{-1}A,(b_t/a_t)B),\quad0\le t\le1.
\tag{RP.15}
\]
It preserves q and both chart transitions and is used at BOTH normal
references. At \(t=0\) the quaternion equality is exact. Higher terms of
the q-units have zero normal linear part because \(q=ZAB\) is quadratic
in the normal directions. The closed rotation is constant along F
and extends over its \(D^3\), with its actual Spin endpoint retained.

\item Take the exact regular binding coordinate
\(q_b=(T-\tau)(T/\tau)^{-\epsilon h n_A}\), derivative1 at \(\tau\). Its full
phase, from the actual polynomial page, is \(\alpha_b=\phi-M\) \(\beta\).
RP.13 gives \(\beta=2\theta+\epsilon(h_0-2g)\), so
\[
\alpha_b=\phi-2M\theta-\epsilon M(h_0-2g).
\tag{RP.16}
\]
The source term is circle-valued and retained. The actual quotient
normal map from the collar tangent(cos \(\alpha\),sin \(\alpha\),\(\rho\),0),\(\rho\)\textgreater0,
sends \((v_q,v_r,v_\eta)\) \(to(v_q-v_r(\cos \alpha,\sin \alpha)/\rho,v_\eta)\).
Projecting the two Cartesian parameter vectors gives
\[
G(\alpha,\theta)=\begin{pmatrix}
-\sin\alpha&\cos\theta\cos\alpha&\sin\theta\cos\alpha\\
\cos\alpha&\cos\theta\sin\alpha&\sin\theta\sin\alpha\\
0&\sin\theta&-\cos\theta
\end{pmatrix}\in SO(3).
\tag{RP.17}
\]
The same signed/reflected frames are used at both ends. An actual
frame on the compression transported by the K flow differs from the
Cox frame by R(p), independent of \(\theta\); its \(S^2\) boundary map extends
over \(D^3\). The compression is a solid torus minus a ball, and its normal
three-bundle has zero \(w_2\) in the stably parallelizable actual orbit \(S^5\),
so such a frame exists. The full clutch factors exactly as
\[
[R(p)^{-1}R_z(-\epsilon M(h_0-2g))]
       [R_z(\phi-2M\theta)G(0,\theta)].\tag{RP.18}
\]
The phase \(\phi\) is fixed on the actual page \(\eta^m-\tau\xi^h=1\). The first factor comes from the actual source \(T^2\), the second from
the parameter circle. Expanding the Maurer-\/-Cartan cubic of their
product gives the sum of factor cubics plus an exact form. Each
factor cubic vanishes by dimension. With the carried collar connection,
the actual normal Pontryagin pairing is therefore 0. No lower-degree
normal data are discarded.

\par

\medskip\noindent\textbf{RP.6. The full component line, regular bands and both cap markings}

\par\nobreak

\item The actual line subfan has rays \((j,\lceil j/2\rceil,1)\), with source component
degrees \((-1,+1)\). The strict contraction supplies M ordered norm shells.
After the first zero, the exact positive-iterate height defects prevent
reentry, so these shells embed. The radial annular comparison is
\(v\mapsto c(q)^{F(s)-s}v\) with \(F(s+1)=F(s)+1\). Its positive multiplier is
periodic, so it commutes with the original contracting lift and \(\delta\),
and fixes the source/base coordinates. Thus it carries the actual
circle line, including the root characters, into the compact page.

\item The graph-normal pushoff is chosen in those same base-fixed annular
and binding charts. Outside the projected regular bands its radial
homotopy avoids the entire compression. In each band the exact boundary
equation is \(\rho=-\eta_0\cot\alpha\) with \(\sin\alpha>0,\ \cos\alpha\le0\). The
three boundary arcs are slightly larger than the finite-collar
intersections, with disjoint gaps. The inner semicircle includes into
its half-disc by a homotopy equivalence. After the source-circle product,
the relative \(H^2\) of each band vanishes; restriction is injective. Hence
the normal meridian line equals the carried page line on the entire
comparison, with every nonzero \(I_2\) degree retained.

\item The fixed \(d_A,d_B\) terms are actual source-trivial finite characters.
Their full sheet products are one. Their compact extensions differ
only by an integral fibre/base class, whose square and pairing with
N vanish. This proves the same relative square \(N^{2}\) in RP.8 for every
retained character, rather than deleting its boundary gauge.

\item The incoming marked basis is \((u,\gamma_B)=(2H+K,-3H-K)\), determinant 1.
The compression cap is \((D^2_w\times S^1_H)\times S^1_K\); its \(H^2\) injects into
\(H^2(T^3)\), and the actual boundary line is flat, so its line class is 0.
The regular cap \(T^2_{u,w}\times D^2_\gamma\) has zero line class and signature.
They kill different primitive cycles; no identification of the caps
is asserted. Novikov additivity and the actual line restrictions give
the same tangent Pontryagin pairing and \(b^{2}=4N^{2}\) despite all odd gauges.

\par

\medskip\noindent\textbf{RP.7. Exact physical/Jin arrow and final numbers}

\par\nobreak

\item With \(U=M^*\) the positive Hopf meridian has fibre integral \(+1\) and the
actual page restriction \(b_U=-2c_1(N)\). Thus \(b_U^{2}=-8kh\), while
\(p_1(T\overline X)\)\(=-24\chi\) and \(p_1(\nu)\)\(=0\). The sphere-bundle pushforward is
\[
(k_U,\ell_U)=((p_1(\nu)+3b_U^2)/4, p_1(T\bar X)+p_1(\nu))
                  =(-6kh,-24\chi).
\tag{RP.19}
\]
K68\textquotesingle s actual Hermitian dual/spinor map \(Q(a,b)=(\overline b,-\overline a)\) sends
U to the physical line M and reverses both characteristic integrals.
Its piecewise gluing gives \(\lambda_M(F)=-\Sigma(F^{-1})=\Sigma(F)\), with F
the actual forward map RP.11. Hence the physical pair is
\[
(k_M,\ell_M)=(6kh,24\chi),\qquad
A_M=(4k_M-\ell_M)/24=kh-\chi=3kl.
\tag{RP.20}
\]
Both arrow reversals RP.4 and RP.11 were needed; changing just one
would falsely give determinant -1. For the old one-cap \(h=1,l=0\) case the
same maps give the retained physical pair \((2m,8m)\), not its negative.
This supports K76/K77 through the corrected actual local maps.

\item The fresh reconstructed certificate passes 1522 exact identities over
274 signed cases. The lost original run reported 2636 identities over
470 cases; those historical counts remain separate from the new
execution. The independent reconstructed page-sign certificate passes
122 additional exact checks.

\qedhere

\end{enumerate}
\end{proof}




\begin{theorem}\label{night:P3K.86}
\StatementRecord{night:P3K.86}{proved}{S6\_06\_BLIND: GPT-6 Astra (OpenAI, ChatGPT Pro mode, effort Pro), recorded re-derivation in the checkpoint dated 18 November 2026}{\cref{night:P3K.68}, \cref{night:P3K.80}, \cref{night:P3K.83}}
\par\medskip\noindent\textbf{Statement.}\par\nopagebreak[4]

Retain the actual fixed-scalar source-zero families of P3K.80 and P3K.83,
\[
 Y^2=X^3+(1-2v)X^2+TX+v^2,\quad G=(0,v),\quad
 q=T-(v^2-2v).
\]
The good branch has \((a,b,m)=(1,2,4p)\); the reduced branch has
\((a,b,m)=(2,1,4p-1)\). Retain either norm sign, the actual Cartier line, every admitted finite character, and the complete integral ordinary-to-cap cochain.

For the initial negative-square fan \(T_0\), the actual total normal clutch is
\[
 \Gamma=(w_--w_+)+c_2(E_{S^4})=0.
\]
The physical characteristic pairs are
\[
 (k_M,\ell_M)=(1,4)\quad(a=1),\qquad
 (k_M,\ell_M)=(0,0)\quad(a=2).
\]
The actual native free pairs are respectively \((1,0)\) and \((0,0)\). The physical Jin integer is zero in both cases. The proof determines the full clutch carried by the source, Cartier and Hopf cap maps.

\end{theorem}
\begin{proof}
\begin{enumerate}[label=\arabic*.]
\item \par\medskip\noindent\textbf{NC.1. The fixed trace and its two actual caps.}\par\nopagebreak[4]

\item The fixed trace is
\[
 F=\operatorname{Bl}_{(s=-1,v=0)}(\mathbb P^1_s\times D_v).
\]
Put \(c=\sqrt{1+v}\), using the branch equal to one at zero, and
\(x=s-c,\ y=s+c\). The two branch divisors are \(D_+=\{x=0\}\) and
\(D_-=\{y=0\}\); the exceptional divisor is \(E\). The full source normalization is
\[
 X=s^2-1,\qquad Y=s(s^2-1-v).
\]
The actual positive-weight normal spinor is
\[
 E_{\rm spin}=
 [M\otimes O(D_+)]\oplus[M\otimes O(-D_-)],
 \qquad M|_F=O(-aE).
\]
This last identification is the actual Cartier map retained in P3K.80/83, including the degree-two source and target \(r\)-arrows for \(a=2\).

\item Let \(Y_F=\partial F\). The actual outside cap is \(C_0=D^3\times S^1\), with the Cox normal product described below. The second comparison cap is \(C=\mathbb P^1\times D^2\), oriented so that \(\partial C=Y_F\); equivalently, \(-C\) has the complex orientation of the second parameter chart. Thus \(F-C\) is the complex blowup of \(\mathbb P^1\times\mathbb P^1\). Write \(H,B,E\) for its two factor classes and exceptional class:
\[
 H^2=B^2=0,\quad HB=1,\quad E^2=-1,\quad HE=BE=0.
\]

\item Use the actual reference divisor frames on \(C\). If
\(e_{\rm cap}=g_\pm e_F\), define \(w_\pm\) by the winding of \(g_\pm\) on the parameter circle. The full source/Cartier/root maps make these ratios global nonzero functions on \(S^2\times S^1\): at either divisor point one uses the alternate Cartier frame, and the identical point transition cancels. The capped spinor has the actual line classes
\[
 H-aE+w_+B,\qquad -H+(1-a)E+w_-B.
\]
Consequently
\[
 c_1=(1-2a)E+(w_++w_-)B,\quad
 c_1^2=-(1-2a)^2,\quad
 c_2=-a(a-1)+(w_--w_+).
\]
The two caps have closed difference
\(S^4=C-C_0\). Its actual normal bundle is glued by their full coefficient maps, rather than by a choice of its Chern classes.

\par\medskip\noindent\textbf{NC.2. The typed normal coefficient and residue arrow.}\par\nopagebreak[4]

\item Choose a local nonzero Cartier frame \(e_M\). The positive spinor frames are
\[
 e_+=x^{-1}e_M,\qquad e_-=y e_M.
\]
The negative physical normal is conjugated. In the associated node chart its actual coefficients are
\[
 U=-x\overline{\alpha},\qquad \xi=y\beta,\qquad
 U\xi=q.
\]
All full normal units of the source and Cartier maps remain in these coefficient arrows. In particular, if the associated resolution writes
\(\xi=S z\), then \(z=(y/S)\beta\); this formula is used with its actual overlap, not with a replacement of \(y/S\) by its divisor.

\item The following bundle map fixes the type of the product. There is an exact morphism
\[
 J=O(-D_+-D_-)\longrightarrow K_F,\qquad
 f\longmapsto f\,\frac{ds\wedge dv}{xy}.
\]
Its meromorphic factor has exactly the displayed polar divisor. In the infinity chart \(z_0=1/s\) it is
\[
 -\frac{dz_0\wedge dv}{1-(1+v)z_0^2}.
\]
In the blowup chart \(s=x_0-1,\ v=x_0r\), it is
\[
 \frac{dx_0\wedge dr}{x_0-2-r}.
\]
The exceptional factor cancels literally. Hence the local frame \(xy\) maps to \(ds\wedge dv\).

\item The reciprocal Hermitian map from the conjugate of \(E_+\) to \(E_+^*\) sends the conjugate of \(x^{-1}e_M\) to \(x e_M^{-1}\), up to its explicit positive norm. Tensoring with \(y e_M\) gives the frame \(xy\), with coefficient \(\overline{\alpha}\beta\). The retained minus in \(U\) therefore gives the actual negative product coefficient. This establishes the normal-to-residue arrow without losing an \(xy\), \(S\), or parameter phase.

\par\medskip\noindent\textbf{NC.3. The exact marked-point calculation and the global winding.}\par\nopagebreak[4]

\item At the regular point \(s=0\), away from the exceptional divisor and both branch divisors,
\[
 (X,Y)=(-1,0),\qquad Z=\frac{s-c}{s+c}=-1.
\]
The section \(G\) has \(s=-1\), so its literal Tate point is
\[
 P=\frac{-1-c}{-1+c}=-\frac{(1+c)^2}{v}.
\]
At \(q=Q=0\), the rigidified theta ratio is
\[
 f_P(-1)=\frac{P+1}{2P}.
\]
The full TL Cartier coefficient is
\[
 \xi=\left(\frac{1-S}{2}\right)^a\frac{V}{f_P^a},
 \qquad S=P/Z.
\]
Since \(S=-P\) at this point, this is exactly \(V=P^{-a}\xi\), with every unit of \(P\) retained.

\item In the native fixed-star chart \(q=ZAB\), the actual coefficient map is
\[
 A=-P^aU,\qquad B=P^{-a}\xi.
\]
Here \(U=c\overline{\alpha}\), \(\xi=c\beta\). The positive-weight coefficients \((\bar A,B)\) are consequently
\[
 \begin{pmatrix}
 -\bar P^{\,a}\bar c&0\\0&cP^{-a}
 \end{pmatrix}
 \binom{\alpha}{\beta}.
\]
The ratio of its two diagonal entries is
\[
 -|P|^{2a}\frac{\bar c}{c}.
\]
The first factor is positive up to the retained constant minus sign, while \(c\) has its extending disc logarithm. Thus the two coefficient phase windings are equal. Each is \(+a\) in this displayed normalization because \(P\) has winding \(-1\). With the frame convention in NC.1 the coefficient arrow is inverted; both frame windings are correspondingly equal as well. In particular,
\[
 w_--w_+=0.
\]
This equality holds for the actual global ratios: their winding is independent of the regular sphere point, since they are global nonzero functions on the connected product. Any further retained common physical scalar multiplies both coefficient phases by the same factor and is carried in the next cap map.

\par\medskip\noindent\textbf{NC.4. The actual Cox/Hopf product on \(C_0\).}\par\nopagebreak[4]

\item On the two-cone fixed star use homogeneous unit Cox factors
\[
 u=a_0+ja_1,\qquad v_N=A+j\bar B,\qquad q_H=u v_N,
 \qquad h=u i u^{-1}.
\]
The auxiliary circle cancels. Choosing a Hopf representative of \(h\) gives the inverse \(v_N=u^{-1}q_H\). In the finite \(Z\)-chart the coefficient matrix is
\[
 \mathcal H_Z=\frac1{\sqrt{1+|Z|^2}}
 \begin{pmatrix}1&-\bar Z\\ Z&1\end{pmatrix},
 \qquad q_H=\mathcal H_Z(A,\bar B)^t .
\]
The physical action is the negative right phase
\(q_H\mapsto q_H e^{-i\phi}\). Therefore the positive cap spinor has coefficients \(\overline{q_H}\), and its actual coefficient arrow is
\(\overline{\mathcal H_Z}(\bar A,B)^t\). This is the full anti-linear physical-normal comparison of P3K.68, followed by the displayed complex-linear map to the positive cap spinor.

\item The sphere coordinate change is the literal map
\[
 \mathcal B_c(P_0)=
 \frac{A_cP_0A_c^*}{\operatorname{tr}(A_cP_0A_c^*)},
 \qquad A_c=\begin{pmatrix}1&-c\\1&c\end{pmatrix},
\]
on trace-one positive Hermitian matrices. Its inverse uses \(A_c^{-1}\); its boundary is \(Z=(s-c)/(s+c)\). Holding the actual \(q_H\) coordinate fixed gives the cap map
\((P_0,q_H,v)\mapsto(\mathcal B_c(P_0),q_H,v)\).
The tangent derivative of the sphere map is not substituted for this normal product.

\item Remove the common phase found in NC.3 by the corresponding simultaneous rotation of both positive coefficient planes on \(C_0=D^3\times S^1\). It extends because it depends on the retained parameter circle. For an integer winding \(d\), its total Spin lift is the product of the two plane lifts; their endpoints are both \((-1)^d\), so the total endpoint is \(+1\). This is a periodic Spin lift of the actual common rotation. The physical circle and its line commute with this coefficient map.

\item The two remaining coefficient ratios have zero winding. They therefore have global smooth logarithms on \(S^2\times S^1\). Retain their exact boundary values, extend those logarithms over the comparison cap using a collar cutoff, and use the exponential frame changes. This carries all original normal units and finite characters instead of replacing them by their winding numbers.

\par\medskip\noindent\textbf{NC.5. Extension on \(C\) and the \(S^4\) clutch.}\par\nopagebreak[4]

\item After those extending scalar maps, the remaining matrix is precisely the conjugate Hopf product from the two branch coordinates. In homogeneous point frames it can be written
\[
 \mathcal G(x,y)=
 \frac1{\sqrt{|x|^2+|y|^2}}
 \begin{pmatrix}-\bar y&-x\\-\bar x&y\end{pmatrix}.
\]
Its determinant is \(-1\). The two columns have the opposite point-frame transitions of \(O(H)\) and \(O(-H)\); thus this formula is a bundle map in both sphere charts. The nonzero positive metric factors in the actual Cox coefficients are retained by their positive normal rescaling. At \(s=0\) it gives exactly the matrix of NC.3 after the two displayed Cartier factors. This is the coefficient computation established in NC.2, not an identification inferred from equal Chern classes.

\item On the second cap use its actual boundary parameter, rescaled so that its boundary is the original parameter circle. In that parameter \(w\), the smooth extension is
\[
 c_C(w)=\sqrt{1+\bar w}
\]
with the same small-disc branch. Its boundary value is the original \(c(v)\) under the specified reversed cap parameter. Put \(x_C=s-c_C(w)\), \(y_C=s+c_C(w)\) in the displayed matrix and in \(\mathcal B_c\). The determinant \(2c_C(w)\) of \(A_{c_C}\) never vanishes; \(x_C,y_C\) never vanish together. Hence both sphere charts and the entire three-ball product extend, with the actual \(q_H\) held fixed. The original radius is retained by the corresponding constant rescaling of \(w\).

\item Combine this extension with the zero-period logarithm extensions of NC.4 and its common physical phase on \(C_0\). On their common boundary the resulting maps are the original full source/Cartier/Hopf maps, including all point transitions. They provide a global complex frame on the normal bundle over \(C-C_0\). Its actual physical associated line is the tautological \(O(-1)\), carried by the same positive spinor coefficients. Consequently
\[
 c_2(E_{S^4})=0.
\]
Together with NC.3 this proves \(\Gamma=0\).

\par\medskip\noindent\textbf{NC.6. The two complete characteristic pairs and the actual return.}\par\nopagebreak[4]

\item For the projectivized normal sphere bundle over \(F-C\), put \(h=c_1(O(1))\). The physical line is \(x=-h\), with
\[
 h^2+c_1h+c_2=0,\quad \pi_*h=1,\quad
 \pi_*h^2=-c_1,\quad \pi_*h^3=c_1^2-c_2 .
\]
The exact tangent splitting gives
\[
 k_B=c_2-c_1^2,\qquad
 \ell_B=-p_1(F-C)-c_1^2+4c_2.
\]
Since \(p_1(F-C)=-3\), these are
\[
 k_B=3a^2-3a+1+(w_--w_+),\qquad
 \ell_B=4+4(w_--w_+).
\]
The \(S^4\) correction contributes
\((c_2(E_{S^4}),4c_2(E_{S^4}))\).

\item The complete reduced finite attachment of P3K.80 has one charge-two boundary and one isolated charge-one boundary, with its proved orientations. The physical projective pair is
\((-r^3,-4r)\) for charge \(r\). Its actual boundary identity therefore gives
\[
 (k_M,\ell_M)=(-\Gamma,-4\Gamma)\qquad(a=2).
\]
The complete good two-step attachment of P3K.83 has two distinct charge-one primitive faces, the whole fixed surface only once, and the free upper trace with its trivial physical line. Its actual boundary identity gives
\[
 (k_M,\ell_M)=(1-\Gamma,4-4\Gamma)\qquad(a=1).
\]
No local fixed-\(S^3\) handle is inserted a second time. Applying \(\Gamma=0\) proves the statement.

\item Choose the delta-averaged connection product on the selected tube in the actual coordinates of NC.4. The induced return on the conifold quotient preserves the physical Euler generator and the orientation, hence both integral homology generators. Its Hopf lift on the complement and its identity extension on the selected tube are exactly those used in P3K.68. The actual directed transfer is
\(\lambda_M(f)=-\Sigma(f^{-1})=\Sigma(f)\).
Thus the above physical pair belongs to the specified source loop and its native free pair is \((k_M,0)\). The integer \((4k_M-\ell_M)/24\) vanishes.

\par\medskip\noindent\textbf{MD.1. The full initial three-curve data.}\par\nopagebreak[4]

\item For completeness the two subsequent native fan paths are calculated from the actual doubled source products. Let
\[
 Q_0=\begin{pmatrix}1&-1\\0&1\end{pmatrix},\qquad
 Q_1=\begin{pmatrix}1&0\\-1&1\end{pmatrix},\qquad
 W=\begin{pmatrix}1&b\\-a&-2\end{pmatrix}.
\]
The four actual period units have winding matrix \(W\), and their common product trace is \(t=-1\). For a basis \(P\), retain every coefficient through the integral map
\[
 U'_{is}=\prod_{h,j}U_{hj}^{(P^{-1})_{ih}P_{js}}.
\]
It follows directly that their doubled winding matrix is \(B=2P^{-1}WP\). The point angles and their physical orientations give
\[
 l=(B_{22}-B_{21},B_{11}-B_{12},B_{22}),\quad
 (x,y)=P^{-1}e_2,\quad w=(x,y,-x-y),
\]
\[
 d=(-x-2y,-2x-y,x-y),\quad
 I_i=-w_i+d_i(l_i+1),\quad J_i=-w_i-d_i(l_i+1).
\]
These are the actual circle and stable-derivative trace equations. For the symmetric variation \(V\) and derivative \(\delta\), the full integral prisms are
\[
 V(w_i e_i-w_0e_0)=I_i e_i-I_0e_0,\qquad
 w_i\delta_i-w_0\delta_0=J_i-J_0,\quad i=1,2.
\]
The initial native fixed pair is \((\mu,0)\). Its actual inverse cylinder contributes left degree \(-2t=2\); hence its doubled ordinary data have \(V_{00}=\delta_0=-2-2\mu\). Solving the full prism gives
\[
 V=\begin{pmatrix}
 -2-2\mu&4a-6&6-4a\\
 4a-6&2-2a&0\\6-4a&0&-4
 \end{pmatrix},\qquad
 \delta=(-2-2\mu,\,2-2a,\,4a-4).
\]

\par\medskip\noindent\textbf{MD.2. The actual native map and its selected reset.}\par\nopagebreak[4]

\item The dense \(Q_0\) character map is \((Z,V)\mapsto(ZV,V)\), with periods \((H,H^{-1}B)\); the \(Q_1\) map is \((Z,V)\mapsto(Z,ZV)\), with periods \((HB^{-1},B)\). The full coefficient formula above supplies their unit maps and inverse maps, preserving every \(q\)-term and endpoint seam.

\item In a selected native chart \((z,n_0,n_1)\), the two resolutions have overlap
\(z'=z^{-1}\), \(n'=zn\), with common rank-one matrix
\[
 \begin{pmatrix}n_0&n_1\\zn_0&zn_1\end{pmatrix},
 \qquad q=zn_0n_1.
\]
The actual bounded torus lift is
\[
 (a_0,a_1,c_0,c_1)\mapsto
 (a_0,Aa_1,(AD)^{-1}c_0,Dc_1).
\]
It acts on both rulings and on their full normal germs. The \(C_1\) ordering interchange is the unitary path
\[
 \frac12
 \begin{pmatrix}1+e^{\pi is}&1-e^{\pi is}\\
                 1-e^{\pi is}&1+e^{\pi is}\end{pmatrix},
 \quad 0\le s\le1 ,
\]
which commutes with the scalar overlap \(zI\). Transporting the link metric and its positive radial normalization keeps the full normal endpoint maps.

\item The common link and oriented exterior maps are
\[
 (h,q_H)\mapsto(q_H^{-1}hq_H,q_H),\qquad
 (y,q_H)\mapsto(q_H^{-1}yq_H,q_H).
\]
The actual complex-meridian reset gives
\[
 (\mu_i,\nu_i,s)\mapsto(\mu_i+2\nu_i,-\nu_i,s+\nu_i).
\]
In doubled ordinary coordinates \(c=V_{ii}\), \(e=\delta_i\), this sends the selected new \(C_2\) data to
\[
 V'_{22}=c,\qquad \delta'_2=2c-e,\qquad
 s'=s+\frac{e-c}{4}.
\]
The quarter factor is derived from the actual doubled native products and their full source seam, not from division of a sphere class.

\item The common exterior inverse character map is the identity in the original source and Cartier coordinates; it retains the finite-root maps. The two projective ruling charts replace exactly the selected tube. The actual common blowup and its integral Gysin maps retain sphere cohomology, while the same primitive period relations retain the fundamental group. Thus both finite paths below are families of the actual completed spheres.

\par\medskip\noindent\textbf{MD.3. Fresh complete target systems.}\par\nopagebreak[4]

\item The first \(Q_0\) step has \(c=e=-2-2\mu\), and hence central increment zero. Its target weights are \((1,1,-2)\). All six symmetric variation and three derivative unknowns are fixed by the full prism and the two selected equations.

\item For the good branch \(a=1\), the target data are
\[
 V'=\begin{pmatrix}
 -8\mu-8&-8\mu-4&4\mu+2\\
 -8\mu-4&-8\mu-8&4\mu+2\\
 4\mu+2&4\mu+2&-2\mu-2
 \end{pmatrix},\quad
 \delta'=(4\mu+4,4\mu+4,-2\mu-2).
\]
For the reduced branch \(a=2\), they are
\[
 V'=\begin{pmatrix}
 -8\mu-20&-8\mu-4&4\mu+2\\
 -8\mu-4&-8\mu-2&4\mu+2\\
 4\mu+2&4\mu+2&-2\mu-2
 \end{pmatrix},\quad
 \delta'=(4\mu+16,4\mu-2,-2\mu-2).
\]
Direct substitution verifies every full prism and both selected equations. The second \(Q_0\) and \(Q_1\) increments are therefore respectively
\[
 (3\mu+3,3\mu+3)\quad(a=1),\qquad
 (3\mu+9,3\mu)\quad(a=2).
\]

\par\medskip\noindent\textbf{MD.4. The two actual nonvertical returns.}\par\nopagebreak[4]

\item NC.6 gives \(\mu=1\) in the good branch and \(\mu=0\) in the reduced branch. Both paths start at central class zero. Thus
\[
 \bigl([\Sigma_{Q_0^2}],[\Sigma_{Q_0Q_1}]\bigr)
 =(6,6)\quad(a=1),\qquad (9,0)\quad(a=2)
 \quad\text{in }\Theta_7.
\]
Every matrix here comes from the actual unit, normal, cap and selected-reset arrows above. This conclusion treats these two finite words and the initial fixed-scalar seed.

\end{enumerate}
\end{proof}




\begin{theorem}\label{night:P3K.87}
\StatementRecord{night:P3K.87}{proved}{S6\_06\_BLIND: GPT-6 Astra (OpenAI, ChatGPT Pro mode, effort Pro), recorded re-derivation in the checkpoint dated 18 November 2026}{\cref{night:P3F.35}, \cref{night:P3K.68}, \cref{night:P3K.70}, \cref{night:P3K.71}, \cref{night:P3K.76}, \cref{night:P3K.84}, \cref{night:P3K.85}}
For the named two-cap OS49 sphere data of P3K.84--85, set \(M=mh\), \(\chi=k(3l+1)\), \(A_0=3kl\). On the vertical fan \(T_n\), the canonical full \(r\)-loop has
\[(k_{M,n},\ell_{M,n})=(2M-6M(n-1)^2,24\chi),\qquad [\Phi_{T_n}]=A_0-M(n-1)^2\pmod {28}.\]
The exact full-unit Cox algorithm in the proof computes its physical return on every lattice-basis cusp fan. The smallest source/cap tuple \((3,7,1,-2,0,0)\) realizes every element of \(\Theta_7\) through the listed explicit fans. The statement retains the canonical full \(r\)-loop; its actual shorter endpoints are treated in P3K.91.
\end{theorem}
\begin{proof}
\begin{enumerate}[label=\arabic*.]
\item \par

\medskip\noindent\textbf{RIF.1. Named family and actual closed cusp piece}

\par\nobreak

\item The exact source and line are
\[
y^2=x^3-3t^3(t-1)x+2t^4(t-1)(t-\tfrac12),\quad
G=(t^2,t^2),\quad\mathcal M=O(G-O),\quad Q=3G.
\]
Use \(\lambda=r^M\) and actual cap actions
\(A_{A,r}=r^{hn_A}A_{A,0}\), \(A_{B,r}=r^{mn_B}A_{B,0}\).
Their powers are \(B^{n_A},B^{n_B}\). The actual closing is
\[
f=t^{hn_A}(t-1)^{mn_B}=q^{-\epsilon}(1-q)^{mn_B},\quad q=1/t.
\tag{RIF.2}
\]
The forward map is \(q^\epsilon u(q)\), \(u=(1-q)^{-mn_B}\), with its full
extending logarithm. Keep periods
\[
u_0=(2s,s)^t+U(q),\quad\gamma_B=(-3s,-s+Mv)^t+C(q),
\quad H=-u_0-\gamma_B,\quad K=3u_0+2\gamma_B.
\]
Here \(s=\log q/(2\pi i),v=\log r/(2\pi i)\). In physical characters
\[
H:(Z,V)\mapsto(\kappa qZ,\alpha V),\qquad
K:(Z,V)\mapsto(gZ,q\eta V),
\tag{RIF.3}
\]
with windings \((\kappa,g,\alpha,\eta)=(0,0,-M,2M)\). All zero-winding
factors are the actual holomorphic q-units. Reindex the forward map
by \(K^{-\epsilon}\); its bounded endpoint is \(E=(g^{-\epsilon},u \eta^{-\epsilon})\).
For an endpoint-flat cut \(a(\tau)\), the full multiplier
\[
C_\tau=(\exp[\epsilon a\log g],\exp[a(-\log u+\epsilon\log\eta)])
\tag{RIF.4}
\]
satisfies \(C_1E=C_0=1\) on entire collars with all derivatives. It commutes
with every period and \(\delta\). On the two triple charts it acts
\[
(a,b,c)\mapsto((AD)^{-1}a,Ab,Dc),\quad
(u_1,v_1,w_1)\mapsto(ADu_1,D^{-1}v_1,A^{-1}w_1).
\tag{RIF.5}
\]
Both products remain q. These are the actual component source and
line of K76. Inverting the character change on the common exterior
is the identity source/Cartier map; at both finite caps it carries
the full K84 annular vector and hence the displayed root action,
including every d character. K85\textquotesingle s corrected plus-product square
has cusp-to-ordinary \(o-\sigma\) and forward \(-\epsilon K\). Its quadratic
factor \(e^{-2\pi i\epsilon M\theta(\theta-1)}\) is the actual global
physical action, carried on all charts and normal references. On a
double parameter circle define the product by literal pullback:
on lap j its seam is \(C_\tau\) \(E^{-j}\), with actual parameter arguments.
All free-coordinate doubling refers to this same product.

\par

\medskip\noindent\textbf{RIF.2. Full vertical line and normal comparison}

\par\nobreak

\item The positive K85 fan is \(T_1\). Its actual fixed-star charts are
\[
(A,B,Z)=(q/V,V/Z,Z),\quad(A',B',Z')=(ZA,ZB,Z^{-1}),\quad q=ZAB.
\]
Their weights are (-1,+1,0) and \(q_H=(1+jZ)(A+j\overline B)\). The shear
\(V_n=Z^{-n}V,H_n=HK^n\) gives the exact real-period coordinates
\(u_n=u,\ v_n=v,\ \gamma_n=\gamma-nv,\ d_n=d-nu\), with all unit phases retained.
For \(U=M^*\) the actual inverse-line connections for bottom windings
\(L=-M,J=2M\) are
\[
\begin{gathered}A_{L,J}=dd+(\alpha_q+J\theta)d\gamma+L\theta dv,\\
A_n=dd_n+(\alpha_q+J\theta)d\gamma_n+(L+nJ)\theta dv,\\
\beta_n=A_n-A_{ext}=-n(du+\alpha_qdv).\end{gathered}
\tag{RIF.6}
\]
The \(\theta\)-cut lines \(T_b,T_\gamma\), glued by \(e^{2\pi ib},e^{2\pi i\gamma}\),
extend over the fixed-star cap using supported real lifts on its \(S^2\).
Their curvatures \(\Xi=d\theta\wedge dv,\ \Psi=d\theta\wedge d\gamma\) obey
\(\Xi^2=\Psi^2=\Xi\Psi=0\). The actual coefficient clutch gives
\(U_{L,J}=U_{0,1}\otimes T_b^L\otimes T_\gamma^{J-1}\). The sole unextended
boundary winding is the angle \(\alpha_q\); it is zero on the actual free
disc \(|Z_*|=1,V=1,\gamma=v=0\). The other b,gamma,\(\theta\) gauges extend,
and Cech partition extends the residual logarithm. The surgery-cap
arrow is the same actual Cartier/Hopf adjoint map as K85.

\item The BW calibration uses the SAME closing \(e=\epsilon\) and SAME u(q):
\(E_s=(g_s^{-e},u_s \eta_s^{-e})\), \((C_s)_1=(g_s^e,u_s^{-1} \eta_s^e)\).
It changes only protected local factors and fixes the outer real
coefficients. At its endpoint all zero-winding factors are q-only.
The actual C2 Cox bordism has first outer map identity and second
base \(\kappa_*^{-1}\), normal \(\operatorname{diag}(\alpha_*^{-1},\kappa_*\alpha_*)\), both parameter-
independent. Its common inner physical gauge extends through the
joint cap. Thus both integer characteristic corrections vanish,
giving integral \(F_{0,1}^2\Xi=1\) and hence \(F_{L,J}^2\Xi=1\). No global
finite-cap datum is varied in this relative calibration.

\item The exact tangent boundary matrix is \(I+n\phi_gE_{12}-nE_{32}\). Its Maurer
form squares to zero. Make the actual regular-level b-map constant
on a tube containing all surgery and connection changes by its real
lift, relative to the outer boundary. The modified level misses the
cap. Its cobordism to the original source-product level proves
\(\langle p_1,\Xi\rangle=0\) with zero collar transgression. Thus \(\ell_M\) is unchanged.
For the cube collar,
\[
F_{ext,n}=(d\alpha_q+Jd\theta)\wedge(d\gamma+n dv)+L d\theta\wedge dv,
\quad d\beta_n=-n d\alpha_q\wedge dv.
\]
In orientation \(\Omega=d\alpha_q\wedge du\wedge dv\wedge d\gamma\wedge d\theta\),
\[
3\int_0^1\beta_n(F_{ext,n}+s d\beta_n)^2ds=(6Ln+3Jn^2)\Omega.
\]
The actual outward collar runs 1 to 0, so contributes \(-6Ln-3Jn^{2}\);
the relative interior contributes \(3nJ\). Carrying K68\textquotesingle s physical arrow gives
\(k_{M,n}-k_{M,0}=-6Ln-3Jn(n-1)\), \(\ell_{M,n}=\ell_{M,0}\).
The common global exterior contributes the identical form with the
same root/line/annular identifications RIF.2-\/-5. K85 gives reference
\((2M,24\chi)\) at \(n=1\). Subtraction proves RIF.1 and the actual native pair
\[
                       (\mu_n,\rho_n)=(6\chi,2M-6\chi-6M(n-1)^2).
\tag{RIF.7}
\]
The left entry is computed on the two-cap page; it is not replaced by 2M.

\par

\medskip\noindent\textbf{RIF.3. Complete doubled ordinary seed}

\par\nobreak

\item The native C0 full multiplier lifts as
\((a_0,a_1,c_0,c_1)\mapsto(a_0,Aa_1,(AD)^{-1}c_0,Dc_1)\), preserving q and both
resolutions. On the literal double product its Hopf arrow is
\[
q_H\mapsto(UH_{2M}U^{-1})q_H H_{\lambda_i-2M}.
\tag{RIF.8}
\]
Here \(H_d\) is GF\textquotesingle s endpoint-compatible normal contraction and \(\lambda_i\)
is the full doubled-period coefficient. K70\textquotesingle s explicit polar-disc
degree gives left degree \(4M\), right degree \(0\), with inverse-cylinder endpoint \(L^{-4M}\).
All zero-winding factors are retained by protected logarithms/radial
homotopies; RIF.4 carries the same endpoint on both references.
Thus the doubled ordinary pair is
\((12\chi-4M,4M-12\chi-12M(n-1)^{2})\). Since ordinary variation and stable
derivative are \((-p-q,-p+q)\),
\(V_{00}=12M(n-1)^{2},\delta_0=24A_0-12M(n-1)^{2}\).

\item The full doubled matrix is \(B_n=2M\begin{pmatrix}0&0\\2n-1&2\end{pmatrix}\), with \(t_0=2M\),
\(w=(0,1,-1),d=(-2,-1,1)\). The literal circle and derivative prisms
give \(I=(8M(n-1),4M,-4M),W=(8M(1-n),0,0)\), so their integral equations
fix every remaining entry:
\[
V=\begin{pmatrix}12M(n-1)^2&8M(1-n)&8M(n-1)\\
8M(1-n)&4M&0\\8M(n-1)&0&4M\end{pmatrix},\qquad
\delta=\begin{pmatrix}24A_0-12M(n-1)^2\\8M(n-1)\\8M(1-n)\end{pmatrix}.
\tag{RIF.9}
\]
The single initial integer is \(S=A_0-M(n-1)^{2}\).

\par

\medskip\noindent\textbf{RIF.4. Every fan and the unique target reconstruction}

\par\nobreak

\item Use K71\textquotesingle s exact identical-fan maps and terminating Euclidean word.
Column swap and \(I,R,R^{2},-I,-R,-R^{2}\), \(R=\begin{pmatrix}-1&-1\\1&0\end{pmatrix}\), identify
the complete infinite triangles by integral translations. The dense
torus identity extends on identical semigroup rings, with original
periods/caps/line. The representative is vertical or
\(P=\begin{pmatrix}p&-r\\q&s\end{pmatrix},\ p,r>0,\ ps+rq=1\). Parent moves by \(Q_0^{-1}\) when \(r\ge p\)
and \(Q_1^{-1}\) when p\textgreater r decrease \(p+r\), with
\[
Q_0=\begin{pmatrix}1&-1\\0&1\end{pmatrix},\quad
Q_1=\begin{pmatrix}1&0\\-1&1\end{pmatrix}.
\tag{RIF.10}
\]
Reversing gives only C0/C1 steps and all nonvertical target weights
are nonzero. The full character map is \((Z,V)\mapsto (Z,V)^{P^{-1}}\), with
period generators given by columns P. For full multiplier \(a_{\rm st}\),
the transformed entry is \(\prod_{s,t}a_{st}^{(P^{-1})_{is}P_{tj}}\).
The valuation identity remains, so residual factors stay holomorphic
q-units. This carries the entire endpoint, plus-origin term and
both root overlaps. Composing its inverse on the exterior is the
actual source/Cartier identity; no finite cap is changed.

\item The native other-ruling charts are \((n_0/n_1,zn_1,n_1)\) and
\((n_1/n_0,zn_0,n_0)\), preserving \(q=zn_0n_1\). The C1 normal interchange
is the actual K71 unitary path with unrestricted four Cox weights.
Both rulings receive the full multipliers, source/line/cap and normal
maps. The common blowup supplies F35\textquotesingle s integral homology transfer,
so each step constructs the compact sphere family.

\item On the actual double parameter circle and target P\textquotesingle{} put
\[
B=2M(P')^{-1}\begin{pmatrix}0&0\\-1&2\end{pmatrix}P',\quad
t_0=2M,\quad\lambda=(b_{22}-b_{21},b_{11}-b_{12},b_{22}),
\]
\[
\begin{gathered}(x,y)^t=(P')^{-1}e_2,\quad w=(x,y,-x-y),\\
d=(-x-2y,-2x-y,x-y),\\
I_i=t_0w_i+d_i(\lambda_i-t_0),\quad W_i=t_0w_i-d_i(\lambda_i-t_0).\end{gathered}
\tag{RIF.11}
\]
The transformed zero-winding logarithms enter the actual GF normal
contractions. Their exponentials and protected normal rescaling supply
the full circle/derivative prisms, whose degrees are exactly I,W.
At selected i retain \(a=V_{ii},\ b=\delta_i\). Its actual single native pair
satisfies \(a=4M-2\mu_i-2\rho_i,b=4M-2\mu_i+2\rho_i\). Hence \(\rho_i=(b-a)/4\).
The unrestricted four-weight Cox link and complex reset send
\((\mu,\rho,S)\mapsto (\mu+2\rho,-\rho,S+\rho)\). Target selected curve is C2, so
\[
V'_{22}=a,\quad\delta'_2=2a-b,\quad S'=S+(b-a)/4.\tag{RIF.12}
\]
This divides free coordinates of the actual doubled product, not \(\Theta_7\).
The complete target equations are
\(V\prime(w_je_j-w_0e_0)=I_je_j-I_0e_0\) and
\(w_j\delta\prime_j-w_0\delta\prime_0=W_j-W_0\) \(\text{for }j=1,2\).
Their symmetric homogeneous solution is \(C(1/w_i)(1/w_j)^t\) and their
derivative homogeneous solution is \(E/w_i\). The selected entries fix
\[
C=w_2^2a-w_2I_2,\quad E=w_2(2a-b)-W_2,\quad
V'_{ij}=(I_i/w_i)\delta_{ij}+C/(w_iw_j),\quad
\delta'_i=(W_i+E)/w_i.\tag{RIF.13}
\]
The constructed maps supply an integral solution, hence this unique
rational expression is integral. Induction proves every computed S.

\par

\medskip\noindent\textbf{RIF.5. Exact witnesses and independent execution}

\par\nobreak

\item The reconstructed certificate runs 6504 unimodular inputs bounded 18,
representing 814 fans, for \((k,l)=(1,1),(1,2),(2,1),(5,1),(7,2):32520\)
executions. Separate three-edge traversals include actual inverse steps;
each reaches 439 states to depth 8. Separate fixed target stars are left
to the proved vertical seed. It checks vertical n from -30 through 30.
For \(k=l=1\), \(P=\begin{pmatrix}2&-1\\-1&1\end{pmatrix}\) has seed \(n=0\)\(,S=-18\) and word (0,1),
selected doubled pairs (252,-180),(1512,-1872), increments -108,-846.
Thus \(S=-972=8\) modulo 28, with
\[
V=\begin{pmatrix}9912&6300&-4200\\6300&3864&-2100\\-4200&-2100&1512\end{pmatrix},
\quad\delta=(-12000,-8016,4896)^t.\tag{RIF.14}
\]
The output retains one actual full basis/word and every intermediate
invariant for each of 28 residues. Fresh SymPy linsolve builds all
target equations independently of RIF.13 at every witness step.
The fresh reconstructed run passes 116 such solves and 26016 full integer
homogeneity checks. Its complete executed coverage is 32520 unimodular
inputs, with all intermediate witness data saved in the certificate JSON.
The earlier run reported the same counts; its historical outputs remain
separate from this fresh execution.
These 28 explicit witnesses establish surjectivity for the fixed tuple
(3,7,1,-2,0,0). The unrestricted domain follows from the proved actual
Euclidean/Cox construction, not from bounded enumeration.

\end{enumerate}
\end{proof}




\begin{theorem}\label{night:P3K.88}
\StatementRecord{night:P3K.88}{proved}{S6\_06\_BLIND: GPT-6 Astra (OpenAI, ChatGPT Pro mode, effort Pro), recorded re-derivation in the checkpoint dated 18 November 2026}{\cref{night:P3F.35}, \cref{night:P3K.68}, \cref{night:P3K.80}, \cref{night:P3K.83}, \cref{night:P3K.86}, \cref{night:P3K.90}}
For the actual E7 source-zero loop with fixed scalar, both marked
branches \((a,b)=(1,2),(2,1)\), every retained cap datum and every fan
obtained from \(T_0\) by a word beginning \(Q_0\) and continuing in
\(Q_0,Q_1\), equations (ZB.5)--(ZB.10) compute its exact physical
return in the K68 convention. Each branch attains every class of
\(\Theta_7=\mathbb Z/28\); the two tables below give all 28 actual
fan witnesses for each branch. The initial free normal correction
is the complete \(\Gamma=0\) of K86.
\end{theorem}
\begin{proof}
\begin{enumerate}[label=\arabic*.]
\item \par

\medskip\noindent\textbf{ZB.1. The specified family and its initial integer pair}

\par\nopagebreak[4]

\item We retain the actual source
\[
Y^2=X^3+(1-2v)X^2+TX+v^2,\qquad G=(0,v),\qquad
q=T-(v^2-2v).
\]
The positive source loop is a small counterclockwise circle about
\(v=0\). Its scalar is fixed in the compact admissible annulus. The
marked datum is \((a,b)=(1,2)\), with good cap order \(4p\), or
\((a,b)=(2,1)\), with reduced cap order \(4p-1\), \(p\ge1\).
Both cap norm signs, every finite character and the full annular cochain
are retained. The orientation and physical line \(M\) are those of K68.

\item K80 and K83 identify the actual initial return with physical integer
zero. K86 computes the complete relative normal correction over the
horizontal normalization, the removed four-ball and their common
boundary. It gives \(\Gamma=0\), including both full Cartier factors.
Consequently its initial physical pair is \((1,4)\) for \(a=1\) and
\((0,0)\) for \(a=2\). Write \(\mu=2-a\). The initial native free
coordinates are \((\mu,0)\); these are integers supplied by the full
normal comparison. They are the inputs to the computation below.

\par

\medskip\noindent\textbf{ZB.2. The full source, period and flop arrows}

\par\nopagebreak[4]

\item The actual marked periods have the form
\[
h=(\kappa qZ,\alpha V),\qquad B=(gZ,q\eta V),
\qquad
\operatorname{wind}_v(\kappa,\alpha,g,\eta)=(1,-a,b,-2).
\]
All four functions are the original nonzero holomorphic units. In
particular their bounded dependence on \(q\) is retained. Their winding
matrix is
\[
W=\begin{pmatrix}1&b\\-a&-2\end{pmatrix}.
\]
For a basis change \(Q\in SL_2(\mathbb Z)\), character substitution
and the inverse period basis give the complete multiplier formula
\[
U'_{is}=\prod_{h,j}U_{hj}^{(Q^{-1})_{ih}Q_{js}}.             \tag{ZB.1}
\]
This equality holds for the units themselves. Its valuation part leaves
the \(q\)-valuation matrix equal to the identity.

\item Put
\[
Q_0=\begin{pmatrix}1&-1\\0&1\end{pmatrix},\qquad
Q_1=\begin{pmatrix}1&0\\-1&1\end{pmatrix}.
\]
We start with \(T_0\), take \(Q_0\) first, and then use any finite
word in these two matrices. In a native source chart
\((z,n_0,n_1)\), the two other-ruling charts are
\[
(n_0/n_1,zn_1,n_1),\qquad(n_1/n_0,zn_0,n_0).              \tag{ZB.2}
\]
They have the same rank-one matrix and the same function
\(q=zn_0n_1\). The second type also interchanges the two normals by
the actual unitary path
\[
\frac12\begin{pmatrix}
1+e^{\pi is}&1-e^{\pi is}\\1-e^{\pi is}&1+e^{\pi is}
\end{pmatrix},\qquad 0\le s\le1.                         \tag{ZB.3}
\]
This path commutes with the scalar normal transition \(zI\).
For full multipliers \((A,D)\), the Cox lift is
\[
(a_0,a_1,c_0,c_1)\longmapsto
(a_0,Aa_1,(AD)^{-1}c_0,Dc_1).                            \tag{ZB.4}
\]
The quotient is the required holomorphic change on each stable chart.
On its exterior it is the same inverse character substitution as
(ZB.1). Thus the actual Cartier line, all finite roots and characters,
and the prescribed closing arrow commute with the comparison there.
The source and finite cap maps are unchanged on that common exterior.

\item We transport the Hermitian link metric by these actual chart maps and
normalize its positive radius. On a protected normal tube the rescaling
\(R_s(z,n)=(z,sn)\), followed by the cutoff flow of its velocity,
carries every higher normal term, as in GF.10a--12. The period
translates commute with these transported choices, so the flows descend
and retain the full normal jets and their collars. The two-ruling
replacement is compact and has the same compact exterior.
Its common blowup has exceptional \(\mathbb P^1\times\mathbb P^1\);
either contraction adds the primitive Gysin summands
\(H^{d-2}(\mathbb P^1;\mathbb Z)\), with exceptional ruling degree
\(-1\), to the same integral cohomology group. The two summands are
identified by interchanging the rulings. Removing the real
codimension-four centre preserves the fundamental group; the normal
sphere and disk bundles are simply connected. Thus the full common
blowup and van Kampen maps of F35 carry the integral sphere groups
through every step. The inverse integral period substitutions preserve
the original primitive relations.

\item The based logarithms of the full unit ratios extend through a radial
collar. Multiplying those logarithms by the fixed collar cutoff gives
the explicit unit homotopy relative to its two boundary ends. The
remaining winding factors act on the native quaternion by
\(H_{-1}q_HH_{\ell_i+1}\). The Hopf comparison has left degree
\(-2\) and right degree zero; its inverse cylinder has left degree
\(+2\). This identifies the ordinary double used below with the
same full normal framing as K86 on each selected star.

\par

\medskip\noindent\textbf{ZB.3. The full integer prism seed and target solution}

\par\nopagebreak[4]

\item For any current basis \(P\), let
\[
\begin{aligned}
B_P&=2P^{-1}WP,&
\ell&=(B_{22}-B_{21},B_{11}-B_{12},B_{22}),\\
(x,y)^t&=P^{-1}e_2,& w&=(x,y,-x-y),\\
d&=(-x-2y,-2x-y,x-y),\\
I_i&=-w_i+d_i(\ell_i+1),&
J_i&=-w_i-d_i(\ell_i+1).
\end{aligned}                                           \tag{ZB.5}
\]
The doubled ordinary product has variation matrix \(V\), stable
derivative \(\delta\), and the integral prism equations
\[
\begin{aligned}
V(w_i e_i-w_0e_0)&=I_ie_i-I_0e_0,\\
w_i\delta_i-w_0\delta_0&=J_i-J_0,\qquad i=1,2.
\end{aligned}                                           \tag{ZB.6}
\]
The complete initial values, from the native pair and inverse cylinder
above, are
\[
\begin{aligned}
c&=-2\mu-2,\\
V&=\begin{pmatrix}
c&4a-6&6-4a\\4a-6&2b-4&0\\6-4a&0&-4
\end{pmatrix},\\
\delta&=(c,8-4a-2b,4a-4),\qquad S=0.
\end{aligned}                                           \tag{ZB.7}
\]
In particular the fixed entry is specified as an integer.

\item At a selected old curve \(i=0,1\), put
\(c=V_{ii}\), \(e=\delta_i\). Its actual single native right
coordinate is \(\rho=(e-c)/4\). The divisor four comes from
the literal doubled product and its already fixed native comparison.
The common Cox link is
\((h,q_H)\mapsto(q_H^{-1}hq_H,q_H)\). It acts by
\[
(\mu,\rho,S)\longmapsto(\mu+2\rho,-\rho,S+\rho).
                                                               \tag{ZB.8}
\]
The sign of the last coordinate is the physical K68 orientation.
After the same ordinary cylinder on both sides the target selected
curve is labelled two, and
\[
V'_{22}=c,\qquad\delta'_2=2c-e,\qquad S'=S+(e-c)/4.
                                                               \tag{ZB.9}
\]
Every target in the specified word has three nonzero weights. Use its
own \(I,J,w\), and set
\[
C=w_2^2c-w_2I_2,\qquad E=w_2(2c-e)-J_2.
\]
The complete target solution is
\[
V'_{rs}=\frac{I_r}{w_r}\delta_{rs}+\frac{C}{w_rw_s},
\qquad\delta'_r=\frac{J_r+E}{w_r}.                       \tag{ZB.10}
\]
Indeed substitution verifies every prism equation and the two selected
entries. A symmetric difference annihilates the plane spanned by
\(w_i e_i-w_0e_0\), hence has the form
\(C'(1/w_r)(1/w_s)\); the selected entry makes \(C'=0\).
The derivative difference is a multiple of \((1/w_r)\), and its
selected entry also vanishes. These arguments prove uniqueness. The
actual target product constructed in ZB.2 supplies an integral
solution; therefore this rational expression equals its integral
invariants. The coordinate \(\rho\) in (ZB.8) is its actual free
integer. No finite-group residue enters this reconstruction.

\par

\medskip\noindent\textbf{ZB.4. Every residue is exhibited by a complete actual word}

\par\nopagebreak[4]

\item We enumerate every word beginning with \(Q_0\) and having length one
through ten: 1,023 words on each branch. Their positive first row has
the form \((u,-r)\), \(u,r>0\), so all target weights are nonzero.
The reverse Euclidean construction terminates at \(T_0\) for each
word, as verified independently by its exact integer factorization.
Equations (ZB.7)--(ZB.10) evaluate every word without a residue lift.
The first two length-two examples \(Q_0^2,Q_0Q_1\) give classes
\(6,6\) for \(a=1\) and \(9,0\) for \(a=2\).

\item For each residue we choose a witness by the lexicographic score
consisting of maximum entry size, sum of entry sizes, word length,
and the four matrix entries. The following two tables give all 28
residues, their actual bases and their computed integer sums. A basis
entry \((u,v;x,y)\) denotes the matrix with those two rows.

\par

\medskip\noindent\textbf{Good branch, \(a=1,b=2\).}
\input{notes/part3_source_zero_base_sector_witnesses_a1.tex}

\clearpage
\par

\medskip\noindent\textbf{Reduced branch, \(a=2,b=1\).}
\input{notes/part3_source_zero_base_sector_witnesses_a2.tex}

\item The executable certificate is
\path{code/part3_root/source_zero_base_sector.py}. Its JSON gives each
full variation matrix, derivative, finite word and successive native
right increments. The original independent audit solved 282 full
target linear systems for these 56 witnesses and checked every
coefficient-conjugation equation. The recovered independent K90
Gaussian audit is separately saved with its exact deterministic
checksum. The finite witness list proves surjectivity onto
\(\Theta_7=\mathbb Z/28\) for each of these two specified
fixed-scalar source-zero fan families.

\end{enumerate}
\end{proof}




\begin{theorem}\label{night:P3K.89}
\StatementRecord{night:P3K.89}{proved}{S6\_06\_BLIND: GPT-6 Astra (OpenAI, ChatGPT Pro mode, effort Pro), recorded re-derivation in the checkpoint dated 18 November 2026}{\cref{night:P3F.9}, \cref{night:P3K.23}, \cref{night:P3K.30}, \cref{night:P3K.62}, \cref{night:P3K.68}, \cref{night:P3K.72}, \cref{night:P3K.82}, \cref{night:P3K.86}}
Retain the full source-zero families, their actual finite-cap
characters and the physical vertical fans \(T_j\), \(j\in\mathbb Z\).
Set \((a,b)=(1,2),m=4p,\mu_0=1\) in the good branch and
\((a,b)=(2,1),m=4p-1,\mu_0=0\) in the reduced branch. Put
\[
 P_b(j)=2b\binom{j+1}{3}+(3-b)\binom{j+1}{2}.
\]
The fixed-scalar source loop has physical Jin pair, native free
pair and integer index
\[
 (k_j,\ell_j)=(\mu_0+6P_b(j),4\mu_0),\qquad
 (\mu_j,\nu_j)=(\mu_0,6P_b(j)),\qquad A_j=P_b(j).
\]
For the actual coupled loop \(\lambda(v)=\lambda_0v^{cm}\), with
its realized norm sign, full Cartier and closing maps, the pair is
\[
\begin{aligned}
 k_j&=\mu_0+6P_b(j)+c[3H-3m j(j-1)],\\
 \ell_j&=4\mu_0+24cp,\qquad \mu_j=\mu_0+6cp,\\
 A_j&=P_b(j)+c\left((a-1)(3p-1)-m\binom j2\right),
\end{aligned}
\]
where \(H=2p\) in the good branch and \(H=8p-2\) in the reduced
branch. Its oriented sphere class is \(A_j\bmod28\).
With \(t=cm-1\), the actual doubled ordinary fixed entries are
\[
 V_{00}=2t-2\mu_j-12A_j,\qquad
 \delta_0=2t-2\mu_j+12A_j.
\]
The proof derives these integers by the complete relative source,
line, tangent, cap and parameter maps, before reducing the index.
\end{theorem}
\begin{proof}
\begin{enumerate}[label=\arabic*.]
\item \par

\medskip\noindent\textbf{VZ.1. Actual source, period and normal maps}

\par\nobreak

\item Retain the two actual source-zero families and all their realized cap
characters. The source parameter is \(v\), its positive angle is
\(\theta\), and \(q\) is the selected cusp coordinate. The fixed-scalar
periods, with every analytic unit retained, are
\[
 h:(Z,V)\longmapsto(\kappa qZ,\alpha V),\qquad
 B:(Z,V)\longmapsto(gZ,q\eta V),\qquad
 \operatorname{wind}(\kappa,\alpha,g,\eta)=(1,-a,b,-2),
\]
where \((a,b)=(1,2)\) or \((2,1)\), so \(ab=2\) and \(a+b=3\).
The columns are \(h,B\); in particular \(\operatorname{wind}g=b\).

\item For the literal vertical fan \(T_j\), use \(V_j=Z^{-j}V\) and
\(h_j=hB^j\). The full period characters are
\[
 Q_j=\kappa qg^j,\quad
 \alpha_j=\alpha\eta^j\kappa^{-j}g^{-j^2},\quad
 g_j=g,\quad \eta_j=\eta g^{-j}.                       \tag{VZ.1}
\]
Consequently
\[
 W_j=
 \begin{pmatrix}1+bj&b\\-a-3j-bj^2&-2-bj\end{pmatrix}
 =\binom1{-a-j}(1+bj,\ b).                            \tag{VZ.2}
\]
No periodic logarithm of \(g\) is used.

\item Put \(Q(q,v)=\kappa(q,v)q\). The actual source-disc map is defined by
\[
 Q(q',v)=Q(q,v)g(q,v)^j,\qquad
 q'=a_j(q,v)q,\qquad a_j(0,v)=g(0,v)^j.                \tag{VZ.3}
\]
The analytic inverse function theorem, uniform shrinking over the compact
parameter annulus, and uniqueness give a single periodic holomorphic map
and its inverse. Identity on the marked Tate coordinate identifies the
source periods. For the remaining suspension translation put
\[
 c(q,v)=g(q',v)/g(q,v),\qquad
 (Z,\gamma_j)\longmapsto(c(q,v)^{\gamma_j}Z,\gamma_j).
                                                               \tag{VZ.4}
\]
The logarithm of \(c\) is the unique one zero on \(q=0\). The identity
\(c^{\gamma_j+1}g(q)=g(q')c^{\gamma_j}\) is the complete suspension
square. On the two node charts this acts by reciprocal normal units,
so it extends and is the identity on the fixed locus.

\item Use the marked contracting-\(B\) coefficient in the F9 proper height
coordinate on both sides. The two node-preimage lifts have the prescribed
strict decreasing order in one short real arc. The actual \(h_j^{-1}\)
identifying the second chart has zero \(B\) coefficient. Thus the
source map carries the whole fixed graph and its distinct endpoint
heights; interpolation of the two proper decreasing coordinates has
no integer-height crossing. The full line norms are carried by this
same map.

\item In the native fixed star, \(q=ZAB\), the source-normal map is
\[
 (Z,A,B)\longmapsto(Z,a_j(q,v)A,B).                    \tag{VZ.5}
\]
It intertwines the second chart \(A_1=ZA,B_1=ZB\) exactly; its inverse
uses the specified inverse \(q\)-map. The higher \(q\) terms have no
normal linear contribution. The character shear is the typed Cartier
tensor map, multiplication by \(Z^j\) from the new line to the old.
It preserves the adjoint normal bundle. The remaining normal coefficient
is therefore exactly \(g(0,v)^j\), of winding \(bj\).

\par

\medskip\noindent\textbf{VZ.2. The actual Hopf boundary map}

\par\nobreak

\item Write
\[
 u=(1+\mathbf jZ)/\sqrt{1+|Z|^2},\quad h=ui u^{-1},\quad
 Q_N=A+\mathbf j\bar B,\quad q_H=uQ_N.
\]
Carry the initial normal map into both references and remove its positive
modulus by the same positive Cox path. On a cut parameter interval write
\(g(0,v)^j/|g(0,v)^j|=e^{i\alpha}\), with \(\alpha(0)=0\) and
increment \(2\pi bj\). For \(H=e^{i\alpha/2}\), the literal map is
\[
 Q_N'=H Q_N H,\qquad
 q_H'=e^{h\alpha/2}q_H e^{i\alpha/2}.                 \tag{VZ.6}
\]
The two half-phase endpoints are both \((-1)^{bj}\); their product is
a closed physical-circle-equivariant map. Its orbit-sphere action is
\(\eta\mapsto\operatorname{Ad}(e^{h\alpha/2})\eta\).

\item The constant-in-\(h\) map \(C(q_H)=Hq_HH\) extends over the actual
\(D^3_h\) cap. Compose the map \(F\) of (VZ.6) on the right with
\(C^{-1}\). The resulting coordinate arrow is
\[
 q_H\longmapsto A_j(h,\theta)q_H,\qquad
 A_j=e^{h\alpha/2}e^{-i\alpha/2}.                     \tag{VZ.7}
\]
Its zero-winding factor contracts by its specified real logarithm,
retaining both endpoints. Set \(N=bj\). The resulting map is
\(e^{\pi N\theta h}e^{-\pi N\theta i}\). Its preimages of \(-1\)
have \(h=-i\) and \(\theta=(2r+1)/(2N)\). In the outward domain
basis \((j,-k,\partial_\theta)\) its differential is
\((k,j,2\pi Ni)\). The positive target basis at \(-1\) is
\((i,j,-k)\); the determinant is \(2\pi N\). Thus
\[
              \deg A_j=bj                           \tag{VZ.8}
\]
in outward \(S^2_h\) times positive parameter orientation. The complex
orientation of the source fixed curve is \(-S^2_h\). That sign is
retained in the cap computation below.

\par

\medskip\noindent\textbf{VZ.3. The complete line comparison}

\par\nobreak

\item Use \(u,t,d,\gamma\) for the real coefficients of the source vanishing
period, \(h\), physical delta and \(B\), respectively; \(t\) here is not
the source parameter \(v\). The two actual character maps give
\[
 u_j=u,\quad t_j=t,\quad\gamma_j=\gamma-jt,\quad d_j=d-ju.
                                                               \tag{VZ.9}
\]
For the meridian-one dual line \(U=M^*\), the integral parts of the
actual full coefficient connections are
\[
 A_0=dd+(\alpha_q-2\theta)d\gamma-a\theta\,dt,
\]
\[
 A_j=dd_j+(\alpha_q-(2+bj)\theta)d\gamma_j
                     -(a+3j+bj^2)\theta\,dt.        \tag{VZ.10}
\]
Their difference in original coefficients is
\[
 A_j-A_0=-j\{du+(\alpha_q+\theta)dt+b\theta\,d\gamma\}.
                                                               \tag{VZ.11}
\]
Including the actual unit terms, the braces are
\(du+\alpha_Qdt+\phi_gd\gamma\). The \(q\) deck changes \(u\) by
\(-t\), and the parameter deck changes it by \(-t-b\gamma\);
the invariant form retains both changes.

\item Under (VZ.3), the target phase is
\(\alpha_q'=\alpha_q+bj\theta\) plus its full zero-winding unit term.
In the new coefficients the target reference connection is
\[
 A_{0,q'}=dd_j+(\alpha_q+(bj-2)\theta)d\gamma_j-a\theta\,dt.
\]
Hence the two residual parameter terms are
\[
 A_j-A_{0,q'}=-(3j+bj^2)\theta\,dt-2bj\theta\,d\gamma_j
                    +\text{the full unit correction}.          \tag{VZ.12}
\]

\item Make the actual graph maps \(b,\gamma\), with outer values \(t,\gamma_j\),
constant near the fixed sphere using their real lifts. Use nested tubes
\(U_0\subset U_1\subset U_2\), with every surgery and connection change
inside \(U_1\) and the level maps constant on \(U_2\). They extend over
the same three-ball cap. The theta-cut coefficient arrows
\(w\mapsto e^{2\pi ib}w\) and \(w\mapsto e^{2\pi i\gamma}w\)
define lines with curvatures
\[
 \Xi=d\theta\wedge db,\qquad \Psi=d\theta\wedge d\gamma,\qquad
 \Xi^2=\Xi\Psi=\Psi^2=0.                              \tag{VZ.13}
\]
The full integral evaluations give
\[
 U_j\simeq U_{0,q'}\otimes
 T_b^{-(3j+bj^2)}\otimes T_\gamma^{-2bj}.              \tag{VZ.14}
\]

\item This isomorphism retains the outer gauge. After the displayed period
and meridian evaluations, the sole boundary obstruction is the \(q\)
angle: \(q\)-monodromy kills \(u\) primitively, whereas theta, \(b\)
and gamma extend. At one theta choose the actual free section disc
\(|Z_*|=1,\gamma=t=0\), avoiding the marked sections and fixed graph.
The prescribed \(Z_*^{-j}\) arrow and both tensor frames extend there,
so the boundary ratio has zero \(q\)-angle degree. In vector frames
\(e_l=e^{2\pi ir_{il}}e_i\), remove the zero integer two-cocycle;
a partition of unity gives \(r_{il}=a_l-a_i\).
Replacing \(e_i\) by \(e^{-2\pi ia_i}e_i\) trivializes the ratio.
Remove its extending theta, \(b\) and gamma characters, then extend
its remaining real boundary logarithm by a collar cutoff. The inverse
trivialization is exactly (VZ.14), with the prescribed outer arrow and
the carried inner cap (VZ.7).

\par

\medskip\noindent\textbf{VZ.4. The exact tangent collar}

\par\nobreak

\item Put \(\tau=\log Q/(2\pi i)\), \(G=\log g/(2\pi i)\), and
\(\delta G=\log c/(2\pi i)\). Only the last logarithm is periodic.
In the full invariant source coframe
\(\omega=du+\operatorname{Re}\tau\,dt+
\operatorname{Re}G\,d\gamma\), the actual pullback is
\[
 \omega'=\omega+\operatorname{Re}\delta G(d\gamma-jdt),\quad
 dt'=dt,\quad d\gamma'=d\gamma-jdt.                   \tag{VZ.15}
\]
The full base map is
\[
 r'=r+\log|a_j|,\quad
 \alpha'=\alpha+\arg(a_j)/(2\pi),\quad\theta'=\theta.  \tag{VZ.16}
\]
Its derivatives are single-valued and retain the winding \(bj\).
The full coframe matrix is block triangular with this base Jacobian
and the source block
\[
 \begin{pmatrix}
 1&-j\operatorname{Re}\delta G&\operatorname{Re}\delta G\\
 0&1&0\\0&-j&1
 \end{pmatrix}.                                     \tag{VZ.17}
\]
In the source order \((\omega,\gamma,t)\) it is upper triangular with
diagonal one. All base-to-source entries lie in the off-diagonal block,
so the curvature-square trace is the sum of the two diagonal-block
traces. The unipotent source trace is zero. The base interpolation lives
on the boundary variables \(\alpha,\theta\) and one collar coordinate,
so its curvature square is zero by dimension. This supplies a connection
with \(\operatorname{tr}(F^2)=0\) on the whole comparison collar, not
merely flat endpoints. The line phases stay in (VZ.10)--(VZ.14).
Use the actual image cut under (VZ.16); no radial coordinate map is
claimed to preserve the source functions.

\par

\medskip\noindent\textbf{VZ.5. The four relative interior pairings}

\par\nobreak

\item For the carried dual-line curvature \(F\), use
\[
 \Omega_+=dr\wedge d\alpha\wedge du\wedge dt
                  \wedge d\gamma\wedge d\theta.      \tag{VZ.18}
\]
The required pairings, with their prescribed boundary connections, are
\[
 \int F^2\Xi=1,\quad \int F^2\Psi=0,\quad
 \int p_1\Xi=0,\quad \int p_1\Psi=2.                 \tag{VZ.19}
\]

\item For the first pairing freeze theta. All four actual units then have
holomorphic logarithms on the protected \(q\) disc. The exact BW.2
source, graph, line and Hopf-cap flows, with the same closing exponent
and full coefficient boundary frame, give its relative calibration.
The actual SCN.9/K62 comparison has both outer normal maps and its
common inner physical gauge extending; its separate characteristic
terms vanish, and its whole collar is \(3\int F^2\Xi-3\).
Thus the pairing is one. No logarithm of \(g\) around the parameter
circle is taken. The carried boundary connection makes this slice
pairing constant.

\item For the Psi Euler pairing choose gamma zero outside the fixed graph's
short arc. This level is the actual \(I_1\) source neighbourhood \(E\),
disjoint from the surgery and cap. At fixed theta its boundary
connection is \(dd-a\theta\,dt\), which is closed. The \(q\) deck
changes \(d\) by minus gamma, so its \(q\) holonomy is zero. The line
has degree zero on the sole normalization sphere, exactly its integral
Chern class since \(E\) retracts to the nodal fibre.
The graph map \(b\) extends over the vanishing-cycle handle by its
constant attaching-circle lift and collar cutoff. The flat line with
graph holonomy \(e^{-2\pi ia\theta}\) has boundary connection
\(dd-a\theta\,db\). Its comparison with the actual line has zero
\(q\)-angle winding by the literal disc \(Z_*=-1,\gamma=0\), where
the Cartier frame, \(Z_*^{-j}\) arrow and complete \(q'\) map extend.
A Cech trivialization, the extending \(b\) character, and the final
real boundary logarithm make its relative boundary value exact.
Transport the flat connection through this isomorphism. Its square
vanishes, giving the second pairing. Terms containing \(d\theta\)
vanish after multiplication by Psi, so all parameter-dependent gauges
remain included.

\item For the Xi Pontryagin pairing, choose a regular graph level missing
\(U_2\). It is relatively cobordant to the original source-three-level
times gamma, with the source-pulled connection. The four-form integrates
to zero there; the comparison collar has zero transgression by
(VZ.15)--(VZ.17). This uses nested tubes and does not extend a product
description through the surgery.

\item For the last pairing keep the actual nodal source frame. In \(xy=q\),
with \(\Omega=dx\wedge dy\), put
\[
 V=(x,-y),\qquad H=(\bar y,\bar x)/R^2,\qquad R^2=|x|^2+|y|^2.
\]
Then \(dq(H)=\Omega(V,H)=1\). The fixed-\(x\) and fixed-\(y\) Tate
horizontals are literally
\[
 H_x=H-\frac{\bar y}{xR^2}V=(0,1/x),\qquad
 H_y=H+\frac{\bar x}{yR^2}V=(1/y,0).                  \tag{VZ.20}
\]
Scaling their shear coefficients gives determinant-one homotopies.
The normalized frame
\[
 (V/R,RH)=R^{-1}
 \begin{pmatrix}x&\bar y\\-y&\bar x\end{pmatrix}
                                                               \tag{VZ.21}
\]
has degree one, so the relative second Chern number is one.
The actual \(q\)-polar frame \((V,qH)\) has determinant \(q\);
its relative first Chern class is the fibre, whose square is zero
by disjoint parallel fibres. Hence the source relative \(p_1\) is
\(-2\). The real period frame differs by (VZ.20) and the positive
matrix of its marked period lattice; the polar rotation factors
through its angle circle and has no third transgression.
Finally \(\Psi=d\theta\wedge d\gamma\) reverses the last two factors
of (VZ.18), giving the asserted \(+2\).

\par

\medskip\noindent\textbf{VZ.6. The cap and the original-boundary Euler collar}

\par\nobreak

\item The complex fixed curve has orientation \(-S^2_h\). The physical
quotient near the fixed tube has orientation
\(-S^2_h\wedge dr\wedge S^2_\eta\); its inner boundary is
\(+S^2_h\times S^2_\eta\), so the inserted cap is
\(-D^3_h\times S^2_\eta\). With coordinate clutch
\(v_{\rm ref}=A_jv_j\), its difference base is
\[
 -(D^3_j\cup-D^3_{\rm ref})\times S^1_\theta.         \tag{VZ.22}
\]
In the positive-spinor representation,
\(\int\operatorname{tr}(A^{-1}dA)^3=-24\pi^2\deg A\).
Use the radial connection \(f(r)A^{-1}dA\) in the first frame,
equivalently \(-f(r)dA A^{-1}\) in the second. The stated clutch
and boundary direction give
\[
 c_1(E)=0,\quad c_2(E)=-bj,\quad
 p_1(\operatorname{ad}E)=4bj.                        \tag{VZ.23}
\]
The physical tautological line has fibre value minus one.
The projective relation yields positive-parameter pair
\((-bj,-4bj)\); K68's literal parameter reversal gives
\[
                 (\Delta k,\Delta\ell)_{\rm cap}=(bj,4bj).
                                                               \tag{VZ.24}
\]
The constant map \(C\), positive factors and unit logarithms extend
as already specified, so this is the cap of the actual line arrow.

\item In the original exterior expressed in the piece coefficients
\(\gamma=\gamma_j\), the complete integral connection and primitive are
\[
 A_{\rm ext}=dd_j+(\alpha-2\theta)d\gamma+j\,du+
                        [j\alpha-(2j+a)\theta]dt,
\]
\[
 \beta=-j\{du+[\alpha+(1+bj)\theta]dt+b\theta\,d\gamma\}.
                                                               \tag{VZ.25}
\]
The full unit terms are the derivatives of the carried coefficient
maps. Their relative Cech boundary arrow is fixed; the zero-period
factors deform together with that endpoint primitive retained.
In \(\Omega_B=d\alpha\wedge du\wedge dt\wedge d\gamma\wedge d\theta\),
direct exterior differentiation gives
\[
 3\int_0^1\beta(F_{\rm ext}+s\,d\beta)^2\,ds
          =(-6aj-9j^2-2bj^3)\Omega_B.                \tag{VZ.26}
\]
The actual outward interpolation is from piece one to exterior zero,
so its contribution is \(6aj+9j^2+2bj^3\). Its tangent transgression
is zero by VZ.4.

\par

\medskip\noindent\textbf{VZ.7. Independent fixed-scalar integer assembly}

\par\nobreak

\item The K68 physical/dual line map and parameter reversal give
\(k=\int_{\Omega_+}F_U^3\), \(\ell=\int_{\Omega_+}p_1F_U\)
in this comparison. Equations (VZ.13), (VZ.14) and (VZ.19) give
\[
 \Delta k_{\rm int}=-9j-3bj^2,\qquad
 \Delta\ell_{\rm int}=-4bj.                          \tag{VZ.27}
\]
Adding the actual cap and original-boundary collar yields
\[
 \Delta\ell=0,\qquad
 \Delta k=2bj^3+(9-3b)j^2+(6a-9+b)j=6P_b(j),
\]
\[
 P_b(j)=2b\binom{j+1}{3}+(3-b)\binom{j+1}{2}.        \tag{VZ.28}
\]
The final equality uses \(a+b=3\). NC/K86 supplies reference pair
\((\mu_0,4\mu_0)\), where \(\mu_0=1\) in the good branch and zero
in the reduced branch. Thus
\[
 (k_j,\ell_j)=(\mu_0+6P_b(j),4\mu_0),\quad
 (\mu_j,\nu_j)=(\mu_0,6P_b(j)),\quad A_j=P_b(j).       \tag{VZ.29}
\]
This is an integer computation before reduction modulo 28.
Explicitly \(P_2(j)=j(j+1)(4j-1)/6\) and
\(P_1(j)=j(j+1)(j+2)/3\). K82 supplies only an independent
residue check, not an integer selection.

\par

\medskip\noindent\textbf{VZ.8. All four native phase-square edges}

\par\nobreak

\item Use the actual RT.12 independent parameters
\(v=v_*e^{2\pi it}\), \(r=r_{\rm base}e^{2\pi i\rho}\),
where \(r_{\rm base}=v_*^c r_0\), \(r_0^m=\lambda_0\).
Retain the full RCZ map and norm
\[
 B=r^mB_0(v),\quad A=r^nA_0(v),\quad
 F=F_0(v)+n\frac{\log r}{2\pi i}s_\infty\delta,\quad A_0^m=B_0^n,
                                                               \tag{VZ.30}
\]
with either realized \(n=\pm1\). Here \(F\) is ordinary-to-cap.
Since \(q=q_\infty^{-1}h_\infty\), increasing \(\log r/(2\pi i)\)
by one gives the literal directed square
\[
 F_{\rm end}R_{\rm ord}=R_{\rm cap}F_{\rm start},\qquad
 R_{\rm ord}=q^n,\quad R_{\rm cap}=h_\infty^n.         \tag{VZ.31}
\]
Both \(s_\infty\) terms cancel. This preserves the actual finite
characters, reduced Cartier correction and complete \(F_0\).

\item At the selected cusp, include \(r_{\rm base}^m\) in \(\eta\), but
exclude its variable phase. On the source cut use full logarithms
\(G=\log g,H=\log\eta\), with exact jumps
\[
 G(t+1)-G(t)=2\pi ib,\qquad H(t+1)-H(t)=-4\pi i.
\]
After the forward \(q^n\) arrow and the actual \(B^{-n}\) reindex,
the scalar endpoint is \(E_r=(g^{-n},\eta^{-n})\).
The current-to-reference cut product is
\[
 C(t,\rho)=(A,D),\quad
 A=e^{n\rho G(t)},\quad
 D=e^{n\rho H(t)+\pi inm\rho(\rho-1)}.                \tag{VZ.32}
\]
It is the inverse of the K30 reference-to-current product and
\(C(t,1)E_r(t)=1\) on full collars. For \(T_j\) put
\(D_j=DA^{-j}\). Its exact native chart factors are
\[
 (z,c_0,c_1)\longmapsto
 (Az,(AD_j)^{-1}c_0,D_jc_1).                         \tag{VZ.33}
\]
The Cox lift is
\((a_0,a_1,c_0,c_1)\mapsto
(a_0,Aa_1,(AD_j)^{-1}c_0,D_jc_1)\).
The second chart factors are \((A^{-1},D_j^{-1},AD_j)\);
both preserve the physical circle and \(q\).

\item The four edges are explicit. At \(\rho=0\) the product is the identity,
the fixed-scalar source product. At \(t=0\) it is the exact inverse
K30 scalar product with the actual amplitude. At \(\rho=1\) it is
\((g^n,\eta^n)=E_r^{-1}\). On \(t=1\) relative to \(t=0\), its full
mismatch in these characters is
\[
 K_j(\rho)=
 (e^{2\pi inb\rho},e^{-2\pi in(2+bj)\rho}).           \tag{VZ.34}
\]
All zero-winding units have equal endpoint values; their chosen
logarithm differences are the two integers just displayed.

\item For native factors \((a,p,q_1)\) with \(apq_1=1\), the literal
quaternion lift has \(L^2=pq_1,\ R^2=p/q_1\).
Substitution gives
\[
 L=e^{-\pi inb\rho},\quad
 R=e^{\pi in(4+2bj-b)\rho},\quad q_H\longmapsto Lq_HR.
                                                               \tag{VZ.35}
\]
Both endpoint half-phases equal \((-1)^{nb}\); the full map closes
even in the odd case. The base \(h\) and orbit sphere \(\eta\) both
rotate by \(\operatorname{Ad}L\), so this is the literal extending
linear rotation on \(D^3_h\times S^2_\eta\), with principal lift
(VZ.35). The base and normal Spin endpoints are both \((-1)^{nb}\);
their product is one. The normal map is independent of \(h\), hence
its primary third classes vanish by its factorization through the
parameter circle. Positive moduli use the same positive Cox maps at
both references with their explicit triangular derivative homotopy.

\par

\medskip\noindent\textbf{VZ.9. Every scalar seam and the coupled integer pair}

\par\nobreak

\item The same formula is defined for \(0\le\rho\le c\). At integer \(\rho=c\)
the quadratic factor is one because \(c(c-1)\) is even, and
\[
 C(t,c)=(g^{nc},\eta^{nc}),\qquad C(t,c)E_r^{(c)}(t)=1.
                                                               \tag{VZ.36}
\]
Every intervening seam is the accumulated product of the actual
endpoint maps. The ordinary and cap endpoints are \(q^{nc}\) and
\(h_\infty^{nc}\), with the same \(F_0\) and finite-character square.
The two half-phases at this endpoint are \((-1)^{nbc}\).

\item On the coupled diagonal \(\rho=ct\), the earlier DZ product has
current-to-reference logarithm
\[
 (nctG(t),\ nctH(t)+2\pi inm c^2t^2).
\]
Its difference from (VZ.32) is precisely
\[
                  \pi inm(c^2t^2+ct)\delta.          \tag{VZ.37}
\]
Its exponential is one at both endpoints because \(c(c+1)\) is even.
It is the actual physical delta loop, with its full winding retained;
it is identity on the orbit normal three-plane and extends over the
same physical line and Hopf cap. Flatten both products on the common
endpoint collars using (VZ.36). The remaining logarithm interpolation
carries every normal jet. Thus this is the actual DZ product comparison,
not a new product chosen on the fixed tube.

\item Apply the RT.12 based phase homotopy from \((t,ct)\) to the broken
source-then-scalar path at the retained amplitude. The four edge maps,
(VZ.35) and (VZ.37), identify its actual native caps; their only
mismatches are the displayed cap-extending maps. Thus free relative
coordinates add in the integral K23 group, or equivalently by the
exact normal-section and derivative composition. Its central
commutator does not alter these free coordinates. The same completed
pair-of-pants carries the physical line, so both characteristic
numbers add as integers.

\item At fixed source the actual scalar page has \(\chi=p\) and
\(H=ma^2/2\). K72 in this same native star gives scalar free pair
\[
 (\mu_{\rm sc},\nu_{\rm sc})=
 \left(6p,\ 6\left((a-1)(3p-1)-m\binom j2\right)\right).
                                                               \tag{VZ.38}
\]
The literal cap orders are \(m=4p\) in the good branch and \(m=4p-1\)
in the reduced branch. Combining with (VZ.29), and using
\(k_M=\mu+\nu,\ell_M=4\mu,A_M=\nu/6\), gives
\[
 \begin{aligned}
 k_j&=\mu_0+6P_b(j)+c[3H-3m j(j-1)],\\
 \ell_j&=4\mu_0+24cp,\\
 \mu_j&=\mu_0+6cp,\\
 A_j&=P_b(j)+c\left((a-1)(3p-1)-m\binom j2\right).
 \end{aligned}                                       \tag{VZ.39}
\]
The original coupled source construction has \(c\ge1\). The written
phase-square maps also carry every integer winding in the retained
contracting annulus. The full data table is
\[
 \begin{array}{c|c|c|c|c|c}
 (a,b)&m&\mu_0&\chi&H&A_{\rm scalar,0}\\\hline
 (1,2)&4p&1&p&2p&0\\
 (2,1)&4p-1&0&p&8p-2&3p-1
 \end{array}.                                        \tag{VZ.40}
\]
For the actual doubled ordinary product put \(t=cm-1\).
Its native-to-ordinary left correction is \(2t\), with the existing
DZ/K82 Cox arrows, and hence
\[
 V_{00}=2t-2\mu_j-12A_j,\qquad
 \delta_0=2t-2\mu_j+12A_j.                            \tag{VZ.41}
\]
These are full integer seeds; no lift is selected from a central
residue and no new global ordinary-product recurrence is asserted.

\item The executed local and assembly certificate is
\path{code/part3/source_zero_vertical_pair.py}. The independent component
audits were PS.1--7 and NP.1--4 by fourfold; their complete maps are
reproduced here. Root and fourfold passed the full proof in the
conversation before the environment loss.

\end{enumerate}
\end{proof}




\begin{theorem}\label{night:P3K.90}
\StatementRecord{night:P3K.90}{proved}{S6\_06\_BLIND: GPT-6 Astra (OpenAI, ChatGPT Pro mode, effort Pro), recorded re-derivation in the checkpoint dated 18 November 2026}{\cref{night:P3F.35}, \cref{night:P3K.68}, \cref{night:P3K.88}, \cref{night:P3K.89}}
For the actual E7 source-zero family and marked data of ZA.1,
\(a=1,2\), \(b=3-a\), \(p\ge1\), \(k\ge0\), and every
\(P\in GL_2(\mathbb Z)\), the terminating construction
(ZA.4)--(ZA.12) gives the physical positive-loop return
\(S\bmod28\) with its full source, Cartier, cap and closing maps.
The vertical seed has physical pair
\((\mu+6A_j,4\mu)\), where
\[
\mu=2-a+6pk,\qquad
A_j=2b\binom{j+1}{3}+(3-b)\binom{j+1}{2}
+k(a-1)(3p-1)-k(4p-a+1)\binom j2.
\]
All free integral matrices and derivatives are computed before
reduction. Both fixed-scalar branches attain every element of
\(\Theta_7\) by the full K88 witnesses.
\end{theorem}
\begin{proof}
\begin{enumerate}[label=\arabic*.]
\item \par

\medskip\noindent\textbf{ZA.1. The actual family, its full vertical integer seed and its closing}

\par\nopagebreak[4]

\item We use the positive source-zero loop in
\[
Y^2=X^3+(1-2v)X^2+TX+v^2,\qquad G=(0,v),\qquad
q=T-(v^2-2v),\qquad\lambda=v^{km}\lambda_0.
\]
The two marked branches are
\[
(a,b,m)=(1,2,4p),\qquad(2,1,4p-1),\qquad
p\ge1,\quad k\ge0.
\]
We retain the actual Cartier line, both cap norm signs, every finite
character, the complete annular cochain and its prescribed endpoint.
The physical circle line and the forward return convention are K68.
For every integer \(j\) put
\[
\begin{aligned}
N&=km,& t&=N-1,& \mu&=2-a+6pk,\\
P_b(j)&=2b\binom{j+1}{3}+(3-b)\binom{j+1}{2},\\
A_j&=P_b(j)+k(a-1)(3p-1)-N\binom j2.
\end{aligned}                                           \tag{ZA.1}
\]
The binomials are their integer-valued falling-factorial polynomials.
K89 proves on the actual vertical fan
\(T_j=T(e_1+je_2,e_2)\) the full physical pair
\[
(k_M,\ell_M)=(\mu+6A_j,4\mu),
\qquad(\mu,\nu)_{\rm native}=(\mu,6A_j).                \tag{ZA.2}
\]
Its proof includes the complete relative normal calculation, the
source-dependent collar and the native parameter composition square.
For clarity the latter retains
\[
B=r^mB_0(v),\qquad A=r^nA_0(v),\qquad
F=F_0+\frac{n\log r}{2\pi i}s_\infty\delta.
\]
Here \(r_{\rm base}=v_*^k r_0\), \(r_0^m=\lambda_0\), and the
actual chosen finite-cap roots obey \(A_0^m=B_0^n\). These
equalities fix the full base and root arrows before the phase square
is formed.
On the scalar/source phase square the current-to-reference arrow is
\[
C(t,\rho)=\left(
e^{n\rho G},\ e^{n\rho H+\pi i nm\rho(\rho-1)}
\right),                                               \tag{ZA.3}
\]
where \(G=\log g\) and \(H=\log\eta\) include the full base units.
The ordinary and cap arrows are \(q^n\) and \(h_\infty^n\), with
\(q=q_\infty^{-1}h_\infty\); they satisfy
\(F_{\rm end}R_{\rm ord}=R_{\rm cap}F_{\rm start}\).
The source-edge ratio is
\[
K_j(\rho)=\left(e^{2\pi inb\rho},
e^{-2\pi in(2+bj)\rho}\right).
\]
In the complete native charts its factors are
\((A,(AD_j)^{-1},D_j)\) and
\((A^{-1},D_j^{-1},AD_j)\), where \(D_j=DA^{-j}\).
The quaternion edge has
\[
L=e^{-\pi inb\rho},\qquad
R=e^{\pi in(4+2bj-b)\rho}.
\]
Their endpoints agree as \((-1)^{nb}\); the paired Spin endpoint
is positive one. The conjugation by \(L\) extends over the fixed
three-ball and its Hopf line. The remaining scalar/source diagonal
comparison is the actual physical circle
\(e^{\pi inm(k^2t^2+kt)}\); it closes because \(k(k+1)\) is even,
extends over that same cap and factors through a circle. These explicit
maps are the attachment carried through every subsequent cone chart.

\par

\medskip\noindent\textbf{ZA.2. Identical-fan maps and terminating construction from a vertical fan}

\par\nopagebreak[4]

\item For columns \(v,w\) of \(P\in GL_2(\mathbb Z)\), the fan
\(T(P)\) consists of the translates of
\(\{0,v,w\}\) and \(\{v,w,v+w\}\). Put
\[
R=\begin{pmatrix}-1&-1\\1&0\end{pmatrix},\qquad R^3=I.
\]
The first and second \(PR\) triangles are respectively the first
and second \(P\) triangles translated by \(-v\) and \(-2v\).
Negating \(P\) interchanges the triangle types after translation
by \(-(v+w)\); swapping columns interchanges their named vertices.
These equalities specify the comparison: the identity on the absolute
dense-torus characters has the same affine semigroup on each matched
cone. It therefore extends holomorphically on every cone and agrees
on overlaps. It commutes with the full periods and the physical circle,
and it is the identity on the exterior source and finite caps. The
full closing (ZA.3) is carried by these same maps.

\item Swap the columns for negative determinant. The six choices
\(P,PR,PR^2,-P,-PR,-PR^2\) have exactly one representative of the form
\[
T_j=\begin{pmatrix}1&0\\j&1\end{pmatrix}
\quad\hbox{or}\quad
P_c=\begin{pmatrix}u&-r\\s&z\end{pmatrix},
\qquad u,r>0,\quad uz+rs=1.                              \tag{ZA.4}
\]
The six origin cones have rays \(\pm v,\pm w,\pm(v-w)\).
The physical primitive vector \(e_2\) lies on one ray or in the
interior of one cone. In the first case its positively oriented
adjacent triangle gives \(T_j\). In the second, the two boundary
vectors ordered from positive to negative first coordinate give
\(P_c\), with coordinates \((r,u)\) for \(e_2\). This proves the
asserted unique choice using the actual fan maps above.

\item Let \(Q_0,Q_1\) have the values in ZB.2. Starting from \(P_c\),
right multiplication by \(Q_0^{-1}\) for \(r\ge u\), or by
\(Q_1^{-1}\) for \(u>r\), replaces the positive pair by
\((u,r-u)\) or \((u-r,r)\). Its sum decreases. The determinant
gives \(\gcd(u,r)=1\). Thus equality occurs at \(u=r=1\), and
the next step is a unique \(T_j\). Reversing the recorded choices
gives
\[
P_c=T_jQ_{i_1}\cdots Q_{i_d},\qquad i_s\in\{0,1\}.
                                                               \tag{ZA.5}
\]
The first forward step is \(Q_0\), and every later first row has
the form \((u,-r)\) with positive entries \(u,r\). Its three
physical weights \((r,u,-r-u)\) are all nonzero.

\item Each factor has the actual other-ruling charts
\[
(n_0/n_1,zn_1,n_1),\qquad(n_1/n_0,zn_0,n_0),
\qquad q=zn_0n_1.
\]
The common rank-one matrix is
\(\left(\begin{smallmatrix}n_0&n_1\\zn_0&zn_1\end{smallmatrix}\right)\).
For \(Q_1\) the normal interchange is the unitary path (ZB.3).
The full multipliers and Cox coefficients are exactly (ZB.1) and
(ZB.4), before winding is taken. They preserve the product defining
\(q\), agree with the inverse character change on the source
exterior, and commute with the finite-root, line and closing squares
of ZA.1. This supplies the actual fan family for each matrix in the
finite construction.

\item For the relative smooth comparison we transport the Hermitian link
metric through these charts and apply its positive radial normalization.
On each protected normal tube the exact rescaling
\(R_s(z,n)=(z,sn)\), followed by the cutoff flow of its velocity,
carries the higher normal terms. This is the GF.10a--12 comparison;
the period translates commute with the transported maps and the
protected vertex collars remain fixed. Thus it carries the full
normal jets, as well as the displayed leading factors.
The compact two-ruling replacement has its original compact exterior.
The common blowup adds, for either contraction, the integral Gysin
summands \(H^{d-2}(\mathbb P^1;\mathbb Z)\). Integration along the
exceptional ruling has degree \(-1\), so these summands are primitive;
interchanging the two rulings identifies them. The normal sphere and
disk bundles are simply connected, and removing the centre of real
codimension four preserves \(\pi_1\). The actual common-blowup and
van Kampen maps of F35 therefore carry integral sphere cohomology
and the trivial fundamental group through each step. The inverse
integral period substitutions retain the primitive exterior relations.

\par

\medskip\noindent\textbf{ZA.3. The complete integer seed at every vertical fan}

\par\nopagebreak[4]

\item The full unit matrix has winding
\[
W=\begin{pmatrix}1&b\\-a&N-2\end{pmatrix}.
\]
For any current basis \(P\), define
\[
\begin{aligned}
B_P&=2P^{-1}WP,&
\ell&=(B_{22}-B_{21},B_{11}-B_{12},B_{22}),\\
(x,y)^t&=P^{-1}e_2,&w&=(x,y,-x-y),\\
d&=(-x-2y,-2x-y,x-y),\\
I_i&=tw_i+d_i(\ell_i-t),&
J_i&=tw_i-d_i(\ell_i-t).
\end{aligned}                                           \tag{ZA.6}
\]
The full normal factors are \(H_tq_HH_{\ell_i-t}\). Their based
unit logarithm at a transformed multiplier is
\[
(P^{-1})_{i1}(P_{1s}f_\kappa+P_{2s}f_g)
+(P^{-1})_{i2}(P_{1s}f_\alpha+P_{2s}f_\eta).
                                                               \tag{ZA.7}
\]
These are logarithms of the actual periodic unit ratios after their
specified winding factors have been removed. A fixed radial cutoff
contracts them relative to the common exterior. The native-to-ordinary
Hopf expression is
\((UH_tU^{-1})q_HH_{\ell_i-t}\). Its left degree is \(2t\) and
right degree zero; the inverse cylinder has left degree \(-2t\).
The same full comparison applies on each selected star in ZA.2.

\item At \(P=T_j\), the physical weights are \((0,1,-1)\).
Using the complete native pair (ZA.2) and this cylinder gives
\[
\begin{aligned}
V_{00}&=2t-2\mu-12A_j,&
\delta_0&=2t-2\mu+12A_j,\\
V&=\begin{pmatrix}
V_{00}&-I_0&I_0\\-I_0&I_1&0\\I_0&0&-I_2
\end{pmatrix},&
\delta&=(\delta_0,J_1-J_0,J_0-J_2),\qquad S=A_j.
\end{aligned}                                           \tag{ZA.8}
\]
Substitution verifies all six variation and both derivative prism
equations
\[
V(w_i e_i-w_0e_0)=I_ie_i-I_0e_0,\qquad
w_i\delta_i-w_0\delta_0=J_i-J_0,\qquad i=1,2.
                                                               \tag{ZA.9}
\]
The fixed entry and its derivative come from the complete free
integer calculation in K89. Their values are determined before
reduction modulo 28.

\par

\medskip\noindent\textbf{ZA.4. The actual single flop and its unique full target solution}

\par\nopagebreak[4]

\item Select the old curve \(i=0\) or \(i=1\), and set
\(c=V_{ii}\), \(e=\delta_i\). Its actual single native right
coordinate is \(\rho=(e-c)/4\). The common link map
\((h,q_H)\mapsto(q_H^{-1}hq_H,q_H)\), with the complex normal
reset, acts on the native free coordinates by
\[
(\mu,\rho,S)\longmapsto(\mu+2\rho,-\rho,S+\rho).
                                                               \tag{ZA.10}
\]
The normal interchange conjugates its two Spin factors by a constant
positive lift. With the same ordinary cylinder on both sides, the
selected target is curve two, and
\[
V'_{22}=c,\qquad\delta'_2=2c-e,\qquad S'=S+(e-c)/4.
                                                               \tag{ZA.11}
\]
Use the target coefficients of (ZA.6). All its weights are nonzero.
The target has the unique symmetric variation and derivative
\[
\begin{aligned}
C&=w_2^2c-w_2I_2,& E&=w_2(2c-e)-J_2,\\
V'_{rs}&=\frac{I_r}{w_r}\delta_{rs}+\frac{C}{w_rw_s},&
\delta'_r&=\frac{J_r+E}{w_r}.
\end{aligned}                                           \tag{ZA.12}
\]
Substitution proves the target prism equations and (ZA.11). To prove
uniqueness, subtract two symmetric solutions. The difference
annihilates the plane spanned by \(w_i e_i-w_0e_0\); its orthogonal
line is generated by \((1/w_0,1/w_1,1/w_2)\). It is therefore a
multiple of the outer square of that vector. The selected diagonal
entry makes that multiple zero. The same plane argument for the
derivative gives a multiple of the vector itself, also zero at the
selected entry. This proves uniqueness of the complete solution.

\item The actual family maps of ZA.2 and the full ordinary product of ZA.3
provide an integral target solution. Its unique rational expression
is (ZA.12), so every entry there is integral. Moreover \((e-c)/4\)
is the native free integer already identified by the literal double.
The computation never divides an element of \(\mathbb Z/28\).
Induction through the terminating word (ZA.5) now computes every
ordinary invariant and the physical return \(S\bmod28\) of the
specified original fan. The identical-fan arrows in ZA.2 carry the
answer back to the input basis, with its complete closing map.

\par

\medskip\noindent\textbf{ZA.5. Executed certificate, independent audit and residue witnesses}

\par\nopagebreak[4]

\item The program \path{code/part3_root/source_zero_all_fans.py} uses integers
and exact fractions. It evaluates all 6,504 unimodular matrices with
entries of absolute value at most 18, representing 814 actual fans,
on each of the ten inputs
\[
a=1,2,\qquad(p,k)=(1,0),(1,1),(1,2),(2,1),(3,2).
\]
All 65,040 inputs pass canonicalization, terminating factorization,
full prism equations, target selected entries, reverse reset and
integer-coordinate checks. Signed vertical seeds from \(-30\)
through \(30\) are evaluated on each input. A separate three-edge
traversal checks 210 states through depth seven for each case; four
transitions needing a separate fixed seed are explicitly recorded
and are not used as premises. Their target fans are calculated by
the proved factorization from (ZA.8).

\item The program also checks 6,504 inverse conjugations on each of the four
independent full multiplier exponents, and 224 exact linearity
identities in the scalar winding. The 56 full K88 witnesses are
recovered with their complete matrices, derivatives and path sums.
They prove that both fixed-scalar source-zero fan families attain
every element of \(\Theta_7\). The eight displayed coupled inputs
also exhibit all 28 residues within the executed range. We make no
universal coupled-image assertion from that finite observation.

\item The independent recovered JavaScript audit solves the full linear
systems by fresh rational Gaussian elimination: 60 seed systems and
1,020 target systems along 420 paths. It performs 1,020 full
coefficient/winding comparisons and 1,020 reverse resets, with
\(a=1,2\), \(p=1,2\), \(k=0,1,3\), \(-2\le j\le2\).
Its target solver does not use (ZA.12). Its complete result has
98,996 JavaScript string code units and deterministic FNV-1a-64 checksum
\texttt{f097696153f172a2}, equal to the record before recovery. The full
ZA.1--5 proof was also independently audited from its pasted text
during the outage. Its recovered scripts and JSON are preserved in
\path{code/part3_fourfold/}. The finite calculations verify implementations
and supply witnesses; the all-input conclusion follows from the
complete family maps, terminating construction and uniqueness proof.

\end{enumerate}
\end{proof}




\begin{theorem}\label{night:P3K.91}
\StatementRecord{night:P3K.91}{proved}{S6\_06\_BLIND: GPT-6 Astra (OpenAI, ChatGPT Pro mode, effort Pro), recorded re-derivation in the checkpoint dated 18 November 2026}{\cref{night:P3K.76}, \cref{night:P3K.78}, \cref{night:P3K.79}, \cref{night:P3K.84}, \cref{night:P3K.87}}
Retain a fixed OS49 source and the two actual finite caps of P3K.84, on any lattice-basis fan, together with finitely many marked smooth order-one clutches. Put \(m=3k\), \(h=6l+1\), \(M=mh\), and \(D=hn_A+mn_B+M\sigma=\epsilon=\pm1\). Their full integral root-groupoid class is \((D,d_A\bmod k,d_B\bmod h)\) in \(\mathbb Z\oplus\mathbb Z/k\oplus\mathbb Z/h\). The explicit family isomorphism redistributes the ordinary clutches to the canonical two-cap datum, retaining both finite characters and all annular integers. Every holomorphic endpoint is covered by the actual physical scalar action; the minimum positive scalar winding is \(kh\). The constructed parameter \(w=r^3\), \(\lambda=w^{kh}\), has positive-fan physical pair \((2kh,8k(3l+1))\) and return \(kl\). On every fan, use the P3K.87 algorithm with \(M=kh\), \(A_0=kl\); an endpoint of winding \(Nkh\) has \(N\) times that return. The proof constructs the endpoint and native cover maps before using integral characteristic naturality.
\end{theorem}
\begin{proof}
\begin{enumerate}[label=\arabic*.]
\item \par

\medskip\noindent\textbf{RE.1. The full two-root integral groupoid}

\par\nobreak

\item In homology order \((\gamma,u,w,\delta)\), put
\(\ell_0=\gamma-u\), \(\psi=(1,0,0,0)\), and retain
\[
A=\begin{pmatrix}1&0&0&0\\0&1&-1&0\\3&3&-2&0\\1&1&-1&1\end{pmatrix},\quad
B=\begin{pmatrix}1&0&0&0\\0&0&1&0\\0&-1&1&0\\0&0&0&1\end{pmatrix},\quad
C=\begin{pmatrix}1&0&0&0\\0&1&0&0\\-3&-2&1&0\\-1&-1&0&1\end{pmatrix}.
\tag{RE.3}
\]
They satisfy \(ABC=1\), \(A^{3}=1\) and preserve \(\psi\). At IV-star the entire
cyclic group \(C_m\) acts by A on the constant root lattice \(\Lambda=\mathbb Z^4\),
including its ineffective subgroup. At the reduced-II root the period
sheaf is the invariant extension of the punctured lattice. Its elliptic
monodromy is the rank-two block \(B_{\rm ell}\), since \(h=1\) modulo 6. The exact K84
flow \(e^{2\pi i\delta}B_\gamma(O,1)\) identifies the extending period stalk
with \(I_B=\mathbb Z\gamma\oplus\mathbb Z\delta\). The root deck fixes these two actual periods.
B-I is inverted only on its rank-two elliptic summand, where its
determinant is 1; it is never inverted on the full rank-four lattice.

\item The constant invariant sections \(I_B\) map equivariantly to the root-disc
period sheaf. Using two slit sectors and the centre, the punctured
circle complex is \(\Lambda\xrightarrow{\,B-I\,}\Lambda\). Its degree-one coinvariant
term cancels the identical radial-extension term, leaving \(I_B\) in degree
zero and zero higher disc cohomology. This integral cancellation commutes
with the actual root rotation. Thus the map from the constant invariant
sections is an equivariant quasi-isomorphism of disc section complexes.
The finite action on \(I_B\) is the identity, so its groupoid resolution has
differentials alternating zero and multiplication by h. Finite invariants
over C have not been applied to this lattice complex.

\item The outer-annulus restriction sends a degree-one cyclic cochain \(a_B\) in
\(I_B\) to the SAME vector in \(\Lambda\). Indeed a positive quotient loop lifts
from one root sector to the next; h successive sectors make the root
circle. Its cyclic norm is h \(a_B\), while the one-sector restriction is
\(a_B\). This supplies the actual index of the restriction map. At the
selected cusp the invariant stalk is \(I_C=\mathbb Zw\oplus\mathbb Z\delta\), with K78\textquotesingle s integral
sector cancellation and full determinant-one toric valuation map.

\item On the pair of pants, a one-cochain (p,q) restricts to \(p,q,-C(p+Aq)\).
The full degree-two cocycle is \((\eta_A,\eta_B,b_A,b_B,b_C)\), with \(\eta_A\) in
\(\Lambda^A=\mathbb Z\ell_0\oplus\mathbb Z\delta\), \(\eta_B\) in \(I_B\) and the three b\textquotesingle s in \(\Lambda\).
Its complete degree-one relations are
\[
\begin{aligned}
\eta_A&=N_m a_A,&\eta_B&=h a_B,\\
b_A&=p-a_A-(A-I)r,&b_B&=q-a_B-(B-I)s,\\
b_C&=-C(p+Aq)-(C-I)t,
\end{aligned}\qquad
N_m=k\begin{pmatrix}3&0&0&0\\-3&0&0&0\\0&0&0&0\\0&0&-1&3\end{pmatrix}.
\tag{RE.4}
\]
Here \(p,q,a_A,r,s,t\) lie in \(\Lambda\) and \(a_B\) lies in \(I_B\). This is the
16-by-26 restriction-cone matrix. Extra ordinary discs add their full
transported vectors to \(v_*=b_A+Ab_B+C^{-1}b_C\) through actual identity
blocks, as in K78. After the unimodular elimination of p,q, its relation is
\[
v_*=-a_A-Aa_B-(A-I)r-A(B-I)s-(I-C^{-1})t.\tag{RE.5}
\]
The last three images are precisely ker \(\psi\). Ordered as (r,s,t), explicit
preimages of u,w,\(\delta\) are
\[
(-u,w,0),\qquad(-w,-u-w,0),\qquad(2w,2(u+w),-Cu).\tag{RE.6}
\]
Hence the elimination is primitive over Z. In the remaining basis
\((\ell_A,\delta_A,\gamma_B,\delta_B,c)\) the complete relation matrix is
\[
R=\begin{pmatrix}m&0&0&0&0&0\\0&0&-k&m&0&0\\0&0&0&0&h&0\\
0&0&0&0&0&h\\-1&0&0&0&-1&0\end{pmatrix}.\tag{RE.7}
\]
It says \(c=m\ell_A=h\gamma_B\), \(k\delta_A=0\), \(h\delta_B=0\). Choose integers
a,b with \(ah+bm=1\). Then \(g=a\ell_A+b\gamma_B\) satisfies \(\ell_A=hg\),
\(\gamma_B=mg,\ c=Mg\), proving the exact full answer
\[
H^2(\mathcal H)=\mathbb Zg\oplus(\mathbb Z/k)\delta_A
 \oplus(\mathbb Z/h)\delta_B,\qquad
[X]=(D,d_A\bmod k,d_B\bmod h).\tag{RE.8}
\]

\par

\medskip\noindent\textbf{RE.2. Actual ordinary-clutch and parameter comparisons}

\par\nobreak

\item The Lie sheaf has coarse terms O \(\delta\) and O(-1). At the reduced-II root
the exact coordinate \(\tau\prime=z^l\tau\) sends the original source vector to
\(z^l\partial_{\tau\prime}\); \(0\le l<h\) proves that these are precisely the invariant
holomorphic vectors, with no new integral coarse twist. At IV-star the
corresponding factor is \(z^{2k}\), with \(0\le 2k<m\). The extension splits by
the Laurent solution of \(H^1(\mathcal O(1))=0\). Hence \(H^1(V)=H^2(V)=0\).

\item The exponential sequence is on the full atlas/groupoid. At reduced II
the central group is additive and the deck derivative is \(\zeta_h^l\ne 1\),
so 1-D has its actual inverse there. Subtract an extending inverse image
of the central displacement; a remaining zero-central norm-zero section
Z is solved by \(-[h]^{-1}\sum_jjD^jZ\). The {[}h{]} inverse is the extending
formal-group inverse. The line lift passes through the actual K84
Cartier identity \(\beta^*N=\pi^*M\). Hom evaluation, followed by division by
its whole-fibre order, makes it invertible; the finite logarithm average
removes only the zero-central norm-one unit. The prescribed constant
character remains. At IV-star the original central three-kernel is
retained. This proves local atlas exponentiation, without asserting
invariant logarithms on a coarse root disc.

\item Redistribute the ordinary vectors by
\[
n_A'=n_A+m\sigma,\quad n_B'=n_B,\quad d_A'=d_A,\quad d_B'=d_B.
\tag{RE.9}
\]
The difference cocycle is \((-m\sigma\ell_0,0,L)\). Set \(a_A=-\sigma\ell_0\),
\(a_B=0\) in RE.4. Its contribution is \((-m\sigma\ell_0,0,\sigma\ell_0)\).
The remaining vector \(L-\sigma\ell_0\) lies in ker \(\psi\); use the negative
of its exact coefficients in RE.6, and recover p,q by the inverse
boundary-coordinate change. This gives the entire integral cochain,
including all \(\delta\) components. On the root the actual map is
\[
                         A_A'=B_\lambda^\sigma A_A.\tag{RE.10}
\]
The generated quotient groups coincide. Its scalar is exactly
\(r^{h(n_A+m\sigma)}=r^{hn_A}\lambda^\sigma\). The II root is unchanged.

\item After this integral correction all differences have Lie logarithms
with zero lattice discrepancies. Laurent/Cauchy splitting of the actual
V-cocycle and finite averaging only on V give local holomorphic
comparisons. They extend through both roots by the preceding maps and
through the cusp by its literal reindexing
\[
                    (j,v)=(3v-j)h_{\rm src}+(2v-j)B.\tag{RE.11}
\]
Its remaining multipliers and their inverses are extending units on
every same-fan cone, as proved by evaluating the actual open torus.

\item The global Isom line is handled by K78\textquotesingle s explicit construction. A
connection in local frames \(F_\alpha\) obeys the transformation law for
\(s=z_\alpha F_\alpha\). Parallel transport around an initial circle has
nonzero holonomy c; multiply by the specified \(e^{-\theta\operatorname{Log}c}\), with
the same endpoint-flat cut, to obtain a periodic nonzero initial section.
Along radial segments solve \(\dot z_\alpha=-\omega_\alpha(\dot r)z_\alpha\).
The transformation law makes the solutions agree. On a thin annulus
write \(\bar\partial s=\eta s\) and use the fixed-contour Cauchy-\/-Green solve
\(\bar\partial u=\eta\); \(e^{-u}\) s is a nonzero holomorphic Isom section. All
parameter period cochains and both finite characters remain in these
maps. On the cut write \(F_t:X_t\mapsto Y_t\). For the actual forward endpoints
\(E_X,E_Y\), the global scalar group proved below gives the typed square
\[
F_1E_XF_0^{-1}=C E_Y,\qquad
F_0(E_X^{-1}F_1^{-1}P_YF_0)F_0^{-1}=E_Y^{-1}C^{-1}P_Y.
\tag{RE.12}
\]
The scalar isotopy removes exactly C. This is an actual inverse-endpoint
return comparison, retaining literal parameter-period monodromy.

\par

\medskip\noindent\textbf{RE.3. Full holomorphic endpoints and the minimum scalar winding}

\par\nobreak

\item K79\textquotesingle s intrinsic base argument applies to these actual spheres. A horizontal
effective divisor would restrict to a nonzero effective curve on a generic
Kähler torus, with positive area, whereas \(H^2(X)=0\) makes its Chern class
zero. Thus every polar divisor is vertical; a meromorphic function is
constant on each generic compact torus. At a special base value use a
transverse disc \(\pi=u^e\), after the extending unit root. Its finite-principal-
part Laurent series is invariant under \(u\mapsto\zeta_eu\), since it is constant
on generic fibres. All exponents are divisible by e. This gives the actual
meromorphic base extension and \(\mathbb C(X)=\mathbb C(t)\). The distinct characteristic
polynomials of A,B,C force every holomorphic endpoint to fix the three
base values and hence the base. Its torus map is affine. The complete
integral centralizer is \(aI+c\delta\psi\), \(a=\pm 1\). A nonzero c would make
the difference from aI complex-linear of real rank one, impossible.

\item The reduced-II inverse retains the section collision. For \(N=O(G'-O')\)
and \(i(X,Y,z)=(X,-Y,z)\), its actual Cartier pairing is
\[
e_N\otimes i^*e_N\longmapsto f_N=z^{2l}X-(1+z^h)^2.
\tag{RE.13}
\]
Near the sections the regular frame is \(e_0=\tau'/(\tau'+z^l)e_N\).
The pairing coefficient is
\[
p(\tau',z)=\frac{\tau'^2}{\tau'^2-z^{2l}}
                 [z^{2l}X-(1+z^h)^2].\tag{RE.14}
\]
The second factor has exactly the divisor \([\tau\prime=-z^l]+[\tau\prime=z^l]-2[\tau\prime=0]\), with no vertical divisor. Thus p is an extending unit,
\(p(0,0)=-1\). Raw inversion sends the actual origin coefficient \(v_i\) to
\(1/(p(0,z)v_i)\); the extending scalar \(p(0,z)v_i^{2}\) makes it fix that
exact origin. This constructs the root inverse through all Cartier
charts. The good-root inverse and full fan inversion are K79\textquotesingle s actual
maps. The latter is torus inversion together with the central-symmetry
translation of the infinite triangle set.

\item For linear part -I the literal ordinary and finite-deck equations are
\[
s_i-L_{ij}s_j=\theta^{\rm end}_{ij}+\theta^{\rm start}_{ij},\qquad
s_{Dz}-L_Ds_z=a_D^{\rm end}+a_D^{\rm start}.\tag{RE.15}
\]
Their integral free obstruction is \(2D\ne0\) in RE.8. Such an endpoint does
not exist. A positive-linear automorphism lifts to the actual source/line
cover: a root discrepancy \(B^a\) has m-th or h-th power 1, and contraction
forces \(a=0\). The remaining suspension powers patch because the affine
base with the two roots has group \(C_m*C_h\) and \(\operatorname{Hom}(C_m*C_h,\mathbb Z)=0\).
Its source section extends by RE.11 and equals jG in the actual
Mordell-\/-Weil group. The rational line lift is \(f(t)U^j\).

\item Write \(U=z_i^{e_i}u_iL_i\) at every noncusp site, with \(L_i\) invertible.
The exact characters give \(e_A=m/3\) modulo m and \(e_B=0\) modulo h.
For \(\rho=\)\(B_0\)\(^{-1}/U^{3}\) the actual divisor identities are
\[
\operatorname{ord}_c\rho=1-3e_c,\quad
\operatorname{ord}_i\rho=-3e_i/r_i,\quad
                         e_c+\sum e_i/r_i=1/3.\tag{RE.16}
\]
Fixed-scalar invertibility gives \(\operatorname{ord}_if=-je_i/r_i\) and cusp order
\(v=j/3\). Hence \(j=3a\) and \(fU^{3a}=cB_0^{-a}\), literally the physical scalar
c \(\lambda^a\) on the quotient. Conversely that scalar extends everywhere.
Thus the complete holomorphic automorphism group is physical \(\mathbb C^*\).

\item For scalar winding \(\ell\) the actual period transport is \(I+\ell\delta\psi\).
It changes the full groupoid class by
\[
                     (0,\ell n_A\bmod k,\ell n_B\bmod h).
\tag{RE.17}
\]
Every ordinary \(\delta\) change is retained and removed by its specified
RE.6 cochain. Since the negative linear part is excluded, the affine
square requires both torsion entries zero. Sphere primitivity gives
\(kh|\ell\). Independently, the actual forward root square introduces
\(z_A^{\ell n_A}\) and \(z_B^{\ell n_B}\); hence
\[
m\operatorname{ord}_A f+je_A=-\ell n_A,\qquad
h\operatorname{ord}_B f+je_B=-\ell n_B.\tag{RE.18}
\]
The finite characters read \(kj+\ell n_A=0\) modulo m and \(\ell n_B=0\) modulo h.
They again give \(k|\ell\) and \(h|\ell\). Ordinary clutch i contributes
\(-\ell\psi(\ell_i)\) to its extending order. The rational degree equation is
\[
                v=j/3+\ell D/M,\qquad 3v-j=\ell\epsilon/(kh).
\tag{RE.19}
\]
This is the same constraint from the full cusp and both cap maps.

\par

\medskip\noindent\textbf{RE.4. The actual endpoint of winding \(kh\)}

\par\nobreak

\item Apply RE.9, so \(hn_A+mn_B=\epsilon\). Keep \(\lambda=r^M\) and put
\(f=t^{hn_A}(t-1)^{mn_B}\), \(\rho=\)\(B_0\)\(^{-1}/U^{3}\). The divisor of
f \(\rho^{-2\epsilon}\) is divisible by 3: at IV-star the residue is
\(hn_A+2\epsilon=0\) modulo 3, at II it is \(mn_B=0\) modulo 3, and at the
cusp its order is \(-3\epsilon\) modulo 3. The other \(\rho\) orders are whole-
fibre multiples of3. Degree-zero divisors on \(\mathbb P^1\) are principal, so choose
the actual rational cube root and endpoint-to-initial arrow
\[
g_0^3=f\rho^{-2\epsilon},\qquad
C_r=r^{2\epsilon kh}g_0(t)U^{-2\epsilon}:X_{\zeta_3r}\to X_r.
\tag{RE.20}
\]
At IV-star the actual normalized lift has \(D_AL_AD_A^{-1}=\zeta_3L_A\),
\(L_A^{3}=\)\(B_0\)\(^{-1}\). At reduced II the central section is \(G'_0=O\) and the
actual K84 line is trivial with deck weight 0. Its commutator character
is 1. Finite averaging of the zero-central norm-one logarithm gives a
commuting lift. Its cubic ratio to \(B_0\)\(^{-1}\) is a deck-invariant nonzero
base unit, and its invariant cube root with the specified central value
gives simultaneously
\[
                       [D_B,L_B]=1,\qquad L_B^3=B_0^{-1}.\tag{RE.21}
\]
All source translation lifts commute in the actual degree-zero line
extension on generic compact fibres, and equality extends through the
root. These operations retain both norm lifts and the given finite
characters.

\item Substituting \(U=z_i^{e_i}u_iL_i\) gives
\(g_0U^{-2\epsilon}=f^{1/3}L_i^{-2\epsilon}\), with every analytic unit root
retained. Its two root powers are \(z_A^{M n_A/3}\), \(z_B^{M n_B/3}\).
Under \(r\mapsto\)\(\zeta_3r\) the plus-log ordinary-to-cap map changes by exactly
the logarithm of that power, because \(\gamma\) changes by M \(\delta\)/3.
Thus the literal squares are
\[
\Phi_{i,r}C_r=C_{i,r}\Phi_{i,\zeta_3r},\qquad
C_{i,r}=(\text{retained extending unit})L_i^{-2\epsilon}.
\tag{RE.22}
\]
The unit includes \(r^{2\epsilon kh}\) and the specified root of the other
base factor. Conjugation at A gives \(\zeta_3\)\(^{hn_A+2\epsilon}=1\) and at B
\(\zeta_3\)\(^{mn_B}=1\). The fixed \(d_A,d_B\) commute and remain. Three consecutive
arrows compose exactly to f \(B_\lambda^{2\epsilon}\), including the successive
r phases. The second factor is an actual deck. This constructs the
proper holomorphic family \(w=r^{3}\), \(\lambda=w^{kh}\); its parameter action
is free since \(r\ne 0\).

\item For the full cusp closure put \(s=\log q/(2\pi i)\), \(v=\log r/(2\pi i)\),
\(a(q)=\log(1-q)^{mn_B}/(2\pi i)\), and \(K=\mathcal D+(s+2Mv)\delta\). Coherent
logs satisfy \(3L_U=-\gamma_B+(Mv-R)\delta\) and
\(3g_*=-\epsilon s+a-2\epsilon R\). Reindex C by \(h_{\rm src}^\epsilon\), whose
valuation is \((2\epsilon,\epsilon)\). Its full bounded logarithm is
\(\epsilon\mathcal D/3+(2\epsilon Mv/3+a/3)\delta\). Multiply the cube-root constant
by \(\exp(-4\pi i \epsilon kh/3)\). The closing is then exactly
\[
c(v)=\sigma(v+1/3)-\sigma(v),\qquad
\sigma(v)=\epsilon v\mathcal D+
                [\epsilon Mv(v-1)+v a(q)]\delta.\tag{RE.23}
\]
These are PLUS origin products; cusp-to-ordinary is \(o-\sigma\), and the
forward map is the inverse closing. The four successive seam logs are
\[
\begin{array}{c|l}
0&0\\
1&\epsilon\mathcal D/3+(2\epsilon Mv/3-2\epsilon M/9+a/3)\delta\\
2&2\epsilon\mathcal D/3+(4\epsilon Mv/3-2\epsilon M/9+2a/3)\delta\\
3&\epsilon\mathcal D+(2\epsilon Mv+a)\delta .
\end{array}\tag{RE.24}
\]
They are the actual products of the inverse forward maps. The last row
is the full canonical bounded closing, with RE.22 identifying both caps.

\par

\medskip\noindent\textbf{RE.5. Native pullbacks and every endpoint return}

\par\nobreak

\item Define the closed native products by RE.23, and their covered products
as the literal triple pullbacks RE.24. Every full-unit seam (A,D) has
native C0 lift \((a_0,a_1,c_0,c_1)\mapsto(a_0,Aa_1,(AD)^{-1}c_0,Dc_1)\), and the actual
K87 C1/C2 four-weight maps. They preserve q, both resolutions, physical
\(\delta\) and all normal jets. On any fan use the actual integral character
and period maps; therefore the native flop square commutes with all
three seams and both root squares.

\item The canonical full-r comparison is the exact K76 polynomial-origin map
with M; the retained term v a(q) occurs on both sides and cancels in
their difference. The remaining physical factor must be split literally:
\[
e^{-2\pi i\epsilon M(\theta^2+\theta)}
=e^{-2\pi i\epsilon M\theta(\theta-1)}
                          e^{-4\pi i\epsilon M\theta}.\tag{RE.25}
\]
Only the first logarithm is zero at both endpoints; multiplication by
1-s contracts it as a based loop. The second retains winding \(-2\epsilon M\).
Its actual fixed-star map is
\[
                         q_H\longmapsto q_H e^{4\pi i\epsilon M\theta}.
\tag{RE.26}
\]
Its left \(SU(2)\) factor is 1 and its right factor depends only on \(\theta\).
Thus both pullbacks of the invariant three-form to the native \(S^2\times S^1\)
are zero, giving zero degree-three free corrections. Its \(\operatorname{Spin}(4)\) endpoint
is 1; the paired base \(\operatorname{Spin}(2)\) endpoint is 1 since the circle winding is
even. This is a closed spin comparison. It acts along the actual Hopf
fibres and is exactly the physical-line scalar, so extends on that
Euler line, both finite cap lines and the fixed-star surgery cap.
Every later Cox square commutes with the same physical action. Carry
it at BOTH normal references and their connections; the common-link
and complex-normal reset transport the same zero degree-three comparison
to each target. No independently chosen target frame or omitted winding
is used. Hence K87 computes the actual covered native products.

\item Characteristic integrals multiply by the degree 3 of this actual cover;
the free native coordinates multiply by 3 in \(\mathbb Z^2\) before taking a sphere
class. Thus the shorter positive physical pair and class are
\[
(k_{M,w},\ell_{M,w})=(2kh,8k(3l+1)),\qquad A_w=kl.\tag{RE.27}
\]
On \(T_n\) its pair is \((2kh-6kh(n-1)^{2},8k(3l+1))\). The complete K87
algorithm uses \(M_w=kh\) and \(A_{0,w}=kl\); all full seed/target invariants
are the actual covered ones divided by 3 through the constructed
mapping-torus cover. No class in \(\Theta_7\) is divided. Denote its resulting
integer by \(S_{k,h}(P)\).

\item Necessity RE.17-\/-19 and existence RE.20-\/-24 prove minimum positive
scalar winding \(kh\). Iterate the constructed endpoint to realize every
\(\ell=Nkh\), using its actual inverse for negative N. Any other holomorphic
endpoint of that winding has positive linear part and differs by the
proved physical scalar. RE.12 and its actual scalar isotopy give return
\[
                              N S_{k,h}(P)\pmod{28}.\tag{RE.28}
\]
This carries every marked ordinary order-one clutch and every actual
fixed-source cap/line variation in the stated scope, retaining both
torsion entries, every integer representative and every finite kernel.

\item The fresh reconstructed certificate passes 10735 exact identities, 610
signed norm instances and 6390 scalar-obstruction tests, including the
full 16-by-26 Smith reduction, both cap characters, coherent seam
identities and the retained native cover maps. The lost original run
reported 10830 identities on the same 610 and 6390 instance counts;
that historical count remains separate. Fourfold\textquotesingle s independent
orchestration certificate passed 26666 checks over 398 norm tuples
for RE.4-\/-5. The fresh JSON records its own inputs and output hash.

\end{enumerate}
\end{proof}




\begin{theorem}\label{night:P3K.92}
\StatementRecord{night:P3K.92}{proved}{S6\_06\_BLIND: GPT-6 Astra (OpenAI, ChatGPT Pro mode, effort Pro), recorded re-derivation in the checkpoint dated 18 November 2026}{\cref{night:P3K.49}, \cref{night:P3K.53}, \cref{night:P3K.66}, \cref{night:P3K.71}, \cref{night:P3K.89}, \cref{night:P3K.90}}
\par\medskip\noindent\textbf{Statement.}\par\nopagebreak[4]

Retain either of the actual E7 one-cap families
\[
 (a,b,m)=(1,2,4p)\quad\hbox{or}\quad(2,1,4p-1),\qquad p\ge1,
\]
with either norm sign \(n=\pm1\), its prescribed finite character and
its full integral ordinary-to-cap reference cochain. Use the physical
circle and orientation of P3K.68. For \(P\in\mathrm{GL}_2(\mathbb Z)\),
let \(\mathcal T(P)\) be the two-triangle periodic fan with triangles
\(\{r,r+v,r+w\}\) and \(\{r+v,r+w,r+v+w\}\), where \(v,w\) are the
columns of \(P\). Put \(J=\operatorname{diag}(-1,1)\).

The marked infinity family is
\[
 y^2=x^3+(-2+3\epsilon^2)x^2+t x+
                  \epsilon^2(1-\epsilon^2)^2,
 \qquad G=(0,\epsilon(1-\epsilon^2)),                 \tag{IFan.1}
\]
with selected height
\[
 q_*=t-(1-\epsilon^2)(1-3\epsilon^2).
\]
The return uses one positive \(\epsilon\)-circle, a fixed sufficiently
small suspension scalar, and the actual source/line/root closing
specified in IFan.4 below. Let its oriented clutched sphere class be
\(I_a(P)\in\Theta_7\cong\mathbb Z/28\).

Let \(Z_a(P)\) be P3K.90's fixed-scalar source-zero class at the same
fan, and let \(C(P)\) be P3K.71's selected-cusp source class. Then
\[
 \boxed{\ I_a(P)=-Z_a(P)-Z_a(JP)-C(P)\pmod {28}.\ } \tag{IFan.2}
\]
This is a complete finite algorithm for every unimodular input, using
the actual algorithms of those two theorems. It is independent of the
cap order \(p\), norm sign, and retained fixed characters within the
stated branches because the two actual input algorithms have that
scope. No characteristic integer or free relative coordinate is
recovered from this equation in \(\Theta_7\).

For every vertical fan \(T_j=\left(\begin{smallmatrix}1&0\\j&1\end{smallmatrix}\right)\),
\[
 I_a(T_j)=1-a\pmod {28}.                            \tag{IFan.3}
\]
In particular these vertical classes are zero in the good branch and
\(-1\) in the reduced branch. On the nonvertical fan \(Q_0^2\), where
\(Q_0=\left(\begin{smallmatrix}1&-1\\0&1\end{smallmatrix}\right)\),
the formula gives respectively \(10\) and \(7\) modulo28. These are
classes of the actual marked infinity loop.

\end{theorem}
\begin{proof}
\begin{enumerate}[label=\arabic*.]
\item \par\medskip\noindent\textbf{IFan.1. A compact parameter family with the complete cap atlas.}\par\nopagebreak[4]

\item Start with the sectioned source
\[
 E_v:\quad Y^2=X^3+(1-2v)X^2+TX+v^2,\qquad G=(0,v),
 \quad q_1=T-(v^2-2v),\quad v=\rho^2-1.              \tag{IFan.4}
\]
Take a large finite \(R>2\) and remove disjoint small discs about
\(\rho=1,0,-1\) from \(|\rho|\le R\). This gives a compact surface
\(B_R\) with four boundary components. The two points above \(v=1/8\)
are filled by P3K.53's actual source/Cartier/root family. Its source
critical point is away from \(G\), its actual Hom map is a unit there,
and its selected-node charts persist. Thus these two discs introduce
no return.

\item The parameter discriminant is \(16v^2(v+1)^6(1-8v)^3\). On \(B_R\)
there are consequently no other excluded source degenerations. Choose
one cap annulus in the \(T\)-plane outside every finite critical value
for this compact parameter range. Its exact root equations are
\[
 y'^2=x'^3+(1-2v)z^{2p}x'^2+
 \begin{cases}x',&m=4p,\\zx',&m=4p-1,\end{cases}
 +v^2z^{6p},\qquad G'=(0,vz^{3p}).                   \tag{IFan.5}
\]
All resolution and Cartier corrections are the actual ones in P3K.49.
The central norm-ratio is the retained constant \(-2\) or \(+2\), so
its root and its logarithm with that reference extend on the root
disc times the bounded source-parameter disc. Finite averaging acts
on the line cochains while retaining their prescribed characters;
it preserves the literal equation \(A_0^m=B_0^n\).

\item Here is the full independent scalar family used in the argument:
\[
 B=r^mB_0(v),\qquad A=r^nA_0(v),\qquad A_0^m=B_0^n,
 \qquad
 F=F_0(v)+n\frac{\log r}{2\pi i}s_\infty\delta.      \tag{IFan.6}
\]
The map \(F\) is ordinary-to-cap. The fixed reference \(F_0\) includes
the actual integral annular period, every finite character and the
reduced doubled period. It is continued on the chosen cap annulus by
the parameter-disc exponential lifting and the fixed-contour Cauchy
splitting of P3K.49. The annular source is regular throughout the
bounded parameter disc. Its full source and line maps therefore
continue, rather than only their generic Picard classes.

\item On the overlap put \(q_1=q_\infty^{-1}h_\infty\),
\(h_\infty=1-(v^2-2v)q_\infty\). A positive scalar-root turn has the
actual forward maps \(R_{\rm ord}=q_1^n\) and
\(R_{\rm cap}=h_\infty^n\) on the physical line. Their exact square is
\[
 F_{\rm end}R_{\rm ord}=R_{\rm cap}F_{\rm start}.     \tag{IFan.7}
\]
The cap unit has its full extending logarithm; it is not replaced by
its central value. The actual parameter return uses the inverse of
the forward closing map after transport. In this theorem we take
\(r=r_0\) constant, \(r_0^m=\lambda_0\). Both scalar forward maps for
a loop in \(B_R\) are then the identity. The nonconstant source and
finite-root maps in \(F_0(v)\) remain in (IFan.6).

\item Choose \(\lambda_0\) in the common contracting range over compact
\(B_R\). All coefficients and their inverses have finite bounds in
the source/cap atlases. The norm shells and finite root quotients of
P3K.49 give the corresponding uniform proper quotient. This argument
uses a compact parameter surface; it does not claim one scalar works
at all arbitrarily small infinity radii. For two radii, the compact
annulus between them admits a common smaller scalar, and the actual
independent \(r\)-family gives the connecting zero-winding amplitude
path. This identifies the two loop returns by actual family transport.

\par\medskip\noindent\textbf{IFan.2. Every cone, full coefficient, and physical circle.}\par\nopagebreak[4]

\item The selected normalization on (IFan.4) is
\[
 X=s^2-1,\quad Y=s(s^2-\rho^2),\qquad
 Z=\frac{s-\rho}{s+\rho},\qquad
 P_G=\frac{1+\rho}{1-\rho}.                          \tag{IFan.8}
\]
The actual Tate period is \(Q=q_1\kappa\), with
\[
 \kappa(0,\rho)=-\frac{\rho^2-1}{64\rho^6}.
\]
Both \(P_G\) and the full \(\kappa\) are nonzero on the compact
parameter surface after taking a uniform selected disc. The chosen
nodal fourth root is \(2\rho\). The source branch marking is thus
single valued. A different lift of \(G\) by \(Q^j\) would have nodal
limit zero or infinity for \(j\ne0\), contrary to (IFan.8), so the
integer lift is fixed. The rigidified theta/Cartier line map supplies
the actual physical coordinate \(V\), with its degree-two correction
when \(a=2\).

\item For an absolute cone \(\sigma\) of the periodic fan, its affine
coordinate ring is generated by the characters
\(Z^uV^v q_1^h\) in the dual integral semigroup. Every change of period
basis carries the four full coefficients by
\[
 U'_{is}=\prod_{h,j}U_{hj}^{(P^{-1})_{ih}P_{js}}.     \tag{IFan.9}
\]
The inverse substitution uses \(P^{-1}\). This carries all higher
\(q_1\) terms, every source-parameter unit and the exact closing.
Multiplying an absolute character by a nowhere-zero base unit to an
integer power defines an automorphism on this actual affine chart;
its inverse retains the reciprocal unit. Therefore all these maps
extend across the toric divisors, and they agree on intersecting
cones as the same character maps. No unit logarithm is used to erase
parameter winding.

\item The resulting family has exactly the P3K.90 periodic chart atlas at
each source-zero boundary and the P3K.71 atlas at the selected-cusp
boundary. A finite set of cone orbits and the uniform norm bounds
give a compact fundamental region. The ordinary exterior and finite
caps are glued by (IFan.6), unchanged outside the selected disc. The
physical action is always \(V\mapsto e^{i\phi}V\). Thus the family is
smooth and proper over \(B_R\), with that actual action. Each fibre
is the completed oriented sphere of P3K.90: its proof constructs every
fan by the actual two-ruling replacements with the common-blowup
integral Gysin and primitive fundamental-group maps. The same full
coefficient maps (IFan.9) make those replacements relative to this
parameter surface.

\item The branch reflection is an especially explicit cone map:
\[
 \mathcal R_\rho:(Z,V,q_1)\longmapsto(Q/Z,V,q_1),
 \qquad P_G\longmapsto P_G^{-1}.                    \tag{IFan.10}
\]
It is the identity on the underlying marked source and Cartier
divisor, expressed in its two branch coordinates. Its line square
is the literal theta identity
\[
 \frac{\Theta_Q(QP_G/Z)}{\Theta_Q(Q/Z)}
       =\frac{\Theta_Q(Z/P_G)}{\Theta_Q(Z)}.
\]
For the actual \(M_a\), its canonical tensor/Cartier correction is
carried as well. On each character the pullback is exactly
\[
 Z^uV^v q_1^h\longmapsto
          \kappa^u Z^{-u}V^v q_1^{h+u}.             \tag{IFan.11}
\]
This is the integral dual-semigroup map of
\((x,y,h)\mapsto(-x+h,y,h)\), including its height translation. Its
square is the identity. It maps the entire fan \(\mathcal T(P)\) to
\(\mathcal T(JP)\), with the indicated integer translate. The ordinary,
cap and \(F_0\) maps depend on \(v\) and are carried identically.
It preserves the physical circle and complex orientation, including
the actual normal and line maps on every cone. More explicitly, write
the full periods as
\[
 H(Z,V)=(QZ,\alpha V),\qquad
 B(Z,V)=(gZ,q_1\eta V).
\]
The target coefficients are \((Q,\alpha^{-1},g^{-1},\eta)\), so the
literal maps satisfy
\[
 \mathcal R_\rho H=H_{\rm target}^{-1}\mathcal R_\rho,
 \qquad \mathcal R_\rho B=B_{\rm target}\mathcal R_\rho.
\]
These identities retain the entire coefficients, not only their
leading valuations, and prove equivariance with the full period
quotient and its physical line.

\par\medskip\noindent\textbf{IFan.3. The three finite local returns with their labels.}\par\nopagebreak[4]

\item At \(\rho=1\) the parameter \(v=\rho^2-1\) has positive complex
local degree one. The complete source, Cartier, root and \(F_0\)
markings above are those of P3K.90, so the positive return is
\(Z_a(P)\). At \(\rho=-1\) the derivative is \(-2\), again positive
as a real oriented complex local map. Equation (IFan.10), whose
ordinary/cap map is literally the identity, identifies that return
with \(Z_a(JP)\). It makes no identification of these two classes.

\item Near \(\rho=0\), choose the holomorphic cube root \(u^3=v\) with
\(u(0)^3=-1\), and put
\[
 X=u^2x,\quad Y=u^3y,\quad T=u^4t,\qquad
 \xi=x-u,\quad \sigma=\rho/u.
\]
The exact selected equation and its actual base map are
\[
 y^2=\xi^2(\xi+\sigma^2),\qquad
 q_{\rm old}=\frac{y^2-\xi^2(\xi+\sigma^2)}{\xi+u},
 \qquad q_1=u^4q_{\rm old}.                         \tag{IFan.12}
\]
The normalized source coordinate is \(s/u\), so its marked Tate
coordinate is exactly \((s-\rho)/(s+\rho)\). There is no branch
reflection at this interface. The local parameter \(\sigma\) has
nonzero complex derivative \(u(0)^{-1}\), hence preserves the
positive source loop. The divisor-defined line map scales the
primitive lift by \(u^{2a}\); the full suspension contains the
height \(u^{-4}\), so \(2ab-4=0\) keeps its scalar unchanged.
All cap-root factors of \(u\) have logarithms on this disc and carry
the full finite character. The degree-two Cartier map is conjugated
on its source and target, with both factors retained. These are the
actual local source/line/root maps of P3K.49 and the selected-source
maps of P3K.71. Their coefficient substitutions (IFan.9) extend on
every cone, so the third finite return is exactly \(C(P)\).

\item The unselected points over \(v=1/8\) have the actual P3K.53 disc
extension. Their source, Cartier and root maps are the identity on
the selected fan overlap; thus this disc extension also carries
every fan in (IFan.9), and its return is zero.

\par\medskip\noindent\textbf{IFan.4. Infinity: the full rescaling, root endpoint and annular lift.}\par\nopagebreak[4]

\item On the outer boundary put \(\epsilon=1/\rho\). The literal source
map is
\[
 X=\epsilon^{-2}x,\quad Y=\epsilon^{-3}y,\quad
 T=\epsilon^{-4}t,\quad v=\epsilon^{-2}-1,
 \qquad q_1=\epsilon^{-4}q_*.                       \tag{IFan.13}
\]
It gives (IFan.1). With \(s_{\rm new}=\epsilon s_{\rm old}\),
\[
 Z_{\rm old}=\frac{s_{\rm old}-\rho}{s_{\rm old}+\rho}
       =\frac{s_{\rm new}-1}{s_{\rm new}+1}
       =Z_{\rm new}.                               \tag{IFan.14}
\]
Consequently \(Q\) and \(P_G\) agree, and the selected fan remains
\(\mathcal T(P)\). The unit relation is
\(\kappa_{\rm old}=\epsilon^4\kappa_{\rm new}\), with
\(P_G=-(1+\epsilon)/(1-\epsilon)\) and
\(\kappa_{\rm new}(0)=-(1-\epsilon^2)/64\).
All integer \(\epsilon\)-powers are retained as actual units on the
parameter annulus. In particular their winding is not discarded.
In the displayed period notation the exact coefficient relations are
\(\kappa_{\rm new}=\epsilon^{-4}\kappa_{\rm old}\),
\(\eta_{\rm new}=\epsilon^{-4}\eta_{\rm old}\), while
\(\alpha\) and \(g\) are unchanged. Therefore both full periods
\((QZ,\alpha V)\) and \((gZ,q\eta V)\) agree in the transported
Tate/Cartier coordinates.

\item The line map is the canonical divisor map: on a local equation
\(f_D\) it carries the actual section \(f/f_D\) to its pullback in
the target divisor chart. Its coordinate factor is precisely the
ratio of the two pulled-back divisor equations. This definition
carries the degree-two Cartier correction and supplies the line
unit on every source chart. In the rigidified Tate frames its source
point and period are those in (IFan.14). Its primitive lift satisfies
\[
 I U_{a,\rm old} I^{-1}=\epsilon^{2a}U_{a,\rm new}.
\]
Therefore the complete suspension obeys
\[
 I\bigl(2(-1)^{b+1}\lambda q_1U_{a,\rm old}^{\,b}\bigr)I^{-1}
 =2(-1)^{b+1}\lambda q_*U_{a,\rm new}^{\,b}.          \tag{IFan.15}
\]
The actual Hom section checks this without omitting the height:
\(F_{ab,\rm old}=\epsilon^{-1}F_{ab,\rm new}\) and
\(v=\epsilon^{-3}g\), where \(g=\epsilon(1-\epsilon^2)\), give
\[
 \frac{q_1}{vF_{ab,\rm old}}
       =\frac{q_*}{gF_{ab,\rm new}}.                \tag{IFan.16}
\]
Any local Cartier unit is used with its reciprocal in the inverse
map and with its actual integer powers in all cone characters.
The full morphism has no additional fan reflection.

\item On the parameter logarithm cover the finite-root map is
\[
 z_{\rm old}=\epsilon^{4/m}z_{\rm new}.              \tag{IFan.17}
\]
The good minimal source coordinates are unchanged. In the reduced
root the old-to-new factors are \(\epsilon^{-2/m}\) and
\(\epsilon^{-3/m}\). Thus its positive endpoint is the actual
\(D^{-4}\), including \((\zeta_m^{-4},\zeta_m^{-2},\zeta_m^{-3})\)
on \((z,x',y')\) in the reduced case. The resolution ideals and
Cartier factors are carried by these coordinate maps.

\item Write \(A_0=L D\) for the actual signed division lift, with \(L\)
covering translation by \(-nR\). The good
central division point is \(O\). In the reduced root the actual line
is trivial on the compact III fibre, and both commuting source maps
have constant line multipliers there. The finite commutator character
is thus one. The prescribed eigenspace averaging retains its
conjugating unit and gives \([D,L]=1\) exactly. Hence
\[
 A_0^4D^{-4}=L^4.                                  \tag{IFan.18}
\]
The full finite action in (IFan.6) is \(A=r_0^n A_0\), so the
actual representative of this endpoint in the quotient is
\[
 A^4D^{-4}=S_{r_0^{4n}}L^4,                         \tag{IFan.18a}
\]
where \(S_c\) is multiplication by \(c\) on the physical line. This
constant scalar is retained. If the comparison along the logarithmic
parameter cut has endpoint \(H_1=D^{-4}H_0\), its endpoint-to-initial
correction is
\[
 G_1=S_{r_0^{-4n}}L^{-4},\qquad
 G_1(A^4D^{-4})=\mathrm{Id}.                        \tag{IFan.18b}
\]
The displayed equality is literal before passing to the quotient;
\(A^4\) is its specified deck map. Fix the same reference logarithm
\(\log r_0\) as in (IFan.6). The scalar part of \(G_1\) is carried
by the explicit based path \(S_{\exp(-4n\theta\log r_0)}\),
\(0\le\theta\le1\), on both ordinary and cap charts. This is a path
from \(1\) to \(r_0^{-4n}\), not a closed loop; these endpoints
implement the prescribed inverse closing. Its square is
\(F S_c=S_c F\) for every \(c\), including the full root/line and
normal maps. On the parameter logarithm cover the corresponding
holomorphic scalar cochain is
\(S_{\exp(-4n t_\epsilon\log r_0)}\), with one-step ratio
\(S_{r_0^{-4n}}\). It is part of the endpoint atlas, not asserted to
be a single-valued scalar on an unglued annulus.

\item The source translation of \(L^4\) specializes to the identity. Its
own central line scalar is likewise retained before taking the small
extending Lie logarithm. The inverse of that logarithm supplies the
remaining \(L^{-4}\) in (IFan.18b); its ordinary-to-cap comparison
is the full relative cochain described below. In the reduced branch the \emph{annular} cochain still
contains the nonzero marked period
\[
 4\log R-\log(4R)=2\omega_1+2\omega_2,              \tag{IFan.19}
\]
where \(\omega_1,\omega_2\) are the actual transported source period
sections, with central square-lattice values \((1,i)\). Thus its
integral vector is \((2,2)\) and its central value is \(2(1+i)\);
the nearby period functions are retained. Its signed source part is
multiplied by \(-n\), since \(L\) covers \(t_{-nR}\), and the full Cartier factor is
carried at the same time. No logarithm of \(L^4\)
replaces that lattice cochain.

\item Fix a nonzero \(\epsilon_*\) and transport the entire original
ordinary-to-cap map there through (IFan.13--17). This defines the
reference infinity overlap, including its integral \(d\delta\).
Continue it by the parameter-disc exponential lifting and fixed
source-annulus Cauchy contours. Every integral or finite discrepancy
is locally constant and is zero at the full reference map; the two
atlases carry the same term (IFan.19). Remove an actual suspension
power before taking a Lie logarithm at a reduced component. Average
only the vector cochains over the finite group, and use their actual
Cartier exponentials. This constructs the same full \(F_0\) on both
sides, rather than choosing another annular clutch.

\item For completeness the scalar cochain of a general coupled reference
is also retained. With \(t_\epsilon=\log\epsilon/(2\pi i)\), the
old and new forward arrows for \(\lambda=v^{km}\lambda_0\) differ
by \(\epsilon^{8nk}\). Their compensating cochain is
\[
 c_0(t_\epsilon)=\exp[-2\pi i\,4nk\,t_\epsilon(t_\epsilon-1)],
 \qquad c_0(t_\epsilon+1)/c_0(t_\epsilon)=\epsilon^{-8nk}.
                                                               \tag{IFan.20}
\]
The full \((1-\epsilon^2)^k\) unit and the cap \(h_\infty\) unit
remain in the same square (IFan.7). In the fixed-scalar theorem
\(k=0\), this displayed correction is exactly one.

\item Finally the local comparisons on two angular annulus charts differ
by an actual extending translation. The source/line valuation
argument in P3K.49 identifies it with the physical scalar. For its
retained overlap logarithm \(\ell\), take
\(b_+=\chi_-\ell\), \(b_-=-\chi_+\ell\), solve
\(\bar\partial w=\bar\partial b_\pm\) by the fixed Cauchy kernel,
and use \(e^{b_\pm-w}\). The two maps glue on both overlap
components, with the integral winding and finite cocycle retained.
This includes both the scalar cochain in (IFan.18b) and the lattice
cochain in (IFan.19). It gives an actual single-valued annular family morphism, preserves
\(V\)'s physical circle, and intertwines the inverse endpoint
closings. It identifies the infinity return with \(I_a(P)\) of the
statement.

\item The separate half-turn map is
\(\mathcal I_\epsilon(Z,V)=(1/Z,P_G^aV)\). At its target
\(P_G\) is inverted, so \(h_{\rm target}\mathcal I_\epsilon=(Q/Z,V)\).
This is precisely (IFan.10). It does not modify (IFan.14); confusing
these two arrows would insert an erroneous \(J\) in (IFan.2).

\par\medskip\noindent\textbf{IFan.5. The actual loop product and the sphere class.}\par\nopagebreak[4]

\item Choose the lower-right corner of a rectangle containing the three
holes as basepoint and use disjoint vertical slits to their lower
points. The boundary of the cut surface gives the based positive
meridians, traversed in the displayed word order,
\[
 \gamma_+\gamma_0\gamma_-\gamma_\infty=1.            \tag{IFan.21}
\]
The positive infinity meridian is clockwise in \(\rho\) and positive
in \(\epsilon=1/\rho\). The coordinate maps above all have positive
complex local degree. Thus no boundary parameter is silently reversed.

\item Use a smooth connection in the actual proper sphere family, averaged
over the physical circle. Retain its slit transports as maps, not
only fibrewise conjugacies. On a lifted path the actual return is
\(R_{\rm end}^{-1}P_\gamma\), using the forward maps in (IFan.7).
For fixed \(r\) their total scalar winding is zero. More generally
the coupled reference has windings \(1,0,1,-2\), whose full forward
maps multiply to \(R_0=1\). The contraction of the cut-disc boundary
lifts to an isotopy
\[
 \Phi_\infty\Phi_-\Phi_0\Phi_+\simeq\mathrm{Id}.    \tag{IFan.22}
\]
This isotopy uses every cap and selected-cone endpoint already
constructed. It is an isotopy of the actual oriented sphere returns.

\item The oriented clutch identification with \(\Theta_7\) is additive
under composition. Its group is abelian, so the retained slit-path
conjugations have no effect after taking the class. The three finite
terms in IFan.3 and the infinity term in IFan.4 give
\[
 Z_a(P)+C(P)+Z_a(JP)+I_a(P)=0.
\]
This proves (IFan.2). It invokes the actual family morphisms before
passing to the group; it does not divide a covered sphere class.

\par\medskip\noindent\textbf{IFan.6. Vertical simplification and exact examples.}\par\nopagebreak[4]

\item The basis identity
\[
 P_{1-j}=J P_j S,\qquad
 S=\begin{pmatrix}-1&0\\1&1\end{pmatrix}
\]
is an identical-fan arrow. Its two triangles are translates of the
matched original triangles, and the identity on absolute characters
extends on their identical semigroups. This is one of the exact
column-swap/rotation arrows of P3K.71 and P3K.90. Thus the reflected
vertical fan is exactly \(T_{1-j}\), including its height translate.

\item Put
\[
 P_b(j)=2b\binom{j+1}{3}+(3-b)\binom{j+1}{2},\qquad b=3-a.
\]
P3K.89 gives \(Z_a(T_j)=P_b(j)\), and P3K.66 gives
\(C(T_j)=-1-3j(j-1)\), as actual classes with their full source
and cap arrows. Polynomial expansion gives
\[
 P_b(j)+P_b(1-j)=3j(j-1)+a.
\]
Substitution in the proved groupoid relation yields (IFan.3).
This is a statement in \(\Theta_7\); no individual infinity
characteristic pair is claimed.

\item For \(P=Q_0^2\), the two actual finite algorithms give
\[
 \begin{array}{c|rrr|r}
 a&Z_a(P)&Z_a(JP)&C(P)&I_a(P)\pmod{28}\\\hline
 1&6&31&-19&10\\
 2&9&59&-19&7
 \end{array}
\]
The middle entries are the algorithms' integer path representatives
before taking residues. They are not asserted to be infinity
characteristic integers. All other inputs are evaluated by the same
finite exact source-zero and selected-cusp algorithms in (IFan.2).

\end{enumerate}
\end{proof}




\begin{theorem}\label{night:P3K.93}
\StatementRecord{night:P3K.93}{proved}{S6\_06\_BLIND: GPT-6 Astra (OpenAI, ChatGPT Pro mode, effort Pro), recorded re-derivation in the checkpoint dated 18 November 2026; also checked in this writing by \CurrentChecker}{\cref{night:P3K.71}, \cref{night:P3K.89}, \cref{night:P3K.90}, \cref{night:P3K.92}}
\par\medskip\noindent\textbf{Statement.}\par\nopagebreak[4]

Retain all source, Cartier, cap, root, scalar and endpoint markings of
P3K.92. For any unimodular fan basis \(P\), put
\[
 (x,y)^t=P^{-1}(0,1)^t,\qquad S=x^2+xy+y^2.
\]
Then its actual fixed-scalar positive infinity return is
\[
 \boxed{\displaystyle
 I_a(P)=f_a(S):=
 1-a-3(a-1)(S-1)-\frac{2a+1}{4}(S-1)(S-3)
 \pmod {28}.}                                      \tag{IC.1}
\]
Here \(a=1\) is the good \(m=4p\) branch and \(a=2\) the reduced
\(m=4p-1\) branch, for every \(p\ge1\), either norm sign and every
fixed admissible finite character and integral annular cochain. The
integer \(S\) is odd, so the displayed expression is integral.

The complete images over all lattice-basis fans are
\[
 \begin{array}{c|l}
 a=1&\{0,6,8,10,14,20,22,24\}\\
 a=2&\{5,7,9,13,19,21,23,27\}.
 \end{array}                                       \tag{IC.2}
\]
The proof below evaluates the sphere class in the actual P3K.92
family. It assigns no infinity characteristic pair or free relative
coordinate from that class.

\end{theorem}
\begin{proof}
\begin{enumerate}[label=\arabic*.]
\item \par\medskip\noindent\textbf{IC.1. Canonical fans, their norm, and the reflected word.}\par\nopagebreak[4]

\item Use the literal identical-fan arrows of P3K.71 and P3K.90: a column
swap, multiplication on the right by
\(R=\left(\begin{smallmatrix}-1&-1\\1&0\end{smallmatrix}\right)\),
and central negation match the actual translated triangles. The
identity on their absolute character semigroups is the family map;
it carries all full period coefficients, original source/Cartier
charts, finite caps and endpoint closings. On \((x,y)\) these
operations are respectively
\[
 (y,x),\qquad (y,-x-y),\qquad(-x,-y).
\]
Each preserves \(x^2+xy+y^2\) by direct expansion. Thus \(S\) is a
function of the actual fan.

\item The canonical representative is either
\(T_j=\left(\begin{smallmatrix}1&0\\j&1\end{smallmatrix}\right)\),
or
\[
 P=\begin{pmatrix}u&-r\\s&z\end{pmatrix},\qquad
 r,u>0,\quad uz+rs=1.                              \tag{IC.3}
\]
In the second case \((x,y)=(r,u)\). Its actual construction is the
finite word from a vertical seed using
\[
 Q_0=\begin{pmatrix}1&-1\\0&1\end{pmatrix},\qquad
 Q_1=\begin{pmatrix}1&0\\-1&1\end{pmatrix}.
\]
The first step is \(Q_0\). The subsequent Euclidean word is the one
proved in P3K.90; all its three physical weights are nonzero.

\item Let \(J=\operatorname{diag}(-1,1)\) and
\(E=\left(\begin{smallmatrix}0&1\\1&0\end{smallmatrix}\right)\).
For nonvertical \(P\), the canonical reflected basis is literally
\[
 P^*=JPE=\begin{pmatrix}r&-u\\z&s\end{pmatrix}.
\]
The column swap is an identical-fan map, while \(J\) is exactly the
source/Cartier reflection with full height translate proved in
P3K.92. Thus the second source-zero class in that theorem is
\(Z_a(P^*)\). Exact matrix multiplication gives
\[
 (PQ_0)^*=P^*Q_1,\qquad (PQ_1)^*=P^*Q_0,             \tag{IC.4}
\]
and at the first step
\[
 (T_jQ_0)^*=T_{1-j}Q_0.                             \tag{IC.5}
\]
Both reflected algorithms therefore start with the actual common
\(Q_0\) step, and thereafter interchange the two generators. These
are identities of full period-basis/character maps; their coefficient
substitution is the same integral monomial formula as in P3K.92.
No new normal identification is introduced in evaluating the arrays
below.

\par\medskip\noindent\textbf{IC.2. The full paired prism equations.}\par\nopagebreak[4]

\item Write \(b=3-a\) and
\[
 W=\begin{pmatrix}1&b\\-a&-2\end{pmatrix}.
\]
For the source-zero algorithm at \(P\), the full winding matrix is
\(B_Z=2P^{-1}WP\), with product trace \(t_Z=-1\). For the
selected-cusp algorithm it is
\(B_C=-6P^{-1}E_{11}P\), with \(t_C=-3\). These are the actual
integer data of P3K.90 and P3K.71, whose complete coefficient maps
and native products are retained in P3K.92.

\item For any such \((B,t)\), their common prism notation is
\[
 \begin{split}
 \ell&=(B_{22}-B_{21},B_{11}-B_{12},B_{22}),\\
 w&=(r,u,-r-u),\qquad d=(-r-2u,-2r-u,r-u),\\
 I_i&=t w_i+d_i(\ell_i-t),\qquad
 J_i=t w_i-d_i(\ell_i-t).
 \end{split}                                      \tag{IC.6}
\]
The source-zero and selected-cusp variation/derivative arrays satisfy
\[
 V(w_i e_i-w_0e_0)=I_i e_i-I_0e_0,\qquad
 w_i\delta_i-w_0\delta_0=J_i-J_0,\quad i=1,2.        \tag{IC.7}
\]
Let \(K\) exchange indices 0 and 1 and fix index 2. Form the following
explicit algebraic sums of the already computed arrays:
\[
 \mathcal V=V_Z(P)+K V_Z(P^*)K+2V_C(P),\qquad
 \mathcal D=\delta_Z(P)+K\delta_Z(P^*)+2\delta_C(P).
                                                               \tag{IC.8}
\]
Their role is to evaluate the right-hand side of P3K.92's formula.

\item The full winding matrices obey
\[
 B_Z+E B_Z(P^*)E+2B_C
   =2P^{-1}(W+JWJ)P-12P^{-1}E_{11}P=-8\operatorname{Id}.
                                                               \tag{IC.9}
\]
The traces in (IC.6) add to \(-1-1+2(-3)=-8\). For indices 0,1,
reflection interchanges \(w,d,\ell\) in their displayed order.
For index 2 it changes the sign of \(d_2\); the corresponding
\(\ell_2-t\) also changes sign because
\(\operatorname{tr}B_Z=-2=2t_Z\). Consequently the full combined
prism right-hand sides are exactly
\[
 \mathcal I_i=\mathcal J_i=-8w_i.                   \tag{IC.10}
\]
This includes index 2; no unlabelled permutation of its sign is used.

\item All three weights are nonzero after the first step. The complete
symmetric solution of (IC.7) with right-hand side (IC.10) is
\[
 \mathcal V_{ij}=-8\delta_{ij}+
                \frac{\mathfrak c}{w_iw_j},\qquad
 \mathcal D_i=-8+\frac{\mathfrak e}{w_i}.            \tag{IC.11}
\]
Indeed \(\mathcal V+8\operatorname{Id}\) annihilates the plane
spanned by \(w_i e_i-w_0e_0\). A symmetric matrix annihilating that
plane is a scalar multiple of the outer product of
\((1/w_0,1/w_1,1/w_2)\). The derivative equations similarly give
one scalar. Thus their selected diagonal and selected derivative
entries determine \(\mathfrak c,\mathfrak e\) uniquely.

\par\medskip\noindent\textbf{IC.3. The complete common first-step seed.}\par\nopagebreak[4]

\item Put
\[
 P_b(j)=2b\binom{j+1}{3}+(3-b)\binom{j+1}{2},\qquad
 \mu=2-a,\qquad \nu_j=6P_b(j).
\]
P3K.89 and P3K.90 give the actual source-zero fixed ordinary entries
on \(T_j\)
\[
 c_Z(j)=-2-2\mu-2\nu_j,\qquad
 e_Z(j)=-2-2\mu+2\nu_j.
\]
P3K.71 gives the selected-cusp entries
\[
 c_C(j)=18j(j-1),\qquad e_C(j)=-12-18j(j-1).
\]
The actual \(Q_0\) Cox arrow and complex-meridian reset send each
selected old pair \((c,e)\) to target entries \((c,2c-e)\) on
\(C_2\). This is the same typed two-ruling map in all three
algorithms. The source-zero class increment is \((e-c)/4\), and
the selected-cusp increment is \((e-c)/2\), exactly as proved by
their respective actual doubled products.

\item The identity
\(P_b(j)+P_b(1-j)=3j(j-1)+a\) now gives the selected combined
entries after this common first step:
\[
 \begin{split}
 \mathcal V_{22}
   &=c_Z(j)+c_Z(1-j)+2c_C(j)=-4(2a+3),\\
 \mathcal D_2
   &=(2c_Z-e_Z)(j)+(2c_Z-e_Z)(1-j)
                    +2(2c_C-e_C)(j)=-4(8a-3).
 \end{split}                                      \tag{IC.12}
\]
The target weights are \((1,1,-2)\). Substitution into the unique
full prism solution (IC.11) gives the complete arrays
\[
 \mathcal V=
 \begin{pmatrix}
 -8(4a+3)&-16(2a+1)&8(2a+1)\\
 -16(2a+1)&-8(4a+3)&8(2a+1)\\
 8(2a+1)&8(2a+1)&-4(2a+3)
 \end{pmatrix},\qquad
 \mathcal D=
 \begin{pmatrix}16(4a-3)\\16(4a-3)\\-4(8a-3)\end{pmatrix}.
                                                               \tag{IC.13}
\]
In particular this is a full common seed, independent of \(j\),
obtained from the actual three target systems.

\item The vertical groupoid class is \(1-a\) by P3K.92. The total of the
three finite increments at the first step is
\[
 \nu_j+\nu_{1-j}+\frac{e_C(j)-c_C(j)}2=6(a-1).
\]
Therefore the actual infinity class after this step is represented
by \(7(1-a)\). This agrees with \(f_a(3)\); the vertical seed agrees
with \(f_a(1)\).

\par\medskip\noindent\textbf{IC.4. Both generator identities and the complete induction.}\par\nopagebreak[4]

\item For positive \(r,u\) define
\[
 D=ru(r+u),\qquad S=r^2+ru+u^2,
\]
and set
\[
 \mathfrak c=-4(2a+1)D^2,\qquad
 \mathfrak e=4D\bigl((2a+1)S+2a-8\bigr).            \tag{IC.14}
\]
At \((r,u)=(1,1)\), these give exactly (IC.13).

\item The two actual forward generators change the weights as follows:
\[
 \begin{array}{c|c|c|c}
 i&(r',u')&w_i&w'_2\\\hline
 0&(r+u,u)&r&-r-2u\\
 1&(r,r+u)&u&-2r-u.
 \end{array}
\]
Direct multiplication gives both identities
\[
 S'=S+\frac{2D}{w_i},\qquad
 \frac{D'}{w'_2}=-\frac{D}{w_i}.                   \tag{IC.15}
\]
By (IC.4) the reflected source-zero step uses the other generator,
while the selected-cusp step uses the original one. Each has the
same actual selected-entry reset. Hence the algebraic sums obey
\[
 \mathcal V'_{22}=\mathcal V_{ii},\qquad
 \mathcal D'_2=2\mathcal V_{ii}-\mathcal D_i.         \tag{IC.16}
\]
The target full prism equations remain (IC.10).

\item Substitute (IC.14--15) into (IC.11) at the target. Its selected
variation is
\[
 -8-4(2a+1)\frac{D'^2}{(w'_2)^2}
   =-8-4(2a+1)\frac{D^2}{w_i^2}=\mathcal V_{ii}.
\]
Its selected derivative is
\[
 -8-\frac{4D}{w_i}
       \left((2a+1)\left(S+\frac{2D}{w_i}\right)+2a-8\right)
   =2\mathcal V_{ii}-\mathcal D_i.
\]
The uniqueness proved in IC.2 therefore gives the entire target
arrays, with (IC.14), for either generator. This proves their formula
along every finite Euclidean word.

\item The change of the \emph{actual class sum} in P3K.92 is
\[
 \begin{split}
 \Delta I_a
  &=-\frac{(e_Z-c_Z)+(e_{Z^*}-c_{Z^*})+2(e_C-c_C)}4\\
  &=-\frac{\mathcal D_i-\mathcal V_{ii}}4
    =-\frac{(2a+1)D^2}{w_i^2}
      -\frac{D\bigl((2a+1)S+2a-8\bigr)}{w_i}.
 \end{split}                                      \tag{IC.17}
\]
Every quotient by 4 here is the sum of the already proved integral
native increments. No sphere class is divided by 4. Expanding the
polynomial in (IC.1) and using (IC.15) gives exactly
\[
                 \Delta I_a=f_a(S')-f_a(S).         \tag{IC.18}
\]
The first-step and vertical seeds in IC.3 agree with \(f_a\), and
every canonical nonvertical fan is reached by the terminating word.
Thus (IC.1) follows for all canonical fans. The identical-fan arrows
and norm invariance in IC.1 give the result for every unimodular
input, including negative determinants.

\par\medskip\noindent\textbf{IC.5. The complete image and explicit actual fans.}\par\nopagebreak[4]

\item Unimodularity implies \(\gcd(x,y)=1\). The three nonzero parity
vectors \((1,0),(0,1),(1,1)\) all give \(x^2+xy+y^2=1\pmod2\).
Thus \(S=2t+1\) for an integer \(t\), and the two polynomials are
\[
 f_1(S)=-3t(t-1),\qquad f_2(S)=-5t^2-t-1.           \tag{IC.19}
\]
The first is even and has residues \(\{0,1,3,6\}\) modulo 7; the
second is odd and has residues \(\{0,2,5,6\}\) modulo 7. Direct
substitution for \(t\) modulo 4 gives respectively residues
\(\{0,2\}\) and \(\{1,3\}\). Equivalently, both polynomials have
period 28 modulo 28, with differences
\[
 f_1(2(t+28)+1)-f_1(2t+1)=-28(6t+81),
\]
\[
 f_2(2(t+28)+1)-f_2(2t+1)=-28(10t+141).
\]
These exact residue computations give the two eight-element upper
sets in (IC.2).

\item Every entry is realized by an actual primitive fan below. The notation
\([u,-r;s,z]\) denotes its full basis matrix; each determinant is 1
and its physical coordinates are \((r,u)\). The construction and all
source/cap endpoints for each such matrix are those of P3K.92.
\[
 \begin{array}{c|r|c|r|c}
 a&I_a\pmod{28}&(r,u)&S&P=[u,-r;s,z]\\\hline
 1&0 &(0,1)&1   &[1,0;0,1]\\
 1&6 &(3,4)&37  &[4,-3;-1,1]\\
 1&8 &(2,3)&19  &[3,-2;-1,1]\\
 1&10&(1,2)&7   &[2,-1;1,0]\\
 1&14&(1,5)&31  &[5,-1;1,0]\\
 1&20&(7,10)&219&[10,-7;3,-2]\\
 1&22&(1,3)&13  &[3,-1;1,0]\\
 1&24&(3,5)&49  &[5,-3;2,-1]\\\hline
 2&5 &(2,3)&19  &[3,-2;-1,1]\\
 2&7 &(1,2)&7   &[2,-1;1,0]\\
 2&9 &(1,3)&13  &[3,-1;1,0]\\
 2&13&(3,4)&37  &[4,-3;-1,1]\\
 2&19&(2,9)&103 &[9,-2;-4,1]\\
 2&21&(1,1)&3   &[1,-1;1,0]\\
 2&23&(2,5)&39  &[5,-2;-2,1]\\
 2&27&(0,1)&1   &[1,0;0,1]
 \end{array}                                       \tag{IC.20}
\]
Substituting their \(S\) values in (IC.19) verifies all entries.
The exact certificate also reexecutes both source-zero algorithms
and the selected-cusp algorithm, with their complete witness
histories, at every displayed matrix. Thus the upper sets are the
complete images of the actual infinity families.

\par\medskip\noindent\textbf{IC.6. Exact certificate scope.}\par\nopagebreak[4]

\texttt{code/part3\_fourfold/source\_infinity\_closed\_formula.py} executes the
symbolic reflected winding/prism identity, all nine unknowns in each
of the three complete first-target systems, their common seed, both
selected reset identities, both class-increment identities and norm
invariance. It checks the period 28 residue calculation and all sixteen
primitive witnesses by the original actual algorithms. A compact
additional replay covers 2,768 full algorithm evaluations, including
negative bases and different vertical seeds. The complete count is
2,885 exact checks. The formula and exhaustiveness follow from the
finite-word proof and the realized residue sets above; the replay
range supplies an additional check of that proof's implementation.

\end{enumerate}
\end{proof}




\begin{theorem}\label{night:P3K.94}
\StatementRecord{night:P3K.94}{proved}{S6\_06\_BLIND: GPT-6 Astra (OpenAI, ChatGPT Pro mode, effort Pro), recorded re-derivation in the checkpoint dated 18 November 2026; also checked in this writing by \CurrentChecker}{\cref{night:P3F.35}, \cref{night:P3F.41}, \cref{night:P3F.46}}
Retain the actual source family
\[
 y^2=x^3+r x^2+t x+t^2,\qquad G=(0,t),\qquad M_6=\mathcal O(6G-O),
\]
on a small disc about \(r=0\), with the analytic branch \(\rho(0)=-1\), \(\rho^2=1-3r\), and selected whole \(I_1\) zero at \(t_c=(1-3\rho^2+2\rho^3)/27\). Use the exact rational Cartier suspension \(B_{\rm sel}=(t-t_c)B_6/t^3\) defined in the proof, with fixed sufficiently small nonzero scalar. At infinity retain the reduced quadratic \(IV^*\) cap, either primitive norm \(n=\pm1\), the prescribed character \((-1)^d\), and the entire reference annular map. There is no extra logarithmic disc at either colliding fibre. For every prescribed lattice-basis fan at the selected cusp, these actual source, line, root, cap and toric maps define a smooth proper oriented sphere family over the full \(r\)-disc, with identity parameter closing. Its positive boundary return is \(0\in\Theta_7\). Under the exact marked-cover rescaling \(D=w^2-3\), \(r=(3-2w^2)/D^2\), the same conclusion holds at both points \(w^2=3/2\), with the scalar unchanged and the same transported fan.
\end{theorem}
\begin{proof}
\begin{enumerate}[label=\arabic*.]
\item \par

\medskip\noindent\textbf{EU.0. Actual source resolution and translation charts}

\par\nobreak

\item Set \(U=y-t-x/2,V=y+t+x/2\). Exactly
\[
 UV=x^2(x+r-1/4),\quad t=(V-U-x)/2,\quad y=(U+V)/2.
\]
On the bidisc at \(r=0\) the factor is a unit. Choose
\(c=\sqrt{x+r-1/4}\) and \(X=xc\). This is a parameter-preserving
analytic isomorphism to \(UV=X^2\); x is its analytic inverse
\(x(X,r)\), and t is the displayed actual base function. The
two minimal-resolution charts are \(X=Uz,V=Uz^2\) and
\(X=Vw,U=Vw^2\), with \(z=1/w\). They give a simultaneous smooth
source resolution over the entire r-disc.

\item The +G formula has slope \(\kappa=(y-t)/x\),
\(x'=\kappa^2-r-x\), \(y'=-t-\kappa x'\).
In the x-blowup chart \(y=xv,\ t=xu\), set \(\kappa=v-u\). Then
\[
 x=\kappa^2-r+(2\kappa-1)u,\quad
 u'=x/(1-2\kappa),\quad \kappa'=-\kappa-2u',\quad
 x'=u(1-2\kappa).
\]
Near the III component intersection the denominator is a unit.
At \(r=0\) the two component maps are
\(\kappa_C\mapsto-\kappa_C/(1-2\kappa_C)\) into E and
\(\kappa_E\mapsto-\kappa_E\) into C; both first derivatives
are -1. At \(\kappa=1/2\), \(x=1/4-r\) is a unit and the target
t-chart is instead
\[
 x'/t=(1-2\kappa)/x,\qquad y'/t=-1-\kappa(1-2\kappa)/x.
\]
This is the actual strict -G chart and has no pole.

\item On the source t-chart \(x=ta,\ y=tb\), the equation is
\(b^2=ta^3+ra^2+a+1\). Near G, use
\(\kappa=(ta^2+ra+1)/(b+1)\), which follows by cancelling
the literal factor \(a\) in \((b-1)(b+1)\). The output \(x',y'\) above is
ordinary and regular, with \(x'(G)=1/4-r\) at \(t=0\). Near -G put
\(N=(b-1)^2-ra^2-ta^3\); then
\[
 x'=N/a^2,\quad y'=-[ta^3+(b-1)N]/a^3,
 \quad \tau'=-x'/y'=aN/[ta^3+(b-1)N].
\]
The last denominator is a unit at \(a=0,\ b=-1\), so this is the actual
O chart. At source O use \(\tau=-x/y\), \(p=1/x=\tau^2v\)
with v an analytic unit satisfying
\(v=1+r\tau^2v+t\tau^4v^2+t^2\tau^6v^3\). Its target
t-chart has
\[
 a'=p(1+2/\tau+2tp),\qquad
 b'=-1+(1/\tau+tp)a',
\]
which extends with \((a',b')=(0,1)\). These charts cover the central
fibre; the inverse is the same construction for -G, or conjugation
by \(y\mapsto-y\). The two maps compose to the identity on a dense open
set and hence everywhere. Thus the actual translation extends as
a holomorphic source automorphism through \(r=0\), fixing t and r.

\item For the line, \(M_6\) is the actual section divisor \(6G-O\). At III, 2G
lies on the identity component at \((1/4,-1/8)\); its additive coordinate
\(-x/y\) is 2, so that of 6G is 6. Thus 6G and O remain
distinct smooth points there, on the same component, and \(M_6\) has
degree zero on each component. On the conductor of the two tangent
components a degree-zero line has gluing jet \(1+c\kappa\) modulo
\(\kappa^2\). Translation swaps the components, inverting this jet, and its actual derivative is -1 on both branches. These two
operations retain c, proving that the central Hom line is trivial.
On smooth, \(I_1\) and \(I_2\) fibres the same Hom line is trivial by the
actual group translation. Proper direct-image evaluation therefore
gives the invertible relative line lift on this whole bidisc;
rigidification at a fixed ordinary source section fixes its base
unit. Its relation to a globally specified suspension multiplier
must retain that unit and its parameter order.

\par

\medskip\noindent\textbf{EU.1. The precise completed datum and the retained earlier stage}

\par\nobreak

\item The preceding source charts and the following full line and cap maps are used together. The initial local calculation remains in the preserved attempt record.

\item Take a small closed complex parameter disc \(\Delta_r\) about zero.
Choose the analytic branch \(\rho^2=1-3r\) with \(\rho(0)=-1\), and put
\[
 t_c(r)=\frac{1-3\rho(r)^2+2\rho(r)^3}{27},\qquad
 q=t-t_c(r),\qquad s=\tfrac14-r.
 \tag{EU.1}
\]
The source is the simultaneous minimal regular model of \(S_r\)
constructed above, with its literal map to the \((t,r)\)-bidisc. Its
line is \(M_6=\mathcal O(6G-O)\). The source map of the suspension is
\(T:p\mapsto p-G\). Fix the quadratic reduced \(IV^*\) cap at infinity,
either signed primitive norm \(n=\pm1\), its prescribed character
\((-1)^d\), and the entire ordinary-to-cap reference map at a nonzero
\(r_*\in\Delta_r\). The selected cusp is the whole \(I_1\) fibre at
\(t_c(r)\), with any specified lattice-basis fan and the physical line
circle. There is no extra logarithmic disc at the colliding fibres.

\item The exact global suspension section is specified next. This defines a
marked scalar family rather than a generic Picard class. The completed
family extends over \(\Delta_r\) with this actual scalar held fixed and
small. Therefore its positive boundary loop has return class zero on
every such fan. The proof below constructs that extension and its
identity parameter closing before using smooth bundle triviality.
The preserved working history records the earlier local-only stage; the present proof retains the complete attachment.

\par

\medskip\noindent\textbf{EU.2. The complete rational Cartier multiplier at the collision}

\par\nobreak

\item Write \(D_j=(jG)-(O)\). In the actual Weierstrass rational frames let
\(f_{2,G}\) be the tangent-over-vertical Miller function, and construct
\(f_{3,G}\), then \(f_{6,G}=f_{3,G}^2f_{2,3G}\), by their actual chord
and tangent functions. The convention is
\[
 \operatorname{div}_{\rm generic}(f_{j,G})=jD_1-D_j,
 \qquad r_6=f_{6,G}^{-1}:M_6\longrightarrow M_1^{\otimes6}.
\]
The arrow includes both source and target factors at every chart; it
does not identify the two line bundles by their generic classes. With
\(U=y-t-x/2\), the full rational functions and transported lift are
\[
 f_{2,G}=\frac{U}{x-s},\qquad
 B_6(p)=f_{2,G}(p)^6\frac{f_{6,G}(Tp)}{f_{6,G}(p)},\qquad
 B_{\rm sel}(p)=\frac{q}{t^3}B_6(p).
 \tag{EU.2}
\]
Thus \(B_{\rm sel}:M_{6,p}\to M_{6,Tp}\) is a typed rational line map.
The denominator \(t^3\) is retained, not absorbed into a source-
parameter scalar.

\item Here is its complete divisor calculation near the present parameter
disc. For \(r\ne0\), the generic identity component of the \(I_2\)
fibre is described by \(w=z+1\) and
\(x_0=s+w^2-w\). The actual inverse-translation chart gives
\[
 \frac{B_6}{t^3}\bigg|_{t=0}
       =\frac{8((1-s)w-s)}{s^3}
 \tag{EU.3}
\]
at its generic point, in the rational source/target frames away from
their horizontal divisors. Here are the full factors establishing it.
Put \(\lambda_3=-2t/s-1/2\),
\(L_\pm=y\mp t\mp\lambda_3x\), and
\(x_3=2t(2t+s)/s^2\). The source equation gives
\[
 L_+L_-=x(x-s)(x-x_3),\qquad
 f_3=\frac{UL_+}{(x-s)(x-x_3)}
     =\frac{s xU}{sV-(s+2t)x}.
\]
The actual inverse translation has
\(x(Tp)/t\to2w/x_0\), \(y(Tp)/t\to(s-w^2)/x_0\).
The tangent of \(3G\) has limiting slope
\((1-3s)/(2(1-s))\) and \(x(6G)\to s^3/(1-s)^2\).
Direct substitution therefore gives
\[
\begin{aligned}
 f_2(p)&\to x_0/w,& f_3(p)&\to x_0,\\
 f_3(Tp)/t&\to\frac{2w^3}{x_0(w-s)},&
 f_{2,3G}(p)&\to\frac{x_0(1-s)}{(1-s)w-s},\\
 f_{2,3G}(Tp)/t&\to
       \frac{2(1-s)(w-s)^2}{s^3x_0}.
\end{aligned}
\]
Their literal product in EU.2 is EU.3. The denominators
\(s\) and \(1-s\) are units on the present disc. The line \(M_6\)
has degree zero on both \(I_2\) components. The Hom divisor consequently
has zero component intersections, and the primitive kernel of the
\(I_2\) matrix is the whole fibre. Hence its exact vertical divisor
there is \(3F_0\), on both components.

\item At a finite irreducible fibre \(F_c\), \(c\ne0\), write the actual
function \(f_{6,G}=(t-c)^e u\) in its component DVR. The displayed
base-preserving translation fixes \(t-c\), so the ratio in EU.2 is
literally \(u(Tp)/u(p)\). Neither \(U\) nor \(x-s\) contains an
entire such fibre. This proves the absence of every other finite
vertical divisor, including a fibre where a horizontal section
meets \(O\). At infinity the line has degree zero on every \(IV^*\)
component and its Hom degree is \(6h(G)=1\). Thus the exact remaining
vertical order is \(-2F_\infty\). Multiplication by \(q/t^3\), whose
divisor is \(F_c-3F_0+2F_\infty\), leaves exactly \(F_c\). This is
the selected whole-zero-one map of F41, with its rational scalar now
specified by EU.2.

\item The same map extends across \(r=0\) without a parameter factor. The
resolved source threefold is smooth by its exact \(A_1\) product
coordinates. The actual \(T\) and its inverse extend by the preceding
addition charts. Hence the relevant Hom sheaf is an invertible sheaf
on that smooth threefold. Outside the curve \(\{r=t=0\}\), EU.2 is
already invertible, except for its prescribed selected zero. To check
the divisor \(\{r=0\}\), evaluate on one ordinary local source
section with \(t\ne0,t_c\) avoiding the finitely many horizontal
divisors of the rational functions. There \(s=1/4\) and every actual
factor and its inverse is holomorphic and nonzero. Thus no component
of the parameter divisor occurs. The remaining omitted curve has
complex codimension two. In regular local rings a rational section
of an invertible sheaf that has valuation zero at every height-one
prime is a unit. Apply this to EU.2 and its inverse. It supplies the
actual extending line morphism on the whole collision neighbourhood.
No new rigidification or scalar depending on a punctured parameter
disc has been inserted.

\item For comparison with a specified local Hom frame, take the value of
EU.2 on that same ordinary section. Its ratio to the frame is an
explicit holomorphic unit on \(\Delta_r\); multiplying the frame by
that ratio makes it exactly EU.2. Its logarithm on the disc is fixed
by its value at \(r_*\), so this typed comparison introduces no
parameter winding. The central conductor-jet calculation above gives
the same fibrewise Hom line and independently checks this extension.

\par

\medskip\noindent\textbf{EU.3. The whole quadratic root, translation and line charts}

\par\nobreak

\item The literal source coordinates at infinity are
\(q_\infty=1/t,x_0=q_\infty^2x,y_0=q_\infty^3y\). The quadratic
root has \(q_\infty=z^2,x_0=z^2X,y_0=z^3Y\), and hence
\[
 Y^2=X^3+z^2(rX^2+X+1),\quad G'=(0,z),\quad
 D(z,X,Y)=(-z,X,-Y).
\]
With \(c(X,r)=\sqrt{1+X+rX^2}\), put \(U=Y-zc,V=Y+zc\).
Then \(UV=X^3\) and \(z=(V-U)/(2c)\); this is an exact
parameter-preserving \(A_2\) product. Near its common central
component point use the middle toric chart
\[
 X=ab,\quad U=a^2b,\quad V=ab^2,\quad
 c=\sqrt{1+ab+ra^2b^2},\quad
 z=\frac{ab(b-a)}{2c},\quad Y=\frac{ab(a+b)}2.
 \tag{EU.4}
\]
The deck is \((a,b,z)\mapsto(-a,-b,-z)\). Define the following
analytic units, all with value one at \(a=b=0\):
\[
\begin{aligned}
 K&=1+\frac{b(b-a)(1+rab)}{2c(c+1)},\\
 F&=1-\frac{ab(1+rab)(K+1)}{2c(c+1)}
           -\frac{ra b^2(a-b)}{4c^2},\\
 X'&=a(a-b)F,\qquad c'=\sqrt{1+X'+rX'^2},\\
 H&=K+\frac{b(a-b)(1+rX')}{2c(1+c')}.
\end{aligned}
 \tag{EU.5}
\]
The literal chord formula for translation by \(G'=(0,z)\) is
\(\kappa=(Y-z)/X=aK\),
\(X'=\kappa^2-rz^2-X\), \(Y'=-z-\kappa X'\).
Substitution into the target toric chart gives exactly
\[
                 a'=(b-a)F/H,\qquad b'=-aH.
 \tag{EU.6}
\]
Indeed \(V'=Y'+zc'=-aX'H\), so \(b'=V'/X'\) and
\(a'=X'^2/V'\). This retains \(z\), hence the actual base
\(q_\infty\), and its derivative at the central point is
\(\left(\begin{smallmatrix}-1&1\\-1&0\end{smallmatrix}\right)\),
of determinant one. It carries the three actual branches cyclically.
The units in EU.5 are invariant under \((a,b)\mapsto(-a,-b)\), so
this translation commutes with the full root deck.

\item The complement of this chart in the central root fibre is its smooth
Neron group locus. There the chord law in the actual resolved
Weierstrass charts is holomorphic in \((z,r)\). The displayed map and
the analogous inverse, obtained by \(Y\mapsto-Y\), agree with that
law on their overlaps. They therefore construct the complete source
automorphism through the central root, not only its tangent action.

\item The unique root division section can be constructed by an explicit
identity-component inverse:
\[
             R=2G'-[2]^{-1}(3G'),\qquad 2R=G'.
 \tag{EU.7}
\]
The section \(3G'\) lies in the small formal group at the central
fibre. The inverse of \([2]\) there is its convergent formal-group
inverse, with derivative \(1/2\). Translation by that section is
the time map of the actual invariant vertical vector field on the
resolved root, with its analytic Lie coefficient. That vector field
extends through EU.4; it is the derivative of the identity-component
addition law and vanishes at the component intersection. Its small
flow exists uniformly on a compact root tube. Together with the
twice-iterated map EU.6 it gives translation by \(R\) and its inverse.
The central component of \(R\) is the nonidentity label two. Formula
EU.7 is deck-invariant, by uniqueness of the identity-component
inverse and the actual deck action on \(G'\).

\item The actual section-line comparison also varies holomorphically over
the whole parameter disc. To specify its complete common graph,
blow up the four literal deck fixed points on the resolved root:
the triple point, the two outer toric-chart origins, and \(O\).
Each has derivative \((-1,-1)\). The deck fixes each exceptional
curve and negates its normal coordinate, so its quotient is smooth.
Each old component meets two centres and has strict square \(-4\)
before quotient, and square \(-2\) afterward. The four exceptional
images also have square \(-2\). Pulling back the actual equation
\(q_\infty=z^2\) gives multiplicities two on the three old
components, three on the triple exceptional image, and one on the
three outer images. The generic map
\((z,X,Y)\mapsto(q_\infty,x_0,y_0)=(z^2,z^2X,z^3Y)\)
therefore yields the minimal regular resolution of the original
Weierstrass equation. Denote the root blowup by \(\beta\) and this
quotient/source map by \(\pi\). These are the actual morphisms
in EU.9, with every exceptional curve retained.

\item In source coordinates at infinity let
\(\tau_0=-x_0/y_0\); in the root let \(\tau'=-X/Y=z\tau_0\).
The exact chord calculation gives
\[
 x_0(3G)\big|_{q_\infty=0}=4/s^2,\quad
 y_0(3G)\big|_{q_\infty=0}=8/s^3,\quad
 \tau_0(6G)\big|_{q_\infty=0}=-s.
 \tag{EU.8}
\]
Thus \(\tau'(6G')=z[-s+O(z)]\). In the actual graph blowup
\(\tau'=z u\), the source point is \(\tau_0=u\); \(O\) and
\(6G'\) have distinct directions \(0\) and \(-s\). Their two
pulled-back divisor equations both contain the same factor \(z\),
which cancels in their difference. Every further resolution centre
is away from these two marked sections. Consequently on the full
common resolved graph the actual equality is
\[
       \beta^*\mathcal O(6G'-O)=\pi^*\mathcal O(6G-O).
 \tag{EU.9}
\]
This is equivariant for the source and deck maps already constructed;
both rational \(r_6\) factors are carried on source and target. It
transports precisely EU.2 to the root line \(N\).

\item For completeness, the pulled-back multiplier is an actual unit before
the norm is chosen. The root scaling is
\(x=z^{-2}X,y=z^{-3}Y,t=z^{-2}\); hence the literal Miller factors
have weights \(z^{-1},z^{-2},z^{-5}\) for \(f_2,f_3,f_6\), respectively.
Writing \(B'_6\) for the same full Miller ratio on the root, this gives
\[
 B_0=B'_{\rm sel}=h z^{-2}B'_6,\qquad
 h=1-t_c(r)z^2.
 \tag{EU.9a}
\]
The graph Cartier map EU.9 and the already computed whole-fibre divisor
give order two for \(B'_6\) on each reduced-IV component. Its leading
coefficient can also be read directly. At a generic point on the
identity component write \(X=w^2,Y=w^3+O(z^2)\). The exact root formula
\[
 f'_3=\frac{s zX(Y-z-zX/2)}
              {s z(Y+z+zX/2)-(s z^2+2)X}
\]
and the actual inverse translation yield
\[
 f'_2\to w,\quad f'_3(p)/z\to-s w^3/2,\quad
 f'_3(Tp)/z^2\to s,\quad
 z f'_{2,3G'}(p),\ z f'_{2,3G'}(Tp)\to-4/s.
\]
Indeed \(z^2X(3G')\to4/s^2\),
\(z^3Y(3G')\to8/s^3\), and \(z^2X(6G')\to1/s^2\),
which give the last two tangent-over-vertical limits. Their product
is \(B'_6/z^2\to4\). In the trivial central line of EU.9 this is
the same constant on the connected reduced fibre. Thus EU.9a has
central coefficient \(4h\), and its full rational coefficient and
inverse extend as units by the height-one-prime argument used in EU.2.

\item On the central reduced IV fibre the two horizontal divisors of
\(N=\mathcal O(6G'-O)\) coincide and its component degrees vanish.
Thus \(N\) is actually trivial there. Choose a lift \(F\) of
translation by \(-R\) using the transported Hom evaluation. Compact
connected fibres make every ratio of two such lifts a base unit.
Both central commuting source maps act by constant multipliers on
the trivial central line, so the finite commutator character is one.
More explicitly,
\[
 DFD^{-1}=c_F(z,r)F,\qquad
 c_F(z,r)c_F(-z,r)=1,\qquad c_F(0,r)=1.
\]
Choose \(\ell=\log c_F\) with \(\ell(0,r)=0\). Then
\(\ell(-z,r)=-\ell(z,r)\), so the literal replacement
\(F_1=\exp(\ell/2)F\) satisfies
\[
 DF_1D^{-1}=\exp(-\ell/2)c_FF=F_1.
\]
The actual root suspension lift is
\(B_0=\beta(z,r)F_1^2\), where \(\beta\) is a nonzero
deck-invariant base unit by EU.9a. Its analytic square root with
the prescribed value at \(r_*\) exists on the full root disc times
\(\Delta_r\) and is deck invariant: the ratio of its deck transform
to itself is a constant sign, equal to one at \(z=0\).
Put \(\widetilde F=\beta^{1/2}F_1\). For either sign and every retained
character the actual cap action is
\[
 A_0=(-1)^d\widetilde F^{\,n}D,\qquad
 A_0^2=B_0^n,\qquad [A_0,B_0]=1.
 \tag{EU.10}
\]
All operations are equalities of the source/Cartier maps, not just
of their central values.

\par

\medskip\noindent\textbf{EU.4. The complete annular reference, scalar and closing}

\par\nobreak

\item Use an auxiliary scalar root \(R_\lambda\) with
\(\lambda=R_\lambda^2\), fixed nonzero and small. It is independent
of the source parameter \(r\). The actual actions are
\[
 B_\lambda=\lambda B_{\rm sel},\qquad
 A_{\lambda}=R_\lambda^{\,n}A_0,\qquad
 A_\lambda^2=B_\lambda^n.
 \tag{EU.11}
\]
Fix a cap annulus \(C\) in \(q_\infty\), uniform over
\(\Delta_r\), whose source fibres are smooth. Cover it by two angular
sectors \(C_\alpha\); their overlap has two components and no triple
overlap is needed. Write
\[
 s_\alpha=\frac{\log_\alpha q_\infty}{2\pi i},\quad
 z_\alpha=\exp(\pi i s_\alpha),\quad
 c_\lambda=\frac{n\log R_\lambda}{2\pi i}+\frac d2.
 \tag{EU.11a}
\]
Let \(I_{z_\alpha}\) be the actual source/root and Cartier map
EU.9 on this smooth annulus. In particular
\(I_{-z}=D I_z\). The rigidified degree-zero line gives its actual
commutative extension of the elliptic group by the line scalar group.
Use its Lie cover and period lattice, with
\(\delta=(0,1)\) exponentiating to one turn of the physical line.
In this cover choose the logarithm
\[
 f_\alpha=(-a_\alpha(R),\ell_\alpha)
       =\operatorname{Log}_\alpha\widetilde F,
 \qquad b_\alpha=\operatorname{Log}_\alpha B_\lambda.
 \tag{EU.11b}
\]
Here \(a_\alpha(R)\) is the actual Abel integral from \(O\) to
\(R\); \(\ell_\alpha\) is the logarithm of the lifted line action
in the same rigidified extension, including both \(r_6\) factors,
\(h\), and the chosen square root of \(\beta\). The displayed
ordered pair is a Lie-cover coordinate, so its exponential is the
full source/line map, not independent multiplication in an arbitrary
Cartier chart. Compactness of the elliptic fibre makes the commutator
of rigidified translation lifts a constant scalar; its value at the
identity is one. These are consequently commuting Lie translations,
and their exponentials commute with \(B_\lambda\).

\item The literal sector formula is
\[
 \Phi_\alpha
   =\operatorname{Exp}\!\left(
       n s_\alpha f_\alpha+s_\alpha c_\lambda\delta
                         +h_\alpha\right)I_{z_\alpha}.
 \tag{EU.11c}
\]
Thus its source Abel displacement is
\(-n s_\alpha a_\alpha(R)\), its scalar factor includes
\(R_\lambda^{n s_\alpha}\exp(\pi i d s_\alpha)\), and its
root changes by the displayed \(D\). The correction \(h_\alpha\)
is now specified relative to the entire reference map, with no
invariant logarithm of a central nonidentity component presumed.

\item On an overlap let \(s_\beta=s_\alpha+k_{\alpha\beta}\) and write
\[
 f_\beta=D_*^{k_{\alpha\beta}}f_\alpha+
                         \Pi_{\alpha\beta}.
\]
The vector \(\Pi_{\alpha\beta}\) is an actual integral combination
of the transported source/line periods. Its coefficients, the
reference suspension deck integers \(j_{\alpha\beta}\), and the
reference full period vectors \(\Lambda_{\alpha\beta}\) are fixed
by the given reference map at \(r_*\). The period functions themselves
vary holomorphically in \(r\). Set
\[
 h_\beta-D_*^{k_{\alpha\beta}}h_\alpha
  =-n s_\alpha\Pi_{\alpha\beta}
          +j_{\alpha\beta}b_\beta+\Lambda_{\alpha\beta}.
 \tag{EU.11d}
\]
The actual reference logarithms satisfy exactly this equation at
\(r_*\). It retains the scalar period \(d\delta\), every elliptic
period and every suspension deck used by that reference. In a
reference lifted with no additional suspension deck one has
\(j_{\alpha\beta}=0\); the formula does not impose that choice.

\item Here is an explicit holomorphic solution retaining that reference.
Use \(I_z\) to express the Lie vectors in the ordinary Lie bundle
over \(C\), and fix a holomorphic frame of that bundle, including
the scalar vector \(\delta\). Such a frame is obtained from the
actual invariant differential and the rigidified line splitting;
the latter is the same additive annulus splitting described next.
Extend the reference \(h_\alpha(r_*)\) in this frame to initial
vectors \(h^0_\alpha\). Subtract their difference from the right
side of EU.11d, obtaining a holomorphic additive overlap vector
\(E_{\alpha\beta}\) that is zero at \(r_*\).
Choose a smooth partition \(\chi_\beta\) subordinate to the fixed
sectors and put \(k_\alpha=-\sum_\beta\chi_\beta
E_{\alpha\beta}\). Then \(k_\beta-k_\alpha=E_{\alpha\beta}\),
and \(\bar\partial k_\alpha\) is a single vector-valued form
\(\omega\). After a fixed cutoff on a slightly larger annulus,
the literal Cauchy-\/-Green operator is
\[
 v(q_\infty,r)=\frac1\pi\int
            \frac{\omega(w,r)}{q_\infty-w}\,dA(w),\qquad
 h_\alpha=h^0_\alpha+k_\alpha-v.
 \tag{EU.11e}
\]
On the smaller fixed annulus \(\bar\partial v=\omega\).
The formula is holomorphic in \(r\), and its corrections are zero
at \(r_*\). It solves EU.11d exactly. Finite averaging, where
needed in the root atlas, is applied to these vector cochains only;
the lattice and finite-character cocycles remain the prescribed
ones. Their cyclic relation is the literal norm
\(A_\lambda^2=B_\lambda^n\), already proved in EU.10-\/-11.

\item Substitute EU.11d into EU.11c. The difference of its two Lie
arguments after subtracting the \(k_{\alpha\beta}\)-fold cap
action is \(j_{\alpha\beta}b_\beta+\Lambda_{\alpha\beta}\).
Exponentiating therefore gives the exact squares
\[
 \Phi_\beta=A_\lambda^{k_{\alpha\beta}}
                   B_\lambda^{j_{\alpha\beta}}\Phi_\alpha,
 \qquad \Phi_\alpha B_\lambda=B_\lambda\Phi_\alpha.
 \tag{EU.12}
\]
In the chosen single-sheet lift this is
\(\Phi(s+1,r)=A_\lambda\Phi(s,r)\), with the full root deck
included. The inverse is given by the negative Lie argument and
\(I_z^{-1}\), so these are actual annular isomorphisms. The finite
cap charts and their actions extend over its central fibre by
EU.4-\/-10; EU.12 is their exact overlap with the ordinary model.
No extension of an annular logarithm across that central fibre is
asserted or needed. This supplies the complete source, Cartier,
norm, scalar and reference-cochain attachment over \(\Delta_r\).

\item All maps are single valued in \(r\) on \(\Delta_r\). Traversing
its positive boundary therefore has source, line, finite-root and
annular forward closing equal to the identity. There is no scalar
winding: the scalar in EU.11 was held fixed, and every additional
base unit has its fixed disc logarithm. This closing assertion is
made for the complete quotient atlas, not inferred from a
fibrewise conjugacy.

\par

\medskip\noindent\textbf{EU.5. Every selected fan and the proper disc family}

\par\nobreak

\item The selected node remains disjoint from the collision and from both
marked sections. Its branch parameter, inverse-j coordinate, Tate
period and actual theta/Cartier multipliers are holomorphic on the
selected disc times \(\Delta_r\). Their leading units are nonzero
there. The selected branch marking is fixed by its value at \(r_*\);
the parameter disc admits no branch exchange. For every specified
periodic fan, each absolute cone character has the form
\(Z^uV^v q^h\). The full source/line maps multiply it by the actual
holomorphic base units to their integral powers. Those factors and
their inverses extend on the affine cone. The resulting chart maps
agree on faces as the same character maps. All periods and the
physical circle are retained. Hence this construction works for
every prescribed lattice-basis fan with the identical closing.

\item On the compact source parameter disc and finitely many source/cap
charts, the actual line multipliers and their inverses are bounded
away from the selected zero. The source translation maps have the
same compact bounds. Choose \(|\lambda|\) smaller than their common
contracting threshold. The ordinary norm shells, root norm shells
and selected periodic toric fundamental regions are then compact
uniformly in \(r\). In the finite cap, choose integers \(u,v\) with
\(nu+2v=1\); its actual generator
\(H=A_\lambda^uB_\lambda^v\) satisfies \(H^2=B_\lambda\), and its
scalar factor is \(R_\lambda\). Strict contraction proves freeness
and properness of the quotient. These are the same full quotient
maps as F41/F46. The explicit product resolutions above and the
source/cap charts show that projection to \(\Delta_r\) is a smooth
submersion. The compact fundamental regions show it is proper.

\item At \(r=0\) its source is the E6-B source of F46. Its two cap orders
are \((m_6,m_2)=(2,1)\), with \((n_6,n_2)=(n,0)\), so the actual
determinant is \(n=\pm1\). F46\textquotesingle s integral maps prove that this
central reference-fan fibre is the constructed oriented sphere.
For a prescribed lattice-basis fan \(P\), take its finite canonical
fan word. At each step the actual common blowup and its two ruling
maps of F35 identify the unchanged complement and the two attaching
boundaries. F35\textquotesingle s primitive integral Gysin maps and its van Kampen
maps preserve the sphere homology and simple connectivity through
that literal surgery. Following this finite word therefore supplies
the sphere fibre on \(P\), with the period and physical-circle
markings carried by those maps. The toric parameter maps above are
defined on that same \(P\) throughout the disc. All its other
fibres are consequently the same smooth sphere type by the proper
disc submersion. Choose a smooth horizontal distribution on this actual
proper family, and use radial paths on \(\Delta_r\) to trivialize
it. The identity atlas closing in EU.12 gives a boundary return
isotopic to the identity. Its clutched class is thus zero in
\(\Theta_7\), on every fan just constructed. No integer division
or source-moduli winding formula is used.

\par

\medskip\noindent\textbf{EU.6. Exact dictionary to the marked w-cover}

\par\nobreak

\item The global marked-source cover found by root is
\[
 D=w^2-3,\quad r=\frac{3-2w^2}{D^2},\quad
 X=D^2x,\ Y=D^3y,\ T=D^3t,
\]
\[
 Y^2=X^3+(3-2w^2)X^2+DTX+T^2,\quad
 G=(0,T),\quad T_c=w^2-1.
 \tag{EU.13}
\]
At either \(w_*^2=3/2\), \(D=-3/2\) is a unit and
\(dr/dw=-16w_*/9\ne0\). These are positive complex local
parameter coordinates on the present \(r\)-disc. The canonical
source/Cartier map carries every section and its actual divisor.
Its Miller functions have the exact homogeneous factors
\[
 f_{2,\rm old}=D^{-1}f_{2,\rm new},\quad
 f_{3,\rm old}=D^{-2}f_{3,\rm new},\quad
 f_{6,\rm old}=D^{-5}f_{6,\rm new}.
 \tag{EU.14}
\]
For example, the target tangent at \(G\) has slope \(D/2\), so
\(U_{\rm old}=D^{-3}[Y-T-(D/2)X]\) and
\((x-s)_{\rm old}=D^{-2}[X-D^2s]\), proving the first identity.
The chord-over-vertical factor has weight one; its actual Miller
recursion proves the other two identities, with all constants kept.
Their source and target \(f_6\) factors cancel in the ratio, so
\[
 B_{6,\rm old}=D^{-6}B_{6,\rm new},\qquad
 (q/t^3)_{\rm old}=D^6\frac{T-T_c}{T^3}.
\]
Consequently \(B_{\rm sel,old}=B_{\rm sel,new}\) under the
canonical divisor line map. This fixes the exact suspension scalar
in the global marked cover. The finite-root coordinate changes use
analytic roots of the unit \(D\) on this local disc; their reference
values and full finite characters are retained. They close as the
identity. The selected Tate branch and full theta/Cartier map are
transported by the same source isomorphism, so the fan is unchanged.
Thus the two unselected-III generators on this marked cover have
the zero return just proved, with their actual root and line closing.

\par

\medskip\noindent\textbf{EU.7. Executed exact certificate}

\par\nobreak

\item The independent symbolic chart certificate executes 96 exact identities for the source addition, root translations, full Miller leading factors, finite averaging, signed annular cochains, and marked-cover weight cancellation. Its source is \path{code/part3/e6_unselected_iii.py}; the complete output is \path{code/part3/e6_unselected_iii.json}. The certificate supplements the complete geometric quotient and properness proof above.

\end{enumerate}
\end{proof}




\begin{theorem}\label{night:P3K.95}
\StatementRecord{night:P3K.95}{proved}{S6\_06\_BLIND: GPT-6 Astra (OpenAI, ChatGPT Pro mode, effort Pro), recorded re-derivation in the checkpoint dated 18 November 2026; also checked in this writing by \CurrentChecker}{\cref{night:P3F.41}, \cref{night:P3F.55}, \cref{night:P3F.56}, \cref{night:P3F.7}, \cref{night:P3F.8}, \cref{night:P3F.9}, \cref{night:P3K.42}, \cref{night:P3K.63}, \cref{night:P3K.66}, \cref{night:P3K.68}, \cref{night:P3K.71}, \cref{night:P3K.94}}
\par\medskip\noindent\textbf{Statement.}\par\nopagebreak[4]

For the actual source family
\[
 Y^2=X^3+(3-2w^2)X^2+(w^2-3)TX+T^2,\qquad
 G=(0,T),\qquad q=T-(w^2-1),
\]
retain \(M_6=\mathcal O(6G-O)\), source suspension by \(-G\),
the complete selected-zero Hom section, and the quadratic reduced-IV
cap with either primitive norm sign \(n=\pm1\), its prescribed
finite character and full integral annular reference cochain.
Fix the sufficiently small nonzero scalar, including its chosen
root and phase. The actual positive \(w\)-loop about zero has,
on every selected lattice-basis fan, the same oriented sphere return
as the complete K71 selected-cusp algorithm. This equality is
induced by the full source, Cartier, finite-root, exterior, normal
and Cox maps below.

On the standard vertical fan \(T_j\), \(j\in\mathbb Z\), its
physical characteristic pair and oriented class are
\[
 (k_M,\ell_M)=(-6-18j(j-1),0),\qquad
 [\Sigma_j]=-1-3j(j-1)\pmod {28}.
\]
The corresponding native free pair is
\((0,-6-18j(j-1))\). At \(j=0\) the complete native triple is
\((0,-6,-1)\). For every other selected lattice-basis fan the
actual unit matrices and native trace maps in EC.6 give the same
K71 calculation. The marked parameter comparison to the original
source is
\[
 r(w)=\frac{3-2w^2}{(w^2-3)^2},\qquad
 r(w)-\frac13=-\frac{w^4}{3(w^2-3)^2};
\]
the \(w\)-loop therefore has actual local degree four over the
\(r-\tfrac13\) parameter. All conclusions refer to this specified
ordered loop and the retained \(M_6,\ b=1,\ m=2\) completion.

\end{theorem}
\begin{proof}
\begin{enumerate}[label=\arabic*.]
\item \par\medskip\noindent\textbf{EC.1. The full physical-height equation and the ordered normalization.}\par\nopagebreak[4]

\item Set
\[
 \begin{gathered}
 v=w^2-1,\quad D=w^2-3,\quad G=(0,T),\quad q=T-v,\\
 E_w:\ Y^2=X^3+(3-2w^2)X^2+DTX+T^2 .
 \end{gathered}
 \tag{EC.1}
\]
The marked-node coordinates are (EC.2)--(EC.4), and the full source/Cartier rescaling is (EC.15) and (EC.20) below. At the selected cusp put
\(\xi=X-v\). The exact equation, including its full height, is
\[
 K=Y^2-\xi^2(\xi+w^2)
       =q\{q+v^2+D\xi\}.
 \tag{EC.2}
\]
Write \(a=v^2+D\xi\). On a sufficiently small common polydisc the
square root equal to \(a\) at \(K=0\) gives the literal inverse
\[
 q=\frac{2K}{a+\sqrt{a^2+4K}},\qquad
 T=v+q,\quad X=v+\xi.
 \tag{EC.3}
\]
Its denominator has value two at the origin. Thus \((\xi,Y,w)\)
are smooth coordinates on the actual source threefold; (EC.3) retains
the original elliptic-base function. In these coordinates the entire
selected rational fibre has normalization
\[
 f(S,w)=\bigl(S^2-w^2,S(S^2-w^2),w\bigr),\quad
 D_+=\{S=w\},\quad D_-=\{S=-w\},\quad
 Z=\frac{S-w}{S+w}.
 \tag{EC.4}
\]
The point at infinity is \(O\), and \(G\) has \(S=-1\). This is
the literal ordered germ used in F56, with \(w\) in place of its
ordered parameter. It is an immersion off \((S,w)=(0,0)\); its three
derivative minors generate \((S,w^2)\), and its only distinct double
pairs there are \((w,w),(-w,w)\). These statements follow directly
from (EC.4), with no base reparametrization.

\item The second selected critical value is
\[
 q_-=-4w^6/27.
 \tag{EC.5}
\]
Indeed the full discriminant is
\(-432T^2(T-v)(T-T_-)\), with \(v-T_-=4w^6/27\).
The unselected \(I_2\) remains at \(T=0\), outside a common small
disc about \(T=-1\). Hence the positive \(w\)-loop has exactly the
ordered two-node braid of F56, including its twelve positive half
twists and the actual boundary twist of its lifted once-punctured
torus. The orientation of that loop is preserved by the identity
\(w=\rho_{\rm F56}\), not inferred from the even original \(r(w)\).
In the original source parameter the exact identity is
\(r(w)-1/3=-w^4/(3D^2)\). Thus the present positive ordered loop
has local degree four over \(r-1/3\). Its actual source, line and
cap closing is the full periodic \(w\)-atlas below; no shorter
endpoint is assigned by this calculation.

\item For the full normal arrow put \(x=S-w,y=S+w\), so \(\xi=xy\), and
retain the nonzero covector
\[
 n_f=(3S^2-w^2)d\xi-2S\,dY+2w(S^2-w^2)dw.
\]
Restriction of the differential of (EC.2) to the normalization gives
\[
 dq=-\frac{xy}{v^2+Dxy}\,n_f.
 \tag{EC.6}
\]
In the two native fixed-edge charts \(q=Zc_0c_1\) and
\(q=Z^{-1}c'_0c'_1\), \(c'_i=Zc_i\), this is exactly
\[
 d(c_0c_1)=-\frac{y^2}{v^2+Dxy}n_f,
 \qquad d(c'_0c'_1)=-\frac{x^2}{v^2+Dxy}n_f.
 \tag{EC.7}
\]
The common factor has an analytic logarithm and square root on the
finite inner ball. These are the actual two Hom arrows of K42 SP.8,
with their full extending source unit. They also show explicitly why
the physical height cannot be replaced by its leading Tate term.

\par\medskip\noindent\textbf{EC.2. The actual Cartier lift and all selected period units.}\par\nopagebreak[4]

\item Put \(s_0=v(w^2+3)/4\) and
\(U=Y-T-DX/2,\ V_s=Y+T+DX/2\). Then
\(UV_s=X^2(X-s_0)\). With the following actual Miller factors define
\[
 f_2=\frac{U}{X-s_0},\qquad
 f_3=\frac{s_0 XU}{s_0V_s-(Ds_0+2T)X},\qquad
 f_6=f_3^2f_{2,3G},
\]
\[
 B_6(p)=f_2(p)^6\frac{f_6(p-G)}{f_6(p)},\qquad
 B_{\rm sel}=\frac q{T^3}B_6.
 \tag{EC.8}
\]
These are coefficient arrows on the actual point line \(M_6\).
The point-line/tensor-line comparison is multiplication by
\(r_6=f_6^{-1}\); both its source and target factors occur in
(EC.8). At \(w=0,q=0,\xi=Y=0\), direct substitution gives
\[
 f_2=2,\quad f_3=3,\quad f_6=6,\quad
 B_6=64,\quad B_{\rm sel}/q=-64.
 \tag{EC.9}
\]
For the tangent factor at \(3G\), its cusp coordinates are
\((X,Y)=(-8/9,-1/27)\); its value at the cusp is \(2/3\),
which supplies the displayed value of \(f_6\). All horizontal
divisors of these factors avoid this cusp point. Thus these are
values of extending units, rather than cancellations of leading
terms at an indeterminate point.

\item There is a full selected-normalization check of its actual theta
frame. Put \(P=(1+w)/(1-w)\) and
\[
 S_6=w\frac{1+P^6}{1-P^6},\qquad
 c_6=w-S_6=\frac{(1+w)^6}{6+20w^2+6w^4}.
 \tag{EC.9a}
\]
The point \(6G\) is \(S=S_6\), with \(S_6(0)=-1/6\).
The exact divisor equations obey
\[
 S-S_6=c_6\,\frac{1-Z/P^6}{1-Z},\qquad
 f_6\big|_{q=0}=\frac{(S+1)^6}{S-S_6}.
 \tag{EC.9b}
\]
For the second identity both sides have the prescribed divisor
and the same leading coefficient \(S^5\) at \(O\); their ratio
is therefore the constant one on the normalization. Both branches
of the node have the same value
\[
 f_6\big|_{\rm node}=6+20w^2+6w^4,\qquad
 f_2\big|_{\rm node}=2,\qquad
 (B_{\rm sel}/q)\big|_{\rm node}=64/v^3.
 \tag{EC.9c}
\]
The reciprocal of the unit \(c_6\) is the coefficient arrow between
the displayed point-divisor and theta frames. It extends on the
entire parameter disc. This fixes the point-line theta gauge with
its full normalization constant, and shows explicitly that it
contributes no parameter winding.

\item Translation and inverse translation are also explicit here. For
\(+G\) use \(k=(Y-T)/X\), and for \(-G\) use \(k=(Y+T)/X\).
The respective target formulae are
\[
 X'=k^2-(3-2w^2)-X,\quad
 Y'=-T-kX'\quad\hbox{or}\quad Y'=T-kX'.
 \tag{EC.10}
\]
Since \(X=-1\) at the cusp, both are analytic there, and the two
maps are inverse by the actual chord law on their dense smooth
fibres. Their common fixed cusp therefore has the full line unit
(EC.9), including its constant sign. The scalar is
\(B_\lambda=\lambda B_{\rm sel}\), with the originally fixed
\(\lambda\); the factor \(-64\) is not absorbed into it.

\item On the selected normalization the marked source multiplier is
\[
 P_G=\frac{1+w}{1-w},\qquad g=P_G^{-1}.
 \tag{EC.11}
\]
The period and line charts give the exact commuting periods
\[
 H:(Z,V,q)\mapsto(\kappa qZ,\alpha V,q),\qquad
 B:(Z,V,q)\mapsto(gZ,\eta qV,q),\quad\eta=\lambda\beta.
 \tag{EC.12}
\]
Here \(\alpha\) and \(\beta\) are the coefficients of the actual
theta/Cartier trivialization of (EC.8), not new independently chosen
units. Its source point is \(6G\), so the marked theta-period
coefficient is \(P_G^6\); the actual rational function \(f_6\)
is elliptic and hence preserves that period on both copies. At the
node the invertible Hom coefficient is the one in (EC.9).

\item There is a joint domain for all the required logarithms. Substitute
\(\xi=w^2\widetilde X,Y=w^3\widetilde Y,q=w^6Q_c\) into (EC.2):
\[
 \widetilde Y^2=\widetilde X^2(\widetilde X+1)
       +Q_c(v^2+Dw^2\widetilde X)+w^6Q_c^2.
 \tag{EC.13}
\]
At \(w=0\) this is the marked compact pencil
\(\widetilde Y^2=\widetilde X^2(\widetilde X+1)+Q_c\).
Choose a sufficiently small fixed \(Q_c\)-disc containing just its
selected node. The inverse-j series, with its actual marked fourth
root, gives
\[
 \kappa=w^{-6}\widehat\kappa(Q_c,w),\qquad
 \widehat\kappa(0,w)=-v^2/64.
 \tag{EC.14}
\]
The factor \(\widehat\kappa\) and its inverse extend on this disc
times the full \(w\)-disc. The original section, viewed in the
weighted source chart at its section at infinity, has a formal
parameter
\[
 -\widetilde X(G)/\widetilde Y(G)
       =\frac{wv}{v+w^6Q_c},
 \tag{EC.14a}
\]
divisible by \(w\). Its denominator is a unit on the full
weighted bidisc. The marked Abel/theta point
therefore has an analytic formal-group logarithm tending to zero.
This proves the extension of \(P_G\), \(g\), and \(\alpha\) as
units in that same domain. Finally (EC.8)--(EC.9) supplies the
extending source and target Cartier units for \(\beta\); its
central value is \(-64\). Thus the based periodic logarithms of
all four units are literal, after retaining the sole angular factor
\(w^{-6}\) in \(\kappa\). Their winding matrix is
\(\operatorname{diag}(-6,0)\).

\par\medskip\noindent\textbf{EC.3. The unselected exterior and the actual finite cap.}\par\nopagebreak[4]

\item The original E6 source coordinate is
\[
 r=(3-2w^2)/D^2,\quad (x,y,t)=(D^{-2}X,D^{-3}Y,D^{-3}T).
 \tag{EC.15}
\]
On a small \(w\)-disc, \(D\) is a unit with a fixed analytic
logarithm, \(r\) lies in a disc about \(1/3\), and
\(s=1/4-r\) is a unit with value \(-1/12\). At the unselected
\(I_2\), the original coordinates satisfy
\[
 (y-t-x/2)(y+t+x/2)=x^2(x-s).
\]
The analytic map \(\widehat x=x\sqrt{x-s}\) identifies this with
\(UV=\widehat x^2\), retaining
\(t=(V-U-x)/2\). The two source resolution charts are
\(\widehat x=Uz,V=Uz^2\) and
\(\widehat x=V\zeta,U=V\zeta^2\), with \(z\zeta=1\).
Translation is the actual addition map of EU.1, including its
\(x\)-chart, the \(t\)-charts at \(\pm G\), and the zero-section
chart. The only denominator singled out in the exceptional
\(k=1/2\) chart is \(x=s\), which is a unit on this disc.
Consequently those maps and their inverses extend jointly here.

\item The full divisor computation in \cref{night:P3K.94}, (EU.2)--(EU.3) and the following component calculation, gives the complete vertical divisor
\(\operatorname{div}_{\rm vert}B_6=3F_0-2F_\infty\).
The factor \(q/T^3\) in (EC.8) cancels it except at the selected
whole zero. At the parameter divisor \(w=0\), its value on any
ordinary reference section avoiding the stated horizontal divisors
is a unit. The resolved total source is smooth. The rational Hom
map and its inverse therefore extend at the remaining
codimension-two loci by their zero valuations at every height-one
prime. This carries the actual Hom coefficient on the entire
ordinary exterior; no parameter scalar is inserted.

\item At infinity put \(q_\infty=z^2\). The original quadratic root is
\[
 Y'^2=X'^3+z^2(rX'^2+X'+1),\quad G'=(0,z),
 \quad\mathfrak D(z,X',Y')=(-z,X',-Y').
 \tag{EC.16}
\]
With \(c=\sqrt{1+X'+rX'^2}\), the coordinates
\(U'=Y'-zc,V'=Y'+zc\) obey \(U'V'=X'^3\) and
\(z=(V'-U')/(2c)\). Use the three exact $A_2$ charts in \cref{night:P3K.94}, (EU.4).
In the middle chart they are
\[
 X'=ab,\quad U'=a^2b,\quad V'=ab^2,\quad
 z=ab(b-a)/(2c),\quad Y'=ab(a+b)/2.
 \tag{EC.17}
\]
The full translation formula EU.5--6 is
\(a'=(b-a)F/H,b'=-aH\), with its actual analytic units
\(F,H\) and value-one denominators. These formulae, their inverse,
and the smooth group charts are uniform on this \(r\)-disc; none
requires \(r=0\). The half-division is the actual section
\(R=2G'-[2]^{-1}(3G')\); the inverse is the convergent formal-group
inverse at the identity. Hence \(2R=G'\) and \(\mathfrak D R=R\),
with its nonidentity central component retained.

\item Blow up the four literal deck-fixed points constructed in \cref{night:P3K.94}, (EU.8). On that common
graph the two section equations at the zero-section blowup are
\(\tau'=zv_0\) and \(\tau'-\tau'(6G')=z(v_0-L_6(s,q_\infty))\),
where \(L_6(s,0)=-s\). The actual point-line arrow is
\[
 \frac{\tau'}{\tau'-\tau'(6G')}=
          \frac{v_0}{v_0-L_6(s,q_\infty)}.
 \tag{EC.18}
\]
It proves \(\beta^*\mathcal O(6G'-O)=\pi^*M_6\), including every
exceptional coefficient. On the root the full selected coefficient
is \(B_0=h_\infty z^{-2}B'_6\),
\(h_\infty=1-t_c(r)z^2\). This is invertible on every reduced-IV
component by the full root divisor computation in \cref{night:P3K.94}, (EU.9)--(EU.9a).

\item For clarity retain the norm normalization. If a lift \(F\) of
\(t_{-R}\) has \(\mathfrak D F\mathfrak D^{-1}=c_FF\), then
\(c_F(-z)c_F(z)=1\), \(c_F(0)=1\). With its odd logarithm put
\(F_1=e^{(\log c_F)/2}F\). It commutes with \(\mathfrak D\).
The invariant unit \(\beta_0=B_0/F_1^2\) has its analytic square
root fixed at the full reference map. Then
\[
 F_2=\beta_0^{1/2}F_1,\quad
 A_0=(-1)^dF_2^n\mathfrak D,\quad A_0^2=B_0^n,
 \qquad B=R_\lambda^2 B_0,\quad A=R_\lambda^n A_0.
 \tag{EC.19}
\]
Here \(R_\lambda^2=\lambda\) is fixed, including its complex
phase. All source, line, character and norm arrows are equalities.

\item Under (EC.15), the root coordinates have the literal comparison
\[
 z_{\rm old}=D^{3/2}z_{\rm new},\quad
 X'_{\rm old}=D X'_{\rm new},\quad
 Y'_{\rm old}=D^{3/2}Y'_{\rm new}.
 \tag{EC.20}
\]
The square root uses the fixed logarithm of the unit \(D\).
The transformed root is the equation obtained from (EC.16) by the displayed substitution (EC.20). The full
Miller weights are \(D^{-1},D^{-2},D^{-5}\), and
\(B_{6,\rm old}=D^{-6}B_{6,\rm new}\). Since
\((t-t_c)/t^3=D^6 q/T^3\), the complete suspension coefficient
is unchanged. The Cartier ratio (EC.18) carries both numerator
and denominator: explicitly
\(\tau'_{\rm old}=D^{-1/2}\tau'_{\rm new}\) and
\(\tau_{0,\rm old}=D^{-2}\tau_{0,\rm new}\), consistently with
\(\tau'=z\tau_0\). Also
\(1-t_{c,\rm old}z_{\rm old}^2=1-vz_{\rm new}^2\).
Thus (EC.20) closes literally over the \(w\)-disc;
no finite root rotation or scalar winding is concealed by the
source rescaling.

\item Here is the literal application of EU.11a--12 to this family. Work in
the current height \(T\), put \(q_\infty=T^{-1}\), and retain
\(t_c(w)=v=w^2-1\) and \(h_\infty=1-t_c(w)z^2\). On a fixed
smooth cap annulus choose the same two angular sectors and reference
ordinary-to-cap map as in that construction. Its actual
\(I_z\) is the source/root and Cartier map (EC.18) followed by
(EC.20); in particular \(I_{-z}=\mathfrak D I_z\). In the
rigidified extension-group Lie cover set
\[
 \begin{gathered}
 s_\alpha=\frac{\log_\alpha q_\infty}{2\pi i},\qquad
 z_\alpha=e^{\pi i s_\alpha},\qquad
 c_\lambda=\frac{n\log R_\lambda}{2\pi i}+\frac d2,\\
 f_\alpha=\operatorname{Log}_\alpha F_2
       =(-a_\alpha(R),\ell_\alpha),\qquad
 b_\alpha=\operatorname{Log}_\alpha B .
 \end{gathered}
 \tag{EC.21a}
\]
The source Abel integral \(a_\alpha(R)\) and the line logarithm
\(\ell_\alpha\) belong to the same extension group. The latter
retains the two \(r_6\) factors, \(h_\infty\), and the selected
square root of \(\beta_0\) in (EC.19). Thus exponentiation is
the full source/line arrow. It is not a choice of independent
scalar factors in unrelated Cartier charts.

\item Write \(s_\beta=s_\alpha+k_{\alpha\beta}\) and
\[
 f_\beta=\mathfrak D_*^{k_{\alpha\beta}}f_\alpha+
                         \Pi_{\alpha\beta}.
\]
Here \(\Pi_{\alpha\beta}\) is the actual transported source/line
period, with the integer coefficients fixed at the given reference
parameter \(w_*\). Retain also its reference suspension integers
\(j_{\alpha\beta}\) and full integral period vectors
\(\Lambda_{\alpha\beta}\). The coefficients are fixed while the
period functions themselves vary holomorphically with \(w\).
The ordinary-to-cap map is
\[
 \begin{split}
 \Phi_\alpha
   &=\operatorname{Exp}\bigl(
       n s_\alpha f_\alpha+s_\alpha c_\lambda\delta+
                         h_\alpha\bigr)I_{z_\alpha},\\
 h_\beta-\mathfrak D_*^{k_{\alpha\beta}}h_\alpha
   &=-n s_\alpha\Pi_{\alpha\beta}
          +j_{\alpha\beta}b_\beta+\Lambda_{\alpha\beta}.
 \end{split}
 \tag{EC.21b}
\]
This is EU.11c--d with \(r\) replaced by \(w\), with every displayed
root, Cartier, period and character map replaced by (EC.18)--(EC.20).
In particular its source displacement is
\(-n s_\alpha a_\alpha(R)\), and its scalar includes
\(R_\lambda^{n s_\alpha}\exp(\pi i d s_\alpha)\).

\item For completeness the correction in (EC.21b) is an actual
parameter-dependent map. Use \(I_z\) to put its Lie vectors into
the ordinary Lie bundle on the fixed annulus and its invariant
differential/rigidified-line frame. Extend the entire reference
\(h_\alpha(w_*)\) to \(h_\alpha^0\) in that frame. Subtract its
difference from the right side of (EC.21b), obtaining the additive
overlap vector \(E_{\alpha\beta}\), zero at \(w_*\).
For a fixed sector partition of unity put
\(k_\alpha=-\sum_\beta\chi_\beta E_{\alpha\beta}\).
Then \(k_\beta-k_\alpha=E_{\alpha\beta}\), and
\(\bar\partial k_\alpha\) is a single vector-valued form \(\omega\).
With a fixed cutoff on a slightly larger annulus use
\[
 v_{\rm CG}(q_\infty,w)=\frac1\pi\int
      \frac{\omega(\zeta,w)}{q_\infty-\zeta}\,dA(\zeta),
 \qquad h_\alpha=h_\alpha^0+k_\alpha-v_{\rm CG}.
 \tag{EC.21c}
\]
On the protected annulus \(\bar\partial v_{\rm CG}=\omega\).
This is holomorphic in \(w\), zero at \(w_*\), and solves the
specified overlap equation. Finite averaging is applied only to
these vector cochains; it preserves the fixed lattice cocycle,
character and complete reference map. Substitution in (EC.21b)
leaves, after subtracting the \(k_{\alpha\beta}\)-fold actual cap
action, exactly \(j_{\alpha\beta}b_\beta+\Lambda_{\alpha\beta}\).
Hence the full quotient arrows satisfy
\[
 \Phi_\beta=A^{k_{\alpha\beta}}B^{j_{\alpha\beta}}\Phi_\alpha,
 \qquad \Phi_\alpha B=B\Phi_\alpha .
 \tag{EC.21}
\]
Their inverse is \(I_z^{-1}\) followed by the negative Lie
argument. Both norm signs retain the fixed factor \(R_\lambda^n\)
and their actual finite character. All these maps are single valued
over the \(w\)-disc, so their continuation has identity closing.
This specifies the complete \(F_0\), rather than only its norm.
For the reflection used in EC.6, perform this exterior/root and
annular construction first on the disc \(\zeta=w^2\), with
reference \(\zeta_*=w_*^2\), and then pull it back by that exact
map. The source equation, \(G\), \(t_c=v\), \(D\), \(I_z\), and
all root/Cartier units above are holomorphic in \(\zeta\).
Choose the reference extensions \(h_\alpha^0\), period sections,
and Cauchy operator on that disc. The resulting complete
ordinary-to-cap arrow is literally even in \(w\) in these
original source coordinates. Its expression in the two ordered
selected Tate markings is then related by the exact reflection
arrow of EC.6; no additional \(w\)-dependent cochain is chosen.

\item The ordinary and cap quotients are proper uniformly on the closed
parameter disc. On their finite source charts away from the selected
zero the actual line maps and their inverses are bounded in fixed
Hermitian metrics. Choose the original fixed \(|R_\lambda|\) below
their common strict-contraction threshold. In the cap the literal
generator \(H=A^n\) satisfies \(H^2=B\), \(H^n=A\), and has scalar
factor \(R_\lambda\). Its norm shell is compact and has the
specified character; this proves freeness and properness for both
norm signs. The selected toric family on the compact punctured
w-circle has its finite union of protected period fundamental
regions. On the exterior of the selected disc the smooth charts,
the root resolution and these compact shells give a proper smooth
family over the entire w-disc, with the actual source/line boundary
marking (EC.21).

\par\medskip\noindent\textbf{EC.4. The real orbit product and the full fixed trace.}\par\nopagebreak[4]

\item On the small source region containing the selected pair the source
is the smooth threefold (EC.3). The section of its group defining
translation is in the smooth locus at the central cusp. Divide it
there by a sufficiently large integer, extend that division section
by the inverse function theorem, take its small-group logarithm,
and multiply the logarithm by that integer. Its exponential is the
actual section by uniqueness and the chord law. The invariant
vertical field, in the original source coordinates,
\[
 2Y\partial_X+\{3X^2+2(3-2w^2)X+DT\}\partial_Y
\]
extends across the cusp and vanishes there. Its indicated time maps
give the actual source translation isotopy. The nonzero line lift
\(B_{\rm sel}/q\) has the same extension-group construction; its
central scalar is the retained \(-64\). Its logarithm is chosen
at that value, not at one. This untwists the actual suspension on
the source product, with its outside collar fixed.

\item At the node use the exact increasing orbit coordinates of F9,
\(2b=r-|y_0|^2/r\), \(2b'=|x_0|^2r-r^{-1}\), \(db'/db>0\).
They give the real gamma lift on the common short arc; the branch
\(S=-w\) is lower and \(S=w\) upper in the marked coordinate
\(Z\). Their full-period closing is the identity on this ordered
\(w\)-cover. A nonzero gamma deck would change the norm of the
strictly contracting line and is absent from that identity.
All source-line logarithms on the outside and the cap extend on
the disc by EC.3. A horizontal field chosen there and on their
common boundary therefore yields the actual ambient orbit-body
product \(Y\times D_w\). Its fibre is the known orbit body of the
constructed E6 sphere. The F9 source-product arrow supplies this
extension also at the cusp, without extending the complex selected
toric quotient there.

\item Fractional source translations act on the normalization by the exact
formula
\[
 S\longmapsto
 \frac{(1+g_t)S+(1-g_t)w}
 {((1-g_t)/w)S+(1+g_t)},\qquad
 g_t=\exp(t\log g),
 \tag{EC.22}
\]
where \(g(0)=1\). If \(g_t=1+wh_t\), its difference from \(S\)
is \(h_t(S^2-w^2)/(2+h_t(w-S))\). The full derivative is
\(I+O(|S|+|w|)\) uniformly on a common sufficiently small ball.
Thus the written branch-fixing isotopy removes this actual source
shear. In the universal real gamma cover interpolate the actual
heights, within the common short arc, to
\(\operatorname{Re}(S\overline w)/\sqrt{|S|^2+|w|^2}\).
The two double points remain strictly ordered, so this is an
isotopy through embeddings. Its outer collar is unchanged.

\item Remove the star-shaped four-ball
\(f^{-1}(B^6_\epsilon)\) from \(\mathbf P^1\times D_w\).
The graph is the same actual embedded cobordism as F56 on the
normalization (EC.4). In the rounded ambient product ball
\(B^6_\epsilon\times[-2\eta\epsilon,2\eta\epsilon]\), its inner
link is literally \((f,\eta H)\), since \(|H|\le |w|\).
This is the F55 lifted-link map, with the same complex domain
orientation and branch-height order. Its outside maps have been
specified in EC.3 and (EC.22), rather than inferred from the cubic
germ alone.

\par\medskip\noindent\textbf{EC.5. The actual circle line, both caps and the reference pair.}\par\nopagebreak[4]

\item The pullback of the actual \(M_6\) to
\(\mathbf P^1\times D_w\) has degree zero on every normalization
fibre. Evaluation in its zero-section Cartier frame gives the
relative nonzero section and hence its literal product line.
Its value at the reference parameter fixes the remaining disc
unit; the full Hom isotopy of EC.4 is retained in that frame.
The normal determinant has the same covector and actual divisor
isomorphism as (EC.6):
\[
 L=\mathcal O(-D_+-D_-),\qquad
 E=\mathcal O(D_+)\oplus\mathcal O(-D_-),\qquad
 \det E\xrightarrow{\ (S-w)/(S+w)\ }\mathcal O.
 \tag{EC.23}
\]
In the positive spinor the physical nonzero-vector line is the
tautological line \(M\); its meridian-one Hermitian dual is \(U\).
The actual coefficient arrow of K42 is
\[
 (a,b)\longmapsto
 \frac1{\sqrt{|x|^2+|y|^2}}
 \begin{pmatrix}y&-\overline x\\x&\overline y\end{pmatrix}
 \overline{(a,b)}.
 \tag{EC.24}
\]
Its source-unit multiplier from (EC.7) has an extending logarithm
and square root. After matching the product of the two normal
coordinates, a residual gauge is \((c_0,c_1)\mapsto(gc_0,g^{-1}c_1)\).
Its unitary part gives the literal Hopf arrow
\(q_H\mapsto q_H g/|g|\), a physical-circle map. Its possible
parameter winding extends over \(D^3\times S^1\). This retains
all separate factors of the actual line/normal frame.

\item The outer cap map uses
\[
 A_w=\begin{pmatrix}1&-w\\1&w\end{pmatrix},\qquad
 (P,q_H,w)\mapsto
 \left(\frac{A_wPA_w^*}{\operatorname{tr}(A_wPA_w^*)},q_H,w\right).
 \tag{EC.25}
\]
Its inverse is the displayed formula with \(A_w^{-1}\). It holds
the actual Cox quaternion fixed, not its tangent derivative, and
is periodic on the full parameter circle. The determinant factor
\(2w\) is retained in (EC.24). The inner cap is the identical
projection-normal four-ball, with the extending unit in (EC.7).
The anti-linear dual map of K68 commutes with the SU(2) clutches
and extends over both normal discs; it carries \(M=U^{-1}\).

\item We next verify the global orbit hypotheses used by the Thom step.
The reference completed E6 member is the F41 sphere with datum
\((a,b,m,n)=(6,1,2,n)\), \(n=\pm1\), and the actual selected-zero
Hom section of EC.2. Its application to the present source is
literal: join the present marked value
\(\rho=w^2/(w^2-3)\) to \(\rho=-1\) by an arc whose interior
avoids \(0,1,-1,\tfrac12,-\tfrac12\). Choose a contractible
tubular neighbourhood and its ordered-node square root. The
Miller/Cartier formulas, (EC.18)--(EC.21), and their inverses give
the completed family on that tube. At the endpoint \(\rho=-1\),
\(r=0\), the K94 unselected-III extension is the actual E6-B
source of F41. The fixed sufficiently small scalar makes this a
proper compact family along the arc, carrying the reference
cochain and character. Its Ehresmann transport identifies the
present reference member \(X\) with that sphere. No loop or
additional parameter return is inserted by this interval map.

\item The physical circle is free on the ordinary part. In the quadratic
cap the actual generator \(H=A^n\) satisfies \(H^2=B\) and
\(H^n=A\), and its strictly contracting norm prevents a nonzero
\(H\)-power from agreeing with a unit physical-circle rotation.
The standard cusp fan has edge directions
\(\pm e_1,\pm e_2,\pm(e_1-e_2)\). The determinant with the
physical cocharacter \(e_2\) is \(\pm1\) off the vertical
edges, whose period quotient is one fixed curve \(F=\mathbf P^1\).
The two charts
\[
 c_0'=Zc_0,\quad c_1'=Zc_1,\quad Z'=Z^{-1},
 \qquad q=Zc_0c_1=Z'c_0'c_1'
 \tag{EC.25a}
\]
have weights \((-1,+1,0)\). Thus the fixed normal bundle and
semifree action are precisely the displayed local input of
F7--8, including the physical circle generator.

\item Put \(E=X\setminus\operatorname{int}(F\times D^4)\) and
\(P=E/S^1\). Since \(X\cong S^6\), Alexander duality gives the
integral homology of \(S^3\) for \(E\), and codimension four
general position gives \(\pi_1(E)=0\). Conjugating the negative
normal line in (EC.25a) gives
\(\mathcal O(1)\oplus\mathcal O(-1)\). Its clutch
\(\operatorname{diag}(Z,Z^{-1})\) bounds the SU(2) disc with
coordinates \((r e^{i\theta},\sqrt{1-r^2})\); adjoining a path
from its centre to the identity gives the actual unitary
trivialization. In that trivialization the boundary quotient is
the product Hopf map
\[
 \partial E=S^2_F\times S^3
       \longrightarrow S^2_F\times S^2_{\rm mer}.
 \tag{EC.25b}
\]
Mayer--Vietoris for \(X=E\cup(F\times D^4)\) identifies
\(H_3(\partial E)\to H_3(E)\) as an isomorphism. Its restriction
to a normal \(S^3\) is the Hopf bundle, whose fibre integration
is the primitive meridian Euler generator. Naturality in the
integral Gysin sequence consequently makes
\(H^3(E)\to H^2(P)=\mathbb Z e\) an isomorphism. The same
sequence gives \(H^3(P)=H^4(P)=0\); the homotopy sequence gives
\(\pi_1(P)=0\). The Euler class restricts to the meridian
generator and has zero \(S^2_F\) coefficient, by (EC.25b).

\item For \(Y=P\cup(F\times D^3)\), the integral Mayer--Vietoris
map in degree two, with boundary basis \((F,\mathrm{meridian})\)
and target \((P,\mathrm{tube})\), is
\[
 \begin{pmatrix}0&1\\-1&0\end{pmatrix}.
 \tag{EC.25c}
\]
Its determinant is one. The boundary \(H_4=\mathbb Z\) maps to
zero and supplies \(H_5(Y)=\mathbb Z\); the other intermediate
groups vanish. Van Kampen gives \(\pi_1(Y)=0\).
The verified \(\Theta_5=0\) and Haefliger applications in F7--8
therefore identify the orbit pair as \((S^5,\text{unknotted }S^2)\).
On its free complement the actual Hopf bundle has Euler
generator. Its comparison with the fixed normal product differs
on \(S^2\times S^2\) by a circle gauge; that gauge has a global
logarithm because \(H^1(S^2\times S^2;\mathbb Z)=0\), and the
logarithm extends over the collar. These are the actual
equivariant reference maps used in (EC.24)--(EC.25).

\item The ambient meridian line extends on the excision complement by
the same integral Thom map as F56: the product orbit body has
\(H^2=H^3=0\), and the properly embedded connected normalization
has meridian evaluation \(H^2\cong\mathbb Z\). The actual outer
line has the specified value, and its ratio gauge extends over
the written outer cap. Thus the extra normal-sphere boundary is
the actual bundle over
\((S^2\times D^2)\cup(D^3\times S^1)=S^4\), with both caps and
their normal clutch fixed by (EC.24)--(EC.25).

\item Consequently every step of the F56/K42 excision uses the same
oriented link, coefficient matrix and cap morphisms. The relative
second Chern number is \(-1\), computed from (EC.24) with cap
boundary orientation negative to \(\mathbf P^1\times S^1\).
The normal-sphere pair is \((1,4)\) for \(U\); the inner lifted
link pair is \((5,-4)\) in its retained convention. Their actual
boundary signs and the Jin parameter reversal give the reference
\(U\)-pair \((6,0)\). The Hermitian dual map gives the physical
pair \((-6,0)\). The K68 first-to-second piece clutch is
\(\lambda_M(F)=-\Sigma_{F^{-1}}=\Sigma_F\); the return is the
one obtained from the delta-averaged product connection, identical
on the selected tube, so this is the actual clutched sphere. Its
reference class and native coordinates are therefore the K71 seed
\(-1\) and \((0,-6,-1)\), respectively. These conclusions use the
complete exterior and cap extension above.

\par\medskip\noindent\textbf{EC.6. The actual vertical seeds and complete fanwise trace comparison.}\par\nopagebreak[4]

\item First retain the complete vertical seed; the initial seed alone does
not specify a different vertical star. On \(T_j\) the literal
characters and periods give
\[
 V_j=Z^{-j}V,\quad H_j=HB^j,\quad B_j=B,\quad
 Q_j=\kappa qg^j,
\]
\[
 \alpha_j=\alpha\eta^j\kappa^{-j}g^{-j^2},
 \qquad\eta_j=\eta g^{-j}.
 \tag{EC.26a}
\]
For each fixed integer \(j\), the exact source interpolation is
\(Q_t=Q\exp(tj\log g)\). Near its node \(xy=Q_t\), the vector
field \(\partial_t+j\log g\,y\partial_y\) is tangent and keeps
\(q,w\) fixed. On the outside torus use the marked-period field
\(u+\tau_t b\mapsto u+\tau_0b\). Their difference on the
overlap is vertical. A partition of unity in this affine space
of horizontal fields, fixed on the two protected collars, supplies
the actual periodic flow. It transports the embedded normalization
graph by the ordered real-height isotopy of EC.4. The spinor arrow
is multiplication by \(Z^j\):
\[
 E\otimes\mathcal O(j(D_+-D_-))\longrightarrow E.
\]
Its two Cartier generators \(x^{-j-1}y^j,x^{-j}y^{j+1}\) go
exactly to \(1/x,y\). Its adjoint and the product cap (EC.25)
therefore carry the full normal frame, including (EC.7).

\item In the full real period coefficients \((u,b,d,\gamma)\), the
actual outer arrow is
\(u_j=u,b_j=b,\gamma_j=\gamma-jb,d_j=d-ju\).
The q-deck is \((u,d)\mapsto(u-b,d-\gamma)\), and the positive
w-deck is \((u,d)\mapsto(u+6b,d)\). These follow from (EC.12)
with every actual unit logarithm retained. In those coefficient
frames the connections are
\[
 A_0=dd+\alpha_q d\gamma,\qquad
 A_j=dd_j+\alpha_q d\gamma_j+6j\theta\,db,
\]
\[
 A_j-A_0=-j\{du+(\alpha_q-6\theta)db\}.
 \tag{EC.26b}
\]
The differentials of the zero-winding logarithms occur in the
definition of \(d\); using these coefficients retains them.
The only new integral theta-clutch is the dual-line factor
\(\mathcal T^{6j}\), where the coefficient of \(\mathcal T\)
changes by \(e^{2\pi i b}\). Its curvature is
\(\Xi=d\theta\wedge db\), with \(\Xi^2=0\).

\item Here the exact relative extension test of K66 applies to the
present source, with the following concrete maps. The graph map
\(b\) is constant on the attaching vanishing circle and extends
over its two-handle; it is made constant near the fixed graph by
a supported real lift and then extends over the product cap.
The integral boundary has
\(H_1=\mathbb Z\langle\alpha_q,\theta,b,\gamma\rangle\);
the q commutator kills \(u\) primitively. The restrictions from
the completed piece contain precisely \(\theta,b,\gamma\).
The remaining q-angle obstruction vanishes on the literal disc
\[
 (q/Z_*,Z_*,1)\mapsto(qZ_*^{j-1},Z_*,Z_*^{-j}),
 \qquad |Z_*|=1,\quad Z_*\text{ near }-1,\quad\gamma=0.
 \tag{EC.26c}
\]
For small \(w\), the marked Cartier sections have Tate points
near one, so this disc avoids them and the fixed tube. Its actual
line arrow extends with no zero. Subtract the integral one-cochain
from the logarithms of the ratio line, solve its real Cech
one-cocycle by a partition of unity, and correct its boundary value
by the extending \(\theta,b,\gamma\) characters and the remaining
real logarithm. This gives the prescribed relative isomorphism
\(U_j\to U_0\otimes\mathcal T^{6j}\), not just its Chern class.

\item The full tangent comparison is the source flow just written and
its actual positive triangular graph shear. The source-frame
matrix is upper triangular after ordering its coframe as in
K63 SQ.4; its strictly triangular curvature has zero trace-square.
The base connection and its interpolation depend on only the
two boundary angles and one collar variable, so its curvature
square vanishes as well. For the relative \(p_1\) pairing use the actual level
cobordism of K63. Choose \(U_1\Subset U_2\Subset U_3\)
around the fixed tube, with all surgery and connection changes
in \(U_1\), and take the real lift \(\widetilde b\) of the
original source-product circle map \(b_{\rm ang}\) on \(U_3\).
For a cutoff \(\chi=1\) on \(U_2\), supported in \(U_3\),
and a fixed real constant \(c\), define the full relative homotopy
\[
 b_t=b_{\rm ang}
       \exp\!\bigl(2\pi i\,t\chi(c-\widetilde b)\bigr),
 \qquad 0\le t\le1.
 \tag{EC.26c}
\]
Outside \(U_3\) it is literally \(b_{\rm ang}\); at \(t=1\)
it is constant on \(U_2\). Choose a common regular value different
from \(\exp(2\pi i c)\). Its preimage under
\((x,t)\mapsto b_t(x)\) is the relative level cobordism, fixed on
the outer boundary. The terminal level misses \(U_1\), so it is
the identical cycle before and after surgery. The original level
is the actual source three-level times \(S^1_\gamma\); its
source-pulled \(p_1\) integral is zero. Stokes on the displayed
cobordism, with the retained zero outer \(p_1\) comparison and
transgression, gives zero on the terminal level. This proves the
required relative pairing without identifying the terminal level
itself with a product.


\item The relative Euler cube has an affine interior term because
\(\Xi^2=0\). On the outward collar the exact difference (EC.26b)
has the same fifth-form coefficient \(18j^2\) as K66 EC.29;
the actual piece-to-exterior interpolation runs from one to zero,
so its physical contribution is \(-18j^2\). Thus the physical
cube has the form \(-6+Lj-18j^2\), and its Pontryagin pairing
stays zero. The two literal endpoint evaluations fix \(L\).
The E6 equation, the section, and all Miller factors depend only
on \(w^2\). The full reflection line arrow can be read directly
from the retained theta frame. One has
\[
 S_6(-w)=S_6(w),\qquad c_6(-w)=P^{-6}c_6(w),\qquad
 f_\theta(Z^{-1},P^{-6})=P^6f_\theta(Z,P^6).
\]
The last identity holds for the full theta product at the same
period \(Q\): use
\(\Theta(Z^{-1})=-Z^{-1}\Theta(Z)\) in its numerator and
denominator. In the first two identities \(P=P_G(0,w)\).
In the full theta identity and the following period arrows,
\(P=P_G(q,w)\) is the actual extending marked lift. The source
identity reverses the marked invariant differential and fixes
\(O\), so it sends this lift to \(P^{-1}\); the integral
Tate reindex is zero on the formal neighbourhood of \(O\)
specified by (EC.14a). Therefore the source/Cartier identity first has
the literal ordered-coordinate arrow
\[
 I:(Z,V)\mapsto(Z^{-1},P^6V).
\]
The actual target source period has coefficient
\(\alpha_{\rm target}=P^{-6}\). Composing that period with \(I\)
gives exactly
\[
 H_{\rm target} I:(Z,V)\mapsto(Q/Z,V).
\]
This retains the full height-one source-period reindex and all
theta factors. This composed half-turn arrow commutes with the
exterior and the finite cap (EC.19)--(EC.21), whose coefficients
are pulled back from the same source disc. This arrow sends
\(T_0\) to its literal reflected \(T_1\). It carries (EC.24)--(EC.25)
and the physical line, so both reference cubes are \(-6\), with
Pontryagin pairing zero. Hence \(L=18\), and on every \(T_j\)
the actual physical pair, class and native pair are
\[
 (k_M,\ell_M)=(-6-18j(j-1),0),\quad
 A_M=-1-3j(j-1),\quad(m_{\rm native},n_{\rm native})
       =(0,-6-18j(j-1)).
 \tag{EC.26d}
\]
All scalar/root and higher-unit terms in these maps are the ones
of EC.2--3. This establishes the complete K71 seed by the actual
E6 line and collar maps.

\item For a chosen lattice-basis fan \(P\), put \(R=P^{-1}\) and use
the literal characters \(Z'=Z^{R_{11}}V^{R_{12}}\),
\(V'=Z^{R_{21}}V^{R_{22}}\), with \(q\) unchanged. The new period
basis is \(H'_j=H^{P_{1j}}B^{P_{2j}}\). Its full coefficients are
\[
 A_{ij}=(\kappa^{P_{1j}}g^{P_{2j}})^{R_{i1}}
              (\alpha^{P_{1j}}\eta^{P_{2j}})^{R_{i2}},
 \qquad B_P=-6P^{-1}e_1e_1^tP.
 \tag{EC.26}
\]
The q-exponents are \(RP=I\). Every cone map and its inverse is
the corresponding integral semigroup map, carrying all full unit
powers. The outside arrow is its inverse character map in the
original source and Cartier coordinates, so the cap squares of
EC.3 are unchanged.

\item The three exact edge transitions, in the two vertex charts, are
\[
 (u,v_2,w_2)=
 (A_{11}A_{21}bc,A_{21}^{-1}ab,A_{11}^{-1}b^{-1}),
\]
\[
 (A_{12}A_{22}bc,A_{22}^{-1}c^{-1},A_{12}^{-1}ac),
 \qquad(a^{-1},ab,ac).
 \tag{EC.27}
\]
Their full normal matrices are
\(\operatorname{diag}(A_{22}/A_{21},A_{11}A_{21})\),
\(\operatorname{diag}(A_{11}/A_{12},A_{12}A_{22})\), and
\(\operatorname{diag}(A_{22},A_{11})\). Their determinant winding
is \(-6\). Retain their first winding numbers
\(l=(B_{22}-B_{21},B_{11}-B_{12},B_{22})\).
The complete unit contractions are the GF.9 maps
\(H_d=(\sin\vartheta e^{id\theta}+j\cos\vartheta)
(\sin\vartheta+j\cos\vartheta)^{-1}\), with left \(H_{-3}\)
and right \(H_{l_i+3}\), multiplied by the based logarithms
specified in EC.2 and their actual radial cutoffs. The exact tube
identity \(zD_{i,*}K_i=zD_i\) matches both vertex collars. Normal
rescaling and the cutoff flow of its velocity carry the full
higher normal jets and preserve their overlap maps.

\item Thus the actual ordinary/native boundary comparison is precisely
GF.12, with left degree \(-6\), right degree zero, and inverse
cylinder correction \(L^6\). Physical weights, odd Spin endpoints,
the capped edge traces and all nine prism equations are the same
integer equations of K71, derived from these actual matrices.
The local Cox ruling change is the same map on its rank-one matrix,
with common link \((h,q_H)\mapsto(q_H^{-1}hq_H,q_H)\); it is
independent of the source parameter and intertwines every full
coefficient and closing map in (EC.26). Its final native reset
retains the physical circle and central generator orientation.
Hence finite Euclidean words construct the same complete K71
algorithm on the present E6 family, starting from the actual seed
of EC.5. This proves equality with the K71 actual return on each
specified fan by these source/exterior/cap/Cox arrows. It assigns
no return to any other E6 puncture.

\par\medskip\noindent\textbf{EC.7. Audit scope and exact certificate.}\par\nopagebreak[4]

\item The certificate \path{code/part3_fourfold/e6_selected_cusp.py} checks
the full source-height equation and inverse polynomial relation,
normalization and normal covector, critical separation, actual
Miller values, source translation at the cusp, original/root
coordinate scaling, the full selected scalar cancellation, and
the exact period-conjugation and representative Cox identities.
The exterior, integral cochain, both normal-cap and physical-clutch
arguments are written above for geometric audit. The source
equation alone is not used to replace those arrows.

\end{enumerate}
\end{proof}




\begin{theorem}\label{night:P3K.96}
\StatementRecord{night:P3K.96}{proved}{S6\_06\_BLIND: GPT-6 Astra (OpenAI, ChatGPT Pro mode, effort Pro), recorded re-derivation in the checkpoint dated 18 November 2026; also checked in this writing by \CurrentChecker}{\cref{night:P3F.35}, \cref{night:P3F.41}, \cref{night:P3F.46}, \cref{night:P3K.94}}
Retain the actual source
\[
 y^2=x^3+r x^2+t x+t^2,\qquad G=(0,t),\qquad
 M_a=\mathcal O(aG-O),\qquad Q=bG,
\]
on a sufficiently small disc about \(r=0\). Choose the analytic branch
\(\rho^2=1-3r\), \(\rho(0)=-1\), and the selected whole
\(I_1\) zero \(t_c=(1-3\rho^2+2\rho^3)/27\). For either of the
complete cap progressions
\[
\begin{array}{c|c|c|c}
(a,b)&m&\text{root fibre}&\text{range}\\\hline
(6,1)&3l+2&IV&l\ge0\\
(2,3)&3l&\text{smooth elliptic}&l\ge1
\end{array}
\]
use the exact Cartier suspension
\(B_{\rm sel}=(t-t_c)B_{a,b}/t^3\) of the proof, either primitive
norm \(n=\pm1\), every prescribed character \(\zeta_m^d\), and
the entire specified reference annular map. Retain a fixed sufficiently
small nonzero scalar for the chosen fan, and no extra logarithmic disc
at the colliding fibres. For every prescribed lattice-basis fan at the
selected cusp, these actual maps extend to a smooth proper oriented
sphere family over the full source disc, with identity parameter
closing. Its positive boundary return is \(0\in\Theta_7\).
The literal marked-cover map \(D=w^2-3\),
\(r=(3-2w^2)/D^2\), \((X,Y,T)=(D^2x,D^3y,D^3t)\)
carries the same conclusion to both points \(w^2=3/2\), with the
full scalar, line, cap and fan marking retained.
\end{theorem}
\begin{proof}
\begin{enumerate}[label=\arabic*.]
\item \par

\medskip\noindent\textbf{EX.1. Full point-line multiplier and its finite vertical divisor}

\par\nobreak

\item Use the same actual source
\[
 y^2=x^3+r x^2+t x+t^2,\quad G=(0,t),\quad
 s=\tfrac14-r,\quad U=y-t-x/2,\quad
 f_2=\frac U{x-s},\quad T(p)=p-G.
\]
The typed point-to-tensor arrow is
\(r_2=f_2^{-1}:M_2\longrightarrow M_1^{\otimes2}\).
It gives the literal single-step lift
\[
 B_{2,1}(p)=f_2(p)^2\frac{f_2(Tp)}{f_2(p)}
          =f_2(p)f_2(Tp).
\]
Composing the three actual source/line arrows, without cancelling
their source or target Cartier factors separately, gives
\[
 B_{2,3}(p)=f_2(p)f_2(T^3p)f_2(Tp)^2f_2(T^2p)^2.
 \tag{EX.1}
\]
Its source map is exactly \(T^3=p-3G\).

\item At a generic I2 identity component \(C_0\), the functions \(U\)
and \(x-s\) are generic units, so \(\operatorname{ord}_{C_0}f_2=0\).
At the other component use the actual chart \(x=t\alpha,y=t\beta\):
\[
 f_2=t\frac{\beta-1-\alpha/2}{-s+t\alpha}.
\]
The displayed numerator and denominator are generic units and
\(s\ne0\), hence \(\operatorname{ord}_{C_1}f_2=1\).
The exact inverse translation of K94 EU.0 exchanges the two
components. Therefore the two valuations of \(B_{2,1}\) are
both one, and those of EX.1 are both three. The whole finite
vertical divisor at zero is exactly \(3F_0\).

\item At any other finite irreducible fibre, neither \(U\) nor \(x-s\)
contains that fibre: substituting \(U=0\) in the source gives
\(x^2(x-s)=0\), and the fibre is not one of these point divisors.
Every power of \(T\) preserves the generic component DVR and its
base uniformizer. Thus each factor in EX.1 is a DVR unit there.
At infinity \(3G\) fixes each component, so the Hom multidegrees
vanish. The primitive affine component kernel gives a whole fibre
there; its total Hom degree is \(2\cdot3\cdot(1/6)=1\).
Consequently its remaining vertical divisor is \(-2F_\infty\).
The exact selected-zero map is therefore
\[
 B^{(2,3)}_{\rm sel}=\frac{t-t_c(r)}{t^3}B_{2,3}.
 \tag{EX.2}
\]
The actual simultaneous A1 source, the maps \(T,T^{-1}\), and
the section \(2G=(s,-t-s/2)\) extend on the r-disc of K94.
EX.2 and its inverse have zero order at the parameter divisor,
as seen on an ordinary source section avoiding the horizontal
divisors. The only remaining locus is the codimension-two central
curve. The same height-one-prime argument as K94 extends this
actual Hom map and its inverse across it, with its prescribed
selected zero retained. No parameter scalar is changed.

\par

\medskip\noindent\textbf{EX.2. The literal good root and its central multiplier}

\par\nobreak

\item For \(m=3l\), \(l\ge1\), the actual good root is
\[
 q_\infty=z^{3l},\quad x_0=z^{4l}X,\quad y_0=z^{6l}Y,
 \qquad Y^2=X^3+r z^{2l}X^2+z^lX+1.
 \tag{EX.3}
\]
It has \(G'=(0,1)\) and
\[
 2G'=(s z^{2l},-1-(s/2)z^{3l}).
\]
The deck acts by \((z,X,Y)\mapsto
(\zeta_m z,\zeta_m^{-4l}X,Y)\), retaining the two marked
three-torsion limits and the actual line \(M_2\).
The original coordinates are \(x=z^{-2l}X,y=z^{-3l}Y,t=z^{-3l}\).
Hence each tangent-over-vertical factor has exact weight
\(f_{2,\rm old}=z^{-l}f'_2\). All six factors in EX.1 give
\[
 B_{2,3,\rm old}=z^{-6l}B'_{2,3},\qquad
 B^{(2,3)}_{\rm sel}=h B'_{2,3},\qquad
 h=1-t_c(r)z^{3l}.
 \tag{EX.4}
\]
This is the full coefficient scaling in the actual rational
source and target Cartier frames.

\item On the central good elliptic curve \(Y^2=X^3+1\), put
\(f=(Y-1)/X\), \(T'=p-G'\). Since \(G'\) has order three,
\(T'^3=1\), and the product of the three translated divisors
of \(f\) is zero. Its norm is therefore a nonzero constant.
The actual orbit of \((2,3)\) is
\((2,3),(2,-3),(-1,0)\); the three values of \(f\) are
\(1,-2,1\). Thus its norm is exactly \(-2\), and
\[
 B'_{2,3}|_{z=0}
   =[f(p)f(T'p)f(T'^2p)]^2=4.
 \tag{EX.5}
\]
The central selected coefficient in EX.4 is consequently \(4h\),
with its full value retained.
The same constant follows from the entire actual central map:
\[
 T'(X,Y)=\left(\frac{2X}{Y-1},-\frac{Y+3}{Y-1}\right),\qquad
 T'^2(X,Y)=\left(-\frac{2X}{Y+1},\frac{Y-3}{Y+1}\right).
\]
These formulae have \(T'^3=1\), and the three factor values as
rational functions are \((Y-1)/X,-(Y+1)/X,2/X\). Their product
is \(-2\) by the exact curve equation \(Y^2-1=X^3\).

\item The section \(Q'=3G'\) is in the identity formal group, with
\(\tau(Q')=-s z^{2l}/2+O(z^{5l})\). Its small analytic
\(m\)-division exists and is deck invariant by uniqueness.
The full good-root point-line correspondence, its invariant toric
correction and norm lift are constructed in EX.5. EX.6 carries
the prescribed finite character and entire annular reference
through the actual disc attachment.

\par

\medskip\noindent\textbf{EX.3. Marked-cover coefficient check}

\par\nobreak

\item Under the source map \(X=D^2x,Y=D^3y,T=D^3t\), every factor
\(f_2\) has weight \(D^{-1}\), so EX.1 has weight \(D^{-6}\).
The selected-height factor has weight \(D^6\), exactly as in
K94 EU.14. Therefore EX.2 is carried with scalar unchanged.
This actual Miller/Cartier coefficient equality carries the
constructed disc family to the marked w-cover without changing
its scalar or selected toroidal fan.

\par

\medskip\noindent\textbf{EX.4. Every reduced root: actual middle charts and full line graph}

\par\nobreak

\item The initial local stage is preserved in \path{records/attempts/e6_unselected_iii_extension/initial_local.md}. The following maps complete both cap progressions.
Keep the same source-parameter disc, selected whole zero, fixed
small scalar and full reference annular map as in K94. The two
named data are
\[
 (a,b;m)=(6,1;3l+2),\quad l\ge0,
 \qquad (a,b;m)=(2,3;3l),\quad l\ge1,
\]
with either signed sphere norm \(n=\pm1\) and the retained finite
character \(\zeta_m^d\). There are no logarithmic discs at the
colliding fibres. The following maps construct the full marked disc family.

\item For the reduced branch put \(m=3l+2,k=2l+1\). The literal root is
\[
 q_\infty=z^m,\quad x_0=z^{2k}X,\quad y_0=z^{3k}Y,
\]
\[
 Y^2=X^3+r z^{2l+2}X^2+z^{l+2}X+z^2,\qquad G'=(0,z).
 \tag{EX.6}
\]
Write \(b_0=z^l,c_0=r z^{2l}\) and
\(c=\sqrt{1+b_0X+c_0X^2}\). The exact coordinates
\(U=Y-zc,V=Y+zc\) satisfy \(UV=X^3\). The equation
\(2zc=V-U\) has a unique analytic inverse for \(z\) near the
origin, since its derivative there is two. Thus the same three
minimal A2 resolution charts give a smooth family over the entire
source disc, retaining the actual function \(q_\infty=z^m\).

\item In the middle chart \(X=ab,U=a^2b,V=ab^2\), set
\[
\begin{aligned}
 K&=1+\frac{b(b-a)(b_0+c_0ab)}{2c(c+1)},\\
 F&=1-\frac{ab(b_0+c_0ab)(K+1)}{2c(c+1)}
           -\frac{c_0ab^2(a-b)}{4c^2},\\
 X'&=a(a-b)F,\qquad c'=\sqrt{1+b_0X'+c_0X'^2},\\
 H&=K+\frac{b(a-b)(b_0+c_0X')}{2c(1+c')}.
\end{aligned}
\]
The actual \(+G'\) chord has \(\kappa=(Y-z)/X=aK\),
\(X'=\kappa^2-c_0z^2-X\), \(Y'=-z-\kappa X'\).
Its complete toric chart arrow is
\[
                  a'=(b-a)F/H,\qquad b'=-aH.
 \tag{EX.7}
\]
As in K94, \(V'=Y'+zc'=-aX'H\) proves the formula, and the
analytic inverse of \(2zc=V-U\) proves that it retains exactly
the same \(z\). All displayed denominators are units and
\(K=F=H=1\) at the triple point. The inverse is obtained from
the actual \(-G'\) chord, equivalently by \(Y\mapsto-Y\).
The smooth group charts at the other central points complete
this actual automorphism.

\item The full deck is
\[
 D(z,X,Y)=(\zeta_mz,\zeta_m^{-2k}X,\zeta_m^{-3k}Y).
\]
Since \(l-2k=-m\), \(c\) is deck invariant; \(U,V\) have
weight \(\zeta_m\). In the middle chart both \(a,b\) have
weight \(\zeta_m^{l+1}\). The identities
\(3(l+1)=m+1\) and \(\gcd(l+1,m)=1\) verify the actual
order and inverse-base relation. Every term of \(K,F,H\) has
weight zero. Formula EX.7 therefore commutes with this whole
deck, including its parameter coordinate.

\item The invariant vertical field has an explicit regular lift there.
Put \(J_0=b_0+2c_0ab\). The usual Weierstrass vector field gives
\[
 W=a\left(2b-a-\frac{zJ_0}{c}\right)\partial_a
   +b\left(2a-b+\frac{zJ_0}{c}\right)\partial_b.
 \tag{EX.8}
\]
This follows from
\(W(U)=3X^2-zJ_0U/c\), \(W(V)=3X^2+zJ_0V/c\).
It is holomorphic and vanishes at the triple point. Its flow
agrees with the identity-component group flow on the smooth
locus. Thus the exact division section
\[
 R=2G'-[m]^{-1}((2m-1)G'),\qquad mR=G'
 \tag{EX.9}
\]
extends as an actual source translation on the full resolved
root. Here \(2m-1=3(2l+1)\), so the argument of the inverse is
in the small identity formal group; the inverse has derivative
\(1/m\). The central component of \(R\) is two. Both EX.9
and the resulting translation commute with \(D\), by uniqueness
of that formal inverse and the displayed deck arrows.

\item For the line, retain the exact source coordinate
\(\tau_0=-x_0/y_0\) and root coordinate \(\tau'=-X/Y\).
Their relation and the full section equation are
\[
 \tau'=z^k\tau_0,\qquad
 \tau'(6G')=z^k L(s,z^m),\qquad L(s,0)=-s\ne0.
 \tag{EX.10}
\]
Use the normalized closure of the graph of the displayed generic
map to the actual minimal regular source. Its two projections
are the morphisms denoted \(\beta\) to the root and \(\pi\)
to the source. Near the marked sections this is the graph of
\(u=\tau'/z^k\), constructed by the actual ideal
\((\tau',z^k)\). The root section equations are
\(\tau'\) and \(\tau'-z^kL(s,z^m)\). The normalized graph adds the ray \((k,1)\). Its cone with
\((1,0)\) is regular. Subdivide the other cone by the literal
rays \((j,1)\), \(1\le j\le k\), followed by \((0,1)\).
Adjacent ordered determinants are
\(j\cdot1-(j-1)\cdot1=1\); these are exactly the exceptional
rays of this chosen regular graph. On the ray \((j,1)\),
\(\operatorname{ord}(\tau')=j\) and
\(\operatorname{ord}(\tau'-z^kL)=\min(j,k)=j\).
At the equality ray \(j=k\),
the graph coordinate \(u\) retains their distinct directions
\(0\) and \(L(s,q_\infty)\); its two strict section equations
are carried separately. The literal coefficient identity is
\[
 \frac{\tau'}{\tau'-z^kL(s,z^m)}
       =\frac{u}{u-L(s,q_\infty)}.
 \tag{EX.11}
\]
On the other graph chart, \(\tau'v=z^k\), the same actual
coefficient is \(1/(1-vL(s,z^m))\). It is a unit at \(v=0\).
The two equations agree by \(uv=1\); the remaining cyclic
singularity of \(\tau'v=z^k\) is resolved by its ordinary
monomial subdivision, and both coefficient arrows pull back
unchanged to every resulting chart. This explicitly retains
the direction at equal valuation as well as both outer charts.
Both marked divisors are absent from the other source/root
quotient centres. Resolving those remaining cyclic charts by
their invariant monomial ideals does not add a coefficient to
either divisor. Thus on the full common graph, and on each of
its such refinements,
\[
                \beta^*N_6=\pi^*M_6.
 \tag{EX.12}
\]
This equality is an equality of Cartier line maps, not only of
generic fibre classes. The source and root translations preserve the graph by the
following actual square. Since the root translation commutes with
\(D\), it and its inverse descend to the finite quotient. An automorphism
of this quotient lifts to its minimal resolution: the local cyclic
resolution is given by the invariant monomial ideals already
specified, and their exceptional valuations are preserved. The
subsequent contraction to the unique minimal regular source is
also preserved, since the product carries each contractible
vertical curve and its intersection data to that of its image.
Consequently there are actual mutually inverse source maps over
\(q_\infty\), extending the descended translation. The product of the
source and root maps carries the generic graph to itself, as does
its inverse. Taking its closure and normalization therefore gives
both inverse graph maps. The marked-section graph ideals and the
displayed cyclic subdivisions carry these maps, retaining the
actual Cartier coefficient square on every chart.

\item The exact root multiplier is also uniform in \(l\). The original
coordinates satisfy \(x=z^{-2(l+1)}X,y=z^{-3(l+1)}Y,t=z^{-m}\).
The Miller weights therefore give
\[
 B_{\rm sel}'=h z^{-2}B'_6,\qquad h=1-t_c(r)z^m.
 \tag{EX.13}
\]
For a generic root identity-component point \(X=w^2,Y=w^3+O(z^2)\),
the full three-point factor is
\[
 f'_3=\frac{s z^k X(Y-z-z^{l+1}X/2)}
              {s z^k(Y+z+z^{l+1}X/2)-(s z^m+2)X}.
\]
The actual inverse chord gives
\[
 f'_2\to w,\quad f'_3(p)/z^k\to-s w^3/2,\quad
 f'_3(Tp)/z^{k+1}\to s,
\]
\[
 z^k f'_{2,3G'}(p),\ z^k f'_{2,3G'}(Tp)\to-4/s.
\]
Consequently \(B'_6/z^2\to4\). The full divisor calculation
and EX.12 give order two on every reduced central component;
normality then extends EX.13 and its inverse as actual units.
Their central coefficient is \(4h\), with no omitted base
factor. The case \(l=0\) is the actual quadratic map of K94.

\par

\medskip\noindent\textbf{EX.5. The full good-root Cartier correction and actual translation lift}

\par\nobreak

\item For the good branch EX.3 write \(u=z^l\). Its quotient to the
source factors through the actual minimal cubic root
\[
 Y^2=X^3+r u^2X^2+uX+1,\qquad q_\infty=u^3.
\]
The \(\mu_l\) kernel acts only on \(z\), and the remaining
\(\mu_3\) action has weights \((2,1)\) at each of its three
fixed sections. At \(2G'\) one can use
\(w=X-su^2\); at \(O\) use \(w=-X/Y\). They have exactly
weight two, with \(u\) of weight one. Their quotient lattice is
\(\mathbb Z^2+\mathbb Z(2,1)/3\), with the two primitive
resolution rays \((2/3,1/3),(1/3,2/3)\).

\item Here is the actual strict-section Cartier calculation. On that
resolved quotient its support function is
\(h(v_w,v_u)=\max(v_w-2v_u,0)\): it is \(v_w-2v_u\)
on the cone adjacent to the section axis and zero on the other
cones. Pulling the section to the common root fan has coefficient
\(v_w\). Their difference is therefore
\(\min(v_w,2v_u)\). In the full \(z\)-root it is
\(\min(v_w,2l v_z)\). If \(E_P\) and \(E_O\) denote
these explicit sums at the two marked sections, the full equality is
\[
          \beta^*N_2=\pi^*M_2+E_P-E_O.
 \tag{EX.14}
\]
The coefficient arrow is the quotient of the actual monomial
Cartier generators on each cone, so it includes their units.
The same projective charts and monomials work holomorphically
over the whole \(r\)-disc; their marked centres are the actual
sections just displayed.

\item The section \(Q'=3G'\) has
\(\tau(Q')=-s u^2/2+O(u^5)\). Its selected small
\(m\)-division \(R=[m]^{-1}Q'\) has
\(\tau(R)=-s u^2/(2m)+O(u^5)\), and is deck invariant.
At each marked fixed point its translation has the exact form
\(w\mapsto w\varepsilon(w,u,r)+u^2 a(w,u,r)\), with
\(\varepsilon\) an analytic unit. It preserves the actual
ideals \((w,u^j)\) and \((u,w^j)\) for \(j=1,2\), as
does its inverse. Their invariant product and normalized blowup
give the common toric refinement above. Thus translation by
\(R\), as well as by \(Q'\), lifts to that graph and preserves
both corrections in EX.14. This proves the full equivariance of
the point-line correspondence. It retains the translated
correction through its actual canonical divisor section.

\item Both branches admit a particularly explicit starting line lift.
For \(N=\mathcal O(P-O)\), define the actual chord-over-vertical
function
\[
 g_{P,R}(p)=\frac{\ell_{P,R}(p)}{x(p)-x(P+R)}.
 \tag{EX.15}
\]
Its divisor is \((P)+(R)-(P+R)-(O)=D_P-t_{-R}^*D_P\),
so it is the coefficient of a typed line map covering
\(p\mapsto p-R\). The points used here are distinct on a
sufficiently small punctured root disc. The line functions are
the literal chord functions; no leading coefficient is removed.
On good roots take \(P=2G'\) and
\(F=\tau(R)g_{P,R}\). On reduced roots take \(P=6G'\) and
\(F=\tau(P)g_{P,R}\). The formal coordinate factors have
orders \(2l\) and \(k\), respectively, with nonzero leading
coefficients \(-s/(2m)\) and \(-s\).

\item To check the full central normalization, a point tending to
\(O\) has \(x=\tau^{-2}+O(1)\),
\(y=-\tau^{-3}+O(\tau^{-1})\). The actual chord slope is
\(-\tau^{-1}+O(1)\). At a generic test point, multiplying
the chord by \(\tau\) tends to its vertical denominator.
Thus the displayed \(F\) has central coefficient exactly one.
In the reduced case the Hom multidegrees are zero, so the
primitive central component kernel makes the chord\textquotesingle s vertical
order a whole-fibre order; the single generic-component limit
fixes it on every component. In the good case there is one
smooth component. Horizontal divisors are exactly those in
EX.15. The height-one-prime argument therefore extends \(F\)
and its inverse through the whole root family, retaining its
full normalizing factor.

\item This particular lift commutes with the deck literally. The
root Weierstrass coordinate weights are \((-2k,-3k)\), with
\(k=2l\) in the good case and \(k=2l+1\) in the reduced case.
Consequently the actual chord-over-vertical function has weight
\(\zeta_m^{-k}\), whereas each of \(\tau(P),\tau(R)\)
has weight \(\zeta_m^k\). Their displayed product has weight
one. These are identities of rational coefficient maps, so
their extending maps commute on every root chart.

\item The ratio \(\beta=B'_{{\rm sel}}/F^m\) covers the identity
on the same compact connected source fibre. It is therefore
an invertible base function, with central value four. It is
deck invariant by the preceding exact arrows. Its analytic
\(m\)-th root, fixed by the complete reference lift, is deck
invariant as well: the ratio of its transform to itself is a
constant \(m\)-th root of unity, equal to one on the central
disc. Define
\[
 \widetilde F=\beta^{1/m}F,\qquad
 A_0=\zeta_m^d\widetilde F^{\,n}D.
\]
Then, as literal source/Cartier maps,
\[
 \widetilde F^m=B'_{{\rm sel}},\qquad
 A_0^m=(B'_{{\rm sel}})^n,\qquad
 [A_0,B'_{{\rm sel}}]=1.
 \tag{EX.16}
\]
The reference finite character is still \(\zeta_m^d\).
The normalized source/tensor factors in EX.1 and the full
correction EX.14 are carried through this equality. No generic
Picard-class comparison replaces their coefficient maps.

\par

\medskip\noindent\textbf{EX.6. General-order annular square, actual disc closing and return}

\par\nobreak

\item Fix \(R_\lambda\ne0\) with \(R_\lambda^m=\lambda\),
independent of the source parameter. Put
\[
 B_\lambda=\lambda B'_{{\rm sel}},\qquad
 A_\lambda=R_\lambda^n A_0.
\]
On the two fixed annulus sectors use
\[
 s_\alpha=\frac{\log_\alpha q_\infty}{2\pi i},\quad
 z_\alpha=\exp(2\pi i s_\alpha/m),\quad
 c_\lambda=\frac{n\log R_\lambda}{2\pi i}+\frac dm.
\]
The actual map in the source/line Lie cover is
\[
 \Phi_\alpha=\operatorname{Exp}
       (n s_\alpha f_\alpha+s_\alpha c_\lambda\delta
                   +h_\alpha)I_{z_\alpha},
 \qquad f_\alpha=\operatorname{Log}_\alpha\widetilde F.
 \tag{EX.17}
\]
The source displacement is the signed \(-n s_\alpha\) Abel
integral of the actual \(R\). The scalar is
\(R_\lambda^{n s_\alpha}\exp(2\pi i d s_\alpha/m)\),
and the root change is the actual \(D\). Thus both the full
norm and every finite character appear in this formula.

\item On an overlap, with \(s_\beta=s_\alpha+k\) and
\(f_\beta=D_*^k f_\alpha+\Pi_{\alpha\beta}\), prescribe
\[
 h_\beta-D_*^k h_\alpha
   =-n s_\alpha\Pi_{\alpha\beta}
       +j_{\alpha\beta}b_\beta+\Lambda_{\alpha\beta},
 \qquad b_\beta=\operatorname{Log}_\beta B_\lambda.
 \tag{EX.18}
\]
All period columns, suspension decks and integral coefficients
are those of the full reference map. Their functions are
transported holomorphically over the source disc. Here is the reference-zero construction in the same fixed
ordinary Lie frame. Extend the prescribed reference cochains
\(h_\alpha^*\) to \(h_\alpha^0\), retaining the full
transported period functions. Let \(E_{\alpha\beta}\) be
the residual of EX.18; it vanishes at the reference parameter.
For the fixed annulus partition put
\[
 k_\alpha=-\sum_\beta\chi_\beta E_{\alpha\beta},
 \qquad \omega=\bar\partial k_\alpha,
 \qquad v(q,r)=\frac1\pi\int
                \frac{\omega(w,r)}{q-w}\,dA(w).
\]
The compactly supported residual and fixed contours are exactly
those of K94 EU.11e. The common \(\omega\) and its Cauchy
transform depend holomorphically on the source parameter and
vanish at the reference. Then
\(h_\alpha=h_\alpha^0+k_\alpha-v\) solves EX.18,
retains every integral period and suspension column, and equals
the whole prescribed reference cochain. Subtracting the target-sheet logarithm
\(kn f_\beta+kc_\lambda\delta\) from the two sides of
EX.17 leaves exactly the prescribed period and suspension deck.
Hence the actual squares are
\[
 \Phi_\beta=A_\lambda^k B_\lambda^{j_{\alpha\beta}}
                  \Phi_\alpha,\qquad
 \Phi_\alpha B_\lambda=B_\lambda\Phi_\alpha,
 \qquad A_\lambda^m=B_\lambda^n.
 \tag{EX.19}
\]
They define the ordinary/cap gluing with its actual inverse
and its full reference value. The finite cap itself extends
by EX.3, EX.6-\/-16, without extending an annular logarithm
across its centre.

\item All these maps are single valued on the whole source disc.
Their positive boundary therefore has identity source, line,
root and annular closing; the scalar is fixed. The selected
cusp remains disjoint from the collision, and its full Tate,
Cartier and cone-character maps extend exactly as in K94 EU.5
on every prescribed lattice-basis fan. On the finite charts
their actual line multipliers have compact uniform bounds.
Choosing the original fixed scalar sufficiently small gives
strict contraction. For both norm signs \(H=A_\lambda^n\)
has \(H^m=B_\lambda\) and scalar factor \(R_\lambda\).
Its compact norm shell and the protected toric fundamental
regions give the same uniform properness, freeness and smooth
disc submersion as the explicit quotient proof of K94.

\item At the central source the two data are exactly the E6-B
entries of F41/F46, with the retained ordinary III cap of
order one and norm zero. Their determinant is \(n=\pm1\),
so the full integral maps give a sphere on the reference fan.
The actual finite fan word, common blowups, two ruling maps
and primitive Gysin/van Kampen maps of F35 carry that sphere
to each prescribed fan. Radial transport in the resulting
proper disc family gives a product with its actual identity
closing. Thus the positive unselected-III boundary return is
zero for both entire cap progressions just constructed.
The quotient, line, annular and topology maps above prove the stated return for the full marked family.

\par

\medskip\noindent\textbf{Executed certificate and retained attempts}

\par\nobreak

\item The exact program \path{code/part3/e6_unselected_iii_extension.py}
checks the full central translation orbit, generic three-step
Cartier multiplier, symbolic root exponents, universal A2 chart
identities, regular graph rays and section coefficients, deck and
line weights, both signed norms and the target-sheet overlap
polynomial. Its current output records 5,241 exact identities
and finite character instances; the complete geometric proof above has the recorded root and independent peer audits.
The initial local stage and the symbolic-power simplification
failure are preserved under
\path{records/attempts/e6_unselected_iii_extension/}.

\end{enumerate}
\end{proof}




\begin{theorem}\label{night:P3K.97}
\StatementRecord{night:P3K.97}{proved}{S6\_06\_BLIND: GPT-6 Astra (OpenAI, ChatGPT Pro mode, effort Pro), recorded re-derivation in the checkpoint dated 18 November 2026; also checked in this writing by \CurrentChecker}{Definitions and the displayed constructions}
Let \(w\) vary on a sufficiently small disc about either point \(w^2=3\), and set
\[
 D=w^2-3,\quad R=3-2w^2,\quad v=w^2-1,\quad
 r=R/D^2,\quad t_c=v/D^3.
\]
On its punctured disc retain the actual original source
\[
 y^2=x^3+r x^2+t x+t^2,\qquad G=(0,t),\qquad
 M_a=\mathcal O(aG-O),\qquad Q=bG,
\]
with the whole selected zero \(t=t_c\), the full Cartier suspension
\(B_{\rm sel}=(t-t_c)B_{a,b}/t^3\), and either complete cap progression
\[
 (a,b;m)=(6,1;3l+2),\ l\ge0,\qquad
 (a,b;m)=(2,3;3l),\ l\ge1.
\]
Use either signed primitive norm \(n=\pm1\), the prescribed finite
character \(\zeta_m^d\), a fixed sufficiently small nonzero scalar,
the entire reference annular map, and no additional logarithmic
sites. For every prescribed lattice-basis fan at the selected cusp,
this punctured family is isomorphic, by the actual source, Cartier,
root, physical-circle and annular maps of the proof, to the restriction
of a smooth proper oriented sphere family on the full native disc
\[
 Y^2=X^3+R X^2+DTX+T^2,\qquad G=(0,T),\qquad q=T-v.
\]
The comparison retains both point and tensor Tate frames, the full
finite-root character, every period and all parameter scalar factors.
It intertwines the specified forward endpoints and hence their inverse
physical closings. The actual positive marked \(w\)-loop has return
\(0\in\Theta_7\).

The original infinity parameter is exactly
\(\epsilon=1/r=D^2/(3-2w^2)\). Each such positive \(w\)-loop
has degree two over the original positive infinity meridian.
No shorter original-parameter endpoint or return is inferred.
\end{theorem}
\begin{proof}
\begin{enumerate}[label=\arabic*.]
\item \par

\medskip\noindent\textbf{DZ6.1. The native disc and actual point-line maps}

\par\nobreak

\item Put
\[
 D=w^2-3,\qquad R=3-2w^2,\qquad v=w^2-1,\qquad
 s_0=D^2/4-R=(w^2-1)(w^2+3)/4.
\]
On a small disc about either of the points \(w^2=3\), use
\[
 Y^2=X^3+R X^2+DTX+T^2,\qquad G=(0,T),\qquad q=T-v.
 \tag{DZ6.1}
\]
Here \(s_0=3\), \(v=2\) at the centre. The selected nodal
fibre remains disjoint from the I2 fibre at \(T=0\). The
actual source A1 coordinates are
\[
 U=Y-T-DX/2,\quad V=Y+T+DX/2,\quad UV=X^2(X-s_0).
\]
Thus the same simultaneous A1 resolution, using the extending
unit square root of \(X-s_0\), is defined on the whole disc.
The exact Miller functions are
\[
 f_2=\frac{U}{X-s_0},\qquad
 f_3=\frac{s_0 XU}{s_0V-(Ds_0+2T)X},\qquad
 f_6=f_3^2 f_{2,3G}.
 \tag{DZ6.2}
\]
For \((a,b)=(6,1)\) retain \(B_6\); for \((a,b)=(2,3)\)
retain the full three-step product of EX.1. In both cases the
selected whole-zero coefficient is exactly
\[
                         B_{\rm sel}=\frac{q}{T^3}B_{a,b}.
 \tag{DZ6.3}
\]
The generic component valuations are unchanged because \(s_0\)
is a unit. The explicit source translations and line maps extend
on the native disc by those same actual resolution charts.

\par

\medskip\noindent\textbf{DZ6.2. All actual roots and their old/native coordinate arrows}

\par\nobreak

\item For good roots \(m=3l\), \(l\ge1\), the native root is
\[
 Y'^2=X'^3+Rz^{2l}X'^2+Dz^lX'+1.
\]
For reduced roots \(m=3l+2\), \(l\ge0\), it is
\[
 Y'^2=X'^3+Rz^{2l+2}X'^2+Dz^{l+2}X'+z^2.
 \tag{DZ6.4}
\]
The full EX middle-chart formula applies with
\(b_0=Dz^l,c_0=Rz^{2l}\); its exact universal polynomial
identities therefore retain these parameter functions. The marked
sections, division sections and all graph ideals use \(s_0\).
Both central selected root multipliers are four. The actual
starting line map is the same \(F=\tau g_{P,R_{\rm div}}\)
from EX.15, with its entire chord and coordinate factor.

\item On \(D\ne0\), the original source coordinates are related by
\[
 X=D^2x,\quad Y=D^3y,\quad T=D^3t,\quad
 r=R/D^2,\quad s=s_0/D^2,\quad q_{\infty,\rm old}
       =D^3q_{\infty,\rm new}.
 \tag{DZ6.5}
\]
Every original \(f_2\) has weight \(D^{-1}\), hence the
six-factor coefficient has weight \(D^{-6}\). The factor
\((t-t_c)/t^3\) has weight \(D^6\). Consequently the full
suspension scalar is unchanged by DZ6.5.

\item The selected toroidal arrow is also literal. On its nodal
normalization \(S_{\rm new}=D S_{\rm old}\) and
\(w_{\rm old}=w/D\), so
\[
 Z_{\rm old}=\frac{S_{\rm old}-w/D}{S_{\rm old}+w/D}
            =Z_{\rm new},\qquad q_{\rm old}=D^{-3}q_{\rm new}.
 \tag{DZ6.5a}
\]
Away from the node the same equality is obtained by integrating
the actual invariant differential divided by its transported
vanishing period; both scale by the same factor under DZ6.5.
Thus the full Tate period and marked point agree.
Let \(\sigma_a\) be the canonical meromorphic section of
\(\mathcal O(aG-O)\), and use the actual theta frame
\[
 e_\theta=\sigma_a/F_a,\qquad
 F_a(Z)=\Theta(Z/P_G^a)/\Theta(Z).
\]
The Cartier pullback sends \(\sigma_{a,\rm old}\) to
\(\sigma_{a,\rm new}\), while the two complete theta products
have identical arguments and periods. Hence this exact frame
map gives \(V_{\rm old}=V_{\rm new}\).
For example, a local Cartier equation
\(S_{\rm old}-S_{a,\rm old}=D^{-1}(S_{\rm new}-S_{a,\rm new})\)
changes its regular frame by \(D\); the two factors cancel in
the displayed \(e_\theta\). This is the actual frame arrow,
including its coefficient, rather than an omitted unit.
Consequently each toric character has the complete pullback
\[
 Z_{\rm old}^pV_{\rm old}^j q_{\rm old}^h
       =D^{-3h}Z_{\rm new}^pV_{\rm new}^j q_{\rm new}^h.
 \tag{DZ6.5b}
\]
The factor \(D^{-3h}\) and its parameter winding are retained.
On the parameter annulus it is an extending chart unit, so the
same fan and physical circle are carried by both inverse maps.

\item This point frame has the following exact dictionary with the
earlier point-to-tensor Tate frame. Write \(v_M\) for the
coefficient in the canonical meromorphic section \(\sigma_a\),
and set \(f_P=\Theta(Z/P_G)/\Theta(Z)\),
\(f_{P^a}=\Theta(Z/P_G^a)/\Theta(Z)\). Then
\[
 V_{\rm pt}=f_{P^a}(Z)v_M,\qquad
 V_{\rm tensor}=f_P(Z)^a f_a(p)^{-1}v_M,\qquad
 \frac{V_{\rm tensor}}{V_{\rm pt}}
       =\frac{f_P^a}{f_a f_{P^a}}.
 \tag{DZ6.5c}
\]
The last ratio is the full nonzero base constant: the numerator
and denominator have the same divisor and the same Tate
automorphy factor, so their quotient is a holomorphic unit on
the compact elliptic fibre. Its actual value, supplied by the
displayed functions, is retained. Since the Miller arrow has
the exact weight \(f_{a,\rm old}=D^{1-a}f_{a,\rm new}\),
the point coordinate is unchanged whereas
\[
                 V_{{\rm tensor},\rm old}
                      =D^{a-1}V_{{\rm tensor},\rm new}.
 \tag{DZ6.5d}
\]
Thus the complete Cox-character map in this tensor-reference
frame is
\[
 Z_{\rm old}^pV_{{\rm tensor},\rm old}^j q_{\rm old}^h
  =D^{(a-1)j-3h}
       Z_{\rm new}^pV_{{\rm tensor},\rm new}^j q_{\rm new}^h.
 \tag{DZ6.5e}
\]
The physical-line factors are \(D^5\) for \(M_6\) and \(D\)
for \(M_2\). They are closed nonzero parameter units on the
punctured comparison, with their full winding retained.

\item Write \(\eta=\log D/(2\pi i)\) on its parameter cover.
Let \(J_\eta\) denote the \textbf{new-to-old} root coordinate map.
For good roots it is
\[
 z_{\rm old}=D^{3/m}z_{\rm new},\quad
 X'_{\rm old}=X'_{\rm new},\quad Y'_{\rm old}=Y'_{\rm new}.
\]
For reduced roots it is
\[
 z_{\rm old}=D^{3/m}z_{\rm new},\quad
 X'_{\rm old}=D^{2/m}X'_{\rm new},\quad
 Y'_{\rm old}=D^{3/m}Y'_{\rm new},
\]
\[
 \tau'_{\rm old}=D^{-1/m}\tau'_{\rm new},\qquad
 \tau_{0,\rm old}=D^{-2}\tau_{0,\rm new}.
 \tag{DZ6.6}
\]
These maps are the actual base-change and minimal-coordinate maps,
not equalities of fibre types. In the reduced branch the root
\(B'_6\) has weight \(D^{6/m}\), canceled by
\(z_{\rm old}^{-2}\). The chord-over-vertical function has
weight \(D^{1/m}\), canceled by \(\tau'\). Thus
\(F\), its full norm ratio and its reference-normalized root
are literally carried to the corresponding native maps.
The good-root native point-line \(N_2\) coordinates already
agree. The common-graph correction still carries its full
monomial factor: \(u_{\rm old}=D u_{\rm new}\), so an
actual correction generator \(w_0^r u^s\) pulls back with
the coefficient \(D^s\). These factors remain on both
projections of the Cartier graph. They do not alter the
displayed point-line coordinate equality.

\item Let \(\mathcal D\) be the actual finite root deck. For the
reduced branch the identities \(6k\equiv-2\pmod m\),
\(9k\equiv-3\pmod m\), \(k=2l+1\), show that
\[
 J_{\eta+1}=\mathcal D_{\rm old}^{\,3}J_\eta,
 \qquad J_{\eta+1}^{-1}
       =\mathcal D_{\rm new}^{-3}J_\eta^{-1}.
 \tag{DZ6.7}
\]
The same typed equations hold for the good branch. The inverse
arrow therefore has the inverse deck; these are two directions
of the same coordinate square.

\par

\medskip\noindent\textbf{DZ6.3. Full cap endpoint, including scalar and finite character}

\par\nobreak

\item Keep \(R_\lambda^m=\lambda\), the actual signed norm \(n\),
finite character \(\zeta_m^d\), and the full EX norm lift
\(\widetilde F\). On the native root put
\[
 A=R_\lambda^n\zeta_m^d\widetilde F^{\,n}\mathcal D,
 \qquad B=\lambda B_{\rm sel},\qquad A^m=B^n.
\]
All displayed factors commute by the exact EX root/Cartier maps.
The base-preserving correction is consequently the literal map
\[
 C=A^3\mathcal D^{-3}
    =R_\lambda^{3n}\zeta_m^{3d}\widetilde F^{\,3n},
 \qquad
 C^{-1}=R_\lambda^{-3n}\zeta_m^{-3d}\widetilde F^{-3n}.
 \tag{DZ6.8}
\]
Its source translation is \(-3nR_{\rm div}\). On the cap
quotient, DZ6.7 therefore has endpoint \(C^{-1}\) in the
new-to-old direction and \(C\) in the old-to-new direction:
\(\mathcal D^3=A^3C^{-1}\). The scalar
\(R_\lambda^{3n}\) and the finite character are retained.

\item In the good branch \(R_{\rm div}\) lies in the identity formal
group. In the reduced branch it has component two, so
\(3R_{\rm div}\) lies in the identity formal group. The
central source map of \(\widetilde F^3\) is consequently the
identity in both cases, and its central line action is a nonzero
constant. The explicit root field and line map provide a formal
logarithm after this fixed central scalar is retained. This
logarithm extends through the native cap.

\item There is still an actual integral difference between this cap
logarithm and three annular logarithms. For the reduced branch,
with the reference Abel lifts retained,
\[
 3\log R_{\rm div}-\log(3R_{\rm div})
     =2\bigl(3\log G'-\log(3G')\bigr).
 \tag{DZ6.9}
\]
Indeed \(R_{\rm div}=2G'-[m]^{-1}((2m-1)G')\),
\(2m-1=3(2l+1)\), and the remaining terms lie in the
identity formal group, where their logarithms are additive.
For this formula take the literal sector lift
\(\log R_{\rm div}=2\log G'-m^{-1}\log((2m-1)G')\), using
the small formal logarithm in its second term. If the prescribed
annular lift differs by a period \(\Pi\), the exact simultaneous
replacement \(f\mapsto f+\Pi,\ h\mapsto h-ns\Pi\) preserves
the entire map EX.17. This carries that prescribed difference;
it does not change the actual annular marking.
The parenthesized quantity is the actual transported lattice
period fixed by this marking. Its period functions are retained;
it is not replaced by a constant complex number. For the complete
line logarithm its simultaneous physical-circle column is retained
by the same exact cochain replacement.

\item The ordinary root-annulus logarithm obeys
\[
 s_{\rm old}=s_{\rm new}+3\eta.
 \tag{DZ6.10}
\]
Thus the exact sector arrow EX.17 changes by the full
\(3\eta(nf+c_\lambda\delta)\) term, with the corresponding
reference-relative cochain. The following construction keeps
DZ6.9 and its line column in that square.

\par

\medskip\noindent\textbf{DZ6.4. The native full disc and the exact cap correction}

\par\nobreak

\item The native source addition has explicit regular charts. In the
exceptional chart \(X=T\alpha,Y=T\beta\),
\[
 \beta^2=T\alpha^3+R\alpha^2+D\alpha+1,\qquad
 \kappa=\frac{\beta-1}{\alpha}
        =\frac{D+R\alpha+T\alpha^2}{\beta+1}.
\]
The second expression is regular at \((\alpha,\beta)=(0,1)\);
the opposite chart uses the formal coordinate at \(O\).
The actual addition is
\[
 X'=\kappa^2-R-T\alpha,\qquad Y'=-T-\kappa X'.
\]
In the other source chart \(T=Xu,Y=Xv_0\), use
\[
 X=v_0^2-R-Du-u^2,\quad
 X'=u(D-2v_0+2u),\quad
 u'=\frac{X}{D-2v_0+2u},\quad v_0'=-u'-v_0+u.
 \tag{DZ6.11}
\]
At the two intersection points the denominator is a unit because
\((D-2v_0)(D+2v_0)=4s_0\ne0\). The inverse is the actual
opposite chord. These maps retain \(T\), \(w\) and the two
source resolution charts, throughout the native disc.

\item On the selected nodal fibre use
\[
 X=S^2-1,\qquad Y=S(S^2-w^2),\qquad G:\ S=-1.
\]
The full tangent-over-vertical factor becomes
\[
 f_2=\frac{(S+1)^2}{S+(1+w^2)/2}.
\]
It has value two on each branch \(S=\pm w\). Translation fixes
each of those node branches. Thus all six factors of \(B_{2,3}\)
give 64 there. For \(B_6=f_2^6 f_6(p-G)/f_6(p)\), the retained
Miller values on the node are
\(f_6=6+20w^2+6w^4\), which is 120 at \(w^2=3\);
the ratio is one. Hence, for both actual lines,
\[
                       (B_{\rm sel}/q)_{\rm node}
                            =64/(w^2-1)^3.
 \tag{DZ6.12}
\]
The actual Tate point \(P_G=(1+w)/(1-w)\), Tate-height unit
\(\kappa(0,w)=-(w^2-1)^2/(64w^6)\), and the two marked
Cartier factors
\[
 c_2=(1+w)^2/2,\qquad
 c_6=(1+w)^6/(6+20w^2+6w^4)
\]
are all nonzero analytic units on the chosen disc. Their full
Tate/Cartier maps therefore extend on every prescribed fan.
These statements use the actual parameters \(w\) near each
chosen square root of three, not an identification by fibre type.

\item The EX root, chord, invariant norm and full-reference annular
construction applies to DZ6.4 with \(b_0=Dz^l,c_0=Rz^{2l}\).
The certificate checks its universal identities before this
substitution. The selected source circle is disjoint from the
root and the I2 resolution. The fixed small scalar yields the
same compact contraction shells as EX.6. Thus the native model
is a proper smooth sphere family on the full disc, with its
literal identity closing.

\item Here is the cap map that implements the old/native comparison.
Let \(\mathcal G\) be the actual identity component of the
degree-zero line-translation group. Its exponential near the
identity is the exponential of the source dualizing vector and
physical \(\delta\) field. The regular EX vector and line map
extend this exponential through the resolved root. The map
\(\widetilde F^3\) has identity central source action and a
nonzero constant central line action. Retain the chosen central
line logarithm, divide by that central scalar to use the inverse
exponential near the identity, and then add it back. This gives
an extending logarithm \(\ell_3=\operatorname{Log}(\widetilde F^3)\).
It is fixed by the whole root deck, since its actual map and
central logarithm are fixed and the near-identity inverse is
unique. All these maps vary holomorphically on the native disc.

\item Set
\[
 \mathcal L_C=n\ell_3+3c_\lambda\delta,\qquad
 K_\eta=\operatorname{Exp}(\eta\mathcal L_C),\qquad
 \widehat J_\eta=J_\eta K_\eta.
 \tag{DZ6.13}
\]
The exponential is evaluated for the bounded complex times in
each fixed angular parameter chart; the source action is the
actual extending group flow. Since \(\operatorname{Exp}\mathcal
L_C=C\), the complete endpoint is
\[
 \widehat J_{\eta+1}
  =\mathcal D_{\rm old}^3 J_\eta C_{\rm new}K_\eta
  =A_{\rm old}^3\widehat J_\eta.
 \tag{DZ6.14}
\]
Its inverse has endpoint \(A_{\rm new}^{-3}\).
Consequently these are single-valued maps on the cap quotients.
They preserve the root function, the actual line, and physical
\(\delta\), including the full scalar in DZ6.8.

\par

\medskip\noindent\textbf{DZ6.5. Literal annular square and relative groupoid solve}

\par\nobreak

\item Choose one nonzero reference \(D_*\). Transport the entire old
ordinary-to-cap map at \(D_*\) by DZ6.5-\/-6, including its root
cut, character, source periods and every integral \(\delta\)
column. Construct the native reference map over the full native
disc from this transported map by the fixed-sector construction
EX.17-\/-19. The root norm branch uses the transported old value.
This specifies the whole native annular marking.

\item On a parameter angular chart and a base-annulus sector let
\(f=\operatorname{Log}\widetilde F\), and put
\[
                 \Pi_F=3f-\ell_3.
 \tag{DZ6.15}
\]
This is the actual period of the complete source/line Lie group,
because both exponentials are \(\widetilde F^3\).
Its source part in the reduced branch is the negative of DZ6.9;
its line part is the same retained period of the actual chord
lift. Pull the two annular cochains into the common native frame
and denote them by \(h_{\rm old}\), \(h_{\rm new}\).
Here this frame change includes DZ6.5c-\/-e: a tensor-reference
map is first conjugated by its complete ratio
\(f_P^a/(f_a f_{P^a})\). The new-to-old tensor arrow has the
literal \(+(a-1)\eta\delta\) contribution; its old-to-new
inverse \(I_D\) has the negative contribution. On adjacent
parameter charts these give the respective integral columns
\(+(a-1)j\delta\) and \(-(a-1)j\delta\). The transported
cochains use these directed frame arrows in full. The next
display is in the common point frame, after precisely these
conjugations.
Direct substitution of \(s_{\rm old}=s_{\rm new}+3\eta\)
and DZ6.13 in EX.17 gives
\[
 \widehat J_\eta^{-1}\Phi_{\rm old}
  =\operatorname{Exp}\!
     \left(\eta n\Pi_F+h_{\rm old}-h_{\rm new}\right)
      \Phi_{\rm new}I_D .
 \tag{DZ6.16}
\]
Here \(I_D\) is the actual old-to-new ordinary source/Cartier
map DZ6.5. This is a square of quotient maps with the same
ordinary input. The term \(3c_\lambda\delta\), including
both the scalar and character, cancels exactly between the two
displayed arrows. On a parameter-chart change
\(\eta_\beta=\eta_\alpha+j\), the indicated lifted translation
changes by \(jn\Pi_F\), together with the transported cochain
columns. This is its actual integral groupoid discrepancy.
It is retained rather than regarded as an extending root
logarithm.

\item For completeness, the relative solution uses the following
explicit sheaf and operators. On the full root atlas the
extending line-translation Lie bundle has
\[
 0\longrightarrow\mathcal O\delta\longrightarrow\mathcal V
       \longrightarrow\mathcal O(-1)\longrightarrow0.
 \tag{DZ6.17}
\]
Indeed the source dualizing vector pulls to \(z^k\) times the
root vector, with its exact character and \(0\le k<m\).
Taking its invariant holomorphic coefficients gives precisely
the original source Lie line. The invariant Cartier graph
arrows already proved identify its line extension; the physical
line is \(\mathcal O\delta\). Here are the actual splitting
and solution operators. Choose the source frames
\(e_\infty=T^{-1}e_0\) and write \(q_\infty=T^{-1}\),
retaining \(q=T-v\) for the selected node. On a fixed
overlap annulus expand the adjusted coefficient as
\[
 f(T,D)=\sum_{j\in\mathbb Z}f_j(D)T^j,\qquad
 f_j(D)=\frac1{2\pi i}\int_{|T|=\rho}
                     f(T,D)T^{-j-1}\,dT .
 \tag{DZ6.17c}
\]
For the source \(\mathcal O(-1)\) component take
\[
 u_0=-\sum_{j\ge0}f_jT^j,\qquad
 u_\infty=\sum_{j\le-1}f_j q_\infty^{-j-1}.
\]
Then \(T^{-1}u_\infty-u_0=f\) literally. For the
\(\mathcal O\delta\) component use the same \(u_0\) and
\(u_\infty=\sum_{j\le-1}f_jq_\infty^{-j}\); here
\(u_\infty-u_0=f\).
These series converge on the two chosen source-base charts
by their Laurent coefficient integrals.

\item The upper triangular transition of the actual Lie frames
is \(\widehat e_\infty=T^{-1}\widehat e_0+
a(T,D)\delta\). Its coefficient is the derivative of the
displayed Cartier arrow: in line-log coordinates it is the
source invariant vector applied to the logarithm of that
arrow\textquotesingle s full line multiplier, divided by \(2\pi i\).
Thus the tensor factor, root monomials and all Cartier units
are the functions being differentiated. Changing these
lifts to \(\widehat e_0+v_0\delta\) and
\(\widehat e_\infty+v_\infty\delta\) gives the exact equation
\[
              T v_\infty-v_0=-T a(T,D).
 \tag{DZ6.17d}
\]
For \(g=-Ta=\sum g_jT^j\), take
\(v_0=-\sum_{j\ge0}g_jT^j\) and
\(v_\infty=\sum_{j\le-1}g_j q_\infty^{1-j}\).
This is the explicit \(\mathcal O(1)\) Laurent solve, and it
splits DZ6.17 before the two component solves above.
All contours and Laurent projections are fixed, so these
operators preserve holomorphic parameter dependence and
a zero reference value.

\item Finally, after retaining the integral finite-group cochains,
let \(a_g\) be the remaining vector one-cocycle, with
\(a_{gh}=a_g+g a_h\). Its literal finite-group solution is
\[
                  v=-\frac1m\sum_h a_h,\qquad
                         g v-v=a_g .
 \tag{DZ6.17e}
\]
The equality follows by reindexing the entire sum using the
displayed cocycle identity. Averaging is applied only to these
complex vector cochains; every integral period and character
has been kept in the preceding adjustment.

\item Use the exponential sequence on this atlas and its entire finite
groupoid. At reduced central components take the actual
suspension power that moves the component to the identity before
the local logarithm. Every finite discrepancy retains
\(A^m=B^n\), its character and the line period in DZ6.15.
After the entire reference map has been matched at \(D_*\),
the difference of the integral base-direction two-cocycles is
zero there. It stays zero on each connected parameter chart:
all its values are continuous lattice-valued functions under
the actual source, root and Cartier identifications. On changing
parameter charts, DZ6.14 and DZ6.16 give exactly the integral
one-cochain \(jn\Pi_F\); its coboundary is the corresponding
change of representatives of those same two-cocycles.
Thus the zero class is transported around the whole annulus.
This records the complete mixed parameter/base discrepancy.

\item Subtract this integral cochain adjustment, apply the fixed
Laurent/Cauchy solution of DZ6.17 to the remaining vector
one-cocycle, and exponentiate. On each parameter chart this
gives an actual old/native family isomorphism. All its local
maps are the source addition/line maps above, the finite-root
maps, or the selected toric translations. At the cusp an
integral component reindexing is made first; the remaining
two characters are bounded holomorphic units. They act on
every prescribed Cox chart by its literal character map.
The inverses use the inverse group elements and the same
proper chart domains.

\par

\medskip\noindent\textbf{DZ6.6. Global physical-scalar factorization and closing}

\par\nobreak

\item The ratio of the local parameter-chart isomorphisms is a
global extending translation of the common marked family.
Here is its exact global group. Lift the marked zero section
through the nonzero-line covering away from the selected cusp.
Its source projection extends over every root; the regular-locus
criterion puts it on a smooth central component. The actual
suspension power moves it to the identity component, and the
degree-zero Hom-line evaluation gives its full line lift.
On generic fibres, constant factors of automorphy show that
this lift commutes with the actual suspension. Repeating the
finite-cap conjugacy \(m\) times and using strict contraction
forces the residual deck exponent to vanish. Thus all finite
characters and component choices have been retained.

\item The resulting source section is \(jG\) with \(j\in\mathbb Z\)
locally constant: these native surfaces have the same marked
E6 plus A1 lattice and primitive generator of height \(1/6\).
More explicitly, the actual simultaneous source resolution above
transports the unimodular rational-source intersection lattice
from a nearby \(D\ne0\) fibre. The E6 and A1 component lattice
has determinant six, and \(G\) meets the same nonidentity
components with corrections \(4/3\) and \(1/2\), while
\(G\cdot O=0\). Its height is \(2-4/3-1/2=1/6\).
Write \(\iota\) for its index in the free Mordell-\/-Weil lattice
and \(t_{\rm MW}\) for the torsion order. The integral
determinant formula gives
\[
 1=\frac{6\,h(G_{\rm prim})}{t_{\rm MW}^2},\qquad
 h(G)=\iota^2 h(G_{\rm prim})=\frac16,\qquad
                 \iota^2 t_{\rm MW}^2=1.
\]
Thus both positive integers are one. This identifies the
actual marked generator on the entire native disc.
The local covering lifts patch to this section: their differences
are integral suspension powers on the source base with its
selected point deleted. Its orbifold fundamental group is a
free product of the finite cap groups, and a homomorphism from
it to \(\mathbb Z\) is zero. The integer one-cocycle is
therefore removed by its actual deck powers, including at the
finite root.

\item The full valuation computation can be made directly in the
present Miller frames. Let \(U_G=B_{a,1}^{-1}\), the actual
meromorphic line lift covering \(+G\). Because \(B_{a,b}\)
is its literal inverse \(b\)-fold composite, DZ6.3 gives
\[
                 \varrho=B_{\rm sel}^{-1}/U_G^b=T^3/q.
 \tag{DZ6.17a}
\]
The ordinary I2 order of \(U_G\) is \(-3/b\).
Its finite-root order is \(2m/b\), by the full root coordinate
maps and the nonzero central coefficient already proved.
Its selected-cusp order is zero; its source Tate coordinate
there is a unit. There are no other fibre orders.
An endpoint lift has the form \(h(T)U_G^j\). Its invertibility
gives
\[
 \operatorname{ord}_0h=3j/b,\qquad
 \operatorname{ord}_\infty h=-2j/b.
\]
At the selected cusp the literal integer reindexing by \(B\)
gives \(\operatorname{ord}_{q=0}h=-k\): composing the
endpoint with \(B^k\) removes that line order. Source period
reindexing changes it by a unit only. The degree of this exact
rational divisor is zero, so \(j=bk\). All three orders of
\(h\) now equal those of \(\varrho^k\); their quotient is a
nonzero constant \(c\). The complete map is consequently
\[
              hU_G^{bk}=cB_{\rm sel}^{-k}.
 \tag{DZ6.17b}
\]
Since \(B=\lambda B_{\rm sel}\) is the actual quotient deck,
this acts by physical scalar \(c\lambda^k\).
Evaluation at an ordinary marked origin makes that scalar
holomorphic in the parameter. This computation includes the
finite integrality and character conditions before the global
equality, rather than erasing a finite central kernel.

\item Let \(F_+=cF_-\) be the actual two parameter-chart comparison
maps, with \(c\) this nonzero holomorphic scalar. On the two
simply connected overlap components choose \(\ell=\log c\),
and a fixed angular partition \(\chi_++\chi_-=1\). Put
\[
 b_+=\chi_-\ell,\qquad b_-=-\chi_+\ell,\qquad
 \omega=\bar\partial b_\pm,\qquad
 v(D)=\frac1\pi\int\frac{\omega(E)}{D-E}\,dA(E).
\]
A radial cutoff is one on the smaller annulus carrying the loop
and supported in the larger one. Then
\(g_\pm=\exp(b_\pm-v)\) are nonzero holomorphic functions,
\(g_+/g_-=c\), and
\[
                         g_+^{-1}F_+=g_-^{-1}F_- .
 \tag{DZ6.18}
\]
This is a single-valued actual family isomorphism over that
annulus. It includes all windings of \(c\) and the literal
physical action on the line and normal charts.

\item Both families are treated with their specified forward endpoint
maps \(E_{\rm old},E_{\rm new}\); for the native disc the latter
is its coordinate identity. Restrict the closed family map to
the cut parameter interval, calling it \(F_t\). Equation
DZ6.18 gives the actual square
\[
                 F_1E_{\rm old}=E_{\rm new}F_0 .
 \tag{DZ6.19}
\]
The physical return uses the inverse endpoint on both sides.
Transporting a native-disc connection through \(F_t\) therefore
conjugates the complete oriented old return to that of the
native disc. The latter extends over the disc and is zero.
This proves the zero return for the specified old family with
the whole transported reference marking at either \(w^2=3\).
The exact original-moduli coordinate at infinity is
\[
                 \epsilon=\frac1r
                    =\frac{D^2}{3-2w^2}.
 \tag{DZ6.20}
\]
Since \(D\) is a simple local coordinate at either point and
the denominator has value \(-3\), a positive \(w\)-loop has
degree two over the positive original infinity meridian.
The constructed family map concerns this actual marked
\(w\)-loop, with the source and cap arrows already specified.
No shorter original-parameter endpoint or return is assigned.
The full source, line, cap and endpoint maps above prove the stated conclusion with every prescribed annular integer retained.

\par

\medskip\noindent\textbf{Executed exact certificate}

\par\nobreak

\item The exact program \path{code/part3/e6_d_zero_disc.py} and output \path{code/part3/e6_d_zero_disc.json} execute 521 identities and finite signed instances. They include the source and nodal Cartier charts, both root and tensor weights, complete cap endpoint, retained period, all three Laurent operators, finite vector averaging, global valuation equation and exact degree-two parameter. The complete geometric family proof is above. The preceding coordinate-only and frame-clarification stages remain under \path{records/attempts/e6_d_zero_disc/}.

\end{enumerate}
\end{proof}




\begin{theorem}\label{night:P3K.98}
\StatementRecord{night:P3K.98}{proved}{S6\_06\_BLIND: GPT-6 Astra (OpenAI, ChatGPT Pro mode, effort Pro), recorded re-derivation in the checkpoint dated 18 November 2026; also checked in this writing by \CurrentChecker}{\cref{night:P3F.41}, \cref{night:P3F.56}, \cref{night:P3F.7}, \cref{night:P3F.8}, \cref{night:P3F.9}, \cref{night:P3K.63}, \cref{night:P3K.68}, \cref{night:P3K.71}, \cref{night:P3K.95}, \cref{night:P3K.96}}
\par\medskip\noindent\textbf{Statement.}\par\nopagebreak[4]

On the source
\[
 y^2=x^3+r x^2+t x+t^2,\qquad G=(0,t),\qquad
 r=\frac{3-2w^2}{(w^2-3)^2},
\]
take the selected zero at
\(t_c=(w^2-1)/(w^2-3)^3\), the actual point line
\(M_a=\mathcal O(aG-O)\), and the whole source/line suspension
over translation by \(-bG\). Retain the exact finite cap and
complete ordinary-to-cap reference map in either row
\[
 (a,b;m)=(6,1;3l+2),\quad l\ge0,\qquad
 (a,b;m)=(2,3;3l),\quad l\ge1 .
 \tag{EG.1}
\]
In both rows \(ab=6\), the primitive norm is \(n=\pm1\), and the scalar
\(\lambda=R_\lambda^m\ne0\), including its chosen root and
phase, is fixed and sufficiently small for the actual compact
contraction shells. Every finite character and integral annular
reference column is retained. The positive ordered \(w\)-loop
defines the completed sphere family and its actual return.

On the initial vertical fan the physical characteristic pair is
\((k_M,\ell_M)=(-6,0)\), its native free pair is \((0,-6)\),
and its clutched sphere has class \(-1\) in \(\Theta_7\cong
\mathbb Z/28\). On every vertical fan \(T_j\), \(j\in\mathbb Z\),
the exact values are
\[
 \begin{split}
 (k_M,\ell_M)&=(-6-18j(j-1),0),\\
 (m_{\rm native},n_{\rm native})&=(0,-6-18j(j-1)),\\
 [\Sigma_j]&=-1-3j(j-1)\pmod {28}.
 \end{split}
\]
For every selected lattice-basis fan, the full actual return is
given by the proved K71 fan algorithm, with these vertical seeds
and the complete coefficient and Cox maps (EG.45)--(EG.46).
This is a comparison of the actual completed families, including
their physical circle, both normal caps, source/Cartier maps and
finite-root closing. The ordered loop has local degree four over
the original \(r=1/3\) puncture; the statement assigns no shorter
source endpoint.

\end{theorem}
\begin{proof}
\begin{enumerate}[label=\arabic*.]
\item \par\medskip\noindent\textbf{EG.1. The full source and the two actual point-line maps.}\par\nopagebreak[4]

\item Retain
\[
 \begin{gathered}
 v=w^2-1,\quad D=w^2-3,\quad A=3-2w^2,\quad
 \\
 Y^2=X^3+AX^2+DTX+T^2,\qquad G=(0,T),\quad q=T-v .
 \end{gathered}
 \tag{EG.2}
\]
The source map to the original family is
\((x,y,t)=(D^{-2}X,D^{-3}Y,D^{-3}T)\), with
\(r=A/D^2\). Set \(s_0=v(w^2+3)/4\) and
\[
 U=Y-T-DX/2,\qquad
 f_2=\frac{U}{X-s_0},\qquad {\cal T}(p)=p-G .
 \tag{EG.3}
\]
The two inverse chord maps are
\[
 k_\pm=(Y\mp T)/X,\quad
 X_\pm=k_\pm^2-A-X,\quad Y_\pm=\mp T-k_\pm X_\pm .
 \tag{EG.4}
\]
Their denominators are units at the selected cusp
\((w,X,Y,T)=(0,-1,0,-1)\). Thus all four translations
\({\cal T}^j\), \(0\le j\le3\), used below extend on a common
source neighbourhood and keep the full \(q,w\).

\item Write \(M_a=\mathcal O(aG-O)\). For \(a=2\), the actual
point-to-tensor map is \(r_2=f_2^{-1}:M_2\to M_1^{\otimes2}\).
For \(a=6\), it is \(r_6=f_6^{-1}:M_6\to M_1^{\otimes6}\), where
\[
 f_3=\frac{s_0XU}{s_0(Y+T+DX/2)-(Ds_0+2T)X},
 \qquad f_6=f_3^2 f_{2,3G}.
 \tag{EG.5}
\]
Here \(f_{2,3G}\) is the whole tangent-over-vertical function.
Its section and tangent coordinates are explicitly
\[
 \begin{split}
 x_3&=4T^2/s_0^2+2DT/s_0,\\
 y_3&=T(8T^2+6Ds_0T+D^2s_0^2-s_0^3)/s_0^3,\\
 \ell_3&=(3x_3^2+2Ax_3+DT)/(2y_3),\qquad
 x_6=\ell_3^2-A-2x_3,\\
 f_{2,3G}&=\frac{Y-y_3-\ell_3(X-x_3)}{X-x_6}.
 \end{split}
 \tag{EG.6}
\]
Every source and target factor is retained in the two suspension
coefficients
\[
 \begin{split}
 B_{6,1}(p)&=f_2(p)^6\,\frac{f_6({\cal T}p)}{f_6(p)},\\
 B_{2,3}(p)&=f_2(p)f_2({\cal T}^3p)
                     f_2({\cal T}p)^2f_2({\cal T}^2p)^2,\\
 B_{\rm sel}^{(a,b)}(p)&=\frac q{T^3}B_{a,b}(p).
 \end{split}
 \tag{EG.7}
\]
Indeed the \(M_1^{\otimes a}\) lift over \({\cal T}^b\)
has coefficient \(\prod_{j=0}^{b-1}f_2({\cal T}^jp)^a\).
Multiplying it by \(r_a(p)/r_a({\cal T}^bp)\) gives precisely
(EG.7), including the uncancelled \(f_2({\cal T}^3p)\) in the
second row. Thus the source translations are exactly \(-G\)
and \(-3G\), respectively.

\item On the selected fibre put
\[
 X=S^2-1,\quad Y=S(S^2-w^2),\quad
 Z=\frac{S-w}{S+w},\quad P=\frac{1+w}{1-w}.
 \tag{EG.8}
\]
The section \(G\) is \(S=-1\), and \(O\) is \(S=\infty\).
For \(a=2,6\) the actual point \(aG\), divisor equation and
normalizing unit are
\[
 S_a=w\frac{1+P^a}{1-P^a},\qquad c_a=w-S_a,\qquad
 S-S_a=c_a\,\frac{1-Z/P^a}{1-Z}.
 \tag{EG.9}
\]
They extend holomorphically over \(w=0\), with
\(S_a(0)=-1/a\), \(c_a(0)=1/a\). In particular
\[
 \begin{aligned}
 S_2&=-(1+w^2)/2,&c_2&=(1+w)^2/2,\\
 c_6&=(1+w)^6/(6+20w^2+6w^4),&
 f_a|_{q=0}&=(S+1)^a/(S-S_a).
 \end{aligned}
 \tag{EG.10}
\]
The identities use the actual Miller functions (EG.3)--(EG.6):
their divisors and their coefficient \(S^{a-1}\) at infinity
agree. For \(a=2\) it also follows by direct substitution in
(EG.3). This retains the whole point-divisor-to-theta coefficient,
not just the point \(P^a\).

\item Let \(f_P(Z)=\Theta_Q(Z/P)/\Theta_Q(Z)\) be the actual rigidified
theta section of \(M_1\), with \(f_P(QZ)=P f_P(Z)\).
Here and below \(P=P_G(q,w)\) is the full marked Tate lift;
(EG.8)--(EG.10) are its restriction to \(q=0\).
The two actual point-line theta coefficients are
\[
 V_a=f_P(Z)^a r_a(p)\,v_a .
 \tag{EG.11}
\]
On the selected normalization their comparison with the strict
point theta factor is the full even unit
\[
 f_P^a/f_a=\frac{c_a}{(1+w)^a}
                   \frac{1-Z/P^a}{1-Z},\qquad
 \frac{c_2}{(1+w)^2}=\frac12,\quad
 \frac{c_6}{(1+w)^6}=\frac1{6+20w^2+6w^4}.
 \tag{EG.12}
\]

\item The exact constant of the \(M_1\) translation is
\[
 \eta_1=f_2(p)\frac{f_P(P^{-1}Z)}{f_P(Z)} .
 \tag{EG.13}
\]
Its divisor is zero on the compact smooth elliptic fibre, so it
is independent of \(p\). Substitution of (EG.7) into (EG.11)
then telescopes literally, giving in both rows
\[
 H:(Z,V_a)\mapsto(\kappa q Z,P^aV_a),\qquad
 B:(Z,V_a)\mapsto(P^{-b}Z,q\eta V_a),\qquad
 \eta=\lambda\,\eta_1^6/T^3 .
 \tag{EG.14}
\]
This is an equality of the full coefficients: \(ab=6\) is used
only after the actual tensor/point maps in (EG.7) and (EG.11).
Both selected node branches have \(f_2=2\). They are fixed by
every \({\cal T}^j\), and therefore
\[
 \eta_1|_{q=0}=2,\qquad
 B_{a,b}|_{\rm node}=64,\qquad
 B_{\rm sel}^{(a,b)}/q=64/v^3 .
 \tag{EG.15}
\]
The coefficient at the central cusp is the retained \(-64\).

\item The weighted node coordinates
\(\xi=X-v=w^2\widetilde X,\ Y=w^3\widetilde Y,\
q=w^6Q_c\) give the full equation
\[
 \widetilde Y^2=\widetilde X^2(\widetilde X+1)
       +Q_c(v^2+Dw^2\widetilde X)+w^6Q_c^2 .
 \tag{EG.16}
\]
The inverse \(j\)-series gives
\(\kappa=w^{-6}\widehat\kappa(Q_c,w)\),
\(\widehat\kappa(0,w)=-v^2/64\).
The actual \(G\) has formal parameter
\(wv/(v+w^6Q_c)\). Thus \(P\), \(P^{-b}\), \(P^a\), and
\(\eta/\lambda\) have extending nonzero logarithms on this
weighted bidisc, with their specified constant values. The only
parameter winding in (EG.14) is the full \(-6\) in \(\kappa\).

\par\medskip\noindent\textbf{EG.2. The complete exterior and both finite-root families.}\par\nopagebreak[4]

\item On the ordinary exterior the actual source maps (EG.4) and
their compositions extend over the unselected \(I_2\), using
the two A1 charts of
\[
 (y-t-x/2)(y+t+x/2)=x^2(x-s),\qquad s=1/4-r .
 \tag{EG.17}
\]
Here \(s\) stays near \(-1/12\) and is a unit. The line
divisor is \(3F_0-2F_\infty\) in both rows: in the second
row this follows from the six retained \(f_2\) factors of
(EG.7), with component orders \(0,1\) and the actual odd
translation exchanging the two \(I_2\) components. Thus
\(q/T^3\) leaves exactly the selected whole zero. The maps
and their inverses have zero order along the source-parameter
divisor, as evaluated on an ordinary section avoiding the
horizontal divisors. Normality extends the typed Hom map
through the remaining codimension-two source locus.

\item We now use the complete K96 root maps, with the following
literal substitutions and units. For the reduced branch put
\(m=3l+2,\ k=2l+1\). In the original source coordinate the root is
\[
 q_\infty=z^m,\quad x_0=z^{2k}X',\quad y_0=z^{3k}Y',
 \qquad
 Y'^2=X'^3+r z^{2l+2}X'^2+z^{l+2}X'+z^2 .
 \tag{EG.18}
\]
With \(c=\sqrt{1+z^lX'+r z^{2l}X'^2}\),
\(U'=Y'-zc,\ V'=Y'+zc\), its equation is \(U'V'=X'^3\).
The inverse of \(2zc=V'-U'\) and all three A2 charts retain
the actual \(z\). In the middle chart \(X'=ab,\ U'=a^2b,\
V'=ab^2\), use \(b_0=z^l,c_0=r z^{2l}\) and the full functions
\[
 \begin{split}
 K&=1+\frac{b(b-a)(b_0+c_0ab)}{2c(c+1)},\\
 F&=1-\frac{ab(b_0+c_0ab)(K+1)}{2c(c+1)}
             -\frac{c_0ab^2(a-b)}{4c^2},\\
 X''&=a(a-b)F,\quad c'=\sqrt{1+b_0X''+c_0X''^2},\\
 H_c&=K+\frac{b(a-b)(b_0+c_0X'')}{2c(1+c')}.
 \end{split}
 \tag{EG.19}
\]
The actual \(+G'\) arrow is
\((a,b)\mapsto((b-a)F/H_c,-aH_c)\); its inverse is the
actual negative chord. It preserves the same \(z\) and
commutes with
\[
 {\mathfrak D}(z,X',Y')
   =(\zeta_mz,\zeta_m^{-2k}X',\zeta_m^{-3k}Y').
 \tag{EG.20}
\]
The exact division is
\(R=2G'-[m]^{-1}((2m-1)G')\), with \(mR=G'\).
The small inverse is taken only on the identity formal
group; \(2m-1=3k\). The nonidentity central component of
\(R\) is two and is retained.

\item The root line graph is the actual ideal \((\tau',z^k)\),
where \(\tau'=-X'/Y'=z^k\tau_0\). The two section equations
are \(\tau'\) and \(\tau'-z^kL(s,z^m)\), \(L(s,0)=-s\ne0\).
Its normalized graph has the ray \((k,1)\), and its regular
subdivision retains every ray \((j,1)\), \(1\le j\le k\).
Both exceptional orders are \(j\); at \(j=k\) the graph
coordinate retains the distinct directions \(0,L\). The
literal coefficient maps on its two charts are
\[
 \frac{\tau'}{\tau'-z^kL}=\frac{u}{u-L},\qquad
 \frac1{1-v_{\rm gr}L},\qquad uv_{\rm gr}=1 .
 \tag{EG.21}
\]
The other quotient centres carry neither marked divisor.
The deck-commuting source translation and its inverse lift
to the quotient resolution and then to the normalized
product graph, with these same coefficient maps. This
gives the full equality \(\beta^*N_6=\pi^*M_6\).
The root selected coefficient is exactly
\[
 B'_0=h_\infty z^{-2}B'_6,\qquad
 h_\infty=1-t_c(r)z^m,\qquad (B'_6/z^2)|_{\rm central}=4 .
 \tag{EG.22}
\]

\item For the good branch \(m=3l\), the actual root is
\[
 q_\infty=z^{3l},\quad x_0=z^{4l}X',\quad y_0=z^{6l}Y',
 \qquad
 Y'^2=X'^3+r z^{2l}X'^2+z^lX'+1 .
 \tag{EG.23}
\]
Its source point is \(G'=(0,1)\), and the small formal
division is \(R=[m]^{-1}(3G')\), with \(mR=3G'\).
Put \(u=z^l\). The quotient first retains its \(\mu_l\)
kernel and then the \(\mu_3\) action of weights \((2,1)\)
on each marked chart \((w_P,u)\). At \(2G'\) take
\(w_P=X'-su^2\); at \(O\), take \(w_O=-X'/Y'\).
The resolution rays are \((2/3,1/3),(1/3,2/3)\).
The strict-section support function
\(\max(v_w-2v_u,0)\) gives the full correction
\[
 \beta^*N_2=\pi^*M_2+E_P-E_O,\qquad
 \operatorname{coeff}_{(v_w,v_z)}E_*=\min(v_w,2l v_z).
 \tag{EG.24}
\]
Every arrow is the quotient of its actual monomial
Cartier generators. The translations have the exact
form \(w\mapsto w\varepsilon+u^2a\) and preserve the
ideals \((w,u^j),(u,w^j)\), \(j=1,2\), together with
their inverses. They therefore lift to this full graph
and carry both correction divisors. The complete root
Hom coefficient is
\[
 B'_0=h_\infty B'_{2,3},\qquad
 B'_{2,3}|_{\rm central}=4 .
 \tag{EG.25}
\]
The last constant is the square of the exact norm \(-2\)
of \(f=(Y'-1)/X'\) under the three central translations,
not a discarded base factor.

\item For both roots use the literal chord-over-vertical function
\(g_{P_*,R}=\ell_{P_*,R}/(x-x(P_*+R))\).
Its divisor is \(D_{P_*}-t_{-R}^*D_{P_*}\).
The actual starting line lift is
\[
 F_{\rm line}=
 \begin{cases}
 \tau(6G')\,g_{6G',R},& a=6,\\
 \tau(R)\,g_{2G',R},& a=2 .
 \end{cases}
 \tag{EG.26}
\]
The formal factor and chord have opposite deck weights.
Their product commutes with \(\mathfrak D\) and has
central coefficient one. In the reduced case the zero
Hom multidegrees and the primitive whole-fibre kernel
carry this limit to every component. The exact ratio
\(\beta=B'_0/F_{\rm line}^m\) is an invariant base unit,
central four. Its analytic \(m\)-th root is fixed by the
complete reference map. Put
\[
 \widetilde F=\beta^{1/m}F_{\rm line},\quad
 A_0=\zeta_m^d\widetilde F^{\,n}\mathfrak D,\quad
 B=\lambda B'_0,\quad A=R_\lambda^nA_0 .
 \tag{EG.27}
\]
Then \(\widetilde F^m=B'_0\), \(A_0^m=(B'_0)^n\),
\(A^m=B^n\), and all source, correction, character and
line arrows commute literally as specified.

\par\medskip\noindent\textbf{EG.3. The root rescaling and the whole annular reference.}\par\nopagebreak[4]

\item All preceding formulas can be written directly in the
current \(T\) coordinate. Choose the analytic unit
\(L_D=D^{1/m}\) on the source disc. The exact new-to-old
root maps are
\[
 \begin{array}{c|ccc}
 &z_{\rm old}&X'_{\rm old}&Y'_{\rm old}\\ \hline
 a=6&L_D^3 z_{\rm new}&L_D^2X'_{\rm new}&L_D^3Y'_{\rm new}\\
 a=2&L_D^3 z_{\rm new}&X'_{\rm new}&Y'_{\rm new}.
 \end{array}
 \tag{EG.28}
\]
Their inverse divides by the same nonzero units. The
new reduced equation has coefficients
\(A z^{2l+2}X'^2+D z^{l+2}X'+z^2\);
the new good equation has
\(A z^{2l}X'^2+D z^lX'+1\).
Substitution uses \(L_D^m=D\) and gives these full
equations, not just their central fibres. In both rows
\(q_{\infty,\rm old}=D^3 q_{\infty,\rm new}\) and
\[
 h_{\infty,\rm old}=1-vz_{\rm new}^m .
 \tag{EG.29}
\]

\item On the reduced graph,
\(\tau'_{\rm old}=L_D^{-1}\tau'_{\rm new}\) and
\(\tau_{0,\rm old}=D^{-2}\tau_{0,\rm new}\);
its \(u=\tau'/z^k\) coordinate changes by \(D^{-2}\).
Both numerator and denominator of (EG.21) are carried.
The root tangent factor has weight \(L_D\), so
\(B'_{6,\rm old}=L_D^6 B'_{6,\rm new}\).
The \(z^{-2}\) factor contributes \(L_D^{-6}\);
the complete root \(B'_0\) is unchanged.
In (EG.26) the formal coordinate has weight
\(L_D^{-1}\) and the full chord has weight \(L_D\);
their product is unchanged as well.

\item In the good graph \(u_{\rm old}=D u_{\rm new}\),
while \(w_P,w_O\) are unchanged. On every correction
cone the actual monomial \(w^r u^s\) consequently carries
the full factor \(D^s\); this is the line arrow in
(EG.24). The root source coordinates and chord functions
are unchanged by (EG.28), including their \(D z^l\)
tangent coefficient. Thus the complete \(B'_0\) and
\(F_{\rm line}\) are again unchanged. In both rows the
reference-fixed norm root and the actual finite
character in (EG.27) are carried by that same map.

\item The ordinary full Miller product has weight \(D^{-6}\)
in both rows of (EG.7). The old factor
\((t-t_c)/t^3\) has weight \(D^6\), so the entire
ordinary suspension scalar is unchanged. All of
(EG.28)--(EG.29), including their fractional unit roots,
are holomorphic on the disc \(\zeta=w^2\), with their
values fixed once at the full reference parameter.

\item The canonical point frame and the tensor-reference frame have the
following directed comparison. With the same canonical meromorphic
section coefficient \(v_M\), they are
\[
 V_{\rm pt}=f_{P^a}v_M,\qquad
 V_{\rm tensor}=f_P^a f_a^{-1}v_M.
 \tag{EG.29a}
\]
The source map carries the canonical section literally and preserves
the full \(Q,P,Z\), so \(V_{{\rm pt},{\rm old}}=V_{{\rm pt},{\rm new}}\).
The whole Miller function has weight
\(f_{a,{\rm old}}=D^{1-a}f_{a,{\rm new}}\). Consequently
\[
 V_{{\rm tensor},{\rm old}}=D^{a-1}V_{{\rm tensor},{\rm new}},\qquad
 q_{\rm old}^{h}Z_{\rm old}^{u}V_{{\rm tensor},{\rm old}}^{j}
 =D^{-3h+(a-1)j}q_{\rm new}^{h}Z_{\rm new}^{u}
                     V_{{\rm tensor},{\rm new}}^{j}.
 \tag{EG.29b}
\]
These exact units enter the annular map below. On this selected-cusp
disc \(D\) is nonzero and its fixed logarithm is holomorphic in
\(\zeta=w^2\); thus its full parameter phase, as well as its norm,
is retained. The root point-line coefficient in the good row is
unchanged, while its graph monomials retain the \(D^s\) just displayed
above. This specifies both sides of each Cartier arrow.

\item Here is the whole annular square. First work over the
\(\zeta\)-disc. Retain the actual source, root and Cartier map
\(I_z\) from (EG.18)--(EG.29). On the fixed angular sectors put
\[
 s_\alpha=\frac{\log_\alpha q_\infty}{2\pi i},\quad
 z_\alpha=e^{2\pi i s_\alpha/m},\quad
 c_\lambda=\frac{n\log R_\lambda}{2\pi i}+\frac dm,\quad
 f_\alpha=\operatorname{Log}_\alpha\widetilde F .
 \tag{EG.30}
\]
Use the actual extension-group Lie coordinates: the
source component is the Abel integral of \(-R\), and
the line component contains (EG.21) or (EG.24), the
full \(h_\infty\), and the norm root. On an overlap
\(s_\beta=s_\alpha+k\), write
\(f_\beta=\mathfrak D_*^k f_\alpha+\Pi_{\alpha\beta}\).
Retain the entire reference period columns, suspension
integers \(j_{\alpha\beta}\), and integral vectors
\(\Lambda_{\alpha\beta}\). Then
\[
 \begin{split}
 \Phi_\alpha&=\operatorname{Exp}
     (n s_\alpha f_\alpha+s_\alpha c_\lambda\delta
                              +h_\alpha)I_{z_\alpha},\\
 h_\beta-\mathfrak D_*^k h_\alpha
   &=-n s_\alpha\Pi_{\alpha\beta}
              +j_{\alpha\beta}b_\beta+\Lambda_{\alpha\beta},
 \qquad b_\beta=\operatorname{Log}_\beta B .
 \end{split}
 \tag{EG.31}
\]
The period functions are transported over the disc,
with their original integer coefficients. In particular
the nonidentity reduced division branch is retained in
these periods; no small logarithm replaces it.

\item Use the invariant differential and rigidified-line
frame to express the vectors on the ordinary annulus.
Extend the whole reference \(h_\alpha\) to \(h_\alpha^0\).
For the residual additive cocycle \(E_{\alpha\beta}\),
zero at the reference, put
\[
 k_\alpha=-\sum_\beta\chi_\beta E_{\alpha\beta},\quad
 \omega=\bar\partial k_\alpha,\quad
 v_{\rm CG}(q,\zeta)=\frac1\pi\int
                   \frac{\omega(t,\zeta)}{q-t}\,dA(t),\quad
 h_\alpha=h_\alpha^0+k_\alpha-v_{\rm CG}.
 \tag{EG.32}
\]
The fixed larger-annulus cutoff and contours are those
of the actual EX construction. This solves (EG.31)
holomorphically in \(\zeta\), with the entire prescribed
reference value. Finite averaging applies only to the
vector cochains. Subtracting the target-sheet
\(knf_\beta+kc_\lambda\delta\) leaves exactly
\(j_{\alpha\beta}b_\beta+\Lambda_{\alpha\beta}\), hence
\[
 \Phi_\beta=A^k B^{j_{\alpha\beta}}\Phi_\alpha,\qquad
 \Phi_\alpha B=B\Phi_\alpha .
 \tag{EG.33}
\]
The inverse is \(I_z^{-1}\) followed by the negative
Lie argument. Pull this whole construction back by
\(\zeta=w^2\). It is literally even in \(w\), includes
both signed norms and all finite characters, and has
identity source, line, root and annular closing.

\par\medskip\noindent\textbf{EG.4. Actual orbit topology, fixed trace and both normal caps.}\par\nopagebreak[4]

\item For each fixed row and root order, choose the fixed
scalar below the common bounds on the finite exterior
and root charts. The actual cap generator \(H=A^n\)
satisfies \(H^m=B\), \(H^n=A\), and has scalar factor
\(R_\lambda\). Its strictly contracting norm gives
the compact shell and excludes an additional physical
circle stabilizer. Together with the protected
selected-fan regions these are the proper outside
products over the source disc.

\item The reference sphere is the corresponding F41 datum.
Its application to the present source uses the
contractible marked arc from the present
\(\rho=w^2/(w^2-3)\) to \(\rho=-1\), avoiding all
other collisions in its interior. The full K96
unselected-III disc attachment at the endpoint is
the actual E6-B model. The source/line/root maps
and their inverses above give the completed
family on a tube about that arc, with the same
reference and character. Proper interval
transport therefore supplies the F41 sphere here.

\item On the standard selected fan the only fixed set
is the vertical \(\mathbf P^1\), whose chart is
\[
 c_i'=Zc_i,\quad Z'=Z^{-1},\quad
 q=Zc_0c_1=Z'c_0'c_1',\qquad
 \operatorname{wt}_\delta(c_0,c_1)=(-1,+1).
 \tag{EG.34}
\]
Every other edge has determinant \(\pm1\) with
the physical cocharacter. On the free complement
\(E\) of its normal tube, the F41 sphere gives
the integral homology of \(S^3\) and \(\pi_1E=0\).
Conjugating the negative line gives
\(\mathcal O(1)\oplus\mathcal O(-1)\); its SU(2)
clutch \(\operatorname{diag}(Z,Z^{-1})\) bounds
the explicit disc \((r e^{i\theta},\sqrt{1-r^2})\).
Thus the boundary is the product Hopf map
\(S^2\times S^3\to S^2\times S^2\).
The primitive normal \(H_3(\partial E)\to H_3(E)\)
map and natural Gysin integration give
\(H^2(E/S^1)=\mathbb Z e\), \(H^3=H^4=0\).
The orbit union with \(S^2\times D^3\) has
degree-two MV matrix \(\left(\begin{smallmatrix}0&1\\-1&0\end{smallmatrix}\right)\).
It is simply connected and has the homology
of \(S^5\); the already verified F7--8
classification identifies its fixed pair with
\((S^5,\text{unknotted }S^2)\). The boundary
Hopf gauges match by their global logarithm
on \(S^2\times S^2\).

\item The full local height is unchanged from K95:
with \(\xi=X-v\),
\[
 Y^2-\xi^2(\xi+w^2)=q(q+v^2+D\xi),\qquad
 q=\frac{2\{Y^2-\xi^2(\xi+w^2)\}}
 {v^2+D\xi+\sqrt{(v^2+D\xi)^2+
                    4\{Y^2-\xi^2(\xi+w^2)\}}}.
 \tag{EG.35}
\]
The square root has value one at the cusp.
Its ordered normalization is
\(f(S,w)=(S^2-w^2,S(S^2-w^2),w)\).
Put \(x=S-w,y=S+w\) and
\[
 n_f=(3S^2-w^2)d\xi-2S\,dY+2w(S^2-w^2)dw .
 \]
The actual native normal arrows are
\[
 dq=-xy\,n_f/(v^2+Dxy),\quad
 d(c_0c_1)=-y^2n_f/(v^2+Dxy),\quad
 d(c'_0c'_1)=-x^2n_f/(v^2+Dxy).
 \tag{EG.36}
\]
The common denominator is an extending unit.
The source translation \(-bG\) and its line
lift are untwisted by their actual extension-group
logarithms on this source neighbourhood. Their
normalization action is
\[
 S\longmapsto
 \frac{(1+g_t)S+(1-g_t)w}{((1-g_t)/w)S+(1+g_t)},
 \qquad g_t=\exp(t\log P^{-b}).
 \tag{EG.37}
\]
Writing \(g_t=1+wh_t\), its difference from
\(S\) is \(h_t(S^2-w^2)/(2+h_t(w-S))\).
Its full derivative is uniformly \(I+O(|S|+|w|)\).
Together with the exact increasing F9 orbit
coordinate this gives the same actual ordered
embedded trace: \(S=-w\) is lower, \(S=w\)
upper. The exterior/root extension (EG.33)
gives its unchanged outer product. Interpolation
on the common real gamma lift gives the literal
height \(\operatorname{Re}(S\overline w)/
\sqrt{|S|^2+|w|^2}\), preserving the branch order.

\item The point-line arrow on this trace is also
explicit. Let \(\sigma_a\) be its canonical
rational section, with divisor \(aG-O\).
The map \(\sigma_a\mapsto S-S_a\) is a
nonzero holomorphic trivialization on
\(\mathbf P^1\times D_w\): at \(aG\) the
two zeros cancel, and at \(O\) their simple
poles cancel. Equivalently the actual
comparison to the \(a=6\) trace line sends
\[
 \sigma_a\longmapsto
       \frac{S-S_a}{S-S_6}\,\sigma_6 .
 \tag{EG.38}
\]
It is a Cartier line isomorphism, including
both marked points, and is holomorphic on
the \(\zeta=w^2\) disc.
For a local coefficient \(g\) of this map,
the normal variables are carried by
\((U,V)\mapsto(g^{-1}U,gV)\). The actual
positive-spinor coefficients consequently
change by \((\overline g^{-1},g)\).
Their phases agree and all full positive
norms are retained. The scalar logarithms
extend on the inner source ball and,
after a collar cutoff, on the outer
\(D^3\times S^1\) cap. This gives an
actual line/normal comparison, not an
identification from degree zero alone.

\item The remaining coefficient map is precisely
the K68 anti-linear map
\[
 (a_+,a_-)\longmapsto
 \frac1{\sqrt{|x|^2+|y|^2}}
 \begin{pmatrix}y&-\overline x\\x&\overline y\end{pmatrix}
 \overline{(a_+,a_-)} .
 \tag{EG.39}
\]
The outer cap is
\((P,q_H,w)\mapsto(A_wPA_w^*/\operatorname{tr}(A_wPA_w^*),
q_H,w)\), \(A_w=\left(\begin{smallmatrix}1&-w\\1&w\end{smallmatrix}\right)\).
It keeps the actual quaternion \(q_H\).
The inner cap is the same projection-normal
four-ball, with the extending factor in
(EG.36). The inverse line \(U=M^{-1}\)
and physical line \(M\) are carried over
both caps. The orbit calculation above
gives the actual integral Thom extension
and the same \(S^4\) normal-sphere boundary.
Thus the complete K95/F56 excision,
with relative second Chern number \(-1\),
has \(U\)-pair \((6,0)\) and physical
\(M\)-pair \((-6,0)\) in both rows.
The delta-averaged product connection
selects the return equal to the identity
on the fixed tube. Its literal K68
Hopf lift satisfies
\(\lambda_M(F)=-\Sigma_{F^{-1}}=\Sigma_F\).
This identifies the actual reference
native triple \((0,-6,-1)\).

\par\medskip\noindent\textbf{EG.5. Every vertical seed and every selected fan.}\par\nopagebreak[4]

\item Keep all coefficients (EG.14). On the
vertical fan \(T_j\) the full characters
and periods are
\[
 V_j=Z^{-j}V,\quad H_j=HB^j,\quad
 Q_j=\kappa qg^j,\quad
 \alpha_j=\alpha\eta^j\kappa^{-j}g^{-j^2},\quad
 \eta_j=\eta g^{-j}.
 \tag{EG.40}
\]
Here \(g=P^{-b}\), \(\alpha=P^a\).
The source interpolation
\(Q_t=Q\exp(tj\log g)\) has the exact
nodal vector field
\(\partial_t+j\log g\,y\partial_y\).
It is joined to the actual marked
period flow on the outside by a
vertical partition of unity, fixed
on the two protected collars.
Its Cartier arrow is \(Z^j\),
carrying the two generators
\(x^{-j-1}y^j,x^{-j}y^{j+1}\)
to \(1/x,y\). It retains (EG.38)
and both caps.

\item In the full real coefficient frame,
the q-deck is
\((u,d)\mapsto(u-b_{\rm ang},d-\gamma)\)
and the \(w\)-deck is
\((u,d)\mapsto(u+6b_{\rm ang},d)\).
All remaining unit logarithms occur
in these coefficients. The exact
relative connection difference is
\[
 A_j-A_0=-j\{du+(\alpha_q-6\theta)
                                     db_{\rm ang}\}.
 \tag{EG.41}
\]
The dual-line tensor has curvature
\(6j\,d\theta\wedge db_{\rm ang}\),
whose square is zero. Its relative
isomorphism is fixed by the actual
disc
\((q/Z_*,Z_*,1)\mapsto(qZ_*^{j-1},Z_*,Z_*^{-j})\)
with \(Z_*\) near \(-1\), away
from both marked point divisors.
After the integral cochain is
removed, its real Cech logarithm,
the extending \(\theta,b_{\rm ang},
\gamma\) characters and the
boundary logarithm give exactly
the specified outer line map.
The actual source/graph tangent
matrix is upper triangular in
the K63 frame order; its entire
comparison collar has zero
\(\operatorname{tr}(F^2)\).
For the relative \(p_1\) pairing use the actual level
cobordism of K63. Choose \(U_1\Subset U_2\Subset U_3\)
around the fixed tube, with all surgery and connection changes
in \(U_1\), and take the real lift \(\widetilde b\) of the
original source-product circle map \(b_{\rm ang}\) on \(U_3\).
For a cutoff \(\chi=1\) on \(U_2\), supported in \(U_3\),
and a fixed real constant \(c\), define the full relative homotopy
\[
 b_t=b_{\rm ang}
       \exp\!\bigl(2\pi i\,t\chi(c-\widetilde b)\bigr),
 \qquad 0\le t\le1.
 \tag{EG.41a}
\]
Outside \(U_3\) it is literally \(b_{\rm ang}\); at \(t=1\)
it is constant on \(U_2\). Choose a common regular value different
from \(\exp(2\pi i c)\). Its preimage under
\((x,t)\mapsto b_t(x)\) is the relative level cobordism, fixed on
the outer boundary. The terminal level misses \(U_1\), so it is
the identical cycle before and after surgery. The original level
is the actual source three-level times \(S^1_\gamma\); its
source-pulled \(p_1\) integral is zero. Stokes on the displayed
cobordism, with the retained zero outer \(p_1\) comparison and
transgression, gives zero on the terminal level. This proves the
required relative pairing without identifying the terminal level
itself with a product.
The outward Euler collar from
(EG.41) contributes \(-18j^2\);
the interior term is affine in
\(j\).

\item Both linear-endpoint values are
fixed by a literal family arrow.
From (EG.9)--(EG.12),
\[
 S_a(-w)=S_a(w),\quad
 c_a(-w)=P^{-a}c_a(w),\quad
 f_\theta(Z^{-1},P^{-a})
                  =P^a f_\theta(Z,P^a).
 \tag{EG.42}
\]
The last identity is the full
theta inversion at the same
\(Q\), with the actual lift
\(P=P_G(q,w)\). Thus the actual
source/line identity can also be checked directly on the full smooth
fibre. The source is the identity in \((X,Y,T)\); every coefficient
of the actual Miller functions (EG.3)--(EG.6) is holomorphic in
\(\zeta=w^2\). The canonical point-section coefficient is unchanged.
Since
\(\Theta_Q(Z^{-1})=-Z^{-1}\Theta_Q(Z)\), the exact tensor-coordinate
map is
\[
 \begin{split}
 V_a(-w)
 &=f_{P^{-1}}(Z^{-1})^a f_a(X,Y,T,-w)^{-1}v_a\\
 &=P^a f_P(Z)^a f_a(X,Y,T,w)^{-1}v_a
  =P^a V_a(w).
 \end{split}
 \tag{EG.42a}
\]
This gives the full source/line map
\((Z,V)\mapsto(Z^{-1},P^aV)\).
The target source period has
coefficient \(P^{-a}\); composing
it gives precisely
\[
 (Z,V)\longmapsto(Q/Z,V).
 \tag{EG.43}
\]
It carries every full root and
annular map because those were
constructed on the \(\zeta\)
disc before pullback. Its height
translation is one, so it sends
\(T_0\) to \(T_1\) with the
same physical characteristic
pair. Equations (EG.39)--(EG.43)
therefore give, for every integer
\(j\),
\[
 \begin{aligned}
 (k_M,\ell_M)&=(-6-18j(j-1),0),\\
 A_M&=-1-3j(j-1),\\
 (m_{\rm native},n_{\rm native})&=(0,-6-18j(j-1)).
 \end{aligned}
 \tag{EG.44}
\]

\item For an arbitrary selected lattice
basis \(P_{\rm bas}\), put
\(R_{\rm bas}=P_{\rm bas}^{-1}\).
In the following display \(P_{ij}\) and \(R_{ij}\) denote the
entries of these two matrices.
The actual character powers are
its two rows and its new period
powers are the two columns of
\(P_{\rm bas}\). Their complete
unit matrix is
\[
 A_{ij}=
 (\kappa^{P_{1j}}g^{P_{2j}})^{R_{i1}}
 (\alpha^{P_{1j}}\eta^{P_{2j}})^{R_{i2}},
 \qquad
 B_P=-6P_{\rm bas}^{-1}e_1e_1^tP_{\rm bas}.
 \tag{EG.45}
\]
The q-exponents are \(RP=I\).
The three full Cox edge rows are
\[
 \begin{gathered}
 (A_{11}A_{21}bc,A_{21}^{-1}ab,A_{11}^{-1}b^{-1}),\\
 (A_{12}A_{22}bc,A_{22}^{-1}c^{-1},A_{12}^{-1}ac),\\
 (a^{-1},ab,ac).
 \end{gathered}
 \tag{EG.46}
\]
Their normal determinants have
winding \(-6\), with first
windings
\((B_{22}-B_{21},B_{11}-B_{12},B_{22})\).
The actual based quaternion
contractions \(H_{-3}\) and
\(H_{l_i+3}\), together with the
full extending logarithms of
(EG.14), satisfy the exact
two-vertex tube identity
\(zD_{i,*}K_i=zD_i\).
Normal rescaling and the cutoff
flow of its velocity carry all
higher normal jets. Hence the
ordinary/native cylinder has
left degree \(-6\), right zero,
and inverse correction \(L^6\),
with the same odd Spin
endpoints and nine integral
prism equations as K71.

\item Each Euclidean fan step uses
the literal rank-one Cox ruling
map, common link
\((h,q_H)\mapsto(q_H^{-1}hq_H,q_H)\),
and its actual native reset.
They commute with every
coefficient in (EG.45).
The inverse character map on
the common exterior is the
identity in the original
source/Cartier coordinates,
so (EG.33) and its full finite
characters remain the same
at every step. The two compact
ruling charts, the unchanged
compact exterior and their
blowup/Gysin/primitive period
maps carry the sphere topology.
This identifies the entire
actual fanwise return algorithm
with K71, starting from all
the proved vertical seeds
(EG.44), for both full
progressions (EG.1).

\item The original parameter equation
\(r-1/3=-w^4/[3(w^2-3)^2]\)
is retained. This comparison
concerns the specified ordered
\(w\)-loop with fixed scalar.
No shorter source endpoint or
other E6 boundary is assigned
by these maps.

\end{enumerate}
\end{proof}


\section{Corrections with their located counterexamples}



\begin{theorem}\label{night:P3F.25}
\StatementRecord{night:P3F.25}{false as stated; correction proved}{S6\_06\_BLIND: GPT-6 Astra (OpenAI, ChatGPT Pro mode, effort Pro), recorded re-derivation in the checkpoint dated 18 November 2026}{\cref{night:P3F.10}, \cref{night:P3F.13}, \cref{night:P3F.16}, \cref{night:P3F.17}, \cref{night:P3F.2}, \cref{night:P3F.21}, \cref{night:P3F.22}, \cref{night:P3F.23}, \cref{night:P3F.24}, \cref{night:P3F.7}, \cref{night:P3F.8}, \cref{night:P3K.33}, \cref{night:P3K.35}, \cref{night:P3K.36}, \cref{night:P3K.6}, \cref{night:P3K.7}}The meridian-one real trace gives $A_U=(12s_\delta-p_{\rm rel})/24=-R\cdot O$, and the physical return is $A_M=R\cdot O$, with both mapping-torus and bordism boundary reversals included. The original and coprime higher reference-fan returns are zero. For the Viaclovsky and Liu--Rollenske source data, the positive loop of the explicit cover $\lambda=r^4$ has reference-fan return zero; P3K.35 also constructs the closed cover $\lambda=w^2$. In the Liu--Rollenske notation $u=\lambda$, P3K.36 proves that one $u$-loop has no biholomorphic endpoint identification. The positive weighted OS47 loop has $A_M=4$, oriented class $4\pmod{28}$, and a mapping torus not diffeomorphic to the product.
\end{theorem}
\begin{proof}
\begin{enumerate}[label=\arabic*.]
\item We retain the complete source, line, cap and frame maps below.
The meridian-one graph line is (U=M^{-1}): the fibre map
(D_M(v)(w)=\langle v,w\rangle/\langle v,v\rangle), with the
Hermitian form linear in its second argument, satisfies
(D_M(e^{i\theta}v)=e^{-i\theta}D_M(v)).
Its (SU(2))-equivariant spinor extension is
((a,b)\mapsto(\overline b,-\overline a)).
The first-to-second cap maps in \cref{night:P3K.68} give
(\lambda_M(F)=-\Sigma_{F^{-1}}) and
(\lambda_U(F)=\Sigma_{F^{-1}}); cancelling the opposite
product collars gives (A_M=-A_U). We therefore distinguish the
intermediate meridian-one number (A_U=(4k_U-\ell_U)/24)
from the physical return (A_M) throughout this calculation.

\item Fix the positive Hopf, product-boundary and Eells--Kuiper conventions of P3F.22, and the positive gamma turn of P3K.6--7. The actual return invariant for such a turn is \[
 A=\frac{12s_\delta-p_{\rm rel}}{24}
   =-R\cdot O
   =-\frac{d\,h(R_{\rm source})-2\chi(\mathcal O)
                  +\sum_v\operatorname{contr}_v(R)}2. \tag{P3F.25.1}
\] Here \(d\) is the actual regular-binding cover degree and \(R\) is the section specifying the line, rather than the section specifying the suspension translation. Thus the sign \(\epsilon\) in the prompt's expression is \(-1\) in these conventions. The local height contributions have the positive sign inside the last numerator.

\item The following table applies the actual trace to the pages already constructed in P3F.10, P3F.17--18 and P3F.21, with their specified reference fan. In the two one-log rows the displayed return is the positive \(r\)-loop of the explicit family \(\lambda=r^4\). For Liu--Rollenske the source scalar is \(u=\lambda\), so that row uses \(u=r^4\). P3K.35 constructs the further actual descent \(u=w^2\), whose reference-fan return is also zero. P3K.36 proves that continuation around one \(u\)-loop has no biholomorphic endpoint identification; this table assigns it no return class.

\begin{longtable}[]{@{}
  >{\raggedright\arraybackslash}p{(\columnwidth - 12\tabcolsep) * \real{0.46}}
  >{\raggedleft\arraybackslash}p{(\columnwidth - 12\tabcolsep) * \real{0.06}}
  >{\raggedright\arraybackslash}p{(\columnwidth - 12\tabcolsep) * \real{0.18}}
  >{\raggedleft\arraybackslash}p{(\columnwidth - 12\tabcolsep) * \real{0.08}}
  >{\raggedleft\arraybackslash}p{(\columnwidth - 12\tabcolsep) * \real{0.07}}
  >{\raggedleft\arraybackslash}p{(\columnwidth - 12\tabcolsep) * \real{0.05}}
  >{\raggedleft\arraybackslash}p{(\columnwidth - 12\tabcolsep) * \real{0.10}}@{}}
\toprule\noalign{}
\begin{minipage}[b]{\linewidth}\raggedright
Actual family and chosen line
\end{minipage} & \begin{minipage}[b]{\linewidth}\raggedleft
\(d\)
\end{minipage} & \begin{minipage}[b]{\linewidth}\raggedright
Page fibres
\end{minipage} & \begin{minipage}[b]{\linewidth}\raggedleft
\(\chi(\mathcal O)\)
\end{minipage} & \begin{minipage}[b]{\linewidth}\raggedleft
\(R\cdot O\)
\end{minipage} & \begin{minipage}[b]{\linewidth}\raggedleft
\(A\)
\end{minipage} & \begin{minipage}[b]{\linewidth}\raggedleft
\(\Theta _7\)
\end{minipage} \\
\midrule\noalign{}
\endhead
\bottomrule\noalign{}
\endlastfoot
Original \((3,4)\), \(M=P-O\) & 12 & \(12I_1\) & 1 & 0 & 0 & 0 \\
Coprime \((3k,4l)\), \(M=P-O\) & \(12kl\) & \(12klI_1\) & \(kl\) & 0 & 0 & 0 \\
Viaclovsky \(III^*+3I_1\), \(M=P-O\), order 4 & 4 & \(12I_1\) & 1 & 0 & 0 & 0 \\
Liu--Rollenske \(III^*+II+I_1\), \(M=P-O\), order 4 & 4 & \(4II+4I_1\) & 1 & 0 & 0 & 0 \\
Weighted OS47, \(M=7G-O\), \(Q=2G\), order 4 & 4 & \(4I_7+8I_1\) & 3 & 4 & \(-4\) & 24 \\
\end{longtable}

\item The OS47 row refers to its specified positive \(r\)-loop, \(\lambda=r^4\). It has order seven in \(\Theta _7\), and its mapping torus is not diffeomorphic to \(S^1\times S^6\), including without an orientation requirement.

\item \textbf{Proof: the actual circle actions.} Away from the scalar-zero fibre, use the actual complement of the zero section of the source line bundle, divided by its contracting lift of translation. A Hermitian norm and a sufficiently small scalar make every positive iterate decrease the norm strictly. A unit-circle rotation can equal an iterate only for the zero iterate, hence is free. This argument works at every smooth point of the original elliptic surface, including the nodal charts and the unchanged type \(II\) chart of Liu--Rollenske. For OS47 it applies to the whole \(I_7\) neighbourhood, with the cyclic component permutation retained, as proved in P3F.21. All these nonzero-linearization local models therefore have the required free circle action.

\item At a finite good-reduction cap of order \(m\), any putative line circle stabilizer must equal an affine cyclic deck transformation. Its invariant gamma displacement is \(jn/m\), where the logarithmic coefficient satisfies \(\gcd(n,m)=1\). The line circle has zero gamma displacement, so \(m\mid j\). The residual circle action on the root torus is free. This excludes all extra stabilizers at the finite caps, including the higher coprime multiplicities.

\item At the unique scalar-zero \(I_1\) fibre the valuation vectors form an integral basis. The unit-square triangulation gives exactly one fixed \(\mathbb P^1\) with normal \(\mathcal O(-1)\oplus\mathcal O(-1)\), of weights \((-1,+1)\), and is free off that curve. This is the exact local model in P3F.7. Thus each action in the table is semifree with precisely this fixed curve. The local orbit product and the real compression in P3K.6 apply at its chosen \(I_1\) cusp; the other singular fibres lie in the graph page and require no cusp replacement in that construction. The line matching in P3K.6 uses the integral meridian class on the entire complement, so it also retains their original circle line.

\item \textbf{Proof: characteristic numbers and the parameter direction.} The actual spin bordism P3K.7 compares the positive geometric family, with its circle Euler line, with the normal sphere bundle of the capped real trace. Its boundary directions are part of the comparison. Let \(f\) denote the quotient of the physical positive return in P3F.8. Write \(Q^+\) for its geometric mapping torus, identifying \((z,1)\sim(f(z),0)\). Then \(Q^+=M_{f^{-1}}\) in Jin's convention. With the circle-first quotient orientation and the complex orientation on the graph page, the capped complement of that graph has oriented boundary \(-Q^+\sqcup(-B)\). Indeed \(\partial(Y^5\times D^2)=-Y\times S^1\), and the inner normal sphere is a deleted-tube boundary. Hence the characteristic number on \(Q^+\) is the negative of that on \(B\), whereas it is also \(-A(f)\) by the reverse mapping-torus convention. These two signs cancel. Thus the meridian-one dual completion satisfies \[
 A_U=\frac{3b^2-p_1(T\bar X)}{24},\qquad
 b^2=4s_\delta,\quad
 p_1(T\bar X)=p_{\rm rel}.
\] The comparison is through that explicit bordism, not an identification of its cap with the usual elliptic cap. The relative square argument allows the numerical evaluation on the latter. On the compact elliptic page, the section line is \(R-O\) up to an integral fibre class. Since \(R^2=O^2=-\chi(\mathcal O)\), its square is \(-2\chi(\mathcal O)-2R\cdot O\). The tangent comparison gives \(p_{\rm rel}=3\sigma=-24\chi(\mathcal O)\). Substitution proves the first two equalities in \(P3F.25.1\). The third is exactly the height formula and degree multiplication in P3F.2.

\item For the original family its explicitly oriented sixfold parameter cover makes the positive gamma turn by P3F.13. P3K.33 supplies the actual simultaneous Cox descent to one original parameter loop, and its integer characteristic numbers give the zero reference-fan return. The same statement supplies the exact descended loop for the coprime higher multiplicities; its page data are P3F.17. The one-log families close on the scalar cover \(\lambda=r^4\): the same finite-cap equation as in P3F.23 gives the ordinary deck \(-\log(t-t_c)\delta/(2\pi i)\), and their nodal periods have multiplier \(q r^4\). The two-origin subtraction of P3F.23, with the corresponding holomorphic \(\Gamma,C\), gives their actual gamma turn. Their values are zero. The explicit quotient family of P3K.35 has \(\lambda=w^2\); its reference-fan return is zero by the characteristic numbers of that actual quotient. For the literal Liu--Rollenske scalar, P3K.36 proves that the once-around \(u\)-continuation has nonisomorphic multiple fibres at its endpoints, so it supplies no closed family loop. The locally labelled surface-deformation parameter of Viaclovsky ranges over a disc, so its loops are contractible. The separate labelled comb-smoothing loop has the explicit identity trivialization of P3F.16.

\item For OS47 the exact subtraction P3F.23 establishes one positive gamma turn for the positive \(r\)-loop. Its actual page has \(s_\delta=-14\) and \(p_{\rm rel}=-72\). Hence \[
 A_U=\frac{-168+72}{24}=-4,
 \qquad A_M=4,\qquad\mu(\Sigma_\Phi)=\frac4{28}.
\] P3F.22 gives the displayed oriented homotopy-sphere class, and P3F.24 gives its nonproduct mapping-torus obstruction for both spin structures.

\item For the primitive one-log extensions of the two \(III^*\) inputs with order \(4k\), the same actual page calculation gives respectively \(12kI_1\) or \(4kII+4kI_1\), degree \(4k\), \(\chi=k\), height \(2k\) and \(R\cdot O=0\). Thus their gamma-return integer is zero for the canonical family \(\lambda=r^{4k}\) by the identical trace calculation. P3K.35 gives the explicit further family \(\lambda=w^{2k}\), with zero reference-fan return. This page computation does not replace the separate integral topology check selecting sphere members among all logarithmic data.

\end{enumerate}
\end{proof}




\begin{theorem}\label{night:P3K.6}
\StatementRecord{night:P3K.6}{false as stated; correction proved}{S6\_06\_BLIND: GPT-6 Astra (OpenAI, ChatGPT Pro mode, effort Pro), recorded re-derivation in the checkpoint dated 18 November 2026}{\cref{night:P3F.10}, \cref{night:P3F.13}, \cref{night:P3F.9}}Let one positive gamma turn denote the actual parameter cover whose
angular clutch is P3F.13.  The orbit family of this cover is the product
\(Y\times S^1_\theta\), with \(Y=S^5\), and its fixed trace is

\[
 T_F=\{(R_\theta(p),\theta):p\in F\},
\]

where \(R_\theta\) is the gamma rotation on the cusp orbit product of
P3F.9.  Use its product trivialization preserving the Seifert gamma
coordinate.  In these coordinates \(R_\theta\) is pure gamma
rotation, including the explicitly trivialized source-translation
suspension.  There is a compact properly embedded oriented four-manifold
\(X\subset Z=Y\times D^2\), with boundary \(T_F\), assembled from the
actual page \(V\) and a product compression trace.  Its complement
has a unique integral class \(x\) with meridian value one.  On the
outer boundary this class is the Euler class of the Hermitian-dual delta line $U=M^{-1}$.
On the graph copy of \(V\), its restriction to the framed normal
section is exactly \(-c_1(M)\).
\end{theorem}

\begin{proof}
\begin{enumerate}[label=\arabic*.]
\item We retain the complete source, line, cap and frame maps below.
The meridian-one graph line is (U=M^{-1}): the fibre map
(D_M(v)(w)=\langle v,w\rangle/\langle v,v\rangle), with the
Hermitian form linear in its second argument, satisfies
(D_M(e^{i\theta}v)=e^{-i\theta}D_M(v)).
Its (SU(2))-equivariant spinor extension is
((a,b)\mapsto(\overline b,-\overline a)).
The first-to-second cap maps in \cref{night:P3K.68} give
(\lambda_M(F)=-\Sigma_{F^{-1}}) and
(\lambda_U(F)=\Sigma_{F^{-1}}); cancelling the opposite
product collars gives (A_M=-A_U). We therefore distinguish the
intermediate meridian-one number (A_U=(4k_U-\ell_U)/24)
from the physical return (A_M) throughout this calculation.

\item \par\medskip\noindent\textbf{The actual orbit trivialization.}\enspace The ordinary and finite
orbit charts are independent of the parameter angle: the variation
of their period and affine-twist columns is in the real delta
direction.  The exact subtraction in P3F.13 leaves, on the cusp
overlap, the gamma translation \(R_\theta\); the residual delta term
is a periodic unitary gauge.  Apply the identity to the exterior
charts and \(R_\theta\) to the cusp chart.  Their overlap maps now
agree pointwise.  Since \(R_1=R_0\), these identifications give a
periodic trivialization of the orbit family.  The fixed sphere in
the cusp chart becomes precisely \(R_\theta F\).  This verifies the
product and the displayed trace for the actual angular family.

\par\medskip\noindent\textbf{The compression.}\enspace Choose the binding at a regular value
\(q=\epsilon>0\) in a small cusp disc.  First write the local
regular-binding Milnor function as
\(u(q)(q-\epsilon)e^{2\pi i12\gamma_{\rm old}}\), with \(u\ne0\),
choose its unit argument, and set
\(\gamma=\gamma_{\rm old}+\arg u(q)/(2\pi\,12)\).
This is an explicit base-dependent circle coordinate change,
commuting with rotation and preserving the integral gamma marking.
Make this unit correction before choosing the short arc and the
handle below.  At its node the source has
coordinates \(xy=q\).  The two normalization discs of the fixed
sphere project to \(x=0\) and \(y=0\) and are separated in the gamma
coordinate.  Keep their values and the rest of the fixed sphere
inside an arc \(J\) of length less than \(1/24\), as allowed by the
relative product construction P3F.9.  Write their actual real gamma
functions as \(g_x(z)\) and \(g_y(z)\), without changing the orbit
map or its page.

\item Use the following exact real handle in the node chart:

\[
 (z,\phi)\longmapsto
 \bigl(x=z\cos\phi,\quad y=\bar z\sin\phi,\quad
       \gamma=g(z,\phi)\bigr),
 \qquad 0\leq\phi\leq\pi/2,
\]

\item Here \(g\) interpolates from \(g_x(z)\) to \(g_y(\bar z)\), with
values in \(J\); for example use their convex linear interpolation
in \(2\phi/\pi\).  The two center values are distinct, so
\(g(0,\phi)\) is strictly monotone.  For \(z\ne0\), the pair
\((x,y)\) determines \((z,\phi)\) by its moduli and phases.  At
\(z=0\), gamma determines \(\phi\), and at least one of
\(\sin\phi,\cos\phi\) is nonzero.  These facts also give full
rank to the derivative, so the map is an embedding at the core.
Its exact base value is

\[
 q=xy=|z|^2\cos\phi\sin\phi\geq0.
\]

\item Restrict to \(|x|,|y|\leq a\) and \(q\leq\epsilon\).
The faces \(\phi=0,\pi/2\) are the two normalization discs at
separated gamma values.  The face \(q=\epsilon\) is the smoothing
annulus, and the remaining faces \(|x|=a\), \(|y|=a\) match the
transported exterior of the source family along the real path.
Their gamma functions extend over that exterior with values in
\(J\), since \(J\) is an interval.  Rounding only these product
corners gives the oriented compression body.  Its two disc-face
orientations are the complex normalization orientations: the
conjugation on the second face is canceled by its opposite boundary
orientation.  Equivalently, reversing it attaches the positive
smoothing one-handle to \(F\).  Its regular boundary is the source
torus embedded as a graph \(\gamma=g_\epsilon(p)\), with the
globally defined real-valued function \(g_\epsilon\) taking values
in \(J\).  Choose a collar-preserving height function to obtain
an embedded three-dimensional compression cobordism

\[
 W\hookrightarrow Y\times[r_0,1],\qquad
 \partial W=F\sqcup(-E_\epsilon).
\]

\item Its projection lies over \(0\leq q\leq\epsilon\) and inside \(J\),
by the exact displayed formula.  The trace is used in the product
with its height interval.  Topologically \(W\) is a solid torus
with an open three-ball removed.  Its meridian on the incoming
torus is the source vanishing cycle \(a\).  The explicit handle
also fixes the framing and avoids an uncontrolled normal
perturbation of the base coordinate.

\par\medskip\noindent\textbf{Choose the page and embed the graph.}\enspace The weighted circle
has degree twelve on the regular-binding Milnor function.  In the
already chosen unit-corrected coordinate, the page equation is

\[
 \arg(q-\epsilon)+2\pi\,12\gamma=\varphi.
\]

\item Along the real path \(0\leq q<\epsilon\), its twelve gamma values
are constant.  Choose \(\varphi\) so that all these values miss
\(J\), and so that the short arc \(12J\) lies strictly on one side
of the page near the binding.  This is possible because of the
chosen length of \(J\).  The interior of the actual page \(V\)
then misses the entire projected compression, not just \(F\).
At the binding the initial torus meets \(\partial V\); pushing the
page a small distance to the opposite page-normal side removes this
intersection as well.  Denote this pushoff by \(V^+\).  It is
isotopic to \(V\) in \(Y\setminus F\) and is disjoint from the
projected compression.

\item On \(\partial V=E_\epsilon\times S^1_\gamma\), define
\(f(p,\gamma)=r_0e^{2\pi i(\gamma-g_\epsilon(p))}\).
This is the exact inverse angular coordinate of the rotating
incoming graph.  Its integral marking is the gamma marking,
because \(g_\epsilon\) is a globally defined real-valued function.
Extend this to a smooth map
\(f:V\to D^2_{r_0}\): on a boundary collar multiply the displayed
value by a cutoff that decreases to zero and is flat at zero,
then set \(f=0\) on the remaining interior.  Its graph in
\(Y\times D^2\) is embedded because its projection to \(Y\) is the
original embedded page.  In the outer annulus attach

\[
 W\times S^1\longrightarrow Y\times D^2,
 \qquad ((y,r),\theta)\longmapsto
 (R_\theta y,r e^{2\pi i\theta}).
\]

\item The polar coordinate records both \(r\) and \(\theta\), so this
map is an embedding.  Its incoming boundary is exactly the graph
of \(f|_{\partial V}\), and its outgoing boundary is \(T_F\).
Choose the radial derivative of the graph collar nonzero.  Extend
that collar outward and adjust only the outer-annulus collar of
the compression to meet it.  In a normal chart of the common
binding these are graphs over the same radial collar, so smooth
interpolation gives an embedded union.  All this adjustment is
arbitrarily small and leaves the graph inside the closed inner
disc unchanged.  At \(\theta=0\) it is chosen on the compression
side of the page; the fixed opposite pushoff \(V^+\) remains
disjoint.  This is the asserted \(X\).  Novikov additivity gives
\(\sigma(X)=\sigma(V)\): the added product \(W\times S^1\) has zero
signature, as is also seen by reversing its circle factor.

\par\medskip\noindent\textbf{The euler line, including its clutch.}\enspace The integral Thom
sequence for the proper codimension-three embedding \(X\subset Z\)
contains

\[
 H^2(Z;\mathbf Z)=0\longrightarrow H^2(Z\setminus X;\mathbf Z)
 \longrightarrow H^0(X;\mathbf Z)=\mathbf Z
 \longrightarrow H^3(Z;\mathbf Z)=0.
\]

\item It gives exactly one class \(x\) with the chosen meridian value one.
The identical calculation for
\(T_F\subset Y\times S^1\) identifies
\(H^2((Y\times S^1)\setminus T_F;\mathbf Z)=\mathbf Z\)
by the same meridian.  The Hermitian-dual delta bundle has that value,
so its first Chern class is the restriction of \(x\).  Classification
of principal circle bundles by integral first Chern class supplies
an isomorphism of these specified bundles.  This is an integral
identification on the whole complement; it retains the angular
clutch of degree minus twelve, rather than deleting its gauge.

\item The normal bundle of the graph copy of \(V\) is the sum of its
oriented page-normal real line and the real tangent plane of the
parameter disc.  The graph derivative gives an upper-triangular
identification with this sum, hence a framing.  Take the section
pointing toward \(V^+\).  Its pushoff first moves its disc coordinate
to zero inside the radius \(r_0\); its \(Y\)-coordinate is off the
graph throughout.  It then moves radially from zero to the outer
parameter point \(\theta=0\).  In the outer annulus the trace at
\(\theta=0\) is exactly \(W\), whose projection misses \(V^+\).
This is a homotopy in \(Z\setminus X\) to the pushoff of the actual
page in the original parameter-zero exterior.  The restriction of
\(x\) along that normal section is therefore

\[
 s^*x=c_1(L_\delta^{-1}|_V)=-c_1(M).
\]

\item The last equality is P3F.10 followed by the explicit Hermitian dual map of \cref{night:P3K.68}.
Boundary collar choices change only its boundary gauge, and their
relative square is zero by the relative cup-product formula.
This completes the construction and the meridian comparison.

\end{enumerate}
\end{proof}




\begin{theorem}\label{night:P3K.10}
\StatementRecord{night:P3K.10}{false as stated; correction proved}{S6\_06\_BLIND: GPT-6 Astra (OpenAI, ChatGPT Pro mode, effort Pro), recorded re-derivation in the checkpoint dated 18 November 2026}{\cref{night:P1.11}, \cref{night:P3F.10}, \cref{night:P3F.21}, \cref{night:P3F.22}, \cref{night:P3F.24}, \cref{night:P3F.7}, \cref{night:P3F.9}, \cref{night:P3K.6}, \cref{night:P3K.7}}Let \(S\) be the minimal regular elliptic surface with section \(O\)
whose affine equation is

\[
 y^2=x^3-3f(t)x+2g(t),\qquad
 f=(t^2+1)^2+t,\qquad
 g=(t^2+1)^3+\frac32t(t^2+1)+\frac38.
\]

Put \(G=(t^2+1,\sqrt3/2)\), \(R=8G\),
\(M=\mathcal O_S(R-O)\), and \(Q=G\).  Choose any one of the four
nodal fibres as the simple zero of the linearization.  At the
ordinary fibre \(t=0\), perform the multiplicity-one logarithmic
gluing whose integral circle clutch is \(\widehat\gamma\).
The resulting construction of Engel's Remark 8.5 gives a proper
holomorphic family of compact complex manifolds diffeomorphic to
the standard \(S^6\), over a punctured scalar disc
\(0<|\lambda|<c\).

For the positive \(\lambda\)-loop, with the complex orientation and
the Eells--Kuiper convention of P3F.22, the return satisfies

\[
 A_M=3,\qquad \mu(\Sigma_\Phi)=\frac3{28},\qquad
 [\Phi]=3\in\mathbf Z/28.
\]

In particular this is a generator.  The positive finite parameter
pullbacks of this same family realize every oriented class in
\(\Theta_7\).  For every nonzero such class the mapping torus is
not diffeomorphic to \(S^1\times S^6\).
\end{theorem}

\begin{proof}
\begin{enumerate}[label=\arabic*.]
\item We retain the complete source, line, cap and frame maps below.
The meridian-one graph line is (U=M^{-1}): the fibre map
(D_M(v)(w)=\langle v,w\rangle/\langle v,v\rangle), with the
Hermitian form linear in its second argument, satisfies
(D_M(e^{i\theta}v)=e^{-i\theta}D_M(v)).
Its (SU(2))-equivariant spinor extension is
((a,b)\mapsto(\overline b,-\overline a)).
The first-to-second cap maps in \cref{night:P3K.68} give
(\lambda_M(F)=-\Sigma_{F^{-1}}) and
(\lambda_U(F)=\Sigma_{F^{-1}}); cancelling the opposite
product collars gives (A_M=-A_U). We therefore distinguish the
intermediate meridian-one number (A_U=(4k_U-\ell_U)/24)
from the physical return (A_M) throughout this calculation.

\item \par\medskip\noindent\textbf{The source and its section.}\enspace More generally put
\(R_0=t^2+a\),
\(f=R_0^2+bt\), and
\(g=R_0^3+\frac32btR_0+\frac38b^2\).  Exact expansion gives

\[
 g^2-f^3=
 \frac34b^2a(t^2+a)^2+
 \frac18b^3t^3+\frac98b^3at+\frac9{64}b^4.
\]

\item At \(a=b=1\), this is
\((48t^4+8t^3+96t^2+72t+57)/64\).  Its four roots are simple
and none is a root of \(f\).  At infinity the minimal coefficients
are units and the discriminant has order eight.  Thus the fibres
are exactly \(I_8+4I_1\).  The point \(G\) satisfies the equation
by direct substitution.

\item Rationality can be checked without identifying the surface by a
table entry.  Set \(u=x-(t^2+a)\).  The equation becomes

\[
 y^2=3(ut-b/2)^2+u^2(u+3a).
\]

\item On \(u\ne0\), the two rational functions

\[
 A_0=\frac{y-\sqrt3(ut-b/2)}u,
 \qquad B_0=\frac{y+\sqrt3(ut-b/2)}u
\]

satisfy \(A_0B_0=u+3a\).  Conversely,

\[
 u=A_0B_0-3a,\qquad
 t=\frac{B_0-A_0}{2\sqrt3}+\frac b{2u},\qquad
 y=\frac u2(A_0+B_0),\qquad x=t^2+a+u.
\]

\item These are inverse rational maps on nonempty open sets.  The smooth
minimal model is therefore a rational elliptic surface.

\item The section \(G\) is disjoint from \(O\).  At infinity, with
\(z=1/t\), its minimal coordinates are
\(X=1+az^2\), \(Y=(\sqrt3b/2)z^3\).
The two local nodal branches have orders three and five when
\(a=b=1\): expanding the square root of
\((1+az^2)^2+bz^3\) cancels the order-three term on exactly
one branch, whose next nonzero term has order five.  Their product
has order eight.  Hence \(G\) meets component three of \(I_8\),
up to reversing the cyclic numbering.  Its height is

\[
 h(G)=2-\frac{3(8-3)}8=\frac18.
\]

\item The root lattice has rank seven and determinant eight.  Shioda--Tate
therefore gives Mordell--Weil rank one.  The unimodular surface
lattice and the height discriminant formula give
\(h(G)=m^2|T|^2/8\), where \(m\) is its index in the free
Mordell--Weil group and \(T\) its torsion subgroup.  Thus
\(m=|T|=1\).  This identifies the actual generator and the
torsion-free \(I_8\) source, including the component marking.

\item The section \(R=8G\) meets the identity component at \(I_8\).
Its height is eight, so the same intersection formula gives
\(R\cdot O=3\).  Exact doubling locates all three intersections:
the reduced denominator of \(x(8G)\) is

\[
 (6t+1)^2(360t^2+282t+199)^2.
\]

\item The three roots are distinct and this denominator is coprime to
the discriminant.  The numerator is coprime to it as well, as
checked in the saved exact group-law calculation.  Consequently
all three intersections are ordinary and simple.  In particular
the Picard point of \(M\) is nonidentity at every \(I_1\) and
at the smooth identity component of \(I_8\).

\par\medskip\noindent\textbf{The actual marked local system.}\enspace We give a comparison
with the displayed source, rather than merely exhibit matrices
having the correct local spectra.  At \(a=0,b=1\) the source has
\(I_9+3I_1\).  Its \(j\)-map has degree twelve and ramification
profiles \(3^4\), \(2^6\), and \((9,1,1,1)\) over
\(0,1728,\infty\).  Indeed \(f=t(t^3+1)\) has four simple
roots, \(g=t^6+\frac32t^3+\frac38\) has six simple roots,
and \(g^2-f^3=(8t^3+9)/64\); these polynomials are pairwise
coprime.  Their indicated ramification already totals twenty-two,
so Riemann--Hurwitz leaves no further branch points.

\item The associated trivalent modular graph is unique.  It has four
vertices and six edges, with three monogon faces.  Each monogon
uses a loop at a vertex and consumes two of its three half-edges.
The three loops therefore occupy three distinct vertices, whose
remaining half-edges must join the fourth vertex.  This is the
three-star with a loop on each leaf.  Its cyclic ordering is unique
up to orientation-preserving graph isomorphism.  This proves
uniqueness of the permutation cover with these profiles by the
usual path-lifting classification of covers of the three-punctured
sphere.

\item For \(v=(p,q)^t\), write

\[
 D_v=\begin{pmatrix}1-pq&p^2\\-q^2&1+pq\end{pmatrix},
 \qquad T=D_{(1,0)}.
\]

\item The subgroup generated by
\(T^9,D_{(-3,1)},D_{(0,1)},D_{(3,1)}\) has the same cover:
its twelve-coset permutation table, exported by the program, has
exactly the three profiles above.  The products of the displayed
matrices equal the identity.  Thus the modular covering comparison
identifies this tuple with the actual \(I_9+3I_1\) source, up to
the explicitly permitted choice of marked vanishing paths and a
single integral symplectic basis.  The lift from projective
monodromy is fixed by the trace-two local \(I_n\) matrices.

\item Near infinity the deformation in \(a\) is analytically
\(xy=z^8(z+c a+O(a^2,az))\), with \(c\ne0\), after absorbing
units.  Hence it splits the \(I_9\) into an \(I_8\) and an
\(I_1\) with the same vanishing cycle.  Transport to \(a=1\)
along a path avoiding the finitely many other discriminant
collisions gives the actual five-point tuple

\[
 E_0=T^8,\quad E_1=T,\quad E_2=D_{(-3,1)},\quad
 E_3=D_{(0,1)},\quad E_4=D_{(3,1)}.
\]

\item Such transport is an isomorphism of the integral local systems
over the punctured surfaces and carries the section \(G\).
The path choice only changes the marking by the corresponding
Hurwitz moves.  No assertion about the topology depends on choosing
a particular branch of this transport: the program verifies all
four choices of the nodal zero below.

\item In the cohomology convention \((\psi,u,w,\delta)\) of P3F.21,
the weighted Jacobi lift for \(M=8G-O,Q=G\) is

\[
 J(E,k,z)=
 \begin{pmatrix}
 1&cj-ah&dj-bh&z\\
 0&a&b&8j\\
 0&c&d&8h\\
 0&0&0&1
 \end{pmatrix},
 \quad E=\begin{pmatrix}a&b\\c&d\end{pmatrix},
 \quad k=\binom jh.
\]

\item The component of \(G\) fixes \(k_0=(3,0)^t\).  An integral
change of Abel--Jacobi lift also sets \(k_4=0\).  Here is the
simultaneous gauge explicitly.  Adding the common integral lift
\((a,b)^t\) changes the first cocycle by
\((E_0-I)(a,b)^t=8b(1,0)^t\), so \(b\) fixes its component-three
representative.  The fourth coefficient changes by \(-a+3b\);
varying \(a\) then sets it to zero without changing the first.
The remaining
cocycles have the forms
\(k_1=r(1,0)^t\), \(k_2=j(-3,1)^t\),
\(k_3=h(0,1)^t\).  The nonzero-linearization local fillings
have rank-one nilpotent image; this forces their central entries
to be \(9,8r^2,8j^2,8h^2\), respectively.  At the scalar-zero
\(I_1\), the central entry is \(-1\).  The actual global product
must be the identity.  Its source columns give
\(j=2h\), \(r=-3h-3\), and its central entry gives

\[
 16(h+1)(4h+5)=0.
\]

\item The unique integral solution is \(h=-1\).  This proves the
actual section-lift identification; it is not a bounded search
for a convenient tuple.  The resulting matrices are

\[
\begin{aligned}
 T_0&=\begin{pmatrix}1&0&3&9\\0&1&8&24\\0&0&1&0\\0&0&0&1\end{pmatrix},
 &T_1&=\begin{pmatrix}1&0&0&0\\0&1&1&0\\0&0&1&0\\0&0&0&1\end{pmatrix},\\
 T_2&=\begin{pmatrix}1&2&6&32\\0&4&9&48\\0&-1&-2&-16\\0&0&0&1\end{pmatrix},
 &T_3&=\begin{pmatrix}1&1&0&8\\0&1&0&0\\0&-1&1&-8\\0&0&0&1\end{pmatrix},\\
 T_4&=\begin{pmatrix}1&0&0&-1\\0&-2&9&0\\0&-1&4&0\\0&0&0&1\end{pmatrix}.
\end{aligned}
\]

\item Moving the linearization zero to any of the other three \(I_1\)
fibres adds one to the central entry of \(T_4\) and subtracts
one at the chosen fibre.  The central matrices commute, so the
global product is still the identity.  Every integer calculation
below has been executed for all four choices.

\par\medskip\noindent\textbf{Line, fillings, and the ordinary logarithmic gluing.}\enspace
The line \(M\) has multidegree zero on every component of
\(I_8\).  Translation by \(Q\) therefore preserves its
multidegree.  The vertical line
\(\operatorname{Hom}(t_Q^*M,M)\) has zero intersection with
every fibre component and is a whole-fibre pullback.  Its degree,
computed by restriction to \(O\), is

\[
 (8G)\cdot O+1-(7G)\cdot O+(-G)\cdot O
 =3+1-3+0=1.
\]

\item Thus it is \(\mathcal O(F)\), and its section may have its
sole simple zero at any chosen nodal fibre.  Scaling the induced
lift by sufficiently small \(\lambda\) makes it uniformly
contracting.  On the complement of that fibre its quotient is a
smooth Hausdorff complex threefold, as follows directly from the
norm contraction and a compact annular fundamental domain.

\item At \(I_8\), the Tate and suspension periods, in a multidegree-zero
line frame, have valuation vectors \((8,0)\) and \((-3,0)\).
If they are denoted by \(h\) and \(g\), respectively, take

\[
 H=h^{-1}g^{-3},\qquad K=h^3g^8.
\]

\item The change has determinant one and their first valuations are
\(1,0\).  The second multiplier of \(K\) is \(\lambda^8\)
times a nonvanishing unit.  Thus the ray fan with consecutive
vertices \((j,1)\) is the actual primitive rank-one filling.
Equivalently it is the contracting quotient of the smooth source
line complement; \(Q\) cyclically permutes all eight components.
This equivalence retains the source elliptic map and its section
line.  The same primitive rank-one model applies at the other
three nonzero-linearization nodal fibres.  At the chosen zero
the two valuations are \((1,0),(0,1)\).  The square triangulation
has unimodular cones and gives the smooth rank-two cusp with
charts \(q=z_0z_1z_2\), its toric period identifications, and
proper reduced central fibre.  These are the specified local
models, with all valuation changes integral.

\item The base point \(p=0\) is ordinary.  On its smooth torus disc,
and on the universal logarithm cover of the scalar parameter,
let \(\gamma(z,v)\) be the holomorphic period of the primitive
suspension cycle. Glue a second copy of this smooth torus disc
by the following punctured-disc isomorphism, directed from the
cap copy to the ordinary copy:

\[
 \Psi_v(\zeta,z)=
 \left(\zeta+\frac{\log z}{2\pi i}\gamma(z,v),z\right),
 \qquad \lambda=e^{2\pi i v}.
\]

\item A logarithm branch change adds an actual lattice period, so this
is a single-valued holomorphic map on the torus family, with
inverse given by subtraction.  The cap is smooth, its central
fibre has multiplicity one, and its boundary clutch is exactly
\(\widehat\gamma\). Indeed a cap meridian with fixed cap coordinate
is carried to the ordinary base meridian plus \(\gamma\), the
\(m=1,n=1\) marking used below. This verifies the allowed logarithmic
transformation at an ordinary fibre even though its multiplicity
is one.  All pieces are proper over their closed base discs and
agree by the displayed holomorphic maps on collars; their finite
gluing is a compact Hausdorff complex threefold.

\par\medskip\noindent\textbf{Integral topology.}\enspace Put \(H_i=T_i^{-t}\) on the
homology lattice \(\Lambda\).  The images of \(H_i-I\) have
Smith factors \((1),(1),(1),(1),(1,1)\).  The primitive toric
collapses kill exactly these vanishing lattices.  We first fix
an actual global origin for the auxiliary untwisted family, before
the ordinary logarithmic transformation.  The restriction
\(M|_O\) has degree \(R\cdot O-O^2=4\).  Choose its section
with divisor \(4t_c\).  Its nonzero values over the other base
points give an origin in the quotient.  Near \(t_c\), its line
character is \(V=q^4\) times a holomorphic unit.  Subtracting
four gamma periods changes this to the nonvanishing bounded
character \(\lambda^{-4}e^{-8\pi iC}\) times a unit and changes
the source coordinate by the bounded translation \(-4\Gamma\).
It therefore extends in the dense torus of the cusp component.
This proves that the origin is global, including all five singular
fibres.  Every base meridian based at this origin contracts in
its local filling.

\item Van Kampen begins with the affine torus-bundle group over the
punctured sphere.  After these base meridians are killed, all
generators lie in the abelian fibre lattice, and the conjugation
relations impose the images of \(H_i-I\).  Their stacked relation
matrix has Smith factors \((1,1,1)\) and quotient \(\mathbf Z\)
measured by \(\psi\).  The ordinary log cap adds a vector with
\(\psi\)-value one.  The augmented Smith factors are
\((1,1,1,1)\), proving \(\pi_1=1\), with every affine origin
correction retained in the kernel already killed.  In the chosen
global-origin marking the complete relations are the abelian
fibre relations, \((H_i-I)\Lambda=0\) for \(0\leq i\leq4\),
and \(\widehat\gamma=0\).  Changing an origin only adds a vector
of \(\ker\psi\), which is precisely the primitive subgroup
generated by those five image lattices.

\item The complete cohomology calculation is integral.  For degree \(q\)
let \(W_i=\bigwedge^qT_i\), let \(L_i\) be the integer column
lattice of \(W_i-I\), and use its Hermite basis.  The marked
parabolic complex is

\[
 \mathbf Z^{\binom4q}\xrightarrow{d^0}
 \bigoplus_{i=0}^4 L_i\xrightarrow{d^1}
 \mathbf Z^{\binom4q},
\]

where \(d^0(a)=((W_i-I)a)_i\) and
\(d^1((b_i))=\sum_i(W_0\cdots W_{i-1})b_i\).
This is obtained by cutting the marked punctured sphere and
attaching each invariant stalk; the product relation gives
\(d^1d^0=0\).  The code exports every matrix in every degree.
For degrees one, two, and three, their ranks are respectively
\((3,3),(5,5),(3,3)\), and every nonzero Smith factor of
both differentials is one.  The middle groups vanish integrally;
the end groups are free of rank one.

\item There are no finite quotient caps in this example.  At each
primitive rank-one filling the collapse is the cellular pinching
of one circle in \(T^2\), crossed with the other elliptic factor;
it identifies cohomology with the full invariant lattice.  At
the rank-two filling the two valuation vectors and triangulation
identify the model integrally with the cusp of Engel's Lemma 7.4.
For an independent local chain verification, use the actual
dual-moment-cell quotient and blowup-collapse comparison of P1.11
in rank two, with \(G=3I_2-\mathbf1\mathbf1^t\).  Its translated
simplex counts are \((1,3,2)\), with twelve codimension-one
incidences.  For a simplex \(\sigma\), put
\(Q_\sigma=\mathbf Z^2/L_\sigma\).  The differential from
\(\bigoplus_{\dim\sigma=p}\bigwedge^q Q_\sigma\) is precisely
the signed integral quotient matrix (1.11.1), now with
\(0\leq p,q\leq2\).  The program
\path{code/part3_fourfold/rank_two_cusp.py} exports all these matrices
and all their Smith transformations.  Their total homology is
free of ranks \((1,2,4,2,1)\).  The collapse is the block sum
of \(\bigwedge^q(\mathbf Z^2\to Q_\sigma)\); its induced
homology matrices are onto with only unit Smith factors.  The
only nontrivial nonidentity block on the associated graded groups,
in total degree two, is

\[
 \begin{pmatrix}0&0&-1&1\\1&0&0&0\\0&-1&0&1\end{pmatrix}.
\]

\item Support leaves no possible higher differential.  The finite
filtration therefore makes the actual collapse onto and primitive
in every degree.  Its dual image is monodromy invariant and
primitive, with the same rank as the full invariant lattice, so
the two lattices agree.  This separately verifies every local
specialization assertion of the cited lemma needed here.
These exact local comparisons
give \(R^qf_*\mathbf Z=j_*\bigwedge^q\Lambda^*\) everywhere.

\item The surviving invariant generators, in degrees zero through four,
are

\[
 1,\quad\psi,\quad\omega=8uw+\psi\delta,
 \quad\eta_3=\psi uw,\quad\eta_4=\psi uw\delta.
\]

\item The corresponding coinvariant generators are
\(1,[\delta],\eta_2,[uw\delta],[\eta_4]\), with
\([\psi\delta]=8\eta_2\), \([uw]=\eta_2\), all other
degree-two basis classes zero.  The only nonzero Leray columns
are zero and two.  There is no unaccounted intrinsic cusp
differential.  Inversion in the untwisted torus family about the
global origin just constructed extends over every local filling:
on primitive nodal charts
\(X_i=Z/q^i\), \(Y_i=q^{i+1}/Z\), it is the actual holomorphic
map \((X_i,Y_i)\mapsto(Y_i,X_i)\) in chart \(-i-1\).
The other elliptic coordinate is inverted as well.  On the
rank-two square fan the lattice map \((v,h)\mapsto(-v,h)\)
sends each cone to its centrally inverted cone; its dual monomial
maps are holomorphic chart isomorphisms and agree on overlaps.
These are integral automorphisms of the specified fans.
The bounded origin characters only add a toric unit translation,
which extends too.  Thus this is an involution of the actual
untwisted compact family.  It acts by \((-1)^q\) on \(R^q\).
Naturality of Leray then makes every untwisted \(d_2\) satisfy
\(d_2=-d_2\).  Its targets are the free groups just computed,
so every one is zero integrally.

\item Replacing the ordinary cap changes only its boundary translation.
The cellular prism for this actual translation changes \(d_2\)
by contraction with its integral clutch.  Consequently the
translation prism of the ordinary
logarithmic clutch gives
\(d_2(a)=H[\iota_{\widehat\gamma}a]\), where \(H\) is the
positive base class.  The four contractions are

\[
 [\iota_\gamma\psi]=1,\quad
 [\iota_\gamma\omega]=[\delta],\quad
 [\iota_\gamma\eta_3]=\eta_2,\quad
 [\iota_\gamma\eta_4]=[uw\delta].
\]

\item For an arbitrary integral affine origin correction, the same
calculation is exactly its \(\psi\)-coordinate; source and line
circle coordinates contribute zero to all four projections.
Thus every displayed differential is multiplication by one for
the actual glued family.  There are no higher differentials,
because the base has dimension two.  Only
\(H^0=H^6=\mathbf Z\) remains, without torsion or extension
ambiguity.  Together with simple connectivity this proves that
the total manifold is a homotopy six-sphere, and \(\Theta_6=0\)
identifies its smooth oriented type with the standard sphere.

\par\medskip\noindent\textbf{The closed positive parameter loop.}\enspace Write \(t_c\) for
the chosen nodal zero and put
\(h(t)=(t-t_c)/t\).  Since
\(\gamma(z,v+1)=\gamma(z,v)+\delta\), the smooth logarithmic
gluing changes by \((\log t)/(2\pi i)\delta\) under the
positive parameter deck.  The global ordinary deck is
\(-\log h(t)\delta/(2\pi i)\): near \(t=0\) it has exactly
that value up to a holomorphic unit logarithm, which extends on
the cap. More explicitly, write
\[
 D_{\rm ord}=-\frac{\log h(t)}{2\pi i}\delta,\qquad
 D_{\rm cap}=-\frac{\log(t-t_c)}{2\pi i}\delta.
\]
The second logarithm is holomorphic on the chosen smooth cap.
The directed gluing satisfies the exact compatibility
\[
 D_{\rm ord}\circ\Psi_v=\Psi_{v+1}\circ D_{\rm cap},
 \qquad
 D_{\rm ord}=D_{\rm cap}+\frac{\log t}{2\pi i}\delta .
\]
Here the \(D\)'s denote the respective translations from \(v\)
to \(v+1\). At the cusp use \(q=h(t)\), \(s=\log q/(2\pi i)\).
Then the actual periods are

\[
 \gamma(s,v)=\binom{\Gamma(q)}{C(q)+s+v},\qquad
 \delta=\binom01,
\]

with \(\Gamma,C\) holomorphic at zero, and the deck translates
by \(-s\delta\).  Its line-character multiplier \(q^{-1}\)
extends through the cusp: compose an inverse gamma lattice shift
with the bounded source translation \(\Gamma\) and the
nonvanishing multiplier \(\lambda e^{2\pi iC}\).  At every
other original fibre \(h\) is a unit.  This proves that the
glued family descends holomorphically over the punctured
\(\lambda\)-disc, with the stated positive closing map.

\item Here \(\gamma\) is the inward period of the actual contracting
suspension. In the physical Abel and line coordinates it is
\[
 B(a,V)=
 \bigl(a+\Gamma(q),\,e^{2\pi iC(q)}q\lambda V\bigr),
 \qquad \Gamma(q)=\log_{\rm Tate}(-G).
\]
For small \(q,\lambda\), its positive real direction decreases
\(\log|V|\). This fixes its orientation relative to the source
page. With \(\delta\) the positive phase of \(V\), the complex
orientation gives
\[
 \operatorname{or}(Y)
   =-\operatorname{or}(S)\,d\log|V|
   =\operatorname{or}(S)\,d\gamma .
\]
Thus the positive gamma in the Hopf page below is this inward
period. The normalized smooth norm coordinate increases in
this direction.

\item For a completely marked angular comparison use the auxiliary
logarithmic cover \(q=Q_0^2\), and the two origins

\[
 o=\binom{-\Gamma s}{-Cs-s^2/2-sv},\qquad
 \sigma=\binom{\Gamma v}{Cv+(v^2-v)/2}.
\]

\item The first has parameter difference \(-s\delta\); the second
has difference \(-s\delta+\gamma\).  Under \(s\mapsto s+2\),
their differences are \(-2\gamma-2\delta\) and zero.  Thus
they descend on this auxiliary cover.  In contrast, the ordinary
origin has one-step difference \(-\gamma-\delta/2\); we do
not impose a false one-step periodicity.  The cusp origin extends
by its bounded unit characters.  Subtracting the ordinary origin
removes the parameter deck and gives the one-node parameter chain
\(X_i=V/\lambda^i\), \(Y_i=\lambda^{i+1}/V\),
\(X_iY_i=\lambda\), over every nonzero \(q\).  Its origin
extends.  These are the two extending collar origins.  No proper
two-variable corner or descent of such a corner is used.

\item Their exact difference is

\[
 F=\binom{\Gamma(s+v)}{C(s+v)+(s+v)^2/2-v/2},
\]

and hence

\[
 F(s,v+\theta)-F(s,v)
 =\theta\gamma(s,v)+\frac{\theta(\theta-1)}2\delta.
\]

\item The additional line gauge is periodic, equals one at both
endpoints, and vanishes in the circle quotient.  The double cusp
deck changes this angular difference by \(2\theta\delta\),
so the retained line gauge has winding two.  In particular the
actual positive \(\lambda\)-loop makes one positive gamma turn,
with no parameter cover.  Every identity is executed in
\path{code/part3_fourfold/os45_origins.py}.

\item More explicitly, the ordinary and cusp coordinates are
\(\zeta-o\) and \(\zeta-\sigma\), so the cusp-to-ordinary
transition is translation by \(F=\sigma-o\). Its angular
increment is the positive translation displayed above. This
identifies the moving fixed sphere in the constant orbit with
the positive inward gamma motion used by the real trace.

\par\medskip\noindent\textbf{The actual page and the return index.}\enspace The line circle is
free away from the selected rank-two cusp.  Indeed its equality
with an iterate of the contracting lift forces that iterate to
be zero by the norm inequality, and then forces the circle
element to be the identity.  This includes the whole \(I_8\)
neighbourhood, where translation permutes the components.  The
ordinary logarithmic cap is free as well.  At the single cusp
the primitive fan gives the one fixed rational curve with normal
weights \((-1,+1)\).  The orbit-pair argument P3F.7 therefore
applies to this proved sphere: its quotient is the standard
\((S^5,S^2)\), and its dual meridian line has positive meridian one.

\item The gamma circle over the base has the one ordinary clutch of
degree one, so its total space is the standard Hopf \(S^3\).
For precision, in the two standard Hopf charts the normalized
sections \((1,t)/\sqrt{1+|t|^2}\) and
\((1/t,1)/\sqrt{1+|1/t|^2}\) have angular coordinates related
by \(\theta_\infty=\theta_0+\arg t\), exactly the clutch
above.  A regular Hopf binding has a complex-oriented disc page
whose projection to \(\mathbf P^1\) has degree one.  Its inverse
image in the actual orbit is the elliptic page \(V\), with
boundary \(T^3\); its regular cap is the displayed rational
elliptic surface \(S\).  The affine origin obstruction on the
disc with marked singular values vanishes by the free-meridian
cochain argument of P3F.10.  Thus this comparison is with the
actual page, including the retained \(I_8\), rather than with
an unrelated elliptic surface.

\item The actual circle line restricts to \(M=8G-O\), up to a base
pullback, by its fibrewise Picard point.  At \(I_8\) both lines
have multidegree zero, so the same normalized fibrewise comparison
extends across its cycle of components.  On the base with boundary
the pullback is topologically trivial.  Therefore the actual
relative page invariants are

\[
 s_\delta=M^2=-2-2(3)=-8,
 \qquad p_{\rm rel}=3\sigma(S)=-24.
\]

\item Near the selected \(I_1\) the source translation is in the smooth
identity component.  The product coordinate of P3F.9 and the real
phase-preserving compression of P3K.6 consequently apply with
page degree \(d=1\).  The page phase can be chosen to avoid the
short gamma arc of that compression.  Its exact graph boundary,
integral Euler-line restriction, and oriented complement bordism
are therefore the ones proved in P3K.6–7.  The angular computation
above identifies this bordism with the actual positive scalar
return.  With the two orientation reversals of P3K.7 retained,

\[
 A_U=\frac{12s_\delta-p_{\rm rel}}{24}
   =\frac{-96+24}{24}=-3,\qquad A_M=-A_U=3.
\]

\item The typed Hopf transfer and Eells--Kuiper normalization P3F.22 give
the asserted oriented class \(3\).  Pullback by
\(\lambda=z^n\) is a proper holomorphic family on the punctured
disc and its positive loop is exactly the \(n\)-fold iterate.
The classes \(3n\pmod{28}\), for positive \(n\), exhaust
\(\mathbf Z/28\).  The mapping-torus obstruction P3F.24 applies
to each nonzero class and to both spin structures.  This proves
all the assertions of the statement.

\end{enumerate}
\end{proof}


The scope of this reduction is the actual whole-fibre cusp
projection specified at the start.  It does not extend the old
gamma coordinate through the component-supported OS49 cusp;
that construction uses its separate mixed \(K\) marking.
The Seifert projection also does not recover the finite-kernel
source quotient or the line character from \(B\) alone.
Those data remain part of every cap in the enumeration.

An arbitrary periodic unimodular fan need not have this fixed-sphere
model.  The present map and its use of P3F.8 are not asserted for
such a fan without its separate orbit comparison.

\begin{theorem}\label{night:P3K.17}
\StatementRecord{night:P3K.17}{false as stated; correction proved}{S6\_06\_BLIND: GPT-6 Astra (OpenAI, ChatGPT Pro mode, effort Pro), recorded re-derivation in the checkpoint dated 18 November 2026}{\cref{night:P3F.22}, \cref{night:P3F.24}, \cref{night:P3F.32}, \cref{night:P3K.7}}
The real master-space resolution, its two moment-face maps,
inner free-orbit line and charged joint-cap maps below are defined
for the actual OS49 local square. Their common inner correction
has the displayed vanishing characteristic pair. The historical
assertion that the two complete OS49 returns are equal is false:
for cap order (m=3k) their physical whole-loop classes are
(-3k) on the negative fan and (0) on the positive fan.
For their actual shorter loop the classes are (-k) and (0),
by the complete marked outer comparison in \cref{night:P3K.65}.
\end{theorem}

\begin{proof}
\begin{enumerate}[label=\arabic*.]
\item \par\medskip\noindent\textbf{The real master space and its actual resolution maps.}\enspace

\item Order the Cox coordinates as \((z_0,z_4,z_1,z_2)\), corresponding
to rays \(u_0,u_4,u_1,u_2\).  Let an auxiliary circle act with
weights \((1,1,-1,-1)\), and let the physical delta circle have
weights \((-1,0,0,+1)\).  These actions commute.  Write
\(a=(z_0,z_4)\), \(c=(z_1,z_2)\), and
\[
       \mu=\|a\|^2-\|c\|^2.
\]
The auxiliary circle is free on \(\mathbb C^4\setminus\{0\}\).
Consequently its quotient is a smooth oriented real seven-manifold.
The residual delta action is semifree.  Indeed every nonzero
two-by-two minor of the matrix with rows
\((1,1,-1,-1)\) and \((-1,0,0,1)\) is \(\pm1\).
On a coordinate support of rank two this proves that the
two-torus has trivial stabilizer.  A rank-one support lies in
one of the two pairs \((z_4,z_1)\), \((z_0,z_2)\),
and the residual circle fixes it pointwise.  Subtracting the
auxiliary lift on these two supports gives normal weights
\((-1,+1)\) in each case.  Thus the circle orbit has the
ordinary rank-three Hopf normal model along its fixed locus,
with no additional finite-quotient strata.
Its positive moment levels are the negative-diagonal small resolution,
and its negative moment levels are the positive-diagonal small
resolution.  Here are the actual maps proving these identifications.
On the positive level use the complex quotient
\[
   \{a\ne0\}\big/\bigl((a,c)\sim(ta,t^{-1}c)\bigr)
       \longrightarrow
   \{([a],ac^t):\operatorname{im}(ac^t)\subset\mathbb Ca\}.
\]
Its target is the total space of \(\mathcal O(-1)^{\oplus2}\)
on \(\mathbb P^1_a\).  Every complex orbit meets a prescribed
positive moment level in exactly one auxiliary-circle orbit:
the required positive real factor solves
\[
 r^2\|a\|^2-r^{-2}\|c\|^2=\tau,
 \qquad
 r^2=\frac{\tau+\sqrt{\tau^2+4\|a\|^2\|c\|^2}}
                 {2\|a\|^2}.
\]
This is smooth for \(a\ne0\) and \(\tau>0\), including
\(c=0\).  The negative-level map is the same formula with
the two pairs interchanged.  Both maps commute with delta.
The positive stability condition excludes \(z_0=z_4=0\),
so its two fan triangles have common edge \(u_1u_2\);
the negative condition has common edge \(u_0u_4\).
Thus these are the specified two fans with their original
coordinate identifications on the common rank-one matrix locus.

\item Restrict to \(|\mu|\le\epsilon\) and a large fixed radius,
and remove a ball of radius \(\rho\) with \(\rho^2<\epsilon\).
The two moment faces are the truncated small resolutions, the
outer face is their common exterior times an interval, and the
new inner face is
\[
             S^7/S^1_{\rm aux}=\mathbb {CP}^3.
\]
This is an actual real bordism relative to the common exterior.
Conjugating the two auxiliary-negative coordinates makes the
auxiliary action scalar-positive; two complex conjugations
preserve the real orientation.  In the resulting homogeneous
coordinates, put
\[
       A=(z_0,\overline z_2),\qquad
       B=(z_4,\overline z_1).
\]
Delta acts by \((A,B)\mapsto(e^{-it}A,B)\).

\par\medskip\noindent\textbf{The circle orbit, its line, and the explicit joint cap.}\enspace

\item The orbit map on this inner face is
\[
 [A,B]\longmapsto
 \left([A],[B],\frac{\|A\|^2}{\|A\|^2+\|B\|^2}\right),
\]
with the irrelevant factor collapsed at either endpoint.
Thus the actual orbit is \(\mathbb {CP}^1*\mathbb {CP}^1\),
identified with \(S^5\subset\mathbb R^3\oplus\mathbb R^3\)
by \((p,q,\phi)\mapsto(\cos\phi\,p,\sin\phi\,q)\).
The two fixed spheres are its two join ends.  On the free orbit
\(S^2_A\times S^2_B\times(0,1)\), the circle line is exactly
\[
           L=\mathcal O(1)_A\boxtimes\mathcal O(-1)_B,
           \qquad x=h_A-h_B.                         \tag{P3K.17.1}
\]
Indeed a nonzero pair determines the relative tensor of its
two tautological lines; delta acts with negative weight on
that tensor.  Taking the associated line with the physical
positive delta character gives the displayed inverse tensor.
This fixes the sign rather than choosing it from the two
normal models independently.

\item The two base charges are therefore \(+1,-1\).  In particular
the line does not extend over either individual conifold cap:
the restriction image of
\(H^2(D^3_A\times S^2_B)\) in the first boundary has no
\(h_A\) term, and the other cap has no \(h_B\) term.
We retain this obstruction and instead join the fixed spheres.

\item Take \(C=S^2_P\times I\), with meridian sphere \(S^2_Q\),
and the line \(\mathcal O(1)_P\boxtimes\mathcal O(1)_Q\)
on \(C\times S^2_Q\).  At its first end use
\[
              (A,B)=(P,\iota Q),
\]
and at its second end use
\[
              (A,B)=(Q,\iota P),                    \tag{P3K.17.2}
\]
where \(\iota\) is the antipodal map of \(\mathbb {CP}^1\).
Both maps pull (P3K.17.1) to \(p+q\).  These are actual
line isomorphisms: the orthogonal-complement map on lines in
\(\mathbb C^2\), followed by its determinant pairing,
identifies \(\iota^*\mathcal O(1)\) complex-linearly with
\(\mathcal O(-1)\); the second end then uses the canonical
interchange of the two tensor factors.  All maps preserve
the physical principal-circle character.

\item The resulting closed five-dimensional orbit correction is the
mapping torus
\[
 Q_{\rm corr}=T_g(S^2\times S^2),\qquad
           g(A,B)=(\iota B,\iota A).
\]
The map \(g\) preserves \(h_A-h_B\) and the product
orientation.  Its line is glued by the specified tensor
isomorphisms, so this statement retains the actual line,
not merely an invariant cohomology class.

\par\medskip\noindent\textbf{The actual parameter gauge on the master space.}\enspace

\item For the closed OS49 parameter \(\lambda=r^m\), with
\(m=3k\) and signed primitive twist \(n=\pm1\), the
extra twist \(d\delta/m\) contributes only a parameter-independent
scalar to the finite-cap line map.  Its source coordinate is unchanged;
the same scalar extends in the Cox lift and across the joint cap.
It therefore introduces no parameter-crossing line class.  The
period-origin calculation in P3F.32 gives the primitive
\(K\)-turn and the angular line winding \(2mn\).
The \(K\)-period on the ordinary Tate coordinates is
\[
      Z\longmapsto a(q)Z,\qquad
      V\longmapsto q^n r^{2mn}b(q)V,
\]
after incorporating the signed traversal (so the bounded units
include the corresponding signed powers), where those units
have parameter-independent logarithms.  Reindexing
the common fan by its integral vertical period removes the
displayed \(q^n\) by the signed integral vertical translation
\(n\).  On the resulting common Cox charts the
bounded map has the explicit lift
\[
 \begin{aligned}
 z_0&\longmapsto\bigl(a(q)b(q)r^{2mn}\bigr)^{-1}z_0,
       &z_1&\longmapsto a(q)z_1,\\
 z_2&\longmapsto b(q)r^{2mn}z_2,
       &z_4&\longmapsto z_4.
 \end{aligned}                                      \tag{P3K.17.3}
\]
The product of its four multipliers is one, so it fixes the
base monomial \(q=z_0z_1z_2z_4\).  Its torus vector is
\(\log a\,(u_1-u_0)+\log(br^{2mn})\,(u_2-u_0)\),
which is exactly the marked \((Z,V)\) action.  Thus its
parameter-dependent factor is one global physical delta
action; it is not two separately chosen gauges at the two
ends of the exceptional geometry.

\item For the exact quotient, label the unit-square flop cells by
\(v\in\mathbb Z^2\).  Their period lattice is
\(\Gamma=\langle(2,1),(-3,-1)\rangle=\mathbb Z^2\),
and \(3(2,1)+2(-3,-1)=(0,1)=K_{\rm trop}\).
It acts freely and transitively on these cells.  Their diagonal
curves are disjoint, since each has the two maximal cones
belonging to its own square.  Choose disjoint equivariant
tubular neighbourhoods and insert one master chart in each.
The quotient of their disjoint union by \(\Gamma\) is
exactly one representative chart with trivial stabilizer.
There is no residual \(K\)-circle on this correction.
The signed parameter return shifts its cell label by
\(n(0,1)=n(3h+2g)\).  Returning to the chosen representative
by that unique period gives (P3K.17.3).  Reversing the
representative convention inverts the one common gauge at
both ends and leaves the comparison unchanged.

\item Fixing \(|r|\), its angular part is
\(e^{2mn i\theta}\) in that physical circle.  The remaining
unit multipliers are independent of \(\theta\); their
chosen logarithms give equivariant isotopies on the Cox
charts.  The positive real rescalings in Section 1 transport
these isotopies to the moment levels and to the inner face.
Explicitly, follow a Cox map by its unique positive auxiliary
rescaling back to the original value of \(\mu\).  Its
implicit derivative is
\(2r^2\|a\|^2+2r^{-2}\|c\|^2>0\) away from the
removed origin, so this is smooth through the real bordism.
The bounded units can be changed to their values at \(q=0\)
using \(a(tq),b(tq)\), \(0\le t\le1\), with an
exterior cutoff.  The maps stay invertible and commute with
both circles; they are independent of the parameter angle.
This change introduces no angular character on the inner face.
These maps need not preserve a fixed Euclidean radius.
Use transported truncations: under the moment-preserving
equivariant isotopy \(G_t\), take the inner and outer
faces to be \(G_t\) of the original faces.  They stay
embedded and disjoint, while the moment faces stay at
\(\mu=\pm\epsilon\).  Glue these transported charts
and collars by their exact maps.  For the remaining
constant-unit clutch \(G_0\), its chosen logarithms give
\(G_0^t\), and
\((x,t)\mapsto(G_0^t x,e^{2\pi it})\) trivializes
the mapping torus.  The angular delta factor preserves
the round face and is trivial on its circle orbit.
Consequently the two maps (P3K.17.2) have the same angular
gauge, which extends by that constant circle multiplication
over all of \(C\times S^2_Q\).  No difference of two
angular gauges is introduced on its ends.  The closed
correction over the parameter circle therefore carries
a line pulled back from \(Q_{\rm corr}\), after this
explicit gauge identification.  Its first Pontryagin class
is pulled back from the same five-dimensional factor.
The potential mixed line class can also be checked directly.
Cover the correction by its free-orbit cylinder and joint-cap
cylinder; their overlap has two components.  The actual
transition functions there are
\(\varphi_1(\theta)=g(\theta)\varphi_{1,0}\) and
\(\varphi_2(\theta)=g(\theta)\varphi_{2,0}\), with
the same scalar from (P3K.17.3).  Change the cap frame by
\(g(\theta)^{-1}\).  Both angular factors disappear
simultaneously, leaving parameter-independent transitions.
This explicitly excludes a line class on
\(S^1_{\rm joint}\times S^1_{\rm parameter}\).
Both six-dimensional characteristic numbers vanish:
\[
  \langle x^3,[Q_{\rm corr}\times S^1]\rangle=0,
  \qquad
  \langle p_1x,[Q_{\rm corr}\times S^1]\rangle=0.       \tag{P3K.17.4}
\]

\par\medskip\noindent\textbf{The additional normal-sphere boundary in the surgery bordism.}\enspace

\item The fixed locus in the master bordism must also be retained.
There are exactly two relevant auxiliary-quotient fixed charts:
\((z_4,z_1)\) and \((z_0,z_2)\).  Each has auxiliary
weights \((1,-1)\), and its quotient is \(\mathbb R^3\)
by the explicit invariant map
\[
      (z_+,z_-)\longmapsto
       \bigl(2\Re(z_+z_-),2\Im(z_+z_-),
                         |z_+|^2-|z_-|^2\bigr).
\]
The inverse orbit is determined by
\(|z_\pm|^2=(\sqrt{\mu^2+4|z_+z_-|^2}\pm\mu)/2\).
The two moment faces in each chart are discs.  Removing the
inner master-space ball removes an interior three-ball in
each of these fixed charts.  Away from the flop curve the
fixed sphere minus its two endpoint discs is unchanged.
Gluing this annulus times the moment interval proves that
the fixed three-dimensional cobordism is exactly
\[
                 W_F=(S^2\times I)\setminus
                             (\mathring D^3\sqcup\mathring D^3).
\]
The Cox reindexing preserves the two charts; it does not
permute or identify their boundary spheres.

\item The oriented normal three-plane bundle is trivial over this
three-dimensional cobordism, with its boundary product
structures retained.  On the two Cox fixed charts its
positive-weight complex normal bundle is respectively
\(L^{\oplus2}\) and \((L^{-1})^{\oplus2}\), so the
associated three-plane bundle has \(w_2=0\).  The other
two boundary spheres have the original \(SU(2)\) normal
model.  Restriction of \(H^2(W_F;\mathbb Z/2)\) to
the four boundary spheres is injective, so \(w_2\) vanishes
globally.  An oriented three-plane bundle on a three-dimensional
complex with this class zero is trivial: the obstruction on
two-cells is \(\pi_1(SO(3))=\mathbb Z/2\), and that
on three-cells is \(\pi_2(SO(3))=0\).  Boundary
trivializations differ by maps \(S^2\to SO(3)\), all
null-homotopic.  Collar homotopies therefore match the chosen
product trivializations at every boundary.  The actual circle
line restricted to the constant normal section has boundary
values \(0,0,+h_A,-h_B\).

\item Cap the two original-family sphere boundaries by \(D^3\),
as in their actual conifold construction, and join the two
inner-face spheres by the charged cap \(C\) above.  This
closes \(W_F\) to \(S^2\times S^1\).  Crossing with
the parameter circle closes the fixed four-dimensional base
to \(\bar X\).  The described local trivialization makes
it the product with that parameter circle, retaining the
possible parameter-independent clutch along the first
\(S^1\).  In particular it is a mapping torus over the
first circle of \(S^2\), crossed with the parameter circle.
All gluings preserve the sphere orientation.  Its signature
is zero, and its tangent first Pontryagin number is zero.
For the exact orientation comparison, capping the first two
boundaries leaves \(S^2\times I\).  Its line has degrees
\(+1\) and \(-1\) in the two complex boundary markings.
Its product transport consequently has degree \(-1\)
between those markings.  The joint cap has exactly that
antipodal base change at its second end.  The composite has
degree \(+1\), giving the oriented sphere mapping torus
above.  On the meridian sphere the first end map is also
antipodal, so both meridian orientations give the same
physical Euler value one.


\item The preceding construction is the inner master-space and charged-cap
comparison. Its actual outer attachment in the OS49 family is the
component-line map, with its unequal two-node defects, carried in
\cref{night:P3K.65}. That proof gives native diagonal pair
((-k,k)) for the shorter negative-to-positive comparison. The
literal Cox conjugacy adds its right integer (k), so starting
from the proved positive return (0) gives the negative return
(-k). The map (w=r^3) pulls the shorter family back to the
whole family and composes three actual positive returns, giving
(-3k). At (k=1) the two classes are (25) and (0).
Thus the claimed equality of the complete returns is refuted,
while all the local master-space and joint-cap arrows above remain
available with their displayed source and target.

\end{enumerate}
\end{proof}


The 26 pairs are a finite necessary source reduction in
the declared orbit and cap scope.  The full cap data still
include their exact cyclic norm, logarithmic site, integral
delta collar, Seifert determinant, and integral topology.
No unverified pair in this table is assigned a sphere or
a return index solely by this arithmetic list.

\begin{theorem}\label{night:P3K.22}
\StatementRecord{night:P3K.22}{false as stated; correction proved}{S6\_06\_BLIND: GPT-6 Astra (OpenAI, ChatGPT Pro mode, effort Pro), recorded re-derivation in the checkpoint dated 18 November 2026}{\cref{night:P3F.1}, \cref{night:P3F.2}, \cref{night:P3F.25}, \cref{night:P3F.9}, \cref{night:P3K.10}, \cref{night:P3K.13}, \cref{night:P3K.15}, \cref{night:P3K.16}, \cref{night:P3K.6}, \cref{night:P3K.7}}For the following ten source profiles, let \(G\) be the marked
primitive Mordell--Weil generator, put \(d=1/h(G)\), and choose
\(M=\mathcal O(dG-O)\), translation \(Q=G\), a simple whole-fibre
zero at an \(I_1\), and one order-one logarithmic gluing with gamma
coefficient \(+1\) at a smooth source fibre.  Use the standard
vertical-edge unimodular cusp fan.  These are compact complex
threefolds diffeomorphic to \(S^6\).  Their positive scalar-parameter
loops have the indicated Jin integer and oriented \(\Theta_7\) class.


\begin{center}\small\begin{tabular}{c p{0.48\textwidth} c c c}\hline
OS row & Source profile & d & Jin integer A & Class mod28 \\
\hline
45 & \(I_8+4I_1\), \(I_8+II+2I_1\) & 8 & \(3\) & 3 \\
47 & \(I_7+I_2+3I_1\), \(I_7+I_2+II+I_1\), \(I_7+III+2I_1\) & 14 & \(6\) & 6 \\
55 & \(I_5+I_4+3I_1\), \(I_5+I_4+II+I_1\) & 20 & \(9\) & 9 \\
56 & \(I_5+I_3+I_2+2I_1\), \(I_5+I_3+III+I_1\), \(I_5+IV+I_2+I_1\) & 30 & \(14\) & 14 \\
\hline\end{tabular}\end{center}

The compact page is the indicated rational elliptic surface itself.
Its line is \(M\), with square \(-d\), and
\((dG)\cdot O=(d-2)/2\).  Every displayed mapping torus is therefore
not diffeomorphic to \(S^1\times S^6\), by the same oriented
mapping-torus detection used in P3K.10 and P3F.25.

There is also a restriction on additional smooth-site caps for these
primitive sources.  After the actual finite-kernel cancellation of
P3K.19, let their primitive cap data be
\((m_i,n_i,\lambda_{ai},\lambda_{bi})\), with
\(\gcd(m_i,n_i)=1\).  The exact orbit Mayer--Vietoris computation
below gives, whenever its first-homology gluing matrix is unimodular,
\[
        H_2(Y;\mathbb Z)=\bigoplus_i(\mathbb Z/m_i)^2.       \tag{P3K.22.1}
\]
Consequently a sphere member with only such smooth-site modifications
has \(m_i=1\) at every site.  For order-one sites the total gamma
coefficient \(\sum n_i=\pm1\) gives sphere topology, including
arbitrary integral source and delta collar coefficients.  The
individual index table above specifies the one-site positive-gamma
family; it does not identify uncomputed singular-source cap choices
with this family.
\end{theorem}

\begin{proof}
\begin{enumerate}[label=\arabic*.]
\item We retain the complete source, line, cap and frame maps below.
The meridian-one graph line is (U=M^{-1}): the fibre map
(D_M(v)(w)=\langle v,w\rangle/\langle v,v\rangle), with the
Hermitian form linear in its second argument, satisfies
(D_M(e^{i\theta}v)=e^{-i\theta}D_M(v)).
Its (SU(2))-equivariant spinor extension is
((a,b)\mapsto(\overline b,-\overline a)).
The first-to-second cap maps in \cref{night:P3K.68} give
(\lambda_M(F)=-\Sigma_{F^{-1}}) and
(\lambda_U(F)=\Sigma_{F^{-1}}); cancelling the opposite
product collars gives (A_M=-A_U). We therefore distinguish the
intermediate meridian-one number (A_U=(4k_U-\ell_U)/24)
from the physical return (A_M) throughout this calculation.

\item \par\medskip\noindent\textbf{The source automorphisms and the line lift.}\enspace
P3F.1 constructs the source surfaces, and P3F.2 fixes their component
markings and proves \(G\cdot O=0\).  Translation and inversion
extend over their resolved singular fibres.  One direct proof keeps
the actual rational maps.  On the generic Weierstrass fibre use its
chord-and-tangent addition and \((x,y)\mapsto(x,-y)\).  Both preserve
the invariant differential up to the constant sign of inversion.
On a rational minimal elliptic surface the form
\((dx/y)\wedge dt\) has canonical divisor \(-F_\infty\), after a
base frame is chosen; the local minimal resolutions preserve this
relative differential.  Resolve the graph of either rational map,
with projections \(p,q:Z\to S\).  They are maps over the same base.
Pulling back that differential gives
\[
 -p^*F_\infty+E_p=-q^*F_\infty+E_q,
 \qquad p^*F_\infty=q^*F_\infty.
\]
Here the Jacobian discrepancy divisors have strictly positive
coefficient on every exceptional curve of their respective
compositions of point blowups.  Thus \(E_p=E_q\), so both
projections contract exactly the same curves and connected fibres.
The graph maps descend in both directions, giving mutually inverse
holomorphic automorphisms of \(S\).  This proves the required
extension, including II, III and IV fibres.  On an \(I_N\) chain
the actual Tate map is multiplication by the section unit and a
shift by its component label.  On III and IV, translation acts on
the dense smooth parts by addition in the component groups of
orders two and three; it therefore permutes the resolved components
by those same labels.  Inversion conjugates translation by \(G\)
to translation by \(-G\) on the whole surface.

\item The point \(dG\) is in the identity component at every reducible
fibre in the table.  Thus \(M\) has multidegree zero there.
The height identity gives
\((dG)\cdot O=(d-2)/2\), and its actual divisor class satisfies
\[
 (t_G)_*[M]-[M]=[F],\qquad t_G^*M\cong M(-F).           \tag{P3K.22.2}
\]
These equalities are also checked directly in the integer matrices
below.  The numerical equality is an equality of line bundles
because the rational surface has \(\operatorname{Pic}^0(S)=0\).
The resulting Hom line is \(\pi^*\mathcal O_{\mathbb P^1}(1)\).
Choose its section with its unique zero at the selected nodal fibre.
This gives the actual holomorphic lift with a simple whole-fibre
zero.

\item For an explicit meromorphic frame of that lift, put \(P_j=jG\).
On the generic cubic the quotient of the chord through
\(P_d,-G\) by the vertical line through \(P_{d-1}\) has divisor
\(P_d+(-G)-P_{d-1}-O\).  It is exactly the fibre divisor of
\(M\otimes(t_G^*M)^{-1}\).  Its remaining vertical orders form
whole fibres, since the component degrees vanish.  Multiplying by
the rational base function that moves this degree-one divisor to the
chosen \(I_1\) supplies the preceding section.  Thus its frame and
all local multipliers are obtained from the stated rational addition
map, rather than from an unspecified linearization.

\item For sufficiently small \(|\lambda|\), multiply this lift by
\(\lambda\).  In a fixed Hermitian metric it is uniformly strictly
contracting, away from its prescribed zero where the Mumford charts
are used.  Quotienting the nonzero line by its iterates is therefore
free and proper.  At the zero the actual valuation periods are
\((1,0),(0,1)\); the standard vertical-edge fan gives the smooth
index-one filling with its one fixed \(\mathbb P^1\), of normal
weights \((-1,+1)\).  The chart and line-gluing formulas are the
same ones used in P3K.10, with the present holomorphic unit logarithms
\(\Gamma,C\).  In particular, on a universal scalar-parameter
chart its cusp period and two origins are
\[
 \begin{aligned}
 \gamma&=(\Gamma,C+s+v),\\
 \sigma&=(\Gamma v,Cv+(v^2-v)/2),\\
 o&=(-\Gamma s,-Cs-s^2/2-sv),
 \end{aligned}                                                     \tag{P3K.22.3}
\]
and the closing ordinary line gauge is \(-s\delta\).  The exact
identities
\[
 \sigma(v+1)-\sigma(v)=\gamma-s\delta,
 \quad (\sigma-o)(v+\theta)-(\sigma-o)(v)
      =\theta\gamma+\tfrac12\theta(\theta-1)\delta
\]
hold for these actual units.  The degree-two cusp root used in
P3K.10 retains the half-integral ordinary-origin term and its integral
deck.  The smooth order-one cap is glued by
\((\log q)/(2\pi i)\,\gamma\); its scalar closing gauge is the
same actual base function used above.  Thus the gluing is holomorphic
in the scalar parameter and closes on its positive loop.  All filled
pieces are proper over their base discs; the contraction shell gives
properness on their common complement.  The compact base gives
compactness.  This establishes the analytic construction and the
positive-gamma parameter direction needed for the index calculation.

\par\medskip\noindent\textbf{Every integer matrix of the source mapping torus.}\enspace
Use the divisor basis
\[
 O,F,\{C_{i,j}:1\le j<N_i\},G,
\]
where \((N_i,k_i)\) is respectively
\((8,3)\), \((7,2),(2,1)\), \((5,2),(4,1)\), or
\((5,1),(3,1),(2,1)\).  A root \(A_1\) may denote III and
a root \(A_2\) may denote IV; their displayed component intersection
and permutation matrices are unchanged.  The Gram matrix has
\(O^2=G^2=-1\), \(O\cdot F=G\cdot F=1\),
\(G\cdot C_{i,j}=\delta_{j,k_i}\), root blocks
\(-A_{N_i-1}\), and all other entries zero.  Its determinant is
minus one in every case.  Since \(H_2(S;\mathbb Z)\) is free
unimodular of rank ten, these classes are an integral basis.

\item Let \(R_j\) be the vector whose \(i\)-th root block is column
\(jk_i\bmod N_i\) of \(A_{N_i-1}^{-1}\), or zero for component
zero.  Put
\[
 \varphi=G-O-F+R_1,\quad
 P_j=O+(P_j\cdot O+1)F+j\varphi-R_j,
 \quad
 P_j\cdot O=\frac{j^2/d-2+\sum_i\operatorname{contr}_i(jG)}2.
\]
The actual pushforward matrix \(T=(t_G)_*\) is now specified
column by column: \(O\mapsto G\), \(F\mapsto F\),
\(G\mapsto P_2\), and
\(C_{i,j}\mapsto C_{i,j+k_i}\), with
\(C_{i,0}=F-\sum_{j>0}C_{i,j}\).  All columns are integral.
The executed program verifies \(T^t\mathcal G T=\mathcal G\) and
\[
 \operatorname{Smith}(T-I)=(1,1,1,1,1,1,1,1,1,0),
 \quad\ker(T-I)=\mathbb ZF,
 \quad\operatorname{coker}(T-I)=\mathbb Z[O].             \tag{P3K.22.4}
\]
It exports every ten-by-ten matrix in
\path{code/part3_fourfold/primitive_source_translation.json}.

\item There is also a direct primitive-presentation check.  In the
coinvariants the components of each fibre become one class \(c_i\),
with \(N_i c_i=F\), and \(G=O\).  The remaining section relation
has coefficient \(P_2\cdot O-1\) on \(F\) and coefficient
\(\sum_j(2R_1-R_2)_{i,j}\) on \(c_i\).  Its Schur complement is
\(1/d\), while \(\prod_iN_i=d\).  The resulting determinant is
therefore \(\pm1\), not merely nonzero rationally.  For the four
rows the determinants are respectively \(-1,1,1,-1\).

\item The actual contracting lift covers \(t_{-G}\).  Its source matrix
has the same image and kernel, since
\(T^{-1}-I=-T^{-1}(T-I)\).  The mapping torus
\(W=T_{-G}(S)\) has the cellular cone complex with groups
\(\mathbb Z,\mathbb Z,\mathbb Z^{10},\mathbb Z^{10},
\mathbb Z,\mathbb Z\) in degrees zero through five; the sole
nonzero differential is \(T^{-1}-I:C_3\to C_2\).
Thus each \(H_i(W)\), \(0\le i\le5\), is \(\mathbb Z\),
and \(\pi_1(W)=\mathbb Z\gamma\).

\item The actual delta quotient before logarithmic gluing is this marked
mapping torus.  On its ordinary part the Hermitian logarithmic norm
is the suspension coordinate, as in P3K.13.  At the chosen \(I_1\),
the section \(G\) has a nonzero holomorphic Tate-unit representative
\(g(q)\).  Choose its logarithm on the disc.  The isotopy
\(Z\mapsto g(q)^u Z\) has node-chart formula
\((x,y)\mapsto(g(q)^u x,g(q)^{-u}y)\), so it extends through the
node for every real \(u\).  It identifies the source mapping torus
with the same \(I_1\)-surface-times-circle product of P3F.9, on
the actual boundary collar.  The bounded line units alter only the
logarithmic-norm coordinate and give the same product map.  This
supplies the comparison across the cusp, including its source
translation, rather than replacing that piece by a different fill.

\par\medskip\noindent\textbf{The primitive rim bases and the decisive involution.}\enspace
Remove \(r\) disjoint regular-fibre tubes
\(N_i=T^2\times D^2\times S^1_\gamma\) from \(W\), and call
the complement \(V\).  Write their boundary bases
\((a_i,b_i,\alpha_i,\gamma)\), with \(\alpha_i\) the positive
base meridian.  The Thom maps
\(H_j(W)\to\bigoplus_iH_{j-2}(T^3)\) are, in descending degrees,
\[
 \begin{array}{c|c}
 j&\text{map}\\\hline
 5&1\longmapsto(1,\ldots,1)\text{ in the }ab\gamma\text{ columns}\\
 4&1\longmapsto(ab)_1+\cdots+(ab)_r\\
 3&0\\
 2&1\longmapsto(1,\ldots,1).
 \end{array}                                                     \tag{P3K.22.5}
\]
For the zero map, the dual inclusion \(H_2(N_i)\to H_2(W)\)
is zero: the invariant class \(F\), which pairs one with \([O]\),
is pulled back from the source base and restricts to zero on the
whole tube core.  The other maps follow from the oriented fibre and
section intersections.  They are primitive diagonal maps.

\item The Thom sequence now gives free groups
\[
 H_1(V)=\mathbb Z^r,\quad H_2(V)=\mathbb Z^{3r},
 \quad H_3(V)=\mathbb Z^{3r},\quad H_4(V)=\mathbb Z^{r-1}.
\]
In degree two an integral basis is exactly
\(A_i=a_i\alpha_i,B_i=b_i\alpha_i,C_i=\alpha_i\gamma\).
In degree one use \(\gamma,\alpha_1,\ldots,\alpha_{r-1}\), with
\(\alpha_r=-\sum_{i<r}\alpha_i\).  In degree three use
\(z=ab\gamma\), the \(r-1\) independent classes \(ab\alpha_i\),
and all \(a_i\alpha_i\gamma,b_i\alpha_i\gamma\).  In degree four
use \(r-1\) boundary fundamental classes with their one sum relation.

\item The remaining boundary coefficients are fixed by an actual smooth
involution.  Inversion on the source conjugates the suspension map
to its inverse, so
\[
                         I(p,u)=(-p,-u)                         \tag{P3K.22.6}
\]
is a well-defined diffeomorphism of its mapping torus and preserves
every removed tube.  In a tube choose a real lift of the local
section and write its mapping-torus trivialization as
\(p'=p+uQ\), where \(Q\) is its signed suspension translation.
Then (P3K.22.6) is literally \((p',\alpha,u)\mapsto(-p',\alpha,-u)\).
Thus \((a,b,\alpha,\gamma)\mapsto(-a,-b,\alpha,-\gamma)\).
Every degree-two rim basis vector is odd, whereas
\(ab,a\gamma,b\gamma\) are even.  Naturality forces each image
of an even class to be two-torsion in \(H_2(V)\).  This group is
free by the primitive Thom calculation, so every such image is zero.
In particular the previously undetermined boundary linking integers
are exactly zero integrally.

\item The same argument gives the full degree-three boundary map.
The three rim classes \(ab\alpha,a\alpha\gamma,b\alpha\gamma\)
are even.  Each \(ab\gamma\) is odd and projects to the same
generator of \(H_3(W)\).  Subtracting two of them gives both an
odd class and a class in the even rim subgroup; freeness again
forces equality.  This proves the basis description and every
boundary coefficient used next.

\par\medskip\noindent\textbf{The actual cap maps and all mayer--vietoris matrices.}\enspace
On the orbit boundary a primitive smooth cap has meridian
\[
 \mu_i=m_i\alpha_i+n_i\gamma+\lambda_{ai}a_i+\lambda_{bi}b_i.
\]
The coefficient vector is the full cyclic norm from P3K.15; no
source-period term is discarded.  The smooth root source, followed
by its extending translation and base rotation, is a mapping torus
of a disc-times-torus diffeomorphism isotopic to the identity by
those explicit rotations and translations.  It is therefore
\(T^2\times D^2\times S^1\), with the displayed marked meridian.
Choose integers \(r_i,s_i\) with \(r_im_i+s_in_i=1\).
The four columns
\[
 a_i,\ b_i,\ \mu_i,\ \kappa_i=-s_i\alpha_i+r_i\gamma
\]
form an integral determinant-one basis.  Let \(Q_i\) be the
matrix of its inverse with the meridian row deleted.  Then
\(Q_i:\mathbb Z^4\to\mathbb Z^3\) is the actual cap inclusion
on first homology, and its higher maps are exactly
\(\bigwedge^jQ_i\).

\item For the exterior, in boundary wedge order
\((ab,a\alpha,a\gamma,b\alpha,b\gamma,\alpha\gamma)\), the
degree-two block is
\[
 i_2=\begin{pmatrix}
 0&1&0&0&0&0\\0&0&0&1&0&0\\0&0&0&0&0&1
 \end{pmatrix}.
\]
The degree-one map sends \(a_i,b_i\) to zero, \(\gamma\) to
the common gamma generator, and \(\alpha_i\) to the meridians with
their stated one sum relation.  The degree-three map in wedge order
\((ab\alpha,ab\gamma,a\alpha\gamma,b\alpha\gamma)\) sends
these columns to the corresponding rim class, the common \(z\),
and the two displayed independent rim classes.  The degree-four
map is \((I_{r-1}\mid-\mathbf1)\).  These descriptions specify
every entry of the integral matrices \(i_j\).

\item The Mayer--Vietoris map in degree \(j\) is the block matrix
\(D_j=\begin{pmatrix}i_j\\-\bigoplus_i\bigwedge^jQ_i\end{pmatrix}\); in degree zero it
is the column-incidence matrix \((\mathbf1,-I_r)^t\).
For one positive order-one cap with zero source collar coefficients,
the three nontrivial square matrices are explicitly
\[
 D_1=\begin{pmatrix}
 0&0&0&1\\-1&0&0&0\\0&-1&0&0\\0&0&1&-1
 \end{pmatrix},\qquad
 D_3=\begin{pmatrix}
 0&1&0&0\\0&0&1&0\\0&0&0&1\\1&-1&0&0
 \end{pmatrix},
\]
\[
 D_2=\begin{pmatrix}
 0&1&0&0&0&0\\0&0&0&1&0&0\\0&0&0&0&0&1\\
 -1&0&0&0&0&0\\0&1&-1&0&0&0\\0&0&0&1&-1&0
 \end{pmatrix}.
\]
All three are unimodular.  The degree-four domain is \(\mathbb Z\)
with zero target, giving the oriented top class of \(Y\).  Thus
\(H_*(Y;\mathbb Z)=H_*(S^5;\mathbb Z)\).

\item For general smooth-site data, eliminating the cap degree-two groups
leaves the exact block presentation
\[
 \begin{pmatrix}
 m_i&0&0&\lambda_{ai}\\
 0&m_i&0&\lambda_{bi}\\
 0&0&m_i&-n_i
 \end{pmatrix}                                                   \tag{P3K.22.7}
\]
on \((A_i,B_i,C_i)\).  Its cokernel is
\((\mathbb Z/m_i)^2\), since \(n_i\) is a unit modulo \(m_i\).
The first-homology determinant is the signed Seifert determinant
\(D=\sum_i n_i\prod_{j\ne i}m_j\), and the degree-three
determinant has the same absolute value.  These identities, every
matrix \(D_0,\ldots,D_4\), and 196 primitive two-cap cases with
nonzero source collars are executed and exported by
\path{code/part3_fourfold/primitive_orbit_filling.py} and its JSON.
This proves (P3K.22.1), including its torsion.

\item For order-one sites the fundamental-group calculation is equally
explicit.  The complement of the source fibres has fundamental
group generated by their base meridians: all elliptic cycles die
as in the rational source, and the section supplies the only sum
relation.  Thus
\(\pi_1(V)=\langle\alpha_i,\gamma\mid
\alpha_1\cdots\alpha_r=1,[\alpha_i,\gamma]=1\rangle\).
The cap imposes \(\alpha_i\gamma^{n_i}=1\); its source loops
already die.  Hence \(\pi_1(Y)=\mathbb Z/|\sum_i n_i|\), which
is zero for the stated sphere data.  The source meridian assertion
can also be checked directly: the missing meridian of a single
regular fibre bounds the punctured zero section, and van Kampen
recovers \(\pi_1(S)=0\).

\par\medskip\noindent\textbf{The actual hopf filling and its six-dimensional homology.}\enspace
In the one-site construction \(Y\) is simply connected with the
integral homology of \(S^5\), hence is the standard oriented
five-sphere by the same smooth classification used for P3F.9.
The fixed curve is the actual embedded sphere furnished by the
index-one cusp chart.  Its normal Hopf meridian has Euler value one.
For \(P=Y\setminus\operatorname{int}\nu F\), Alexander duality
gives \(H^*(P;\mathbb Z)=H^*(S^2;\mathbb Z)\); the circle
bundle on \(P\) therefore has \(H_*(E)=H_*(S^3)\) by the
integral Gysin sequence.  The meridian value one makes
\(\pi_2(P)\to\pi_1(S^1)\) surjective, so \(\pi_1(E)=0\).

\item The other piece of \(X\) is \(F\times D^4\), with boundary
\(F\times S^3\).  The positive-weight SU(2) normal trivialization
of the actual \(\mathcal O(-1)\oplus\mathcal O(-1)\) chart gives
this product and its Hopf map.  Its Mayer--Vietoris matrices are
\((1,-1)^t\) in degree zero, \((-1)\) in degree two,
\((1)\) in degree three, and the zero map from \(\mathbb Z\)
in degree five.  All omitted groups are zero.  Therefore
\(H_*(X;\mathbb Z)=H_*(S^6;\mathbb Z)\) and van Kampen gives
\(\pi_1(X)=0\).  The already used smooth six-sphere classification
and \(\Theta_6=0\) identify this actual smooth manifold with
\(S^6\).  The same calculation excludes the nonzero torsion in
(P3K.22.1) from a sphere member; the quotient of a semifree sphere
with this fixed sphere must have zero such torsion.

\par\medskip\noindent\textbf{The actual page and the signed return.}\enspace
The order-one gamma gluing makes the orbit base the ordinary Hopf
\(S^3\), by the signed meridian matrix of P3K.16.  Its regular
binding page has degree one over the source base.  The norm shell
of \(M^\times\) maps injectively into the orbit: a positive iterate
strictly decreases its norm, and a negative iterate increases it.
At the whole-zero cusp the line graph is horizontal and lies in
the chosen standard fan; its successive suspension translates have
strictly increasing graph defect one, so no new identification of
this shell occurs there.  It is disjoint from the delta fixed
sphere.  On the smooth logarithmic collar the displayed log-gamma
map is the exact inverse angular coordinate, so the shell closes
as the page \(S\setminus\operatorname{int}\nu F\), with its
original source line \(M\).  This is the same actual norm-shell
and binding-collar comparison proved for OS45 in P3K.10; all its
maps here use the present \(\Gamma,C\), and no finite cap or
component-supported zero is introduced.

\item Thus the capped page is \(S\), of \(\chi(\mathcal O)=1\) and
signature \(-8\), while its line square is
\(M^2=-2-2(dG\cdot O)=-d\).  The real compression, meridian-line
comparison and signed characteristic calculation of P3K.6--7 apply
to this actual page and give
\[
 A_U=\frac{12M^2-p_1(S)}{24}
   =1-\frac d2,\qquad A_M=\frac d2-1=\frac{d-2}{2}.                        \tag{P3K.22.8}
\]
The identities (P3K.22.3) identify one positive scalar-parameter
loop with one positive gamma turn, retaining the integral line
gauge and the degree-two cusp-origin root.  Hence (P3K.22.8)
computes the stated actual loop, with no unproved division by a
parameter cover.  The explicit line and cap dictionary of \cref{night:P3K.68}
reduces this integer to the classes in the table.  This proves all
the asserted source, topology, page and return claims.

\end{enumerate}
\end{proof}


The statement supplies ten explicit sphere families and an exact
obstruction to additional smooth finite caps in this primitive-source
scope.  Singular-source cap replacements and arbitrary cusp-fan return
changes retain their separate actual comparisons; they have not been
identified with these examples by the source-lattice calculation.

\begin{theorem}\label{night:P3K.23}
\StatementRecord{night:P3K.23}{false as stated; correction proved}{S6\_06\_BLIND: GPT-6 Astra (OpenAI, ChatGPT Pro mode, effort Pro), recorded re-derivation in the checkpoint dated 18 November 2026}{\cref{night:P3F.22}}Consider one of the actual sphere families constructed above with a
semifree fixed curve \(F\simeq\mathbb P^1\), whose complex normal bundle
is \(\mathcal O(-1)\oplus\mathcal O(-1)\). Perform the toric flop of
this fixed edge in the periodic cusp model, retaining the actual
parameter clutch and all exterior gluings. Write the original positive
return's physical index and meridian-one dual pair as
\[
             A_M=-A_U,\qquad A_U=\frac{4k_U-\ell_U}{24},
\]
The positive return of the flopped family has oriented class
\[
                         7A_M\pmod {28}                 \tag{K23.1}
\]
in the run's normalization of \(\Theta _7\). The flopped curve has two
isolated circle fixed points. Formula (K23.1) specifies its homotopy
sphere class; it does not assign a new Jin integer to that circle
action.

In particular, the four groups of primitive-source examples in
P3K.22 give the following flopped classes:


\begin{center}\small\begin{tabular}{c l l}\hline
\(d=1/h(G)\) & Original physical index \(A_M\) & Flopped class in \(\mathbb Z/28\) \\
\hline
\(8\) & \(3\) & \(21\) \\
\(14\) & \(6\) & \(14\) \\
\(20\) & \(9\) & \(7\) \\
\(30\) & \(14\) & \(14\) \\
\hline\end{tabular}\end{center}

The earlier degree-four OS47 example of physical index \(4\) gives class
\(0\). Every original zero-index example remains class zero after
this fixed-edge flop. No formula for a flop at another toric edge is
asserted here.
\end{theorem}

\begin{proof}
\begin{enumerate}[label=\arabic*.]
\item We retain the complete source, line, cap and frame maps below.
The meridian-one graph line is (U=M^{-1}): the fibre map
(D_M(v)(w)=\langle v,w\rangle/\langle v,v\rangle), with the
Hermitian form linear in its second argument, satisfies
(D_M(e^{i\theta}v)=e^{-i\theta}D_M(v)).
Its (SU(2))-equivariant spinor extension is
((a,b)\mapsto(\overline b,-\overline a)).
The first-to-second cap maps in \cref{night:P3K.68} give
(\lambda_M(F)=-\Sigma_{F^{-1}}) and
(\lambda_U(F)=\Sigma_{F^{-1}}); cancelling the opposite
product collars gives (A_M=-A_U). We therefore distinguish the
intermediate meridian-one number (A_U=(4k_U-\ell_U)/24)
from the physical return (A_M) throughout this calculation. In the intermediate Hopf-quotient calculation below, unadorned (k,\ell) denote (k_U,\ell_U).

\item \par\medskip\noindent\textbf{The actual two local products and the parameter clutch.}\enspace
Write the conifold link as
\[
 K=(S^3_a\times S^3_c)/S^1,\qquad
 (a,c)\sim(e^{i\phi}a,e^{-i\phi}c).
\]
With \(\mathbb H\) oriented by \((1,i,j,k)\), put
\[
 u=a_0+j a_1,\quad v=c_0+j\overline{c_1},\quad
 q=uv,\quad h=u i u^{-1},\quad k=v^{-1} i v.
                                                        \tag{K23.2}
\]
The first resolution has the actual product coordinates \((h,q)\)
on its real normal sphere bundle; the second has \((k,q)\).
To invert the first map, choose a unit representative \(u\) of \(h\)
and put \(v=u^{-1}q\). Changing the representative changes only the
displayed auxiliary-circle orbit. For the second choose \(v\) and
put \(u=qv^{-1}\). These inverse formulas are smooth in the standard
quotient charts. The common boundary map is consequently
\[
             \tau(h,q)=(q^{-1}h q,q).                  \tag{K23.3}
\]
The fixed sphere is unknotted in the smooth six-sphere:
Haefliger's Theorem 1.2 and Corollary 6.6, with \(n=2,q=4\),
give the trivial embedding class \cite{Haefliger1966}.
Its exterior is therefore the actual product
\(E=D^3_y\times S^3_q\), with its given boundary marking extended
by the normal framing. The map (K23.3) extends over this exterior:
\[
 T(y,q)=(q^{-1}yq,q),\qquad
 T^{-1}(y,q)=(q yq^{-1},q).                           \tag{K23.4}
\]
These are smooth also at \(y=0\). Thus replacing the local resolution
conjugates the exterior return by \(T\).

\item This last sentence includes the actual parameter family. For a
local Cox angular action of weights
\((-r-s,0,r,s)\), write
\[
 a_\theta=e^{-i(r+s)\theta/2},\qquad
 b_\theta=e^{i(r-s)\theta/2}.
\]
The first product action and second product action are respectively
\[
 (h,q)\longmapsto
 (a_\theta h a_\theta^{-1},a_\theta q b_\theta),\qquad
 (k,q)\longmapsto
 (b_\theta^{-1}k b_\theta,a_\theta q b_\theta).
\]
Substitution in (K23.3) proves exact equivariance for every \(\theta\).
At \(2\pi\), both scalar endpoints have the same central sign;
the displayed product actions close even when \(r+s\) is odd.
Therefore the comparison of the two local family products is the
same map \(\tau\), independent of the parameter angle. The period
reindexing acts on the periodic collection of common conifold cells.
That collection has trivial cell stabilizer; choose one quotient
representative and extend every chart change equivariantly to its
translates. Extracting the integer winding from each bounded
nonzero toric unit leaves a logarithm on this representative. Contract
its positive modulus and its remaining logarithm in both resolution
charts at once, transporting the truncation collars by the resulting
common conifold isotopy. The integer winding is exactly the Cox
action just displayed. This retains the actual exterior
identifications, including their bounded units. No null-homotopy of
an odd Cox loop, or additional disc-supported correction, is inserted.

\par\medskip\noindent\textbf{Exact relative generators and conjugation.}\enspace
Choose a smooth degree-one map
\(\zeta:(D^3,S^2)\to(S^3,1)\), constant on a boundary collar,
which is equivariant for quaternionic conjugation and commutes
with \(y\). Explicitly one can use
\[
 \zeta(y)=\cos\phi(|y|)-\sin\phi(|y|)\frac{y}{|y|},
\]
where \(\phi\) decreases from \(\pi\) to zero, has the smooth odd
radial expansion at zero, and is flat near the boundary. Define
\[
 L(y,q)=(y,\zeta(y)q),\qquad R(y,q)=(y,q\zeta(y)).
                                                        \tag{K23.5}
\]
They commute and extend over \(D^3\times D^4\) by the same quaternion
multiplications. Their induced six-sphere classes are therefore zero.
They represent the left and right generators of
\(\pi_3SO(4)=\mathbb Z^2\). Extension by the identity to the double
injects the relative pseudo-isotopy group into the doubled group;
Krylov's Proposition 1 and Lemma 3 give, in this dimension, the
relative group \(\Theta _7\oplus\mathbb Z^2\)
\cite{Krylov2004}. Pseudo-isotopy is sufficient here, since extending
the resulting collar gives the same oriented sphere gluing.

\item On the double write
\[
 {\cal A}(s,t)=(tst^{-1},t),\quad
 {\cal B}(s,t)=(s,sts^{-1}),\quad
 {\cal C}=[{\cal A},{\cal B}].
\]
Krylov's Theorems 2--3 establish centrality and order \(28\) of
\({\cal C}\), and the relation
\({\cal A}Y=Y{\cal A}{\cal B}\), where \(Y(s,t)=(s,st)\)
\cite{Krylov2001}. In the marked relative products
\[
       T={\cal A}^{-1},\qquad L=Y,\qquad
       R=Y{\cal B}^{-1}.
\]
The first identity below also follows by literal substitution in
(K23.4)--(K23.5); the second then follows from the displayed central
relation:
\[
 TLT^{-1}=R,\qquad
 TRT^{-1}=L^{-1}R^2{\cal C}.
\]
Consequently
\[
 T(L^mR^n{\cal C}^s)T^{-1}
       =L^{-n}R^{m+2n}{\cal C}^{s+n}.                \tag{K23.6}
\]
All maps on the left are relative to the boundary. The injection
just cited retains this equality in the relative group. Transport
of a supporting oriented disc shows that \(T\) fixes the central
disc-supported subgroup.

\par\medskip\noindent\textbf{The actual integer coefficient from the hopf quotient.}\enspace
The orientations in (K23.2) are essential. At \(a_1=x+iy=0\),
the derivatives of \(h\) are \((-2k,-2j)\), so \(h\) reverses the
complex curve orientation. At \(u=1\),
\[
 q=\Re c_0+i\Im c_0+j\Re c_1+k\Im c_1
\]
preserves the complex normal orientation. Thus the complex
orientation of the exterior is \(D^3_y\times S^3_q\).
For the second resolution the derivatives of \(k\) are \((2k,-2j)\);
its base preserves and its normal coordinate reverses the respective
complex orientations. Its exterior has the same product orientation.

\item The physical delta action is \(q\mapsto q e^{-i\theta}\), with
Hopf quotient \(z=q i q^{-1}\). At \(q=1\), its circle vector is
\(-i\), and the quotient derivatives in the \(j,k\) directions are
\(-2k,2j\). Hence
\[
 [S^3_q]=-[S^1_\delta][S^2_z],\qquad
 [E]=[S^1_\delta][D^3_y][S^2_z].                    \tag{K23.7}
\]
The closed conifold quotient therefore has orientation
\(S^3_y\times S^2_z\), equivalently \(S^2_z\times S^3_y\).

\item Let \(t\) be the relative variation of the actual exterior return in
this product. The integral circle Gysin argument can be made
explicit. Put \(P=Y\setminus\operatorname{int}\nu F\), the orbit
complement before its conifold cap. On its physical mapping torus,
whose fourth cohomology vanishes, let \(\eta\)
be a circle connection with \(d\eta=x\), and take
\(d\beta=x^2\). The integral degree-three generator of the circle
total space is represented by
\(-\eta x+\pi^*\beta\); fibre integration and the meridian-one
condition fix its sign. Its boundary restriction in the complex
curve orientation is \(s+k_{\rm Jin}h_F\theta\):
Stokes gives \(k_{\rm phys}=-\int_{F\times S^1}\beta\), and the
physical mapping torus is the inverse of Jin's convention.
Since \(h=-h_F\), its restriction in (K23.2) is
\[
                         t=-k_{\rm Jin}.             \tag{K23.8}
\]
Equivalently, the boundary of a relative \(D^3\times[0,1]\) is
\(S^2\times[0,1]-D^3\times\{1\}+D^3\times\{0\}\), giving precisely
this variation coefficient. The integral Gysin sequence supplies
the integral class; choosing differential representatives does not
replace the integral assertion.

\item The generator \(L\) has variation \(+1\) and quotient
\((y,z)\mapsto(y,\operatorname{Ad}_{\zeta(y)}z)\). In Jin's
identification \((v,0)\sim(g(y)v,1)\), a connection is
\(A_\theta=-\theta\,dg\,g^{-1}\). With \(\omega=dg\,g^{-1}\),
\[
 F=-d\theta\,\omega+(\theta^2-\theta)\omega^2,\qquad
 \int_{S^3}\operatorname{tr}(\omega^3)
                      =-24\pi^2\deg g.
\]
The orientation is \(S^3\times S^1\), so direct integration gives
\(c_2=\deg g\) and \(p_1(\operatorname{ad}g)=-4\deg g\).
The vertical Euler class of this sphere bundle is \(2x\), and
the base \(S^3\times S^1\) has \(H^2=0\). Therefore
\[
                       (k_L,\ell_L)=(-1,-4).
\]
Composing the actual equivariant return with \(L^{-t}\) kills its
variation and gives
\[
                      k_0=0,\qquad \ell_0=\ell-4k.    \tag{K23.9}
\]
Additivity follows from the mapping-torus pair-of-pants bordism
carrying the actual Euler line.

\item Put \({\cal B}=LR^{-1}\). The relative class after (K23.9) is
\({\cal B}^b{\cal C}^s\). Its double is the actual principal circle
bundle over the closed conifold mapping torus. Its tangent bundle
is the pullback of the quotient tangent bundle plus a trivial
line. Since \(H^2=\mathbb Zx\), \(H^4=\mathbb Z[y\theta]\), and
\(x^2=0\), the integral Gysin sequence makes the pullback on \(H^4\)
injective.

\item Represent \({\cal C}^s\) on a disc disjoint from
\(S^3_y\times\{q=1\}\). The representative
\({\cal B}^b{\cal C}^s\) fixes that sphere pointwise and has
normal clutch \(\operatorname{Ad}_{\zeta(y)^b}\) on
\(S^3_y\times S^1\). Its Pontryagin number is \(-4b\), by the
preceding connection calculation. The pullback of \([y\theta]\)
evaluates \(1\) on this cycle. This retains the actual marking:
an isotopy of representatives is the identity on an initial fibre,
and the possible \(x\theta\) ambiguity in a quotient degree-three
lift pulls back to zero. Thus
\[
 \ell_U-4k_U=-4b,\qquad b=6A_U,\qquad n=-6A_U=6A_M.              \tag{K23.10}
\]
Here \(n\) is the \(R\)-exponent after rewriting the original
relative class in \(L,R\). This calculation uses the actual Hopf
quotient before the flop; it makes no corresponding quotient claim
after the fixed curve has become two isolated points.

\par\medskip\noindent\textbf{The sign of the central sphere by an oriented plumbing.}\enspace
Let \(\Sigma_f\) denote \(D^4_s\times S^3_t\) glued to
\(S^3_s\times D^4_t\), right coordinate equal to \(f\) of the left
coordinate, with the orientations from
\(\partial(D^4_s\times D^4_t)\). Take three products
\(D_i=D^4_s\times D^4_t\), with orientations \(+,-,+\).
Glue their faces by
\[
 (s_1,t_1)\longmapsto(s_2,t_2)
       =(s_1,s_1t_1s_1^{-1}),\qquad
 (s_2,t_2)\longmapsto(s_3,t_3)
       =(t_2^{-1}s_2t_2,t_2).
\]
The first is on the \(S^3_s\times D^4_t\) faces, the second on the
\(D^4_s\times S^3_t\) faces. Their coordinate orientations agree,
while the adjacent eight-ball signs differ, as required for oriented
gluing. After rounding, the two remaining faces have boundary word
\({\cal A}^{-1}{\cal B}\). The map \({\cal A}\) extends the left
face and \({\cal B}^{-1}\) the right face, so the oriented boundary is
\[
 \Sigma_{{\cal B}^{-1}{\cal A}^{-1}{\cal B}{\cal A}}
                       =\Sigma_{{\cal C}^{-1}}.
\]

\item Orient the first zero-section sphere by \(ds\) in \(D_1\) and
\(-ds\) in \(D_2\), and the second by \(dt\) in \(D_2\) and
\(-dt\) in \(D_3\). At their unique intersection in \(D_2\) their
concatenated orientation is \(-ds\,dt\), the ambient orientation.
Each normal clutch fixes a real line. Thus
\[
 Q=\begin{pmatrix}0&1\\1&0\end{pmatrix},\qquad \sigma(W)=0.
\]
The normal clutches are \(\operatorname{Ad}_s\) and
\(\operatorname{Ad}_{t^{-1}}\). Their normal orientations on their
two hemispheres are respectively \(dt\) and \(-ds\).
Quaternionic conjugation reverses the four-plane orientation and
commutes with every adjoint action, so this second normal-frame
orientation introduces no extra Pontryagin sign. The two
Pontryagin evaluations are \((p,-p)\), with \(|p|=4\).
The eight-manifold has only two four-handles besides its zero-handle;
in particular it is spin. Its unimodular intersection form gives
\[
 p_1^2=(p,-p)Q^{-1}(p,-p)^{\mathsf T}=-32.
\]
The spin Eells--Kuiper formula consequently gives
\[
 \mu(\Sigma_{{\cal C}^{-1}})=-\frac1{28},\qquad
 \mu(\Sigma_{\cal C})=+\frac1{28}.                   \tag{K23.11}
\]

\item The physical comparison uses the explicit cap maps of
\cref{night:P3K.68}, with the oppositely oriented (y)-ball
implemented by the reflection (
ho(y_0,y_1,y_2,y_3)=(-y_0,y_1,y_2,y_3)).
For (L=D^4_y\times S^3_q), (R=S^3_y\times D^4_q), and
first-to-second clutch (F), the physical pieces are (-R,-L).
Writing (L) first changes the arrow to (F^{-1}), hence
(lambda_M(F)=-Sigma_{F^{-1}}).
A disc-supported clutch and its inverse bound the cylinder obtained
by cancelling their opposite product collars, so
(Sigma_{F^{-1}}=-Sigma_F). Therefore
(lambda_M(mathcal C)=Sigma_{mathcal C}) with the positive
value in (K23.11). This fixes the central generator independently
of the dual Euler line used in the characteristic calculation.

\item Finally \(L\) and \(R\) extend the seven-disc, so the central exponent
of the original relative class is its actual physical class \(A_M\bmod28\).
Equation (K23.6) adds \(n\) to that exponent, and (K23.10) gives
\(n=6A_M\). This proves (K23.1) and the displayed examples.

\item All local orientation matrices and the Chern--Weil integral are
checked exactly in \path{code/part3_fourfold/fixed_edge_signs.py}. The
central word, hyperbolic pairing and rational Eells--Kuiper
arithmetic are checked in \path{code/part3_fourfold/commutator_orientation.py}.
These programs certify those algebraic steps; the actual family,
circle, collar and oriented gluing maps are the constructions in
this proof.

\end{enumerate}
\end{proof}




\begin{theorem}\label{night:P3K.24}
\StatementRecord{night:P3K.24}{false as stated; correction proved}{S6\_06\_BLIND: GPT-6 Astra (OpenAI, ChatGPT Pro mode, effort Pro), recorded re-derivation in the checkpoint dated 18 November 2026}{\cref{night:P3F.22}, \cref{night:P3F.24}, \cref{night:P3F.37}, \cref{night:P3F.9}, \cref{night:P3K.5}, \cref{night:P3K.6}, \cref{night:P3K.7}}For the family of P3F.37, put $m=4l+3$, $l\ge0$, and retain
its actual parameter $\lambda=r^m$. Its positive $r$-return
satisfies
\[
 A_M(\Phi)=3l+2,\qquad
 \mu(\Sigma_\Phi)=\frac{3l+2}{28}\pmod{\mathbb Z},
 \qquad [\Phi]=3l+2\pmod{28}.                         \tag{K24.1}
\]
The same construction and formula hold for the explicit source
$y^2=x^3+tx-t$, $G=(1,1)$, with fibres III*$+$II$+I_1$,
using its $I_1$ as the sole simple whole-fibre zero. As $l$ runs
from zero to twenty-seven, either set of families realizes all
oriented classes in $\Theta_7$. For a nonzero residue its mapping
torus is not diffeomorphic to $S^1\times S^6$.
\end{theorem}
\begin{proof}
\begin{enumerate}[label=\arabic*.]
\item We retain the complete source, line, cap and frame maps below.
The meridian-one graph line is (U=M^{-1}): the fibre map
(D_M(v)(w)=\langle v,w\rangle/\langle v,v\rangle), with the
Hermitian form linear in its second argument, satisfies
(D_M(e^{i\theta}v)=e^{-i\theta}D_M(v)).
Its (SU(2))-equivariant spinor extension is
((a,b)\mapsto(\overline b,-\overline a)).
The first-to-second cap maps in \cref{night:P3K.68} give
(\lambda_M(F)=-\Sigma_{F^{-1}}) and
(\lambda_U(F)=\Sigma_{F^{-1}}); cancelling the opposite
product collars gives (A_M=-A_U). We therefore distinguish the
intermediate meridian-one number (A_U=(4k_U-\ell_U)/24)
from the physical return (A_M) throughout this calculation.

\item \par\medskip\noindent\textbf{Physical characters and the closing arrow.}
First use the source of P3F.37. Write $t_c$ for the selected cusp,
$q=t-t_c$, $s=\log q/(2\pi i)$, and $v=\log r/(2\pi i)$.
Use the physical source Abel coordinate $a=\zeta_1$ and the
physical line character $V=e^{2\pi i\zeta_2}$. Let $g(q)$
be the holomorphic Tate logarithm of $+G$, and let $C(q)$ be
the bounded unit logarithm in the linearization. The actual maps
are
\[
 \begin{aligned}
 B(a,V)&=(a-g,e^{2\pi iC}qr^mV),\\
 B^{-1}(a,V)&=(a+g,e^{-2\pi iC}q^{-1}r^{-m}V).
 \end{aligned}
\]
Thus the positive outward suspension loop is the period
\[
 \gamma_{\rm out}=\binom{g}{-C-s-mv},\qquad
 \delta=\binom01.                                    \tag{K24.2}
\]
Delta is positive physical scalar rotation. In particular the
period (K24.2) is $B^{-1}$, the same outward loop as in the
cap meridian and the Seifert marking; neither coordinate has
been silently inverted.

\item The annular arrow in P3F.37 is
$F_c:\text{ordinary}\to\text{cap}$, where
$c=\log\lambda/(2\pi i)$. Its Lie translation satisfies
$g_c(s_\infty+1)=Dg_c(s_\infty)+a_c$, so its scalar variation
is $(c-c_0)s_\infty\delta/m$. The cap is identical at $c+m$.
Consequently the inverse closing is $F_c^{-1}F_{c+m}$, with
multiplier $q_\infty$ at infinity, and globally
$(t-t_c)^{-1}$. The positive parameter deck is the inverse:
\[
              (v,\zeta)\longmapsto(v+1,\zeta+s\delta). \tag{K24.3}
\]
This square fixes its direction, including its bounded unit;
it does not infer the direction from a finite-difference identity.

\item On the auxiliary cusp cover $q=Q_0^{2m}$ take the extending origins
\[
 o(s,v)=\binom{-gs/m}{Cs/m+s^2/(2m)+sv},\qquad
 \sigma(s,v)=\binom{gv}{-Cv-m(v^2-v)/2}.
\]
Their exact differences are
\[
 \begin{aligned}
 o(s,v+1)-o(s,v)&=s\delta,\\
 \sigma(s,v+1)-\sigma(s,v)&=s\delta+\gamma_{\rm out},\\
 o(s+2m,v)-o(s,v)&=-2\gamma_{\rm out}+2m\delta,\\
 \sigma(s+2m,v)-\sigma(s,v)&=0.
 \end{aligned}                                                    \tag{K24.4}
\]
The factor two in this auxiliary cover is necessary for these
particular origins: $m$ is odd, and the one-$m$ shift would
contain $m\delta/2$. Every discrepancy in (K24.4) is an actual
integral period. The cusp origin has bounded characters in $q$
and extends in the unimodular cusp charts. Subtracting $o$
removes the parameter deck on the ordinary collar. For fixed
nonzero $q$, its remaining multiplier is a unit times $r^m$;
the actual $m$-node chain
$X_i=V/r^i$, $Y_i=r^{i+1}/V$, $X_iY_i=r$ gives its extending
origin. These are the two specified collar comparisons, without
an asserted two-variable proper analytic corner.

\item Their difference and angular change are
\[
 F=\sigma-o=
 \binom{g(s/m+v)}{-C(s/m+v)-(s+mv)^2/(2m)+mv/2},
\]
\[
 F(s,v+\theta)-F(s,v)
 =\theta\gamma_{\rm out}
          -\frac m2\theta(\theta-1)\delta.             \tag{K24.5}
\]
The last term is an endpoint-fixed periodic unitary gauge.
It vanishes in the circle quotient. The whole displayed angular
difference changes under $s\mapsto s+2m$ by $-2m\theta\delta$, so its retained
line transition has winding $-2m$. The actual positive $r$-loop
therefore makes one positive outward gamma turn. The auxiliary
cusp cover has not changed the parameter-loop count.

\par\medskip\noindent\textbf{The actual Seifert space and page.}
The cap marking is
\[
 (z,u)\sim(\zeta_mz,u-1),\qquad q_\infty=z^m,
 \qquad u=m\gamma_{\rm out}-\arg(q_\infty)/(2\pi).
\]
Let $R(z)>0$ solve $R^2+|z|^2R^{2/m}=1$. The explicit map
\[
 (z,u)\longmapsto
 (z_1,z_2)=
 \bigl(zR^{1/m}e^{2\pi iu/m},R e^{2\pi iu}\bigr)       \tag{K24.6}
\]
identifies the marked three-dimensional quotient with the
weighted circle fibration on $S^3$ of weights $(1,m)$.
The root and phase factors cancel under its indicated equivalence.
Conversely a real argument of $z_2$ recovers $u$ and then $z$;
another argument gives exactly that equivalence. Its base
coordinate is $t=z_2/z_1^m$, and its positive gamma rotation
is $(z_1,z_2)\mapsto(e^{2\pi i\theta}z_1,
e^{2\pi im\theta}z_2)$. In particular it retains
$q\mapsto-h$, $\gamma\mapsto mh$.

\item Choose a regular binding $t=t_*$ near, but distinct from,
$t_c$. Its positive Milnor function is $z_2-t_*z_1^m$.
The affine page $b-t_*a^m=1$ embeds by
$(a,b)\mapsto(\rho a,\rho^mb)$, where the unique positive
$\rho$ solves $\rho^2|a|^2+\rho^{2m}|b|^2=1$.
The inverse uses the positive weighted root of the Milnor
function. Its projection is exactly
\[
                         t=t_*+a^{-m}.                 \tag{K24.7}
\]
The compactification is $\mathbb P^1_a$, branched to order
$m$ at $a=0$ over the additive cap and at $a=\infty$ over
the regular binding. Its open page $C^\circ$ is a disc with
one boundary circle. The chosen cusp and the other finite
critical values each have $m$ unramified preimages. This is
the regular-binding cover; no cusp-branched map is substituted.

\par\medskip\noindent\textbf{The root slice, torsor and line.}
The full cap orbit is the actual mapping torus of $a_0$ on
the resolved III source. In (K24.6) the page equation is
\[
 2\pi u+\arg(1-t_*z^m)=\varphi\pmod{2\pi}.
\]
On a small cap disc choose its real argument. The lifted
graphs are $u=u_0(z)+j$, $j\in\mathbb Z$, where
$u_0=\varphi/(2\pi)-\arg(1-t_*z^m)/(2\pi)$.
The actual $A_r$ moves adjacent graphs and applies $a_0$ to
the source point. Every orbit has a unique representative
on $j=0$: $u_0(\zeta_mz)=u_0(z)$, so an $A_r^j$ equality
on that graph forces $j=0$. Thus
\[
                         p\longmapsto[p,u_0(z(p))]
\]
identifies the entire resolved III source with the actual page
near that core, including both components and their tangency
point. No residual stabilizer quotient remains. The delta
bundle on the graph is the norm circle of the actual line
$N=\mathcal O(2G'-O)$ with a smooth positive radius. Radial
normalization identifies it with the unit circle of $N$.

\item The other finite page pieces are unramified source pullbacks.
At the cusp use P3F.9 relative to its gamma and source boundary
marking. The degree-$m$ page can avoid the short gamma arc of
the fixed-sphere compression exactly as in P3K.6. This gives
an actual affine elliptic fibration over $C^\circ$, locally
translation-isomorphic to the minimal pullback under (K24.7).

\item The comparison can fix the binding collar. In one sector of
the cover the three original nodal vanishing paths lift with
their source marking. Their primitive vanishing cycles span
the full elliptic lattice, by P3F.37's trace-modulo-$d^2$
argument. Thus the global lattice coinvariant group is zero.
For the period local system $\Lambda$ on the punctured page,
compact-support duality gives $H_c^2(\Lambda)=H_0(\Lambda)=0$.
The quotient $j_*\Lambda/j_!\Lambda$ is a finite skyscraper
sheaf with no $H^1$, so
\[
             H^2(C^\circ,\partial C^\circ;j_*\Lambda)=0. \tag{K24.8}
\]
The smooth extending Neron translation sheaf has an identity
part with exponential kernel $j_*\Lambda$ and a fine
Lie-algebra sheaf. At III the identity group is additive;
the local exponential surjectivity follows from its derivative.
At a node it is the multiplicative exponential. Its component
quotient is the skyscraper $\mathbb Z/2$ at III, with zero
$H^1$ relative to the ordinary binding collar. Remove the
component cocycle by local component sections, lift the remainder
to the Lie bundle, remove its integral two-cocycle using (K24.8),
and solve the remaining Lie cocycle by a partition of unity
vanishing on that collar. The resulting translations give
mutually inverse smooth maps to the sectioned pullback, equal
to the prescribed comparison at the binding. Their extensions
over III are the actual extending source translations.

\item For the line comparison use the rigidified Picard map before
pulling back by these smooth translations. On a holomorphic
translation chart form $\operatorname{Hom}(N,T^*N)$ using
the universal extending source translation $T$. Its restriction
to each compact fibre is trivial. On regular elliptic fibres
this is the degree-zero group law; on the III fibre $N$ is
trivial, as the doubling calculation below shows; on a node
the normalization's multiplicative translation preserves its
gluing unit. The retained II case below has normalization
coordinate $u$ translated by a constant: its cusp coordinate
$w=1/u$ changes to $w/(1+tw)$, preserving the first jet
classifying its degree-zero Picard line. Thus the fibre $H^0$
is always one-dimensional. Cohomology and base change makes
the direct image a line, and its evaluation is an isomorphism
on each fibre. Evaluation along $O$ rigidifies it and identifies
it with a base line. Pulling this actual isomorphism back along
the smooth translation gives a smooth, fibrewise holomorphic
line comparison; it does not assign a holomorphic divisor to
a merely smooth total-space map.

\item The actual page line consequently has Chern class that of
$P-O$, $P=2G'$, up to a base pullback. Its absolute base class
vanishes on $C^\circ$. Retaining the prescribed boundary
gauge can leave an integral relative class $nF$, whose square
and product with $M$ vanish because $F^2=FM=0$. On the marked
boundary $T^3$ the degree-zero line and its commuting
translation lift admit a flat representative, so the boundary
Chern class is zero. The relative cup-product calculation
P3K.5 preserves the square for all such retained gauges.

\par\medskip\noindent\textbf{Characteristic numbers and the signed Thom map.}
The compact pullback has one III and $3m$ nodal fibres.
Hence
\[
 \chi=\frac{m+1}{4}=l+1,\qquad \sigma=-8\chi.
\]
The pulled section $P=2G'$ has height $2m$, identity component
at III, and no other reducible-fibre contribution. Thus
\[
 P\cdot O=m-\chi=3l+2,\qquad
 (P-O)^2=-2\chi-2P\cdot O=-2m.                         \tag{K24.9}
\]
These intersections have an explicit location. For (F37.1)
the doubling slope at $G'$ is $1/(2z^{3l+2})$, and the
zero-section parameter of $2G'$ is $2z^{3l+2}$ times a unit.
All $3l+2$ intersections occur at this actual cap; none is
inserted at the regular binding. In particular the central
degree-zero line there is trivial, as used above.

\item The regular elliptic cap and the real compression cap of
P3K.6 have zero signature. The latter kills the nodal vanishing
cycle, and its $H^2$ restriction to the marked $T^3$ is
injective. Thus the actual page gives
\[
                 s_\delta=-2m,\qquad p_{\rm rel}=-24\chi.
\]
We now retain the physical orientation reversal as well as
the Thom meridian. With $V=\rho e^{i\phi}$ and positive
delta $\partial_\phi$, the complex orientation is
$\operatorname{or}(\text{source}_4)\wedge d\rho\wedge d\phi$.
The circle-first quotient therefore has
$\operatorname{or}(Y)=-\operatorname{or}(\text{source}_4)
\wedge d\log\rho$. In particular its identity with the
weighted positive-circle product reverses that product
orientation. This observation alone does not assign the sign
of the return, since it also changes the normal meridian.

\item Let $\kappa$ be the actual sign in
$\partial V=\kappa E\times S^1_\theta$. The compression
has boundary $F-E$, so the glued trace has
$\partial X=\kappa F\times S^1_\theta$. In its outer collar,
\[
 \Omega_Z=\Omega_F\wedge\Omega_{\nu,\mathrm{Hopf}}
                  \wedge dr\wedge d\theta,
 \qquad
 \Omega_X=\kappa\Omega_F\wedge dr\wedge d\theta.
\]
The oriented normal sphere is consequently $\kappa$ times
the Hopf meridian, and the extended actual Euler class has
$\pi_*x=\kappa$. Put $\widehat x=\kappa x$ and
$e_v=2\widehat x-\pi^*b$. The framed graph section gives
$b|_V=-2\kappa c_1(M)$, with the same retained boundary
correction as P3K.5. Therefore $b^2=4M_{\rm rel}^2$.
Fibre integration and both outer/Jin orientation reversals
of P3K.7 give
\[
 A_U(\Phi)=\kappa\frac{3\langle b^2,[\bar X]\rangle
                   -\langle p_1(T\bar X),[\bar X]\rangle}{24}.
                                                               \tag{K24.10}
\]
In the actual affine page the complex boundary of a large
$a$-disc is $+\arg a$. Formula (K24.6) identifies it with
outward gamma, and (K24.5) identifies the positive parameter
with that same outward gamma. Hence $\kappa=+1$. The actual
line, Thom sign and parameter loop are now all fixed. Equations
(K24.9)--(K24.10) give
\[
 A_U(\Phi)=\frac{12(-2m)+24\chi}{24}
        =\chi-m=-(3l+2),\qquad A_M=3l+2.
\]
The Hopf transfer P3F.22 gives (K24.1). The residues
$3l+2$ exhaust $\mathbb Z/28$ because $\gcd(3,28)=1$;
the mapping-torus assertion follows from P3F.24.

\par\medskip\noindent\textbf{The explicit III*$+$II$+I_1$ source.}
For $y^2=x^3+tx-t$ the discriminant is a nonzero constant
times $t^2(4t+27)$. At zero its coefficient orders are
$(1,1,2)$, giving II, and $-27/4$ gives $I_1$. At infinity,
\[
 y_0^2=x_0^3+q^3x_0-q^5,\qquad G=(q^2,q^3),
\]
giving III*. The source is rational by
$t=(y^2-x^3)/(x-1)$. Its fourth good root sends $G$ to the
nonzero fixed two-torsion point, so its $E_7$ component,
height $1/2$, primitivity and absence of torsion are the same
marked data as in P3F.37. In particular $2G\cdot O=0$ and
its Hom line for $Q=G$ is the whole-fibre line of degree one.
Use its sole $I_1$ as the zero. The actual reduced root is
\[
 Y^2=X^3+zX-z^{2l+3},\qquad
 G'=(z^{2l+2},z^{3l+3}).
\]
The quadratic part is again $Y^2-zX$; the same resolution,
deck, unique $m$-division, line norm and meridian construct its
III cap. The basis $O,F,E_7,G$ has the identical unimodular
Gram matrix and $TM=M-F$. The two vanishing cycles of its
actual II fibre meet once and generate the elliptic lattice.
This gives the source-complement simple connectivity and every
unchanged integral gluing map in P3F.37. Thus this source too
gives the stated proper family of standard six-spheres.

\item Its actual page has one III, $m$ II and $m$ nodal fibres,
again of Euler number $3+3m$. The same lifted vanishing
relations prove (K24.8) relative to the binding. Its doubling
slope is $(1+3z^{4l+3})/(2z^{3l+2})$, with the same leading
zero-section parameter. Every line, signature and oriented
trace calculation above therefore has the same value. The
independent certificates are
\path{code/part3/reduced_iii_return.py} and
\path{code/part3_root/reduced_iii_return_audit.py}; the working
note retains the superseded character convention verbatim in
its dated historical record.

\end{enumerate}
\end{proof}




\begin{theorem}\label{night:P3K.25}
\StatementRecord{night:P3K.25}{false as stated; correction proved}{S6\_06\_BLIND: GPT-6 Astra (OpenAI, ChatGPT Pro mode, effort Pro), recorded re-derivation in the checkpoint dated 18 November 2026}{\cref{night:P3F.22}, \cref{night:P3K.10}, \cref{night:P3K.17}, \cref{night:P3K.22}}
The literal square charts, period reindexing, scalar Cox lift,
common origin difference and inner charged-cap maps below exist
for each primitive scalar source of \cref{night:P3K.22}.
The assertion that every vertical shear has the same complete
return is false. For the actual OS45 family the reference fan
has physical class (3), whereas (T_2) has class (2).
The missing outer determinant is calculated explicitly below.
\end{theorem}
\begin{proof}
\begin{enumerate}[label=\arabic*.]
\item \par\medskip\noindent\textbf{The literal square chart and the unchanged period group.}\enspace

\item First take the horizontal edge from \(0\) to \(e_1\). Its two
adjacent triangles have the other vertices \(e_2,e_1-e_2\).
Order the four rays, including their height-one coordinate, as
\[
\begin{aligned}
 u_0&=(0,0,1),&u_4&=(1,0,1),\\
 u_1&=(1,-1,1),&u_2&=(0,1,1).
\end{aligned}
                                                        \tag{P3K.25.2}
\]
Their relation is \(u_0+u_4=u_1+u_2\). This is exactly the Cox
square of P3K.17 under
\[
 B=\begin{pmatrix}1&0\\-1&1\end{pmatrix}.
\]
The auxiliary weights are \((1,1,-1,-1)\), and the physical
delta weights are \((-1,0,0,1)\): their image in the height-zero
lattice is \(u_2-u_0=e_2\). Hence every nonzero two-by-two
minor of the two-circle weight matrix is \(1\) or \(-1\).
The seven-dimensional real master space has the same semifree
residual circle, fixed normal pairs, and inner \(\mathbb{CP}^3\)
face as P3K.17.

\item The torus coordinates in this square are the literal functions
\[
                  Z_{\mathrm{loc}}=Z,\qquad
                  V_{\mathrm{loc}}=ZV.                 \tag{P3K.25.3}
\]
They are used on the common dense torus to name the Cox maps;
they do not change the global source line. Write the actual
whole-I1 periods in its multidegree-zero source frame as
\[
\begin{aligned}
 h:(Z,V)&\longmapsto(qZ,\alpha(q)V),\\
 \gamma:(Z,V)&\longmapsto(g(q)Z,q\lambda\beta(q)V),
\end{aligned}
                                                        \tag{P3K.25.4}
\]
where the nonzero units have parameter-independent logarithms.
These are the periods of P3K.22.3, with
\(g=\exp(2\pi i\Gamma)\), \(\beta=\exp(2\pi iC)\).
In (P3K.25.3) they are exactly
\[
\begin{aligned}
 h:(Z_{\mathrm{loc}},V_{\mathrm{loc}})
       &\longmapsto(qZ_{\mathrm{loc}},q\alpha V_{\mathrm{loc}}),\\
 \gamma:(Z_{\mathrm{loc}},V_{\mathrm{loc}})
       &\longmapsto(gZ_{\mathrm{loc}},
                          qg\beta\lambda V_{\mathrm{loc}}).
\end{aligned}
                                                        \tag{P3K.25.5}
\]
Their integer labels are \((1,1)\) and \((0,1)\).
They generate all of \(\mathbb Z^2\). Thus the disjoint periodic
collection of square-diagonal tubes is one free period orbit,
with no stabilizer on a chosen master-chart representative.
The quotient of the disjoint collection is that single chart.
It has no residual source or gamma circle.

\item The original fan uses the \(u_0u_4\) diagonal, and the new fan
uses \(u_1u_2\). These are the two moment faces of P3K.17 in
the reverse order from that statement. All its actual quotient
maps on the common rank-one matrix locus remain the identity
maps in the named global torus coordinates.

\par\medskip\noindent\textbf{The actual scalar loop and the common endpoint gauge.}\enspace

\item The positive parameter deck on the ordinary cusp collar is
the actual line multiplication \(q^{-1}\). It is the deck
derived from the directed cap-to-ordinary gluing in P3K.10
and P3K.22. Returning its shifted square representative using
the unique gamma period cancels the integral \(q\)-translation.
The remaining bounded torus map in (P3K.25.3) is
\[
       (Z_{\mathrm{loc}},V_{\mathrm{loc}})
           \longmapsto(gZ_{\mathrm{loc}},
                             g\beta\lambda V_{\mathrm{loc}}).
                                                        \tag{P3K.25.6}
\]
Its Cox lift is the literal formula
\[
\begin{aligned}
 z_0&\longmapsto(g^2\beta\lambda)^{-1}z_0,
 & z_4&\longmapsto z_4,\\
 z_1&\longmapsto gz_1,
 & z_2&\longmapsto g\beta\lambda z_2 .
\end{aligned}                                          \tag{P3K.25.7}
\]
The product of the four multipliers is one, so the base
monomial is fixed. The parameter-dependent part has weights
\((-1,0,0,1)\); it is one global physical delta multiplication.
The same formula acts at both fixed charts of the master
space. Its remaining unit multipliers are independent of the
parameter and have the stated logarithms.

\item This comparison also retains the actual extending origins,
rather than choosing a new loop from (P3K.25.7).
With \(s=\log q/(2\pi i)\), \(v=\log\lambda/(2\pi i)\),
they are
\[
 o=(-\Gamma s,-Cs-s^2/2-sv),\qquad
 \sigma=(\Gamma v,Cv+(v^2-v)/2).
\]
Their difference changes by
\[
 (\sigma-o)(v+\theta)-(\sigma-o)(v)
     =\theta\gamma+\frac{\theta(\theta-1)}2\delta.
                                                        \tag{P3K.25.8}
\]
The first term is the one positive gamma traversal of the
actual family. The second is one common delta gauge. On the
auxiliary source cover \(q=Q^2\), the source deck changes this
angular expression by \(2\theta\delta\). This is likewise the
same physical circle scalar on every master chart and both
cap endpoints. The half-integral one-step ordinary-origin
term is retained exactly as in P3K.10; no false one-step
descent is imposed.

\item To pass through the real moment bordism, follow each bounded
Cox unit map by its unique positive auxiliary rescaling back
to the prescribed moment value. The derivative of the moment
equation is positive away from the removed origin. Use the
images of the inner and outer truncation faces under these
equivariant isotopies, exactly as in P3K.17. Contract the
parameter-independent units through their chosen logarithms,
with the exterior cutoff and transported collars. The
remaining angular map is physical delta and therefore acts
trivially on the circle orbit and on its oriented normal
three-plane bundle.

\par\medskip\noindent\textbf{All extra boundaries and their actual lines.}\enspace

\item The inner face has orbit \(S^2_A*S^2_B\). Its free part carries
the actual circle line
\[
                       x=h_A-h_B.
                                                        \tag{P3K.25.9}
\]
The two charged cap maps are those of P3K.17:
\((A,B)=(P,\iota Q)\) at the first end and
\((A,B)=(Q,\iota P)\) at the second. Both pull the line to
\(p+q\), by the explicit orthogonal-complement line
isomorphism. The resulting closed five-dimensional correction
is \(T_g(S^2\times S^2)\), \(g(A,B)=(\iota B,\iota A)\).
Here this \(g\) denotes the conjugate-swap map, not the
source-period unit in (P3K.25.4).

\item On the two overlap components of this joint cap, (P3K.25.7)
and (P3K.25.8) give transition functions
\(c(\theta)\varphi_{1,0}\) and
\(c(\theta)\varphi_{2,0}\), with the identical scalar
\(c(\theta)\). One change of cap frame by \(c(\theta)^{-1}\)
removes both angular factors. Thus its line and tangent
Pontryagin class pull back from the five-dimensional
correction. Its two six-dimensional characteristic numbers
\(\int x^3,\int p_1x\) vanish.

\item The fixed locus of the master space is the same actual
three-dimensional cobordism
\[
 (S^2\times I)\setminus
                   (\mathring D^3\sqcup\mathring D^3).
\]
It is obtained by keeping the annulus in the original fixed
curve and the two Cox fixed charts. Its four boundary
line charges are \(0,0,+h_A,-h_B\).
The two zero endpoint charges follow from the literal
\(\mathcal O(-1)\oplus\mathcal O(-1)\) normal charts with
weights \((-1,+1)\), independently of the source line \(M\).
Conjugating the negative-weight summand gives the oriented
Hopf normal model with determinant class
\(-h+h=0\); the adjacent ray relation is unchanged by
every shear fixing \(e_2\).
The original endpoint conifold caps and the joint charged
cap close it to the oriented \(S^2\)-mapping torus of P3K.17.
Its parameter circle is a separate factor: the parameter
gauge is literally physical delta, and the common scalar was
removed by the preceding single frame change. Consequently
the class \(b\) of the additional normal-sphere boundary
satisfies \(b^2=0\), and its base tangent Pontryagin number
is zero. Retaining that entire sphere boundary gives
\[
 \frac{4\int x^3-\int p_1x}{24}
       =\frac{3\int b^2-\int p_1(T\bar X)}{24}=0.
                                                        \tag{P3K.25.10}
\]
This uses no value for the normal bundle's \(p_1\), and no
spin extension across the comparison bordism.


\item The actual outer attaching map has to be retained as well.
For full periods (h:(Z,V)\mapsto(\kappa qZ,\alpha V)) and
(B:(Z,V)\mapsto(gZ,q\eta V)), the second outer star has
(z_1=(g/\kappa)z_0) and normal coefficient map
[
 (N_1^-,N_1^+)=
 ((\eta/\alpha)N_0^-,\;\kappa\alpha/(g\eta)N_0^+).
]
Conjugating the negative-weight normal coordinate makes its
complex determinant ((\kappa/g)(\alpha/\overline\alpha)
(\overline\eta/\eta)), of parameter winding (-2J).
The simultaneous marked equivariant product equations would give
(-J=0), whereas the scalar family has (J\ne0).
These are the full coefficient arrows of \cref{night:P3K.61}.
The inner charged-cap calculation above therefore cannot be glued
with an unchanged outer product to deduce equality of the returns.
For OS45, (J=1,H=8,\chi=1). The exact vertical source and period
arrows (Z_n=Z, V_n=Z^{-n}V, h_n=hB^n, B_n=B) in
\cref{night:P3K.72} give
(A_M(n)=H/2-\chi-Jn(n-1)/2=3-n(n-1)/2).
Thus (A_M(0)=3) and (A_M(2)=2), proving the stated counterexample.

\end{enumerate}
\end{proof}




\begin{theorem}\label{night:P3K.27}
\StatementRecord{night:P3K.27}{false as stated; correction proved}{S6\_06\_BLIND: GPT-6 Astra (OpenAI, ChatGPT Pro mode, effort Pro), recorded re-derivation in the checkpoint dated 18 November 2026}{\cref{night:P3F.9}, \cref{night:P3K.22}}For the ten primitive scalar families of K22, use the actual
positive scalar loop, the standard vertical-edge cusp, and the
normal product transported by the extending cusp origin. Its
free normal-framing gauge agrees with the transported rotating
framing of K9. A possible constant rotation loop extends over
the fixed three-disc and changes neither integer. Thus the
meridian-one dual characteristic pair and the corresponding physical relative
coordinates may be taken to be
\[
       k=-3d,\qquad \ell=-24,\qquad
       (m_{\rm rel},n_{\rm rel},s)
             =(6,\ 3d-6,\ d/2-1\bmod28).
                                                        \tag{P3K.27.1}
\]
Here \(d=8,14,20,30\) is the height denominator in K22. The
coordinate convention is exactly K23, including its complex
orientation and the positively oriented Eells--Kuiper generator.
This does not identify a further edge's relative four-cycle.
\end{theorem}
\begin{proof}
\begin{enumerate}[label=\arabic*.]
\item We retain the complete source, line, cap and frame maps below.
The meridian-one graph line is (U=M^{-1}): the fibre map
(D_M(v)(w)=\langle v,w\rangle/\langle v,v\rangle), with the
Hermitian form linear in its second argument, satisfies
(D_M(e^{i\theta}v)=e^{-i\theta}D_M(v)).
Its (SU(2))-equivariant spinor extension is
((a,b)\mapsto(\overline b,-\overline a)).
The first-to-second cap maps in \cref{night:P3K.68} give
(\lambda_M(F)=-\Sigma_{F^{-1}}) and
(\lambda_U(F)=\Sigma_{F^{-1}}); cancelling the opposite
product collars gives (A_M=-A_U). We therefore distinguish the
intermediate meridian-one number (A_U=(4k_U-\ell_U)/24)
from the physical return (A_M) throughout this calculation.

\item \par\medskip\noindent\textbf{1. The actual two products and their literal transition.}\enspace

\item Write the actual periods as
\[
 h:(Z,V)\mapsto(qZ,\alpha(q)V),\qquad
 \gamma:(Z,V)\mapsto(g(q)Z,q\lambda\beta(q)V).
\]
On the positive parameter circle, \(|\lambda|\) is fixed.
After quotienting the physical delta circle, the parameter
phase in this gamma period disappears: the orbit period is
\((Z,|V|)\mapsto(g(q)Z,|q\lambda\beta(q)|\,|V|)\).
The same assertion holds in each actual Cox chart: its
parameter phase has the physical-delta weights, so it is the
identity on that circle orbit. The chosen fan and its
period-cell quotient therefore supply a fixed smooth orbit
body along this parameter circle, before applying the
closing translation.

\item The actual ordinary closing translation is \(q^{-1}\) in the
line character. Returning its period cell by gamma leaves
the bounded unit map \((g,\lambda\beta)\). It is removed by
the actual extending cusp origin
\[
 \sigma(v)=(\Gamma v,Cv+(v^2-v)/2),\qquad
 \lambda=e^{2\pi iv},\quad
 g=e^{2\pi i\Gamma},\quad\beta=e^{2\pi iC}.
                                                        \tag{P3K.27.2}
\]
Indeed
\[
 \sigma(v+1)-\sigma(v)=\gamma-s\delta,\qquad
             s=\frac{\log q}{2\pi i}.
\]
Thus the ordinary deck minus the cusp-origin difference is
the actual negative gamma period. In the period quotient
it is the identity. This proves the closed cusp product,
including the representative-cell identification; it is
not a choice of a new parameter loop.

\item On the same universal parameter interval the bounded
coordinate change is
\[
 C_\theta=
 \left(g^\theta,\,
       (\beta\lambda_0)^\theta
                       e^{\pi i\theta(\theta-1)}\right).
                                                        \tag{P3K.27.3}
\]
All powers use the given logarithms. These are literal
bounded torus maps and extend through both adjacent
unimodular charts. They also extend in the Cox master
coordinates used in K23. At \(\theta=1\) they equal the
bounded period-return map, as required. The last factor
is physical delta and acts trivially on the orbit normal
three-plane. Thus the product from (P3K.27.2) is precisely
the normal product obtained by carrying the Cox charts
through their actual bounded clutch.

\item The ordinary extending origin is
\[
                  o=(-\Gamma s,-Cs-s^2/2-sv).
\]
The actual cusp-to-ordinary transition in the two product
coordinates is \(\sigma-o\). Its angular change is
\[
        \theta\gamma+\frac{\theta(\theta-1)}2\delta.
                                                        \tag{P3K.27.4}
\]
On the free orbit annulus, the second term is the identity
and the first is the marked gamma translation. In the
source-preserving Seifert product of P3F.9, that translation
is exactly \(R_\theta\), rotation of its gamma circle.
This retains the source-translation trivialization
\(Z\mapsto\exp(t\log g(q))Z\) used there.

\par\medskip\noindent\textbf{2. The real extension at the node and its derivative.}\enspace

\item The radial part of (P3K.27.4) contains \(|q|^\theta\).
It is used as an annular expression only. The actual smooth
extension in the orbit is the following fixed product map,
which agrees with that annular expression.

\item Use P3F.9's two exact small-resolution charts for
\(xy=uv\):
\[
 (x,v,z),\ u=xz,\ y=vz;\qquad
 (u,y,w),\ x=uw,\ v=yw.
\]
The physical circle has weights \((-1,+1)\) on \((v,z)\)
in the first chart, and \((+1,-1)\) on \((u,w)\) in the
second. Its actual normal Hopf coordinates are
\[
 (y,b),\quad b=(|z|^2-|v|^2)/2;
 \qquad
 (x,b'),\quad b'=(|u|^2-|w|^2)/2.
                                                        \tag{P3K.27.5}
\]
The fixed curves in these charts are \(v=z=0\) and
\(u=w=0\), respectively. On their overlap, with \(r=|z|^2\),
\[
 2b=r-|y|^2/r,\qquad
 2b'=|x|^2r-r^{-1},\qquad
 \frac{db'}{db}
       =\frac{|x|^2+r^{-2}}{1+|y|^2r^{-2}}>0.
                                                        \tag{P3K.27.6}
\]
Choose the relative, increasing proper fibre coordinate
\(\chi(x,y,b)\) constructed in P3F.9, agreeing with the
marked gamma coordinate on the outside collar. Its
derivative \(\chi_b\) is strictly positive. It gives the
fixed smooth map
\[
         \Phi:(x,y,b)\longmapsto(x,y,\chi(x,y,b))
\]
to the nodal source neighbourhood times the gamma circle,
with the corresponding second-chart formula.
The construction is independent of the parameter angle
in the closed cusp product (P3K.27.2).

\item The actual quotient family, in the two extending products,
is glued on its outside annulus by (P3K.27.4). It therefore
has the following explicit smooth description in the
constant orbit:
\[
          J_\theta=R_\theta\circ\Phi
             \quad\hbox{on its cusp piece},             \tag{P3K.27.7}
\]
with the fixed ordinary product on its other piece.
Equation (P3K.27.4) proves that these maps agree on the
actual annular identification. Equation (P3K.27.6) proves
smoothness through the two node charts. Thus (P3K.27.7)
is an actual comparison map of the orbit family, rather
than a formal differentiation of a fractional monomial
at \(q=0\).

\item At the first fixed chart, quotient the derivative by the
tangent to its fixed curve, including the derivative of
its graph gamma coordinate. On the normal coordinates
\((\Re y,\Im y,b)\), the resulting map is
\[
 Q_x=
 \begin{pmatrix}
 1&0&0\\
 0&1&0\\
 \chi_{\Re y}&\chi_{\Im y}&\chi_b
 \end{pmatrix},\qquad \chi_b>0.                         \tag{P3K.27.8}
\]
There is the same matrix in the second chart with \(x,b'\)
in place of \(y,b\). These are the actual normal derivatives
at the node endpoints as well: (P3K.27.6) has no excluded
nonzero-source-coordinate hypothesis in either smooth
Hopf chart.

\item The matrices (P3K.27.8) are independent of \(\theta\).
Their off-diagonal entries contract linearly to zero,
and their last diagonal entry contracts through positive
numbers to one. The overlap derivative in (P3K.27.6) is
positive, so these contractions are the coordinate
expressions of the global normal-bundle comparison.
Intrinsically, the free gamma vector gives its oriented
normal line, and the quotient is the normal two-plane
bundle of the immersed nodal normalization in the source.
The derivative is the identity on that quotient and is
positive on the gamma line. Choosing a splitting is
contractible; the preceding contractions therefore glue
globally even though the source normal two-plane bundle
has nonzero Euler class.
The source tangent shear, already removed in obtaining
(P3K.27.8), also contracts through invertible triangular
matrices. In particular the \(|q|^\theta\) annular factor
introduces no omitted normal derivative: it has been
extended by the explicitly smooth map (P3K.27.7).

\par\medskip\noindent\textbf{3. The fixed-frame comparison.}\enspace

\item The initial Cox product is explicit on the whole fixed
sphere. Write its two affine coordinates as
\(\zeta,\zeta'=\zeta^{-1}\), and its two complex normal
coordinates as \(c_0,c_1\), with
\(c_i'=\zeta c_i\). In the quaternion convention of K23 put
\[
\begin{aligned}
 q_{\rm H}
   &=(1+j\zeta)(c_0+j\overline{c_1})\\
   &=(\zeta^{-1}+j)(c_0'+j\overline{c_1'}).
\end{aligned}                                          \tag{P3K.27.10}
\]
These are literally equal. To check the overlap, substitute
\(c_i'=\zeta c_i\) and use
\(zj=j\bar z\) for a complex number \(z\).
The first product expands as
\[
       c_0+j\bar c_1+j\zeta c_0-\bar\zeta\bar c_1,
\]
and the second has exactly these four terms.
The formula is real linear and invertible on every normal
fibre, including either chart origin. It therefore gives
the global real normal product, without an unnamed
trivialization of \(\mathcal O(-1)^{\oplus2}\).

\item The physical weights \((-1,+1)\) give
\(q_{\rm H}\mapsto q_{\rm H}e^{-i\theta}\).
Consequently its global Hopf normal coordinate is
\(q_{\rm H}i\overline{q_{\rm H}}\).
For \(y_{\rm H}=c_0c_1\) and
\(b_{\rm H}=(|c_1|^2-|c_0|^2)/2\), the inner Hopf vector is
\[
 (c_0+j\bar c_1)i(\bar c_0-j\bar c_1)
   =-2b_{\rm H}i+2\Im(y_{\rm H})j-2\Re(y_{\rm H})k.
                                                        \tag{P3K.27.11}
\]
Conjugation by \(1+j\zeta\) gives the desired global
three-plane frame. The same overlap identity proves that
the second-chart expression agrees. The local ordered
coordinates \((\Re y_{\rm H},\Im y_{\rm H},b_{\rm H})\)
have the opposite orientation to the positive physical
Hopf normal orientation, as retained in K23; the displayed
matrix has negative determinant in those local coordinates
and hence preserves the prescribed normal orientation.

\item Choose the global positive-Hopf normal frame supplied by
this actual Cox quaternion product at parameter zero.
Push it to the constant orbit by \(\Phi\). K9 chooses
exactly such a fixed initial frame and transports it by
the actual \(R_\theta\) flow. By (P3K.27.7), the Cox
frame in the closed cusp product is carried to that same
frame:
\[
        dJ_\theta(\text{Cox frame})
           =dR_\theta\,d\Phi(\text{initial Cox frame}).
                                                        \tag{P3K.27.9}
\]
The derivative comparison in (P3K.27.8) is independent
of the parameter and cannot contribute a three-dimensional
normal-framing class on \(F\times S^1_\theta\).
Changing the initial frame changes a map
\(S^2\to SO(3)\), which is null-homotopic; it extends
over the fixed three-disc and is transported by the same
flow. This retains the actual initial Cox marking.

\item For completeness, the bounded coordinate maps in
(P3K.27.3) are applied to both products before this
comparison. Their positive parts cancel as the same
normal-bundle isomorphism. Their unitary part in the
K23 quaternion coordinate \(q_{\rm H}\) is
\[
                   q_{\rm H}\longmapsto
                         a_\theta q_{\rm H}b_\theta,
\]
with \(a_\theta,b_\theta\) independent of the fixed-curve
point. Here is the literal bounded unitary identity fixing
this assertion. Write the first Cox chart as
\(q=\zeta c_0c_1\). A unitary torus map has the form
\[
  (\zeta,c_0,c_1)\longmapsto
                  (a\zeta,p c_0,q_1c_1),\qquad apq_1=1.
\]
Choose its continuous square roots along the parameter interval,
\(A^2=pq_1\), \(B^2=p/q_1\), with \(AB=p\) initially
and hence throughout. Substitution in (P3K.27.10) gives
\[
                      q_{\rm H}\longmapsto A q_{\rm H}B.
                                                        \tag{P3K.27.12}
\]
For example the two complex coefficients of \(q_{\rm H}\),
in the form \(a_{\rm H}+jb_{\rm H}\), are
\(c_0-\bar\zeta\bar c_1\) and
\(\bar c_1+\zeta c_0\). They are multiplied by \(AB\)
and \(B/A\), respectively, on both sides of
(P3K.27.12). This proves the identity without an implicit
normal-frame convention. The induced Hopf normal map is
\(\eta\mapsto A\eta\bar A\), a constant
three-dimensional rotation. Dependence of \(g,\beta\)
on the base produces only the source tangent shears:
in the fixed normal charts the base monomial is
quadratic in the two complex normal variables, so its
higher unit terms have zero normal linear part.
Any residual constant rotation loop is the
\(\mathbb Z/2\) summand of
\(\pi_1{\rm Map}(S^2,SO(3))\); it extends over the fixed
three-disc and contributes neither \(k\) nor \(\ell\).
There is no free \(\pi_3(SO(3))\) discrepancy between
the two actual framings.

\par\medskip\noindent\textbf{4. The integer coordinates.}\enspace

\item We include the complete degree-one normal-clutch calculation in the actual matched framing.

\item P3K.22 identifies the page as the original rational elliptic
source minus a regular fibre, with its original line \(M\).
The page boundary has binding degree one along the positive
inward gamma circle. The actual K6 graph boundary is
\(\gamma=\theta+g_\epsilon(p)\), where \(p\in T^2\) and
\(\theta\) is the positive physical scalar-parameter direction.
The cap-to-ordinary arrow and the two extending origins in K22.3
give exactly this positive direction; no parameter cover is used.

\item Use the graph normal consisting of its oriented normal in the
orbit \(Y\) and the two Cartesian parameter directions.
Use on the compression the actual rank-three normal frame
transported by \(R_\theta\). The compression is the same solid
torus minus a ball as in K6. Its normal bundle has \(w_2=0\):
its tangent bundle is trivial, and \(TY\oplus\mathbb R\) is
trivial for the actual orbit \(Y=S^5\). Obstruction theory on
this three-dimensional compression gives a normal framing.
Its restriction at the fixed-sphere end differs from the
positive-Hopf frame by a map \(S^2\to SO(3)\), which extends
over the fixed \(D^3\). Transport that extension by the same
\(R_\theta\). This defines the rotating compression-cap
framing, including its fixed-sphere attachment.

\item Let \(\phi\) be the fixed outward page-collar phase in radians.
The actual degree-one binding gives
\[
                  \alpha=\phi-\gamma
                         =\phi-\theta-g_\epsilon(p).
\]
The four-dimensional quotient used to compare the graph and
compression normals has the actual collar tangent
\[
             u=(\cos\alpha,\sin\alpha,\rho,0),
             \qquad \rho>0.
\]
Its quotient map is
\[
 (v_q,v_r,v_\eta)\longmapsto
       \bigl(v_q-\rho^{-1}v_r(\cos\alpha,\sin\alpha),v_\eta\bigr).
\]
It kills \(u\). Project the Cartesian parameter vectors
with radial-angular coefficients
\((\cos\theta,-\sin\theta)\), \((\sin\theta,\cos\theta)\),
reverse both of these vectors, and apply the positive
rescalings of K9. The resulting actual graph normal matrix is
\[
 G(\alpha,\theta)=
 \begin{pmatrix}
 -\sin\alpha&\cos\theta\cos\alpha&\sin\theta\cos\alpha\\
 \cos\alpha&\cos\theta\sin\alpha&\sin\theta\sin\alpha\\
 0&\sin\theta&-\cos\theta
 \end{pmatrix}\in SO(3).
\]
Source derivatives of \(g_\epsilon\) are removed by the
explicit triangular shear before taking this quotient.
The outward-page and inward-compression tangent remain
in the same positive radial half-space during the actual
rounding; that half-space is contractible.

\item Write the compression frame in the same transported product
coordinates as \(h(p)\). It is independent of \(\theta\).
The full clutch is therefore homotopic, through these actual
invertible quotient maps and collar homotopies, to
\[
\begin{aligned}
 h(p)^{-1}G(\phi-\theta-g_\epsilon(p),\theta)
   &=[h(p)^{-1}R_z(-g_\epsilon(p))]\\
   &\quad\cdot[R_z(\phi-\theta)G(0,\theta)].
\end{aligned}
\]
The first factor is pulled back from \(T^2\), and the second
from \(S^1_\theta\). The cubic Maurer--Cartan form of their
product differs from the sum of the two cubic forms by an
exact form, by direct expansion of
\((ab)^{-1}d(ab)=b^{-1}a^{-1}(da)b+b^{-1}db\).
Both factor cubic forms vanish by dimension. Integrating
the collar Chern--Weil transgression over the marked \(T^3\)
therefore proves
\[
                     \int_{\bar X}p_1(\bar\nu)=0.
\]
This retains all one-dimensional frame classes and does not
replace the actual normal bundle with a trivial bundle.

\item The actual K22 page has
\(b^2=-4d\) and \(p_1(T\bar X)=-24\), with positive
meridian integral. Thus \(k=-3d\), \(\ell=-24\).
The exact K23 dictionary now gives
\[
      m_{\rm rel}=-k+6A_U=-\ell/4=6,\qquad
      n_{\rm rel}=-6A_U=3d-6.
\]
The central class is \(A_M=-A_U=d/2-1\pmod{28}\) by the explicit dual cap transfer. This proves the
tuple (P3K.27.1) in the explicitly compared products.

\item The actual different-edge four-cycle is still to be
transported through the common exterior in
\path{notes/part3_three_curve_exterior.md.} No finite-isotropy
fan value is inferred from this frame comparison.

\end{enumerate}
\end{proof}




\begin{theorem}\label{night:P3K.29}
\StatementRecord{night:P3K.29}{false as stated; correction proved}{S6\_06\_BLIND: GPT-6 Astra (OpenAI, ChatGPT Pro mode, effort Pro), recorded re-derivation in the checkpoint dated 18 November 2026}{Definitions and the displayed constructions}
The integral cellular bases and reciprocal-section common-exterior
maps below give
[
 m'_0=-m_0+2m_1+2m_2,\quad m'_j=m_j\ (j=1,2),\qquad
 D'_0=-D_0,\quad D'_j=D_j+2D_0.
]
The symmetry of variation and the exact stable frame-change pairing
also hold for the actual relative return maps to which those
constructions are applied. The old scalar initialization
(V=\operatorname{diag}(v,0,0)) and vanishing nonfixed right
integers are refuted by the full outer coefficient maps of
\cref{night:P3K.61,night:P3K.62}.
\end{theorem}
\begin{proof}
\begin{enumerate}[label=\arabic*.]
\item \par\medskip\noindent\textbf{1. The common exterior and its oriented bases.}\enspace

\item Let \(C=C_0\cup C_1\cup C_2\) be the three double curves of the
index-one cusp, in its actual smooth six-sphere. Each curve is a
complex projective line. They meet at the same two triple points,
where the actual charts have \(q=x_0x_1x_2\) and the three curves
are the coordinate axes. Put
\[
 E_C=S^6\setminus\mathring N(C).
\]
The neighbourhood is formed from small coordinate balls at those
points and complex normal tubes along the curves, with all faces
rounded equivariantly.

\item Here is the complete marked cellular calculation. Give each
\(C_i\) two oriented edges \(a_i,b_i\) from the first vertex
to the second and two oriented faces with boundaries
\(a_i-b_i\) and \(b_i-a_i\). In the order
\((a_0,b_0,a_1,b_1,a_2,b_2)\), its chain complex has
\[
 C_2=\mathbb Z^6\xrightarrow{D_2}C_1=\mathbb Z^6
                      \xrightarrow{D_1}C_0=\mathbb Z^2,
 \qquad
 D_1=\begin{pmatrix}-1&-1&-1&-1&-1&-1\\
                     1&1&1&1&1&1\end{pmatrix},
\]
\[
 D_2=\operatorname{diag}(B,B,B),\qquad
               B=\begin{pmatrix}1&-1\\-1&1\end{pmatrix}.
\]
The kernel of \(D_2\) is freely generated by the three
sums of paired faces, which are the positive complex
spheres. Modulo the image of \(D_2\), the two edges
\(a_i,b_i\) become the same edge with coefficient one.
The remaining differential is
\((-1,1)\) on each of these three edges. Consequently
\[
 H_2(C)=\bigoplus_i\mathbb Z[C_i],\qquad
 H_1(C)=\mathbb Z[a_1-a_0]\oplus\mathbb Z[a_2-a_0].
\]
All the identifications used unit coefficients; these
are primitive integral bases.

\item Triangulate the two coordinate-ball models and their
normal tubes compatibly with this cell structure.
In the oriented six-sphere, assign to each cell its
oriented transverse dual cell, with intersection
coefficient \(+1\). Passing to the complement pairs the
reduced complement chains with the cochains of \(C\),
shifted by five, and reverses the boundary incidence
in the usual dual-cell order. The map is integral:
each displayed cell is met once by its own normal
dual cell and not by the others. The cells of the
ambient six-sphere supply only the degree-zero
constant and the ambient orientation class, so the
resulting reduced complement homology is precisely
\(\widetilde H_j(E_C)=\widetilde H^{5-j}(C)\).
This gives
\[
 H_3(E_C;\mathbb Z)=\bigoplus_{i=0}^2\mathbb Z m_i,\qquad
 H_4(E_C;\mathbb Z)=\mathbb Z^2,\qquad H_2(E_C;\mathbb Z)=0.
                                                        \tag{P3K.29.1}
\]
Here \(m_i\) is the boundary of a small positively oriented complex
normal four-disc at a smooth point of \(C_i\). This fixes the
orientation by a local chart, independently of a quaternion
trivialization of a different curve's exterior. The complement is
simply connected, by relative general position. All positive
homology groups not displayed in (P3K.29.1) vanish.

\item Write \(D_i\in H_3(E_C,\partial E_C;\mathbb Z)\) for the
Poincare--Lefschetz dual basis:
\[
                    m_j\cdot D_i=\delta_{ji}.
                                                        \tag{P3K.29.2}
\]
For completeness, its relative duality is also integral.
Small representatives of the three normal meridians
are pairwise disjoint, so the absolute middle
intersection homomorphism is zero. The relative
duality chain map of the same triangulation identifies
\(H_3(E_C,\partial E_C)\) with \(H^3(E_C)\), whose
evaluation on the three primitive meridians is the
full \(\mathbb Z^3\). Thus (P3K.29.2) defines its entire
relative group. The boundary sequence further gives
\[
 H_j(\partial E_C)=
 \begin{cases}
  \mathbb Z^2,&j=1,4,\\
  \mathbb Z^3,&j=2,3,
 \end{cases}
\]
with the degree-three and degree-four maps to \(E_C\)
isomorphisms. In particular no division by two occurs
at a triple point.

\item For any chosen \(C_i\), filling the other two curve tubes identifies
its single-curve exterior \(E_i\) with \(D^3\times S^3\), in the
marking already established in K23. Its map on absolute third
homology is \(m_j\mapsto\delta_{ij}m_i\). A relative three-disc
representing the \(D^3\) factor can be chosen to miss the other two
curves: first choose its boundary normal section away from their
four specified directions at the two vertices, then apply relative
general position in the interior, since \(3+2<6\).
Its class in the common exterior is \(D_i\). Small normal
meridians of the other two curves miss this compact disc, while
the marked meridian of \(C_i\) meets it once. This is the actual
map used below to compare a common-exterior variation with K23's
single-curve variation.

\par\medskip\noindent\textbf{2. The reciprocal-section transfer chain.}\enspace

\item Flop \(C_0\), and denote the three new curves by \(C'_0,C'_1,C'_2\),
where the latter two are the strict transforms. The common blowup
has exceptional divisor
\[
          E=\mathbb P^1_a\times\mathbb P^1_c,\qquad
          N_E=L=\mathcal O_E(-1,-1).
                                                        \tag{P3K.29.3}
\]
This follows from the actual blowup of the zero section in
\(\mathcal O(-1)\otimes\mathbb C^2\): \(a\) is its old base
coordinate, \(c\) is the projective direction of its two normal
coordinates, and the exceptional scalar is a vector in \(L\).
The two blowdowns are the two projections of (P3K.29.3).

\item The strict transforms of \(C_1,C_2\) meet \(E\) in four distinct
points. The two local coordinate axes at each endpoint of \(C_0\)
give its two distinct normal directions. Each of \(C_1,C_2\)
therefore supplies two of these points. Remove four disjoint
small complex coordinate balls from \(E\), and write \(E^\circ\)
for the resulting four-manifold.

\item Let \(T\) denote oriented circle transfer in \(S(L)|_{E^\circ}\),
with the fibre oriented by increasing complex argument of the
exceptional scalar. At a fixed \(a\), the old complex normal
meridian is exactly \(T[c]\): it is the unit sphere of the two
normal complex coordinates before blowup. At fixed \(c\), the new
complex normal meridian is exactly \(T[a]\), by the other blowdown.
Thus both signs follow from the actual complex normal four-discs.

\item Choose a smooth divisor \(D\) of bidegree \((1,1)\), avoiding the
four balls, as the zero locus of a transverse section \(\sigma\)
of \(L^\vee\). On the complement of a small tube about \(D\),
the normalized reciprocal \(\sigma^{-1}/|\sigma^{-1}|\) is an
actual section of \(S(L)\). Its graph is an oriented four-chain.
Along \(D\), let \(\sigma=re^{i\phi}\) in a positive local normal
coordinate. The boundary orientation of the complement of its
tube is
\(-\operatorname{or}(D)\wedge d\phi\). The reciprocal section
has fibre argument \(-\phi\), so this boundary maps to
\(+T[D]\). At each of the four deleted balls the boundary is
minus its positive three-sphere. Consequently the graph has
the literal oriented boundary
\[
           T[a]+T[c]-\mu_{00}-\mu_{01}-\mu_{10}-\mu_{11}.
                                                        \tag{P3K.29.4}
\]
Here \(T[D]=T[a]+T[c]\) follows by the homology transfer of
\([D]=[a]+[c]\) in \(E^\circ\).

\item Each small section lift \(\mu\) is the positive complex normal
meridian of the strict curve meeting that point. Indeed, in the
blowup chart that strict curve is transverse to \(E\); the two
coordinates on \(E\) are its two complex normal coordinates.
Fixing a nonzero exceptional scalar identifies its small normal
three-sphere with the indicated section lift. A change of unit
frame on that three-sphere has degree zero in \(S^1\), so it
does not alter this homology class. Transport along either
complex curve, away from its endpoints, preserves this positive
normal orientation. The two lifts belonging to each curve
therefore have the same sign.

\item Push (P3K.29.4) into the common exterior. It proves
\[
        m'_0=-m_0+2m_1+2m_2,\qquad
        m'_1=m_1,\qquad m'_2=m_2.                        \tag{P3K.29.5}
\]
The identification of common exteriors is the literal blowdown
isomorphism off the three curves; take the same compact
complement in that open set. The other two meridians can be
chosen away from both endpoint charts, where this map is the
identity. The new ambient orientation is the complex orientation,
so this open-set map is orientation preserving.

\item Thus the columns of the new absolute basis and of its relative
dual are, respectively,
\[
 M=
 \begin{pmatrix}-1&0&0\\2&1&0\\2&0&1\end{pmatrix},
 \qquad
 Q=M^{-T}=
 \begin{pmatrix}-1&2&2\\0&1&0\\0&0&1\end{pmatrix}.
                                                        \tag{P3K.29.6}
\]
In particular
\[
              D'_0=-D_0,\qquad
              D'_1=D_1+2D_0,\qquad
              D'_2=D_2+2D_0.                           \tag{P3K.29.7}
\]
All coefficients are integral, and \(M^2=I\). Relabelling
the selected edge gives the two other matrices recorded in
the exact certificate.

\item For a relative map (F) equal to the identity on the boundary,
its variation is the integral chain map
(operatorname{Var}_F(D)=F(D)-D).
Double the common exterior and extend (F) by the identity.
The meridians and doubled relative cycles have pairings
(m_i\cdot D_j=\delta_{ij}) and (D_j\cdot m_i=-\delta_{ij}).
The extended map fixes the meridians and sends
(D_j\mapsto D_j+\sum_i V_{ij}m_i). Preservation of
(D_i\cdot D_j) gives (-V_{ij}+V_{ji}=0), proving symmetry.
This applies to the closed actual ordinary products constructed
in \cref{night:P3K.71,night:P3K.73}; it makes no scalar
initialization from the static meridian map.

\item For stable frames related by (sigma'=sigma g), the exact
matrix identity is
[
 c_{sigma'}(F)(p)=g(Fp)^{-1}c_sigma(F)(p)g(p).
]
The primitive third class of Spin changes on a relative cycle (D)
by (-\langle g^*\eta,F(D)-D\rangle).
Indeed Spin is two-connected, so choose a cellular representative
of (g) constant on the two-skeleton. On each three-cell the
product obstruction is the sum of the three factor integers;
the boundary terms cancel because (F) is the identity there.
The remaining chain is exactly (F(D)-D), with the displayed minus
sign. This proves the frame-change formula integrally.

\item The actual outer attaching map has to be retained as well.
For full periods (h:(Z,V)\mapsto(\kappa qZ,\alpha V)) and
(B:(Z,V)\mapsto(gZ,q\eta V)), the second outer star has
(z_1=(g/\kappa)z_0) and normal coefficient map
[
 (N_1^-,N_1^+)=
 ((\eta/\alpha)N_0^-,\;\kappa\alpha/(g\eta)N_0^+).
]
Conjugating the negative-weight normal coordinate makes its
complex determinant ((\kappa/g)(\alpha/\overline\alpha)
(\overline\eta/\eta)), of parameter winding (-2J).
The simultaneous marked equivariant product equations would give
(-J=0), whereas the scalar family has (J\ne0).
These are the full coefficient arrows of \cref{night:P3K.61}.
The inner charged-cap calculation above therefore cannot be glued
with an unchanged outer product to deduce equality of the returns.
On the actual native (C_1) star \cref{night:P3K.62} gives the
free pair ((J,-J)). At (J=1) its right integer is (-1),
which directly refutes the old vanishing claim. The static cell
and reciprocal-section maps proved above remain the input maps
for the corrected marked-prism calculation.
\end{enumerate}
\end{proof}




\begin{theorem}\label{night:P3K.30}
\StatementRecord{night:P3K.30}{false as stated; correction proved}{S6\_06\_BLIND: GPT-6 Astra (OpenAI, ChatGPT Pro mode, effort Pro), recorded re-derivation in the checkpoint dated 18 November 2026}{\cref{night:P3F.9}}For every actual page of P3K.26 and P3K.28, let \(M\) be its page-cover degree,
\(h(G)=1/d\), and \(aG\) its marked line point. Write
\(H=Ma^2/d\) and \(\chi=\chi(\mathcal O_{S_{\rm page}})\).
Use the actual vertical-edge Cox normal product transported by the
extending signed cusp origin. For one positive loop in the actual
parameter \(r\), with \(\lambda=r^M\), the individual meridian-one dual pair and
the K23 relative coordinates are
\[
 \begin{aligned}
 k&=-3H,&\ell&=-24\chi,\\
 m_{\rm rel}&=6\chi,&n_{\rm rel}&=3H-6\chi,\\
 s&=H/2-\chi\pmod{28}.
 \end{aligned}
\]
This holds for both one-cap norms \(n=\pm1\) and both joint determinants
\(D=\pm1\). The comparison fixes the actual initial Cox quaternion
frame; a possible extending constant rotation loop changes neither
integer. No different fan edge's relative four-cycle is assigned here.

\end{theorem}
\begin{proof}
\begin{enumerate}[label=\arabic*.]
\item We retain the complete source, line, cap and frame maps below.
The meridian-one graph line is (U=M^{-1}): the fibre map
(D_M(v)(w)=\langle v,w\rangle/\langle v,v\rangle), with the
Hermitian form linear in its second argument, satisfies
(D_M(e^{i\theta}v)=e^{-i\theta}D_M(v)).
Its (SU(2))-equivariant spinor extension is
((a,b)\mapsto(\overline b,-\overline a)).
The first-to-second cap maps in \cref{night:P3K.68} give
(\lambda_M(F)=-\Sigma_{F^{-1}}) and
(\lambda_U(F)=\Sigma_{F^{-1}}); cancelling the opposite
product collars gives (A_M=-A_U). We therefore distinguish the
intermediate meridian-one number (A_U=(4k_U-\ell_U)/24)
from the physical return (A_M) throughout this calculation.

\item \par

\medskip\noindent\textbf{1. The two actual cusp products for the signed parameter}

\par\nobreak

\item Use the physical coordinates of K26 and K28:
\[
 V=e^{2\pi i\zeta_2},\qquad
 B_\lambda(Z,V)=(U(q)^{-1}Z,q\lambda W(q)^{-1}V),
\]
where \(U=e^{2\pi i g}\), \(W=e^{-2\pi i C}\) are the actual
nonvanishing cusp units. Thus the positive outward period is
\[
 \gamma:(Z,V)\longmapsto
      (U(q)Z,q^{-1}\lambda^{-1}W(q)V).
 \tag{K30.2}
\]
At fixed \(|r|\), variation of the phase of \(\lambda=r^M\)
is a physical delta rotation. It disappears in the circle orbit,
including in the two actual unimodular Cox charts. Consequently
the orbit body and its marked period-cell quotient are fixed
before applying the ordinary closing map.

\item Set \(s=\log q/(2\pi i)\) and \(v=\log r/(2\pi i)\).
The actual extending cusp origin is
\[
 \sigma_\varepsilon(v)
       =\varepsilon\left(gv,-Cv-\frac{M(v^2-v)}2\right).
 \tag{K30.3}
\]
The ordinary parameter deck is \(+\varepsilon s\delta\).
The identity
\[
 \sigma_\varepsilon(v+1)-\sigma_\varepsilon(v)
            =\varepsilon s\delta+\varepsilon\gamma
 \tag{K30.4}
\]
shows that the ordinary deck minus the cusp-origin difference
is the actual integral period \(-\varepsilon\gamma\).
This closes the cusp product, with the same positive \(r\)-loop
as in the construction proofs.

\item For a fixed initial parameter \(\lambda_0\), the bounded
coordinate map on its universal angular interval is literally
\[
 C_\theta=
 \left(
 U^{\varepsilon\theta},
 (W\lambda_0^{-1})^{\varepsilon\theta}
       e^{-\pi i\varepsilon M\theta(\theta-1)}
 \right).
 \tag{K30.5}
\]
All powers use the specified logarithms. These nonzero bounded
torus maps extend in the actual two adjacent cusp charts and
their Cox master coordinates. At \(\theta=1\) they are exactly
the bounded period-return map from (K30.4). The quadratic
factor is physical delta; it acts trivially on the orbit normal
three-plane. An additional unit in the global ordinary closing
map adds the same bounded origin term to both products and
is carried by the same extending torus map.

\item The actual ordinary origin is
\[
 o_\varepsilon(s,v)
       =\varepsilon\left(
          -\frac{gs}{M},
          \frac{Cs}{M}+\frac{s^2}{2M}+sv\right).
 \tag{K30.6}
\]
The source and parameter descents of these origins, including
the auxiliary \(2M\)-root gauge, were proved in K26 and K28.
Their literal transition has angular difference
\[
 (\sigma_\varepsilon-o_\varepsilon)(v+\theta)
       -(\sigma_\varepsilon-o_\varepsilon)(v)
   =\varepsilon\theta\gamma
       -\frac{\varepsilon M\theta(\theta-1)}2\,\delta .
 \tag{K30.7}
\]
On the circle orbit this is the marked rotation
\(R_{\varepsilon\theta}\) of the gamma product. Equivalently,
it is the positive rotation \(R^\beta_\theta\) for
\(\beta=\varepsilon\gamma\).

\par

\medskip\noindent\textbf{2. The exact smooth extension and Cox normal comparison}

\par\nobreak

\item We retain the actual map, rather than differentiate the annular
expression involving \(|q|^{\varepsilon\theta}\) at the node.
In the first small-resolution chart
\[
 (x,v,z),\qquad u=xz,\quad y=vz,
\]
the circle weights on \((v,z)\) are \((-1,+1)\), and its
Hopf coordinates are
\[
                (y,b),\qquad b=(|z|^2-|v|^2)/2.
\]
In the second chart \((u,y,w)\), with \(x=uw,v=yw\), use
\[
                (x,b'),\qquad b'=(|u|^2-|w|^2)/2.
\]
On the overlap, putting \(r=|z|^2\), the exact derivatives are
\[
 2b=r-|y|^2/r,\qquad
 2b'=|x|^2r-r^{-1},\qquad
 \frac{db'}{db}
   =\frac{|x|^2+r^{-2}}{1+|y|^2r^{-2}}>0.
 \tag{K30.8}
\]
The proper increasing orbit coordinate \(\chi_0(x,y,b)\)
of P3F.9, fixed relative to its marked outside collar,
has \((\chi_0)_b>0\). It gives the smooth product map
\(\Phi(x,y,b)=(x,y,\chi_0(x,y,b))\), with the corresponding
second-chart formula.

\item In the two actual extending products the orbit comparison is
\[
                       J_\theta=R_{\varepsilon\theta}\circ\Phi.
 \tag{K30.9}
\]
Equation (K30.7) proves agreement on the original annular
identification; (K30.8) proves the smooth extension through the
two node charts. The gamma rotation is an actual flow on the
fixed source-preserving product, so its negative direction for
\(\varepsilon=-1\) is equally defined.

\item Quotienting the derivative by the tangent of the fixed curve,
including the derivative of its graph gamma coordinate, gives
\[
 Q_x=
 \begin{pmatrix}
 1&0&0\\0&1&0\\
 (\chi_0)_{\Re y}&(\chi_0)_{\Im y}&(\chi_0)_b
 \end{pmatrix},
 \qquad(\chi_0)_b>0 .
 \tag{K30.10}
\]
The second chart has the same formula with \(x,b'\).
These maps are independent of \(\theta\). Their source shears
contract linearly, and the positive last entry contracts to one.
The intrinsic normal line is the free gamma direction; the
quotient is the source normal two-plane. Its splittings form
a contractible space, so the displayed contractions agree
globally. This remains true for either signed direction of
the transported rotation in (K30.9).

\item The sign is also retained explicitly in beta coordinates.
There the product map is
\(\Phi_\beta(x,y,b)=(x,y,\varepsilon\chi_0(x,y,b))\);
its raw normal derivative is
\[
 Q_x^\beta=T_\varepsilon Q_x,\qquad
 T_\varepsilon=\operatorname{diag}(1,1,\varepsilon).
 \tag{K30.10a}
\]
For the negative sign its last raw derivative is indeed negative.
Apply this same constant coordinate and orientation change to
both initial normal frames, as prescribed by the positive-Hopf
Thom orientation in K26 and K28. In the transported reference,
\(T_\varepsilon^{-1}Q_x^\beta=Q_x\), so the remaining comparison
is the positive triangular, theta-independent map just displayed.
Equivalently the full relative clutch is conjugated by
\(T_\varepsilon\); this is also conjugation by
\(\varepsilon T_\varepsilon\in SO(3)\).
This gives the actual placement of the signed coordinate change
at both ends. It contributes no free three-dimensional
framing class.

\item The initial actual Cox frame is the quaternion frame of K27,
which we retain explicitly. If the fixed-sphere charts satisfy
\(\zeta'=\zeta^{-1}\), \(c_i'=\zeta c_i\), then
\[
 q_{\rm H}
  =(1+j\zeta)(c_0+j\overline{c_1})
  =(\zeta^{-1}+j)(c_0'+j\overline{c_1'}).
 \tag{K30.11}
\]
The two products are literally equal, using \(zj=j\bar z\).
The physical delta action is right multiplication by
\(e^{-i\theta}\), so the positive Hopf normal coordinate is
\(q_{\rm H}i\overline{q_{\rm H}}\), with exactly the K23
orientation. Transport this initial frame by \(J_\theta\).
By (K30.9), it is the same frame obtained by applying
\(dR_{\varepsilon\theta}d\Phi\) to the initial Cox frame.
Equation (K30.10) supplies no parameter-dependent free
normal-framing class.

\item The bounded units in (K30.5) also retain their actual normal
map. In a Cox chart \(q=\zeta c_0c_1\), a unitary torus map
has form
\[
 (\zeta,c_0,c_1)\longmapsto
       (a\zeta,p c_0,q_1c_1),\qquad apq_1=1.
\]
Choose continuous square roots along the actual interval,
\[
             A^2=pq_1,\qquad B^2=p/q_1,\qquad AB=p.
\]
Substitution in (K30.11) gives the literal equality
\[
                         q_{\rm H}\longmapsto A q_{\rm H}B .
 \tag{K30.12}
\]
Its Hopf normal map is the constant rotation
\(\eta\mapsto A\eta\bar A\), independent of the fixed-sphere
point. The positive factors are the same normal isomorphism
on both sides of the comparison and cancel. The higher
base-unit terms have zero normal linear part, because \(q\)
is quadratic in \(c_0,c_1\). Thus any remaining loop is a
constant rotation loop. It extends over the fixed three-disc
and contributes neither \(k\) nor \(\ell\). There is no
free \(\pi_3(SO(3))\) discrepancy between the actual closed
Cox product and the transported rotating product.

\par

\medskip\noindent\textbf{3. The degree-\(M\) binding coordinate is an actual holomorphic chart}

\par\nobreak

\item For one cap the map in K26 has weighted polynomial
\(z_2-\tau z_1^M\), and its source-base coordinate is
\(t=z_2/z_1^M\). Its page ray in the binding coordinate
\(q_b=t-\tau\) has exactly
\[
                            \alpha=\phi-M\beta,
 \tag{K30.13}
\]
where \(\phi\) is the fixed outward page phase in radians.

\item For the two-cap page of K28, its exact ordinary phases give
\[
 \arg(z_2^{m_0}-\tau z_1^{m_\infty})
       =M\beta+\arg(t-\tau)-\varepsilon m_\infty n_0\arg t .
 \tag{K30.14}
\]
Near the regular binding \(t=\tau\ne0\), define the literal
holomorphic coordinate
\[
              q_b=(t-\tau)(t/\tau)^{-\varepsilon m_\infty n_0}.
 \tag{K30.15}
\]
Its derivative at \(\tau\) is one, so it is a local
biholomorphism. The exponent is an integer and the factor is
a specified nonvanishing holomorphic unit. Equations
(K30.14)--(K30.15), after retaining their constant \(\tau\)
phase in \(\phi\), prove exactly (K30.13).
This retains the actual two-cap source coordinate; it does
not replace a variable base phase by its limiting value.

\item The actual graph of K26 or K28 has
\[
                         \beta=\theta+g(p)
 \tag{K30.16}
\]
on the binding, with \(p\in T^2\) the source point and
\(\theta\) the positive physical parameter angle.
Here \(g:T^2\to\mathbb R/2\pi\mathbb Z\) is the graph phase
fixed by the source-preserving orbit product. Local radian
lifts suffice; \(R_z(-Mg(p))\) is a global map because \(M\)
is an integer. No source-dependent modification of the
page phase is made.

\par

\medskip\noindent\textbf{4. The full normal clutch and its Pontryagin number}

\par\nobreak

\item Give the graph its oriented page-normal vector and its two
Cartesian parameter vectors. The actual graph derivative
contributes only a contractible triangular shear to this
frame. On the compression choose a normal frame at parameter
zero and transport it by \(R_{\varepsilon\theta}\).
The compression is the same solid torus minus a ball used
in K6. Its tangent bundle is trivial, and
\(TY\oplus\mathbb R\) is trivial for the actual orbit
\(Y=S^5\). Hence its rank-three oriented normal bundle has
\(w_2=0\), and the three-dimensional obstruction groups
give a frame. At the fixed-sphere end its difference from
the specified positive Hopf frame is a map
\(S^2\to SO(3)\), which extends across the fixed three-disc.
Transport this actual extension by the same signed flow.

\item After quotienting by the source tangent directions and
the graph tangent, use the four coordinates consisting of
the real and imaginary \(q_b\) directions, the parameter
radius, and the relative angle \(\beta-\theta\).
The outward page collar and inward compression tangent
point in the same positive radial half-space. Their actual
common tangent is
\[
           u=(\cos\alpha,\sin\alpha,\rho,0),\qquad\rho>0.
\]
The literal quotient map is
\[
 (v_q,v_r,v_\eta)\longmapsto
       \left(v_q-\frac{v_r}{\rho}(\cos\alpha,\sin\alpha),
               v_\eta\right).
 \tag{K30.17}
\]
It kills \(u\). Project the Cartesian parameter vectors,
whose radial-relative coefficients are
\((\cos\theta,-\sin\theta)\), \((\sin\theta,\cos\theta)\);
reverse both and apply the positive reference rescalings.
The resulting actual oriented graph normal matrix is
\[
 G(\alpha,\theta)=
 \begin{pmatrix}
 -\sin\alpha&\cos\theta\cos\alpha&\sin\theta\cos\alpha\\
 \cos\alpha&\cos\theta\sin\alpha&\sin\theta\sin\alpha\\
 0&\sin\theta&-\cos\theta
 \end{pmatrix}\in SO(3).
 \tag{K30.18}
\]
The signed Thom calculation of K26 and K28 fixes the normal
orientation here as positive Hopf for both signs.
Equivalently, changing from the gamma relative coordinate
to the beta relative coordinate multiplies the reference
three-space by the constant matrix
\(T_\varepsilon=\operatorname{diag}(1,1,\varepsilon)\).
It conjugates both frame comparisons. This conjugation is
also conjugation by the oriented constant matrix
\(\varepsilon T_\varepsilon\in SO(3)\), so it introduces no
additional clutch or Pontryagin term.

\item The derivatives of \(g(p)\) are exactly the source tangent
shears already removed in (K30.17). The positive \(\rho\)
rescaling is contractible, and the actual collar rounding
stays in the same contractible positive radial half-space.
These specify a homotopy of the actual transition, including
all source tangent terms.

\item In the transported product coordinates the compression
frame is \(h(p)\), independent of \(\theta\). Since
\(G(\alpha,\theta)=R_z(\alpha)G(0,\theta)\), the full
clutch on the actual \(T^3\) is homotopic to
\[
 \begin{split}
 &h(p)^{-1}G(\phi-M(\theta+g(p)),\theta)\\
 &\quad=
  [h(p)^{-1}R_z(-Mg(p))]
  [R_z(\phi-M\theta)G(0,\theta)] .
 \end{split}
 \tag{K30.19}
\]
The first factor is pulled back from \(T^2\), the second
from the actual parameter circle. Both are globally
defined, including any nonzero one-dimensional frame
classes or graph-phase winding. Expanding
\[
 (ab)^{-1}d(ab)=b^{-1}(a^{-1}da)b+b^{-1}db
\]
and using the Maurer--Cartan identities shows that the
cubic trace form of the product differs from the sum of
the two cubic forms by an exact form. Each factor cubic
form vanishes by dimension. Therefore its integral over
the marked \(T^3\) is zero. The interpolating connection
\(t\,g^{-1}dg\) on the actual collar identifies this integral
with the normal Pontryagin clutching number, and gives
\[
                    \langle p_1(\bar\nu),[\bar X]\rangle=0.
 \tag{K30.20}
\]
This retains the full normal bundle and all lower-degree
frame data; it does not assert that the bundle is trivial.

\par

\medskip\noindent\textbf{5. The individual integers and the relative coordinates}

\par\nobreak

\item The already proved actual page-line comparison gives
\[
 b^2=4c_1(N)^2=-4H,\qquad
 p_1(T\bar X)=-24\chi .
 \tag{K30.21}
\]
The signed trace bordism has positive meridian integral
for each norm sign. Its sphere-bundle characteristic
formulas, in the matched normal framing just constructed,
are
\[
 k=\frac{p_1(\bar\nu)+3b^2}{4},\qquad
 \ell=p_1(T\bar X)+p_1(\bar\nu).
 \tag{K30.22}
\]
Substitution of (K30.20)--(K30.21) proves
\(k=-3H\), \(\ell=-24\chi\).
The actual line and height equality
\(H=2\chi+2(P\cdot O)\) makes both integers integral.
The full K23 coordinate morphism is
\[
 m_{\rm rel}=-k+6A_U=-\ell/4,\qquad
 n_{\rm rel}=-6A_U,\qquad s=-A_U\bmod28 .
 \tag{K30.23}
\]
Together with the independently proved
\(A_U=\chi-H/2\), this gives (K30.1).

\item The construction fixes the initial Cox quaternion product,
the signed cusp product, and its actual transport. The
possible constant rotation loop extends over the fixed
three-disc and changes neither number. The formula is for
one positive loop in the stated \(r\) parameter; it has not
been divided by the page degree, by a source-root order,
or by the number of scalar turns in \(\lambda\).

\par

\medskip\noindent\textbf{Certificate}

\par\nobreak

\item The exact source-phase and normal-matrix identities are
executed in \path{code/part3/general_cap_framing.py}. It also
combines the already certified K26 and K28 page arithmetic
with (K30.22)--(K30.23). The actual normal comparison is
the explicit map and clutch proof above.

\end{enumerate}
\end{proof}




\begin{theorem}\label{night:P3K.31}
\StatementRecord{night:P3K.31}{false as stated; correction proved}{S6\_06\_BLIND: GPT-6 Astra (OpenAI, ChatGPT Pro mode, effort Pro), recorded re-derivation in the checkpoint dated 18 November 2026}{\cref{night:P3F.35}, \cref{night:P3K.22}, \cref{night:P3K.23}, \cref{night:P3K.24}, \cref{night:P3K.25}, \cref{night:P3K.26}, \cref{night:P3K.27}, \cref{night:P3K.29}, \cref{night:P3K.30}, \cref{night:P3K.33}}
The actual single-star product frame differs from the ambient
stable frame by its displayed degree-one normal generator.
The literal common-link flop changes the ambient third frame
class by (2e_i), and the explicit orientation-preserving reset
((y,q)\mapsto(Jy,Jq)) swaps the free generators.
Consequently the native directed update is
((m,n,s)\mapsto(m+2n,-n,s+n)).
The historical universal norm polynomial is false: the actual
OS45 fans (T_0,T_2) both have physical cocharacter norm (1),
but their physical classes are (3,2).
\end{theorem}
\begin{proof}
\begin{enumerate}[label=\arabic*.]
\item \par\medskip\noindent\textbf{1. The three exact objects that must be transported.}\enspace

\item Keep K29's positive complex normal meridians \(m_i\) and relative
dual cycles \(m_j\cdot D_i=\delta_{ji}\). Let \(V\) denote the
variation matrix of the actual relative return
\(\operatorname{Var}_F(D)=F(D)-D\). Choose a global stable tangent
frame of the actual six-sphere and restrict it to the common exterior.
The derivative of \(F\), equal to the identity near its boundary,
lifts to the stable Spin group. Write
\[
 \delta_i=\langle c(F)^*\eta,D_i\rangle,\qquad
       \delta=(\delta_0,\delta_1,\delta_2)^T,
                                                        \tag{P3K.31.1}
\]
where the integral third generator \(\eta\) is normalized so that
the stable map \(q\mapsto L_q\), left quaternion multiplication,
has degree \(+1\).

\item The third object is the physical central return class
(s\in\mathbb Z/28). It is carried together with the actual
variation and derivative data; no diagonal initialization of the
full scalar plumbing is made here.

\item A different global stable frame gives a map from the whole six-sphere
to Spin. Its third class is zero because \(H^3(S^6)=0\).
K29's exact frame-change pairing therefore shows that (P3K.31.1) is
independent of that global choice at the degree being used.
No assertion about a higher framing homotopy group is needed.

\par\medskip\noindent\textbf{2. The first-resolution product differs by one normal generator.}\enspace

\item Use the actual first small resolution and the K23 quaternion product
\[
 u=a_0+j a_1,\qquad
 c=c_0+j\overline{c_1},\qquad
 q=uc,\qquad h=u i u^{-1}.
                                                        \tag{P3K.31.3}
\]
The positive complex meridian is the \(q\)-sphere. The base
\(h\)-sphere has the opposite of the complex curve orientation.
The exterior product is oriented \(D^3_h\times S^3_q\).
At a fixed base point, an outward radial vector of \(D^3_h\)
corresponds to the negative of the outward normal radial
vector \(q\) of the curve tube: the two are the collar
directions on opposite sides of their common boundary.

\item The derivative has a source-to-normal shear that is
retained before making the following frame comparison.
The inverse product map is literally \(c=u^{-1}q\).
At fixed \(q\), a source derivative therefore gives
\[
                   dc=-u^{-1}(du)u^{-1}q.
\]
Together with the source derivative and the radial
normal column, the full frame matrix has block form
\[
       \begin{pmatrix}
         B_2&0\\ A(q)&L_qR_4
       \end{pmatrix}.
\]
Here \(A(q)\) is precisely the normal component just
displayed, after the fixed source basis change.
Every matrix
\[
       \begin{pmatrix}
         B_2&0\\ tA(q)&L_qR_4
       \end{pmatrix},\qquad 0\le t\le1,
\]
is invertible and has the same positive determinant.
This is the actual continuous triangular homotopy
removing that derivative term. In the second product
the inverse is \(u=qc^{-1}\), so at fixed \(q\)
\(du=-qc^{-1}(dc)c^{-1}\); the identical block homotopy,
with the fixed reflection \(J\) in its normal block,
retains and removes its corresponding shear.

\item A fixed linear change of the three Cartesian \(h\)-vectors
therefore gives the two tangent vectors of the curve and
the vector \(-q\). The tangent curve pair has orientation
opposite to its complex orientation. The three left
quaternion vectors of the normal sphere are \(qi,qj,qk\).
Thus, in a local complex tangent frame of the ambient
six-manifold, the actual product frame along the meridian
has the form
\[
       \operatorname{diag}(B_2,L_q R_4),\qquad
       \det B_2=-1,\quad
       R_4=\operatorname{diag}(-1,1,1,1).
                                                        \tag{P3K.31.4}
\]
Both determinants are negative, so the total matrix is
orientation preserving. After multiplying by its constant
value at \(q=1\), its class is exactly
\(\operatorname{diag}(I_2,L_q)\).
The ambient stable frame extends over the small normal
four-disc, so its restriction has zero third class.
Consequently the product frame differs from the ambient
frame by degree \(+1\) on its positive complex meridian.
It has degree zero on the other two meridians, since its
frame extends over their normal discs in the selected
single-curve exterior.

\item Let \(g_i\) be this change from ambient to the selected
first-resolution product frame. Its absolute third
cohomology vector in the common exterior is exactly
\[
                              [g_i]=e_i.
                                                        \tag{P3K.31.5}
\]
By K29, the actual derivative in the product frame
differs from the ambient derivative by
\[
           -\langle[g_i],\operatorname{Var}_F(D_i)\rangle=-V_{ii}.
\]
The ordered product disc and the relative dual
disc have opposite signs:
\[
              D_i=-[D^3_h],\qquad
              m_i\cdot[D^3_h]=-1.
                                                        \tag{P3K.31.6}
\]
K23's derivative on \([D^3_h]\) is \(-2n_i\). Equations
(P3K.31.5)--(P3K.31.6) give the exact first-resolution formula
\[
                    2n_i=\delta_i-V_{ii},\qquad
                    m_i^{\rm rel}=-V_{ii}-n_i.
                                                        \tag{P3K.31.7}
\]
Here \(m_i^{\rm rel}\) is a relative group integer, not a
normal meridian. In the initial fan K29 makes all three
quantities zero for \(i=1,2\), and the fixed edge has
\(\delta_0=v+2n=n-m\), proving (P3K.31.2).

\par\medskip\noindent\textbf{3. The literal ambient-frame change under the flop is \(+2e_i\).}\enspace

\item First keep the unchanged \(q\) coordinate of the actual
two small resolutions. The second product has coordinates
\[
             k=c^{-1}ic,\qquad q=uc,
\]
and the common boundary map is
\[
                 T(h,q)=(q^{-1}h q,q).
                                                        \tag{P3K.31.8}
\]
K23 proves that this is the literal actual-family map,
independent of the parameter angle, including odd Cox
loops and the transported bounded-unit truncation collars.
Its displayed extension over \(D^3\times S^3\) gives a
normal collar representative of the same boundary map.

\item The derivative in Cartesian base and left quaternion
normal frames is triangular. Its base block is
\(\operatorname{Ad}_{q^{-1}}\), its normal tangent block
is the identity, and its off-diagonal block contracts
linearly to zero. Thus its stable third class on the
old positive \(q\)-meridian is
\[
                         [c_T]=-2.
                                                        \tag{P3K.31.9}
\]
This sign and factor follow from the actual quaternion
action: \(\operatorname{Ad}_q=L_qR_{q^{-1}}\) has stable
degree \(+2\) in the normalization (P3K.31.1), while inversion
of \(q\) reverses its degree. Equivalently the marked
adjoint clutch has \(p_1=-4c_2\), with the same direct
orientation calculation as K23; the left quaternion
clutch is its degree-one stable normal generator.

\item In the second product, the \(k\)-sphere preserves the
complex curve orientation and the unchanged \(q\)
coordinate reverses the complex normal orientation.
Let \(J\) be the fixed reflection of that normal real
four-plane. The analogue of (P3K.31.4), parametrized by
the unchanged \(q\)-sphere, is
\(\operatorname{diag}(B_2',J L_qR_4)\), with
\(\det B_2'=+1\).
After its constant value is removed its stable class
is \(J L_qJ^{-1}\), still degree \(+1\).
Indeed conjugation by \(J\) becomes conjugation by
\(\operatorname{diag}(J,-1)\) after one stabilization;
the latter lies in the connected special orthogonal
group and hence induces the identity on the third
homotopy group. The new positive complex meridian is
\(-[S^3_q]\), so this same product-frame difference has
degree \(-1\) on the new complex meridian, exactly as
required by the old/new normal orientations.

\item Write \(P_o,P_n\) for the two product frames and
\(\sigma_o,\sigma_n\) for the two global ambient stable
frames. Put \(P_o=\sigma_og_o\), \(P_n=\sigma_ng_n\).
Pulling the new frame back by the actual map gives
\[
 (dT)^{-1}\sigma_n
       =\sigma_o\,g_o\,c_T^{-1}\,(g_n\circ T)^{-1}.
                                                        \tag{P3K.31.10}
\]
On the old positive complex meridian the three displayed
degrees are \(+1\), \(-(-2)\), and \(-1\), respectively.
The exact ambient-frame change is therefore
\[
                         [G]=1+2-1=2.
                                                        \tag{P3K.31.11}
\]
On either other curve's normal meridian the birational
map is a smooth complex isomorphism on the whole normal
four-disc, so the change has degree zero. Hence for a
flop of edge \(i\), its absolute common-exterior third
class is exactly \(2e_i\). This computes the relative
normal Thom correction; it is not chosen from the
vanishing of an absolute characteristic class.

\par\medskip\noindent\textbf{4. Resetting the first-resolution complex convention.}\enspace

\item The unchanged-\(q\) second product in (P3K.31.8) is the one
used in K23's conjugation formula
\[
             (m,n,s)\longmapsto(-n,m+2n,s+n).
                                                        \tag{P3K.31.12}
\]
For the next selected edge we must return to the first
product whose normal quaternion preserves its complex
orientation. This reset has an explicit map.

\item Write a quaternion as \(x_0+ix_1+jx_2+kx_3\), and set
\[
 J(x_0+ix_1+jx_2+kx_3)=x_0+ix_1+jx_2-kx_3.
                                                        \tag{P3K.31.13}
\]
It is the anti-automorphism \(J(q)=k\bar q k^{-1}\).
For the new first resolution its base quaternion is
\(J(c)=c_0+jc_1\), its normal quaternion is
\(J(u)=a_0+j\bar a_1\), and their product is
\[
                      q_{\rm new}=J(c)J(u)=J(q).
                                                        \tag{P3K.31.14}
\]
Its Hopf base is \(J(k)\). Thus the exact reset is
\[
                       U(y,q)=(J(y),J(q)).
                                                        \tag{P3K.31.15}
\]
The base three-reflection has determinant \(-1\) on
\(\operatorname{Im}\mathbb H\), and the normal
four-reflection has determinant \(-1\) on
\(\mathbb H\). Their product is \(+1\), so the map
preserves both the six-dimensional exterior and
the seven-dimensional filling orientation.
The map extends over \(D^3\times D^4\).
For K23's conjugation-equivariant radial \(\zeta\),
\(J\zeta(Jy)=\zeta(y)\), so literal substitution gives
\[
                    ULU^{-1}=R,\qquad URU^{-1}=L.
                                                        \tag{P3K.31.16}
\]
Its orientation-preserving extension fixes the central
disc-supported class. Hence the selected first-convention
pair after the flop is
\[
                        (m+2n,-n),\qquad s'=s+n.
                                                        \tag{P3K.31.17}
\]
In particular the reverse flop changes the central class
by \(-n\), and restores the original class.
This reset also retains the actual parameter family.
For the Cox product action \(q\mapsto a_\theta q b_\theta\),
the same literal anti-automorphism gives
\(q_{\rm new}\mapsto J(b_\theta)q_{\rm new}J(a_\theta)\).
The corresponding base action follows by substituting
in its Hopf coordinate. The map \(U\) is independent of
\(\theta\); it therefore intertwines the two existing
family products, including their endpoint signs, and
does not add a parameter-dependent frame loop.


\item The maps just proved are local comparison maps with their exact
frames and source/target markings. The remaining ordinary-product
transport is supplied, with all outer units, by
\cref{night:P3K.71,night:P3K.73}. For the counterexample use OS45
(J=1,H=8,\chi=1) and the literal vertical changes
(Z_n=Z,V_n=Z^{-n}V,h_n=hB^n,B_n=B).
The physical cocharacter is still (e_2), so its coordinates are
((0,1)) and the proposed norm equals (1) for both (n=0,2).
The actual class (3-n(n-1)/2) is (3) and (2), respectively.
Thus no function of that norm alone gives these scalar returns;
the rejected recurrence and its residue table are preserved only
in the historical audit file.

\end{enumerate}
\end{proof}




\begin{theorem}\label{night:P3K.32}
\StatementRecord{night:P3K.32}{false as stated; correction proved}{S6\_06\_BLIND: GPT-6 Astra (OpenAI, ChatGPT Pro mode, effort Pro), recorded re-derivation in the checkpoint dated 18 November 2026}{\cref{night:P3F.35}, \cref{night:P3F.36}, \cref{night:P3F.38}, \cref{night:P3F.40}, \cref{night:P3F.41}, \cref{night:P3F.44}, \cref{night:P3F.45}, \cref{night:P3F.48}, \cref{night:P3F.49}, \cref{night:P3F.9}, \cref{night:P3K.15}, \cref{night:P3K.22}, \cref{night:P3K.26}, \cref{night:P3K.27}, \cref{night:P3K.28}, \cref{night:P3K.30}, \cref{night:P3K.5}, \cref{night:P3K.7}}Consider the actual source and root-cap data of P3F.36, P3F.38,
P3F.41--43, P3F.46--47, and the primitive-source data of P3K.22.
Retain any finite collection of fixed order-one elliptic, gamma,
and delta clutches at additional smooth or multiplicative source
values, including their annular units and finite line characters.
Put the surviving core values at $\infty,0$; a missing core is
an ordinary disc of order one and norm zero. Write their orders
and signed norms as $(m_\infty,n_\infty),(m_0,n_0)$, the added
signed norms as $n_j$, and the retained base relation as
$\sum\alpha_i+b_0\gamma=0$. Define
\[
 N=\sum_jn_j-b_0,\qquad
 n_\infty'=n_\infty+m_\infty N,\qquad
 D=m_0n_\infty'+m_\infty n_0.                         \tag{K32.1}
\]
The sphere locus is exactly $D=\pm1$. This topology conclusion
holds for every allowed cusp fan.

On that locus put $M=m_\infty m_0$, and use the actual
parameter $\lambda=r^M$. The vertical-edge reference member
has corrected compact page line $N_{\rm page}$, with
\[
 H=-c_1(N_{\rm page})^2=\frac{Ma^2}{d},\qquad
 \chi=\chi(\mathcal O_{S_{\rm page}}),\qquad h(G)=1/d.
\]
Its full characteristic and relative data are
\[
 \begin{aligned}
 k&=-3H,&\ell&=-24\chi,& A_U&=\chi-H/2,\\
 m_{\rm rel}&=6\chi,&n_{\rm rel}&=3H-6\chi,
 &s&=-A_U=H/2-\chi\pmod {28}.                           
 \end{aligned}\tag{K32.2}
\]
Thus the single- and joint-cap tables of P3K.26 and P3K.28
apply with the actual norm $n_\infty'$ and all their root-order
restrictions retained. In the primitive sector without a
distinguished additive cap, the page is the original rational
elliptic source and
\[
 (k,\ell)=(-3d,-24),\qquad
 (m_{\rm rel},n_{\rm rel})=(6,3d-6),\qquad
 A_M=d/2-1.                                   \tag{K32.3}
\]
The return data in these formulas use the reference fan;
the complete marked-prism transport of \cref{night:P3K.73} computes the other whole-zero scalar fans. The individual reference fixed-star product alone supplies no simultaneous outer plumbing product.
\end{theorem}

\begin{proof}
\begin{enumerate}[label=\arabic*.]
\item We retain the complete source, line, cap and frame maps below.
The meridian-one graph line is (U=M^{-1}): the fibre map
(D_M(v)(w)=\langle v,w\rangle/\langle v,v\rangle), with the
Hermitian form linear in its second argument, satisfies
(D_M(e^{i\theta}v)=e^{-i\theta}D_M(v)).
Its (SU(2))-equivariant spinor extension is
((a,b)\mapsto(\overline b,-\overline a)).
The first-to-second cap maps in \cref{night:P3K.68} give
(\lambda_M(F)=-\Sigma_{F^{-1}}) and
(\lambda_U(F)=\Sigma_{F^{-1}}); cancelling the opposite
product collars gives (A_M=-A_U). We therefore distinguish the
intermediate meridian-one number (A_U=(4k_U-\ell_U)/24)
from the physical return (A_M) throughout this calculation.

\item \par\medskip\noindent\textbf{The holomorphic source cochain.}
We first construct the source arrows for all retained elliptic
clutches. Let $f:S\to\mathbb P^1$ be the marked rational
relatively minimal elliptic source. Choose disjoint small
closed discs at the clutch sites, slightly enlarge them to
$U_i$, and let $U_0$ be the complement of smaller closed discs.
The overlaps $U_0\cap U_i$ are annuli with smooth fibres.
Let $u_{0i}$ be the actual holomorphic translation section
specified by the elliptic part of the clutch. The gamma
component is retained separately.

\item At a smooth site with period frame $\Pi(q)$ and marked
elliptic period $\lambda$, this section is literally
\[
       u_{0i}(q)=\Pi(q)\lambda\frac{\log q}{2\pi i}
                         \pmod {\Lambda(q)}.
\]
Changing a logarithm branch adds its displayed integral
period, so it descends. At a singular site we use the actual
annular descent law. For example, on a Tate annulus set
$s=\log q/(2\pi i)$ and $\tau=Ns+\tau_0(q)$. For the
marked elliptic integers $(c,j)$ the exact section is
\[
 \begin{split}
 u(s)&=\frac{jN}{2}s(s-1)+j\tau_0(q)s+cs,\\
 u(s+1)-u(s)&=j\tau(q)+c.                             
 \end{split}\tag{K32.4}
\]
In a line frame with $1$-period character one and
$\tau$-period character $P(q)$, the multiplier
$\exp(j\log(P(q))s)$ has quotient $P(q)^j$ under
$s\mapsto s+1$ and compensates that lattice shift. The
independent delta multiplier is retained.

\item Glue the local source surfaces by
$p_0=t_{u_{0i}}p_i$. This gives a smooth complex surface
$S_u\to\mathbb P^1$. It is Hausdorff: points over distinct
base values are separated on the base, while points over
one value lie in a common original proper local model.
Properness is local on the target and holds on each such
model. Hence $S_u$ is compact. Its fibres and their local
analytic models are the original ones, with no multiple
fibres.

\item On a local minimal model choose its invariant dualizing
differential $\eta$ and base coordinate $q$. The local
two-form $dq\wedge\eta$ is nowhere zero. A base-dependent
translation satisfies $t_u^*\eta=\eta+h\,dq$ on smooth
fibres, and therefore
\[
             t_u^*(dq\wedge\eta)=dq\wedge\eta.
\]
Both sides extend across the local model. The canonical
transition functions of the glued surface are consequently
the original ones, giving the actual line isomorphism
\[
                    \omega_{S_u}=f_u^*\mathcal O(-1).
                                                               \tag{K32.5}
\]
Translation acts identically on $H^1(E,\mathcal O_E)$:
its Serre-dual action fixes the invariant differential.
The local proper-direct-image maps therefore also glue to
\[
 f_{u*}\mathcal O_{S_u}=\mathcal O_{\mathbb P^1},\qquad
 R^1f_{u*}\mathcal O_{S_u}=\mathcal O_{\mathbb P^1}(-1).
                                                               \tag{K32.6}
\]
Leray gives $q(S_u)=p_g(S_u)=0$ and $\chi(\mathcal O_{S_u})=1$.
Equation (K32.5) makes every positive plurigenus zero.
The compact-surface classification, with the precise alternatives
in \cite[VI.1.1 and Table~10, pp.~243--244]{P4_BHPV2004},
makes $S_u$ rational: class VII has holomorphic Euler
characteristic zero, and a ruled surface over a positive-genus
curve has positive irregularity. Since $K_{S_u}^2=0$,
Noether's formula gives $e(S_u)=12$. A minimal rational
surface has Euler number three or four, so $S_u$ contains
a smooth rational $(-1)$-curve $E$. Adjunction and (K32.5)
give $F\cdot E=1$. Thus $f_u|_E$ has degree one and is
an isomorphism to $\mathbb P^1$. We have constructed a
holomorphic section $\sigma$.

\item Let $s_i$ be this section in the original local source chart.
The actual gluing equation gives
$s_0=s_i+u_{0i}$. The maps
\[
                 \Psi_i(p_i)=p_i-s_i(f(p_i))            \tag{K32.7}
\]
therefore agree on the glued overlaps and give an isomorphism
$\Psi:S_u\to S$ over the base.
We record its extension over singular fibres. A section
meets the smooth multiplicity-one part of a Kodaira fibre.
Resolve the graph of its meromorphic translation with
projections $p,q:Z\to S$ over a small disc. For the local
nowhere-zero two-form $\Omega$ above, $p^*\Omega=q^*\Omega$.
Each point blowup gives a positive canonical discrepancy,
so these zero divisors have support exactly the exceptional
curves of the respective projections. Their supports are
equal. Each connected fibre of $p$ is consequently sent to
a point by $q$. Normal descent through $p_*\mathcal O_Z
=\mathcal O_S$ factors $q$ holomorphically through $p$.
The same argument for the inverse proves that (K32.7)
extends as an isomorphism on the entire local minimal model.

\item These translations commute with $t_{-Q}$. Their pullbacks
to a root are equivariant for its canonical deck, so they
also commute with the affine source norm maps there.
The same graph argument gives their minimal-root extensions.
The cochain is chosen once, independently of $r$, and fixes
the actual source gamma generator.

\item Keep now the actual line lifts in these clutches. They glue
to a line $\mathcal M_u$ on $S_u$. Transport it by $\Psi$
and call the result $\widetilde{\mathcal M}$. Translation
preserves the degree-zero Picard point on every smooth fibre,
so $Q_{\rm line}=\widetilde{\mathcal M}\otimes\mathcal M^{-1}$
is trivial there. The proper direct image
$L=f_*Q_{\rm line}$ is torsion-free of rank one on the
smooth base, hence is a line bundle. Evaluation is an
isomorphism on every smooth fibre. Its divisor is effective
and vertical, giving the exact isomorphism
\[
 \widetilde{\mathcal M}
   =\mathcal M\otimes f^*L\otimes\mathcal O(D_{\rm vert}).
                                                               \tag{K32.8}
\]
We retain this evaluation map, its base transition units,
and every coefficient of $D_{\rm vert}$. On each minimal
root, the multidegree-zero line comparison of P3K.15 gives
the actual lift of (K32.7). Its generic constant-factor
uniformization makes it commute with the suspension lift;
extension gives the equality on the root. The finite
deck character is retained by P3F.44 and P3F.49.

\par\medskip\noindent\textbf{The exact map on the orbit base.}
We next handle gamma on the three-dimensional base. Put
the added sites at $c_j$ and write
$\alpha_j=\arg(t-c_j)/(2\pi)$. The actual order-one
cap coordinate is
\[
                      h_j=\gamma-n_j\alpha_j.           \tag{K32.9}
\]
It gives the physical images $\alpha_j\mapsto-n_jh_j$,
$\gamma\mapsto h_j$ and killed meridian
$\alpha_j+n_j\gamma$.
The retained framing $\sum\alpha_i+b_0\gamma=0$ has a
literal base-atlas description: introduce an ordinary
disc with transition $h=\gamma+b_0\alpha$. This chart
has order one and signed norm $-b_0$. Eliminating its
meridian gives the displayed framing relation. This
identifies the marked circle-base atlases without adding
an elliptic or line clutch on the total space.

\item First use the ordinary annular atlas with $b_0=0$. Set
\[
 u(t)=\prod_j(t-c_j)^{n_j},\qquad
 \gamma'=\gamma-\frac{\arg u(t)}{2\pi}.                 \tag{K32.10}
\]
The phase is circle valued and well defined. Near $c_j$
write $u=(t-c_j)^{n_j}v_j$, where $v_j$ is a nonzero
holomorphic unit. Then
$\gamma'=h_j-\arg(v_j)/(2\pi)$ extends across the actual
solid torus, with inverse given by adding that same unit
phase. At infinity, $q=1/t$ gives $u=q^{-N}v_\infty(q)$,
so
\[
 \gamma'=\gamma+N\alpha_\infty
                         -\frac{\arg v_\infty}{2\pi}.
\]
Retain the unit as an extending change of cap coordinate.
On the marked loops this is
\[
 \begin{split}
 \alpha_\infty'&=\alpha_\infty-N\gamma,\\
 m_\infty\alpha_\infty+n_\infty\gamma
  &=m_\infty\alpha_\infty'+n_\infty'\gamma.             
 \end{split}\tag{K32.11}
\]
At zero $u$ is an extending unit and transports its
marking by the same formula. Apply this calculation
also to the framing chart with exponent $-b_0$.
It gives exactly $N=\sum n_j-b_0$ in (K32.1).
These local maps and inverses give an actual diffeomorphism
of the filled three-dimensional bases to the two-cap base
with norms $(n_\infty',n_0)$.

\item There is also a literal local generator identity:
\[
 A_\infty'=A_\infty B^N,\qquad
 (A_\infty')^{m_\infty}=B^{n_\infty'},\qquad
 \langle A_\infty',B\rangle=\langle A_\infty,B\rangle.
                                                               \tag{K32.12}
\]
It holds in the actual family:
$B_r=r^LB_0$ and
$A_{\infty,r}=r^{Ln_\infty/m_\infty}A_{\infty,0}$
give the new exponent $Ln_\infty'/m_\infty$.
Thus the local quotient and all its finite characters
are preserved. We use (K32.10) only on the base;
the source and line maps on its page are constructed below.

\item Eliminating the order-one base relations gives
$\alpha_0+\alpha_\infty+(b_0-\sum n_j)\gamma=0$,
which also verifies the sign in (K32.1).
The base abelianization has determinant $D$. It is a
quotient of the orbit group, so every sphere member has
$D=\varepsilon\in\{1,-1\}$.

\par\medskip\noindent\textbf{Integral topology and sufficiency.}
We establish the converse before using the orbit as a
five-sphere. Let $W_0$ be the actual source exterior with
only its chosen core tubes removed. Remove the additional
ordinary tubes $U_j$ to obtain $W$, and put
$E_j=\partial U_j$. First use their original ordinary
fillings. The maps $H_2(E_j)\to H_2(U_j)$ are onto, with
primitive kernels $K_j$: the three smooth rims, or the
two multiplicative classes $R_j,L_j$ of P3F.40.
Every old degree-three representative avoids the new
tubes. Old source-boundary and local-rim representatives
remain in old collars, and the common generic-fibre
trace lies over a regular value outside them. Hence
$H_3(W)\to H_3(W_0)$ is onto.

\item The next ordinary boundary map from $H_2(W_0)$ lands
in the free meridian kernels in $\bigoplus H_1(E_j)$.
That target is inversion-even. The actual source
involution is inversion-odd on $H_2(W_0;\mathbb Q)$,
by the source calculations in P3F.45 and P3F.48.
Thus the rational map is zero, and a homomorphism into
this free target has no nonzero torsion image. The map
is zero integrally. Ordinary Mayer--Vietoris, using
the cap surjections just displayed, gives
\[
 0\longrightarrow K=\bigoplus_jK_j
  \longrightarrow H_2(W)
  \longrightarrow H_2(W_0)\longrightarrow0.             \tag{K32.13}
\]
The involution is odd on $K$ and on the last group
rationally. An involution over $\mathbb Q$ is semisimple,
so it is odd on $H_2(W;\mathbb Q)$ as well. Therefore
every even boundary core column has torsion image in
$H_2(W)$. This argument has not divided by two.

\item Now use the actual old core cap maps. Their combined
degree-two map projects isomorphically onto
$H_2(W_0)\oplus H_2(\text{old caps})$, by the complete
maps in P3F.36, P3F.38, P3F.41--47, or by P3F.45 and
P3K.22 for the regular primitive core. Its image is
therefore an exact graph over this whole target, with
zero intersection with $K$. Quotienting by that image
leaves precisely the free group $K$ and the new cap
targets. All the preceding torsion core columns vanish
there. This proves the integral assertion also at
primes dividing the weighted source integer; no
involution has been extended over a replacement cap.

\item For a smooth order-one cap retain its actual meridian
\[
          \mu=\alpha+n\gamma+\lambda_a a+\lambda_b b.
\]
The cap map is the exterior square of the actual
primitive quotient on $H_1$. After eliminating its
target, use the kernel columns
\[
 (a\wedge\mu,\ b\wedge\mu,\ \mu\wedge\gamma,
                                      \alpha\wedge\mu)
\]
and the ordered rims
$(a\wedge\alpha,b\wedge\alpha,\alpha\wedge\gamma)$.
The exact map is
\[
 \begin{pmatrix}
 1&0&0&-\lambda_a\\
 0&1&0&-\lambda_b\\
 0&0&1&n
 \end{pmatrix}.                                      \tag{K32.14}
\]
Its primitive relation is
$\lambda_a\mathrm{col}_1+\lambda_b\mathrm{col}_2
-n\mathrm{col}_3+\mathrm{col}_4=0$.
Consequently its first three columns form an actual
integral basis of the kernel and map to the rim basis.

\item For a multiplicative order-one cap retain the Tate
integer $A$ and all the P3F.40 coefficients. Its
isotropy relation and actual target basis change give
the cap rows
\[
                 \begin{pmatrix}1&-A&0&0\\0&0&1&A\end{pmatrix}
\]
in the marked boundary order $(z,R,B,L)$. Together
with the exterior $R,L$ rows this reduces integrally
to $\operatorname{diag}(1,1,-1,-1)$, exactly as in
P3F.48 with effective order one. The $K_j$ are
independent summands of (K32.13), so the complete
degree-two Mayer--Vietoris map is an isomorphism.
All elliptic and delta collar coefficients have been
retained in these maps.

\item We next compute the fundamental group. The holomorphic
cochain (K32.7) identifies the elliptic-only clutches
with the same source, commutes with its suspension,
and transports the root diagrams. Equivalently, the
actual van Kampen maps keep every local vanishing
cycle at an order-one cap and change its meridian
only by the specified elliptic and gamma terms.
Translations act identically on the fibre lattice.
Thus every vanishing relation and its conjugates by
the marked linear monodromy are exactly the source
relations of P3F.36, P3F.38, P3F.45--47. Their normal
closure kills the entire elliptic lattice, including
the graph generator of the quadratic $I_2$ root.
Adding an elliptic term to a meridian does not change
that normal closure: conjugation by a fibre translation
is trivial on the abelian fibre lattice. The actual
finite component maps at the old caps remain their
proved isomorphisms.

\item After these exact source identifications the remaining
group is the marked Seifert group. Eliminate each
added order-one meridian and use (K32.12); its
presentation is
\[
 \left\langle h_\infty,h_0\ \middle|\
 h_\infty^{m_\infty}=h_0^{m_0},\quad
 h_\infty^{n_\infty'}h_0^{n_0}=1\right\rangle.          \tag{K32.15}
\]
The common power is $\gamma$. The original root norm
is primitive, and
$\gcd(m_\infty,n_\infty')=
\gcd(m_\infty,n_\infty)=1$. Choose integers $a,c$
with $am_\infty+cn_\infty'=1$. Then
\[
          h_\infty=h_0^{am_0-cn_0},
\]
so (K32.15) is cyclic, of order $|D|$; at $D=0$
it is infinite cyclic. In particular it is trivial
at $D=\pm1$.

\item The degree-one maps retain the same source and finite
component identifications. Their free marked relation
matrix has rows
\[
 \begin{gathered}
 m_\infty\alpha_\infty+n_\infty\gamma,\qquad
 m_0\alpha_0+n_0\gamma,\qquad
 \alpha_j+n_j\gamma,\qquad
 \sum\alpha_i+b_0\gamma.
 \end{gathered}                                      \tag{K32.16}
\]
The column changes
$\alpha_\infty=\alpha_\infty'+N\gamma$,
$\alpha_j=\alpha_j'-n_j\gamma$, followed by subtracting
the order-one relation rows from the last row, leave
the two-core matrix and identity blocks. Both changes
are unimodular, and the determinant is $-D$.
The full matrices are exported by the certificate.

\item We also verify injectivity directly, retaining possible
torsion rather than relying only on that free matrix.
As reference use the already proved core sphere with
norm $n_\infty'$ and all new tubes filled ordinarily.
Its exterior is the same $W$. Each new order-one cap
has the same individual quotient as its ordinary cap;
(K32.12) gives the identical old local quotient.
Thus the domain and target of our degree-one map have
the same free ranks as that reference, and those ranks
are equal. Its cokernel is zero by (K32.15), so its
kernel has rank zero. The only torsion in its domain
is in old boundary blocks; their maps into the
corresponding cap torsion groups are the unchanged
actual isomorphisms of the cap proofs. New smooth
boundaries are free, and transitive multiplicative
boundaries are free because their Euler pair is
primitive. The torsion kernel is consequently zero.
This proves the full integral degree-one isomorphism.

\item Mayer--Vietoris now gives $H_1(Y)=H_2(Y)=0$, and
(K32.15) gives $\pi_1(Y)=0$. Integral Poincare duality
and the universal coefficient theorem give
$H_3(Y)=H_4(Y)=0$. The connected closed oriented orbit
is a simply connected integral homology five-sphere;
the smooth five-dimensional sphere theorem already
used in the construction identifies it with $S^5$.
Apply the actual fixed-sphere Hopf gluing and circle
Gysin calculation of P3F.36, P3F.38 and P3K.22, with
the retained boundary line. The principal-circle part,
primitive normal Hopf charge, and fixed-sphere tube
are carried by the maps above. They give a simply
connected integral homology six-sphere, hence $S^6$
by the proved vanishing of the smooth homotopy
six-sphere group. P3F.35 transfers this topology
through the actual finite flop chain to every
allowed fan. This proves the asserted sufficiency
before any normal-frame argument uses $Y=S^5$.

\par\medskip\noindent\textbf{The physical parameter cochain and page.}
We now work on the sphere locus $D=\varepsilon=\pm1$.
The orders are coprime, so $M=m_\infty m_0$ is their
least common multiple and the parameter is
$\lambda=r^M$. The fixed elliptic and delta clutch
coefficients are independent of $r$. The gamma part
has the exact ordinary-to-cap scalar variation
$n_i(c-c_0)s_i\delta/m_i$, where
$c=\log\lambda/(2\pi i)$ and $s_i=\log q_i/(2\pi i)$.
For one positive $r$-loop put
\[
 k_\infty=m_0n_\infty,\quad
 k_0=m_\infty n_0,\quad k_j=Mn_j.
\]
In the zero-framing atlas the actual inverse closing
multiplier and its redistributed counterpart are
\[
 \begin{split}
 f_{\rm old}(t)
   &=t^{k_0}\prod_j(t-c_j)^{Mn_j}(t-t_c)^{-D},\\
 f_{\rm new}(t)&=t^{k_0}(t-t_c)^{-D},\qquad
 \frac{f_{\rm old}}{f_{\rm new}}=u(t)^M.                
 \end{split}\tag{K32.17}
\]
Their orders at all capped values, including infinity,
are exactly the marked integers. For the retained
framing include exponent $-Mb_0$ on its ordinary
chart. This gives the same identity with the $u$
of that extended atlas.

\item Choose the selected cusp and the regular binding away
from the additional sites. There $u$ is a unit.
Add its local logarithm to both compared cusp origins,
retaining it in each. It cancels in their difference.
With $v=\log r/(2\pi i)$, the actual K28 origin map
therefore remains
\[
 F_\varepsilon(v+\theta)-F_\varepsilon(v)
   =\varepsilon\theta\gamma
        -\frac{\varepsilon M\theta(\theta-1)}2\delta.
                                                               \tag{K32.18}
\]
The quadratic gauge is zero at both endpoints.
The weighted base angle is $\beta=\varepsilon\gamma'$;
its vector field is the same as that of
$\varepsilon\gamma$, since (K32.10) changes only its
origin by a fixed base function. Thus the physical
positive return has the signed angular coordinate
used in P3K.28. All other fixed clutch maps glue the
product comparison over the parameter interval with
the same endpoint maps. Equation (K32.17) records
the entire additional parameter closing term.

\item Apply the explicit weighted three-sphere map of
P3K.28 to the redistributed orbit base. Its regular
page has complex base
\[
 \mathcal C^\circ:\quad
 \eta^{m_0}-\tau\xi^{m_\infty}=1,\qquad
 t=\tau+\xi^{-m_\infty}.                              \tag{K32.19}
\]
It has degree $M$, the ramification at zero, infinity,
and the regular binding computed there, and no
ramification at the added order-one sites.
Define the actual total-space page by pulling this
page back through the original orbit projection and
the displayed diffeomorphism of bases. This definition
retains every original source and line transition.

\item On an ordinary simply connected chart choose a real
lift of its prescribed gamma graph. The source arrow
is literally
\[
                 p\longmapsto[p,\gamma_{\rm graph}(t(p))].
                                                               \tag{K32.20}
\]
On overlaps the lifts differ by an integer; the
source maps differ by the corresponding actual
power of $t_{-Q}$ and by the given elliptic
translation. These are holomorphic fibre translations
with their specified annular descent laws, so they
define the source translation torsor on (K32.19).
At an added site (K32.10) gives an extending graph
of $h_j$. The order is one, and its sheet map is
the original local source. A relation preserving
one representative level has $n_jk+h=0$; the actual
identity $A_j^kB^h=1$ proves that no residual quotient
remains. At the surviving roots use P3K.28's sheet
map with (K32.12). There the relation is
$(k,h)=(m_i l,-n_i' l)$, and the same norm identity
proves uniqueness. These maps extend across all
resolved components. At the selected nodes use the
actual orbit/source charts of P3F.9, retaining their
Hopf charge in the compression.

\item Translations do not change integral linear monodromy,
and the cover (K32.19) is unchanged. In the joint
cases the two closed lifted nodal loops of P3K.28
still have images $\mathbb Zv$ and $\mathbb ZAv$,
with $|\langle v,Av\rangle|=1$. In the single-cap
case every retained source meridian lifts with
trivial sheet permutation. The proved torsion-free
source Mordell--Weil group gives zero integral
coinvariants by the invariant-torsion-section
argument of P3K.26. The added translations change
none of these integral maps.

\item Let $j:\mathcal C_{\rm reg}\to\mathcal C$ and let
$\Lambda$ be the actual period lattice. The
coinvariant calculation gives
\[
       H^2(\mathcal C,\partial\mathcal C;j_*\Lambda)=0.
                                                               \tag{K32.21}
\]
Indeed compact-support duality gives the coinvariants
for $j_!\Lambda$, and $j_*\Lambda/j_!\Lambda$ is a
skyscraper with zero first cohomology. The finite
component quotient of extending translations is
also a skyscraper. Remove its cocycle, lift the
remaining translations to the smooth Lie sheaf,
and use (K32.21) to remove the lattice two-cocycle.
A partition of unity vanishing on the prescribed
binding collar removes the remaining Lie cocycle.
Exponentiation gives the actual smooth translation
isomorphism with the sectioned elliptic page, fixed
on that collar. The $I_2$ invariant lattice and
component group are included as in P3K.28.
The holomorphic source cochain (K32.7) can also be
used before this page construction; neither procedure
requires a global total-space lift of (K32.10).

\par\medskip\noindent\textbf{The corrected page line and its evaluation map.}
On each surviving root the norm-sphere map sends a
vector in the actual line
$N_i=\mathcal O(aG_i'-O_i)$ to its specified graph
level. This delta-equivariant map identifies its
associated complex line. We use all the actual
Cartier maps of the cap constructions. In particular
the quadratic root uses
\[
 \beta^*N_*
  =\pi^*\mathcal O(aG-O-L_0)+E_0+E_1,                 \tag{K32.22}
\]
including the four-point blowup correspondence
transported by P3F.49. The good-starred toric
corrections, good III/IV exceptional differences,
cancelling good-II coefficients, and reduced-root
whole-fibre factors are likewise retained.

\item At an added order-one site the norm-sphere map has
no root quotient and carries its actual local line.
Every component degree there is zero. For a retained
multiplicative fibre, the actual line norm preserves
the component-degree vector and its translation is
transitive on components. All component degrees are
therefore equal and their total is zero. Smooth
fibres have degree zero by construction. The
additive and quadratic root degrees are the
individually verified zero degrees in the cap proofs.
Thus the actual page line has degree zero on every
component of every fibre.

\item Perform the preceding source comparison through
the rigidified relative Picard map of P3K.26--28.
We describe its remaining line explicitly. On a
holomorphic translation-parameter chart form the
ratio of the two lines before pulling back by the
smooth translating section. This ratio is trivial
on every smooth elliptic fibre. Proper direct image
and evaluation give a base line and an effective
vertical evaluation divisor. Both compared lines
have zero component degrees, so the divisor
intersects every fibre component by zero. The
connected Kodaira intersection matrix has kernel
the whole-fibre multiplicity vector. Hence this
divisor is an integral sum of whole fibres.
Dividing by their base equations gives the
nowhere-zero Hom section. This proves the
evaluation isomorphism at singular and smooth
fibres alike. Rigidify at the zero section and
pull it back by the smooth translating section.
The result is the actual smooth line map,
holomorphic along fibres.

\item Glue these evaluations with every delta transition
unit and finite character retained. The remaining
ratio is pulled back from the page base; on its
compactification its class is $kF$. This processes
the exact source output (K32.8) through the corrected
root diagrams. It is the corrected page line whose
vertical remainder has been proved to be a sum of
whole fibres. The original divisor $D_{\rm vert}$
on the unrooted source has remained in the
comparison until this step. Finite root characters
may leave base holonomy, retained in the same
pulled-back line. Thus
\[
 (c_1(N_{\rm page})+kF)^2=c_1(N_{\rm page})^2,\qquad
               c_1(N_{\rm page})F=F^2=0.               \tag{K32.23}
\]
On the regular binding collar the comparison is
the prescribed one up to this base line. That
base line extends over the ordinary disc and
the fibre degree is zero, so the boundary first
Chern class is zero. The relative boundary-gauge
formula of P3K.5 gives zero square ambiguity,
including odd gauges. The source diffeomorphism
also carries the stable tangent pairing.

\item The fixed source and delta clutches glue the page
and its rotating trace by the same maps throughout
the parameter interval. Source maps commute with
the suspension translation. Their degree-zero
line lifts commute with the suspension lift on
smooth fibres by constant-factor uniformization;
extension gives that identity on the resolved
roots. The dependent parameter cochain is exactly
(K32.17). At the selected nodes and binding its
ratio is the extending unit $u^M$. Transport the
source and line charts by this unit over the
whole local parameter trace. Both origins receive
the same logarithm, leaving (K32.18) and its
positive Hopf normal sphere. These are the actual
boundary arrows of P3K.7 and P3K.28's signed Thom
calculation; our relative page comparison is
fixed on that boundary. Their characteristic
bordism therefore gives
\[
 A^{\rm Jin}
 =\frac{12c_1(N_{\rm page})^2-p_1(TS_{\rm page})}{24}
 =\chi-\frac{Ma^2}{2d}.                               \tag{K32.24}
\]
Every added site is unramified on the page and
retains its original local source model. The
Euler count is therefore exactly
\[
 \chi=\frac{M(12-e_\infty-e_0)
                +m_0e_\infty'+m_\infty e_0'}{12},
 \qquad c_1(N_{\rm page})^2=-\frac{Ma^2}{d},             \tag{K32.25}
\]
with the root Euler numbers and corrected height
line in P3K.26--28. No extra Euler replacement
comes from an order-one clutch.

\par\medskip\noindent\textbf{The full normal frame and the primitive sector.}
We finish with the actual P3K.30 normal comparison,
retaining the added unit. In radian coordinates,
the page phase on the ordinary part is
\[
 \arg P=M\varepsilon\gamma+\arg(t-\tau)
       -\varepsilon m_\infty n_0\arg t
       -\varepsilon M\arg u(t).
\]
Near the regular binding choose the literal
holomorphic coordinate
\[
 q_b=(t-\tau)(t/\tau)^{-\varepsilon m_\infty n_0}
                  (u(t)/u(\tau))^{-\varepsilon M}.      \tag{K32.26}
\]
Its derivative at $\tau$ is one. Every exponent
is an integer, and each other factor is a
specified extending unit. Retain the constant
phases of $\tau$ and $u(\tau)$ in the page
angle $\phi$. Then (K32.26) gives exactly
$\alpha=\phi-M\beta$, with
$\beta=\varepsilon\gamma$. In a one-cap chart
the $t/\tau$ factor is absent. This coordinate
retains the variable added phase throughout
the binding neighbourhood.

\item The actual source-preserving graph has
$\beta=\theta+g(p)$ on the binding, where
$p\in T^2$ and $g$ is its fixed circle-valued
graph phase. It includes the retained source,
delta, and norm-unit gauges. Apply P3K.30's
literal quotient by the common tangent
$(\cos\alpha,\sin\alpha,\rho,0)$, $\rho>0$.
The explicit triangular shear removes the
source derivatives and gives the same oriented
normal matrix
\[
 G(\alpha,\theta)=
 \begin{pmatrix}
 -\sin\alpha&\cos\theta\cos\alpha&\sin\theta\cos\alpha\\
 \cos\alpha&\cos\theta\sin\alpha&\sin\theta\sin\alpha\\
 0&\sin\theta&-\cos\theta
 \end{pmatrix}.                                      \tag{K32.27}
\]
The transported compression frame is $h(p)$,
independent of $\theta$. The actual clutch
therefore factors as
\[
 \begin{split}
 &h(p)^{-1}G(\phi-M(\theta+g(p)),\theta)\\
 &\quad=[h(p)^{-1}R_z(-Mg(p))]
           [R_z(\phi-M\theta)G(0,\theta)].             
 \end{split}\tag{K32.28}
\]
Both factors are global. The first comes from
the source two-torus and the second from the
parameter circle; the integer $M$ retains every
graph-phase winding. The normal Pontryagin
transgression vanishes by P3K.30's exact
two-factor calculation, including its
one-dimensional frame classes. The bounded
Cox comparison of P3K.27 and P3K.30 is
transported by the fixed source and line maps
above. Its remaining constant rotation loop
extends over the fixed three-disc. Hence the
actual normal pairing is zero here as well.

\item Write $H=Ma^2/d$. The signed sphere-bundle
formulas of P3K.30 now give
\[
 k=-3H,\qquad\ell=-24\chi,\qquad
 (m_{\rm rel},n_{\rm rel})=(6\chi,3H-6\chi).
\]
They satisfy $-(m_{\rm rel}+n_{\rm rel})=k$
and $A_U=(4k-\ell)/24=\chi-H/2$, while the explicit physical cap comparison gives $A_M=-A_U$. This proves (K32.2).

\item For the primitive sector without a distinguished
additive cap, choose a regular infinity disc
of order one and an ordinary zero disc with
$(m_0,n_0)=(1,0)$. P3F.45 makes all the effective
additional orders one. The preceding base
map puts their actual signed sum at infinity
with $m_\infty=1$, $n_\infty'=\varepsilon$.
The page has degree one and is the original
rational elliptic source of P3K.22. Its line
is the actual $\mathcal O(dG-O)$ after the
evaluation remainder (K32.23). P3K.22 and
P3K.27 give $\chi=1$, $H=d$, and the same
period and rotating-frame comparisons.
Equation (K32.3) follows. Both signs of the
total determinant are covered by the positive
physical angle (K32.18) and the signed normal
frame. Every original source fibre remains
in this page, including any ordinary II,
III, IV, or multiplicative fibre.

\item The certificate
\path{code/part3/order_one_pages.py} and its
executed JSON give the complete integral
base-relation changes, with $b_0$ retained,
the divisor identity (K32.17), the phase
matrix, the root-generator identity, the
binding derivative, and the actual order-one
kernel maps. The geometric source, line,
page, and boundary arrows used by these
calculations are the maps proved above.

\end{enumerate}
\end{proof}




\begin{theorem}\label{night:P3K.33}
\StatementRecord{night:P3K.33}{false as stated; correction proved}{S6\_06\_BLIND: GPT-6 Astra (OpenAI, ChatGPT Pro mode, effort Pro), recorded re-derivation in the checkpoint dated 18 November 2026}{\cref{night:P3K.11}, \cref{night:P3K.23}, \cref{night:P3K.27}, \cref{night:P3K.28}, \cref{night:P3K.30}, \cref{night:P3K.31}, \cref{night:P5.5}, \cref{night:P5.6}}
For logarithmic orders (3K,4L), \(\gcd(3K,4L)=1\), and
sphere equation (4Ln_0+3Kn_1=\varepsilon=\pm1), the explicit
integral cochain and chartwise translations below construct the
single endpoint for
(w_h=e^{-\pi i c/(KL)},\ u=\kappa w_h^{2KL},\ w_h=r^6).
They retain every finite character and give the reference
fixed-star dual pair ((-KL,-4KL)), native free pair ((KL,0))
and physical reference class (0).
The all-fan polynomial formerly attached to this endpoint is
refuted: for (K=L=1), the actual (T_2) shorter return is
(26\pmod{28}), while that polynomial is zero.
\end{theorem}
\begin{proof}
\begin{enumerate}[label=\arabic*.]
\item \par\medskip\noindent\textbf{1. The actual integral one-period primitive.}\enspace

\item Use exactly the homology marking of P5.6 and K11:
\[
\begin{split}
 A_1&=\begin{pmatrix}
 1&0&0&0\\6&0&1&0\\-6&-1&-1&0\\-2&1&0&1
 \end{pmatrix},\\
 A_2&=\begin{pmatrix}
 1&0&0&0\\0&0&-1&0\\-6&1&0&0\\3&0&1&1
 \end{pmatrix},\qquad
 A_0=(A_1A_2)^{-1}
 =\begin{pmatrix}
 1&0&0&0\\0&1&0&0\\0&1&1&0\\-1&0&0&1
 \end{pmatrix}.
\end{split}                                                     \tag{P3K.33.2}
\]
Here \(e_1\) is the marked gamma period and \(\delta=e_4\).
Let \(J\) have columns \(e_3,\delta\), and put
\(N_1=1+A_1+A_1^2\), \(N_2=1+A_2+A_2^2+A_2^3\).
For clarity, retain the full degree-zero and degree-one maps,
including the finite cyclic group arrows and the cusp invariant stalk:
\[
\begin{split}
 d^0(a,b_0,b_1,z)={}&((A_1-I)a,(A_2-I)a,
 (A_1-I)b_0,(A_2-I)b_1,a-b_0,a-b_1,a-Jz),\\
 d^1(p,q,r,s,u_0,u_1,w)={}&(KN_1r,LN_2s,
 p-r-(A_1-I)u_0,q-s-(A_2-I)u_1,\\
 &\hspace{30mm}-A_0p-A_0A_1q-(A_0-I)w).
\end{split}                                                     \tag{P3K.33.3}
\]
These are K11's actual root-disc and annulus restriction maps.
The first two target summands record the finite-root characteristic
classes, and the last three record the literal annular matches.

\item Define
\[
\begin{split}
 p=r&=\frac{P\delta-\varepsilon(A_1-I)e_1}{6}
       =(0,-\varepsilon,\varepsilon,(P+2\varepsilon)/6)^t,\\
 q=s&=\frac{Q\delta-\varepsilon(A_2-I)e_1}{6}
       =(0,0,\varepsilon,(Q-3\varepsilon)/6)^t,\\
 u_0&=u_1=w=0,
 \qquad x=(p,q,p,q,0,0,0).
\end{split}                                                     \tag{P3K.33.4}
\]
This is an integral cochain, with zero cusp-annulus component.
Indeed \(P+Q=\varepsilon\), \(Q\) is divisible by three, and \(P\)
is divisible by four. Reducing the sphere equation modulo three
gives \(P\equiv\varepsilon\pmod3\), hence
\(P+2\varepsilon\equiv0\pmod6\). The odd sphere equation
gives \(Q-3\varepsilon\equiv0\pmod6\).

\item The finite norm equations are literal:
\[
 N_1p=\frac P2\delta,\qquad
 N_2q=\frac{2Q}3\delta.
\]
Moreover
\((A_1A_2-I)e_1=\delta\), so
\[
 p+A_1q=\frac{(P+Q)\delta-\varepsilon(A_1A_2-I)e_1}{6}=0.
\]
Consequently
\[
 d^1x=(bn_0\delta,bn_1\delta,0,0,0),                 \tag{P3K.33.5}
\]
exactly the characteristic change of the actual \(c\mapsto c-b\)
shear. This proves the full local-to-global integral match; it
does not replace a finite norm by an invariant projector.
Every block in \(x\) has first coordinate zero. Hence the actual
successive parameter shears \(S_b=I+bE_{41}\) fix \(x\).

\item The sixfold endpoint used in K11 has the logarithmic primitive
\[
 x_{\rm cov}=(P\delta,Q\delta,P\delta,Q\delta,0,0,\varepsilon e_1).
\]
There is the following exact integral comparison:
\[
 x_{\rm cov}-6x=d^0y,\qquad
 y=(\varepsilon e_1,\varepsilon e_1,\varepsilon e_1,
       0\in\mathbb Z^2).                             \tag{P3K.33.6}
\]
In particular the cusp component of \(y\) is zero. Formula
(P3K.33.6) follows by substituting (P3K.33.4) in every block; its last
block is \(\varepsilon e_1\), exactly the cusp component of
\(x_{\rm cov}\).

\par\medskip\noindent\textbf{2. Chartwise holomorphic endpoint maps.}\enspace

\item This cochain is realized by an actual translation on the family.
Use the same integral Jacobian identification
\(\Pi_c=\Pi_{c-b}S_b\) as K11. On the ordinary base, let
\[
 H(t)=-\frac{P\log t+Q\log(t-1)}{2\pi i}.
\]
The cover's ordinary endpoint primitive is \(H\delta\).
At the cusp, with \(q=1/t\) and
\(s=\log q/(2\pi i)\), K11 gives
\[
 H=\varepsilon s+U(q),\qquad
 \gamma_c=(\Gamma(q),C(q)-s-dv)^t,
\]
and its actual bounded cusp primitive is
\(H\delta+\varepsilon\gamma_c\).
All logarithms and the marked period \(\Gamma\) are those of
the original family; a sixth root of a different source
translation is not substituted.

\item Let \(\tau_{{\rm cov},i}\) denote those actual cover primitives
on the ordinary, two finite-root, and cusp charts. The finite-root
primitives are the holomorphic vectors in K11's extending cover
map. Define the one-step primitives by the explicit formula
\[
       \tau_i=\frac{\tau_{{\rm cov},i}+\Pi_cy_i}{6}.
                                                               \tag{P3K.33.7}
\]
On a finite-root disc all periods in \(\Pi_c y_i\) are
holomorphic in the actual smooth root model, so (P3K.33.7) is
holomorphic there. The change of the annular gluing under the
parameter shift is linear in \(b\) and is in the fixed delta
direction. Thus its sixfold value is exactly six times its
one-step value. Taking the lattice-valued coboundary of (P3K.33.7)
and using (P3K.33.6) gives the integral cochain \(x\), rather
than a fractional cochain. This verifies equivariance for every
finite group arrow as well as equality on every annular overlap.
Exponentiation therefore gives an actual holomorphic isomorphism
\(D_c:X_c\longrightarrow X_{c-b}\). Its inverse is given
chartwise by the negative translation in the target marking.

\item One can see the finite-disc cancellation directly. The coefficients
of the potentially unbounded logarithms in \(H/6\) are
\(-P/6=-bn_0/(3K)\) and
\(-Q/6=-bn_1/(4L)\). The change of the respective logarithmic
gluing adds their negatives. The remaining unit logarithm and
root-model period in (P3K.33.7) are holomorphic. The finite cyclic
norms in (P3K.33.5) are the exact integral compatibility at the
centres of those discs.

\item The cusp formula is especially simple because \(y_{\rm cusp}=0\):
\[
 \tau_{\rm cusp}
     =\frac{H\delta+\varepsilon\gamma_c}{6}
     =\binom{\varepsilon\Gamma/6}
              {(\varepsilon C+U-\varepsilon dv)/6}.    \tag{P3K.33.8}
\]
It is a bounded holomorphic vector. In the actual cusp torus
characters it is literally
\[
 (Z,V)\longmapsto
 \left(e^{2\pi i\varepsilon\Gamma(q)/6}Z,
       e^{2\pi i(\varepsilon C(q)+U(q)-\varepsilon dv)/6}V\right).
                                                               \tag{P3K.33.9}
\]
Every Cox monomial \(q^\alpha Z^\beta V^\gamma\) is multiplied
by the corresponding character of these two units. This formula
extends simultaneously to every cone of the chosen fan, fixes
its ray and curve labels, and has the same expression on the
common Cox refinement. It therefore includes all three curves
and both triple points, not just the fixed sphere.

\item There is also an exact norm check. Under the next parameter
step the marked gamma period changes by \(-b\delta\), so
\[
 \tau(c-b)-\tau(c)=-\frac{\varepsilon b}{6}\delta.
\]
All other cochain blocks remain fixed by \(S_b\). Hence
\[
 \sum_{j=0}^5\tau(c-jb)=6\tau(c)-\frac{5\varepsilon b}{2}\delta
                      =6\tau(c)-5\varepsilon KL\delta. \tag{P3K.33.10}
\]
The last summand is an actual integral delta period. Together
with (P3K.33.6), this shows that \(D^6\) is the previously specified
cover endpoint map in the torus quotient.

\item The construction also fixes any possible endpoint ambiguity.
For the higher root base, \(H^1(\widehat B,\mathcal H)=0\)
by the same integral calculation as P5.5: multiplication of the
finite norm maps by \(K,L\) does not change their kernels in
the free lattice, and all annular and cusp maps are unchanged.
The source-period principal-cocycle matrix has determinant
\(-1\), as displayed there. Thus an extending translation has
a global Lie logarithm, and its only period is \(\mathbb Z\delta\).
The resulting translation group is \(\mathbb C/\mathbb Z\),
the global physical-delta flow. Therefore changing an admissible
endpoint primitive changes it by an actual extending delta
isotopy. The return diffeomorphism class and the free normal
comparison below are unaffected.

\par\medskip\noindent\textbf{3. The closed endpoint atlas and the reference fixed-star framing.}\enspace

\item K11's bounded cusp origin on the universal parameter cover is
\[
 \sigma(v)=
 \binom{\varepsilon\Gamma v}
       {v(\varepsilon C+U)-\varepsilon dv^2/2+\varepsilon dv/2}.
\]
Direct subtraction, retaining the constant term, gives
\[
 \sigma(v+1/6)-\sigma(v)
      =\tau_{\rm cusp}+\frac{5\varepsilon KL}{6}\delta. \tag{P3K.33.11}
\]
Define instead the explicit origin
\[
             \sigma_{\rm desc}(v)=\sigma(v)+\varepsilon KLv\delta.
                                                               \tag{P3K.33.12}
\]
Its one-step difference is
\[
 \sigma_{\rm desc}(v+1/6)-\sigma_{\rm desc}(v)
                    =\tau_{\rm cusp}+\varepsilon KL\delta. \tag{P3K.33.13}
\]
In coordinates translated by \(\sigma_{\rm desc}\), the actual
endpoint translation (P3K.33.8) is therefore \(-\varepsilon KL\delta\).
It is literally the identity on every integral torus character.
Thus (P3K.33.12) closes the toric endpoint characters for the single
positive \(w_h\)-loop. It supplies the reference fixed-star framing; it does not identify the complete marked outer plumbing with a product. No residual source translation, curve
permutation, base power, or unrecorded constant remains.

\item At intermediate parameter angles the two origins differ by the
global physical-delta action. On the sixfold cover the difference
is the closed loop \(\varepsilon KLv\delta\). Its action on the actual K23
quaternion normal coordinate is
\[
             q_{\rm H}\longmapsto
                       q_{\rm H}e^{-2\pi i\varepsilon KLv}.
\]
Consequently the Hopf normal vector
\(q_{\rm H}i\overline{q_{\rm H}}\) is unchanged at every
point of the fixed sphere and every parameter angle. The same
statement on the other two curves is the literal Cox-character
formula (P3K.33.9); those charts are already closed by (P3K.33.13).
This proves the reference fixed-star normal endpoint comparison; the other displayed character maps are retained as endpoint maps.

\item The remaining bounded torus units have the exact normal action
computed in K27: in a first chart \(q=\zeta c_0c_1\), a unitary
map \((\zeta,c_0,c_1)\mapsto(a\zeta,pc_0,q_1c_1)\) with
\(apq_1=1\) acts on \(q_{\rm H}\) by
\(q_{\rm H}\mapsto A q_{\rm H}B\), where
\(A^2=pq_1\), \(B^2=p/q_1\). These factors are independent
of the fixed-curve point. Their positive parts are applied to
both compared products. The induced Hopf-normal map is a
constant rotation; its possible endpoint loop extends over the
fixed three-disc and has zero free framing component.

\item On the ordinary annulus the known K11 difference of origins
has gamma part \(\varepsilon\theta\gamma\). The change (P3K.33.12) is pure
physical delta and is the identity on the orbit. Thus the
smooth node comparison is still exactly K27's
\(R_{\varepsilon\theta}\Phi\), with its positive triangular normal
derivative. This retains the extension of the annular
\(|q|^\theta\) factor through both node charts. On the
sixfold cover it is the signed rotating comparison typed in the
next section.

\item The closed single-loop Cox framing and the descended signed rotating
framing therefore have the same sixfold free normal class.
Their difference is an integer in the free summand of
\(\pi_1\operatorname{Map}(S^2,SO(3))\); sixfold pullback
multiplies that integer by six, which is injective. The possible
constant rotation is the torsion summand already shown to extend
over the fixed disc. Hence the free classes agree on the single
loop itself. The signed pair proved in the next section is therefore
\[
                       (k,\ell)=(-KL,-4KL)
\]
in the actual reference fixed-star product carried by (P3K.33.12).

\par\medskip\noindent\textbf{4. The signed original-family dictionary and its normal pair.}\enspace

\item We identify the original family with K28's actual E6 row, rather
than reversing a different sphere datum. Use the literal holomorphic
base coordinate
\[
             T=\frac{t-1}{t}.
\]
It sends the original \(IV^*\) fibre at \(t=0\) to infinity,
the original III fibre at \(t=1\) to zero, and the selected nodal
fibre at \(t=\infty\) to \(T=1\). The marked data become
\[
\begin{array}{c|c|c}
 &\text{infinity in }T&\text{zero in }T\\\hline
 \text{root order}&m_\infty=3K&m_0=4L\\
 \text{norm coefficient}&n_\infty=n_0&n_0^{\rm new}=n_1
\end{array}
\]
Thus \(m_0n_\infty+m_\infty n_0^{\rm new}=\varepsilon\)
and the actual page degree is \(M=m_\infty m_0=12KL\).
On each root chart this coordinate map multiplies the old base
coordinate by a holomorphic unit. Taking its specified local root
gives the exact root-chart map; it commutes with the cyclic cut.
It has no integral angular winding, so retains the marked norm
coefficients and all finite affine line characters.

\item The positive endpoint multiplier has exactly the K28 arrow. Indeed
\[
 H=-Q\frac{\log T}{2\pi i}
       +\varepsilon\frac{\log(1-T)}{2\pi i}.
                                                               \tag{P3K.33.14}
\]
K28's positive deck is
\(T^{-Q}(T-1)^{\varepsilon}\), and (P3K.33.14) differs only by
the specified constant factor \((-1)^{\varepsilon}\). Its chosen
logarithm is carried in the common bounded unit of both origins.
At the cusp \(q_T=T-1=-q\); the corresponding constant logarithm
changes \(C\) by the same bounded constant. Therefore K30's
physical coordinates are literally obtained by
\[
 g=\Gamma,\qquad -C_{30}=C+\frac{\log(q_T/q)}{2\pi i},
 \qquad \lambda=u=\kappa r^M.
\]
Its outward gamma is the same marked period \(\gamma_c\).
The line and translation data are exactly the E6 row
\((a_{\rm line},b_{\rm source})=(1,6)\): the source line is
\(\mathcal O(G-O)\), the contracting line map covers
translation by \(-6G\), and its inverse outward gamma period
covers translation by \(+6G\),
and the two good roots are of orders \(3K,4L\). The strict
section/root line maps used for this row in K28 are the same
ones used in K11 and the original source-page construction.

\item To retain the sign, write \(\alpha=\arg T/(2\pi)\). The actual
solid-torus maps of K28 have ordinary phases
\[
 \phi_1=\varepsilon(m_0\gamma-n_0^{\rm new}\alpha),\qquad
 \phi_2=\varepsilon(m_\infty\gamma+n_\infty\alpha).
                                                               \tag{P3K.33.15}
\]
Their determinant is \(\varepsilon\), and
\(m_0\phi_2-m_\infty\phi_1=\alpha\). They extend on the
actual norm cut \((z,u_i)\mapsto(\zeta_{m_i}z,u_i-n_i)\)
by K28's explicit root-solid-torus formulas. In particular the
positive weighted Hopf angle is \(\beta=\varepsilon\gamma\).
The positive \(r\)-loop has exactly that angular motion by
(P3K.33.14) and K11's displayed origins. Its complex page is the
actual sheet graph
\[
                 \eta^{m_0}-\tau\xi^{m_\infty}=1,
\]
with precisely the root source and line, because the cut changes
the sheet index by \(-\varepsilon m_jn_i\). The sheet-zero
maps of K28 are independent of this sign and retain every
finite character. Thus this is the same marked complex page
and source line for either determinant sign.

\item At its binding, if \(w\) is the complex parameter at the single
point at infinity, the actual boundary angle is
\(-\arg w=\beta=\theta\), where \(\theta\) is the positive
physical parameter angle. Consequently
\(\partial V=+E\times S^1_\theta\). Orient the actual compression
by \(\partial W=F-E\); the glued trace then has
\(\partial X=+F\times S^1_\theta\). In its actual Hopf chart
the ambient and trace orientations are, respectively,
\[
 \operatorname{or}(F)\operatorname{or}(\nu_{\rm Hopf})\,dr\,d\theta,
 \qquad \operatorname{or}(F)\,dr\,d\theta.
\]
Passing the three normal directions past the last two gives a
positive sign. Hence its meridian sphere is the positive Hopf
sphere and \(\pi_*x=+1\) for both signs. This is the actual
Thom comparison, including the orientation change in (P3K.33.15).

\item The initial Cox/rotating normal map also retains that sign.
In the beta product its raw matrix is
\(T_\varepsilon Q_x\), where
\(T_\varepsilon=\operatorname{diag}(1,1,\varepsilon)\) and
\(Q_x\) is K30's positive triangular node derivative. The same
\(T_\varepsilon\) is applied to both initial frames. Returning
to their common reference leaves \(Q_x\); conjugation by
\(T_\varepsilon\) is also conjugation by
\(\varepsilon T_\varepsilon\in SO(3)\). Thus neither a hidden
orientation reversal nor a free framing term remains.

\item For this actual page the good roots contribute no singular Euler
term. The source Euler numbers are \(8,3,1\), so
\[
       \chi=\frac{M(12-8-3)}{12}=KL,
       \qquad H_{\rm page}=M h(G)=M/6=2KL.
\]
These are the line and tangent numbers already proved in K11.
The full degree-\(M\) normal clutch of K30 applies through
the explicit maps (P3K.33.14)--(P3K.33.15). Its two factors are pulled
back from \(T^2\) and \(S^1\), so its cubic clutch integral
is zero. The positive-meridian sphere calculation gives
\[
              k_{\rm cov}=-3H_{\rm page}=-6KL,
              \qquad \ell_{\rm cov}=-24\chi=-24KL,
                                                               \tag{P3K.33.16}
\]
for both \(\varepsilon\). The preceding simultaneous closed
single-loop product proves that these are the pullbacks of its
individual integers. Characteristic-number naturality for the
degree-six parameter cover gives
\(6k=-6KL\), \(6\ell=-24KL\). These are equalities of
ordinary integers in the explicitly compared normal product;
they do not divide a class in \(\Theta_7\). Thus the single-loop
pair is \((-KL,-4KL)\) for either sphere sign. Its integer
index and its initial oriented central class are zero.


\item The endpoint and reference free invariants just constructed are the
inputs to \cref{night:P3K.73}, with shorter scalar winding
(J=2KL), \(\widehat H=KL/3\), and
\(\widehat\chi=KL/6\).
Their actual vertical class is
(\widehat H/2-\widehat\chi-Jn(n-1)/2=-KL,n(n-1)).
For (K=L=1,n=2) this is (-2=26), whereas the physical
cocharacter still has coordinates ((0,1)) and norm (S=1).
The rejected expression (KL(S-1)(S-3)/8) is zero there.
The complete all-fan answer is the actual marked-prism algorithm,
whose original-family witness (P=\left(\begin{smallmatrix}1&-2\\0&1\end{smallmatrix}\right))
has class (3), and whose full image is all of \(\mathbb Z/28\),
as proved in \cref{night:P3K.73,night:P3K.74}.

\end{enumerate}
\end{proof}




\begin{theorem}\label{night:P3K.34}
\StatementRecord{night:P3K.34}{false as stated; correction proved}{S6\_06\_BLIND: GPT-6 Astra (OpenAI, ChatGPT Pro mode, effort Pro), recorded re-derivation in the checkpoint dated 18 November 2026}{\cref{night:P3F.26}, \cref{night:P3F.29}, \cref{night:P3F.32}, \cref{night:P3K.65}, \cref{night:P3K.68}, \cref{night:P3K.76}, \cref{night:P3K.77}}
For the component-supported OS49 families with
\(\lambda=r^{3k}\), \(k>0\), signed primitive cap norm
\(n=\pm1\), and retained character \(d\delta\), the positive-fan
source, line, solid-torus and normal-frame maps below give
[
 (k_U,\ell_U)=(-6k,-24k),\qquad
 (k_M,\ell_M)=(6k,24k),\qquad
 (m_{\rm rel},n_{\rm rel},s)=(6k,0,0).
]
The actual component line retains its \(I_2\) multidegree
\((-1,+1)\) and relative square \(-2k\).
The old universal fan polynomial is refuted by the negative fan:
its whole-loop class is \(-3k\), and its actual shorter-loop
class is \(-k\), by \cref{night:P3K.65}. The complete all-fan
calculation is \cref{night:P3K.77}.
\end{theorem}
\begin{proof}
\begin{enumerate}[label=\arabic*.]
\item \par

\medskip\noindent\textbf{1. The actual one-cap specialization and both solid-torus maps}

\par\nobreak

\item Retain the actual source, section and component line
\[
y^2=x^3-3t^3(t-1)x+2t^4(t-1)(t-\tfrac12),\quad
G=(t^2,t^2),\quad\mathcal M=\mathcal O(G-O),\quad Q=3G.
\]
The finite \(IV^*\) cap at \(t=0\) has order \(m=3k\),
full vector \(\nu\ell_0+d\delta\), \(\nu=\pm1\), and the literal
Cartier/norm maps of F26/F29/F32. The selected \(I_2\) cusp is at
\(t=\infty\); the original type-II fibre at \(t=1\) is retained.
The actual scalar family is \(\lambda=r^m\). Its endpoint-to-initial
ordinary closing is \(t^\nu\).

\item Here is the complete specialization morphism from RP.1-\/-7. Set
\[
h=1,\quad l=0,\quad n_A=\nu,\quad n_B=0,\quad
 d_A=d,\quad d_B=0,\quad M=m,\quad\epsilon=\nu.
\]
At the second site the root cover is the identity, its root deck and
norm action are the identity, and its ordinary-to-cap logarithmic
cochain is zero. Thus its quotient is exactly the original type-II
source/Cartier chart. In the coordinate \(q=t-1\), the source section
is \(((1+q)^2,(1+q)^2)\). At \(q=0\) it is the smooth point
\((1,1)\) of \(y^2=x^3\), distinct from \(O\); its coordinate
\(\tau=-x/y\) is \(-1\). Therefore the \(l=0\) terminal Cartier
ratio is \((\tau+1)/\tau\), the ratio for the original horizontal
divisor \(G-O\), with no exceptional graph blowup and no extra
vertical divisor. This is the actual source and line arrow at this
site. At the first site every root, Cartier and finite-character map
is unchanged. All the following maps are therefore maps of the
one-cap family itself.

\item Let \(y_B\) denote the coefficient of the actual contracting period
\(B\). On the first root \(q_A=z^m\), write
\(a=\arg z/(2\pi)\), \(\alpha_A=ma\). The plus-log arrow is
\[
y_{\rm cap}=y_B+\nu a,\qquad
(a,y_{\rm cap})\longmapsto(a+1/m,y_{\rm cap}+\nu/m).
\]
Its invariant cap circle is \(u_B=my_B+\nu\alpha_A\). The filled
ordinary meridian is \(m\alpha_A-\nu\gamma_B\). Define the exact
outward coordinates \(\gamma_{\rm out}=-y_B\),
\(u_{\rm out}=-u_B\). Then
\(u_{\rm out}=m\gamma_{\rm out}-\nu\alpha_A\), and the root cut
changes \(u_{\rm out}\) by \(-\nu\). The positive weighted Hopf
angle is consequently \(\beta=-\nu y_B\).

\item Put \(T=t/(t-1)\). On the ordinary annulus the phases are
\[
\phi_1=\nu(m\gamma_{\rm out}-\nu\arg T/(2\pi)),\qquad
\phi_2=\nu\gamma_{\rm out},\qquad
m\phi_2-\phi_1=\arg T/(2\pi).
\]
Choose the unique positive radii with
\(r_1^2+r_2^2=1\) and \(r_2^m/r_1=|T|\); then
\((z_1,z_2)=(r_1e^{2\pi i\phi_1},r_2e^{2\pi i\phi_2})\)
has \(T=z_2^m/z_1\). The extending maps at the two solid tori are
\[
\begin{aligned}
T^{-1}=z:\quad &(z,u_\infty)\longmapsto
(zR^m e^{2\pi i\nu m u_\infty},R e^{2\pi i\nu u_\infty}),
&&R^2+|z|^2R^{2m}=1,\\
T=z^m:\quad &(z,u_0)\longmapsto
(R e^{2\pi i\nu u_0},zR^{1/m}e^{2\pi i\nu u_0/m}),
&&R^2+|z|^2R^{2/m}=1.
\end{aligned}
\]
The root cut changes their phases by integral amounts, using
\(\nu^2=1\). Each inverse first recovers the positive radius from
the nonzero coordinate, then an argument lift and the remaining
root coordinate. Different lifts differ precisely by the actual
root cut. The maps have positive radial derivative and agree with
the ordinary overlap after the specified analytic coordinate units
are retained. On the actual source/line sheet a stabilizer satisfies
\(\nu a+mb=0\), hence \((a,b)=(mj,-\nu j)\); its action is
\(A^{mj}B^{-\nu j}=1\) by the literal norm. Thus the sheet-zero
map introduces no residual quotient.

\item For a regular \(\tau\notin\{0,1,\infty\}\), the page is
\[
\eta^m-\tau\xi=1,\qquad T=\tau+\xi^{-1}.
\]
The positive weighted radial map is
\((\xi,\eta)\mapsto(\rho^m\xi,\rho\eta)\), with
\(\rho^{2m}|\xi|^2+\rho^2|\eta|^2=1\). This is the actual complex
page of the polynomial \(z_2^m-\tau z_1\). Its degree over the
source base is \(m\); it has one point of ramification index \(m\)
over the finite root and the binding, and \(m\) unramified points
over the retained type-II fibre. Its compact base has genus zero.

\par

\medskip\noindent\textbf{2. The plus-origin square and the mixed boundary map}

\par\nobreak

\item Keep F32\textquotesingle s actual periods and prescribed unit logarithms:
\[
u=(2s,s)^t+U(q),\quad
\gamma_B=(-3s,-s+mv)^t+C(q),\quad\mathcal D=3U+2C,
\]
\[
H=-u-\gamma_B,\qquad
K=3u+2\gamma_B=\mathcal D+(s+2mv)\delta,
\qquad s=\frac{\log q}{2\pi i},\quad v=\frac{\log r}{2\pi i}.
\tag{OC.3}
\]
The valuation matrix has columns \((2,1),(-3,-1)\) and determinant
one. The two actual products have PLUS origins
\[
\xi_{\rm ord}=\zeta+o,\quad\xi_{\rm cusp}=\zeta+\sigma,\quad
 o=-\frac{\nu s\mathcal D}{2m}-\frac{\nu s^2}{4m}\delta-\nu sv\delta,
\quad\sigma=\nu v\mathcal D+\nu mv(v-1)\delta.
\tag{OC.4}
\]
Thus cusp-to-ordinary is \(o-\sigma\). The ordinary closing logarithm
is \(c=-\nu s\delta\), whereas
\(\sigma(v+1)-\sigma(v)=c+\nu K\). The bounded FORWARD endpoint
is therefore \(-\nu K\). Its full angular path is
\[
(Z,V)\longmapsto
\left(U_K^{-\nu\theta}Z,
(W_K\lambda_0^2)^{-\nu\theta}
 e^{-2\pi i\nu m\theta(\theta-1)}V\right).
\tag{OC.5}
\]
The source-root \(s\mapsto s+2m\) difference retains its integral
boundary gauge. All powers use the actual logarithms of the bounded
units. The quadratic scalar factor is the global physical-circle
map on the source, line and finite cap, with its endpoint retained.
For full multiplier \((A,D)\), its local Cox lift is
\[
(z_0,z_4,z_1,z_2)\longmapsto
((AD)^{-1}z_0,z_4,Az_1,Dz_2).
\tag{OC.6}
\]
It preserves \(q=z_0z_4z_1z_2\) and the actual period and cap arrows.
Here it is used for the positive fixed-star product and its carried
normal references; no simultaneous marked product for all curve
neighbourhoods is asserted.

\item Let \((x,y_B)\) be point angles of \((u,\gamma_B)\) and
\((h_0,\kappa)\) those of \((H,K)\). Their literal maps are
\[
\binom{h_0}{\kappa}=\begin{pmatrix}2&-3\\1&-1\end{pmatrix}
\binom{x}{y_B},\qquad
\binom{x}{y_B}=\begin{pmatrix}-1&3\\-1&2\end{pmatrix}
\binom{h_0}{\kappa}.
\tag{OC.7}
\]
The incoming compression graph \(\kappa=g(h_0,w)\) moves by
\(-\nu\theta\), by OC.5. Hence
\[
\theta=-\nu(\kappa-g),\qquad\beta=-\nu y_B,
\tag{OC.8}
\]
\[
\binom{h_0}{\theta}=\begin{pmatrix}2&3\nu\\-\nu&-1\end{pmatrix}
\binom{x}{\beta}+\binom0{\nu g},\qquad\det=1.
\tag{OC.9}
\]
Both signs have the same positive degree. The source cycle \(w\) is
unchanged. The Tate lengths of \(u,H\) have positive leading terms
\(2\operatorname{Im}s,\operatorname{Im}s\), so their complex source
orientations agree on the chosen collar. Consequently
\[
\partial V_m\longrightarrow E_H\times S^1_\theta
\quad\text{has degree }+1.
\tag{OC.10}
\]
Orient the compression by \(\partial W=F-E_H\). Its product trace
has boundary \(F\times S^1_\theta\). The following actual Hopf
coordinate fixes the normal orientation; its positive meridian
line is \(U=\mathcal M^*\). The physical line is carried by K68\textquotesingle s
explicit dual arrow at the end of the computation.

\par

\medskip\noindent\textbf{3. The full Cox normal frame and its signed reference}

\par\nobreak

\item The positive fan has charts
\[
(A,B,Z)=(q/V,V/Z,Z),\qquad
(A',B',Z')=(ZA,ZB,Z^{-1}),\qquad q=ZAB,
\tag{OC.11}
\]
with physical weights \((-1,+1,0)\). Put
\(y=AB\), \(b=(|B|^2-|A|^2)/2\). The positive \(K\) height
\(\chi_K\) decreases with \(\log|V|\), so
\(\partial_b\chi_K<0\). Retain this reflection at both references.
The global quaternion is
\[
q_H=(1+jZ)(A+j\overline B)
       =(Z^{-1}+j)(A'+j\overline{B'}).
\tag{OC.12}
\]
Physical \(\delta\) acts by \(q_H\mapsto q_He^{-i\phi}\).
At \(Z=0\),
\[
q_Hi\overline{q_H}=-2bi+2\operatorname{Im}(y)j
                               -2\operatorname{Re}(y)k.
\tag{OC.13}
\]
Its determinant relative to \((\operatorname{Re}y,\operatorname{Im}y,b)\)
is \(-8\); the circle-first quotient determinant at \((A,B)=(1,0)\)
is \(-1\). Together with OC.10 these give the positive normal
Hopf meridian. The initial comparison in
\((\operatorname{Re}y,\operatorname{Im}y,-b)\) is positive triangular.
Its shear and positive diagonal contract. For the signed coordinate
\(\beta\), apply
\[
T_{-\nu}=\operatorname{diag}(1,1,-\nu)
\tag{OC.14}
\]
to BOTH references. The residual comparison is the same positive
triangular matrix; conjugation by \((-\nu)T_{-\nu}\in SO(3)\)
adds no free framing class.

\item For a bounded complex multiplier \((Z,V)\mapsto(aZ,b_1V)\), the
actual chart map is
\((Z,A,B)\mapsto(aZ,b_1^{-1}A,(b_1/a)B)\). The quaternion formula
uses its UNITARY phases. With \(p=b_1^{-1}\), \(q_1=b_1/a\), choose
\(L^2=pq_1\), \(R^2=p/q_1\), \(LR=p\) on the actual interval;
then \(q_H\mapsto Lq_HR\). Carry the positive moduli by
\[
a_t=(a/|a|)|a|^t,\quad b_t=(b_1/|b_1|)|b_1|^t,\qquad
(Z,A,B)\mapsto(a_tZ,b_t^{-1}A,(b_t/a_t)B),\quad0\le t\le1.
\tag{OC.15}
\]
This preserves \(q\), both chart transitions and both normal
references. At \(t=0\) the quaternion identity is literal. Higher
\(q\)-unit terms have zero normal linear part because \(q=ZAB\)
is quadratic in the normal directions. The closed rotation is
constant along \(F\), extends over its \(D^3\), and retains its
actual Spin endpoint.

\par

\medskip\noindent\textbf{4. The binding phase and the complete normal clutch}

\par\nobreak

\item The actual regular binding coordinate is
\(q_b=(T-\tau)(T/\tau)^{-1}\), with derivative one at \(T=\tau\).
The polynomial phase gives \(\alpha_b=\phi-m\beta\). By OC.9,
\[
\beta=2\theta+\nu(h_0-2g),\qquad
\alpha_b=\phi-2m\theta-\nu m(h_0-2g).
\tag{OC.16}
\]
The source circle term is retained. In the quotient of the collar
normal tangent \((\cos\alpha,\sin\alpha,\rho,0)\), \(\rho>0\),
the actual map sends
\((v_q,v_r,v_\eta)\) to
\((v_q-v_r(\cos\alpha,\sin\alpha)/\rho,v_\eta)\).
The two Cartesian parameter vectors give
\[
G(\alpha,\theta)=\begin{pmatrix}
-\sin\alpha&\cos\theta\cos\alpha&\sin\theta\cos\alpha\\
\cos\alpha&\cos\theta\sin\alpha&\sin\theta\sin\alpha\\
0&\sin\theta&-\cos\theta
\end{pmatrix}\in SO(3).
\tag{OC.17}
\]
A compression frame transported by the actual \(K\) flow differs
from the Cox frame by \(R(p)\), independent of \(\theta\). Its
\(S^2\) boundary map extends over \(D^3\). The compression is a
solid torus minus a ball, and its oriented normal three-bundle has
zero \(w_2\) in the actual stably parallelizable orbit \(S^5\).
The complete quotient matrix obeys the literal factorization
\[
G(\alpha,\theta)=R_z(\alpha)G(0,\theta),
\tag{OC.18}
\]
so the full clutch is
\[
[R(p)^{-1}R_z(-\nu m(h_0-2g))]
[R_z(\phi-2m\theta)G(0,\theta)].
\tag{OC.19}
\]
The page is the fixed polynomial level \(\eta^m-\tau\xi=1\), so its phase \(\phi\) is fixed throughout this clutch. The first factor is pulled from the actual source two-torus; the
second is pulled from the parameter circle. Expansion of the
Maurer-\/-Cartan cubic for their product gives the sum of factor
cubics plus an exact form. Both factor cubics vanish by dimension.
With the carried collar connection this proves
\[
\langle p_1(\overline\nu),[\overline X]\rangle=0.
\tag{OC.20}
\]
All one-dimensional normal data and both signed references remain
in this calculation.

\par

\medskip\noindent\textbf{5. The actual component line, cap markings and physical pair}

\par\nobreak

\item The sheet-zero morphism above carries the exact root line to
\(N=\mathcal O(G_{\rm page}-O)\). At the unmodified type-II points
\(G\) and \(O\) stay distinct, so their intersection contribution
is zero. The source page has \(m\) unchanged \(I_2\) fibres,
\(m\) unchanged type-II fibres and one good root. Thus
\[
\chi(\mathcal O_{S_{\rm page}})=m/3=k,\quad
\sigma(S_{\rm page})=-8k,\quad N^2=-2k,
\quad p_1(TS_{\rm page})=-24k.
\tag{OC.21}
\]
The relative period torsor is trivialized with its exact binding
collar: a lifted type-II meridian has elliptic monodromy with
\(B_{\rm ell}-I=\left(\begin{smallmatrix}0&1\\-1&-1\end{smallmatrix}\right)\),
which is unimodular. Hence the integral monodromy coinvariants
vanish. The component terms remain finite skyscraper sheaves.
The rigidified fibrewise Picard/Poincare Hom map carries the line
under the resulting translations, retaining every component degree
and finite character.

\item The literal line subfan has rays \((j,\lceil j/2\rceil,1)\) and
component degrees \((-1,+1)\). Strict contraction gives \(m\)
ordered, separated norm shells; the positive-iterate height defects
prevent reentry after the first zero. The exact radial comparison
\(v\mapsto c(q)^{F(s)-s}v\), \(F(s+1)=F(s)+1\), is periodic,
commutes with the contracting lift and physical circle, and fixes
the source. It carries the actual line into the compact page.

\item The graph-normal pushoff uses these same annular/binding charts.
Outside the regular bands its radial homotopy misses the compression.
Inside each band the exact equation is
\(\rho=-\eta_0\cot\alpha\), \(\sin\alpha>0\), \(\cos\alpha\le0\).
The slightly enlarged boundary arcs have disjoint gaps, and the
inner semicircle includes into its half-disc by a homotopy
equivalence. After the source-circle product, each relative
\(H^2\) is zero, so restriction is injective. The normal meridian
line is therefore the carried page line throughout the comparison.
The fixed \(d\) character is retained on every sheet; its full
product is one. Its compact extension changes only an integral
fibre/base class with zero square and zero pairing with \(N\).

\item The incoming basis is \((u,\gamma_B)=(2H+K,-3H-K)\), determinant
one. The compression cap \((D_w^2\times S_H^1)\times S_K^1\)
has \(H^2\) injecting into the boundary and a flat actual line.
The regular cap \(T^2_{u,w}\times D^2_{\gamma_B}\) also has zero
line square and signature. They retain their different primitive
killed cycles. Their actual line restrictions and Novikov
additivity give \(b_U^2=4N^2=-8k\) and
\(p_1(T\overline X)=-24k\).

\item With the positive Hopf meridian \(U=\mathcal M^*\), the
sphere-bundle pushforward gives
\[
(k_U,\ell_U)=(-6k,-24k).
\]
K68\textquotesingle s actual dual/spinor arrow
\(Q(a,b)=(\overline b,-\overline a)\), together with the inverse
closing in OC.5, gives
\(\lambda_{\mathcal M}(F)=-\Sigma(F^{-1})=\Sigma(F)\).
It reverses both characteristic integrals, so the physical pair is
\[
(k_{\mathcal M},\ell_{\mathcal M})=(6k,24k)=(2m,8m),\qquad
A_{\mathcal M}=\frac{4k_{\mathcal M}-\ell_{\mathcal M}}{24}=0.
\tag{OC.22}
\]
Both cap and origin arrow reversals were used. The one-cap proof
therefore retains the physical positive-fan reference used in
K76/K77. The actual negative-fan/shorter comparison is K65; the
all-fan computation is K77. No old OC all-fan polynomial or
simultaneous marked-plumbing premise is reinstated.

\end{enumerate}
\end{proof}




\begin{theorem}\label{night:P3F.53}
\StatementRecord{night:P3F.53}{false as stated; correction proved}{S6\_06\_BLIND: GPT-6 Astra (OpenAI, ChatGPT Pro mode, effort Pro), recorded re-derivation in the checkpoint dated 18 November 2026}{\cref{night:P3F.2}, \cref{night:P3F.22}, \cref{night:P3F.23}, \cref{night:P3F.24}, \cref{night:P3F.35}, \cref{night:P3F.36}, \cref{night:P3F.37}, \cref{night:P3F.38}, \cref{night:P3F.39}, \cref{night:P3F.40}, \cref{night:P3F.41}, \cref{night:P3F.42}, \cref{night:P3F.43}, \cref{night:P3F.44}, \cref{night:P3F.45}, \cref{night:P3F.46}, \cref{night:P3F.47}, \cref{night:P3F.48}, \cref{night:P3F.49}, \cref{night:P3F.50}, \cref{night:P3F.51}, \cref{night:P3F.52}, \cref{night:P3K.19}, \cref{night:P3K.20}, \cref{night:P3K.21}, \cref{night:P3K.22}, \cref{night:P3K.26}, \cref{night:P3K.28}, \cref{night:P3K.30}, \cref{night:P3K.31}, \cref{night:P3K.32}, \cref{night:P3K.33}, \cref{night:P3K.34}, \cref{night:P3K.35}}
The fixed-source canonical scalar sphere families in the literal
whole-fibre construction of E26 Remark~8.5, with rational elliptic
source and Mordell--Weil group \(\mathbb Z\), are exactly the
sixteen profiles and twenty-six source pairs in the two tables below.
Their root table has thirty-one branches; the good/good
\(III^*+II\) branch has no sphere solution, and the other thirty
have precisely the displayed coprimality and determinant conditions.
All fixed source coefficients, allowed vertical source lines,
primitive effective cap actions and order-one clutches are included
through the comparison maps in the proof.

For the actual canonical scalar parameter \(\lambda=r^M\), put
\[
 M=m_\infty m_0,\qquad H=Ma^2/d,\qquad
 \chi=\frac{M(12-e_\infty-e_0)+m_0e'_\infty+m_\infty e'_0}{12}.
\]
The specified complex page, source section and relative line have
height \(H\), relative square \(-H\), and tangent number
\(-24\chi\). Both determinant signs give the reference characteristic
dual pair \((-3H,-24\chi)\) and physical return index \(A_M=H/2-\chi\).
For every actual transported fan, the complete units and endpoint
maps enter the marked-prism calculation of \cref{night:P3K.73}.
The historical polynomial in the physical cocharacter norm is
refuted by the OS45 reference value (3), for which it gave (25).
The actual shorter original loop, component-supported OS49 loop,
and general weighted shorter scalar loop are accounted for through
K33--35 and their explicit endpoint maps.

\end{theorem}
\begin{proof}
\begin{enumerate}[label=\arabic*.]
\item \par

\medskip\noindent\textbf{C1. Data and the actual scalar family}

\par\nobreak

\item We consider the literal whole-fibre construction in E26 Remark 8.5
on a fixed rational elliptic surface with Mordell--Weil group
\(\mathbb ZG\). All source coefficients, the selected zero, the
finite root actions, the line lifts, the collar data and the
Mumford fan remain fixed while the scalar multiplying the chosen
Hom section varies. F50 supplies the actual inverse-action map
for the expanding family, and we first use its contracting
coordinate \(\lambda\). Replacing the marked generator by its
negative through F50's displayed source map puts the generic
line point and translating point in the form \(aG,bG\) with
\(a,b>0\).

\item The complete source list in K20--21 consists of the following
sixteen retained profiles and twenty-six pairs. Here \(d\)
denotes the denominator of \(h(G)=1/d\). The symbols in each
profile are the actual Kodaira fibres of the source.

\begin{longtable}[]{@{}lrll@{}}
\toprule\noalign{}
OS entry & \(d\) & Source profile & \((a,b)\) \\
\midrule\noalign{}
\endhead
\bottomrule\noalign{}
\endlastfoot
43 & 2 & \(III^*+3I_1\) & \((1,2),(2,1)\) \\
43 & 2 & \(III^*+II+I_1\) & \((1,2),(2,1)\) \\
45 & 8 & \(I_8+4I_1\) & \((8,1)\) \\
45 & 8 & \(I_8+II+2I_1\) & \((8,1)\) \\
47 & 14 & \(I_7+I_2+3I_1\) & \((14,1)\) \\
47 & 14 & \(I_7+I_2+II+I_1\) & \((14,1)\) \\
47 & 14 & \(I_7+III+2I_1\) & \((7,2),(14,1)\) \\
49 & 6 & \(IV^*+I_2+2I_1\) & \((2,3),(6,1)\) \\
49 & 6 & \(IV^*+III+I_1\) & \((1,6),(2,3),(3,2),(6,1)\) \\
50 & 12 & \(I_1^*+I_3+2I_1\) & \((6,2)\) \\
50 & 12 & \(I_1^*+IV+I_1\) & \((2,6),(6,2)\) \\
55 & 20 & \(I_5+I_4+3I_1\) & \((20,1)\) \\
55 & 20 & \(I_5+I_4+II+I_1\) & \((20,1)\) \\
56 & 30 & \(I_5+I_3+I_2+2I_1\) & \((30,1)\) \\
56 & 30 & \(I_5+I_3+III+I_1\) & \((15,2),(30,1)\) \\
56 & 30 & \(I_5+IV+I_2+I_1\) & \((10,3),(30,1)\) \\
\end{longtable}

\item For each row choose one of its actual \(I_1\) fibres as the simple
whole zero. The generic Picard line is represented by
\(\mathcal O(aG-O)\). At the \(I_1^*\) source use the actual
corrected divisor \(\mathcal O(aG-O-L_0)\) of F36/F38, with its
specified component \(L_0\). All additional allowed vertical
source divisors are carried by the actual maps of F52.

\item Put every additive fibre in its specified root disc. Its
effective order and reduced minimal root type are given below.
The variable \(j\) in each entry is independent.

\begin{longtable}[]{@{}
  >{\raggedright\arraybackslash}p{(\columnwidth - 4\tabcolsep) * \real{0.3333}}
  >{\raggedright\arraybackslash}p{(\columnwidth - 4\tabcolsep) * \real{0.3333}}
  >{\raggedright\arraybackslash}p{(\columnwidth - 4\tabcolsep) * \real{0.3333}}@{}}
\toprule\noalign{}
\begin{minipage}[b]{\linewidth}\raggedright
Source fibre and allowed translating coefficient
\end{minipage} & \begin{minipage}[b]{\linewidth}\raggedright
Effective order
\end{minipage} & \begin{minipage}[b]{\linewidth}\raggedright
Minimal root fibre
\end{minipage} \\
\midrule\noalign{}
\endhead
\bottomrule\noalign{}
\endlastfoot
\(III^*,\ b=1\) & \(4j+3,\ j\geq0\) & \(III\) \\
\(III^*,\ b=2\) & \(4j,\ j\geq1\) & smooth \\
\(IV^*,\ b=1,2\) & \(3j+2,\ j\geq0\) & \(IV\) \\
\(IV^*,\ b=3,6\) & \(3j,\ j\geq1\) & smooth \\
\(I_1^*,\ b=2,6\) & \(2\) & \(I_2\) \\
\(II\) & \(6j+1,\ j\geq0\), or \(6j,\ j\geq1\) & \(II\), or smooth \\
\(III,\ b=1,3\) & \(4j+1,\ j\geq0\) & \(III\) \\
\(III,\ b=2,6\) & \(4j,\ j\geq1\) & smooth \\
\(IV,\ b=1,2\) & \(3j+1,\ j\geq0\) & \(IV\) \\
\(IV,\ b=3,6\) & \(3j,\ j\geq1\) & smooth \\
\end{longtable}

\item There are at most two additive discs in these profiles. Label
their orders \(m_\infty,m_0\), placing a starred fibre at infinity.
For a single additive disc take \(m_0=1\) at an actual regular
disc; with no additive disc use two actual regular discs of
order one. These regular discs are the marked base charts in
K32's construction. The order-one reduced \(II,III,IV\) entries
are included in the table.

\item Keep only coprime \(m_\infty,m_0\). This removes the
\(III^*+II\) good/good branch, whose orders are both even.
At \(I_1^*+IV\), it selects the reduced orders \(6j+1\),
and the good orders \(3j\) with \(j\) odd. In every other
two-disc row it imposes the displayed integer gcd condition
on the two independently chosen orders.

\item Every other smooth or multiplicative logarithmic disc has
effective order one. Its integer norm \(n_j\), its actual
elliptic translation and its finite line character are
retained. Write \(b_0\) for the marked ordinary base framing,
and put
\[
 N=\sum_j n_j-b_0,\qquad
 M=m_\infty m_0,\qquad
 D=m_0n_\infty+m_\infty n_0+MN.                 \tag{SC.1}
\]
The complete sphere condition is \(D=+1\) or \(D=-1\).
All its integral solutions are allowed with their actual
cap and collar data. In particular it forces
\(\gcd(m_i,n_i)=1\). The parameter used in this catalogue is
the actual closed canonical cover
\[
                         \lambda=r^M.          \tag{SC.2}
\]

\item There is an exact description of every norm solution. Choose
integers \(u,v\) with \(m_0u+m_\infty v=1\). For every
\(N,j\in\mathbb Z\) and \(\varepsilon\in\{1,-1\}\), put
\[
 n_\infty=\varepsilon u+m_\infty j,\qquad
 n_0=\varepsilon v-m_0(j+N).
\]
Substitution gives \(D=\varepsilon\). Conversely the same
equation modulo \(m_\infty\) gives
\(n_\infty\equiv\varepsilon u\pmod {m_\infty}\), so the
displayed \(j\) is integral and determines the displayed
\(n_0\). With \(m_0=1\), retain the arbitrary regular-disc
norm \(n_0\) and use
\(n_\infty=\varepsilon-m_\infty(n_0+N)\).
With both orders one the equation is
\(n_\infty+n_0+N=\varepsilon\).
These are actual integer norm formulas with every
order-one slot present.

\par

\medskip\noindent\textbf{C2. The page, line and return for each entry}

\par\nobreak

\item Let \(e_\infty,e_0\) be the source Euler numbers at the two
marked discs, and \(e'_\infty,e'_0\) those of their reduced
minimal root models. A regular disc has both Euler numbers
zero. The table uses
\[
 e(II)=2,\quad e(III)=3,\quad e(IV)=4,\quad
 e(I_1^*)=7,\quad e(IV^*)=8,\quad e(III^*)=9,
 \quad e(I_2)=2.
\]
K32's exact base map absorbs the order-one norms into
\(n'_\infty=n_\infty+m_\infty N\). Its actual page is the
compactification of
\[
 \eta^{m_0}-\tau\xi^{m_\infty}=1,\qquad
 t=\tau+\xi^{-m_\infty},\qquad
 g(C)=\frac{(m_\infty-1)(m_0-1)}2.               \tag{SC.3}
\]
Its degree over the source base is \(M\). Pull back the
source section and the actual root line through the
sheet-zero maps, retaining the divisor corrections
at every root resolution. The resulting compact page
surface and its relative line satisfy
\[
 \begin{aligned}
 \chi&=\frac{M(12-e_\infty-e_0)
                 +m_0e'_\infty+m_\infty e'_0}{12},
 &H&=\frac{Ma^2}{d},\\
 c_1(N_{\rm page})_{\rm rel}^2&=-H,
 &p_{\rm rel}&=-24\chi.
 \end{aligned}                                  \tag{SC.4}
\]
The reduced root line has zero degree on each component.
Its marked section has height \(H\); the full intersection
formula, including every reducible-fibre correction,
is the proved formula of F2. Thus (SC.4) determines the
actual page and line even when a section intersection
alone does not specify its relative square.

\item For one positive \(r\)-circle, the actual reference Cox
product has
\[
 \begin{split}
 (k_{\rm Jin},\ell_{\rm Jin})&=(-3H,-24\chi),\\
 (m_{\rm rel},n_{\rm rel},A_M)&=
           (6\chi,3H-6\chi,H/2-\chi).
 \end{split}                                   \tag{SC.5}
\]
These formulas hold for both values of \(D\); the positive
Hopf angle is the signed physical suspension angle in
K28/K32, and their sheet and Thom maps give the same
complex orientation in (SC.5).

\item For an allowed fan, retain the full period functions and
apply \cref{night:P3K.73} with (J=M), reference data
(H,\chi), and the displayed source/cap endpoint.
For the actual shorter family use (J=M/b),
(widehat H=H/b,widehat\chi=\chi/b), through its explicit
covering morphism. The physical reference is
(widehat H/2-widehat\chi).
Pullback (r=z^h) composes (h) actual positive returns;
F50's displayed inverse-action map sends a large-scalar expanding
circle to the inverse return. The old norm polynomial is not used.

\par

\medskip\noindent\textbf{C3. Necessity and construction through exact maps}

\par\nobreak

\item K20 obtains one simple whole \(I_1\) zero directly from
the integral cusp and cap maps. The height equation
for its Hom line is \(ab/d=1\). K21 applies the actual
component permutations and the starred intersection
matrices to that equation. Its executable source
enumeration leaves exactly the sixteen rows above.
The excluded \(I_3^*\) rows have no permitted reduced
cap, and the three excluded profiles with no \(I_1\)
have no selected simple whole zero. Every exclusion
is a statement about the actual component or cusp
map specified in those proofs.

\item F37, F41--43 and F46--47 derive the root table from
the minimal equations, actual section positions,
finite deck weights and line divisors. Their maps
construct the reduced \(III,IV,II\) roots and the
good roots for every displayed order. F36/F38 supply
the quadratic \(I_1^*\) root and its corrected
line, including the exceptional divisor and
finite character. F44/F49 compare every permitted
affine root action by a holomorphic translation
conjugation and the rigidified line evaluation.
In the quadratic case that comparison carries
the blowup centres and the entire divisor diagram.
K19's literal quotient removes a normal root
kernel and gives precisely these primitive
effective root data.

\item F45 proves the simultaneous smooth and multiplicative
order obstruction for primitive translation. F48
proves it for each weighted additive core, including
all primes dividing its source translation degree.
Their integral rim groups and local matrices force
every added smooth or multiplicative order to be one.
The topology computations in F36--47 and K32 then
give exactly (SC.1). The source vanishing cycles
kill the original elliptic lattice; the remaining
Seifert group is cyclic of order \(|D|\).
The full integral degree-two and degree-three maps
in those computations are isomorphisms at
\(D=\pm1\). Their Hopf/Gysin comparison and the
oriented six-sphere classification give the
actual sphere. F35's chartwise fan blowups and
blowdowns retain this integral topology for
every allowed fan.

\item Conversely, take a row, orders, and norm solution
in (SC.1). Use the exact local models and
linearizations just specified and glue by their
displayed annular morphisms. Their finite actions
are free because the primitive norm conditions
have just been verified. K32 includes every
fixed order-one elliptic clutch: its translated
source is identified with the original source
through the constructed section, and its line
is carried by the proper direct-image evaluation
with the residual vertical divisor retained.
These are holomorphic maps over the source base.
F51 proves the displayed root and line diagrams
for every fixed source in these profiles through
its marked Weierstrass equations and actual
coordinate transformations. F52 supplies the
invariant-divisor line map for every source line
with the same generic point. Its full-fibre
factor at the selected zero induces the actual
cone-height map
\((v,h)\mapsto(v+kh\delta,h)\); its factors at
the caps become the recorded characters and
integer collar transitions. Therefore these
constructions realize all the fixed source,
line and cap data under consideration.

\item The local toric charts are smooth by their
unimodular cones; the finite caps are smooth
by their free actions. The overlap maps above
are holomorphic and coincide with the given
line suspension on the common open set.
The proper Mumford quotient and the compact
finite cap pieces give the compact complex
threefold in the sphere calculation. The
closing cochain in K32 uses the actual
degree-\(M\) base function and all its retained
units. It constructs the holomorphic family
on (SC.2), including its actual closed cusp endpoint atlas and reference fixed-star product. Its further fan values use the complete marked-prism construction of \cref{night:P3K.73}.

\item For the page, K26/K28 give the actual
sheet-zero root maps in (SC.3). K32 carries
the order-one clutches and the exact
binding coordinate through those maps.
Their integral coinvariant calculation and
rigidified Hom evaluation identify the
actual relative line, giving (SC.4).
K30 computes its full normal clutch,
including both determinant signs, and
gives (SC.5). Finally K31's common-exterior
meridians, stable frame transitions and
explicit flop recurrence apply to that
simultaneous product and give (SC.6).
This derivation retains the maps used at
each source, cap, page and fan comparison.

\par

\medskip\noindent\textbf{C4. Other specified loops and scope}

\par\nobreak

\item K33 supplies the explicit shorter closed
loop for the original and coprime higher
\((3K,4L)\) families. Its exact parameter
is \(w_h=r^6\), and its integral cochain
and closed Cox product give
\[
 [\Sigma_{\Phi,T}]
       =KL\frac{(S-1)(S-3)}8 \pmod {28}.
\]
The original \(K=L=1\) example with
\(v_1=(3,1),v_2=(-1,0)\) has
\(\delta=(0,1)=v_1+3v_2\), determinant
one and \(S=7\). Its return class is
three. The exact clutch transfer and
mapping-torus argument of F22--24 apply
to this explicit holomorphic family.

\item K34 handles the named OS49
component-supported line through its
actual source component degrees. For
\(\lambda=r^{3k}\), its pair is
\((-6k,-24k)\), and its all-fan class
is \(3k(S-1)(S-3)/4\). The line,
period basis and mixed-angle maps
in that proof specify the comparison
with its page and each fan.

\item This catalogue fixes the source and
all auxiliary data along the scalar
circle. A loop varying the marked
source, a finite-cap lift or an
endpoint identification needs its
own transport maps. The separate
source-loop working proof records
the computed nodal branch exchange
and the uncomputed fan and cap
endpoint transport. K35 now constructs the shorter
weighted scalar loop by its actual
meromorphic translation lift, finite
cap characters and cyclic scalar
correction. Its parameter is
\(w=r^b,\lambda=w^{M/b}\); its
characteristic pair is
\((-3H/b,-24\chi/b)\), and its
physical reference index is \(A_M/b\).
The proof uses the same closed Cox
origin after the exact descended
endpoint map. These are integer
characteristic numbers of that
constructed family. The varying-source
loops still require the recorded
relative transport calculation.

\item The executable catalogue is
\(\texttt{code/part3/canonical\_scalar\_catalogue.py}\),
with its complete JSON and captured output beside it.
It reads the K21 source reduction and independently
derives all thirty-one root branches from the local
table. Each branch records its bounds, gcd domain,
source and root Euler numbers, symbolic
\(\chi,H,A\), individual pair, all-fan expression,
and a historical primitive norm witness for each
determinant sign. The sole empty branch is recorded
with its universal even-determinant obstruction;
its formal arithmetic is assigned no sphere return.
The search box of twenty-five consecutive values
per root index supplies witnesses only; the exact
gcd and Bezout equations above specify the full
unbounded parameter sets.

\item The certificate verifies all sixteen K26 rows in
their twenty compatible source-profile instances,
all nine K28 joint rows, both one-\(I_1^*\)
specializations, and the ten primitive K22
representatives. It also proves polynomial
integrality of \(\chi,H/2,M/b\), verifies
the reference and flopped fan substitutions, and
checks the full norm formulas just proved. Its 437 recorded checks are historical; the former all-fan polynomial fields in that output are refuted and are not used here. In particular every
one of the twenty-six source pairs has a supplied
coprime sphere witness, and exactly thirty
root branches are nonempty.

\end{enumerate}
\end{proof}




\begin{theorem}\label{night:P3K.35}
\StatementRecord{night:P3K.35}{false as stated; correction proved}{S6\_06\_BLIND: GPT-6 Astra (OpenAI, ChatGPT Pro mode, effort Pro), recorded re-derivation in the checkpoint dated 18 November 2026}{\cref{night:P3F.36}, \cref{night:P3F.38}, \cref{night:P3F.41}, \cref{night:P3F.42}, \cref{night:P3F.43}, \cref{night:P3F.44}, \cref{night:P3F.45}, \cref{night:P3F.46}, \cref{night:P3F.47}, \cref{night:P3F.48}, \cref{night:P3F.49}, \cref{night:P3F.51}, \cref{night:P3F.52}, \cref{night:P3F.53}, \cref{night:P3K.14}, \cref{night:P3K.15}, \cref{night:P3K.30}, \cref{night:P3K.31}, \cref{night:P3K.32}, \cref{night:P3K.33}}
Write the actual source line as
\(\mathcal M=\mathcal O(aG-O+D)\), with the prescribed vertical
correction when it occurs, and the contracting source translation as
\(-bG\). The whole-zero catalogue has \(ab=d\), where \(h(G)=1/d\).
Let the actual nontrivial cap orders be \(m_i\), put
\(M=\prod_i m_i\), and include any order-one caps in the signed
norm equation. The sphere datum is
\[
           \sum_i P_i=\varepsilon,\qquad
           P_i=\frac{M n_i}{m_i},\qquad \varepsilon=\pm1.
                                                               \tag{P3K.35.1}
\]
This sum includes K32's actual ordinary framing chart. If its marked
base relation is \(\sum\alpha_i+b_0\gamma=0\), that chart has
transition \(h=\gamma+b_0\alpha\), effective order one, norm
\(-b_0\), and \(P_{\rm frame}=-Mb_0\). Thus (P3K.35.1) is exactly
\[
 m_0n_\infty+m_\infty n_0
                +M\left(\sum_jn_j-b_0\right)=\varepsilon.
\]
The framing chart is the actual base-atlas map of K32, with no
additional elliptic or line clutch. Its exponent is retained in
the closing rational function below.
For the already constructed one-gamma family the parameter is
\(\lambda=r^M\), and the actual finite-root maps satisfy
\[
 A_{i,r}=r^{P_i}A_{i,0},\qquad B_\lambda=r^M B_0,
 \qquad A_{i,0}^{m_i}=B_0^{n_i},\qquad[A_{i,0},B_0]=1.
                                                               \tag{P3K.35.2}
\]
Every weighted row has \(b\mid M\). The construction below gives
an actual free deck action over \(r\mapsto\zeta_b r\), with all
cap, line, and cusp maps retained. It descends the family to
\[
                     w=r^b,\qquad \lambda=w^{M/b}.       \tag{P3K.35.3}
\]
The single positive \(w\)-loop has the closed endpoint atlas and reference fixed-star framing constructed below.
Put \(H=Ma^2/d\) and let \(\chi\) be the page holomorphic Euler
characteristic from K30/K32. Its integers and index are
\[
 k_w=-\frac{3H}{b},\qquad
 \ell_w=-\frac{24\chi}{b},\qquad
 A_{U,w}=\frac{\chi-H/2}{b},\qquad A_{M,w}=\frac{H/2-\chi}{b}.                              \tag{P3K.35.4}
\]
These are characteristic-number equalities on an explicitly descended
family. For every fan the return is computed from these actual endpoint
maps and all full period units by \cref{night:P3K.73}, with
(J=M/b,widehat H=H/b,widehat\chi=\chi/b).
The old all-fan polynomial is refuted by the OS45 (b=1)
reference return (3). No homotopy-sphere class is divided.

\end{theorem}
\begin{proof}
\begin{enumerate}[label=\arabic*.]
\item \par\medskip\noindent\textbf{2. The meromorphic source lift and its actual norm.}\enspace

\item The generic-fibre source translation by \(+G\) fixes every degree-zero
Picard class. Here its isomorphism is literal. Put \(P=aG\) and
\(D_P=P-O\). For \(a\ne1\), take the ratio of the two chord functions
through \((P-G,O)\) and \((P,-G)\). Their third intersection is the
same, and the ratio \(f\) has divisor
\[
             \operatorname{div}(f)
                 =(P-G)+O-P-(-G)=t_G^*D_P-D_P.         \tag{P3K.35.6}
\]
For \(a=1\), the exact function is \(f=1/(x-x(G))\), with the
same divisor. In the rational section frame, the isomorphism
\(t_G^*\mathcal M\to\mathcal M\) sends \(t_G^*s\) to \(fs\);
its inverse supplies a meromorphic bundle map \(U\) covering \(+G\).
The prescribed vertical divisor \(D\) is retained in its rational
frame. Its zeros and poles are vertical, because (P3K.35.6) is already
an isomorphism on every generic fibre.

\item The actual map \(B_0^{-1}\) covers \(+bG\). Thus
\[
                        \rho(t)=B_0^{-1}/U^b            \tag{P3K.35.7}
\]
is a meromorphic base function. To verify this assertion, compare the
two line maps on a generic compact elliptic fibre: their ratio has
neither a zero nor a pole there and is therefore constant. The resulting
function lies in the rational function field of the base, because both
maps are the displayed rational divisor maps. At the selected I1, if
the divisor order of \(U\) is \(e\), then
\[
                      \operatorname{ord}\rho=-1-be.     \tag{P3K.35.8}
\]
The \(-1\) is the actual simple pole of the inverse of the given
linearization section. At an uncapped retained fibre the line has zero
degree on every component. The vertical divisor of \(U\) has zero
intersection with every component, so it is an integral whole-fibre
multiple; the same local order is then \(-be\). This is the integral
kernel calculation already proved for the actual reduced fibres in
K14/K15. A local whole-fibre power therefore makes the lift invertible
there, and introduces no fractional source-root character.

\item This retains every vertical line in F52, not only the chosen
representative divisor. Its canonical rational section gives the
actual ordinary comparison. On a cap root F52's direct-image
evaluation gives \(N_i\simeq N_{i,0}\otimes f_i^*L_i^{\rm base}\),
after dividing by the displayed integral whole-fibre monomial.
Translation by \(G_i'\) fixes this base line. Any extra deck
character multiplies every fibre weight by the same scalar and
therefore cancels in the commutator ratios below. At the selected
zero a coefficient \(kF\) has the literal frame map
\((Z,V,q)\mapsto(Z,q^kV,q)\); its height-one lattice map translates
every vertex by \(k\delta\). This is the actual period reindexing
of F52, so it preserves all edge differences, the physical
cocharacter, and the integer \(S\) in the transported fan.

\par\medskip\noindent\textbf{3. Actual extending lifts on the cap root lines.}\enspace

\item Use the actual minimal root source \(S_i'\), its deck \(D_i\), and
the line \(N_i\) established by the source-line diagrams of F41--F49.
Its generic Picard point is \(aG_i'\). The source translation by
\(+G_i'\) extends on that actual minimal regular source. The generic
chord lift of (P3K.35.6) differs from an everywhere invertible line map
only by a whole-fibre divisor: on a smooth root this follows from
the single reduced fibre, and on a reduced III, IV, II or I2 root
it follows from the zero component degrees and the primitive integral
kernel of its component intersection matrix. Remove that integer
base power. Denote the resulting extending map by \(L_i\).

\item All translation lifts on a degree-zero elliptic line commute. This
has an explicit generic-fibre model: write that line as
\((\mathbb C\times\mathbb C^*)/\Lambda\), with lattice action
\((z,v)\mapsto(z+\omega,\chi(\omega)v)\). The lifts are translations
in the first coordinate and scalar multiplications in the second,
which commute. Equality on the generic fibres extends over the actual
root line. In particular \(L_i\) commutes with the translation lift
\(F_i\) used in the cap's \(A_{i,0}\).

\item The deck commutator is a base unit:
\[
                      D_iL_iD_i^{-1}=\chi_i(z)L_i.
\]
Its cyclic norm is one. Divide by its constant value at zero, take
the logarithm \(g\) vanishing at zero, and put
\(\mathcal D f(z)=f(\zeta_{m_i}^{-1}z)\). The norm identity
gives \(\sum_j\mathcal D^jg=0\). The explicit function
\[
 h=-\frac1{m_i}\sum_{j=0}^{m_i-1}j\mathcal D^jg
 \quad\text{satisfies}\quad (\mathcal D-I)h=-g.
\]
Rescaling \(L_i\) by the target scalar \(e^h\) therefore makes
the commutator the constant
\(\chi_i=\chi_i(0)\). Once \(\chi_i^b=1\), the unit
\(L_i^b/B_0^{-1}\) is deck invariant. An invariant analytic b-th
root rescales \(L_i\) further so that
\[
 L_i^b=B_0^{-1},\qquad
 D_iL_iD_i^{-1}=\chi_i L_i,
 \qquad L_iA_{i,0}L_i^{-1}=\chi_i^{-1}A_{i,0}.          \tag{P3K.35.9}
\]
The last equality uses the actual form \(A_{i,0}=\text{unit}\cdot
F_i^{n_i}D_i\) and the just-proved commutativity, including its
retained finite line character.

\par\smallskip\noindent\textit{The good-root characters.}\enspace

\item On a good root let \(\rho_i\) be the deck derivative at a marked
fixed point of the central elliptic curve. The natural line
\(N_i=\mathcal O(aG_i'-O)\) has fibre weights \(\rho_i\) at
\(aG_i'\), \(\rho_i^{-1}\) at \(O\), and one at every other
fixed point. This is precisely the Cartier support-function line
map in F41/F42, including its exceptional correction, so it is
the actual cap line.

\item Translation by \(G_i'\) cycles the marked fixed points. Comparing
the fibre weights at two successive points gives its deck
commutator. For a good III or III* root the marked point has order
two, \(a\) is odd, and \(\rho_i\) has order four. The two successive
ratios are both \(\rho_i^2=-1\). For a good IV or IV* root the
marked point has order three. The three successive ratios are
all \(\rho_i^{-a}\). Thus
\[
 \chi_{\mathrm{III},\mathrm{III}^*}=-1,\qquad
 \chi_{\mathrm{IV},\mathrm{IV}^*}=\rho_i^{-a}.           \tag{P3K.35.10}
\]
The actual minimal-root substitutions give
\(\rho_{\mathrm{IV}^*}=\zeta_3^2\) and
\(\rho_{\mathrm{IV}}=\zeta_3\). Consequently the single-cap
weighted E6 row \((a,b)=(2,3)\) and the ordinary IV row
\((a,b)=(10,3)\) both have \(\chi=\zeta_3^2\).
In the original E6 two-cap row \((1,6)\), the IV* character is
\(\zeta_3\) and the III character is \(-1\).

\par\smallskip\noindent\textit{Reduced partners with the component class killed by a.}\enspace

\item In the weighted two-cap rows, a reduced cap is the partner on
which \(aG_i'\) has zero component class. Its central line is
trivial. One can check this on the actual subsequent good root:
after multiplication by \(a\), the marked section has reached
the identity point in the original minimal good model. Its
formal parameter vanishes to positive order. Substitution of a
root order \(m\ge2\) multiplies that order by \(m\). A nonzero
point in the additive identity group of a reduced II, III or IV
fibre has order exactly one in this good-root formal parameter,
as follows from \(x=w^2X,y=w^3Y\) and
\(-X/Y=w(-x/y)\). Therefore the root section specializes to
the zero element of the additive group. For the order-one root
the deck itself is the identity. In every case the central line
is trivial, or the deck is identity. A nonvanishing global
section of a trivial line on the connected proper central fibre
has a single deck eigencharacter, so comparison after translation
gives
\[
                              \chi_i=1.                 \tag{P3K.35.11}
\]
This covers the reduced III partner of the E6 \((2,3)\) row,
the reduced IV partner of its \((3,2)\) row, the reduced II
partner of the E7 \((1,2)\) row, and the reduced IV partner
of the starred \((6,2)\) row.

\par\smallskip\noindent\textit{The quadratic I1* root.}\enspace

\item Here retain the actual Tate marking
\[
 q_T=z^2u(z^2),\qquad Z_G=zv(z),\qquad
 v(z)v(-z)=-u(z^2).
\]
Thus \(+G\) exchanges the two central I2 components, whereas
the root deck fixes their labels and reverses each normalized
component. The actual root line is \(\mathcal O(aG'-O)\),
\(a=2\) or \(6\), with identity-component Picard parameter
\(Z_{aG'}(0)=-1\), by the exact F36/F38 line diagram.

\item Trivialize this multidegree-zero line on the two normalizations.
Its node-gluing scalars can be chosen \((\lambda_0,\lambda_1)
=(1,-1)\). Label the node from the infinity of \(C_j\) to
the zero of \(C_{j+1}\) by \(j\), with indices modulo two.
The translation exchanges the nodes in their cyclic order;
the deck exchanges their ends. If their frame factors are
\(t_j,d_j\), compatibility is literally
\[
 t_{j+1}\lambda_j=\lambda_{j+1}t_j,\qquad
 d_j=\lambda_j\lambda_{j+1}d_{j+1}.
\]
Hence \(t_1=-t_0\) and \(d_1=-d_0\). Their commutator is
\(d_{j+1}/d_j=-1\), independently of the nonzero coordinate
units. This proves on the actual corrected root line
\[
                         \chi_{I_1^*}=-1.              \tag{P3K.35.12}
\]
It retains both components and both node gluings.

\item All these calculations retain arbitrary marked source coefficients.
F51 supplies the exact minimal Weierstrass root substitutions,
the linearization fixing the two marked sections, and the same
Cartier corrections with their holomorphic unit coefficients.
For the starred quadratic root it supplies precisely the units
\(u,v\) and the identity displayed above. Thus the character
calculation takes place on each actual root line under the F51/F52
maps, rather than identifying source members only by their fibre
profiles.

\par\medskip\noindent\textbf{4. The global rational root and exact finite-cap map.}\enspace

\item Pull the global meromorphic \(U\) to a root line by its already
proved rational line correspondence. Comparison with \(L_i\)
is a meromorphic scalar constant along generic elliptic fibres,
so it has the form
\[
                          U=z^{e_i}h_i(z)L_i
\]
with \(e_i\in\mathbb Z\) and \(h_i\) a nonzero holomorphic
unit after shrinking the disc. Since \(U\) is defined over
the original base, it is invariant under the deck conjugation.
For a target scalar the exact rule is
\[
 D_i m_f D_i^{-1}=m_{f(\zeta_{m_i}^{-1}z)}.
\]
Inserting the displayed comparison gives
\[
 \chi_i=\zeta_{m_i}^{e_i},\qquad
 h_i(\zeta_{m_i}z)=h_i(z),\qquad
 \operatorname{ord}_{q_i}\rho=-\frac{be_i}{m_i}.        \tag{P3K.35.13}
\]
The last number is an integer, because \(\rho\) is a base
rational function. Write \(\chi_i=\zeta_b^{z_i}\); then
\(\operatorname{ord}_{q_i}\rho\equiv-z_i\pmod b\).

\item The actual one-gamma positive closing multiplier is the base
rational function \(h_0(t)=e^{2\pi iH(t)}\), with order
\(-P_i\) at cap \(i\), order \(\varepsilon\) at the
selected I1, and no other zeros or poles.
In an affine chart containing the finite sites it is explicitly
\[
 h_0(t)=c_0(t-t_c)^\varepsilon
                   \prod_{i\text{ finite}}(t-t_i)^{-P_i},
 \qquad c_0\ne0,
\]
with its recorded infinity order. The framing chart contributes
\((t-t_{\rm frame})^{Mb_0}\), exactly K32's inverse-closing
exponent with its arrow reversed. The constant is the fixed
bounded unit of the actual closing map. Therefore
\[
                h_0(t)\rho(t)^\varepsilon
\]
has a divisor divisible by \(b\) exactly as follows. At the
selected I1, (P3K.35.8) cancels the order \(\varepsilon\).
At any retained fibre its order is a multiple of \(b\).
At a cap its order modulo \(b\) is
\(-P_i-\varepsilon z_i\). The following actual character
table and the sphere equation prove this is zero in every row:

\begin{center}\small\begin{tabular}{p{0.32\linewidth}p{0.04\linewidth}p{0.22\linewidth}p{0.27\linewidth}}\toprule
Weighted sector & b & Nontrivial root orders & Character exponents \(z_i\) modulo b\\
\midrule
E7 good III*, one cap & 2 & 4l & 1\\
E6 good IV*, one cap & 3 & 3l & 2\\
OS47 good III & 2 & 4l & 1\\
OS56 good III & 2 & 4l & 1\\
OS56 good IV & 3 & 3l & 2\\
quadratic I1*, one cap & 2 & 2 & 1\\
E6, (a,b)=(1,6) & 6 & 3l,4k & 2,3\\
E6, (a,b)=(2,3) & 3 & 3l,4k+1 & 2,0\\
E6, (a,b)=(3,2) & 2 & 3l+2,4k & 0,1\\
E7, (a,b)=(1,2) & 2 & 4l,6k+1 & 1,0\\
I1*+IV, (a,b)=(6,2) & 2 & 2,6l+1 & 1,0\\
I1*+IV, (a,b)=(2,6) & 6 & 2,3k, k odd & 3,2\\
\bottomrule\end{tabular}\end{center}

\item All root orders retain the catalogue's positivity and primitive
sphere restrictions. For one cap, \(P_i\equiv\varepsilon\pmod b\)
and the character is \(\zeta_b^{-1}\). For a two-cap row
the condition is the reduction of
\(m_jn_i+m_in_j=\varepsilon\) modulo two and/or three.
For example in the last row, \(3kn_*+2n_4=\varepsilon\)
gives \(3kn_*+3\varepsilon\equiv0\pmod6\) and
\(2n_4+2\varepsilon\equiv0\pmod6\). The other rows have
the identical displayed reductions. Every order-one cap has
\(\chi_i=1\) and \(P_i=Mn_i\), divisible by \(b\), so it
does not change these reductions. Their total contribution to
the sphere equation is also divisible by \(b\).

\item A rational function on \(\mathbb P^1\) whose divisor is
divisible by \(b\) has a rational b-th root: factor its zeros
and poles, divide their integer orders by \(b\), and choose a
b-th root of its nonzero constant. Choose this root \(g_0\)
with the specified cusp value. Then the actual ordinary endpoint
map is
\[
                    D_r=r^{-\varepsilon M/b}g_0(t)U^\varepsilon.
                                                               \tag{P3K.35.14}
\]
The exponent is an integer. On the cap root, its apparent power
of \(z\) is \(-Mn_i/b\). The actual ordinary-to-cap parameter
gluing changes by \(+Mn_i\log z\,\delta/(2\pi ib)\)
under \(r\mapsto\zeta_b r\). This cancels that power exactly.
The remaining map is an invariant unit times \(L_i^\varepsilon\),
and is therefore invertible through every root component. Its
conjugation is exactly
\[
 L_i^\varepsilon A_{i,r}L_i^{-\varepsilon}
      =\chi_i^{-\varepsilon}A_{i,r}
      =\zeta_b^{P_i}A_{i,r}=A_{i,\zeta_b r}.            \tag{P3K.35.15}
\]
The same map commutes with \(B_\lambda\), since
\(\lambda=(\zeta_b r)^M=r^M\). This is the exact quotient
map on the actual cap, with no lost finite character.

\item The weighted Jacobi marking checks the same map integrally.
For a source monodromy \(E\) and its marked normal-function
column \(k\), the actual cohomology lift is
\[
\resizebox{\linewidth}{!}{$\displaystyle
 T(E,k,z)=\begin{pmatrix}
 1&b(E_{21}k_1-E_{11}k_2)&b(E_{22}k_1-E_{12}k_2)&z\\
 0&E_{11}&E_{12}&ak_1\\
 0&E_{21}&E_{22}&ak_2\\
 0&0&0&1
 \end{pmatrix}.
$}
\]
For \(A=T^{-t}\), literal multiplication gives
\[
 (A-I)e_\gamma=(0,bk_2,-bk_1,-z)^t.
\]
Thus the one-step lattice cochain is
\[
 \frac{P_i\delta-\varepsilon(A_i-I)e_\gamma}{b}
    =(0,-\varepsilon k_2,\varepsilon k_1,
                     (P_i+\varepsilon z_i)/b)^t.
                                                               \tag{P3K.35.16}
\]
Its integrality is precisely (P3K.35.15). The source columns in
this formula are retained integral source periods; the last
entry is supplied by the actual line commutator above.

\par\medskip\noindent\textbf{5. Exact cyclic norm and closed cusp endpoint origin.}\enspace

\item At the selected I1 choose the actual bounded Tate lift of \(G\).
The outward \(bG\) period then has source unit its b-th power.
This fixes its logarithm \(\Gamma=b\mu_G\) in the marked source
cover. In the actual cusp notation
\[
 H=\varepsilon s+U(q),\qquad
 \gamma=(\Gamma,C-s-Mv),\qquad v=\frac{\log r}{2\pi i}.
\]
The endpoint (P3K.35.14) has the bounded lift
\[
       \tau=\frac{H\delta+\varepsilon\gamma}{b}
          =\binom{\varepsilon\Gamma/b}
                  {(U+\varepsilon C-\varepsilon Mv)/b}.
                                                               \tag{P3K.35.17}
\]
It multiplies every cusp Cox character by its literal bounded
unit, so it extends simultaneously through all three curves,
both triple points, and every allowed periodic fan.

\item There is a harmless but necessary scalar norm correction. Set
\[
 \alpha=\frac{\varepsilon M(b-1)}{2b^2},\qquad
 \widetilde D_r=e^{2\pi i\alpha}D_r.
                                                               \tag{P3K.35.18}
\]
This is an actual constant physical-delta translation and changes
none of the cap conjugations. It must be retained even when
\(b=2\) and \(M/b\) is odd. Since \(v\) advances by \(1/b\),
direct summation gives
\[
 \sum_{j=0}^{b-1}(\tau(v+j/b)+\alpha\delta)
                         =b\tau(v)=H\delta+\varepsilon\gamma.
                                                               \tag{P3K.35.19}
\]
On the ordinary part the same equality follows directly from
\(U^b\rho=B_0^{-1}\) in (P3K.35.7) and the parameter factor
\(r^{-\varepsilon M/b}\). Hence the b-th iterate of the
endpoint is exactly the already specified one-gamma cover
endpoint, including its line scalar. The cyclic quotient of the
parameter family is taken in that order: first work over the
universal \(v\)-plane, where the proposed map covers
\(v\mapsto v+1/b\). Its b-th iterate equals the prescribed
\(v\mapsto v+1\) gluing map just calculated. Quotient by this
latter map to obtain the already constructed r-family. The first
map then induces an order-b action over \(r\mapsto\zeta_b r\).
In particular the nonconstant multiplier \(h_0\) is retained
as part of the r-family gluing; it is not declared the identity
in an unglued ordinary chart. The quotient of the
proper r-family by this induced deck action is therefore an actual
proper holomorphic family over the punctured w-disc. The action
is free because it rotates the nonzero parameter.

\item Use the already specified cover cusp origin
\[
 \sigma(v)=\binom{\varepsilon\Gamma v}
                {v(U+\varepsilon C)-\varepsilon M(v^2-v)/2}.
\]
Its exact difference is
\[
               \sigma(v+1/b)-\sigma(v)=\tau+\alpha\delta.
                                                               \tag{P3K.35.20}
\]
Thus the same origin, without changing the fan or any ray label,
is a closed simultaneous Cox product for the w-family. Its
pullback is exactly the closed r-product of K30/K32. The scalar
correction is physical delta; it acts on the actual quaternion
coordinate by right scalar multiplication and fixes the Hopf
normal vector. The bounded torus-unit comparison is K30's
literal \(q_{\rm H}\mapsto A q_{\rm H}B\); it has no free
normal-frame component. All initial normal and exterior markings
are therefore the ones used by K31, not just abstractly
isomorphic bundles.

\par\medskip\noindent\textbf{6. Integer characteristic naturality and the fan class.}\enspace

\item The single w-loop has the actual vertical-edge semifree circle,
its standard fixed sphere, and the closed normal framing just
constructed. Its b-fold parameter pullback is the prescribed
r-loop with exactly the K30/K32 framing. Consequently its
ordinary characteristic numbers satisfy
\[
 b k_w=-3H,\qquad b\ell_w=-24\chi.
\]
The left sides are ordinary integers of this actual mapping
torus. Its dual index is \(A_{U,w}=(4k_w-\ell_w)/24\), and the explicit cap duality of \cref{night:P3K.68} assigns the physical central class \(A_{M,w}=-A_{U,w}\pmod{28}\). This establishes (P3K.35.4) by naturality on
the constructed covering, not by choosing a residue-class
quotient. The full outer coefficient maps are retained in \cref{night:P3K.73}; applying its marked-prism algorithm to this constructed shorter endpoint gives the remaining fan returns.

\item The exact certificate \path{code/part3_fourfold/weighted_shorter_loop.py}
checks the full Jacobi gamma-column identity, the two I2 frame
compatibility equations, all twelve finite character rows over
their complete residue domains, and the exact origin/norm
identities (P3K.35.19)--(P3K.35.20). Its JSON and captured output are
retained alongside it. The explicit root-line and global
meromorphic maps above are the geometric proof; the modular
computation does not replace them.

\item The same certificate reads F53's full thirty nonempty branch
polynomials and proves, over \(\mathbb Z[l,k]\), that
\(M/b\), \(-3H/b\), \(-24\chi/b\), and
\((\chi-H/2)/b\) all have integer coefficients. It also checks
\(4k_w-\ell_w=24A_{U,w}\) identically on every branch. The
determinant-excluded good/good branch remains excluded. Thus
integrality is checked on the complete symbolic catalogue,
rather than on a bounded selection of root orders.

\par\medskip\noindent\textbf{Audit scope.}\enspace

\item The component-supported OS49 cusp has its separate mixed-circle
marking and is not included here. Source-moduli loops are also
not assigned by this fixed-source scalar endpoint construction.
The proved original K33 b=6 descent is recovered after its
constant physical-delta normalization; this does not introduce
a different source-family or fan marking.

\end{enumerate}
\end{proof}




\begin{theorem}\label{night:P3K.37}
\StatementRecord{night:P3K.37}{false as stated; correction proved}{S6\_06\_BLIND: GPT-6 Astra (OpenAI, ChatGPT Pro mode, effort Pro), recorded re-derivation in the checkpoint dated 18 November 2026}{\cref{night:P3F.26}, \cref{night:P3F.29}, \cref{night:P3F.31}, \cref{night:P3K.22}, \cref{night:P3K.31}, \cref{night:P3K.34}, \cref{night:P3K.35}}
Keep the exact source, component line, cap and fans of F26/F29/K34:
\(M=O(G-O)\), suspension by \(-3G\), source \(I_2\) at
\(q=1/t=0\), cap order \(m=3k\), \(k>0\), norm
\(n=\pm1\), and arbitrary fixed delta character \(d\).
The already constructed parameter is \(\lambda=r^{3k}\).
The construction below produces an actual family over
\[
 w=r^3,\qquad \lambda=w^k.                                      \tag{P3K.37.1}
\]
Its initial positive-diagonal closed Cox pair is
\[
 (k_{\rm Jin},\ell_{\rm Jin})=(-2k,-8k),\qquad
 (m_{\rm rel},n_{\rm rel},s)=(2k,0,0).
                                                                    \tag{P3K.37.2}
\]
Its physical reference value on the positive fan is zero; on the
negative fan the same actual shorter loop has class (-k), by
\cref{night:P3K.65}. The full fan transport is
\cref{night:P3K.77}. The old polynomial is refuted already at
(k=1) on the negative fan, where it gives zero and the actual
class is (27). These comparisons use the constructed covering
and endpoint maps, with no division in \(\Theta_7\).

\end{theorem}
\begin{proof}
\begin{enumerate}[label=\arabic*.]
\item \par\medskip\noindent\textbf{1. The actual meromorphic +G lift and its component norm.}\enspace

\item Let \(B_0\) be the actual contracting line lift over \(-3G\), before
its scalar parameter is inserted, so \(B_\lambda=\lambda B_0\).
The target-point linearization has divisor \(C_0\); as a map of
source-point line fibres, \(B_0\) has its zero on source \(C_1\),
which is mapped to \(C_0\). This is the exact graph calculation in F31.
Consequently \(B_0^{-1}\), over \(+3G\), has its simple pole on
source \(C_0\).

\item There is an actual rational lift \(U\) over \(+G\). In the rational
section frame of \(M=O(G-O)\), the function
\[
 f=\frac1{x-x(G)},\qquad
 \operatorname{div}(f)=2O-G-(-G)
       =t_G^*(G-O)-(G-O)
\]
defines the isomorphism \(t_G^*M\to M\); its inverse gives \(U\).
All its remaining zeros and poles are vertical. The source translation
extends on the actual relatively minimal regular elliptic source, by the
canonical-form graph argument already proved in K22.

\item At the I2 cusp use F31's literal line graph
\(h_j=\lceil j/2\rceil\), \(h_{j+2}=h_j+1\).
The forward +G source shift is \(j\mapsto j+1\). After multiplying
\(U\) by an integral base monomial, its line multiplier has valuation
zero. The two actual graph defects are
\[
 h_j-h_{j+1}=-1\quad(j\text{ even}),\qquad
 h_j-h_{j+1}=0\quad(j\text{ odd}).                              \tag{P3K.37.4}
\]
Thus the pole divisor is \(-C_0\). This includes both node charts,
because their endpoint graph defects determine the two fibre-coordinate
monomials. The third translated norm has divisor \(-2C_0-C_1\).
The maps \(B_0^{-1}\) and \(U^3\) cover the same translation,
so their ratio is a rational base function
\[
 \rho=B_0^{-1}/U^3,\qquad \operatorname{ord}_{q=0}\rho=1.         \tag{P3K.37.5}
\]
The ratio is constant on each compact generic elliptic fibre. With another
whole-fibre normalization of \(U\), the last order is \(1-3e\),
which is the same exact congruence needed below.

\item In particular the component divisor has not been replaced by a whole fibre.
It is precisely the unequal translated coefficients in (P3K.37.4) that produce
the whole-fibre ratio in (P3K.37.5).

\par\medskip\noindent\textbf{2. Actual cap root lift and its character.}\enspace

\item On the actual order-\(m\) good root at IV*, the line is the degree-zero
root line with Picard point \(G'\). Its marked central point has order
three. In the explicit root \(t=z^3,x=z^4X,y=z^6Y\), it is
\((X,Y)=(0,1)\) on \(Y^2=X^3+1\); the deck derivative on the
elliptic curve is \(\varrho=\zeta_3^2\). For the order-\(3k\)
root the same effective deck acts on the central curve.

\item The translation by +G' has an invertible line lift \(L\) on this actual
root line. One may construct it in the universal flat-line model
\((\mathbb C\times\mathbb C^*)/\Lambda\), and then compare it to
the rational pulled-back \(U\). Their ratio is a meromorphic base scalar;
on the smooth central fibre its only divisor is a whole fibre, so removing
that base power gives the extending invertible lift. Translation lifts on
this degree-zero line commute, as their universal representatives are a
source translation and a constant fibre scalar.

\item The exact Cartier line has deck weights \(\varrho\) at G',
\(\varrho^{-1}\) at O, and 1 at the remaining point of the order-three
orbit. Translation by +G' cyclically permutes them. Their successive ratios
are all \(\varrho^{-1}=\zeta_3\), giving
\[
 DLD^{-1}=\zeta_3 L.                                             \tag{P3K.37.6}
\]
The finite cyclic averaging of K35 removes any nonconstant unit in this
commutator; an invariant analytic cube root then fixes the actual norm
\(L^3=B_0^{-1}\). The fixed \(d\)-character is scalar and commutes
with this map.

\item Write the global rational lift on the root as
\(U=z^e a(z)L\). Global deck invariance and the target-scalar identity
\(D m_f D^{-1}=m_{f(\zeta_m^{-1}z)}\) give
\[
 \zeta_m^e=\zeta_3,\qquad a(\zeta_mz)=a(z),\qquad
 \operatorname{ord}_{t=0}\rho=-\frac{3e}{m}\equiv-1\pmod3.
                                                                    \tag{P3K.37.7}
\]
At the retained type-II fibre, and at every other ordinary fibre, the
irreducible vertical divisor of \(U\) is an integral whole-fibre divisor;
the order of \(\rho\) is therefore a multiple of three. No cap or line
character has been dropped in these calculations.

\par\medskip\noindent\textbf{3. The global rational endpoint.}\enspace

\item K34's one-r-turn ordinary CLOSING map is multiplication by
\(h_0=t^n=q^{-n}\). Its arrow is from the continued endpoint to the
initial chart. Equations (P3K.37.5)--(P3K.37.7) prove that
\[
                    h_0\rho^{-2n}
\]
has divisor divisible by three at every source fibre. At the I2 cusp its
order is \(-3n\) modulo multiples of three; at IV* its order is
\(n+2n=3n\) modulo three; elsewhere it is divisible by three.
Factor the rational function on \(\mathbb P^1\), and fix a cube root
of its nonzero constant. This gives an actual rational function
\[
                    g_0^3=h_0\rho^{-2n}.                         \tag{P3K.37.8}
\]
The proposed ordinary closing map is
\[
 C_r=r^{2nk}g_0(t)U^{-2n}: X_{\zeta_3r}\longrightarrow X_r.
                                                                    \tag{P3K.37.9}
\]
Its inverse is the forward parameter deck. At an ordinary fibre its base
power cancels the whole-fibre power of \(U^{-2n}\), so it is an
invertible map of the actual line-complement quotients. The exact maps
commute with \(B_\lambda\); the scalar \(\lambda\) is unchanged
under \(r\mapsto\zeta_3r\).

\item On the IV* root, inserting \(U=z^eaL\) and (P3K.37.7) leaves the precise
base power \(z^{mn/3}\) times an invariant unit and \(L^{-2n}\).
This is the power used in the cap comparison below, not a discarded pole.

\par\medskip\noindent\textbf{4. The typed finite-cap commuting square.}\enspace

\item Retain F26/F29's physical convention \(V=e^{2\pi i\zeta_2}\).
Its root action is
\[
 A_r=r^n A_0,
\]
because the actual cap period is
\(\ell=n(e_\gamma-e_u)+d\delta\) in F26's homology marking and the
gamma period has line multiplier \(\lambda=r^m\).
The displayed plus-log map is ordinary to cap:
\[
 \Phi_r(z,s)=z+\frac{\log s}{2\pi i}\Pi_r(s)\ell.
\]
Under \(r\mapsto\zeta_3r\), its parameter difference is
\((mn/3)\log s\,\delta/(2\pi i)\). It is exactly the positive
power left in (P3K.37.9). Thus the actual punctured maps satisfy
\[
 \Phi_r C_r=C_{\rm cap}\Phi_{\zeta_3r},\qquad
 C_{\rm cap}=\text{invariant unit}\cdot L^{-2n}.                 \tag{P3K.37.10}
\]
This bounded cap map extends across the root fibre. Its conjugation is
\[
 C_{\rm cap} A_{\zeta_3r} C_{\rm cap}^{-1}
   =\zeta_3^n\,(\zeta_3)^{2n} A_r=A_r.                         \tag{P3K.37.11}
\]
It therefore descends to the exact cap, with the original \(d\)-character.
The opposite direction would reverse both the parameter deck and this
conjugation; (P3K.37.10) specifies the actual direction used here.

\par\medskip\noindent\textbf{5. The cusp charts, cyclic quotient and identical Cox origin.}\enspace

\item Use K34's actual periods
\[
 u=(2s,s)+U(q),\qquad
 \gamma=(-3s,-s+mv)+C(q),\qquad
 K=3u+2\gamma=D+(s+2mv)\delta,
\]
where \(D=3U+2C\), \(v=\log r/(2\pi i)\), and
\(s=\log q/(2\pi i)\). The source map of \(U\) has Tate
valuation +1 and, after the chosen base normalization, line valuation 0.
Hence (P3K.37.9) has tropical vector \((-2n,-n)\). Reindexing by the
actual source period \(h^n=(2n,n)\) makes it bounded on every toric
chart, not only on the fixed curve. This period preserves every allowed
index-one fan and every cell label. Choose K34's logarithms coherently with
the actual bounded Tate unit of G: the source unit of \(B_0\) is the
inverse cube of that unit, because both are the same marked source translation
and their valuations already agree. This specifies the first-coordinate
logarithm, including its integral period. Changing a previous logarithm by an
integer supplies the literal bounded unit torus loop; K34's quaternion
identity carries that loop as a normal rotation constant along the fixed
sphere, so its free framing class is unchanged. Retain every line unit as
follows. Put \(R=\log\rho/(2\pi i)\), let \(L_U\) be the coherent
Tate logarithm of \(U\), and put \(g=\log g_0/(2\pi i)\). The
actual norm and rational root equations give, with the constant choices
retained,
\[
 3L_U=-\gamma+mv\delta-R\delta,\qquad
 3g=-ns-2nR.
\]
Consequently the logarithm of \(C_r\), after the actual \(h^n\)
reindexing, is literally
\[
 nu-2nL_U+(2nkv+g)\delta
   =\frac{nK}{3}-\frac{ns}{3}\delta
   =\frac{nD}{3}+\frac{2nmv}{3}\delta.
\]
Thus all nonconstant units cancel by their actual norm identity, not by an
omitted change of frame. The bounded logarithm is exactly
\[
 \tau(v)=\frac n3\bigl(D+2mv\delta\bigr).                      \tag{P3K.37.12}
\]
Changing the cube root in (P3K.37.8) by the actual cube root of unity
\(e^{-4\pi i nk/3}\) gives the closing map
\(\widetilde C_r=e^{2\pi i\alpha}C_r\),
\(\alpha=-2nk/3\). This constant is physical delta.

\item The already specified K34 origin, with no new framing choice, satisfies
\[
 \sigma(v)=nvD+nmv(v-1)\delta,
\]
\[
 \sigma(v+1/3)-\sigma(v)=\tau(v)+\alpha\delta,
\qquad
 \sum_{j=0}^{2}\bigl(\tau(v+j/3)+\alpha\delta\bigr)
       =nD+2nmv\delta.                                         \tag{P3K.37.13}
\]
The threefold ordinary composition is exactly
\[
 \widetilde C_r\widetilde C_{\zeta_3r}
       \widetilde C_{\zeta_3^2r}
     =h_0 B_\lambda^{2n};                                      \tag{P3K.37.14}
\]
the parameter-root phase is \(e^{2\pi i2nk}=1\), and the constant
correction has cube one. The period \(B_\lambda^{2n}\) is already
divided out in the ordinary family. Thus this is exactly its prescribed
one-r-turn closing, with the cap square (P3K.37.10), not just its action on
cohomology.

\item First use these maps over the universal v-plane. Their third iterate is
the existing v-to-v+1 family gluing. Quotient first by that gluing, obtaining
the already constructed r-family. The inverse of (P3K.37.9) then induces an
order-three deck over \(r\mapsto\zeta_3r\). It is free because
\(r\ne0\). Its quotient is an actual proper holomorphic w-family with
(P3K.37.1). The toric bounded units in (P3K.37.12) act on every Cox character,
so this quotient retains all three curves, both triple points, and every
allowed fan simultaneously.

\item In K34's product convention the coordinate translation by \(\sigma\)
matches the endpoint-to-initial closing (P3K.37.13). Thus the same Cox product
closes on the w-loop, and its pullback is literally the original r-product.
The constant correction is physical delta. In K34's actual quaternion frame
it is right multiplication by a scalar and fixes the Hopf normal vector.
The bounded unit comparison is the same \(q_H\mapsto Lq_HR\),
independent of the fixed-sphere point. The decreasing-K reflection and both
norm signs remain at both endpoints exactly as in K34; no free normal-frame
term is introduced.

\par\medskip\noindent\textbf{6. The actual framed covering and the return.}\enspace

\item For the initial positive fan, the w-loop now has an actual closed normal
framing. Its threefold parameter pullback is the exact K34 loop and framing.
The induced mapping-torus covering preserves the physical circle Euler line,
the complex orientation and the prescribed spin structures. Ordinary
characteristic-number naturality therefore gives
\[
 3k_{\rm Jin}=-6k,\qquad 3\ell_{\rm Jin}=-24k.
\]
This proves (P3K.37.2) as equations of integers of the constructed mapping torus.
Its physical reference index is zero. The complete fan transport uses the actual outer maps in \cref{night:P3K.77}; the reference fixed-star covering does not imply the rejected norm polynomial.

\item The executed exact certificate is
\path{code/part3_fourfold/os49_shorter_loop.py}. It retains the homology-to-positive
Tate sign, rather than identifying those two u coordinates silently. In the
displayed F26 homology basis,
\(u_{\rm Tate}=-e_u\), \(K=2e_\gamma-3e_u\), and
\[
 (A-I)K=-3e_w-\delta,\quad
 (B-I)K=3e_u+3e_w,\quad (C-I)K=\delta.
\]
The integral one-step cochain is explicitly
\[
 p=\frac{-n\delta-n(A-I)K}{3}=ne_w,\qquad
 q=-\frac{n(B-I)K}{3}=-n(e_u+e_w),
\]
\[
 p+Aq=0,\qquad q=(B-I)(ne_u),\qquad
 k(1+A+A^2)p=-nk\delta.
\]
Moreover \(-n\delta-3p=n(A-I)K\) and \(-3q=n(B-I)K\),
so the threefold comparison retains the exact integer period primitive.
These agree with the literal local periods and the finite-root logarithmic
power. The component defects, character/divisor identities, cyclic norm,
bounded origins and characteristic substitutions all pass in the same
certificate; its JSON and captured output are retained alongside the script.


\item The constructed endpoint, finite-root commuting square and
framed covering above remain the geometric descent proof.
The positive fan has the displayed dual pair ((-2k,-8k)),
whose physical index is zero. The explicit component-normal
comparison of \cref{night:P3K.65} gives native pair ((-k,k))
for the negative-to-positive flop, hence negative class (-k).
At (k=1) this is (27), refuting the old zero at norm (S=1).
Every further fan is computed by \cref{night:P3K.77}; the old
polynomial witnesses and claimed even-only range are archived
with the rejected compound assertion.

\end{enumerate}
\end{proof}




\begin{theorem}\label{night:P3K.40}
\StatementRecord{night:P3K.40}{false as stated; correction proved}{S6\_06\_BLIND: GPT-6 Astra (OpenAI, ChatGPT Pro mode, effort Pro), recorded re-derivation in the checkpoint dated 18 November 2026}{\cref{night:P3F.56}, \cref{night:P3K.23}, \cref{night:P3K.31}}Retain the actual positive selected-cusp source loop of P3F.54 and
P3F.56, with fixed sufficiently small scalar, the III-star cap, no
logarithmic modification at either member of the coalescing finite
pair, and the initial physical-delta-edge fan. Its parameter is one
positive \(\rho\)-loop with \(a(u)=\rho^2\), the actual twofold
marked closure of the selected source puncture. Flop that fixed edge
in the periodic cusp completion, carrying its actual source, line,
cap and ordinary-annulus maps. The positive return of the resulting
family has oriented homotopy-sphere class
\[
                 [\Sigma_{\Phi,\mathrm{flop}}]
                     = -7=21\quad\hbox{in }\mathbb Z/28.
\]
In the actual first-resolution relative product, its initial right
exponent is \(n=-6\). The other free exponent \(m\) is retained
without assigning its value. This statement assigns neither the
separate Jin characteristic pair nor a return for another edge or
an arbitrary fan.
\end{theorem}
\begin{proof}
\begin{enumerate}[label=\arabic*.]
\item We retain the complete source, line, cap and frame maps below.
The meridian-one graph line is (U=M^{-1}): the fibre map
(D_M(v)(w)=\langle v,w\rangle/\langle v,v\rangle), with the
Hermitian form linear in its second argument, satisfies
(D_M(e^{i\theta}v)=e^{-i\theta}D_M(v)).
Its (SU(2))-equivariant spinor extension is
((a,b)\mapsto(\overline b,-\overline a)).
The first-to-second cap maps in \cref{night:P3K.68} give
(\lambda_M(F)=-\Sigma_{F^{-1}}) and
(\lambda_U(F)=\Sigma_{F^{-1}}); cancelling the opposite
product collars gives (A_M=-A_U). We therefore distinguish the
intermediate meridian-one number (A_U=(4k_U-\ell_U)/24)
from the physical return (A_M) throughout this calculation.

\item \par

\medskip\noindent\textbf{SF.1. The closed product is supplied by the actual source coordinates}

\par\nobreak

\item Write the trace of the fixed curve, over a small punctured rho-disc, as
\(F=\mathbf P^1_s\times D_\rho^*\). Its two distinguished sections are
\(D_+=\{s=\rho\}\) and \(D_-=\{s=-\rho\}\). The literal ordered
coordinate and inverse are
\[
 Z=\frac{s-\rho}{s+\rho},\qquad
 s=\rho\frac{1+Z}{1-Z}.                                      \tag{SF.1}
\]
This is an actual biholomorphism from \(F\) to
\(\mathbf P^1_Z\times D_\rho^*\). Its determinant is \(2\rho\), so
it has no zero on this base. In particular it is single-valued after
one rho turn, and retains the two ordered endpoint sections.

\item The relative small-resolution charts used in P3F.56 have normal
coordinates \(c_0,c_1\), with the actual overlap
\[
             c_i'=Zc_i,\qquad i=0,1,                       \tag{SF.2}
\]
and physical delta weights \((-1,+1)\). Thus their normal bundle is the
two copies of \(\mathcal O_{\mathbf P^1}(-1)\) in precisely the
coordinate (SF.1). This is also the Cartier identification used in
P3F.56: after conjugating the negative-weight normal coordinate its
spinor is
\(\mathcal O(D_+)\oplus\mathcal O(-D_-)\), and the determinant
trivialization is the ratio \(Z\). We use that same identification.

\item Here is why no source-scaling loop is discarded in using (SF.2).
The source chart and its actual Tate parameter satisfy
\[
 \xi\eta=q,\qquad
 \frac{q_T}{q}=-\frac{u(\rho)}{64\rho^6}+O(q).               \tag{SF.3}
\]
The coefficient in (SF.3) has winding \(-6\) on the rho-circle.
Choose the source collar uniformly smaller than its distance from the
second moving critical value; then this coefficient is an actual
nonzero holomorphic function on the whole parameter annulus and
normal coordinate polydisc. Every remaining nonzero toric chart unit
\(U\) on such a representative has an exact factorization
\[
                    U=\rho^{e}\exp H,\qquad e\in\mathbb Z. \tag{SF.4}
\]
Indeed \(e\) is its integral winding on the rho generator; subtracting
\(e\,d\rho/\rho\) from \(dU/U\) kills its only possible period,
and integration gives the single-valued holomorphic \(H\). This
argument is on the annulus, not across its missing centre. The integer
\(-6\) in (SF.3) is retained as one of these exponents.

\item The source and line units in P3F.56 extend on the original source
polydisc. Their changes to the Tate cell use (SF.3), so the preceding
factorization is applied only after that exact coordinate change.
The period action has trivial stabilizer on an individual conifold
cell. We choose one representative of the cell containing the fixed
edge and transport its charts to all translates. Thus (SF.4) and its
chosen logarithm give the same maps on both resolution charts and
on their actual ordinary annulus. Contracting \(H\) and the positive
modulus transports their common truncation collar by the resulting
isotopy; it leaves the integral angular powers in place. The source
coordinate (SF.1) is carried along, rather than identified with the
old source coordinate by an unmarked homotopy.

\item Equivalently, for the linear normal maps this assertion can be checked
over the whole compact fixed curve. Delta-equivariance preserves its
two distinct weight lines. After (SF.1)--(SF.2), an automorphism of
either line is multiplication by a holomorphic function on
\(\mathbf P^1\), hence by a function of rho alone. Its integer winding
is still retained by (SF.4). Higher normal terms are removed by the
literal equivariant normal rescaling and isotopy extension on the
same collar. The linearization homotopy has invertible diagonal
blocks throughout and starts with the actual map. No change of an
outer gluing or of an initial normal frame is made.

\par

\medskip\noindent\textbf{SF.2. The two relative family products have the same boundary arrow}

\par\nobreak

\item The common conifold link is
\[
 (S^3_a\times S^3_c)/S^1,\qquad
 (a,c)\sim(e^{i\phi}a,e^{-i\phi}c).
\]
Put, in the quaternion convention of P3K.23,
\[
 u=a_0+ja_1,\qquad v=c_0+j\overline{c_1},\qquad
 q_H=uv,\quad h=uiu^{-1},\quad k=v^{-1}iv.                 \tag{SF.5}
\]
The map to the first product is \([a,c]\mapsto(h,q_H)\), and that to
the second product is \([a,c]\mapsto(k,q_H)\). Their inverses choose
one Hopf representative and use respectively \(v=u^{-1}q_H\) or
\(u=q_Hv^{-1}\); changing the representative is exactly the displayed
auxiliary circle. These are the actual normal-sphere products of the
two small resolutions.

\item On the first source chart (SF.5) is precisely P3F.56's retained map
\[
              q_H=(1+jZ)(c_0+j\overline{c_1}),             \tag{SF.6}
\]
after the common positive normalization of its norm. The positive
normalizing factor extends across both charts and is carried along
in both resolutions. Thus the local identification in this proof is
the one in the proved trace comparison, including its normal frame.

\item The exact common boundary map, computed from the same two factors,
is
\[
               \tau(h,q_H)=(q_H^{-1}hq_H,q_H).             \tag{SF.7}
\]
It has no rho dependence. To verify this for all retained integer
chart powers, give the four Cox coordinates \((a_0,a_1,c_0,c_1)\)
their actual integer weights \((A_0,A_1,C_0,C_1)\). Such integer
lifts exist because the auxiliary-circle vector \((1,1,-1,-1)\)
is primitive; the quotient cocharacter lattice is free. We impose
no condition on their sum. Put
\[
 \begin{aligned}
 a_\theta&=e^{i(A_0-A_1)\theta/2},&
 d_\theta&=e^{i(A_0+A_1)\theta/2},\\
 e_\theta&=e^{i(C_0+C_1)\theta/2},&
 b_\theta&=e^{i(C_0-C_1)\theta/2},\\
 S&=A_0+A_1+C_0+C_1.
 \end{aligned}
\]
The literal quaternion changes are
\(u'=a_\theta u d_\theta\) and
\(v'=e_\theta v b_\theta\). Therefore
\[
 \begin{split}
 h'&=a_\theta h a_\theta^{-1},\qquad
 k'=b_\theta^{-1}k b_\theta,\\
 q_H'&=a_\theta u e^{iS\theta/2}v b_\theta\\
     &=a_\theta e^{hS\theta/2}q_H b_\theta
       =a_\theta q_H e^{kS\theta/2}b_\theta .
 \end{split}                                             \tag{SF.8}
\]
Here \(h=u i u^{-1}\), \(k=v^{-1}iv\), and both displayed
exponential identities follow by conjugating the same middle
complex scalar. Substitution into (SF.7) gives
\[
 (q_H')^{-1}h'q_H'
 =b_\theta^{-1}e^{-kS\theta/2}
              (q_H^{-1}h q_H)e^{kS\theta/2}b_\theta
 =b_\theta^{-1}k b_\theta=k'.
\]
This proves equivariance even when \(S\ne0\). At \(\theta=2\pi\)
the factors \(a_\theta,d_\theta\) are central signs whose product
is \((-1)^{2A_0}=1\), and the product of the two signs
\(e_\theta,b_\theta\) is \((-1)^{2C_0}=1\). Thus both actual Cox
factors, and consequently both product actions, close. The possible
nonzero middle exponential has been retained.

\item Formula (SF.8) consequently retains all integer windings from
(SF.3)--(SF.4), including the actual sixth source scaling. It does not
require evaluating or discarding any of them. The same comparison
holds during the common unit isotopy in SF.1, so its initial and
final ordinary-annulus maps are the original ones. The earlier
working passage with \(S=0\) is preserved as superseded in
\path{code/part3/attempts/source_cusp_fixed_edge_sum_zero_restriction/};
that restricted passage was not promoted as a theorem.

\item Choose the closed product connection just described on the fixed
curve tube and equivariant connections on its complement; an
invariant partition of unity patches them while retaining the fixed
collar. The resulting physical return is the identity near that
curve. The same choice in the flopped family gives the identical
common-exterior return, conjugated only by the extension of (SF.7),
\[
 T(y,q_H)=(q_H^{-1}yq_H,q_H),\qquad
 T^{-1}(y,q_H)=(q_Hyq_H^{-1},q_H).                         \tag{SF.9}
\]
It is smooth at \(y=0\). The first and second normal conventions,
their complex orientations and their physical Hopf circle are exactly
those in P3K.23. This verifies the family/framing hypothesis required
for that statement, for the source loop itself.

\par

\medskip\noindent\textbf{SF.3. The actual relative integer and the flopped class}

\par\nobreak

\item Write the actual first relative return in the proved relative group
as \(L^mR^n\mathcal C^s\), using P3K.23's maps
\[
 L(y,q_H)=(y,\zeta(y)q_H),\qquad
 R(y,q_H)=(y,q_H\zeta(y)),
\]
and its positively oriented central generator \(\mathcal C\).
We retain the actual, presently unevaluated integer \(m\).
The central coordinate is the oriented closed six-sphere return,
so the explicit physical transfer of \cref{night:P3K.68} gives \(s=-1\pmod {28}\).

\item The coefficient \(n\) is fixed without a separate normal
Pontryagin calculation. We spell out the comparison being used.
The first quaternion normal product has physical circle
\(q_H\mapsto q_He^{-i\theta}\), so the exact Hopf quotient has
orientation \(S^2_z\times S^3_y\). Its relative variation is
\(-k_{\rm Jin}\). The actual generator \(L\) has quotient
characteristic pair \((-1,-4)\), computed from its adjoint clutch
with \(p_1=-4\). Removing its variation therefore changes the
characteristic pair to \((0,\ell-4k)\). The remaining free relative
map is \((LR^{-1})^b\). On its fixed three-sphere its normal clutch
is \(\operatorname{Ad}_{\zeta^b}\), hence its Pontryagin evaluation
is \(-4b\). The integral Gysin pullback is injective on the relevant
fourth cohomology and the actual marked cycle evaluates to one.
Consequently, in these same products,
\[
           \ell-4k=-4b,
           b=6\frac{4k-\ell}{24},\qquad
           n=-b=-6A_U.                                    \tag{SF.10}
\]
Every arrow and orientation in this argument is the one verified in
P3K.23, now applicable by SF.1--SF.2; it does not import the scalar-loop
value of either \(k\) or \(\ell\). Since the actual source loop has
\(A_U=1\) in the dual-line computation of P3F.56, it gives \(n=-6\).

\item The literal conjugation in (SF.9) is
\[
 T(L^mR^n\mathcal C^s)T^{-1}
       =L^{-n}R^{m+2n}\mathcal C^{s+n}.                   \tag{SF.11}
\]
It follows that the actual positive source-loop return after its
first fixed-edge flop has class
\[
                          -1-6=-7=21\pmod {28}.           \tag{SF.12}
\]
This computation retains the actual source parameter, fixed scalar,
cap norm, physical delta character, and ordinary annular maps. It
does not replace the source loop by a scalar loop.

\item The transported complex first-resolution reset of P3K.31 changes the
relative pair to \((m-12,6)\), while retaining the same central class
\(-7\). To compute a later flop at another edge one still needs the
source loop's full common-exterior variation and derivative. Neither
\(m\), the separate Jin pair, nor an arbitrary-fan formula is assigned
by the present proof.

\end{enumerate}
\end{proof}




\begin{theorem}\label{night:P3K.41}
\StatementRecord{night:P3K.41}{false as stated; correction proved}{S6\_06\_BLIND: GPT-6 Astra (OpenAI, ChatGPT Pro mode, effort Pro), recorded re-derivation in the checkpoint dated 18 November 2026}{\cref{night:P3F.41}, \cref{night:P3F.42}, \cref{night:P3F.43}, \cref{night:P3F.44}, \cref{night:P3F.45}, \cref{night:P3F.46}, \cref{night:P3F.47}, \cref{night:P3F.48}, \cref{night:P3F.49}, \cref{night:P3F.51}, \cref{night:P3K.15}, \cref{night:P3K.35}, \cref{night:P5.5}}
Fix a marked rational elliptic source in the whole-zero selected
\(I_1\) sphere catalogue, its actual source line, translation
\(Q=bG\), logarithmic sites, root orders \(m_i\), norm integers
\(n_i\), and an index-one periodic fan. Retain every finite line
character and every annular map. Consider loops of the actual affine
source maps and line units of P3F.44/P3F.49 with their prescribed
linear source/lattice marking, together with the scalar \(\lambda\).
Put \(M=\prod_i m_i\), including the order-one and framing-chart
norms in
\[
           \sum_i\frac{n_i}{m_i}=\frac{\varepsilon}{M},
                \qquad\varepsilon\in\{1,-1\}.
\]
The global group of extending fibrewise translations is exactly
the physical \(\mathbb C^*\) action. In particular every such
fixed-scalar cap-data loop has trivial sphere mapping class.

More generally, a closed loop in this marked scope with scalar
winding \(\ell\) satisfies
\[
                       \ell=N\frac{M}{b},\qquad N\in\mathbb Z.
\]
Its physical return is isotopic to the \(N\)-th power of the
actual shorter scalar return of P3K.35, whose parameter is
\(\lambda=w^{M/b}\). Thus the local finite kernels of \(1-D_i\)
produce no additional return class in this specified cap-data scope.
For the transported fan (P), write (A_{
m short}(P))
for the actual shorter return computed by \cref{night:P3K.73}.
The oriented return is (N A_{
m short}(P)).
The proof retains the inverse closing arrow in the
physical monodromy. It concerns the stated affine cap/line/annular
variations; moving logarithmic sites, source-moduli loops, arbitrary
nontranslation endpoint maps, and the component-supported OS49
model are outside this statement.


\end{theorem}
\begin{proof}
\begin{enumerate}[label=\arabic*.]
\item \par

\medskip\noindent\textbf{CP.1. The exact local group and its Lie sheaf}

\par\nobreak

\item Let \(B\) be the orbifold \(\mathbf P^1\) with the actual cap orders. Its selected \(I_1\)
point is ordinary. Let \(A\) be the sheaf of fibrewise translations of the
compactified suspension that extend on its actual root and cusp charts.
All linear lattice maps are fixed. The physical delta field gives an
exact sequence of orbifold holomorphic vector bundles

\[
 0\longrightarrow\mathcal O_B\delta\longrightarrow V
       \longrightarrow L\longrightarrow0.                 \tag{CP.1}
\]

\item The quotient is the source translation Lie line. Here is the exact
coarse extension, including the changed additive root model. In the
minimal Weierstrass substitution

\[
 q=z^m,\qquad x=z^{2k}X,\quad y=z^{3k}Y,
\]

the source relative differential is \(z^{-k}dX/(2Y)\). Its dual source
translation vector therefore maps to \(z^k\) times the root translation
vector. The latter has deck character \(\zeta_m^k\). An invariant holomorphic
section has the form \(z^k h(q)\) times that root vector, because \(0\leq k<m\).
It is thus exactly an original source Lie section, not a different
coarse line. The allowed \(k\) values are \(\lfloor3m/4\rfloor\), \(\lfloor2m/3\rfloor\),
\(\lfloor m/4\rfloor\), \(\lfloor m/3\rfloor\), \(\lfloor m/6\rfloor\), and \(k=1\) for the quadratic I1-star
root. These are respectively the III-star, IV-star, III, IV, II, and
I1-star substitutions already proved in F41--F49. Every one has
\(0\leq k<m\). The root character and this equality are retained under the
actual line diagrams and coefficient-generic maps of F51.

\item On the selected I1 the source vector is the logarithmic Tate vector,
which extends as the relative dualizing vector through the actual
toric charts. The quotient line in (CP.1) has the original coarse
extension \(\mathcal O(-1)\), since the source is rational with \(\chi(\mathcal O)=1\). The
physical delta line has coarse extension \(\mathcal O\). Taking finite invariants
is exact over \(\mathbb C\). Hence the coarse direct image of \(V\) is an extension
of \(\mathcal O(-1)\) by \(\mathcal O\), and the elementary two-chart Laurent calculation gives

\[
                       H^1(B,V)=0.                         \tag{CP.2}
\]

\item At a reduced retained/root fibre the source component group can be
removed from a translation lift by a suspension power: the actual
source translation \(Q\) acts transitively on those components, by the
already proved sphere restriction K15. The remaining central point
lies in the additive or multiplicative identity group and has a local
logarithm. On a good root ordinary elliptic local logarithms apply.
At the selected I1 cusp the index-one component lattice gives the
unique integral toric reindexing making an extending translation
preserve the chosen component; the remaining two torus multipliers
are holomorphic units. Their logarithms give the local Lie lift.
This is the actual local argument of P5.5, which uses the toric cover
and its deck lattice, not a claim based on the dimension of an
automorphism group. In particular

\[
 0\longrightarrow H\longrightarrow V\xrightarrow{\exp}A
       \longrightarrow0                                   \tag{CP.3}
\]

is exact as a sequence of sheaves on the orbifold atlas and its
full groupoid, where \(H\) is the extending integral period sheaf. It
is not an assertion that exponentiation is surjective on invariant
sections of a coarse root disc. A fixed torsion translation can lack
an invariant Lie logarithm; its non-equivariant atlas logarithm
has a finite-group lattice discrepancy, which remains in (CP.3).
Taking finite invariants was used only for the complex vector sheaf
\(V\) in (CP.1)--(CP.2), where averaging is exact.

\par

\medskip\noindent\textbf{CP.2. Parametric conjugacies and the discrete torsor class}

\par\nobreak

\item Work first over a small simply connected parameter disc. F44 and F49
give actual local source conjugations, with the line lifts and their
transported annular maps. Their constructions are analytic in that
parameter after shrinking: a central isogeny has a local inverse at
each chosen inverse image; the multiplicative square root has a local
analytic branch; the residual zero-central logarithm is solved by the
displayed finite averaging sum. The residual line unit has its
central finite character retained and its zero-central logarithm
solved by the same sum. The component choices and this finite
character are locally constant. No choice has yet been made around
the entire parameter circle.

\item At the cusp the scalar and the source line are fixed, so the same
toric translation model is used. Comparing all these local models
with one fixed member gives a cocycle in \(A\) on the full Cech
groupoid of the actual orbifold atlas. This includes the finite-root
group cochains as well as ordinary overlaps. Its image under (CP.3)
is a lattice-valued total two-cocycle. Local atlas logarithms can be
chosen continuously on the parameter disc; they are not presumed
invariant. Their finite-group discrepancies and ordinary overlap
discrepancies are continuous lattice-valued functions and hence are
locally constant. Subtract the reference-parameter cocycle and use
the corresponding integral total-cochain adjustment. The result is
a \(V\)-valued total one-cocycle, including its finite-group terms.
Equation (CP.2) solves this cocycle; finite averaging is valid here
because \(V\) is a complex vector sheaf. Thus no fixed torsion point is
removed by an unjustified invariant-logarithm choice.

\item This solution can be taken analytic in the parameter. On the two
standard source-base charts, the solution is the Cauchy/Laurent
splitting of the negative and nonnegative powers for \(\mathcal O\) and \(\mathcal O(-1)\);
finite averaging gives its orbifold version. The integrals are over
fixed compact circles, so holomorphic parameter dependence is
preserved. A splitting of the vector extension in (CP.1) also exists,
since its extension group has coarse \(H^1(\mathcal O(1))=0\); the same Laurent
calculation constructs it. Exponentiating the resulting sections
gives the actual global family isomorphism. On every local chart it
is the specified extending translation, so it transports the
annular, finite-character, line, and cusp maps literally.

\item Consequently a fixed-scalar cap-data family is locally isomorphic to
the fixed compactified suspension by translations. Continuing these
actual isomorphisms around a parameter interval leaves a global
extending translation at the endpoint. The finite kernel of \(1-D\)
has not been discarded: it is included in that endpoint. The next
step computes this global endpoint group.

\par

\medskip\noindent\textbf{CP.3. Lifting an endpoint away from the selected I1}

\par\nobreak

\item Put \(U=\mathbf P^1\setminus\{c\}\), where \(c\) is the selected
I1. We give the actual local lift, including at every reduced root or
retained fibre. On such a root disc the construction is literally
\[
              X'=N^\times/\langle B_\lambda\rangle .
                                                               \tag{CP.3a}
\]
Here \(N\) is the actual corrected multidegree-zero root line, and
\(B_\lambda\) is an invertible line map with strict uniform norm
contraction. Thus (CP.3a) is an unramified regular covering of smooth
total spaces. The same description holds at every uncapped retained
fibre, since the only zero of the Hom section is at \(c\).

\item Choose a holomorphic nonzero vector \(v_O(z)\) in \(N|_O\), possible
on this disc. For an extending translation \(\tau\), the map
\(z\mapsto\tau([v_O(z)])\) is an actual holomorphic section of
\(X'\). Pulling the covering (CP.3a) back to the simply connected
root disc gives a trivial covering; lifting this one section gives
a holomorphic map \(p(z)\in N^\times\). Its projection \(T(z)\)
to the source extends through \(z=0\). It meets the smooth locus
of the reduced central fibre: \([v_O(0)]\) is a regular point of
the map to the disc, a base-preserving automorphism preserves this
property, and the quotient covering is locally biholomorphic. For
the good root this assertion is immediate. For reduced II, III,
IV or \(I_n\) the regular locus of that map is exactly the smooth locus
of the reduced fibre. Compose the lifted point with a suspension
power if necessary to put its central component in the identity
component. The required power exists because the actual Q-action
is transitive on the components. This is a definite deck map in
(CP.3a), not an identification of distinct central components by hand.

\item The resulting section translation \(t_T\) has an invertible line
lift \(L_T\) on the entire root surface. To see the extension,
\(t_T^*N\otimes N^{-1}\) is generically trivial. Proper direct
image gives a generator whose evaluation divisor is vertical. Its
intersection with every component is zero, since both lines have
multidegree zero. The explicit component intersection matrix has
primitive kernel the whole reduced fibre. Dividing by the
corresponding integral power of its base equation therefore leaves
a nowhere-zero section. This is the actual construction of \(L_T\)
and applies equally to the II, III, IV, \(I_2\) and retained \(I_n\) charts.
For a good fibre the same argument has its single reduced component.
The finite quotient/resolution and transported-centre diagrams are
those of F41--F49; the line being used here is their actual N.

\item There is exact commutation with the suspension, including its line
scalar. On a smooth generic source elliptic fibre write the
specified degree-zero line as
\[
 (\mathbf C\times\mathbf C^*)/\Lambda,\qquad
 (z,v)\sim(z+\omega,\chi(\omega)v).
\]
The constant factors of automorphy give a literal lift of a source
translation, \((z,v)\mapsto(z+t,av)\). Every other holomorphic
line lift of that translation differs by a nowhere-zero holomorphic
function on the compact elliptic curve, hence by a constant. Thus
any two such lifts commute. In particular \([L_T,B_\lambda]=1\)
on every generic fibre, and holomorphic extension gives equality
on the entire reduced/root chart. Compare \(L_T(v_O)\) with the
chosen lifted point \(p\). Their ratio is a holomorphic base unit
\(c(z)\). The actual map \(c(z)L_T\) therefore has the same value
as \(\tau\) at the chosen origin of every regular compact torus.
Two translations with the same origin value agree. Equality on the
dense regular locus proves that \(c(z)L_T\) is an extending lift of
\(\tau\) on all of (CP.3a). This constructs the lift without
assuming a covering-lift theorem for an arbitrary singular fibre.

\item It is exactly equivariant under the finite cap action. Pullback of
a base-preserving automorphism to the actual root cover is defined
by \((x,z)\mapsto(\tau x,z)\); it commutes with its deck action.
The two line-cover maps \(A_i L A_i^{-1}\) and \(L\) therefore
differ by \(B_\lambda^{j_i}\), where \(j_i\) is an integer constant
on the connected root disc. Iterating \(m_i\) times gives
\(B_\lambda^{m_i j_i}\). But \(A_i^{m_i}=B_\lambda^{n_i}\),
and the just-proved commutation with \(B_\lambda\) makes this
iterate the identity. Strict contraction on the covering line forces
\(j_i=0\). The lift is consequently exactly equivariant for the
actual affine action; its component and finite line character have
not been omitted.

\item The overlaps of these lifts differ by suspension powers and give
an integral one-cocycle on \(U\) with its finite orbifold points.
Its cohomology is zero. Explicitly the orbifold fundamental group
is the free product of the finite cyclic cap groups, and every
homomorphism from it to \(\mathbf Z\) is zero. Changing the local
lifts by the resulting integral powers of \(B_\lambda\) patches
them to one lift on \(U\). Exact finite-root equivariance descends
its source point to an original source section. The actual
root/quotient/resolution diagrams give its meromorphic line map
through the original additive fibre. On an overlap the scalar
ratio is a base function, since it is constant on the compact
smooth source elliptic fibre.

\item Here is the selected-cusp comparison, with both integer reindexings
written. In its actual Tate and line coordinates the two period maps
have the form
\[
 \begin{split}
 h:(Z,V)&\longmapsto(q_T Z,P_T(q)V),\\
 B_\lambda:(Z,V)&\longmapsto(Q_T(q)Z,
                              q\lambda C(q)V),
 \end{split}                                                \tag{CP.3b}
\]
where \(q_T/q,P_T,Q_T,C\) are holomorphic units. The source
sections and the line divisor sections avoid the selected node;
this gives the asserted unit Tate coordinates. In the height-one
component lattice the period shifts of (CP.3b) are exactly
\((1,0)\) and \((0,1)\). This is the whole-zero I1, index-one
calculation, and is why the component-supported OS49 case is kept
outside this note.

\item An extending translation on the toric covering sends its chosen
component to the component with a definite index \((r,s)\).
It induces the identity on the nearby marked lattice, hence commutes
with the component deck action. Compose its lift with
\(h^{-r}B_\lambda^{-s}\). It now preserves the chosen component
and each of its boundary strata. On the open torus times the base
its map is literally
\[
                  (Z,V,q)\longmapsto(a(q)Z,b(q)V,q).
                                                               \tag{CP.3c}
\]
Evaluate it and its inverse at \((Z,V)=(1,1)\). Since both preserve
the open central torus, \(a\) and \(b\) are holomorphic units
at zero. This proves the bounded multiplier and its inverse, not
only a bound on their valuations. Conversely (CP.3c) extends in
every toric chart. The original lift is therefore
\(h^rB_\lambda^s\) times this unit map.

\item Its source Tate coordinate is \(q_T^r Q_T^s a(q)\). Dividing by
the actual elliptic period \(q_T^r\) leaves the same source point
and a holomorphic unit, so the source section extends through the
smooth Neron group at the selected I1. Its line multiplier is
\(P_T^r(q\lambda C)^s b(q)\); the first factor is a unit and
its exact q-order is s. Thus no line order was lost by the source
period reindexing. The global source section has meromorphic
Weierstrass coordinates on \(\mathbf P^1\), which are rational
functions. The already proved Mordell--Weil group is exactly
\(\mathbf ZG\); hence it is \(kG\) for an integer \(k\). These actual
cusp orders are the input to (CP.6), where \(j=-s\).

\par

\medskip\noindent\textbf{CP.4. The global valuation equation retains the finite kernel}

\par\nobreak

\item Use K35's actual meromorphic line lift \(U_G\) covering \(+G\),
and its actual rational norm ratio
\[
                   \rho=B_0^{-1}/U_G^b.                   \tag{CP.4}
\]
At a cap \(i\) write the literal local line comparison
\(U_G=z_i^{e_i}h_iL_i\), where \(h_i\) is a nonzero unit and
\(L_i\) extends invertibly. The quotient/resolution map defining
\(L_i\) is retained. At an uncapped retained fibre use its
corresponding integer order \(e_i\) and \(m_i=1\). At the
selected \(I_1\) put \(\operatorname{ord}U_G=e_c\) in its
ordinary rational frame. K35 gives the exact rational divisor orders
\[
 \operatorname{ord}_i\rho=-be_i/m_i,
 \qquad\operatorname{ord}_c\rho=-1-be_c.
\]
Their sum is zero, so
\[
                   e_c+\sum_{i\ne c}e_i/m_i=-1/b.          \tag{CP.5}
\]
The lifted global endpoint has the form \(h(t)U_G^k\). On every
root and retained fibre its invertibility gives
\(\operatorname{ord}_i h=-ke_i/m_i\). This includes the full
finite-character condition: the right hand side must be integral.
At the selected cusp apply the actual integer reindexing of CP.3.
Its source-period power has a unit line multiplier. Composing the
endpoint with \(B_\lambda^j\) removes its line order; that map
has source translation \(-jbG\) and line zero of order \(j\).
Therefore the original endpoint has line order \(-j\), with
\(j=-s\) in CP.3, and
\[
                      \operatorname{ord}_c h+ke_c=-j.
                                                               \tag{CP.6}
\]
Taking the degree of the actual rational divisor of \(h\) and
using (CP.5) gives \(j=k/b\). Hence \(k=bj\). All finite-kernel
endpoints have now been tested against the global ordinary, root
and cusp maps; no central kernel was set to zero locally.

\item The functions \(h\) and \(\rho^j\) have exactly the same divisor,
by (CP.4)--(CP.6), and therefore differ by a nonzero constant \(c\).
Thus the actual global line map is
\[
                hU_G^{bj}=cB_0^{-j}.
\]
On the compactified suspension \(B_\lambda=\lambda B_0\) is its
literal deck generator. The displayed map consequently acts by
physical delta multiplication \(c\lambda^j\). Conversely every
nonzero scalar is an extending physical-delta automorphism in the
source, cap and cusp charts. The scalar action is faithful: equality
with a suspension power would make a nonzero multiple of \(bG\)
the zero section, contradicting the torsion-free Mordell--Weil group.
The global extending translation group is therefore precisely
\(\mathbb C^*\).

\par

\medskip\noindent\textbf{CP.5. The actual fixed-scalar loop endpoint}

\par\nobreak

\item The endpoint of CP.2 is the physical scalar just computed, with every
finite central translation tested in CP.4. Choose a logarithm of
that one nonzero number \(c_{\rm end}\). Multiplication by
\(\exp(-t\log c_{\rm end})\), \(0\leq t\leq1\), is an actual
isotopy through global delta automorphisms; compose it with the
interval trivialization. Its endpoint closes literally. It is a
smooth bundle trivialization over the parameter circle, with the
same source, cap and annular identifications and the same transported
toric fan. Thus the fixed-scalar cap-data loop has zero sphere
mapping class in this scope. The scalar path acts in the actual
Cox normal coordinates by the prescribed physical circle weights;
no additional normal framing is inserted.

\item For a varying scalar we retain the full endpoint and compare it
with K35's actual shorter coordinate \(w\), where
\(\lambda=w^{M/b}\). The next step determines the allowed winding
by the literal cap and cusp maps.

\par

\medskip\noindent\textbf{CP.6. Exact scalar-winding obstruction and the shorter endpoint}

\par\nobreak

\item The arrow in this step is fixed to K35.
Let the scalar have signed winding \(\ell\in\mathbb Z\), and use
the canonical universal parameter
\(v=\log\lambda/(2\pi iM)\). Its change is
\(\alpha=\ell/M\). The comparison in CP.2 works with the actual
scalar varying as well: on each parameter disc compare to the
canonical scalar family at that same scalar. Its endpoint is a
translation isomorphism from the canonical model at \(v\) to the one
at \(v+\alpha\). All finite characters are carried with that map.

\item On root \(i\) the actual ordinary-to-cap parameter cochain changes by
\(M\alpha n_i\log z_i\,\delta/(2\pi i)\). Its exponential is
exactly \(z_i^{\ell n_i}\), in the arrow of WS.14--15. Hence a
source lift \(hU_G^k\) extends as an endpoint precisely when its
ordinary root multiplier has the compensating order
\[
             m_i\operatorname{ord}_i h+k e_i=-\ell n_i.
                                                               \tag{CP.7}
\]
This is an equality of actual line orders. Reducing it modulo \(m_i\)
retains the finite commutator character; it is not a rational-rank
condition. The ordinary order-one and framing charts have exactly
this formula with \(m_i=1\) and their recorded signed norm.

\item At the selected I1 the two endpoint source and line charts have the
same literal scalar \(\lambda\). The integral component reindexing of
CP.3 is still required; it gives
\(\operatorname{ord}_c h+ke_c=-j\), \(j\) an integer. Taking the degree
of \(h\), using (CP.5) and the actual sphere norm equation
\(\sum_i n_i/m_i=\varepsilon/M\), gives
\[
                 j=\frac{k}{b}-\frac{\varepsilon\ell}{M}.
                                                               \tag{CP.8}
\]
Consequently
\[
                  k-bj=\frac{\varepsilon b\ell}{M}
                          \quad\hbox{is an integer}.           \tag{CP.9}
\]
Since \(b\) divides \(M\) in the proved catalogue, (CP.9) forces
\((M/b)\mid\ell\). This derives the minimal closed scalar winding
from the actual cap and cusp maps even when the intermediate cap
conjugators have nontrivial finite kernels.

\item Conversely K35 constructs a literal endpoint for \(\ell=M/b\),
with its rational norm root, finite characters, constant delta
normalization and simultaneous Cox product. Iterate that actual
endpoint \(q=\ell/(M/b)\) times, reversing its arrow for negative
\(q\). Comparing an arbitrary cap-data endpoint of this scalar winding
with that particular endpoint now gives a global extending
translation of the same compactified family. CP.3--4 identify that
ratio with physical delta multiplication. The path in CP.5 removes
it as an actual global isotopy. Thus the return of the given
cap-data loop is the \(q\)-th power of K35's specified shorter scalar
return. No finite-kernel contribution remains untested, and no
residue-class division is used.

\item For completeness, the physical return uses the inverse closing
arrow. Let \(I_t:X_t\to Y_t\) be the actual interval comparison to
the canonical scalar family. Write \(E_X:X_0\to X_1\) and
\(E_Y:Y_0\to Y_1\) for their forward endpoint identifications;
\(E_Y\) is the corresponding iterate of K35's explicit endpoint.
The comparison just proved is
\[
             I_1E_XI_0^{-1}=C E_Y,                         \tag{CP.10}
\]
where \(C\) is a physical delta automorphism of \(Y_1\). Choose
parallel transport \(P_X=I_1^{-1}P_YI_0\). The physical return
is \(\Phi_X=E_X^{-1}P_X\), so the literal comparison is
\[
 I_0\Phi_XI_0^{-1}=E_Y^{-1}C^{-1}P_Y
                    \simeq E_Y^{-1}P_Y.                  \tag{CP.11}
\]
The isotopy is the global scalar path of CP.5. The right-hand map
is exactly the \(q\)-th power of the positive shorter scalar
return, with the inverse endpoint and the same parallel-transport
direction. This proves the claimed oriented sphere mapping class
without changing a closing map into a physical return by notation.

\item The whole argument fixes the marked source, root orders, norm
integers, logarithmic sites and fan. Its parameter transitions have
the prescribed linear source/lattice marking and vary the affine
translations, line units and transported annular data of F44/F49.
It makes no assertion about adding arbitrary nontranslation endpoint
automorphisms, moving the source or logarithmic sites, or changing
the component-supported OS49 model.

\item The exact certificate is \path{code/part3/fixed_source_cap_loops.py},
with its JSON and captured output. It checks every root Lie weight,
the actual valuation identities, all thirty nonempty catalogue
branches, and the free-group closing-arrow identity (CP.11).

\item The actual OS45 datum has (b=M=1), (H=8),
(chi=1), and physical reference return (3).
Taking this cap-data loop with (N=1) gives (3) through
(CP.10)--(CP.11), while the old numerical substitution gave
(25). The family-morphism reduction is retained; only that
substitution is refuted.

\end{enumerate}
\end{proof}




\begin{theorem}\label{night:P3K.43}
\StatementRecord{night:P3K.43}{false as stated; correction proved}{S6\_06\_BLIND: GPT-6 Astra (OpenAI, ChatGPT Pro mode, effort Pro), recorded re-derivation in the checkpoint dated 18 November 2026}{\cref{night:P3K.32}, \cref{night:P3K.35}, \cref{night:P3K.41}}
Fix the marked rational elliptic source, actual line, translation
\(Q=bG\), whole-zero selected \(I_1\), effective additive caps,
and transported periodic index-one fan of P3K.41. Allow the ordinary
order-one logarithmic sites to move through the regular-source locus,
away from every fixed singular value and effective cap. Retain their
signed norms, all elliptic and delta period coefficients, finite
characters, annular maps, and prescribed linear source/lattice marking.
Consider an actual closed marked family, with its specified endpoint;
site permutations must carry the complete cap data.

There is an actual local family isomorphism from a moving-site family
to the reference-site family at the same scalar, extending through every
root and cusp chart and carrying the physical circle and normal collar.
These isomorphisms identify every fixed-scalar closed loop, including
closed braids around the fixed singular values, with a smoothly trivial
sphere bundle. More generally, if the scalar has winding \(\ell\),
then
\[
                    \ell=N\frac{M}{b},\qquad N\in\mathbb Z,
\]
and its physical return is the \(N\)-th power of the constructed
shorter scalar return of P3K.35, with parameter
\(\lambda=w^{M/b}\). In particular moving these sites contributes
no additional return class. For the transported fan (P), its oriented class is
(N A_{
m short}(P)), where \cref{night:P3K.73} computes the
actual shorter return from the full period and endpoint maps.
The statement concerns the stated closed marked family scope; it does
not assign a return to source-moduli loops, collisions with source
singular values, arbitrary nontranslation endpoints, or the
component-supported OS49 cusp.


\end{theorem}
\begin{proof}
\begin{enumerate}[label=\arabic*.]
\item \par

\medskip\noindent\textbf{MS.1. The full regular-disc logarithmic map}

\par\nobreak

\item Fix a simply connected regular-source disc \(D\) and a smaller disc
containing the moving site \(p\). The local suspension is the actual
compact complex two-torus family
\[
 \mathcal X_D=N_D^\times/\langle B_\lambda\rangle
       =V_D/\Pi(t)\mathbb Z^4.                         \tag{MS.1}
\]
The first presentation retains the source degree-zero line and its
contracting lift. In the second presentation a holomorphic marking of
the full period lattice is chosen on \(D\), including the source-line
characters. Its period map \(\Pi(t)\) takes values in the actual Lie
bundle \(V\) of P3K.41. It is not a projection to just the gamma
coordinate. On a small scalar-parameter disc the same marking is
holomorphic in that parameter as well.

\item Let \(v\in\mathbb Z^4\) be the complete marked period of the
order-one clutch. With the ordinary-to-cap direction fixed, its map is
\[
 \Phi_p(t)=\operatorname{trans}\left(
       \Pi(t)v\,\frac{\log(t-p)}{2\pi i}\right).        \tag{MS.2}
\]
The gamma coefficient of \(v\) is \(-n\) for the norm \(n\).
This sign is the literal equation
\(h=\gamma-n\arg(t-p)/(2\pi)\) of P3K.32, not a change to its
inverse arrow. The remaining entries of \(v\) are the specified
elliptic and delta coefficients. A branch change of the logarithm adds
\(\Pi(t)v\), an actual full period in (MS.1). Thus (MS.2) is a
single-valued holomorphic translation on the punctured disc, including
its line automorphy. Any prescribed extending unit or translation is
also retained; its independent variation is the actual cap/line
comparison already proved in P3K.41.

\item There is a literal local model on the whole common disc \(D\).
For order one the norm identity is \(A_p=B_\lambda^n\), hence
\[
 N_D^\times/\langle A_p,B_\lambda\rangle
              =N_D^\times/\langle B_\lambda\rangle .    \tag{MS.3}
\]
Inside the original small cap disc use this quotient. On its overlap
with \(D\setminus\{p\}\), use (MS.2) to carry the exterior
coordinate into that same quotient. These maps agree by the original
cap gluing. They identify the entire piece over \(D\) with the
fixed model (MS.1), using the actual source line and suspension, while
its outer annular transition remains exactly \(\Phi_p\). This
constructs the enlarged chart; it does not presume a lift of a
circle-valued base phase or move the source coordinate \(t\).

\item The resulting charts form an actual holomorphic family as \(p\)
varies locally. Their interior models and root/cusp models are fixed,
and their outer annular maps are holomorphic in \(t,p\). Each local model is proper over its source-base chart, and the
models are glued by the displayed fibre isomorphisms. Compactness of
the fixed source base therefore gives properness over the parameter,
after restricting to the chosen parameter disc.

\par

\medskip\noindent\textbf{MS.2. The exact Lie logarithm on a fixed outer annulus}

\par\nobreak

\item Choose the outer annulus once, with
\(r<|t-p_0|<R\), and restrict \(|p-p_0|<r/2\). The ratio of
the two actual maps in (MS.2) has the following honest Lie lift:
\[
 \begin{split}
 L(t,p)&=\frac{\Pi(t)v}{2\pi i}
                   \operatorname{Log}\frac{t-p}{t-p_0},\\
 \operatorname{Log}\frac{t-p}{t-p_0}
       &=-\sum_{j\geq1}\frac{(p-p_0)^j}{j(t-p_0)^j}.
 \end{split}                                             \tag{MS.4}
\]
The series converges uniformly on compact subsets of the fixed annulus
and parameter disc, and vanishes at \(p=p_0\). Its exponential is
the displayed ratio, with no residual period. In particular
\[
       \Phi_p\Phi_{p_0}^{-1}=\operatorname{trans}(L).      \tag{MS.5}
\]
All four entries of \(v\) participate in this identity. For several
distinct moving sites choose disjoint discs for this small parameter
step and use the sum of their Lie cocycles on the corresponding
annuli. The translations commute by the actual degree-zero-line group
law of P3K.41 CP.3, including the suspension scalar. Thus summing the
lifts is the actual composition of the gluing maps.

\item On a change of regular-source marking, \(\Pi' =\Pi C\) and
\(v'=C^{-1}v\) for the same integral monodromy matrix \(C\).
Consequently \(\Pi'v'=\Pi v\) literally. This is how (MS.4)
is transported to the next disc along a site path. No global
single-valued source period basis is assumed, even when the path
winds around a fixed singular fibre. The complete period is carried
by the actual local system; it is never reset after a turn.

\par

\medskip\noindent\textbf{MS.3. The actual global family comparison}

\par\nobreak

\item Use the fixed orbifold base atlas of P3K.41, adding the discs and fixed
outer annuli just constructed. Its Lie bundle has
\(H^1(B,V)=0\), proved there from the exact coarse extension of
\(\mathcal O(-1)\) by \(\mathcal O\), with all root weights
retained. Equation (MS.4) is a \(V\)-valued total Cech one-cocycle
on this atlas and groupoid. It is supported on the regular annular
overlaps and is zero on unchanged overlaps; on a refinement the
indicated bundle transition maps are included. Unlike a merely
torus-valued logarithm, it has no new integral two-cocycle. The
pre-existing finite-root group cocycle remains the one of P3K.41.

\item The Laurent/Cauchy solution from CP.2 gives holomorphic sections
\(L_0,L_i\), analytic in \(p\), satisfying
\[
                    L_0-L_i=L=\Psi_p-\Psi_{p_0},          \tag{MS.6}
\]
where \(\Psi\) denotes the Lie expression in (MS.2) on the
overlap. The solution is zero at \(p=p_0\). At root charts this
is the full groupoid solution; finite averaging is applied to the
complex vector sheaf, not to the integral lattice. Thus no invariant
logarithm of a nonzero fixed torsion point is being assumed.

\item Set \(T_0=\operatorname{trans}(L_0)\) on the ordinary chart and
\(T_i=\operatorname{trans}(L_i)\) on the cap chart. With the
ordinary-to-cap direction of (MS.2), equation (MS.6) gives
\[
                  \Phi_{p_0}T_0=T_i\Phi_p.                \tag{MS.7}
\]
Hence these maps glue to an actual isomorphism from the moving-site
family to the reference-site family. Formula (MS.7) fixes its direction;
it is not obtained by reversing a closing map informally.

\item Every local map is an extending translation of the sheaf used in
P3K.41. On finite roots it is the actual lifted source and line
translation commuting with \(A_i,B_\lambda\), so the finite
character and norm remain unchanged. At the selected cusp its Lie
section exponentiates, after the same integral period reindexing, to
the actual bounded toric unit map
\[
                 (Z,V,q)\longmapsto(a(q)Z,b(q)V,q),       \tag{MS.8}
\]
with both units and their inverses holomorphic. This is CP.3's
toric extension on the fixed fan. It transports the actual physical
delta action, fixed-curve tube, and Cox normal coordinates. Thus the
comparison includes the line and normal maps on the cusp collar;
there is no separate uncomputed framing change.

\item The same construction applies when affine cap data, line units, and
annular units vary at the same time. First use the actual local
conjugacies of P3K.41, retaining their finite kernel cocycle, then
apply (MS.4)--(MS.7). All maps are analytic on a small parameter disc
and keep the marked source itself fixed.

\par

\medskip\noindent\textbf{MS.4. A closed braid and its exact endpoint}

\par\nobreak

\item Cover the compact parameter interval of a site motion by finitely many
of the steps above. At each step keep its chosen outer annuli fixed.
Compose the actual isomorphisms (MS.7), always carrying the current
period marking and every annular map. This gives an interval
comparison to the reference-site family. It applies to any actual
closed braid in this marked scope. If sites are permuted, their full
cap data and the specified endpoint are carried by that permutation;
no additional endpoint identification is inserted.

\item For fixed scalar, both endpoint models are the same fixed compactified
suspension. Their discrepancy is a global extending translation,
because every map in the construction is such a translation and
has identity linear lattice part. P3K.41 CP.3--4 computes this group
as exactly the physical \(\mathbb C^*\) action, using the root
line orders and the selected-cusp valuation. That computation includes
all finite kernels of \(1-D_i\). In particular a braid around
a source singular value has not produced an untested finite endpoint.

\item Write the discrepancy as multiplication by \(c\ne0\). The
path \(c_t=\exp(t\log c)\) is a global physical delta isotopy
on all source, root, and toric charts. Composing the interval
comparison with its inverse closes it over the parameter circle.
This is a smooth isomorphism of the actual closed bundle with the
fixed bundle. Its physical return is therefore zero. The assertion
uses the actual family and endpoint; it does not infer loop
triviality merely from isomorphism of individual fibres.

\par

\medskip\noindent\textbf{MS.5. Varying scalar and the physical return arrow}

\par\nobreak

\item When the scalar varies, make the comparison at the same scalar on
each parameter disc. The marked period map \(\Pi(t,\lambda)\)
and (MS.4) are holomorphic there. CP.2's simultaneous Laurent splitting
therefore supplies the same actual family isomorphisms. Continuing
them yields a comparison to the fixed-site scalar/cap-data family of
P3K.41. Its exact root/cusp valuation gives
\[
                     \ell=N\frac{M}{b},\qquad N\in\mathbb Z,
                                                               \tag{MS.9}
\]
for the scalar winding of any closed marked family in this scope.
The reference endpoint is the \(N\)-fold iterate of P3K.35's
constructed shorter endpoint, for \(\lambda=w^{M/b}\).

\item For clarity retain the full arrows. Let \(I_t:X_t\to Y_t\)
be the moving-site to fixed-site comparison, let
\(E_X:X_0\to X_1\), \(E_Y:Y_0\to Y_1\) be the actual
forward endpoints, and let \(P_Y:Y_0\to Y_1\) be parallel
transport. Their discrepancy is a physical scalar \(C\) on
\(Y_1\), so
\[
 I_1E_XI_0^{-1}=C E_Y,\qquad
 P_X=I_1^{-1}P_YI_0.
\]
The physical return is the inverse endpoint followed by parallel
transport. Consequently
\[
 I_0(E_X^{-1}P_X)I_0^{-1}
       =E_Y^{-1}C^{-1}P_Y
       \simeq E_Y^{-1}P_Y.                                \tag{MS.10}
\]
The isotopy is the actual global delta path in MS.4. Thus the return
is precisely the \(N\)-th power of P3K.35's shorter scalar return,
with no additional contribution from the moving ordinary sites.
For the fixed transported fan (P), the actual class is
(N A_{
m short}(P)), where the endpoint just compared is the
input to \cref{night:P3K.73}. In the allowed OS45 constant-site
case (N=1,P=T_0), the result is (3), whereas the rejected
polynomial gave (25). This is the located counterexample.

\item This proves the stated extension for regular order-one sites, including
their closed braids around the fixed singular values. It retains the
fixed marked source and effective caps. Source-moduli loops,
collisions with source singular values, arbitrary nontranslation
endpoints, and the component-supported OS49 cusp are not assigned a
return by this argument.

\end{enumerate}
\end{proof}




\begin{theorem}\label{night:P3K.44}
\StatementRecord{night:P3K.44}{false as stated; correction proved}{S6\_06\_BLIND: GPT-6 Astra (OpenAI, ChatGPT Pro mode, effort Pro), recorded re-derivation in the checkpoint dated 18 November 2026}{\cref{night:P3F.26}, \cref{night:P3F.29}, \cref{night:P3F.44}, \cref{night:P3K.37}, \cref{night:P3K.41}}
Fix the actual OS49 source and component line
\(M=\mathcal O(G-O)\), suspension by \(-3G\), selected
\(I_2\) cusp, good IV-star cap of order \(m=3k\), \(k>0\),
norm \(n=\pm1\), finite delta character \(d\), and an actual
periodic index-one fan. Retain the prescribed linear source/lattice
marking and every annular map. For the actual affine source and line
variations on these fixed charts, the global group of extending fibre
translations is precisely the physical \(\mathbb C^*\) action.
Every closed fixed-scalar cap-data family in this marked scope has
trivial sphere return.

More generally, the scalar winding of a closed family satisfies
\[
                          \ell=Nk,\qquad N\in\mathbb Z.
\]
Its physical return is the \(N\)-th power of the actual shorter
OS49 return of P3K.37, whose parameter is \(\lambda=w^k\).
For the two fixed component-fan objects its classes are
(N(-k)) on the negative fan and (0) on the positive fan,
by \cref{night:P3K.65}. For a general transported fan (P)
the exact expression is (N A_{
m short}(P)), with
(A_{
m short}(P)) computed by \cref{night:P3K.77}.
The statement keeps the component divisor, both mixed cusp period
orders, the finite root kernel and character, and the inverse endpoint
in the physical monodromy. It concerns the fixed source and marked
cap/line/annular data above; no added ordinary logarithmic clutches,
source-moduli loops, moving cap sites, or arbitrary nontranslation
endpoints are assigned a return here.


\end{theorem}
\begin{proof}
\begin{enumerate}[label=\arabic*.]
\item \par

\medskip\noindent\textbf{OC.1. The actual local groups and the Lie extension}

\par\nobreak

\item Away from the selected I2 the source line is the actual degree-zero
line on every reduced fibre, and the suspension lift is invertible.
On the good IV-star root use the corrected root line of F26/F29 and
K37, with its original finite character. Its central Picard point
\(G'\) has order three, and the effective deck has order three.
The suspension point \(Q'=3G'\) is zero on that central fibre.
For any actual affine source root \(a=t_S D\) the exact equation is
\[
       a^m=t_{-nQ'},\qquad
       \sum_{j=0}^{m-1}D^jS=-nQ'.                         \tag{OC.1}
\]
The branch \(R\) with \(mR=Q'\), \(R(0)=O\), is
well defined and deck invariant. Put \(Z=S+nR\). On the central
elliptic curve \(1-D_0\) is a surjective isogeny. Choose a local
inverse image of \(Z_0\), extend it to a section, and subtract
its \(1-D\) image. Write the chosen extending central inverse-image section as
\(\widetilde T\), and put
\(Z_1=Z-(1-D)\widetilde T\). This remainder has central value zero
and cyclic norm zero. In that zero-central section group the analytic
inverse of multiplication by \(m\) gives the literal solution
\[
 T=\widetilde T-[m]^{-1}\sum_{j=0}^{m-1}jD^jZ_1,
                 \qquad (1-D)T=Z.                        \tag{OC.2}
\] Substitution
gives \(t_T(t_{-nR}D)t_T^{-1}=a\). This is the good-root
construction of F44, now on this actual OS49 root. Its central kernel
of order three is retained; no inverse image is silently identified
with another.

\item There is an extending line lift of \(t_T\). On the smooth
generic elliptic fibre the Hom line is trivial by its constant
degree-zero factors of automorphy. Proper direct-image evaluation
extends it over the good root; any vertical divisor is a multiple of
the reduced fibre, and division by the corresponding base power gives
an invertible line map. At the retained type-II fibre the same
argument uses its one reduced component. On a generic flat-line
model every translation lift is \((z,V)\mapsto(z+t,cV)\),
because the ratio of any two such lifts is a holomorphic function on
a compact elliptic curve and is constant. Thus these lifts commute
with each other and with the actual suspension. Equality extends
over the root and retained fibre. The remaining line-unit norm is
the exact cyclic product; its central \(m\)-th root of unity is
retained and its zero-central logarithm is removed by the finite sum
in F44. This supplies analytic local conjugacies, including line and
annular maps, on a small parameter disc.

\item Let \(A\) be the sheaf of extending fibre translations of the
actual compactified OS49 suspension and \(V\) its Lie sheaf.
On the orbifold base with the order-\(m\) root there is the exact
vector-bundle sequence
\[
 0\longrightarrow\mathcal O\delta\longrightarrow V
          \longrightarrow L\longrightarrow0.             \tag{OC.3}
\]
At the good IV-star root the actual Weierstrass substitution is
\(t=z^{3k}\), \(x=z^{4k}X\), \(y=z^{6k}Y\). Hence the
original relative dualizing differential is
\(z^{-2k}dX/(2Y)\), and the original source translation vector
is \(z^{2k}\) times the root vector. The invariant holomorphic
Lie sections are exactly these original source sections, since
\(0\le2k<3k\). At the retained II and selected I2 the
logarithmic source vector is the original dualizing vector. At I2
this assertion follows directly from the unit torus maps in OC.2
below, with source coordinate \(Z\) and line coordinate \(V\).
Thus the coarse extension of \(L\) is the source
\(\mathcal O(-1)\), and that of the physical delta line is
\(\mathcal O\). Finite averaging is exact for this complex
vector sheaf, so
\[
                         H^1(B,V)=0.                      \tag{OC.4}
\]

\item As in the actual atlas construction of K41, the exponential sequence
\(0\to H\to V\to A\to0\) is exact on the full root
atlas and groupoid. This does not assert an invariant logarithm for
every fixed torsion point of the coarse root disc. The finite-group
lattice discrepancies remain part of the total Cech two-cocycle.
They are locally constant under parameter variation. Subtracting the
reference parameter gives a Lie one-cocycle, and (OC.4), solved by
the analytic Laurent/Cauchy splitting and finite vector averaging,
gives an actual interval comparison. Its endpoint is a global
extending translation. We compute that group using this cusp's
actual period matrix.

\par

\medskip\noindent\textbf{OC.2. The mixed cusp, with both integer reindexings}

\par\nobreak

\item Use K37's actual Tate line coordinates. The source period \(h\),
suspension \(B_\lambda\), and rational forward lift \(U\)
of \(+G\) have the form
\[
 \begin{aligned}
 h:(Z,V)&\longmapsto(q^2H(q)Z,qP(q)V),\\
 B_\lambda:(Z,V)&\longmapsto(q^{-3}J(q)Z,
                                      q^{-1}\lambda C(q)V),\\
 U:(Z,V)&\longmapsto(qg(q)Z,q^{e_c}u(q)V),
 \end{aligned}                                             \tag{OC.5}
\]
where all displayed capital/lowercase coefficient functions are
holomorphic units. K37 chooses \(e_c=0\); we retain an arbitrary
integral whole-fibre change of that rational lift. Its exact graph is
\(h_j=\lceil j/2\rceil\), and the two \(+G\) defects are
\(-1,0\), so \(U\)'s actual pole divisor is \(-C_0\)
in the chosen normalization. The maps and source/line characters
in (OC.5) are those graph maps; no whole-fibre replacement of this
component divisor is made.

\item The valuation vectors of the first two maps are
\[
                   h=(2,1),\qquad B=(-3,-1),\qquad
             \det\begin{pmatrix}2&-3\\1&-1\end{pmatrix}=1.
                                                               \tag{OC.6}
\]
They generate the entire height-one component lattice of the
periodic toric covering. Consider an extending translation with
identity linear lattice part. It sends a central component and its
open torus to another such component and open torus: these are
components of the central fibre minus its intersection strata, which
an automorphism of the fibration preserves. Choose that image
component in the toric cover and reindex by its actual periods. On
the open torus the reindexed map is
\[
                    (Z,V,q)\longmapsto(a(q)Z,b(q)V,q).     \tag{OC.7}
\]
Evaluate the map and its inverse on the extending section
\((Z,V)=(1,1)\) of this chart. Both images remain in the open
central torus, so \(a,b\) and their inverses extend as
holomorphic units. Conversely (OC.7) extends on every chart of
the same fan. This proves the bounded map, not just an integral
valuation test.

\item In particular a meromorphic translation with valuation vector
\((j,v)\in\mathbb Z^2\) is reindexed by precisely
\[
       (j,v)=(3v-j)h+(2v-j)B.                             \tag{OC.8}
\]
After this period map its remaining coefficients are the units in
(OC.7). This identity is the mixed-cusp replacement for K41's
whole-zero calculation. Since the toric fan is periodic under the
actual lattice (OC.6), it holds for every transported index-one fan,
including the non-semifree choices. Its unit maps carry all Cox
characters and the physical delta action.

\item On the source I2, remove source-period powers from the first
coordinate in (OC.5). A source coordinate \(q^j a(q)\) is
then \(q^\epsilon a'(q)\), \(\epsilon=0\) or \(1\),
by division by the literal \(q^2H(q)\) period. These are the
smooth points on the two actual Neron components, respectively.
This gives the source extension of an endpoint through I2, with
its component retained. Its line multiplier still records the
second coordinate in (OC.8); the line order of \(h\) is one
and cannot be omitted.

\par

\medskip\noindent\textbf{OC.3. The global rational endpoint and its exact divisor}

\par\nobreak

\item On the complement of the selected I2 the local line quotient is an
unramified regular covering by the actual line complement. To lift
a global extending translation, choose a nonzero vector along the
source zero section on each simply connected root or retained disc
and lift its image through that covering. The source point is a
holomorphic section meeting the regular locus of the reduced or
good fibre. Its line translation extends by the Hom evaluation
argument in OC.1. Matching its value at the chosen zero vector
makes its quotient exactly the given translation on every generic
compact torus, and hence on the full chart.

\item The lift commutes with \(B_\lambda\) by OC.1. Its conjugate
by the finite cap root differs by a deck power \(B_\lambda^a\).
Iterating the order \(m\) and using \(A^m=B^n\) gives
\(B^{ma}=1\); strict contraction forces \(a=0\). Thus the
lift retains the actual finite character. Overlap lifts differ by
integer suspension powers. The base complement is the affine line
with one finite orbifold point; its orbifold fundamental group is
cyclic of order \(m\), so it has no homomorphism to \(\mathbb Z\).
These powers can be patched to one lift. OC.2 gives its source
extension at the missing I2. The resulting original source section
is rational, and the proved Mordell--Weil group \(\mathbb ZG\)
makes it \(jG\) for an integer \(j\). The lifted line map is
therefore
\[
                            f(t)U^j                       \tag{OC.9}
\]
for a rational base function \(f\): its ratio with \(U^j\)
is constant on each compact generic elliptic fibre, and the root
and selected-cusp maps prove its meromorphic extension.

\item Retain K37's literal rational norm ratio
\[
                   \rho=B_0^{-1}/U^3.                    \tag{OC.10}
\]
At the I2 cusp the graph defects of \(U^3\) are
\(-2C_0-C_1\), whereas those of \(B_0^{-1}\) are
\(-C_0\). Thus their ratio has the whole-fibre divisor
\(+C_0+C_1\), and
\[
                  \operatorname{ord}_c\rho=1-3e_c.       \tag{OC.11}
\]
At every other site \(i\), write the actual line comparison
\(U=z_i^{e_i}u_i L_i\), with \(L_i\) invertible and
\(m_i=m\) at the good root, \(m_i=1\) elsewhere. The
source/line norm diagram gives
\[
 \operatorname{ord}_i\rho=-3e_i/m_i,\qquad
             e_c+\sum_{i\ne c}e_i/m_i=\frac13.           \tag{OC.12}
\]
The second equality is the degree of the actual rational divisor
of \(\rho\). At the good root its finite character is still
\(DLD^{-1}=\zeta_3L\), so \(e_i\equiv m/3\pmod m\);
the order in (OC.12) is integral. This retains the finite character,
not just its rational average.

\item For a fixed scalar, invertibility of (OC.9) away from the cusp
gives \(\operatorname{ord}_i f=-je_i/m_i\), including the
integrality at the root. Write
\(v=\operatorname{ord}_c f+je_c\) for its actual cusp line
valuation. Degree zero of \(\operatorname{div}f\) and
(OC.12) give
\[
                              v=j/3.                     \tag{OC.13}
\]
Both \(j\) and \(v\) are integers, so \(j=3a\).
Equation (OC.8) now has \(3v-j=0\), \(2v-j=-a\);
the exact cusp reindexing is \(B^{-a}\). Moreover \(f\)
and \(\rho^a\) have the same divisor at every site by
(OC.11)--(OC.13), so \(f=c_0\rho^a\) for a nonzero
constant. Therefore the global endpoint is literally
\[
                       fU^{3a}=c_0B_0^{-a}.              \tag{OC.14}
\]
On the quotient by \(B_\lambda=\lambda B_0\) this is
physical multiplication by \(c_0\lambda^a\). Every physical
nonzero scalar extends on every actual root and toric chart.
It acts faithfully because a nonzero suspension power has nonzero
source translation in the torsion-free Mordell--Weil group. Thus
the global extending translation group is exactly \(\mathbb C^*\).
This proof has tested the central order-three kernel and every finite
line character against the actual component cusp.

\par

\medskip\noindent\textbf{OC.4. Scalar winding, its root square, and the shorter endpoint}

\par\nobreak

\item For the varying scalar use the universal parameter
\(v_\lambda=\log\lambda/(2\pi i m)\). A scalar winding
\(\ell\) changes it by \(\ell/m\). The actual OS49
ordinary-to-cap map is K37's plus-log map, with period
\(n(e_\gamma-e_u)+d\delta\). The source-period and fixed
delta terms have no scalar change. Consequently its exact change
on the root is
\[
                  \ell n\frac{\log z}{2\pi i}\delta.    \tag{OC.15}
\]
Let \(E\) be a forward endpoint from the initial model to the
continued final model, with ordinary expression \(fU^j\).
The lifting proof in OC.3 applies with the exact relation
\(A_{\rm final}E=EA_{\rm initial}\) in place of commutation:
a possible difference \(B^a\) has \(m\)-fold iterate
\(B^{ma}\), while both root norms are \(B^n\). It is
therefore zero by contraction. The two canonical root source maps
are identical; only their scalar factors differ. Their source
section and rational-base ratio are consequently those of OC.9.
The actual square is
\[
          \Phi_{\rm final}E=E_{\rm cap}\Phi_{\rm initial}.
                                                               \tag{OC.16}
\]
Thus the cap expression contains the extra factor \(z^{\ell n}\).
Its extending unit condition is exactly
\[
        m\operatorname{ord}_{i}f+je_i=-\ell n.            \tag{OC.17}
\]
At all other noncusp sites the condition is the fixed-scalar one.
The source and line characters in this square are the same as K37
OS.10--11, with the arrow reversed there because OS.9 is a closing
map. Taking degrees as before now gives
\[
             v=\frac j3+\frac{\ell n}{m},\qquad
                      3v-j=\frac{\ell n}{k}.             \tag{OC.18}
\]
The left side is an integer and \(n=\pm1\), so a closed
marked family necessarily has
\[
                              \ell=Nk,\quad N\in\mathbb Z.
                                                               \tag{OC.19}
\]
This is an obstruction from the actual mixed-cusp lattice and the
root square, not divisibility of an already computed sphere class.

\item P3K.37 supplies the actual endpoint for every such winding. For one
positive shorter turn its closing map is
\(\widetilde C_r:X_{\zeta_3r}\to X_r\) with source
translation \(-2nG\). Its inverse is the forward endpoint with
\(j=2n\). Equation (OC.18) then gives \(v=n\), and its
valuation vector is
\[
                           (2n,n)=n h.                  \tag{OC.20}
\]
This is exactly the inverse of K37's bounded closing reindexing.
At the root, \(j\equiv-n\pmod3\) gives the full character
equation with \(DLD^{-1}=\zeta_3L\) and
\(A_{\zeta_3r}=\zeta_3^n A_r\); the fixed \(d\)
commutes and is retained. K37's rational cube root, actual norm,
constant delta adjustment, and common Cox origin therefore give a
literal closed endpoint. Iterating it \(N\) times, or using its
inverse for negative \(N\), realizes (OC.19) with the same fan.

\par

\medskip\noindent\textbf{OC.5. The actual interval comparison and physical return}

\par\nobreak

\item OC.1's parameter-dependent source and line maps and the full
groupoid Lie solution give an actual interval comparison to that
canonical scalar family. The cusp maps are precisely OC.2's bounded
units after its integer period reindexings. They preserve the fan
and transport its Cox normal maps, so no extra framing is inserted.
The ratio of any given endpoint of winding (OC.19) to K37's
constructed endpoint is a global extending translation of the same
family. OC.3 identifies it with a physical scalar, including every
finite kernel. Its chosen logarithm gives a global scalar isotopy.

\item Write the actual interval comparison as \(I_t:X_t\to Y_t\),
the forward endpoints as \(E_X,E_Y\), and the reference parallel
transport as \(P_Y\). Then for a physical scalar \(C\),
\[
 I_1E_XI_0^{-1}=C E_Y,\qquad
 I_0(E_X^{-1}I_1^{-1}P_YI_0)I_0^{-1}
                  =E_Y^{-1}C^{-1}P_Y.                   \tag{OC.21}
\]
Removing \(C\) by the actual scalar isotopy leaves the inverse
reference endpoint followed by parallel transport. This is exactly
the physical return of K37, iterated \(N\) times. In particular fixed scalar gives the trivial sphere bundle.
For the negative and positive component fans, respectively,
\cref{night:P3K.65} gives (N(-k)) and (0); for a transported
fan (P) the same actual arrow gives (N A_{
m short}(P)),
using \cref{night:P3K.77}. The case (N=k=1) on the negative
fan has class (27), contradicting the old norm polynomial's zero.
The scalar reduction has been proved by the literal forward
endpoint square and inverse-endpoint physical return above.

\end{enumerate}
\end{proof}




\begin{theorem}\label{night:P3K.45}
\StatementRecord{night:P3K.45}{false as stated; correction proved}{S6\_06\_BLIND: GPT-6 Astra (OpenAI, ChatGPT Pro mode, effort Pro), recorded re-derivation in the checkpoint dated 18 November 2026}{\cref{night:P3F.54}, \cref{night:P3F.56}, \cref{night:P3K.40}}Retain the selected-cusp source family and positive \(\rho\)-loop
of P3F.54/P3F.56, with fixed scalar, the actual III-star cap, and no
logarithmic modification on the coalescing finite pair. For every
allowed cusp fan \(T\), put \(J=\operatorname{diag}(-1,1)\)
on its marked two-dimensional lattice. The actual returns for
\(T\) and \(J(T)\) are conjugate by an orientation-preserving
diffeomorphism, through the literal half-source family map.
For the initial fan of P3F.56, \(J(T)\) is its nonfixed
diagonal-edge flop, and its positive return has oriented class
\[
                 [\Sigma_{\Phi,J(T)}]=-1\in\mathbb Z/28.
\]
This gives neither a once-u endpoint on one fixed fan nor the
individual characteristic pair or integer derivative of the
reflected fan. The different fixed-edge flop has class21 by
P3K.40.
\end{theorem}
\begin{proof}
\begin{enumerate}[label=\arabic*.]
\item \par

\medskip\noindent\textbf{RF.1. The literal source, line and toric endpoint}

\par\nobreak

\item Retain the exact physical source equation and ordered node maps of
P3F.54. A positive half rho turn is one positive turn of the selected
u puncture. The square root in
\[
 \xi=\frac{y-X\sqrt{X+a(u)}}{\sqrt{X+u}},\qquad
 \eta=\frac{y+X\sqrt{X+a(u)}}{\sqrt{X+u}}
\]
changes sign, while the denominator root remains on its extending
branch. Thus the actual nodal endpoint exchanges \(\xi,\eta\).
It has the Tate source arrow \(Z\mapsto Z^{-1}\), with \(q_T\)
unchanged. The marked Picard point changes from \(P\) to
\(P^{-1}\).

\item We retain the line map. For the canonical theta ratio
\(f_P(Z)=\Theta(Z/P)/\Theta(Z)\), the exact identity in F54 is
\[
                   f_{P^{-1}}(Z^{-1})=P f_P(Z).
                                                               \tag{RF.1}
\]
The divisor-compatible arrow is consequently
\((Z,V)\mapsto(Z^{-1},PV)\). Reindex once by the new source
period \((Z',V')\mapsto(q_TZ',P^{-1}V')\). The two line factors
cancel, giving the actual annular endpoint
\[
                   (Z,V,q_T)\longmapsto(q_T/Z,V,q_T).
                                                               \tag{RF.2}
\]
Its square is literally the identity. The source translation and
its Hom map are intertwined by this same arrow: their source
multiplier changes to its inverse, while their scalar line factor
is retained. These are precisely the source and line transitions
of F54, not a replacement of the actual parameter family.

\item The III-star root charts, division section, unit root and finite
character extend over the full small u-disc. Their ordinary-to-cap
annuli continue with those coefficients. The half transport retains
these maps; its second iterate returns their literal charts. The
ordinary exterior is carried by the corresponding source and line
isomorphisms. Thus (RF.2) is compatible with every part of the actual
family, including its finite cap.

\item Here is the chart-ring map at the cusp. Write \(q_T=\kappa q\),
where \(q\) is the physical height character and \(\kappa\)
is the retained nonzero unit on the parameter annulus. For a target
character \(Z^mV^nq^\ell\), pullback by (RF.2) is
\[
          Z^mV^nq^\ell\longmapsto
             \kappa^m Z^{-m}V^nq^{\ell+m}.             \tag{RF.3}
\]
The integral character map is
\((m,n,\ell)\mapsto(-m,n,\ell+m)\). Its dual cone-height
map is
\[
                 (v_Z,v_V,h)\mapsto(-v_Z+h,v_V,h).
                                                               \tag{RF.4}
\]
On the height-one plane this is the affine reflection
\(p\mapsto Jp+e_1\). Periodicity identifies its translated fan
with \(J(T)\). Equation (RF.3) maps each target dual-cone
semigroup into the source semigroup with equality after the inverse
map. The extra factor \(\kappa^m\) is an actual holomorphic unit
on the full parameter annulus. Hence the monomial map and its
inverse extend over every actual unimodular chart. The period
actions are intertwined by the already displayed source and line
maps, so these chart maps descend to the compact models.

\item The physical delta action is multiplication in \(V\). It is
preserved literally by (RF.2), including its orientation. The
holomorphic chart maps preserve the complex orientation of the
six-dimensional fibres. No source-circle inversion or new line
character is substituted.

\par

\medskip\noindent\textbf{RF.2. The actual two-component family and its return arrows}

\par\nobreak

\item On the rho cover consider the two actual families with fans
\(T\) and \(J(T)\). The deck change \(\rho\mapsto-\rho\)
and (RF.2)--(RF.4) interchange them. Applying that arrow twice is
the identity on the source, line and toric charts, and the retained
finite-cap maps also return. Their disjoint union therefore descends
to a proper smooth family over the punctured u-disc whose fibre has
two connected components, labelled \(T,J(T)\) at the initial
point. This construction does not assert a once-u endpoint on either
one fixed fan.

\item Choose a smooth horizontal field on the disjoint cover family,
tangent to its prescribed collars, and average it under the actual
involution that exchanges the two families while taking rho to minus
rho. It descends to a horizontal field on the displayed u-family.
The physical delta action can be retained by an additional compact
circle average. Its transport along one positive u-loop has arrows
\[
           \alpha:X_T\longrightarrow X_{J(T)},\qquad
           \beta:X_{J(T)}\longrightarrow X_T.
                                                               \tag{RF.5}
\]
The domains and targets in (RF.5) are the actual components of that
same family. Restricting its second iterate to either component is
the physical return along one positive rho-loop in the corresponding
pulled-back family. Thus, with compatible representatives,
\[
                \Phi_T=\beta\alpha,\qquad
                \Phi_{J(T)}=\alpha\beta
                       =\alpha\Phi_T\alpha^{-1}.         \tag{RF.6}
\]
Any change of the horizontal field gives the usual isotopy through
its interpolated horizontal fields, with these same endpoint maps.
Equation (RF.6) therefore proves conjugacy of the actual positive
return classes. Since \(\alpha\) preserves the fibre orientation,
the oriented class in \(\Theta_7\) is the same for the two fans.
Explicitly, the already fixed identification of the oriented sphere
mapping-class group with \(\Theta_7=\mathbb Z/28\) is a group
homomorphism, and conjugation acts trivially in this abelian group.

\par

\medskip\noindent\textbf{RF.3. Identification of the particular reflected diagonal}

\par\nobreak

\item Use the exact two triangles of the initial physical-height atlas:
\[
       (0,e_1,e_2),\qquad(e_1,e_2,e_1+e_2),
                                                               \tag{RF.7}
\]
and all their lattice translates. Its three edge directions are
\(e_1,e_2,e_1-e_2\); the physical delta edge is \(e_2\).
The reflection \(J\) retains the two axis edge directions and
changes the third to \(e_1+e_2\), up to sign. It therefore replaces
exactly the diagonal \([e_1,e_2]\) of the unit parallelogram by
\([0,e_1+e_2]\). This is the periodic nonfixed diagonal-edge
flop, denoted \(C_2\) in the physical-height six-chart atlas.

\item The physical transfer of \cref{night:P3K.68} gives the original positive rho return class $-1$.
Equations (RF.3)--(RF.7) consequently give the actual reflected-fan
return
\[
                    [\Sigma_{\Phi,J(T)}]=-1
                             \quad\text{in }\mathbb Z/28.
                                                               \tag{RF.8}
\]
The individual derivative integer on \(C_2\) is still not computed
by this conjugacy. The separate fixed-edge flop of P3K.40 has class
21; its edge is \(C_0\), not the edge in (RF.7).

\item The certificate \path{code/part3/source_cusp_reflection.py} checks the
character/cone involution, physical delta retention, exact theta
line-factor cancellation, and the two labelled triangle sets.

\end{enumerate}
\end{proof}




\begin{theorem}\label{night:P3K.46}
\StatementRecord{night:P3K.46}{false as stated; correction proved}{S6\_06\_BLIND: GPT-6 Astra (OpenAI, ChatGPT Pro mode, effort Pro), recorded re-derivation in the checkpoint dated 18 November 2026}{\cref{night:P3F.54}, \cref{night:P3K.23}, \cref{night:P3K.40}, \cref{night:P3K.42}, \cref{night:P3K.45}}Retain the actual selected-cusp positive \(\rho\)-loop of
P3F.54/P3F.56 and P3K.45: the scalar is fixed, the III-star cap
is unchanged, and the coalescing finite pair has no logarithmic
modification. Let \(T'=J(T)\) be the reflected initial fan, obtained
by flopping its nonfixed diagonal edge \(C_2\). Its meridian-one dual
characteristic pair and relative coordinates at its retained fixed
edge \(C_0\), in the specified Cox product, are
\[
 (k,\ell)=(6,0),\qquad
 (m_{C_0},n_{C_0},s)=(0,-6,-1).
\]
The physical positive-parameter completion has pair \((-6,0)\).
The actual cap map below transports its meridian line, and the common
parameter reversal gives the displayed Jin pair. This statement
does not assign the separate relative integers at \(C_2\) or an
arbitrary-fan return formula.
\end{theorem}
\begin{proof}
\begin{enumerate}[label=\arabic*.]
\item We retain the complete source, line, cap and frame maps below.
The meridian-one graph line is (U=M^{-1}): the fibre map
(D_M(v)(w)=\langle v,w\rangle/\langle v,v\rangle), with the
Hermitian form linear in its second argument, satisfies
(D_M(e^{i\theta}v)=e^{-i\theta}D_M(v)).
Its (SU(2))-equivariant spinor extension is
((a,b)\mapsto(\overline b,-\overline a)).
The first-to-second cap maps in \cref{night:P3K.68} give
(\lambda_M(F)=-\Sigma_{F^{-1}}) and
(\lambda_U(F)=\Sigma_{F^{-1}}); cancelling the opposite
product collars gives (A_M=-A_U). We therefore distinguish the
intermediate meridian-one number (A_U=(4k_U-\ell_U)/24)
from the physical return (A_M) throughout this calculation.

\item \par

\medskip\noindent\textbf{RC.1. The fixed-edge chart map after the actual source period}

\par\nobreak

\item We use the physical height coordinate \(q\), as in P3K.42. Write
\(q_T=\kappa(q,\rho)q\), retaining the actual unit \(\kappa\) and its
rho winding. In the first chart of the initial fixed edge, put
\[
 Z=z,\qquad V=c_1,\qquad q=z c_0c_1.                 \tag{RC.1}
\]
Its second chart has coordinate \(z^{-1}\) and normal coordinates
\(z c_0,z c_1\), after the actual source-period reindex and its
retained units. These are precisely the two fixed-edge products used
in P3K.40 and P3K.42.

\item The reflected triangles are
\((0,e_1,e_1+e_2)\) and \((0,e_2,e_1+e_2)\). Their affine toric
coordinates are, respectively,
\[
 \left(\frac qZ,\frac ZV,V\right),\qquad
 \left(\frac qV,\frac VZ,Z\right).                  \tag{RC.2}
\]
Thus the reflected fixed edge at \([0,e_2]\) has first chart
\[
 z'=Z',\qquad d_0=\frac q{V'},\qquad
 d_1=\frac{V'}{Z'},\qquad q=z'd_0d_1.              \tag{RC.3}
\]
The adjacent translate of the first triangle has tangent coordinate
\((z')^{-1}\) and normal coordinates \(z'd_0,z'd_1\).
In particular its actual normal transition is again multiplication
by its displayed fixed-curve coordinate.

\item The half-source family map proved in P3K.45 is
\[
 (Z,V,q_T,\rho)\longmapsto(q_T/Z,V,q_T,-\rho).
                                                               \tag{RC.4}
\]
It carries the vertical edge at \(0\) to the vertical edge at \(e_1\).
Let the target source period, in the target marking, be
\[
 h_{\rm tar}(Z',V',q)
       =(\kappa qZ',\alpha(q,\rho)V',q).             \tag{RC.5}
\]
Here \(\alpha\) means the target line multiplier, including the inverse
Picard marking produced by the half-source continuation. It is the
actual nonzero source-period unit from P3F.54. We do not replace it by
the source multiplier with its old marking. Composing (RC.4) with the
literal inverse of (RC.5) gives
\[
       (Z,V,q)\longmapsto(Z^{-1},\alpha^{-1}V,q).
                                                               \tag{RC.6}
\]
The cancellation of \(\kappa\) in (RC.6) is the equality
\((q_T/Z)/(\kappa q)=Z^{-1}\). Its winding has been retained in both
maps before this composition.

\item Substitution in (RC.3), using (RC.1), gives the fixed-edge arrow
\[
 z'=z^{-1},\qquad
 d_0=\alpha z c_0,\qquad d_1=\alpha^{-1}z c_1.
                                                               \tag{RC.7}
\]
In the second target chart this is
\[
 (z')^{-1}=z,\qquad z'd_0=\alpha c_0,
                   \qquad z'd_1=\alpha^{-1}c_1.       \tag{RC.8}
\]
Consequently the map extends over both endpoints of the fixed curve.
It preserves the physical height exactly:
\(z'd_0d_1=z c_0c_1\). It also preserves the physical delta action,
whose normal weights are \((-1,+1)\).

\item The coordinates in (RC.1) on the fixed curve agree with the ordered
normalization coordinate \(z=(s-\rho)/(s+\rho)\) of P3K.40/P3K.42:
both use the same two ordered branches and send the zero section to
\(z=1\). The target uses its continued branch order. Thus (RC.7) is
an arrow between the specified fixed-edge products, including their
source coordinate; it is not an unidentified bundle isomorphism.

\par

\medskip\noindent\textbf{RC.2. Its literal Cox product and extension over the cap}

\par\nobreak

\item Only the normal derivative of (RC.7) enters its normal-sphere product.
We retain the actual nonlinear unit by the explicit homotopy
\[
 \alpha_t(q,\rho)=\alpha(t^2q,\rho),\qquad 0\le t\le1.
                                                               \tag{RC.9}
\]
On a uniformly small height disc all these functions are nonzero.
Equations (RC.7)--(RC.8) with \(\alpha_t\) are invertible,
delta-equivariant and preserve \(q\); at \(t=0\) the unit is
\(a(\rho)=\alpha(0,\rho)\). We next use the exact diagonal homotopy
\(a_s=a|a|^{-s}\), \(0\le s\le1\), in those same equations.
It ends at the closed circle-valued function
\[
                         \tau(\rho)=a(\rho)/|a(\rho)|. \tag{RC.10}
\]
No logarithm or square root of this function is chosen. Any winding
is retained.

\item Use the homogeneous Cox coordinates \((a_0,a_1;b_0,b_1)\) with
auxiliary equivalence
\((a,b)\sim(e^{i\phi}a,e^{-i\phi}b)\). On the chart \(a_0\ne0\),
the local source coordinate is \(z=a_1/a_0\), and the local normal
coordinates are \(c_i=a_0b_i\). The unitary version of (RC.7) is
therefore exactly
\[
 (a_0,a_1;b_0,b_1)
       \longmapsto(a_1,a_0;\tau b_0,\tau^{-1}b_1).
                                                               \tag{RC.11}
\]
Indeed its target local normal coordinates are
\(a_1\tau b_0=\tau zc_0\) and
\(a_1\tau^{-1}b_1=\tau^{-1}zc_1\), as required.

\item Keep the same quaternion convention and the same ordered normal
coordinates as in P3K.40/P3K.42:
\[
 u=a_0+ja_1,\qquad v=b_0+j\overline{b_1},\qquad
 h=uiu^{-1},\qquad q_H=uv.
\]
Set \(r=ji\) and \(d=-i\). Literal multiplication gives
\[
 r(a_0+ja_1)d=a_1+ja_0,
 \qquad
 \tau b_0+j\overline{\tau^{-1}b_1}=v\tau.
                                                               \tag{RC.12}
\]
Hence (RC.11) acts on the actual product as
\[
 (h,q_H)\longmapsto
 \left(rhr^{-1},\;r\exp(-\pi h/2)q_H\tau(\rho)\right).
                                                               \tag{RC.13}
\]
To obtain the second entry we use the exact equality
\(udu^{-1}=\exp(-\pi h/2)\), with \(d=\exp(-\pi i/2)\).
This retains the determinant of the constant interchange of the two
Cox source coordinates. It is not replaced by a different lift of
that projective map.

\item Formula (RC.13) has the following explicit extension over the
actual product cap:
\[
 \begin{aligned}
 D^3\times S^3\times S^1_\rho&\longrightarrow
                  D^3\times S^3\times S^1_\rho,\\
 (y,q_H,\rho)&\longmapsto
 \left(ryr^{-1},\;r\exp(-\pi y/2)q_H\tau(\rho),\;-\rho\right).
 \end{aligned}                                               \tag{RC.14}
\]
Its inverse first takes \(y=r^{-1}y'r\), and then takes
\(q_H=\exp(\pi y/2)r^{-1}q'_H\tau(\rho)^{-1}\), with
\(\rho=-\rho'\). Thus it is a smooth diffeomorphism, including
\(y=0\). The base map is an orientation-preserving real rotation;
the fibre map is left and right multiplication by unit quaternions.
The positive rho-circle is carried to its positive half rotation.
It follows that (RC.14) preserves all the product orientations used
in the completion.

\item The physical Hopf action is \(q_H\mapsto q_He^{-i\theta}\).
Since \(\tau(\rho)\) is complex, (RC.14) commutes with this action
with its given orientation. It therefore descends to a map of
\(D^3\times S^2\times S^1_\rho\) and is an isomorphism of the
actual principal circle bundles above that map. In particular the
meridian Euler class is preserved, with sign \(+1\).

\item The homotopies (RC.9)--(RC.10) also give maps of the original
completions, not just their linear models. More explicitly, for each
homotopy of the boundary maps we take its velocity
\((\partial_tF_t)F_t^{-1}\), multiply it by a collar cutoff which is
one at the inner boundary and zero at the outer boundary, and
integrate this time-dependent vector field. It is periodic in rho
and delta-equivariant. Its flow supplies the required collar
diffeomorphism and transports the same principal circle bundle.
Thus composing (RC.14) with that flow extends the original boundary
map (RC.7), with all its units.

\par

\medskip\noindent\textbf{RC.3. The completed quotient and its characteristic numbers}

\par\nobreak

\item Let \(E_T\) and \(E_{T'}\) denote the complements of the fixed tubes
in the two actual families over the rho-circle. P3K.45 supplies the
literal isomorphism
\[
       E_{T,\rho}\longrightarrow E_{T',-\rho},       \tag{RC.15}
\]
including its finite cap and every ordinary annulus. It is
delta-equivariant and orientation-preserving. On the tube boundary,
(RC.15) is exactly the map calculated in RC.1. Use the collar flow
just constructed so that it matches (RC.14) on the attached cap.
We thereby obtain a map of the actual glued total spaces, not only
of a single fibre or of their homology groups.

\item Divide by the physical circle and attach
\(D^3\times S^2\times S^1_\rho\) in the same Cox product. Denote
this physical positive-parameter completion by \(P_T\), and denote
the corresponding reflected completion by \(P_{T'}\). The preceding
maps glue to an oriented diffeomorphism
\[
              \widehat I:P_T\longrightarrow P_{T'},
       \qquad \widehat I^*x^P_{T'}=x^P_T.            \tag{RC.16}
\]
Here \(x^P_T,x^P_{T'}\) are the Euler classes of the actual meridian
lines. The pullback identity follows from the explicitly equivariant
principal-bundle map in RC.2 and its identical map on the exterior.
Periodicity of (RC.14) is what makes (RC.16) a map of the completed
mapping-torus quotients. Equivalently, one may choose the mutually
transported connections of P3K.45; the fibre arrows then give the
same glued map with the displayed positive half rotation of the
parameter circle.

\item For the positive physical returns \(f_T,f_{T'}\), P3K.42 gives
\(P_T=M_{f_T^{-1}}\) and \(P_{T'}=M_{f_{T'}^{-1}}\) in Jin's
endpoint convention. In particular the proved pair of \(P_T\) is
\((-6,0)\). Naturality applied directly to (RC.16) first gives that
same physical pair on \(P_{T'}\).

\item Now put \(N_T=M_{f_T}\) and \(N_{T'}=M_{f_{T'}}\). These are obtained
from their respective physical completions by the same reversal of
the parameter circle, carrying the meridian line unchanged. Each
reversal changes the six-manifold orientation and both characteristic
numbers. Conjugating (RC.16) by these two actual parameter reversals
gives an orientation-preserving map \(N_T\to N_{T'}\), still with
pullback of the meridian Euler class equal to the original class.
Thus, writing those classes as \(x_T,x_{T'}\), we obtain the exact
integer equalities
\[
 \begin{aligned}
 \langle x_{T'}^3,[N_{T'}]\rangle
       &=\langle x_T^3,[N_T]\rangle,\\
 \langle p_1(N_{T'})x_{T'},[N_{T'}]\rangle
       &=\langle p_1(N_T)x_T,[N_T]\rangle.
 \end{aligned}                                               \tag{RC.17}
\]
P3K.42 evaluates the right-hand sides as \((6,0)\). This proves the
stated reflected pair. Its dual index is \(A_U=(4k-\ell)/24=1\), and the explicit cap comparison gives \(A_M=-1\), in agreement with the oriented conjugacy of \cref{night:P3K.45}.

\item Finally the actual fixed-edge product is the one retained throughout
RC.1--RC.3. The exact fixed-edge identities of P3K.23/P3K.42 therefore
give
\[
        n_{C_0}=-6A_U=-6,\qquad
        m_{C_0}=6A_U-k=0,\qquad s=-1\pmod {28}.         \tag{RC.18}
\]
These are the relative coordinates for \(C_0\) in the reflected fan.
No relative integer for its flopped nonfixed edge \(C_2\) is obtained
by this calculation.

\item The exact certificate is \path{code/part3/source_cusp_reflected_pair.py},
with its JSON and captured output. It verifies the chart maps,
source-period cancellation, quaternion interchange, physical-circle
commutation and the stated characteristic-number arithmetic.

\end{enumerate}
\end{proof}




\begin{theorem}\label{night:P3K.50}
\StatementRecord{night:P3K.50}{false as stated; correction proved}{S6\_06\_BLIND: GPT-6 Astra (OpenAI, ChatGPT Pro mode, effort Pro), recorded re-derivation in the checkpoint dated 18 November 2026}{\cref{night:P3K.23}, \cref{night:P3K.31}, \cref{night:P3K.40}, \cref{night:P3K.42}, \cref{night:P3K.48}}In the actual ordinary plumbing product of P3K.48, the positive
selected-cusp source return restricted to the exterior of the
original fixed curve has relative coordinates
\[
 (m_0,n_0,s)=(6,-6,-1),\qquad V_{00}=0,\qquad\delta_0=-12.
\]
The derivative is \(+12\) on the product cycle \([D^3_h]\) and
\(-12\) on its ambient relative dual \(D_0=-[D^3_h]\).
The native-to-ordinary product change at each of the two nonfixed
stars likewise adds the explicit pure-left relative class \(L^6\)
and leaves its right integer unchanged. This statement computes
those comparison maps; it assigns no native nonfixed integer or
further fan return.
\end{theorem}
\begin{proof}
\begin{enumerate}[label=\arabic*.]
\item \par

\medskip\noindent\textbf{OR.1. The exact direction of the two boundary products}

\par\nobreak

\item Let \(E_\theta\) be the exterior of the actual fixed curve, with the
parameter running positively from \(0\) to \(2\pi\). Write
\(b^H_\theta:\partial E_\theta\to\partial E_0\) for the actual
Hopf/Cox product of P3K.40/P3K.42 and
\(b^O_\theta:\partial E_\theta\to\partial E_0\) for the restriction
of P3K.48's ordinary plumbing product. These are actual-to-reference
coordinate arrows, both based at the identity.

\item We verify their local coordinates before taking a homotopy. In the
first fixed-edge chart of OE.4,
\[
 Z=b,\qquad(c_0,c_1)=(a,c),\qquad q=bac,\qquad
 q_H=U(b)(a+j\overline c),\quad
 U(b)=\frac{1+jb}{\sqrt{1+|b|^2}}.                 \tag{OR.1}
\]
Here \(q\) is the original source-height coordinate. The retained
period multiplier \(\kappa=q_T/q\) is not substituted for this
height. The second standard Hopf chart has coordinates
\[
 (Z',c'_0,c'_1)=(\kappa w,\alpha v,(\kappa\alpha)^{-1}u).
                                                               \tag{OR.2}
\]
Substitution of OE.4 gives exactly
\(Z'=b^{-1}\), \(c'_0=ba\), \(c'_1=bc\).
The ordinary product instead uses the second vertex coordinates
\((u,\eta v,\kappa w)\), with the same transformation at the
initial reference parameter. Its first-chart normal arrow is the
literal \(\mathcal K_0=\mathcal D_0 T_-\) of OP.9. Thus, after the
specified based homotopy of the actual zero-winding unit factors,
\[
 \Gamma_\theta=b^O_\theta(b^H_\theta)^{-1},\qquad
 \Gamma_\theta(h,q_H)=
       (h,L(h,\theta)q_H R(h,\theta)),
 \quad L=UHU^{-1},\quad R=H^{-1}.                     \tag{OR.3}
\]
This is the direction of OP.20, rather than its inverse. The full
\(\mathcal D_0\) is retained before that based homotopy. Its
positive normal metric and its based periodic logarithms extend the
homotopy on the protected collars, so it changes neither endpoint
relative class. The higher normal terms are carried by the same
normal-rescaling collar isotopy as OP.3. No additional source
inversion, Tate-height substitution or line gauge is made.

\item P3K.48 OP.5 computed the degrees of the two closed maps in (OR.3):
\[
 (\deg L,\deg R)=(-6,0)
 \quad\text{on }S^2_h\times S^1_\theta.                \tag{OR.4}
\]
The complex-curve orientation has the opposite sign; (OR.4) uses
exactly the \(h\)-sphere boundary orientation of P3K.23.

\par

\medskip\noindent\textbf{OR.2. Extension of the inverse gauge and its endpoint degree}

\par\nobreak

\item Use the actual oriented exterior product
\(E_0=D^3_y\times S^3_{q_H}\) of P3K.23. Each of the based
homotopies \(L^{-1},R^{-1}:S^2\times[0,2\pi]\to SU(2)\)
extends to \(D^3\times[0,2\pi]\), with value one at the initial
end. To see this as an exact extension problem, the restriction
\[
 D^3\times\{0\}\cup S^2\times[0,2\pi]
                    \longrightarrow D^3\times[0,2\pi]
\]
is a cofibration with an explicit retraction. Rescale time to
\(t\in[0,1]\), set
\[
 \lambda=\min\{2/(2-t),1/|y|\},\qquad
 r(y,t)=(\lambda y,\,2+\lambda(t-2)),
\]
with \(1/0=\infty\). The image is on the bottom or side, and the
map fixes that subspace. Straight interpolation between the identity
and this map is a deformation through the cylinder. Composing the
prescribed bottom-and-side map with \(r\) gives a continuous
extension. Relative smoothing in \(\mathbb R^4\), followed by
radial normalization to its unit sphere, gives smooth extensions
while retaining the protected side and bottom collars, denoted
\(A(y,\theta)\), \(B(y,\theta)\), without changing their prescribed
boundary maps. Set
\[
 \psi_\theta(y,q_H)=
       (y,A(y,\theta)q_HB(y,\theta)).                \tag{OR.5}
\]
This is a smooth orientation-preserving diffeomorphism for every
parameter, with inverse obtained by replacing the two multipliers
by their inverses. On the boundary it is precisely
\(\Gamma_\theta^{-1}\). Its endpoint is relative to the boundary.

\item We fix the connecting sign by integration. Orient
\(D^3\times[0,2\pi]\) by \(\Omega_{D^3}\wedge d\theta\).
Its side has orientation \(\Omega_{S^2}\wedge d\theta\), while
its top has orientation \(-\Omega_{D^3}\). For the normalized
closed degree-three form \(\omega\) of OP.21, Stokes gives
\[
 0=\int_{S^2\times[0,2\pi]}A^*\omega
                    -\int_{D^3}A(\cdot,2\pi)^*\omega.
\]
The bottom integral is zero. Inverting an \(SU(2)\) map reverses
its third class, so (OR.4) yields
\[
 \deg A(\cdot,2\pi)=6,\qquad
 \deg B(\cdot,2\pi)=0.                               \tag{OR.6}
\]
These are degrees of maps \((D^3,S^2)\to(S^3,1)\), with exactly
P3K.23's orientation.

\item The endpoint map (OR.5) is a product of the actual left and right
multiplications of K23.5. Homotopies of the two coefficient maps
relative to their boundary identify its relative class as
\[
                         [\psi_{2\pi}]=L^6.           \tag{OR.7}
\]
There is no disc-supported central factor: the same formula extends
over \(D^3\times D^4\), and its coefficient maps are classified by
the two degrees in (OR.6). Changing either extension changes neither
of these degrees nor the endpoint relative class represented by
these fibrewise multiplications.

\par

\medskip\noindent\textbf{OR.3. The actual physical return and its three coordinates}

\par\nobreak

\item Let \(P^H_\theta:E_0\to E_\theta\) be parallel transport chosen
so that \(b^H_\theta P^H_\theta\) is the identity on the boundary.
Equation (OR.3) gives the ordinary-product transport
\[
                       P^O_\theta=P^H_\theta\psi_\theta,
 \qquad b^O_\theta P^O_\theta=\operatorname{Id}
                       \text{ on }\partial E_0.      \tag{OR.8}
\]
Its endpoint is therefore
\[
                         \Phi_O=\Phi_H\psi_{2\pi}.   \tag{OR.9}
\]
These are the actual forward returns of the positive parameter loop;
we do not identify their physical mapping tori with Jin's oppositely
parameterized completion. That sign was already retained when
P3K.42 computed \([\Phi_H]=L^0R^{-6}\mathcal C^{-1}\).
Connections compatible with the same fixed boundary product form
an affine space; an interpolation with fixed boundary gives a
relative isotopy of their endpoints. Thus (OR.9) also computes the
return from the connection fixed on all three components of the
ordinary plumbing in OP.3.

\item The generators \(L,R\) commute literally, and \(\mathcal C\) is
central. Substituting (OR.7) consequently proves
\[
       [\Phi_O]=L^6R^{-6}\mathcal C^{-1},
       \qquad (m,n,s)=(6,-6,-1).                      \tag{OR.10}
\]
The fixed-curve variation is \(m+n=0\). In the Cartesian/left-
quaternion stable frame of P3K.23 its relative stable derivative is
\(-2n=12\) on the product-oriented relative cycle \([D^3_h]\).
In the ambient-relative-dual convention used for the common exterior,
\(D_0=-[D^3_h]\), so the same derivative evaluates \(-12\) on
\(D_0\). The identification with the ambient stable frame also retains the
actual frame-change map of P3K.31: its third class is
\([g_0]=e_0\), and its derivative correction on this cycle is
\(-V_{00}=0\), because the fixed-curve variation was just
computed to vanish. Hence the ambient value is exactly
\(\delta_0=-12\). These statements concern this one exterior restriction
in the explicitly constructed ordinary product. The two nonfixed
edge restrictions, their mutual variations, and any further fan
return require additional actual maps and are not assigned here.

\par

\medskip\noindent\textbf{OR.4. The two exact nonfixed single-star frame comparisons}

\par\nobreak

\item The same calculation identifies, without evaluating their native
returns, the two remaining changes of boundary product. For \(C_1\),
the native equivariant star of OE.7--8 has first coordinates
\((z,n_0,n_1)=(c,a,b)\) and second coordinates
\[
 (z',n'_0,n'_1)=(\eta v,g w,(g\eta)^{-1}u).
\]
The transition is exactly \(z'=z^{-1}\), \(n'=zn\).
For \(C_2\), the native star has first coordinates
\((z,n_0,n_1)=(a,b,c)\) and second \((u,v,w)\), with the
same standard transition. Thus both native products use
\(q_i=U(z)(n_0+j\overline{n_1})\), in the identical source/normal
orientation convention. The complete parameter units in the first
star have their original extending logarithms, while the second
star is already parameter-free in these charts.

\item The ordinary product of OP.1--3 consequently gives the actual
native-to-ordinary boundary arrows, after their carried based unit
homotopies,
\[
 \begin{array}{ll}
 C_1:&q_1\longmapsto (UHU^{-1})q_1 H,\\
 C_2:&q_2\longmapsto (UHU^{-1})q_2 H^{-1}.
 \end{array}                                             \tag{OR.11}
\]
The left factor is the identical closed map of OP.20 on the
corresponding source sphere. Each right factor depends on only the
radial coordinate and parameter. Therefore each arrow has factor
degrees \((-6,0)\) in its own \(h\)-sphere orientation. The full
extension and direction proof of OR.2--3 applies to these displayed
arrows: the ordinary relative return is the native single-star
relative return composed with the explicit pure-left endpoint of
degree six. In particular this boundary change leaves the native
right relative integer unchanged. This computes the comparison
map, not the native nonfixed return. Those two actual returns and
their mutual common-exterior entries are still required.

\item The exact certificate is \path{code/part3/source_cusp_ordinary_relative.py}, with its JSON and captured output.

\end{enumerate}
\end{proof}




\begin{theorem}\label{night:P3K.52}
\StatementRecord{night:P3K.52}{false as stated; correction proved}{S6\_06\_BLIND: GPT-6 Astra (OpenAI, ChatGPT Pro mode, effort Pro), recorded re-derivation in the checkpoint dated 18 November 2026}{\cref{night:P3F.56}, \cref{night:P3K.31}, \cref{night:P3K.50}, \cref{night:P3K.51}}In the actual ordinary plumbing product of P3K.48, use the ambient
stable frame and integral relative dual cycles of P3K.51. The
positive selected-cusp source return has stable-derivative vector
\[
 \delta=(-12,-12,0)^T.
\]
Its three ordinary single-curve relative pairs are
\[
 (m_0,n_0)=(6,-6),\qquad(m_1,n_1)=(6,-6),\qquad
 (m_2,n_2)=(0,0).
\]
Their physical central coordinate is $-1$. The actual native nonfixed products
have pairs \((0,-6)\) at \(C_1\) and \((-6,0)\) at \(C_2\).
These are computed from the literal derivative action and its
closed Spin lifts; no subsequent fan-product reset or general fan
return is asserted here.
\end{theorem}
\begin{proof}
\begin{enumerate}[label=\arabic*.]
\item \par

\medskip\noindent\textbf{ST.1. The exact derivative prism and its sign}

\par\nobreak

\item Let \(E\) be the actual common exterior, \(P_t:E\to E_t\) the
ordinary-product parallel transport, and \(F=P_1\) its positive
return. On the reference exterior put
\(a_t=P_t^{-1}\delta_tP_t\), with \(t\in[0,1]\). The actions
preserve the common boundary, and there are exact identities
\[
 a_1(\phi)=F^{-1}a_0(\phi)F,\qquad
 \Delta_1=F^{-1}\Delta_0F.                            \tag{ST.1}
\]
Fix an oriented global ambient stable frame of the reference
six-sphere and restrict it to \(E\). Write \(c_F\) for the stable
derivative of \(F\), and \(d_t(\phi,p)\) for the derivative of
\(a_t(\phi)\), both in this frame. Their matrices lie in the stable group
\(GL^+\), not necessarily in its orthogonal subgroup. We lift them
to the universal covering group \(\widetilde{GL^+}\), retaining
its actual group multiplication. The following endpoints are
prescribed. The lift \(c_F\) is one on the
boundary, because \(F\) is the identity on a boundary collar and
\(E\) is simply connected. Lift \(d_t\) starting at one for
\(\phi=0\). Its endpoint at \(\phi=2\pi\) is one: the actual
circle action extends to the vertices, whose three complex weights
are \((-1,0,1)\) or \((1,-1,0)\). Their sum is zero, so the
stable rotation loop has a closed lift there. Continuity on the
connected exterior fixes the same endpoint everywhere and for all
\(t\). No unclosed half-angle lift is used.

\item The chain rule, with every source and target retained, now gives
the literal identity in the covering group
\[
 d_1(\phi,p)=c_F(a_1(\phi)p)^{-1}
                    d_0(\phi,Fp)c_F(p).              \tag{ST.2}
\]
Both sides have lift value one at \(\phi=0\), so no central sign is
missing. The explicit polar deformation
\[
 A\longmapsto A(A^T A)^{-s/2},\qquad 0\leq s\leq1,
\]
lifts from the identity to a deformation retraction of this covering
group onto Spin. We use this deformation to identify its third
class; we do not apply it to a product identity. Normalize the
integral third class \(\omega\) by degree one on left quaternion
multiplication. In the
stable real vector representation we can use
\[
 \omega(g)=-\frac{\operatorname{tr}_{\mathbb R}
                         (g^{-1}dg)^3}{48\pi^2},
\]
whose restriction to the left quaternion subgroup has integral one.
The Maurer form is closed on the full linear group, and the displayed
deformation proves that it represents the same integral class on
its cover. The local orthogonal blocks below lie in the Spin
subgroup of this cover.
Its Maurer--Cartan product identity has exact boundary form
\(\operatorname{tr}_{\mathbb R}(g^{-1}dg\wedge dhh^{-1})/(16\pi^2)\).
Put \(\Omega=c_F^*\omega\), a relative closed three-form, and
\(\delta(D)=\langle\Omega,D\rangle\).

\item For a relative two-chain \(x\) with boundary \(\gamma\), integrate
(ST.2) on \(S^1_\phi\times x\). Every product boundary form vanishes:
\(c_F\) is one and its derivative is zero on the collar containing
\(S^1\times\gamma\). The term pulled back from \(c_F(p)\) factors
through a two-chain. The remaining difference involving \(d_0\)
is its integral on \(S^1\times(Fx-x)\). Since \(F\) fixes the
boundary, \(Fx-x\) is an absolute two-cycle; it bounds a three-chain
because \(H_2(E)=0\). The integral of that closed form is therefore
zero. We obtain
\[
 \int_{S^1\times x}(d_1^*\omega-d_0^*\omega)
                   =-\langle\Omega,\Delta_1x\rangle. \tag{ST.3}
\]
The actual action prism \(K\) of CT.8 has
\(\Delta_1x-\Delta_0x=\partial Kx-K\gamma\).
The first term pairs to zero with the closed form \(\Omega\), and
the second lies in the boundary where \(\Omega=0\). Thus the
right side of (ST.3) is \(-\delta(\Delta_0x)\).

\item Apply Stokes to the derivative map on
\([0,1]_t\times S^1_\phi\times x\), with this displayed product
orientation. Its boundary is top minus bottom plus
\([0,1]\times S^1\times\gamma\). Consequently
\[
 \int_{S^1\times x}(d_1^*\omega-d_0^*\omega)
   =-W(\gamma),\qquad
 W(\gamma)=\int_{[0,1]\times S^1\times\gamma}d_t^*\omega.
                                                               \tag{ST.4}
\]
Combining the two exact signs proves
\[
                       \delta(\Delta_0x)=W(\gamma).    \tag{ST.5}
\]

\par

\medskip\noindent\textbf{ST.2. Actual local frames and the closed Spin blocks}

\par\nobreak

\item We compute (ST.4) using the same generic normal sections and
vertex-link paths as CT.3. This is an actual frame comparison.
For a fixed change of stable frame \(g\), the derivative action
changes by
\[
 d_t'(\phi,p)=g(a_t(\phi)p)^{-1}d_t(\phi,p)g(p).
                                                               \tag{ST.6}
\]
On the closed graph trace, the last term factors through its path,
and the exact Maurer boundary terms cancel. Hence its third
integral changes by minus the pairing of \(g^*\omega\) with the
actual orbit trace \(K\gamma\).

\item On each edge choose the reference complex coordinate frame on the
annulus times its normal four-disc, retaining its stable orientation.
Its comparison with the global ambient frame extends through that
four-disc. This domain has zero third cohomology. Close the edge
trace with the rank-one \(D^2_t\times S^1_\phi\) endpoint caps of
CT.3; the preceding pairing with \(g^*\omega\) is then zero.
Thus this frame change contributes zero to its integer. The vertex
coordinate frames extend over their actual balls. The vertex traces
and the same caps are independent of \(t\), hence their third
integrals vanish as well. If a determinant adjustment is used to
choose a Spin-compatible complex frame, its scalar factor has no
third class and is retained in (ST.6). This verifies the passage to
local frames on closed cycles, rather than identifying two frames
by their names.

\item The full physical action in these coordinates is OP.19. Its
zero-winding unit factors have the actual based logarithm homotopy
already displayed there. The higher normal terms are removed by
the literal normal rescaling of OP.3. To compute the derivative
trace, move its normal evaluation point to the zero section through
that four-disc, where the coordinate frame and action both extend.
At the zero section the derivative is block triangular, with base
rotation and normal block. Its off-diagonal block vanishes there;
the rescaling is a homotopy through invertible derivatives. All end
faces remain parameter-independent. Thus these operations preserve
the closed third integrals and retain the actual full-unit action.

\item In the actual quaternion coordinate
\(Q=n_0+j\overline{n_1}\), the three normal blocks are exactly
\[
 \begin{array}{c|c|c|c}
 &\text{base rotation}& A& B\quad (Q\mapsto AQB)\\ \hline
 C_0&1&1&H e^{-i\phi}H^{-1}\\
 C_1&e^{i\phi}&H e^{-i\phi/2}H^{-1}&H^{-1}e^{-i\phi/2}H\\
 C_2&e^{-i\phi}&H e^{i\phi/2}H^{-1}&H e^{-i\phi/2}H^{-1}.
 \end{array}                                             \tag{ST.7}
\]
For the stable class, a left factor has degree \(+1\) and a right
factor has degree \(-1\). This is the same normalization as
P3K.31: \(\operatorname{Ad}_q=L_qR_{q^{-1}}\) has stable degree
two. Equivalently the real quaternion trace restricts to the
positive normalized form on the left factor and its negative on
the right factor. The degree of (ST.7)'s normal block is therefore
\(\deg A-\deg B\).

\item At the two nonfixed curves, the normal half-angle pair ends at
\((-1,-1)\), representing the central minus sign in the normal
Spin lift. The source \(SO(2)\) rotation has Spin endpoint minus
one. Their images in Spin of the full six-dimensional tangent
space multiply to plus one. The full lift therefore closes exactly.
The source factor depends only on \(\phi\), so its own third
form is zero. Product-form boundary terms vanish on the same
parameter-independent end faces and constant central-sign angle
faces used in CT.2. Thus the half-angle block calculation evaluates
the closed full stable lift, not an artificially closed normal loop.

\par

\medskip\noindent\textbf{ST.3. The full derivative vector and the three relative pairs}

\par\nobreak

\item CT.3's axis calculations, with the orientation
path--parameter--circle, give
\[
 \begin{array}{c|rr|r}
 &\deg A&\deg B&W_i=\deg A-\deg B\\ \hline
 C_0&0&-6&6\\
 C_1&-3&3&-6\\
 C_2&3&-3&6.
 \end{array}                                             \tag{ST.8}
\]
The vertex connectors contribute zero by ST.2. The graph loops are
\(\gamma_j=a_j-a_0\); hence
\[
                  W(\gamma_1)=-12,\qquad W(\gamma_2)=0.
\]
CT.7 supplies the exact integral signs
\(\Delta x_1=D_1\), \(\Delta x_2=-D_2\). Equation (ST.5)
therefore gives \(\delta_1=-12\) and \(\delta_2=0\).
K50 computed the remaining entry in the same ambient frame and
relative-dual orientation. Thus
\[
                       \delta=(-12,-12,0)^T.           \tag{ST.9}
\]

\item For clarity the frame relation to each single-curve exterior is also
retained. K31's actual product-frame comparison has third class
\(e_i\). Its derivative correction is \(-V_{ii}\); all three
diagonal entries in K51 are zero. K31.7 consequently gives
\[
 2n_i=\delta_i-V_{ii},\qquad m_i=-V_{ii}-n_i.
\]
The actual ordinary relative pairs are therefore
\[
 (m_0,n_0)=(6,-6),\qquad
 (m_1,n_1)=(6,-6),\qquad
 (m_2,n_2)=(0,0).                                     \tag{ST.10}
\]
Each exterior's central coordinate is the same oriented closed
return, namely $-1$ by the explicit F56-to-physical cap map of \cref{night:P3K.68}. The native nonfixed products differ by
the explicit pure-left correction in OR.4: subtracting that \(L^6\)
gives native pairs \((0,-6)\) at \(C_1\) and \((-6,0)\) at
\(C_2\), with the central coordinate unchanged. No additional fan
return is assigned by this calculation; the actual flop maps and
subsequent basis/frame transports must still be applied.

\item The exact certificate is \path{code/part3/source_cusp_stable_trace.py}, with its JSON and captured output.

\end{enumerate}
\end{proof}




\begin{theorem}\label{night:P3K.54}
\StatementRecord{night:P3K.54}{false as stated; correction proved}{S6\_06\_BLIND: GPT-6 Astra (OpenAI, ChatGPT Pro mode, effort Pro), recorded re-derivation in the checkpoint dated 18 November 2026}{\cref{night:P3K.23}, \cref{night:P3K.31}, \cref{night:P3K.40}, \cref{night:P3K.42}, \cref{night:P3K.45}, \cref{night:P3K.50}, \cref{night:P3K.52}}Retain exactly the positive selected-cusp source loop, scalar, finite cap,
source periods and line maps of P3F.54/P3F.56. In its initial fan,
the actual first periodic flops at the fixed vertical curve \(C_0\),
the horizontal curve \(C_1\), and the diagonal curve \(C_2\) have
oriented return classes
\[
 21,\qquad21,\qquad27\quad\text{in }\mathbb Z/28,
\]
respectively. The proof retains each native family product and its
actual exterior boundary arrow. No subsequent ordinary-product
reset or general fan formula is assigned by this statement.
\end{theorem}
\begin{proof}
\begin{enumerate}[label=\arabic*.]
\item \par

\medskip\noindent\textbf{FF.1. The three native family stars and their actual maps}

\par\nobreak

\item For the fixed star \(C_0\) use P3K.40/P3K.42's literal source
coordinate and physical normal chart:
\[
 z=Z=\frac{s-\rho}{s+\rho},\qquad
 q=zc_0c_1,\qquad z'=z^{-1},\quad c_i'=zc_i.
                                                               \tag{FF.1}
\]
Its complete normal framing, physical delta action and retained
\(q_T/q\) winding are those proved in P3K.40 and P3K.42.

\item The two other native family stars are also actual coordinate maps,
not only normal-bundle identifications. In the six charts OE.4,
for \(C_1\) take
\[
 (z,c_0,c_1)=(c,a,b),\qquad
 (z',c'_0,c'_1)=(\eta v,g w,(g\eta)^{-1}u),           \tag{FF.2}
\]
where \(g\) is the actual source translation multiplier and
\(\eta=\lambda\beta\) the actual line/suspension multiplier.
Direct substitution of OE.4 gives
\(z'=z^{-1}\), \(c'_i=zc_i\). For \(C_2\) use
\[
 (z,c_0,c_1)=(a,b,c),\qquad (z',c'_0,c'_1)=(u,v,w),  \tag{FF.3}
\]
with the identical standard transition. The functions in (FF.2)
are the original nonzero holomorphic units on the annular source
representative, including their initial constants and their given
extending logarithms. The transformations are invertible on those
charts. They are applied as coordinate arrows to the actual family,
so all source-period, finite-cap, line and annular maps are conjugated
by those same arrows. None of those data is replaced by a marking
chosen only on homology.

\item In particular, every chosen conifold cell has trivial stabilizer
under the period lattice. Choose its quotient representative,
make (FF.1), (FF.2) or (FF.3) there, and extend to all translates by
the original period maps. This is the same equivariant construction
of the periodic flop as P3K.40. The finite cap is outside this cusp
neighbourhood and its actual annular identification is retained.
The full nonconstant units are transported by their literal common
normal-rescaling isotopies on the two small-resolution charts and
on their shared punctured conifold. Their endpoint maps are the
ones in (FF.1)--(FF.3); no new endpoint translation is inserted.

\item Each of these native stars therefore has the actual Cox link
\[
 (S^3_a\times S^3_c)/S^1,\qquad
 (a,c)\sim(e^{i\phi}a,e^{-i\phi}c),
\]
with \(u=a_0+ja_1\), \(v=c_0+j\overline{c_1}\), and
\[
 h=uiu^{-1},\qquad k=v^{-1}iv,\qquad q_H=uv.
\]
The old and new products are respectively \((h,q_H)\) and
\((k,q_H)\). Their exact common boundary arrow and exterior
extension are
\[
 \tau(h,q_H)=(q_H^{-1}hq_H,q_H),\qquad
 T(y,q_H)=(q_H^{-1}yq_H,q_H).                         \tag{FF.4}
\]
The inverse uses \(q_Hyq_H^{-1}\), so the extension is smooth at
\(y=0\) and orientation preserving. This is the actual family arrow
as well: all four integer Cox weights obey the unrestricted
identity SF.8, which intertwines (FF.4) even when their sum is
nonzero. The complete unit isotopies just described carry that
identity with their actual boundaries. Thus \(T\) compares the
actual positive returns, with no discarded parameter winding.

\par

\medskip\noindent\textbf{FF.2. The native relative pairs and the resulting central classes}

\par\nobreak

\item Use exactly K23's relative generators
\[
 L(y,q_H)=(y,\zeta(y)q_H),\qquad
 R(y,q_H)=(y,q_H\zeta(y)),
\]
with its degree-one \(\zeta\) and its oriented central generator
\(\mathcal C\). Their six-sphere classes are zero, while
\(\mathcal C\) is the positive generator in this run's convention.
The native triples of the original positive source return are
\[
 \begin{array}{c|rrr}
 &m&n&s\\ \hline
 C_0&0&-6&-1\\
 C_1&0&-6&-1\\
 C_2&-6&0&-1.
 \end{array}                                             \tag{FF.5}
\]
The first is P3K.42; the other two are the actual native-star
comparison and derivative computation of P3K.50/P3K.52. In
particular, the pure-left ordinary-product changes have already
been removed by the exact arrows OR.11. The central coordinate is
the oriented closed source-loop return $-1$ in each case by \cref{night:P3K.68}.

\item K23's exact relative conjugation for the same map (FF.4) is
\[
 T(L^mR^n\mathcal C^s)T^{-1}
       =L^{-n}R^{m+2n}\mathcal C^{s+n}.              \tag{FF.6}
\]
There is no circle-quotient hypothesis in this identity. We use
only the literal local products, the actual exterior extension and
the already established native integers. Substitution into (FF.6)
gives central coordinates \((-7,-7,-1)\), hence precisely
\((21,21,27)\) modulo28. If one records the first-resolution complex
normal convention after the flop, K31's explicit fixed reflection
\(J(q)=k\overline qk^{-1}\), acting by
\((y,q_H)\mapsto(Jy,Jq_H)\), swaps the two free integers and
preserves the central coordinate. The three resulting native pairs
in that convention are \((-12,6),(-12,6),(-6,0)\).

\item The diagonal value $-1$ agrees with the independently constructed
half-transport/conjugacy of P3K.45; the proof above did not infer its
integer from that agreement. It computes all three first flops
from their actual family products. The extension of an arbitrary
ordinary plumbing product through a later flop is a separate map
and is not asserted here.

\item The exact certificate is \path{code/part3/source_cusp_first_flops.py}, with its JSON and captured output.

\end{enumerate}
\end{proof}




\begin{theorem}\label{night:P3K.56}
\StatementRecord{night:P3K.56}{false as stated; correction proved}{S6\_06\_BLIND: GPT-6 Astra (OpenAI, ChatGPT Pro mode, effort Pro), recorded re-derivation in the checkpoint dated 18 November 2026}{\cref{night:P3K.31}}Retain the actual positive selected-cusp source loop, scalar, finite
cap, source periods and line maps of P3F.54/P3F.56. For every integer
\(n\), use the physical-delta-retaining periodic lattice-basis fan
\(T_n=T(e_1+ne_2,e_2)\). Its actual oriented return class is
\[
                  s(n)=-1-3n(n-1)\pmod{28}.
\]
These fans, with the exact periodic relabellings proved below,
exhaust all allowed lattice-basis fans containing a physical vertical
edge. Their return classes have exact image
\[
                     \{5,7,9,13,19,21,23,27\}.
\]
The native right relative integers of the two nonfixed stars are
\(6n-6\) and \(-6n\), respectively. The proof retains both actual
adjacent-fan character and period arrows. It does not assign a full
fixed-star characteristic pair, a return on a fan with no physical
vertical edge, or a shorter source-loop endpoint.
\end{theorem}
\begin{proof}
\begin{enumerate}[label=\arabic*.]
\item \par

\medskip\noindent\textbf{VF.1. Actual shear fans, characters, periods and coverage}

\par\nobreak

\item Let \(v_n=e_1+ne_2\), \(w=e_2\), for any integer \(n\), and let
\(T_n=T(v_n,w)\) be the actual periodic fan with triangles
\[
 \{a,a+v_n,a+w\},\qquad
 \{a+v_n,a+w,a+v_n+w\},\quad a\in\mathbb Z^2.
                                                               \tag{VF.1}
\]
The original source-height character \(q\), the physical delta
cocharacter \(e_2\), the source exterior, finite cap and line maps
are retained. Its actual coordinates and periods are
\[
 Z_n=Z,\qquad V_n=Z^{-n}V,\qquad h_n=hB^n,\quad B_n=B.
                                                               \tag{VF.2}
\]
Their inverses use \(V=Z_n^nV_n\), \(h=h_nB^{-n}\), so they
are integral coordinate and period arrows on the same actual
quotient. In these characters the complete bounded multipliers are
\[
 A_{11}=\kappa g^n,\quad A_{12}=g,\quad
 A_{21}=\frac{\alpha\eta^n}{\kappa^n g^{n^2}},\quad
 A_{22}=\frac{\eta}{g^n}.                              \tag{VF.3}
\]
These follow by substitution into the actual OE.2 periods; no
bounded factor or rho winding is discarded.

\item These fans exhaust the lattice-basis fans containing a physical
vertical edge. To see this with exact triangles, the three edge
vectors of \(T(v,w)\) are \(v,w,v-w\). Swapping \(v,w\) preserves
the same two sets of triangles. The replacement
\((v,w)\mapsto(v,v-w)\) also preserves the periodic triangulation:
its first triangle is the old second triangle translated by \(-w\),
and the reverse comparison is the same identity. Thus a vertical
edge can be made the second basis vector. If this second vector is \(-e_2\), replace both basis vectors by
their negatives. This preserves the fan exactly: the two triangle
types for \((-v,-w)\), translated by \(v+w\), are the two
old types in reverse order. The second vector is now \(e_2\). Since the basis is unimodular, the first coordinate of
the other vector is \(\pm1\). For \(v=-e_1+ne_2\), the actual
triangle identity
\[
                         T(v,e_2)=T(e_2-v,e_2)
\]
follows by translating the two triangle types by \(v\) and
\(v-e_2\), respectively. The latter basis has first vector
\(e_1+(1-n)e_2\). This proves coverage by (VF.1), using only
periodic relabellings and their integral source-period arrows.

\par

\medskip\noindent\textbf{VF.2. The exact native right integers on each shear fan}

\par\nobreak

\item For \(P_n=\begin{pmatrix}1&0\\n&1\end{pmatrix}\), the full
winding matrix, normal first exponents, and physical weights in
GF.4, GF.8 and GF.13 are
\[
 B_{P_n}=\begin{pmatrix}-6&0\\6n&0\end{pmatrix},\quad
 l=(-6n,-6,0),\quad
 w=(0,1,-1),\quad d=(-2,-1,-1).                       \tag{VF.4}
\]
The actual local orbit and stable-derivative traces are consequently
\[
 I=(12n-6,0,0),\qquad W=(6-12n,-6,6).                 \tag{VF.5}
\]
Apply the integral chain identities GF.19, with their actual
positive complex meridians and relative dual cycles. They give
\[
 V_{11}=V_{22}=0,\qquad
             \delta_1=12n-12,\quad\delta_2=-12n.      \tag{VF.6}
\]
No value of \(V_{00}\) or \(\delta_0\) is used. K31's local
ambient-to-native-product frame map is the same explicit complex
normal map on each of these stars; its derivative correction is
\(-V_{ii}\). Therefore the ordinary right integers are
\[
              n_{C_1}=6n-6,\qquad n_{C_2}=-6n.         \tag{VF.7}
\]
GF.12 and its explicit cylinder extension compare the ordinary
star with its own actual native product by the pure-left map
\(L^6\). Thus (VF.7) are also the native right integers, with no
additional period or frame choice. Their native left integers and
the transverse fixed-curve data are unnecessary for the comparison
that follows.

\par

\medskip\noindent\textbf{VF.3. The two literal adjacent-fan flops and the closed formula}

\par\nobreak

\item A periodic \(C_2\) flop replaces the diagonal joining \(v_n\) to
\(w\) in the parallelogram \(0,v_n,w,v_n+w\) by the diagonal
joining \(0\) to \(v_n+w=v_{n+1}\). Its new triangles are exactly
\[
 \{0,w,v_{n+1}\},\qquad\{0,v_n,v_{n+1}\}.
\]
These are the first triangle of \(T_{n+1}\) and its second triangle
translated by \(-w\). The new labels are: \(C_0\) is the same
vertical edge, \(C_1\) is the newly flopped edge, and \(C_2\) is
the old \(C_1\). The actual character and period arrows are
\[
 Z_{n+1}=Z_n,\quad V_{n+1}=Z_n^{-1}V_n,\qquad
 h_{n+1}=h_nB,\quad B_{n+1}=B.                        \tag{VF.8}
\]
They retain (VF.3), the physical height and the original source
exterior. The \(C_1\) flop is the inverse operation: its adjacent
quadrilateral has vertices \(w,0,v_n,v_n-w\), and its new diagonal
joins \(w\) to \(v_n-w=v_{n-1}\). The resulting fan is \(T_{n-1}\),
with the new diagonal labelled \(C_2\). These are identities of the
actual height-one cones and their character rings, not a comparison
of numerical invariants of unrelated fillings.

\item For every such flop the native star has the two actual charts
GF.11. In the \(C_1\) chart its units are \(A_{12}=g\) and
\(A_{22}=\eta/g^n\), which carry their original extending logarithms;
the \(C_2\) chart is already the literal standard product. The
common Cox arrow is therefore the exact \(\tau\) of FF.4, with
its full unit, line, annular and finite-cap maps transported as in
FF.1. Its exterior extension conjugates the actual positive return
and changes its central coordinate by the native right integer,
by FF.6. This argument is applied afresh on each fan; it makes no
assertion that a previously chosen ordinary product extends through
the new zero section.

\item Write \(s(n)\in\mathbb Z/28\) for the oriented source-loop return
on \(T_n\). The \(C_2\) arrow and (VF.7) give
\[
                     s(n+1)=s(n)-6n.                 \tag{VF.9}
\]
The \(C_1\) arrow gives \(s(n-1)=s(n)+6(n-1)\), the identical
inverse relation. The proved initial physical source return is \(s(0)=-1\).
Summing these exact changes in either direction proves
\[
                    s(n)=-1-3n(n-1)\pmod{28}
                    \qquad\text{for every }n\in\mathbb Z.
                                                               \tag{VF.10}
\]
This covers every physical-delta-retaining lattice-basis fan by
VF.1. It computes the actual specified positive rho loop, and
claims neither a shorter endpoint nor a characteristic pair for
these other circle quotients. Fans with no physical vertical edge
still require the remaining transverse comparison.

\item The formula has period 28: its change under \(n\mapsto n+28\) is
\(28(-6n-81)\). To give its range by a finite exact calculation,
the values for \(n=0,\ldots,14\) are
\[
 27,27,21,9,19,23,21,13,27,7,9,5,23,7,13.
\]
The identity \(s(1-n)=s(n)\), together with period 28, supplies
all remaining residues. Hence the exact image over all integers is
\[
                  \{5,7,9,13,19,21,23,27\}.
\]
The exact certificate is \path{code/part3/source_cusp_vertical_fans.py},
with its complete JSON and captured output.

\end{enumerate}
\end{proof}




\begin{theorem}\label{night:P3K.57}
\StatementRecord{night:P3K.57}{false as stated; correction proved}{S6\_06\_BLIND: GPT-6 Astra (OpenAI, ChatGPT Pro mode, effort Pro), recorded re-derivation in the checkpoint dated 18 November 2026}{\cref{night:P3F.22}, \cref{night:P3K.49}, \cref{night:P3K.62}}For the marked source
\[
 y^2=x^3+(-2+3\epsilon^2)x^2+t x
       +\epsilon^2(1-\epsilon^2)^2,\qquad
 G=(0,\epsilon(1-\epsilon^2)),
\]
retain the selected whole-zero fibre at
\(t_*=(1-\epsilon^2)(1-3\epsilon^2)\), the actual Cartier line
\(M_a=\mathcal O([aG]-[O])\), and the suspension over \(-bG\),
where \(ab=2\). Use either the initial vertical-edge cusp fan or its
literal reflected fan, and the admitted cap data
\[
 \begin{array}{c|c|c}
 (a,b)&\text{cap order}&\text{root type}\\ \hline
 (1,2)&m=4p&\text{good III-star root}\\
 (2,1)&m=4p-1&\text{reduced III root}
 \end{array}\qquad(p\ge1),
\]
with either norm sign \(n=\pm1\) and every retained integral
\(d\delta\) twist. These are the completed sphere families of
P3F.41 and P3K.24, with their prescribed source, line and cap maps.
For a fixed sufficiently small nonzero suspension scalar, the actual
positive \(\epsilon\)-loop has the meridian-one dual characteristic integer and physical oriented
clutched-sphere class
\[
 \begin{array}{c|c|c}
 (a,b)&\displaystyle\int\frac{4x^3-p_1x}{24}
       &\Theta_7\simeq\mathbb Z/28\\ \hline
 (1,2)&0&0\\
 (2,1)&1&-1.
 \end{array}
\]
The actual source/Cartier/cap morphism under
\(v=\epsilon^{-2}-1\) preserves the suspension scalar and retains
the full \(D^{-4}\) root endpoint, the reduced doubled lattice period,
the annular reference overlap and the quadratic scalar cochain.
The theorem assigns no source-zero return and makes no further fan comparison.
\end{theorem}
\begin{proof}
\begin{enumerate}[label=\arabic*.]
\item We retain the complete source, line, cap and frame maps below.
The meridian-one graph line is (U=M^{-1}): the fibre map
(D_M(v)(w)=\langle v,w\rangle/\langle v,v\rangle), with the
Hermitian form linear in its second argument, satisfies
(D_M(e^{i\theta}v)=e^{-i\theta}D_M(v)).
Its (SU(2))-equivariant spinor extension is
((a,b)\mapsto(\overline b,-\overline a)).
The first-to-second cap maps in \cref{night:P3K.68} give
(\lambda_M(F)=-\Sigma_{F^{-1}}) and
(\lambda_U(F)=\Sigma_{F^{-1}}); cancelling the opposite
product collars gives (A_M=-A_U). We therefore distinguish the
intermediate meridian-one number (A_U=(4k_U-\ell_U)/24)
from the physical return (A_M) throughout this calculation.

\item \par\medskip\noindent\textbf{IF.1. The marked infinity chart and selected-node units.}\par\nopagebreak[4]

\item On the cover \(\rho^2=1+v\), put \(\epsilon=\rho^{-1}\). The source map is
\[
 v=\epsilon^{-2}-1,\qquad
 X=\epsilon^{-2}x,\quad Y=\epsilon^{-3}y,\quad T=\epsilon^{-4}t.
\]
It carries the actual marked equation and section to
\[
 y^2=x^3+(-2+3\epsilon^2)x^2+t x
                     +\epsilon^2(1-\epsilon^2)^2,\qquad
 G=(0,g),\quad g=\epsilon(1-\epsilon^2).                 \tag{P3K.57.IF.1}
\]
Write \(A=1-\epsilon^2\), \(t_*=A(1-3\epsilon^2)\), and \(q_*=t-t_*\).
The selected fibre factors and normalizes by the literal formulas
\[
 (x-A)^2(x+\epsilon^2),\qquad
 x=s^2-\epsilon^2,\quad y=s(s^2-1).                     \tag{P3K.57.IF.2}
\]
Its double preimages are \(s=\pm1\), independent of the parameter. The marked
section has \(s=-\epsilon\). The normalization derivative has first coordinate
\(2s\); at \(s=0\) its second coordinate is \(-1\). It is therefore an immersion
through the full parameter disc, with its two double sections disjoint.

\item Choose the nodal multiplicative coordinate \(Z=(s-1)/(s+1)\). Its marked point
is
\[
 P_G=-\frac{1+\epsilon}{1-\epsilon},\qquad
 P_G(-\epsilon)=P_G(\epsilon)^{-1}.                    \tag{P3K.57.IF.3}
\]
The selected Tate parameter satisfies
\[
 \left.\frac{Q}{q_*}\right|_{q_*=0}
       =-\frac{1-\epsilon^2}{64}.                       \tag{P3K.57.IF.4}
\]
Indeed the cubic discriminant has \(t\)-derivative \(-4A\) at \(t_*\), its
elliptic discriminant is sixteen times that cubic discriminant, and \(c_4=16\)
there. The first term of \(j(Q)^{-1}\) is \(Q\), which gives (P3K.57.IF.4).

\item Retain the actual two chord factors
\[
 F_{21}=\frac{y+g+t x/(2g)}x,\qquad
 F_{12}=\frac{y-g-t x/(2g)}x .
\]
At the selected node they are \(A/(2\epsilon)\) and
\(-A/(2\epsilon)\). The transported actual Hom map is
\[
 \phi_{ab}=\frac{q_*}{g F_{ab}},\qquad
 \left.\frac{\phi_{21}}{q_*}\right|_{\rm node}=\frac2{A^2},
 \quad
 \left.\frac{\phi_{12}}{q_*}\right|_{\rm node}=-\frac2{A^2}.       \tag{P3K.57.IF.5}
\]
All three units (P3K.57.IF.3)--\allowbreak (P3K.57.IF.5) extend and are nonzero on a sufficiently small
full \(\epsilon\)-disc. Their full-disc parameter windings are zero. This is
an actual selected-cusp chart statement; it does not trivialize the exterior
suspension at the other singular fibres.

\par\medskip\noindent\textbf{IF.2. The complete suspension keeps the scalar.}\par\nopagebreak[4]

\item Use the canonical meromorphic divisor sections to identify the old Cartier
line with the line of (P3K.57.IF.1). Translation by \(G\) commutes with the displayed
source isomorphism. The primitive chord lift \(U_a\), covering \(-G\), has
target multiplier \(X(p-G)^{-a}\); consequently
\[
 I_\epsilon U_{a,\rm old}I_\epsilon^{-1}
       =\epsilon^{2a}U_{a,\rm new},\qquad
 q_{\rm old}=\epsilon^{-4}q_* .                         \tag{P3K.57.IF.6}
\]
The actual suspension is
\[
 B_{a,b}=2(-1)^{b+1}\lambda q_{\rm old}U_{a,\rm old}^{\,b},
 \qquad ab=2 .
\]
Thus its transported map is exactly
\[
 I_\epsilon B_{a,b}I_\epsilon^{-1}
       =2(-1)^{b+1}\lambda q_*U_{a,\rm new}^{\,b}.       \tag{P3K.57.IF.7}
\]
The exponent \(2ab-4\) is zero. In degree two the Cartier multiplier has the
same scalar factor at source and target; these factors cancel in the
conjugation as well. Therefore \(\lambda_{\rm new}=\lambda_{\rm old}\).

\par\medskip\noindent\textbf{IF.3. The finite root comparison retains its endpoint and lattice period.}\par\nopagebreak[4]

\item Let the old cap use \(z_{\rm old}^m=T^{-1}\), and the new cap use
\(z_{\rm new}^m=t^{-1}\). On the universal \(\epsilon\)-cover the literal arrow is
\[
 z_{\rm old}=\epsilon^{4/m}z_{\rm new}.                 \tag{P3K.57.IF.8}
\]
The admitted branches are the existing good III-star root
\((a,b,m)=(1,2,4p)\), and reduced III root
\((a,b,m)=(2,1,4p-1)\), with their actual norm sign \(n=\pm1\).
The comparison of minimal-root coordinates is
\[
 \begin{array}{c|cc}
  m&X'_{\rm new}&Y'_{\rm new}\\ \hline
 4p&x'_{\rm old}&y'_{\rm old}\\
 4p-1&\epsilon^{-2/m}x'_{\rm old}&\epsilon^{-3/m}y'_{\rm old}.
 \end{array}                                           \tag{P3K.57.IF.9}
\]
This follows directly by writing \(X'=z_{\rm new}^{2p}x\) and
\(Y'=z_{\rm new}^{3p}y\) and substituting (P3K.57.IF.8). A full positive
\(\epsilon\)-turn has endpoint \(D_{\rm new}^{-4}\). In the reduced case its
factors on \((z_{\rm new},X',Y')\) are
\((\zeta_m^{-4},\zeta_m^{-2},\zeta_m^{-3})\), since \(4p=m+1\).
In the good case the last two factors are one.

\item The reduced branch retains the actual central nonidentity division section.
On the transported square good-root lattice its chosen logarithm tends to
\((1+i)/2\), and
\[
                   4\log R-\log(4R)=2(1+i).             \tag{P3K.57.IF.10}
\]
The second logarithm is the small identity logarithm. Equation (P3K.57.IF.10)
remains a lattice term in the annular affine cochain. It is not replaced by
zero or multiplied by \(\log\epsilon\) with an assertion of central
holomorphic extension.

\item The actual full reference overlap is fixed by transporting the prescribed
old ordinary-to-cap overlap at one nonzero parameter, including its integral
delta term and finite line character. The K49 relative source/Hom atlas
comparison then continues that full cocycle on the cap annulus. Its finite
averaging retains the central character and its Cauchy splitting retains the
annular integer. These maps identify the two punctured parameter families.
They do not yet attach a central collision model.

\par\medskip\noindent\textbf{IF.4. Half-loop endpoint at the selected cusp.}\par\nopagebreak[4]

\item The source equation is even in \(\epsilon\), while \(G\) changes sign. Its
actual endpoint isomorphism is
\[
                  \iota(x,y)=(x,-y).                    \tag{P3K.57.IF.11}
\]
For \(\Theta(QZ)=-Z^{-1}\Theta(Z)\) and
\(\Theta(Z^{-1})=-Z^{-1}\Theta(Z)\), write
\(f_P(Z)=\Theta(Z/P)/\Theta(Z)\). The exact identity
\[
 f_{P^{-a}}(Z^{-1})=P^a f_{P^a}(Z)
\]
gives the selected-cusp line map
\[
                  (Z,V)\longmapsto(Z^{-1},P^a V).       \tag{P3K.57.IF.12}
\]
Here \(P\) is a bounded unit and \(P(-\epsilon)=P^{-1}\). This map squares
to the identity. It sends \(h\) to \(h_{\rm target}^{-1}\), \(g\) to
\(g_{\rm target}\), and preserves the physical delta circle.
Its selected-cusp valuation map is \(\operatorname{diag}(-1,1)\); it sends
the two specified diagonal fans to one another by that actual map.

\item The cap endpoint of a half \(\epsilon\)-turn is
\(\iota D_{\rm new}^{-2}\). In the reduced case its literal factors from
(P3K.57.IF.9) are \(\zeta_m^{-1}\) on \(X'\) and
\(\exp(-3\pi i/m)\) on \(Y'\); the latter includes the \(y\)-sign of
\(\iota\). Thus the source inversion and cap root rotation are both retained.
The full \(\epsilon\)-turn closes the selected source reflection twice.
In the global rho marking the same arrow is
\(\mathcal R_\rho(Z,V)=(Q/Z,V)\). This is literally
\(h_{\rm target}\circ\mathcal I_\epsilon\), because the target
source period is \((Q,P^{-a})\). Thus the integral height
translation and the factor \(P^a\) in (P3K.57.IF.12) are carried by an
actual period map; neither is dropped from a six-chart closing.

\item The root coordinate used in the other infinity chart is
\(\zeta=\epsilon/\sqrt{1-\epsilon^2}\), with the square root equal to one
at zero. The ratio is a nonzero unit and has zero winding. A positive
\(\epsilon\)-turn covers a negative double \(v\)-turn. The original
\(u^3=v\) endpoint is further discrete data; it is not closed by
discarding this covering degree.

\par\medskip\noindent\textbf{IF.5. The unselected collision and its actual period orders.}\par\nopagebreak[4]

\item At \(\epsilon=0\) the cubic discriminant is \(4t^2(1-t)\). The selected
simple zero remains at \(t=1\); the unselected pair meets at \(t=0\).
The exact other two critical values are the analytic branches
\[
 t_\pm=\frac{\epsilon}{8}
 \left\{\epsilon(4-3\epsilon^2)
          \pm\sqrt{(9\epsilon^2-8)^3}\right\}.           \tag{P3K.57.IF.13}
\]
Choose the square root and the labels once. They have distinct first
derivatives \(\pm2i\sqrt2\) and are exchanged by \(\epsilon\mapsto-\epsilon\).
Set \(q_A=t-t_+\), \(q_B=t-t_-\), with the choice agreeing with the marked
sectional small resolution. In its retained theta/Cartier gauge,
\[
 Q=q_Aq_Bu_Q,\qquad P=q_B^{-1}u_P,\qquad
 c_X=q_Bu_X,                                           \tag{P3K.57.IF.14}
\]
where the units are retained as actual functions. The source and
Cartier derivation is the literal chord/theta calculation reproduced
in IF.5a through the unit Weierstrass coordinate map described below.
It retains the degree-two Cartier divisor and its central component.
The leading coefficients in the transported marking are
\[
               u_Q(0)=1/4096,\qquad u_P(0)=64,\qquad u_X(0)=1/8 .
\]

\item For completeness, a literal source map to the marked local model exists.
Choose either unit root \(\kappa\) of
\[
 \kappa^3-(-2+3\epsilon^2)\kappa-2\epsilon(1-\epsilon^2)
 =(\kappa-\epsilon)(\kappa^2+\epsilon\kappa+2-2\epsilon^2).
\]
Its central value is \(i\sqrt2\) or \(-i\sqrt2\).
Then \(x=\kappa^2 X_0\), \(y=\kappa^3Y_0\),
\(t=\kappa^4T_0\), \(v_0=g/\kappa^3\) sends (P3K.57.IF.1) to
\[
 Y_0^2=X_0^3+(1-2v_0)X_0^2+T_0X_0+v_0^2,\qquad G_0=(0,v_0).
\]
This map also carries the marked small resolution and the canonical
Cartier sections. The whole-zero site is carried to the other, distant
simple discriminant value. It is not changed to one of the colliding sites.
More explicitly, the two source heights are divided by \(\kappa^4\);
the earlier units \(u_Q^0,u_P^0,u_X^0\) become
\(\kappa^{-8}u_Q^0,\kappa^4u_P^0,\kappa^{-2}u_X^0\).
Their previous central values \(1/256,16,-1/4\), together with
\(\kappa^2=-2\), give the displayed constants. This also fixes the signs
of the two node labels by the chosen \(\kappa\) branch.

\item The actual periods for \(\lambda=\epsilon^N\lambda_0\) are
\[
 h=(Q,P^a),\qquad
 g=(P^{-b},\,2(-1)^{b+1}\lambda_0\epsilon^N q_*c_X^{-2}).
                                                               \tag{P3K.57.IF.15}
\]
Because \(q_*(0,0)=-1\), its collision order is zero. Their divisor matrices
at the two strict transforms, the parameter direction, and their sum are
\[
 P_A=\begin{pmatrix}1&0\\0&0\end{pmatrix},\quad
 P_B=\begin{pmatrix}1&b\\-a&-2\end{pmatrix},\quad
 E=\begin{pmatrix}0&0\\0&N\end{pmatrix},\quad
 H_\infty=\begin{pmatrix}2&b\\-a&N-2\end{pmatrix}.
                                                               \tag{P3K.57.IF.16}
\]
In particular \(\det H_\infty=2N-2\). These are period-order identities,
not a count of fixed strata, links, or characteristic numbers.

\par\smallskip\noindent\textit{IF.5a. Exact source, theta and Cartier derivation of the periods.}\par\nopagebreak[4]

\item We give the local derivation used in (P3K.57.IF.14)--\allowbreak (P3K.57.IF.15), so no
unpromoted source-zero construction is needed. Work first in the
unit source coordinates \((X_0,Y_0,T_0,v_0)\) just displayed and
write \(a_2=1-2v_0\). Its discriminant and \(c_4\) give
\[
 j^{-1}=\frac{a_2^2T_0^2-4T_0^3-4a_2^3v_0^2-27v_0^4
                         +18a_2T_0v_0^2}
                   {256(a_2^2-3T_0)^3}.                 \tag{P3K.57.IF.16a}
\]
The convergent inverse of the \(j\)-series therefore gives
\(Q=q_1q_2u_Q^0\), with \(u_Q^0(0)=1/256\), on a
smaller two-node polydisc. Here \(q_1,q_2\) are its two
actual discriminant heights, with the chosen small-resolution
label. All the following unit logarithms are fixed on that
full polydisc.

\item The node map is explicit. Put
\[
 x_c=\frac{-a_2+\sqrt{a_2^2-3T_0}}3,\qquad
 R=(X_0-x_c)\sqrt{a_2+3x_c+X_0-x_c},
\]
using the square roots tending to one. If
\(f(X)=X^3+a_2X^2+T_0X+v_0^2\), then
\(f(X_0)=f(x_c)+R^2\). Thus
\((Y_0-R)(Y_0+R)=f(x_c)\), with
\(f(x_c)=q_1q_2\upsilon\), \(\upsilon(0)=-1/4\).
Dividing one factor by this actual unit gives the literal
node equation \(AB=q_1q_2\). The derivative in \(X_0\)
of \(R\) at the origin is one, so the displayed map is
invertible by the holomorphic inverse theorem.

\item The convergent Tate sums extend over their node chart
\(Q=Z_0W_0\), with leading terms
\(x_T=Z_0+W_0+\cdots\),
\(y_T=-W_0+\cdots\). Rewriting negative-index terms in
\(W_0\) and \(Q\), and positive-index terms in \(Z_0\)
and \(Q\), proves normal convergence there. After the
specified \(c_4\)-root scaling and Weierstrass translation,
this is the invertible node map above. The actual section
\(G_0=(0,v_0)\) has node factors with linear terms
\(v_0+T_0/2\) and \(v_0-T_0/2\). Its two Tate
coordinates therefore vanish simply on opposite
discriminant branches. Analytic division gives
\[
 Z_0(G_0)=q_1u_G,\quad W_0(G_0)=q_2u'_G,\qquad
 P=Z_0(G_0)/Q=q_2^{-1}u_P^0.                         \tag{P3K.57.IF.16b}
\]
The selected nodal normalization gives the exact value
\(P=(1+\sqrt{1+v_0})/(1-\sqrt{1+v_0})\) on
\(q_1=0\). Since \(q_2=-4v_0+O(v_0^2)\) there,
\(u_P^0(0)=16\). This selects the actual period lift
of the section, including its integral reindexing.

\item Write \(f_P(Z)=\Theta(Z/P)/\Theta(Z)\). Its divisor
and exact multiplier \(f_P(QZ)=P f_P(Z)\) identify
it with the canonical meromorphic section of the
actual degree-zero line. The source function \(X_0\)
has divisor \([G_0]+[-G_0]-2[O]\), and consequently
\[
 X_0(Z)=c_X^0 f_P(Z)f_{P^{-1}}(Z),\qquad
                     c_X^0=q_2u_X^0.                   \tag{P3K.57.IF.16c}
\]
The coefficient is fixed by evaluating at the ordinary
Tate point \(Z=-1\). The theta shift identity shows
that \(q_2f_P(-1)f_{P^{-1}}(-1)\) is a nonzero
holomorphic unit; so is the actual \(X_0(-1)\).
At the central nodal cubic the retained coordinate
is \(X_0=4Z/(1-Z)^2\), hence \(c_X^0\sim-4/P\)
and \(u_X^0(0)=-1/4\). The unit map by \(\kappa\)
then gives exactly the three constants and functions
in (P3K.57.IF.14).

\item The primitive source lift covering \(-G\) has target
Cartier multiplier \(X(p-G)^{-1}\). The identity
\(f_P(Z)f_{P^{-1}}(Z/P)=1\) cancels every
\(Z\)-dependent factor in its source and target
Cartier coefficients. On the tensor-power line the
literal resulting map is
\[
          U_a:(Z,V)\longmapsto(P^{-1}Z,c_X^{-a}V).
                                                               \tag{P3K.57.IF.16d}
\]
For \(a=2\), retain the actual Cartier comparison
before using that tensor-power coordinate. With
\(F_2=1-Bz^3-(1-2v_0)z^2\) in the small-resolution
chart of IE.3, the addition formula gives
\[
 r(p)=X_0(p-G_0)=\frac{2v_0F_{21}(p)}{X_0(p)}
                         =\frac{F_2}{z^2}.             \tag{P3K.57.IF.16e}
\]
Its divisor is \([2G_0]+[O]-2[G_0]\). Multiplication
by \(r\) is therefore the actual map
\(\mathcal O([2G_0]-[O])\to
\mathcal O([G_0]-[O])^{\otimes2}\). The last
expression retains its central component and has no
extra factor of \(v_0\). In these unit coordinates the
actual selected height is \(q_0=\kappa^{-4}q_*\).
At \(p'=p-G_0\),
\(F_{21}(p')=r(p')r(p)/(2v_0)\), so its conjugation
of the suspension is the exact identity
\[
 \frac{r(p')}{r(p)}\frac{\lambda q_0}{v_0F_{21}(p')}
                   =\frac{2\lambda q_0}{X_0(p')^2}.
                                                               \tag{P3K.57.IF.16f}
\]
The scalar unit source change is carried on both sides.
This proves (P3K.57.IF.15) for the actual degree-two divisor
line, as well as the degree-one line. It does not
replace that Cartier line by its generic Picard class.

\par\medskip\noindent\textbf{IF.6. The half-loop Cartier arrow through the collision.}\par\nopagebreak[4]

\item The representative \(P\sim q_B^{-1}\) in (P3K.57.IF.14) requires a different
integral Tate reindexing after the two branches are exchanged. The actual
target representative is
\[
                       P'=(QP)^{-1}.                   \tag{P3K.57.IF.17}
\]
It represents \(-G\), has the required \(q_A^{-1}\) leading order, and
retains the integral period omitted by the bounded selected-cusp expression.
Applying the theta shift identity \(a\) times gives
\[
 f_{(P')^a}(Z^{-1})
   =c_a Z^a f_{P^a}(Z),\qquad
 c_a=(-1)^aQ^{-a(a-1)/2}P^{a(1-a)}.                    \tag{P3K.57.IF.18}
\]
Thus the actual collision line map is
\[
              \mathcal I(Z,V)=(Z^{-1},c_a Z^aV).        \tag{P3K.57.IF.19}
\]
It preserves delta and squares literally to the identity:
\(c_a(P)c_a(P')=1\).
Applying the same two theta identities to the divisor of \(x\) gives
\[
                   c_X'=(QP^2)c_X .                    \tag{P3K.57.IF.20}
\]
Consequently the complete source and suspension periods satisfy
\[
 \mathcal I h=h'^{-1}\mathcal I,\qquad
 \mathcal I g=h'^{-b}g'\mathcal I.                      \tag{P3K.57.IF.21}
\]
For the second equality, its source coordinate has factor
\(P^b=Q^{-b}(P')^{-b}\). Its line coordinate has factor
\(P^{-ab}g_V=(P')^{-ab}(QP^2)^{-2}g_V\), since \(ab=2\).
This proves the full line equality including the scalar, rather than only
the generic Picard-point equality.

\item Let \(c=a(a-1)/2\). In the ordered valuation lattice
\((Z,V,q_A,q_B)\), (P3K.57.IF.19) and the branch exchange give
\[
 S=\begin{pmatrix}
 -1&0&0&0\\ a&1&-c&c\\0&0&0&1\\0&0&1&0
 \end{pmatrix},\qquad S^2=1,\qquad S\delta=\delta.       \tag{P3K.57.IF.22}
\]
Its fibre block and period relabelling are
\[
 J=\begin{pmatrix}-1&0\\a&1\end{pmatrix},\qquad
 R=\begin{pmatrix}-1&-b\\0&1\end{pmatrix}.
\]
The exact identities
\[
 JP_A=P_BR,\quad JP_B=P_AR,\quad JE=ER,\quad
 JH_\infty=H_\infty R                                  \tag{P3K.57.IF.23}
\]
prove compatibility with every retained period order. The nonconstant
unit factors in \(c_a\) stay in (P3K.57.IF.18); (P3K.57.IF.22) records only their divisors.
Equations (P3K.57.IF.18)--\allowbreak (P3K.57.IF.21) are the actual morphism relating the
selected-cusp and collision half-loop descriptions.

\par\medskip\noindent\textbf{IF.7. The actual origin and the attachment sequence.}\par\nopagebreak[4]

\item The scalar origin on the ordinary source uses \(\log q_*\). Near the
unselected collision \(q_*\) is a nonzero unit on the full bidisc, so this
logarithm is single valued there. Its \(q_A\) and \(q_B\) continuations do
not insert the extra \(v^{-nk}\) deck found at the selected source-zero
collision. Its parameter endpoint is the extending unit \(q_*^{nk}\).
The finite cap still carries its full norm generator and root arrow
(P3K.57.IF.8)--\allowbreak (P3K.57.IF.10).

\item The following sections construct the rank-one quotient atlas, its actual finite-cap
attachment, and the real selected-cusp comparison, retaining the physical Euler line
and oriented boundary. The analytic collision filling uses the specified positive
scalar rescaling; IE.6 removes that scalar winding by an actual parameter bordism.

\item The executed script
\(\texttt{code/part3\_fourfold/source\_infinity\_exact.py}\) checks the
source equations, discriminant branches, selected units, full scalar
cancellation, all root-coordinate exponents, and the complete half-loop
source/line/period identities. Its JSON retains the explicit cap rows
and has global-bordism and return fields false.

\par\medskip\noindent\textbf{IF.8. Literal rank-one sector fans and protected boundary chains.}\par\nopagebreak[4]

\item All three lower matrices in (P3K.57.IF.16) have rank one. The following is an
actual construction of their cone collections. It is kept in the original
integer lattice, with the physical cocharacter \(\delta=(0,1,0,0)\).
For \(P_0=P_A,P_B,E\), put
\[
 R_0(n)=(P_0n,1,0),\quad R_1(n)=(H_\infty n,0,1),\qquad
 T_j(x,r,s)=(x+(rP_0+sH_\infty)j,r,s).                 \tag{P3K.57.IF.24}
\]
Use the basis \(u=e_1,v=-e_2\) on A, the basis
\(u=e_1,v=e_2-be_1\) on B, and \(u=-e_2,v=e_1\) on E.
In each case \(P_0v=0\). Start with the two labelled
triangles \((0,u,u+v)\), \((0,v,u+v)\). The three
nondegenerate maximal cones in one prism cell are
\[
 \begin{split}
 &\operatorname{cone}(R_0(0),R_1(0),R_1(u),R_1(u+v)),\\
 &\operatorname{cone}(R_0(0),R_1(0),R_1(v),R_1(u+v)),\\
 &\operatorname{cone}(R_0(0),R_0(u),R_1(u),R_1(u+v)).
 \end{split}                                             \tag{P3K.57.IF.24a}
\]
Take all their literal \(T_j\)-translates. The other
three tetrahedra of the two triangular prisms collapse
to the lower-dimensional faces identified below.
In each case \(P_0v=0\). Start with the two triangles
\((0,u,u+v),(0,v,u+v)\) and their three usual staircase tetrahedra
between the \(R_0\) and \(R_1\) copies. Omit their collapsed
four-dimensional cones and retain every ordinary face.

\item Here is a direct common-face and support proof. Set
\(p=P_0u,\ A_1=H_\infty u,\ B_1=H_\infty v\).
The rational inspection map with basis \((A_1,B_1)\) sends \(p\) to
\((c,d)\), with
\[
 \begin{array}{c|cc}
 &c&d\\ \hline
 A&(N-2)/(2N-2)&-a/(2N-2)\\
 B&(N-2)/(2N-2)&-a/(2N-2)\\
 E&N/(N-1)&bN/(2N-2).
 \end{array}                                           \tag{P3K.57.IF.25}
\]
Every \(c\) is positive for the admitted \(N>2\). At fixed \(r,s>0\)
the two top cones give the square \([0,s]^2\), split along its
diagonal. The third surviving cone gives
\[
          (s,0)+[0,r](c,d)+[0,s](0,1).
\]
Their union is the strip
\[
 0\le x\le L=s+rc,\qquad f(x)\le y\le f(x)+s,\quad
 f(x)=0\ (x\le s),\quad f(x)=(d/c)(x-s)\ (x\ge s).
\]
The actual period vectors are \((L,rd)\) and \((0,s)\).
Reducing \(x\) uniquely modulo \(L\), then reducing
\(y-i rd-f(x-iL)\) uniquely modulo \(s\), assigns every point
off a boundary to exactly one of these strip cells.
The square and parallelogram share their entire vertical edge;
the far edge is the next square's near edge; and the upper
piecewise linear boundary is the lower one translated by \((0,s)\).
Thus all intersections are full common faces. The two collapsed
bottom tetrahedra have the already existing face
\(\operatorname{cone}(R_0(0),R_0(u),R_1(u+v))\), and the collapsed
mixed tetrahedron has
\(\operatorname{cone}(R_0(0),R_1(v),R_1(u+v))\).
No image of a collapsed cone is omitted from the face complex.
This proves the fan and support statements with the exact translations
(P3K.57.IF.24), including the rank-one height boundary.

\item On the H boundary the A and E constructions give the negative
unit-square triangulation in \(H_\infty^{-1}\) coordinates; B gives
the shear with basis \((e_1,e_2-be_1)\). Their literal polygonal
overlay is a common rational periodic refinement. Intersections
are common faces because their supporting equalities are those of
the two original complexes. It has finitely many cells modulo
\(\mathbb Z^2\). All rays are then replaced by their primitive
integer vectors on the same ray, preserving their true heights.
For \(b=2\) the added H-edge midpoint ray is the primitive multiple
of \((2,N-2,0,2)\).

\item The A boundary retains every primitive ray \((j,0,1,0)\); the B
boundary retains \((j,-aj,1,0)\). They are exactly the actual source
line chains: on B the literal coordinate \(Y=Z^aV\) has no vanishing
line character along that lower face. The E boundary retains all
primitive rays \((0,j,1,0)\), not only the raw step-\(N\) rays.
Its adjacent coordinates are
\[
       X_j=V\xi_E^{-j},\quad Y_j=\xi_E^{j+1}/V,\qquad
                       X_jY_j=\xi_E .                    \tag{P3K.57.IF.26}
\]
The actual scalar period shifts \(j\) by \(N\), retaining its
nonzero multiplier in the two chart maps.

\item Extend these boundary refinements by coning already subdivided
faces from their primitive ray sum, face by face in increasing
dimension. This gives common faces and commutes with every
unimodular map (P3K.57.IF.24). The relative regularization then uses a
nonzero lattice point in each singular cone's half-open
parallelepiped, subdividing all occurrences of its minimal face
simultaneously. The incident determinants strictly decrease;
the relative separator prevents two translates of that centre
in one cone. The protected boundary faces are already regular
and hence never contain a resolution centre.

\item The executable \(\texttt{source\_infinity\_regular\_fan.py}\) is a
separate saved adaptation of the earlier complete algorithm.
Its derivation file records every changed matrix, rank-one
canonical representative, stabilizer and protected-face dimension.
It exports every actual primitive ray, determinant-one cone,
period, face permutation and higher parameter order. Its legacy
tags I,W,E mean the present A,B,E faces; they do not identify A
with the earlier full-rank I face.

\par\smallskip\noindent\textit{IF.8a. Relative separator and terminating integral resolution.}\par\nopagebreak[4]

\item Here is the complete relative regularization argument for the
present infinity matrices. Protect all the actual A, B and E
lower-face chains listed above. The adjacent primitive rays
have a unit two-dimensional minor; their span lattices are
saturated, with equations respectively \(V=0\),
\(V=-aZ\), and \(Z=0\). Hence every protected cone is
regular in its actual lattice. In each sector the protected
subfan is exactly its lower face \(s=0\). A nonprotected
ray has positive H height. Its period stabilizer is zero:
\[
 \begin{aligned}
 \det(rP_A+sH_\infty)&=\det(rP_B+sH_\infty)\\
                    &=(N-2)rs+(2N-2)s^2,\\
 \det(rE+sH_\infty)&=2Nrs+(2N-2)s^2.
 \end{aligned}                                         \tag{P3K.57.IF.26a}
\]
Both are positive for \(r\ge0,s>0,N>2\). Thus
\(T_\nu(x,r,s)=(x,r,s)\) on such a ray implies
\(\nu=0\). The lower-face kernel stabilizers fix those
faces pointwise; their actual analytic unit actions are
retained in IF.9 and are not replaced by identity maps.

\item First take the relative barycentric subdivision outside
the protected subfan. Process old faces in increasing
dimension, insert the primitive ray on the sum of their
primitive generators, and cone the already subdivided
boundary from that ray. A convex transverse section gives
the literal radial coverage and common-face map. The
same choices are made on the common H faces in all three
sectors. Every period map is unimodular, preserves face
dimension, and takes a generator sum to the corresponding
generator sum; hence this subdivision is exactly equivariant.

\item Each resulting cone consists of a protected face and a
strictly increasing chain of nonprotected old faces. It
therefore contains at most one ray in any nonprotected
period orbit: two rays in an orbit would arise from old
faces of the same dimension, whereas a strict face chain
contains only one of that dimension. An inserted ray is
interior to its named face, so equality of rays also gives
equality of their named faces. Every singular face contains
a nonprotected ray, because a cone generated entirely by
protected rays is a face of one of the regular protected
cones. Two distinct period translates of a singular face
cannot lie in one cone: their nonprotected rays would
violate this orbit separation, and the ray stabilizer is
trivial by (P3K.57.IF.26a).

\item For a singular maximal cone let \(B\) be its primitive
column matrix and \(D=|\det B|>1\). Some column
\(B^{-1}e_j\) is nonintegral, since otherwise \(B^{-1}\)
would be integral and its determinant would force \(D=1\).
Take its coordinatewise fractional part \(\alpha\) and set
\[
 p=B\alpha=e_j-B\lfloor B^{-1}e_j\rfloor,
 \qquad 0\le\alpha_i<1.                              \tag{P3K.57.IF.26b}
\]
This is a nonzero integral point. Divide \(p\) and every
\(\alpha_i\) by the positive gcd of its coordinates, so
\(p\) is primitive and the same strict inequalities hold.
Its minimal face \(F\), generated by the rays with
\(\alpha_i>0\), is singular. Indeed on a regular face
an integral vector with all these coefficients in
\([0,1)\) is zero in the saturated span lattice, which
contradicts \(p\ne0\). Thus \(F\) is outside the
protected subfan and has trivial period stabilizer.

\item Insert \(T_\nu p\) simultaneously on every translated
face \(T_\nu F\). Orbit separation proves that no old
cone receives two distinct centres from this orbit.
For a point \(x=\sum t_i r_i\) of a containing cone,
subtract \(\tau p\), where
\(\tau=\min_{\alpha_i>0}t_i/\alpha_i\). The remaining
coefficients are nonnegative and at least one is zero,
which gives precisely a cone of the star subdivision.
Ties are precisely the shared faces; the same formula
on an old shared face proves compatibility. For each
affected maximal cone and each replaced ray,
\[
                        D_{
m new}=\alpha_i D_{
m old},
                 \qquad 0<\alpha_i<1.                  \tag{P3K.57.IF.26c}
\]
These are positive integers. Every affected determinant
strictly decreases, and no protected cone is subdivided.

\item Orbit separation survives the star: a new cone lies in
an old cone and contains at most one new ray of the chosen
orbit. This new orbit was not an old ray orbit, since its
point was interior to a face of dimension at least two.
Its positive H height gives trivial stabilizer again.
There are initially finitely many cone orbits, and every
star introduces finitely many new ones. At each stage
choose a maximal-determinant cone orbit. The pair
\((D_{\max},\#\{\text{maximal orbits of determinant }D_{\max}\})\)
strictly decreases lexicographically: no new orbit has
the old maximal determinant. There are finitely many
stages at each first entry and finitely many decreases
of that entry. Hence the process terminates with every
maximal determinant one. If the centre lies on H, perform
the identical star in all incident sectors; formula
(P3K.57.IF.26c) applies to each. This proves a compatible
period-equivariant regular refinement relative to all
three actual chains, in the original integral lattice,
with every higher parameter height retained.

\par\medskip\noindent\textbf{IF.9. Exact face kernels and the analytic quotient.}\par\nopagebreak[4]

\item Put \(C=2(-1)^{b+1}\lambda_0\). On the A face the kernel period is
\(g\). Its transverse line coordinate is \(V\), and its multiplier is
\[
                  C\epsilon^N q_*q_B^{-2}u_X^{-2}.       \tag{P3K.57.IF.27}
\]
On the B face the kernel is \(k=g h^{-b}\). Its source multiplier
and its multiplier on the literal transverse coordinate \(Y=Z^aV\)
are
\[
 (q_Au_Qu_P)^{-b},\qquad
          C\epsilon^N q_*q_A^{-2}u_Q^{-2}u_P^{-4}u_X^{-2}.       \tag{P3K.57.IF.28}
\]
These equalities follow directly from (P3K.57.IF.15): the second line is
\(g_V/(QP^2)^2\). Since the two critical values differ by
\(\epsilon\) times a nonzero unit, both raw transverse kernel
multipliers have order \(N-2\) at their corresponding strict face.
All constants and source multipliers remain in (P3K.57.IF.27)--\allowbreak (P3K.57.IF.28).
They are not replaced by a formal valuation.

\item The raw multiplier order \(N-2\) must be combined with the actual nodal-line metric and source reindexing. The complete operator bounds below retain losses three and four, respectively, and the geometric proof uses the safe orders \(N=8\) and \(N=6\).

\item At E the source period \(h\) retains the actual compact elliptic
source fibre and its line holonomy. The scalar period has order \(N\)
and acts on the primitive chain (P3K.57.IF.26). Near H, both previously
kernel directions have nonzero image under the full-rank matrix
\(H_\infty\). Thus the period separation in the full-rank cones
controls the closures of both face norm shells; no additional
kernel is lost at that overlap.

\item Here is the complete analytic quotient argument, including the
norm bound for the actual nodal lines. The analytic charts are
the affine semigroup charts of the regular fan, pulled back by
the stated base-blowup maps. The full period indexed by
\(\nu\in\mathbb Z^2\) has lattice map \(T_\nu\) from
(P3K.57.IF.24), followed by the actual extending base units from
(P3K.57.IF.15). Multiplication by these units acts on every character
by its literal integer powers, and therefore extends to the
corresponding semigroup chart. The units form a homomorphism
in \(\nu\), since they are products of the two commuting
periods. This proves the group law, its inverse, and the
common-H overlap law on the analytic atlas.

\item For a nonzero base value, let \(r,s\) be the positive logarithmic
radii of the two sector coordinates. The full real matrix of
the logarithmic magnitudes of its two actual periods is
\[
                       L=-rP_0-sH_\infty+C,             \tag{P3K.57.IF.28a}
\]
where \(C\) is bounded; the fixed \(\lambda_0\) belongs
to that bound. By (P3K.57.IF.26a) the determinant of the leading
matrix is bounded below by a positive multiple of
\(s(r+s)\), and its norm is bounded above by a multiple
of \(r+s\). Its least singular value is therefore at
least a positive multiple of \(s\). Shrinking the corner
polydisc makes the bounded perturbation harmless. Use the
actual parallelogram \(L([0,1]^2)\), with its compact
angular torus. It differs from the leading parallelogram
by a uniformly bounded set.

\item The strip construction in IF.8 places that set in finitely
many adjacent chart types, uniformly in \(r/s\). Indeed
its width is \(s+rc\), its height is \(s\), and its
periods are \((s+rc,rd)\), \((0,s)\). Compare the
linear function \(rd\,x/(s+rc)\) with the two pieces
of \(f(x)\). Their difference divided by \(s\) is
bounded uniformly, so only a bounded number of vertical
translates is required. A bounded perturbation changes
both assigned strip integers by a bounded amount.
Consequently one fixed finite collection of adjacent
cone charts covers the quotient. Its monomial coordinate
magnitudes are uniformly bounded above and its angles
lie in compact tori. Closing those regions, including
their toric faces, gives compact chart sets over a
closed smaller base polydisc. Away from the three
marked H points the same argument uses the full-rank
matrix \(H_\infty\) alone.

\item At the strict B face, restrict to
\(|q_B|\le\eta|q_A|\), with a fixed small \(\eta\).
Then \(|\epsilon|\) is comparable with \(|q_A|\).
The nodal line in the character \(Y=Z^aV\) has source
period multiplier \(c=(QP)^a=(q_Au_Qu_P)^a\).
On \(|Q|\le|Z|\le1\), interpolate a positive metric
between its two actual endpoint values. Its ratio is
bounded by \(C|q_A|^{-a}\). The same interpolation on
the plumbing annuli extends to the nodal normalization
with the prescribed gluing at both node preimages.
The source multiplier of the kernel is bounded by
\(C|q_A|^{-b}\). Reindexing its image to the Tate
fundamental annulus uses at most two source periods
for \((a,b)=(1,2)\), and at most one for \((2,1)\).
Each period changes the line norm by at most
\(C|q_A|^{-a}\). Combining these bounds with the
raw multiplier of order \(N-2\) gives
\[
 \|k\|\le
 \begin{cases}
 C|\lambda_0|\,|\epsilon|^{N-5},&(a,b)=(1,2),\\
 C|\lambda_0|\,|\epsilon|^{N-6},&(a,b)=(2,1).
 \end{cases}                                            \tag{P3K.57.IF.28b}
\]
At A restrict instead to \(|q_A|\le\eta|q_B|\).
Here \(|\epsilon|\) is comparable with \(|q_B|\),
the literal line character is \(V\), and its source
period multiplier is \(P^a\), of norm comparable
with \(|q_B|^{-a}\). Choose its endpoint metric
with the prescribed ratio \(|P|^a\), and interpolate
on the source annulus and plumbing charts. Its
maximum-to-minimum norm ratio is at most
\(C|q_B|^{-a}\), including both preimages of the
nodal point. The source part of \(g\) is
\(P^{-b}\), of norm comparable with \(|q_B|^b\).
Since \(|Q|\le C|q_B|^2\), returning the image of
\(|Q|\le|Z|\le1\) to that same annulus uses at most
two source periods for \(b=2\) and at most one for
\(b=1\), after fixing \(\eta\) and shrinking.
Each positive or negative reindex changes the line
norm by at most \(C|q_B|^{-a}\). The raw multiplier
has order \(N-2\); adding this metric loss and those
reindex losses gives exactly (P3K.57.IF.28b) for \(g\).

\item These estimates include all their face overlaps.
For fixed \(\eta<1\), the strict A and B regions
are disjoint; the intervening sector is covered by
the full-rank H parallelograms already constructed.
On each strict face its middle annulus and two node
charts use the same endpoint metric and transition
multiplier, so no further norm ratio is introduced
on their overlaps. Away from H the relevant
source/line maps lie in compact families and have
a common finite norm bound. At E the kernel is
the actual source Tate period, with its complete
elliptic fundamental annulus and its line holonomy.
The common H charts retain the nonzero images of
all these kernel vectors, which controls the
closures of their norm shells as described below.

\item We use the analytic construction at the two executed
orders \(N=8\) and \(N=6\), respectively. In the first
case (P3K.57.IF.28b) tends to zero; in the second it is
uniformly less than one after choosing \(|\lambda_0|\)
small. For an arbitrary admitted cap, take a sufficiently
large multiple of its order; IE.6 compares these orders
by an actual scalar pair of pants. The additional
\(N=3,4\) fan executions are finite arithmetic checks,
not an assertion of this norm bound at those orders.
The actual line-norm shells prove freeness and proper
discontinuity of the retained kernels on the nonzero
nodal lines. On E use the actual Tate shell for
\(Z\mapsto QZ\) together with the primitive chain.
Near H the two kernel vectors have nonzero H images,
so an unbounded sequence of their translates cannot
accumulate in a fixed H-cone chart. The full-rank
separation controls the closures of the face shells.

\item Finally a nonzero period cannot stabilize a bounded
cell on a full-rank face, and only finitely many
translates of a compact chart meet it. On a rank-one
face the only possible stabilizer is the specified
kernel, whose strict norm contraction has the same
finite-intersection property on a compact shell.
The equivalence relation is locally a finite union
of closed graphs on these compact covers. Its quotient
is Hausdorff, and inverse images of compact smaller
base sets are compact. On A and B its boundary charts
are precisely the actual Tate source chain and
Cartier line with their retained kernel maps; on E
they are (P3K.57.IF.26). Thus the punctured-parameter part
is the prescribed family under the displayed theta
and Cartier maps, with the stated boundary charts.

\item The global E-chain monodromy still carries its actual source-base
data. A small loop around the unselected collision has zero winding
of \(q_*\), so it fixes all \(N\) chain labels. A loop around the
separate selected value shifts labels by \(-nk\), and the inverse
cap cut at infinity shifts them by \(+nk\), for \(N=km\).
The latter two have \(k\) cycles of length \(m\). This is a
compatible list of branch locations; it does not itself construct
the missing selected-cusp attachment.

\item The real joining of this partial filling to the selected scalar trace is
constructed in IE.1--6, with its physical Euler line and normal orientation.

\par\medskip\noindent\textbf{IF.10. The changed-origin partial attachment away from the selected disc.}\par\nopagebreak[4]

\item Fix disjoint source-base discs around the unselected collision at \(t=0\),
the selected value \(t=t_*(\epsilon)\), and the cap at infinity. The first
and last closures are disjoint from the selected disc for all sufficiently
small \(\epsilon\). Put \(N=km\), retain the actual \(n=\pm1\), and use
\[
 \tau=\frac{\log\epsilon}{2\pi i},\quad
 s_*=\frac{\log q_*}{2\pi i},\quad
 s_\infty=\frac{\log q_\infty}{2\pi i},\quad
 h=1-t_*(\epsilon)q_\infty,\qquad
 s_*=-s_\infty+\frac{\log h}{2\pi i}.                   \tag{P3K.57.IF.29}
\]
On the unselected collision bidisc \(q_*\) is a nonzero unit, so \(s_*\)
has one holomorphic branch on that full bidisc. On the cap polydisc \(h\)
is likewise a nonzero unit, with its chosen reference logarithm.

\item On the universal parameter cover the actual ordinary-to-cap map and
the actual positive parameter-deck square are
\[
 F_\tau=F_0+nk\,\tau s_\infty\delta,\qquad
 R_{\rm ord}=q_*^{nk},\quad R_{\rm cap}=h^{nk},\qquad
 F_{\tau+1}R_{\rm ord}=R_{\rm cap}F_\tau.               \tag{P3K.57.IF.30}
\]
These are the admitted cap's norm and cochain formulas, with the current
source and the transported full \(F_0\) of IF.3. The source and line
characters are not changed. Define
\[
 o_{\rm ord}=nk\,\tau s_*\delta,\qquad
 o_{\rm cap}=nk\,\tau\frac{\log h}{2\pi i}\delta,\qquad
 \eta=\zeta-o_{\rm ord}.
\]
The exact overlap is
\[
                F_\tau+o_{\rm ord}-o_{\rm cap}=F_0.    \tag{P3K.57.IF.31}
\]
After a positive parameter turn the old ordinary chart is identified
by \(R_{\rm ord}\). Its change is exactly the change of \(o_{\rm ord}\);
therefore \(\eta\) has identity parameter deck. The cap coordinates
have the same property. This retains the closing arrow: neither
\(R_{\rm ord}\) nor \(R_{\rm cap}\) is silently set equal to one in
the old chart.

\item Every source period and line lift commutes with these base-only scalar
translations. Hence in the new ordinary coordinate the full periods
are exactly (P3K.57.IF.15), including all their units. Moreover positive
\(q_A\) or \(q_B\) continuation in the collision bidisc leaves \(s_*\)
unchanged. Thus the new collision coordinate has no additional
source-base chain deck there. The fan of IF.8 may therefore be
constructed directly in \(\eta\), and its quotient of IF.9 has the
pure parameter deck just proved. On its punctured parameter part
this is an actual isomorphism to the prescribed family, using
(P3K.57.IF.30), not an assertion that a fractional scalar power extends in
the old \(V\) coordinate.

\item The cap chain is constructed in \(\zeta_{\rm cap}-o_{\rm cap}\).
The actual primitive generator is
\[
 H=A^uB^w=\epsilon^k r_0H_0,\qquad nu+mw=1,
 \qquad H^m=B,\quad H^n=A ,
\]
with the actual admitted source and line maps. The displayed Bezout
equation is the literal exponent identity for \(A=H^n,B=H^m\).
Thus its node chart shift is
\(k\), and \(A\) has signed shift \(nk\). The source and finite
root maps in \(H_0\) are retained. The chain equations and all line
frame transitions are (P3K.57.IF.26), with \(H\) acting by that shift and
by its actual source/line map. A nonzero iterate changes the index,
so the action is free and properly discontinuous also at
\(\epsilon=0\).

\item On the cap annulus positive \(q_\infty\) continuation changes
\(s_*\) by \(-1\); the new ordinary vertical coordinate consequently
has the factor \(\epsilon^{nk}\). This is exactly \(A\)'s chain
shift and its retained \(A_0\) arrow. The cap unit \(h^{nk}\) acts
on a node chart by
\((X_j,Y_j)\mapsto(h^{nk}X_j,h^{-nk}Y_j)\); it is an extending
unit even though its central value is now \(1-q_\infty\), rather
than the constant one of the source-zero calculation. Its origin
change has been retained in (P3K.57.IF.30)--\allowbreak (P3K.57.IF.31).

\item On each remaining ordinary source-base chart, away from the selected
disc and these two special pieces, the source and Cartier line are
smooth over the full parameter disc. Use their full chain with the
actual source/line period. On angular overlaps the affine cochains
are exactly those already prescribed by \(F_0\); the only nonzero
integral base continuation is the signed cap shift just displayed.
The collision overlap has zero such continuation. Consequently
the ordinary chains, the collision fan in \(\eta\), and the cap
chain glue by literal extending source/line maps. This constructs
the partial family over the complement of the selected source-base
disc.

\item The relative regularization is by primitive determinant-one cones,
and the protected face charts are the actual smooth source/chain
charts. The total complex fourfold is therefore smooth. The map
to the epsilon disc may have a singular central fibre; smoothness
here refers to the total space and does not replace that map by a
submersion.

\item Properness over a closed smaller parameter disc follows from the
finite compact fundamental covers already described: collision
cones and face shells in IF.9, finitely many node/middle chain
charts on every ordinary compact source patch, and one finite
fundamental chain segment for \(H\) on the cap. A finite source-base
cover suffices. Every overlap is the actual biholomorphic map
(P3K.57.IF.31), so these compact covers produce a compact inverse image.
The only retained real boundary is over the boundary circle of the
deleted selected disc.

\item This partial attachment supplies the exact source-base boundary used in
IE.1--6. Those sections construct its real comparison with the selected
scalar trace, with the physical Euler line and its normal orientation.

\par\medskip\noindent\textbf{IF.11. Executed local fixed-stratum data.}\par\nopagebreak[4]

\item The two independently saved regular fans have respectively
\[
 \begin{array}{c|r|r|rrr}
 (a,b,N)&\text{checks}&\text{stars}&A&B&E\\ \hline
 (1,2,8)&220943&1289&1838&1838&1818\\
 (2,1,6)&222609&1381&1748&1768&2021
 \end{array}
\]
maximal cone orbits. Every exported ordered cone has determinant \(+1\).
The fixed components are determined by the minimal face containing
the actual \(\delta\), with every recorded period and vertex
permutation used. In particular neither rank-one A nor rank-one B
contains a fixed component on its protected boundary.

\item The compact interior and H-normal calculations use the local
normal-link convention and formula (P3K.57.IE.16). On H the full unit
factors, with \(w=q_B/q_A\), are
\[
 h:(Z,V)\mapsto(wZ,w^{-a}V),\qquad
 g:(Z,V)\mapsto(w^bZ,w^{-2}(1-w)^NV)
\]
after the \(H_\infty\) normal powers are separated; the nonzero
constants remain part of the atlas. Thus the three normal-degree
orders at \(w=0,1,\infty\) are the same literal expressions
(P3K.57.IE.22), with the actual \(H_\infty\) used when unwrapping
the recorded periods.

\item The executed values, verified component by component using two
independent auxiliary circles, are
\[
 \begin{array}{c|rrr}
 (a,b,N)&\text{compact interior}&H&\text{sum}\\ \hline
 (1,2,8)&9333/16&-9301/16&2\\
 (2,1,6)&739&-734&5 .
 \end{array}                                           \tag{P3K.57.IF.32}
\]
All H surface tangent Euler numbers, canonical squares,
signatures and holomorphic Euler numbers satisfy their exact
integrality checks.

\item Each E fixed surface has the actual unmodified local I2 cycle
with intersection matrix \(\left(\begin{smallmatrix}-2&2\\2&-2\end{smallmatrix}\right)\).
The relative determinant marking \(Z^aV^2/\xi_E^{2j+1}\)
has divisor zero on both components, as the unwrapped cone
characters show. The local tangent marking
\(\Omega=d\log Z\wedge dt\) gives relative
\(p_1=-4\). Indeed the actual source-period change adds
\(d\log Q(t)\wedge dt=0\), so \(\Omega\) descends.
On a node chart \(xy=t\), it is \(dx\wedge dy\).
The vertical vector is \(x\partial_x-y\partial_y\);
it has one simple complex zero at each of the two nodes.
Together with the boundary vector normalized by
\(\Omega(W,H)=1\), this gives relative \(c_1=0\)
and \(c_2=2\), hence \(p_1=c_1^2-2c_2=-4\).
Thus the annular normal-link residue is \(-1/6\) per
E surface. These local numbers are retained with their
boundary frames. IE.3 constructs the actual E double,
and IE.4--6 assemble its completed-family contribution.

\item All programs and their full exact outputs are saved with the prefix
\path{code/part3_fourfold/source_infinity_}.
The source-family return uses the complete IE.1--7 comparison,
including its actual boundary frames and parameter orientation.

\par\medskip\noindent\textbf{IF.12. Full old-v / new-epsilon annular family comparison.}\par\nopagebreak[4]

\item Fix an integer root-scalar winding \(j\), and use
\[
 r(\epsilon)=\epsilon^j u(\epsilon)r_0,\qquad
 \lambda=r^m ,
\]
where the specified nonzero unit \(u\) has an extending
logarithm. The comparison in this section is over the
punctured epsilon annulus carrying this actual path.
It does not assume a trivial Isom torsor over a product
of two independent punctured parameter discs.

\item The source map (P3K.57.IF.1), the canonical divisor-line map,
and (P3K.57.IF.6)--\allowbreak (P3K.57.IF.7) identify the full ordinary suspension.
The selected Tate \(Q,P\) agree under this source
isomorphism. Their physical heights differ by
\(q_{\rm old}=\epsilon^{-4}q_*\), an actual single-valued
integer-power unit on this annulus. Its powers act on
the toric character charts and preserve the chosen
periodic fan and physical delta. Its winding is retained.

\item On each simply connected parameter chart use the actual
root map (P3K.57.IF.8)--\allowbreak (P3K.57.IF.9). It carries the marked source
section, the resolution ideals and the corrected Cartier
line. The finite source/line map has the form \(A=LD\),
where \(L\) covers \(-nR\). Its central commutator
character is one: in the good branch \(R=O\) centrally
and \(L\) is a scalar; in the reduced branch the
corrected root line is trivial on the compact reduced
III fibre, so the multipliers of \(D,L\) are constants,
and their source maps commute. The finite averaging
conjugation of the actual norm-one unit gives
\[
 DLD^{-1}=L,\qquad A^4D^{-4}=L^4 .
\]
The source part of \(L^4\) specializes to the identity
in both branches. Divide its central scalar and take
its small extending Lie logarithm. The actual annular
difference from four times the logarithm of \(L\)
retains the period (P3K.57.IF.10), with its full line factor
and signed \(n\). No multiple of that period is
asserted to be an extending reduced-root logarithm.

\item Choose one reference \(\epsilon_*\) and transport the
entire old ordinary-to-cap overlap through these
source/root/Cartier arrows, including \(d\delta\),
the chosen root cut, norm and finite character.
Use this full map as the reference of the new
\(F_0\). On a fixed new cap annulus the source and
root families are smooth through epsilon zero.
Restriction to a chosen root sheet identifies
the ordinary line quotient with the cap quotient:
an element \(A^hB^\ell\) fixing that sheet has
\(h=mk\), so is exactly \(B^{nk+\ell}\).
Thus the local comparisons are literal source
and line maps, and their sheet changes are the
retained powers of \(A\).

\item Correct those maps at \(\epsilon_*\) to the full
reference overlap. On two angular base charts
their overlap translations have Lie logarithms
zero at that reference. Lift along the simply
connected epsilon disc. Every ordinary and
finite-group discrepancy is an integral period
and is zero at the reference, hence stays zero.
The fixed-contour Cauchy splitting in the base
variable therefore glues the actual local maps
to \(F_0\), holomorphically in epsilon. This
retains its full annular integer and the
doubled reduced period. It is the explicit
Hom/source-period torsor construction of
P3K.49, RCZ.5, with the present displayed root maps
as its input.

\item Now compare the two completed families on a
small simply connected parameter chart.
Their total ordinary/root/cusp lattice
two-cocycles agree under the actual source
period identification: they agree at the
reference, and their integral discrepancies
are continuous in the parameter. The reduced
component translations are first removed by
their actual suspension powers. The remaining
Lie bundle has the literal coarse-base exact
sequence
\[
     0\longrightarrow\mathcal O\delta
       \longrightarrow V\longrightarrow\mathcal O(-1)
       \longrightarrow0 .
\]
The source dualizing vector and physical delta
give its relative frames. Split the extension
and its one-cocycles by the fixed-circle
Cauchy/Laurent operators for \(\mathcal O(1)\),
\(\mathcal O\), and \(\mathcal O(-1)\).
Finite averaging preserves the full root
cocycles. Exponentiation then gives actual
local family isomorphisms extending on the
chosen cusp and root charts.

\item On two parameter charts their ratio is a
global extending translation. The actual
source/rational-divisor endpoint argument
of P3K.49, RCZ.4 identifies that group with the
physical \(\mathcal O^*\), including all
finite-kernel choices. To glue the local
maps, take logarithms \(\ell\) of their
scalar ratio on the two angular overlap
components, choose \(\chi_++\chi_-=1\),
and set \(b_+=\chi_-\ell,b_-=-\chi_+\ell\).
Their common \(\bar\partial\) is solved on
a smaller annulus by
\[
 v(\epsilon)=\frac1\pi\int_{\mathbb C}
       \frac{\widetilde f(w)}{\epsilon-w}\,dA(w).
\]
Then \(h_\pm=\exp(b_\pm-v)\) have ratio
equal to the actual scalar transition,
and \(h_\pm^{-1}I_\pm\) glue literally.
The cutoff is supported in the larger
annulus and equals one on the loop.
This construction retains both overlap
components and any integral scalar winding.

\item There is a concrete forward-endpoint
factor before this last scalar splitting.
For the root winding \(j\), the old and
new ordinary forward decks have ratio
\[
      q_{\rm old}^{nj}/q_*^{nj}
                         =\epsilon^{-4nj}.
\]
With \(\tau=\log\epsilon/(2\pi i)\), retain
the actual compensating cochain
\[
 c_0(\tau)=\exp\{2\pi i\,2nj\,\tau(\tau-1)\},
 \qquad
 c_0(\tau+1)/c_0(\tau)=\epsilon^{4nj}.                 \tag{P3K.57.IF.33}
\]
This is not declared single valued by
itself. It belongs to the forward
parameter-deck square, alongside the
\(D^{-4}\) root endpoint and the
retained period above.

\item All constructed maps and inverses are
source/Cartier coordinate maps,
extending translations on their stated
charts, or physical scalars. Thus they
give a holomorphic delta-equivariant
family isomorphism on the epsilon
annulus, carrying the actual line,
fan, normal circle and closing arrows.
A connection transported through
this family map identifies the physical
returns by conjugacy; there is no
loop reversal or numerical class division.

\item For the old global scalar
\(r=\epsilon^{-2k}(1-\epsilon^2)^kr_0\),
one has \(j=-2k\). Formula (P3K.57.IF.33) is
\(\exp\{-2\pi i\,4nk\,\tau(\tau-1)\}\),
whose one-step ratio is
\(\epsilon^{-8nk}\). The nonzero unit
\((1-\epsilon^2)^k\) is retained and
can be homotoped to its reference value
through its extending logarithm.
The local rescaling
\(r=\epsilon^{k_\infty}r_\infty\)
is instead \(j=k_\infty\). Their
root-winding difference is
\(-2k-k_\infty\), as in RG.5.
These are two explicitly compared
loops, not two notations for one scalar.

\par\medskip\noindent\textbf{IE.1. Two actual collision fillings and their boundary families.}\par\nopagebreak[4]

\item Choose a small closed source-base disc \(U\) around \(t=0\), disjoint
from the selected value and the cap. Its boundary annulus contains no
singular source fibre for any small \(\epsilon\). Fix a nonzero
reference \(\epsilon_*\). Let \(S_\epsilon\) be the actual minimal
source over \(U\), with its marked small resolution at \(\epsilon=0\),
and retain \(M_a=\mathcal O(aG-O)\), \(a=1,2\).

\item The first filling \(W_{\rm c}(N)\) is the regularized collision family
IF.8--10, with \(\lambda=\epsilon^N\lambda_0\). After its
determinant-one regularization the total fourfold is smooth; the
function \(\epsilon\) may have a singular central fibre. Its actual
source/line periods are (P3K.57.IF.15), its physical circle is \(\delta\),
and its protected boundary is the primitive source-line/parameter
chain.

\item The reference filling \(W_{\rm r}(N)\) uses the fixed source
\(S_{\epsilon_*}\), its fixed actual line \(M_a\), and the same scalar
\(\lambda=\epsilon^N\lambda_0\). No selected zero lies in \(U\).
Its actual suspension lift \(B_0\) is therefore invertible over
\(U\), including both ordinary I1 fibres. Use the full chain
\[
 X_j=V\epsilon^{-j},\qquad
 Y_j=\epsilon^{j+1}/V,\qquad X_jY_j=\epsilon ,
 \quad j\in\mathbb Z.                                  \tag{P3K.57.IE.1}
\]
The suspension carries chart \(j\) to \(j+N\), with \(B_0\) on
its source and line frames. A nonzero iterate changes the chain
index; a segment of \(N\) node/middle charts is a compact
fundamental cover over the compact source patch. Thus
\(W_{\rm r}(N)\) is a smooth fourfold with exactly \(N\) fixed
source surfaces in its central chain, of normal weights \(+1,-1\).

\item On the boundary annulus the varying source, Cartier line,
translation and line lift all extend over the full
\(\epsilon\)-disc. They are the full source/line maps, not just
their limiting Picard points. The changed-origin coordinate
of IF.10 has pure parameter deck, and its source-base
continuation on this annulus is zero because \(q_*\) is a unit
on \(U\).

\item Here is the actual smooth boundary transport. Introduce
\(s\in[0,1]\), use source parameter
\(\epsilon_s=(1-s)\epsilon+s\epsilon_*\), and keep the
degeneration coordinate \(\epsilon\) and scalar
\(\epsilon^N\lambda_0\) fixed. The boundary source fibres
stay smooth. Continue their actual line and suspension
using the source/Cartier expressions. On each chain chart
the equation is always \(XY=\epsilon\). There are local
lifts of \(\partial_s\) keeping \(\epsilon\) and the
source-base boundary coordinate fixed: take the
source-family chart lift and its actual line frame,
with zero \(X,Y\) derivatives. Under a unit frame change,
carry its derivative into the two opposite normal
coordinates. These give lifts on the actual quotient
charts.

\item Use a finite partition of unity on the compact quotient
collar to combine the lifts, then average over the physical
circle. The resulting field still has \(ds=1\),
\(d\epsilon=0\), and fixed source-base boundary coordinate.
Its compact fundamental cover makes the flow complete for
\(0\le s\le1\). Its time-one map is an actual equivariant
boundary diffeomorphism
\[
        \Psi:\partial_U W_{\rm c}(N)
                   \longrightarrow\partial_U W_{\rm r}(N).   \tag{P3K.57.IE.2}
\]
It retains the full source/line units and \(j\)-label
transport. The flow is in the chain atlas, not a
translation logarithm across the interior collision.

\item Glue the two fillings along (P3K.57.IE.2), reversing the reference
orientation and rounding the common corners. The result
\(W_{\rm rel}(N)\) is a compact oriented smooth
eight-dimensional delta-manifold. Its outer boundary is
the closed relative mapping-torus piece. Its fixed set
consists of the compact interior/H components of
\(W_{\rm c}\), and \(N\) closed doubles of its E patches
with the oppositely oriented reference patches. Their
normal bundles are glued by the actual derivative of
\(\Psi\).

\par\medskip\noindent\textbf{IE.2. Exterior transport and selected-cusp cancellation.}\par\nopagebreak[4]

\item The same source deformation extends on the complement
of \(U\), including the actual cap root. Its division
section, finite character, norm lift and full annular
integral cochain are those of IF.3 and IF.10. On the
selected I1 collar the actual normalization is (P3K.57.IF.2);
the Tate, Picard and Hom units (P3K.57.IF.3)--\allowbreak (P3K.57.IF.5) have
extending logarithms and zero source-parameter winding.
Its Cox products continue through \(\epsilon_s\), with
the same scalar parameter and physical circle.

\item On the outer circle the completed threefolds are smooth
and proper. Lift \(\partial_s\) as above, prescribing
(P3K.57.IE.2) on a collar of \(\partial U\). Extend the lift
using a partition of unity and circle averaging.
Near the selected fixed sphere prescribe the literal
Cox-coordinate unit map. Its normal derivative extends
through \(s\) and the source-parameter disc, so it
transports the actual Hopf Euler line and normal
framing. The resulting exterior diffeomorphism agrees
with (P3K.57.IE.2).

\item The two selected fixed-sphere traces and their normal
products are thereby identified, including the positive
physical orbit. Remove the product tubes and perform
the same conifold replacement on both sides. The
identification extends over \(D^3\times S^2\), with
its transported meridian and Hopf Euler line. The
remaining exterior is a product bordism. Therefore
the difference of the two completed Jin six-manifolds
is the relative collision boundary of \(W_{\rm rel}\),
with the parameter-convention sign in IE.5.
This cancels actual selected scalar traces; it does
not supply a holomorphic full selected corner.

\par\medskip\noindent\textbf{IE.3. The E doubles and actual Cartier class.}\par\nopagebreak[4]

\item The marked source small resolution in the unit
Weierstrass coordinates of IF.5 is
\[
 A=Dz,\quad C=Bz,\quad X=Bz,\quad
 Y=(B+Dz)/2,\quad v_0=(B-Dz)/2 .
\]
Here \(A=Y-v_0\), \(B=Y+v_0\), \(C=X\), and
\(D=T_0+X^2+(1-2v_0)X\). The complementary chart is
\[
 D=Aw,\quad B=Cw,\quad w=z^{-1},\quad
 v_0=(Cw-A)/2.
\]
Thus \(dv_0/dB=1/2\) in the first chart and
\(dv_0/dA=-1/2\) in the second. These are actual
submersion charts on both sides of the exceptional
curve. Choose the source-base boundary outside the
two moving critical values. Its restriction is then
a proper submersion as well; the boundary transports
used below are transports of this resolved family.
At \(v_0=0\) the I2 has components \(C_0,C_1\);
the marked section meets \(C_1\), and \(O\) meets
\(C_0\). The actual degree-two section closure has
Cartier equation
\[
                  1-Bz^3-(1-2v_0)z^2=0 .
\]
To see this directly, the tangent slope at
\(G=(0,v_0)\) is \(T_0/(2v_0)\), so the actual
double has \(Y+v_0=-T_0X/(2v_0)\). In the first
resolution chart this gives \(T_0=-2v_0/z\).
Substitute \(2v_0=B-Dz\), \(X=Bz\) and
\(D=T_0+X^2+(1-2v_0)X\); canceling on the dense
ordinary chart gives the displayed Cartier equation,
which defines its closure. At \(v_0=0\), one has
\(B=Dz\) and
\(T_0=D(1-z^2-Dz^4)\). The two central components
are therefore \(C_1=\{D=0\}\) and
\(C_0=\{1-z^2-Dz^4=0\}\), with the latter exactly
the displayed double-section closure.
It restricts to \(C_0\) with multiplicity one and
does not contain \(C_1\). Consequently
\[
 c_1(M_1)[C_1]=1,\qquad c_1(M_2)[C_1]=2.               \tag{P3K.57.IE.3}
\]
The second equality retains the actual vertical
Cartier specialization \(2G-O=C_0\).

\item The local source patch retracts onto its I2
plumbing: use the radial collapse of \(xy=t\)
to the axes in the two node charts and the
component normal-disc retractions elsewhere.
Its integral two-cycle lattice has primitive
basis \(F=C_0+C_1,C_1\), with pairing
\[
                  \begin{pmatrix}0&0\\0&-2\end{pmatrix}.       \tag{P3K.57.IE.4}
\]
The normalized component two-cells give integral
\(C_0,C_1\), and the zero-section relative disc
meets \(F\) once. The smooth source resolution
is a submersion over \(v_0\), also with boundary,
so this marking transports to \(S_{\epsilon_*}\).

\item On the reference chain (P3K.57.IE.1) the positive normal
line is literally \(M_a\), and the negative one
is \(M_a^{-1}\). Conjugating the negative-weight
line gives determinant
\[
                          b_{\rm r}=2c_1(M_a).          \tag{P3K.57.IE.5}
\]
Indeed the actual fibre-linear change \(V\mapsto gV\)
on \(M_a^\times\) induces
\((X_j,Y_j)\mapsto(gX_j,g^{-1}Y_j)\) in the chain.
After conjugating the negative-weight coordinate,
the determinant transition is
\(g/\overline g=g^2/|g|^2\). A Hermitian metric
identifies the conjugate of \(M_a^{-1}\) with
\(M_a\) by this explicit positive real factor.
The boundary flow \(\Psi\) of IE.1 preserves the
physical circle, so its normal derivative becomes
a complex linear map after this conjugation. Its
full determinant is transported, including any
off-diagonal entries; no diagonal derivative is
assumed.
It evaluates to zero on \(F\) and \(2a\) on \(C_1\).
Its boundary class is zero integrally: in the
central Tate marking \(QP^2=1\), so \(Z^aV^2\)
is a nonvanishing boundary frame. Transport that
actual frame through the smooth source/Cartier
family. Equivalently the boundary component class
has order two and (P3K.57.IE.5) kills it. A relative lift
is \(-aC_1\), up to the null fibre and a boundary
connecting class. Thus
\[
                              b_{\rm r}^2=-2a^2.        \tag{P3K.57.IE.6}
\]

\item Here is the chart map identifying the collision E
patch with that same minimal I2 patch. The E face
between the primitive rays
\((0,j,1,0)\), \((0,j+1,1,0)\) has quotient lattice
\((Z,s)\), star rays \((0,1),(1,1)\), and primitive
source period \((Z,s)\mapsto(Z+2s,s)\). On the central
source write \(Q=t^2u_Q(t)\) and use the actual
square root with transported leading coefficient,
\(\widetilde q=t\sqrt{u_Q(t)}\). The two-node chart
chain is
\[
 \xi_\ell=Z\widetilde q^{-\ell},\qquad
 \eta_\ell=\widetilde q^{\ell+1}/Z,\qquad
 \xi_\ell\eta_\ell=\widetilde q,
 \quad \ell\pmod2.                                  \tag{P3K.57.IE.5a}
\]
Its literal overlaps are
\(\xi_{\ell+1}=\eta_\ell^{-1}\),
\(\eta_{\ell+1}=\xi_\ell\eta_\ell^2\).
These are the E-star characters with the retained
Tate units; the primitive stabilizer shifts
\(\ell\) by two. In particular the map has no
additional base ramification or exceptional curve.

\item For completeness, compose this map with the actual
Tate uniformization into the unit source
\(Y_0^2=X_0^3+X_0^2+T_0X_0\). In its marked Tate
coordinate the convergent functions are
\[
 X_T=\sum_{\nu\in\mathbb Z}
       \frac{Q^\nu Z}{(1-Q^\nu Z)^2}
       -2\sum_{\nu\ge1}\frac{\nu Q^\nu}{1-Q^\nu},
 \qquad
 Y_T=\sum_{\nu\in\mathbb Z}
       \frac{(Q^\nu Z)^2}{(1-Q^\nu Z)^3}
       +\sum_{\nu\ge1}\frac{\nu Q^\nu}{1-Q^\nu}.
\]
The retained minimal Weierstrass coordinate change
has leading terms
\(X_0=4X_T-T_0/2+O(T_0^2,T_0X_T)\) and
\(Y_0=8Y_T+4X_T+O(T_0Y_T,T_0X_T)\), while
\(T_0=16\widetilde q+O(\widetilde q^2)\).
On the chart \(Z=\xi\),
\(\widetilde q=\xi\eta\), this gives
\[
 T_0=16\xi\eta+O(\xi^2),\quad
 X_0=4\xi(1-\eta)^2+O(\xi^2),\quad
 Y_0=4\xi(1-\eta^2)+O(\xi^2).                       \tag{P3K.57.IE.5b}
\]
For every negative summation index, rewrite the
summand in \(Q^k/Z\); its exponents are nonnegative
in this node chart. The resulting series converge
normally on a smaller chart, so these are actual
holomorphic maps, including the axes.

\item The map to the two small-resolution charts is
\[
 z=X_0/Y_0,\quad D=T_0+X_0+X_0^2,\quad B=Y_0,
 \qquad
 w=Y_0/X_0,\quad C=X_0,\quad A=Y_0.                 \tag{P3K.57.IE.5c}
\]
The source equation gives \(A=Dz,C=Bz\) exactly.
At the marked two-torsion point \(\eta=1\), the
double zero of \(X_0\) and simple zero of \(Y_0\)
make the first quotient extend. On \(C_1\),
\[
 z=\frac{1-\eta}{1+\eta},\qquad
 D=4\xi(1+\eta)^2+O(\xi^2).
\]
The Jacobian of \((\xi,\eta)\mapsto(D,z)\) on
\(\xi=0\) is \(-8\). At \(\eta=-1\), use
\((C,w)\): its Jacobian has magnitude eight there.
The second node is covered by the literal chain
overlap in (P3K.57.IE.5a). This proves the biholomorphism
on the whole component, its two nodes, and the
marked section, with the source orientation and
boundary marking retained.

\item The actual normal determinant is marked by
\[
                         Z^aV^2/\xi_E^{\,2j+1}.
\]
After extracting its weight-two normal factor,
this character is a meromorphic section of
\(N_+\otimes N_-^{-1}\). It is invariant under the
source period, since \(QP^2=1\), and is nonzero
away from the two central components. Its
dual-character order is zero on each component.
The actual node charts above therefore make the
section invertible also at their intersections.
The Hermitian identification of this line with
the positive-weight spinor determinant is the
same one used in (P3K.57.IE.5). This proves zero
integral first Chern class, rather than merely
a rational vanishing of its degree. One can also
read the integral obstruction from the plumbing
retraction: \(H_1=\mathbb Z\), \(H_2=\mathbb Z^2\)
with primitive basis \(C_0,C_1\), so
\(\operatorname{Ext}(H_1,\mathbb Z)=0\); vanishing
on both components determines the zero class
in \(H^2\). Hence in this boundary marking
\[
                               b_{\rm c}^2=0.           \tag{P3K.57.IE.7}
\]
Replacing \(\xi_E=\epsilon/t\) by \(\epsilon\)
adds the explicit boundary gauge \(t^{2j+1}\).
It may be odd and is retained. A boundary
gauge changes the glued class by the
Mayer--Vietoris image of \(H^1(\partial S;\mathbb Z)\).
Its square and its product with the glued
determinant class vanish. Indeed represent
the connecting class as \(d\chi\wedge\alpha\)
in the collar; its square repeats \(d\chi\),
and its product integrates the zero boundary
restriction of the other class. Thus the
actual derivative of (P3K.57.IE.2) has exactly the
square difference (P3K.57.IE.7) minus (P3K.57.IE.6).

\item The source patches are oriented diffeomorphic
with boundary. Their closed double has
signature zero and hence tangent Pontryagin
integral zero. This uses only the ordinary
four-dimensional fixed-surface signature
identity. With weights \(+1,-1\), its
normal-link contribution is
\[
 \frac1{24}\int_{S_{\rm c}\cup(-S_{\rm r})}
            \{p_1(TS)-3b^2\}=-\frac{a^2}{4}.           \tag{P3K.57.IE.8}
\]
There are exactly \(N\) E doubles: the actual
scalar period shifts the primitive chain by
\(N\), and the collision-boundary source
loop fixes these labels.

\par\medskip\noindent\textbf{IE.4. Finite stabilizers and the actual Stokes identity.}\par\nopagebreak[4]

\item Remove invariant normal tubes of every fixed
component from \(W_{\rm rel}\). The physical
circle is locally free on the complement.
Its quotient is a compact oriented
seven-dimensional orbifold \(Z\), carrying
the actual orbifold Euler line with class \(x\).

\item A slice with stabilizer \(\mu_d\) gives a
finite chart \(\widetilde U/\mu_d\). Average
metrics, partitions of unity and connections
over \(\mu_d\) and over the physical circle.
Normalize the connection by
\(\alpha(\delta)=1\). Its curvature represents
the same physical \(x\) on every chart. The
actual tangent splitting is
\[
 TW_{\rm rel}|_{\rm locally\ free}
                    \cong\pi^*TZ\oplus\mathbb R_\delta .      \tag{P3K.57.IE.9}
\]
Choose its connection to match the actual
boundary tangent connections after a collar
interpolation. The basic Chern--Weil form
\[
                      \Omega=\frac{4x^3-p_1(TZ)x}{24}         \tag{P3K.57.IE.10}
\]
is closed. Integrate each finite chart as
\(d^{-1}\int_{\widetilde U}\). Ordinary Stokes
on the finite covers and their averaged
partition of unity cancels every interior
face with exactly these multiplicities.
The finite-isotropy strata are not boundary
faces. Consequently
\(\int_{\partial Z}\Omega=0\).

\item On a fixed normal link, (P3K.57.IE.9) restricts
to the fixed tangent and weighted normal
bundles. Its equivariant Thom denominator
is \(\prod_i(x_i+w_i u)\). The finite
chart integral, equivalently the weighted
projective power map, gives
\(1/\prod_iw_i\), with every negative
normal conjugation's orientation sign.
Thus its actual contribution is (P3K.57.IE.16), proved below.
The compact interior/H values and (P3K.57.IE.8)
are exactly the terms in this Stokes
identity, including finite stabilizers.
No spin structure on this comparison
bordism or eight-dimensional signature
formula is used.

\par\medskip\noindent\textbf{IE.5. Outer boundary and Jin conventions.}\par\nopagebreak[4]

\item Orient \(W_{\rm c}\) complexly and the
reference piece oppositely. Define
\(\delta\wedge\operatorname{or}(Z)
=\operatorname{or}(W_{\rm rel})\).
Near the outer parameter boundary,
\[
 \operatorname{or}(W)
 =\delta\wedge\operatorname{or}(Y)
                  \wedge d\rho\wedge d\theta .
\]
Since \(\dim Y=5\), its quotient boundary
orientation is
\[
                         -\operatorname{or}(Y)\wedge d\theta.       \tag{P3K.57.IE.11}
\]
At an inner fixed-normal boundary the
minus sign from removing the normal disc
and the minus sign from quotienting an
oriented boundary by the odd-dimensional
orbit cancel. Its orientation is the
positive complex-normal link convention
of (P3K.57.IE.16). Therefore Stokes gives
\[
                 A_{\rm physical,difference}
                          =\sum_F\mathcal R(F).          \tag{P3K.57.IE.12}
\]

\item Jin identifies \((z,0)\sim(f(z),1)\).
The positive physical family with return
\(f\) has the inverse parameter marking.
The map to the Jin six-manifold reverses
that circle and preserves the physical
Euler line. Hence
\[
                  A_{\rm Jin,difference}
                           =-\sum_F\mathcal R(F).         \tag{P3K.57.IE.13}
\]
The common conifold replacement in IE.2
uses the same transported positive
meridian, so adds no line or orientation
change.

\par\medskip\noindent\textbf{IE.6. Scalar-order independence and exact local values.}\par\nopagebreak[4]

\item The selected-cusp fan in this argument is the initial
vertical-edge fan or its literal reflected fan under
\((Z,V)\mapsto(Z^{-1},P^aV)\). In either marking the
physical delta circle has the already specified
\((-1,+1)\) fixed-curve model, and the selected-unit
transport in IE.2 closes that same fan. The admitted
source/cap sectors are \((a,b)=(1,2)\), \(m=4p\),
and \((a,b)=(2,1)\), \(m=4p-1\), with \(p\ge1\),
their actual \(n=\pm1\) norm markings, and the
retained integral \(d\delta\) twist. No comparison
with a further flopped fan is used here.

\item The fixed-source scalar integers in these two particular markings
are equal by the actual C2 comparison of P3K.62. Its first outer
branch map is the identity; its second has base multiplier
\(\kappa^{-1}\) and normal matrix
\(\operatorname{diag}(\alpha^{-1},\kappa\alpha)\), both independent
of the scalar parameter. These are exactly the outer maps used
on the selected collar here. Together with the equal inner
charged-cap gauges, they give the line-preserving integer
comparison in P3K.62. This invocation retains the two outer
maps and uses no assertion of equality for a further fan.

\item Over \(U\) the ordinary suspension, with
no linearization zero, is an actual family
over the product of a source annulus and
a small scalar annulus:
\[
                   B_\lambda=2(-1)^{b+1}\lambda q_*U_a^b.
\]
There is no finite cap in this local
family, so every integer scalar winding
is a closed local loop. Its source-base
boundary family extends over the source
disc, with the actual periods and Cartier
frames. Choose IE.1's transport on this
entire product.

\item A pair of pants mapping into this
product realizes the based-loop identity
\[
 (\epsilon\text{ turn},N\text{ scalar turns})
   =(\epsilon\text{ turn},0)+(0,N\text{ scalar turns}).    \tag{P3K.57.IE.14}
\]
Pull back the actual local family and
quotient its free physical circle.
This is a seven-dimensional bordism
of the mapping-torus pieces. Glue the
source-base boundary collars by the
full-product transport. The fixed-source
scalar-loop boundary is exactly the
reference piece of IE.1 and cancels.
Thus the closed relative collision
boundaries \(Q_{\rm rel}(N)\) and
\(Q_{\rm rel}(0)\) are bordant with
their actual Euler line. Their
characteristic difference is independent
of \(N\). This is a literal parameter
product and collar construction; no
division of a covered sphere class occurs.

\item For any admitted cap, its exterior
source/root/line family extends over
the source disc. IE.2 cancels it with
the constant exterior. Hence this
same local relative boundary represents
the fixed-scalar source loop for that
completed family. It suffices to use
the two fully executed orders \(N=8\)
for \((a,b)=(1,2)\), and \(N=6\) for
\((a,b)=(2,1)\):
\[
 \begin{array}{c|rrr}
 (a,b,N)&\text{interior/H}&N\text{ E doubles}&\text{total}\\ \hline
 (1,2,8)&2&-2&0\\
 (2,1,6)&5&-6&-1 .
 \end{array}                                           \tag{P3K.57.IE.15}
\]
Independent executions at \(N=4,12\)
in the first case and \(N=3,9\) in
the second give the same relative
sums. They check the arithmetic;
the independence proof is (P3K.57.IE.14).

\item By (P3K.57.IE.13), the fixed-scalar marked infinity
meridian-one dual differences are respectively \(0\) and \(+1\).
Apply P3F.22's proved Jin-to-\(\Theta_7\) transfer
to the two actual completed boundary mapping
tori: on each, IE.2 has made the specified conifold
replacement, transported the same positive meridian,
and retained the actual Euler line of its free circle
bundle. To select their actual return maps, prescribe
the product connection on the common selected tube
in IE.2's closed Cox product, extend it to a connection
on the actual parameter family, and average over
the physical delta circle. The prescribed product
connection remains unchanged under that averaging.
After the conifold replacement its return is a
diffeomorphism \(f\) of \(Y=S^2\times S^3\), equal
to the identity on the prescribed tube. The actual
Euler line extends around the family, so \(f\)
preserves its primitive class \(x\in H^2(Y;\mathbb Z)\).
The return preserves orientation; the perfect
pairing \(H^2\otimes H^3\to H^5\) then shows that
it preserves \(H^3\) as well. Thus \(f\in I(Y)\).

\item On the free complement, apply the Hermitian line map
(D_M:M^\times\to U^\times) to the actual equivariant return.
It conjugates the lifted physical map to its meridian-one dual
lift and extends across the tube by
((a,b)\mapsto(\overline b,-\overline a)), with the full
(SU(2)) transition functions retained.
On the two cap pieces the first-to-second maps are
(lambda_M(F)=-Sigma_{F^{-1}}) and
(lambda_U(F)=Sigma_{F^{-1}}), as proved in
\cref{night:P3K.68}. Thus the computed characteristic difference
(A_U=\int(4x^3-p_1x)/24) gives physical class (-A_U).
The pair-of-pants endpoint maps (P3K.57.IE.14) identify this
relative difference with the fixed-scalar positive epsilon loop.
Its physical classes are therefore (0) for ((a,b)=(1,2))
and (-1) for ((a,b)=(2,1)), on both specified vertical fan
objects. Every residue and Euler denominator above is retained;
the change is the explicitly proved cap/line transfer.

\item The old
\(v\)/new \(\epsilon\) scalar path,
\(D^{-4}\) root endpoint and
half-loop fan map remain the explicit
dictionary of IF.1--6. Their winding
corrections are not absorbed into
(P3K.57.IE.15).

\par\medskip\noindent\textbf{IE.7. Complete normal-link calculation for the two executed fans.}\par\nopagebreak[4]

\item This section supplies the local formulas used in IE.4 and
the exact finite sums in (P3K.57.IE.15), independently of the
uncompleted source-zero attachment. Let a compact fixed
component \(F\) have complex normal roots \(x_i\) and
nonzero physical delta weights \(w_i\). Retain its complex
orientation and the original complex normal-sphere
orientation. Its actual normal-link quotient contribution is
\[
 \mathcal R(F)=[u^{-1}]\int_F
 \frac{4u^3-u\{p_1(TF)+\sum_i(x_i+w_i u)^2\}}
      {24\prod_i(x_i+w_i u)}.                         \tag{P3K.57.IE.16}
\]
The coefficient is the Laurent coefficient at infinity.
On the normal sphere, adding its radial real line gives
the pullback of \(TF\oplus\nu_F\); quotienting the
physical orbit removes its oriented real line. Thus the
stable quotient Pontryagin class is the displayed
equivariant class. The equivariant Thom class restricts
to \(\prod_i(x_i+w_i u)\). The localized pushforward
therefore gives (P3K.57.IE.16), with the Euler and integration
identities of Atiyah--Bott \cite[pp.~8--9, Theorem~(3.5), equations~(3.7)--(3.8)]{AtiyahBott1984}.

\item The fibre normalization can be verified directly. For
positive weights, the weighted projective power map
\([z_i]\mapsto[z_i^{w_i}]\) has generic degree
\(\prod_iw_i\) and pulls its rational hyperplane class
back to the ordinary hyperplane class. Consequently
\(\int u^{r-1}=1/\prod_iw_i\). The relation
\(\prod_i(x_i+w_i u)=0\) gives all higher pushforwards
recursively. Conjugating a negative-weight line changes
\((x_i,w_i)\) to \((-x_i,-w_i)\) and reverses the
original normal orientation. This is precisely the
denominator sign in (P3K.57.IE.16); every finite stabilizer
factor remains.

\item For a point and a curve the respective formulas are
\[
 \frac{4-\sum_iw_i^2}{24\prod_iw_i},\qquad
 -\frac1{24\prod_iw_i}\int_F\sum_i
 \left\{\frac{4-\sum_jw_j^2}{w_i}+2w_i\right\}x_i.
                                                               \tag{P3K.57.IE.17}
\]
For a surface, put
\[
 P=w_1w_2,\quad c=4-w_1^2-w_2^2,\quad
 A=\frac{x_1}{w_1}+\frac{x_2}{w_2},\quad
 B=\frac{x_1^2}{w_1^2}+\frac{x_1x_2}{w_1w_2}
                         +\frac{x_2^2}{w_2^2}.
\]
Its contribution is
\[
 \int_F\frac{cB+2(w_1x_1+w_2x_2)A
                 -p_1(TF)-x_1^2-x_2^2}{24P}.             \tag{P3K.57.IE.18}
\]
Weights \((1,-1)\) give
\((p_1(TF)-3(x_1-x_2)^2)/24\), the formula used
in (P3K.57.IE.8), in the same original normal orientation.

\item For each exported regular cone let \(C\) be the
ordered matrix of its four primitive rays in the
original \((Z,V,\text{base}_0,\text{base}_1)\)
lattice. The coordinate characters are the rows of
\(C^{-1}\), and their actual physical weights are
\[
                    C^{-1}(0,1,0,0)^{\mathsf T}.        \tag{P3K.57.IE.19}
\]
The vanishing of the nonzero-weight coordinates
specifies every local fixed stratum. The recorded
face permutations and literal period matrices glue
these strata, so no component is counted twice or
omitted. There is no fixed stratum in the open
protected A or B face: their span is respectively
\(\langle(1,0,0,0),(0,0,1,0)\rangle\) or
\(\langle(1,-a,0,0),(0,0,1,0)\rangle\), neither
of which contains the physical delta vector.
Their remaining boundary strata belong to the
already recorded H/interior cones. The protected
E face does contain delta; its fixed surfaces are
exactly the N chain components treated in IE.3.
Thus the interior, H and E contributions exhaust
the fixed set of the relative comparison.
On an interior fixed component, unwrap
each chart by its recorded period \(n\). An
auxiliary cocharacter \(e\) then has tangent and
normal weights
\[
                            C^{-1}T_n e.                 \tag{P3K.57.IE.20}
\]
This is the dual of the actual map
\(C_{\rm unwrapped}=T_{-n}C\). Applying the same
pushforward to the fixed component divides (P3K.57.IE.17)
or (P3K.57.IE.18), evaluated at its auxiliary roots, by
the product of its tangent weights. The two
independently executed choices are
\((1,4,16,64)\) and \((1,5,25,125)\); the code
verifies agreement separately on every component.

\item On the exceptional base H use \(w=q_B/q_A\) and
normal coordinate \(q_A\). The exact periods have
normal-order matrix
\[
 H_\infty=\begin{pmatrix}2&b\\-a&N-2\end{pmatrix},
\]
and remaining nonconstant factors
\[
 h:(Z,V)\mapsto(wZ,w^{-a}V),\qquad
 g:(Z,V)\mapsto(w^bZ,w^{-2}(1-w)^N V).                 \tag{P3K.57.IE.21}
\]
All other factors are the retained nonzero source,
Cartier and suspension units of IF.5. The three
marked base points are \(w=0,1,\infty\), on the
B,E,A faces respectively. For each regular rank-three
H face, drop its zero first base coordinate and
form the unimodular matrix \(B_H\). Its orbit
closure has the actual base coordinate \(w\), so
is \(\mathbb P^1\). If
\((p_j,q_j,s_j)\) is row \(j\) of \(B_H^{-1}\),
its generic normal coordinate is
\(\eta_j=Z^{p_j}V^{q_j}q_A^{s_j}\).

\item For an occurrence at period \(n=(n_1,n_2)\), let
\((p,q,r,s)\) be its corresponding full chart row.
The exact recorded checks are
\((p,q)=(p_j,q_j)\) and
\(s+(p,q)H_\infty n=s_j\). Set
\[
 e_w=p(n_1+bn_2)+q(-an_1-2n_2),\qquad e_1=Nn_2q.
\]
The three actual coefficient orders of the generic
normal vector, in its holomorphic local frames, are
\[
 \begin{array}{c|cc|c}
  \text{sector}&x&y&\text{order at its marked point}\\ \hline
  B&w&q_A&e_w+r\\
  E&(1-w)/u_\epsilon&q_A&e_1+r\\
  A&w^{-1}&wq_A&-e_w+r-s-e_1.
 \end{array}                                             \tag{P3K.57.IE.22}
\]
Their sum is the actual normal-line degree. In
particular the \((1-w)^N\) term, the higher ray
heights, and every recorded period remain in the
degree computation.

\item The two auxiliary fibre circles \((1,2,0,0)\) and
\((1,3,0,0)\) commute with (P3K.57.IE.21) and the full
periods. On each compact H fixed surface, their
fixed sets consist of these rank-three curves and
the isolated chart fixed points. A sector-interior
three-face has base rank two and only the physical
one-dimensional fibre kernel, so supplies no
additional auxiliary fixed curve. If an auxiliary
fixed curve has normal weight \(\eta\), normal
degree \(d_t\), and ambient normal weights/degrees
\((\beta_i,d_i)\), substitute
\[
 x_i=\beta_i u+d_i h,\qquad
 p_1(TF)=\eta^2u^2+2\eta d_t hu,\qquad h^2=0
\]
into (P3K.57.IE.18). Writing the numerator polynomial
there as \(R_0u^2+R_hhu\), its exact localized
contribution is
\[
                         R_h/\eta-R_0d_t/\eta^2.         \tag{P3K.57.IE.23}
\]
This divides by the actual Euler class
\(\eta u+d_t h\). Isolated fixed points use the
two nonzero tangent weights. Both auxiliary circles
agree component by component. Independent outputs
also compute \(c_1(TF)^2,c_2(TF)\), and verify the
integrality of \((c_1^2-2c_2)/3\) and
\((c_1^2+c_2)/12\) on every H surface.

\item The complete finite outputs are
\[
 \begin{array}{c|rr|rr|r}
 (a,b,N)&\#\text{ interior}&\sum_{\rm int}\mathcal R
 &\#H&\sum_H\mathcal R&\text{combined}\\ \hline
 (1,2,8)&2153&9333/16&125&-9301/16&2\\
 (2,1,6)&2147&739&111&-734&5.
 \end{array}                                             \tag{P3K.57.IE.24}
\]
There are 194 and 152 full-fibre auxiliary fixed
curves on H, respectively. Every primitive ray,
ordered determinant-one cone, period, face
permutation, normal degree and rational summand
is retained in the corresponding
\(\texttt{source\_infinity\_regular\_fan}\),
\(\texttt{fixed\_charts}\),
\(\texttt{interior\_residues}\), and
\(\texttt{h\_residues}\) JSON/certificate files
under \(\texttt{code/part3\_fourfold/}\).
Their scope is the actual infinity matrices and
period maps just displayed. Combining (P3K.57.IE.24)
with the E doubles proved in IE.3 gives (P3K.57.IE.15).

\end{enumerate}
\end{proof}




\begin{theorem}\label{night:P3K.58}
\StatementRecord{night:P3K.58}{false as stated; correction proved}{S6\_06\_BLIND: GPT-6 Astra (OpenAI, ChatGPT Pro mode, effort Pro), recorded re-derivation in the checkpoint dated 18 November 2026}{\cref{night:P3F.51}, \cref{night:P3F.56}, \cref{night:P3K.26}, \cref{night:P3K.45}, \cref{night:P3K.49}, \cref{night:P3K.53}, \cref{night:P3K.57}, \cref{night:P3K.62}, \cref{night:P3K.92}}Use the actual marked E7 source and section
\[
 Y^2=X^3+(1-2v)X^2+TX+v^2,\qquad
 G=(0,v),\qquad T_0(v)=v^2-2v.
\]
Retain the selected whole-zero \(I_1\) cusp, one III-star cap,
and the complete Cartier line, finite character and reference
annular data of P3K.49. The cap data are
\[
 (a,b,m)=(1,2,4p)\quad\hbox{or}\quad(2,1,4p-1),
 \qquad p\ge1,\quad n=\pm1,\quad
 N=km,\quad k\ge1.
\]
The literal scalar is \(\lambda(v)=v^N\lambda_0\), in the
common contracting range of each bounded parameter surface
used below. No additional logarithmic disc meets the two
finite sites that coalesce. The cusp fan is the initial
physical-vertical-edge fan \(\mathcal T_0\) or its reflection
\(J\mathcal T_0\), where \(J=\operatorname{diag}(-1,1)\).

On the cover \(v=\rho^2-1\), the constructed four-holed
parameter surface has the actual cut-surface endpoint relation
\[
 [Z_{\mathcal T}]+[Z_{J\mathcal T}]-1+[I_{\mathcal T}]=0
                    \quad\text{in }\Theta_7\cong\mathbb Z/28.
\]
Here \(Z_{\mathcal T}\) is the positive source-zero return
near \(\rho=1\), and \(I_{\mathcal T}\) is the positive
\(\epsilon=1/\rho\) return with the full original scalar
path. The reflected source-zero object and every closing
arrow are supplied by the literal source, theta/Cartier,
root and annular maps in the proof.

The full infinity comparison retains \(D^{-4}\), the reduced
doubled lattice period, and the scalar quadratic cochain.
Together with the independently proved infinity excision
P3K.57 and the actual independent scalar parameter, it gives
\[
 \begin{array}{c|c|c}
 (a,b,m)&I_{\mathcal T}
                &[Z_{\mathcal T}]+[Z_{J\mathcal T}]\\ \hline
 (1,2,4p)&0&1\\
 (2,1,4p-1)&-1-2k(3p-1)&2+2k(3p-1).
 \end{array}
\]
These formulas retain both norm signs and the stated two
fan objects. P3K.49 carries each original positive \(u\)-loop
to three positive \(v\)-loops by an actual family morphism.
No individual source-zero class or division of a class is
assigned by this statement.
\end{theorem}
\begin{proof}
\begin{enumerate}[label=\arabic*.]
\item \par

\medskip\noindent\textbf{RG.1. A bounded parameter surface and its actual charts}

\par\nobreak

\item Take a large finite \(R>2\), and remove small disjoint discs around
\(\rho=-1,0,1\) from the disc \(|\rho|\le R\). Denote this compact
four-holed surface by \(B_R\). Its two points over \(v=1/8\) are
filled using the actual family of K53. We choose \(\lambda_0\) in
the common contracting range on this compact parameter surface.
This does not presume that the same fixed scalar remains
contracting throughout a punctured neighbourhood of infinity.

\item All source equations and the canonical Cartier divisors are the
ones displayed above. The discriminant of their cubic in \(X\),
viewed as a polynomial in \(T\), has parameter discriminant
\[
 16v^2(v+1)^6(1-8v)^3.
\]
Thus on this bounded parameter surface the only excluded source
degenerations are the two source-zero lifts and the selected
collision. The unselected type-II point is precisely the disc
already constructed in K53.

\item Choose a single cap-overlap annulus sufficiently far out in the
\(T\)-plane to avoid all finite critical values for the compact
range of \(v\). The III-star root equations are
\[
 y'^2=x'^3+(1-2v)z^{2p}x'^2+
 \begin{cases}x',&m=4p,\\zx',&m=4p-1,\end{cases}
 +v^2z^{6p},\qquad G'=(0,vz^{3p}).
\]
Their central models, finite kernel branch and root-line
corrections are the ones in K49 and K53. They are uniform for
bounded \(v\) after shrinking the root disc. The norm-ratio has
the actual nonzero central value \(-2\) or \(+2\), respectively.
Consequently its logarithm and chosen root extend on the root
disc times the bounded \(v\)-disc with that reference value.
The same finite averaging preserves the literal root equation
\(A_0^m=B_0^n\) and every finite character. No cube root of \(v\)
is used to define these global cap charts.

\item The full reference annular map \(F_0\) of K49 is continued on
this fixed cap annulus. Its regular elliptic source, root,
Cartier and norm maps are holomorphic for every \(v\) in that
bounded disc, including zero. The parameter-disc exponential
lifting and the fixed-contour Cauchy splitting of RCZ.5 therefore
apply on the entire bounded disc, with the same full reference
map. One may choose the fixed cap annulus this large at the
initial reference; transport to it is through the literal source
and cap charts, carrying every integral annular period. The
construction retains the entire reference \(d\delta\) and the
reduced doubled period. Its restriction near \(v=0\) is the
marked continuation used in K49.

\item The selected cusp has the exact normalization
\[
 X=s^2-1,\qquad Y=s(s^2-\rho^2),\qquad
 Z=\frac{s-\rho}{s+\rho}.
\]
The marked point \(G\) is \(s=-1\), so
\[
 P_G=\frac{1+\rho}{1-\rho}.
\]
It is a nonzero finite function on \(B_R\). The unique small
Tate period and the physical height \(q_1=T-T_0\) satisfy
\[
 \left.\frac{Q}{q_1}\right|_{q_1=0}
       =-\frac{v}{64(v+1)^3}
       =-\frac{\rho^2-1}{64\rho^6}.                 \tag{RG.1}
\]
The actual unit \(Q/q_1\), not just its leading coefficient, is
retained in the charts. The chosen fourth root of the nodal
\(c_4\) is \(2\rho\). Hence the Tate source coordinate is fixed
over the rho surface; there is no unrecorded branch interchange.
The lift of \(G\) is fixed by its displayed finite nodal limit.
Any alternative continuation differing by \(Q^j\) would have
limit zero or infinity for \(j\ne0\), so that integral ambiguity
does not occur. The theta/Cartier line map then fixes the line
coordinate as well.

\item A chosen periodic cusp fan \(\mathcal T\) therefore gives actual
toric charts over this entire bounded parameter surface.
Nonzero unit windings in (RG.1) and \(P_G\) remain in their
transition maps; no logarithm is used to discard them.
The uniform norm shells, finite caps and compact-base gluing
give a smooth proper sphere family over \(B_R\), with the
prescribed physical circle. The cap and selected-cusp charts
just specified identify the actual endpoint groupoid below.

\par

\medskip\noindent\textbf{RG.2. The full scalar cut and its forward endpoint arrows}

\par\nobreak

\item Choose \(t=\log v/(2\pi i)\) on parameter cuts and
\[
 s_\infty=\frac{\log q_\infty}{2\pi i},\qquad
 h=1-T_0(v)q_\infty,\qquad q_1=q_\infty^{-1}h.
\]
The actual ordinary-to-cap cochain is
\[
 F_t=F_0+nk\,t\,s_\infty\delta.
\]
For an integer change \(t\mapsto t+w\), its forward arrows are
\[
 R^{\rm ord}_w=q_1^{nkw},\qquad
 R^{\rm cap}_w=h^{nkw}.                              \tag{RG.2}
\]
Multiplication here denotes the actual line maps. Their square is
\[
 F_{t+w}\circ R^{\rm ord}_w
        =R^{\rm cap}_w\circ F_t,                     \tag{RG.3}
\]
because the additional factor in \(F\) is \(q_\infty^{nkw}\).
This retains the direction: \(R_w\) is the map from the initial
lifted fibre to the terminal lifted fibre. The physical forward
return uses its inverse as the closing arrow.

\item These are maps of every retained local model. On the cap, \(h\)
is an invariant nonzero unit, so its multiplication commutes
with the root, source, line and norm maps. At the selected cusp,
\(q_1^{nkw}\) is its integral line-height translation. Reindexing
by the corresponding suspension period leaves a bounded toric
multiplier, with its full source and line factors. Thus it
extends through the periodic fan as the actual integral
reindexing. Its inverse extends too. The finite characters and
the total cyclic lattice cocycle stay in \(A_0^m=B_0^n\).

\item The scalar-root parameter itself is single valued:
\[
 r=v^k r_0,\qquad r_0^m=\lambda_0,\qquad
 A=r^nA_0,\quad B=r^mB_0.
\]
The cut in (RG.2) is consequently the retained annular gluing
cut, not a newly inserted root character. The arrows obey
\(R_{w+w'}=R_{w'}R_w\) literally. For the positive local rho
loops the integers are
\[
 \begin{array}{c|rrrr}
 \text{puncture}&-1&0&1&\infty\\ \hline
 \operatorname{wind}v&1&0&1&-2\\
 \operatorname{wind}r&k&0&k&-2k .
 \end{array}                                           \tag{RG.4}
\]
The infinity orientation in this table is the positive
\(\epsilon=1/\rho\) orientation.

\par

\medskip\noindent\textbf{RG.2a. The actual source-zero closing on every cusp chart}

\par\nobreak

\item Near the lift \(\rho=1\), retain the two-node Tate/Cartier
coordinates of SZ.39--44. With \(q=q_1\), their periods are
\[
 h:(Z,V)\mapsto(\kappa qZ,\alpha V),\qquad
 g_{\rm per}:(Z,V)\mapsto(gZ,\eta qV),
\]
where
\[
 \kappa=q_2u_Q,\quad \alpha=P^a,\quad
 g=P^{-b},\quad
 \eta=2(-1)^{b+1}\lambda_0v^Nc_X^{-2},\quad
 P=q_2^{-1}u_P,\quad c_X=q_2u_X.
                                                        \tag{RG.4a}
\]
All the displayed \(u\)'s are the actual nonzero two-node
units with their extending logarithms. At \(q_1=0\), the
function \(q_2\) has a simple zero in \(v\). A positive \(v\)
loop therefore has the complete period-unit winding matrix
\[
 \begin{pmatrix}
  \operatorname{wind}\kappa&\operatorname{wind}g\\
  \operatorname{wind}\alpha&\operatorname{wind}\eta
 \end{pmatrix}
 =
 \begin{pmatrix}1&b\\-a&N-2\end{pmatrix}.              \tag{RG.4b}
\]
Its trace is \(N-1\). In particular the three raw normal
determinants in the labelled six-chart plumbing have that
winding; for even \(N\) their underlying real normal loops
are nontrivial in \(\pi_1(SO(4))\). The fixed-labelled-vertex
ordinary product used for the selected rho loop is not used
in this source-zero comparison.

\item The endpoint itself is also retained. Put \(L=nk\). The
forward arrow of (RG.2) at this puncture is
\[
             R_{\rm ord}:(Z,V,q)\mapsto(Z,q^LV,q).
                                                        \tag{RG.4c}
\]
For the two triangle charts
\[
 \begin{aligned}
 T^+_{i,j}:\quad &q=x_0x_1x_2,\quad Z=q^ix_1,\quad V=q^jx_2,\\
 T^-_{i,j}:\quad &q=y_0y_1y_2,\quad
 Z=q^iy_0y_1,\quad V=q^jy_0y_2.
 \end{aligned}
\]
this sends the index \((i,j)\) to \((i,j+L)\), with the
three chart coordinates unchanged. To return to the chosen
index representatives apply the actual period
\(g_{\rm per}^{-L}\). On their dense tori the resulting map is
\[
             (Z,V)\mapsto(g^{-L}Z,\eta^{-L}V).        \tag{RG.4d}
\]
Consequently its literal maps in the named charts are
\[
 \begin{split}
 T^+:\ &(x_0,x_1,x_2)\mapsto
       ((g\eta)^Lx_0,g^{-L}x_1,\eta^{-L}x_2),\\
 T^-:\ &(y_0,y_1,y_2)\mapsto
       ((g\eta)^{-L}y_0,\eta^Ly_1,g^Ly_2).
 \end{split}                                         \tag{RG.4e}
\]
These formulas retain the complete bounded multipliers at the
endpoint, rather than replacing the endpoint by identity in
the six named charts. Every exponent is integral. They glue
by the actual period maps and extend to the toric boundaries.
For another periodic fan the same proof uses its character
matrix: (RG.4c) translates every height-one vertex by
\((0,L)\), and (RG.4d) is the remaining bounded torus map.
The physical return uses the inverse of (RG.4e) after path
transport. The cap endpoint is the unchanged-index map
\[
       (X_j,Y_j)\mapsto(h^LX_j,h^{-L}Y_j)
\]
with its actual finite root action retained. Equations
(RG.3) and (RG.4c--e) give their full common overlap.

\item The opposite source-zero lift is obtained by the literal
reflection in RG.3; its fan is \(J\mathcal T\), and its
closing is the conjugate of these full maps. No identity
closing, labelled product, or omitted chain reindexing is
assumed in (RG.8).

\par

\medskip\noindent\textbf{RG.3. Reflection as an actual arrow between the two fan objects}

\par\nobreak

\item At \(-\rho\) the source and cap data depend on the same \(v\).
The selected normalization exchanges its two branches. The
literal Tate and line map is
\[
 \mathcal R_\rho:\ (Z,V)\longmapsto(Q/Z,V),
 \qquad P\longmapsto P^{-1}.                          \tag{RG.5}
\]
It is the identity on the underlying marked source and its
Cartier divisor, expressed in the two actual Tate coordinate
systems. The line arrow follows from the theta identities,
without a suppressed multiplier:
\[
 \frac{\Theta_Q(QP/Z)}{\Theta_Q(Q/Z)}
       =\frac{\Theta_Q(Z/P)}{\Theta_Q(Z)}.
\]
For \(M_a\), apply the same identity to \(P^a\), retaining its
actual Cartier representative. The ordinary source, cap and
\(F_0\) maps are carried identically because they depend on
\(v\). Equation (RG.3) is also intertwined, since \(q_1,h\) are
unchanged and the line coordinate in (RG.5) is unchanged.

\item The fan map in the original lattice is
\[
 (x_Z,x_V,h)\longmapsto(-x_Z+h,x_V,h).
\]
Its linear two-dimensional part is
\(J=\operatorname{diag}(-1,1)\); the integral height translation
and the nonzero unit \(Q/q_1\) are retained. This gives the
actual arrow
\[
 \mathcal X_{\mathcal T,\rho}\longrightarrow
             \mathcal X_{J\mathcal T,-\rho}.
\]
Its square is the identity on all charts, including the cap.
It preserves the physical circle and the complex orientation.
The parameter map \(\rho\mapsto-\rho\) has positive local
complex degree one at both source-zero lifts.

\item Let \(Z_{\mathcal T}\) denote the positive return near
\(\rho=1\), with its actual full marked data, and let
\(Z^-_{\mathcal T}\) denote that near \(\rho=-1\). Then the
actual family arrow (RG.5) proves
\[
                  [Z^-_{\mathcal T}]=[Z_{J\mathcal T}]. \tag{RG.6}
\]
This does not identify \(Z_{\mathcal T}\) with
\(Z_{J\mathcal T}\). The two fan objects remain distinct inputs
to the loop relation.

\par

\medskip\noindent\textbf{RG.4. The cut-surface groupoid relation and the selected loop}

\par\nobreak

\item Choose a rectangular outer boundary containing the three holes,
with disjoint vertical slits from their lower points to its
lower edge. Use its lower-right point as basepoint. Walking the
boundary of the cut-open disc, and then reversing that boundary
word, gives based positive meridians in the order
\[
             \gamma_+\gamma_0\gamma_-\gamma_\infty=1. \tag{RG.7}
\]
Concatenation in this equation means traversal from left to
right. The outer positive infinity meridian is clockwise in
the rho plane. The slit paths give actual transport maps from
the reference fibre to the four local fibres.

\item Choose a smooth connection for the actual family over the
cut surface, averaged over the physical circle. On a lifted
parameter path, let \(P_\gamma\) be its transport. Its endpoint
in the reference fibre is
\[
             \Phi_\gamma=R_{w(\gamma)}^{-1}P_\gamma,
\]
with (RG.2), and with the slit-path conjugations included.
The total integer in (RG.7) is \(1+0+1-2=0\), so its total
closing arrow is literally \(R_0=1\). A contraction of the
cut-disc boundary gives, by lifting its paths through this
connection, an isotopy of the product of its four actual
physical returns to the identity. Thus, as arrows on the
reference oriented sphere,
\[
       \Phi_\infty\Phi_-\Phi_0\Phi_+\simeq\operatorname{Id}.
                                                               \tag{RG.8}
\]
No fibrewise conjugacy alone is used to deduce this relation.
It comes from the actual family on the cut surface and its
literal endpoint maps.

\item Near \(\rho=0\), take the analytic cube root \(u^3=v\) with
\(u(0)^3=-1\). In the \(G=(0,1)\) source coordinates put
\[
 \xi=x-u,\qquad \sigma=\rho/u.
\]
The selected curve is exactly
\[
 y^2=\xi^2(\xi+\sigma^2),\qquad
 q=\frac{y^2-\xi^2(\xi+\sigma^2)}{\xi+u}.
\]
These are the actual F56 source coordinates. The parameter
\(\sigma=\rho/u\) has nonzero complex derivative \(u(0)^{-1}\).
The normalized source point has
\(s_{\rm norm}=s/u\), so its Tate coordinate is exactly
\((s-\rho)/(s+\rho)\), with no further reflection. The
Cartier and root maps are the analytic local coordinate maps
of K49/F51; their parameter units have extending logarithms.
The scalar ratio
\(\lambda(v)/\lambda(-1)\) has an extending logarithm on this
disc. Its explicit homotopy through that logarithm, with the
same root/line cochains, identifies the positive local return
with F56's fixed-scalar return. For the initial fan and its
reflected fan, K45/F56 therefore give
\[
                         [\Phi_0]=-1.
\]
The two unselected points over \(v=1/8\) have the actual
disc extensions of K53 and give zero returns.

\item Write \(I_{\mathcal T}\) for the positive infinity return in
this actual bounded-parameter family. Applying the oriented
mapping-class identification with \(\Theta_7\), and using the
actual slit-path conjugations in (RG.8), gives
\[
 [Z_{\mathcal T}]+[Z_{J\mathcal T}]-1+[I_{\mathcal T}]
                  =0\quad\hbox{in }\Theta_7.          \tag{RG.9}
\]
Conjugation disappears here because this oriented group is
abelian; the underlying family arrows were retained before
passing to it. K49 separately identifies each original
positive \(u\)-return with three positive \(v\)-returns.
No value or division is assigned to either source-zero term.
The infinity term is evaluated by the separately proved excision
in RG.7 below.

\par

\medskip\noindent\textbf{RG.5. The infinity scalar path and the independent scalar parameter}

\par\nobreak

\item The term \(I_{\mathcal T}\) is defined by the actual positive
\(\epsilon=1/\rho\) loop, before any change to the infinity
source coordinates. Its literal scalar is
\[
 \lambda=(\epsilon^{-2}-1)^N\lambda_0,\qquad
 r=\epsilon^{-2k}(1-\epsilon^2)^k r_0.                \tag{RG.10}
\]
The unit \((1-\epsilon^2)^k\) has an extending logarithm;
using that logarithm gives an actual based homotopy to a
root-scalar loop of winding \(-2k\). At a fixed sufficiently
small epsilon circle, all amplitudes are chosen in a common
contracting range.

\item For clarity, one can use the actual source construction with
independent root-scalar variable \(r\). Its maps are
\[
 B=r^mB_0,\qquad A=r^nA_0,\qquad
 F=F_0+n\frac{\log r}{2\pi i}s_\infty\delta.
\]
The positive \(r\)-loop has the full forward arrows
\(q_1^n\) and \(h^n\). These commute with the retained source
charts and finite groupoid maps, so this is an actual family
over the epsilon annulus times a small \(r\)-annulus. The
parameter square \((\epsilon,r)=(\epsilon_*e^{2\pi is},
r_*e^{2\pi i(ps+qt)})\) gives the usual product-parameter
homotopies with all those endpoint arrows retained.

\item Denote by \(I_{\mathcal T}(j)\) the actual coupled loop of this
family with one positive epsilon turn and \(r\)-winding \(j\).
Then
\[
 I_{\mathcal T}(j)-I_{\mathcal T}(j')
       =(j-j')A^{\rm scalar}_{\mathcal T}.             \tag{RG.11}
\]
This is a comparison of actual loops in the same constructed
two-parameter family. In particular \(I_{\mathcal T}\) in
(RG.9) is \(I_{\mathcal T}(-2k)\). A different local rescaling
\(\lambda=\epsilon^{N_\infty}\lambda_\infty\), with
\(N_\infty=k_\infty m\), is the distinct loop
\(I_{\mathcal T}(k_\infty)\), with difference
\[
 I_{\mathcal T}(-2k)-I_{\mathcal T}(k_\infty)
       =-(2k+k_\infty)A^{\rm scalar}_{\mathcal T}.
\]
An identification of the latter loop with another infinity
model must carry that model's complete source, root,
Cartier, annular and normal arrows. RG.6 supplies this
annular dictionary; only afterward does RG.7 use the
separately proved infinity excision.

\item For the initial and reflected vertical-edge fans the actual
fixed-source page calculation K26 and the literal C2 native
comparison K62 give
\[
 A^{\rm scalar}=0\quad(a,b,m)=(1,2,4p),\qquad
 A^{\rm scalar}=3p-1\quad(a,b,m)=(2,1,4p-1).
\]
This scalar input does not identify the two source-zero fan
returns in (RG.9).

\par

\medskip\noindent\textbf{RG.6. Full comparison with the marked infinity source}

\par\nobreak

\item This paragraph supplies the punctured-family arrow to the infinity
presentation IF.1--6. It uses that presentation's full reference
overlap, defined by transport of the global reference at one
nonzero epsilon. It does not identify a central collision attachment
or infer the characteristic value of its boundary.

\item The source, section and actual selected height map are
\[
 v=\epsilon^{-2}-1,\quad
 X=\epsilon^{-2}x,\quad Y=\epsilon^{-3}y,\quad
 T=\epsilon^{-4}t,\quad
 g=\epsilon(1-\epsilon^2),\quad
 q_{1,\rm old}=\epsilon^{-4}q_{*,\rm new}.
                                                        \tag{RG.12}
\]
They give exactly
\[
 y^2=x^3+(-2+3\epsilon^2)x^2+tx+
              \epsilon^2(1-\epsilon^2)^2,\quad G=(0,g),
\qquad
 t_*=(1-\epsilon^2)(1-3\epsilon^2).
\]
The selected normalization coordinate changes by
\(s_{\rm new}=\epsilon s_{\rm old}\). Therefore
\[
 Z_{\rm old}=\frac{s_{\rm old}-\rho}{s_{\rm old}+\rho}
            =\frac{s_{\rm new}-1}{s_{\rm new}+1}
            =Z_{\rm new}.
                                                        \tag{RG.13}
\]
The Tate period \(Q\) and marked \(P\) agree literally.
The physical-height units satisfy
\[
 \kappa_{\rm old}=\epsilon^4\kappa_{\rm new},\qquad
 P=-\frac{1+\epsilon}{1-\epsilon},\qquad
 \left.\kappa_{\rm new}\right|_{q_*=0}
                       =-\frac{1-\epsilon^2}{64}.
\]
Every nonconstant unit in these equations is retained. The
integer epsilon power is a toric unit on the punctured parameter
annulus, with its winding; the selected fan is unchanged by
(RG.13).

\item The Cartier map is the canonical divisor map, not an identification
only of generic Picard classes. The primitive lift satisfies
\[
 I_\epsilon U_{a,\rm old}I_\epsilon^{-1}
                  =\epsilon^{2a}U_{a,\rm new}.
                                                        \tag{RG.14}
\]
Since \(ab=2\), its \(b\)-th power cancels exactly the
\(\epsilon^{-4}\) in the selected height. Thus
\[
 I_\epsilon\bigl(2(-1)^{b+1}\lambda
             q_{1,\rm old}U_{a,\rm old}^{\,b}\bigr)
 I_\epsilon^{-1}
       =2(-1)^{b+1}\lambda q_{*,\rm new}U_{a,\rm new}^{\,b}.
                                                        \tag{RG.15}
\]
The Hom section itself provides a second literal check. With
\[
 F_{21}=\frac{Y+v+TX/(2v)}X,\qquad
 F_{12}=\frac{Y-v-TX/(2v)}X
\]
in the old presentation, substitution gives
\(F_{ab,\rm old}=\epsilon^{-1}F_{ab,\rm new}\) and
\(v=\epsilon^{-3}g\). Hence
\[
       \frac{q_{1,\rm old}}{vF_{ab,\rm old}}
               =\frac{q_{*,\rm new}}{gF_{ab,\rm new}}. \tag{RG.16}
\]
This retains the actual degree-two Cartier correction as well.
In particular the scalar in (RG.10) is unchanged.

\item On the logarithmic epsilon cover, the cap comparison is
\[
 z_{\rm old}=\epsilon^{4/m}z_{\rm new}.
                                                        \tag{RG.17}
\]
For \(m=4p\), its minimal source coordinates \(X',Y'\) are
unchanged. For \(m=4p-1\), the old-to-new arrow multiplies
them by \(\epsilon^{-2/m},\epsilon^{-3/m}\). It carries the
marked section, source resolution ideals and every Cartier
correction by their actual coordinate maps. Its full positive
epsilon endpoint is
\[
 D_{\rm new}^{-4}:\quad
 (z,X',Y')\mapsto
 (\zeta_m^{-4}z,\zeta_m^{-2}X',\zeta_m^{-3}Y')
                                                        \tag{RG.18}
\]
in the reduced case; the last two factors are one in the good
case. The ordinary and cap source arrows commute with the
specified finite root maps on each parameter chart.

\item Here is the completion of that endpoint on the affine line
atlas. Write \(A_0=L D\), with \(L\) covering the specified
signed division translation. The central finite commutator
character is one: in the good root the division point is \(O\);
in the reduced root the actual line
\(\mathcal O(2G'-O)\) restricts to the trivial line on the
compact reduced III fibre. Its commuting source maps have
constant line multipliers there. The finite eigenspace
average, retaining its actual conjugating unit, therefore
gives \(DLD^{-1}=L\). In this actual atlas
\[
                         A_0^4D^{-4}=L^4.
                                                        \tag{RG.19}
\]
For the actual fixed scalar \(r_0\), the full quotient
representative is \(A^4D^{-4}=S_{r_0^{4n}}L^4\).
P3K.92, equations (IFan.18a--b), supplies its exact
endpoint-to-initial arrow on both ordinary and cap charts,
including the retained scalar path and its endpoints.
The latter extends and has identity central source translation
in both branches; after its central scalar is retained it has
the small extending Lie logarithm. For the reduced branch the
annular comparison still contains
\[
                    4\log R-\log(4R)=2\omega_1+2\omega_2.
                                                        \tag{RG.20}
\]
Here \(\omega_1,\omega_2\) are the actual transported period
sections, with central values \((1,i)\). The signed source part
for \(L=t_{-nR}\) is multiplied by \(-n\), and its line
part is the actual Cartier factor of automorphy. This entire
period is included in the total atlas cochain. It is not
replaced by the small logarithm or suppressed in (RG.19).

\item To specify the annular marking, fix one nonzero
\(\epsilon_*\), transport the full global ordinary-to-cap
map through (RG.12--17), and use that entire map as the
infinity reference \(F_*\). It includes the given integral
\(d\delta\) and all finite characters. The cap-overlap source
and Cartier family is smooth over the full epsilon disc:
its base annulus avoids the unselected collision and the
selected fibre. The actual root division branches, norm
units and their equivariant line lifts extend on the cap.
Their nonzero norm-ratio has the fixed central value \(+2\)
or \(-2\), including at \(\epsilon=0\).

\item Continue \(F_*\) by the parameter-disc exponential lifting
on every ordinary and finite groupoid overlap, and solve the
remaining vector cochains by the fixed source-base Cauchy
contours of RCZ.3 and RCZ.5. At the reference fibre the full
maps agree, so every integral overlap and cyclic-group
discrepancy is zero there. It is locally constant and hence
zero throughout the parameter disc. The nonzero doubled
period (RG.20) is kept when changing logarithmic parameter
charts; the two compared total cocycles carry the same
term. Before taking a Lie logarithm at the reduced root,
use the actual suspension power to remove its component
translation. Finite averaging is applied to vector
cochains, not to lattice components. The root line
exponentials extend by the actual resolved Cartier maps.
This constructs precisely IF.3's full \(F_0\), with its
reference, rather than choosing an unmarked overlap.

\item For clarity, the scalar closing defect is explicit. Put
\[
 t_\epsilon=\frac{\log\epsilon}{2\pi i},\qquad
 \ell(\epsilon)=\frac{\log(1-\epsilon^2)}{2\pi i}.
\]
Then \(t_v=-2t_\epsilon+\ell(\epsilon)\).
The old and new ordinary forward arrows for the positive
epsilon loop are \(q_{1,\rm old}^{-2nk}\) and
\(q_{*,\rm new}^{-2nk}\). Their ratio is
\(\epsilon^{8nk}\). The obligatory scalar cochain on the
logarithmic parameter cover is
\[
 c_0(t_\epsilon)=
       \exp\{-2\pi i\,4nk\,t_\epsilon(t_\epsilon-1)\},
 \qquad
 \frac{c_0(t_\epsilon+1)}{c_0(t_\epsilon)}
                              =\epsilon^{-8nk}.
                                                        \tag{RG.21}
\]
The unit \((1-\epsilon^2)^k\), its actual logarithm and
the corresponding cap unit \(h\) remain in the full
ordinary-to-cap square. Equation (RG.21) is a cochain in
that square, not by itself a single-valued global scalar.
The cap correction (RG.19), the source and line maps, and
the remaining vector cochain are carried simultaneously.

\item The two local family comparisons on angular parameter
charts now differ by an actual global extending
translation. The rational source-section and line
valuation argument of RCZ.4 applies to these marked
punctured E7 sources and identifies this map with a
holomorphic physical scalar. Factor that scalar on the
two-chart epsilon annulus by RCZ.4's literal
Cauchy--Green solution: for its overlap logarithm
\(\ell_c\), set \(b_+=\chi_-\ell_c\),
\(b_-=-\chi_+\ell_c\), solve
\(\bar\partial v=\bar\partial b_\pm\) using
\(1/(\pi(\epsilon-w))\), and use
\(e^{b_\pm-v}\). The two adjusted actual maps then glue
on both overlap components. This retains the scalar
winding, (RG.21), the full finite cocycle and (RG.20).

\item Thus on the punctured parameter annulus the global
infinity restriction is carried to IF's actual source,
Cartier, cap and \(F_0\) family with the literal scalar
\((\epsilon^{-2}-1)^N\lambda_0\). Its positive return
is \(I_{\mathcal T}(-2k)\) in (RG.11), using the inverse
closing arrows on both sides. This is an annular family
morphism; it asserts neither a central filling nor
any equality between the two source-zero fan returns.

\item Finally the reflected source descriptions match through
one actual period. IF's half-turn arrow is
\(\mathcal I_\epsilon(Z,V)=(1/Z,P^aV)\).
At the target \(P_{\rm target}=P^{-1}\), so
\[
        h_{\rm target}\circ\mathcal I_\epsilon
                  =(Q/Z,V)=\mathcal R_\rho .
                                                        \tag{RG.22}
\]
This is the exact quotient arrow, with its integral
height reindexing and all its unit factors. It is the
reflection used in (RG.6), rather than an additional
unmarked endpoint identification.

\par

\medskip\noindent\textbf{RG.7. The actual paired source-zero returns}

\par\nobreak

\item This conclusion has the initial fan \(\mathcal T_0\) and
its reflected fan \(J\mathcal T_0\) as its precise scope.
It retains the one III-star cap, its prescribed finite
character, the full \(F_0\), the norm sign \(n=\pm1\),
and the original scalar
\(\lambda=v^N\lambda_0\) with \(N=km\), \(k>0\).
No additional logarithmic disc is introduced at either
colliding finite site.

\item The actual infinity collision excision IE.1--6 of
P3K.57 computes the fixed-scalar positive epsilon
return in the marked IF source/cap family:
\[
 \begin{array}{c|c}
 (a,b,m)&I^{\rm fixed}_{\mathcal T}\\ \hline
 (1,2,4p)&0\\
 (2,1,4p-1)&-1 .
 \end{array}                                            \tag{RG.23}
\]
The two entries apply to each of the two stated fan
objects, with the actual positive physical orbit and
dual-to-physical cap transfer of that excision, made explicit in \cref{night:P3K.68}. Its local
Euler-line bordism and cap cancellation establish this
input independently of the groupoid relation.

\item To pass from (RG.23) to the actual global infinity loop,
use the family with independent root-scalar variable
constructed in RG.5. On a fixed small epsilon circle,
the unit \((1-\epsilon^2)^k\) in (RG.10) has the actual
extending logarithm. Multiplying its argument by
\(s\in[0,1]\) gives a based parameter homotopy that
retains the cap root and the full annular cochain.
The resulting loop has one positive epsilon turn and
exactly \(-2k\) positive \(r\)-turns. In the actual
two-parameter family its based loop is the product of
the fixed-scalar source loop and \(-2k\) scalar loops.
The endpoint squares are (RG.2--3) and (RG.21), with
the finite root and Cartier arrows already matched
by RG.6. Thus the relation is one of the actual
completed family returns, not a substitution into a
characteristic formula:
\[
             I_{\mathcal T}
                  =I^{\rm fixed}_{\mathcal T}
                         -2kA^{\rm scalar}_{\mathcal T}.
                                                        \tag{RG.24}
\]
K26 and the actual C2 characteristic comparison K62 give
\(A^{\rm scalar}=0\) in the good branch
and \(A^{\rm scalar}=3p-1\) in the reduced branch,
for these two actual fan objects. Therefore
\[
 \begin{array}{c|c|c}
 (a,b,m)&I_{\mathcal T}
                 &[Z_{\mathcal T}]+[Z_{J\mathcal T}]\\ \hline
 (1,2,4p)&0&1\\
 (2,1,4p-1)&-1-2k(3p-1)&2+2k(3p-1)
 \end{array}
 \qquad\text{in }\Theta_7\cong\mathbb Z/28.             \tag{RG.25}
\]
The last column is (RG.9), with its already proved
selected-cusp term \(-1\), actual zero unselected
term, and the two separate source-zero fan objects.
It is valid for both retained norm signs. The original
positive \(u\)-returns are separately the cubes of
these positive \(v\)-returns by K49; that is an actual
covering-family statement, not division of a class.

\item Equation (RG.25) determines only the paired sum.
Neither individual source-zero return, nor an equality
between those two fan returns, is assigned here.
The ordinary-product obstruction and the full
source-zero closing of RG.2a remain part of the
construction.

\item Root has audited RG.1--6 and independently checked
the two paired formulas. Fourfold has independently
audited RG.1--3, RG.5 and the complete RG.6 dictionary.
The source/cut/closing identities and the paired
arithmetic are executed in
\path{code/part3/source_global_groupoid.py}.

\end{enumerate}
\end{proof}




\begin{theorem}\label{night:P3K.59}
\StatementRecord{night:P3K.59}{false as stated; correction proved}{S6\_06\_BLIND: GPT-6 Astra (OpenAI, ChatGPT Pro mode, effort Pro), recorded re-derivation in the checkpoint dated 18 November 2026}{\cref{night:P3F.56}, \cref{night:P3K.23}, \cref{night:P3K.31}, \cref{night:P3K.51}, \cref{night:P3K.52}, \cref{night:P3K.54}, \cref{night:P3K.55}}Retain the actual positive selected-cusp source loop of
P3F.54 and P3F.56, with its fixed scalar, prescribed finite cap,
physical height, Cartier line and complete period units. For every
integer \(j\geq0\), consider the two specified periodic fans with bases
\[
 P_j=\begin{pmatrix}1&-j\\0&1\end{pmatrix},\qquad
 R_j=\begin{pmatrix}j+1&-j\\-1&1\end{pmatrix}.
\]
Their actual oriented clutched-sphere classes are
\[
 [\Phi_{P_j}]=-1-3j(j+1),\qquad
 [\Phi_{R_j}]=-7-24j(j+1)\quad\text{in }\Theta_7=\mathbb Z/28.
\]
On \(P_j\), the complete ordinary-product variation and ambient
stable-derivative invariants in the marked three-curve basis are
\[
 V_j=\begin{pmatrix}
 0&6(j+1)&-6\\
 6(j+1)&18j(j+1)&-6j\\
 -6&-6j&0
 \end{pmatrix},\qquad
 \delta_j=\begin{pmatrix}-12(j+1)\\-24j(j+1)-12\\12j\end{pmatrix}.
\]
The corresponding native relative pairs, in the order \(C_0,C_1,C_2\), are
\[
 (6j,-6(j+1)),\qquad
 (3j(j+1),-21j(j+1)-6),\qquad
 (-6j-6,6j).
\]
For each selected periodic flop used in these constructions, write
\(a=V_{ii}\), \(b=\delta_i\), and let \(i'\) be the new label of that
same flopped curve. The actual native Cox arrow and the carried
ordinary-product comparison give exactly
\[
 V'_{i'i'}=a,\qquad \delta'_{i'}=2a-b,\qquad
 s'=s+\frac{b-a}{2}\pmod {28}.
\]
The \(R_j\) return sequence has least positive period seven and exact
range \(\{1,13,17,21\}\). In particular \(R_1\) has class one.
The statement concerns these explicitly constructed families and their
selected updates; it assigns no return to an additional fan.
\end{theorem}
\begin{proof}
\begin{enumerate}[label=\arabic*.]
\item \par\medskip\noindent\textbf{CF.1. Three exact lattice and native-star arrows.}\par\nopagebreak[4]

\item Write the ordered fan basis as \(P=(v\ w)\). Its triangles are
\((0,v,w)\) and \((v,w,v+w)\), with their lattice translates.
In the first triangle use \((a,b,c)\), with
\(q=abc, Z=b,V=c\). The three curves have native first charts
\[
 C_0:(z,n_0,n_1)=(b,a,c),\qquad
 C_1:(z,n_0,n_1)=(c,a,b),\qquad
 C_2:(z,n_0,n_1)=(a,b,c).                              \tag{P3K.59.CF.1}
\]
Their native second charts are exactly GF.11 and satisfy
\(z'=z^{-1},n'_r=zn_r\), including every complete period unit.

\item The three periodic flips have the following actual new bases and
new labels for the flopped curve:
\[
 Q_0=\begin{pmatrix}1&-1\\0&1\end{pmatrix},\qquad
 Q_1=\begin{pmatrix}1&0\\-1&1\end{pmatrix},\qquad
 Q_2=\begin{pmatrix}1&1\\0&1\end{pmatrix},
 \qquad P'=PQ_i,\qquad i'=(2,2,0)_i.                  \tag{P3K.59.CF.2}
\]
For \(C_0\), the old quadrilateral is
\((0,v,w,w-v)\); its old diagonal is \((0,w)\), and its new
diagonal joins \(v\) to \(w-v\). The new two triangles are those
of \((v,w-v)\). For \(C_1\), use
\((0,w,v,v-w)\), giving the basis \((v-w,w)\).
For \(C_2\), the quadrilateral is \((0,v,w,v+w)\), and its
new diagonal is \((0,v+w)\). The second triangle of
\((v,v+w)\), translated by \(-v\), is \((0,w,v+w)\).
These are literal triangle equalities, repeated over the period
lattice. Each basis change and its inverse is integral.

\item The actual torus characters are, respectively,
\[
 (Z',V')=(ZV,V),\qquad (Z',V')=(Z,ZV),\qquad
 (Z',V')=(Z/V,V),\qquad q'=q.                         \tag{P3K.59.CF.3}
\]
Their new first-triangle coordinates are
\[
 (a/c,bc,c),\qquad (a/b,b,bc),\qquad (ac,b/c,c).        \tag{P3K.59.CF.4}
\]
Thus in the native first charts of the flopped curves, (P3K.59.CF.3)
is exactly the standard other-ruling arrow
\[
 (z,n_0,n_1)\longmapsto
  (\zeta,u,v)=(n_0/n_1,zn_1,n_1)                     \tag{P3K.59.CF.5}
\]
for \(i=0,2\). For \(i=1\) the last two coordinates are
\((v,u)\). In each case \(\zeta uv=zn_0n_1=q\).
The second new chart uses
\((\zeta',u',v')=(n_1/n_0,zn_0,n_0)\), with the same
constant interchange for \(i=1\). It has
\(\zeta'=\zeta^{-1},(u',v')=\zeta(u,v)\).
These formulas are the two quotient charts of the rank-one matrix
\[
 \begin{pmatrix}n_0&n_1\\zn_0&zn_1\end{pmatrix}.       \tag{P3K.59.CF.6}
\]
They identify the common complement and both native normal bundles
by the actual Cox construction. For the normal interchange at
\(i=1\), use the explicit path
\[
 S_t=\frac12
 \begin{pmatrix}
 1+e^{\pi it}&1-e^{\pi it}\\
 1-e^{\pi it}&1+e^{\pi it}
 \end{pmatrix},\qquad 0\le t\le1.                    \tag{P3K.59.CF.6a}
\]
Its eigenspaces are the two constant lines spanned by
\((1,1)\) and \((1,-1)\), with eigenvalues \(1,e^{\pi it}\).
Thus \(S_t^*S_t=1\), \(S_0=1\), and
\(S_1=(\begin{smallmatrix}0&1\\1&0\end{smallmatrix})\).
It commutes with the transition \(z\,1\) and therefore acts on
the whole native normal bundle. It is independent of the source
loop parameter. Its oriented real lift conjugates the two
Spin factors through constant conjugations, which preserve
both integral third degrees.

\item Here are the full unit maps. Write the old two periods as GF.2,
with unit matrix \(A=(A_{rs})\). The new generators are the
columns of \(Q_i\) in that period basis. The actual new units are
\[
 A'_{rs}=\prod_{h,j}A_{hj}^{(Q_i^{-1})_{rh}(Q_i)_{js}}. \tag{P3K.59.CF.7}
\]
In the three cases their matrices are
\[
 A'_0=\begin{pmatrix}
 A_{11}A_{21}&A_{12}A_{22}/(A_{11}A_{21})\\
 A_{21}&A_{22}/A_{21}
 \end{pmatrix},
\]
\[
 A'_1=\begin{pmatrix}
 A_{11}/A_{12}&A_{12}\\
 A_{11}A_{21}/(A_{12}A_{22})&A_{12}A_{22}
 \end{pmatrix},
\]
\[
 A'_2=\begin{pmatrix}
 A_{11}/A_{21}&A_{11}A_{12}/(A_{21}A_{22})\\
 A_{21}&A_{21}A_{22}
 \end{pmatrix}.                                      \tag{P3K.59.CF.8}
\]
Substitution in the new GF.11 second charts gives the second
native chart of (P3K.59.CF.5) exactly. For example in case \(i=2\),
the second vertex coordinates before applying GF.11 are
\[
 (u',v',w')=\bigl(A_{11}b,ab/A_{21},A_{21}c/(A_{11}b)\bigr).
\]
GF.11 sends these to \((c/b,ab,b)\), precisely the second
chart of (P3K.59.CF.5). For \(i=0,1\), the new curve is \(C_2\),
and the second native charts are directly
\((c/a,ab,a)\) and \((b/a,a,ac)\), respectively.
These equalities retain all unit factors; their parameter winding
is included in (P3K.59.CF.7). In particular
\(B'=Q_i^{-1}BQ_i\), without replacing a unit by its winding.

\item The characters (P3K.59.CF.3), inverse characters and full period changes
give the same original source and line maps on their dense overlap.
The physical height is literally \(q\). The retained finite cap
and source exterior therefore use their original maps there.
Every periodic curve orbit has the same free representative charts
as K55, so (P3K.59.CF.5)--\allowbreak (P3K.59.CF.8) descend to its whole periodic flop.
The unrestricted four-weight Cox identity of K54 intertwines
the physical circle with these maps. Transporting the Hermitian
metric and applying the positive radial normalization gives its
actual common-link arrow; the normal rescaling of K54 carries
the full nonlinear units. This supplies the actual family flop,
including its native normal conventions.

\par\medskip\noindent\textbf{CF.2. Exact selected-coordinate update.}\par\nopagebreak[4]

\item Let the old ordinary relative pair at \(C_i\) be \((m_i,n_i)\),
and let its central coordinate be the actual sphere return \(s\).
K55's explicit cylinder comparison removes the pure-left
\(L^6\), giving the native pair \((m_i-6,n_i)\).
The actual common-link arrow (P3K.59.CF.5), in the K23 Hopf coordinates,
is \(\tau(h,q_H)=(q_H^{-1}hq_H,q_H)\), with the exterior
extension \(T(y,q_H)=(q_H^{-1}yq_H,q_H)\). The complex-normal
reset is exactly K54's fixed \(J\). Applying their proved
relative maps gives the new native pair and central coordinate
\[
 (m_i-6+2n_i,-n_i),\qquad s'=s+n_i.                   \tag{P3K.59.CF.9}
\]
The constant normal interchange in (P3K.59.CF.5) acts by a constant
oriented four-plane change. Its lift conjugates the two Spin
factors separately, through the displayed \(U(2)\) path, and
preserves their degrees. Thus the same formula applies to
all three cases of (P3K.59.CF.2).

\item Add the actual new \(L^6\) comparison on the new native star.
Its ordinary pair is \((m_i+2n_i,-n_i)\). In terms of the
actual ordinary variation and ambient derivative, put
\(a=V_{ii}\), \(b=\delta_i\). K31.7's carried single-star
frame map gives
\[
 a=-(m_i+n_i),\qquad b=2n_i+a.
\]
Consequently the new selected entries and the actual return obey
\[
 V'_{i'i'}=a,\qquad \delta'_{i'}=2a-b,\qquad
 s'=s+\frac{b-a}{2}\pmod {28}.                       \tag{P3K.59.CF.10}
\]
This derivation uses only the actual selected native star on
each side. It does not use a simultaneous-product matrix reset.

\par\medskip\noindent\textbf{CF.3. Recovering the other entries from the actual trace maps.}\par\nopagebreak[4]

\item For every target with nonzero tangent weights \(w_0,w_1,w_2\),
write \(r_i=1/w_i\). The exact GF trace identities and symmetry
have the complete solutions
\[
 V=\operatorname{diag}(I_0/w_0,I_1/w_1,I_2/w_2)+c\,rr^T,
 \qquad \delta_j=(W_j+e)/w_j.                         \tag{P3K.59.CF.11}
\]
Indeed the displayed diagonal is a particular solution. The
two vectors \(w_je_j-w_0e_0\) span the kernel of the functional
with coefficients \(r\). A symmetric matrix annihilating that
plane has image in its orthogonal line, and is exactly
\(c rr^T\). The same argument for the derivative leaves
the single constant \(e\). All these statements are rational
linear algebra applied to the actual integral identities.

\item At the selected new curve the already proved entries (P3K.59.CF.10)
therefore fix
\[
 c=w_{i'}^2a-w_{i'}I_{i'},\qquad
 e=w_{i'}(2a-b)-W_{i'}.                               \tag{P3K.59.CF.12}
\]
Equations (P3K.59.CF.11)--\allowbreak (P3K.59.CF.12) give the entire actual new matrix
and derivative. Their values are integral because they equal
the actual integral invariants uniquely, whose existence is
the explicit family construction above.

\item For completeness a target has at most one zero tangent weight.
For the zero index \(r\), the other two weights are \(+1,-1\),
by primitivity. The same trace equations give
\[
 \begin{aligned}
 V_{jj}&=I_j/w_j,\qquad V_{rj}=V_{jr}=-I_r/w_j,\\
 V_{pq}&=0\qquad(p,q\ne r,\ p\ne q),\\
 \delta_j&=(W_j-W_r)/w_j\qquad(j\ne r).
 \end{aligned}                                      \tag{P3K.59.CF.13}
\]
Their two remaining entries are precisely \(V_{rr},\delta_r\).
For a flop whose selected target is that fixed curve, (P3K.59.CF.10)
supplies both. A target with a different retained fixed curve
requires its own actual fixed-star data. This describes the
exact rank of the trace map and leaves that specific data open.

\par\medskip\noindent\textbf{CF.4. A complete horizontal staircase.}\par\nopagebreak[4]

\item For every integer \(j\ge0\), set
\[
 P_j=\begin{pmatrix}1&-j\\0&1\end{pmatrix}.             \tag{P3K.59.CF.14}
\]
The actual return and complete ordinary invariants on this fan are
\[
 s_j=-1-3j(j+1)\pmod {28},
\]
\[
 V_j=\begin{pmatrix}
 0&6(j+1)&-6\\
 6(j+1)&18j(j+1)&-6j\\
 -6&-6j&0
 \end{pmatrix},\qquad
 \delta_j=\begin{pmatrix}
 -12(j+1)\\-24j(j+1)-12\\12j
 \end{pmatrix}.                                      \tag{P3K.59.CF.15}
\]
The native pairs on its three curves are
\[
 (6j,-6(j+1)),\quad
 (3j(j+1),-21j(j+1)-6),\quad
 (-6j-6,6j).                                         \tag{P3K.59.CF.16}
\]

\item Here is an induction by actual family maps. At \(j=0\), (P3K.59.CF.15)
is exactly K51/K52 and the return is F56's class one. GF.10a
supplies the explicit based homotopy identifying those initial
products with K55's chosen product. The \(C_0\) flip of \(P_j\)
is precisely \(P_jQ_0=P_{j+1}\); its selected target is \(C_2\).
For \(P_j\) the literal complete windings and trace degrees are
\[
 B_j=\begin{pmatrix}-6&6j\\0&0\end{pmatrix},\quad
 l=(0,-6(j+1),0),\quad w=(j,1,-j-1),
\]
\[
 I=(-6(j+1),12j(j+1),6j),\qquad
 W=(6,-12j(j+1)-6,6).                                \tag{P3K.59.CF.17}
\]
These follow by substituting the actual (P3K.59.CF.3)/(P3K.59.CF.7) maps in GF.
The target weights \((j+1,1,-j-2)\) are all nonzero.
Equation (P3K.59.CF.10) sets its selected entries to
\(V'_{22}=0,\delta'_2=12(j+1)\). The unique solution
(P3K.59.CF.11)--\allowbreak (P3K.59.CF.12) is (P3K.59.CF.15) with \(j\) replaced by \(j+1\).
Thus the full matrix and derivative have been obtained from the
actual target, including their off-diagonal entries.
The old right integer is \(-6(j+1)\), so
\[
 s_{j+1}=s_j-6(j+1).
\]
Summing this actual consecutive-flop identity gives (P3K.59.CF.15).
Finally the native comparison gives
\(n_i=(\delta_i-V_{ii})/2\) and
\(m_i=-(\delta_i+V_{ii})/2-6\), which proves (P3K.59.CF.16).

\par\medskip\noindent\textbf{CF.5. The adjacent complete family.}\par\nopagebreak[4]

\item Flop \(C_1\) on each of the just constructed actual families.
Its new basis is
\[
 R_j=P_jQ_1=\begin{pmatrix}j+1&-j\\-1&1\end{pmatrix},
 \qquad j\ge0.                                       \tag{P3K.59.CF.18}
\]
Equation (P3K.59.CF.16) gives its actual native right integer.
The same Cox family arrow (P3K.59.CF.5)--\allowbreak (P3K.59.CF.10) therefore gives
\[
 [\Phi_{R_j}]=-1-3j(j+1)-21j(j+1)-6
             =-7-24j(j+1)\pmod {28}.                 \tag{P3K.59.CF.19}
\]
This computes the actual return on every fan (P3K.59.CF.18), including
\(j=0\). Its sequence has period seven: modulo 28 it equals
\(21+4j(j+1)\), and shifting \(j\) by seven changes this
by \(56j+224\). The seven consecutive values are
\[
                         21,1,17,13,17,1,21 .
\]
They give the exact range \(\{1,13,17,21\}\); the sequence
is not constant, so seven is its least positive period.
In particular
\(R_1=(\begin{smallmatrix}2&-1\\-1&1\end{smallmatrix})\)
has the actual positive return class one.

\item All original physical source and cap maps are carried through
the specified finite sequence for each \(j\). This proof assigns
no value to a fan outside its computed domain. The exact algebra
and repeated-basis/fan consistency checks are retained under
\path{code/part3_root/source_cusp_mutation_diagnostic.*}; their
first unsimplified symbolic-matrix equality attempt is preserved.
The written Cox/period maps, native-frame comparison and induction
are the proof for all \(j\). The active finite search starts only
from the proved \(P_0=1\) seed: at depth seven it checks
210 basis states and 128 fan states, with no conflicting state
or failed integrality check. Four transitions to a target with
a separately retained fixed curve are left unevaluated.
The earlier exploratory two-seed output is retained separately;
its second fixed-star seed supplies no input to this proof.

\end{enumerate}
\end{proof}




\begin{theorem}\label{night:P3K.60}
\StatementRecord{night:P3K.60}{false as stated; correction proved}{S6\_06\_BLIND: GPT-6 Astra (OpenAI, ChatGPT Pro mode, effort Pro), recorded re-derivation in the checkpoint dated 18 November 2026}{\cref{night:P3K.23}, \cref{night:P3K.29}, \cref{night:P3K.31}, \cref{night:P3K.57}, \cref{night:P3K.58}}Retain P3K.58's actual marked source, section, scalar and selected height
\[
 Y^2=X^3+(1-2v)X^2+TX+v^2,\qquad G=(0,v),\qquad
 q=T-(v^2-2v),\qquad \lambda=v^{km}\lambda_0 .
\]
The one III-star cap has
\[
 (a,b,m)=(1,2,4p)\quad\text{or}\quad(2,1,4p-1),
 \qquad p,k\ge1,\qquad n=\pm1 .
\]
Keep its complete root, Cartier line, finite character, reference
annular map and integer cochain. No additional logarithmic disc
meets either colliding finite site. Let \(\mathcal T\) be the initial
physical-vertical-edge fan, and \(J\mathcal T\) its opposite-diagonal
reflection. The positive source-zero loop is the positive \(v\)-loop
with the literal scalar above.

In the single native diagonal-star product constructed below, its
actual relative return has free coordinates
\[
                              (m_{\rm rel},n_{\rm rel})=(1-a,a).
\]
The actual single-loop Cox flop therefore gives
\[
                   [Z_{J\mathcal T}]-[Z_{\mathcal T}]
                                  =a\quad\text{in }\Theta_7\cong\mathbb Z/28.
\]
The proof constructs the single native product, its literal pullback
to the double parameter cover, and a closed ordinary product of
the doubled three-curve plumbing. It identifies the actual relative
square before using the two torsion-free integer coordinates.

Together with P3K.58's proved paired sum and its P3K.57 dependency,
the pair of actual returns is constrained exactly to
\[
 \bigl([Z_{\mathcal T}],[Z_{J\mathcal T}]\bigr)
   =\bigl(kA_{\rm scalar}+14e,\ a+kA_{\rm scalar}+14e\bigr),
       \qquad e\in\{0,1\},
\]
where
\[
 A_{\rm scalar}=
 \begin{cases}
 0,&(a,b,m)=(1,2,4p),\\
 3p-1,&(a,b,m)=(2,1,4p-1).
 \end{cases}
\]
The same \(e\) occurs in the two entries of a given pair. This particular sum-and-difference argument does not select its value; the completed absolute calculation is \cref{night:P3K.88}. No class in \(\Theta_7\) is divided by two.
\end{theorem}
\begin{proof}
\begin{enumerate}[label=\arabic*.]
\item \par

\medskip\noindent\textbf{DZ.1. The literal parameter seam on all three edges}

\par\nobreak

\item Take a small positive \(v_*\), and first parameterize one positive
turn by \(v=v_*e^{2\pi i\tau}\), \(0\leq\tau\leq1\). Choose one
\(q\)-disc which excludes the other moving nodal value for this
compact parameter circle. Its full period maps are
\[
 h:(Z,V,q)\longmapsto(\kappa qZ,\alpha V,q),\qquad
 B:(Z,V,q)\longmapsto(gZ,\eta qV,q),                 \tag{DZ.1}
\]
where all four coefficients are the actual nowhere-zero holomorphic
functions of \(q\) and the cut parameter. Their windings on the
double loop, later parameterized by \(v=v_*e^{2i\theta}\), are
\[
 \begin{pmatrix}\operatorname{wind}\kappa&
                 \operatorname{wind}g\\
                 \operatorname{wind}\alpha&
                 \operatorname{wind}\eta\end{pmatrix}
       =2\begin{pmatrix}1&b\\-a&N-2\end{pmatrix}.    \tag{DZ.2}
\]
The functions themselves, rather than only these integers, are kept
below.

\item Put \(L_0=nk\). The actual forward parameter arrow for one turn is
\[
 R_{\rm ord}^{(1)}:(Z,V,q)\longmapsto(Z,q^{L_0}V,q). \tag{DZ.3}
\]
It translates the fan indices by \((0,L_0)\). Applying the actual
period \(B^{-L_0}\) returns each chosen cone representative to its
original index and leaves the bounded map
\[
 E_v:(Z,V,q)\longmapsto(g(q,v)^{-L_0}Z,
                         \eta(q,v)^{-L_0}V,q).     \tag{DZ.4}
\]
This is an equality of quotient arrows, not a substitution of a
bounded multiplier for the original period.

\item There are holomorphic-in-\(q\), smooth-in-\(\tau\) logarithms
\(G=\log g\), \(H=\log\eta\) on the cut disc times interval. Choose a
smooth real function \(\chi\), zero on an initial parameter collar and
one on a terminal parameter collar. Define the actual local maps
\[
 C^{(1)}_\tau:(Z,V,q)\longmapsto
       (A_\tau Z,D_\tau V,q),\qquad
 A_\tau=e^{L_0\chi G},\quad D_\tau=e^{L_0\chi H}.    \tag{DZ.5}
\]
They have explicit inverses obtained by negating the exponents.
On the terminal collar they are exactly
\((Z,V)\mapsto(g^{L_0}Z,\eta^{L_0}V)\), independently of the chosen logarithm
branches. Thus, on the two parameter collars identified by the
original forward arrow,
\[
                     C_{\rm end}E_v=C_{\rm start}=1. \tag{DZ.6}
\]
This equality holds as a map in \(q\) and the collar parameter, so
all parameter and spatial derivatives match at the seam. Using only
the reference values \(g_*,\eta_*\) would not supply this collar
identity; (DZ.5) uses the varying functions.

\item Every multiplier in (DZ.1) and (DZ.5) depends only on the unchanged
coordinates \(q,v\). Consequently \(C^{(1)}_\tau\) commutes literally
with \(h,B\), and with the physical circle. It therefore descends
to the actual quotient neighbourhood. On the two triangle charts its
maps are
\[
 \begin{split}
 T^+:(x_0,x_1,x_2)&\longmapsto
       ((AD)^{-1}x_0,Ax_1,Dx_2),\\
 T^-:(y_0,y_1,y_2)&\longmapsto
       (ADy_0,D^{-1}y_1,A^{-1}y_2).
 \end{split}                                          \tag{DZ.7}
\]
Their three-coordinate products remain \(q\), so the maps extend
through both vertices and every toric boundary. In the six named
charts \((a_0,b_0,c_0)\), \((u_0,v_0,w_0)\), the three transitions are
\[
 \begin{array}{c|ccc}
 &u_0&v_0&w_0\\ \hline
 C_0&\kappa\alpha b_0c_0&\alpha^{-1}a_0b_0&
                                      \kappa^{-1}b_0^{-1}\\
 C_1&g\eta b_0c_0&\eta^{-1}c_0^{-1}&g^{-1}a_0c_0\\
 C_2&a_0^{-1}&a_0b_0&a_0c_0 .
 \end{array}                                          \tag{DZ.8}
\]
Substitution of (DZ.7) on either side gives the same factors
\((AD,D^{-1},A^{-1})\) in all three rows. This verifies the common
edge and corner maps, retaining all four complete units. In
particular (DZ.5) changes none of the coefficients in (DZ.8).
The vertex charts defined by (DZ.5) are now periodic by (DZ.6).
In particular the diagonal native star is already a literal product
for this single turn: its two charts satisfy
\(z'=z^{-1},c'_i=zc_i\), the \(C_2\) row of (DZ.8), with no
parameter-dependent transition. This assertion uses no
contraction of the other two curve tubes.

\item We now specify, rather than independently choose, the doubled
seam. On \(0\leq u\leq2\), write \(u=j+\tau\), \(j=0,1\), and put
\[
              C^{(2)}_{j+\tau}
                    =C^{(1)}_\tau E_{v(\tau)}^{-j}. \tag{DZ.8a}
\]
Near the intermediate seam the two formulas both equal \(E_v^{-1}\);
they agree on collars with every derivative by (DZ.6). At the last
seam the product is \(E_v^{-2}\), which cancels the doubled
forward arrow after the literal period \(B^{-2L_0}\).
Every factor in (DZ.8a) commutes with the full periods, the physical
circle, and all rows of (DZ.8). Thus these closed doubled charts
retain exactly (DZ.2). The double native-star product so obtained
is precisely the pullback of the single product under
\(u\bmod2\mapsto u\bmod1\): the factor \(E_v^{-j}\) is the
transition from the \(j\)-th raw sheet to its first-sheet
representative. There is no independently selected second-loop
normal framing.

\item Put \(L=2L_0=2nk\). The double cap retains the actual forward map
\(h_\infty^L\) of RG.2,
where \(h_\infty=1-T_0(v)q_\infty\), and the full root/norm
cochain. Its overlap square remains
\[
                   F_{\rm end}R_{\rm ord}
                         =R_{\rm cap}F_{\rm start}.  \tag{DZ.9}
\]
The local arrows (DZ.5) and (DZ.8a) are coordinate maps on the cusp
neighbourhood, not replacement global source or cap maps. To choose
parallel transport compatible with their boundary product, extend
its time-dependent velocity from a collar using a cutoff in a
slightly larger cusp neighbourhood. This is a smooth vector field
on the already constructed compact family and has a complete flow
for the compact parameter interval. Outside that neighbourhood use
the original family connection and its literal cap arrow.
The two vector fields are joined in their affine space. This
constructs actual global transport with the prescribed cusp
boundary; its endpoint is the physical return formed with the
inverse of the unchanged global closing arrow.

\par

\medskip\noindent\textbf{DZ.2. A closed ordinary plumbing product}

\par\nobreak

\item Work in the doubled closed charts of DZ.1, and suppress their
tildes only after the explicit pullback map (DZ.8a). At the second vertex use
\((u_0,\eta v_0,\kappa w_0)\), with the same initial constant
transformation at the reference fibre. The ordered normal matrices
are literally
\[
 D_0=\operatorname{diag}(\eta/\alpha,\kappa\alpha),\quad
 D_1=\operatorname{diag}(\kappa/g,g\eta),\quad
 D_2=\operatorname{diag}(\eta,\kappa).                \tag{DZ.10}
\]
Their common determinant winding is \(2t\), \(t=N-1\), and their
first diagonal windings are
\[
 l_0=2(N+a-2),\qquad l_1=2(1-b),\qquad
 l_2=2(N-2).                                        \tag{DZ.11}
\]
The based ratio of each diagonal entry is its displayed angular
factor times the exponential of its actual based periodic
logarithm. This follows by subtracting its winding from the full
nonzero loop, on the same \(q\)-disc.

\item For every integer \(d\) use the actual based quaternion contraction
\[
 H_d(s,\theta)=
  (\sin\vartheta(s)e^{id\theta}+j\cos\vartheta(s))
  (\sin\vartheta(s)+j\cos\vartheta(s))^{-1},          \tag{DZ.12}
\]
where \(\vartheta\) is smooth, constant at \(0,\pi/2\) near the two
ends. Thus \(H_d(0,\theta)=1\), \(H_d(1,\theta)=e^{id\theta}\), and
both parameter endpoints are one. For
\(Q=n_0+j\overline{n_1}\), set
\[
              T_i(Q)=H_t Q H_{l_i-t}.               \tag{DZ.13}
\]
At its second radial end this has the two complex factors
\((e^{il_i\theta},e^{i(2t-l_i)\theta})\), exactly the angular part
of \(D_{i,*}^{-1}D_i\). Multiply by the diagonal exponentials of
the retained periodic unit logarithms, with a radial cutoff zero
at the first end and one at the second. The resulting
\(\mathcal K_i\) is based and periodic and satisfies
\[
 \mathcal K_i(0,\theta)=1,\quad
 \mathcal K_i(1,\theta)=D_{i,*}^{-1}D_i,\quad
                   zD_{i,*}\mathcal K_i(1,\theta)=zD_i. \tag{DZ.14}
\]
On the tube of the \(i\)-th curve use
\((z,n)\mapsto(z,\mathcal K_i(s(|z|),\theta)n)\).
The last equality is its exact second-vertex collar equality;
near the first vertex it is the identity. The three curve interiors
are disjoint, so the three maps glue to the two vertex maps.

\item For the original, nonlinear germs retain the full normal rescaling
\(R_\tau(z,n)=(z,\tau n)\). The conjugation
\(R_\tau^{-1}FR_\tau\) extends smoothly to \(\tau=0\) with the
zero-section map and its invertible normal derivative. A common
small tube works on the compact parameter interval. Integrating
the cutoff velocity of this isotopy restores the actual higher
normal terms, with the vertex collars protected. These isotopies
are periodic in the closed charts, and the maps (DZ.8a) carry them
back to the original family with its seam (DZ.6).
Finally carry the rounded boundary by the same collar isotopy.
This is an actual based oriented ordinary product of the full
three-curve plumbing, not merely of its normal bundles.

\par

\medskip\noindent\textbf{DZ.3. Actual circle traces, derivative and the diagonal right integer}

\par\nobreak

\item The maps (DZ.8a) commute with the physical circle. Therefore that
circle has precisely the original vertex weights
\((-1,0,1)\), \((1,-1,0)\) in the closed charts. In the ordered
curve normal coordinates its tangent weights and normal differences
are
\[
                  w=(0,1,-1),\qquad d=(-2,-1,-1).
\]
The ordinary product (DZ.13) carries its normal action to
\[
 Q\longmapsto
 (H_t e^{-iw_i\phi/2}H_t^{-1})Q
 (H_{l_i-t}^{-1}e^{id_i\phi/2}H_{l_i-t}).            \tag{DZ.15}
\]
The retained diagonal unit exponentials conjugate this literal
formula and give a based homotopy back to it, with the normal metric
transported. Normal rescaling carries the higher derivatives as in
DZ.2. Both parameter ends and every radial-end physical orbit are
closed; the vertex orbit maps are independent of the parameter.

\item Here is the actual degree calculation with its arrows. On the
positive source path put \(U=\cos\chi+j\sin\chi\) and
\(E_i=e^{iw_i\phi/2}\). The moved source Hopf representative is
\(E_i^{-1}UE_i\). At the normal section \(q_H=1\), the trace is
\[
 E_i^{-1}UE_i H_t E_i^{-1}H_t^{-1}U^{-1}
              H_{l_i-t}^{-1}e^{id_i\phi/2}H_{l_i-t}. \tag{DZ.16}
\]
For (DZ.12), direct differentiation at the inverse images of
\(-i\) gives
\(\deg(H_diH_d^{-1})=-d\) and
\(\deg(H_d^{-1}iH_d)=d\), in the path--parameter orientation.
The factor \(UiU^{-1}\) has degree zero. The positive quaternion
volume gives degree \(c e\) for an axis of degree \(e\) and angle
\(c\phi/2\). The exact Maurer product boundary terms vanish on
the radial ends and the constant central-sign angle ends; the
parameter faces cancel. Thus (DZ.16) has normal-meridian degree
\[
                    I_i=t w_i+d_i(l_i-t).            \tag{DZ.17}
\]
For the derivative use the universal covering group of stable
\(GL^+\), with its integral third class identified by the lifted
polar deformation. The chain-rule identity is retained in that
group. A left quaternion factor contributes its degree and a right
factor its negative. The base rotation has only the \(\phi\)
variable. At \(2\pi\) the base and normal Spin endpoints are both
\((-1)^{w_i}\), so their product is one. The resulting full
closed stable derivative has degree
\[
                    W_i=t w_i-d_i(l_i-t).            \tag{DZ.18}
\]
These computations use the same actual endpoint caps as CT.3:
take generic normal sections with all vertex coordinates nonzero,
connect them by paths interpolating squared moduli and phases in
the vertex links, and cap each parameter-independent end by
\(D^2_\theta\times S^1_\phi\), constant in its disc variable.
The circle changes only phases and preserves the axis-tube
inequalities. The vertex traces factor through dimension two.
Thus no extra vertex meridian or third derivative class occurs.
The change from the coordinate tangent frame to the ambient
stable frame extends through each annulus times its normal
four-disc, whose third cohomology is zero; its pairing on each
such capped trace is zero. This is the same exact frame arrow
used in ST.2, now on the closed charts constructed in DZ.1.

\item Let \(V\) be the actual ordinary-product variation, \(\delta\)
the actual ambient stable derivative, and \(D_i,m_i\) the
oriented relative cycles and normal meridians. For
\(\gamma_j=a_j-a_0=\partial x_j\), the positive-circle boundary
identity \(\partial\Delta=-\Delta\partial\) gives
\[
             \Delta x_j=w_jD_j-w_0D_0 .
\]
The literal action prism and the chain-rule/Stokes prism then give
\[
 V(w_je_j-w_0e_0)=I_je_j-I_0e_0,\qquad
 w_j\delta_j-w_0\delta_0=W_j-W_0.                    \tag{DZ.19}
\]
The return fixes the common boundary by the actual product in
DZ.2; its integral meridians generate \(H_3\), and \(H_2\) of the
common exterior is zero. These are exactly the groups in the
three-curve complement of K29, since the curves and their
neighbourhood have been carried by the explicit ordinary product.
They justify both prism identities as identities of the actual
return.

\item Substituting (DZ.11) gives
\[
 \begin{split}
 I&=(-2(N+2a-3),\,2N+2b-4,\,-2N+4),\\
 W&=(2(N+2a-3),\,2-2b,\,-2).
 \end{split}                                          \tag{DZ.20}
\]
In particular \(w_0=0,w_2=-1\) yield
\[
                  V_{22}=2N-4,\qquad
                  \delta_2=2N+4a-4.                 \tag{DZ.21}
\]
K31's actual ambient-to-single-star product frame map has third
class \(e_2\). Carrying that map by the based product above retains
its derivative correction \(-V_{22}\); it is not an identification
of two unmarked frames. Hence the actual ordinary relative
integers satisfy
\[
 2n_2=\delta_2-V_{22},\qquad m_2=-V_{22}-n_2,
 \quad\text{so}\quad
           (m_2,n_2)=(-2N+4-2a,\,2a).               \tag{DZ.22}
\]

\par

\medskip\noindent\textbf{DZ.4. The native Cox product and the actual flop}

\par\nobreak

\item The diagonal native star in the closed charts is literally
\[
 (z,c_0,c_1)=(a_0,b_0,c_0),\qquad
 (z',c'_0,c'_1)=(u_0,v_0,w_0),\qquad
                  z'=z^{-1},\quad c'_i=zc_i.
\]
Its Hopf coordinate is
\(q_H=U(z)(c_0+j\overline{c_1})\).
The native-to-ordinary boundary arrow is, with the full unit
homotopy retained,
\[
 q_H\longmapsto
       (U H_t U^{-1})q_H H_{l_2-t}.                 \tag{DZ.23}
\]
Its left factor has degree \(2t\) on the Hopf-oriented sphere
times the parameter, by the exact Maurer boundary integral of
OP.21--23 with exponent \(t\). Its right factor depends on only
the radial coordinate and parameter and has third degree zero.
Extending the inverse factors over \(D^3\times I\) by the
literal cylinder retraction of OR.2 gives endpoint degrees
\((-2t,0)\). Consequently the ordinary relative return is the
native return followed by the actual pure-left map \(L^{-2t}\).
There is no disc-supported factor: these coefficient maps extend
over \(D^3\times D^4\). In particular the native right integer is
unchanged, and the complete native pair is
\[
                        (m_2^{\rm nat},n_2^{\rm nat})
                                      =(2-2a,\,2a). \tag{DZ.24}
\]

\item The single and doubled seam maps themselves also extend through
the flop. In Cox
coordinates with auxiliary weights \((1,1,-1,-1)\), take the
chart \(z=a_1/a_0,c_i^{\rm chart}=a_0c_i\).
For the multipliers \(A,D\) in (DZ.5) or (DZ.8a), the literal Cox lift is
\[
 (a_0,a_1,c_0,c_1)\longmapsto
       (a_0,(AD)^{-1}a_1,Ac_0,Dc_1).               \tag{DZ.25}
\]
Here the actual height is \(q=a_0a_1c_0c_1\), which this map
preserves. Thus even for the full \(q\)-dependent multipliers its
inverse is obtained by replacing \(A,D\) by \(A^{-1},D^{-1}\)
at the same \(q\); no implicit inverse or constant-unit
replacement is needed.
It commutes with the auxiliary action, the physical circle and
all period translates. It gives (DZ.7) on the original star and
extends to the other small resolution by its other quotient
chart. Transport the Hermitian metric by this map, then use its
positive radial normalization on the common link. Thus the
two native product links and their parameter seams are carried
by the same actual Cox map, with all its nonunit moduli retained.
On each free period representative this is a literal map;
equivariance extends it to the whole periodic orbit of curves.

\item The common-link arrow and its oriented exterior extension are
therefore the same explicit maps
\[
 \tau(h,q_H)=(q_H^{-1}hq_H,q_H),\qquad
 T(y,q_H)=(q_H^{-1}yq_H,q_H) .
\]
The unrestricted four-weight identity of FF.1/SF.8 intertwines
them with (DZ.25), and normal rescaling carries the full
\(q\)-dependent unit germs on both small resolutions. The
physical cap and exterior family are unchanged outside these
matched cusp charts. The diagonal flip gives exactly
\(J\mathcal T\) from \(\mathcal T\), including the integral
period indexing.

\item We first identify the actual relative square, before applying the
flop formula. Let \(f\) be the positive one-turn return on the
exterior of the single native star, using precisely the single
product following (DZ.8). Choose a global family connection
compatible with this product on its boundary. Its literal pullback
to the degree-two parameter cover has relative return \(f^2\):
the first sheet has endpoint \(f\), and the second uses the same
connection and boundary product under the transition \(E_v^{-1}\)
in (DZ.8a). The full forward arrows compose as
\((R_{\rm ord}^{(1)})^2=R_{\rm ord}^{(2)}\) and likewise on the
cap, by RG.3. Thus the closed parameter and cap square is exact.

\item The ordinary-product calculation (DZ.22), followed by the inverse
of the actual boundary change (DZ.23), computes this same
pulled-back native return. Indeed both resulting connections use
the literal native boundary product (DZ.8a); their interpolation
with that boundary fixed supplies a relative isotopy of endpoints.
Hence the native pair (DZ.24) is the free pair of \(f^2\).
K23's relative generators \(L,R\) commute literally, and its
disc-supported central factor is central. The free-coordinate
homomorphism to \(\mathbb Z^2\) consequently sends
\(f^2\) to twice the free coordinates of \(f\).
Multiplication by two is injective on each of these two integral
coordinates. We therefore obtain the actual single native pair
\[
                         (m(f),n(f))=(1-a,a).       \tag{DZ.26}
\]
This argument takes an actual square in the relative group and
uses its torsion-free coordinates. It does not divide a class in
\(\Theta_7\).

\item K23's already proved actual relative conjugation for the same
single-loop Cox map is
\[
 T(L^mR^n\mathcal C^s)T^{-1}
          =L^{-n}R^{m+2n}\mathcal C^{s+n}.
\]
Here the native central coordinate is the actual single return
\(s=[Z_{\mathcal T}]\), and the single Cox family arrow is the
one carried by (DZ.5) and (DZ.25). Applying the identity to
(DZ.26) proves
\[
                 [Z_{J\mathcal T}]-[Z_{\mathcal T}]
                         =a\quad\text{in }\mathbb Z/28. \tag{DZ.27}
\]
Every source/cap/scalar datum is the same one as RG.25. This
relation is derived from the actual single-loop relative
conjugation. Combining it with the already proved paired sum
still leaves a two-torsion ambiguity for either individual
source-zero class; neither is selected here.

\par

\medskip\noindent\textbf{DZ.5. The paired returns with their exact remaining ambiguity}

\par\nobreak

\item Keep the actual paired sum of P3K.58, with its P3K.57 infinity
comparison. Write
\[
 A_{\rm scalar}=0\quad(a,b,m)=(1,2,4p),\qquad
 A_{\rm scalar}=3p-1\quad(a,b,m)=(2,1,4p-1).
\]
In both rows P3K.58 is exactly
\[
       [Z_{\mathcal T}]+[Z_{J\mathcal T}]
                          =a+2kA_{\rm scalar}.
\]
Together with (DZ.27) it implies
\[
             2\bigl([Z_{\mathcal T}]-kA_{\rm scalar}\bigr)=0
                              \quad\text{in }\mathbb Z/28.
\]
The kernel of multiplication by two in this specified group is
exactly \(\{0,14\}\): for an integer representative \(r\),
\(2r\equiv0\pmod{28}\) is equivalent to \(r\equiv0\pmod{14}\).
Consequently the actual pair has precisely the form
\[
 \boxed{\;
 ([Z_{\mathcal T}],[Z_{J\mathcal T}])
   =(kA_{\rm scalar}+14e,\ a+kA_{\rm scalar}+14e),
       \qquad e\in\{0,1\}.\;}                       \tag{DZ.28}
\]
This is a determination of the two-element set allowed by the
proved sum and the actual single-loop flop morphism. No value of
\(e\) is selected. The common \(e\) belongs to the particular
retained family data; this proof makes no assertion that it is
constant across distinct norm or source data.

\par

\medskip\noindent\textbf{DZ.6. Scope and exact verification}

\par\nobreak

\item DZ.27 concerns only the two stated initial/reflected fan objects,
the one III-star cap and the literal scalar \(v^{km}\lambda_0\).
The single endpoint \(q^{nk}\), its literal pullback
\(q^{2nk}\), their period reindexings, the bounded seam
multipliers, and the cap square have all been carried.
The earlier diagnostic's algebra is recovered, but its formerly
missing plumbing/seam map is now (DZ.5)--(DZ.14), with the actual
single-to-double square in DZ.4.
All symbolic chart, seam, quaternion-exponent, trace and pair
identities are checked in
\path{code/part3/source_zero_double_loop_product.py}.
The complete geometric proof has passed the independent checks and
root audits, including the literal pullback and actual relative square.
The scalar \(C_1\) consistency diagnostic
is separate and is not used to extend this exact \(C_2\) comparison.

\end{enumerate}
\end{proof}


\section{Results not yet established: the remaining attachments}



\begin{theorem}\label{night:P3K.75}
\StatementRecord{night:P3K.75}{open}{S6\_06\_BLIND: GPT-6 Astra (OpenAI, ChatGPT Pro mode, effort Pro), recorded re-derivation in the checkpoint dated 18 November 2026}{\cref{night:P3K.58}, \cref{night:P3K.59}, \cref{night:P3K.60}, \cref{night:P3K.65}, \cref{night:P3K.68}, \cref{night:P3K.69}, \cref{night:P3K.71}, \cref{night:P3K.72}, \cref{night:P3K.73}, \cref{night:P3K.74}, \cref{night:P3K.76}, \cref{night:P3K.77}, \cref{night:P3K.78}, \cref{night:P3K.79}, \cref{night:P3K.80}, \cref{night:P3K.82}, \cref{night:P3K.83}, \cref{night:P3K.85}, \cref{night:P3K.92}}
P3K.71 computes the selected-cusp source return on every lattice-basis
fan. P3K.72--74 give the full corrected scalar-fan algorithms in their
specified whole-zero source scope. P3K.85--91 complete the following
additional sectors through their actual source, Cartier, root, normal
and inverse-closing maps.

For the source-zero E7 branches put
\[
 (a,b,m)=(1,2,4p)\quad\hbox{or}\quad(2,1,4p-1),\qquad
 N=km,\quad \mu=2-a+6pk.
\]
The scalar is \(\lambda=v^N\lambda_0\); the fixed-scalar case is
\(k=0\). For every vertical fan \(T_j\), P3K.89 proves
\[
\begin{aligned}
 A_j={}&2b\binom{j+1}{3}+(3-b)\binom{j+1}{2}
       +k(a-1)(3p-1)-km\binom j2,\\
 (k_M,\ell_M)={}&(\mu+6A_j,4\mu).
\end{aligned}
\]
The full normal clutch used here is computed in P3K.86.
P3K.90 transports these free integer data to every
\(P\in\mathrm{GL}_2(\mathbb Z)\) by the actual two-ruling maps and
the uniquely solved integral target-prism system. P3K.88 supplies
explicit fixed-scalar fan witnesses for every element of
\(\Theta_7\) on each branch.

For the component-supported OS49 source, P3K.77--79 complete all
one-cap lattice-basis fans, marked smooth order-one clutches and
fixed-source holomorphic endpoints. P3K.81 supplies the exact
integral obstructions for the additional smooth caps and good-II
caps treated there. P3K.84--85 construct the two actual finite caps
\(m=3k\), \(h=6l+1\) and their positive canonical return.
With the full ordinary-clutch integer \(\sigma\), P3K.91 proves
\[
 D=hn_A+mn_B+mh\sigma=\epsilon=\pm1,\qquad
 H^2(\mathcal H)=\mathbb Z\oplus\mathbb Z/k\oplus\mathbb Z/h.
\]
The class is \((D,d_A\bmod k,d_B\bmod h)\). Its complete
holomorphic endpoint group is the physical \(\mathbb C^*\).
The actual shortest parameter \(w=r^3\), \(\lambda=w^{kh}\), has
positive-fan physical pair \((2kh,8k(3l+1))\) and integer \(kl\).
P3K.87 and P3K.91 compute every fan and every admitted winding
\(Nkh\), with the endpoint and native covering maps constructed
before integral characteristic naturality is applied.

P3K.92 completes the infinity boundary of the marked E7
four-boundary parameter family on every fan. In its exact marking,
\[
 I_a(P)=-Z_a(P)-Z_a(JP)-C(P)\pmod {28},\qquad
 J=\operatorname{diag}(-1,1).
\]
The branch reflection and the infinity rescaling are given by
their full character maps. The infinity root endpoint retains
\(A^4D^{-4}=S_{r_0^{4n}}L^4\), its inverse scalar path on both
ordinary and cap charts, and the actual transported lattice
period \(-n(2\omega_1+2\omega_2)\). The two unselected
\(v=1/8\) discs have the P3K.53 extension.

P3K.93 evaluates that complete algorithm in closed form. Put
\((x,y)^t=P^{-1}(0,1)^t\) and \(S=x^2+xy+y^2\). Then
\[
 I_a(P)=1-a-3(a-1)(S-1)-\frac{2a+1}{4}(S-1)(S-3)
          \pmod {28}.
\]
The proof supplies the full paired prism equations, the common
first-step seed and both generator identities. Its complete images
are \(\{0,6,8,10,14,20,22,24\}\) for \(a=1\) and
\(\{5,7,9,13,19,21,23,27\}\) for \(a=2\), with an actual
unimodular fan witness for every entry.

The outstanding work is the remaining marked source-moduli
generators, singular-site sectors and endpoint scopes outside
the completed families just listed. In the E6 family
\[
 y^2=x^3+rx^2+tx+t^2,\qquad G=(0,t),
\]
P3K.94 completes the unselected \(r=0\) collision for the actual
line \(M_6=\mathcal O(6G-O)\), suspension by \(-G\), and the
quadratic reduced \(IV^*\) cap, with either signed primitive norm,
each retained character, the full reference annular map and every
selected lattice-basis fan. Its complete proper disc family has
identity closing, so each of the corresponding marked-cover
generators \(w^2=3/2\) has return zero. Its proof retains the full
section \((t-t_{\rm sel})B_6/t^3\), the root coefficient \(4h\),
the signed norm and every source/line period cochain.
P3K.96 extends this disc construction to both full progressions
\[
 (a,b;m)=(6,1;3l+2),\quad l\ge0,
 \qquad (a,b;m)=(2,3;3l),\quad l\ge1.
\]
The root fibres are respectively reduced type \(IV\) and smooth
elliptic. The proof gives the actual higher A2 translation, every
regular graph ray \((j,1)\), both section coefficient maps, the
good-root exceptional correction, the literal chord lift with central
coefficient one, its full norm unit of value four, every finite
character, and the complete reference-relative annular cochain.
For both primitive norm signs and every selected lattice-basis fan,
the proper source-disc family has identity closing and return zero.
P3K.97 completes the native E6 infinity-coordinate disc at either
\(w^2=3\), for both cap progressions of P3K.96. The actual source,
Cartier, root, physical-circle and annular maps carry the punctured
original family to a proper smooth sphere family on the full native
disc. Both point and tensor Tate frames, every finite character,
the transported period discrepancy and the full scalar closing are
retained. Each prescribed lattice-basis fan and both primitive norm
signs have marked positive \(w\)-loop return zero. The exact original
parameter is
\[
 \epsilon=1/r=(w^2-3)^2/(3-2w^2).
\]
This marked loop has degree two over the original infinity meridian;
no shorter original-parameter endpoint or return is inferred.
The \(r=1/4\) order-three collision
has an exact central line divisor \(3F-C_1+C_2\) and the actual
two-ruling multiplier \(8\lambda q w^{-3}\); its global oriented
attachment and return are still being calculated.

P3K.95 completes the E6 selected-cusp generator on the ordered
\(w\)-cover for \(M_6\), suspension by \(-G\), and the quadratic
cap, with both primitive norm signs, all retained characters and
full reference annular maps, and every selected lattice-basis fan.
The exact source, Cartier, root, exterior, normal and Cox arrows
identify its full all-fan algorithm with P3K.71. On the vertical fan
\(T_j\) its physical pair and class are
\[
 (k_M,\ell_M)=(-6-18j(j-1),0),\qquad
 [\Sigma_j]=-1-3j(j-1)\pmod {28}.
\]
Its native free pair is \((0,-6-18j(j-1))\). The actual parameter
map satisfies \(r-1/3=-w^4/(3(w^2-3)^2)\), so the stated
ordered loop has degree four over the original source parameter.
P3K.98 extends this full selected-cusp calculation to both complete
cap progressions of P3K.96, with the same fixed scalar, both norm signs,
every finite character and complete annular reference. Its actual
point-to-tensor line arrow, all higher root graphs, even-parameter
annular map and both normal caps identify the exact vertical free
pair above and every K71 fan step. The full smooth-fibre theta
reflection and its target period give the actual height-one arrow;
the relative \(p_1\) calculation carries the explicit regular-level
cobordism and its boundary transgression.
Other selected-node generators and the remaining singular-site and
endpoint scopes still require their full comparison maps. The working notes assign no
return to an unfinished attachment. Part 3.4 remains incomplete.
\end{theorem}
\begin{proof}[Current reductions and preserved earlier attachment history]
\begin{enumerate}[label=\arabic*.]
\item The current reductions are the full proofs P3K.71--74 and
P3K.76--98 cited in the statement. The recovered P3K.85--91
sources have fresh mathematical and visual audits; their
reconstruction records identify the lost earlier bytes and the
newly verified files. P3K.92--98 are new complete proofs.

\par\medskip\noindent\textbf{Historical working record from checkpoint 19.}\par\nopagebreak[4]
The following retained paragraphs document the successive earlier
attachment attempts and their status at that time. The current
status and completed replacements are the statement above and
P3K.85--93. In particular their earlier unresolved source-zero
and reduced-II stages are preserved here as history. The exact
pre-update file is in
\path{records/checkpoint22_prior_master68/notes/part3_remaining_return_obligations.tex}.

\item K80 now supplies the complete reduced period-strip bordism:
\[
 \partial\mathcal W=-N-\mathbb {CP}^{3}_{2}-B_F-B_4
                                      +\mathbb {CP}^{3}_{1}.
\]
The respective auxiliary physical indices are \(-1,1,0,0\),
giving \(A(N)=0\). Its full inverse period and mixed-face
charts identify this outside boundary with the actual parameter
return through K68. The whole fixed trace occurs once; its
normal clutch remains in the proof. K82 carries the resulting
initial class through every actual vertical \(C_1\) flop.

\item The following paragraphs retain the earlier attempted routes
and the exact intermediate obstructions. The prior full status
text is preserved in
\path{records/checkpoint19_prior_master64/notes/part3_remaining_return_obligations.tex}.
The unresolved good branch has its own chronological two-step
construction; no result on it is inferred from the reduced proof.

\par\medskip\noindent\textbf{The earlier paired source-zero constraint.}\par\nopagebreak[4]
The physical/dual line map and both Hopf-gluing orientations are
proved in K68. K69 carries them through the actual global
parameter groupoid and the full source-infinity cap comparison.
They give the paired sum. K60's retained single native diagonal
pair \((1-a,a)\), with the literal single-to-double pullback
product, gives the difference \(Z_{J\mathcal T}-Z_{\mathcal T}=a\).
Solving these two equations in the stated group leaves exactly
the kernel \(\{0,14\}\) of multiplication by two. These actual
maps prove the displayed two possibilities; choosing either
element requires further information.

\item The completed selected-cusp and vertical-source calculations
do not supply that further integer: their loops and endpoints
are the explicitly different paths recorded in the global
groupoid. Their composition yields the paired relation just
used. The precise groupoid arrows, rather than an identification
of individual loops, are retained in K58--60 and K69.

\par\medskip\noindent\textbf{The attempted source completion and its actual obstruction.}\par\nopagebreak[4]
The marked small-resolution source charts and both Cartier
spinor maps have been constructed in the source-zero completion
notes. On the unresolved central contraction the cyclic
generator fixes the cone origin; every power fixes it. Thus
the literal origin stabilizer is infinite and that proposed
central quotient does not meet the proper-discontinuity
condition used for the punctured construction. Deleting or
resolving this point changes the required gluing problem.
No index for a smooth isolated quotient has been assigned
to that infinite-stabilizer point.

\item For the earlier specified fractional-period root quotient
of SZ.27, the root map has local base equation \(q=z^m\);
its central
divisor therefore has multiplicity \(m\), rather than the
specified simple selected zero. Moving that branching to a
regular value inserts its actual additional cap. Those maps
are retained with these exact changes of the requested data.
The subsequently constructed global noninvariant norm sections
have their proved source identities, but their full line and
root action has not been constructed. The multiplicity
calculation just given applies to the earlier SZ.27 quotient;
no multiplicity is assigned to the new source-section route.
Its full source-line-cap comparison is still required.

\par\medskip\noindent\textbf{The real trace and its extra end.}\par\nopagebreak[4]
The current actual normalization trace is
\(F=\operatorname{Bl}_{(-1,0)}(\mathbb P^1_s\times D_v)\).
Its two blowup-chart maps are written in
\path{notes/part3_source_zero_real_trace_working.md}; their
Jacobians are units on the exceptional curve. The actual
source line restricts to \(\mathcal O(-aE)\), with
\(E^2=-1\). The two transported spinor determinants are
\((1-2a)E\) and \(-(1+2a)E\), giving local intersection
indices \(a(a-1)/2\) and \(a(a+1)/2\). Their difference
agrees with the proved native difference \(a\).

\item The exact Cox coordinates also give
\[
 k=s_1-s_0,\qquad l=t_0+t_1,\qquad
 q=kl-\tfrac14(s_0l-t_1k)^2.
\]
The transverse quadratic form along \(k=l=0\) has unit
Hessian. Its equivariant factorization produces the associated
node chart \(UV=q\), with the actual auxiliary line character.
At \(q=0\) this compactification carries an additional
vertical fixed surface. Therefore identifying the complete
fixed trace with the horizontal blowup alone has not been
proved. The missing surface and its end must be included
before those local indices can be used as global values.

\item The residual ruling end is the actual
\(S^2\times S^3\) clutch
\(\tau(h,q_H)=(q_H^{-1}h q_H,q_H)\), with its typed inverse
in the forward period direction. Its mapping torus is the
projective line bundle over \(S^3\times S^1\) of the
quaternion clutch, with \(c_2=\pm1\) according to that
unresolved boundary placement. The charge-\(a\) line gives
the exactly computed pair \((-a^3c_2,-4ac_2)\) and index
\(a(a^2-1)c_2/6\). These computations record the dependence
on the actual boundary orientation; they do not select its
placement in the complete bordism.

\item The full finite cap arrow is now explicit. In its root
Cartier frame,
\[
 U_{\rm cap}=z^{-ma/2}U_a,\qquad
 B=2(-1)^{b+1}\lambda h U_{\rm cap}^{b},\qquad
 h=1-T_0q_\infty.
\]
For \((a,b)=(2,1)\) it commutes with the actual finite cap
map and retains one forward \(t\)-to-\(s\) ruling arrow.
For \((a,b)=(1,2)\), its exact commutator changes the finite
character by \(m/2\), and the second step closes it. These
intermediate objects and the complete annular factors are
retained in
\path{notes/part3_source_zero_annular_attachment_working.md}.

\item The active construction cuts the compact suspension outside
the short fixed-normalization graph, removes a fixed compact
Cox-origin neighbourhood, and uses the finite Cox moment
cobordism on that cut. The source translation, its finite-cap
arrow above and the associated-node chart have been written.
On the stable outside portion the newly constructed exact
metric arrow is
\[
 N_t=|\xi|\|t\|^a=R^aN_s,\quad R=\|s\|\|t\|,\qquad
 (\xi,\xi_{\rm dual})\longmapsto
 (R^a\xi,R^{-a}\xi_{\rm dual}).
\]
It preserves \(UV=q\), both node branches and the physical
circle on \(R\ge\varepsilon\). The Tate index change
\(j\mapsto j+1\) absorbs exactly the original \(q\)
factor. At the inner boundary the fixed trace is the actual
\(S^3\), with coordinate \(p=e^{i\phi}(S_0,S_1)\) and
positive spinor frame
\((e^{i(a-1)\phi}z,e^{ia\phi}\overline V)\).
The local framed surgery is now constructed in RT.10. In the
reduced case the positive node model is the explicitly identified
total space of \(\mathcal O(-1)\oplus\mathcal O(-2)\) over
\(\mathbb P^2\), with its fixed surface and isolated charge-one
point retained. RT.13 connects its radial boundary to the
product-radius cut by the actual compact region
\(\{\rho\le\rho_0,\ R\ge\varepsilon\}\), with two small-\(R\)
arms having cores \(S^3\) and \(\mathbb{RP}^3\).

\item The earlier \(j=-1\) designation used the SZ chart index.
TL supplies the full source and Cartier map to the current Tate
fan, retaining
\[
 Z_0W_0=q_1q_2u_Q,\quad P=q_2^{-1}u_P,\quad
 Z=P/S=1/W_0,\quad Z_0=u_Qu_PT_0,\quad W_0=q_2S/u_P.
\]
Its valuation arrow is \((x,y,h)\mapsto(h-x,y,h)\).
The associated positive square maps to the current negative
fan \(\mathcal T_0\), and the opposite resolution to
\(\mathcal T_1\). The actual source period and its full theta
line factor occur before that valuation. The two resolution
charts and their unit inverses extend across both branches.

\item EA.17 gives the full two-node refinement using the normalized
Morse functions \(\alpha\beta=q_1q_2\) and exact factors
\(F_+F_-=q_1\):
\[
 S_0=F_+,\quad T_0=F_-,\quad
 S_1=-\beta/(2F_-),\quad T_1=-2\alpha/F_+.
\]
The branch-divisor argument makes both quotients holomorphic;
their Jacobian at the origin is \(4\). They give literally
\(q_1=S_0T_0\), \(q_2=S_1T_1\), and
\(\widetilde v=q_1-q_2\). The marked coefficient is multiplied
by \(\kappa^a=((S_1-S_0)/(2s_0))^a\), an extending unit;
the two original Cartier \(r\) maps are retained for \(a=2\).

\item There is also an exact obstruction to replacing the present
end by a radial sphere link. In the charge-one chart,
\[
 R^2=(|S_0|^2+|S_1|^2)(|zV|^2+|T_1|^2).
\]
Thus \(R\) vanishes on the entire fibre over
\([0:0:1]\), not only on its zero section. The verified
local fixed-boundary surgery therefore requires an actual
comparison to this product-radius end. Its closed
projective spinor cap has pair
\((3a^2-3a+1,4)\) and index \(a(a-1)/2\);
the charge-one point has pair \((-1,-4)\) and index zero.
The oriented local surgered end on the reduced branch
has pair \((-8,-8)\) and index \(-1\), as proved in
RT.10 of the real-trace working record and independently
checked there. This remains a local statement. Removing
the capped exceptional disk from the global fixed trace
also removes its intersection class; it cannot be counted
again in the complementary capped trace. The lower arm has an actual charged-shape comparison. In its
charge-two cover put \(\sigma^2=\|x\|^2+|y|^2\) and
\(u=|y|^2/\sigma^2\). Near that arm,
\(\rho=\sigma^2\sqrt{1-u}\). Extend this positive factor
to an invariant \(f\) with \(1/\sqrt2\le f\le1\).
The ambient annulus
\[
 \{\rho_0/2\le\sigma^2\le\rho_0/f\}/\mu_2
\]
has positive width and retains the source functions by inclusion.
Its physical line is carried by the explicit Hopf map
\([u,z]\mapsto z u^{\otimes2}\) into \(\mathcal O(-2)\).
Its full quaternion charts have \(w'=w^{-1}\) and
\(u'=(w/|w|)u\), retaining the hemisphere inversion.
EA.15 gives their native ruling-quaternion comparison.
RT.15 carries the actual canonical spin lift \(+\omega_8\).

\item After the actual node reindex and source reset, the common
coefficient is \(U_j\). The gauge \(U_j=1\), transported to
\(V_{j-1}=1\), gives \(N_t=\|T\|^a\), \(N_s=\|S\|^{-a}\).
The transported coefficient is \(V_{j-1}^{-1}\), retaining
the \(q\) factor in the other node coordinate. Thus the entire
finite faces are
\[
 H_t=(D^4_S(r)\times S^3_T(r))/\mu_a,\qquad
 H_s=(S^3_S(r)\times D^4_T(r))/\mu_a.
\]
On \(R>0\), the actual physical factor \(d=(R/r^2)^a\)
gives \((S,T)\mapsto(d^{-1/a}S,d^{1/a}T)\), preserving
\(q_1,q_2\) and the original base functions.
The forward contracting field \(-\eta\), where
\(\eta=-S\partial_S+T\partial_T\), enters \(H_t\) and
exits \(H_s\). The oriented target-minus-source chain is
the positive product-ball boundary. Both cores are retained;
the normalization is not extended across \(R=0\).

\item These maps completed the local face and spin comparisons.
The subsequent reduced global period-strip equation, its whole
fixed cap and every retained or excised face are proved in K80.
K83 supplies the good branch's separate two-primitive finite
strip, its exact upper free source trace and final actual-return
comparison. Its complete oriented physical quotient boundary is
\[
              -N-P_0-P_1-B_F-B_4+Q_{\rm up}.
\]
Both charged physical lines are \(\mathcal O(-1)\); the upper
line has its explicit global nonzero coefficient frame. The whole
fixed and four-sphere normal clutches remain in the calculation.
Every auxiliary characteristic index is zero. The same selected
Hopf tube and K68 inverse-clutch map identify the resulting
\(A(N)=0\) with the actual good parameter return.

\item The alternative anti-holomorphic Tate filling retains its
boundary parameter \(\kappa=4\lambda^2/w\) and replaces
it on the filling disc by
\(4\lambda^2\overline w/R^2\), with a positive radial
collar interpolation. This matches the displayed boundary
parameter and stays in the contraction domain. Its full
source-line-cap square and global orientation have not
been supplied, so it also has not assigned \(e\).

\par\medskip\noindent\textbf{The other remaining return sectors.}\par\nopagebreak[4]
For OS49 the complete actual primitive periods in K65 have
windings \((0,-m,0,2m)\) on the positive \(r\)-loop.
Its two diagonal values and shorter endpoint are proved.
K76 extends the relative-line proof of K72 to the two
bottom-row windings \((L,J)=(-m,2m)\).
The actual boundary connection difference remains
\(-n(du+\alpha\,dv)\); the additional \(L\)-term enters
the complete collar cube. The completed integer comparison,
with its exact same-closing-exponent local calibration, is in
\path{notes/part3_bottom_row_winding_working.md}.
It gives the actual vertical-fan integers
\(-m(n-1)^2\), or \(-k(n-1)^2\) on the shorter loop.
It does not choose a free integer from the earlier residue.
K77 completes the nonvertical construction with its full
period and Cox maps, its integer target-prism solution,
and all intermediate cap and shorter-endpoint arrows.
It also proves the exact image on every cap order and
the retained marked cap-data loop calculation. The
fixed smooth order-one clutch sectors are now completed
by K78, including its signed norm and actual endpoint conjugacy.
The remaining scopes are stated above.

\item The source-loop catalogue, its discriminant and braid maps,
and the endpoint constraints remain in the named-loop and
source-moduli working records. K71 now computes all fans of
the specified selected-cusp loop; K73 and K74 compute all
fans of the specified whole-zero scalar loops and their
retained marked cap/site variations. The remaining source
generators require their own closed source and cap arrows.
K79 now carries every holomorphic endpoint of the stated fixed
OS49 sphere data through its full lattice, local inverse, physical
circle and finite-cap groupoid. It gives the same return for all
specified endpoints at each allowed scalar winding. For the other
remaining source data, those endpoint maps still require their
own such comparison; no value is assigned by changing a marking.

\end{enumerate}
\end{proof}


\end{document}
